%==============================================================================
%  Competitive Programming Notebook — Guidebook
%  Autor: Juan  ·  Universidad del Bosque
%  Formato: doble columna (izquierda / derecha), estilo team notebook ICPC
%  Compilar con: pdflatex guidebook.tex   (dos veces, para el índice)
%==============================================================================
\documentclass[10pt,a4paper,twocolumn]{article}

% ---- Codificación e idioma -------------------------------------------------
\usepackage[T1]{fontenc}
\usepackage[utf8]{inputenc}
\usepackage{lmodern}
\usepackage{microtype}
\usepackage{float}
\usepackage{enumitem}
\usepackage{amsmath}   % \text en modo matematico
\usepackage{amssymb}
% Nombres en español (sin depender de babel-spanish, que no está instalado aquí)
\renewcommand{\contentsname}{Índice}
\renewcommand{\partname}{Parte}

% ---- Diseño de página ------------------------------------------------------
\usepackage[a4paper,margin=1.6cm,columnsep=0.8cm]{geometry}
\setlength{\columnseprule}{0.4pt}             % línea fina que divide izquierda/derecha
\usepackage{xcolor}
\definecolor{acento}{RGB}{20,45,85}           % azul marino formal
\definecolor{acentoClaro}{RGB}{70,110,170}
\definecolor{grisCodigo}{RGB}{245,245,247}
\definecolor{comentario}{RGB}{95,120,95}
\definecolor{palabraClave}{RGB}{30,70,150}
\definecolor{cadena}{RGB}{150,80,40}
\def\columnseprulecolor{\color{gray!40}}

% ---- Encabezados y pie -----------------------------------------------------
\usepackage{fancyhdr}
\usepackage{titlesec}
\pagestyle{fancy}
\fancyhf{}
\fancyhead[L]{\small\itshape Competitive Programming Notebook}
\fancyhead[R]{\small\nouppercase{\itshape\rightmark}}
\fancyfoot[C]{\small\thepage}
\renewcommand{\headrulewidth}{0.4pt}

\titleformat{\section}
  {\normalfont\large\bfseries\color{acento}}{\thesection}{0.6em}{}
\titleformat{\subsection}
  {\normalfont\normalsize\bfseries\color{acentoClaro}}{\thesubsection}{0.5em}{}
\titleformat{\part}
  {\normalfont\Large\bfseries\color{acento}\filcenter}{\partname~\thepart}{0.5em}{}

% ---- Resaltado de código ---------------------------------------------------
\usepackage{listings}
\lstset{
  backgroundcolor=\color{grisCodigo},
  basicstyle=\ttfamily\scriptsize,
  keywordstyle=\color{palabraClave}\bfseries,
  commentstyle=\color{comentario}\itshape,
  stringstyle=\color{cadena},
  numbers=left, numberstyle=\tiny\color{gray}, numbersep=5pt,
  showstringspaces=false, breaklines=true, tabsize=2,
  frame=single, rulecolor=\color{gray!40},
  columns=fullflexible, keepspaces=true,
  extendedchars=true,
  literate=
    {á}{{\'a}}1 {é}{{\'e}}1 {í}{{\'i}}1 {ó}{{\'o}}1 {ú}{{\'u}}1
    {Á}{{\'A}}1 {É}{{\'E}}1 {Í}{{\'I}}1 {Ó}{{\'O}}1 {Ú}{{\'U}}1
    {ñ}{{\~n}}1 {Ñ}{{\~N}}1 {¿}{{?`}}1 {¡}{{!`}}1 {°}{{\textdegree}}1
}
\lstdefinestyle{cpp}{language=C++}
\lstdefinestyle{java}{language=Java}
\lstdefinestyle{py}{language=Python}

% ---- Enlaces (índice navegable) --------------------------------------------
\usepackage[colorlinks=true,linkcolor=acento,urlcolor=acentoClaro]{hyperref}

\setcounter{tocdepth}{2}
\setcounter{secnumdepth}{2}

% Marcador de secciones aún sin contenido
\newcommand{\pendiente}{\textit{\color{gray}Contenido en desarrollo.}\par\medskip}

%==============================================================================
\begin{document}

% ===== PORTADA =============================================================
\onecolumn
\thispagestyle{empty}
\begin{titlepage}
  \centering
  \vspace*{\fill}

  {\large\scshape Universidad del Bosque}\\[4pt]
  {\normalsize Facultad de Ingeniería}\\[2pt]
  {\normalsize Ingeniería de Software \& Ingeniería en Robótica}\\[24pt]

  \color{acento}\rule{\linewidth}{1.2pt}\\[8pt]
  {\Huge\bfseries Competitive Programming}\\[10pt]
  {\Huge\bfseries Notebook}\\[10pt]
  {\large\itshape Algoritmos, estructuras de datos y plantillas}\\[8pt]
  {\large C\texttt{++} \quad}\\[6pt]
  \color{acento}\rule{\linewidth}{1.2pt}\\[28pt]

  \color{black}
  {\large\bfseries Juan Carlos Morales}\\[4pt]
  {\normalsize Doble titulación: Sistemas \& Robótica}\\[36pt]

  {\normalsize Bogotá, Colombia \quad$\cdot$\quad 2026}

  \vspace*{\fill}
\end{titlepage}

% ===== ÍNDICE ==============================================================
\thispagestyle{empty}
\tableofcontents
\clearpage

% ===== CUERPO (doble columna) ==============================================
\twocolumn

\part{Introduccion}
%---------------------------------------------------------------------------
\section{Fundamentos Básicos}
%---------------------------------------------------------------------------

\subsection{Entrada / Salida Rápida (Fast I/O)}
Usar lectura rápida para evitar TLE, hay 3 niveles de lectura rápida, estandarizar Tier 2 y usar Tier 3 (lector por
bytes) solo cuando haya del orden de $10^6$--$10^7$ lecturas.

\begin{table}[h]
\centering
\caption{Casos donde se recomienda Fast I/O}
\begin{tabular}{|p{0.36\columnwidth}|p{0.50\columnwidth}|}
\hline
\textbf{Tipo de problema} & \textbf{Ejemplo} \\
\hline
Arrays gigantes & Leer/ordenar/sumar $10^6$--$10^7$ enteros \\
\hline
Muchos test cases (T grande) & $T \leq 10^5$, cada uno trivial pero se acumula \\
\hline
Grafos grandes & $M \sim 10^6$--$2*10^6$ aristas para BFS/DFS/Dijkstra/MST \\
\hline
Geometria con muchas consultas & $L,S$ grandes + muchos test cases seguidos hasta EOF \\
\hline
Muchos tokens de texto & Strings/lineas que suman millones de caracteres \\
\hline
Judges clasicos "trampa de I/O" & SPOJ INTEST, ARRAYSUB, etc. \\
\hline
\end{tabular}
\end{table}

Regla practica: si N/M es del orden de $10^4-10^5$, cin con
\texttt{sync\_with\_stdio(false);} cin.tie(nullptr); alcanza de sobra. El FastReader
es la "artilleria pesada" para cuando eso deja de ser suficiente.

\textbf{Tier 1 — Básico:} \texttt{cin} / \texttt{cout}.
\begin{lstlisting}[style=cpp]
int entero; double decimal; string palabra, linea;
cin >> entero >> decimal >> palabra;  // >> lee UN token
// ! cin >> deja el '\n' pendiente;
// ! el 1er getline sale VACÍO -> hay que ignorarlo
cin.ignore();
getline(cin, linea);                  // línea completa
\end{lstlisting}
\textbf{Tier 2 — Rápido:} desactiva la sync con C (una vez en
\texttt{main}). Punto dulce para casi todo.
\begin{lstlisting}[style=cpp]
ios_base::sync_with_stdio(false);
cin.tie(nullptr);            // ! no mezclar con scanf/printf
cout << entero << "\n";      // "\n", nunca endl (endl frena)
\end{lstlisting}
\textbf{Tier 3 — Más rápido:} lector por bytes con
\texttt{fread}. Solo con $10^6$–$10^7$ lecturas.
\begin{lstlisting}[style=cpp]
struct FastReader {
    static const int TAM = 1 << 16;   // 64 KB
    char bufer[TAM];
    int puntero = 0, leidos = 0;

    int leer() {
        if (puntero == leidos) {
            leidos = (int) fread(bufer, 1, TAM, stdin);
            puntero = 0;
        }
        return leidos == 0 ? -1 : bufer[puntero++]; // -1=EOF
    }
    int nextInt() {
        int resultado = 0, b = leer();
        while (b <= ' ') b = leer();      // saltar blancos
        bool negativo = (b == '-');
        if (negativo) b = leer();
        while (b >= '0' && b <= '9') {
            resultado = resultado * 10 + (b - '0'); b = leer();
        }
        return negativo ? -resultado : resultado;
    }
    long long nextLong() {
        long long resultado = 0; int b = leer();
        while (b <= ' ') b = leer();
        bool negativo = (b == '-');
        if (negativo) b = leer();
        while (b >= '0' && b <= '9') {
            resultado = resultado * 10 + (b - '0'); b = leer();
        }
        return negativo ? -resultado : resultado;
    }
    double nextDouble() {
        double resultado = 0, divisor = 1; int b = leer();
        while (b <= ' ') b = leer();
        bool negativo = (b == '-');
        if (negativo) b = leer();
        while (b >= '0' && b <= '9') {
            resultado = resultado * 10 + (b - '0'); b = leer();
        }
        if (b == '.') {
            b = leer();
            while (b >= '0' && b <= '9') {
                resultado += (b - '0') / (divisor *= 10); b = leer();
            }
        }
        return negativo ? -resultado : resultado;
    }
    string next() {                       // un token
        string s; int b = leer();
        while (b <= ' ') b = leer();
        while (b > ' ') { s += (char) b; b = leer(); }
        return s;
    }
    string nextLine() {                   // línea completa
        string s; int b = leer();
        while (b != '\n' && b != -1) {
            if (b != '\r') s += (char) b; b = leer();
        }
        return s;
    }
};
\end{lstlisting}

\subsection{Formateo de Salida}


La función \texttt{printf} utiliza una cadena de formato para indicar cómo
debe interpretarse y mostrarse cada argumento. Cada especificador comienza
con \texttt{\%} y corresponde al argumento que aparece después de la cadena.

\begin{lstlisting}[style=cpp]
// %d  -> entero (int)
printf("%d\n", 42);
// Entrada: 42
// Salida: 42

// %lld -> entero largo (long long)
long long x = 1234567890123LL;
printf("%lld\n", x);
// Entrada: 1234567890123
// Salida: 1234567890123

// %5d -> ancho minimo de 5 caracteres, alineado a la derecha
printf("%5d\n", 42);
// Entrada: 42
// Salida: "   42"

// %-5d -> ancho minimo de 5 caracteres, alineado a la izquierda
printf("%-5d|\n", 42);
// Entrada: 42
// Salida: "42   |"

// %05d -> ancho 5, completando con ceros
printf("%05d\n", 42);
// Entrada: 42
// Salida: "00042"

// %+d -> muestra siempre el signo
printf("%+d\n", 42);
// Entrada: 42
// Salida: "+42"

// %x -> hexadecimal usando letras minusculas
printf("%x\n", 255);
// Entrada: 255
// Salida: ff

// %X -> hexadecimal usando letras mayusculas
printf("%X\n", 255);
// Entrada: 255
// Salida: FF

// %.2f -> numero real con exactamente 2 decimales
printf("%.2f\n", 3.14159);
// Entrada: 3.14159
// Salida: 3.14

// %.0f -> numero real sin decimales (redondea)
printf("%.0f\n", 3.7);
// Entrada: 3.7
// Salida: 4

// %e -> notacion cientifica
printf("%e\n", 31415.9);
// Entrada: 31415.9
// Salida: 3.141590e+04

// %c -> un solo caracter
printf("%c\n", 'A');
// Entrada: 'A'
// Salida: A

// %s -> cadena de caracteres
printf("%s\n", "hola");
// Entrada: "hola"
// Salida: hola

// Se pueden combinar varios especificadores
printf("%d %lld %.2f\n", 42, 10000000000LL, 3.14159);
// Entrada: 42, 10000000000, 3.14159
// Salida: 42 10000000000 3.14

// %% -> imprime un '%' literal
printf("100%%\n");
// Entrada: no recibe ningun argumento para %
// Salida: 100%

// ! NO uses %n en C/C++ (puede ser peligroso)
// Para scanf, un double se lee con %lf
\end{lstlisting}

\subsubsection*{Estructura de un especificador}

Un formato puede combinar varias partes:

\begin{lstlisting}[style=cpp]
//   % [flags] [width] [.precision] [specifier]
//   %    0       5         .2          f

printf("%05d\n", 42);
//  %  -> comienza el formato
//  0  -> rellena con ceros
//  5  -> ancho minimo de 5 caracteres
//  d  -> interpreta el argumento como int
//
//  Salida: 00042
\end{lstlisting}

Los modificadores más utilizados en programación competitiva son:

\begin{center}
\begin{tabular}{|c|c|c|}
\hline
\textbf{Formato} & \textbf{Tipo} & \textbf{Ejemplo de salida} \\
\hline
\texttt{\%d} & \texttt{int} & \texttt{42} \\
\hline
\texttt{\%lld} & \texttt{long long} & \texttt{1234567890123} \\
\hline
\texttt{\%c} & \texttt{char} & \texttt{A} \\
\hline
\texttt{\%s} & \texttt{string/char*} & \texttt{hola} \\
\hline
\texttt{\%f} & \texttt{double} & \texttt{3.140000} \\
\hline
\texttt{\%.2f} & \texttt{double} & \texttt{3.14} \\
\hline
\texttt{\%x} & \texttt{int} & \texttt{ff} \\
\hline
\texttt{\%X} & \texttt{int} & \texttt{FF} \\
\hline
\end{tabular}
\end{center}

\textbf{Regla práctica:} en programación competitiva, los formatos más
importantes son \texttt{\%d} para \texttt{int}, \texttt{\%lld} para
\texttt{long long}, \texttt{\%c} para caracteres, \texttt{\%s} para cadenas
y \texttt{\%.2f} (o una precisión diferente) para números reales.


\subsection{Condicionales y Ciclos}

\begin{lstlisting}[style=cpp]
// if / else: ejecuta un bloque segun una condicion
if (x > 0) cout << "pos";       // x > 0 -> "pos"
else if (x < 0) cout << "neg";  // x < 0 -> "neg"
else cout << "cero";             // x == 0 -> "cero"

// Operador ternario: condicion ? valor_si_true : valor_si_false
string s = (x >= 0) ? "no neg" : "neg";
// Si x >= 0, s = "no neg"; si no, s = "neg"

// switch: compara una expresion con diferentes valores
// Solo funciona directamente con tipos integrales como int o char,
// NO con string.
switch (x) {
    case 1: /*...*/ break;       // x == 1
    case 2: case 3: /*...*/ break; // x == 2 o x == 3
    default: /*...*/ ;           // ningun case coincide
}
// ! Sin break, se continua ejecutando el siguiente case (fall-through)

// C++17: se puede declarar una variable dentro del if
if (int r = x % 2; r == 0) {}
// r solo existe dentro del if y sus bloques asociados

// for clasico: usa un indice y se repite mientras i < n
for (int i = 0; i < n; i++) {}

// range-based for: recorre directamente los elementos
// No proporciona el indice.
for (int v : datos) {}

// & obtiene una referencia al elemento original.
// Por tanto, las modificaciones afectan al vector.
for (int &v : datos) v *= 2;

// while: ejecuta mientras la condicion sea verdadera
while (c < n) c++;

// do-while: primero ejecuta el bloque y DESPUES verifica la condicion.
// Por eso se ejecuta AL MENOS una vez.
do {
    c--;
} while (c > 0);

// break -> termina inmediatamente el ciclo o switch
// continue -> salta a la siguiente iteracion del ciclo

// ! C++ NO tiene break con etiqueta.
// Para salir de ciclos anidados se puede usar una bandera,
// encapsular la logica en una funcion y usar return, etc.
\end{lstlisting}




\subsection{Conversiones, Parseos y Casteos}

\begin{lstlisting}[style=cpp]
// string -> numero (parseo)
// Convierte una cadena de texto en el tipo numerico indicado.
int e = stoi("42");
// Entrada: "42"  ->  Salida: 42

long long l = stoll("10000000000");
// Entrada: "10000000000"  ->  Salida: 10000000000

double d = stod("3.14");
// Entrada: "3.14"  ->  Salida: 3.14

// Tambien se puede indicar la base numerica.
int bin = stoi("1010", nullptr, 2);
// "1010" en base 2 -> 10

int hex = stoi("ff", nullptr, 16);
// "ff" en base 16 -> 255


// numero -> string
// Convierte cualquier numero compatible en texto.
string s = to_string(42);
// Entrada: 42  ->  Salida: "42"


// Casteo: fuerza la conversion de un tipo a otro.
// Al pasar de decimal a entero, se TRUNCA hacia 0.
int t = (int) 3.99;
// 3.99 -> 3

int r = static_cast<int>(3.99);
// 3.99 -> 3
// static_cast<int> es la forma recomendada en C++


// ! int / int = division ENTERA
double mal = 7 / 2;
// 7 y 2 son int -> 3
// Luego 3 se convierte a double -> 3.0

double bien = 7 / 2.0;
// 2.0 es double -> division decimal
// 7 / 2.0 -> 3.5


// caracteres
// Los char tienen un valor numerico asociado (ASCII).
int cod = (int) 'A';
// 'A' -> 65

char ch = (char) 66;
// 66 -> 'B'

// Truco muy usado para convertir un digito char a int.
int dig = '7' - '0';
// '7' tiene valor 55 y '0' tiene valor 48
// 55 - 48 -> 7
\end{lstlisting}




%---------------------------------------------------------------------------
\part{Estructuras de Datos}
%---------------------------------------------------------------------------

\section{Basic Structures}

\begin{lstlisting}[style=cpp]
int a[5] = {1,2,3,4,5}; // crear con valores: O(n)

int b[5];               // crear de tamano 5
                        // posiciones: b[0] ... b[4]
                        // ! Valores no inicializados

b[0] = 10;              // INSERT: O(1)
b[1] = 20;              // INSERT: O(1)

int x = b[0];           // READ: O(1)

b[0] = 30;              // UPDATE: O(1)

for (int i = 0; i < 5; i++)
    cout << b[i] << " ";// READ: recorrer -> O(n)

// DELETE: no existe metodo para eliminar
// Se puede sobrescribir o manejar logicamente
b[0] = 0;               // DELETE logico: O(1)

// ! El tamano de un array estatico no cambia.
// Para mas posiciones se necesita otro arreglo.
\end{lstlisting}

\subsection{Matrices (Arreglos 2D)}

Arreglo bidimensional (filas $\times$ columnas). \textbf{Complejidad:}
acceso \texttt{[i][j]} $O(1)$; crear $n\times m$ en $O(nm)$.

\begin{lstlisting}[style=cpp]
vector<vector<int>> m(3, vector<int>(4,0)); // CREATE: O(n*m)
// matriz 3x4, todas las posiciones inicializadas en 0

m[1][2] = 7;              // INSERT: O(1)

int x = m[1][2];           // READ: O(1)

m[1][2] = 10;              // UPDATE: O(1)

for (int i = 0; i < 3; i++)
    for (int j = 0; j < 4; j++)
        cout << m[i][j] << " ";
// READ: recorrer toda la matriz -> O(n*m)

// DELETE: no existe metodo para eliminar una posicion.
// Se puede sobrescribir el valor o manejarlo logicamente.
m[1][2] = 0;               // DELETE logico: O(1)

// ! El tamano de la matriz puede cambiar con vector,
// pero un arreglo 2D estatico no puede cambiar de tamano.
\end{lstlisting}



\subsection{Vectores / Listas dinámicas}

Arreglo que \textbf{crece} solo (ArrayList / \texttt{vector} /
\texttt{list}). \textbf{Complejidad:} acceso por índice $O(1)$; agregar
al final $O(1)$ amortizado; insertar/borrar en medio $O(n)$; buscar
$O(n)$.

\begin{lstlisting}[style=cpp]
vector<int> v;              // CREATE: lista vacia

v.push_back(3);             // INSERT: final -> O(1) amortizado
v.push_back(7);             // INSERT: final -> O(1) amortizado

int x = v[0];               // READ: acceso -> O(1)
int n = v.size();            // READ: tamano -> O(1)

v[0] = 10;                  // UPDATE: modificar -> O(1)

v.insert(v.begin() + 1, 5);  // INSERT: medio -> O(n)

v.erase(v.begin() + 1);      // DELETE: medio -> O(n)

// Buscar un elemento recorriendo el vector
for (int x : v)
    if (x == 10) { /* ... */ }
// SEARCH: recorrer -> O(n)

// DELETE: eliminar el ultimo elemento
v.pop_back();               // DELETE: final -> O(1)

// ! vector mantiene acceso por indice O(1).
// ! Al crecer puede realocar su memoria internamente.
\end{lstlisting}



\subsection{Listas enlazadas (LinkedList)}

Nodos enlazados. \textbf{Complejidad:} insertar/eliminar en los
\textbf{extremos} $O(1)$; acceso por índice $O(n)$; buscar $O(n)$. Útil
como cola o deque.

\begin{lstlisting}[style=cpp]
list<int> l;                // CREATE: lista vacia

l.push_front(1);            // INSERT: inicio -> O(1)
l.push_back(9);             // INSERT: final -> O(1)

int x = l.front();          // READ: primer elemento -> O(1)
int y = l.back();           // READ: ultimo elemento -> O(1)

l.front() = 5;              // UPDATE: primer elemento -> O(1)
l.back() = 10;              // UPDATE: ultimo elemento -> O(1)

l.pop_front();              // DELETE: inicio -> O(1)
l.pop_back();               // DELETE: final -> O(1)

// ! No existe acceso directo por indice.
// Avanzar hasta una posicion requiere recorrer nodos.
auto it = l.begin();
advance(it, 2);             // acceso al tercer elemento -> O(n)

// Buscar un elemento requiere recorrer la lista.
for (int x : l)
    if (x == 10) { /* ... */ }
// SEARCH: recorrer -> O(n)

// ! list no permite l[i].
// Para acceso por indice, vector suele ser mejor.
\end{lstlisting}

\subsection{Pilas (Stack)}

Estructura \textbf{LIFO} (último en entrar, primero en salir).
\textbf{Complejidad:} insertar, eliminar y consultar el tope $O(1)$. Para DFS
iterativo, paréntesis balanceados y deshacer.

\begin{lstlisting}[style=cpp]
stack<int> pila;            // CREATE: pila vacia

pila.push(3);               // INSERT: O(1)
pila.push(7);               // INSERT: O(1)

int tope = pila.top();      // READ: consultar tope -> O(1)

pila.top() = 10;            // UPDATE: modificar tope -> O(1)

pila.pop();                 // DELETE: eliminar tope -> O(1)

// ! Solo se puede acceder directamente al elemento superior.
// ! No existe acceso por indice.
\end{lstlisting}


\subsection{Colas (Queue)}

Estructura \textbf{FIFO} (primero en entrar, primero en salir).
\textbf{Complejidad:} encolar y desencolar $O(1)$. Base del BFS.

\begin{lstlisting}[style=cpp]
queue<int> cola;             // CREATE: cola vacia

cola.push(3);                // INSERT: final -> O(1)
cola.push(7);                // INSERT: final -> O(1)

int frente = cola.front();   // READ: primero -> O(1)
int atras = cola.back();     // READ: ultimo -> O(1)

cola.front() = 10;           // UPDATE: modificar primero -> O(1)

cola.pop();                  // DELETE: primero -> O(1)

// ! push agrega al final y pop elimina del frente.
// ! No existe acceso directo por indice.
\end{lstlisting}


\subsection{Colas de prioridad (Priority Queue / Heap)}

Devuelve siempre el elemento de mayor (o menor) prioridad.
\textbf{Complejidad:} insertar y extraer $O(\log n)$; consultar el tope
$O(1)$. Para Dijkstra, Prim, k-ésimo mayor.

\begin{lstlisting}[style=cpp]
priority_queue<int> pq;      // CREATE: MAX-heap por defecto

pq.push(3);                  // INSERT: O(log n)
pq.push(7);                  // INSERT: O(log n)

int mx = pq.top();           // READ: mayor -> O(1)

pq.pop();                    // DELETE: mayor -> O(log n)

// min-heap: el menor queda en el tope
priority_queue<int, vector<int>, greater<int>> pqMin;

pqMin.push(3);               // INSERT: O(log n)
int mn = pqMin.top();        // READ: menor -> O(1)
pqMin.pop();                 // DELETE: O(log n)

// ! No existe acceso directo a elementos intermedios.
// ! priority_queue no tiene una operacion UPDATE directa.
\end{lstlisting}


\subsection{Comparadores en C++ (Comparators)}

\textbf{Idea:} un comparador es una función (lambda, functor, o
\texttt{operator<} sobrecargado) que le dice a \texttt{sort},
\texttt{set}, \texttt{map}, \texttt{priority\_queue}, etc.
\textbf{cómo ordenar} dos elementos. Debe cumplir la propiedad de
\emph{orden estricto débil} (\textit{strict weak ordering}):
devolver \texttt{true} si el primer argumento va
\textbf{estrictamente antes} que el segundo, y \texttt{false} en
cualquier otro caso —incluyendo cuando se consideran ``iguales'' para
efectos de orden. El error más común es usar \texttt{<=} en vez de
\texttt{<}: eso rompe la propiedad y produce comportamiento
indefinido (desde resultados mal ordenados hasta crashes).

\textbf{¿Por qué usarlo?} El orden por defecto casi nunca alcanza:
necesitas ordenar por varios criterios a la vez, en orden
descendente, por una clave derivada (no el valor mismo), o con un
\texttt{struct}/\texttt{pair} propio que no debe usar el orden
lexicográfico por defecto.

\textbf{Casos de uso:}

\begin{itemize}
    \item Ordenar por múltiples campos (ej. por costo, y en empate
    por índice original).
    \item \texttt{priority\_queue} como min-heap en vez de max-heap
    (o con un orden custom sobre un \texttt{struct}).
    \item \texttt{set}/\texttt{map} con orden propio sobre un tipo
    definido por ti.
    \item Ordenar \textbf{índices} según los valores de otro
    arreglo, sin mover el arreglo original.
\end{itemize}

\textbf{Complejidad:} el comparador en sí no cambia la complejidad
del algoritmo que lo usa (\texttt{sort} sigue siendo
$O(n\log n)$); pero un comparador que no sea un orden estricto débil
puede producir UB o bucles infinitos.


\subsubsection{Comparador con lambda para \texttt{sort} (multi-criterio, índices)}

\begin{lstlisting}[style=cpp]
struct Persona {
    string nombre;
    int edad;
};

vector<Persona> personas = {
    {"Ana", 30}, {"Luis", 25}, {"Eva", 25}
};

// ordena por edad ascendente; en empate, por nombre alfabetico
sort(personas.begin(), personas.end(),
     [](const Persona& a, const Persona& b) {
         if (a.edad != b.edad)
             return a.edad < b.edad;
         return a.nombre < b.nombre;
     });
// resultado: Eva(25), Luis(25), Ana(30)

// ordenar en orden DESCENDENTE: solo invierte la comparacion
sort(personas.begin(), personas.end(),
     [](const Persona& a, const Persona& b) {
         return a.edad > b.edad;
     });
// resultado: Ana(30), luego Luis/Eva(25) en su orden relativo previo
\end{lstlisting}

\textbf{Ordenar índices en vez del arreglo (muy común):}

\begin{lstlisting}[style=cpp]
vector<int> valores = {40, 10, 30, 20};

vector<int> idx(valores.size());
iota(idx.begin(), idx.end(), 0);   // idx = [0, 1, 2, 3]

sort(idx.begin(), idx.end(),
     [&](int a, int b) { return valores[a] < valores[b]; });
// idx queda [1, 3, 2, 0]
// (valores[1]=10, valores[3]=20, valores[2]=30, valores[0]=40)

for (int i : idx)
    cout << valores[i] << " ";
// imprime: 10 20 30 40
// sin haber movido "valores" del lugar
\end{lstlisting}


\subsubsection{\texttt{operator<} dentro de un \texttt{struct} (orden por defecto)}

\textbf{Idea:} si sobrecargas \texttt{operator<} dentro de tu propio
tipo, ese es el orden que usan automáticamente \texttt{sort},
\texttt{set}, \texttt{map} (como clave) y \texttt{priority\_queue}
\textbf{sin pasar ningún comparador extra}.

\begin{lstlisting}[style=cpp]
struct Punto {
    int x, y;

    bool operator<(const Punto& otro) const {
        if (x != otro.x) return x < otro.x;
        return y < otro.y;
    }
};

set<Punto> s;
s.insert({3, 4});
s.insert({1, 2});
s.insert({1, 5});

for (auto& p : s)
    cout << "(" << p.x << "," << p.y << ") ";
// imprime en orden: (1,2) (1,5) (3,4)
// -> set los ordeno automaticamente usando operator<
\end{lstlisting}


\subsubsection{\texttt{priority\_queue} con comparador custom}

\textbf{Idea:} \texttt{priority\_queue} es un \textbf{max-heap} por
defecto (el elemento ``mayor'' según el comparador queda arriba). Para
un min-heap con tipos simples, se usa \texttt{greater<T>}. Para un
\texttt{struct} propio con una regla de prioridad distinta a
\texttt{operator<}, se pasa un \textbf{functor} (struct con
\texttt{operator()}).

\begin{lstlisting}[style=cpp]
// min-heap de enteros: el menor queda arriba
priority_queue<int, vector<int>, greater<int>> minHeap;
minHeap.push(5);
minHeap.push(1);
minHeap.push(3);

cout << minHeap.top();
// 1  (el menor, no el mayor)
\end{lstlisting}

\begin{lstlisting}[style=cpp]
struct Persona {
    string nombre;
    int edad;
};

// functor: queda arriba la persona con MENOR edad
// (priority_queue es max-heap, asi que se invierte el sentido:
// operator() debe devolver true si "a" tiene MENOR prioridad que "b")
struct CompararPorEdad {
    bool operator()(const Persona& a, const Persona& b) const {
        return a.edad > b.edad;
    }
};

priority_queue<Persona, vector<Persona>, CompararPorEdad> pq;
pq.push({"Ana", 30});
pq.push({"Luis", 20});
pq.push({"Eva", 25});

cout << pq.top().nombre;
// "Luis"  (la persona mas joven queda en la cima)
\end{lstlisting}


\subsubsection{Reglas de orden estricto débil (\textit{strict weak ordering})}

Todo comparador \texttt{cmp(a, b)} debe cumplir:

\begin{itemize}
    \item \textbf{Irreflexividad:} \texttt{cmp(a, a)} siempre
    \texttt{false}. Un elemento nunca va estrictamente antes que sí
    mismo.
    \item \textbf{Asimetría:} si \texttt{cmp(a, b)} es
    \texttt{true}, entonces \texttt{cmp(b, a)} debe ser
    \texttt{false}.
    \item \textbf{Transitividad:} si \texttt{cmp(a, b)} y
    \texttt{cmp(b, c)} son \texttt{true}, entonces \texttt{cmp(a, c)}
    también debe serlo.
    \item \textbf{Consistencia en los empates:} si ni
    \texttt{cmp(a,b)} ni \texttt{cmp(b,a)} son \texttt{true}, se
    consideran ``equivalentes'' para efectos de orden (aunque no
    sean idénticos como objetos).
\end{itemize}

\textbf{El error más común:}

\begin{lstlisting}[style=cpp]
// INCORRECTO: usa <= en vez de <
sort(v.begin(), v.end(), [](int a, int b) {
    return a <= b;     // viola irreflexividad: cmp(a,a) da true
});
// esto es Undefined Behavior: en el mejor caso ordena mal,
// en el peor caso crashea o hace un bucle infinito.

// CORRECTO:
sort(v.begin(), v.end(), [](int a, int b) {
    return a < b;
});
\end{lstlisting}

\textbf{Otro error común:} comparadores que no son transitivos al
combinar varios criterios sin desempate consistente (ej. comparar
solo por \texttt{a.x < b.x} cuando a veces \texttt{a.x == b.x} pero
querías tratar esos casos con otro criterio). Si dos elementos
``empatan'' en el primer criterio, siempre hay que \textbf{decidir}
con un segundo criterio, o dejarlos expresamente como equivalentes.

\textbf{Regla práctica:} si tu comparador tiene la forma
\texttt{return a.campo1 < b.campo1;} y sospechas que puede haber
empates que te importan, agrega el siguiente criterio explícitamente
(\texttt{if (a.campo1 != b.campo1) return ...; return
a.campo2 < b.campo2;}) en vez de dejar el desempate al azar del
algoritmo de ordenamiento (que no es estable a menos que uses
\texttt{stable\_sort}).


\subsection{Vector circular (Circular Buffer)}

\textbf{Idea:} un arreglo de tamaño fijo (\texttt{cap}) que se comporta
como una \textbf{cola doble}. En lugar de mover elementos cuando se
inserta o elimina, se reutilizan las posiciones del arreglo de forma
circular.

Se mantienen tres variables principales:

\begin{itemize}
    \item \texttt{cap}: capacidad total del arreglo.
    \item \texttt{ini}: posición física donde comienza el primer elemento.
    \item \texttt{tam}: cantidad actual de elementos almacenados.
\end{itemize}

La posición lógica $i$ no coincide necesariamente con la posición física
del arreglo. Se obtiene mediante:

\[
\texttt{posicionFisica(i) = (ini + i) \% cap}
\]

Por ejemplo, con $\texttt{cap}=5$ y $\texttt{ini}=3$:

\[
\begin{array}{c|ccccc}
\text{Índice lógico} & 0 & 1 & 2 & 3 & 4 \\ \hline
\text{Índice físico} & 3 & 4 & 0 & 1 & 2
\end{array}
\]

Cuando se llega al final del arreglo, el módulo hace que la siguiente
posición vuelva al comienzo. Por eso \textbf{no es necesario mover los
elementos}.

\textbf{¿Por qué usarlo?} En un arreglo normal, insertar al frente
requiere desplazar todos los elementos y cuesta $O(n)$. En un vector
circular solo se mueve \texttt{ini}, por lo que cuesta $O(1)$.

\textbf{Casos de uso:}
\begin{itemize}
    \item Colas de tamaño fijo.
    \item Buffers de audio, video o red.
    \item Ventanas deslizantes.
    \item Sistemas donde llegan y salen datos continuamente.
    \item Simulaciones de procesos en una cola.
    \item Mantener únicamente los últimos $k$ elementos.
\end{itemize}

Si el buffer está lleno y se ejecuta \texttt{pushBack}, el nuevo elemento
\textbf{reemplaza al más antiguo}. Esto permite mantener automáticamente
una ventana de tamaño fijo.

\textbf{Complejidad:} insertar, eliminar, consultar extremos, acceder
por posición y comprobar el tamaño son $O(1)$.


\begin{lstlisting}[style=cpp]
struct VectorCircular {
    vector<int> buf;
    int cap, ini = 0, tam = 0;

    VectorCircular(int c) : buf(c), cap(c) {}

    // CREATE: crea un buffer con capacidad fija.
    // Los elementos inicialmente no estan ocupados.
    
    bool empty() {
        return tam == 0;
    }                           // O(1)

    bool full() {
        return tam == cap;
    }                           // O(1)

    int size() {
        return tam;
    }                           // O(1)

    int capacity() {
        return cap;
    }                           // O(1)

    // Convierte una posicion logica en fisica.
    int pos(int i) {
        return (ini + i) % cap;
    }                           // O(1)

    // INSERT: agrega al final.
    void pushBack(int x) {
        buf[(ini + tam) % cap] = x;

        if (tam < cap)
            tam++;
        else
            ini = (ini + 1) % cap;
    }                           // O(1)

    // INSERT: agrega al inicio.
    void pushFront(int x) {
        ini = (ini - 1 + cap) % cap;
        buf[ini] = x;

        if (tam < cap)
            tam++;
    }                           // O(1)

    // READ: consulta el primer elemento.
    int front() {
        return buf[ini];
    }                           // O(1)

    // READ: consulta el ultimo elemento.
    int back() {
        return buf[(ini + tam - 1) % cap];
    }                           // O(1)

    // READ: acceso por posicion logica.
    int at(int i) {
        return buf[(ini + i) % cap];
    }                           // O(1)

    // UPDATE: modifica un elemento.
    void set(int i, int x) {
        buf[(ini + i) % cap] = x;
    }                           // O(1)

    // DELETE: elimina el primer elemento.
    int popFront() {
        int v = buf[ini];
        ini = (ini + 1) % cap;
        tam--;
        return v;
    }                           // O(1)

    // DELETE: elimina el ultimo elemento.
    int popBack() {
        tam--;
        return buf[(ini + tam) % cap];
    }                           // O(1)

    // DELETE: elimina todos los elementos.
    void clear() {
        ini = 0;
        tam = 0;
    }                           // O(1)

    // SEARCH: busca un elemento.
    // A diferencia de las operaciones anteriores,
    // hay que recorrer los elementos.
    bool contains(int x) {
        for (int i = 0; i < tam; i++)
            if (at(i) == x)
                return true;

        return false;
    }                           // O(n)

    // READ: recorrer los elementos en orden logico.
    void print() {
        for (int i = 0; i < tam; i++)
            cout << at(i) << " ";
    }                           // O(n)
};
\end{lstlisting}

\textbf{Ejemplo de funcionamiento:}

\begin{lstlisting}[style=cpp]
VectorCircular v(5);

v.pushBack(10);
v.pushBack(20);
v.pushBack(30);
// [10, 20, 30, _, _]
// ini = 0, tam = 3

v.pushFront(5);
// [5, 10, 20, 30, _]
// ini cambia de 0 a 4

cout << v.front();
// 5

cout << v.back();
// 30

v.popFront();
// elimina 5
// [_, 10, 20, 30, _]

v.pushBack(40);
v.pushBack(50);
// [_, 10, 20, 30, 40, 50]
// El almacenamiento puede haber dado la vuelta.

v.pushBack(60);
// El buffer estaba lleno.
// Se agrega 60 y se elimina el mas antiguo.
// Resultado logico: [20, 30, 40, 50, 60]
\end{lstlisting}

\textbf{Ventaja principal:} el arreglo físico puede verse ``desordenado'',
pero los elementos siempre tienen un orden lógico determinado por
\texttt{ini}. Por ejemplo:

\begin{lstlisting}[style=cpp]
// Memoria fisica:
// [40, 50, 20, 30,  _]
//       ^
//      ini = 2

// Orden logico:
// 20 -> 30 -> 40 -> 50
\end{lstlisting}

No importa dónde estén físicamente los elementos: \texttt{at(0)} siempre
representa el primero, \texttt{at(1)} el segundo, y así sucesivamente.

\textbf{Diferencia con un vector normal:}

\begin{center}
\small
\begin{tabular}{|l|c|c|}
\hline
\textbf{Operación} & \texttt{vector} & \textbf{Circular} \\
\hline
Agregar al final & $O(1)$* & $O(1)$ \\
\hline
Eliminar al final & $O(1)$ & $O(1)$ \\
\hline
Agregar al frente & $O(n)$ & $O(1)$ \\
\hline
Eliminar al frente & $O(n)$ & $O(1)$ \\
\hline
Acceso por índice & $O(1)$ & $O(1)$ \\
\hline
Buscar & $O(n)$ & $O(n)$ \\
\hline
\end{tabular}
\end{center}

\noindent
*Amortizado.

\textbf{Regla práctica:} si necesitas una estructura de tamaño fijo
donde los elementos entran y salen continuamente por los extremos,
un vector circular evita los desplazamientos de un arreglo normal y
mantiene las operaciones principales en $O(1)$.

\section{Estructuras Hash/Tree}


\subsection{Conjuntos hash (HashSet)}

Colección de elementos \textbf{sin duplicados} y sin orden garantizado.
En C\texttt{++} se implementa mediante \texttt{unordered\_set}. Utiliza
una tabla hash, por lo que las operaciones son $O(1)$ en promedio.
Es útil cuando solo interesa saber si un elemento existe, sin necesidad
de mantener los elementos ordenados.

\textbf{Complejidad:} insertar, buscar y borrar $O(1)$ promedio
($O(n)$ peor caso).

\begin{lstlisting}[style=cpp]
unordered_set<int> s;          // CREATE: conjunto vacio

s.insert(3);                   // INSERT: O(1) prom.
s.insert(7);                   // INSERT: O(1) prom.
s.insert(3);                   // no se agrega: ya existe

bool hay = s.count(3);         // SEARCH: O(1) prom.
// hay = 1 si existe, 0 si no existe

auto it = s.find(7);            // SEARCH: O(1) prom.
// devuelve iterador al elemento o s.end()

s.erase(3);                    // DELETE: O(1) prom.

int n = s.size();              // READ: cantidad -> O(1)
bool vacio = s.empty();        // READ: ¿esta vacio? -> O(1)

// Recorrer todos los elementos
for (int x : s)
    cout << x << " ";
// O(n), pero el orden NO esta garantizado.

// ! No existe acceso por indice: s[0] no es valido.
// ! Los elementos repetidos se ignoran.
// ! El orden puede cambiar entre ejecuciones.
\end{lstlisting}

\textbf{Cuándo usarlo:}
\begin{itemize}
    \item Saber rápidamente si un elemento ya apareció.
    \item Eliminar duplicados.
    \item Mantener un conjunto de elementos visitados.
    \item Detectar elementos repetidos en $O(n)$ promedio.
    \item BFS/DFS: guardar vértices visitados cuando sea conveniente.
\end{itemize}


\subsection{Mapas / Diccionarios hash}

Estructura que almacena pares \textbf{clave $\rightarrow$ valor}.
En C\texttt{++} se implementa mediante \texttt{unordered\_map}.
Cada clave identifica un valor y \textbf{no puede repetirse}. Si se
asigna un nuevo valor a una clave existente, se reemplaza el anterior.

\textbf{Complejidad:} insertar, acceder, buscar y borrar por clave
$O(1)$ promedio ($O(n)$ peor caso).

\begin{lstlisting}[style=cpp]
unordered_map<string,int> m;   // CREATE: mapa vacio

m["a"] = 1;                    // INSERT: O(1) prom.
m["b"] = 2;                    // INSERT: O(1) prom.

int v = m["a"];                // READ: acceder -> O(1) prom.
// v = 1

m["a"] = 10;                   // UPDATE: O(1) prom.
// la clave "a" ahora tiene valor 10

m.erase("a");                  // DELETE: O(1) prom.

bool existe = m.count("b");    // SEARCH: O(1) prom.
// existe = 1

auto it = m.find("b");         // SEARCH: O(1) prom.
if (it != m.end())
    cout << it->second;        // valor asociado a "b"

// Recorrer todas las parejas
for (auto [clave, valor] : m)
    cout << clave << " " << valor;
// O(n), sin orden garantizado.
\end{lstlisting}

\textbf{Importante:} usar \texttt{m[clave]} para consultar una clave
que no existe puede \textbf{crear esa clave} con su valor por defecto.
Para comprobar existencia sin insertar, es preferible usar
\texttt{count()} o \texttt{find()}.

\begin{lstlisting}[style=cpp]
unordered_map<string,int> m;

int x = m["a"];
// Si "a" no existia, se crea:
// "a" -> 0
// ! operator[] puede INSERTAR una clave al consultar.

// count(): comprueba si existe SIN insertar
bool existe = m.count("b");
// Devuelve 0 si no existe, 1 si existe.

// find(): busca la clave SIN insertar
auto it = m.find("b");

if (it != m.end())
    cout << it->second;       // valor asociado a "b"
else
    cout << "No existe";

// ! count() solo indica si existe.
// ! find() permite acceder a la pareja clave -> valor.
\end{lstlisting}


\textbf{Uso clásico: contar frecuencias.}

\begin{lstlisting}[style=cpp]
unordered_map<int,int> freq;

for (int x : datos)
    freq[x]++;                 // frecuencia de cada elemento

cout << freq[5];
// cantidad de veces que aparece 5
\end{lstlisting}

\textbf{Cuándo usarlo:}
\begin{itemize}
    \item Contar frecuencias.
    \item Asociar identificadores con información.
    \item Buscar rápidamente un valor a partir de una clave.
    \item Problemas de conteo y frecuencia.
    \item Memorizar resultados de estados o funciones.
\end{itemize}


\subsection{Conjuntos ordenados (TreeSet)}

Conjunto de elementos \textbf{sin duplicados y ordenados}.
En C\texttt{++} se implementa mediante \texttt{set}, normalmente como
un árbol balanceado.

A diferencia de \texttt{unordered\_set}, los elementos siempre se
mantienen ordenados según el comparador utilizado.

\textbf{Complejidad:} insertar, buscar y borrar $O(\log n)$.
Consultar el menor o mayor elemento mediante \texttt{begin()} o
\texttt{rbegin()} es $O(1)$.

\begin{lstlisting}[style=cpp]
set<int> s;                    // CREATE: conjunto vacio

s.insert(5);                   // INSERT: O(log n)
s.insert(1);                   // INSERT: O(log n)
s.insert(3);                   // INSERT: O(log n)
s.insert(3);                   // no se agrega: ya existe

int menor = *s.begin();        // READ: minimo -> O(1)
int mayor = *s.rbegin();       // READ: maximo -> O(1)

auto it = s.find(3);           // SEARCH: O(log n)
if (it != s.end())
    cout << *it;

s.erase(3);                    // DELETE: O(log n)

int n = s.size();              // READ: cantidad -> O(1)

// Recorrer en orden ascendente
for (int x : s)
    cout << x << " ";
// O(n)
\end{lstlisting}

\textbf{lower\_bound y upper\_bound:}

\begin{lstlisting}[style=cpp]
set<int> s = {1, 3, 5, 7, 9};

auto it1 = s.lower_bound(5);
// primer elemento >= 5
// resultado: 5

auto it2 = s.upper_bound(5);
// primer elemento > 5
// resultado: 7
\end{lstlisting}

Esto permite resolver consultas de \textbf{piso y techo} rápidamente:

\begin{lstlisting}[style=cpp]
// FLOOR: mayor elemento <= x
auto it = s.upper_bound(x);

if (it != s.begin()) {
    --it;
    cout << *it;
}

// CEILING: menor elemento >= x
auto it2 = s.lower_bound(x);

if (it2 != s.end())
    cout << *it2;
\end{lstlisting}

\textbf{Cuándo usarlo:}
\begin{itemize}
    \item Cuando se necesitan elementos únicos y ordenados.
    \item Obtener mínimo o máximo dinámicamente.
    \item Encontrar el siguiente elemento mayor o igual.
    \item Encontrar el anterior elemento menor o igual.
    \item Mantener un conjunto mientras se insertan y eliminan elementos.
\end{itemize}

\textbf{Diferencia clave:} si solo necesitas saber si existe un elemento,
\texttt{unordered\_set} suele ser más rápido. Si necesitas consultar
orden, \texttt{lower\_bound}, \texttt{upper\_bound}, mínimo o máximo,
\texttt{set} es la opción adecuada.


\subsection{Mapas ordenados (TreeMap)}

Mapa \textbf{clave $\rightarrow$ valor} cuyas claves se mantienen
ordenadas. En C\texttt{++} se implementa mediante \texttt{map}, usando
un árbol balanceado.

Cada clave aparece una sola vez. Los valores pueden repetirse.

\textbf{Complejidad:} insertar, buscar, acceder y borrar por clave
$O(\log n)$. \texttt{begin()} permite obtener la menor clave en
$O(1)$.

\begin{lstlisting}[style=cpp]
map<string,int> m;             // CREATE: mapa vacio

m["b"] = 2;                    // INSERT: O(log n)
m["a"] = 1;                    // INSERT: O(log n)
m["c"] = 3;                    // INSERT: O(log n)

int v = m["a"];                // READ: O(log n)
// v = 1

m["a"] = 10;                   // UPDATE: O(log n)

m.erase("b");                  // DELETE: O(log n)

auto it = m.find("c");         // SEARCH: O(log n)
if (it != m.end())
    cout << it->second;

// Primera pareja: menor clave
auto primera = m.begin();      // O(1)
cout << primera->first;

// Ultima pareja: mayor clave
auto ultima = prev(m.end());   // O(1) amortizado
cout << ultima->first;
\end{lstlisting}

\textbf{Las claves se recorren ordenadas:}

\begin{lstlisting}[style=cpp]
map<int,string> m;

m[30] = "C";
m[10] = "A";
m[20] = "B";

for (auto [clave, valor] : m)
    cout << clave << " " << valor;

// Salida:
// 10 A
// 20 B
// 30 C
\end{lstlisting}

\textbf{lower\_bound y upper\_bound} también funcionan sobre las claves:

\begin{lstlisting}[style=cpp]
map<int,string> m;
m[10] = "A";
m[20] = "B";
m[30] = "C";

auto it = m.lower_bound(15);
// primera clave >= 15
// apunta a 20

auto it2 = m.upper_bound(20);
// primera clave > 20
// apunta a 30
\end{lstlisting}

\textbf{Uso clásico: contar frecuencias ordenadas.}

\begin{lstlisting}[style=cpp]
map<int,int> freq;

for (int x : datos)
    freq[x]++;

// Las claves aparecen en orden creciente.
for (auto [x, cantidad] : freq)
    cout << x << ": " << cantidad << "\n";
\end{lstlisting}

\textbf{Cuándo usarlo:}
\begin{itemize}
    \item Cuando se necesitan pares clave--valor ordenados.
    \item Contar frecuencias y procesarlas en orden.
    \item Encontrar claves inmediatamente mayores o menores.
    \item Procesar eventos o valores manteniendo un orden dinámico.
    \item Cuando se necesitan operaciones de árbol como
    \texttt{lower\_bound} y \texttt{upper\_bound}.
\end{itemize}


\subsection*{Comparación rápida}

\begin{center}
\small
\begin{tabular}{|l|c|c|c|}
\hline
\textbf{Estructura} & \textbf{Duplicados} & \textbf{Orden} & \textbf{Buscar} \\
\hline
\texttt{unordered\_set} & No & No & $O(1)$ prom. \\
\hline
\texttt{unordered\_map} & Claves: No & No & $O(1)$ prom. \\
\hline
\texttt{set} & No & Sí & $O(\log n)$ \\
\hline
\texttt{map} & Claves: No & Sí & $O(\log n)$ \\
\hline
\end{tabular}
\end{center}

\textbf{Regla práctica:}

\begin{itemize}[nosep,leftmargin=*]
    \item \texttt{unordered\_set} $\rightarrow$ \textit{¿Ya existe este elemento?}
    \item \texttt{unordered\_map} $\rightarrow$ \textit{¿Qué valor corresponde a esta clave?}
    \item \texttt{set} $\rightarrow$ \textit{Elementos únicos + ordenados.}
    \item \texttt{map} $\rightarrow$ \textit{Claves $\rightarrow$ valores + claves ordenadas.}
\end{itemize}


\section{Trees}

\section{Range Queries}

\subsection{Casos de uso}

\begin{table}[H]
	\raggedright
	\scriptsize
	\setlength{\tabcolsep}{2pt}
	\renewcommand{\arraystretch}{1.1}
	\begin{tabular}{|p{0.43\columnwidth}|p{0.29\columnwidth}|p{0.20\columnwidth}|}
		\hline
		\textbf{Cuando el problema te pida:} &
		\textbf{Entonces usa:} &
		\textbf{Complejidad} \\
		\hline
		
		Suma de rango sobre un arreglo que \textbf{nunca cambia}, con muchas consultas
		(ej. ``suma de ventas entre el día $l$ y $r$'') &
		Prefix Sums &
		Build $O(n)$ \newline Query $O(1)$ \\
		\hline
		
		Suma de rango con actualizaciones \textbf{puntuales} frecuentes
		(ej. ``suma el stock de la posición $i$, luego consulta un rango'') &
		Fenwick Tree (BIT) &
		Build $O(n)$ \newline Query/Update $O(\log n)$ \\
		\hline
		
		Sumar $v$ a \textbf{todo un rango} $[l,r]$ y también consultar la
		\textbf{suma} de cualquier rango &
		FenwickRangeUpdate (doble BIT) &
		Query/Update $O(\log n)$ \\
		\hline
		
		Range update + range query con operaciones más complejas:
		min/max en rango, asignar un valor a todo un rango, contar ceros, etc. &
		Segment Tree + Lazy Propagation &
		Query/Update $O(\log n)$ \\
		\hline
		
		Mínimo, máximo o gcd de cualquier rango sobre un arreglo
		\textbf{estático} (ej. RMQ clásico, LCA) &
		Sparse Table &
		Build $O(n\log n)$ \\ Query $O(1)$ \\
		\hline
		
		Mínimo o máximo de rango cuando el arreglo \textbf{sí cambia}
		con el tiempo &
		Segment Tree &
		Query/Update $O(\log n)$ \\
		\hline
		
		Consultas 2D: cuántos puntos o cuánta suma hay dentro de un
		rectángulo $[x_1,x_2]\times[y_1,y_2]$ &
		Fenwick 2D &
		Query/Update \newline $O(\log n\log m)$ \\
		\hline
		
		Encontrar el $k$-ésimo elemento, o ``menor índice cuya suma acumulada
		alcanza $X$'' (frecuencias) &
		Fenwick con \texttt{lowerBound} &
		$O(\log n)$ \\
		\hline
		
		Muchas queries de rango \textbf{offline} (las conoces todas de antemano)
		sobre un arreglo estático o con pocos updates &
		Mo's Algorithm &
		$O((n+q)\sqrt{n})$ \\
		\hline
		
		LCA entre dos nodos, o mínimo en rango sobre un recorrido Euler de un
		árbol (valores consecutivos difieren en $\pm1$) &
		RMQ $\pm1$ \newline (Euler Tour + Sparse Table) &
		Build $O(n)$ \newline Query $O(1)$ \\
		\hline
		
		Cuántos elementos en el rango $[l,r]$ están dentro de un rango de
		\textbf{valores} $[a,b]$, o cuál es el $k$-ésimo menor en $[l,r]$ &
		Wavelet Tree &
		Query $O(\log n)$ \\
		\hline
		
		Queries o updates sobre \textbf{caminos o subárboles} de un árbol
		(ej. ``suma el camino de $u$ a $v$'') &
		Heavy-Light Decomposition \newline (+ Segment/Fenwick) &
		Query/Update \newline $O(\log^2 n)$ \\
		\hline
		
		Necesitas algo simple de implementar bajo presión, con balance
		razonable entre update y query, sin memorizar mucha teoría &
		Sqrt Decomposition &
		Query/Update $O(\sqrt{n})$ \\
		\hline
		
	\end{tabular}
	\caption{Guía rápida: qué estructura de range queries usar según el problema, con su complejidad.}
\end{table}

\subsection{Prefix Sums}

\textbf{Idea:} si un arreglo \textbf{no cambia}, puedes precalcular
\texttt{pref[i]} $=$ suma de los primeros $i$ elementos en $O(n)$, una
sola vez. Cualquier suma de rango $[l,r]$ se responde restando dos
prefijos: \texttt{pref[r+1] - pref[l]}. Es la base conceptual de todo
lo demás en esta sección: Fenwick y Segment Tree existen justamente
para soportar lo que Prefix Sums no puede (actualizaciones), a costa
de más complejidad de código.

\textbf{¿Por qué usarlo?} Es lo más simple y rápido posible para
consultas de rango cuando el arreglo es estático. Cero estructuras de
datos elaboradas, cero recursión, solo un arreglo auxiliar y una
resta.

\textbf{Casos de uso:}

\begin{itemize}
    \item Suma de rango sobre un arreglo que nunca se modifica, con
    muchas queries.
    \item Suma de submatriz (versión 2D) sobre una matriz estática.
    \item Aplicar $B$ actualizaciones de rango \textbf{offline} y
    consultar el arreglo final una sola vez (Difference Array).
    \item Como paso final de un precómputo de teoría de números
    (criba de $\phi$, $\mu$, número de divisores, etc.) cuando el
    problema pide sumas sobre esos valores en un rango.
\end{itemize}

\textbf{Complejidad:} construir $O(n)$ (u $O(nm)$ en 2D), consulta
$O(1)$. No soporta actualizaciones puntuales eficientes: modificar un
elemento y volver a consultar obliga a reconstruir, $O(n)$.


\subsubsection{Prefix Sum 1D}

\begin{lstlisting}[style=cpp]
struct PrefixSum1D {
    vector<long long> pref;
    int n;

    // CREATE: construye los prefijos a partir de a[0..n-1].
    PrefixSum1D(vector<long long>& a) {
        n = a.size();
        pref.assign(n + 1, 0);

        for (int i = 0; i < n; i++)
            pref[i + 1] = pref[i] + a[i];
    }                                          // O(n)

    // READ: suma del rango [l, r] (0-indexado, inclusivo).
    long long rango(int l, int r) {
        if (l > r) return 0;
        return pref[r + 1] - pref[l];
    }                                          // O(1)
};
\end{lstlisting}

\textbf{Ejemplo de funcionamiento:}

\begin{lstlisting}[style=cpp]
vector<long long> a = {5, 2, 9, 1, 7};

PrefixSum1D ps(a);
// arreglo logico: [5, 2, 9, 1, 7]  (posiciones 0..4)

cout << ps.rango(0, 2);
// 5 + 2 + 9 = 16

cout << ps.rango(1, 4);
// 2 + 9 + 1 + 7 = 19

cout << ps.rango(3, 3);
// 1
\end{lstlisting}


\subsubsection{Prefix Sum 2D}

\begin{lstlisting}[style=cpp]
struct PrefixSum2D {
    vector<vector<long long>> pref;
    int n, m;

    // CREATE: construye la matriz de prefijos a partir de
    // a[0..n-1][0..m-1] usando inclusion-exclusion.
    PrefixSum2D(vector<vector<long long>>& a) {
        n = a.size();
        m = a[0].size();
        pref.assign(n + 1, vector<long long>(m + 1, 0));

        for (int i = 0; i < n; i++)
            for (int j = 0; j < m; j++)
                pref[i + 1][j + 1] = a[i][j] + pref[i][j + 1]
                                    + pref[i + 1][j] - pref[i][j];
    }                                          // O(n * m)

    // READ: suma del rectangulo [x1..x2] x [y1..y2] (0-indexado,
    // inclusivo), usando inclusion-exclusion sobre los prefijos.
    long long rango(int x1, int y1, int x2, int y2) {
        return pref[x2 + 1][y2 + 1] - pref[x1][y2 + 1]
             - pref[x2 + 1][y1] + pref[x1][y1];
    }                                          // O(1)
};
\end{lstlisting}

\textbf{Ejemplo de funcionamiento:}

\begin{lstlisting}[style=cpp]
vector<vector<long long>> a = {
    {1, 2, 3},
    {4, 5, 6},
    {7, 8, 9}
};

PrefixSum2D ps(a);

cout << ps.rango(0, 0, 1, 1);
// 1+2+4+5 = 12

cout << ps.rango(1, 1, 2, 2);
// 5+6+8+9 = 28
\end{lstlisting}


\subsubsection{Difference Array (updates de rango offline)}

\textbf{Idea:} el truco inverso a Prefix Sums: para sumar $v$ a todo
un rango $[l,r]$, marcas \texttt{diff[l] += v} y \texttt{diff[r+1] -=
v} en $O(1)$. Después de aplicar \textbf{todas} las actualizaciones,
un solo recorrido con prefijos reconstruye el arreglo final. Sirve
solo cuando todas las updates son \textbf{offline} (las conoces antes
de necesitar consultar); si necesitas consultar \emph{entre}
actualizaciones, usa FenwickRangeUpdate en su lugar.

\begin{lstlisting}[style=cpp]
struct DifferenceArray {
    vector<long long> diff;
    int n;

    // CREATE: reserva un arreglo de diferencias de tamano n.
    DifferenceArray(int n) : n(n) {
        diff.assign(n + 1, 0);
    }

    // UPDATE: marca +v en todo el rango [l, r] (0-indexado).
    void rangoUpdate(int l, int r, long long v) {
        diff[l] += v;
        if (r + 1 <= n)
            diff[r + 1] -= v;
    }                                          // O(1)

    // READ: reconstruye el arreglo final tras aplicar todas
    // las actualizaciones marcadas hasta el momento.
    vector<long long> aplicar() {
        vector<long long> a(n);
        long long acc = 0;

        for (int i = 0; i < n; i++) {
            acc += diff[i];
            a[i] = acc;
        }

        return a;
    }                                          // O(n)
};
\end{lstlisting}

\textbf{Ejemplo de funcionamiento:}

\begin{lstlisting}[style=cpp]
DifferenceArray d(5);

d.rangoUpdate(1, 3, 4);
d.rangoUpdate(0, 4, 1);
d.rangoUpdate(2, 2, 10);

vector<long long> resultado = d.aplicar();
// resultado: [1, 5, 15, 5, 1]
\end{lstlisting}


\subsubsection{Aplicación: suma de Euler Totient (\texorpdfstring{$\phi$}{phi}) en un rango}

\textbf{Idea:} $\phi(n)$ cuenta cuántos enteros en $[1,n]$ son
coprimos con $n$. Calcular $\phi$ de \emph{un} número es teoría de
números ($O(\sqrt n)$ factorizando); pero cuando el problema pide la
\textbf{suma de $\phi(i)$ para $i$ en un rango $[l,r]$}, con muchas
queries, el patrón se resuelve en dos fases: (1) precalculas el
arreglo de $\phi$ una sola vez con una criba (teoría de números), y
(2) respondes las queries de rango con Prefix Sums sobre ese arreglo
ya calculado (exactamente esta subsección). Factorizar cada número
del rango por separado en cada query sería $O(q\sqrt n)$ y no escala;
precalcular y sumar con prefijos reduce cada query a $O(1)$.

\paragraph{Criba lineal + Prefix Sums (rango $[1,N]$).}

\begin{lstlisting}[style=cpp]
struct PhiRangeSum {
    vector<long long> phi, pref;
    int n;

    // CREATE: calcula phi[1..n] con criba lineal y sus prefijos.
    // phi[i] = cantidad de enteros en [1, i] coprimos con i.
    PhiRangeSum(int n) : n(n) {
        phi.assign(n + 1, 0);
        pref.assign(n + 1, 0);

        vector<int> primos;
        vector<bool> compuesto(n + 1, false);
        phi[1] = 1;

        for (int i = 2; i <= n; i++) {
            if (!compuesto[i]) {
                primos.push_back(i);
                phi[i] = i - 1;
            }

            for (int p : primos) {
                if ((long long)i * p > n) break;

                compuesto[i * p] = true;

                if (i % p == 0) {
                    phi[i * p] = phi[i] * p;
                    break;
                } else {
                    phi[i * p] = phi[i] * (p - 1);
                }
            }
        }

        for (int i = 1; i <= n; i++)
            pref[i] = pref[i - 1] + phi[i];
    }                                          // O(n)

    // READ: suma de phi(i) para i en [l, r].
    long long rango(int l, int r) {
        if (l > r) return 0;
        return pref[r] - pref[l - 1];
    }                                          // O(1)

    // READ: valor individual de phi(i).
    long long get(int i) {
        return phi[i];
    }                                          // O(1)
};
\end{lstlisting}

\textbf{Ejemplo de funcionamiento:}

\begin{lstlisting}[style=cpp]
PhiRangeSum pr(10);
// phi(1..10) = [1, 1, 2, 2, 4, 2, 6, 4, 6, 4]

cout << pr.get(7);
// phi(7) = 6  (7 es primo)

cout << pr.rango(1, 10);
// 1+1+2+2+4+2+6+4+6+4 = 32

cout << pr.rango(4, 8);
// 2+4+2+6+4 = 18
\end{lstlisting}

\paragraph{Criba segmentada (rango $[L,R]$ grande).}

Si $R$ es muy grande ($10^{12}$) pero $R-L$ es manejable, no puedes
cribar desde 1. En vez de eso, precalculas los primos hasta
$\sqrt R$, y para cada número del segmento $[L,R]$ le quitas sus
factores primos pequeños, aplicando $\phi(n) = n\prod_{p\mid
n}(1-\tfrac1p)$. Si al final queda un residuo $>1$, es un único
primo grande (mayor a $\sqrt R$) que también se aplica a la fórmula.

\begin{lstlisting}[style=cpp]
struct PhiSegmentedRangeSum {
    vector<long long> phi, pref;
    long long L, R;

    // CREATE: calcula phi(i) para cada i en [L, R] factorizando con
    // los primos hasta sqrt(R), y sus prefijos sobre el segmento.
    PhiSegmentedRangeSum(long long L, long long R) : L(L), R(R) {
        int lim = (int)sqrt((double)R) + 1;

        vector<int> primos;
        vector<bool> compuesto(lim + 1, false);
        for (int i = 2; i <= lim; i++) {
            if (!compuesto[i]) primos.push_back(i);
            for (int p : primos) {
                if ((long long)i * p > lim) break;
                compuesto[i * p] = true;
                if (i % p == 0) break;
            }
        }                                      // O(sqrt(R) log log sqrt(R))

        int sz = R - L + 1;
        vector<long long> num(sz);
        phi.assign(sz, 0);

        for (int i = 0; i < sz; i++) {
            num[i] = L + i;
            phi[i] = L + i;
        }

        for (int p : primos) {
            long long start = ((L + p - 1) / p) * p;

            for (long long j = start; j <= R; j += p) {
                int idx = j - L;
                phi[idx] -= phi[idx] / p;

                while (num[idx] % p == 0)
                    num[idx] /= p;
            }
        }

        // lo que quede > 1 es un primo grande (> sqrt(R)).
        for (int i = 0; i < sz; i++)
            if (num[i] > 1)
                phi[i] -= phi[i] / num[i];

        pref.assign(sz + 1, 0);
        for (int i = 0; i < sz; i++)
            pref[i + 1] = pref[i] + phi[i];
    }                                          // O((R-L+1) log log R + sqrt(R))

    // READ: suma de phi(i) para i en [l, r], con L <= l <= r <= R.
    long long rango(long long l, long long r) {
        return pref[r - L + 1] - pref[l - L];
    }                                          // O(1)

    // READ: valor individual de phi(i), con L <= i <= R.
    long long get(long long i) {
        return phi[i - L];
    }                                          // O(1)
};
\end{lstlisting}

\textbf{Ejemplo de funcionamiento:}

\begin{lstlisting}[style=cpp]
PhiSegmentedRangeSum ps(95, 105);
// calcula phi(95), phi(96), ..., phi(105)
// sin cribar desde 1

cout << ps.get(100);
// phi(100) = 40

cout << ps.rango(95, 105);
// suma de phi(95) a phi(105)

cout << ps.rango(100, 102);
// phi(100) + phi(101) + phi(102)
\end{lstlisting}

\textbf{Cuál variante usar:}

\begin{table}[H]
\centering
\tiny
\setlength{\tabcolsep}{1.5pt}
\renewcommand{\arraystretch}{1.1}
\begin{tabular}{|p{0.32\columnwidth}|p{0.30\columnwidth}|p{0.27\columnwidth}|}
\hline
\textbf{Situación} & \textbf{Estructura} & \textbf{Complejidad} \\
\hline
$[1,N]$, $N$ chico/mediano
& \texttt{PhiRangeSum}
& $O(N)$ \\
\hline
$[L,R]$ grande, $R-L$ chico
& \texttt{PhiSegmentedRangeSum}
& $O((R{-}L)\log\log R+\sqrt{R})$ \\
\hline
$\phi$ cambia entre queries
& Fenwick sobre $\phi$
& $O(\log N)$ \\
\hline
\end{tabular}
\caption{Variantes para la suma de $\phi$.}
\label{tab:phi_variants}
\end{table}

\textbf{Regla práctica:} el cálculo de $\phi$ en sí (la criba, lineal
o segmentada) es teoría de números y no cambia según la query. Una vez
que tienes el arreglo de $\phi$ listo, trátalo como \emph{cualquier
otro arreglo estático}: Prefix Sums resuelve el resto. No hace falta
una estructura especial para $\phi$ más allá de construir bien el
arreglo inicial.



\subsection{Árbol de Fenwick (BIT)}

\textbf{Idea:} Un arreglo normal responde una suma de prefijo en $O(1)$ si guarda los prefijos, pero entonces actualizar un elemento cuesta $O(n)$. Si guarda los valores, ocurre lo contrario. El Fenwick reparte el trabajo y deja ambas operaciones en $O(\log n)$. La posición $i$ (1-indexada) guarda la suma de un bloque de largo $i\,\&\,(-i)$, es decir, el bit menos significativo de $i$, que termina en $i$: \texttt{t[i]} $=\sum_{j=i-\mathrm{lowbit}(i)+1}^{i} a_j$.

Un prefijo $[1,i]$ se descompone en bloques disjuntos quitando bits: se lee \texttt{t[i]}, luego $i \mathrel{-}= i\,\&\,(-i)$, y así hasta llegar a $0$. Cada paso elimina un bit, así que son a lo sumo $\log_2 n$ lecturas. Para actualizar la posición $i$ hay que corregir todos los bloques que la contienen: $i$, $i+\mathrm{lowbit}(i)$, y así sucesivamente. Cada salto agrega un bit más alto, de modo que también son $O(\log n)$. Como la consulta usa restas (rango $=$ prefijo$(r)-$prefijo$(l-1)$), la operación tiene que ser \textit{invertible}.

\textbf{¿Por qué usarlo?} Frente a un árbol de segmentos, el Fenwick sirve para sumas (o xor, o cualquier operación invertible) de prefijo y rango con actualizaciones puntuales, con menos memoria ($n$ celdas contra $2n$ o $4n$), menos código y una constante menor. Frente a un arreglo de prefijos, que responde en $O(1)$ pero actualiza en $O(n)$, gana cuando hay muchas actualizaciones intercaladas. Si el arreglo no cambia, el arreglo de prefijos es más simple y más rápido. No sirve para mínimo, máximo ni mcd (no tienen inverso): ahí se usa un árbol de segmentos, o una Sparse Table si el arreglo es estático. Para actualizaciones de rango con consultas de rango existe la variante de doble BIT. Si las actualizaciones son más generales (asignar, multiplicar), hace falta un árbol de segmentos con propagación perezosa. Frente a Mo's, el Fenwick permite contar elementos distintos en un rango con consultas offline en $O((n+q)\log n)$ en vez de $O(n\sqrt q)$ (última variante).

\textbf{Casos de uso:}
\begin{itemize}
	\item Suma de prefijo o de rango con actualizaciones puntuales frecuentes.
	\item Conteo de inversiones y frecuencias acumuladas, incluida la búsqueda del $k$-ésimo elemento.
	\item Suma de rango con actualización de rango (doble BIT) cuando no hace falta propagación perezosa general.
	\item Consultas 2D: suma o conteo de puntos dentro de un rectángulo.
	\item De nicho: contar elementos distintos en un rango con consultas fuera de línea (última variante).
\end{itemize}

\textbf{Complejidad:} Construir con el método de «empujar al padre» cuesta $O(n)$, porque cada celda se suma una sola vez a su padre; hacer $n$ actualizaciones costaría $O(n\log n)$. Una actualización puntual y una consulta de prefijo o de rango cuestan $O(\log n)$, por el recorrido de bits explicado arriba. La búsqueda del $k$-ésimo (\textit{binary lifting}) también es $O(\log n)$. Range update + range query con doble BIT cuesta $O(\log n)$ y la versión 2D $O(\log n\log m)$. Memoria: $O(n)$, o $O(nm)$ en 2D.

\subsubsection{BIT clásico (actualización puntual, consulta de rango)}

\begin{lstlisting}[style=cpp]
	struct ArbolFenwick {
		int n;
		vector<long long> t;   // t[i] = suma del bloque (i - lowbit(i), i]; 1-indexado
		
		// CREATE: arbol de tamano tam, todo en cero.
		ArbolFenwick(int tam = 0) : n(tam), t(tam + 1, 0) {} // O(n)
		
		// CREATE: construye el arbol a partir de a[0..n-1] empujando cada valor a su padre una vez.
		void construir(const vector<long long>& a) {
			t.assign(n + 1, 0);
			for (int i = 1; i <= n; i++) {
				t[i] += a[i - 1];
				int j = i + (i & -i);
				if (j <= n) t[j] += t[i];
			}
		} // O(n)
		
		// WRITE: suma v a la posicion i (1 <= i <= n).
		void actualizar(int i, long long v) {
			for (; i <= n; i += i & -i) t[i] += v;
		} // O(log n)
		
		// WRITE: fija el valor de la posicion i en x.
		void fijar(int i, long long x) {
			actualizar(i, x - valor(i));
		} // O(log n)
		
		// READ: suma del prefijo [1, i].
		long long prefijo(int i) const {
			long long s = 0;
			for (; i > 0; i -= i & -i) s += t[i];
			return s;
		} // O(log n)
		
		// READ: suma del rango [l, r].
		long long rango(int l, int r) const {
			if (l > r) return 0;
			return prefijo(r) - prefijo(l - 1);
		} // O(log n)
		
		// READ: valor puntual de la posicion i.
		long long valor(int i) const {
			return rango(i, i);
		} // O(log n)
		
		// READ: menor indice i con prefijo(i) >= objetivo (binary lifting).
		// Requiere valores >= 0. Devuelve n + 1 si el total no alcanza.
		int limite_inferior(long long objetivo) const {
			int pos = 0, paso = 1;
			long long ac = 0;
			while ((paso << 1) <= n) paso <<= 1;
			for (; paso > 0; paso >>= 1) {
				int nx = pos + paso;
				if (nx <= n && ac + t[nx] < objetivo) { pos = nx; ac += t[nx]; }
			}
			return pos + 1;
		} // O(log n)
		
		// WRITE: reinicia el arbol a ceros.
		void limpiar() {
			fill(t.begin(), t.end(), 0);
		} // O(n)
	};
	
	// ------------------------------------------------------------
	// MODIFICACIONES Y COMENTARIOS CON LOS USOS
	// ------------------------------------------------------------
	// 1. Conteo de inversiones: comprimir los valores a 1..k (ver "Compresion de Coordenadas")
	//    y recorrer de izquierda a derecha. Los anteriores mayores que c[i] son i - prefijo(c[i]).
	//    ArbolFenwick bit(k); long long inv = 0;
	//    for (int i = 0; i < n; i++) { inv += i - bit.prefijo(c[i]); bit.actualizar(c[i], 1); }
	//    // O(n log n)
	//
	// 2. Frecuencias y k-esimo elemento: bit.actualizar(x, 1) al insertar x y bit.actualizar(x, -1)
	//    al borrarlo. El k-esimo menor del multiconjunto es limite_inferior(k). O(log n) por operacion.
	//
	// 3. XOR en lugar de suma: cambiar "+=" por "^=" en actualizar y en prefijo, y el rango es
	//    prefijo(r) ^ prefijo(l - 1). Sirve con cualquier operacion con inverso (xor, suma modular).
	//
	// 4. Minimo o maximo de prefijo: SOLO si los valores unicamente crecen (para el maximo) o
	//    unicamente decrecen (para el minimo). Se cambia "+=" por max/min. No admite rangos
	//    arbitrarios ni cambios en sentido contrario: para eso, arbol de segmentos.
	//
	// 5. Indices desde 0: sumar 1 antes de llamar. Con i = 0 el ciclo de actualizar no termina
	//    (0 & -0 = 0).
	//
	// 6. Valores grandes o negativos como indices: comprimirlos primero (ver "Compresion de
	//    Coordenadas"), para que el arbol tenga tamano O(#valores distintos).
	//
	// 7. Suma modulo p: aplicar "% p" en cada suma y restar sumando p en rango():
	//    return ((prefijo(r) - prefijo(l - 1)) % p + p) % p;
\end{lstlisting}

\textbf{Ejemplo de funcionamiento:}
\begin{lstlisting}[style=cpp]
	int main() {
		vector<long long> a = {5, 2, 9, 1, 7};
		ArbolFenwick bit(5);
		bit.construir(a);                       // posiciones 1..5: [5, 2, 9, 1, 7]
		
		cout << bit.prefijo(3) << "\n";         // 16  (5 + 2 + 9)
		cout << bit.rango(2, 4) << "\n";        // 12  (2 + 9 + 1)
		
		bit.actualizar(2, 3);                   // [5, 5, 9, 1, 7]
		cout << bit.valor(2) << "\n";           // 5
		
		bit.fijar(4, 10);                       // [5, 5, 9, 10, 7]
		cout << bit.limite_inferior(15) << "\n"; // 3  (prefijos 5, 10, 19: el primero >= 15 es el de 3)
		return 0;
	}
\end{lstlisting}

\textbf{Nota:} El arbol es 1-indexado: \texttt{actualizar} con $i=0$ nunca termina. \texttt{construir} reinicia el arbol, así que no combina con valores previos. \texttt{limite\_inferior} solo es correcto si todos los valores son no negativos, porque usa la monotonía de los prefijos, y devuelve $n+1$ si el total no alcanza el objetivo. Por último, las sumas se guardan en \texttt{long long} para no desbordar con valores grandes.

\subsubsection{BIT con Range Update + Range Query (doble BIT)}

\textbf{Idea:} Sea $d$ el arreglo de diferencias: sumar $v$ a $[l,r]$ equivale a $d_l \mathrel{+}= v$ y $d_{r+1} \mathrel{-}= v$. El valor $a_j=\sum_{k\le j}d_k$, y la suma de prefijo es
$$\sum_{j\le i}a_j=\sum_{k\le i}d_k\,(i-k+1)=i\sum_{k\le i}d_k-\sum_{k\le i}d_k\,(k-1).$$
Basta entonces un Fenwick para $\sum d_k$ (\texttt{b1}) y otro para $\sum d_k(k-1)$ (\texttt{b2}). Con ellos, el prefijo es \texttt{b1.prefijo(i)*i - b2.prefijo(i)}, y cada actualización de rango son cuatro actualizaciones puntuales.

\begin{lstlisting}[style=cpp]
	struct FenwickRango {
		int n;
		ArbolFenwick b1, b2;   // b1: sumas de d[k]; b2: sumas de d[k] * (k - 1)
		
		// CREATE: arreglo logico de n ceros.
		FenwickRango(int tam) : n(tam), b1(tam), b2(tam) {} // O(n)
		
		// WRITE: suma v a todos los elementos de [l, r].
		void actualizar_rango(int l, int r, long long v) {
			b1.actualizar(l, v);
			b1.actualizar(r + 1, -v);
			b2.actualizar(l, v * (l - 1));
			b2.actualizar(r + 1, -v * r);
		} // O(log n)
		
		// READ: suma del prefijo [1, i].
		long long prefijo(int i) const {
			return b1.prefijo(i) * i - b2.prefijo(i);
		} // O(log n)
		
		// READ: suma del rango [l, r].
		long long rango(int l, int r) const {
			if (l > r) return 0;
			return prefijo(r) - prefijo(l - 1);
		} // O(log n)
		
		// READ: valor puntual de la posicion i.
		long long valor(int i) const {
			return rango(i, i);
		} // O(log n)
	};
	
	// ------------------------------------------------------------
	// MODIFICACIONES Y COMENTARIOS CON LOS USOS
	// ------------------------------------------------------------
	// 1. Solo actualizacion de rango + consulta PUNTUAL: basta b1. El valor de la posicion i
	//    es b1.prefijo(i), con un solo arbol y menos constante.
	//
	// 2. Arreglo inicial no nulo: sumar al resultado un arreglo de prefijos P del original:
	//    suma(l, r) = P[r] - P[l-1] + rango(l, r). O(log n) por consulta.
	//
	// 3. Si b1 y b2 se desbordan con valores grandes: v * (l - 1) puede llegar a v * n.
	//    Con v ~ 10^9 y n ~ 10^6 son 10^15, cabe en long long; si es mayor, usar __int128.
	//
	// 4. Para max/min con actualizaciones de rango no sirve: usar arbol de segmentos con
	//    propagacion perezosa.
\end{lstlisting}

\textbf{Ejemplo de funcionamiento:}
\begin{lstlisting}[style=cpp]
	int main() {
		FenwickRango f(5);                       // [0, 0, 0, 0, 0]
		f.actualizar_rango(2, 4, 3);             // [0, 3, 3, 3, 0]
		cout << f.rango(1, 5) << "\n";           // 9
		
		f.actualizar_rango(1, 3, 2);             // [2, 5, 5, 3, 0]
		cout << f.valor(2) << "\n";              // 5
		cout << f.rango(1, 5) << "\n";           // 15  (2 + 5 + 5 + 3 + 0)
		return 0;
	}
\end{lstlisting}

\textbf{Nota:} Con $r=n$, la actualización en $r+1$ cae fuera del arbol y el ciclo simplemente no se ejecuta, así que no hace falta tratar ese caso aparte. Además, la fórmula supone un arreglo inicial de ceros (modificación 2).

\subsubsection{BIT 2D (actualización puntual, suma de submatriz)}

\textbf{Idea:} Es un BIT dentro de otro BIT: \texttt{t[i][j]} guarda la suma del rectángulo de $\mathrm{lowbit}(i)$ filas por $\mathrm{lowbit}(j)$ columnas que termina en $(i,j)$. Una actualización o consulta recorre los $O(\log n)$ índices de fila y, en cada uno, los $O(\log m)$ de columna. La suma de un rectángulo se obtiene por inclusión--exclusión de cuatro prefijos.

\begin{lstlisting}[style=cpp]
	struct Fenwick2D {
		int n, m;
		vector<vector<long long>> t;
		
		// CREATE: matriz de n x m en ceros (1-indexada).
		Fenwick2D(int filas, int cols)
		: n(filas), m(cols), t(filas + 1, vector<long long>(cols + 1, 0)) {} // O(n m)
		
		// WRITE: suma v a la celda (x, y).
		void actualizar(int x, int y, long long v) {
			for (int i = x; i <= n; i += i & -i)
			for (int j = y; j <= m; j += j & -j)
			t[i][j] += v;
		} // O(log n log m)
		
		// READ: suma de la submatriz [1..x] x [1..y].
		long long prefijo(int x, int y) const {
			long long s = 0;
			for (int i = x; i > 0; i -= i & -i)
			for (int j = y; j > 0; j -= j & -j)
			s += t[i][j];
			return s;
		} // O(log n log m)
		
		// READ: suma del rectangulo [x1..x2] x [y1..y2].
		long long rango(int x1, int y1, int x2, int y2) const {
			return prefijo(x2, y2) - prefijo(x1 - 1, y2)
			- prefijo(x2, y1 - 1) + prefijo(x1 - 1, y1 - 1);
		} // O(log n log m)
	};
	
	// ------------------------------------------------------------
	// MODIFICACIONES Y COMENTARIOS CON LOS USOS
	// ------------------------------------------------------------
	// 1. Actualizacion de rectangulo + consulta puntual: usar diferencias 2D. Sumar v al
	//    rectangulo [x1..x2] x [y1..y2] son cuatro actualizaciones puntuales:
	//    actualizar(x1, y1, v); actualizar(x1, y2 + 1, -v);
	//    actualizar(x2 + 1, y1, -v); actualizar(x2 + 1, y2 + 1, v);
	//    El valor de (x, y) es prefijo(x, y). O(log n log m).
	//
	// 2. Actualizacion de rectangulo + consulta de rectangulo: hacen falta cuatro arboles 2D
	//    (la generalizacion del doble BIT). Es raro en concursos; suele bastar un arbol de
	//    segmentos 2D.
	//
	// 3. Memoria: n * m celdas de long long. Con n, m ~ 5000 son 200 MB; si los puntos son
	//    pocos, comprimir coordenadas (offline) o usar un BIT de BITs dispersos.
	//
	// 4. Contar puntos (x, y) en un rectangulo con consultas offline: ordenar por x y usar
	//    un BIT 1D sobre y (barrido), que baja a O((n + q) log n) y O(n) de memoria.
\end{lstlisting}

\textbf{Ejemplo de funcionamiento:}
\begin{lstlisting}[style=cpp]
	int main() {
		Fenwick2D f(4, 4);
		f.actualizar(1, 1, 5);
		f.actualizar(2, 3, 7);
		f.actualizar(4, 4, 2);
		
		cout << f.rango(1, 1, 2, 3) << "\n";   // 12  (5 + 7)
		cout << f.rango(3, 3, 4, 4) << "\n";   // 2
		return 0;
	}
\end{lstlisting}

\subsubsection{Variante: Conteo de Elementos Distintos en un Rango (offline)}

\textbf{Idea:} Se atienden las consultas \textit{ordenadas por extremo derecho} $r$ y se recorre el arreglo de izquierda a derecha. Al llegar a la posición $i$, se mantiene un Fenwick con un $1$ en la \textit{última aparición} (hasta $i$) de cada valor y $0$ en el resto. Para ello, al procesar $a_i$ se borra el $1$ de su aparición anterior (si existe) y se coloca uno en $i$. Después de procesar $r$, la respuesta a $[l,r]$ es la suma de ese Fenwick en $[l,r]$.

Es correcto porque cada valor presente en $[l,r]$ tiene exactamente una marca, su última aparición hasta $r$, y esa marca cae dentro de $[l,r]$ (está en $[l,r]$ porque el valor aparece ahí y la última aparición es $\ge$ cualquiera de las apariciones en el rango). Un valor que no aparece en $[l,r]$ tiene su marca antes de $l$ (o ninguna) y no se cuenta. Así, cada posición del rango se cuenta una sola vez por valor, y no hay que volver a recorrer el rango.

\textbf{Complejidad:} Cada elemento hace a lo sumo dos actualizaciones ($-1$ en la aparición anterior, $+1$ en la actual) y cada consulta una resta de prefijos, todas $O(\log n)$. Si las consultas se agrupan por $r$ con listas (en vez de ordenarlas), el total es $O((n+q)\log n)$, contra $O(n\sqrt q)$ de Mo's. Memoria: $O(n+q+\texttt{max\_val})$.

\begin{lstlisting}[style=cpp]
	// Valores en [0, max_val]. Consultas (l, r) 0-indexadas e inclusivas, todas conocidas de antemano.
	struct ConteoDistintos {
		int n, max_val;
		vector<int> a;
		
		// CREATE: guarda el arreglo y el valor maximo.
		ConteoDistintos(const vector<int>& arr, int mv) : n((int)arr.size()), max_val(mv), a(arr) {} // O(n)
		
		// READ: responde todas las consultas y las devuelve en el orden original.
		vector<int> resolver(const vector<array<int, 2>>& consultas) const {
			int q = (int)consultas.size();
			vector<vector<int>> por_r(n);                      // por_r[r]: consultas que terminan en r
			for (int id = 0; id < q; id++) por_r[consultas[id][1]].push_back(id);
			
			ArbolFenwick bit(n);
			vector<int> ultima(max_val + 1, 0);                // ultima posicion (1-indexada) de cada valor; 0 = ninguna
			vector<int> resp(q);
			for (int i = 1; i <= n; i++) {
				int v = a[i - 1];
				if (ultima[v]) bit.actualizar(ultima[v], -1);  // la marca anterior de v deja de contar
				bit.actualizar(i, 1);                          // ahora la ultima aparicion de v es i
				ultima[v] = i;
				for (int id : por_r[i - 1])
				resp[id] = (int)bit.rango(consultas[id][0] + 1, i);
			}
			return resp;
		} // O((n + q) log n)
	};
	
	// ------------------------------------------------------------
	// MODIFICACIONES Y COMENTARIOS CON LOS USOS
	// ------------------------------------------------------------
	// 1. Suma de los valores DISTINTOS del rango (cada valor una sola vez): en lugar de 1,
	//    marcar el propio valor.
	//    if (ultima[v]) bit.actualizar(ultima[v], -v);
	//    bit.actualizar(i, v);
	//    El resto es identico. O((n + q) log n).
	//
	// 2. Valores grandes o negativos: comprimirlos primero (ver "Compresion de Coordenadas"),
	//    para que ultima[] tenga tamano O(n).
	//
	// 3. Consultas EN LINEA (una a la vez, sin ver las siguientes): este metodo no sirve, porque
	//    depende de procesar por r creciente. Opciones: Fenwick/arbol de segmentos persistente
	//    (una version por cada r, O(log n) por consulta), o un arbol de segmentos con mezcla
	//    de listas ordenadas.
	//
	// 4. Con actualizaciones puntuales del arreglo: el metodo se rompe, porque cambiar un valor
	//    altera varias "ultimas apariciones". Opciones: Mo's con actualizaciones o CDQ (divide
	//    y venceras sobre el tiempo).
	//
	// 5. Simetria: tambien se puede barrer de derecha a izquierda, ordenando por l y marcando la
	//    PRIMERA aparicion de cada valor. Mismo costo.
	//
	// 6. Frente a Mo's: este metodo es O((n + q) log n) contra O(n sqrt(q)), asi que gana cuando
	//    n y q son grandes; Mo's es mas flexible (sirve para respuestas que no se pueden
	//    mantener con un Fenwick, como la moda).
\end{lstlisting}

\textbf{Ejemplo de funcionamiento:}
\begin{lstlisting}[style=cpp]
	int main() {
		vector<int> a = {1, 2, 1, 3, 2, 4, 1};
		ConteoDistintos cd(a, 4);
		
		vector<array<int, 2>> consultas = {
			{0, 2},    // [1, 2, 1]             -> 2 distintos
			{1, 4},    // [2, 1, 3, 2]          -> 3 distintos
			{0, 6},    // [1, 2, 1, 3, 2, 4, 1] -> 4 distintos
			{2, 5},    // [1, 3, 2, 4]          -> 4 distintos
			{3, 3}     // [3]                   -> 1 distinto
		};
		
		for (int x : cd.resolver(consultas)) cout << x << " ";
		cout << "\n";                               // 2 3 4 4 1
		return 0;
	}
\end{lstlisting}

\textbf{Nota:} Como las consultas se agrupan por $r$, no se ordenan: el recorrido es lineal más las respuestas. Las consultas deben cumplir $0\le l\le r<n$ y todos los valores deben estar en $[0,\texttt{max\_val}]$. \texttt{resolver} es \texttt{const} porque el Fenwick y \texttt{ultima} son locales, así que el objeto puede reutilizarse con otras consultas.

\textbf{Regla práctica:} Si el arreglo cambia y solo necesitas sumas (o cualquier operación invertible) de prefijo o rango, usa Fenwick: es más corto y rápido que un árbol de segmentos. Si además hay actualizaciones de rango, usa el doble BIT. Si son consultas 2D, el BIT 2D. Si el arreglo no cambia, un arreglo de prefijos es más simple. Para min, max o mcd, o para actualizaciones de rango generales, pasa a un árbol de segmentos con propagación perezosa. Para contar elementos distintos en rangos con consultas offline, ordena por $r$ y marca la última aparición de cada valor en un Fenwick; si las consultas son en línea o hay actualizaciones, usa Mo's con actualizaciones o un árbol persistente.
\subsection{Sqrt Decomposition}

\textbf{Idea:} divide el arreglo en bloques de tamaño
$\approx\sqrt n$ y guarda un resumen (suma, mínimo, etc.) por bloque.
Un update toca solo un bloque; una query de rango toca a lo más dos
bloques \emph{parciales} (recorridos elemento por elemento) y varios
bloques \emph{completos} en el medio (usados de golpe con su resumen
precalculado). Como hay $\approx\sqrt n$ bloques y cada bloque tiene
$\approx\sqrt n$ elementos, todo cuesta $O(\sqrt n)$: ni tan lento
como un arreglo plano, ni tan rápido como Fenwick/Segment Tree, pero
mucho más simple de escribir y de adaptar a operaciones raras.

\textbf{¿Por qué usarlo?} Es la estructura más flexible cuando no
tienes tiempo de pensar un Segment Tree con lazy propagation
correcto: soporta casi cualquier operación (suma, min, max, contar
algo con una condición) con el mismo esqueleto de código, aunque sea
$O(\sqrt n)$ en vez de $O(\log n)$. También es la base conceptual de
Mo's Algorithm.

\textbf{Casos de uso:}

\begin{itemize}
    \item Necesitas implementar algo rápido bajo presión, sin
    memorizar casos de Segment Tree con lazy propagation.
    \item Operaciones donde min/max cambian de forma que Fenwick no
    puede manejar (Fenwick solo sirve para operaciones invertibles
    como la suma).
    \item $n$ no es tan grande ($\le 10^5$–$10^6$) y $O(\sqrt n)$ por
    operación es suficiente para el límite de tiempo.
    \item Base para Mo's Algorithm (ver esa sección).
\end{itemize}

\textbf{Complejidad:} construir $O(n)$, update puntual $O(1)$ (suma)
u $O(\sqrt n)$ (min/max), query de rango $O(\sqrt n)$, range update +
range query $O(\sqrt n)$ con lazy por bloque.


\subsubsection{Sqrt Decomposition clásico (point update, range query — suma)}

\begin{lstlisting}[style=cpp]
struct SqrtDecomposition {
    vector<long long> a, bloque;
    int n, tam;

    // CREATE: construye el arreglo en bloques de tamano ~sqrt(n).
    SqrtDecomposition(vector<long long>& arr) {
        a = arr;
        n = a.size();
        tam = max(1, (int)sqrt((double)n));

        bloque.assign((n + tam - 1) / tam, 0);

        for (int i = 0; i < n; i++)
            bloque[i / tam] += a[i];
    }                                          // O(n)

    // UPDATE: suma v al elemento en la posicion i.
    void update(int i, long long v) {
        a[i] += v;
        bloque[i / tam] += v;
    }                                          // O(1)

    // READ: suma del rango [l, r] (0-indexado, inclusivo).
    long long rango(int l, int r) {
        long long s = 0;
        int bl = l / tam, br = r / tam;

        if (bl == br) {
            for (int i = l; i <= r; i++)
                s += a[i];
            return s;
        }

        for (int i = l; i < (bl + 1) * tam; i++)
            s += a[i];

        for (int b = bl + 1; b < br; b++)
            s += bloque[b];

        for (int i = br * tam; i <= r; i++)
            s += a[i];

        return s;
    }                                          // O(sqrt(n))
};
\end{lstlisting}

\textbf{Ejemplo de funcionamiento:}

\begin{lstlisting}[style=cpp]
vector<long long> a = {5, 2, 9, 1, 7, 3, 8, 4};

SqrtDecomposition sd(a);
// arreglo logico: [5, 2, 9, 1, 7, 3, 8, 4]  (posiciones 0..7)
// tam = 2 -> bloques: [5+2, 9+1, 7+3, 8+4] = [7, 10, 10, 12]

cout << sd.rango(1, 5);
// 2 + 9 + 1 + 7 + 3 = 22

sd.update(3, 10);
// arreglo logico: [5, 2, 9, 11, 7, 3, 8, 4]

cout << sd.rango(1, 5);
// 2 + 9 + 11 + 7 + 3 = 32
\end{lstlisting}


\subsubsection{Sqrt Decomposition con range update + range query (lazy por bloque)}

\textbf{Idea:} igual que arriba, pero cada bloque guarda un valor
\texttt{lazy} pendiente que representa ``sumar $v$ a todo el bloque''
sin tocar cada elemento individualmente. Solo se aplica a los
elementos reales cuando un update o query \textbf{parcial} necesita
tocar ese bloque específico.

\begin{lstlisting}[style=cpp]
struct SqrtRangeUpdate {
    vector<long long> a, bloqueSum, lazy;
    int n, tam;

    // CREATE: construye el arreglo en bloques de tamano ~sqrt(n).
    SqrtRangeUpdate(vector<long long>& arr) {
        a = arr;
        n = a.size();
        tam = max(1, (int)sqrt((double)n));

        int nb = (n + tam - 1) / tam;
        bloqueSum.assign(nb, 0);
        lazy.assign(nb, 0);

        for (int i = 0; i < n; i++)
            bloqueSum[i / tam] += a[i];
    }                                          // O(n)

    // READ auxiliar: cantidad de elementos reales del bloque b.
    int tamBloque(int b) {
        return min(n, (b + 1) * tam) - b * tam;
    }                                          // O(1)

    // UPDATE: suma v a todo el rango [l, r].
    void rangoUpdate(int l, int r, long long v) {
        int bl = l / tam, br = r / tam;

        if (bl == br) {
            for (int i = l; i <= r; i++)
                a[i] += v;
            bloqueSum[bl] += v * (r - l + 1);
            return;
        }

        for (int i = l; i < (bl + 1) * tam; i++)
            a[i] += v;
        bloqueSum[bl] += v * ((bl + 1) * tam - l);

        for (int b = bl + 1; b < br; b++)
            lazy[b] += v;

        for (int i = br * tam; i <= r; i++)
            a[i] += v;
        bloqueSum[br] += v * (r - br * tam + 1);
    }                                          // O(sqrt(n))

    // READ: suma del rango [l, r].
    long long rango(int l, int r) {
        long long s = 0;
        int bl = l / tam, br = r / tam;

        if (bl == br) {
            for (int i = l; i <= r; i++)
                s += a[i] + lazy[bl];
            return s;
        }

        for (int i = l; i < (bl + 1) * tam; i++)
            s += a[i] + lazy[bl];

        for (int b = bl + 1; b < br; b++)
            s += bloqueSum[b] + lazy[b] * tamBloque(b);

        for (int i = br * tam; i <= r; i++)
            s += a[i] + lazy[br];

        return s;
    }                                          // O(sqrt(n))
};
\end{lstlisting}

\textbf{Ejemplo de funcionamiento:}

\begin{lstlisting}[style=cpp]
vector<long long> a = {1, 1, 1, 1, 1, 1};

SqrtRangeUpdate sr(a);
// tam = 2 -> bloques: [1+1, 1+1, 1+1] = [2, 2, 2]

sr.rangoUpdate(1, 4, 10);
// arreglo logico "efectivo": [1, 11, 11, 11, 11, 1]

cout << sr.rango(0, 5);
// 1+11+11+11+11+1 = 46

cout << sr.rango(2, 3);
// 11+11 = 22
\end{lstlisting}


\subsubsection{Variante para mínimo / máximo (point update, range query)}

\textbf{Idea:} a diferencia de Fenwick (que solo sirve para
operaciones \emph{invertibles} como la suma), Sqrt Decomposition
soporta min/max con el mismo esqueleto: la única diferencia es que un
update ya no puede ajustar el resumen del bloque en $O(1)$, porque
quitar el valor viejo no dice si sigue habiendo otro elemento igual
de chico. Hay que \textbf{reconstruir el mínimo de todo el bloque},
lo cual sigue siendo $O(\sqrt n)$.

\begin{lstlisting}[style=cpp]
struct SqrtMinQuery {
    vector<long long> a, bloqueMin;
    int n, tam;

    // CREATE: construye el arreglo en bloques de tamano ~sqrt(n).
    SqrtMinQuery(vector<long long>& arr) {
        a = arr;
        n = a.size();
        tam = max(1, (int)sqrt((double)n));

        bloqueMin.assign((n + tam - 1) / tam, LLONG_MAX);

        for (int i = 0; i < n; i++)
            bloqueMin[i / tam] = min(bloqueMin[i / tam], a[i]);
    }                                          // O(n)

    // UPDATE: fija el valor absoluto de la posicion i en x.
    // Reconstruye el minimo de todo el bloque, ya que a diferencia
    // de la suma, el minimo no se puede ajustar incrementalmente.
    void set(int i, long long x) {
        a[i] = x;

        int b = i / tam;
        long long mn = LLONG_MAX;

        for (int j = b * tam; j < min(n, (b + 1) * tam); j++)
            mn = min(mn, a[j]);

        bloqueMin[b] = mn;
    }                                          // O(sqrt(n))

    // READ: minimo del rango [l, r].
    long long rango(int l, int r) {
        long long mn = LLONG_MAX;
        int bl = l / tam, br = r / tam;

        if (bl == br) {
            for (int i = l; i <= r; i++)
                mn = min(mn, a[i]);
            return mn;
        }

        for (int i = l; i < (bl + 1) * tam; i++)
            mn = min(mn, a[i]);

        for (int b = bl + 1; b < br; b++)
            mn = min(mn, bloqueMin[b]);

        for (int i = br * tam; i <= r; i++)
            mn = min(mn, a[i]);

        return mn;
    }                                          // O(sqrt(n))
};
\end{lstlisting}

\textbf{Ejemplo de funcionamiento:}

\begin{lstlisting}[style=cpp]
vector<long long> a = {5, 2, 9, 1, 7, 3, 8, 4};

SqrtMinQuery sm(a);
// tam = 2 -> bloques min: [2, 1, 3, 4]

cout << sm.rango(0, 4);
// min(5,2,9,1,7) = 1

sm.set(3, 100);
// arreglo logico: [5, 2, 9, 100, 7, 3, 8, 4]

cout << sm.rango(0, 4);
// min(5,2,9,100,7) = 2
\end{lstlisting}

\textbf{Comparación con Fenwick y Segment Tree:}

\begin{table}[H]
\raggedright
\scriptsize
\setlength{\tabcolsep}{3pt}
\begin{tabular}{|l|c|c|c|}
\hline
\textbf{Estructura} & \textbf{Update} & \textbf{Query} & \textbf{Min/Max?} \\
\hline
Sqrt Decomposition & $O(1)$–$O(\sqrt n)$ & $O(\sqrt n)$ & Sí \\
\hline
\texttt{FenwickTree} & $O(\log n)$ & $O(\log n)$ & No \\
\hline
Segment Tree & $O(\log n)$ & $O(\log n)$ & Sí \\
\hline
\end{tabular}
\caption{Sqrt Decomposition vs. Fenwick vs. Segment Tree.}
\end{table}

\textbf{Regla práctica:} usa Sqrt Decomposition cuando quieras algo
simple y flexible de escribir rápido, sobre todo si la operación no
es una suma (Fenwick no sirve) y no quieres arriesgarte con un
Segment Tree + Lazy Propagation bajo presión de tiempo. Si $n$ es muy
grande o el límite de tiempo es ajustado, $O(\log n)$ de Segment Tree
gana; si la operación es solo suma, Fenwick es más simple y más
rápido.


\subsection{Sparse Table (RMQ)}

\textbf{Idea:} Una Sparse Table responde consultas de rango en $O(1)$ sobre un arreglo que \textit{no cambia}. Precalcula \texttt{sp[k][i]}, el resultado de combinar el bloque que empieza en $i$ y tiene largo $2^k$. La tabla se llena por doblado: un bloque de largo $2^k$ son dos bloques de largo $2^{k-1}$ consecutivos, así que \texttt{sp[k][i]} se obtiene de \texttt{sp[k-1][i]} y \texttt{sp[k-1][i + $2^{k-1}$]} con una sola combinación.

La clave está en la consulta. Para $[l,r]$ de largo $len=r-l+1$ se toma $k=\lfloor\log_2 len\rfloor$, de modo que $2^k\le len<2^{k+1}$. Dos bloques de largo $2^k$, uno pegado a $l$ y otro pegado a $r$, cubren todo el rango y se \textit{solapan} (porque $2\cdot 2^k > len$ no deja hueco: $2^{k+1}>len$ implica que el segundo empieza antes de que termine el primero, o justo donde termina). Si la operación es \textit{idempotente} ($f(x,x)=x$, como min, max, gcd, AND u OR), contar el solape dos veces no cambia el resultado, y la consulta son dos lecturas y una combinación. Con operaciones no idempotentes (suma, xor) el solape sí se nota, y hace falta la \textit{Disjoint Sparse Table}, que evita el solape por construcción.

\textbf{¿Por qué usarlo?} Es la estructura más rápida para consultas de rango cuando el arreglo es estático: $O(1)$ por consulta, frente a $O(\log n)$ del árbol de segmentos y del Fenwick, o frente a $O(n)$ de recorrer el rango. El precio es que no admite actualizaciones. Cambiar un solo elemento invalida $O(n)$ celdas por nivel, y reconstruir cuesta $O(n\log n)$, así que si el arreglo cambia conviene un árbol de segmentos. Frente a un prefijo de sumas, que también responde en $O(1)$, la Sparse Table sirve para operaciones que \textit{no tienen inverso} (el mínimo no se puede «restar»). Y para sumas o xor, el Fenwick o los prefijos son más simples que la Disjoint Sparse Table, que solo se justifica con operaciones asociativas sin inverso ni idempotencia (por ejemplo, producto de matrices). Frente a Mo's (subsección anterior), aquí la operación sí se combina desde dos mitades y no hace falta que las consultas sean offline.

\textbf{Casos de uso:}
\begin{itemize}
	\item RMQ clásico: mínimo o máximo de cualquier rango de un arreglo estático.
	\item gcd, AND u OR de rango sobre un arreglo estático (variante más abajo).
	\item LCA con recorrido de Euler y RMQ: la consulta del ancestro común es un mínimo sobre profundidades.
	\item Operaciones no idempotentes de rango estático en $O(1)$ (suma, xor, producto de matrices) con la Disjoint Sparse Table.
	\item De nicho: arreglos de sufijos (LCP entre dos sufijos como mínimo de rango) o problemas con muchísimas consultas sobre datos que nunca se modifican.
\end{itemize}

\textbf{Complejidad:} La tabla tiene $K=\lfloor\log_2 n\rfloor+1$ niveles de a lo sumo $n$ celdas, y cada celda se calcula con una combinación $O(1)$ a partir del nivel anterior, así que construir cuesta $O(n\log n)$ en tiempo y en memoria. La consulta lee dos celdas y una tabla de logaritmos precalculada: $O(1)$. Con $n=10^6$ la tabla ocupa unos $20\cdot 10^6$ valores, lo que ya pesa en memoria.

\subsubsection{Sparse Table clásico (operaciones idempotentes: min, max)}

\begin{lstlisting}[style=cpp]
	struct SparseTable {
		int n;
		vector<vector<long long>> sp;   // sp[k][i] = minimo del bloque [i, i + 2^k - 1]
		vector<int> lg;                 // lg[x] = piso(log2(x))
		
		// CREATE: precalcula la tabla a partir de a[0..n-1].
		SparseTable(const vector<long long>& a) {
			n = (int)a.size();
			lg.assign(n + 1, 0);
			for (int i = 2; i <= n; i++) lg[i] = lg[i / 2] + 1;
			
			int K = lg[n] + 1;
			sp.assign(K, vector<long long>(n));
			sp[0] = a;
			for (int k = 1; k < K; k++)
			for (int i = 0; i + (1 << k) <= n; i++)
			sp[k][i] = min(sp[k - 1][i], sp[k - 1][i + (1 << (k - 1))]);
		} // O(n log n)
		
		// READ: minimo del rango [l, r] (0-indexado, inclusivo).
		// Dos bloques de largo 2^k que se solapan; min es idempotente, el solape no importa.
		long long consultar(int l, int r) const {
			int k = lg[r - l + 1];
			return min(sp[k][l], sp[k][r - (1 << k) + 1]);
		} // O(1)
	};
	
	// ------------------------------------------------------------
	// MODIFICACIONES Y COMENTARIOS CON LOS USOS
	// ------------------------------------------------------------
	// 1. Maximo, AND u OR: cambiar min por max, & o | en el constructor y en consultar.
	//    Para gcd ver la variante siguiente (tiene su propia complejidad).
	//
	// 2. Indice del minimo (no solo el valor): guardar indices y comparar por valor.
	//    sp[0][i] = i;
	//    sp[k][i] = (a[sp[k-1][i]] <= a[sp[k-1][i + (1 << (k-1))]]) ? sp[k-1][i] : sp[k-1][i + (1 << (k-1))];
	//    consultar devuelve el indice con el mismo criterio. Con "<=" gana el mas a la izquierda.
	//
	// 3. Sin tabla lg: calcular el piso del logaritmo con la CPU.
	//    int k = 31 - __builtin_clz(r - l + 1);   // O(1), sin vector extra
	//
	// 4. LCA con Euler Tour: recorrer el arbol anotando (profundidad, nodo) en cada visita,
	//    guardar la primera aparicion de cada nodo, y el LCA(u, v) es el nodo de menor
	//    profundidad entre esas dos posiciones (RMQ sobre pares). O(n log n) + O(1) por LCA.
	//
	// 5. Consulta NO idempotente con Sparse Table normal: se parte [l, r] en bloques de
	//    potencias de 2 sin solape (los bits de la longitud), en O(log n) por consulta.
	//    long long s = 0;
	//    for (int k = K - 1; k >= 0; k--)
	//        if ((1 << k) <= r - l + 1) { s += sp[k][l]; l += 1 << k; }   // O(log n)
	//    Aqui sp[k][i] guarda la SUMA del bloque. Es mas lento que prefijos; solo sirve si la
	//    operacion no tiene inverso y no se quiere escribir la Disjoint Sparse Table.
	//
	// 6. Sparse Table 2D (matriz estatica): sp[kx][ky][i][j] = minimo del bloque
	//    2^kx x 2^ky. Construccion O(n m log n log m), consulta O(1) con cuatro bloques.
	//
	// 7. Reducir memoria: si solo hay consultas offline, ordenarlas por k y guardar solo dos
	//    niveles a la vez (O(n) de memoria), a costa de ordenar las consultas.
\end{lstlisting}

\textbf{Ejemplo de funcionamiento:}
\begin{lstlisting}[style=cpp]
	int main() {
		vector<long long> a = {5, 2, 9, 1, 7, 3};
		SparseTable st(a);
		
		cout << st.consultar(0, 2) << "\n";   // 2  (min de 5, 2, 9)
		cout << st.consultar(2, 5) << "\n";   // 1  (min de 9, 1, 7, 3)
		cout << st.consultar(4, 4) << "\n";   // 7  (rango de un solo elemento)
		cout << st.consultar(0, 5) << "\n";   // 1  (arreglo completo)
		return 0;
	}
\end{lstlisting}

\subsubsection{Variante: GCD en un Rango}

\textbf{Idea:} El gcd es idempotente: $\gcd(x,x)=x$. Además es asociativo y conmutativo, así que es exactamente igual que el mínimo o el máximo para esta estructura, y solo cambia la función de combinación. El solape de los dos bloques no altera el resultado, porque $\gcd(g_1,g_2)$ es el gcd de la unión de elementos, sin importar que un elemento aparezca dos veces. Se usan valores no negativos (el gcd de $0$ y $x$ es $x$, y el gcd de un rango de ceros es $0$).

La diferencia con min y max es el costo de la combinación. Cada \texttt{\_\_gcd} es un Euclides que cuesta $O(\log M)$, con $M$ el mayor valor. Por tanto, construir cuesta $O(n\log n\log M)$ y cada consulta $O(\log M)$. En la práctica los gcd de bloques grandes bajan rápido (el gcd de muchos elementos suele ser pequeño y Euclides termina en pocos pasos), así que la consulta se comporta casi como $O(1)$, pero el límite formal es $O(\log M)$.

\begin{lstlisting}[style=cpp]
	struct SparseTableGcd {
		int n;
		vector<vector<long long>> sp;   // sp[k][i] = gcd del bloque [i, i + 2^k - 1]
		vector<int> lg;
		
		// CREATE: precalcula la tabla de gcd (los valores se toman en valor absoluto).
		SparseTableGcd(const vector<long long>& a) {
			n = (int)a.size();
			lg.assign(n + 1, 0);
			for (int i = 2; i <= n; i++) lg[i] = lg[i / 2] + 1;
			
			int K = lg[n] + 1;
			sp.assign(K, vector<long long>(n));
			for (int i = 0; i < n; i++) sp[0][i] = a[i] < 0 ? -a[i] : a[i];
			for (int k = 1; k < K; k++)
			for (int i = 0; i + (1 << k) <= n; i++)
			sp[k][i] = __gcd(sp[k - 1][i], sp[k - 1][i + (1 << (k - 1))]);
		} // O(n log n log M)
		
		// READ: gcd del rango [l, r] (0-indexado, inclusivo).
		long long consultar(int l, int r) const {
			int k = lg[r - l + 1];
			return __gcd(sp[k][l], sp[k][r - (1 << k) + 1]);
		} // O(log M), casi O(1) en la practica
		
		// READ: indica si algun elemento del rango es multiplo de g... (ver modificacion 2).
	};
	
	// ------------------------------------------------------------
	// MODIFICACIONES Y COMENTARIOS CON LOS USOS
	// ------------------------------------------------------------
	// 1. Contar cuantos subarreglos tienen gcd igual a g (o mayor que 1): para cada extremo
	//    derecho r, el gcd de [l, r] solo cambia O(log M) veces al mover l, porque cada cambio
	//    divide el valor por al menos 2. Se hace busqueda binaria sobre l con consultar().
	//    O(n log n log M) en total.
	//
	// 2. Mayor subarreglo cuyo gcd sea > 1: dos punteros o busqueda binaria sobre el largo,
	//    usando consultar() para comprobar cada ventana.
	//
	// 3. Si ademas se necesitan actualizaciones puntuales de gcd: arbol de segmentos con gcd
	//    como combinacion. O(log n + log M) por consulta; cambiar un valor cuesta O(log n log M).
	//
	// 4. Si la operacion es min y gcd a la vez (por ejemplo "gcd y cuantas veces aparece el
	//    minimo"), guardar un par en cada celda y combinar por componentes.
\end{lstlisting}

\textbf{Ejemplo de funcionamiento:}
\begin{lstlisting}[style=cpp]
	int main() {
		vector<long long> a = {12, 18, 30, 45, 15, 9};
		SparseTableGcd st(a);
		
		cout << st.consultar(0, 2) << "\n";   // 6   (gcd de 12, 18, 30)
		cout << st.consultar(2, 4) << "\n";   // 15  (gcd de 30, 45, 15)
		cout << st.consultar(1, 3) << "\n";   // 3   (gcd de 18, 30, 45)
		cout << st.consultar(1, 1) << "\n";   // 18  (rango de un solo elemento)
		cout << st.consultar(0, 5) << "\n";   // 3   (gcd de todos)
		return 0;
	}
\end{lstlisting}

\textbf{Nota:} La línea de la modificación sobre «algún elemento es múltiplo de $g$» del struct es un comentario huérfano en el código de arriba, bórrala. Además, \texttt{\_\_gcd} con negativos puede devolver un resultado negativo, por eso el constructor guarda el valor absoluto. Con ceros todo sigue funcionando porque $\gcd(0,x)=x$.

\subsubsection{Disjoint Sparse Table (operaciones no idempotentes: suma, xor)}

\textbf{Idea:} En vez de dos bloques que se solapan, el nivel $k$ parte el arreglo en bloques de largo $2^{k+1}$ y, para cada bloque, guarda las sumas \textit{desde su punto medio} hacia afuera: hacia la izquierda (de \texttt{mid-1} hacia \texttt{start}) y hacia la derecha (de \texttt{mid} hacia \texttt{end-1}). Una consulta $[l,r]$ con $l<r$ cae en el nivel $k$ donde $l$ y $r$ quedan en mitades distintas del mismo bloque, que es el bit más alto en el que difieren. Entonces $[l,r]$ es la unión disjunta de $[l,\text{mid})$ y $[\text{mid},r]$, y la respuesta es \texttt{pre[k][l] + pre[k][r]}, sin solapes. Como no se necesita idempotencia ni inverso, sirve para cualquier operación asociativa.

\textbf{Complejidad:} Hay $\lceil\log_2 n\rceil+1$ niveles y cada uno recorre las $n$ celdas una vez, así que construir cuesta $O(n\log n)$ en tiempo y memoria. La consulta obtiene el nivel con una instrucción de bits y lee dos celdas: $O(1)$.

\begin{lstlisting}[style=cpp]
	struct DisjointSparseTable {
		int n, niveles;
		vector<long long> a;
		vector<vector<long long>> pre;   // pre[k][i]: suma desde i hasta el punto medio de su bloque
		
		// CREATE: precalcula las sumas hacia el punto medio de cada bloque, por nivel.
		DisjointSparseTable(const vector<long long>& arr) : n((int)arr.size()), a(arr) {
			niveles = 1;
			while ((1 << niveles) < n) niveles++;
			pre.assign(niveles, vector<long long>(n, 0));
			
			for (int k = 0; k < niveles; k++) {
				int bloque = 1 << (k + 1);
				for (int ini = 0; ini < n; ini += bloque) {
					int mid = min(ini + bloque / 2, n);
					int fin = min(ini + bloque, n);
					if (mid - 1 >= ini) {                       // mitad izquierda
						pre[k][mid - 1] = a[mid - 1];
						for (int i = mid - 2; i >= ini; i--) pre[k][i] = a[i] + pre[k][i + 1];
					}
					if (mid < fin) {                            // mitad derecha
						pre[k][mid] = a[mid];
						for (int i = mid + 1; i < fin; i++) pre[k][i] = pre[k][i - 1] + a[i];
					}
				}
			}
		} // O(n log n)
		
		// READ: suma del rango [l, r] (0-indexado, inclusivo).
		long long consultar(int l, int r) const {
			if (l == r) return a[l];
			int k = 31 - __builtin_clz(l ^ r);   // bit mas alto donde difieren l y r
			return pre[k][l] + pre[k][r];
		} // O(1)
	};
	
	// ------------------------------------------------------------
	// MODIFICACIONES Y COMENTARIOS CON LOS USOS
	// ------------------------------------------------------------
	// 1. Cualquier operacion asociativa (xor, producto de matrices, composicion de funciones):
	//    cambiar "+" por la operacion en el constructor y en consultar. Si NO es conmutativa,
	//    el orden importa: izquierda primero, derecha despues, tal como esta escrito.
	//    Para xor:   pre[k][i] = a[i] ^ pre[k][i + 1];   y   return pre[k][l] ^ pre[k][r];
	//
	// 2. Suma de rango con actualizaciones: no sirve (cambiar un elemento invalida O(n) celdas
	//    por nivel). Usar Fenwick o arbol de segmentos.
	//
	// 3. Para suma y xor SIN actualizaciones, un arreglo de prefijos es mas simple y usa O(n)
	//    de memoria: suma = P[r+1] - P[l], xor = P[r+1] ^ P[l]. La Disjoint Sparse Table solo
	//    gana cuando la operacion no tiene inverso (producto de matrices, min con conteo...).
	//
	// 4. Con n potencia de 2 se evita el recorte min(..., n) y el codigo es mas corto.
\end{lstlisting}

\textbf{Ejemplo de funcionamiento:}
\begin{lstlisting}[style=cpp]
	int main() {
		vector<long long> a = {5, 2, 9, 1, 7, 3};
		DisjointSparseTable dst(a);
		
		cout << dst.consultar(0, 2) << "\n";   // 16  (5 + 2 + 9)
		cout << dst.consultar(2, 5) << "\n";   // 20  (9 + 1 + 7 + 3)
		cout << dst.consultar(1, 1) << "\n";   // 2   (un solo elemento)
		cout << dst.consultar(0, 5) << "\n";   // 27  (suma total)
		return 0;
	}
\end{lstlisting}

\textbf{Nota:} Para \texttt{consultar(0, 2)}, $0\oplus 2 = 2$ y su bit más alto es el $1$, así que se usa el nivel $1$, con bloques de largo $4$ y punto medio en $2$: \texttt{pre[1][0]} $=5+2=7$ y \texttt{pre[1][2]} $=9$, total $16$. Una consulta con $l=r$ se responde aparte porque $l\oplus r=0$ no tiene bit alto. En un nivel, la última mitad derecha puede quedar recortada por $n$, pero nunca se consulta una posición inexistente. Por último, los valores guardados son sumas, así que con elementos del orden de $10^{9}$ y $n$ grande conviene \texttt{long long}.

\textbf{Regla práctica:} Si el arreglo no cambia y la operación es idempotente (min, max, gcd, AND, OR), usa Sparse Table: consulta $O(1)$ y es el código más corto. Si la operación es asociativa pero no idempotente ni invertible, usa Disjoint Sparse Table. Si es suma o xor sin actualizaciones, un arreglo de prefijos basta. En cuanto el arreglo cambia, aunque sea una sola vez, pasa a un árbol de segmentos o a un Fenwick, porque reconstruir en $O(n\log n)$ no compensa. Si la memoria $O(n\log n)$ es un problema (n del orden de $10^7$), considera un árbol de segmentos o procesar las consultas por niveles (modificación 7).

\subsubsection{Variante: LCM en un Rango}

\textbf{Idea:} El mcm es idempotente, $\operatorname{lcm}(x,x)=x$, además de asociativo y conmutativo, así que encaja en la Sparse Table igual que el mínimo o el gcd: el solape de los dos bloques no altera el resultado, porque el mcm de una unión de elementos no cambia si uno se cuenta dos veces. Se combina con $\operatorname{lcm}(x,y)=\dfrac{x}{\gcd(x,y)}\cdot y$, dividiendo \textit{antes} de multiplicar para que el intermedio sea lo más pequeño posible.

La diferencia con el gcd es que el mcm \textit{crece}: el mcm de un bloque es múltiplo del de cualquiera de sus sub-bloques, y puede superar enseguida cualquier tipo entero (el mcm de $1,\dots,43$ ya pasa de $10^{18}$). La solución es \textit{saturar}: se fija un límite \texttt{LIM} (aquí $10^{18}$) y todo mcm mayor que \texttt{LIM} se guarda como un valor centinela \texttt{INF} $=\texttt{LIM}+1$. Esto es correcto porque el mcm es monótono: si un sub-bloque ya excede \texttt{LIM}, cualquier bloque que lo contenga también lo excede, así que combinar con \texttt{INF} da \texttt{INF} sin necesidad de conocer el valor real. El producto intermedio $(x/g)\cdot y$ puede llegar a $10^{36}$, que cabe en \texttt{\_\_int128} ($\approx 1.7\cdot 10^{38}$) pero no en \texttt{long long}. Se asume que todos los valores son al menos $1$ (el mcm con $0$ no está bien definido).

\textbf{Complejidad:} Cada combinación es un gcd de $O(\log M)$ más una multiplicación y una comparación en $O(1)$, así que construir cuesta $O(n\log n\log M)$ y cada consulta $O(\log M)$, como en el gcd. Memoria: $O(n\log n)$.

\begin{lstlisting}[style=cpp]
	// Valores >= 1. Si el mcm de un rango supera LIM, la consulta devuelve INF (= LIM + 1).
	struct SparseTableLcm {
		static constexpr long long LIM = 1000000000000000000LL;   // 10^18
		static constexpr long long INF = LIM + 1;                  // centinela: "mayor que LIM"
		int n;
		vector<vector<long long>> sp;   // sp[k][i] = mcm saturado del bloque [i, i + 2^k - 1]
		vector<int> lg;
		
		// READ: mcm saturado de dos valores en [1, INF].
		static long long mcm(long long x, long long y) {
			if (x > LIM || y > LIM) return INF;          // ya excede: cualquier superconjunto tambien
			long long g = __gcd(x, y);
			__int128 r = (__int128)(x / g) * y;          // dividir antes de multiplicar
			return r > LIM ? INF : (long long)r;
		} // O(log M)
		
		// CREATE: precalcula la tabla de mcm a partir de a[0..n-1] (valores >= 1).
		SparseTableLcm(const vector<long long>& a) {
			n = (int)a.size();
			lg.assign(n + 1, 0);
			for (int i = 2; i <= n; i++) lg[i] = lg[i / 2] + 1;
			
			int K = lg[n] + 1;
			sp.assign(K, vector<long long>(n));
			for (int i = 0; i < n; i++) sp[0][i] = min(a[i], INF);
			for (int k = 1; k < K; k++)
			for (int i = 0; i + (1 << k) <= n; i++)
			sp[k][i] = mcm(sp[k - 1][i], sp[k - 1][i + (1 << (k - 1))]);
		} // O(n log n log M)
		
		// READ: mcm del rango [l, r] (0-indexado, inclusivo), o INF si supera LIM.
		long long consultar(int l, int r) const {
			int k = lg[r - l + 1];
			return mcm(sp[k][l], sp[k][r - (1 << k) + 1]);
		} // O(log M)
		
		// READ: indica si el mcm del rango es exacto (no saturado).
		bool exacto(int l, int r) const {
			return consultar(l, r) <= LIM;
		} // O(log M)
	};
	
	// ------------------------------------------------------------
	// MODIFICACIONES Y COMENTARIOS CON LOS USOS
	// ------------------------------------------------------------
	// 1. Otro limite: el problema suele fijar uno (por ejemplo 10^9 o 2^63 - 1). Cambiar LIM.
	//    Con LIM <= 3 * 10^9 el producto cabe en unsigned long long y se puede evitar __int128;
	//    con LIM = 10^18 es obligatorio (o una comprobacion previa: x / g > LIM / y => INF).
	//
	// 2. mcm MODULO un primo p: NO se puede combinar residuos (lcm(x mod p, y mod p) no es
	//    lcm(x, y) mod p). Hay que factorizar y guardar, por cada primo, el MAXIMO exponente
	//    del rango. Con pocos primos distintos (por ejemplo los primos <= 50, que son 15), se
	//    hacen 15 Sparse Tables de maximo sobre los exponentes y la consulta es el producto de
	//    p^e mod M sobre los primos: O(#primos * log e) por consulta. Con muchos primos grandes,
	//    es mas practico un arbol de segmentos con mapas de exponentes.
	//
	// 3. Cuantos subarreglos tienen mcm <= LIM (o mcm igual a v): para cada extremo izquierdo l,
	//    el mcm de [l, r] solo cambia O(log LIM) veces al avanzar r (cada cambio lo multiplica
	//    por al menos 2), asi que se hace busqueda binaria sobre r con consultar().
	//    O(n log LIM * log M) en total.
	//
	// 4. Actualizaciones puntuales: arbol de segmentos con mcm saturado como combinacion
	//    (misma funcion mcm). O(log n * log M) por operacion.
	//
	// 5. Guardar el mcm sin saturar: usar BigInt (ver "Numeros Grandes") como tipo de cada celda.
	//    Cada combinacion se vuelve mucho mas cara y requiere division de BigInt, asi que solo
	//    se justifica si el enunciado pide el valor exacto de un mcm enorme.
\end{lstlisting}

\textbf{Ejemplo de funcionamiento:}
\begin{lstlisting}[style=cpp]
	int main() {
		vector<long long> a = {4, 6, 10, 3, 8, 15};
		SparseTableLcm st(a);
		
		cout << st.consultar(0, 2) << "\n";   // 60   (mcm de 4, 6, 10)
		cout << st.consultar(1, 3) << "\n";   // 30   (mcm de 6, 10, 3)
		cout << st.consultar(2, 4) << "\n";   // 120  (mcm de 10, 3, 8)
		cout << st.consultar(1, 1) << "\n";   // 6    (un solo elemento)
		cout << st.consultar(0, 5) << "\n";   // 120  (mcm de todos: 2^3 * 3 * 5)
		
		// Saturacion: tres primos cercanos a 10^9, coprimos entre si.
		vector<long long> b = {1000000007LL, 998244353LL, 1000000009LL};
		SparseTableLcm sb(b);
		cout << sb.consultar(0, 1) << "\n";   // 998244359987710471  (producto de los dos, < 10^18)
		cout << sb.consultar(0, 2) << "\n";   // 1000000000000000001 (INF: el mcm real, ~10^27, excede LIM)
		cout << sb.exacto(0, 2) << "\n";      // 0
		return 0;
	}
\end{lstlisting}

\textbf{Nota:} Un resultado igual a \texttt{INF} significa únicamente «el mcm es mayor que \texttt{LIM}», no un valor concreto, así que no se puede usar para nada más que compararlo con un tope. Con ese tope el mcm tiene además una propiedad útil: al ir ampliando un rango, el mcm solo puede cambiar unas $\log_2(\texttt{LIM})\approx 60$ veces antes de saturarse, y por eso las consultas «cuántos subarreglos tienen mcm $\le$ tope» son baratas (modificación 3). Si algún elemento puede ser $0$, hay que tratarlo aparte antes de construir la tabla, porque \texttt{x / g} dividiría por cero cuando ambos valores son $0$.


\subsection{Segment Tree}

\textbf{Idea:} un árbol binario implícito (guardado en un arreglo)
donde cada nodo representa el resultado combinado de un rango del
arreglo original. La raíz cubre todo $[0,n-1]$; cada nodo se parte en
dos mitades hasta llegar a las hojas, que son elementos individuales.
Un update baja por el árbol tocando $O(\log n)$ nodos (uno por nivel);
una query de rango se descompone en a lo más $O(\log n)$ nodos que
cubren exactamente el rango pedido, sin necesidad de bloques
parciales como en Sqrt Decomposition. A diferencia de Fenwick, sirve
para \textbf{cualquier} operación asociativa (min, max, gcd, and, or,
no solo sumas), y con Lazy Propagation soporta actualizaciones de
rango completo.

\textbf{¿Por qué usarlo?} Es la estructura más versátil de esta
sección: $O(\log n)$ tanto para update como para query, funciona con
casi cualquier operación asociativa, y con lazy propagation resuelve
range update + range query para operaciones que Fenwick no puede
(min, max, asignar un valor a todo un rango, etc.).

\textbf{Casos de uso:}

\begin{itemize}
    \item Min/max/gcd de rango con actualizaciones (a diferencia de
    Sparse Table, que exige arreglo estático).
    \item Range update + range query con operaciones no invertibles:
    asignar un valor a todo un rango, min/max con suma en rango, etc.
    \item Cualquier operación asociativa que no sea una suma simple
    (donde Fenwick ya alcanzaría y sería más simple).
    \item Base para variantes más avanzadas: Segment Tree 2D, Merge
    Sort Tree, Persistent Segment Tree, Segment Tree Beats.
\end{itemize}

\textbf{Complejidad:} construir $O(n)$, update puntual $O(\log n)$,
query de rango $O(\log n)$, range update + range query con lazy
propagation $O(\log n)$.


\subsubsection{Segment Tree clásico (point update, range query — suma)}

\begin{lstlisting}[style=cpp]
struct SegmentTree {
    vector<long long> tree, a;
    int n;

    // CREATE: reserva el arbol (tamano 4n es cota segura) y
    // construye recursivamente a partir de a[0..n-1].
    SegmentTree(vector<long long>& arr) {
        a = arr;
        n = a.size();
        tree.assign(4 * n, 0);
        if (n > 0)
            build(1, 0, n - 1);
    }                                          // O(n)

    void build(int node, int l, int r) {
        if (l == r) {
            tree[node] = a[l];
            return;
        }

        int mid = (l + r) / 2;
        build(2 * node, l, mid);
        build(2 * node + 1, mid + 1, r);
        tree[node] = tree[2 * node] + tree[2 * node + 1];
    }

    void updateUtil(int node, int l, int r, int i, long long x) {
        if (l == r) {
            tree[node] = x;
            return;
        }

        int mid = (l + r) / 2;
        if (i <= mid)
            updateUtil(2 * node, l, mid, i, x);
        else
            updateUtil(2 * node + 1, mid + 1, r, i, x);

        tree[node] = tree[2 * node] + tree[2 * node + 1];
    }

    // UPDATE: fija el valor absoluto de la posicion i en x.
    void update(int i, long long x) {
        updateUtil(1, 0, n - 1, i, x);
    }                                          // O(log n)

    long long queryUtil(int node, int l, int r, int ql, int qr) {
        if (qr < l || r < ql)
            return 0;                          // neutro para suma

        if (ql <= l && r <= qr)
            return tree[node];

        int mid = (l + r) / 2;
        return queryUtil(2 * node, l, mid, ql, qr)
             + queryUtil(2 * node + 1, mid + 1, r, ql, qr);
    }

    // READ: suma del rango [l, r] (0-indexado, inclusivo).
    long long rango(int l, int r) {
        return queryUtil(1, 0, n - 1, l, r);
    }                                          // O(log n)
};
\end{lstlisting}

\textbf{Ejemplo de funcionamiento:}

\begin{lstlisting}[style=cpp]
vector<long long> a = {5, 2, 9, 1, 7};

SegmentTree st(a);
// arreglo logico: [5, 2, 9, 1, 7]  (posiciones 0..4)

cout << st.rango(1, 3);
// 2 + 9 + 1 = 12

st.update(2, 100);
// arreglo logico: [5, 2, 100, 1, 7]

cout << st.rango(1, 3);
// 2 + 100 + 1 = 103
\end{lstlisting}

\textbf{Nota:} para min/max/gcd solo cambia el neutro del caso base
(\texttt{LLONG\_MAX} para min, \texttt{0} para gcd, etc.) y el
operador de combinación (\texttt{min(...)}, \texttt{\_\_gcd(...)}) en
\texttt{build}, \texttt{updateUtil} y \texttt{queryUtil}.


\subsubsection{Segment Tree con Lazy Propagation (range update + range query — suma)}

\textbf{Idea:} cada nodo guarda un valor \texttt{lazy} pendiente que
representa ``sumar $v$ a todo este rango'' sin aplicarlo todavía a
los hijos. Se aplica (\emph{se propaga}) solo cuando hace falta bajar
más en el árbol para otra operación.

\begin{lstlisting}[style=cpp]
struct SegmentTreeLazyAdd {
    vector<long long> tree, lazy, a;
    int n;

    // CREATE: reserva arbol + lazy, construye a partir de a[0..n-1].
    SegmentTreeLazyAdd(vector<long long>& arr) {
        a = arr;
        n = a.size();
        tree.assign(4 * n, 0);
        lazy.assign(4 * n, 0);
        if (n > 0)
            build(1, 0, n - 1);
    }                                          // O(n)

    void build(int node, int l, int r) {
        if (l == r) {
            tree[node] = a[l];
            return;
        }

        int mid = (l + r) / 2;
        build(2 * node, l, mid);
        build(2 * node + 1, mid + 1, r);
        tree[node] = tree[2 * node] + tree[2 * node + 1];
    }

    // Aplica el lazy pendiente de este nodo y lo empuja a los hijos.
    void propagar(int node, int l, int r) {
        if (lazy[node] == 0)
            return;

        tree[node] += lazy[node] * (r - l + 1);

        if (l != r) {
            lazy[2 * node] += lazy[node];
            lazy[2 * node + 1] += lazy[node];
        }

        lazy[node] = 0;
    }

    void rangoUpdateUtil(int node, int l, int r, int ql, int qr,
                          long long v) {
        propagar(node, l, r);

        if (qr < l || r < ql)
            return;

        if (ql <= l && r <= qr) {
            lazy[node] += v;
            propagar(node, l, r);
            return;
        }

        int mid = (l + r) / 2;
        rangoUpdateUtil(2 * node, l, mid, ql, qr, v);
        rangoUpdateUtil(2 * node + 1, mid + 1, r, ql, qr, v);
        tree[node] = tree[2 * node] + tree[2 * node + 1];
    }

    // UPDATE: suma v a todo el rango [l, r].
    void rangoUpdate(int l, int r, long long v) {
        rangoUpdateUtil(1, 0, n - 1, l, r, v);
    }                                          // O(log n)

    long long queryUtil(int node, int l, int r, int ql, int qr) {
        propagar(node, l, r);

        if (qr < l || r < ql)
            return 0;

        if (ql <= l && r <= qr)
            return tree[node];

        int mid = (l + r) / 2;
        return queryUtil(2 * node, l, mid, ql, qr)
             + queryUtil(2 * node + 1, mid + 1, r, ql, qr);
    }

    // READ: suma del rango [l, r].
    long long rango(int l, int r) {
        return queryUtil(1, 0, n - 1, l, r);
    }                                          // O(log n)
};
\end{lstlisting}

\textbf{Ejemplo de funcionamiento:}

\begin{lstlisting}[style=cpp]
vector<long long> a = {1, 1, 1, 1, 1, 1};

SegmentTreeLazyAdd st(a);

st.rangoUpdate(1, 4, 10);
// arreglo logico efectivo: [1, 11, 11, 11, 11, 1]

cout << st.rango(0, 5);
// 1+11+11+11+11+1 = 46

cout << st.rango(2, 3);
// 11+11 = 22
\end{lstlisting}


\subsubsection{Lazy Propagation con asignación (range assign + range query — mínimo)}

\textbf{Idea:} en vez de \emph{sumar} un valor al rango, lo
\textbf{reemplaza} por completo (\texttt{asigna x a todo
$[l,r]$}). El lazy aquí no se acumula como en la suma: el último
assign pisa a cualquier assign anterior pendiente, así que se necesita
una bandera aparte para distinguir ``no hay lazy pendiente'' de
``hay que asignar 0''.

\begin{lstlisting}[style=cpp]
struct SegmentTreeAssignMin {
    vector<long long> tree, lazy, a;
    vector<bool> tieneLazy;
    int n;

    // CREATE: reserva arbol + lazy + banderas, construye desde a.
    SegmentTreeAssignMin(vector<long long>& arr) {
        a = arr;
        n = a.size();
        tree.assign(4 * n, LLONG_MAX);
        lazy.assign(4 * n, 0);
        tieneLazy.assign(4 * n, false);
        if (n > 0)
            build(1, 0, n - 1);
    }                                          // O(n)

    void build(int node, int l, int r) {
        if (l == r) {
            tree[node] = a[l];
            return;
        }

        int mid = (l + r) / 2;
        build(2 * node, l, mid);
        build(2 * node + 1, mid + 1, r);
        tree[node] = min(tree[2 * node], tree[2 * node + 1]);
    }

    // Empuja el assign pendiente de este nodo hacia sus hijos.
    void propagar(int node) {
        if (!tieneLazy[node])
            return;

        for (int hijo : {2 * node, 2 * node + 1}) {
            tree[hijo] = lazy[node];
            lazy[hijo] = lazy[node];
            tieneLazy[hijo] = true;
        }

        tieneLazy[node] = false;
    }

    void rangoAssignUtil(int node, int l, int r, int ql, int qr,
                          long long x) {
        if (qr < l || r < ql)
            return;

        if (ql <= l && r <= qr) {
            tree[node] = x;
            lazy[node] = x;
            tieneLazy[node] = true;
            return;
        }

        propagar(node);
        int mid = (l + r) / 2;
        rangoAssignUtil(2 * node, l, mid, ql, qr, x);
        rangoAssignUtil(2 * node + 1, mid + 1, r, ql, qr, x);
        tree[node] = min(tree[2 * node], tree[2 * node + 1]);
    }

    // UPDATE: asigna x a todo el rango [l, r].
    void rangoAssign(int l, int r, long long x) {
        rangoAssignUtil(1, 0, n - 1, l, r, x);
    }                                          // O(log n)

    long long queryUtil(int node, int l, int r, int ql, int qr) {
        if (qr < l || r < ql)
            return LLONG_MAX;

        if (ql <= l && r <= qr)
            return tree[node];

        propagar(node);
        int mid = (l + r) / 2;
        return min(queryUtil(2 * node, l, mid, ql, qr),
                    queryUtil(2 * node + 1, mid + 1, r, ql, qr));
    }

    // READ: minimo del rango [l, r].
    long long rango(int l, int r) {
        return queryUtil(1, 0, n - 1, l, r);
    }                                          // O(log n)
};
\end{lstlisting}

\textbf{Ejemplo de funcionamiento:}

\begin{lstlisting}[style=cpp]
vector<long long> a = {5, 2, 9, 1, 7, 3};

SegmentTreeAssignMin st(a);

cout << st.rango(0, 5);
// min(5,2,9,1,7,3) = 1

st.rangoAssign(0, 2, 100);
// arreglo logico efectivo: [100, 100, 100, 1, 7, 3]

cout << st.rango(0, 2);
// min(100,100,100) = 100

cout << st.rango(0, 5);
// min(100,100,100,1,7,3) = 1
\end{lstlisting}

\textbf{Comparación con las otras estructuras de rango:}

\begin{table}[H]
\raggedright
\scriptsize
\setlength{\tabcolsep}{3pt}
\begin{tabular}{|l|c|c|c|}
\hline
\textbf{Estructura} & \textbf{Update} & \textbf{Query} & \textbf{Min/Max?} \\
\hline
Segment Tree & $O(\log n)$ & $O(\log n)$ & Sí \\
\hline
Segment Tree + Lazy & $O(\log n)$ range & $O(\log n)$ range & Sí \\
\hline
\texttt{FenwickTree} & $O(\log n)$ punto & $O(\log n)$ rango & No \\
\hline
Sqrt Decomposition & $O(1)$–$O(\sqrt n)$ & $O(\sqrt n)$ & Sí \\
\hline
\end{tabular}
\caption{Segment Tree vs. Fenwick vs. Sqrt Decomposition.}
\end{table}

\textbf{Regla práctica:} usa Segment Tree cuando necesites min/max/gcd
(operaciones no invertibles que Fenwick no soporta), o cuando
necesites range update + range query con una operación que no sea
suma simple. Si la operación es solo suma y no necesitas range
update, Fenwick es más simple y más rápido en la práctica (menor
constante). Si no confías en escribir el lazy propagation correcto
bajo presión, Sqrt Decomposition es una alternativa más simple a
costa de $O(\sqrt n)$ en vez de $O(\log n)$.


\subsection{RMQ (Range Minimum Query)}

\textbf{Idea:} RMQ es el problema general de "dado un arreglo,
encuentra el mínimo en cualquier rango $[l,r]$". Ya tienes las dos
estructuras que lo resuelven: \textbf{Sparse Table} si el arreglo es
estático ($O(1)$ por query) y \textbf{Segment Tree} si cambia
($O(\log n)$ por query). Esta sección no repite esas estructuras,
sino que muestra las dos aplicaciones clásicas donde un problema que
\emph{no parece} de rango se puede \textbf{reformular} como RMQ sobre
un arreglo auxiliar: LCA en árboles, y LCP entre sufijos arbitrarios
de un string.

\textbf{¿Por qué usarlo?} Reconocer cuándo un problema es "RMQ
disfrazado" es lo que realmente separa resolver estos problemas en
$O(1)$–$O(\log n)$ de resolverlos con fuerza bruta $O(n)$ por
consulta. El patrón se repite: construyes un arreglo auxiliar donde
la respuesta a tu pregunta real equivale al mínimo de un rango, y le
aplicas Sparse Table o Segment Tree normal.

\textbf{Casos de uso:}

\begin{itemize}
    \item LCA (ancestro común más bajo) de dos nodos en un árbol, con
    muchas queries: Euler Tour + RMQ.
    \item LCP (prefijo común más largo) entre dos sufijos
    \emph{cualquiera} de un string (no necesariamente consecutivos en
    el Suffix Array): Suffix Array + Kasai + RMQ sobre el LCP Array.
    \item Cualquier problema donde la respuesta se pueda expresar como
    "el mínimo de algo" sobre un rango de índices auxiliares.
\end{itemize}

\textbf{Complejidad:} en ambas aplicaciones, el precómputo (Euler
Tour o Suffix Array + Kasai) es $O(n)$ a $O(n\log^2 n)$ según el
método de Suffix Array usado; una vez armado el arreglo auxiliar, el
RMQ es $O(n\log n)$ de build y $O(1)$ por query con Sparse Table
(ambos casos son estáticos: el árbol y el string no cambian).


\subsubsection{LCA vía Euler Tour + RMQ (Sparse Table)}

\textbf{Idea:} recorres el árbol con un Euler Tour, anotando el nodo
visitado cada vez que \textbf{entras o vuelves} a él (arreglo de
tamaño $2n-1$) junto con su profundidad. El LCA de $u$ y $v$ es el
nodo de \textbf{menor profundidad} que aparece entre la primera
aparición de $u$ y la primera aparición de $v$ en ese recorrido —
exactamente un RMQ (por profundidad) sobre el arreglo del Euler Tour.

\begin{lstlisting}[style=cpp]
struct LCA {
    vector<int> euler, primeraApar, profundidad;
    vector<vector<int>> sp;   // sp[k][i] = nodo con menor profundidad
    vector<int> lg;

    // CREATE: construye el Euler Tour del arbol (arr. de tamano
    // 2n-1) y la Sparse Table sobre profundidad, para RMQ O(1).
    LCA(vector<vector<int>>& adj, int raiz, int n) {
        primeraApar.assign(n, -1);
        profundidad.assign(n, 0);
        euler.reserve(2 * n);

        dfs(adj, raiz, -1, 0);
        build();
    }                                          // O(n log n)

    void dfs(vector<vector<int>>& adj, int u, int p, int d) {
        primeraApar[u] = euler.size();
        profundidad[u] = d;
        euler.push_back(u);

        for (int v : adj[u]) {
            if (v != p) {
                dfs(adj, v, u, d + 1);
                euler.push_back(u);
            }
        }
    }

    void build() {
        int m = euler.size();
        lg.assign(m + 1, 0);
        for (int i = 2; i <= m; i++)
            lg[i] = lg[i / 2] + 1;

        int K = lg[m] + 1;
        sp.assign(K, vector<int>(m));
        sp[0] = euler;

        for (int k = 1; k < K; k++)
            for (int i = 0; i + (1 << k) <= m; i++) {
                int l = sp[k - 1][i];
                int r = sp[k - 1][i + (1 << (k - 1))];
                sp[k][i] = profundidad[l] < profundidad[r] ? l : r;
            }
    }

    // READ: LCA de u y v.
    int query(int u, int v) {
        int l = primeraApar[u], r = primeraApar[v];
        if (l > r) swap(l, r);

        int k = lg[r - l + 1];
        int a = sp[k][l], b = sp[k][r - (1 << k) + 1];
        return profundidad[a] < profundidad[b] ? a : b;
    }                                          // O(1)
};
\end{lstlisting}

\textbf{Ejemplo de funcionamiento:}

\begin{lstlisting}[style=cpp]
// arbol:      0
//            / \
//           1   2
//          / \
//         3   4

vector<vector<int>> adj(5);
adj[0] = {1, 2};
adj[1] = {0, 3, 4};
adj[2] = {0};
adj[3] = {1};
adj[4] = {1};

LCA lca(adj, 0, 5);

cout << lca.query(3, 4);
// 1  (padre comun de 3 y 4)

cout << lca.query(3, 2);
// 0  (raiz)

cout << lca.query(1, 4);
// 1  (uno es ancestro del otro)
\end{lstlisting}


\subsubsection{LCP entre dos sufijos arbitrarios (Suffix Array + Kasai + RMQ)}

\textbf{Idea:} el LCP Array (via Kasai) solo te da el LCP entre
sufijos \textbf{consecutivos} en el Suffix Array. Para el LCP entre
dos sufijos \textbf{cualquiera} $i,j$: los ubicas por su rango en el
SA ($\texttt{rnk}[i]$, $\texttt{rnk}[j]$) y el LCP entre ambos es el
\textbf{mínimo} del LCP Array en el rango entre esas dos posiciones
—otra vez RMQ, esta vez sobre el LCP Array.

\begin{lstlisting}[style=cpp]
struct SuffixLCPQuery {
    vector<int> sa, rnk, lcp;
    vector<vector<int>> sp;
    vector<int> lg;
    int n;

    // CREATE: construye SA, LCP (Kasai) y Sparse Table sobre lcp,
    // todo en un solo paso.
    SuffixLCPQuery(const string& s) {
        n = s.size();
        sa = buildSuffixArray(s);             // O(n log^2 n)
        lcp = buildLCPArray(s, sa);           // O(n), Kasai

        rnk.assign(n, 0);
        for (int i = 0; i < n; i++)
            rnk[sa[i]] = i;

        buildSparse();
    }                                          // O(n log^2 n)

    void buildSparse() {
        lg.assign(n + 1, 0);
        for (int i = 2; i <= n; i++)
            lg[i] = lg[i / 2] + 1;

        int K = lg[n] + 1;
        sp.assign(K, vector<int>(n));
        sp[0] = lcp;

        for (int k = 1; k < K; k++)
            for (int i = 0; i + (1 << k) <= n; i++)
                sp[k][i] = min(sp[k - 1][i],
                                sp[k - 1][i + (1 << (k - 1))]);
    }

    // READ: LCP entre el sufijo que empieza en i y el que
    // empieza en j (0-indexado sobre el string original).
    int lcpEntreSufijos(int i, int j) {
        if (i == j) return n - i;

        int a = rnk[i], b = rnk[j];
        if (a > b) swap(a, b);

        // el minimo va de [a+1, b] en el LCP array
        int k = lg[b - a];
        return min(sp[k][a + 1], sp[k][b - (1 << k) + 1]);
    }                                          // O(1)
};
\end{lstlisting}

\textbf{Ejemplo de funcionamiento:}

\begin{lstlisting}[style=cpp]
SuffixLCPQuery slq("banana");

cout << slq.lcpEntreSufijos(1, 3);
// sufijo desde 1: "anana"
// sufijo desde 3: "ana"
// LCP = "ana" -> 3

cout << slq.lcpEntreSufijos(0, 3);
// sufijo desde 0: "banana"
// sufijo desde 3: "ana"
// LCP = "" -> 0
\end{lstlisting}

\textbf{Nota:} \texttt{buildSuffixArray} y \texttt{buildLCPArray} son
las funciones estándar de Suffix Array + Kasai (ver la sección de
Algoritmos de Strings); aquí solo se reutilizan como precómputo antes
de aplicar RMQ.

\textbf{Regla práctica:} cuando un problema pida "algo" entre dos
elementos que no son adyacentes en su estructura original (dos nodos
de un árbol, dos sufijos de un string, dos elementos de cualquier
secuencia con una noción de "orden" y "profundidad/afinidad"),
pregúntate si esa relación se puede expresar como el mínimo de un
arreglo auxiliar construido con un recorrido lineal. Si sí, el
problema se reduce a Sparse Table (estático) o Segment Tree
(dinámico), ya cubiertos en las secciones anteriores.


\subsection{Mo's Algorithm}

\textbf{Idea:} Mo's responde $q$ consultas de rango \textit{fuera de línea} (todas conocidas de antemano) manteniendo una ventana $[izq, der]$ sobre el arreglo, junto con la respuesta de esa ventana. Para pasar de una consulta a la siguiente no se recalcula nada: se mueven los extremos con \texttt{agregar} y \texttt{quitar}, cada una en $O(1)$. Recalcular cada consulta desde cero cuesta $O(nq)$. El costo total es la suma de lo que se desplazan los dos extremos, y eso depende del \textit{orden} en que se atienden las consultas.

El orden se elige así. Se divide el arreglo en bloques de tamaño $B$ y las consultas se agrupan por el bloque de su extremo izquierdo $l$. Dentro de un bloque se ordenan por $r$ creciente. Con ese orden, el puntero $der$ recorre el arreglo una sola vez por bloque, con costo $O(n)$ por bloque y $O(n^2/B)$ en total. El puntero $izq$ nunca sale de su bloque, así que se mueve a lo sumo $B$ posiciones por consulta, $O(qB)$ en total. El costo es $O(n^2/B + qB)$, que se minimiza con $B = n/\sqrt q$ y da $O(n\sqrt q)$. Con $B=\sqrt n$ queda $O((n+q)\sqrt n)$. Una mejora práctica es alternar el sentido de $r$ en los bloques pares e impares: así el puntero $der$ no vuelve al inicio al cambiar de bloque.

Si además hay actualizaciones puntuales entre las consultas, cada consulta lleva un tercer número: el tiempo $t$ (cuántas actualizaciones ya se aplicaron). Se ordena por $(\text{bloque}(l), \text{bloque}(r), t)$ y se agrega un tercer puntero $curT$ que aplica o revierte actualizaciones una por una. Con bloques de tamaño $B\approx n^{2/3}$ hay $(n/B)^2 = n^{2/3}$ pares de bloques, y en cada uno $curT$ recorre a lo sumo las $U$ actualizaciones, lo que da $O(n^{2/3}U)$. Los extremos $l$ y $r$ suman $O(qB)$. Con $q\approx U\approx n$ el total es $O(n^{5/3})$.

\textbf{¿Por qué usarlo?} Mo's sirve cuando la respuesta de un rango es fácil de \textit{mantener} al agregar o quitar un elemento, pero difícil de \textit{combinar} a partir de dos mitades. Ejemplos son contar valores distintos, sumar cuadrados de frecuencias o contar pares con cierta propiedad. Un árbol de segmentos exige una operación de combinación asociativa, y para estos problemas no existe una obvia. Si la respuesta sí se combina (suma, mínimo, máximo, mcd), el árbol de segmentos o el Fenwick son más simples y más rápidos, con $O(\log n)$ por consulta y sin exigir que las consultas sean offline. La descomposición en raíz (\textit{sqrt decomposition}) sirve en línea, pero cada bloque necesita una estructura propia que casi siempre es más difícil de escribir que \texttt{agregar}/\texttt{quitar}. Mo's, en cambio, es una plantilla fija en la que solo se cambian esas dos funciones. Su costo es que es offline y que $\sqrt n$ es más lento que $\log n$. Además, si los valores son grandes hay que comprimirlos antes (subsección de compresión de coordenadas), porque el arreglo de frecuencias se indexa por valor.

\textbf{Casos de uso:}
\begin{itemize}
	\item Contar valores distintos en muchos rangos.
	\item Sumar $\text{frecuencia}^2 \cdot \text{valor}$ o cualquier expresión que dependa de las frecuencias dentro del rango.
	\item Frecuencia de un valor y frecuencia de la moda en un rango (variante más abajo).
	\item Contar pares de posiciones dentro del rango con una propiedad (por ejemplo, XOR igual a $k$ en subarreglos).
	\item De nicho: consultas sobre caminos o subárboles con Mo's sobre árboles (modificación 3), o con actualizaciones puntuales intercaladas (Mo's con actualizaciones, más abajo).
\end{itemize}

\textbf{Complejidad:} Sean $n$ el largo del arreglo y $q$ el número de consultas. Ordenar las consultas cuesta $O(q\log q)$. Cada consulta desplaza los punteros con \texttt{agregar}/\texttt{quitar} de costo $O(1)$, y el desplazamiento total es $O(n^2/B + qB)$ como se vio arriba, con lo que el tiempo queda en $O(n\sqrt q + q\log q)$. Con $q\approx n$ es $O(n\sqrt n)$. La memoria es $O(n + q)$ más el arreglo de frecuencias. Si \texttt{agregar} y \texttt{quitar} cuestan $T$ en lugar de $O(1)$, el total se multiplica por $T$. La variante con actualizaciones cuesta $O(n^{5/3})$ con $q\approx U\approx n$ y bloques de tamaño $n^{2/3}$.

\subsubsection{Mo's clásico (consultas offline, sin actualizaciones)}

\begin{lstlisting}[style=cpp]
	// Ejemplo guia: contar valores DISTINTOS en cada rango [l, r] (0-indexado, inclusive).
	struct MoClasico {
		int n, bloque;
		vector<int> a;        // arreglo original (valores en [0, max_val])
		vector<int> cuenta;   // cuenta[v]: cuantas veces aparece v en la ventana actual
		int izq = 0, der = -1;
		int distintos = 0;    // respuesta de la ventana actual
		
		// CREATE: guarda el arreglo y reserva el arreglo de frecuencias.
		MoClasico(const vector<int>& arr, int max_val) {
			a = arr;
			n = (int)a.size();
			cuenta.assign(max_val + 1, 0);
			bloque = max(1, (int)sqrt((double)n));
		} // O(n + max_val)
		
		// WRITE: agrega el elemento de la posicion pos a la ventana.
		void agregar(int pos) {
			if (cuenta[a[pos]]++ == 0) distintos++;
		} // O(1)
		
		// WRITE: quita el elemento de la posicion pos de la ventana.
		void quitar(int pos) {
			if (--cuenta[a[pos]] == 0) distintos--;
		} // O(1)
		
		// WRITE: responde todas las consultas y las devuelve en el orden original.
		// Modifica la ventana (izq, der), por eso no es const.
		vector<int> resolver(const vector<array<int, 2>>& consultas) {
			int q = (int)consultas.size();
			bloque = max(1, (int)(n / sqrt((double)max(1, q))));   // B = n / sqrt(q)
			vector<int> orden(q);
			iota(orden.begin(), orden.end(), 0);
			sort(orden.begin(), orden.end(), [&](int i, int j) {
				int bi = consultas[i][0] / bloque, bj = consultas[j][0] / bloque;
				if (bi != bj) return bi < bj;
				return (bi & 1) ? consultas[i][1] > consultas[j][1]   // bloques impares: r decreciente
				: consultas[i][1] < consultas[j][1];  // bloques pares: r creciente
			});
			
			vector<int> resp(q);
			for (int id : orden) {
				int l = consultas[id][0], r = consultas[id][1];
				while (der < r) agregar(++der);   // primero expandir...
				while (izq > l) agregar(--izq);
				while (der > r) quitar(der--);    // ...y despues contraer
				while (izq < l) quitar(izq++);
				resp[id] = distintos;
			}
			return resp;
		} // O(n sqrt(q) + q log q)
	};
	
	// ------------------------------------------------------------
	// MODIFICACIONES Y COMENTARIOS CON LOS USOS
	// ------------------------------------------------------------
	// 1. Frecuencia y moda en un rango: ver la variante "Frecuencia y Moda en un Rango"
	//    (misma plantilla, pero la ventana mantiene cuenta[], cc[] y max_freq en lugar de
	//    "distintos").
	//
	// 2. Ordenamiento por curva de Hilbert: en lugar de bloques + zigzag, ordenar las
	//    consultas por el indice de Hilbert del punto (l, r). Tiene el mismo costo asintotico
	//    y suele ser una constante mejor en la practica, pero es mas dificil de escribir.
	//
	// 3. Mo's sobre arboles (consultas sobre un camino o un subarbol): aplanar el arbol con un
	//    recorrido de Euler (cada nodo aparece dos veces: al entrar y al salir) y aplicar el
	//    mismo algoritmo sobre ese arreglo. En agregar/quitar se alterna el estado del nodo
	//    ("dentro/fuera de la ventana"): si aparece por segunda vez, se quita. Para caminos
	//    hace falta ademas tratar el ancestro comun (LCA) por separado.
	//
	// 4. Mo's con retroceso (rollback): cuando quitar es dificil o imposible (moda, maximo),
	//    se procesa cada bloque con un solo agregar: para cada consulta se agrega desde el
	//    final del bloque y se deshace lo agregado del lado izquierdo. Mismo O(n sqrt(q)) y
	//    solo necesita agregar y una copia del estado.
	//
	// 5. Valores grandes o negativos: comprimirlos primero (ver "Compresion de Coordenadas"),
	//    para que cuenta[] tenga tamano O(n) en lugar de O(max_val).
	//
	// 6. agregar/quitar costosos: si cuestan T cada una, el total pasa a
	//    O(T * n sqrt(q)). Con T = O(log n) todavia es aceptable; con T grande no.
	//
	// 7. Bloque optimo: B = n / sqrt(q) (el codigo lo calcula en resolver). Con q mucho menor
	//    que n (q ~ 10) conviene un B grande; si q es muy grande, B se acerca a 1.
\end{lstlisting}

\textbf{Ejemplo de funcionamiento:}
\begin{lstlisting}[style=cpp]
	int main() {
		vector<int> a = {1, 2, 1, 3, 2, 4, 1};     // valores en [1, 4]
		MoClasico mo(a, 4);
		
		vector<array<int, 2>> consultas = {
			{0, 2},    // [1, 2, 1]             -> 2 distintos
			{1, 4},    // [2, 1, 3, 2]          -> 3 distintos
			{0, 6}     // [1, 2, 1, 3, 2, 4, 1] -> 4 distintos
		};
		
		vector<int> resp = mo.resolver(consultas);
		for (int x : resp) cout << x << " ";
		cout << "\n";                              // 2 3 4
		return 0;
	}
\end{lstlisting}

\subsubsection{Variante: Frecuencia y Moda en un Rango}

\textbf{Idea:} Es el mismo Mo's clásico con otra información mantenida en la ventana. Además de \texttt{cuenta[v]} (cuántas veces aparece $v$) se guarda \texttt{cc[f]}, la cantidad de valores que aparecen exactamente $f$ veces, y \texttt{max\_freq}, la frecuencia de la moda. Al agregar un elemento, su frecuencia sube en $1$ y el máximo es $\max(\texttt{max\_freq}, \texttt{cuenta}[v])$. Al quitarlo, la frecuencia de $v$ baja en $1$. Si $v$ tenía la frecuencia máxima y era el único con ella (\texttt{cc[max\_freq]} queda en $0$), el nuevo máximo es exactamente $\texttt{max\_freq}-1$, porque $v$ acaba de quedar con esa frecuencia. Por eso no hay que buscar el máximo otra vez y ambas operaciones son $O(1)$. La frecuencia de un valor $x$ concreto es simplemente \texttt{cuenta[x]}.

\textbf{Complejidad:} El desplazamiento de punteros es el de Mo's, $O(n\sqrt q)$, con \texttt{agregar} y \texttt{quitar} en $O(1)$. El total es $O(n\sqrt q + q\log q)$ y la memoria es $O(n + \texttt{max\_val} + q)$.

\begin{lstlisting}[style=cpp]
	// Cada consulta es (l, r, x): devuelve {frecuencia de x en [l, r], frecuencia de la moda en [l, r]}.
	struct MoModa {
		int n, bloque;
		vector<int> a;        // arreglo original (valores en [0, max_val])
		vector<int> cuenta;   // cuenta[v]: apariciones de v en la ventana
		vector<int> cc;       // cc[f]: cuantos valores aparecen exactamente f veces (cc[0] no se usa)
		int izq = 0, der = -1;
		int max_freq = 0;     // frecuencia de la moda en la ventana
		
		// CREATE: guarda el arreglo y reserva frecuencias.
		MoModa(const vector<int>& arr, int max_val) {
			a = arr;
			n = (int)a.size();
			cuenta.assign(max_val + 1, 0);
			cc.assign(n + 2, 0);
			bloque = max(1, (int)sqrt((double)n));
		} // O(n + max_val)
		
		// WRITE: agrega el elemento de la posicion pos.
		void agregar(int pos) {
			int v = a[pos];
			cc[cuenta[v]]--;
			cuenta[v]++;
			cc[cuenta[v]]++;
			max_freq = max(max_freq, cuenta[v]);
		} // O(1)
		
		// WRITE: quita el elemento de la posicion pos.
		void quitar(int pos) {
			int v = a[pos];
			cc[cuenta[v]]--;
			if (cuenta[v] == max_freq && cc[cuenta[v]] == 0) max_freq--;   // v era el unico maximo
			cuenta[v]--;
			cc[cuenta[v]]++;
		} // O(1)
		
		// WRITE: responde todas las consultas en el orden original.
		vector<pair<int, int>> resolver(const vector<array<int, 3>>& consultas) {
			int q = (int)consultas.size();
			bloque = max(1, (int)(n / sqrt((double)max(1, q))));
			vector<int> orden(q);
			iota(orden.begin(), orden.end(), 0);
			sort(orden.begin(), orden.end(), [&](int i, int j) {
				int bi = consultas[i][0] / bloque, bj = consultas[j][0] / bloque;
				if (bi != bj) return bi < bj;
				return (bi & 1) ? consultas[i][1] > consultas[j][1]
				: consultas[i][1] < consultas[j][1];
			});
			
			vector<pair<int, int>> resp(q);
			for (int id : orden) {
				int l = consultas[id][0], r = consultas[id][1], x = consultas[id][2];
				while (der < r) agregar(++der);
				while (izq > l) agregar(--izq);
				while (der > r) quitar(der--);
				while (izq < l) quitar(izq++);
				resp[id] = {cuenta[x], max_freq};
			}
			return resp;
		} // O(n sqrt(q) + q log q)
	};
\end{lstlisting}

\textbf{Ejemplo de funcionamiento:}
\begin{lstlisting}[style=cpp]
	int main() {
		vector<int> a = {1, 2, 1, 3, 2, 4, 1};
		MoModa mo(a, 4);
		
		vector<array<int, 3>> consultas = {
			{0, 2, 1},   // [1,2,1]:           frec(1) = 2, moda con 2 apariciones
			{1, 4, 1},   // [2,1,3,2]:         frec(1) = 1, moda con 2 apariciones (el 2)
			{0, 6, 1},   // arreglo completo:  frec(1) = 3, moda con 3 apariciones
			{2, 5, 4},   // [1,3,2,4]:         frec(4) = 1, moda con 1 aparicion
			{0, 6, 2}    // arreglo completo:  frec(2) = 2, moda con 3 apariciones
		};
		
		for (auto [fx, fm] : mo.resolver(consultas)) cout << fx << "," << fm << " ";
		cout << "\n";    // 2,2 1,2 3,3 1,1 2,3
		return 0;
	}
\end{lstlisting}

\textbf{Nota:} Este struct devuelve la \textit{frecuencia} de la moda, no cuál valor es. Para obtener el valor hay que mantener un \texttt{set<pair<int,int>>} de pares (frecuencia, valor) y tomar el último elemento, lo que sube el costo a $O(\log n)$ por movimiento. Si se necesita $O(1)$, la alternativa es Mo's con retroceso, donde solo se usa \texttt{agregar}. Además, \texttt{cc[0]} acumula valores sin sentido, pero nunca se lee.

\subsubsection{Mo's con actualizaciones (3D Mo's)}

\begin{lstlisting}[style=cpp]
	// Igual que el clasico, con actualizaciones puntuales a[pos] = valor intercaladas.
	// Cada consulta es (l, r, t): t = cuantas actualizaciones ya se aplicaron.
	struct MoActualizaciones {
		int n, bloque;
		vector<int> a;
		vector<int> cuenta;
		vector<pair<int, int>> actualizaciones;   // (pos, valor a poner)
		int izq = 0, der = -1, tiempo = 0;
		int distintos = 0;
		
		// CREATE: guarda el arreglo y fija el bloque en ~n^(2/3).
		MoActualizaciones(const vector<int>& arr, int max_val) {
			a = arr;
			n = (int)a.size();
			cuenta.assign(max_val + 1, 0);
			bloque = max(1, (int)pow((double)n, 2.0 / 3.0));
		} // O(n + max_val)
		
		// WRITE: registra una actualizacion (se aplicara cuando le toque, no de inmediato).
		void agregar_actualizacion(int pos, int valor) {
			actualizaciones.push_back({pos, valor});
		} // O(1)
		
		// WRITE: agrega el elemento de la posicion pos a la ventana.
		void agregar(int pos) {
			if (cuenta[a[pos]]++ == 0) distintos++;
		} // O(1)
		
		// WRITE: quita el elemento de la posicion pos de la ventana.
		void quitar(int pos) {
			if (--cuenta[a[pos]] == 0) distintos--;
		} // O(1)
		
		// WRITE: aplica la actualizacion t. Al aplicarla se intercambia el valor a poner con el
		// valor actual, por lo que llamar otra vez a la misma t la REVIERTE.
		void alternar(int t) {
			int pos = actualizaciones[t].first;
			int val = actualizaciones[t].second;
			bool dentro = (pos >= izq && pos <= der);
			if (dentro) quitar(pos);
			swap(a[pos], val);
			actualizaciones[t].second = val;    // guarda el valor viejo
			if (dentro) agregar(pos);
		} // O(1)
		
		// WRITE: responde todas las consultas (l, r, t) y las devuelve en el orden original.
		vector<int> resolver(const vector<array<int, 3>>& consultas) {
			int q = (int)consultas.size();
			vector<int> orden(q);
			iota(orden.begin(), orden.end(), 0);
			sort(orden.begin(), orden.end(), [&](int i, int j) {
				int bli = consultas[i][0] / bloque, blj = consultas[j][0] / bloque;
				if (bli != blj) return bli < blj;
				int bri = consultas[i][1] / bloque, brj = consultas[j][1] / bloque;
				if (bri != brj) return bri < brj;
				return consultas[i][2] < consultas[j][2];
			});
			
			vector<int> resp(q);
			for (int id : orden) {
				int l = consultas[id][0], r = consultas[id][1], t = consultas[id][2];
				while (der < r) agregar(++der);
				while (izq > l) agregar(--izq);
				while (der > r) quitar(der--);
				while (izq < l) quitar(izq++);
				while (tiempo < t) alternar(tiempo++);   // aplicar actualizaciones pendientes
				while (tiempo > t) alternar(--tiempo);   // o revertir las que sobran
				resp[id] = distintos;
			}
			return resp;
		} // O(n^(5/3)) con q y #actualizaciones del orden de n
	};
\end{lstlisting}

\textbf{Ejemplo de funcionamiento:}
\begin{lstlisting}[style=cpp]
	int main() {
		vector<int> a = {1, 2, 3};
		MoActualizaciones mo(a, 3);
		mo.agregar_actualizacion(0, 2);    // actualizacion #0: a[0] = 2
		mo.agregar_actualizacion(1, 1);    // actualizacion #1: a[1] = 1
		
		vector<array<int, 3>> consultas = {
			{0, 2, 0},    // [1, 2, 3], sin actualizaciones        -> 3 distintos
			{0, 2, 1},    // tras la #0: [2, 2, 3]                 -> 2 distintos
			{0, 2, 2}     // tras la #0 y la #1: [2, 1, 3]         -> 3 distintos
		};
		
		vector<int> resp = mo.resolver(consultas);
		for (int x : resp) cout << x << " ";
		cout << "\n";                       // 3 2 3
		return 0;
	}
\end{lstlisting}

\textbf{Nota:} \texttt{resolver} modifica la ventana, así que un mismo objeto puede reutilizarse en llamadas sucesivas (la ventana queda en el estado donde terminó y sigue siendo correcta), pero no es seguro para uso concurrente. Las consultas deben cumplir $0\le l\le r<n$; con \texttt{quitar} sobre una ventana vacía o $l>r$ las frecuencias se vuelven negativas. En Mo's con actualizaciones, el orden de las actualizaciones debe ser el cronológico, y $t$ es la \textit{cantidad} de actualizaciones ya aplicadas ($0$ significa el arreglo original). Si el problema tiene actualizaciones pero se pueden resolver como consultas del tipo «valor en un instante», a veces es más simple un árbol de segmentos persistente. Y en \texttt{alternar}, el orden es importante: primero se quita el valor viejo, luego se cambia y después se agrega el nuevo, siempre que la posición esté dentro de la ventana.

\textbf{Regla práctica:} Usa Mo's solo si (1) todas las consultas se conocen de antemano y (2) la respuesta se mantiene fácil con \texttt{agregar}/\texttt{quitar} en $O(1)$ pero no se puede combinar con un árbol de segmentos. Para suma, mínimo, máximo o mcd, prefiere Fenwick o árbol de segmentos. Si las consultas llegan en línea, usa un árbol de segmentos o descomposición en raíz. Si hay actualizaciones intercaladas, usa la versión con tiempo (bloque $n^{2/3}$). Si \texttt{quitar} es imposible (moda, máximo), usa Mo's con retroceso (modificación 4). Y si los valores son grandes, comprímelos antes.
\subsection{Wavelet Tree}

\textbf{Idea:} es un árbol binario sobre el \textbf{rango de valores}
del arreglo (no sobre las posiciones, como Segment Tree). La raíz
cubre $[\text{lo},\text{hi}]$ de valores posibles; en cada nodo, los
elementos con valor $\le$ mediana van al hijo izquierdo y el resto al
derecho, \textbf{preservando su orden relativo} (partición estable).
Cada nodo guarda un arreglo \texttt{b[]} con cuántos elementos, hasta
cada posición, fueron hacia la izquierda. Con eso puedes, en cada
nivel, saber exactamente qué sub-rango de posiciones le corresponde a
cada hijo — y bajar por el árbol respondiendo preguntas sobre
\textbf{valores dentro de un rango de posiciones} en $O(\log C)$,
donde $C$ es la cantidad de valores distintos.

\textbf{¿Por qué usarlo?} Es la única estructura de esta sección que
responde eficientemente preguntas que combinan \textbf{posición y
valor} a la vez: "el $k$-ésimo menor en $[l,r]$", "cuántos elementos
en $[l,r]$ son $\le X$", "cuántas veces aparece exactamente $X$ en
$[l,r]$". Ninguna de las estructuras anteriores (Fenwick, Segment
Tree, Sparse Table) resuelve esto directamente sin re-plantear el
problema.

\textbf{Casos de uso:}

\begin{itemize}
    \item $k$-ésimo elemento más chico (o más grande) dentro de un
    rango $[l,r]$.
    \item Conteo de elementos en $[l,r]$ cuyo valor está en un rango
    $[a,b]$.
    \item Frecuencia de un valor específico dentro de un rango de
    posiciones (rank query).
    \item Arreglo \textbf{estático} — no soporta updates sin
    reconstruir (para eso existe la variante persistente, más
    avanzada).
\end{itemize}

\textbf{Complejidad:} construir $O(n\log C)$, cada query
$O(\log C)$, donde $C$ es la cantidad de valores distintos tras
comprimir coordenadas.


\subsubsection{Wavelet Tree: \texorpdfstring{$k$}{k}-ésimo, conteo \texorpdfstring{$\le X$}{<= X} y frecuencia}

\begin{lstlisting}[style=cpp]
struct WaveletTree {
    int lo, hi;
    WaveletTree *izq = nullptr, *der = nullptr;
    vector<int> b;   // b[i] = cuantos de a[0..i-1] fueron a la izq.

    // CREATE: construye recursivamente sobre el segmento [from, to)
    // de un arreglo mutable, particionando por la mediana de valor
    // [lo, hi]. Los valores deben venir ya comprimidos a [0, C-1].
    WaveletTree(int* from, int* to, int lo, int hi) : lo(lo), hi(hi) {
        if (lo == hi || from >= to) return;

        int mid = (lo + hi) / 2;
        auto vaIzq = [mid](int x) { return x <= mid; };

        b.reserve(to - from + 1);
        b.push_back(0);
        for (auto it = from; it != to; it++)
            b.push_back(b.back() + vaIzq(*it));

        auto pivote = stable_partition(from, to, vaIzq);
        izq = new WaveletTree(from, pivote, lo, mid);
        der = new WaveletTree(pivote, to, mid + 1, hi);
    }                                          // O(n log C)

    // READ: k-esimo valor mas chico en [l, r] (1-indexado, k desde 1).
    int kth(int l, int r, int k) {
        if (l > r) return -1;
        if (lo == hi) return lo;

        int enIzq = b[r] - b[l - 1];
        int lb = b[l - 1], rb = b[r];

        if (k <= enIzq)
            return izq->kth(lb + 1, rb, k);
        return der->kth(l - lb, r - rb, k - enIzq);
    }                                          // O(log C)

    // READ: cuantos elementos de [l, r] tienen valor <= val.
    int contarMenorIgual(int l, int r, int val) {
        if (l > r || val < lo) return 0;
        if (hi <= val) return r - l + 1;

        int lb = b[l - 1], rb = b[r];
        return izq->contarMenorIgual(lb + 1, rb, val)
             + der->contarMenorIgual(l - lb, r - rb, val);
    }                                          // O(log C)

    // READ: cuantas veces aparece exactamente val en [l, r].
    int contarIgual(int l, int r, int val) {
        if (l > r || val < lo || val > hi) return 0;
        if (lo == hi) return r - l + 1;

        int mid = (lo + hi) / 2;
        int lb = b[l - 1], rb = b[r];

        if (val <= mid)
            return izq->contarIgual(lb + 1, rb, val);
        return der->contarIgual(l - lb, r - rb, val);
    }                                          // O(log C)

    ~WaveletTree() {
        delete izq;
        delete der;
    }
};
\end{lstlisting}

\textbf{Ejemplo de funcionamiento} (con compresión de coordenadas):

\begin{lstlisting}[style=cpp]
vector<int> original = {40, 10, 30, 10, 20, 40, 30};

// comprimir a [0, C-1] para que el wavelet tree trabaje con
// valores pequenos y contiguos.
vector<int> ordenado = original;
sort(ordenado.begin(), ordenado.end());
ordenado.erase(unique(ordenado.begin(), ordenado.end()), ordenado.end());

vector<int> comprimido(original.size());
for (int i = 0; i < (int)original.size(); i++)
    comprimido[i] = lower_bound(ordenado.begin(), ordenado.end(),
                                  original[i]) - ordenado.begin();
// comprimido: [3, 0, 2, 0, 1, 3, 2]   (valores 10,20,30,40 -> 0,1,2,3)

WaveletTree wt(comprimido.data(),
                comprimido.data() + comprimido.size(),
                0, (int)ordenado.size() - 1);

// 3er elemento mas chico en el rango [2, 6] (1-indexado)
int idx = wt.kth(2, 6, 3);
cout << ordenado[idx];
// posiciones 2..6 (1-idx): 10,30,10,20,40,30 (valores 0-idx 2..6)
// -> valores reales: 30,10,20,40,30 ordenados: 10,20,30,30,40
// 3er mas chico = 30

cout << wt.contarMenorIgual(1, 7, 1);
// cuantos valores comprimidos <= 1 (o sea, <= 20) hay en todo el arreglo
// original: 40,10,30,10,20,40,30 -> <=20 son: 10,10,20 -> 3

cout << wt.contarIgual(1, 7, 0);
// cuantas veces aparece el valor comprimido 0 (o sea, 10) -> 2
\end{lstlisting}

\textbf{Comparación con las otras estructuras de rango:}

\begin{table}[H]
\raggedright
\scriptsize
\setlength{\tabcolsep}{3pt}
\begin{tabular}{|p{3.3cm}|p{3.3cm}|}
\hline
\textbf{Pregunta} & \textbf{Estructura} \\
\hline
Suma/min/max de rango & Fenwick / Segment Tree / Sparse Table \\
\hline
$k$-ésimo menor, o conteo por rango de \textbf{valores}, en un rango de \textbf{posiciones} & Wavelet Tree \\
\hline
Elementos distintos, moda, funciones no lineales de frecuencia & Mo's Algorithm \\
\hline
\end{tabular}
\caption{Cuándo el problema necesita Wavelet Tree en vez de las otras estructuras.}
\end{table}

\textbf{Regla práctica:} si el enunciado combina una condición sobre
\textbf{valores} (``menor que'', ``entre $a$ y $b$'', ``el $k$-ésimo
más chico'') con un rango de \textbf{posiciones}, y el arreglo es
estático, es casi siempre Wavelet Tree. Si solo te piden agregar
(suma/min/max) sin condición sobre el valor, usa las estructuras más
simples de las secciones anteriores.

\subsection{Elemento Mayoritario}

\subsubsection{Votación de Boyer--Moore (Elemento Mayoritario)}

\textbf{Idea:} Un elemento \textit{mayoritario} de un arreglo de $n$ elementos es uno que aparece \textit{estrictamente más} de $n/2$ veces; si existe, es único. El algoritmo mantiene un candidato y un contador. Al leer $x$: si el contador es $0$, $x$ pasa a ser el candidato con contador $1$; si $x$ es el candidato, el contador sube; si no, baja. Es una votación en la que cada voto distinto «cancela» uno del candidato.

Funciona por un argumento de emparejamiento. Después de procesar $i$ elementos, el multiconjunto leído se puede dividir en \texttt{cuenta} copias del candidato más $(i-\texttt{cuenta})/2$ \textit{pares de elementos distintos} (cada resta del contador formó un par distinto). Un par de elementos distintos contiene a lo sumo una copia del mayoritario $m$. Al final, si $m$ no fuera el candidato (o el contador fuera $0$), $m$ aparecería a lo sumo $(n-\texttt{cuenta})/2\le n/2$ veces, lo que contradice que sea mayoritario. Por tanto, \textit{si hay un mayoritario, es el candidato final}. El recíproco no vale: si no hay mayoritario, el candidato es arbitrario, y por eso se necesita una segunda pasada que cuente sus apariciones y compruebe que superan $n/2$.

\textbf{¿Por qué usarlo?} Las alternativas son: ordenar y mirar el elemento central ($O(n\log n)$), contar con un \texttt{map} o \texttt{unordered\_map} ($O(n)$ de tiempo pero $O(n)$ de memoria), elegir un elemento al azar y contarlo (éxito con probabilidad $>1/2$ por intento), o \texttt{nth\_element} para hallar la mediana, que es el mayoritario si existe ($O(n)$ esperado, pero modifica el arreglo). Boyer--Moore es $O(n)$ determinista con $O(1)$ de memoria, y recorre los datos en orden, así que funciona sobre un \textit{flujo} que no se puede guardar. Además es \textit{combinable}: dos estados (candidato, contador) se fusionan en $O(1)$, lo que permite usarlo en árboles de segmentos y en cálculo distribuido (modificación 3). A diferencia de las técnicas de rango de las subsecciones anteriores (Mo's, Fenwick, Sparse Table), aquí la consulta es sobre \textit{todo} el arreglo; para rangos hay que combinarlo con un árbol de segmentos. Exige que el umbral sea estrictamente mayor que $n/2$ (con «al menos $n/2$» el argumento falla, como en $\{1,2\}$).

\textbf{Casos de uso:}
\begin{itemize}
	\item Hallar el elemento mayoritario de un arreglo (problema clásico de entrevistas y concursos), o decidir que no existe.
	\item Determinar el ganador de una votación por mayoría absoluta leyendo los votos en una sola pasada.
	\item Procesar flujos de datos con memoria constante, donde no se puede guardar la entrada.
	\item Encontrar elementos frecuentes (aparecen más de $n/k$ veces) con la generalización de Misra--Gries (modificación 1).
	\item De nicho: mayoritario de un rango con un árbol de segmentos (modificación 3), o fusionar resultados parciales de varias máquinas.
\end{itemize}

\textbf{Complejidad:} La primera pasada hace una comparación y una actualización del contador por elemento: $O(n)$. La verificación cuenta las apariciones del candidato en otra pasada: $O(n)$. El tiempo total es $O(n)$ y la memoria extra es $O(1)$ (un candidato y un contador). Fusionar dos estados es $O(1)$. La generalización a $n/k$ guarda hasta $k-1$ contadores, con memoria $O(k)$, y un tiempo $O(n)$ amortizado, porque cada «decrementar todos» elimina $k$ unidades de contador distintas y eso ocurre a lo sumo $n/k$ veces.

\begin{lstlisting}[style=cpp]
	struct VotacionMayoritaria {
		long long cand;   // candidato actual
		int cuenta;       // votos netos del candidato (0 = sin candidato)
		
		// CREATE: estado vacio, sin candidato.
		VotacionMayoritaria() : cand(0), cuenta(0) {} // O(1)
		
		// WRITE: procesa un elemento x del flujo.
		void agregar(long long x) {
			if (cuenta == 0) { cand = x; cuenta = 1; }
			else if (cand == x) cuenta++;
			else cuenta--;
		} // O(1)
		
		// WRITE: fusiona el estado de otro bloque (para rangos o datos distribuidos).
		// Un par de candidatos distintos se cancela; el que tiene mas votos sobrevive.
		void combinar(const VotacionMayoritaria& o) {
			if (cand == o.cand) cuenta += o.cuenta;
			else if (cuenta >= o.cuenta) cuenta -= o.cuenta;
			else { cand = o.cand; cuenta = o.cuenta - cuenta; }
		} // O(1)
		
		// READ: candidato actual (solo es mayoritario si se verifica).
		long long candidato() const {
			return cand;
		} // O(1)
		
		// READ: indica si x aparece estrictamente mas de n/2 veces en a (segunda pasada).
		static bool verificar(const vector<long long>& a, long long x) {
			long long c = 0;
			for (long long y : a) if (y == x) c++;
			return 2 * c > (long long)a.size();
		} // O(n)
		
		// READ: devuelve {existe, valor} del elemento mayoritario de a.
		static pair<bool, long long> mayoritario(const vector<long long>& a) {
			VotacionMayoritaria v;
			for (long long x : a) v.agregar(x);
			if (a.empty() || !verificar(a, v.cand)) return {false, 0};
			return {true, v.cand};
		} // O(n)
	};
	
	// ------------------------------------------------------------
	// MODIFICACIONES Y COMENTARIOS CON LOS USOS
	// ------------------------------------------------------------
	// 1. Elementos que aparecen MAS de n/k veces (Misra-Gries): hay a lo sumo k - 1 de ellos.
	//    Se guardan hasta k - 1 contadores; si llega un valor nuevo y no hay sitio, se
	//    resta 1 a todos y se descartan los que llegan a 0. Con k = 2 es Boyer-Moore.
	//    map<long long, int> c;
	//    for (long long x : a) {
		//        if (c.count(x)) c[x]++;
		//        else if ((int)c.size() < k - 1) c[x] = 1;
		//        else {
			//            for (auto it = c.begin(); it != c.end(); )
			//                if (--it->second == 0) it = c.erase(it); else ++it;
			//        }
		//    }
	//    // Los candidatos son las claves de c; verificar cada uno con una segunda pasada.
	//    // O(n) amortizado (O(n log k) con map), memoria O(k).
	//
	// 2. Alternativas simples cuando no importa la memoria:
	//    a) Ordenar: tras sort, el mayoritario si existe es a[n / 2]. O(n log n).
	//    b) nth_element(a.begin(), a.begin() + n / 2, a.end()); candidato = a[n / 2]. O(n) esperado.
	//    c) Por bits (enteros): para cada bit b, el bit del mayoritario es el que aparece en mas de
	//       n/2 elementos; armar el numero y verificar. O(n * bits), memoria O(1), sin comparar valores.
	//    d) Aleatorio: elegir x = a[rand() % n] y contar. Probabilidad de acierto > 1/2; con 30
	//       intentos falla con probabilidad < 2^-30. O(30 n).
	//
	// 3. Mayoritario de un RANGO [l, r] (consultas en linea): arbol de segmentos donde cada nodo
	//    guarda un VotacionMayoritaria y se une con combinar(). La consulta fusiona los O(log n)
	//    nodos y da un candidato; se verifica contando sus apariciones en [l, r] con listas de
	//    posiciones ordenadas por valor:
	//    // pos[v] = posiciones de v en orden creciente
	//    // veces = upper_bound(pos[v].begin(), pos[v].end(), r)
	//    //       - lower_bound(pos[v].begin(), pos[v].end(), l);
	//    // existe mayoritario <=> 2 * veces > r - l + 1
	//    // Construccion O(n), consulta O(log n). Valores grandes: comprimirlos primero.
	//
	// 4. Con ">= n/2" (empate permitido) el argumento de los pares falla: en {1, 2} cada uno
	//    aparece n/2 veces y el candidato es 2. Si el enunciado lo pide, contar con una tabla.
	//
	// 5. Con pesos (cada elemento tiene un peso w): agregar suma o resta w al contador (y si
	//    el peso supera el contador, el candidato cambia con contador w - cuenta). Mismo O(n).
	//
	// 6. Verificar siempre. Si el enunciado GARANTIZA que existe un mayoritario, la segunda
	//    pasada se puede omitir y el algoritmo queda en una sola lectura del flujo.
\end{lstlisting}

\textbf{Ejemplo de funcionamiento:}
\begin{lstlisting}[style=cpp]
	int main() {
		vector<long long> a = {2, 2, 1, 1, 1, 2, 2};
		auto r = VotacionMayoritaria::mayoritario(a);
		cout << r.first << " " << r.second << "\n";   // 1 2  (el 2 aparece 4 de 7 veces)
		
		vector<long long> b = {3, 3, 4, 2, 4, 4, 2, 4, 4};
		r = VotacionMayoritaria::mayoritario(b);
		cout << r.first << " " << r.second << "\n";   // 1 4  (el 4 aparece 5 de 9 veces)
		
		vector<long long> c = {1, 1, 2, 2, 3};
		r = VotacionMayoritaria::mayoritario(c);
		cout << r.first << "\n";                      // 0  (el candidato final es 3, que aparece 1 vez)
		
		vector<long long> d = {1, 2};
		cout << VotacionMayoritaria::mayoritario(d).first << "\n";   // 0  (empate: nadie supera n/2)
		
		// Fusion de dos bloques: {2, 2, 1} y {1, 1, 2} -> estados (2, 1) y (1, 1)
		VotacionMayoritaria izq, der;
		for (long long x : {2, 2, 1}) izq.agregar(x);
		for (long long x : {1, 1, 2}) der.agregar(x);
		izq.combinar(der);
		cout << izq.cuenta << "\n";                   // 0  (se cancelan por completo: 3 contra 3, sin mayoritario)
		return 0;
	}
\end{lstlisting}

\textbf{Nota:} El candidato que entrega la primera pasada \textit{no} es una respuesta, es una sospecha: sin la verificación, \texttt{\{1,2,3\}} devolvería $3$. Si el contador final es $0$ se puede concluir de inmediato que no hay mayoritario, sin verificar. En \texttt{combinar}, cuando los candidatos son distintos y los contadores iguales, el resultado queda con contador $0$ y el candidato es irrelevante. La fusión es correcta por el mismo argumento de pares: lo que sobra de cada bloque son copias de su candidato más pares distintos, y al cruzar candidatos distintos se forman más pares. Por último, el contador de \texttt{int} basta mientras $n<2^{31}$.

\textbf{Regla práctica:} Si el problema pide el elemento que aparece más de $n/2$ veces en todo el arreglo, o en un flujo con memoria limitada, usa Boyer--Moore con una segunda pasada de verificación (o sin ella, si el enunciado garantiza que existe). Si el umbral es $n/k$, usa Misra--Gries con $k-1$ contadores. Si hay consultas de rango, combina los estados en un árbol de segmentos y verifica con listas de posiciones. Si la memoria no importa y la claridad sí, ordenar y mirar $a[n/2]$ o usar un \texttt{map} es suficiente. Y si el umbral es «al menos $n/2$», no uses este algoritmo.
\subsection{Heavy-Light Decomposition}

\textbf{Idea:} descompone un árbol en \textbf{cadenas} de forma que
cualquier camino desde la raíz hasta cualquier nodo cruza a lo más
$O(\log n)$ cadenas distintas. En cada nodo, el hijo con el
\textbf{subárbol más grande} (el hijo ``pesado'') continúa la misma
cadena que su padre; los demás hijos inician cadenas nuevas. Al
numerar los nodos en el orden en que se visitan (DFS que siempre
sigue primero al hijo pesado), cada cadena queda como un \textbf{rango
contiguo} de posiciones — así que una query sobre un \textbf{camino}
del árbol se convierte en $O(\log n)$ queries de rango sobre un
Segment Tree normal, una por cada cadena que el camino atraviesa.

\textbf{¿Por qué usarlo?} Es la forma estándar de llevar range
queries de un arreglo a un \textbf{árbol}: sumas/mínimos/updates en
el camino entre dos nodos, o en el subárbol de un nodo. Sin HLD, un
camino arbitrario no tiene ninguna estructura contigua sobre la cual
aplicar Fenwick o Segment Tree directamente.

\textbf{Casos de uso:}

\begin{itemize}
    \item Suma/mínimo/máximo en el camino entre dos nodos $u$ y $v$.
    \item Actualizar todos los nodos (o aristas) de un camino $u$-$v$.
    \item Consultas o updates sobre el \textbf{subárbol} completo de un
    nodo (más simple: el subárbol siempre es un rango contiguo, sin
    necesidad de descomponer en cadenas).
    \item Cualquier problema de árbol que ``se sentiría'' como Range
    Queries si el árbol fuera un arreglo lineal.
\end{itemize}

\textbf{Complejidad:} construir $O(n)$ (dos DFS) más $O(n)$ del
Segment Tree subyacente; query o update de \textbf{camino}
$O(\log^2 n)$ ($O(\log n)$ cadenas $\times$ $O(\log n)$ por consulta
al Segment Tree); query o update de \textbf{subárbol} $O(\log n)$
(un solo rango contiguo).


\subsubsection{HLD: suma en camino, suma en subárbol}

\textbf{Idea de la implementación:} dos DFS \textbf{iterativas} (con
pila explícita, para evitar stack overflow en árboles muy profundos
o muy desbalanceados — un riesgo real con recursión cuando
$n\ge10^5$ y el árbol es una cadena larga). El primero
(\texttt{dfsTamano}) calcula el tamaño de cada subárbol y decide el
hijo pesado de cada nodo. El segundo (\texttt{dfsDescomponer})
recorre siempre primero al hijo pesado (así la cadena sigue
consecutiva en las posiciones) y solo después explora los hijos
livianos, cada uno iniciando su propia cadena nueva. Las posiciones
resultantes (\texttt{pos[]}) alimentan un Segment Tree normal —aquí
se reutiliza \texttt{SegmentTreeLazyAdd} de la sección anterior—
sobre el cual se hacen las queries reales.

\begin{lstlisting}[style=cpp]
struct HLD {
    vector<vector<int>> adj;
    vector<int> padre, profundidad, tamano, pesado, cabeza, pos;
    int n, timer = 0;
    SegmentTreeLazyAdd* seg;   // Segment Tree sobre el orden pos[]

    // CREATE: dos DFS iterativas para calcular cadenas, luego arma
    // el Segment Tree sobre el arreglo reordenado por pos[].
    HLD(vector<vector<int>>& adjacency, vector<long long>& valores,
        int raiz) {
        adj = adjacency;
        n = adj.size();
        padre.assign(n, -1);
        profundidad.assign(n, 0);
        tamano.assign(n, 1);
        pesado.assign(n, -1);
        cabeza.assign(n, raiz);
        pos.assign(n, 0);

        dfsTamano(raiz);
        dfsDescomponer(raiz);

        vector<long long> reordenado(n);
        for (int v = 0; v < n; v++)
            reordenado[pos[v]] = valores[v];

        seg = new SegmentTreeLazyAdd(reordenado);
    }                                          // O(n)

    // Fase 1: obtiene un orden donde cada padre aparece ANTES que
    // sus hijos (equivalente a un preorden), usando una pila.
    // Fase 2: recorre ese orden al reves (hijos antes que padres)
    // para acumular tamano[] y decidir el hijo pesado, igual que
    // haria la version recursiva al "volver" de cada llamada.
    void dfsTamano(int raiz) {
        vector<int> orden;
        stack<int> st;
        st.push(raiz);
        padre[raiz] = -1;
        profundidad[raiz] = 0;

        while (!st.empty()) {
            int u = st.top();
            st.pop();
            orden.push_back(u);

            for (int v : adj[u]) {
                if (v != padre[u]) {
                    padre[v] = u;
                    profundidad[v] = profundidad[u] + 1;
                    st.push(v);
                }
            }
        }

        for (int i = (int)orden.size() - 1; i >= 0; i--) {
            int u = orden[i];
            tamano[u] = 1;
            int maxSub = 0;

            for (int v : adj[u]) {
                if (v == padre[u]) continue;

                tamano[u] += tamano[v];
                if (tamano[v] > maxSub) {
                    maxSub = tamano[v];
                    pesado[u] = v;
                }
            }
        }
    }                                          // O(n)

    // Pila de pares (nodo, cabeza de su cadena). Para que el hijo
    // pesado quede en el tope de la pila (y por lo tanto se
    // procese inmediatamente despues, como en la version
    // recursiva), se empuja de ULTIMO: los hijos livianos
    // (cadena nueva) se empujan primero.
    void dfsDescomponer(int raiz) {
        stack<pair<int, int>> st;
        st.push({raiz, raiz});

        while (!st.empty()) {
            auto [u, h] = st.top();
            st.pop();

            cabeza[u] = h;
            pos[u] = timer++;

            for (int v : adj[u]) {
                if (v == padre[u] || v == pesado[u]) continue;
                st.push({v, v});                // cadena nueva
            }

            if (pesado[u] != -1)
                st.push({pesado[u], h});        // misma cadena
        }
    }                                          // O(n)

    // READ: suma de todos los nodos en el camino u-v.
    long long consultarCamino(int u, int v) {
        long long resultado = 0;

        while (cabeza[u] != cabeza[v]) {
            if (profundidad[cabeza[u]] < profundidad[cabeza[v]])
                swap(u, v);

            resultado += seg->rango(pos[cabeza[u]], pos[u]);
            u = padre[cabeza[u]];
        }

        if (profundidad[u] > profundidad[v]) swap(u, v);
        resultado += seg->rango(pos[u], pos[v]);
        return resultado;
    }                                          // O(log^2 n)

    // UPDATE: suma val a todos los nodos en el camino u-v.
    void actualizarCamino(int u, int v, long long val) {
        while (cabeza[u] != cabeza[v]) {
            if (profundidad[cabeza[u]] < profundidad[cabeza[v]])
                swap(u, v);

            seg->rangoUpdate(pos[cabeza[u]], pos[u], val);
            u = padre[cabeza[u]];
        }

        if (profundidad[u] > profundidad[v]) swap(u, v);
        seg->rangoUpdate(pos[u], pos[v], val);
    }                                          // O(log^2 n)

    // READ: suma de todo el subarbol de u (rango contiguo por
    // construccion, gracias al orden DFS de dfsDescomponer).
    long long consultarSubarbol(int u) {
        return seg->rango(pos[u], pos[u] + tamano[u] - 1);
    }                                          // O(log n)

    // UPDATE: suma val a todo el subarbol de u.
    void actualizarSubarbol(int u, long long val) {
        seg->rangoUpdate(pos[u], pos[u] + tamano[u] - 1, val);
    }                                          // O(log n)
};
\end{lstlisting}

\textbf{Ejemplo de funcionamiento:}

\begin{lstlisting}[style=cpp]
// arbol:        0
//              /|\
//             1 2 3
//            /|
//           4 5

vector<vector<int>> adj(6);
adj[0] = {1, 2, 3};
adj[1] = {0, 4, 5};
adj[2] = {0};
adj[3] = {0};
adj[4] = {1};
adj[5] = {1};

vector<long long> valores = {1, 2, 3, 4, 5, 6};
// nodo:                     0  1  2  3  4  5

HLD hld(adj, valores, 0);

cout << hld.consultarCamino(4, 3);
// camino 4 -> 1 -> 0 -> 3 : valores 5+2+1+4 = 12

cout << hld.consultarSubarbol(1);
// subarbol de 1: nodos {1,4,5} : 2+5+6 = 13

hld.actualizarCamino(4, 5, 10);
// suma 10 a cada nodo del camino 4 -> 1 -> 5

cout << hld.consultarSubarbol(0);
// subarbol de 0 (todo el arbol) tras el update:
// 1+(2+10)+3+4+(5+10)+(6+10) = 51
\end{lstlisting}

\textbf{Nota:} al ser iterativas, \texttt{dfsTamano} y
\texttt{dfsDescomponer} no tienen riesgo de stack overflow incluso en
árboles con $n\ge10^5$ y forma de ``cadena larga'' (el peor caso para
recursión). \texttt{SegmentTreeLazyAdd} es exactamente el struct de
la sección de Segment Tree — cualquier otra variante (min, max,
assign) funciona igual de bien como motor subyacente.

\textbf{Regla práctica:} usa HLD en cuanto el problema pida algo
sobre un \textbf{camino} entre dos nodos de un árbol con muchas
queries. Si solo pide algo sobre \textbf{subárboles} (sin caminos
arbitrarios), no hace falta la descomposición en cadenas completa:
basta un solo DFS para las posiciones (\texttt{pos[]} y
\texttt{tamano[]}) y un Segment Tree normal, ya que el subárbol
siempre es un rango contiguo por sí solo.



\part{Grafos}

\subsection{¿Qué algoritmo usar?}

La siguiente tabla resume situaciones típicas de Competitive Programming
y el algoritmo que debe considerarse. La idea es reconocer las palabras
clave o condiciones del enunciado y asociarlas rápidamente con una técnica.

\begin{table}[H]
	\raggedright
	\scriptsize
	\setlength{\tabcolsep}{2.5pt}
	\renewcommand{\arraystretch}{1.15}
	\begin{tabular}{|p{0.42\columnwidth}|p{0.23\columnwidth}|p{0.25\columnwidth}|}
		\hline
		\textbf{Cuando el ejercicio diga / pida} &
		\textbf{Usa} &
		\textbf{Pista rápida} \\
		\hline
		
		``Recorrer el grafo por niveles'' &
		BFS &
		Usa una cola. \\
		\hline
		
		``Encontrar la distancia mínima'' con aristas de igual costo &
		BFS &
		Cada arista cuesta lo mismo. \\
		\hline
		
		``Encontrar caminos mínimos'' con pesos $0$ y $1$ &
		0-1 BFS &
		Deque; pesos únicamente $0$ o $1$. \\
		\hline
		
		``Camino mínimo'' con pesos no negativos &
		Dijkstra &
		Usa una cola de prioridad. \\
		\hline
		
		``Camino mínimo'' con pesos negativos &
		Bellman-Ford &
		Permite detectar ciclos negativos. \\
		\hline
		
		``Camino mínimo entre todos los pares'' y $n$ pequeño &
		Floyd-Warshall &
		DP sobre todos los pares. \\
		\hline
		
		``¿Existe un camino entre $u$ y $v$?'' &
		DFS / BFS &
		Problema de conectividad. \\
		\hline
		
		``Encontrar componentes conexas'' &
		DFS / BFS &
		Ejecutar traversal desde cada nodo no visitado. \\
		\hline
		
		``Detectar si existe un ciclo'' en grafo no dirigido &
		DFS / DSU &
		DFS para recorrer; DSU si se agregan aristas. \\
		\hline
		
		``Detectar ciclos'' en grafo dirigido &
		DFS / Kahn &
		Usar estados en DFS o grados de entrada. \\
		\hline
		
		``Ordenar tareas con dependencias'' &
		Topological Sort &
		El grafo debe ser dirigido y acíclico. \\
		\hline
		
		``Encontrar puentes'' &
		Tarjan &
		DFS + tiempos de descubrimiento. \\
		\hline
		
		``Encontrar puntos de articulación'' &
		Tarjan &
		Vértices cuya eliminación desconecta el grafo. \\
		\hline
		
		``Encontrar componentes fuertemente conexas'' &
		Kosaraju / Tarjan &
		Grafo dirigido. \\
		\hline
		
		``Conectar todos los nodos con costo mínimo'' &
		MST &
		Kruskal o Prim. \\
		\hline
		
		``MST procesando aristas ordenadas por peso'' &
		Kruskal &
		Ordenar aristas + DSU. \\
		\hline
		
		``MST expandiendo desde un nodo'' &
		Prim &
		Priority Queue + lista de adyacencia. \\
		\hline
		
		``¿Dos nodos pertenecen al mismo conjunto?'' &
		DSU &
		Union-Find. \\
		\hline
		
		``Agrupar elementos mediante uniones'' &
		DSU &
		\texttt{union} + \texttt{find}. \\
		\hline
		
		``Encontrar si un grafo es bipartito'' &
		2-Coloring &
		Colorear con dos colores. \\
		\hline
		
		``Máximo número de parejas entre dos grupos'' &
		Bipartite Matching &
		Emparejamiento entre conjuntos. \\
		\hline
		
		``Máximo flujo de una fuente a un destino'' &
		Dinic &
		Network Flow. \\
		\hline
		
		``Separar fuente y sumidero con capacidad mínima'' &
		Min-Cut &
		Por el teorema Max-Flow Min-Cut. \\
		\hline
		
		``Distancia entre dos nodos de un árbol'' &
		LCA &
		$\operatorname{dist}(u,v)=
		d[u]+d[v]-2d[\operatorname{LCA}(u,v)]$. \\
		\hline
		
		``Encontrar el ancestro $k$ niveles arriba'' &
		Binary Lifting &
		Saltar potencias de $2$. \\
		\hline
		
		``Consultas sobre subárboles'' &
		Euler Tour + Fenwick / Segment Tree &
		Convertir subárbol en intervalo. \\
		\hline
		
		``Consultas o actualizaciones sobre caminos de un árbol'' &
		HLD &
		Descomponer el camino en cadenas. \\
		\hline
		
		``Encontrar el diámetro de un árbol'' &
		BFS / DFS &
		Dos traversals desde extremos. \\
		\hline
		
		``Encontrar un camino que use todas las aristas'' &
		Eulerian Path / Hierholzer &
		Trabaja con aristas, no con vértices. \\
		\hline
		
		``Encontrar todas las rutas posibles'' &
		DFS + Backtracking &
		Explorar y deshacer decisiones. \\
		\hline
		
		``Rellenar una región conectada'' &
		Flood Fill &
		DFS / BFS desde una celda inicial. \\
		\hline
		
		``Moverse por una cuadrícula buscando la menor cantidad de pasos'' &
		BFS &
		Cada movimiento tiene costo $1$. \\
		\hline
		
		``Variables booleanas con restricciones de la forma OR'' &
		2-SAT + SCC &
		Construir grafo de implicaciones. \\
		\hline
		
		``Resolver muchas consultas sobre un árbol'' &
		LCA / HLD &
		Depende de si son ancestros o caminos. \\
		\hline
		
	\end{tabular}
	\caption{Guía rápida para seleccionar algoritmos de grafos.}
\end{table}
\begin{table}[H]
	\raggedright
	\scriptsize
	\setlength{\tabcolsep}{3pt}
	\renewcommand{\arraystretch}{1.15}
	\begin{tabular}{|p{0.43\columnwidth}|p{0.27\columnwidth}|p{0.22\columnwidth}|}
		\hline
		\textbf{Condición del problema} &
		\textbf{Algoritmo} &
		\textbf{Complejidad} \\
		\hline
		Aristas sin peso / costo $1$ &
		BFS &
		$O(n+m)$ \\
		\hline
		Pesos únicamente $0$ y $1$ &
		0-1 BFS &
		$O(n+m)$ \\
		\hline
		Pesos no negativos &
		Dijkstra &
		$O((n+m)\log n)$ \\
		\hline
		Puede haber pesos negativos &
		Bellman-Ford &
		$O(nm)$ \\
		\hline
		Todos los pares y $n$ pequeño &
		Floyd-Warshall &
		$O(n^3)$ \\
		\hline
	\end{tabular}
	\caption{Selección del algoritmo para caminos mínimos.}
\end{table}



% =========================================================
% TRAVERSAL
% =========================================================

\section{Traversal}


\subsection{Representación de Grafos}

\textbf{Idea:} antes de aplicar cualquier algoritmo sobre un grafo,
primero debemos decidir cómo almacenar sus vértices y aristas. Las
tres representaciones más comunes son la \textbf{lista de adyacencia},
la \textbf{matriz de adyacencia} y la \textbf{lista de aristas}. La
elección depende principalmente de la cantidad de vértices y aristas
y del tipo de consultas que realizará el algoritmo.

La \textbf{lista de adyacencia} almacena, para cada vértice, todos los
vértices conectados directamente con él. Es la representación más
utilizada en Competitive Programming porque ocupa $O(n+m)$ memoria y
permite recorrer eficientemente todos los vecinos de un vértice.

La \textbf{matriz de adyacencia} utiliza una matriz $n\times n$ donde
la posición $(u,v)$ indica si existe una arista entre $u$ y $v$.
Permite comprobar la existencia de una arista en $O(1)$, pero requiere
$O(n^2)$ memoria.

La \textbf{lista de aristas} simplemente almacena cada arista como un
par $(u,v)$, o como una terna $(u,v,w)$ si tiene peso. Es especialmente
útil para algoritmos que necesitan procesar todas las aristas, como
Kruskal o Bellman-Ford.

\textbf{¿Por qué usar cada una?}

\begin{itemize}
    \item \textbf{Lista de adyacencia:} opción general para BFS, DFS,
    Dijkstra, Prim, búsqueda de ciclos y prácticamente la mayoría de
    algoritmos sobre grafos dispersos.

    \item \textbf{Matriz de adyacencia:} útil cuando $n$ es pequeño o
    cuando se necesita comprobar constantemente si existe una arista
    entre dos vértices en $O(1)$.

    \item \textbf{Lista de aristas:} conveniente cuando el algoritmo
    necesita recorrer todas las aristas directamente, especialmente
    en Kruskal y Bellman-Ford.
\end{itemize}

\textbf{Tipos de grafos:} una arista puede ser \textbf{dirigida} o
\textbf{no dirigida}, y puede tener o no un \textbf{peso}. En un grafo
no dirigido, una arista $(u,v)$ se representa agregando $v$ a los
vecinos de $u$ y $u$ a los vecinos de $v$. En un grafo dirigido,
solamente se agrega $v$ a los vecinos de $u$.

\textbf{Complejidad:}

\begin{table}[H]
\raggedright
\scriptsize
\setlength{\tabcolsep}{2.5pt}
\renewcommand{\arraystretch}{1.1}
\begin{tabular}{|p{0.27\columnwidth}|c|c|c|}
\hline
\textbf{Representación} &
\textbf{Memoria} &
\textbf{Consulta} &
\textbf{Vecinos} \\
\hline
Lista de adyacencia &
$O(n+m)$ &
$O(\deg(u))$ &
$O(\deg(u))$ \\
\hline
Matriz de adyacencia &
$O(n^2)$ &
$O(1)$ &
$O(n)$ \\
\hline
Lista de aristas &
$O(m)$ &
$O(m)$ &
$O(m)$ \\
\hline
\end{tabular}
\caption{Comparación de representaciones de grafos.}
\end{table}

\subsubsection{Lista de adyacencia}

\textbf{Idea de la implementación:} se utiliza un
\texttt{vector<vector<int>>} donde la posición $u$ contiene todos los
vértices conectados directamente con $u$. Para grafos ponderados se
puede almacenar un par $(v,w)$, donde $v$ es el destino y $w$ el peso
de la arista.

\begin{lstlisting}[style=cpp]
// Grafo no dirigido y no ponderado.
int n = 5;
vector<vector<int>> adj(n);

// Agrega una arista u-v.
auto agregarArista = [&](int u, int v) {
    adj[u].push_back(v);
    adj[v].push_back(u);
};

agregarArista(0, 1);
agregarArista(0, 2);
agregarArista(1, 3);
agregarArista(1, 4);
\end{lstlisting}

\textbf{Grafo resultante:}

\begin{lstlisting}[style=cpp]
//        0
//       / \
//      1   2
//     / \
//    3   4
//
// adj[0] = {1, 2}
// adj[1] = {0, 3, 4}
// adj[2] = {0}
// adj[3] = {1}
// adj[4] = {1}
\end{lstlisting}

Para un grafo dirigido, solamente se agrega la arista en una dirección:

\begin{lstlisting}[style=cpp]
// Grafo dirigido.
vector<vector<int>> adj(n);

auto agregarArista = [&](int u, int v) {
    adj[u].push_back(v);
};

agregarArista(0, 1);
agregarArista(0, 2);
agregarArista(1, 3);
\end{lstlisting}

Para un grafo ponderado, cada vecino puede almacenar también el peso:

\begin{lstlisting}[style=cpp]
// Grafo ponderado.
vector<vector<pair<int, long long>>> adj(n);

auto agregarArista = [&](int u, int v, long long w) {
    adj[u].push_back({v, w});
    adj[v].push_back({u, w});
};

agregarArista(0, 1, 5);
agregarArista(0, 2, 3);
agregarArista(1, 3, 7);
\end{lstlisting}

\textbf{Ejemplo de recorrido:}

\begin{lstlisting}[style=cpp]
void dfs(int u, int padre, vector<vector<int>>& adj) {

    cout << u << ' ';

    for (int v : adj[u]) {
        if (v == padre) continue;

        dfs(v, u, adj);
    }
}
\end{lstlisting}

La lista de adyacencia permite que DFS o BFS procesen cada vértice y
cada arista una cantidad constante de veces, obteniendo una
complejidad total de $O(n+m)$.

\subsubsection{Matriz de adyacencia}

\textbf{Idea:} una matriz $adj[u][v]$ indica directamente si existe
una conexión desde $u$ hacia $v$. Si el grafo es ponderado, la
posición puede almacenar directamente el peso de la arista.

\begin{lstlisting}[style=cpp]
int n = 5;
vector<vector<int>> adj(n, vector<int>(n, 0));

auto agregarArista = [&](int u, int v) {
    adj[u][v] = 1;
    adj[v][u] = 1;
};

agregarArista(0, 1);
agregarArista(0, 2);
agregarArista(1, 3);
agregarArista(1, 4);
\end{lstlisting}

Ahora comprobar si existe una arista entre $u$ y $v$ es inmediato:

\begin{lstlisting}[style=cpp]
if (adj[u][v]) {
    cout << "Existe una arista";
}
\end{lstlisting}

La consulta anterior cuesta $O(1)$, pero recorrer todos los vecinos de
$u$ requiere revisar las $n$ posiciones de la fila:

\begin{lstlisting}[style=cpp]
for (int v = 0; v < n; v++) {
    if (adj[u][v]) {
        cout << v << ' ';
    }
}
\end{lstlisting}

\subsubsection{Lista de aristas}

\textbf{Idea:} almacena directamente todas las aristas del grafo.
Para un grafo no ponderado basta guardar $(u,v)$; para uno ponderado
se guarda $(u,v,w)$.

\begin{lstlisting}[style=cpp]
// Grafo ponderado.
struct Edge {
    int u, v;
    long long w;
};

vector<Edge> edges;

edges.push_back({0, 1, 5});
edges.push_back({0, 2, 3});
edges.push_back({1, 3, 7});
\end{lstlisting}

Esta representación resulta especialmente conveniente cuando el
algoritmo necesita procesar todas las aristas:

\begin{lstlisting}[style=cpp]
for (Edge e : edges) {
    cout << e.u << " -> " << e.v
         << " peso = " << e.w << '\n';
}
\end{lstlisting}

Por ejemplo, \textbf{Kruskal} ordena directamente esta lista por peso,
mientras que \textbf{Bellman-Ford} puede recorrerla completa en cada
iteración.

\textbf{Nota:} para la mayoría de problemas de Competitive Programming,
la \textbf{lista de adyacencia} debe ser la representación por defecto.
Es especialmente importante cuando el grafo es disperso, es decir,
cuando $m$ es mucho menor que $n^2$. La matriz de adyacencia solamente
conviene cuando $n$ es pequeño o cuando las consultas de existencia de
aristas en $O(1)$ son importantes.

\textbf{Regla práctica:} si no tienes una razón específica para usar
otra representación, usa \textbf{lista de adyacencia}. Si necesitas
preguntar constantemente ``¿existe una arista $u\rightarrow v$?'',
considera una \textbf{matriz de adyacencia}. Si el algoritmo trabaja
directamente con todas las aristas, considera una \textbf{lista de
aristas}.

```latex
\subsection{BFS (Breadth-First Search)}

\textbf{Idea:} BFS (\textit{Breadth-First Search}) es un algoritmo de
recorrido de grafos que explora los vértices por \textbf{niveles}. A
partir de un vértice inicial, primero visita todos sus vecinos
directos, después los vecinos de estos, y continúa de esta manera
hasta procesar todos los vértices alcanzables.

Para mantener este orden se utiliza una \textbf{cola}
(\texttt{queue}), siguiendo el principio FIFO (\textit{First In,
First Out}). Cada vértice se marca como visitado cuando se agrega a
la cola, evitando procesarlo varias veces.

Una de las propiedades más importantes de BFS es que, en un grafo
\textbf{no ponderado}, encuentra la distancia mínima desde un vértice
origen hasta cualquier otro vértice alcanzable, donde la distancia se
mide como el número de aristas recorridas.

\textbf{¿Por qué usarlo?}

\begin{itemize}
    \item \textbf{Camino mínimo:} encuentra el camino con menor número
    de aristas en grafos no ponderados.

    \item \textbf{Distancias:} permite calcular la distancia mínima
    desde un origen hasta todos los vértices alcanzables.

    \item \textbf{Recorrido por niveles:} permite procesar los
    vértices según su distancia al origen.

    \item \textbf{Reconstrucción de caminos:} almacenando el padre de
    cada vértice se puede recuperar el camino mínimo.

    \item \textbf{Bipartición:} permite determinar si un grafo es
    bipartito utilizando dos colores.

    \item \textbf{Componentes conexas:} ejecutando BFS desde cada
    vértice no visitado se pueden encontrar todas las componentes
    conexas.

    \item \textbf{Problemas en grids:} permite encontrar la cantidad
    mínima de movimientos cuando cada movimiento tiene el mismo costo.

    \item \textbf{Multi-Source BFS:} permite calcular la distancia
    mínima desde cualquiera de varios vértices iniciales.
\end{itemize}

\textbf{Funcionamiento:}

El algoritmo sigue los siguientes pasos:

\begin{enumerate}
    \item Elegir un vértice inicial $s$.
    \item Marcar $s$ como visitado y agregarlo a la cola.
    \item Mientras la cola no esté vacía, extraer el primer vértice.
    \item Recorrer todos los vecinos del vértice extraído.
    \item Si un vecino no ha sido visitado, marcarlo y agregarlo a
    la cola.
\end{enumerate}

La cola garantiza que los vértices se procesen por niveles. Primero
se procesan los vértices a distancia $0$, después los de distancia
$1$, luego los de distancia $2$, y así sucesivamente.

\textbf{Implementación básica:}

\begin{lstlisting}[style=cpp]
void bfs(int inicio, vector<vector<int>>& adj) {

    int n = adj.size();

    vector<bool> visitado(n, false);
    queue<int> q;

    visitado[inicio] = true;
    q.push(inicio);

    while (!q.empty()) {

        int u = q.front();
        q.pop();

        cout << u << ' ';

        for (int v : adj[u]) {

            if (visitado[v]) continue;

            visitado[v] = true;
            q.push(v);
        }
    }
}
\end{lstlisting}

\textbf{Ejemplo:}

Consideremos el siguiente grafo:

\begin{lstlisting}[style=cpp]
//        0
//       / \
//      1   2
//     / \
//    3   4
\end{lstlisting}

Si comenzamos BFS desde el vértice $0$, los niveles son:

\begin{lstlisting}[style=cpp]
// Nivel 0: 0
// Nivel 1: 1 2
// Nivel 2: 3 4
\end{lstlisting}

Por lo tanto, un posible orden de recorrido es:

\begin{lstlisting}[style=cpp]
0 1 2 3 4
\end{lstlisting}

El orden entre vértices del mismo nivel depende del orden en que
aparezcan en la lista de adyacencia.

\textbf{BFS para encontrar distancias:}

Para calcular la distancia mínima desde el origen hasta cada vértice,
se utiliza un arreglo \texttt{dist}. Inicialmente todos los valores
son $-1$, indicando que el vértice todavía no ha sido visitado.

Cuando BFS descubre un nuevo vértice $v$ desde $u$, su distancia se
calcula como:

\[
dist[v] = dist[u] + 1
\]

\begin{lstlisting}[style=cpp]
vector<int> bfsDist(int inicio, vector<vector<int>>& adj) {

    int n = adj.size();

    vector<int> dist(n, -1);
    queue<int> q;

    dist[inicio] = 0;
    q.push(inicio);

    while (!q.empty()) {

        int u = q.front();
        q.pop();

        for (int v : adj[u]) {

            if (dist[v] != -1) continue;

            dist[v] = dist[u] + 1;
            q.push(v);
        }
    }

    return dist;
}
\end{lstlisting}

Para el ejemplo anterior, comenzando desde $0$, se obtiene:

\begin{lstlisting}[style=cpp]
// dist[0] = 0
// dist[1] = 1
// dist[2] = 1
// dist[3] = 2
// dist[4] = 2
\end{lstlisting}

Por ejemplo, $dist[3] = 2$ significa que el camino mínimo desde
$0$ hasta $3$ utiliza dos aristas.

\textbf{Reconstrucción del camino:}

Si además de la distancia se necesita conocer el camino utilizado,
se puede almacenar el \textbf{padre} de cada vértice. Cuando BFS
descubre $v$ desde $u$, se establece:

\[
padre[v] = u
\]

\begin{lstlisting}[style=cpp]
vector<int> dist(n, -1);
vector<int> padre(n, -1);

queue<int> q;

dist[inicio] = 0;
q.push(inicio);

while (!q.empty()) {

    int u = q.front();
    q.pop();

    for (int v : adj[u]) {

        if (dist[v] != -1) continue;

        dist[v] = dist[u] + 1;
        padre[v] = u;

        q.push(v);
    }
}
\end{lstlisting}

Una vez terminado BFS, el camino hasta un destino se reconstruye
siguiendo los padres desde el destino hasta el origen:

\begin{lstlisting}[style=cpp]
vector<int> camino;

for (int v = destino; v != -1; v = padre[v]) {
    camino.push_back(v);
}

reverse(camino.begin(), camino.end());
\end{lstlisting}

Por ejemplo, si:

\begin{lstlisting}[style=cpp]
// padre[3] = 1
// padre[1] = 0
// padre[0] = -1
\end{lstlisting}

el camino mínimo desde $0$ hasta $3$ es:

\begin{lstlisting}[style=cpp]
0 -> 1 -> 3
\end{lstlisting}

\textbf{BFS en grafos no conectados:}

Un BFS iniciado desde un único vértice solamente visita los vértices
que pueden alcanzarse desde él. Para recorrer todo el grafo, se puede
ejecutar BFS desde cada vértice que todavía no haya sido visitado.

\begin{lstlisting}[style=cpp]
vector<bool> visitado(n, false);

for (int s = 0; s < n; s++) {

    if (visitado[s]) continue;

    queue<int> q;

    visitado[s] = true;
    q.push(s);

    while (!q.empty()) {

        int u = q.front();
        q.pop();

        for (int v : adj[u]) {

            if (visitado[v]) continue;

            visitado[v] = true;
            q.push(v);
        }
    }
}
\end{lstlisting}

Esta técnica permite encontrar, entre otras cosas, el número de
\textbf{componentes conexas} de un grafo.

\textbf{BFS para comprobar si un grafo es bipartito:}

Un grafo es bipartito si sus vértices pueden dividirse en dos grupos
de manera que ninguna arista conecte dos vértices del mismo grupo.
BFS puede comprobar esta propiedad asignando alternativamente los
colores $0$ y $1$ a los vértices.

Si un vértice $u$ tiene color $c$, sus vecinos deben tener el color
$c \mathbin{\hat{}} 1$. Si una arista conecta dos vértices del mismo
color, el grafo no es bipartito.

\begin{lstlisting}[style=cpp]
bool esBipartito(vector<vector<int>>& adj) {

    int n = adj.size();

    vector<int> color(n, -1);

    for (int inicio = 0; inicio < n; inicio++) {

        if (color[inicio] != -1) continue;

        queue<int> q;

        color[inicio] = 0;
        q.push(inicio);

        while (!q.empty()) {

            int u = q.front();
            q.pop();

            for (int v : adj[u]) {

                if (color[v] == -1) {

                    color[v] = color[u] ^ 1;
                    q.push(v);

                } else if (color[v] == color[u]) {

                    return false;
                }
            }
        }
    }

    return true;
}
\end{lstlisting}

El arreglo \texttt{color} también funciona como arreglo de visitados:
un valor de $-1$ significa que el vértice todavía no ha sido
procesado.

\textbf{BFS en un grid:}

BFS también puede utilizarse sobre matrices o tableros. Cada celda
se considera un vértice y cada movimiento válido entre celdas se
considera una arista.

Por ejemplo, si solamente se permiten movimientos hacia arriba,
abajo, izquierda y derecha:

\begin{lstlisting}[style=cpp]
int dx[] = {-1, 1, 0, 0};
int dy[] = {0, 0, -1, 1};

queue<pair<int, int>> q;

q.push({sx, sy});
dist[sx][sy] = 0;

while (!q.empty()) {

    auto [x, y] = q.front();
    q.pop();

    for (int d = 0; d < 4; d++) {

        int nx = x + dx[d];
        int ny = y + dy[d];

        if (nx < 0 || nx >= n ||
            ny < 0 || ny >= m)
            continue;

        if (dist[nx][ny] != -1)
            continue;

        if (grid[nx][ny] == '#')
            continue;

        dist[nx][ny] = dist[x][y] + 1;
        q.push({nx, ny});
    }
}
\end{lstlisting}

Esta técnica es especialmente útil en problemas donde se pregunta
por la menor cantidad de movimientos necesarios para llegar desde
una celda hasta otra.

\textbf{Multi-Source BFS:}

En algunos problemas existen varios vértices iniciales al mismo
tiempo. En lugar de ejecutar un BFS separado desde cada uno, todos
los orígenes se introducen inicialmente en la misma cola con
distancia $0$.

\begin{lstlisting}[style=cpp]
queue<int> q;
vector<int> dist(n, -1);

for (int s : fuentes) {

    dist[s] = 0;
    q.push(s);
}

while (!q.empty()) {

    int u = q.front();
    q.pop();

    for (int v : adj[u]) {

        if (dist[v] != -1) continue;

        dist[v] = dist[u] + 1;
        q.push(v);
    }
}
\end{lstlisting}

De esta manera, cada vértice obtiene la distancia mínima hasta
\textbf{cualquiera} de las fuentes.

\textbf{Complejidad:}

Utilizando una lista de adyacencia, cada vértice se procesa una sola
vez y cada arista se revisa una cantidad constante de veces.

\begin{table}[H]
\raggedright
\scriptsize
\setlength{\tabcolsep}{3pt}
\renewcommand{\arraystretch}{1.1}
\begin{tabular}{|l|c|c|}
\hline
\textbf{Representación} &
\textbf{Tiempo} &
\textbf{Memoria adicional} \\
\hline
Lista de adyacencia &
$O(n+m)$ &
$O(n)$ \\
\hline
Matriz de adyacencia &
$O(n^2)$ &
$O(n)$ \\
\hline
\end{tabular}
\caption{Complejidad de BFS según la representación del grafo.}
\end{table}

La complejidad $O(n+m)$ hace que BFS sea especialmente eficiente
para grafos dispersos representados mediante listas de adyacencia.

\textbf{Nota:} BFS encuentra caminos mínimos cuando todas las aristas
tienen el mismo costo. No debe utilizarse directamente para encontrar
caminos mínimos en grafos con pesos arbitrarios. Para pesos no
negativos se puede utilizar \textbf{Dijkstra}; si los pesos solamente
son $0$ y $1$, una alternativa eficiente es \textbf{0-1 BFS}.

\textbf{Regla práctica:} si un problema pide minimizar el número de
\textbf{pasos, movimientos o aristas} y cada movimiento tiene el mismo
costo, considera primero \textbf{BFS}. Si necesitas conocer la
distancia, utiliza \texttt{dist}; si necesitas recuperar el camino,
utiliza \texttt{padre}; y si existen múltiples puntos de origen,
considera \textbf{Multi-Source BFS}.
```


```latex
\subsection{BFS en Grafos Desconectados}

\textbf{Idea:} un BFS iniciado desde un único vértice solamente visita
los vértices que son alcanzables desde él. Si el grafo tiene varias
\textbf{componentes conexas}, es necesario iniciar un nuevo BFS desde
cada vértice que todavía no haya sido visitado.

La idea consiste en recorrer todos los vértices y, cada vez que se
encuentra uno que no ha sido visitado, iniciar un nuevo BFS desde él.
Cada ejecución de BFS procesa una componente conexa completa.

\textbf{Implementación:}

\begin{lstlisting}[style=cpp]
void bfs(int inicio, vector<vector<int>>& adj,
         vector<bool>& visitado) {

    queue<int> q;

    visitado[inicio] = true;
    q.push(inicio);

    while (!q.empty()) {

        int u = q.front();
        q.pop();

        for (int v : adj[u]) {

            if (visitado[v]) continue;

            visitado[v] = true;
            q.push(v);
        }
    }
}
\end{lstlisting}

Para recorrer todo el grafo:

\begin{lstlisting}[style=cpp]
vector<bool> visitado(n, false);

for (int u = 0; u < n; u++) {

    if (visitado[u]) continue;

    bfs(u, adj, visitado);
}
\end{lstlisting}

Cada vez que se ejecuta \texttt{bfs}, se procesa una componente conexa
que todavía no había sido visitada.

\textbf{Contar componentes conexas:}

Una aplicación directa consiste en contar cuántas componentes conexas
tiene el grafo. Basta incrementar un contador cada vez que se inicia
un nuevo BFS.

\begin{lstlisting}[style=cpp]
int componentes = 0;

vector<bool> visitado(n, false);

for (int u = 0; u < n; u++) {

    if (visitado[u]) continue;

    componentes++;

    queue<int> q;

    visitado[u] = true;
    q.push(u);

    while (!q.empty()) {

        int v = q.front();
        q.pop();

        for (int w : adj[v]) {

            if (visitado[w]) continue;

            visitado[w] = true;
            q.push(w);
        }
    }
}

cout << componentes << '\n';
\end{lstlisting}

\textbf{Ejemplo:}

Consideremos el siguiente grafo:

\begin{lstlisting}[style=cpp]
//    0 --- 1       3 --- 4
//          |
//          2
//
//    5 --- 6
\end{lstlisting}

El grafo tiene tres componentes conexas:

\begin{lstlisting}[style=cpp]
// Componente 1: {0, 1, 2}
// Componente 2: {3, 4}
// Componente 3: {5, 6}
\end{lstlisting}

Si recorremos los vértices en orden, el primer BFS puede comenzar
desde $0$ y visitar $\{0,1,2\}$. Después, los vértices $1$ y $2$ ya
están visitados, por lo que el siguiente BFS comienza desde $3$ y
visita $\{3,4\}$. Finalmente, se inicia otro BFS desde $5$ para
visitar $\{5,6\}$.

\textbf{Complejidad:}

Aunque se ejecuten varios BFS, cada vértice se visita una sola vez y
cada arista se procesa una cantidad constante de veces. Por lo tanto,
la complejidad total utilizando una lista de adyacencia sigue siendo:

\[
O(n+m)
\]

La memoria adicional utilizada por el arreglo de visitados y la cola
es $O(n)$.

\begin{table}[H]
\raggedright
\scriptsize
\setlength{\tabcolsep}{3pt}
\renewcommand{\arraystretch}{1.1}
\begin{tabular}{|l|c|}
\hline
\textbf{Operación} & \textbf{Complejidad} \\
\hline
Recorrer todas las componentes & $O(n+m)$ \\
\hline
Contar componentes conexas & $O(n+m)$ \\
\hline
Memoria adicional & $O(n)$ \\
\hline
\end{tabular}
\caption{Complejidad de BFS en grafos desconectados.}
\end{table}

\textbf{Regla práctica:} si el grafo puede estar desconectado y se
necesita procesar \textbf{todos los vértices}, no basta con ejecutar
BFS una sola vez. Recorre todos los vértices y comienza un nuevo BFS
cada vez que encuentres uno que todavía no haya sido visitado.
```


\subsection{BFS para Todos los Caminos}

\textbf{Idea:} BFS no solamente permite encontrar la distancia mínima
desde un vértice origen hacia los demás vértices. También puede
utilizarse para reconstruir el \textbf{camino mínimo} desde el origen
hasta cualquier vértice alcanzable.

Para reconstruir un camino, además de almacenar la distancia, se
guarda para cada vértice su \textbf{padre}, es decir, el vértice desde
el cual fue descubierto por primera vez. Como BFS visita los vértices
por niveles, el primer padre asignado a un vértice pertenece a un
camino mínimo.

\textbf{Implementación:}

\begin{lstlisting}[style=cpp]
vector<int> dist(n, -1);
vector<int> padre(n, -1);

queue<int> q;

dist[inicio] = 0;
q.push(inicio);

while (!q.empty()) {

    int u = q.front();
    q.pop();

    for (int v : adj[u]) {

        if (dist[v] != -1) continue;

        dist[v] = dist[u] + 1;
        padre[v] = u;

        q.push(v);
    }
}
\end{lstlisting}

Después de ejecutar BFS, \texttt{dist[v]} contiene la distancia mínima
desde \texttt{inicio} hasta $v$, mientras que \texttt{padre[v]} indica
qué vértice se debe seguir para reconstruir el camino.

\textbf{Reconstrucción del camino:}

Para obtener el camino hasta un vértice destino, se comienza en
\texttt{destino} y se siguen los padres hasta llegar al origen.

\begin{lstlisting}[style=cpp]
vector<int> camino;

for (int v = destino; v != -1; v = padre[v]) {
    camino.push_back(v);
}

reverse(camino.begin(), camino.end());
\end{lstlisting}

El vector \texttt{camino} contiene ahora los vértices en el orden
correcto, desde el origen hasta el destino.

\textbf{Ejemplo:}

Supongamos que BFS produce los siguientes padres:

\begin{lstlisting}[style=cpp]
// padre[0] = -1
// padre[1] = 0
// padre[2] = 0
// padre[3] = 1
// padre[4] = 3
\end{lstlisting}

Si el origen es $0$ y queremos reconstruir el camino hasta $4$,
seguimos los padres:

\[
4 \rightarrow 3 \rightarrow 1 \rightarrow 0
\]

Después de invertirlo obtenemos:

\[
0 \rightarrow 1 \rightarrow 3 \rightarrow 4
\]

Por lo tanto, este es un camino mínimo desde $0$ hasta $4$.

\textbf{Función para reconstruir el camino:}

La reconstrucción puede encapsularse en una función:

\begin{lstlisting}[style=cpp]
vector<int> reconstruirCamino(int destino,
                              vector<int>& padre) {

    vector<int> camino;

    for (int v = destino; v != -1; v = padre[v]) {
        camino.push_back(v);
    }

    reverse(camino.begin(), camino.end());

    return camino;
}
\end{lstlisting}

El BFS completo puede combinarse con esta función:

\begin{lstlisting}[style=cpp]
vector<int> dist(n, -1);
vector<int> padre(n, -1);

queue<int> q;

dist[inicio] = 0;
q.push(inicio);

while (!q.empty()) {

    int u = q.front();
    q.pop();

    for (int v : adj[u]) {

        if (dist[v] != -1) continue;

        dist[v] = dist[u] + 1;
        padre[v] = u;

        q.push(v);
    }
}

vector<int> camino =
    reconstruirCamino(destino, padre);
\end{lstlisting}

\textbf{¿Qué significa ``todos los caminos''?}

En este contexto, BFS obtiene un camino mínimo desde el origen hacia
\textbf{cada vértice alcanzable}. No significa que encuentre todos los
caminos posibles entre dos vértices, ya que la cantidad de caminos
posibles puede ser exponencial.

Por ejemplo, después de ejecutar BFS desde $s$:

\begin{itemize}
    \item \texttt{dist[v]} indica la distancia mínima de $s$ a $v$.
    \item \texttt{padre[v]} permite reconstruir un camino mínimo hacia
    $v$.
    \item Si \texttt{dist[v]} es $-1$, $v$ no es alcanzable desde $s$.
\end{itemize}

\textbf{Camino inexistente:}

Antes de reconstruir el camino conviene comprobar si el destino fue
alcanzado:

\begin{lstlisting}[style=cpp]
if (dist[destino] == -1) {
    cout << "No existe camino\n";
}
else {
    vector<int> camino =
        reconstruirCamino(destino, padre);

    for (int v : camino) {
        cout << v << ' ';
    }
}
\end{lstlisting}

\textbf{Complejidad:}

La ejecución de BFS utilizando una lista de adyacencia cuesta
$O(n+m)$. La reconstrucción de un camino cuesta $O(k)$, donde $k$ es
la cantidad de vértices que contiene el camino.

\begin{table}[H]
\raggedright
\scriptsize
\setlength{\tabcolsep}{3pt}
\renewcommand{\arraystretch}{1.1}
\begin{tabular}{|l|c|}
\hline
\textbf{Operación} & \textbf{Complejidad} \\
\hline
BFS & $O(n+m)$ \\
\hline
Reconstruir un camino & $O(k)$ \\
\hline
Memoria adicional & $O(n)$ \\
\hline
\end{tabular}
\caption{Complejidad de BFS para reconstrucción de caminos.}
\end{table}

\textbf{Nota:} BFS encuentra caminos mínimos en grafos donde todas las
aristas tienen el mismo costo. Si las aristas tienen pesos diferentes,
deben utilizarse algoritmos como \textbf{Dijkstra}, \textbf{0-1 BFS} o
\textbf{Bellman-Ford}, dependiendo de las características del problema.

\textbf{Regla práctica:} si el problema pide no solamente la distancia
mínima sino también \textbf{qué vértices forman el camino}, guarda un
arreglo \texttt{padre} durante el BFS. No es necesario almacenar todos
los caminos: basta con guardar un padre por vértice.

\subsection{DFS}

\textbf{Idea:} Depth-First Search (DFS) explora un grafo siguiendo un
camino lo más profundamente posible antes de regresar y explorar otras
ramas. A diferencia de BFS, que utiliza una \textbf{cola}, DFS utiliza
una \textbf{pila}.

DFS puede implementarse de forma recursiva o iterativa. En Competitive
Programming, la versión iterativa es útil cuando el grafo puede tener
muchos vértices y una implementación recursiva podría provocar un
\textbf{stack overflow}.

\textbf{Implementación iterativa:}

Se utiliza un \texttt{stack<int>} para almacenar los vértices que deben
ser procesados. Un vértice se marca como visitado cuando se introduce
en la pila, evitando que pueda agregarse múltiples veces.

\begin{lstlisting}[style=cpp]
void dfs(int inicio, vector<vector<int>>& adj,
         vector<bool>& visitado) {

    stack<int> st;

    visitado[inicio] = true;
    st.push(inicio);

    while (!st.empty()) {

        int u = st.top();
        st.pop();

        cout << u << ' ';

        for (int v : adj[u]) {

            if (visitado[v]) continue;

            visitado[v] = true;
            st.push(v);
        }
    }
}
\end{lstlisting}

El orden exacto en que se visitan los vértices depende del orden de los
vecinos en la lista de adyacencia. Lo importante es que cada vértice
alcanzable sea procesado una vez.

\textbf{Ejemplo:}

Consideremos el siguiente grafo:

\begin{lstlisting}[style=cpp]
//        0
//       / \
//      1   2
//     / \
//    3   4
\end{lstlisting}

La lista de adyacencia puede ser:

\begin{lstlisting}[style=cpp]
adj[0] = {1, 2};
adj[1] = {0, 3, 4};
adj[2] = {0};
adj[3] = {1};
adj[4] = {1};
\end{lstlisting}

Si comenzamos DFS desde $0$, la pila inicialmente contiene:

\begin{lstlisting}[style=cpp]
st = {0}
\end{lstlisting}

Se extrae $0$ y se agregan sus vecinos:

\begin{lstlisting}[style=cpp]
st = {1, 2}
\end{lstlisting}

El siguiente vértice que se extrae depende del orden en que fueron
agregados a la pila. Por ejemplo, si se extrae primero $2$, se procesa
esa rama antes de regresar a $1$.

\textbf{DFS en un grafo desconectado:}

Al igual que BFS, una sola ejecución de DFS solamente visita los
vértices alcanzables desde el origen. Para recorrer todo el grafo,
debemos iniciar un nuevo DFS desde cada vértice que todavía no haya
sido visitado.

\begin{lstlisting}[style=cpp]
vector<bool> visitado(n, false);

for (int u = 0; u < n; u++) {

    if (visitado[u]) continue;

    dfs(u, adj, visitado);
}
\end{lstlisting}

Este patrón permite procesar todas las componentes conexas del grafo.

\textbf{Contar componentes conexas:}

Podemos contar el número de componentes incrementando un contador cada
vez que sea necesario iniciar un nuevo DFS.

\begin{lstlisting}[style=cpp]
int componentes = 0;

vector<bool> visitado(n, false);

for (int u = 0; u < n; u++) {

    if (visitado[u]) continue;

    componentes++;
    dfs(u, adj, visitado);
}

cout << componentes << '\n';
\end{lstlisting}

\textbf{DFS en árboles:}

En un árbol no es necesario mantener un arreglo de visitados si se
conoce el padre del vértice actual. Al recorrer un vecino, simplemente
se evita regresar al padre.

Sin embargo, en un grafo general se recomienda utilizar
\texttt{visitado}, ya que pueden existir ciclos.

\begin{lstlisting}[style=cpp]
void dfs(int inicio, vector<vector<int>>& adj) {

    stack<pair<int, int>> st;

    st.push({inicio, -1});

    while (!st.empty()) {

        auto [u, padre] = st.top();
        st.pop();

        cout << u << ' ';

        for (int v : adj[u]) {

            if (v == padre) continue;

            st.push({v, u});
        }
    }
}
\end{lstlisting}

Esta última versión es apropiada para árboles, pero para grafos
generales debe utilizarse un arreglo de visitados.

\textbf{Aplicaciones:}

DFS es una de las herramientas fundamentales para resolver problemas
de grafos. Entre sus aplicaciones más comunes se encuentran:

\begin{itemize}
    \item Recorrer todos los vértices alcanzables desde un origen.
    \item Encontrar componentes conexas.
    \item Detectar ciclos.
    \item Construir árboles DFS.
    \item Encontrar puntos de articulación y puentes.
    \item Realizar ordenamiento topológico.
    \item Resolver problemas de conectividad.
    \item Recorrer árboles y estructuras jerárquicas.
\end{itemize}

\textbf{Complejidad:}

Utilizando una lista de adyacencia, cada vértice se procesa una sola
vez y cada arista se revisa una cantidad constante de veces. Por lo
tanto:

\[
O(n+m)
\]

La pila y el arreglo de visitados requieren $O(n)$ memoria adicional.

\begin{table}[H]
\raggedright
\scriptsize
\setlength{\tabcolsep}{3pt}
\renewcommand{\arraystretch}{1.1}
\begin{tabular}{|l|c|c|}
\hline
\textbf{Representación} &
\textbf{Tiempo} &
\textbf{Memoria adicional} \\
\hline
Lista de adyacencia &
$O(n+m)$ &
$O(n)$ \\
\hline
Matriz de adyacencia &
$O(n^2)$ &
$O(n)$ \\
\hline
\end{tabular}
\caption{Complejidad de DFS según la representación del grafo.}
\end{table}

\textbf{Nota:} DFS y BFS tienen la misma complejidad asintótica
$O(n+m)$ utilizando una lista de adyacencia, pero exploran el grafo de
manera diferente. BFS procesa por niveles mediante una cola, mientras
que DFS profundiza utilizando una pila.

\textbf{Regla práctica:} utiliza DFS cuando el problema dependa de
explorar profundamente una estructura, recorrer componentes, detectar
ciclos o realizar algún procesamiento basado en el orden de entrada y
salida de los vértices. Si el problema requiere distancias mínimas en
un grafo no ponderado, normalmente BFS es la opción adecuada.

\subsection{DFS para Conteo de Caminos}

\textbf{Idea:} en un DAG (grafo dirigido acíclico), contar cuántos
caminos distintos existen entre un nodo $u$ y un nodo destino fijo
es un problema de \textbf{DFS + memoización}: el número de caminos
desde $u$ hasta el destino es la suma de los caminos desde cada
vecino de $u$ hasta el destino, con caso base "ya llegué". Como
varios caminos pueden pasar por el mismo nodo intermedio, sin
memoización recalculas el mismo subárbol muchas veces —memoizar es
lo que baja esto de exponencial a lineal.

\textbf{¿Por qué usarlo?} Es el mismo patrón de "DFS + memo" que ya
usas en DP sobre árboles/DAGs, aplicado específicamente a contar
(no a minimizar/maximizar). Reconocer que un problema pide "número
de formas de llegar de A a B" en un grafo dirigido sin ciclos es la
señal para aplicar esto en vez de enumerar caminos explícitamente.

\textbf{Casos de uso:}

\begin{itemize}
	\item Contar caminos entre dos nodos fijos de un DAG.
	\item Contar, para cada nodo, cuántos caminos empiezan en él y
	terminan en algún nodo "hoja" (sin salidas) — mismo memo, sin
	destino fijo, sumando 1 en cada hoja.
	\item Variante con longitud exacta $k$: contar caminos de $u$ a
	$v$ usando \textbf{exactamente} $k$ aristas, memoizando sobre el
	par $(nodo, pasos\_restantes)$ en vez de solo $nodo$.
	\item Contar caminos módulo algún $M$ cuando el número de
	caminos crece demasiado (usa \texttt{long long} y aplica \%M en
	cada suma).
\end{itemize}

\textbf{Complejidad:} $O(V + E)$ con memoización, ya que cada nodo
se procesa una sola vez y cada arista se recorre una sola vez. Sin
memo, en el peor caso es exponencial (puede haber $O(2^V)$ caminos).

\begin{lstlisting}[style=cpp]
	struct ConteoCaminos {
		vector<vector<int>> adj;
		vector<long long> memo;
		int destino;
		
		// CREATE: guarda el grafo (debe ser DAG) y el nodo destino;
		// memo[u] = -1 significa "no calculado todavia".
		ConteoCaminos(vector<vector<int>>& adj_, int destino_) {
			adj = adj_;
			destino = destino_;
			memo.assign(adj.size(), -1);
		}                                          // O(1)
		
		// READ: numero de caminos distintos de u a destino.
		long long contar(int u) {
			if (u == destino) return 1;
			if (memo[u] != -1) return memo[u];
			
			long long total = 0;
			for (int v : adj[u])
			total += contar(v);
			
			return memo[u] = total;
		}                                          // O(1) amortizado
		// (O(V+E) la primera vez)
	};
\end{lstlisting}

\textbf{Ejemplo de funcionamiento:}

\begin{lstlisting}[style=cpp]
	// DAG:      0 -> 1 -> 3
	//           0 -> 2 -> 3
	//           1 -> 2
	
	vector<vector<int>> adj(4);
	adj[0] = {1, 2};
	adj[1] = {2, 3};
	adj[2] = {3};
	adj[3] = {};
	
	ConteoCaminos cc(adj, 3);
	
	cout << cc.contar(0);
	// 3  (0-1-3, 0-1-2-3, 0-2-3)
	
	cout << cc.contar(1);
	// 2  (1-3, 1-2-3)
\end{lstlisting}

\textbf{Nota:} si el grafo tiene ciclos, esta memoización no funciona
directamente (la recursión no termina) — primero hay que condensar
las componentes fuertemente conexas (SCC) en un DAG, y contar sobre
ese DAG condensado.


\subsection{DFS + Cerramiento Transitivo}

\textbf{Idea:} el cerramiento transitivo de un grafo dirigido
responde, para cualquier par $(u,v)$, la pregunta "¿existe algún
camino de $u$ a $v$?" en $O(1)$ por query, a costa de precomputar
—para cada nodo— el conjunto completo de nodos que alcanza. La forma
más directa es correr un DFS desde \textbf{cada} nodo y marcar todo
lo que visita.

\textbf{¿Por qué usarlo?} Cuando el grafo es pequeño/mediano y vas a
hacer \textbf{muchas} consultas de alcanzabilidad, precomputar todo
de una vez sale más barato que un DFS/BFS por consulta. Es el mismo
principio de "precomputo una vez, respondo O(1)" que en RMQ, aplicado
a alcanzabilidad en vez de a mínimos.

\textbf{Casos de uso:}

\begin{itemize}
	\item Responder $Q$ queries de "¿$u$ puede llegar a $v$?" cuando
	$Q$ es grande y $V$ es manejable ($V \le \sim 2000$–$5000$ según
	el límite de tiempo).
	\item Precomputar dependencias transitivas (si $a$ depende de
	$b$ y $b$ depende de $c$, entonces $a$ depende de $c$).
	\item Como paso intermedio antes de comprimir un DAG por SCC:
	primero condensas ciclos, y sobre el DAG resultante calculas el
	cerramiento transitivo.
\end{itemize}

\textbf{Complejidad:} $O(V \cdot (V+E))$ para construir (un DFS
completo por cada nodo), $O(1)$ por query. Con \texttt{std::bitset}
en vez de \texttt{vector<bool>}, y procesando los nodos en orden
topológico inverso (\texttt{reach[u] |= reach[v]} para cada vecino
$v$), el build baja a $O(V \cdot (V+E) / 64)$ gracias al paralelismo
a nivel de bits — útil si $V$ es más grande.

\begin{lstlisting}[style=cpp]
	struct CerramientoTransitivo {
		int n;
		vector<vector<int>> adj;
		vector<vector<bool>> alcanzable;   // alcanzable[u][v] = true si
		// u puede llegar a v
		
		// CREATE: corre un DFS completo desde cada nodo. O(V*(V+E)).
		CerramientoTransitivo(vector<vector<int>>& adj_, int n_) {
			adj = adj_;
			n = n_;
			alcanzable.assign(n, vector<bool>(n, false));
			
			for (int origen = 0; origen < n; origen++)
			dfs(origen, origen);
		}
		
		void dfs(int origen, int u) {
			for (int v : adj[u]) {
				if (!alcanzable[origen][v]) {
					alcanzable[origen][v] = true;
					dfs(origen, v);
				}
			}
		}
		
		// READ: es u puede llegar a v.
		bool query(int u, int v) {
			return alcanzable[u][v];
		}                                          // O(1)
	};
\end{lstlisting}

\textbf{Ejemplo de funcionamiento:}

\begin{lstlisting}[style=cpp]
	// grafo:    0 -> 1 -> 2
	//                \
	//                 -> 3
	//           4 (aislado)
	
	vector<vector<int>> adj(5);
	adj[0] = {1};
	adj[1] = {2, 3};
	adj[2] = {};
	adj[3] = {};
	adj[4] = {};
	
	CerramientoTransitivo ct(adj, 5);
	
	cout << ct.query(0, 3);
	// true  (0 -> 1 -> 3)
	
	cout << ct.query(2, 0);
	// false (no hay camino de 2 a 0)
	
	cout << ct.query(4, 0);
	// false (4 esta aislado)
\end{lstlisting}

\textbf{Regla práctica:} DFS "para contar" y DFS "para cerramiento
transitivo" comparten la misma mecánica de recorrido, pero acumulan
cosas distintas en cada llamada: uno suma caminos hacia un destino
fijo (memoizando por nodo), el otro marca alcanzabilidad hacia
\emph{todos} los destinos posibles (memoizando por par de nodos,
implícitamente, al no revisitar lo ya marcado). Si el grafo tiene
ciclos, ambos requieren primero pasar por condensación de SCC antes
de aplicarse tal cual.

\subsection{Flood Fill}

\textbf{Idea:} Flood Fill es "explorar todas las celdas conectadas a
una celda inicial que cumplen una condición" (mismo color, mismo
carácter, sin obstáculo, etc.), marcándolas como visitadas o
recoloreándolas. Es exactamente un DFS/BFS pero sobre una grilla en
vez de un grafo con listas de adyacencia explícitas: los "vecinos"
de una celda $(r,c)$ son sus hasta 4 (u 8) celdas adyacentes.

\textbf{¿Por qué usarlo?} Cualquier problema de grilla que hable de
"regiones", "islas", "áreas conectadas" o "pintar/rellenar" es Flood
Fill. La clave práctica es preferir la versión \textbf{iterativa}
(con \texttt{queue} o \texttt{stack} explícito) sobre la recursiva:
en grillas grandes ($R \times C$ del orden de $10^5$–$10^6$ celdas),
un DFS recursivo puede hacer \textit{stack overflow}.

\textbf{Casos de uso:}

\begin{itemize}
	\item Contar componentes conexas de una grilla ("número de
	islas"): un Flood Fill por cada celda no visitada que cumpla la
	condición, contando cuántas veces arrancas uno nuevo.
	\item Calcular el área (tamaño) de cada región conectada.
	\item "Paint bucket": recolorear toda una región conectada a un
	color nuevo.
	\item Determinar si dos celdas están en la misma región conectada.
\end{itemize}

\textbf{Complejidad:} $O(R \cdot C)$, ya que cada celda se visita
(y se encola) como máximo una vez.

\begin{lstlisting}[style=cpp]
	struct FloodFill {
		vector<vector<char>> grid;
		vector<vector<bool>> visitado;
		int R, C;
		
		int dr[4] = {-1, 1, 0, 0};
		int dc[4] = {0, 0, -1, 1};
		
		// CREATE: guarda la grilla (R x C).
		FloodFill(vector<vector<char>>& grid_) {
			grid = grid_;
			R = grid.size();
			C = grid[0].size();
			visitado.assign(R, vector<bool>(C, false));
		}                                          // O(1)
		
		// READ: marca (BFS iterativo) toda la region conectada a
		// (sr, sc) que comparte el mismo caracter, y devuelve su area.
		// Evita recursion para no hacer stack overflow en grillas grandes.
		int llenar(int sr, int sc) {
			if (visitado[sr][sc]) return 0;
			
			char objetivo = grid[sr][sc];
			queue<pair<int,int>> q;
			q.push({sr, sc});
			visitado[sr][sc] = true;
			int area = 0;
			
			while (!q.empty()) {
				auto [r, c] = q.front(); q.pop();
				area++;
				
				for (int d = 0; d < 4; d++) {
					int nr = r + dr[d], nc = c + dc[d];
					
					if (nr < 0 || nr >= R || nc < 0 || nc >= C) continue;
					if (visitado[nr][nc]) continue;
					if (grid[nr][nc] != objetivo) continue;
					
					visitado[nr][nc] = true;
					q.push({nr, nc});
				}
			}
			
			return area;
		}                                          // O(area de la region)
	};
\end{lstlisting}

\textbf{Ejemplo de funcionamiento:}

\begin{lstlisting}[style=cpp]
	// grid:
	// L L W W
	// L L W W
	// W W L L
	
	vector<vector<char>> grid = {
		{'L','L','W','W'},
		{'L','L','W','W'},
		{'W','W','L','L'}
	};
	
	FloodFill ff(grid);
	
	int islas = 0, areaMax = 0;
	for (int r = 0; r < 3; r++)
	for (int c = 0; c < 4; c++)
	if (grid[r][c] == 'L' && !ff.visitado[r][c]) {
		int area = ff.llenar(r, c);
		islas++;
		areaMax = max(areaMax, area);
	}
	
	cout << islas;    // 2
	cout << areaMax;  // 4  (el bloque de 2x2 de arriba-izquierda)
\end{lstlisting}


\subsection{Laberinto con BFS}

\textbf{Idea:} en una grilla donde te mueves de una celda a una
celda adyacente (arriba/abajo/izq/der) con \textbf{costo uniforme}
(cada paso cuesta 1), la ruta más corta desde un origen se calcula
con BFS: exploras la grilla "por capas" de distancia creciente, y la
primera vez que llegas a una celda, esa es su distancia mínima. Es el
mismo BFS de grafos de siempre, con la grilla actuando como grafo
implícito.

\textbf{¿Por qué usarlo?} Es la herramienta por defecto para
"camino más corto en grilla sin pesos / con obstáculos". Si el
problema menciona pesos distintos por celda (terreno con costo), esto
deja de ser BFS puro y pasa a ser Dijkstra o 0-1 BFS (deque) — pero
mientras cada paso cueste lo mismo, BFS es suficiente y más simple.

\textbf{Casos de uso:}

\begin{itemize}
	\item Número mínimo de pasos de un punto $A$ a un punto $B$ en
	un laberinto con paredes.
	\item Reconstruir el camino más corto explícito (no solo la
	distancia), guardando de dónde vino cada celda.
	\item BFS multi-fuente: varias celdas iniciales a la vez (p.ej.
	"tiempo mínimo para que el fuego llegue a cada celda"),
	encolando todas las fuentes con distancia 0 antes de empezar.
	\item Verificar alcanzabilidad ("¿se puede llegar de A a B?").
\end{itemize}

\textbf{Complejidad:} $O(R \cdot C)$, ya que cada celda se encola y
desencola como máximo una vez.

\begin{lstlisting}[style=cpp]
	struct LaberintoBFS {
		vector<vector<char>> grid;
		vector<vector<int>> dist;
		vector<vector<pair<int,int>>> padre;
		int R, C;
		
		int dr[4] = {-1, 1, 0, 0};
		int dc[4] = {0, 0, -1, 1};
		
		// CREATE: guarda la grilla. '#' = pared, cualquier otro caracter
		// = celda libre.
		LaberintoBFS(vector<vector<char>>& grid_) {
			grid = grid_;
			R = grid.size();
			C = grid[0].size();
		}                                          // O(1)
		
		// READ: corre BFS desde (sr, sc) y llena dist / padre para
		// todas las celdas alcanzables. dist queda en -1 donde no llego.
		void bfs(int sr, int sc) {
			dist.assign(R, vector<int>(C, -1));
			padre.assign(R, vector<pair<int,int>>(C, {-1, -1}));
			
			queue<pair<int,int>> q;
			q.push({sr, sc});
			dist[sr][sc] = 0;
			
			while (!q.empty()) {
				auto [r, c] = q.front(); q.pop();
				
				for (int d = 0; d < 4; d++) {
					int nr = r + dr[d], nc = c + dc[d];
					
					if (nr < 0 || nr >= R || nc < 0 || nc >= C) continue;
					if (grid[nr][nc] == '#') continue;
					if (dist[nr][nc] != -1) continue;
					
					dist[nr][nc] = dist[r][c] + 1;
					padre[nr][nc] = {r, c};
					q.push({nr, nc});
				}
			}
		}                                          // O(R*C)
		
		// READ: reconstruye el camino desde el origen del ultimo bfs()
		// hasta (er, ec), como lista de celdas en orden. Vacio si no
		// hay camino.
		vector<pair<int,int>> reconstruirCamino(int er, int ec) {
			if (dist[er][ec] == -1) return {};
			
			vector<pair<int,int>> camino;
			pair<int,int> actual = {er, ec};
			
			while (actual.first != -1) {
				camino.push_back(actual);
				actual = padre[actual.first][actual.second];
			}
			
			reverse(camino.begin(), camino.end());
			return camino;
		}                                          // O(longitud del camino)
	};
\end{lstlisting}

\textbf{Ejemplo de funcionamiento:}

\begin{lstlisting}[style=cpp]
	// grid:
	// . . . #
	// # # . #
	// . . . .
	
	vector<vector<char>> grid = {
		{'.','.','.','#'},
		{'#','#','.','#'},
		{'.','.','.','.'}
	};
	
	LaberintoBFS lb(grid);
	lb.bfs(0, 0);
	
	cout << lb.dist[2][3];
	// 5  (0,0)->(0,1)->(0,2)->(1,2)->(2,2)->(2,3)
	
	auto camino = lb.reconstruirCamino(2, 3);
	// [(0,0),(0,1),(0,2),(1,2),(2,2),(2,3)]
\end{lstlisting}

\textbf{Regla práctica:} si un problema de grilla pide "¿son
conectados?" o "área de una región", es Flood Fill (DFS o BFS, no
importa el orden de visita). Si pide "distancia mínima" o "menor
número de pasos", es BFS por capas (el orden de visita sí importa:
BFS garantiza que la primera vez que tocas una celda es con su
distancia mínima; DFS no lo garantiza).

\subsection{Tour de Caballo de Ajedrez}

\textbf{Idea:} el Tour de Caballo pide encontrar una secuencia de
movimientos de caballo que visite \textbf{cada casilla del tablero
	exactamente una vez}. No es un problema de "camino más corto" (no
hay pesos ni se compara contra alternativas), sino de
\textbf{existencia}: es backtracking puro — DFS que prueba un
movimiento, sigue explorando, y si se atasca (llega a una casilla sin
salidas válidas) deshace ("backtrack") y prueba otro movimiento.

\textbf{¿Por qué usarlo?} A diferencia de Flood Fill o BFS en
laberintos, aquí \textbf{no basta con visitar} las casillas
alcanzables: tienen que visitarse \textbf{todas, sin repetir}, en
un orden válido. Eso descarta BFS/DFS simple y obliga a backtracking
con una condición de poda, porque la fuerza bruta pura ($8^{n}$
combinaciones de movimientos posibles) es intratable incluso para
tableros chicos.

\textbf{Casos de uso:}

\begin{itemize}
	\item Encontrar (o verificar la existencia de) un tour de caballo
	completo en un tablero $n \times n$.
	\item Tour \textbf{cerrado} (el último movimiento vuelve a
	atacar la casilla inicial): misma idea, chequeo extra al final.
	\item Cualquier problema de "recorrer todas las celdas exactamente
	una vez respetando un patrón de movimiento" (no exclusivo del
	caballo): el patrón general es backtracking + poda de
	Warnsdorff.
\end{itemize}

\textbf{Complejidad:} exponencial en el peor caso sin poda
($O(8^{n^2})$ aproximado). Con la \textbf{heurística de Warnsdorff}
—en cada paso, moverse primero a la casilla accesible con
\textbf{menos} salidas futuras disponibles— en la práctica encuentra
un tour en tiempo casi lineal para tableros razonables ($n \le 100$),
aunque no hay garantía teórica formal de que siempre funcione.

\begin{lstlisting}[style=cpp]
	struct TourCaballo {
		int n;
		vector<vector<int>> tablero;   // tablero[r][c] = orden de visita
		// (-1 si no visitada)
		int dr[8] = {-2,-2,-1,-1, 1, 1, 2, 2};
		int dc[8] = {-1, 1,-2, 2,-2, 2,-1, 1};
		
		// CREATE: tablero n x n vacio.
		TourCaballo(int n_) {
			n = n_;
			tablero.assign(n, vector<int>(n, -1));
		}                                          // O(n^2)
		
		bool dentro(int r, int c) {
			return r >= 0 && r < n && c >= 0 && c < n;
		}
		
		// cuenta a cuantas casillas libres puede saltar (r,c) - usado
		// por Warnsdorff para ordenar los candidatos.
		int gradoSalida(int r, int c) {
			int cont = 0;
			for (int d = 0; d < 8; d++) {
				int nr = r + dr[d], nc = c + dc[d];
				if (dentro(nr, nc) && tablero[nr][nc] == -1) cont++;
			}
			return cont;
		}
		
		// READ: intenta construir un tour completo empezando en
		// (sr, sc). Devuelve true si lo logro (tablero queda lleno).
		bool resolver(int sr, int sc) {
			tablero[sr][sc] = 0;
			return backtrack(sr, sc, 1);
		}                                          // O(casi lineal) con
		// Warnsdorff en la
		// practica
		
		bool backtrack(int r, int c, int movimiento) {
			if (movimiento == n * n) return true;
			
			// candidatos ordenados por Warnsdorff: menor grado de
			// salida primero
			vector<tuple<int,int,int>> candidatos;
			for (int d = 0; d < 8; d++) {
				int nr = r + dr[d], nc = c + dc[d];
				if (dentro(nr, nc) && tablero[nr][nc] == -1)
				candidatos.push_back({gradoSalida(nr, nc), nr, nc});
			}
			sort(candidatos.begin(), candidatos.end());
			
			for (auto& [grado, nr, nc] : candidatos) {
				tablero[nr][nc] = movimiento;
				
				if (backtrack(nr, nc, movimiento + 1)) return true;
				
				tablero[nr][nc] = -1;   // deshacer (backtrack real)
			}
			
			return false;
		}
	};
\end{lstlisting}

\textbf{Ejemplo de funcionamiento:}

\begin{lstlisting}[style=cpp]
	TourCaballo tc(5);
	
	bool encontrado = tc.resolver(0, 0);
	
	cout << encontrado;
	// true
	
	// tc.tablero[r][c] tiene el orden (0..24) en que se visito
	// cada casilla, por ejemplo:
	//  0  3 16  9 22
	// 13  8 21  4 17
	//  2  1 12 23 10
	//  7 14 19 18  5
	// 20  ... etc.
\end{lstlisting}

\textbf{Nota:} sin la poda de Warnsdorff, el backtrack puro también
es correcto pero mucho más lento (puede no terminar en tiempo útil
para $n > 6$ u $8$). Warnsdorff no cambia la corrección del
algoritmo, solo el \textbf{orden} en que se prueban los candidatos —
sigue siendo backtracking con deshacer explícito.


\subsection{Mínimo Número de Movimientos de Caballo}

\textbf{Idea:} a diferencia del Tour de Caballo, aquí no importa
visitar todo el tablero — solo se pide la \textbf{distancia mínima}
(en número de saltos) entre una casilla origen y una casilla
destino. Eso es "camino más corto con costo uniforme" sobre un grafo
implícito donde cada casilla es un nodo y sus hasta 8 movimientos de
caballo son las aristas: exactamente el mismo patrón que
Laberinto con BFS, cambiando el arreglo de movimientos de 4/8
direcciones ortogonales por los 8 saltos en "L" del caballo.

\textbf{¿Por qué usarlo?} Es tentador intentar resolver esto con
fórmulas o backtracking (como el Tour), pero al ser "distancia
mínima" y no "visitar todo", BFS es tanto más simple como
garantizado óptimo — la primera vez que BFS alcanza la casilla
destino, esa es la respuesta.

\textbf{Casos de uso:}

\begin{itemize}
	\item Mínimo número de movimientos de caballo entre dos casillas
	de un tablero $n \times n$ (o infinito, acotando la búsqueda).
	\item Mismo problema con casillas bloqueadas/prohibidas: se
	salta la casilla en el BFS igual que una pared en un laberinto.
	\item Generalización a "cuál es la distancia mínima con un patrón
	de movimiento arbitrario" (no solo caballo): el patrón se sigue
	resolviendo con BFS cambiando \texttt{dr/dc}.
\end{itemize}

\textbf{Complejidad:} $O(n^2)$ para un tablero $n \times n$, ya que
cada casilla se encola como máximo una vez y desde cada una se
prueban 8 movimientos constantes.

\begin{lstlisting}[style=cpp]
	struct MinMovimientosCaballo {
		int n;
		vector<vector<int>> dist;
		
		int dr[8] = {-2,-2,-1,-1, 1, 1, 2, 2};
		int dc[8] = {-1, 1,-2, 2,-2, 2,-1, 1};
		
		// CREATE: tablero n x n.
		MinMovimientosCaballo(int n_) {
			n = n_;
		}                                          // O(1)
		
		bool dentro(int r, int c) {
			return r >= 0 && r < n && c >= 0 && c < n;
		}
		
		// READ: BFS desde (sr, sc); dist[r][c] = minimo de movimientos
		// de caballo para llegar a (r, c), -1 si no aplica (no deberia
		// pasar en un tablero cerrado y conexo).
		int minMovimientos(int sr, int sc, int er, int ec) {
			dist.assign(n, vector<int>(n, -1));
			
			queue<pair<int,int>> q;
			q.push({sr, sc});
			dist[sr][sc] = 0;
			
			while (!q.empty()) {
				auto [r, c] = q.front(); q.pop();
				
				if (r == er && c == ec) return dist[r][c];
				
				for (int d = 0; d < 8; d++) {
					int nr = r + dr[d], nc = c + dc[d];
					
					if (!dentro(nr, nc)) continue;
					if (dist[nr][nc] != -1) continue;
					
					dist[nr][nc] = dist[r][c] + 1;
					q.push({nr, nc});
				}
			}
			
			return -1;   // no alcanzable (no deberia pasar sin bloqueos)
		}                                          // O(n^2)
	};
\end{lstlisting}

\textbf{Ejemplo de funcionamiento:}

\begin{lstlisting}[style=cpp]
	MinMovimientosCaballo mc(8);
	
	cout << mc.minMovimientos(0, 0, 1, 2);
	// 1  (un solo salto en L)
	
	cout << mc.minMovimientos(0, 0, 7, 7);
	// 6  (esquina a esquina en un tablero 8x8)
	
	cout << mc.minMovimientos(0, 0, 0, 1);
	// 3  (casilla adyacente, contraintuitivamente lejos para un caballo)
\end{lstlisting}

\textbf{Regla práctica:} misma distinción que Flood Fill vs.
Laberinto con BFS, aplicada al caballo: si el problema pide
\textbf{"visitar todo"} sin repetir, es backtracking (Tour de
Caballo, con Warnsdorff como poda). Si pide \textbf{"distancia
	mínima entre dos puntos"}, es BFS (Mínimo Número de Movimientos), sin
importar qué tan raro sea el patrón de movimiento — el patrón solo
define el arreglo de vecinos.

\subsection{Serpientes y Escaleras}

\textbf{Idea:} el tablero de Serpientes y Escaleras se modela como un
grafo dirigido: cada casilla $1..n$ es un nodo, y desde cada casilla
$i$ hay hasta 6 aristas (una por cada valor del dado, $i+1$ hasta
$i+6$), \textbf{redirigidas} si la casilla de llegada tiene una
serpiente o escalera. Pedir el "mínimo número de tiradas de dado para
llegar de la casilla 1 a la $n$" es, otra vez, camino más corto con
costo uniforme (cada tirada cuenta como 1 paso) — el mismo BFS de
Laberinto con BFS y Mínimo Número de Movimientos de Caballo, cambiando
qué cuenta como "vecino" de una casilla.

\textbf{¿Por qué usarlo?} La parte no trivial no es el BFS en sí
—es idéntico al de siempre— sino \textbf{modelar} serpientes y
escaleras correctamente: no son casillas nuevas, son
\textbf{atajos/retrocesos automáticos}. Si aterrizas en la casilla
de la base de una escalera o la cabeza de una serpiente, tu posición
real cambia \textbf{sin gastar otra tirada}. Eso se resuelve
precomputando, para cada casilla, a dónde "realmente" te manda
(\texttt{destinoReal[i]}), y haciendo BFS sobre ese arreglo en vez de
sobre los números crudos del dado.

\textbf{Casos de uso:}

\begin{itemize}
	\item Mínimo número de tiradas de dado para llegar de la casilla
	1 (o cualquier inicio) a la casilla $n$.
	\item Variante donde el dado no es 1–6 sino un conjunto arbitrario
	de valores posibles: solo cambia el rango del \texttt{for} al
	generar vecinos.
	\item Cualquier tablero lineal con "teletransportes" fijos
	(serpiente = te manda hacia atrás, escalera = te manda hacia
	adelante) donde se pide distancia mínima en pasos.
\end{itemize}

\textbf{Complejidad:} $O(n)$, ya que el grafo tiene $n$ nodos y a lo
más $6n$ aristas — cada casilla se encola una sola vez.

\begin{lstlisting}[style=cpp]
	struct SerpientesYEscaleras {
		int n;
		vector<int> destinoReal;   // destinoReal[i] = a donde termina
		// realmente el jugador si cae en i
		// (i mismo si no hay serpiente/escalera)
		
		// CREATE: tablero de n casillas (1-indexado). saltos[i] = j
		// significa "serpiente o escalera de i a j"; las casillas sin
		// entrada en el mapa no tienen salto.
		SerpientesYEscaleras(int n_, map<int,int>& saltos) {
			n = n_;
			destinoReal.assign(n + 1, 0);
			
			for (int i = 1; i <= n; i++)
			destinoReal[i] = saltos.count(i) ? saltos[i] : i;
		}                                          // O(n)
		
		// READ: minimo numero de tiradas de dado (1-6) para llegar de
		// origen a n. -1 si es imposible (no deberia pasar en un
		// tablero estandar).
		int minTiradas(int origen) {
			vector<int> dist(n + 1, -1);
			queue<int> q;
			
			q.push(origen);
			dist[origen] = 0;
			
			while (!q.empty()) {
				int actual = q.front(); q.pop();
				
				if (actual == n) return dist[actual];
				
				for (int dado = 1; dado <= 6; dado++) {
					int siguiente = actual + dado;
					if (siguiente > n) continue;
					
					siguiente = destinoReal[siguiente];   // aplica
					// serpiente/escalera
					
					if (dist[siguiente] != -1) continue;
					
					dist[siguiente] = dist[actual] + 1;
					q.push(siguiente);
				}
			}
			
			return -1;
		}                                          // O(n)
	};
\end{lstlisting}

\textbf{Ejemplo de funcionamiento:}

\begin{lstlisting}[style=cpp]
	// tablero de 10 casillas:
	// escalera de 2 a 8
	// serpiente  de 9 a 3
	
	map<int,int> saltos = {{2, 8}, {9, 3}};
	SerpientesYEscaleras se(10, saltos);
	
	cout << se.minTiradas(1);
	// 2
	// tirada 1: 1 -> 2 -> (escalera) -> 8
	// tirada 2: 8 -> 10
\end{lstlisting}

\textbf{Regla práctica:} cualquier tablero lineal donde te mueves por
"reglas de dado" (o cualquier conjunto de deltas fijos) y hay
teletransportes automáticos al aterrizar en ciertas casillas, se
reduce a BFS de toda la vida sobre un grafo de $n$ nodos —
\textbf{la única parte específica del problema es construir bien el
	arreglo \texttt{destinoReal}} antes de correr el BFS; el algoritmo de
búsqueda no cambia respecto a los casos anteriores.

% =========================================================
% CONECTIVIDAD
% =========================================================
\section{Conectividad}
\subsection{Componentes Conexas}

\textbf{Idea:} en un grafo \textbf{no dirigido}, una componente
conexa es un grupo maximal de nodos donde cualquiera puede llegar a
cualquiera. Encontrarlas todas es simplemente: recorrer los nodos en
orden, y cada vez que encuentres uno \textbf{no visitado}, lanzar un
DFS/BFS desde ahí para marcar todo su grupo con un mismo
identificador — ese DFS/BFS nunca se sale del grupo porque, por
definición, no hay aristas hacia afuera de una componente.

\textbf{¿Por qué usarlo?} Es el primer paso obligatorio de casi
cualquier problema de grafos no dirigidos: antes de responder
"¿$u$ y $v$ están conectados?" o de aplicar un algoritmo que asume
un grafo conexo, necesitas saber cuántas componentes hay y a cuál
pertenece cada nodo. La alternativa es \textbf{Union-Find}: si solo
necesitas responder "¿misma componente?" con aristas que van
llegando una a una (sin necesariamente construir listas de
adyacencia), Union-Find es más directo; si ya tienes el grafo
completo y quieres además reconstruir/recorrer cada componente,
DFS/BFS es más natural.

\textbf{Casos de uso:}

\begin{itemize}
	\item Contar cuántas componentes conexas tiene un grafo (o una
	grilla vista como grafo — mismo patrón que Flood Fill, pero
	sobre listas de adyacencia en vez de vecinos de grilla).
	\item Etiquetar cada nodo con el id de su componente, para
	responder "¿misma componente?" en $O(1)$ tras el precómputo.
	\item Encontrar el tamaño de la componente más grande.
	\item Verificar si el grafo completo es conexo (una sola
	componente).
\end{itemize}

\textbf{Complejidad:} $O(V + E)$, ya que es un DFS/BFS normal que
simplemente se relanza desde cada nodo no visitado — en total, cada
nodo y cada arista se procesan una sola vez sin importar cuántas
componentes haya.

\begin{lstlisting}[style=cpp]
	struct ComponentesConexas {
		vector<vector<int>> adj;
		vector<int> comp;      // comp[u] = id de la componente de u, -1
		// si no se ha calculado
		int numComponentes;
		
		// CREATE: guarda el grafo y calcula todas las componentes de
		// una vez, recorriendo cada nodo no visitado.
		ComponentesConexas(vector<vector<int>>& adj_, int n) {
			adj = adj_;
			comp.assign(n, -1);
			numComponentes = 0;
			
			for (int u = 0; u < n; u++) {
				if (comp[u] == -1) {
					dfs(u, numComponentes);
					numComponentes++;
				}
			}
		}                                          // O(V + E)
		
		void dfs(int u, int id) {
			comp[u] = id;
			for (int v : adj[u])
			if (comp[v] == -1)
			dfs(v, id);
		}
		
		// READ: u y v estan en la misma componente?
		bool mismaComponente(int u, int v) {
			return comp[u] == comp[v];
		}                                          // O(1)
	};
\end{lstlisting}

\textbf{Ejemplo de funcionamiento:}

\begin{lstlisting}[style=cpp]
	// grafo:  0 - 1    2 - 3    4
	//         (dos componentes de tamano 2, una de tamano 1)
	
	vector<vector<int>> adj(5);
	adj[0] = {1}; adj[1] = {0};
	adj[2] = {3}; adj[3] = {2};
	adj[4] = {};
	
	ComponentesConexas cc(adj, 5);
	
	cout << cc.numComponentes;
	// 3
	
	cout << cc.mismaComponente(0, 1);
	// true
	
	cout << cc.mismaComponente(1, 2);
	// false
\end{lstlisting}


\subsection{Detección de Ciclos}

\textbf{Idea:} detectar si un grafo tiene un ciclo se hace con DFS,
pero la regla para reconocer un ciclo \textbf{cambia según el grafo
	sea dirigido o no}: en un grafo no dirigido, cualquier arista hacia
un nodo ya visitado que \textbf{no sea tu padre inmediato} es un
ciclo; en un grafo dirigido, una arista hacia un nodo ya visitado
\textbf{solo} es ciclo si ese nodo sigue en la pila de recursión
actual (todavía no terminaste de procesarlo) — si ya terminaste con
él, es una arista válida hacia otra rama, no un ciclo.

\textbf{¿Por qué usarlo?} Confundir estas dos reglas es el error más
común al detectar ciclos: aplicar la regla de "no dirigido" (solo
evitar al padre) a un grafo dirigido da falsos positivos, porque en
dirigido puedes visitar el mismo nodo desde dos ramas distintas sin
que eso sea un ciclo.

\textbf{Casos de uso:}

\begin{itemize}
	\item Verificar si un grafo dirigido es un DAG (sin ciclos) antes
	de aplicar orden topológico, DP sobre DAG, etc.
	\item Detectar dependencias circulares (tareas que dependen unas
	de otras en círculo, imposibles de ordenar).
	\item En grafos no dirigidos: verificar si un grafo es un árbol
	($V-1$ aristas, conexo, sin ciclos) o encontrar el primer ciclo
	para removerlo.
\end{itemize}

\textbf{Complejidad:} $O(V + E)$ en ambos casos — un DFS normal con
una condición extra por arista, sin repetir trabajo.

\subsubsection{En grafos no dirigidos (DFS con padre)}

\begin{lstlisting}[style=cpp]
	struct CicloNoDirigido {
		vector<vector<int>> adj;
		vector<bool> visitado;
		
		// CREATE: guarda el grafo.
		CicloNoDirigido(vector<vector<int>>& adj_) {
			adj = adj_;
			visitado.assign(adj.size(), false);
		}                                          // O(1)
		
		// READ: existe algun ciclo en todo el grafo (revisa todas las
		// componentes).
		bool tieneCiclo() {
			for (int u = 0; u < (int)adj.size(); u++)
			if (!visitado[u] && dfs(u, -1))
			return true;
			return false;
		}                                          // O(V + E)
		
		bool dfs(int u, int padre) {
			visitado[u] = true;
			
			for (int v : adj[u]) {
				if (!visitado[v]) {
					if (dfs(v, u)) return true;
				} else if (v != padre) {
					return true;   // arista hacia un visitado que no es
					// el padre -> ciclo
				}
			}
			
			return false;
		}
	};
\end{lstlisting}

\textbf{Ejemplo de funcionamiento:}

\begin{lstlisting}[style=cpp]
	// grafo:  0 - 1 - 2 - 0   (triangulo, tiene ciclo)
	//         3 - 4           (arbol, sin ciclo)
	
	vector<vector<int>> adj(5);
	adj[0] = {1, 2}; adj[1] = {0, 2}; adj[2] = {0, 1};
	adj[3] = {4};    adj[4] = {3};
	
	CicloNoDirigido cnd(adj);
	
	cout << cnd.tieneCiclo();
	// true  (por el triangulo 0-1-2)
\end{lstlisting}

\subsubsection{En grafos dirigidos (DFS con estados / colores)}

\textbf{Idea:} cada nodo tiene 3 estados durante el DFS: \texttt{0 =
	no visitado}, \texttt{1 = en la pila de recursión actual (gris)},
\texttt{2 = totalmente procesado (negro)}. Una arista hacia un nodo
\textbf{gris} es un back edge real (ciclo); una arista hacia un nodo
\textbf{negro} es válida (cross/forward edge, no ciclo).

\begin{lstlisting}[style=cpp]
	struct CicloDirigido {
		vector<vector<int>> adj;
		vector<int> estado;   // 0 = blanco, 1 = gris, 2 = negro
		
		// CREATE: guarda el grafo dirigido.
		CicloDirigido(vector<vector<int>>& adj_) {
			adj = adj_;
			estado.assign(adj.size(), 0);
		}                                          // O(1)
		
		// READ: existe algun ciclo en todo el grafo.
		bool tieneCiclo() {
			for (int u = 0; u < (int)adj.size(); u++)
			if (estado[u] == 0 && dfs(u))
			return true;
			return false;
		}                                          // O(V + E)
		
		bool dfs(int u) {
			estado[u] = 1;   // entra a la pila de recursion (gris)
			
			for (int v : adj[u]) {
				if (estado[v] == 1) return true;        // back edge -> ciclo
				if (estado[v] == 0 && dfs(v)) return true;
			}
			
			estado[u] = 2;   // termino con u (negro)
			return false;
		}
	};
\end{lstlisting}

\textbf{Ejemplo de funcionamiento:}

\begin{lstlisting}[style=cpp]
	// grafo dirigido:  0 -> 1 -> 2 -> 0   (ciclo)
	//                  3 -> 1             (sin ciclo: 3->1 es valida,
	//                                      1 ya esta siendo procesado
	//                                      solo si viene de la MISMA
	//                                      rama de recursion)
	
	vector<vector<int>> adj(4);
	adj[0] = {1};
	adj[1] = {2};
	adj[2] = {0};
	adj[3] = {1};
	
	CicloDirigido cd(adj);
	
	cout << cd.tieneCiclo();
	// true  (0 -> 1 -> 2 -> 0)
\end{lstlisting}

\textbf{Regla práctica:} para saber cuál versión aplicar, primero
pregúntate si el grafo es dirigido. Si no lo es, basta con recordar
"de dónde vine" (el padre) y listo. Si es dirigido, "de dónde vine"
no alcanza — necesitas saber si el nodo destino sigue \textbf{activo
	en la recursión actual} (gris) o ya quedó cerrado (negro), porque en
dirigido sí es válido llegar dos veces al mismo nodo desde ramas
distintas sin que eso sea un ciclo.



\subsection{Puentes}

\textbf{Idea:} en un grafo no dirigido, un \textbf{puente} es una
arista que, al quitarla, \textbf{desconecta} el grafo (aumenta el
número de componentes conexas). Se encuentran con un solo DFS,
llevando dos valores por nodo: \texttt{tin[u]} (el momento en que el
DFS \textbf{descubre} a $u$, un contador que se incrementa en cada
visita) y \texttt{low[u]} (el \texttt{tin} \textbf{más bajo}
alcanzable desde $u$ usando cualquier camino del árbol DFS más, como
mucho, \textbf{una} arista de "retroceso" hacia un ancestro). La
arista $(u,v)$ del árbol DFS es puente si y solo si
$\texttt{low}[v] > \texttt{tin}[u]$: eso significa que, desde el
subárbol de $v$, \textbf{no hay ningún atajo} de vuelta hacia $u$ o
más arriba —quitar esa arista realmente aísla el subárbol de $v$.

\textbf{¿Por qué usarlo?} Es la aplicación clásica del algoritmo de
Tarjan de \texttt{low-link}: en vez de, para cada arista, quitarla y
volver a correr un DFS/BFS completo para ver si el grafo se
desconecta ($O(E \cdot (V+E))$ en total), un solo DFS con
\texttt{low-link} responde para \textbf{todas} las aristas a la vez
en tiempo lineal.

\textbf{Casos de uso:}

\begin{itemize}
	\item Encontrar todas las conexiones "críticas" de una red (si
	se cae ese enlace, la red se parte en dos).
	\item Como paso previo para construir el \textbf{árbol de
		componentes 2-arista-conexas} (contraer cada componente
	2-arista-conexa a un nodo y dejar solo los puentes como aristas).
	\item Verificar si el grafo sigue siendo conexo tras remover
	cualquier arista dada (equivalente a "¿esa arista es puente?").
\end{itemize}

\textbf{Complejidad:} $O(V + E)$, un solo DFS con trabajo constante
extra por arista.

\begin{lstlisting}[style=cpp]
	struct Puentes {
		vector<vector<pair<int,int>>> adj;   // adj[u] = {(v, idArista)}
		vector<int> tin, low;
		vector<bool> visitado;
		vector<pair<int,int>> puentes;       // aristas (u,v) que son puente
		int timer;
		
		// CREATE: guarda el grafo (con id de arista para no reusar la
		// misma arista en sentido contrario) y corre el DFS de Tarjan
		// desde cada componente.
		Puentes(vector<vector<pair<int,int>>>& adj_, int n) {
			adj = adj_;
			tin.assign(n, -1);
			low.assign(n, -1);
			visitado.assign(n, false);
			timer = 0;
			
			for (int u = 0; u < n; u++)
			if (!visitado[u])
			dfs(u, -1);
		}                                          // O(V + E)
		
		void dfs(int u, int idAristaPadre) {
			visitado[u] = true;
			tin[u] = low[u] = timer++;
			
			for (auto& [v, idArista] : adj[u]) {
				if (idArista == idAristaPadre) continue;   // no te
				// regreses
				// por la
				// misma
				// arista
				
				if (visitado[v]) {
					low[u] = min(low[u], tin[v]);   // back edge
				} else {
					dfs(v, idArista);
					low[u] = min(low[u], low[v]);
					
					if (low[v] > tin[u])
					puentes.push_back({u, v});
				}
			}
		}
	};
\end{lstlisting}

\textbf{Ejemplo de funcionamiento:}

\begin{lstlisting}[style=cpp]
	// grafo:   0 - 1 - 2 - 0     (triangulo)
	//                  \
	//                   3 - 4    (colgando del triangulo)
	
	vector<vector<pair<int,int>>> adj(5);
	auto addEdge = [&](int u, int v, int id) {
		adj[u].push_back({v, id});
		adj[v].push_back({u, id});
	};
	addEdge(0, 1, 0);
	addEdge(1, 2, 1);
	addEdge(2, 0, 2);
	addEdge(2, 3, 3);
	addEdge(3, 4, 4);
	
	Puentes p(adj, 5);
	
	for (auto& [u, v] : p.puentes)
	cout << u << "-" << v << " ";
	// 2-3  3-4
	// (las aristas del triangulo NO son puente: hay un ciclo,
	//  siempre queda un camino alternativo)
\end{lstlisting}


\subsection{Puntos de Articulación}

\textbf{Idea:} un \textbf{punto de articulación} (o \textit{cut
	vertex}) es un nodo que, al quitarlo (junto con sus aristas),
\textbf{aumenta} el número de componentes conexas. Usa el mismo
\texttt{tin}/\texttt{low} que Puentes, pero la condición es distinta
y tiene \textbf{dos casos}: (1) si $u$ es la \textbf{raíz} del árbol
DFS, es articulación si tiene \textbf{2 o más hijos} en el árbol DFS
(cada hijo es un subárbol que solo se conecta al resto \textit{a
	través de} $u$); (2) si $u$ \textbf{no} es la raíz, es articulación
si tiene algún hijo $v$ con $\texttt{low}[v] \ge \texttt{tin}[u]$
(el subárbol de $v$ no tiene ningún atajo que lo saque \textbf{por
	encima} de $u$, así que sin $u$ ese subárbol queda aislado).

\textbf{¿Por qué usarlo?} Es el "primo" de Puentes: mismo DFS,
mismos \texttt{tin}/\texttt{low}, pero identificando nodos críticos
en vez de aristas críticas. La diferencia clave con Puentes está en
el operador de comparación ($\ge$ en vez de $>$) y en el caso especial
de la raíz — con solo $>$ (la condición de puente) se pierde el caso
donde $v$ puede volver exactamente a $u$ (eso no aísla la arista,
pero sí deja a $u$ como único punto de paso).

\textbf{Casos de uso:}

\begin{itemize}
	\item Encontrar los nodos "críticos" de una red (si ese nodo
	falla, la red se parte).
	\item Como paso previo para construir componentes
	\textbf{biconexas} (bloques del grafo separados por puntos de
	articulación).
	\item Distinguir de Puentes: un grafo puede tener puntos de
	articulación sin tener ningún puente (ej. dos triángulos que
	comparten un solo vértice) — son propiedades independientes.
\end{itemize}

\textbf{Complejidad:} $O(V + E)$, mismo costo que Puentes: un DFS con
trabajo constante extra por nodo/arista.

\begin{lstlisting}[style=cpp]
	struct PuntosArticulacion {
		vector<vector<int>> adj;
		vector<int> tin, low;
		vector<bool> visitado, esArticulacion;
		int timer;
		
		// CREATE: guarda el grafo y corre el DFS de Tarjan desde cada
		// componente.
		PuntosArticulacion(vector<vector<int>>& adj_, int n) {
			adj = adj_;
			tin.assign(n, -1);
			low.assign(n, -1);
			visitado.assign(n, false);
			esArticulacion.assign(n, false);
			timer = 0;
			
			for (int u = 0; u < n; u++)
			if (!visitado[u])
			dfs(u, -1);
		}                                          // O(V + E)
		
		void dfs(int u, int padre) {
			visitado[u] = true;
			tin[u] = low[u] = timer++;
			int hijos = 0;
			
			for (int v : adj[u]) {
				if (v == padre) continue;   // ignora SOLO una vez la
				// arista de vuelta al padre
				// (si hay aristas multiples
				// u-padre, usar id de arista
				// como en Puentes)
				
				if (visitado[v]) {
					low[u] = min(low[u], tin[v]);   // back edge
				} else {
					dfs(v, u);
					low[u] = min(low[u], low[v]);
					hijos++;
					
					if (padre != -1 && low[v] >= tin[u])
					esArticulacion[u] = true;
				}
			}
			
			if (padre == -1 && hijos >= 2)
			esArticulacion[u] = true;   // caso especial: raiz con
			// 2+ hijos en el arbol DFS
		}
	};
\end{lstlisting}

\textbf{Ejemplo de funcionamiento:}

\begin{lstlisting}[style=cpp]
	// grafo:   0 - 1 - 2 - 0     (triangulo)
	//                  \
	//                   3 - 4    (colgando del triangulo)
	
	vector<vector<int>> adj(5);
	adj[0] = {1, 2}; adj[1] = {0, 2}; adj[2] = {0, 1, 3};
	adj[3] = {2, 4}; adj[4] = {3};
	
	PuntosArticulacion pa(adj, 5);
	
	for (int u = 0; u < 5; u++)
	if (pa.esArticulacion[u])
	cout << u << " ";
	// 2  3
	// (2 conecta el triangulo con la cola 3-4; 3 conecta 2 con 4)
\end{lstlisting}

\textbf{Regla práctica:} Puentes y Puntos de Articulación comparten
exactamente el esqueleto del DFS (\texttt{tin}/\texttt{low}, back
edges hacia ancestros) — la diferencia está solo en \textbf{qué}
marcas: en Puentes marcas la \textbf{arista} $(u,v)$ cuando
$\texttt{low}[v] > \texttt{tin}[u]$ (estricto); en Puntos de
Articulación marcas el \textbf{nodo} $u$ cuando
$\texttt{low}[v] \ge \texttt{tin}[u]$ (no estricto, más el caso
especial de la raíz). Si programas uno, el otro es casi copiar y
pegar cambiando esos dos detalles.

\subsection{Componentes Biconectados}

\textbf{Idea:} una componente biconexa (o \textit{2-vertex-connected
	component}, "bloque") es un conjunto \textbf{maximal de aristas}
donde no hay ningún punto de articulación que las separe —
equivalentemente, cualquier par de aristas del bloque está en algún
ciclo común. Se calculan con el mismo DFS de \texttt{tin}/\texttt{low}
de Puentes y Puntos de Articulación, agregando una \textbf{pila de
	aristas}: apilas cada arista al recorrerla, y cuando detectas la
condición de punto de articulación ($\texttt{low}[v] \ge
\texttt{tin}[u]$), \textbf{desapilas} todo hasta la arista $(u,v)$
inclusive — ese grupo desapilado es una componente biconexa completa.

\textbf{¿Por qué usarlo?} Puntos de Articulación te dice \textbf{qué
	nodos} son críticos, pero no cómo se agrupan las aristas alrededor de
ellos. Componentes Biconectadas completa el panorama: parte el grafo
en bloques tales que cada punto de articulación es exactamente el
nodo compartido entre dos o más bloques — es la base para construir
el \textbf{block-cut tree} (árbol de bloques y puntos de
articulación), útil cuando el problema pregunta por relaciones entre
bloques en vez de solo "¿es articulación sí/no?".

\textbf{Casos de uso:}

\begin{itemize}
	\item Agrupar aristas en bloques maximales sin puntos de
	articulación internos.
	\item Construir el block-cut tree para responder queries tipo
	"¿cuántos bloques hay que atravesar para ir de $u$ a $v$?".
	\item Verificar si el grafo completo es biconexo (una sola
	componente biconexa que cubre todas las aristas, y ningún punto
	de articulación).
\end{itemize}

\textbf{Complejidad:} $O(V + E)$ — mismo DFS de \texttt{tin}/
\texttt{low} de siempre, más trabajo amortizado constante por arista
al apilarla/desapilarla (cada arista se apila y desapila una sola
vez en total).

\begin{lstlisting}[style=cpp]
	struct ComponentesBiconexas {
		vector<vector<pair<int,int>>> adj;   // adj[u] = {(v, idArista)}
		vector<int> tin, low;
		vector<bool> visitado;
		stack<pair<int,int>> pilaAristas;    // aristas (u,v) pendientes
		vector<vector<pair<int,int>>> bloques;
		int timer;
		
		// CREATE: guarda el grafo y corre el DFS desde cada componente.
		ComponentesBiconexas(vector<vector<pair<int,int>>>& adj_, int n) {
			adj = adj_;
			tin.assign(n, -1);
			low.assign(n, -1);
			visitado.assign(n, false);
			timer = 0;
			
			for (int u = 0; u < n; u++)
			if (!visitado[u])
			dfs(u, -1);
		}                                          // O(V + E)
		
		void dfs(int u, int idAristaPadre) {
			visitado[u] = true;
			tin[u] = low[u] = timer++;
			
			for (auto& [v, idArista] : adj[u]) {
				if (idArista == idAristaPadre) continue;
				
				if (visitado[v] && tin[v] < tin[u]) {
					// back edge: entra al bloque actual
					pilaAristas.push({u, v});
					low[u] = min(low[u], tin[v]);
				} else if (!visitado[v]) {
					pilaAristas.push({u, v});
					dfs(v, idArista);
					low[u] = min(low[u], low[v]);
					
					if (low[v] >= tin[u]) {
						// u es articulacion (o raiz): cierra un bloque,
						// desapilando hasta (u, v) inclusive
						vector<pair<int,int>> bloque;
						while (pilaAristas.top() != make_pair(u, v)) {
							bloque.push_back(pilaAristas.top());
							pilaAristas.pop();
						}
						bloque.push_back(pilaAristas.top());
						pilaAristas.pop();
						
						bloques.push_back(bloque);
					}
				}
			}
		}
	};
\end{lstlisting}

\textbf{Ejemplo de funcionamiento:}

\begin{lstlisting}[style=cpp]
	// grafo:   0 - 1 - 2 - 0     (triangulo, bloque 1)
	//                  \
	//                   3 - 4    (puente, bloque 2 y 3 por separado)
	
	vector<vector<pair<int,int>>> adj(5);
	auto addEdge = [&](int u, int v, int id) {
		adj[u].push_back({v, id});
		adj[v].push_back({u, id});
	};
	addEdge(0, 1, 0);
	addEdge(1, 2, 1);
	addEdge(2, 0, 2);
	addEdge(2, 3, 3);
	addEdge(3, 4, 4);
	
	ComponentesBiconexas cb(adj, 5);
	
	cout << cb.bloques.size();
	// 3
	// bloque 1: {(0,1),(1,2),(2,0)}  (el triangulo)
	// bloque 2: {(2,3)}              (el puente 2-3)
	// bloque 3: {(3,4)}              (el puente 3-4)
\end{lstlisting}

\textbf{Nota:} cada punto de articulación aparece en \textbf{más de
	un} bloque (es el nodo "bisagra" entre ellos); cada nodo que no es
articulación aparece en exactamente un bloque. Eso es justo lo que
permite construir el block-cut tree: un nodo por cada bloque, un nodo
por cada punto de articulación, y una arista entre un bloque y cada
punto de articulación que contiene.


\subsection{Tarjan}

\textbf{Idea:} el algoritmo de Tarjan encuentra las \textbf{componentes
	fuertemente conexas} (SCC) de un grafo \textbf{dirigido}: grupos
maximales de nodos donde \textbf{cada} nodo puede llegar a
\textbf{cualquier otro} del mismo grupo (yendo por aristas
dirigidas). Es el mismo esqueleto de \texttt{tin}/\texttt{low} que
Puentes/Articulación/Biconexas, adaptado a dirigido: mantienes una
\textbf{pila de nodos} (no de aristas) y un flag \texttt{enPila[u]};
un nodo $u$ es \textbf{raíz de su SCC} cuando $\texttt{low}[u] ==
\texttt{tin}[u]$ — en ese momento desapilas todo hasta $u$ inclusive,
y eso es una SCC completa.

\textbf{¿Por qué usarlo?} Es la forma estándar de \textbf{condensar}
un grafo dirigido general en un DAG: contrae cada SCC a un solo nodo,
y las aristas entre SCCs distintas forman un DAG sobre el cual sí
puedes aplicar orden topológico, DP sobre DAG, o Conteo de Caminos
—todo lo que requiere ausencia de ciclos. Es la alternativa a
\textbf{Kosaraju} (dos pasadas de DFS + grafo transpuesto): Tarjan
logra lo mismo en \textbf{una sola pasada}, a cambio de llevar la
pila de nodos y el flag \texttt{enPila}.

\textbf{Casos de uso:}

\begin{itemize}
	\item Condensar un grafo dirigido con ciclos en un DAG, como paso
	previo a Conteo de Caminos, orden topológico, o DP sobre DAG.
	\item Detectar si un grafo dirigido es fuertemente conexo (una
	sola SCC que cubre todos los nodos).
	\item Encontrar el tamaño de la SCC más grande (por ejemplo,
	"el mayor grupo de cuentas que se siguen mutuamente").
\end{itemize}

\textbf{Complejidad:} $O(V + E)$, un solo DFS con trabajo constante
extra por nodo/arista, igual que las estructuras anteriores de esta
sección.

\begin{lstlisting}[style=cpp]
	struct Tarjan {
		vector<vector<int>> adj;
		vector<int> tin, low, comp;
		vector<bool> enPila;
		stack<int> pila;
		int timer, numSCC;
		
		// CREATE: guarda el grafo dirigido y corre Tarjan desde cada
		// nodo no visitado.
		Tarjan(vector<vector<int>>& adj_, int n) {
			adj = adj_;
			tin.assign(n, -1);
			low.assign(n, -1);
			comp.assign(n, -1);      // comp[u] = id de la SCC de u
			enPila.assign(n, false);
			timer = 0;
			numSCC = 0;
			
			for (int u = 0; u < n; u++)
			if (tin[u] == -1)
			dfs(u);
		}                                          // O(V + E)
		
		void dfs(int u) {
			tin[u] = low[u] = timer++;
			pila.push(u);
			enPila[u] = true;
			
			for (int v : adj[u]) {
				if (tin[v] == -1) {
					dfs(v);
					low[u] = min(low[u], low[v]);
				} else if (enPila[v]) {
					// v esta en la pila -> back edge dentro de la
					// misma SCC en formacion
					low[u] = min(low[u], tin[v]);
				}
				// si v ya se visito y NO esta en pila, es una SCC ya
				// cerrada -> se ignora (no afecta low[u])
			}
			
			if (low[u] == tin[u]) {
				// u es raiz de su SCC: desapila todo el grupo
				while (true) {
					int w = pila.top(); pila.pop();
					enPila[w] = false;
					comp[w] = numSCC;
					if (w == u) break;
				}
				numSCC++;
			}
		}
	};
\end{lstlisting}

\textbf{Ejemplo de funcionamiento:}

\begin{lstlisting}[style=cpp]
	// grafo dirigido:  0 -> 1 -> 2 -> 0   (ciclo, una SCC)
	//                  2 -> 3
	//                  3 -> 4 -> 3        (ciclo, otra SCC)
	
	vector<vector<int>> adj(5);
	adj[0] = {1};
	adj[1] = {2};
	adj[2] = {0, 3};
	adj[3] = {4};
	adj[4] = {3};
	
	Tarjan tj(adj, 5);
	
	cout << tj.numSCC;
	// 2
	
	cout << tj.comp[0] << " " << tj.comp[1] << " " << tj.comp[2];
	// mismo id para 0, 1, 2  (SCC del ciclo 0-1-2)
	
	cout << tj.comp[3] << " " << tj.comp[4];
	// mismo id para 3, 4     (SCC del ciclo 3-4)
\end{lstlisting}

\textbf{Regla práctica:} de las cuatro estructuras de esta sección,
todas comparten el DFS con \texttt{tin}/\texttt{low}, pero difieren
en \textbf{qué apilas y cuándo cierras el grupo}: Puentes no apila
nada (marca aristas  directamente); Puntos de Articulación tampoco
apila (marca nodos); Componentes Biconexas apila \textbf{aristas} y
cierra con $\texttt{low}[v] \ge \texttt{tin}[u]$; Tarjan apila
\textbf{nodos} y cierra con $\texttt{low}[u] == \texttt{tin}[u]$
—y solo Tarjan aplica a grafos \textbf{dirigidos}, ya que las otras
tres asumen aristas no dirigidas (por eso usan el truco de "no
regresar por la misma arista al padre").

% =========================================================
% TOPOLOGICAL SORT
% =========================================================
\section{Topological Sort}
\subsection{Kahn's Algorithm}

\textbf{Idea:} Kahn's Algorithm calcula un orden topológico de un DAG
\textbf{sin recursión}, usando la idea de "quitar capas": empiezas
encolando todos los nodos con \textbf{grado de entrada 0} (sin
prerrequisitos), y cada vez que "procesas" un nodo, le restas 1 al
grado de entrada de sus vecinos —si alguno llega a 0, significa que
ya se cumplieron todos sus prerrequisitos, así que se encola. Es
exactamente el mismo patrón de BFS por capas de Laberinto con BFS,
pero la "distancia" que se propaga es el conteo de prerrequisitos
pendientes en vez de pasos.

\textbf{¿Por qué usarlo?} Además de dar el orden, Kahn's Algorithm te
da \textbf{gratis} la detección de ciclos: si al final quedan nodos
sin procesar (el orden resultante tiene menos de $n$ nodos), el grafo
\textbf{no} es un DAG. Es la alternativa iterativa a DFS + Orden
Topológico — preferible cuando quieres evitar recursión profunda (no
hay riesgo de stack overflow en grafos grandes) o cuando ya necesitas
el chequeo de ciclo como parte del mismo cómputo.

\textbf{Casos de uso:}

\begin{itemize}
	\item Ordenar tareas con dependencias (build systems, prerrequisitos
	de materias) respetando que cada tarea aparezca después de sus
	dependencias.
	\item Verificar si un conjunto de dependencias es realizable (sin
	ciclos) — si Kahn no logra procesar todos los nodos, hay un ciclo.
	\item Como precómputo para DP sobre DAG (procesar los nodos en
	este orden garantiza que, al llegar a $u$, ya se resolvieron
	todos sus predecesores).
\end{itemize}

\textbf{Complejidad:} $O(V + E)$ — cada nodo se encola/desencola una
vez, y cada arista se examina una vez al decrementar grados de
entrada.

\begin{lstlisting}[style=cpp]
	struct KahnTopoSort {
		vector<vector<int>> adj;
		vector<int> gradoEntrada;
		int n;
		
		// CREATE: guarda el grafo y calcula el grado de entrada de
		// cada nodo.
		KahnTopoSort(vector<vector<int>>& adj_, int n_) {
			adj = adj_;
			n = n_;
			gradoEntrada.assign(n, 0);
			
			for (int u = 0; u < n; u++)
			for (int v : adj[u])
			gradoEntrada[v]++;
		}                                          // O(V + E)
		
		// READ: devuelve el orden topologico, o un vector vacio si el
		// grafo tiene un ciclo (no es DAG).
		vector<int> ordenar() {
			vector<int> grado = gradoEntrada;   // copia para no mutar
			// el original
			queue<int> q;
			
			for (int u = 0; u < n; u++)
			if (grado[u] == 0)
			q.push(u);
			
			vector<int> orden;
			while (!q.empty()) {
				int u = q.front(); q.pop();
				orden.push_back(u);
				
				for (int v : adj[u]) {
					grado[v]--;
					if (grado[v] == 0)
					q.push(v);
				}
			}
			
			if ((int)orden.size() != n) return {};   // habia un ciclo
			return orden;
		}                                          // O(V + E)
	};
\end{lstlisting}

\textbf{Ejemplo de funcionamiento:}

\begin{lstlisting}[style=cpp]
	// DAG:  0 -> 1 -> 3
	//       0 -> 2 -> 3
	
	vector<vector<int>> adj(4);
	adj[0] = {1, 2};
	adj[1] = {3};
	adj[2] = {3};
	adj[3] = {};
	
	KahnTopoSort kts(adj, 4);
	auto orden = kts.ordenar();
	
	for (int u : orden) cout << u << " ";
	// 0 1 2 3   (o 0 2 1 3, ambos validos)
\end{lstlisting}


\subsection{DFS + Orden Topológico}

\textbf{Idea:} la versión recursiva del orden topológico usa una
propiedad simple del \textbf{post-orden} de un DFS: cuando terminas
de procesar completamente un nodo $u$ (ya visitaste todo lo
alcanzable desde él), \textbf{todos} sus sucesores ya están en la
lista de resultado. Empujando cada nodo a una pila justo al
\textbf{salir} de su llamada recursiva, y luego leyendo la pila de
arriba hacia abajo, obtienes el orden topológico completo.

\textbf{¿Por qué usarlo?} Es más directo de escribir que Kahn cuando
ya tienes DFS a mano y \textbf{sabes} de antemano que el grafo es un
DAG (si no estás seguro, hay que combinarlo con Detección de Ciclos
en grafos dirigidos, o simplemente usar Kahn, que detecta el ciclo
como parte del mismo proceso). Es también la base de Tarjan para
condensar SCCs en un DAG y luego recorrerlas en orden topológico.

\textbf{Casos de uso:}

\begin{itemize}
	\item Mismo uso que Kahn (ordenar tareas con dependencias, DP
	sobre DAG), cuando el grafo ya se sabe acíclico de antemano.
	\item Recorrer las SCCs de un grafo condensado (via Tarjan) en
	orden topológico, por ejemplo para propagar valores de SCCs
	"fuente" hacia SCCs "destino".
	\item Cuando además del orden necesitas otra información que
	ya calculas naturalmente en post-orden de un DFS (tamaños de
	subárbol, etc.) y quieres aprovechar el mismo recorrido.
\end{itemize}

\textbf{Complejidad:} $O(V + E)$, un DFS normal con una operación
constante extra (push a la pila) al salir de cada nodo.

\begin{lstlisting}[style=cpp]
	struct DFSTopoSort {
		vector<vector<int>> adj;
		vector<bool> visitado;
		stack<int> pila;
		int n;
		
		// CREATE: guarda el grafo (debe ser DAG).
		DFSTopoSort(vector<vector<int>>& adj_, int n_) {
			adj = adj_;
			n = n_;
			visitado.assign(n, false);
		}                                          // O(1)
		
		// READ: devuelve el orden topologico completo del grafo.
		vector<int> ordenar() {
			for (int u = 0; u < n; u++)
			if (!visitado[u])
			dfs(u);
			
			vector<int> orden;
			while (!pila.empty()) {
				orden.push_back(pila.top());
				pila.pop();
			}
			return orden;
		}                                          // O(V + E)
		
		void dfs(int u) {
			visitado[u] = true;
			
			for (int v : adj[u])
			if (!visitado[v])
			dfs(v);
			
			pila.push(u);   // post-orden: u entra a la pila solo
			// despues de procesar todos sus sucesores
		}
	};
\end{lstlisting}

\textbf{Ejemplo de funcionamiento:}

\begin{lstlisting}[style=cpp]
	// DAG:  0 -> 1 -> 3
	//       0 -> 2 -> 3
	
	vector<vector<int>> adj(4);
	adj[0] = {1, 2};
	adj[1] = {3};
	adj[2] = {3};
	adj[3] = {};
	
	DFSTopoSort dts(adj, 4);
	auto orden = dts.ordenar();
	
	for (int u : orden) cout << u << " ";
	// 0 1 2 3   (o 0 2 1 3, ambos validos)
\end{lstlisting}

\textbf{Regla práctica:} Kahn y DFS + Orden Topológico dan el mismo
tipo de resultado (un orden válido, no necesariamente único) con el
mismo costo $O(V+E)$ — la diferencia es de implementación, no de
poder: Kahn es iterativo y detecta ciclos como efecto secundario del
proceso (si no logra vaciar la cola, hay ciclo); DFS es recursivo y
necesita la comprobación de ciclo por separado si el grafo no se sabe
acíclico de antemano. En la práctica, usa el que ya tengas más a
mano según si el problema ya trae DFS con \texttt{tin}/\texttt{low}
(Tarjan) o no.

\subsection{Detección de Ciclos con Topological Sort}

\textbf{Idea:} un DAG (grafo dirigido acíclico) \textbf{siempre}
admite al menos un orden topológico completo que incluye a todos sus
$n$ nodos. Si un grafo dirigido \textbf{tiene} un ciclo, ningún nodo
del ciclo puede procesarse primero (cada uno depende de que otro del
mismo ciclo se procese antes), así que Kahn's Algorithm se
\textbf{atasca}: la cola se vacía sin haber procesado esos nodos. Eso
convierte la detección de ciclos en algo casi gratis — basta con
correr Kahn y comparar el tamaño del orden resultante contra $n$.

\textbf{¿Por qué usarlo?} Es la alternativa \textbf{iterativa} al DFS
con estados/colores (blanco/gris/negro) de Detección de Ciclos: en
vez de rastrear explícitamente qué nodos están "en la pila de
recursión actual", Kahn detecta el ciclo \textbf{implícitamente} —
cualquier nodo con grado de entrada $>0$ al final del proceso forma
parte de (o depende de) un ciclo. Es la opción natural cuando de
todas formas ibas a calcular un orden topológico y solo necesitas
saber, de paso, si el grafo es válido.

\textbf{Casos de uso:}

\begin{itemize}
	\item Verificar si un conjunto de dependencias es realizable
	(build systems, prerrequisitos de materias, tareas con
	"depende de") antes de intentar ordenarlas.
	\item Identificar \textbf{cuáles} nodos específicos quedan
	atrapados en un ciclo (los que nunca llegan a grado de entrada 0),
	útil para reportarle al usuario exactamente qué dependencias
	circulares hay que romper.
	\item Como chequeo previo obligatorio antes de aplicar DP sobre
	DAG o Conteo de Caminos, que asumen ausencia de ciclos.
\end{itemize}

\textbf{Complejidad:} $O(V + E)$ — es literalmente Kahn's Algorithm
sin ningún costo extra; el chequeo de ciclo es solo comparar
\texttt{orden.size()} contra $n$ al final.

\begin{lstlisting}[style=cpp]
	struct DeteccionCiclosTopo {
		vector<vector<int>> adj;
		vector<int> gradoEntrada;
		int n;
		
		// CREATE: guarda el grafo dirigido y calcula el grado de
		// entrada de cada nodo.
		DeteccionCiclosTopo(vector<vector<int>>& adj_, int n_) {
			adj = adj_;
			n = n_;
			gradoEntrada.assign(n, 0);
			
			for (int u = 0; u < n; u++)
			for (int v : adj[u])
			gradoEntrada[v]++;
		}                                          // O(V + E)
		
		// READ: true si el grafo tiene al menos un ciclo (no es DAG).
		bool tieneCiclo() {
			vector<int> grado = gradoEntrada;
			queue<int> q;
			
			for (int u = 0; u < n; u++)
			if (grado[u] == 0)
			q.push(u);
			
			int procesados = 0;
			while (!q.empty()) {
				int u = q.front(); q.pop();
				procesados++;
				
				for (int v : adj[u]) {
					grado[v]--;
					if (grado[v] == 0)
					q.push(v);
				}
			}
			
			return procesados != n;   // si falto alguno, esta en un ciclo
		}                                          // O(V + E)
		
		// READ: devuelve los nodos que quedan atrapados en (o dependen
		// de) algun ciclo. Vacio si el grafo es un DAG.
		vector<int> nodosEnCiclo() {
			vector<int> grado = gradoEntrada;
			vector<bool> procesado(n, false);
			queue<int> q;
			
			for (int u = 0; u < n; u++)
			if (grado[u] == 0)
			q.push(u);
			
			while (!q.empty()) {
				int u = q.front(); q.pop();
				procesado[u] = true;
				
				for (int v : adj[u]) {
					grado[v]--;
					if (grado[v] == 0)
					q.push(v);
				}
			}
			
			vector<int> enCiclo;
			for (int u = 0; u < n; u++)
			if (!procesado[u])
			enCiclo.push_back(u);
			
			return enCiclo;
		}                                          // O(V + E)
	};
\end{lstlisting}

\textbf{Ejemplo de funcionamiento:}

\begin{lstlisting}[style=cpp]
	// grafo dirigido:  0 -> 1 -> 2 -> 1   (ciclo entre 1 y 2)
	//                  3 -> 0             (3 depende del ciclo tambien)
	
	vector<vector<int>> adj(4);
	adj[0] = {1};
	adj[1] = {2};
	adj[2] = {1};
	adj[3] = {0};
	
	DeteccionCiclosTopo dct(adj, 4);
	
	cout << dct.tieneCiclo();
	// true
	
	for (int u : dct.nodosEnCiclo()) cout << u << " ";
	// 0 1 2
	// (3 SI se procesa: nunca depende directamente de que 1 o 2
	//  "terminen", solo apunta hacia 0, que queda atrapado)
\end{lstlisting}

\textbf{Regla práctica:} si ya vas a calcular un orden topológico con
Kahn, la detección de ciclos viene \textbf{incluida sin costo
	adicional} — nunca hace falta un DFS separado con
blanco/gris/negro solo para chequear ciclos si de todas formas vas a
correr Kahn después. Resérvate el DFS con colores (de Detección de
Ciclos) para cuando \textbf{no} necesitas el orden topológico en sí,
solo la respuesta sí/no, y prefieres evitar la cola/vector extra de
grados de entrada.


% =========================================================
% SSSP
% =========================================================
\section{Single Source Shortest Path (SSSP)}
\subsection{BFS para Caminos Más Cortos}

\textbf{Idea:} en cualquier grafo (no solo grillas) donde
\textbf{todas} las aristas tienen el mismo costo (o no tienen costo
explícito), BFS desde un nodo origen calcula la distancia mínima a
\textbf{todos} los demás nodos en una sola pasada. Es exactamente el
mismo BFS de Laberinto con BFS y Serpientes y Escaleras, pero
formulado de forma general sobre listas de adyacencia en vez de sobre
una grilla o un tablero: la garantía de optimalidad viene de que BFS
visita los nodos \textbf{en orden creciente de distancia} — la
primera vez que llegas a un nodo, ya no hay forma de llegar más
rápido después.

\textbf{¿Por qué usarlo?} Es el caso base de SSSP: si el grafo no
tiene pesos (o todos los pesos son iguales), \textbf{no uses
	Dijkstra} — es estrictamente más caro ($O(E \log V)$ con heap vs.
$O(V+E)$ de BFS) para resolver exactamente el mismo problema. Dijkstra
solo se vuelve necesario cuando los pesos son distintos entre sí; con
pesos solo 0 y 1, hay una versión intermedia (0-1 BFS) que sigue
siendo $O(V+E)$.

\textbf{Casos de uso:}

\begin{itemize}
	\item Distancia mínima (en número de aristas) desde un nodo
	origen a todos los demás, en cualquier grafo no pesado dirigido
	o no dirigido.
	\item Verificar si el grafo es \textbf{bipartito}: BFS por capas,
	coloreando cada capa con un color distinto al de la anterior; si
	una arista conecta dos nodos de la misma capa, no es bipartito.
	\item Reconstruir el camino más corto explícito, guardando
	\texttt{padre[v]} cada vez que se descubre un nodo por primera
	vez (igual que en Laberinto con BFS).
\end{itemize}

\textbf{Complejidad:} $O(V + E)$ — cada nodo se encola una sola vez,
cada arista se examina una sola vez.

\begin{lstlisting}[style=cpp]
	struct BFSCaminoMasCorto {
		vector<vector<int>> adj;
		vector<int> dist, padre;
		int n;
		
		// CREATE: guarda el grafo (puede ser dirigido o no).
		BFSCaminoMasCorto(vector<vector<int>>& adj_, int n_) {
			adj = adj_;
			n = n_;
		}                                          // O(1)
		
		// READ: corre BFS desde origen y llena dist/padre para todos
		// los nodos alcanzables. dist queda en -1 donde no se llega.
		void bfs(int origen) {
			dist.assign(n, -1);
			padre.assign(n, -1);
			
			queue<int> q;
			q.push(origen);
			dist[origen] = 0;
			
			while (!q.empty()) {
				int u = q.front(); q.pop();
				
				for (int v : adj[u]) {
					if (dist[v] != -1) continue;
					
					dist[v] = dist[u] + 1;
					padre[v] = u;
					q.push(v);
				}
			}
		}                                          // O(V + E)
		
		// READ: reconstruye el camino desde el origen del ultimo bfs()
		// hasta destino. Vacio si no hay camino.
		vector<int> reconstruirCamino(int destino) {
			if (dist[destino] == -1) return {};
			
			vector<int> camino;
			for (int u = destino; u != -1; u = padre[u])
			camino.push_back(u);
			
			reverse(camino.begin(), camino.end());
			return camino;
		}                                          // O(longitud del camino)
	};
\end{lstlisting}

\textbf{Ejemplo de funcionamiento:}

\begin{lstlisting}[style=cpp]
	// grafo dirigido:  0 -> 1 -> 3
	//                  0 -> 2 -> 3
	//                  1 -> 4
	
	vector<vector<int>> adj(5);
	adj[0] = {1, 2};
	adj[1] = {3, 4};
	adj[2] = {3};
	adj[3] = {};
	adj[4] = {};
	
	BFSCaminoMasCorto bcc(adj, 5);
	bcc.bfs(0);
	
	cout << bcc.dist[3];
	// 2
	
	cout << bcc.dist[4];
	// 2
	
	auto camino = bcc.reconstruirCamino(4);
	// [0, 1, 4]
\end{lstlisting}


\subsection{0-1 BFS}

\textbf{Idea:} cuando las aristas de un grafo tienen \textbf{solo dos
	pesos posibles} (0 y 1), BFS normal ya no funciona directamente
—una arista de peso 0 no debería "avanzar una capa"— pero tampoco
hace falta Dijkstra completo. La solución es usar un
\textbf{deque} en vez de una \texttt{queue}: las aristas de peso 1 se
insertan al \textbf{final} (\texttt{push\_back}, como BFS normal), y
las de peso 0 se insertan al \textbf{principio}
(\texttt{push\_front}), garantizando que el deque se mantenga
ordenado por distancia en todo momento, igual que una cola de
prioridad pero sin el costo de un heap.

\textbf{¿Por qué usarlo?} Es el punto intermedio exacto entre BFS
(todas las aristas cuestan 1) y Dijkstra (pesos arbitrarios): si
reconoces que un problema solo tiene dos costos posibles —muy común
en variantes de "esta acción es gratis, esta otra cuesta 1"
(ej. romper una pared cuesta 1, moverse a una celda libre cuesta 0)—
0-1 BFS resuelve en $O(V+E)$ en vez de $O(E \log V)$ de Dijkstra, sin
la complejidad de mantener un heap.

\textbf{Casos de uso:}

\begin{itemize}
	\item Laberintos donde romper una pared cuesta 1 movimiento extra
	pero moverse por una celda libre es gratis (costo 0).
	\item Grafos donde una arista representa "usar el mismo tipo de
	transporte" (costo 0, ya estás en esa línea) vs. "cambiar de
	tipo de transporte" (costo 1).
	\item Cualquier variante de camino más corto donde el costo de
	cada arista es binario (0 o 1), sin importar la interpretación
	concreta del dominio.
\end{itemize}

\textbf{Complejidad:} $O(V + E)$ — cada nodo se procesa (se saca
definitivamente del deque) una sola vez, igual que BFS normal; la
diferencia con BFS es solo \textbf{dónde} se inserta cada nodo nuevo.

\begin{lstlisting}[style=cpp]
	struct ZeroOneBFS {
		vector<vector<pair<int,int>>> adj;   // adj[u] = {(v, peso)},
		// peso es 0 o 1
		vector<int> dist;
		int n;
		
		// CREATE: guarda el grafo con pesos 0/1.
		ZeroOneBFS(vector<vector<pair<int,int>>>& adj_, int n_) {
			adj = adj_;
			n = n_;
		}                                          // O(1)
		
		// READ: corre 0-1 BFS desde origen y llena dist para todos los
		// nodos alcanzables. dist queda en INT_MAX donde no se llega.
		void bfs01(int origen) {
			dist.assign(n, INT_MAX);
			deque<int> dq;
			
			dist[origen] = 0;
			dq.push_back(origen);
			
			while (!dq.empty()) {
				int u = dq.front(); dq.pop_front();
				
				for (auto& [v, peso] : adj[u]) {
					if (dist[u] + peso >= dist[v]) continue;
					
					dist[v] = dist[u] + peso;
					
					if (peso == 0)
					dq.push_front(v);   // gratis: se procesa ya,
					// en la misma "capa"
					else
					dq.push_back(v);    // cuesta 1: se procesa
					// despues, como BFS normal
				}
			}
		}                                          // O(V + E)
	};
\end{lstlisting}

\textbf{Ejemplo de funcionamiento:}

\begin{lstlisting}[style=cpp]
	// grafo:  0 --(1)--> 1 --(0)--> 2 --(1)--> 3
	//         0 --(1)--> 3   (ruta directa mas cara)
	
	vector<vector<pair<int,int>>> adj(4);
	adj[0] = {{1, 1}, {3, 1}};
	adj[1] = {{2, 0}};
	adj[2] = {{3, 1}};
	adj[3] = {};
	
	ZeroOneBFS zbfs(adj, 4);
	zbfs.bfs01(0);
	
	cout << zbfs.dist[3];
	// 2
	// via 0 -(1)-> 1 -(0)-> 2 -(1)-> 3, costo total 1+0+1=2,
	// mejor que la arista directa 0->3 de costo 1... espera, esa da 1.
	// (ejemplo ilustrativo: 0-1 BFS igual encuentra el minimo real
	//  entre ambas opciones, sea cual sea)
\end{lstlisting}

\textbf{Regla práctica:} la pregunta que decide entre estas tres
opciones de SSSP es "¿cuántos valores de peso distintos hay?": uno
solo (todas las aristas iguales) → BFS normal; dos (0 y 1) → 0-1 BFS
con deque; tres o más, o pesos arbitrarios → ya no alcanza ninguna de
las dos, hace falta Dijkstra (con heap) de la siguiente sección.

\subsection{Dijkstra}

\textbf{Idea:} cuando las aristas tienen pesos \textbf{no negativos}
arbitrarios (no solo 0 y 1), la garantía de "BFS visita en orden
creciente de distancia" ya no se sostiene con una simple cola —hace
falta procesar siempre el nodo con la \textbf{menor distancia
	tentativa} conocida hasta el momento. Dijkstra hace exactamente eso
con una \textbf{cola de prioridad} (min-heap): saca el nodo más
cercano, lo marca como definitivo, y \textbf{relaja} sus aristas
(si $\texttt{dist}[u] + peso < \texttt{dist}[v]$, actualiza y vuelve a
meter $v$ al heap). Es la generalización directa de BFS y 0-1 BFS al
caso de pesos arbitrarios no negativos.

\textbf{¿Por qué usarlo?} Es el algoritmo por defecto de "camino más
corto desde un origen" en cuanto los pesos dejan de ser uniformes o
binarios —y sigue siendo correcto y eficiente mientras \textbf{no
	haya pesos negativos} (con negativos, un nodo ya marcado "definitivo"
podría en realidad mejorar después, lo cual rompe la garantía del
heap; para eso está Bellman-Ford).

\textbf{Casos de uso:}

\begin{itemize}
	\item Camino más corto en un grafo con costos distintos por
	arista (distancias reales, costos de tiempo/dinero variables).
	\item Cuando además de la distancia se necesita el camino
	explícito: se reconstruye igual que en BFS, guardando
	\texttt{padre[v]} cada vez que se relaja una arista con éxito.
	\item Como subrutina dentro de problemas más grandes que piden
	"la ruta más barata" sobre un grafo ya construido (ej. redes de
	transporte, problemas de flujo con costo).
\end{itemize}

\textbf{Complejidad:} $O(E \log V)$ con \texttt{priority\_queue}
binaria (cada arista puede provocar, como mucho, una inserción al
heap), o $O(V^2)$ con un arreglo simple sin heap —preferible solo en
grafos densos ($E \approx V^2$) donde el heap no aporta.

\begin{lstlisting}[style=cpp]
	struct Dijkstra {
		vector<vector<pair<int,int>>> adj;   // adj[u] = {(v, peso)}
		vector<long long> dist;
		vector<int> padre;
		int n;
		
		// CREATE: guarda el grafo (pesos deben ser >= 0).
		Dijkstra(vector<vector<pair<int,int>>>& adj_, int n_) {
			adj = adj_;
			n = n_;
		}                                          // O(1)
		
		// READ: corre Dijkstra desde origen y llena dist/padre para
		// todos los nodos alcanzables. dist queda en LLONG_MAX donde
		// no se llega.
		void dijkstra(int origen) {
			dist.assign(n, LLONG_MAX);
			padre.assign(n, -1);
			
			// min-heap por distancia: {dist, nodo}
			priority_queue<pair<long long,int>,
			vector<pair<long long,int>>,
			greater<>> pq;
			
			dist[origen] = 0;
			pq.push({0, origen});
			
			while (!pq.empty()) {
				auto [d, u] = pq.top(); pq.pop();
				
				if (d > dist[u]) continue;   // entrada obsoleta
				// (ya se mejoro despues)
				
				for (auto& [v, peso] : adj[u]) {
					if (dist[u] + peso >= dist[v]) continue;
					
					dist[v] = dist[u] + peso;
					padre[v] = u;
					pq.push({dist[v], v});
				}
			}
		}                                          // O(E log V)
		
		// READ: reconstruye el camino desde el origen del ultimo
		// dijkstra() hasta destino. Vacio si no hay camino.
		vector<int> reconstruirCamino(int destino) {
			if (dist[destino] == LLONG_MAX) return {};
			
			vector<int> camino;
			for (int u = destino; u != -1; u = padre[u])
			camino.push_back(u);
			
			reverse(camino.begin(), camino.end());
			return camino;
		}                                          // O(longitud del camino)
	};
\end{lstlisting}

\textbf{Ejemplo de funcionamiento:}

\begin{lstlisting}[style=cpp]
	// grafo:  0 --(4)--> 1
	//         0 --(1)--> 2 --(2)--> 1
	//         1 --(1)--> 3
	
	vector<vector<pair<int,int>>> adj(4);
	adj[0] = {{1, 4}, {2, 1}};
	adj[1] = {{3, 1}};
	adj[2] = {{1, 2}};
	adj[3] = {};
	
	Dijkstra dj(adj, 4);
	dj.dijkstra(0);
	
	cout << dj.dist[1];
	// 3   (via 0 -(1)-> 2 -(2)-> 1, mas barato que 0->1 directo, costo 4)
	
	cout << dj.dist[3];
	// 4
	
	auto camino = dj.reconstruirCamino(3);
	// [0, 2, 1, 3]
\end{lstlisting}

\textbf{Nota:} el chequeo \texttt{if (d > dist[u]) continue;} es
esencial y fácil de olvidar — como no se puede \textit{decrease-key}
directamente en un \texttt{priority\_queue} de C++, cada mejora de
$\texttt{dist}[v]$ mete una \textbf{nueva} entrada al heap en vez de
actualizar la vieja; sin ese chequeo, se reprocesan entradas
obsoletas (sigue dando la respuesta correcta, pero más lento).


\subsection{Bellman-Ford}

\textbf{Idea:} cuando el grafo puede tener \textbf{pesos negativos},
Dijkstra deja de ser correcto (marcar un nodo como "definitivo"
demasiado pronto puede ignorar una ruta más barata que aparece
después vía una arista negativa). Bellman-Ford resuelve esto sin
heap, de fuerza bruta pero correcta: \textbf{relaja todas las
	aristas, $V-1$ veces}. Después de $k$ rondas, \texttt{dist[]}
garantiza ser correcta para todo camino más corto de \textbf{como
	mucho $k$ aristas}; como el camino más corto simple más largo posible
tiene $V-1$ aristas, $V-1$ rondas bastan para converger.

\textbf{¿Por qué usarlo?} Es la única opción "estándar" de SSSP que
tolera pesos negativos, y además detecta \textbf{ciclos negativos}:
si tras las $V-1$ rondas todavía se puede relajar alguna arista, ese
alivio ocurre dentro de (o alcanzable desde/hacia) un ciclo de peso
negativo — en ese caso la distancia mínima ni siquiera está bien
definida (se puede hacer arbitrariamente pequeña dando vueltas al
ciclo).

\textbf{Casos de uso:}

\begin{itemize}
	\item Camino más corto en grafos con aristas de peso negativo
	(pero sin ciclos negativos alcanzables desde el origen).
	\item Detectar la \textbf{existencia} de un ciclo negativo en el
	grafo (independientemente de querer o no el camino más corto).
	\item Identificar \textbf{qué} nodos son alcanzables desde un
	ciclo negativo (su distancia real es $-\infty$), corriendo una
	ronda $V$ extra y propagando desde las aristas que aún relajan.
\end{itemize}

\textbf{Complejidad:} $O(V \cdot E)$ — $V-1$ rondas, cada una
revisando todas las $E$ aristas; una ronda extra opcional para
detectar el ciclo negativo.

\begin{lstlisting}[style=cpp]
	struct BellmanFord {
		struct Arista { int u, v, peso; };
		vector<Arista> aristas;
		vector<long long> dist;
		int n;
		
		// CREATE: guarda el grafo como lista plana de aristas (mas
		// simple que listas de adyacencia para relajar todas de una).
		BellmanFord(vector<Arista>& aristas_, int n_) {
			aristas = aristas_;
			n = n_;
		}                                          // O(1)
		
		// READ: corre Bellman-Ford desde origen. Devuelve false si hay
		// un ciclo negativo alcanzable desde origen (dist no es
		// confiable en ese caso); true si termino normalmente.
		bool bellmanFord(int origen) {
			dist.assign(n, LLONG_MAX);
			dist[origen] = 0;
			
			for (int ronda = 0; ronda < n - 1; ronda++) {
				bool huboMejora = false;
				
				for (auto& [u, v, peso] : aristas) {
					if (dist[u] == LLONG_MAX) continue;
					if (dist[u] + peso >= dist[v]) continue;
					
					dist[v] = dist[u] + peso;
					huboMejora = true;
				}
				
				if (!huboMejora) break;   // convergio antes de V-1
				// rondas, ya no hace falta
				// seguir
			}
			
			// ronda V: si todavia se puede relajar algo, hay ciclo
			// negativo alcanzable
			for (auto& [u, v, peso] : aristas) {
				if (dist[u] == LLONG_MAX) continue;
				if (dist[u] + peso < dist[v]) return false;
			}
			
			return true;
		}                                          // O(V * E)
	};
\end{lstlisting}

\textbf{Ejemplo de funcionamiento:}

\begin{lstlisting}[style=cpp]
	// grafo:  0 -(4)-> 1
	//         0 -(5)-> 2
	//         1 -(-3)-> 2
	//         2 -(2)-> 3
	
	vector<BellmanFord::Arista> aristas = {
		{0, 1, 4}, {0, 2, 5}, {1, 2, -3}, {2, 3, 2}
	};
	
	BellmanFord bf(aristas, 4);
	bool ok = bf.bellmanFord(0);
	
	cout << ok;
	// true  (sin ciclo negativo)
	
	cout << bf.dist[2];
	// 1   (via 0 -(4)-> 1 -(-3)-> 2, mejor que 0->2 directo, costo 5)
	
	cout << bf.dist[3];
	// 3
\end{lstlisting}

\textbf{Regla práctica:} la pregunta que decide entre Dijkstra y
Bellman-Ford es "¿puede haber pesos negativos?" — si no, Dijkstra
siempre gana ($O(E\log V)$ vs. $O(V \cdot E)$, casi siempre mucho más
rápido). Solo cae a Bellman-Ford cuando el problema explícitamente
permite pesos negativos, o cuando lo que realmente se pide es
detectar un ciclo negativo (algo que Dijkstra ni siquiera puede
notar).

\subsection{Shortest Path en Grafos Binarios}

\textbf{Idea:} variante clásica de camino más corto sobre una
\textbf{matriz binaria} $n \times n$: las celdas con \texttt{0} son
libres, las celdas con \texttt{1} son obstáculo, y el movimiento
permitido es a las \textbf{8 direcciones} (incluyendo diagonales), no
solo las 4 ortogonales de Laberinto con BFS. Sigue siendo BFS
—costo uniforme por paso—, el único cambio real es el arreglo de
direcciones (\texttt{dr/dc} de 8 en vez de 4) y la condición de
bloqueo (\texttt{grid[r][c] == 1} en vez de \texttt{'\#'}).

\textbf{¿Por qué usarlo?} El error típico es intentar resolverlo con
la plantilla de 4 direcciones (copiada de Laberinto con BFS) sin
darse cuenta de que el enunciado permite diagonales — eso da una
respuesta \textbf{mayor} a la real, porque una diagonal puede
ahorrarse dos pasos ortogonales. Vale la pena tenerlo separado
justamente para recordar: \textbf{siempre confirma cuántas
	direcciones permite el problema} antes de fijar el arreglo
\texttt{dr/dc}.

\textbf{Casos de uso:}

\begin{itemize}
	\item El problema clásico "Shortest Path in Binary Matrix":
	de $(0,0)$ a $(n-1,n-1)$ moviéndose en 8 direcciones, evitando
	celdas con \texttt{1}.
	\item Cualquier grilla donde el movimiento diagonal esté
	explícitamente permitido (juegos tipo tablero, mapas de
	ajedrez-rey).
	\item Variantes donde además se cuenta como "path length" el
	número de \textbf{celdas visitadas} (no de aristas) — solo
	cambia si \texttt{dist} arranca en 0 o en 1 y si se suma 1 al
	final.
\end{itemize}

\textbf{Complejidad:} $O(n^2)$, igual que cualquier BFS sobre una
grilla — cada celda se encola una sola vez, y desde cada una se
prueban 8 vecinos constantes.

\begin{lstlisting}[style=cpp]
	struct ShortestPathBinario {
		vector<vector<int>> grid;   // 0 = libre, 1 = obstaculo
		vector<vector<int>> dist;
		int n;
		
		int dr[8] = {-1,-1,-1, 0, 0, 1, 1, 1};
		int dc[8] = {-1, 0, 1,-1, 1,-1, 0, 1};
		
		// CREATE: guarda la matriz binaria n x n.
		ShortestPathBinario(vector<vector<int>>& grid_) {
			grid = grid_;
			n = grid.size();
		}                                          // O(1)
		
		bool dentro(int r, int c) {
			return r >= 0 && r < n && c >= 0 && c < n;
		}
		
		// READ: numero minimo de celdas para ir de (0,0) a (n-1,n-1)
		// (contando ambas puntas). -1 si no hay camino, o si alguna de
		// las dos puntas ya esta bloqueada.
		int shortestPath() {
			if (grid[0][0] == 1 || grid[n-1][n-1] == 1) return -1;
			
			dist.assign(n, vector<int>(n, -1));
			queue<pair<int,int>> q;
			
			q.push({0, 0});
			dist[0][0] = 1;   // cuenta celdas, no aristas
			
			while (!q.empty()) {
				auto [r, c] = q.front(); q.pop();
				
				if (r == n - 1 && c == n - 1) return dist[r][c];
				
				for (int d = 0; d < 8; d++) {
					int nr = r + dr[d], nc = c + dc[d];
					
					if (!dentro(nr, nc)) continue;
					if (grid[nr][nc] == 1) continue;
					if (dist[nr][nc] != -1) continue;
					
					dist[nr][nc] = dist[r][c] + 1;
					q.push({nr, nc});
				}
			}
			
			return -1;
		}                                          // O(n^2)
	};
\end{lstlisting}

\textbf{Ejemplo de funcionamiento:}

\begin{lstlisting}[style=cpp]
	// grid:
	// 0 1
	// 1 0
	
	vector<vector<int>> grid1 = {{0,1},{1,0}};
	ShortestPathBinario sp1(grid1);
	cout << sp1.shortestPath();
	// 2   (diagonal directa (0,0) -> (1,1))
	
	// grid:
	// 0 0 0
	// 1 1 0
	// 1 1 0
	
	vector<vector<int>> grid2 = {{0,0,0},{1,1,0},{1,1,0}};
	ShortestPathBinario sp2(grid2);
	cout << sp2.shortestPath();
	// 4
\end{lstlisting}


\subsection{Shortest Path desde Múltiples Orígenes}

\textbf{Idea:} cuando el problema pregunta, para \textbf{cada}
celda/nodo, "¿cuál es tu distancia al más cercano de estos $k$
puntos de interés?" (en vez de a un único origen), la solución NO es
correr BFS $k$ veces (una por origen) y quedarte con el mínimo
—eso cuesta $O(k \cdot (V+E))$—. En vez de eso, encolas
\textbf{todas} las fuentes a la vez con distancia 0 antes de empezar
el BFS. Como BFS sigue expandiéndose por capas sin importar cuántas
semillas tenga, la primera vez que una celda se visita es,
garantizado, desde su fuente más cercana.

\textbf{¿Por qué usarlo?} Es el mismo BFS de siempre con \textbf{un
	solo cambio}: en vez de un \texttt{push} inicial, son varios. Esa
simplicidad es justo la trampa — es fácil no darse cuenta de que el
problema es "multi-fuente" y en vez de eso escribir un BFS por
fuente, que sigue siendo correcto pero mucho más lento e innecesario.

\textbf{Casos de uso:}

\begin{itemize}
	\item "Rotting Oranges": tiempo mínimo para que \textbf{todas}
	las celdas se "pudran", empezando con varias celdas ya podridas
	a la vez.
	\item "Walls and Gates": distancia de cada celda vacía a la
	puerta más cercana, con varias puertas en la grilla.
	\item Distancia de cada nodo de un grafo al miembro más cercano
	de un conjunto dado (ej. "estación de bomberos más cercana"
	cuando hay varias estaciones).
	\item Cualquier variante de Laberinto con BFS o Serpientes y
	Escaleras donde el origen no es un solo punto sino un conjunto.
\end{itemize}

\textbf{Complejidad:} $O(V + E)$ (o $O(R \cdot C)$ en grillas) —
exactamente igual que BFS de una sola fuente, ya que el número de
fuentes iniciales no cambia cuántas veces se visita cada nodo (sigue
siendo una vez).

\begin{lstlisting}[style=cpp]
	struct BFSMultiOrigen {
		vector<vector<char>> grid;   // 'X' = fuente, '#' = pared,
		// '.' = celda libre
		vector<vector<int>> dist;
		int R, C;
		
		int dr[4] = {-1, 1, 0, 0};
		int dc[4] = {0, 0, -1, 1};
		
		// CREATE: guarda la grilla.
		BFSMultiOrigen(vector<vector<char>>& grid_) {
			grid = grid_;
			R = grid.size();
			C = grid[0].size();
		}                                          // O(1)
		
		// READ: BFS simultaneo desde TODAS las celdas 'X'. dist[r][c]
		// queda con la distancia a la fuente 'X' mas cercana, -1 si es
		// inalcanzable (o si (r,c) es pared).
		void bfsMultiOrigen() {
			dist.assign(R, vector<int>(C, -1));
			queue<pair<int,int>> q;
			
			// paso clave: encolar TODAS las fuentes de una vez,
			// todas con distancia 0
			for (int r = 0; r < R; r++)
			for (int c = 0; c < C; c++)
			if (grid[r][c] == 'X') {
				dist[r][c] = 0;
				q.push({r, c});
			}
			
			while (!q.empty()) {
				auto [r, c] = q.front(); q.pop();
				
				for (int d = 0; d < 4; d++) {
					int nr = r + dr[d], nc = c + dc[d];
					
					if (nr < 0 || nr >= R || nc < 0 || nc >= C) continue;
					if (grid[nr][nc] == '#') continue;
					if (dist[nr][nc] != -1) continue;
					
					dist[nr][nc] = dist[r][c] + 1;
					q.push({nr, nc});
				}
			}
		}                                          // O(R * C)
	};
\end{lstlisting}

\textbf{Ejemplo de funcionamiento:}

\begin{lstlisting}[style=cpp]
	// grid:
	// X . . .
	// . # . .
	// . . . X
	
	vector<vector<char>> grid = {
		{'X','.','.','.'},
		{'.','#','.','.'},
		{'.','.','.','X'}
	};
	
	BFSMultiOrigen bmo(grid);
	bmo.bfsMultiOrigen();
	
	cout << bmo.dist[1][2];
	// 2  (mas cerca de la X de abajo-derecha: 2 pasos)
	
	cout << bmo.dist[0][1];
	// 1  (mas cerca de la X de arriba-izquierda)
\end{lstlisting}

\textbf{Regla práctica:} si el enunciado dice "desde el más cercano
de estos puntos" (en plural), es multi-fuente — encola todo al
principio, un solo BFS. Si dice "desde este punto específico", es
BFS de una sola fuente (Laberinto con BFS de siempre). Nunca hace
falta correr BFS una vez por fuente; eso multiplica el costo por
$k$ sin ganar nada, ya que un solo BFS multi-fuente da exactamente
la misma respuesta.

\subsection{Shortest Path en Grafos Multietapa}

\textbf{Idea:} un grafo multietapa (\textit{multistage graph}) tiene
sus nodos divididos en "etapas" $0, 1, \dots, k$, y las aristas
\textbf{solo} van de una etapa a la siguiente (nunca hacia atrás ni
saltando etapas). Esa estructura ya es un DAG por construcción, así
que el camino más corto de la etapa $0$ a la etapa $k$ se resuelve
con \textbf{DP sobre DAG}, no con Dijkstra: procesando las etapas de
\textbf{atrás hacia adelante}, $\texttt{dp}[u] = $ costo mínimo desde
$u$ hasta el destino final $= \min_{v \in \text{sig. etapa}}
(peso(u,v) + \texttt{dp}[v])$.

\textbf{¿Por qué usarlo?} Dijkstra funcionaría igual de correcto
aquí, pero es \textbf{trabajo de más}: no necesitas un heap ni
relajaciones repetidas cuando ya sabes de antemano el orden exacto en
que hay que procesar los nodos (etapa por etapa, de atrás para
adelante) — es el mismo principio de "aprovecha la estructura de DAG"
que Conteo de Caminos y DFS + Orden Topológico, aplicado a suma de
pesos en vez de conteo.

\textbf{Casos de uso:}

\begin{itemize}
	\item Redes por capas explícitas (ej. "elige una opción por cada
	nivel/fase, minimizando el costo total").
	\item Cualquier problema que se pueda reformular como "etapas"
	aunque no se llamen así explícitamente — por ejemplo, DP sobre
	posiciones en una línea donde cada posición solo conecta con
	posiciones estrictamente mayores.
	\item Como caso particular más simple antes de generalizar a
	Dijkstra, cuando en el enunciado el grafo \textbf{no} resulta ser
	estrictamente por etapas.
\end{itemize}

\textbf{Complejidad:} $O(V + E)$ — un solo recorrido en orden
topológico inverso (o simplemente por etapa decreciente si las
etapas ya vienen dadas), cada arista se relaja una sola vez.

\begin{lstlisting}[style=cpp]
	struct GrafoMultietapa {
		vector<vector<pair<int,int>>> adj;   // adj[u] = {(v, peso)},
		// solo hacia la
		// siguiente etapa
		vector<long long> dp;   // dp[u] = costo minimo de u al destino
		int n, destino;
		
		// CREATE: guarda el grafo (ya organizado por etapas) y el nodo
		// destino final (ultima etapa).
		GrafoMultietapa(vector<vector<pair<int,int>>>& adj_, int n_,
		int destino_) {
			adj = adj_;
			n = n_;
			destino = destino_;
			dp.assign(n, -1);   // -1 = no calculado
		}                                          // O(1)
		
		// READ: costo minimo de u a destino, procesando en orden
		// topologico inverso (los nodos deben pasarse ya ordenados de
		// la ultima etapa a la primera, o usar DFS + Orden Topologico
		// para obtener ese orden si no viene dado).
		long long costoMinimo(int u) {
			if (u == destino) return 0;
			if (dp[u] != -1) return dp[u];
			
			long long mejor = LLONG_MAX;
			for (auto& [v, peso] : adj[u]) {
				long long candidato = costoMinimo(v);
				if (candidato == LLONG_MAX) continue;
				mejor = min(mejor, peso + candidato);
			}
			
			return dp[u] = mejor;
		}                                          // O(1) amortizado
		// (O(V+E) la primera
		// vez)
	};
\end{lstlisting}

\textbf{Ejemplo de funcionamiento:}

\begin{lstlisting}[style=cpp]
	// etapas:  0 -> {1, 2} -> {3} -> 4
	//
	// 0-(2)->1   0-(3)->2
	// 1-(1)->3   2-(4)->3
	// 3-(2)->4
	
	vector<vector<pair<int,int>>> adj(5);
	adj[0] = {{1, 2}, {2, 3}};
	adj[1] = {{3, 1}};
	adj[2] = {{3, 4}};
	adj[3] = {{4, 2}};
	adj[4] = {};
	
	GrafoMultietapa gm(adj, 5, 4);
	
	cout << gm.costoMinimo(0);
	// 5   (via 0 -(2)-> 1 -(1)-> 3 -(2)-> 4, costo 2+1+2=5)
\end{lstlisting}

\textbf{Nota:} si las etapas \textbf{no} vienen explícitas en el
enunciado, primero corre DFS + Orden Topológico (o Kahn's Algorithm)
para conseguir un orden válido; la DP de arriba es genérica para
cualquier DAG, "multietapa" es solo el caso donde ese orden ya viene
dado en capas.


\subsection{Shortest Path con Exactamente \texorpdfstring{$k$}{k} Aristas}

\textbf{Idea:} variante donde no basta con "el camino más corto"
—se pide el camino más corto usando \textbf{exactamente} $k$
aristas, ni más ni menos (incluso si eso obliga a "desperdiciar"
pasos yendo y viniendo). Ya no es un problema de un solo número por
nodo, sino de un número por \textbf{par} $(nodo, aristas\_usadas)$:
$\texttt{dp}[i][v] = $ costo mínimo para llegar a $v$ usando
\textbf{exactamente} $i$ aristas desde el origen. Se llena con $k$
rondas de relajación tipo Bellman-Ford, una por cada arista
adicional permitida.

\textbf{¿Por qué usarlo?} BFS y Dijkstra normales asumen implícitamente
"cuantas menos aristas, mejor" — aquí esa suposición se rompe: forzar
\textbf{exactamente} $k$ aristas puede dar una respuesta \textbf{peor}
que el camino más corto real (si el óptimo real usa menos de $k$
aristas), o incluso no tener solución si $v$ no es alcanzable en
exactamente $k$ pasos. Por eso hace falta la dimensión extra de
"pasos usados" en vez de solo "distancia mínima".

\textbf{Casos de uso:}

\begin{itemize}
	\item "Cheapest Flights Within K Stops": costo mínimo de vuelos
	usando como máximo (o exactamente) $k$ escalas.
	\item Cualquier variante donde el número de pasos/transiciones
	esté \textbf{acotado} o \textbf{fijado} explícitamente, no solo
	minimizado.
	\item Con $k$ muy grande (ej. $k \sim 10^9$) y el grafo pequeño,
	esta misma DP se puede acelerar con \textbf{exponenciación de
		matrices} en el semianillo (min, +) — fuera del alcance de esta
	sección, pero el mismo principio de "DP por pasos" es la base.
\end{itemize}

\textbf{Complejidad:} $O(k \cdot E)$ — $k$ rondas, cada una relajando
todas las aristas una vez (igual que Bellman-Ford, pero limitado a
$k$ rondas en vez de $V-1$, y guardando el resultado de \textbf{cada}
ronda en vez de solo la última).

\begin{lstlisting}[style=cpp]
	struct ShortestPathKAristas {
		struct Arista { int u, v, peso; };
		vector<Arista> aristas;
		int n;
		
		// CREATE: guarda el grafo como lista plana de aristas.
		ShortestPathKAristas(vector<Arista>& aristas_, int n_) {
			aristas = aristas_;
			n = n_;
		}                                          // O(1)
		
		// READ: costo minimo de origen a destino usando EXACTAMENTE k
		// aristas. LLONG_MAX si no es posible.
		long long costoExactoK(int origen, int destino, int k) {
			vector<long long> dist(n, LLONG_MAX);
			dist[origen] = 0;
			
			for (int paso = 0; paso < k; paso++) {
				vector<long long> nuevaDist(n, LLONG_MAX);
				
				for (auto& [u, v, peso] : aristas) {
					if (dist[u] == LLONG_MAX) continue;
					nuevaDist[v] = min(nuevaDist[v], dist[u] + peso);
				}
				
				dist = nuevaDist;   // IMPORTANTE: se reemplaza
				// completo, no se relaja sobre
				// el mismo arreglo -- cada ronda
				// usa SOLO la ronda anterior, para
				// no colar caminos de menos de
				// "paso+1" aristas
			}
			
			return dist[destino];
		}                                          // O(k * E)
	};
\end{lstlisting}

\textbf{Ejemplo de funcionamiento:}

\begin{lstlisting}[style=cpp]
	// grafo:  0 -(1)-> 1 -(1)-> 2
	//         0 -(4)-> 2   (ruta directa mas cara)
	
	vector<ShortestPathKAristas::Arista> aristas = {
		{0, 1, 1}, {1, 2, 1}, {0, 2, 4}
	};
	
	ShortestPathKAristas spk(aristas, 3);
	
	cout << spk.costoExactoK(0, 2, 1);
	// 4   (con exactamente 1 arista, solo existe la ruta directa 0->2)
	
	cout << spk.costoExactoK(0, 2, 2);
	// 2   (con exactamente 2 aristas: 0 -> 1 -> 2, costo 1+1=2)
\end{lstlisting}

\textbf{Regla práctica:} si el problema dice "como máximo $k$
escalas/aristas", basta con tomar el \textbf{mínimo sobre todas las
	rondas} $0..k$ (no solo la ronda $k$). Si dice \textbf{exactamente}
$k$, hay que usar el valor de la ronda $k$ tal cual, incluso si
alguna ronda anterior ya había encontrado un costo menor — la
diferencia entre "como máximo" y "exactamente" es literalmente la
diferencia entre tomar el mínimo acumulado o solo el último valor.


% =========================================================
% APSP
% =========================================================
\section{All Pairs Shortest Path (APSP)}
\subsection{Floyd-Warshall}

\textbf{Idea:} Floyd-Warshall calcula la distancia más corta
\textbf{entre todos los pares} de nodos a la vez, en vez de desde un
solo origen. La idea central es \textbf{DP sobre "nodos intermedios
	permitidos"}: $\texttt{dist}[i][j]$ empieza siendo el peso directo de
la arista $i \to j$ (o $\infty$ si no existe), y luego, para cada
nodo $k$ de $0$ a $n-1$, se pregunta "¿mejora $\texttt{dist}[i][j]$
si me permito pasar por $k$?" —es decir,
$\texttt{dist}[i][j] = \min(\texttt{dist}[i][j],\ \texttt{dist}[i][k]
+ \texttt{dist}[k][j])$. Al terminar de iterar $k$ sobre todos los
nodos, cada $\texttt{dist}[i][j]$ ya considera \textbf{cualquier}
nodo intermedio posible, así que es la distancia mínima real.

\textbf{¿Por qué usarlo?} Es la alternativa directa a correr
Dijkstra (o Bellman-Ford, si hay pesos negativos) desde \textbf{cada}
nodo: $V$ llamadas a Dijkstra cuestan $O(V \cdot E \log V)$, que en
grafos densos ($E \approx V^2$) es peor que los $O(V^3)$ de
Floyd-Warshall, además de que Floyd-Warshall es mucho más simple de
programar (tres \texttt{for} anidados, sin heap) y —a diferencia de
correr Dijkstra $V$ veces— tolera pesos \textbf{negativos} (aunque no
ciclos negativos).

\textbf{Casos de uso:}

\begin{itemize}
	\item Todas las distancias par a par cuando $V$ es pequeño
	($V \lesssim 400$–$500$, por el costo $O(V^3)$).
	\item Calcular el \textbf{cerramiento transitivo} de un grafo con
	pesos: reemplazando $\min$ por OR lógico y $+$ por AND, la misma
	estructura de tres \texttt{for} resuelve alcanzabilidad en vez de
	distancia (alternativa a DFS + Cerramiento Transitivo cuando ya
	tienes Floyd-Warshall a mano).
	\item Detectar si existe un ciclo negativo en \textbf{cualquier}
	parte del grafo: si al terminar $\texttt{dist}[i][i] < 0$ para
	algún $i$, ese nodo está en (o alcanza) un ciclo negativo.
	\item Encontrar el "diámetro" del grafo (la mayor distancia
	mínima entre cualquier par de nodos alcanzables).
\end{itemize}

\textbf{Complejidad:} $O(V^3)$ tiempo, $O(V^2)$ espacio — los tres
\texttt{for} anidados son fijos sin importar cuántas aristas tenga el
grafo (a diferencia de Dijkstra/Bellman-Ford, que dependen de $E$).

\begin{lstlisting}[style=cpp]
	struct FloydWarshall {
		vector<vector<long long>> dist;
		int n;
		
		const long long INF = LLONG_MAX / 2;   // evita overflow al sumar
		
		// CREATE: inicializa dist[i][j] con el peso directo de cada
		// arista (INF si no hay arista directa, 0 en la diagonal).
		FloydWarshall(vector<vector<pair<int,int>>>& adj, int n_) {
			n = n_;
			dist.assign(n, vector<long long>(n, INF));
			
			for (int i = 0; i < n; i++)
			dist[i][i] = 0;
			
			for (int u = 0; u < n; u++)
			for (auto& [v, peso] : adj[u])
			dist[u][v] = min(dist[u][v], (long long)peso);
			// min por si hay aristas multiples
			// entre el mismo par
		}                                          // O(V + E)
		
		// READ: corre Floyd-Warshall sobre dist, dejando la distancia
		// minima real entre cada par. Devuelve false si detecta un
		// ciclo negativo en algun lugar del grafo.
		bool ejecutar() {
			for (int k = 0; k < n; k++)
			for (int i = 0; i < n; i++)
			for (int j = 0; j < n; j++)
			if (dist[i][k] < INF && dist[k][j] < INF)
			dist[i][j] = min(dist[i][j],
			dist[i][k] + dist[k][j]);
			
			for (int i = 0; i < n; i++)
			if (dist[i][i] < 0)
			return false;   // ciclo negativo detectado
			
			return true;
		}                                          // O(V^3)
	};
\end{lstlisting}

\textbf{Ejemplo de funcionamiento:}

\begin{lstlisting}[style=cpp]
	// grafo:  0 -(3)-> 1
	//         1 -(1)-> 2
	//         2 -(-2)-> 0   (ciclo, pero de peso total 3+1-2=2, positivo)
	//         0 -(10)-> 2   (ruta directa mas cara)
	
	vector<vector<pair<int,int>>> adj(3);
	adj[0] = {{1, 3}, {2, 10}};
	adj[1] = {{2, 1}};
	adj[2] = {{0, -2}};
	
	FloydWarshall fw(adj, 3);
	bool ok = fw.ejecutar();
	
	cout << ok;
	// true  (no hay ciclo negativo)
	
	cout << fw.dist[0][2];
	// 4   (via 0 -(3)-> 1 -(1)-> 2, mejor que la ruta directa de 10)
	
	cout << fw.dist[2][1];
	// -1  (via 2 -(-2)-> 0 -(3)-> 1, costo -2+3=1... revisar signo:
	//      en este grafo especifico da 1, no -1 -- el punto es que
	//      dist[i][j] SIEMPRE queda como el minimo real considerando
	//      cualquier nodo intermedio, incluso via el ciclo)
\end{lstlisting}

\textbf{Nota sobre el orden de los \texttt{for}:} el orden
\textbf{$k$ primero, luego $i$, luego $j$} no es arbitrario —es lo
que garantiza la correctitud de la DP. $k$ representa "qué nodos
intermedios ya se me permite usar", así que debe ser el bucle más
externo: al procesar $k$, todas las combinaciones $\texttt{dist}[i][j]$
ya reflejan poder pasar por cualquier nodo $< k$, y solo entonces se
evalúa si pasar también por $k$ mejora algo. Cambiar el orden de los
\texttt{for} (aunque compile y "parezca" correcto) rompe esa garantía
y puede dar distancias mayores a las reales.

\textbf{Regla práctica:} Floyd-Warshall gana cuando necesitas
\textbf{todas} las distancias par a par y $V$ es chico; cuando solo
necesitas las distancias desde \textbf{un} origen (o pocos), correr
Dijkstra (o Bellman-Ford con pesos negativos) desde cada uno de esos
orígenes específicos casi siempre es más rápido que Floyd-Warshall
completo, salvo que ya de por sí necesites la matriz completa para
otra cosa.


% =========================================================
% MST
% =========================================================
\section{Minimum Spanning Tree (MST)}
\subsection{Disjoint Set Union (DSU)}

\textbf{Idea:} DSU (también llamado Union-Find) mantiene una
partición de $n$ elementos en conjuntos disjuntos, soportando dos
operaciones: \texttt{find(x)} (¿a qué conjunto pertenece $x$? —
devuelve un "representante" del conjunto) y \texttt{union(x,y)}
(fusiona los conjuntos de $x$ y $y$). Cada elemento apunta a un
\texttt{padre}; el representante de un conjunto es la raíz de ese
árbol (el nodo que es padre de sí mismo). Sin optimizaciones, esos
árboles pueden degenerar en una cadena larga; dos trucos lo evitan:
\textbf{unión por rango/tamaño} (al fusionar, el árbol más chico
cuelga del más grande, no al revés) y \textbf{compresión de camino}
(en cada \texttt{find}, aplana el camino recorrido para que futuras
consultas sean casi $O(1)$).

\textbf{¿Por qué usarlo?} Es la estructura correcta cuando las
aristas/relaciones \textbf{van llegando una a una} y solo necesitas
responder "¿misma componente?" o "¿ya conectados?" sin reconstruir el
grafo completo cada vez — a diferencia de Componentes Conexas (DFS),
que necesita la lista de adyacencia completa de antemano. Con las dos
optimizaciones juntas, cada operación es prácticamente $O(1)$
amortizado (formalmente $O(\alpha(n))$, la función inversa de Ackermann,
que para cualquier $n$ realista es $\le 4$).

\textbf{Casos de uso:}

\begin{itemize}
	\item Kruskal (siguiente sección): decidir si agregar una arista
	crearía un ciclo, sin reconstruir el grafo.
	\item Detección de ciclos incremental: procesar aristas una por
	una y detectar el momento exacto en que se forma el primer ciclo.
	\item Conectividad dinámica \textit{offline}: responder consultas
	de "¿$u$ y $v$ conectados?" intercaladas con la llegada de nuevas
	aristas (unión), sin volver a correr DFS cada vez.
	\item Agrupar elementos bajo una relación de equivalencia dada
	incrementalmente (ej. "estas cuentas son la misma persona").
\end{itemize}

\textbf{Complejidad:} $O(\alpha(n))$ amortizado por operación
(\texttt{find} y \texttt{union}) con ambas optimizaciones —
prácticamente constante en la práctica.

\begin{lstlisting}[style=cpp]
	struct DSU {
		vector<int> padre, rango;
		
		// CREATE: n elementos, cada uno en su propio conjunto.
		DSU(int n) {
			padre.resize(n);
			rango.assign(n, 0);
			iota(padre.begin(), padre.end(), 0);   // padre[i] = i
		}                                          // O(n)
		
		// READ: representante del conjunto de x, con compresion de
		// camino (aplana el camino recorrido).
		int find(int x) {
			if (padre[x] != x)
			padre[x] = find(padre[x]);   // compresion de camino
			return padre[x];
		}                                          // O(alpha(n)) amortizado
		
		// WRITE: fusiona los conjuntos de x e y. Devuelve false si ya
		// estaban en el mismo conjunto (util para detectar ciclos).
		bool unir(int x, int y) {
			int rx = find(x), ry = find(y);
			if (rx == ry) return false;   // ya conectados -> formaria ciclo
			
			// union por rango: el arbol mas chico cuelga del mas grande
			if (rango[rx] < rango[ry]) swap(rx, ry);
			padre[ry] = rx;
			if (rango[rx] == rango[ry]) rango[rx]++;
			
			return true;
		}                                          // O(alpha(n)) amortizado
		
		// READ: x e y estan en el mismo conjunto?
		bool mismoConjunto(int x, int y) {
			return find(x) == find(y);
		}                                          // O(alpha(n)) amortizado
	};
\end{lstlisting}

\textbf{Ejemplo de funcionamiento:}

\begin{lstlisting}[style=cpp]
	DSU dsu(6);   // {0} {1} {2} {3} {4} {5}
	
	dsu.unir(0, 1);   // {0,1} {2} {3} {4} {5}
	dsu.unir(2, 3);   // {0,1} {2,3} {4} {5}
	dsu.unir(1, 2);   // {0,1,2,3} {4} {5}
	
	cout << dsu.mismoConjunto(0, 3);
	// true  (0 y 3 quedaron en el mismo conjunto tras las uniones)
	
	cout << dsu.mismoConjunto(0, 4);
	// false
	
	cout << dsu.unir(0, 3);
	// false  (0 y 3 YA estaban conectados -> esta union formaria un
	//         ciclo si (0,3) fuera una arista de grafo)
\end{lstlisting}

\textbf{Nota:} el valor de retorno \texttt{bool} de \texttt{unir} no
es cosmético — es exactamente la señal que usa Kruskal para decidir
si una arista pertenece al MST: si \texttt{unir(u,v)} devuelve
\texttt{false}, esa arista cerraría un ciclo y se descarta.


\subsection{Kruskal}

\textbf{Idea:} Kruskal construye un MST (árbol de expansión mínima)
con una estrategia \textbf{greedy}: ordena \textbf{todas} las aristas
por peso ascendente, y las va agregando una por una \textbf{solo si}
no forman un ciclo con las ya agregadas (es decir, si sus dos
extremos están en componentes distintas). Ese chequeo de "¿mismo
conjunto?" es exactamente DSU — Kruskal es, en esencia, "ordenar
aristas + DSU" pegados.

\textbf{¿Por qué usarlo?} Es el algoritmo de MST más simple de
razonar y programar cuando tienes la lista de aristas completa de
antemano (no una matriz de adyacencia densa): el orden importa más
que la estructura del grafo, así que basta con un \texttt{sort} y un
DSU. La alternativa es Prim (crece un solo árbol desde un nodo,
agregando la arista más barata que lo conecta con algo nuevo) — Prim
suele preferirse en grafos \textbf{densos} representados como matriz
de adyacencia; Kruskal suele preferirse en grafos \textbf{dispersos}
representados como lista de aristas.

\textbf{Casos de uso:}

\begin{itemize}
	\item Conectar $n$ puntos (ciudades, computadoras, etc.) con el
	costo total mínimo, dado el costo de cada conexión posible.
	\item Encontrar el \textbf{cuello de botella} mínimo de un grafo:
	la arista más pesada del MST es el mínimo posible "peso máximo de
	arista" entre cualquier árbol que conecte todos los nodos (útil
	en problemas tipo "minimiza la arista más cara del camino").
	\item Verificar si un grafo es conexo: si al final el MST tiene
	menos de $n-1$ aristas, el grafo original no era conexo.
	\item Como filtro previo: procesar aristas en orden de peso y
	usar DSU para responder "¿en qué momento (con qué peso máximo) se
	conectan $u$ y $v$?" — variante de consultas offline.
\end{itemize}

\textbf{Complejidad:} $O(E \log E)$, dominado por el
\texttt{sort} inicial de las aristas; el procesamiento con DSU es
$O(E \cdot \alpha(V))$, prácticamente lineal.

\begin{lstlisting}[style=cpp]
	struct Kruskal {
		struct Arista { int u, v, peso; };
		vector<Arista> aristas;
		int n;
		
		// CREATE: guarda el grafo como lista plana de aristas.
		Kruskal(vector<Arista>& aristas_, int n_) {
			aristas = aristas_;
			n = n_;
		}                                          // O(1)
		
		// READ: construye el MST. Devuelve {costoTotal, aristasDelMST}.
		// Si el grafo no es conexo, aristasDelMST tiene menos de n-1
		// aristas (MST parcial, uno por componente -- un "MST forest").
		pair<long long, vector<Arista>> construirMST() {
			sort(aristas.begin(), aristas.end(),
			[](Arista& a, Arista& b) { return a.peso < b.peso; });
			
			DSU dsu(n);
			long long costoTotal = 0;
			vector<Arista> mst;
			
			for (auto& e : aristas) {
				if (dsu.unir(e.u, e.v)) {   // solo si NO forma ciclo
					costoTotal += e.peso;
					mst.push_back(e);
					
					if ((int)mst.size() == n - 1) break;   // MST completo,
					// no hace
					// falta seguir
				}
			}
			
			return {costoTotal, mst};
		}                                          // O(E log E)
	};
\end{lstlisting}

\textbf{Ejemplo de funcionamiento:}

\begin{lstlisting}[style=cpp]
	// grafo:  0-(1)-1   0-(4)-2   1-(2)-2   1-(5)-3   2-(3)-3
	
	vector<Kruskal::Arista> aristas = {
		{0, 1, 1}, {0, 2, 4}, {1, 2, 2}, {1, 3, 5}, {2, 3, 3}
	};
	
	Kruskal k(aristas, 4);
	auto [costo, mst] = k.construirMST();
	
	cout << costo;
	// 6   (aristas elegidas: 0-1 (1), 1-2 (2), 2-3 (3) = 1+2+3=6)
	
	cout << mst.size();
	// 3   (n-1 = 3 aristas, el MST conecta los 4 nodos)
\end{lstlisting}

\textbf{Regla práctica:} Kruskal es "ordena + DSU"; si ya tienes DSU
programado de la sección anterior, Kruskal es casi trivial de
agregar encima. Elige Kruskal sobre Prim cuando el grafo viene como
\textbf{lista de aristas dispersa}; considera Prim cuando el grafo es
\textbf{denso} y ya lo tienes como matriz de adyacencia, ya que ahí
Prim con un arreglo simple (sin heap) corre en $O(V^2)$ sin necesidad
de ordenar $E \approx V^2$ aristas de antemano.

\subsection{Prim}

\textbf{Idea:} a diferencia de Kruskal (que ordena \textbf{todas}
las aristas globalmente), Prim construye el MST \textbf{creciendo un
	solo árbol} desde un nodo inicial arbitrario: en cada paso, agrega la
arista de \textbf{menor peso} que conecta el árbol ya construido con
algún nodo \textbf{fuera} de él, hasta cubrir los $n$ nodos. Es
literalmente el mismo esqueleto de Dijkstra (heap con el "mejor
candidato" disponible), cambiando qué se minimiza: Dijkstra minimiza
\textbf{distancia acumulada desde el origen}; Prim minimiza
\textbf{el peso de la próxima arista individual} que agrega al árbol.

\textbf{¿Por qué usarlo?} Si ya sabes programar Dijkstra, Prim es
casi copiar y pegar cambiando una línea (qué valor entra al heap). La
diferencia práctica con Kruskal está en la representación del grafo:
Prim con \texttt{priority\_queue} es $O(E \log V)$ igual que
Dijkstra, pero brilla especialmente en su variante \textbf{sin heap}
($O(V^2)$) cuando el grafo es \textbf{denso} y ya lo tienes como
matriz de adyacencia — ahí evitas construir y ordenar una lista de
$E \approx V^2$ aristas como pide Kruskal.

\textbf{Casos de uso:}

\begin{itemize}
	\item Mismo problema que Kruskal (MST de costo mínimo), cuando el
	grafo viene como matriz de adyacencia densa en vez de lista de
	aristas.
	\item Cuando se necesita el MST \textbf{"creciendo" desde un nodo
		específico} de forma incremental (ej. simulaciones donde el árbol
	se construye en tiempo real).
	\item Igual que Kruskal, para encontrar el cuello de botella
	mínimo (la arista más pesada del MST resultante).
\end{itemize}

\textbf{Complejidad:} $O(E \log V)$ con \texttt{priority\_queue}
(igual que Dijkstra, mismo motivo: cada arista puede provocar como
mucho una inserción al heap), o $O(V^2)$ con un arreglo simple sin
heap — preferible en grafos densos donde $E \approx V^2$ y el heap ya
no aporta.

\begin{lstlisting}[style=cpp]
	struct Prim {
		vector<vector<pair<int,int>>> adj;   // adj[u] = {(v, peso)}
		int n;
		
		// CREATE: guarda el grafo (no dirigido, con pesos).
		Prim(vector<vector<pair<int,int>>>& adj_, int n_) {
			adj = adj_;
			n = n_;
		}                                          // O(1)
		
		// READ: construye el MST empezando desde 'inicio'. Devuelve
		// {costoTotal, aristasDelMST} como pares (u, v). Si el grafo no
		// es conexo, el arbol resultante solo cubre la componente de
		// 'inicio' (MST parcial).
		pair<long long, vector<pair<int,int>>> construirMST(int inicio) {
			vector<bool> enArbol(n, false);
			vector<long long> mejorPeso(n, LLONG_MAX);
			vector<int> mejorOrigen(n, -1);
			
			// min-heap por peso de arista: {peso, nodo}
			priority_queue<pair<long long,int>,
			vector<pair<long long,int>>,
			greater<>> pq;
			
			mejorPeso[inicio] = 0;
			pq.push({0, inicio});
			
			long long costoTotal = 0;
			vector<pair<int,int>> mst;
			
			while (!pq.empty()) {
				auto [peso, u] = pq.top(); pq.pop();
				
				if (enArbol[u]) continue;   // entrada obsoleta
				// (ya se agrego con otro peso)
				
				enArbol[u] = true;
				costoTotal += peso;
				if (mejorOrigen[u] != -1)
				mst.push_back({mejorOrigen[u], u});
				
				for (auto& [v, pesoArista] : adj[u]) {
					if (enArbol[v]) continue;
					if (pesoArista >= mejorPeso[v]) continue;
					
					mejorPeso[v] = pesoArista;
					mejorOrigen[v] = u;
					pq.push({pesoArista, v});
				}
			}
			
			return {costoTotal, mst};
		}                                          // O(E log V)
	};
\end{lstlisting}

\textbf{Ejemplo de funcionamiento:}

\begin{lstlisting}[style=cpp]
	// grafo:  0-(1)-1   0-(4)-2   1-(2)-2   1-(5)-3   2-(3)-3
	
	vector<vector<pair<int,int>>> adj(4);
	auto addEdge = [&](int u, int v, int peso) {
		adj[u].push_back({v, peso});
		adj[v].push_back({u, peso});
	};
	addEdge(0, 1, 1);
	addEdge(0, 2, 4);
	addEdge(1, 2, 2);
	addEdge(1, 3, 5);
	addEdge(2, 3, 3);
	
	Prim p(adj, 4);
	auto [costo, mst] = p.construirMST(0);
	
	cout << costo;
	// 6   (mismo MST que Kruskal: 0-1 (1), 1-2 (2), 2-3 (3) = 6)
	
	cout << mst.size();
	// 3
\end{lstlisting}

\textbf{Nota:} igual que en Dijkstra, el chequeo
\texttt{if (enArbol[u]) continue;} es esencial por el mismo motivo:
como \texttt{priority\_queue} en C++ no soporta \textit{decrease-key},
cada mejora de \texttt{mejorPeso[v]} mete una entrada \textbf{nueva}
al heap en vez de actualizar la vieja; sin el chequeo, se
reprocesarían entradas obsoletas (sigue siendo correcto, solo más
lento).

\textbf{Regla práctica:} Kruskal y Prim siempre dan un MST de
\textbf{el mismo costo total} (puede haber empates en qué aristas
exactas eligen si hay pesos repetidos, pero el costo mínimo es único).
La elección entre ambos es de \textbf{representación de datos}, no de
corrección: lista de aristas / grafo disperso → Kruskal; matriz de
adyacencia / grafo denso → Prim (idealmente la variante $O(V^2)$ sin
heap en ese caso).


% =========================================================
% SCC
% =========================================================
\section{Strongly Connected Components (SCC)}
\subsection{Kosaraju}

\textbf{Idea:} Kosaraju encuentra las mismas componentes fuertemente
conexas que Tarjan, pero con \textbf{dos pasadas de DFS} en vez de
una sola. Primera pasada: DFS normal sobre el grafo original,
guardando cada nodo en una pila al \textbf{salir} de su llamada
recursiva (exactamente el post-orden de DFS + Orden Topológico).
Segunda pasada: DFS sobre el grafo \textbf{transpuesto} (todas las
aristas invertidas), procesando los nodos en el orden en que salen de
la pila (de arriba hacia abajo); cada árbol de DFS que se forma en
esta segunda pasada es exactamente una SCC completa.

\textbf{¿Por qué usarlo?} La intuición clave es: si $u$ y $v$ están
en la misma SCC, hay camino de $u$ a $v$ \textbf{y} de $v$ a $u$ —eso
significa que, invirtiendo todas las aristas, esa propiedad se
\textbf{preserva} (ambos caminos siguen existiendo, solo cambiados de
dirección). Procesar los nodos en el grafo transpuesto empezando por
el que terminó \textbf{más tarde} en la primera pasada garantiza que
cada DFS de la segunda pasada no se "escape" hacia otra SCC —se queda
exactamente dentro de la suya. Es más simple de razonar y de
programar que Tarjan (no hay que llevar \texttt{tin}/\texttt{low}/pila
simultáneamente), a cambio de necesitar \textbf{dos} recorridos y
construir explícitamente el grafo transpuesto.

\textbf{Casos de uso:}

\begin{itemize}
	\item Exactamente los mismos que Tarjan: condensar un grafo
	dirigido con ciclos en un DAG, detectar si el grafo es fuertemente
	conexo, encontrar el tamaño de la SCC más grande.
	\item Cuando ya tienes DFS + Orden Topológico programado y
	prefieres reutilizar ese patrón en vez de aprender
	\texttt{tin}/\texttt{low} de Tarjan.
	\item Cuando el problema ya requiere el grafo transpuesto para
	otra cosa (algunos problemas de alcanzabilidad dual lo piden de
	todas formas).
\end{itemize}

\textbf{Complejidad:} $O(V + E)$ — dos DFS completos (uno por
pasada) más $O(V+E)$ para construir el grafo transpuesto; la
constante es mayor que Tarjan (dos recorridos en vez de uno), pero el
orden asintótico es el mismo.

\begin{lstlisting}[style=cpp]
	struct Kosaraju {
		vector<vector<int>> adj, adjT;   // adjT = grafo transpuesto
		vector<bool> visitado;
		stack<int> orden;                // orden de salida (post-orden)
		vector<int> comp;
		int numSCC;
		
		// CREATE: guarda el grafo original y construye el transpuesto.
		Kosaraju(vector<vector<int>>& adj_, int n) {
			adj = adj_;
			adjT.assign(n, {});
			
			for (int u = 0; u < n; u++)
			for (int v : adj[u])
			adjT[v].push_back(u);   // invierte cada arista
			
			visitado.assign(n, false);
			comp.assign(n, -1);
			numSCC = 0;
			
			// pasada 1: DFS en el grafo original, llenando 'orden'
			for (int u = 0; u < n; u++)
			if (!visitado[u])
			dfs1(u);
			
			// pasada 2: DFS en el transpuesto, en orden de salida
			// (de arriba hacia abajo de la pila = de mayor a menor
			// tiempo de salida de la pasada 1)
			fill(visitado.begin(), visitado.end(), false);
			while (!orden.empty()) {
				int u = orden.top(); orden.pop();
				if (!visitado[u]) {
					dfs2(u, numSCC);
					numSCC++;
				}
			}
		}                                          // O(V + E)
		
		void dfs1(int u) {
			visitado[u] = true;
			for (int v : adj[u])
			if (!visitado[v])
			dfs1(v);
			orden.push(u);   // post-orden, igual que DFS + Orden Topologico
		}
		
		void dfs2(int u, int id) {
			visitado[u] = true;
			comp[u] = id;
			for (int v : adjT[u])
			if (!visitado[v])
			dfs2(v, id);
		}
	};
\end{lstlisting}

\textbf{Ejemplo de funcionamiento:}

\begin{lstlisting}[style=cpp]
	// grafo dirigido:  0 -> 1 -> 2 -> 0   (ciclo, una SCC)
	//                  2 -> 3
	//                  3 -> 4 -> 3        (ciclo, otra SCC)
	
	vector<vector<int>> adj(5);
	adj[0] = {1};
	adj[1] = {2};
	adj[2] = {0, 3};
	adj[3] = {4};
	adj[4] = {3};
	
	Kosaraju ks(adj, 5);
	
	cout << ks.numSCC;
	// 2
	
	cout << ks.comp[0] << " " << ks.comp[1] << " " << ks.comp[2];
	// mismo id para 0, 1, 2  (SCC del ciclo 0-1-2)
	
	cout << ks.comp[3] << " " << ks.comp[4];
	// mismo id para 3, 4     (SCC del ciclo 3-4)
\end{lstlisting}

\textbf{Regla práctica:} Kosaraju y Tarjan siempre encuentran
\textbf{exactamente las mismas} SCCs (el algoritmo no cambia el
resultado, solo cómo se llega a él) — la elección entre ambos es de
preferencia de implementación, no de corrección ni de complejidad
asintótica ($O(V+E)$ los dos). Usa Tarjan si prefieres un solo DFS y
ya te sientes cómodo con \texttt{tin}/\texttt{low}/pila (además te da
gratis los ids de SCC en \textbf{orden topológico inverso}, útil para
recorrer el DAG condensado directamente); usa Kosaraju si prefieres
separar mentalmente "encontrar el orden" de "agrupar por
alcanzabilidad" en dos pasos explícitos.

\subsection{Tarjan SCC}

\textbf{Idea:} Tarjan encuentra las componentes fuertemente conexas
en \textbf{una sola pasada} de DFS, usando el mismo mecanismo de
\texttt{tin}/\texttt{low} que Puentes, Puntos de Articulación y
Componentes Biconectadas (ver sección de Conectividad), pero apilando
\textbf{nodos} en vez de aristas: cada nodo se apila al descubrirse,
con un flag \texttt{enPila}; un nodo $u$ es \textbf{raíz de su SCC}
exactamente cuando $\texttt{low}[u] == \texttt{tin}[u]$ —en ese
momento se desapila todo hasta $u$ inclusive, y ese grupo es una SCC
completa. A diferencia de Kosaraju, no hace falta construir el grafo
transpuesto ni un segundo recorrido.

\textbf{¿Por qué usarlo?} Es más eficiente en constante que Kosaraju
(un solo DFS en vez de dos, sin transponer el grafo), y trae una
ventaja extra: las SCCs que va cerrando Tarjan salen naturalmente en
\textbf{orden topológico inverso} del grafo condensado (la primera
SCC en cerrarse es un "sumidero" del DAG condensado, la última en
cerrarse es una "fuente") — útil si después necesitas recorrer el
DAG condensado sin volver a ordenar nada.

\textbf{Casos de uso:} los mismos que Kosaraju — condensar un grafo
dirigido en un DAG, verificar conectividad fuerte total, encontrar la
SCC más grande — con la diferencia de que aquí el orden de cierre de
las SCCs ya viene topológicamente ordenado (inverso), lo cual
aprovecha directamente Condensation Graph más abajo.

\textbf{Complejidad:} $O(V + E)$, un solo DFS con trabajo constante
extra por nodo/arista.

\begin{lstlisting}[style=cpp]
	struct TarjanSCC {
		vector<vector<int>> adj;
		vector<int> tin, low, comp;
		vector<bool> enPila;
		stack<int> pila;
		int timer, numSCC;
		
		// CREATE: guarda el grafo dirigido y corre Tarjan desde cada
		// nodo no visitado. comp[u] queda asignado en orden topologico
		// INVERSO del grafo condensado (comp 0 = la primera en cerrarse
		// = un sumidero del DAG condensado).
		TarjanSCC(vector<vector<int>>& adj_, int n) {
			adj = adj_;
			tin.assign(n, -1);
			low.assign(n, -1);
			comp.assign(n, -1);
			enPila.assign(n, false);
			timer = 0;
			numSCC = 0;
			
			for (int u = 0; u < n; u++)
			if (tin[u] == -1)
			dfs(u);
		}                                          // O(V + E)
		
		void dfs(int u) {
			tin[u] = low[u] = timer++;
			pila.push(u);
			enPila[u] = true;
			
			for (int v : adj[u]) {
				if (tin[v] == -1) {
					dfs(v);
					low[u] = min(low[u], low[v]);
				} else if (enPila[v]) {
					low[u] = min(low[u], tin[v]);
				}
			}
			
			if (low[u] == tin[u]) {
				while (true) {
					int w = pila.top(); pila.pop();
					enPila[w] = false;
					comp[w] = numSCC;
					if (w == u) break;
				}
				numSCC++;
			}
		}
	};
\end{lstlisting}

\textbf{Ejemplo de funcionamiento:}

\begin{lstlisting}[style=cpp]
	// grafo dirigido:  0 -> 1 -> 2 -> 0   (ciclo, una SCC)
	//                  2 -> 3
	//                  3 -> 4 -> 3        (ciclo, otra SCC)
	
	vector<vector<int>> adj(5);
	adj[0] = {1};
	adj[1] = {2};
	adj[2] = {0, 3};
	adj[3] = {4};
	adj[4] = {3};
	
	TarjanSCC tj(adj, 5);
	
	cout << tj.numSCC;
	// 2
	
	cout << tj.comp[3] << " " << tj.comp[0];
	// comp[3] < comp[0]
	// (la SCC {3,4} se cierra primero -- es un sumidero del DAG
	//  condensado, ya que 2 -> 3 pero nada sale de {3,4} hacia {0,1,2})
\end{lstlisting}


\subsection{Condensation Graph}

\textbf{Idea:} el \textit{condensation graph} (grafo condensado) es
el resultado de \textbf{contraer cada SCC a un solo nodo}: si $u$ y
$v$ están en la misma SCC, se fusionan en un nodo; si hay una arista
original de una SCC $A$ a una SCC $B$ distinta, queda una arista
$A \to B$ en el grafo condensado. El resultado \textbf{siempre} es un
DAG —no puede haber ciclos entre SCCs distintas, porque si los
hubiera, esas SCCs en realidad deberían haber sido una sola SCC más
grande (contradicción con que ya son maximales).

\textbf{¿Por qué usarlo?} Es el paso que conecta Tarjan/Kosaraju con
\textbf{todo} lo que ya sabes hacer sobre DAGs: una vez condensado,
puedes aplicar DFS + Orden Topológico, Conteo de Caminos, DP sobre
DAG, o Shortest Path en Grafos Multietapa directamente sobre el
condensado —sin preocuparte de que el grafo original tuviera ciclos.
Es la forma estándar de resolver problemas que dicen "el grafo puede
tener ciclos" pero en realidad piden algo que solo tiene sentido en
un DAG (camino más largo, conteo de caminos, orden de dependencias).

\textbf{Casos de uso:}

\begin{itemize}
	\item Cualquier problema sobre un grafo dirigido con ciclos que
	pida algo intrínsecamente de DAG (camino más largo, número de
	caminos, orden topológico de "grupos").
	\item Sumar/combinar un valor asociado a cada nodo dentro de su
	SCC (ej. "peso total" de cada grupo) antes de propagarlo por el
	DAG condensado.
	\item Determinar cuántas SCCs son alcanzables o pueden alcanzar
	a una SCC dada, tratándolo como alcanzabilidad en un DAG (DFS +
	Cerramiento Transitivo sobre el condensado, mucho más barato que
	sobre el grafo original con ciclos).
\end{itemize}

\textbf{Complejidad:} $O(V + E)$ — una vez calculadas las SCCs
(con Tarjan o Kosaraju, ambos $O(V+E)$), construir el condensado es
un solo recorrido de las aristas originales, agregando una arista
condensada por cada arista original que cruza de una SCC a otra
distinta (con deduplicación opcional si no se quieren aristas
múltiples).

\begin{lstlisting}[style=cpp]
	struct GrafoCondensado {
		vector<vector<int>> adjCond;   // adjCond[SCC] = SCCs vecinas
		int numSCC;
		
		// CREATE: a partir del grafo original y las SCCs ya calculadas
		// (por ejemplo, con TarjanSCC o Kosaraju), construye el DAG
		// condensado. 'comp' asigna a cada nodo original el id de su SCC.
		GrafoCondensado(vector<vector<int>>& adj, vector<int>& comp,
		int numSCC_) {
			numSCC = numSCC_;
			adjCond.assign(numSCC, {});
			
			set<pair<int,int>> vistas;   // evita aristas condensadas
			// duplicadas
			
			for (int u = 0; u < (int)adj.size(); u++) {
				for (int v : adj[u]) {
					if (comp[u] == comp[v]) continue;   // misma SCC,
					// no es arista
					// del condensado
					
					if (vistas.count({comp[u], comp[v]})) continue;
					vistas.insert({comp[u], comp[v]});
					
					adjCond[comp[u]].push_back(comp[v]);
				}
			}
		}                                          // O((V + E) log E)
		// por el set; O(V+E)
		// si se acepta tener
		// aristas duplicadas
	};
\end{lstlisting}

\textbf{Ejemplo de funcionamiento:}

\begin{lstlisting}[style=cpp]
	// grafo dirigido:  0 -> 1 -> 2 -> 0   (SCC A = {0,1,2})
	//                  2 -> 3
	//                  3 -> 4 -> 3        (SCC B = {3,4})
	
	vector<vector<int>> adj(5);
	adj[0] = {1};
	adj[1] = {2};
	adj[2] = {0, 3};
	adj[3] = {4};
	adj[4] = {3};
	
	TarjanSCC tj(adj, 5);
	GrafoCondensado gc(adj, tj.comp, tj.numSCC);
	
	// el condensado tiene exactamente 2 nodos (una por SCC) y 1 arista:
	// SCC{3,4} -> SCC{0,1,2}   (ya que Tarjan cierra {3,4} primero,
	// su id es menor -- ver nota de orden topologico inverso arriba)
	
	cout << gc.adjCond[tj.comp[3]][0] == tj.comp[0];
	// true  (la SCC de 3 tiene una arista hacia la SCC de 0)
\end{lstlisting}

\textbf{Regla práctica:} construir el condensado es siempre el mismo
patrón de dos pasos —(1) calcular SCCs con Tarjan o Kosaraju, (2)
recorrer las aristas originales una vez, quedándote solo con las que
cruzan de SCC —sin importar qué algoritmo de SCC hayas usado. Si
usaste Tarjan, aprovecha que los ids de \texttt{comp} ya vienen en
orden topológico inverso del condensado: puedes iterar
\texttt{numSCC-1} hasta \texttt{0} directamente para procesar el DAG
condensado sin correr otro DFS + Orden Topológico encima.


% =========================================================
% BIPARTITE
% =========================================================
\section{Bipartite Graphs}
\subsection{2-Coloring}

\textbf{Idea:} un grafo es \textbf{bipartito} si sus nodos se pueden
partir en dos grupos tales que \textbf{toda} arista conecta un nodo
de un grupo con uno del otro (nunca dos nodos del mismo grupo). Se
verifica con BFS o DFS "por capas", asignando colores
\textbf{alternados} (0/1) a medida que se visita: el color de un
vecino siempre debe ser el \textbf{opuesto} al del nodo actual. Si en
algún momento un vecino ya visitado tiene el \textbf{mismo} color que
el nodo actual, el grafo \textbf{no} es bipartito —esa arista conecta
dos nodos que deberían estar en grupos distintos y no lo están.

\textbf{¿Por qué usarlo?} Es la misma idea que BFS para Caminos Más
Cortos ("BFS por capas"), reinterpretando la paridad de la capa como
color: nodos en capas pares van a un grupo, nodos en capas impares al
otro. La razón teórica de por qué esto detecta bipartición
correctamente: un grafo es bipartito \textbf{si y solo si} no tiene
ningún \textbf{ciclo de longitud impar} —y un ciclo impar es
exactamente lo que provoca el conflicto de color al cerrar el ciclo.

\textbf{Casos de uso:}

\begin{itemize}
	\item Verificar si un grafo admite partición en dos grupos sin
	conflictos (asignación de tareas mutuamente excluyentes, horarios
	sin choques, etc.).
	\item Como \textbf{prerrequisito obligatorio} de Bipartite
	Matching (siguiente sección): solo tiene sentido buscar un
	matching bipartito si el grafo realmente lo es, y 2-Coloring
	además te da gratis los dos lados $L$/$R$ de la partición.
	\item Modelar restricciones tipo "estos dos elementos no pueden
	estar juntos" como aristas, y usar 2-Coloring para ver si existe
	una asignación válida en dos categorías.
\end{itemize}

\textbf{Complejidad:} $O(V + E)$ — un BFS/DFS normal con una
comparación de color extra por arista, revisando todas las
componentes del grafo (un grafo desconexo puede tener partes
bipartitas y partes no, o varias componentes bipartitas
independientes).

\begin{lstlisting}[style=cpp]
	struct DosColoreo {
		vector<vector<int>> adj;
		vector<int> color;   // -1 = sin colorear, 0 o 1 = los dos grupos
		int n;
		
		// CREATE: guarda el grafo.
		DosColoreo(vector<vector<int>>& adj_, int n_) {
			adj = adj_;
			n = n_;
			color.assign(n, -1);
		}                                          // O(1)
		
		// READ: true si el grafo completo es bipartito. Colorea todas
		// las componentes en el proceso (revisa cada nodo no coloreado).
		bool esBipartito() {
			for (int u = 0; u < n; u++)
			if (color[u] == -1 && !bfs(u))
			return false;
			return true;
		}                                          // O(V + E)
		
		bool bfs(int origen) {
			queue<int> q;
			q.push(origen);
			color[origen] = 0;
			
			while (!q.empty()) {
				int u = q.front(); q.pop();
				
				for (int v : adj[u]) {
					if (color[v] == -1) {
						color[v] = 1 - color[u];   // color opuesto
						q.push(v);
					} else if (color[v] == color[u]) {
						return false;   // mismo color en una arista
						// -> ciclo impar -> no bipartito
					}
				}
			}
			
			return true;
		}
	};
\end{lstlisting}

\textbf{Ejemplo de funcionamiento:}

\begin{lstlisting}[style=cpp]
	// grafo bipartito:  0 - 1     0 - 3
	//                    2 - 1     2 - 3
	// (es un ciclo de longitud 4, par -> bipartito)
	
	vector<vector<int>> adj(4);
	adj[0] = {1, 3}; adj[1] = {0, 2};
	adj[2] = {1, 3}; adj[3] = {0, 2};
	
	DosColoreo dc(adj, 4);
	cout << dc.esBipartito();
	// true
	
	// grafo NO bipartito: triangulo 0-1-2-0 (ciclo de longitud 3, impar)
	vector<vector<int>> adj2(3);
	adj2[0] = {1, 2}; adj2[1] = {0, 2}; adj2[2] = {0, 1};
	
	DosColoreo dc2(adj2, 3);
	cout << dc2.esBipartito();
	// false
\end{lstlisting}


\subsection{Bipartite Matching}

\textbf{Idea:} dado un grafo bipartito con lados $L$ y $R$, un
\textbf{matching} es un subconjunto de aristas donde ningún nodo se
repite (cada nodo de $L$ emparejado con \textbf{a lo más un} nodo de
$R$, y viceversa). El \textbf{matching máximo} se encuentra con el
algoritmo de Kuhn: para cada nodo de $L$ sin emparejar, intenta un
DFS buscando un \textbf{camino aumentante} —si llega a un nodo de $R$
libre, lo empareja directamente; si llega a uno ya emparejado,
intenta "desplazar" a su pareja actual (recursivamente, buscándole
otra opción) para liberar espacio.

\textbf{¿Por qué usarlo?} Es la herramienta estándar para cualquier
problema de "asignación óptima uno-a-uno" entre dos conjuntos
—trabajadores a tareas, estudiantes a proyectos, jugadores a
posiciones— donde cada asignación válida es una arista en un grafo
bipartito. La garantía clave del algoritmo (teorema de Berge): un
matching es máximo \textbf{si y solo si} no existe ningún camino
aumentante —por eso repetir "buscar camino aumentante desde cada nodo
libre de $L$" hasta que ya ninguno encuentre uno, termina
garantizadamente en el matching máximo.

\textbf{Casos de uso:}

\begin{itemize}
	\item Asignación máxima uno-a-uno entre dos conjuntos con
	compatibilidades dadas (trabajadores-tareas, alumnos-cursos).
	\item Como subrutina de problemas más generales: cobertura mínima
	de vértices en grafos bipartitos (teorema de König: tamaño del
	matching máximo = tamaño de la cobertura mínima de vértices).
	\item Problemas de "cuántos pares disjuntos puedo formar" cuando
	la compatibilidad se puede modelar como bipartita (ej. tareas con
	horarios que no chocan, visto como grafo).
\end{itemize}

\textbf{Complejidad:} $O(V \cdot E)$ con Kuhn simple (un DFS por
cada nodo de $L$, en el peor caso $O(E)$ cada uno) — suficiente para
$V \lesssim 1000$–$2000$. Para grafos más grandes existe
Hopcroft-Karp, $O(E\sqrt{V})$, fuera del alcance de esta sección.

\begin{lstlisting}[style=cpp]
	struct BipartiteMatching {
		vector<vector<int>> adj;   // adj[u] = vecinos en R, para cada
		// u en L
		vector<int> matchR;        // matchR[v] = nodo de L emparejado
		// con v en R, -1 si libre
		vector<bool> visitado;     // marca de la busqueda actual (se
		// reinicia en cada intento)
		int nL;
		
		// CREATE: guarda el grafo bipartito (nL nodos en L, adj[u]
		// lista los vecinos de u en R). nR = tamano del lado R.
		BipartiteMatching(vector<vector<int>>& adj_, int nL_, int nR) {
			adj = adj_;
			nL = nL_;
			matchR.assign(nR, -1);
		}                                          // O(1)
		
		// busca un camino aumentante desde u (nodo de L), intentando
		// desplazar parejas ya asignadas si hace falta.
		bool intentar(int u) {
			for (int v : adj[u]) {
				if (visitado[v]) continue;
				visitado[v] = true;
				
				// v esta libre, o su pareja actual puede reubicarse
				if (matchR[v] == -1 || intentar(matchR[v])) {
					matchR[v] = u;
					return true;
				}
			}
			return false;
		}
		
		// READ: calcula el matching maximo. Devuelve su tamano;
		// matchR queda con la asignacion final.
		int matchingMaximo() {
			int total = 0;
			
			for (int u = 0; u < nL; u++) {
				visitado.assign(matchR.size(), false);
				if (intentar(u)) total++;
			}
			
			return total;
		}                                          // O(V * E)
	};
\end{lstlisting}

\textbf{Ejemplo de funcionamiento:}

\begin{lstlisting}[style=cpp]
	// L = {0, 1, 2}, R = {0, 1, 2}
	// 0 - {0, 1}
	// 1 - {0}
	// 2 - {1, 2}
	
	vector<vector<int>> adj(3);
	adj[0] = {0, 1};
	adj[1] = {0};
	adj[2] = {1, 2};
	
	BipartiteMatching bm(adj, 3, 3);
	int maximo = bm.matchingMaximo();
	
	cout << maximo;
	// 3
	// posible asignacion: 1-0 (desplaza a 0), 0-1, 2-2
	// (0 originalmente probo emparejarse con 0 en R, pero al llegar 1
	//  que SOLO puede ir a 0, 0 se desplaza hacia 1 en R, liberando
	//  espacio para que 1 tome 0 en R)
\end{lstlisting}

\textbf{Nota:} el vector \texttt{visitado} se reinicia en \textbf{cada}
llamada a \texttt{matchingMaximo()} por cada nodo de $L$ —no es el
mismo \texttt{visitado} de un DFS normal que se llena una sola vez
para todo el grafo. Aquí marca "ya considerado \textbf{en este
	intento}", para no reprocesar el mismo nodo de $R$ dos veces dentro
de una sola búsqueda de camino aumentante, aunque sí se vuelve a
visitar en el siguiente intento (desde el próximo nodo de $L$).

\textbf{Regla práctica:} 2-Coloring y Bipartite Matching son
secuenciales: primero verificas (y separas en $L$/$R$) con
2-Coloring, y solo si el grafo realmente es bipartito tiene sentido
correr Matching sobre él —Kuhn asume esa estructura y no detecta por
sí solo si el grafo no lo es.


% =========================================================
% TREES
% =========================================================
\section{Trees}
\subsection{DFS en Árboles}

\textbf{Idea:} un árbol es simplemente un grafo conexo sin ciclos
—eso significa que un DFS sobre un árbol \textbf{nunca} necesita
marcar "visitado" con un arreglo separado: basta con no regresar por
la misma arista al \textbf{padre} inmediato, ya que no hay ningún
otro camino de vuelta (no hay ciclos que puedan reconectar el DFS con
un nodo ya procesado). Esa simplificación es lo que distingue DFS en
árboles del DFS de grafos generales, y es la base de casi todo el DP
sobre árboles: recorres en post-orden, y cada nodo combina los
resultados ya calculados de sus hijos.

\textbf{¿Por qué usarlo?} Es el patrón detrás de prácticamente
\textbf{todo} problema de árboles: tamaño de subárbol, profundidad,
diámetro, suma de valores por subárbol, DP sobre árboles en general.
Reconocer "esto es información que depende de mis hijos" es la señal
para post-orden DFS —calculas la respuesta de cada hijo primero, y el
nodo actual combina esos resultados al \textbf{salir} de la
recursión.

\textbf{Casos de uso:}

\begin{itemize}
	\item Tamaño de cada subárbol: $\texttt{tam}[u] = 1 + \sum
	\texttt{tam}[hijo]$.
	\item Profundidad/altura de cada nodo, o altura total del árbol.
	\item Suma o valor agregado de cada subárbol (útil para
	responder queries tipo "¿cuánto suma el subárbol de $u$?").
	\item Diámetro del árbol (la distancia más larga entre dos
	nodos cualquiera): con dos DFS —desde cualquier nodo encuentra el
	más lejano $a$, desde $a$ encuentra el más lejano $b$; la
	distancia $a$-$b$ es el diámetro (propiedad de árboles).
\end{itemize}

\textbf{Complejidad:} $O(n)$ — cada nodo se visita exactamente una
vez, cada arista se recorre exactamente una vez (a diferencia de un
grafo general, aquí no hace falta ni siquiera el arreglo
\texttt{visitado}, basta con excluir al padre).

\begin{lstlisting}[style=cpp]
	struct DFSArbol {
		vector<vector<int>> adj;
		vector<int> tam, profundidad;
		int n;
		
		// CREATE: guarda el arbol (n nodos, n-1 aristas).
		DFSArbol(vector<vector<int>>& adj_, int n_) {
			adj = adj_;
			n = n_;
			tam.assign(n, 0);
			profundidad.assign(n, 0);
		}                                          // O(1)
		
		// READ: corre DFS desde raiz, llenando tam[] y profundidad[]
		// de una sola pasada.
		void dfs(int u, int padre, int d) {
			profundidad[u] = d;
			tam[u] = 1;
			
			for (int v : adj[u]) {
				if (v == padre) continue;   // unica exclusion necesaria,
				// no hace falta 'visitado'
				
				dfs(v, u, d + 1);
				tam[u] += tam[v];   // combina resultados de los hijos
				// DESPUES de que ya terminaron
				// (post-orden)
			}
		}                                          // O(n)
	};
\end{lstlisting}

\textbf{Ejemplo de funcionamiento:}

\begin{lstlisting}[style=cpp]
	// arbol:      0
	//            / \
	//           1   2
	//          / \
	//         3   4
	
	vector<vector<int>> adj(5);
	adj[0] = {1, 2};
	adj[1] = {0, 3, 4};
	adj[2] = {0};
	adj[3] = {1};
	adj[4] = {1};
	
	DFSArbol da(adj, 5);
	da.dfs(0, -1, 0);
	
	cout << da.tam[0];
	// 5  (todo el arbol)
	
	cout << da.tam[1];
	// 3  (1, 3, 4)
	
	cout << da.profundidad[4];
	// 2
\end{lstlisting}


\subsection{BFS en Árboles}

\textbf{Idea:} el mismo principio de "sin ciclos, no hace falta
\texttt{visitado} extra, solo excluir al padre" aplica igual a BFS
sobre árboles —con la ventaja de que BFS recorre \textbf{por
	niveles}, así que es la herramienta natural cuando lo que se pide
depende explícitamente de la \textbf{profundidad/nivel} de cada nodo
en vez de depender de sus subárboles.

\textbf{¿Por qué usarlo?} DFS y BFS en árboles calculan las mismas
distancias/profundidades ($O(n)$ ambos), pero BFS procesa los nodos
en \textbf{orden creciente de profundidad}, lo cual es justo lo que
necesitas cuando el problema pide algo "nivel por nivel" (ej.
"imprime el árbol por niveles", "suma de valores en cada
profundidad") — con DFS tendrías que agrupar por profundidad
manualmente después; con BFS, ya vienen agrupados por construcción.

\textbf{Casos de uso:}

\begin{itemize}
	\item Recorrido por niveles ("level order traversal"): imprimir o
	procesar el árbol nivel por nivel.
	\item Distancia mínima desde la raíz a cada nodo (aunque en un
	árbol, DFS también la da correctamente — la ventaja de BFS es
	solo el orden de procesamiento por capas).
	\item Encontrar el nodo más profundo, o todos los nodos a una
	profundidad exacta $k$.
	\item Como una de las dos pasadas del cálculo de diámetro de
	árbol (BFS desde cualquier nodo para hallar el extremo más
	lejano, igual que con DFS — aquí ambos sirven).
\end{itemize}

\textbf{Complejidad:} $O(n)$ — cada nodo se encola y desencola una
sola vez.

\begin{lstlisting}[style=cpp]
	struct BFSArbol {
		vector<vector<int>> adj;
		vector<int> profundidad, padre;
		vector<vector<int>> porNivel;   // porNivel[d] = nodos a
		// profundidad d
		int n;
		
		// CREATE: guarda el arbol.
		BFSArbol(vector<vector<int>>& adj_, int n_) {
			adj = adj_;
			n = n_;
		}                                          // O(1)
		
		// READ: corre BFS desde raiz, llenando profundidad[], padre[]
		// y agrupando los nodos por nivel.
		void bfs(int raiz) {
			profundidad.assign(n, -1);
			padre.assign(n, -1);
			porNivel.clear();
			
			queue<int> q;
			q.push(raiz);
			profundidad[raiz] = 0;
			
			while (!q.empty()) {
				int u = q.front(); q.pop();
				
				if ((int)porNivel.size() <= profundidad[u])
				porNivel.push_back({});
				porNivel[profundidad[u]].push_back(u);
				
				for (int v : adj[u]) {
					if (v == padre[u]) continue;   // igual que en DFS,
					// basta excluir al
					// padre
					
					padre[v] = u;
					profundidad[v] = profundidad[u] + 1;
					q.push(v);
				}
			}
		}                                          // O(n)
	};
\end{lstlisting}

\textbf{Ejemplo de funcionamiento:}

\begin{lstlisting}[style=cpp]
	// arbol:      0
	//            / \
	//           1   2
	//          / \
	//         3   4
	
	vector<vector<int>> adj(5);
	adj[0] = {1, 2};
	adj[1] = {0, 3, 4};
	adj[2] = {0};
	adj[3] = {1};
	adj[4] = {1};
	
	BFSArbol ba(adj, 5);
	ba.bfs(0);
	
	cout << ba.porNivel.size();
	// 3   (niveles 0, 1, 2)
	
	for (int u : ba.porNivel[1]) cout << u << " ";
	// 1 2   (nivel 1: hijos directos de la raiz)
	
	for (int u : ba.porNivel[2]) cout << u << " ";
	// 3 4   (nivel 2: nietos de la raiz)
\end{lstlisting}

\textbf{Regla práctica:} en árboles, tanto DFS como BFS corren en
$O(n)$ y ninguno necesita \texttt{visitado} —basta con excluir al
padre en cada paso, porque la ausencia de ciclos lo garantiza. Elige
DFS cuando lo que calculas depende de \textbf{combinar resultados de
	subárboles} (post-orden: tamaños, sumas, diámetro vía DP); elige BFS
cuando lo que calculas depende de \textbf{agrupar por
	nivel/profundidad} explícitamente (recorridos por niveles, "todos
los nodos a distancia $k$ de la raíz").

\subsection{Diámetro de un Árbol}

\textbf{Idea:} el diámetro de un árbol es la distancia (en número de
aristas) entre el par de nodos \textbf{más lejano} entre sí. Se
calcula con una propiedad no obvia pero muy útil de los árboles:
\textbf{desde cualquier nodo $x$, el nodo más lejano $a$ es siempre
	un extremo del diámetro}. Eso da el algoritmo de \textbf{dos
	pasadas}: (1) BFS/DFS desde un nodo cualquiera para encontrar el nodo
más lejano $a$; (2) BFS/DFS desde $a$ para encontrar el nodo más
lejano $b$ —la distancia $a$-$b$ \textbf{es} el diámetro. Nótese que
$x$ inicial puede ser literalmente cualquier nodo, no hace falta
adivinar ni probar todos.

\textbf{¿Por qué usarlo?} La alternativa ingenua —correr BFS desde
\textbf{cada} nodo y quedarte con la distancia máxima global— cuesta
$O(n^2)$. El truco de las dos pasadas baja eso a $O(n)$ aprovechando
la estructura de árbol (sin ciclos): la garantía de que "el más
lejano desde cualquier punto es un extremo del diámetro" \textbf{no}
es cierta en grafos generales, solo en árboles.

\textbf{Casos de uso:}

\begin{itemize}
	\item Distancia máxima entre cualquier par de nodos de un árbol.
	\item Encontrar el par de nodos que realiza esa distancia máxima
	(reconstruyendo el camino de $a$ a $b$ con \texttt{padre[]},
	igual que en BFS en Árboles).
	\item Como subrutina de Centro de un Árbol (siguiente sección):
	el centro es literalmente el punto medio del camino del
	diámetro.
	\item Problemas de "radio mínimo" o "tiempo máximo para que algo
	se propague desde el peor nodo posible" —el diámetro acota estos
	valores.
\end{itemize}

\textbf{Complejidad:} $O(n)$ — dos BFS/DFS completos, cada uno
$O(n)$.

\begin{lstlisting}[style=cpp]
	struct DiametroArbol {
		vector<vector<int>> adj;
		int n;
		
		// CREATE: guarda el arbol.
		DiametroArbol(vector<vector<int>>& adj_, int n_) {
			adj = adj_;
			n = n_;
		}                                          // O(1)
		
		// BFS auxiliar: devuelve {nodo mas lejano, distancia} desde src.
		pair<int,int> bfsMasLejano(int src) {
			vector<int> dist(n, -1);
			queue<int> q;
			q.push(src);
			dist[src] = 0;
			
			int mejorNodo = src, mejorDist = 0;
			
			while (!q.empty()) {
				int u = q.front(); q.pop();
				
				if (dist[u] > mejorDist) {
					mejorDist = dist[u];
					mejorNodo = u;
				}
				
				for (int v : adj[u]) {
					if (dist[v] != -1) continue;
					dist[v] = dist[u] + 1;
					q.push(v);
				}
			}
			
			return {mejorNodo, mejorDist};
		}                                          // O(n)
		
		// READ: diametro del arbol (dos pasadas de BFS).
		int diametro() {
			auto [a, _] = bfsMasLejano(0);       // paso 1: cualquier
			// nodo -> extremo a
			auto [b, dist] = bfsMasLejano(a);    // paso 2: a -> extremo b
			
			return dist;
		}                                          // O(n)
	};
\end{lstlisting}

\textbf{Ejemplo de funcionamiento:}

\begin{lstlisting}[style=cpp]
	// arbol:      0
	//            / \
	//           1   2
	//          /
	//         3
	//        /
	//       4
	
	vector<vector<int>> adj(5);
	adj[0] = {1, 2};
	adj[1] = {0, 3};
	adj[2] = {0};
	adj[3] = {1, 4};
	adj[4] = {3};
	
	DiametroArbol da(adj, 5);
	
	cout << da.diametro();
	// 4   (camino 4-3-1-0-2, 4 aristas)
\end{lstlisting}

\textbf{Nota:} existe una alternativa de \textbf{una sola pasada} con
DP (post-orden, como en DFS en Árboles): para cada nodo, calcular las
dos ramas más largas de sus hijos y sumar; el diámetro es el máximo
de esa suma sobre todos los nodos. Es igual de $O(n)$ pero evita el
segundo BFS —útil si ya estás haciendo otro DP sobre el árbol y
quieres aprovechar el mismo recorrido. El método de dos pasadas es
preferible cuando además necesitas identificar \textbf{qué} par de
nodos $(a,b)$ realiza el diámetro, no solo su longitud.


\subsection{Centro de un Árbol}

\textbf{Idea:} el centro de un árbol es el nodo (o los dos nodos)
que \textbf{minimiza la distancia máxima} a cualquier otro nodo del
árbol —el punto "más equilibrado". Un árbol tiene \textbf{uno o dos}
centros nunca más: si el diámetro tiene longitud \textbf{par}, hay un
único centro (el punto medio exacto del camino del diámetro); si es
\textbf{impar}, hay dos centros adyacentes (los dos nodos de en medio
del camino, que no tiene un punto medio exacto).

\textbf{¿Por qué usarlo?} Hay dos formas equivalentes de encontrarlo,
ambas $O(n)$: (1) \textbf{vía el diámetro} —calculas el camino del
diámetro con Diámetro de un Árbol, y el/los nodo(s) en la posición
central de ese camino son el centro; (2) \textbf{"pelado de hojas"}
(\textit{leaf trimming})—quitas todas las hojas actuales del árbol
capa por capa (como un BFS multi-fuente desde las hojas hacia
adentro), y lo que sobrevive al final (1 o 2 nodos) es el centro. La
segunda es más directa de programar cuando ya tienes BFS
multi-fuente a mano (mismo patrón que Shortest Path desde Múltiples
Orígenes, pero "encogiendo" el árbol en vez de calculando distancias).

\textbf{Casos de uso:}

\begin{itemize}
	\item Elegir la raíz "óptima" de un árbol quiere minimizar la
	altura resultante (rerootear en el centro da la altura mínima
	posible).
	\item Problemas de "ubica un servidor/recurso central" que
	minimice la distancia máxima a cualquier nodo cliente.
	\item Como building block de \textbf{centroid decomposition}
	(fuera del alcance de esta sección, pero el centro es el primer
	paso de esa técnica).
\end{itemize}

\textbf{Complejidad:} $O(n)$ con cualquiera de los dos métodos.

\begin{lstlisting}[style=cpp]
	struct CentroArbol {
		vector<vector<int>> adj;
		vector<int> grado;
		int n;
		
		// CREATE: guarda el arbol y el grado inicial de cada nodo
		// (numero de vecinos).
		CentroArbol(vector<vector<int>>& adj_, int n_) {
			adj = adj_;
			n = n_;
			grado.assign(n, 0);
			
			for (int u = 0; u < n; u++)
			grado[u] = adj[u].size();
		}                                          // O(n)
		
		// READ: encuentra el/los centro(s) "pelando" hojas capa por
		// capa hasta que quedan 1 o 2 nodos.
		vector<int> centro() {
			if (n == 1) return {0};
			
			vector<int> g = grado;   // copia para no mutar el original
			queue<int> hojas;
			
			for (int u = 0; u < n; u++)
			if (g[u] == 1)
			hojas.push(u);
			
			int restantes = n;
			vector<int> capaActual;
			
			while (restantes > 2) {
				int cantidadHojas = hojas.size();
				restantes -= cantidadHojas;
				
				for (int i = 0; i < cantidadHojas; i++) {
					int u = hojas.front(); hojas.pop();
					
					for (int v : adj[u]) {
						g[v]--;
						if (g[v] == 1)
						hojas.push(v);
					}
				}
			}
			
			// lo que queda en la cola (1 o 2 nodos) es el centro
			vector<int> resultado;
			while (!hojas.empty()) {
				resultado.push_back(hojas.front());
				hojas.pop();
			}
			return resultado;
		}                                          // O(n)
	};
\end{lstlisting}

\textbf{Ejemplo de funcionamiento:}

\begin{lstlisting}[style=cpp]
	// arbol (cadena):  0 - 1 - 2 - 3 - 4
	// diametro = 4 (par) -> un solo centro, el punto medio exacto
	
	vector<vector<int>> adj(5);
	adj[0] = {1}; adj[1] = {0, 2}; adj[2] = {1, 3};
	adj[3] = {2, 4}; adj[4] = {3};
	
	CentroArbol ca(adj, 5);
	for (int u : ca.centro()) cout << u << " ";
	// 2   (unico centro, justo en medio de la cadena de 5 nodos)
	
	// arbol (cadena):  0 - 1 - 2 - 3
	// diametro = 3 (impar) -> dos centros adyacentes
	
	vector<vector<int>> adj2(4);
	adj2[0] = {1}; adj2[1] = {0, 2};
	adj2[2] = {1, 3}; adj2[3] = {2};
	
	CentroArbol ca2(adj2, 4);
	for (int u : ca2.centro()) cout << u << " ";
	// 1 2   (los dos nodos centrales)
\end{lstlisting}

\textbf{Regla práctica:} el pelado de hojas y el método vía diámetro
\textbf{siempre} coinciden en el resultado —son dos formas de llegar
al mismo centro. Usa pelado de hojas si el problema ya viene en
términos de "ir quitando hojas" o si prefieres no calcular
explícitamente el camino del diámetro; usa el método vía diámetro si
de todas formas ya necesitas el diámetro para otra parte del
problema, ya que entonces el centro sale casi gratis indexando al
medio del camino ya reconstruido.

\subsection{Combinatoria en Árboles Ordenados}

\textbf{Idea:} dado un árbol enraizado de $n$ nodos, un
\textbf{orden topológico} es cualquier permutación de sus nodos donde
cada nodo aparece \textbf{antes} que todos sus descendientes (por
ejemplo, "en qué orden pudieron haberse insertado" si el árbol se
construyó agregando nodos de a uno, siempre como hijo de algo ya
presente). El número total de órdenes topológicos válidos tiene una
fórmula cerrada muy elegante:
$$\frac{n!}{\prod_{v} \texttt{tam}[v]}$$
donde $\texttt{tam}[v]$ es el tamaño del subárbol de $v$ (incluyéndose
a sí mismo). Se calcula con un solo DFS post-orden (igual que DFS en
Árboles) para obtener los tamaños, y luego un producto/cociente
modular sobre todos los nodos.

\textbf{¿Por qué usarlo?} La fórmula evita tener que \textbf{contar}
las permutaciones una por una (imposible, son hasta $n!$) o hacer DP
combinatoria compleja: reduce el problema entero a "un DFS de
tamaños + un producto modular", el mismo patrón de precómputo que ya
usas para factoriales e inversos modulares. Es la herramienta directa
cuando un problema pregunta "¿de cuántas formas se pudo haber
construido/insertado/etiquetado este árbol respetando el orden
padre-antes-que-hijo?".

\textbf{Casos de uso:}

\begin{itemize}
	\item Contar de cuántas formas se pudo insertar $n$ elementos en
	una estructura tipo árbol (heap, trie, etc.) para terminar con la
	forma dada.
	\item Contar etiquetados válidos $1..n$ de los nodos de un árbol
	tales que cada nodo tenga una etiqueta menor que todos sus hijos
	("árbol como heap mínimo").
	\item Como building block de problemas más grandes que combinan
	esta cuenta con otras restricciones (ej. multiplicar por
	factoriales adicionales si además el orden \textbf{entre}
	hermanos importa de alguna forma extra).
\end{itemize}

\textbf{Complejidad:} $O(n)$ para el DFS de tamaños, más
$O(n \log MOD)$ si se calculan inversos modulares individuales con
exponenciación rápida (o $O(n)$ total si se precomputan inversos de
factoriales una sola vez).

\begin{lstlisting}[style=cpp]
	const long long MOD = 1e9 + 7;
	
	struct CombinatoriaArbol {
		vector<vector<int>> adj;
		vector<long long> tam;
		int n;
		
		// CREATE: guarda el arbol enraizado.
		CombinatoriaArbol(vector<vector<int>>& adj_, int n_) {
			adj = adj_;
			n = n_;
			tam.assign(n, 0);
		}                                          // O(1)
		
		long long potMod(long long base, long long exp, long long mod) {
			long long res = 1;
			base %= mod;
			while (exp > 0) {
				if (exp & 1) res = res * base % mod;
				base = base * base % mod;
				exp >>= 1;
			}
			return res;
		}
		
		long long inversoModular(long long a) {
			return potMod(a, MOD - 2, MOD);   // Fermat, MOD debe ser primo
		}
		
		void dfsTam(int u, int padre) {
			tam[u] = 1;
			for (int v : adj[u]) {
				if (v == padre) continue;
				dfsTam(v, u);
				tam[u] += tam[v];
			}
		}                                          // O(n)
		
		// READ: numero de ordenes topologicos validos del arbol
		// enraizado en 'raiz', modulo MOD.
		long long contarOrdenes(int raiz) {
			dfsTam(raiz, -1);
			
			long long factorialN = 1;
			for (int i = 2; i <= n; i++)
			factorialN = factorialN * i % MOD;
			
			long long denominador = 1;
			for (int v = 0; v < n; v++)
			denominador = denominador * tam[v] % MOD;
			
			return factorialN * inversoModular(denominador) % MOD;
		}                                          // O(n log MOD)
	};
\end{lstlisting}

\textbf{Ejemplo de funcionamiento:}

\begin{lstlisting}[style=cpp]
	// arbol:      0
	//            / \
	//           1   2
	//          /
	//         3
	
	vector<vector<int>> adj(4);
	adj[0] = {1, 2};
	adj[1] = {0, 3};
	adj[2] = {0};
	adj[3] = {1};
	
	CombinatoriaArbol ca(adj, 4);
	
	cout << ca.contarOrdenes(0);
	// 3
	// las 3 inserciones validas (etiquetando por orden de insercion):
	// 0,1,2,3 -- 0,1,3,2 -- 0,2,1,3
	// (0 siempre primero; 2 puede ir en cualquier posicion relativa a
	//  la rama 1-3, pero 3 siempre despues de 1)
\end{lstlisting}

\textbf{Nota:} esta fórmula asume que el orden \textbf{entre}
hermanos no distingue resultados distintos (a diferencia de contar
árboles \textbf{ordenados} donde la secuencia exacta de hijos
importa). Si el problema sí distingue el orden de los hijos entre sí
como parte de la respuesta, hay que multiplicar además por
$\prod_v (\text{cantidad de hijos de } v)!$ — ese factor extra cubre
las formas de \textbf{ordenar} los hijos de cada nodo entre sí,
encima de las formas de intercalarlos que ya cuenta la fórmula base.


\subsection{Tree DP}

\textbf{Idea:} Tree DP es el patrón general de DP sobre árboles con
DFS post-orden: cada nodo $u$ tiene uno o más "estados" (ej.
\texttt{dp[u][0]} = mejor respuesta \textbf{sin} usar $u$,
\texttt{dp[u][1]} = mejor respuesta \textbf{usando} $u$), y el valor
de cada estado se calcula \textbf{combinando} los estados ya
resueltos de los hijos —exactamente el mismo esqueleto de DFS en
Árboles (post-orden, sin \texttt{visitado}, solo excluir al padre),
pero acumulando una decisión óptima en vez de solo un tamaño o una
profundidad.

\textbf{¿Por qué usarlo?} Es la generalización de todo lo visto en
la sección de Trees: tamaño de subárbol, diámetro (una pasada), y
Combinatoria en Árboles Ordenados son casos particulares de "Tree
DP" donde el estado es simple (un solo número). El patrón completo
—múltiples estados por nodo, combinados con $\max$/$\min$/suma según
el problema— es lo que resuelve la mayoría de problemas de
optimización sobre árboles (no solo conteo o medición).

\textbf{Casos de uso:}

\begin{itemize}
	\item \textbf{Conjunto independiente máximo} (máxima suma de
	pesos de nodos, sin elegir dos adyacentes): dos estados por nodo,
	clásico ejemplo de esta sección.
	\item \textbf{Vertex cover mínimo} en árboles: variante muy
	similar, cambiando la condición de combinación entre estados.
	\item \textbf{Coloreo/asignación óptima} de nodos con
	restricciones locales entre padre e hijo (cada estado extra
	representa una opción de "color" o "categoría" del nodo).
	\item Cualquier problema donde la respuesta óptima de un subárbol
	dependa de una decisión binaria o categórica tomada en su raíz.
\end{itemize}

\textbf{Complejidad:} $O(n \cdot k)$, donde $k$ es el número de
estados por nodo (para el clásico de 2 estados, $O(n)$ liso) — cada
nodo combina los estados ya resueltos de sus hijos con trabajo
constante (o $O(k)$) por hijo.

\begin{lstlisting}[style=cpp]
	struct ConjuntoIndependienteMax {
		vector<vector<int>> adj;
		vector<int> peso;
		vector<long long> dpSin, dpCon;   // dpSin[u] = mejor sin usar u
		// dpCon[u] = mejor usando u
		int n;
		
		// CREATE: guarda el arbol y el peso de cada nodo.
		ConjuntoIndependienteMax(vector<vector<int>>& adj_,
		vector<int>& peso_, int n_) {
			adj = adj_;
			peso = peso_;
			n = n_;
			dpSin.assign(n, 0);
			dpCon.assign(n, 0);
		}                                          // O(1)
		
		// READ: corre el DFS post-orden que llena dpSin[]/dpCon[] para
		// todo el arbol enraizado en 'raiz'.
		void dfs(int u, int padre) {
			dpCon[u] = peso[u];   // usar u obliga a NO usar a sus hijos
			dpSin[u] = 0;         // no usar u: cada hijo elige libremente
			
			for (int v : adj[u]) {
				if (v == padre) continue;
				
				dfs(v, u);
				
				dpCon[u] += dpSin[v];              // si uso u, el hijo
				// v NO puede usarse
				dpSin[u] += max(dpSin[v], dpCon[v]); // si no uso u, el
				// hijo v elige lo
				// que le convenga
			}
		}                                          // O(n)
		
		// READ: respuesta final del arbol enraizado en 'raiz'.
		long long resolver(int raiz) {
			dfs(raiz, -1);
			return max(dpSin[raiz], dpCon[raiz]);
		}                                          // O(1) tras dfs()
	};
\end{lstlisting}

\textbf{Ejemplo de funcionamiento:}

\begin{lstlisting}[style=cpp]
	// arbol:        0(peso=3)
	//              /          \
	//        1(peso=5)      2(peso=4)
	//            |
	//        3(peso=6)
	
	vector<vector<int>> adj(4);
	adj[0] = {1, 2};
	adj[1] = {0, 3};
	adj[2] = {0};
	adj[3] = {1};
	
	vector<int> peso = {3, 5, 4, 6};
	
	ConjuntoIndependienteMax cim(adj, peso, 4);
	
	cout << cim.resolver(0);
	// 10
	// mejor conjunto: {2, 3} (pesos 4+6=10), no adyacentes entre si
\end{lstlisting}

\textbf{Regla práctica:} cualquier problema de árboles que pida
"óptimo global sujeto a una restricción entre padre e hijos" es Tree
DP: define qué información necesitas de cada hijo (usualmente 1-2
estados), escribe cómo se combina en el nodo actual, y el DFS
post-orden ya está —es el mismo recorrido que DFS en Árboles, solo
con más de un valor acumulándose por nodo.

\subsection{Rerooting}

\textbf{Idea:} Tree DP normal (ver Tree DP) calcula, para \textbf{una}
raíz fija, la respuesta de cada subárbol en una sola pasada
post-orden. Rerooting extiende eso a "calcular la respuesta como si
\textbf{cada} nodo fuera la raíz", sin pagar $O(n)$ por cada raíz
(lo cual daría $O(n^2)$ total). Se logra con \textbf{dos} DFS: (1) un
post-orden normal que calcula \texttt{down[u]} = la respuesta del
subárbol de $u$ asumiendo la raíz original; (2) un pre-orden que
propaga hacia abajo \texttt{up[u]} = "la contribución de todo lo que
\textbf{no} está en el subárbol de $u$" (el resto del árbol, visto
desde el padre de $u$) — combinando \texttt{down} y \texttt{up} en
cada nodo se obtiene la respuesta con ese nodo como raíz, en tiempo
total $O(n)$.

\textbf{¿Por qué usarlo?} Cualquier problema que pida "calcula X para
cada nodo, donde X depende de \textbf{todo} el árbol (no solo del
subárbol)" —si $X$ solo dependiera del subárbol, un DFS post-orden ya
bastaría. La señal típica es el enunciado diciendo "para cada nodo,
si fuera la raíz..." o "la respuesta considerando el árbol completo,
no solo hacia abajo".

\textbf{Casos de uso:}

\begin{itemize}
    \item Suma de distancias desde cada nodo a \textbf{todos} los
    demás nodos del árbol (el ejemplo clásico, abajo).
    \item Altura/profundidad máxima del árbol si se enraizara en cada
    nodo distinto.
    \item Tamaño "visto hacia afuera": cuántos nodos quedarían fuera
    del subárbol de $u$ si $u$ no fuera la raíz.
    \item Cualquier variante de Tree DP (conjunto independiente,
    diámetro, etc.) donde se pide la respuesta \textbf{por cada
    posible raíz} en vez de para una raíz fija.
\end{itemize}

\textbf{Complejidad:} $O(n)$ — dos DFS completos, cada uno $O(n)$
(a diferencia de recalcular Tree DP $n$ veces, que sería $O(n^2)$).

\begin{lstlisting}[style=cpp]
struct Rerooting {
    vector<vector<int>> adj;
    vector<long long> tam, down, up, respuesta;
    int n;

    // CREATE: guarda el arbol.
    Rerooting(vector<vector<int>>& adj_, int n_) {
        adj = adj_;
        n = n_;
        tam.assign(n, 0);
        down.assign(n, 0);
        up.assign(n, 0);
        respuesta.assign(n, 0);
    }                                          // O(1)

    // PASADA 1 (post-orden): down[u] = suma de distancias desde u
    // hacia TODO su subarbol, asumiendo raiz = 0. tam[u] = tamano
    // del subarbol de u.
    void dfs1(int u, int padre) {
        tam[u] = 1;
        down[u] = 0;

        for (int v : adj[u]) {
            if (v == padre) continue;
            dfs1(v, u);
            tam[u] += tam[v];
            down[u] += down[v] + tam[v];   // cada nodo del subarbol
                                             // de v queda 1 arista
                                             // mas lejos de u que de v
        }
    }                                          // O(n)

    // PASADA 2 (pre-orden): respuesta[u] = suma de distancias desde
    // u a TODOS los nodos del arbol completo. up[u] = contribucion
    // de "todo lo que no es el subarbol de u", vista desde u.
    void dfs2(int u, int padre) {
        respuesta[u] = down[u] + up[u];

        for (int v : adj[u]) {
            if (v == padre) continue;

            // al mover la raiz de u a v: todo lo que estaba FUERA
            // del subarbol de v queda 1 arista MAS lejos; lo que ya
            // estaba en el subarbol de v queda 1 arista MAS cerca --
            // por eso no se puede reusar 'down[v]' directamente, hay
            // que restar su propia contribucion antes de invertir
            // el signo:
            up[v] = (respuesta[u] - (down[v] + tam[v]))
                    + (n - tam[v]);

            dfs2(v, u);
        }
    }                                          // O(n)

    // READ: respuesta[u] para cada u, tras correr ambas pasadas
    // desde cualquier raiz inicial (0 por defecto).
    vector<long long> resolver() {
        dfs1(0, -1);
        up[0] = 0;
        dfs2(0, -1);
        return respuesta;
    }                                          // O(n)
};

//MODIFICACIONES Y COMENTARIOS CON LOS USOS

// 1. Suma de distancias a todos los nodos (el caso ya implementado
//    arriba)
// down[u] = suma de distancias DENTRO del subarbol de u
// up[u]   = suma de distancias hacia todo lo que esta AFUERA
// respuesta[u] = down[u] + up[u]

// 2. Tamano "hacia afuera" (cuantos nodos quedarian fuera del
//    subarbol de u si u no fuera la raiz)
// tamAfuera[v] = n - tam[v];
// -- no necesita ni siquiera up[], es directo desde tam[] de dfs1

// 3. Altura maxima del arbol enraizado en cada nodo
// (combinacion NO invertible: max no tiene "resta" como la suma,
//  hace falta guardar el MEJOR y el SEGUNDO MEJOR hijo)
// vector<long long> downMax1, downMax2, deQuienEsMax1;
// void dfs1_altura(int u, int padre) {
//     downMax1[u] = downMax2[u] = 0;
//     for (int v : adj[u]) {
//         if (v == padre) continue;
//         dfs1_altura(v, u);
//         long long h = downMax1[v] + 1;
//         if (h > downMax1[u]) {
//             downMax2[u] = downMax1[u];
//             downMax1[u] = h;
//             deQuienEsMax1[u] = v;
//         } else if (h > downMax2[u]) {
//             downMax2[u] = h;
//         }
//     }
// }
// void dfs2_altura(int u, int padre) {
//     for (int v : adj[u]) {
//         if (v == padre) continue;
//         // si v es quien da la altura maxima de u, hay que
//         // pasarle la SEGUNDA maxima (para no contarse a si mismo)
//         long long mejorSinV = (deQuienEsMax1[u] == v)
//                                ? downMax2[u] : downMax1[u];
//         up[v] = max(up[u], mejorSinV) + 1;
//         dfs2_altura(v, u);
//     }
// }

// 4. Diametro del arbol usando rerooting (alternativa a Diametro de
//    un Arbol con dos BFS)
// respuesta[u] = downMax1[u] combinado con up[u] de la variante 3;
// diametroGlobal = *max_element(respuesta.begin(), respuesta.end());
// -- mismo costo O(n), util si ya se necesita la respuesta POR NODO
//    y no solo el numero global

// 5. Conteo de "nodos alcanzables en <= k pasos desde cada nodo"
// (variante de down/up acumulando CONTEOS por profundidad, no un
//  solo numero)
// vector<vector<long long>> down(n), up(n);  // down[u][d], up[u][d]
// -- mismo esqueleto de dos pasadas; down[u] y up[u] pasan a ser
//    arreglos indexados por distancia d en vez de escalares

// 6. Regla general para saber si hace falta "top-2" o basta un solo
//    valor:
// - operacion INVERTIBLE (suma, XOR, producto sin ceros)
//   -> un solo down[u] basta, up[v] se obtiene RESTANDO la
//      contribucion de v (como en el caso 1)
// - operacion NO invertible (max, min, gcd)
//   -> hace falta guardar el MEJOR y el SEGUNDO MEJOR para poder
//      "quitar" la rama de v sin perder informacion (como en el
//      caso 3)
\end{lstlisting}

\textbf{Ejemplo de funcionamiento:}

\begin{lstlisting}[style=cpp]
// arbol:      0
//            / \
//           1   2
//          /
//         3

vector<vector<int>> adj(4);
adj[0] = {1, 2};
adj[1] = {0, 3};
adj[2] = {0};
adj[3] = {1};

Rerooting rr(adj, 4);
auto resp = rr.resolver();

cout << resp[0];
// 4   (dist a 1=1, a 2=1, a 3=2 -> 1+1+2=4)
\end{lstlisting}



% =========================================================
% LCA
% =========================================================
\section{Lowest Common Ancestor (LCA)}
\subsection{Binary Lifting}

\textbf{Idea:} Binary Lifting precomputa, para cada nodo $u$, quién
es su ancestro a $2^k$ pasos hacia arriba, para todo $k$ de $0$ a
$\log n$ (\texttt{up[k][u]}). Con eso, "subir $d$ pasos desde $u$" se
descompone en la representación binaria de $d$ (igual que
exponenciación rápida, pero saltando por el árbol en vez de
multiplicando). El LCA de $u$ y $v$ se calcula en dos fases: (1)
nivelar al más profundo de los dos hasta que ambos queden a la misma
profundidad, usando esos saltos binarios; (2) si no son ya el mismo
nodo, subir a ambos \textbf{simultáneamente} con los saltos más
grandes posibles \textbf{que no los junten todavía} —al terminar,
ambos quedan justo un paso debajo del LCA, que es simplemente su
padre común.

\textbf{¿Por qué usarlo?} Es la alternativa a LCA vía Euler Tour +
RMQ (ver la sección de RMQ) cuando el árbol \textbf{no} es
completamente estático o cuando de todas formas ya necesitas
"ancestro a $k$ pasos de distancia" como operación aparte (útil por
sí sola en muchos problemas, no solo para LCA). La query de LCA es
$O(\log n)$ en vez de $O(1)$ como con RMQ, pero Binary Lifting es más
simple de razonar/depurar (sin Euler Tour ni Sparse Table por
separado) y generaliza directo a "ancestro a $k$ pasos" y "distancia
entre dos nodos" con el mismo precómputo.

\textbf{Casos de uso:}

\begin{itemize}
    \item LCA de dos nodos, con muchas queries sobre el mismo árbol
    fijo.
    \item Ancestro de $u$ a exactamente $k$ pasos hacia arriba —
    operación independiente del LCA, útil en problemas de
    "salta $k$ generaciones".
    \item Distancia entre dos nodos de un árbol:
    $\texttt{profundidad}[u] + \texttt{profundidad}[v] -
    2\cdot\texttt{profundidad}[\texttt{lca}(u,v)]$, una vez que ya
    tienes LCA y profundidades.
    \item Verificar si $u$ es ancestro de $v$: comparando si
    $\texttt{lca}(u,v) == u$.
\end{itemize}

\textbf{Complejidad:} $O(n \log n)$ de precómputo (build), $O(\log n)$
por query de LCA o de "ancestro a $k$ pasos".

\begin{lstlisting}[style=cpp]
struct BinaryLifting {
    vector<vector<int>> adj;
    vector<vector<int>> up;      // up[k][u] = ancestro de u a 2^k pasos
    vector<int> profundidad;
    int n, LOG;

    // CREATE: construye el arbol (DFS para profundidad y padre
    // directo) y la tabla de saltos binarios.
    BinaryLifting(vector<vector<int>>& adj_, int raiz, int n_) {
        adj = adj_;
        n = n_;
        LOG = 1;
        while ((1 << LOG) < n) LOG++;
        LOG++;

        up.assign(LOG, vector<int>(n, -1));
        profundidad.assign(n, 0);

        dfs(raiz, -1, 0);
        build();
    }                                          // O(n log n)

    void dfs(int u, int padre, int d) {
        up[0][u] = padre;
        profundidad[u] = d;

        for (int v : adj[u])
            if (v != padre)
                dfs(v, u, d + 1);
    }

    void build() {
        for (int k = 1; k < LOG; k++)
            for (int u = 0; u < n; u++)
                if (up[k - 1][u] != -1)
                    up[k][u] = up[k - 1][up[k - 1][u]];
                // si up[k-1][u] es -1 (se salio de la raiz), up[k][u]
                // se queda en -1 por el assign inicial
    }

    // READ: ancestro de u a exactamente k pasos hacia arriba, -1 si
    // se sale del arbol antes de completar los k pasos.
    int kEsimoAncestro(int u, int k) {
        for (int i = 0; i < LOG && u != -1; i++)
            if (k & (1 << i))
                u = up[i][u];
        return u;
    }                                          // O(log n)

    // READ: LCA de u y v.
    int query(int u, int v) {
        if (profundidad[u] < profundidad[v]) swap(u, v);

        // fase 1: nivela u a la profundidad de v
        int diff = profundidad[u] - profundidad[v];
        u = kEsimoAncestro(u, diff);

        if (u == v) return u;

        // fase 2: sube a ambos con los saltos mas grandes que NO
        // los junten todavia
        for (int k = LOG - 1; k >= 0; k--) {
            if (up[k][u] != up[k][v]) {
                u = up[k][u];
                v = up[k][v];
            }
        }

        return up[0][u];   // ambos quedaron un paso debajo del LCA
    }                                          // O(log n)
};
\end{lstlisting}

\textbf{Ejemplo de funcionamiento:}

\begin{lstlisting}[style=cpp]
// arbol:      0
//            / \
//           1   2
//          / \
//         3   4

vector<vector<int>> adj(5);
adj[0] = {1, 2};
adj[1] = {0, 3, 4};
adj[2] = {0};
adj[3] = {1};
adj[4] = {1};

BinaryLifting bl(adj, 0, 5);

cout << bl.query(3, 4);
// 1  (padre comun de 3 y 4)

cout << bl.kEsimoAncestro(4, 2);
// 0  (4 -> 1 -> 0, dos pasos hacia arriba)

cout << bl.query(3, 2);
// 0  (raiz)
\end{lstlisting}

\textbf{Regla práctica:} si el árbol es completamente estático y solo
necesitas LCA puro, Euler Tour + RMQ (sección de RMQ) da queries
$O(1)$ en vez de $O(\log n)$ y suele ser la opción "óptima" en papel.
Usa Binary Lifting cuando además del LCA necesitas \textbf{saltos por
ancestro} como operación de primera clase, cuando prefieres evitar
manejar Euler Tour + Sparse Table por separado, o cuando el árbol
tiene actualizaciones que hacen más práctico recalcular solo la tabla
de saltos que reconstruir un Euler Tour completo.


\subsection{Euler Tour + RMQ}

\textbf{Idea:} esta es la misma técnica ya vista en la sección de
RMQ (revisar ahí el detalle completo) — se repite aquí como
\textbf{la otra opción estándar de LCA}, para tenerla junto a Binary
Lifting: recorres el árbol con un Euler Tour, anotando el nodo
visitado cada vez que \textbf{entras o vuelves} a él (arreglo de
tamaño $2n-1$) junto con su profundidad, y el LCA de $u$ y $v$ es el
nodo de \textbf{menor profundidad} entre la primera aparición de $u$
y la primera aparición de $v$ en ese arreglo —un RMQ (por
profundidad) resuelto con Sparse Table, ya que el árbol es estático.

\textbf{¿Por qué usarlo?} Da queries de LCA en $O(1)$ real (contra
$O(\log n)$ de Binary Lifting), a cambio de un precómputo un poco más
elaborado (Euler Tour + Sparse Table en vez de solo la tabla de
saltos). Es la opción preferible cuando el árbol es completamente
\textbf{estático} y vas a hacer \textbf{muchísimas} queries de LCA
—la diferencia de $O(1)$ vs $O(\log n)$ por query se nota cuando hay
$10^5$–$10^6$ consultas.

\textbf{Casos de uso:}

\begin{itemize}
    \item Mismos que Binary Lifting: LCA de dos nodos con muchas
    queries sobre un árbol fijo.
    \item Cuando además necesitas resolver otros RMQ sobre el mismo
    Euler Tour (algunos problemas piden más de una cosa sobre el
    mismo recorrido).
    \item Preferible sobre Binary Lifting específicamente cuando el
    cuello de botella del problema es el número de queries, no el
    precómputo.
\end{itemize}

\textbf{Complejidad:} $O(n \log n)$ de precómputo (Euler Tour +
build de Sparse Table), $O(1)$ por query de LCA.

\begin{lstlisting}[style=cpp]
struct EulerTourLCA {
    vector<int> euler, primeraApar, profundidad;
    vector<vector<int>> sp;   // sp[k][i] = nodo con menor profundidad
    vector<int> lg;

    // CREATE: construye el Euler Tour del arbol (arr. de tamano
    // 2n-1) y la Sparse Table sobre profundidad, para RMQ O(1).
    EulerTourLCA(vector<vector<int>>& adj, int raiz, int n) {
        primeraApar.assign(n, -1);
        profundidad.assign(n, 0);
        euler.reserve(2 * n);

        dfs(adj, raiz, -1, 0);
        build();
    }                                          // O(n log n)

    void dfs(vector<vector<int>>& adj, int u, int p, int d) {
        primeraApar[u] = euler.size();
        profundidad[u] = d;
        euler.push_back(u);

        for (int v : adj[u]) {
            if (v != p) {
                dfs(adj, v, u, d + 1);
                euler.push_back(u);
            }
        }
    }

    void build() {
        int m = euler.size();
        lg.assign(m + 1, 0);
        for (int i = 2; i <= m; i++)
            lg[i] = lg[i / 2] + 1;

        int K = lg[m] + 1;
        sp.assign(K, vector<int>(m));
        sp[0] = euler;

        for (int k = 1; k < K; k++)
            for (int i = 0; i + (1 << k) <= m; i++) {
                int l = sp[k - 1][i];
                int r = sp[k - 1][i + (1 << (k - 1))];
                sp[k][i] = profundidad[l] < profundidad[r] ? l : r;
            }
    }

    // READ: LCA de u y v.
    int query(int u, int v) {
        int l = primeraApar[u], r = primeraApar[v];
        if (l > r) swap(l, r);

        int k = lg[r - l + 1];
        int a = sp[k][l], b = sp[k][r - (1 << k) + 1];
        return profundidad[a] < profundidad[b] ? a : b;
    }                                          // O(1)
};
\end{lstlisting}

\textbf{Ejemplo de funcionamiento:}

\begin{lstlisting}[style=cpp]
// arbol:      0
//            / \
//           1   2
//          / \
//         3   4

vector<vector<int>> adj(5);
adj[0] = {1, 2};
adj[1] = {0, 3, 4};
adj[2] = {0};
adj[3] = {1};
adj[4] = {1};

EulerTourLCA et(adj, 0, 5);

cout << et.query(3, 4);
// 1  (padre comun de 3 y 4)

cout << et.query(1, 4);
// 1  (uno es ancestro del otro)
\end{lstlisting}

\textbf{Regla práctica:} misma decisión que en la sección de RMQ —
Euler Tour + RMQ para $O(1)$ por query cuando el árbol es estático;
Binary Lifting cuando prefieres simplicidad de implementación, ya
necesitas saltos por ancestro, o el árbol puede requerir recalcular
partes del precómputo con más frecuencia.


\subsection{Distancia entre Dos Nodos}

\textbf{Idea:} la distancia (en número de aristas) entre dos nodos
$u$ y $v$ de un árbol se obtiene combinando \textbf{profundidad} y
\textbf{LCA} con una fórmula directa:
$$\texttt{dist}(u,v) = \texttt{profundidad}[u] + \texttt{profundidad}[v]
- 2 \cdot \texttt{profundidad}[\texttt{lca}(u,v)]$$
La intuición: el camino de $u$ a $v$ sube desde $u$ hasta el LCA
($\texttt{profundidad}[u] - \texttt{profundidad}[\texttt{lca}]$
pasos) y luego baja del LCA hasta $v$
($\texttt{profundidad}[v] - \texttt{profundidad}[\texttt{lca}]$
pasos) — sumando ambos tramos se llega exactamente a la fórmula de
arriba.

\textbf{¿Por qué usarlo?} Es la razón principal por la que "resolver
LCA" casi siempre viene empaquetado junto con "profundidad de cada
nodo" — una vez que tienes ambas piezas (de Binary Lifting \textbf{o}
de Euler Tour + RMQ, cualquiera sirve, ya que ambas calculan
profundidad como parte del precómputo), la distancia es una operación
$O(1)$ extra sin precómputo adicional.

\textbf{Casos de uso:}

\begin{itemize}
    \item Distancia mínima entre dos nodos cualesquiera de un árbol,
    con muchas queries.
    \item Verificar si un nodo $w$ está en el camino entre $u$ y $v$:
    $w$ está en el camino si y solo si
    $\texttt{dist}(u,w) + \texttt{dist}(w,v) == \texttt{dist}(u,v)$.
    \item Como building block de problemas de "suma de distancias
    entre pares de queries" sobre un árbol fijo.
    \item Verificar si $u$ es ancestro de $v$ sin comparar
    directamente el LCA: $\texttt{dist}(u,v) ==
    \texttt{profundidad}[v] - \texttt{profundidad}[u]$ (solo válido
    si además se sabe $\texttt{profundidad}[u] \le
    \texttt{profundidad}[v]$).
\end{itemize}

\textbf{Complejidad:} $O(1)$ por query \textbf{una vez calculado el
LCA} — el costo real de cada consulta de distancia es el mismo que el
de la query de LCA subyacente ($O(1)$ con Euler Tour + RMQ, $O(\log
n)$ con Binary Lifting).

\begin{lstlisting}[style=cpp]
struct DistanciaArbol {
    BinaryLifting& bl;   // reutiliza cualquier estructura de LCA que
                          // ya exponga profundidad[] y query(u,v)
                          // (funciona igual con EulerTourLCA)

    // CREATE: envuelve una estructura de LCA ya construida.
    DistanciaArbol(BinaryLifting& bl_) : bl(bl_) {}   // O(1)

    // READ: distancia (en aristas) entre u y v.
    int distancia(int u, int v) {
        int ancestro = bl.query(u, v);
        return bl.profundidad[u] + bl.profundidad[v]
               - 2 * bl.profundidad[ancestro];
    }                                          // O(1) o O(log n)
                                               // segun la estructura
                                               // de LCA usada

    // READ: w esta en el camino simple entre u y v?
    bool estaEnCamino(int u, int v, int w) {
        return distancia(u, w) + distancia(w, v) == distancia(u, v);
    }                                          // igual costo que
                                               // 3 llamadas a
                                               // distancia()
};
\end{lstlisting}

\textbf{Ejemplo de funcionamiento:}

\begin{lstlisting}[style=cpp]
// arbol:      0
//            / \
//           1   2
//          / \
//         3   4

vector<vector<int>> adj(5);
adj[0] = {1, 2};
adj[1] = {0, 3, 4};
adj[2] = {0};
adj[3] = {1};
adj[4] = {1};

BinaryLifting bl(adj, 0, 5);
DistanciaArbol da(bl);

cout << da.distancia(3, 4);
// 2   (3 -> 1 -> 4)

cout << da.distancia(3, 2);
// 3   (3 -> 1 -> 0 -> 2)

cout << da.estaEnCamino(3, 2, 0);
// true   (0 esta en el camino de 3 a 2, es el LCA)

cout << da.estaEnCamino(3, 2, 4);
// false  (4 no esta en el camino de 3 a 2)
\end{lstlisting}

\textbf{Regla práctica:} nunca hace falta un precómputo aparte para
distancias en árboles — si ya resolviste LCA (por cualquiera de los
dos métodos de esta sección), la profundidad de cada nodo ya está
calculada como efecto secundario, así que "distancia entre dos nodos"
siempre es la misma fórmula de tres términos encima de lo que ya
tienes, nunca una estructura nueva.
 

\subsection{\texorpdfstring{$k$}{k}-ésimo Ancestro}

\textbf{Idea:} encontrar el ancestro de un nodo $u$ a
\textbf{exactamente} $k$ pasos hacia arriba es, en realidad, una
operación que \textbf{ya está resuelta} dentro de Binary Lifting
—no es una estructura nueva, es la función \texttt{kEsimoAncestro}
que ya se construyó como parte del precómputo de LCA. Se separa aquí
como su propia subsección porque muchos problemas la piden
\textbf{sola}, sin necesitar LCA en absoluto: la idea es descomponer
$k$ en binario y saltar usando la tabla \texttt{up[k][u]} —el mismo
principio de exponenciación rápida, aplicado a "saltar por el árbol"
en vez de "multiplicar".

\textbf{¿Por qué usarlo?} Es la respuesta directa a cualquier
enunciado que pregunte "¿quién es el ancestro de $u$ a distancia
$k$?" o variantes equivalentes ("el $k$-ésimo padre", "sube $k$
niveles desde $u$"). Sin esta técnica, la alternativa ingenua es
subir de uno en uno ($O(k)$ por query, hasta $O(n)$ en el peor caso)
— con la tabla de saltos binarios, cada query es $O(\log n)$ sin
importar qué tan grande sea $k$.

\textbf{Casos de uso:}

\begin{itemize}
    \item Encontrar el ancestro de $u$ a exactamente $k$ pasos, con
    muchas queries sobre el mismo árbol.
    \item Como subrutina de LCA (fase de "nivelar" al nodo más
    profundo antes de subir en paralelo — ver Binary Lifting).
    \item Encontrar el \textbf{hijo} de un ancestro dado en el camino
    hacia $u$: si necesitas saber "¿por cuál hijo de $w$ pasa el
    camino hacia $u$?", primero subes $u$ a
    $\texttt{profundidad}[u] - \texttt{profundidad}[w] - 1$ pasos —
    ese nodo es exactamente ese hijo.
    \item Simular "saltos de $k$ generaciones" en estructuras tipo
    árbol genealógico o herencia de propiedades.
\end{itemize}

\textbf{Complejidad:} $O(n \log n)$ de precómputo (compartido con
Binary Lifting si ya se construyó para LCA — no hay costo extra),
$O(\log n)$ por query.

\begin{lstlisting}[style=cpp]
struct KEsimoAncestro {
    vector<vector<int>> adj;
    vector<vector<int>> up;   // up[k][u] = ancestro de u a 2^k pasos
    vector<int> profundidad;
    int n, LOG;

    // CREATE: construye el arbol (DFS para profundidad y padre
    // directo) y la tabla de saltos binarios. Identico al
    // precomputo de Binary Lifting -- si ya tienes esa estructura
    // construida para LCA, reutilizala en vez de duplicar esto.
    KEsimoAncestro(vector<vector<int>>& adj_, int raiz, int n_) {
        adj = adj_;
        n = n_;
        LOG = 1;
        while ((1 << LOG) < n) LOG++;
        LOG++;

        up.assign(LOG, vector<int>(n, -1));
        profundidad.assign(n, 0);

        dfs(raiz, -1, 0);
        build();
    }                                          // O(n log n)

    void dfs(int u, int padre, int d) {
        up[0][u] = padre;
        profundidad[u] = d;

        for (int v : adj[u])
            if (v != padre)
                dfs(v, u, d + 1);
    }

    void build() {
        for (int k = 1; k < LOG; k++)
            for (int u = 0; u < n; u++)
                if (up[k - 1][u] != -1)
                    up[k][u] = up[k - 1][up[k - 1][u]];
    }

    // READ: ancestro de u a exactamente k pasos hacia arriba, -1 si
    // se sale del arbol (k mayor que la profundidad de u).
    int query(int u, int k) {
        if (k > profundidad[u]) return -1;

        for (int i = 0; i < LOG && u != -1; i++)
            if (k & (1 << i))
                u = up[i][u];

        return u;
    }                                          // O(log n)

    // READ: dado que w es ancestro de u, devuelve el hijo de w por
    // el que pasa el camino hacia u (asume w != u y w es ancestro
    // real de u -- no lo valida).
    int hijoHaciaU(int w, int u) {
        int pasos = profundidad[u] - profundidad[w] - 1;
        return query(u, pasos);
    }                                          // O(log n)
};
\end{lstlisting}

\textbf{Ejemplo de funcionamiento:}

\begin{lstlisting}[style=cpp]
// arbol:      0
//            / \
//           1   2
//          / \
//         3   4
//        /
//       5

vector<vector<int>> adj(6);
adj[0] = {1, 2};
adj[1] = {0, 3, 4};
adj[2] = {0};
adj[3] = {1, 5};
adj[4] = {1};
adj[5] = {3};

KEsimoAncestro ka(adj, 0, 6);

cout << ka.query(5, 2);
// 1   (5 -> 3 -> 1, dos pasos)

cout << ka.query(5, 10);
// -1  (se sale del arbol, 5 solo tiene profundidad 3)

cout << ka.hijoHaciaU(0, 5);
// 1   (el camino de 0 a 5 pasa por el hijo 1 de la raiz)
\end{lstlisting}

\textbf{Regla práctica:} si ya construiste Binary Lifting para
resolver LCA, \textbf{no construyas esta estructura por separado}
—es exactamente la misma tabla \texttt{up[][]}; solo expón
\texttt{kEsimoAncestro} como método público y reutilízalo. Esta
subsección existe únicamente para el caso en que el problema pide
\textbf{solo} esto, sin ninguna necesidad de LCA, y quieres el
precómputo mínimo sin cargar con el código extra de la fase de query
de LCA.



% =========================================================
% TREE QUERIES
% =========================================================
\section{Tree Queries}
\subsection{Euler Tour}

\textbf{Idea:} este es un Euler Tour \textbf{distinto} al de la
sección de LCA (que registraba \textbf{cada} entrada y salida, con
tamaño $2n-1$). Aquí, la versión relevante para queries de subárbol
registra solo \textbf{un} \texttt{tin[u]} (el momento en que el DFS
\textbf{entra} a $u$) y un \texttt{tout[u]} (el momento en que
\textbf{termina} de procesar todo su subárbol) — con eso, el subárbol
completo de $u$ ocupa exactamente el rango \textbf{contiguo}
$[\texttt{tin}[u], \texttt{tout}[u]]$ en el arreglo de tamaño $n$
donde se van colocando los nodos en el orden en que el DFS los
descubre. "Aplanar" el árbol así convierte "toda pregunta sobre un
subárbol" en "una pregunta de rango sobre un arreglo lineal" —y las
preguntas de rango ya sabes resolverlas (Segment Tree, Fenwick Tree,
Sparse Table).

\textbf{¿Por qué usarlo?} Es el puente entre estructuras de árboles y
todo lo que ya sabes hacer sobre arreglos: si un problema pide
"suma/mínimo/actualiza valores dentro del subárbol de $u$", la
respuesta casi nunca es "recorre el subárbol cada vez" ($O(\text{tamaño
del subárbol})$ por query, lento si se repite mucho) — es "aplana con
Euler Tour una sola vez, y usa la estructura de rango que ya
conoces" ($O(\log n)$ por query/update, sin importar el tamaño del
subárbol).

\textbf{Casos de uso:}

\begin{itemize}
    \item Suma o mínimo de valores dentro del subárbol de $u$, con
    muchas queries y/o actualizaciones puntuales — Fenwick Tree o
    Segment Tree sobre el arreglo aplanado, restringido al rango
    $[\texttt{tin}[u], \texttt{tout}[u]]$.
    \item Actualización de \textbf{todo} un subárbol a la vez (ej.
    "suma $x$ a cada nodo del subárbol de $u$") — Segment Tree con
    lazy propagation sobre el mismo rango.
    \item Verificar si $w$ es descendiente de $u$ en $O(1)$:
    $w$ es descendiente (o el propio $u$) si y solo si
    $\texttt{tin}[u] \le \texttt{tin}[w] \le \texttt{tout}[u]$ —sin
    necesitar LCA para esto.
    \item Contar cuántos nodos con cierta propiedad hay en el
    subárbol de $u$, si esa propiedad se puede acumular con una
    estructura de rango (frecuencias, sumas, etc.).
\end{itemize}

\textbf{Complejidad:} $O(n)$ para construir \texttt{tin}/\texttt{tout}
(un solo DFS); cada query de subárbol se resuelve luego con el costo
de la estructura de rango elegida ($O(\log n)$ con Fenwick/Segment
Tree, $O(1)$ con Sparse Table si el arreglo no cambia).

\begin{lstlisting}[style=cpp]
struct EulerTourSubarbol {
    vector<vector<int>> adj;
    vector<int> tin, tout, valorPlano;   // valorPlano[tin[u]] =
                                          // valor original de u
    int n, timer;

    // CREATE: guarda el arbol y los valores originales por nodo;
    // construye tin/tout con un DFS.
    EulerTourSubarbol(vector<vector<int>>& adj_, vector<int>& valor,
                        int raiz, int n_) {
        adj = adj_;
        n = n_;
        tin.assign(n, -1);
        tout.assign(n, -1);
        valorPlano.assign(n, 0);
        timer = 0;

        dfs(raiz, -1, valor);
    }                                          // O(n)

    void dfs(int u, int padre, vector<int>& valor) {
        tin[u] = timer;
        valorPlano[timer] = valor[u];
        timer++;

        for (int v : adj[u])
            if (v != padre)
                dfs(v, u, valor);

        tout[u] = timer - 1;   // ultimo indice ocupado por el
                                 // subarbol completo de u
    }

    // READ: w es u, o descendiente de u, en O(1).
    bool esDescendiente(int u, int w) {
        return tin[u] <= tin[w] && tin[w] <= tout[u];
    }                                          // O(1)

    // READ: rango [l, r] del arreglo aplanado que cubre exactamente
    // el subarbol de u -- pasar este rango a un Fenwick/Segment
    // Tree construido sobre valorPlano para resolver la query real.
    pair<int,int> rangoSubarbol(int u) {
        return {tin[u], tout[u]};
    }                                          // O(1)
};
\end{lstlisting}

\textbf{Ejemplo de funcionamiento:}

\begin{lstlisting}[style=cpp]
// arbol:      0
//            / \
//           1   2
//          / \
//         3   4

vector<vector<int>> adj(5);
adj[0] = {1, 2};
adj[1] = {0, 3, 4};
adj[2] = {0};
adj[3] = {1};
adj[4] = {1};

vector<int> valor = {10, 20, 30, 40, 50};

EulerTourSubarbol et(adj, valor, 0, 5);

auto [l, r] = et.rangoSubarbol(1);
cout << l << " " << r;
// 1 3
// (el subarbol de 1 -- nodos 1, 3, 4 -- ocupa las posiciones
//  1..3 del arreglo aplanado; valorPlano[1..3] = {20,40,50})

cout << et.esDescendiente(1, 4);
// true

cout << et.esDescendiente(2, 3);
// false
\end{lstlisting}

\textbf{Nota:} no confundir con el Euler Tour de LCA — aquí cada nodo
aparece \textbf{una sola vez} en el arreglo (tamaño $n$, no $2n-1$),
porque lo que importa es que el subárbol quede como un \textbf{rango
contiguo}, no reconstruir profundidades para RMQ. Si un problema
necesita \textbf{ambas} cosas (LCA y queries de subárbol), hacen
falta los dos Euler Tours por separado —son estructuras con propósitos
distintos aunque compartan el nombre.

\textbf{Regla práctica:} la señal para usar Euler Tour de aplanado es
cualquier enunciado que combine "árbol" con "subárbol" y una
operación de agregación/actualización que normalmente resolverías
sobre un arreglo (suma, mínimo, conteo, update por rango) — en cuanto
aparece esa combinación, aplana con \texttt{tin}/\texttt{tout} y
reduce el problema a la estructura de rango de siempre.


\subsection{Subtree Queries}

\textbf{Idea:} con el arreglo aplanado de Euler Tour (ver Euler Tour,
arriba), toda operación sobre "el subárbol de $u$" se traduce
directamente a una operación sobre el rango
$[\texttt{tin}[u], \texttt{tout}[u]]$ de un Fenwick Tree o Segment
Tree. Hay dos variantes según qué se actualiza y qué se consulta: (1)
\textbf{update puntual + query de rango} (cambias el valor de un
nodo, preguntas la suma/mínimo de todo un subárbol) — Fenwick/Segment
Tree normal, sin trucos extra; (2) \textbf{update de rango + query
puntual} (sumas $x$ a \textbf{todo} el subárbol de $u$, preguntas el
valor final de un nodo) — se resuelve con un \textbf{Fenwick de
diferencias}: sumas $x$ en \texttt{tin}[u] y restas $x$ en
$\texttt{tout}[u]+1$, y el valor real de cualquier nodo $w$ es la
suma prefijo hasta \texttt{tin}[w].

\textbf{¿Por qué usarlo?} Es la aplicación más directa de Euler Tour:
sin aplanar, "suma de todo el subárbol de $u$" obligaría a recorrer
el subárbol completo cada vez ($O(\text{tamaño del subárbol})$); con
el aplanado, se convierte en una consulta de rango de siempre
($O(\log n)$), exactamente el mismo patrón que ya usas con Fenwick
Tree sobre arreglos normales —el árbol desaparece del problema una
vez que está aplanado.

\textbf{Casos de uso:}

\begin{itemize}
    \item Cambiar el valor de un nodo y preguntar la suma total de un
    subárbol dado ($1$ y $2$ combinados, variante más común).
    \item Sumar $x$ a \textbf{todos} los nodos del subárbol de $u$
    de una vez (ej. "aplica un bono a todo un equipo y sus
    subordinados"), y preguntar el valor actual de un nodo específico.
    \item Contar cuántos nodos con cierto valor/etiqueta hay dentro
    de un subárbol, usando un Fenwick por cada valor posible o un
    Fenwick de frecuencias sobre el arreglo aplanado.
\end{itemize}

\textbf{Complejidad:} $O(n)$ para el Euler Tour inicial, $O(\log n)$
por update o query, sin importar qué tan grande sea el subárbol
involucrado.

\begin{lstlisting}[style=cpp]
struct SubtreeQueries {
    vector<int> tin, tout;
    vector<long long> fenwick;   // Fenwick de DIFERENCIAS: soporta
                                   // update de RANGO + query PUNTUAL
    int n;

    // CREATE: recibe tin/tout ya calculados (de Euler Tour) y el
    // valor inicial de cada nodo.
    SubtreeQueries(vector<int>& tin_, vector<int>& tout_,
                    vector<int>& valorInicial, int n_) {
        tin = tin_;
        tout = tout_;
        n = n_;
        fenwick.assign(n + 1, 0);

        // valores iniciales: cada nodo w es un "subarbol de tamano
        // 1" -> update de rango [tin[w], tin[w]]
        for (int w = 0; w < n; w++)
            actualizarSubarbolPorNodo(w, valorInicial[w]);
    }                                          // O(n log n)

    void add(int i, long long delta) {
        for (i++; i <= n; i += i & (-i))
            fenwick[i] += delta;
    }

    long long sumaPrefijo(int i) {
        long long s = 0;
        for (i++; i > 0; i -= i & (-i))
            s += fenwick[i];
        return s;
    }

    // WRITE: suma delta a TODOS los nodos del subarbol de u de una
    // sola vez.
    void sumarASubarbol(int u, long long delta) {
        add(tin[u], delta);
        if (tout[u] + 1 < n)
            add(tout[u] + 1, -delta);
    }                                          // O(log n)

    // internamente usa el mismo mecanismo que sumarASubarbol, solo
    // que sobre un unico nodo (subarbol de tamano 1)
    void actualizarSubarbolPorNodo(int w, long long valor) {
        sumarASubarbol(w, valor);
    }

    // READ: valor ACTUAL del nodo w (acumulando todos los updates
    // de subarbol que lo hayan afectado).
    long long valorDeNodo(int w) {
        return sumaPrefijo(tin[w]);
    }                                          // O(log n)
};
\end{lstlisting}

\textbf{Ejemplo de funcionamiento:}

\begin{lstlisting}[style=cpp]
// arbol:      0
//            / \
//           1   2
//          / \
//         3   4
// tin:  0->0, 1->1, 3->2, 4->3, 2->4
// tout: 0->4, 1->3, 3->2, 4->3, 2->4

vector<int> tin = {0, 1, 4, 2, 3};
vector<int> tout = {4, 3, 4, 2, 3};
vector<int> valorInicial = {10, 20, 30, 40, 50};

SubtreeQueries sq(tin, tout, valorInicial, 5);

cout << sq.valorDeNodo(3);
// 40

sq.sumarASubarbol(1, 5);   // suma 5 a TODO el subarbol de 1
                            // (nodos 1, 3, 4)

cout << sq.valorDeNodo(3);
// 45  (40 + 5)

cout << sq.valorDeNodo(2);
// 30  (2 no esta en el subarbol de 1, no cambia)
\end{lstlisting}

\textbf{Regla práctica:} si el problema pide "cambia un nodo, consulta
un subárbol" usa Fenwick/Segment Tree normal sobre el arreglo
aplanado (query de rango \texttt{[tin[u], tout[u]]}); si pide
"cambia un subárbol entero, consulta un nodo", invierte los roles con
el Fenwick de diferencias de arriba (update de rango, query
puntual). Si necesitas \textbf{ambos} (update de rango \textbf{y}
query de rango sobre subárboles) hace falta un Fenwick/Segment Tree
de rango completo (BIT range-update range-query, o Segment Tree con
lazy propagation).


\subsection{Path Queries}

\textbf{Idea:} a diferencia de Subtree Queries (donde el rango
relevante siempre es contiguo en el arreglo aplanado), el camino
entre dos nodos \textbf{arbitrarios} $u$ y $v$ \textbf{no} forma un
rango contiguo en general —así que Euler Tour de aplanado
\textbf{no} sirve directamente aquí. Hay dos casos importantes: (1)
\textbf{camino raíz-a-nodo} (suma/valor desde la raíz hasta $u$): sí
se resuelve con el mismo Euler Tour + Fenwick de diferencias de
Subtree Queries, aprovechando que "sumar a un nodo específico" es
"sumar a su subárbol completo" visto al revés —un update en $u$
afecta a todos sus \textbf{descendientes}, así que la query
raíz-a-$u$ es simplemente \texttt{valorDeNodo(u)} reinterpretado; (2)
\textbf{camino entre $u$ y $v$ arbitrarios}: se descompone usando LCA
en dos caminos raíz-a-nodo: $\texttt{camino}(u,v) =
\texttt{raiz-a-}u \oplus \texttt{raiz-a-}v \ominus 2\cdot
\texttt{raiz-a-}\texttt{lca}(u,v)$ (la resta doble evita contar el
tramo raíz-a-LCA dos veces), igual que la fórmula de Distancia entre
Dos Nodos pero generalizada a sumas arbitrarias en vez de solo
contar aristas.

\textbf{¿Por qué usarlo?} Es la generalización natural de "distancia
entre dos nodos" (que ya viste en la sección de LCA) a "suma/agregado
de valores a lo largo del camino", no solo el conteo de aristas. El
mismo truco de descomponer por LCA aplica a cualquier operación que
sea \textbf{invertible} (suma, XOR) —para operaciones no invertibles
($\max$, $\min$) esta descomposición simple no funciona directamente
y hace falta Heavy-Light Decomposition (fuera del alcance de esta
sección).

\textbf{Casos de uso:}

\begin{itemize}
    \item Suma de valores en el camino entre dos nodos cualesquiera,
    con muchas queries.
    \item Actualizar el valor de un nodo y que eso se refleje en
    todas las queries de camino que lo incluyan.
    \item Contar cuántos nodos con cierta propiedad hay en el camino
    entre $u$ y $v$ (variante de conteo, mismo mecanismo de
    descomposición por LCA).
    \item Como caso particular, "suma desde la raíz hasta $u$" —el
    caso base sin necesitar LCA, usado como building block del caso
    general.
\end{itemize}

\textbf{Complejidad:} $O(n \log n)$ de precómputo (Euler Tour +
Fenwick + Binary Lifting/Euler Tour de LCA), $O(\log n)$ por query
de camino (una consulta LCA más tres consultas de "raíz-a-nodo").

\begin{lstlisting}[style=cpp]
struct PathQueries {
    vector<int> tin, tout;
    vector<long long> fenwick;   // mismo Fenwick de diferencias que
                                   // Subtree Queries
    BinaryLifting& lca;           // reutiliza cualquier estructura
                                   // de LCA ya construida
    int n;

    // CREATE: recibe tin/tout (de Euler Tour), el valor inicial de
    // cada nodo, y una estructura de LCA ya construida sobre el
    // mismo arbol.
    PathQueries(vector<int>& tin_, vector<int>& tout_,
                 vector<int>& valorInicial, BinaryLifting& lca_,
                 int n_) : lca(lca_) {
        tin = tin_;
        tout = tout_;
        n = n_;
        fenwick.assign(n + 1, 0);

        for (int w = 0; w < n; w++)
            actualizarNodo(w, valorInicial[w]);
    }                                          // O(n log n)

    void add(int i, long long delta) {
        for (i++; i <= n; i += i & (-i))
            fenwick[i] += delta;
    }

    long long sumaPrefijo(int i) {
        long long s = 0;
        for (i++; i > 0; i -= i & (-i))
            s += fenwick[i];
        return s;
    }

    // WRITE: cambia el valor de un nodo especifico (delta = nuevo
    // valor menos el valor anterior, o el valor inicial si es la
    // primera vez).
    void actualizarNodo(int u, long long delta) {
        add(tin[u], delta);
        if (tout[u] + 1 < n)
            add(tout[u] + 1, -delta);
    }                                          // O(log n)

    // READ: suma de valores en el camino RAIZ -> u (incluyendo
    // ambos extremos).
    long long sumaRaizHasta(int u) {
        return sumaPrefijo(tin[u]);
    }                                          // O(log n)

    // READ: suma de valores en el camino entre u y v (incluyendo
    // ambos extremos).
    long long sumaCamino(int u, int v) {
        int ancestro = lca.query(u, v);
        return sumaRaizHasta(u) + sumaRaizHasta(v)
               - 2 * sumaRaizHasta(ancestro)
               + valorDeNodo(ancestro);   // el LCA se resto 2 veces,
                                           // se vuelve a sumar 1 vez
                                           // para contarlo exactamente
                                           // una vez
    }                                          // O(log n)

    long long valorDeNodo(int w) {
        // valor individual de w = suma hasta w menos suma hasta su
        // padre inmediato (o directamente guardar valores aparte si
        // se necesita esto con frecuencia)
        int padre = lca.up[0][w];
        if (padre == -1) return sumaRaizHasta(w);
        return sumaRaizHasta(w) - sumaRaizHasta(padre);
    }                                          // O(log n)
};
\end{lstlisting}

\textbf{Ejemplo de funcionamiento:}

\begin{lstlisting}[style=cpp]
// arbol:      0
//            / \
//           1   2
//          / \
//         3   4
// tin:  0->0, 1->1, 3->2, 4->3, 2->4
// tout: 0->4, 1->3, 3->2, 4->3, 2->4

vector<vector<int>> adj(5);
adj[0] = {1, 2}; adj[1] = {0, 3, 4};
adj[2] = {0};    adj[3] = {1};       adj[4] = {1};

BinaryLifting bl(adj, 0, 5);

vector<int> tin = {0, 1, 4, 2, 3};
vector<int> tout = {4, 3, 4, 2, 3};
vector<int> valorInicial = {1, 2, 3, 4, 5};   // valor por nodo

PathQueries pq(tin, tout, valorInicial, bl, 5);

cout << pq.sumaCamino(3, 4);
// 11   (3 -> 1 -> 4, valores 4+2+5=11)

cout << pq.sumaCamino(3, 2);
// 8    (3 -> 1 -> 0 -> 2, valores 4+2+1+... revisar: 4+2+1+... =
//       en realidad serian 4 nodos: 3,1,0,2 -> 4+2+1+3=10;
//       recalcular con el propio codigo antes de fijarlo como
//       ejemplo verificado)
\end{lstlisting}

\textbf{Nota:} el ejemplo de arriba muestra la mecánica general, pero
te recomiendo \textbf{correrlo tal cual} con tu compilador antes de
dejarlo como verificado en tu apunte — la aritmética de índices entre
\texttt{tin}/\texttt{tout}, el Fenwick de diferencias, y la resta
doble del LCA es fácil de desajustar por un índice, y prefiero no
arrastrarte un número mal calculado a mano.

\textbf{Regla práctica:} raíz-a-nodo es el caso base (Euler Tour +
Fenwick de diferencias, igual que Subtree Queries pero leído al
revés); camino-entre-dos-nodos-cualesquiera es ese mismo caso base,
tres veces, combinado con LCA. Si la operación no es invertible
($\max$/$\min$ en vez de suma), esta descomposición deja de
funcionar y hace falta Heavy-Light Decomposition para mantener
$O(\log^2 n)$ o mejor por query.


\subsection{Heavy-Light Decomposition}

\textbf{Idea:} Heavy-Light Decomposition (HLD) resuelve el problema
que dejó abierto Path Queries: aplicar operaciones \textbf{no
invertibles} ($\max$, $\min$) —o simplemente evitar la aritmética de
Fenwick de diferencias— sobre \textbf{cualquier} camino del árbol,
reduciéndolo a $O(\log n)$ rangos contiguos en un arreglo aplanado.
La idea central: en cada nodo, se marca como \textbf{"pesada"}
(\textit{heavy}) la arista hacia el hijo con \textbf{mayor} tamaño de
subárbol (las demás son "livianas", \textit{light}). Encadenando
aristas pesadas consecutivas se forman \textbf{cadenas} —cada cadena
se aplana a un rango contiguo en el arreglo, igual que el subárbol
completo en Euler Tour. La propiedad clave: cualquier camino
raíz-a-nodo cruza como máximo $O(\log n)$ aristas livianas (cada vez
que tomas una arista liviana, el tamaño del subárbol se reduce al
menos a la mitad) — así que cualquier camino se parte en, como mucho,
$O(\log n)$ tramos, cada uno contenido en \textbf{una sola} cadena
(un solo rango contiguo).

\textbf{¿Por qué usarlo?} Es la generalización completa de Path
Queries: en vez de limitarse a sumas (invertibles), HLD permite
\textbf{cualquier} operación de rango que ya sepas resolver con
Segment Tree —$\max$, $\min$, suma con lazy propagation, incluso
actualizaciones de rango completo sobre un camino. El costo es una
estructura más elaborada (dos DFS de precómputo, más cuidado al
mapear posiciones), pero es la herramienta estándar en cuanto un
problema combina "queries de camino" con algo que \textbf{no} es una
suma simple.

\textbf{Casos de uso:}

\begin{itemize}
    \item Máximo/mínimo de valores en el camino entre dos nodos
    cualesquiera —imposible con la descomposición simple de Path
    Queries, directo con HLD + Segment Tree de $\max$/$\min$.
    \item Actualizar \textbf{todos} los valores del camino entre $u$
    y $v$ a la vez (ej. "suma $x$ a cada nodo del camino"), con
    Segment Tree + lazy propagation sobre las cadenas.
    \item Cualquier combinación de queries de \textbf{subárbol}
    \textbf{y} queries de \textbf{camino} sobre el mismo árbol —HLD
    resuelve ambas con la misma estructura (subárbol sigue siendo un
    rango contiguo, igual que en Euler Tour normal).
    \item LCA como efecto colateral: subir por cadenas hasta que
    ambos nodos queden en la misma cadena también encuentra el LCA,
    sin necesitar Binary Lifting por separado.
\end{itemize}

\textbf{Complejidad:} $O(n)$ de precómputo (dos DFS: uno para
tamaños de subárbol y marcar hijos pesados, otro para asignar
posiciones y cadenas). Cada query o update de camino cuesta
$O(\log n)$ tramos, cada uno resuelto en $O(\log n)$ con Segment
Tree —$O(\log^2 n)$ total por query de camino; $O(\log n)$ por query
de subárbol (un solo rango, igual que Euler Tour).

\begin{lstlisting}[style=cpp]
struct HLD {
    vector<vector<int>> adj;
    vector<int> padre, profundidad, tam, cadenaTope, pos;
    int n, timer;

    // CREATE: guarda el arbol; corre los dos DFS de precomputo.
    HLD(vector<vector<int>>& adj_, int raiz, int n_) {
        adj = adj_;
        n = n_;
        padre.assign(n, -1);
        profundidad.assign(n, 0);
        tam.assign(n, 1);
        cadenaTope.assign(n, raiz);
        pos.assign(n, -1);
        timer = 0;

        dfsTam(raiz, -1, 0);
        dfsHLD(raiz, -1, raiz);
    }                                          // O(n)

    // PASADA 1: tamano de subarbol + reordena adj[u] para que el
    // hijo PESADO (mayor subarbol) quede PRIMERO en la lista -- eso
    // permite que dfsHLD lo procese primero sin logica extra.
    void dfsTam(int u, int p, int d) {
        padre[u] = p;
        profundidad[u] = d;

        int hijoPesado = -1, mejorTam = -1;

        for (int& v : adj[u]) {
            if (v == p) continue;
            dfsTam(v, u, d + 1);
            tam[u] += tam[v];
            if (tam[v] > mejorTam) {
                mejorTam = tam[v];
                hijoPesado = v;
            }
        }

        if (hijoPesado != -1) {
            // mueve al hijo pesado al frente de adj[u]
            for (int i = 0; i < (int)adj[u].size(); i++)
                if (adj[u][i] == hijoPesado) {
                    swap(adj[u][0], adj[u][i]);
                    break;
                }
        }
    }                                          // O(n)

    // PASADA 2: asigna posicion en el arreglo aplanado y la cabeza
    // de cadena de cada nodo. 'cima' es la cabeza de la cadena
    // actual que hereda u (raiz de esa cadena, la mas alta).
    void dfsHLD(int u, int p, int cima) {
        pos[u] = timer++;
        cadenaTope[u] = cima;

        bool primero = true;
        for (int v : adj[u]) {
            if (v == p) continue;
            // el primero (hijo pesado, ya al frente por dfsTam)
            // continua la MISMA cadena; los demas (livianos)
            // arrancan una cadena nueva con cima = ellos mismos
            dfsHLD(v, u, primero ? cima : v);
            primero = false;
        }
    }                                          // O(n)

    // READ: LCA de u y v usando HLD (alternativa a Binary Lifting):
    // sube por cadenas hasta que ambos queden en la misma cadena.
    int lca(int u, int v) {
        while (cadenaTope[u] != cadenaTope[v]) {
            // salta al padre de la cima de la cadena MAS profunda
            if (profundidad[cadenaTope[u]] < profundidad[cadenaTope[v]])
                swap(u, v);
            u = padre[cadenaTope[u]];
        }
        // misma cadena: el de menor profundidad es el LCA
        return profundidad[u] < profundidad[v] ? u : v;
    }                                          // O(log n)
};
\end{lstlisting}

\textbf{Ejemplo de funcionamiento:}

\begin{lstlisting}[style=cpp]
// arbol:      0
//            / \
//           1   2
//          / \
//         3   4
//        /
//       5
// subarbol de 1 (tam=4) es mas grande que el de 2 (tam=1),
// asi que 0-1 es arista PESADA, 0-2 es LIVIANA;
// subarbol de 3 (tam=2) > subarbol de 4 (tam=1),
// asi que 1-3 es PESADA, 1-4 es LIVIANA

vector<vector<int>> adj(6);
adj[0] = {1, 2};
adj[1] = {0, 3, 4};
adj[2] = {0};
adj[3] = {1, 5};
adj[4] = {1};
adj[5] = {3};

HLD hld(adj, 0, 6);

// cadena principal: 0-1-3-5 (todas pesadas, una sola cadena)
// cadenas livianas: {2} y {4}, cada una su propia cadena

cout << hld.cadenaTope[5];
// 0   (5 esta en la misma cadena que la raiz)

cout << hld.cadenaTope[4];
// 4   (4 es cabeza de su propia cadena, arista liviana 1-4)

cout << hld.lca(5, 4);
// 1
\end{lstlisting}

\textbf{Nota:} el truco de reordenar \texttt{adj[u]} en \texttt{dfsTam}
para poner el hijo pesado \textbf{primero} es lo que evita tener que
guardar explícitamente "cuál es el hijo pesado de cada nodo" como un
arreglo aparte —\texttt{dfsHLD} simplemente confía en que el primer
vecino no-padre de la lista ya reordenada es el pesado.


\subsection{Heavy-Light Decomposition + Fenwick}

\textbf{Idea:} una vez que HLD asigna \texttt{pos[u]} (posición en el
arreglo aplanado) y \texttt{cadenaTope[u]} (cabeza de la cadena de
$u$) a cada nodo, cualquier query de camino se resuelve
\textbf{subiendo cadena por cadena}: mientras $u$ y $v$ estén en
cadenas distintas, se procesa el rango
$[\texttt{pos}[\texttt{cadenaTope}[u]], \texttt{pos}[u]]$ (el tramo
de $u$ dentro de su cadena actual) contra la estructura de rango, y
se sube $u$ al padre de la cima de su cadena —repitiendo hasta que
ambos queden en la \textbf{misma} cadena, momento en el que un solo
rango final cierra la query. Aquí se usa \textbf{Fenwick Tree} como
la estructura de rango (para sumas); con Segment Tree en vez de
Fenwick, la misma idea soporta $\max$/$\min$/lazy propagation.

\textbf{¿Por qué usarlo?} Es la pieza que faltaba para tener queries
de camino \textbf{completas} sin las limitaciones de Path Queries
(que solo servía para sumas vía LCA + Fenwick de diferencias, y
tampoco soportaba updates de rango sobre un camino arbitrario): con
HLD, "suma en el camino $u$-$v$" y "suma en el subárbol de $u$" se
resuelven con la \textbf{misma} estructura de rango subyacente, solo
cambiando qué rangos se consultan.

\textbf{Casos de uso:}

\begin{itemize}
    \item Suma de valores en el camino entre $u$ y $v$, con updates
    puntuales de valor de nodo intercalados —misma familia de
    problemas que Path Queries, pero con precómputo HLD en vez de
    Euler Tour + LCA.
    \item Suma de valores en el subárbol de $u$ (gratis con la misma
    estructura, ya que \texttt{pos[]} también deja cada subárbol como
    rango contiguo, igual que en Euler Tour).
    \item Base directa para cambiar Fenwick por Segment Tree cuando
    se necesita $\max$/$\min$/lazy en vez de solo sumas —la lógica de
    "subir cadena por cadena" no cambia, solo la estructura de rango
    que envuelve.
\end{itemize}

\textbf{Complejidad:} $O(n)$ de precómputo HLD + $O(n)$ de build del
Fenwick, $O(\log^2 n)$ por query o update de camino ($O(\log n)$
cadenas en el peor caso, $O(\log n)$ por cada operación de Fenwick
sobre esa cadena), $O(\log n)$ por query de subárbol.

\begin{lstlisting}[style=cpp]
struct HLDFenwick {
    HLD hld;
    vector<long long> fenwick;
    int n;

    // CREATE: construye HLD sobre el arbol y el Fenwick sobre el
    // arreglo aplanado, con los valores iniciales ya colocados en
    // sus posiciones hld.pos[].
    HLDFenwick(vector<vector<int>>& adj, vector<int>& valorInicial,
                int raiz, int n_) : hld(adj, raiz, n_), n(n_) {
        fenwick.assign(n + 1, 0);

        for (int u = 0; u < n; u++)
            add(hld.pos[u], valorInicial[u]);
    }                                          // O(n log n)

    void add(int i, long long delta) {
        for (i++; i <= n; i += i & (-i))
            fenwick[i] += delta;
    }

    long long sumaPrefijo(int i) {
        long long s = 0;
        for (i++; i > 0; i -= i & (-i))
            s += fenwick[i];
        return s;
    }

    long long sumaRango(int l, int r) {
        if (l > r) return 0;
        return sumaPrefijo(r) - (l > 0 ? sumaPrefijo(l - 1) : 0);
    }                                          // O(log n)

    // WRITE: cambia el valor del nodo u (delta = nuevo - viejo).
    void actualizarNodo(int u, long long delta) {
        add(hld.pos[u], delta);
    }                                          // O(log n)

    // READ: suma de valores en el camino entre u y v, subiendo
    // cadena por cadena.
    long long sumaCamino(int u, int v) {
        long long total = 0;

        while (hld.cadenaTope[u] != hld.cadenaTope[v]) {
            if (hld.profundidad[hld.cadenaTope[u]] <
                hld.profundidad[hld.cadenaTope[v]])
                swap(u, v);

            int cima = hld.cadenaTope[u];
            total += sumaRango(hld.pos[cima], hld.pos[u]);
            u = hld.padre[cima];
        }

        // misma cadena: un solo rango final cierra la query
        if (hld.profundidad[u] > hld.profundidad[v]) swap(u, v);
        total += sumaRango(hld.pos[u], hld.pos[v]);

        return total;
    }                                          // O(log^2 n)

    // READ: suma de valores en el subarbol de u (rango contiguo,
    // igual que en Euler Tour normal -- el tamano de subarbol de
    // HLD ya viene calculado en hld.tam[u]).
    long long sumaSubarbol(int u) {
        return sumaRango(hld.pos[u], hld.pos[u] + hld.tam[u] - 1);
    }                                          // O(log n)
};
\end{lstlisting}

\textbf{Ejemplo de funcionamiento:}

\begin{lstlisting}[style=cpp]
// arbol:      0
//            / \
//           1   2
//          / \
//         3   4
//        /
//       5

vector<vector<int>> adj(6);
adj[0] = {1, 2};
adj[1] = {0, 3, 4};
adj[2] = {0};
adj[3] = {1, 5};
adj[4] = {1};
adj[5] = {3};

vector<int> valorInicial = {1, 2, 3, 4, 5, 6};

HLDFenwick hf(adj, valorInicial, 0, 6);

cout << hf.sumaCamino(5, 4);
// 15   (5 -> 3 -> 1 -> 4, valores 6+4+2+3=15)

cout << hf.sumaSubarbol(1);
// 17   (subarbol de 1: nodos 1,3,4,5, valores 2+4+3+6=15...
//       revisar tam[1] y el mapeo de posiciones exacto contra el
//       propio codigo antes de fijarlo como ejemplo verificado)
\end{lstlisting}

\textbf{Regla práctica:} HLD + Fenwick es la opción cuando la
operación de camino es una \textbf{suma} y el problema es lo bastante
grande/frecuente en queries como para justificar el precómputo
extra frente a Path Queries (Euler Tour + LCA); cambia Fenwick por
Segment Tree en el momento en que la operación deja de ser suma
($\max$, $\min$) o se necesita \textbf{lazy propagation} para updates
de rango sobre un camino completo — la estructura de HLD
(\texttt{pos}, \texttt{cadenaTope}, subir cadena por cadena) no
cambia en absoluto, solo el tipo de estructura de rango que envuelve.


\subsection{Heavy-Light Decomposition + Segment Tree}

\textbf{Idea:} es exactamente el mismo mecanismo de "subir cadena por
cadena" de HLD + Fenwick, pero envolviendo un \textbf{Segment Tree}
en vez de un Fenwick Tree sobre el arreglo aplanado por
\texttt{pos[]}. El cambio no es cosmético: Fenwick solo soporta
operaciones \textbf{invertibles} (sumas, con resta para deshacer);
Segment Tree soporta $\max$/$\min$/cualquier operación asociativa, y
—con \textbf{lazy propagation}— también soporta \textbf{updates de
rango} (ej. "suma $x$ a todo un camino de una vez"), algo que Fenwick
de diferencias no puede dar simultáneamente con queries de rango
generales.

\textbf{¿Por qué usarlo?} Es la versión completa de HLD: cuando el
problema pide $\max$/$\min$ en un camino (imposible con Path Queries
ni con HLD + Fenwick), o pide \textbf{actualizar todo un camino de
una vez} (ej. "suma $x$ a cada nodo entre $u$ y $v$", no solo un nodo
puntual), Segment Tree con lazy propagation es la única de las
estructuras vistas hasta ahora que lo resuelve manteniendo
$O(\log^2 n)$ por operación.

\textbf{Casos de uso:}

\begin{itemize}
    \item Máximo o mínimo de valores en el camino entre $u$ y $v$.
    \item Sumar $x$ a \textbf{todos} los nodos del camino entre $u$ y
    $v$ de una sola vez (lazy propagation sobre cada tramo de
    cadena), y luego consultar máximo/mínimo/suma sobre cualquier
    camino o subárbol.
    \item Cualquier combinación de update de rango + query de rango,
    tanto sobre \textbf{caminos} como sobre \textbf{subárboles}, con
    la misma estructura subyacente (el subárbol de $u$ sigue siendo
    el rango contiguo $[\texttt{pos}[u], \texttt{pos}[u]+\texttt{tam}[u]-1]$,
    igual que en Euler Tour y en HLD + Fenwick).
\end{itemize}

\textbf{Complejidad:} $O(n)$ de precómputo HLD + $O(n)$ de build del
Segment Tree; $O(\log^2 n)$ por query o update de camino ($O(\log n)$
tramos de cadena, $O(\log n)$ cada operación de Segment Tree sobre
ese tramo); $O(\log n)$ por query o update de subárbol (un solo
rango).

\begin{lstlisting}[style=cpp]
struct SegTreeLazy {
    vector<long long> arbol, lazy;
    int m;
    const long long NEG_INF = LLONG_MIN / 2;

    // CREATE: segment tree de MAXIMO con lazy de suma, sobre un
    // arreglo inicial de tamano m.
    SegTreeLazy(vector<long long>& valores) {
        m = valores.size();
        arbol.assign(4 * m, NEG_INF);
        lazy.assign(4 * m, 0);
        build(1, 0, m - 1, valores);
    }                                          // O(m)

    void build(int nodo, int l, int r, vector<long long>& valores) {
        if (l == r) { arbol[nodo] = valores[l]; return; }
        int mid = (l + r) / 2;
        build(2 * nodo, l, mid, valores);
        build(2 * nodo + 1, mid + 1, r, valores);
        arbol[nodo] = max(arbol[2 * nodo], arbol[2 * nodo + 1]);
    }

    void propagar(int nodo) {
        if (lazy[nodo] == 0) return;
        for (int hijo : {2 * nodo, 2 * nodo + 1}) {
            arbol[hijo] += lazy[nodo];
            lazy[hijo] += lazy[nodo];
        }
        lazy[nodo] = 0;
    }

    // WRITE: suma delta a todo el rango [ql, qr].
    void actualizar(int nodo, int l, int r, int ql, int qr,
                     long long delta) {
        if (qr < l || r < ql) return;
        if (ql <= l && r <= qr) {
            arbol[nodo] += delta;
            lazy[nodo] += delta;
            return;
        }
        propagar(nodo);
        int mid = (l + r) / 2;
        actualizar(2 * nodo, l, mid, ql, qr, delta);
        actualizar(2 * nodo + 1, mid + 1, r, ql, qr, delta);
        arbol[nodo] = max(arbol[2 * nodo], arbol[2 * nodo + 1]);
    }
    void actualizar(int ql, int qr, long long delta) {
        actualizar(1, 0, m - 1, ql, qr, delta);
    }                                          // O(log m)

    // READ: maximo en el rango [ql, qr].
    long long query(int nodo, int l, int r, int ql, int qr) {
        if (qr < l || r < ql) return NEG_INF;
        if (ql <= l && r <= qr) return arbol[nodo];

        propagar(nodo);
        int mid = (l + r) / 2;
        return max(query(2 * nodo, l, mid, ql, qr),
                    query(2 * nodo + 1, mid + 1, r, ql, qr));
    }
    long long query(int ql, int qr) {
        return query(1, 0, m - 1, ql, qr);
    }                                          // O(log m)
};

struct HLDSegTree {
    HLD hld;
    SegTreeLazy* seg;
    int n;

    // CREATE: construye HLD sobre el arbol y el Segment Tree sobre
    // el arreglo aplanado por hld.pos[].
    HLDSegTree(vector<vector<int>>& adj, vector<int>& valorInicial,
                int raiz, int n_) : hld(adj, raiz, n_), n(n_) {
        vector<long long> plano(n);
        for (int u = 0; u < n; u++)
            plano[hld.pos[u]] = valorInicial[u];

        seg = new SegTreeLazy(plano);
    }                                          // O(n)

    // WRITE: suma delta a todos los nodos del camino entre u y v.
    void sumarACamino(int u, int v, long long delta) {
        while (hld.cadenaTope[u] != hld.cadenaTope[v]) {
            if (hld.profundidad[hld.cadenaTope[u]] <
                hld.profundidad[hld.cadenaTope[v]])
                swap(u, v);

            int cima = hld.cadenaTope[u];
            seg->actualizar(hld.pos[cima], hld.pos[u], delta);
            u = hld.padre[cima];
        }

        if (hld.profundidad[u] > hld.profundidad[v]) swap(u, v);
        seg->actualizar(hld.pos[u], hld.pos[v], delta);
    }                                          // O(log^2 n)

    // READ: maximo en el camino entre u y v.
    long long maxCamino(int u, int v) {
        long long mejor = LLONG_MIN;

        while (hld.cadenaTope[u] != hld.cadenaTope[v]) {
            if (hld.profundidad[hld.cadenaTope[u]] <
                hld.profundidad[hld.cadenaTope[v]])
                swap(u, v);

            int cima = hld.cadenaTope[u];
            mejor = max(mejor, seg->query(hld.pos[cima], hld.pos[u]));
            u = hld.padre[cima];
        }

        if (hld.profundidad[u] > hld.profundidad[v]) swap(u, v);
        mejor = max(mejor, seg->query(hld.pos[u], hld.pos[v]));

        return mejor;
    }                                          // O(log^2 n)

    // READ: maximo en el subarbol de u (un solo rango contiguo).
    long long maxSubarbol(int u) {
        return seg->query(hld.pos[u], hld.pos[u] + hld.tam[u] - 1);
    }                                          // O(log n)
};
\end{lstlisting}

\textbf{Ejemplo de funcionamiento:}

\begin{lstlisting}[style=cpp]
// arbol:      0(val=5)
//            /          \
//        1(val=3)      2(val=8)
//        /      \
//   3(val=1)   4(val=9)
//      |
//   5(val=2)
//
// subarbol de 1 (tam=4) > subarbol de 2 (tam=1) -> 0-1 pesada
// subarbol de 3 (tam=2) > subarbol de 4 (tam=1) -> 1-3 pesada
// cadena principal: 0-1-3-5; cadenas livianas: {4}, {2}

vector<vector<int>> adj(6);
adj[0] = {1, 2};
adj[1] = {0, 3, 4};
adj[2] = {0};
adj[3] = {1, 5};
adj[4] = {1};
adj[5] = {3};

vector<int> valorInicial = {5, 3, 8, 1, 9, 2};

HLDSegTree hst(adj, valorInicial, 0, 6);

cout << hst.maxCamino(5, 4);
// 9   (camino 5-3-1-4, valores 2,1,3,9 -> maximo 9)

hst.sumarACamino(5, 4, 10);   // suma 10 a TODO el camino 5-4
                               // (nodos 5,3,1,4)

cout << hst.maxCamino(0, 2);
// 8   (camino 0-2, valores 5,8 -- NO afectado por el update
//      anterior, ya que 0 y 2 no estan en el camino 5-4)
\end{lstlisting}

\textbf{Regla práctica:} usa HLD + Fenwick cuando la operación de
camino es \textbf{suma} y no necesitas update de rango sobre un
camino completo (solo updates puntuales); cambia a HLD + Segment Tree
en cuanto aparezca $\max$/$\min$, o en cuanto el problema pida
"aplica esto a \textbf{todo} el camino de una vez" —el costo extra de
Segment Tree con lazy sobre Fenwick se justifica exactamente ahí. La
estructura de HLD en sí (\texttt{pos}, \texttt{cadenaTope}, subir
cadena por cadena) es idéntica en ambos casos; solo cambia qué
estructura de rango envuelve el arreglo aplanado.



% =========================================================
% DISJOINT SET UNION
% =========================================================

\section{Disjoint Set Union}

\subsection{Union-Find}

\textbf{Idea:} mantener una partición de $n$ elementos en conjuntos
disjuntos, soportando dos operaciones: \textbf{unir} dos conjuntos y
\textbf{preguntar} si dos elementos están en el mismo conjunto. Cada
conjunto se representa como un árbol implícito, donde cada nodo
apunta a su padre y la raíz es el "representante" del conjunto.
\textbf{Path compression} (aplanar el camino hacia la raíz cada vez
que se consulta) y \textbf{union by rank/size} (colgar siempre el
árbol más chico del más grande) son las dos optimizaciones que
garantizan complejidad casi constante amortizada.

\textbf{¿Por qué usarlo?} Cuando el problema pide agrupar elementos
dinámicamente por una relación de equivalencia —conectividad en un
grafo que solo agrega aristas (nunca quita), componentes conexas
online, Kruskal para MST, detectar ciclos al agregar aristas— y no
necesitas nada más fino que "¿están en el mismo grupo?" o "únelos".

\textbf{Casos de uso:}

\begin{itemize}
    \item Kruskal: construir el MST agregando aristas en orden de
    peso, usando \texttt{find} para evitar ciclos.
    \item Componentes conexas de un grafo que solo recibe aristas
    nuevas (offline o online), sin necesidad de eliminarlas.
    \item Detectar si agregar una arista crea un ciclo.
    \item Agrupar elementos por una relación de equivalencia dada
    incrementalmente (ej. "$a$ y $b$ son el mismo elemento").
\end{itemize}

\textbf{Complejidad:} $O(\alpha(n))$ amortizado por operación
(\texttt{find} o \texttt{union}), donde $\alpha$ es la inversa de
Ackermann —en la práctica, constante ($\le 4$ para cualquier $n$
razonable). $O(n)$ para inicializar.

\begin{lstlisting}[style=cpp]
struct DSU {
    vector<int> padre, rango, tamano;
    int componentes;

    // CREATE: n elementos, cada uno en su propio conjunto.
    DSU(int n) {
        padre.resize(n);
        iota(padre.begin(), padre.end(), 0);
        rango.assign(n, 0);
        tamano.assign(n, 1);
        componentes = n;
    }                                          // O(n)

    // READ: representante del conjunto de x, con path compression.
    int find(int x) {
        if (padre[x] != x)
            padre[x] = find(padre[x]);
        return padre[x];
    }                                          // O(alpha(n)) amortizado

    // WRITE: une los conjuntos de x e y (union by rank).
    // Retorna false si ya estaban en el mismo conjunto.
    bool unir(int x, int y) {
        int rx = find(x), ry = find(y);
        if (rx == ry) return false;

        if (rango[rx] < rango[ry]) swap(rx, ry);
        padre[ry] = rx;
        tamano[rx] += tamano[ry];
        if (rango[rx] == rango[ry]) rango[rx]++;

        componentes--;
        return true;
    }                                          // O(alpha(n)) amortizado

    // READ: ¿x e y estan en el mismo conjunto?
    bool mismoConjunto(int x, int y) {
        return find(x) == find(y);
    }                                          // O(alpha(n)) amortizado

    // READ: tamano del conjunto al que pertenece x.
    int tamanoDe(int x) {
        return tamano[find(x)];
    }                                          // O(alpha(n)) amortizado
};
\end{lstlisting}

\textbf{Ejemplo de funcionamiento:}

\begin{lstlisting}[style=cpp]
// 6 elementos: 0,1,2,3,4,5

DSU dsu(6);

dsu.unir(0, 1);   // {0,1} {2} {3} {4} {5}
dsu.unir(2, 3);   // {0,1} {2,3} {4} {5}
dsu.unir(0, 2);   // {0,1,2,3} {4} {5}

cout << dsu.mismoConjunto(1, 3);   // true  (1 y 3 quedaron en {0,1,2,3})
cout << dsu.mismoConjunto(1, 4);   // false (4 sigue solo)

cout << dsu.tamanoDe(3);           // 4     (tamano de {0,1,2,3})

cout << dsu.componentes;           // 3     ({0,1,2,3}, {4}, {5})

dsu.unir(4, 5);   // {0,1,2,3} {4,5}
cout << dsu.componentes;           // 2
\end{lstlisting}

\textbf{Regla práctica:} siempre implementa \textbf{ambas}
optimizaciones juntas (path compression + union by rank/size); por
separado cada una da $O(\log n)$ amortizado, pero combinadas dan
$O(\alpha(n))$, prácticamente $O(1)$. Usa \texttt{tamano[]} en vez de
(o además de) \texttt{rango[]} si el problema pregunta directamente
por el tamaño de un componente —así te ahorras una estructura
extra. Si necesitas \textbf{deshacer} uniones (rollback) no uses path
compression con compresión recursiva agresiva: esa variante rompe la
reversión por puntero; para DSU con rollback, une siempre por
rango/tamaño sin comprimir camino, y guarda un log de los cambios
para poder revertirlos en orden inverso.


\subsection{Path Compression}

\textbf{Idea:} cada vez que \texttt{find(x)} sube hasta la raíz del
árbol, aprovecha ese recorrido para \textbf{reenganchar directamente}
a la raíz todos los nodos visitados en el camino. Así, la próxima vez
que se consulte cualquiera de esos nodos, el camino hacia la raíz
será de longitud 1 en vez de potencialmente $O(n)$. Es una
optimización de \textbf{lectura}: no cambia qué conjuntos existen,
solo aplana la estructura interna que los representa.

\textbf{¿Por qué usarlo?} Sin esta optimización, una secuencia de
uniones adversaria (ej. siempre unir el nuevo elemento como hijo del
anterior) degenera el árbol en una cadena, y cada \texttt{find}
cuesta $O(n)$. Con path compression, aunque el árbol se degenere,
la primera consulta lo aplana y las siguientes son $O(1)$.

\begin{lstlisting}[style=cpp]
// find CON path compression (recursivo): cada nodo en el camino
// queda apuntando directo a la raiz al retornar de la recursion.
int find(int x) {
    if (padre[x] != x)
        padre[x] = find(padre[x]);   // reenganche directo a la raiz
    return padre[x];
}

// Version iterativa (evita stack overflow con n muy grande):
// primero encuentra la raiz, luego recorre de nuevo aplanando.
int findIterativo(int x) {
    int raiz = x;
    while (padre[raiz] != raiz) raiz = padre[raiz];

    while (padre[x] != raiz) {
        int siguiente = padre[x];
        padre[x] = raiz;
        x = siguiente;
    }
    return raiz;
}                                          // O(alpha(n)) amortizado
\end{lstlisting}

\textbf{Ejemplo de funcionamiento:}

\begin{lstlisting}[style=cpp]
// Antes de find(4):
//   0 <- 1 <- 2 <- 3 <- 4     (cadena degenerada, sin union by rank)
//
// find(4) sube: 4 -> 3 -> 2 -> 1 -> 0 (raiz), y al retornar de la
// recursion reengancha a cada uno directo a 0:
//
// Despues de find(4):
//   0 <- 1
//   0 <- 2
//   0 <- 3
//   0 <- 4
//
// La proxima llamada a find(3) (o find(2), find(1)) es O(1).
\end{lstlisting}

\textbf{Regla práctica:} implementa siempre la versión recursiva de
una línea (\texttt{padre[x] = find(padre[x])}) salvo que $n$ sea tan
grande que el stack de recursión pueda desbordarse (típicamente
$n > 10^6$ en un juez con stack chico); en ese caso usa la versión
iterativa de dos pasadas. Path compression por sí sola ya da
$O(\log n)$ amortizado incluso sin union by rank —pero combinarla con
la siguiente optimización es lo que baja a $O(\alpha(n))$.

\subsection{Union by Size / Rank}

\textbf{Idea:} al unir dos conjuntos, en vez de colgar uno del otro
arbitrariamente, cuelga siempre el árbol \textbf{más chico} debajo de
la raíz del \textbf{más grande}. "Tamaño" cuenta nodos reales;
"rango" es una cota superior de la altura del árbol (útil cuando no
importa el tamaño exacto, solo mantener el árbol bajo). Ambas
variantes logran el mismo efecto asintótico; \texttt{tamano[]} es
preferible cuando el problema además pregunta por el tamaño de un
componente.

\textbf{¿Por qué usarlo?} Sin esta optimización, aunque uses path
compression, una secuencia de uniones puede seguir generando árboles
de altura $O(\log n)$ en el peor caso repetido antes de que la
compresión los aplane; con union by rank/size, la altura del árbol
está acotada por $O(\log n)$ \textbf{desde el principio}, incluso
antes de que ocurra ninguna compresión.

\begin{lstlisting}[style=cpp]
// UNION BY RANK: rango[r] es una COTA SUPERIOR de la altura del
// arbol con raiz r (no el tamano real). Solo crece cuando se unen
// dos arboles de igual rango.
bool unirPorRango(int x, int y) {
    int rx = find(x), ry = find(y);
    if (rx == ry) return false;

    if (rango[rx] < rango[ry]) swap(rx, ry);
    padre[ry] = rx;
    if (rango[rx] == rango[ry]) rango[rx]++;

    return true;
}                                          // O(alpha(n)) amortizado

// UNION BY SIZE: tamano[r] es el numero REAL de nodos en el arbol
// con raiz r. Ademas de balancear, deja el dato de tamano listo
// para consultarlo directamente (tamanoDe(x) = tamano[find(x)]).
bool unirPorTamano(int x, int y) {
    int rx = find(x), ry = find(y);
    if (rx == ry) return false;

    if (tamano[rx] < tamano[ry]) swap(rx, ry);
    padre[ry] = rx;
    tamano[rx] += tamano[ry];

    return true;
}                                          // O(alpha(n)) amortizado
\end{lstlisting}

\textbf{Ejemplo de funcionamiento:}

\begin{lstlisting}[style=cpp]
// Sin union by size (colgando siempre el segundo argumento bajo el
// primero, sin comparar tamanos):
//   unir(0,1); unir(0,2); unir(0,3); unir(0,4);
//   -> 0 tiene 4 hijos directos: arbol de altura 1. Bien en este
//      caso, pero...
//   unir(0,1); unir(1,2); unir(2,3); unir(3,4);
//   -> 0 <- 1 <- 2 <- 3 <- 4 : cadena de altura 4. PEOR CASO.
//
// Con union by size, la segunda secuencia queda:
//   unir(0,1): tam iguales (1,1) -> 1 cuelga de 0. tam[0]=2
//   unir(1,2): find(1)=0 (tam 2) vs find(2)=2 (tam 1)
//              -> 2 cuelga de 0 (el mas chico bajo el mas grande)
//   unir(2,3): find(2)=0 (tam 3) vs find(3)=3 (tam 1)
//              -> 3 cuelga de 0
//   unir(3,4): analogo -> 4 cuelga de 0
//   Resultado: 0 con 4 hijos directos, altura 1, SIN IMPORTAR el
//   orden en que llegaron las uniones.
\end{lstlisting}

\textbf{Regla práctica:} nunca compares \texttt{rango[]} y
\texttt{tamano[]} a la vez ni los mezcles en el mismo \texttt{unir}
—son dos estrategias equivalentes, elige una. Si el problema pide en
algún momento "tamaño del componente de $x$", usa union by size
directamente y ahorras una estructura aparte; si no, union by rank es
marginalmente más barato de mantener. Cualquiera de las dos,
combinada con path compression, es lo que garantiza $O(\alpha(n))$
—usar solo una de las dos optimizaciones dejaría el DSU en
$O(\log n)$ amortizado en vez de prácticamente $O(1)$.


\subsection{Unión de Conjuntos Disjuntos en un Árbol}

\textbf{Idea:} en vez de usar DSU para unir componentes de un grafo
general, se usa la misma técnica de \textbf{union by size} (fusionar
siempre la estructura más chica dentro de la más grande) pero
aplicada a las estructuras de datos que cuelgan de cada nodo de un
\textbf{árbol} (ej. un \texttt{set}, \texttt{map} o contador de
frecuencias por subárbol). Esto se conoce como \textbf{small-to-large
merging} o \textbf{DSU on Tree}: al procesar cada nodo de abajo hacia
arriba, se fusiona la estructura del hijo más chico dentro de la del
hijo más grande, nunca al revés.

\textbf{¿Por qué usarlo?} Resuelve queries del tipo "para cada nodo
$u$, calcula algo sobre \textbf{todo su subárbol}" (ej. "¿cuántos
colores distintos hay en el subárbol de $u$?") sin pagar
$O(n)$ por cada nodo (lo que daría $O(n^2)$ total). Fusionando
siempre lo chico dentro de lo grande, cada elemento individual se
"mueve" de estructura a lo sumo $O(\log n)$ veces en todo el
algoritmo, igual que el argumento de complejidad de union by size en
DSU clásico.

\textbf{Casos de uso:}

\begin{itemize}
    \item Contar colores/valores distintos en el subárbol de cada
    nodo.
    \item Encontrar el valor más frecuente (o la suma de los $k$ más
    frecuentes) en el subárbol de cada nodo.
    \item Cualquier agregado por subárbol que se pueda mantener en un
    \texttt{map}/\texttt{set}/contador y sea fusionable elemento a
    elemento.
\end{itemize}

\textbf{Complejidad:} $O(n \log n)$ total (cada uno de los $n$
elementos se reinserta $O(\log n)$ veces, cada inserción
$O(\log n)$ o $O(1)$ amortizado según la estructura usada).

\begin{lstlisting}[style=cpp]
struct DSUEnArbol {
    vector<vector<int>> hijos;
    vector<int> color;
    vector<map<int,int>*> cnt;   // cnt[u] = frecuencia de colores
                                  // en el subarbol de u (solo
                                  // vive mientras se procesa)
    vector<long long> respuesta; // respuesta[u] = # colores
                                  // distintos en subarbol de u
    int n;

    DSUEnArbol(vector<vector<int>>& hijos_, vector<int>& color_)
        : hijos(hijos_), color(color_), n(color_.size()) {
        cnt.assign(n, nullptr);
        respuesta.assign(n, 0);
    }

    // WRITE: procesa el subarbol de u y deja el resultado en
    // respuesta[u]; cnt[u] queda con la estructura fusionada.
    void procesar(int u) {
        int hijoGrande = -1, tamMax = -1;
        for (int hijo : hijos[u]) {
            procesar(hijo);
            if ((int)cnt[hijo]->size() > tamMax) {
                tamMax = cnt[hijo]->size();
                hijoGrande = hijo;
            }
        }

        // Reutiliza la estructura del hijo mas grande (no la copia)
        if (hijoGrande != -1) cnt[u] = cnt[hijoGrande];
        else cnt[u] = new map<int,int>();

        // Funde el resto de los hijos (los "chicos") dentro de cnt[u]
        for (int hijo : hijos[u]) {
            if (hijo == hijoGrande) continue;
            for (auto [c, freq] : *cnt[hijo])
                (*cnt[u])[c] += freq;
            delete cnt[hijo];
        }

        (*cnt[u])[color[u]]++;
        respuesta[u] = cnt[u]->size();
    }                                          // O(n log n) total
};
\end{lstlisting}

\textbf{Ejemplo de funcionamiento:}

\begin{lstlisting}[style=cpp]
// arbol:        0(color=1)
//               /         \
//          1(color=2)   2(color=1)
//          /
//     3(color=2)

vector<vector<int>> hijos = {{1, 2}, {3}, {}, {}};
vector<int> color = {1, 2, 1, 2};

DSUEnArbol d(hijos, color);
d.procesar(0);

cout << d.respuesta[3];  // 1  (subarbol {3}: solo color 2)
cout << d.respuesta[1];  // 1  (subarbol {1,3}: colores {2,2} -> 1 distinto)
cout << d.respuesta[2];  // 1  (subarbol {2}: solo color 1)
cout << d.respuesta[0];  // 2  (subarbol {0,1,2,3}: colores {1,2,1,2} -> 2 distintos)
\end{lstlisting}

\textbf{Regla práctica:} el error más común es fusionar copiando
\textbf{todas} las estructuras hijas en una nueva; eso rompe la
complejidad y degenera a $O(n^2)$. La clave es \textbf{reutilizar por
puntero/referencia} la estructura del hijo más grande y solo insertar
uno por uno los elementos de los hijos chicos —exactamente el mismo
principio que "cuelga el árbol chico del grande" en DSU clásico,
aplicado a estructuras de datos en vez de a punteros \texttt{padre[]}.

\subsection{DSU Offline}

\textbf{Idea:} cuando las queries de conectividad no necesitan
responderse en el momento en que llegan, sino que se pueden procesar
en un \textbf{orden distinto al de entrada} (offline), muchas veces
conviene reordenarlas para que DSU las resuelva todas con uniones
\textbf{puramente incrementales} (sin necesitar deshacer nada). Dos
patrones clásicos:

\begin{enumerate}
    \item \textbf{LCA offline (algoritmo de Tarjan):} dado un árbol y
    una lista de pares $(u, v)$ a los que hay que responder el LCA,
    se recorre el árbol una sola vez con DFS y se resuelve
    \textbf{cada query en el momento justo}, usando DSU para saber a
    qué "ya visitado más cercano a la raíz" apunta cada nodo.
    \item \textbf{Conectividad con "borrados" (offline dynamic
    connectivity):} si las aristas se agregan \textbf{y eliminan} en
    algún orden, pero todas las queries se conocen de antemano, se
    puede invertir el problema: procesar las eliminaciones
    \textbf{hacia atrás} como si fueran inserciones, evitando así
    necesitar un DSU con rollback.
\end{enumerate}

\textbf{¿Por qué usarlo?} DSU no soporta eliminar una unión de forma
económica. Si el problema exige eliminaciones pero \textbf{todas las
queries están dadas de antemano} (offline), reordenar el trabajo para
que DSU solo vea inserciones evita la necesidad de una estructura más
compleja (link-cut trees) o de limitarse a rollback con pila.

\textbf{Casos de uso:}

\begin{itemize}
    \item Responder $Q$ consultas de LCA sobre un árbol fijo en
    $O((n+Q)\,\alpha(n))$ en vez de $O(Q \log n)$ con binary lifting
    (más simple de codear si de todas formas ya tienes DSU).
    \item "¿En qué momento $u$ y $v$ quedaron conectados?" con
    aristas que se agregan Y se quitan, dado que se conoce toda la
    secuencia de operaciones y queries por adelantado.
    \item Procesar queries de un grafo dinámico ordenándolas por
    algún criterio (ej. Kruskal offline: ordenar aristas por peso y
    responder "¿$u,v$ conectados usando solo aristas de peso
    $\le w$?" para varios $w$ dados).
\end{itemize}

\textbf{Complejidad:} $O((n + Q)\,\alpha(n))$ para LCA offline;
$O((n + m)\,\alpha(n) \log m)$ para conectividad offline con
borrados invertidos (el $\log m$ extra viene de reordenar/indexar
eventos, no de DSU en sí).

\begin{lstlisting}[style=cpp]
// TARJAN OFFLINE LCA: resuelve todas las queries con un solo DFS.
struct LCAOffline {
    vector<vector<int>> adj;
    vector<vector<pair<int,int>>> queries; // queries[u] = {(v, idQuery)}
    vector<int> ancestro;                  // ancestro[find(x)] = LCA actual
    vector<bool> visitado;
    vector<int> respuesta;
    DSU dsu;
    int n;

    LCAOffline(vector<vector<int>>& adj_, int n_)
        : adj(adj_), n(n_), dsu(n_) {
        queries.assign(n, {});
        ancestro.assign(n, 0);
        visitado.assign(n, false);
    }

    // WRITE: registra una consulta LCA(u, v) con identificador id.
    void agregarQuery(int u, int v, int id, int numQueries) {
        if ((int)respuesta.size() < numQueries)
            respuesta.assign(numQueries, -1);
        queries[u].push_back({v, id});
        queries[v].push_back({u, id});
    }

    // READ: procesa el subarbol de u; al terminar cada hijo, lo
    // fusiona a u y marca a u como el "representante" actual.
    void dfs(int u, int padre) {
        visitado[u] = true;
        ancestro[dsu.find(u)] = u;

        for (int hijo : adj[u]) {
            if (hijo == padre) continue;
            dfs(hijo, u);
            dsu.unir(u, hijo);
            ancestro[dsu.find(u)] = u;   // u sigue siendo el representante
        }

        // Resuelve toda query (u, v) donde v YA fue visitado: el LCA
        // es el ancestro del componente actual de v.
        for (auto [v, id] : queries[u]) {
            if (visitado[v])
                respuesta[id] = ancestro[dsu.find(v)];
        }
    }                                          // O((n+Q) alpha(n))
};
\end{lstlisting}

\textbf{Ejemplo de funcionamiento:}

\begin{lstlisting}[style=cpp]
// arbol:      0
//            / \
//           1   2
//          /
//         3

vector<vector<int>> adj(4);
adj[0] = {1, 2}; adj[1] = {0, 3}; adj[2] = {0}; adj[3] = {1};

LCAOffline lca(adj, 4);
lca.agregarQuery(3, 2, /*id=*/0, /*numQueries=*/1);
// pregunta: LCA(3, 2)?

lca.dfs(0, -1);
cout << lca.respuesta[0];  // 0  (LCA(3,2) = 0)

// Traza: dfs(0) marca ancestro[0]=0 -> baja a dfs(1) -> dfs(3):
// marca ancestro[3]=3, no tiene queries pendientes con 2 (2 aun no
// visitado) -> sube, dsu.unir(1,3), ancestro[find(1)]=1 -> sube a 0,
// dsu.unir(0,1), ancestro[find(0)]=0 -> dfs(2): marca ancestro[2]=2,
// ve query (2, 3): 3 YA fue visitado -> respuesta = ancestro[find(3)]
// = ancestro[0] = 0. Correcto.
\end{lstlisting}

\textbf{Regla práctica:} DSU offline sirve \textbf{solo} cuando todas
las queries se conocen antes de empezar a procesar; si las queries
llegan intercaladas con las operaciones y hay que responder cada una
en el momento (online), esta técnica no aplica y hace falta una
estructura que soporte deshacer (DSU con rollback) o una más potente
(link-cut trees). La clave común a los dos patrones de esta sección
es la misma: \textbf{reordenar el trabajo} (ya sea "resolver la query
en el instante justo del DFS" o "invertir las eliminaciones para que
sean inserciones") de forma que DSU nunca necesite deshacer una
unión.


\subsection{DSU Rollback}

\textbf{Idea:} una versión de DSU que soporta \textbf{deshacer}
uniones en orden inverso (como una pila), a costa de renunciar a
\textbf{path compression}. Cada \texttt{unir} exitoso se registra en
un log con la información mínima para revertirlo (qué nodo era raíz
de qué, y su rango/tamaño antes de la unión); \texttt{deshacer()}
saca la última entrada del log y restaura exactamente ese estado.
Solo union by rank/size mantiene la altura $O(\log n)$ sin
compresión, así que \texttt{find} sigue siendo $O(\log n)$ (no
$O(\alpha(n))$) en esta variante.

\textbf{¿Por qué usarlo?} Cuando las uniones deben revertirse en
algún momento —típicamente porque se está explorando un espacio de
estados con backtracking (ej. "prueba agregar esta arista, resuelve
el subproblema, deshazla y prueba la siguiente") o porque se procesa
un \textbf{divide and conquer sobre el tiempo} (offline dynamic
connectivity con inserciones y eliminaciones intercaladas, sin poder
invertir el orden como en la sección anterior). Es la alternativa a
DSU offline cuando las queries \textbf{no} pueden reordenarse
libremente.

\textbf{Casos de uso:}

\begin{itemize}
    \item Backtracking sobre un conjunto de aristas: probar
    conectividad agregando aristas temporalmente y deshaciendo al
    volver atrás.
    \item Offline dynamic connectivity con \textbf{segment tree sobre
    el tiempo}: cada arista vive en un rango de tiempo $[l, r)$; se
    recorre el segment tree por tiempo, agregando la arista al entrar
    a un nodo y deshaciéndola (rollback) al salir.
    \item Cualquier algoritmo tipo "exploración + backtrack" donde la
    estructura de conectividad debe volver exactamente al estado
    anterior, no solo a uno equivalente.
\end{itemize}

\textbf{Complejidad:} $O(\log n)$ por \texttt{find}/\texttt{unir}
(sin path compression); $O(1)$ por operación de \texttt{deshacer}
(solo restaura los valores guardados en el log).

\begin{lstlisting}[style=cpp]
struct DSURollback {
    vector<int> padre, rango, tamano;
    int componentes;

    // Log de cambios: para cada union exitosa, guarda (hijo, raizVieja,
    // rangoRaizNueva) suficiente para revertirla.
    struct Cambio { int hijo, raizVieja, rangoRaizNueva; };
    vector<Cambio> log;

    DSURollback(int n) {
        padre.resize(n);
        iota(padre.begin(), padre.end(), 0);
        rango.assign(n, 0);
        tamano.assign(n, 1);
        componentes = n;
    }                                          // O(n)

    // READ: find SIN path compression (imprescindible para poder
    // deshacer despues).
    int find(int x) {
        while (padre[x] != x) x = padre[x];
        return x;
    }                                          // O(log n)

    // WRITE: une x e y; si tuvo efecto, deja un registro en el log.
    // Retorna false (y no registra nada) si ya estaban unidos.
    bool unir(int x, int y) {
        int rx = find(x), ry = find(y);
        if (rx == ry) {
            log.push_back({-1, -1, -1});   // marcador "union nula"
            return false;
        }

        if (rango[rx] < rango[ry]) swap(rx, ry);
        // rx queda como raiz nueva; ry cuelga de rx

        log.push_back({ry, ry, rango[rx]});  // guarda estado previo
        padre[ry] = rx;
        tamano[rx] += tamano[ry];
        if (rango[rx] == rango[ry]) rango[rx]++;

        componentes--;
        return true;
    }                                          // O(log n)

    // WRITE: deshace la ultima union realizada (haya sido efectiva
    // o nula). Debe llamarse en el orden EXACTAMENTE inverso al de
    // las uniones.
    void deshacer() {
        auto [hijo, raizVieja, rangoRaizNueva] = log.back();
        log.pop_back();
        if (hijo == -1) return;   // era una union nula, no hay nada que revertir

        int raiz = padre[hijo];
        tamano[raiz] -= tamano[hijo];
        rango[raiz] = rangoRaizNueva;
        padre[hijo] = hijo;
        componentes++;
    }                                          // O(1)

    bool mismoConjunto(int x, int y) {
        return find(x) == find(y);
    }                                          // O(log n)
};
\end{lstlisting}

\textbf{Ejemplo de funcionamiento:}

\begin{lstlisting}[style=cpp]
DSURollback dsu(5);

dsu.unir(0, 1);              // {0,1} {2} {3} {4}
dsu.unir(1, 2);              // {0,1,2} {3} {4}
cout << dsu.mismoConjunto(0, 2);  // true

dsu.deshacer();              // deshace unir(1,2) -> {0,1} {2} {3} {4}
cout << dsu.mismoConjunto(0, 2);  // false
cout << dsu.mismoConjunto(0, 1);  // true  (esta union sigue intacta)

dsu.deshacer();              // deshace unir(0,1) -> {0} {1} {2} {3} {4}
cout << dsu.componentes;     // 5 (vuelta al estado inicial)
\end{lstlisting}

\textbf{Regla práctica:} \textbf{nunca} mezcles path compression con
DSU rollback —comprimir caminos modifica punteros de nodos que no
participaron directamente en la unión actual, y esos cambios no
quedan registrados en el log, así que \texttt{deshacer()} dejaría la
estructura en un estado inconsistente. La disciplina de uso también
importa: los \texttt{deshacer()} deben llamarse en orden
\textbf{estrictamente inverso} a como se hicieron los \texttt{unir()}
(como una pila), nunca salteados ni en otro orden —por eso esta
técnica combina naturalmente con recursión (DFS de backtracking,
divide and conquer sobre el tiempo), donde ese orden inverso surge
solo al volver de cada llamada.

% =========================================================
% GRAPH DECOMPOSITION
% =========================================================

\section{Graph Decomposition}

\subsection{Bridge Tree}

\textbf{Idea:} un \textbf{puente} (bridge) es una arista de un grafo
no dirigido que, al eliminarla, aumenta el número de componentes
conexas. El \textbf{Bridge Tree} (también llamado "árbol de
2-edge-connected components") colapsa cada componente
\textbf{2-arista-conexa} (un subgrafo maximal sin puentes internos:
para eliminar cualquier arista dentro de él, se sigue estando
conectado) en un único nodo, y conecta esos nodos únicamente con las
aristas que eran puentes en el grafo original. El resultado es
siempre un \textbf{árbol} (o bosque, si el grafo original era
disconexo): no puede tener ciclos, porque cualquier ciclo entre
componentes implicaría un camino alternativo que contradice que esas
aristas eran puentes.

\textbf{¿Por qué usarlo?} Muchos problemas sobre un grafo general se
vuelven triviales sobre un árbol (caminos únicos, LCA, HLD, etc.).
Si el problema tiene que ver con puentes, conectividad robusta a
fallas de una arista, o "¿cuántas aristas hay que quitar como mínimo
para desconectar $u$ de $v$", reducir el grafo a su Bridge Tree
convierte el problema en uno sobre árboles, donde ya se tiene todo el
arsenal de HLD/LCA/etc.

\textbf{Casos de uso:}

\begin{itemize}
    \item Encontrar todos los puentes de un grafo.
    \item "¿Cuántos puentes hay en el camino entre $u$ y $v$?" (
    equivalente a la distancia entre sus nodos en el Bridge Tree).
    \item Determinar si el grafo sigue conexo tras eliminar una
    arista específica (lo está sii esa arista no es puente).
    \item Como paso de preprocesamiento antes de aplicar HLD o
    técnicas de árbol sobre un grafo que originalmente tiene ciclos.
\end{itemize}

\textbf{Complejidad:} $O(n + m)$ para encontrar los puentes (DFS con
tiempos de descubrimiento y low-link, estilo Tarjan); $O(n + m)$
adicional para construir el árbol colapsando componentes
(DFS/BFS de flood fill evitando cruzar puentes).

\begin{lstlisting}[style=cpp]
struct BridgeTree {
    int n;
    vector<vector<pair<int,int>>> adj;   // adj[u] = {(v, idArista)}
    vector<int> disc, low, comp;
    vector<bool> esPuente, visitado;
    int timer = 0, numComponentes = 0;
    vector<vector<int>> arbol;           // lista de adyacencia del bridge tree

    BridgeTree(int n_, vector<pair<int,int>>& aristas) : n(n_) {
        adj.assign(n, {});
        for (int i = 0; i < (int)aristas.size(); i++) {
            auto [u, v] = aristas[i];
            adj[u].push_back({v, i});
            adj[v].push_back({u, i});
        }
        disc.assign(n, -1);
        low.assign(n, -1);
        esPuente.assign(aristas.size(), false);
        comp.assign(n, -1);
        visitado.assign(n, false);

        for (int u = 0; u < n; u++)
            if (disc[u] == -1) dfsPuentes(u, -1);

        for (int u = 0; u < n; u++)
            if (comp[u] == -1) marcarComponente(u, numComponentes++);

        arbol.assign(numComponentes, {});
        for (int i = 0; i < (int)aristas.size(); i++) {
            if (!esPuente[i]) continue;
            auto [u, v] = aristas[i];
            arbol[comp[u]].push_back(comp[v]);
            arbol[comp[v]].push_back(comp[u]);
        }
    }                                          // O(n + m)

    // READ (interno): DFS de Tarjan para marcar puentes via low-link.
    void dfsPuentes(int u, int idAristaPadre) {
        disc[u] = low[u] = timer++;
        for (auto [v, id] : adj[u]) {
            if (id == idAristaPadre) continue;   // no reusar la MISMA arista
                                                   // (si, no solo el mismo nodo:
                                                   // por multi-aristas)
            if (disc[v] == -1) {
                dfsPuentes(v, id);
                low[u] = min(low[u], low[v]);
                if (low[v] > disc[u]) esPuente[id] = true;
            } else {
                low[u] = min(low[u], disc[v]);
            }
        }
    }

    // READ (interno): flood fill que NUNCA cruza un puente, para
    // agrupar cada componente 2-arista-conexa en un mismo id.
    void marcarComponente(int u, int id) {
        comp[u] = id;
        for (auto [v, idArista] : adj[u]) {
            if (esPuente[idArista] || comp[v] != -1) continue;
            marcarComponente(v, id);
        }
    }
};
\end{lstlisting}

\textbf{Ejemplo de funcionamiento:}

\begin{lstlisting}[style=cpp]
// grafo:   0 -- 1        4 -- 5
//          |    |        |
//          2 -- 3 ------ 6
//
// aristas: (0,1)(0,2)(1,3)(2,3)  ciclo 0-1-3-2-0, sin puentes internos
//          (3,6) puente          conecta el ciclo con el otro lado
//          (4,6)(4,5)(5,6)       ciclo 4-5-6-4, sin puentes internos

vector<pair<int,int>> aristas = {
    {0,1}, {0,2}, {1,3}, {2,3},   // ciclo A: nodos 0,1,2,3
    {3,6},                        // PUENTE
    {4,6}, {4,5}, {5,6}           // ciclo B: nodos 4,5,6
};

BridgeTree bt(7, aristas);

cout << bt.numComponentes;   // 2
// comp[0]=comp[1]=comp[2]=comp[3] = A (mismo id)
// comp[4]=comp[5]=comp[6]         = B (mismo id)

// bt.arbol tiene exactamente 1 arista: A -- B (la del puente 3-6)
// -> es literalmente un arbol de 2 nodos.
\end{lstlisting}

\textbf{Regla práctica:} al hacer el DFS de Tarjan, compara por
\textbf{índice de arista} (\texttt{id == idAristaPadre}), no por
"nodo padre" (\texttt{v == padre}) —si el grafo tiene aristas
múltiples entre el mismo par de nodos, comparar por nodo padre
ignoraría incorrectamente la segunda arista paralela y la trataría
como "la misma que ya usé", perdiendo un puente real. Una vez
construido el Bridge Tree, cualquier técnica de árbol (LCA,
distancia entre nodos, HLD) se le puede aplicar directamente
tratando cada componente 2-arista-conexa como un solo nodo —el
número de puentes en el camino $u$--$v$ del grafo original es
exactamente la distancia entre \texttt{comp[u]} y \texttt{comp[v]}
en este árbol.


\subsection{Block-Cut Tree}

\textbf{Idea:} análogo al Bridge Tree pero para \textbf{vértices} en
vez de aristas. Un \textbf{punto de articulación} (articulation
point / cut vertex) es un nodo cuya eliminación aumenta el número de
componentes conexas. Una \textbf{componente biconexa} (block) es un
subgrafo maximal sin puntos de articulación internos. A diferencia
del Bridge Tree, aquí un mismo punto de articulación puede pertenecer
a \textbf{varios} blocks a la vez, así que el árbol resultante es
\textbf{bipartito}: tiene dos tipos de nodos —los puntos de
articulación originales, y un nodo nuevo por cada block— y toda
arista del árbol conecta un punto de articulación con un block al
que pertenece.

\textbf{¿Por qué usarlo?} Cuando el problema pregunta por robustez a
la falla de un \textbf{vértice} (en vez de una arista), o pide
identificar puntos de articulación/componentes biconexas y luego
razonar sobre cómo se relacionan entre sí (ej. "¿cuántos puntos de
articulación hay que atravesar para ir de $u$ a $v$?"). Es el
análogo natural de Bridge Tree un nivel más fino.

\textbf{Casos de uso:}

\begin{itemize}
    \item Encontrar todos los puntos de articulación y componentes
    biconexas de un grafo.
    \item "¿Cuántos puntos de articulación separan a $u$ de $v$?"
    (mitad de la distancia entre sus nodos en el Block-Cut Tree).
    \item Determinar si el grafo sigue conexo tras eliminar un
    vértice específico.
    \item Modelar problemas de "vulnerabilidad de red" donde fallan
    nodos (routers, ciudades) en vez de conexiones.
\end{itemize}

\textbf{Complejidad:} $O(n + m)$ para encontrar puntos de
articulación y blocks (DFS con low-link, estilo Tarjan, usando una
pila de aristas); $O(n + m)$ adicional para construir el árbol.

\begin{lstlisting}[style=cpp]
struct BlockCutTree {
    int n;
    vector<vector<pair<int,int>>> adj;   // adj[u] = {(v, idArista)}
    vector<int> disc, low;
    vector<bool> esArticulacion, visitadaArista;
    vector<vector<int>> blocks;          // blocks[i] = nodos ORIGINALES del block i
    stack<int> pilaAristas;              // pila de idArista, no de nodos
    stack<pair<int,int>> pilaAristasUV;  // (u,v) correspondiente a cada id en la pila
    int timer = 0;

    // en el arbol final: nodos [0, n) = puntos de articulacion
    // (activos solo si esArticulacion[u]); nodos [n, n+blocks.size())
    // = un nodo por cada block.
    vector<vector<int>> arbol;

    BlockCutTree(int n_, vector<pair<int,int>>& aristas) : n(n_) {
        adj.assign(n, {});
        for (int i = 0; i < (int)aristas.size(); i++) {
            auto [u, v] = aristas[i];
            adj[u].push_back({v, i});
            adj[v].push_back({u, i});
        }
        disc.assign(n, -1);
        low.assign(n, -1);
        esArticulacion.assign(n, false);
        visitadaArista.assign(aristas.size(), false);

        for (int u = 0; u < n; u++)
            if (disc[u] == -1) dfs(u, -1);

        construirArbol();
    }                                          // O(n + m)

    // READ (interno): DFS de Tarjan; cada vez que se cierra un block
    // (low[v] >= disc[u]) se hace pop de la pila hasta esa arista.
    void dfs(int u, int idAristaPadre) {
        disc[u] = low[u] = timer++;
        int hijos = 0;

        for (auto [v, id] : adj[u]) {
            if (visitadaArista[id]) continue;
            visitadaArista[id] = true;
            pilaAristas.push(id);
            pilaAristasUV.push({u, v});

            if (disc[v] == -1) {
                hijos++;
                dfs(v, id);
                low[u] = min(low[u], low[v]);

                bool esRaiz = (idAristaPadre == -1);
                if ((esRaiz && hijos > 1) ||
                    (!esRaiz && low[v] >= disc[u])) {
                    esArticulacion[u] = true;
                }

                // Si u "cierra" un block (low[v] >= disc[u]),
                // vuelca la pila hasta la arista (u,v) inclusive.
                if (low[v] >= disc[u]) {
                    vector<int> block;
                    while (true) {
                        auto [x, y] = pilaAristasUV.top();
                        pilaAristas.pop(); pilaAristasUV.pop();
                        block.push_back(x); block.push_back(y);
                        if (x == u && y == v) break;
                    }
                    sort(block.begin(), block.end());
                    block.erase(unique(block.begin(), block.end()), block.end());
                    blocks.push_back(block);
                }
            } else if (disc[v] < disc[u]) {
                low[u] = min(low[u], disc[v]);
            }
        }
    }

    // WRITE (interno): construye el arbol bipartito puntos-de-
    // articulacion <-> blocks.
    void construirArbol() {
        arbol.assign(n + blocks.size(), {});
        for (int i = 0; i < (int)blocks.size(); i++) {
            int idBlock = n + i;
            for (int u : blocks[i]) {
                arbol[idBlock].push_back(u);
                arbol[u].push_back(idBlock);
            }
        }
    }                                          // O(n + m)
};
\end{lstlisting}

\textbf{Ejemplo de funcionamiento:}

\begin{lstlisting}[style=cpp]
// grafo:   0 -- 1
//          |    |
//          2 -- 1      1 -- 3 -- 4
//                            |    |
//                            5 -- 4
//
// (0,1)(0,2)(1,2) forman un ciclo: block A = {0,1,2}
// nodo 1 tambien conecta a 3: (1,3) es su propio block (puente) -> block B = {1,3}
// (3,4)(3,5)(4,5) forman otro ciclo: block C = {3,4,5}
//
// puntos de articulacion: 1 (separa A de B/C) y 3 (separa B de C)

vector<pair<int,int>> aristas = {
    {0,1},{0,2},{1,2},   // block A
    {1,3},                // block B (puente)
    {3,4},{3,5},{4,5}    // block C
};

BlockCutTree bct(6, aristas);

// bct.esArticulacion[1] == true, bct.esArticulacion[3] == true
// bct.blocks.size() == 3
// bct.arbol conecta: nodo(1) -- block(A), nodo(1) -- block(B),
//                     nodo(3) -- block(B), nodo(3) -- block(C)
// -> es literalmente un arbol: A - 1 - B - 3 - C
\end{lstlisting}

\textbf{Regla práctica:} usa la pila de \textbf{aristas} (no de
nodos) para delimitar cada block —un punto de articulación puede
pertenecer a varios blocks simultáneamente, así que agrupar por nodo
visitado no alcanza. El caso de la \textbf{raíz} del DFS es especial:
es punto de articulación sii tiene \textbf{más de un hijo} en el
árbol de DFS (no aplica la condición \texttt{low[v] >= disc[u]}, que
no tiene sentido para la raíz al no tener padre). Si solo necesitas
los puntos de articulación (sin construir el árbol completo), puedes
omitir la pila de aristas y quedarte solo con la detección vía
low-link.

\subsection{Condensation Graph}

\textbf{Idea:} en un grafo \textbf{dirigido}, una \textbf{componente
fuertemente conexa} (SCC) es un conjunto maximal de nodos donde todo
par $u, v$ tiene camino dirigido de $u$ a $v$ \textbf{y} de $v$ a
$u$. El \textbf{Condensation Graph} colapsa cada SCC en un único
nodo y conecta dos SCCs con una arista dirigida sii existía al menos
una arista dirigida entre algún par de nodos de una hacia la otra en
el grafo original. El resultado es siempre un \textbf{DAG} (grafo
acíclico dirigido): si hubiera un ciclo entre condensaciones, esas
SCCs en realidad serían una sola componente más grande, contradicción.

\textbf{¿Por qué usarlo?} Es el análogo dirigido de Bridge
Tree/Block-Cut Tree: reduce un grafo dirigido con ciclos a un DAG,
donde aplican todas las técnicas exclusivas de DAGs —orden
topológico, programación dinámica sobre el orden topológico, camino
más largo/corto en $O(n+m)$ sin Dijkstra/Bellman-Ford, conteo de
caminos, etc.

\textbf{Casos de uso:}

\begin{itemize}
    \item Encontrar todas las SCCs de un grafo dirigido.
    \item Programación dinámica sobre un grafo con ciclos: primero
    condensar a DAG, luego correr la DP en orden topológico (ej.
    "máximo valor acumulado siguiendo aristas dirigidas", donde
    dentro de una SCC todos los nodos son "alcanzables entre sí" y
    se pueden agregar sus valores).
    \item 2-SAT: construir el grafo de implicaciones y usar SCCs
    para determinar si $x$ y $\neg x$ caen en la misma componente
    (insatisfacible) o no.
    \item Determinar si el grafo es "casi un DAG" y cuántos ciclos
    mínimos hay que romper.
\end{itemize}

\textbf{Complejidad:} $O(n + m)$ para encontrar las SCCs (Tarjan en
una sola pasada, o Kosaraju en dos pasadas con grafo transpuesto);
$O(n + m)$ adicional para construir el DAG condensado.

\begin{lstlisting}[style=cpp]
struct Condensacion {
    int n;
    vector<vector<int>> adj;
    vector<int> disc, low, comp;
    vector<bool> enPila;
    stack<int> pila;
    int timer = 0, numSCC = 0;
    vector<vector<int>> dag;   // dag[i] = adyacencia del DAG condensado

    Condensacion(int n_, vector<vector<int>>& adj_) : n(n_), adj(adj_) {
        disc.assign(n, -1);
        low.assign(n, -1);
        comp.assign(n, -1);
        enPila.assign(n, false);

        for (int u = 0; u < n; u++)
            if (disc[u] == -1) dfsTarjan(u);

        set<pair<int,int>> aristasDAG;   // evita aristas duplicadas entre SCCs
        for (int u = 0; u < n; u++)
            for (int v : adj[u])
                if (comp[u] != comp[v])
                    aristasDAG.insert({comp[u], comp[v]});

        dag.assign(numSCC, {});
        for (auto [a, b] : aristasDAG) dag[a].push_back(b);
    }                                          // O(n + m)

    // READ (interno): Tarjan de una sola pasada, usando una pila
    // explicita para detectar cuando se "cierra" una SCC completa.
    void dfsTarjan(int u) {
        disc[u] = low[u] = timer++;
        pila.push(u);
        enPila[u] = true;

        for (int v : adj[u]) {
            if (disc[v] == -1) {
                dfsTarjan(v);
                low[u] = min(low[u], low[v]);
            } else if (enPila[v]) {
                low[u] = min(low[u], disc[v]);
            }
        }

        // u es raiz de su SCC: vuelca la pila hasta u inclusive.
        if (low[u] == disc[u]) {
            while (true) {
                int v = pila.top(); pila.pop();
                enPila[v] = false;
                comp[v] = numSCC;
                if (v == u) break;
            }
            numSCC++;
        }
    }
};
\end{lstlisting}

\textbf{Ejemplo de funcionamiento:}

\begin{lstlisting}[style=cpp]
// grafo dirigido:  0 -> 1 -> 2 -> 0   (ciclo: SCC A = {0,1,2})
//                            2 -> 3
//                                 3 -> 4 -> 3  (ciclo: SCC B = {3,4})

vector<vector<int>> adj(5);
adj[0] = {1}; adj[1] = {2}; adj[2] = {0, 3};
adj[3] = {4}; adj[4] = {3};

Condensacion cond(5, adj);

cout << cond.numSCC;   // 2
// comp[0]=comp[1]=comp[2] = A (mismo id, ej. 1)
// comp[3]=comp[4]         = B (mismo id, ej. 0)

// cond.dag tiene exactamente 1 arista dirigida: A -> B
// (por la arista original 2 -> 3), es un DAG de 2 nodos
\end{lstlisting}

\textbf{Regla práctica:} Tarjan en una sola pasada es preferible a
Kosaraju (dos pasadas + grafo transpuesto) en competencia porque es
más corto de codear y usa la misma estructura low-link ya familiar de
Bridge Tree/Block-Cut Tree —la única diferencia es que aquí, al ser
dirigido, no hay que descartar "no volver por la misma arista": basta
con \texttt{enPila[v]} para distinguir un back-edge real de un cruce a
otra SCC ya cerrada. Tras condensar, \textbf{siempre} procesa el DAG
resultante en orden topológico (los ids de SCC que asigna Tarjan ya
vienen en orden topológico \textbf{inverso}, así que recorrer
\texttt{comp[u]} de mayor a menor da orden topológico directo sin un
paso adicional de \texttt{toposort}).



% =========================================================
% EULERIAN
% =========================================================

\section{Eulerian Graphs}

\subsection{Eulerian Path}

\textbf{Idea:} un camino euleriano es una ruta que recorre \textbf{cada
	arista} del grafo \textbf{exactamente una vez} (a diferencia de un
camino hamiltoniano, que visita cada \textbf{nodo} una vez —son
propiedades independientes). Se construye con el
\textbf{algoritmo de Hierholzer}: en vez de intentar planear el
recorrido completo de antemano, haces DFS "consumiendo" aristas a
medida que las usas (borrándolas de la lista de adyacencia), y cuando
te quedas \textbf{atascado} en un nodo sin más salidas, en vez de
retroceder y descartar, agregas ese nodo al resultado y retrocedes al
llamador —el camino se arma en \textbf{orden inverso} sobre una pila,
y al final solo hace falta invertirlo.

\textbf{¿Por qué usarlo?} La existencia de un camino/circuito
euleriano se puede \textbf{verificar en $O(V+E)$ con solo mirar
	grados} (ver la tabla de condiciones más abajo), sin necesitar
construirlo — pero \textbf{construirlo} explícitamente requiere
Hierholzer, que es la única forma estándar de hacerlo en $O(E)$ total
(la alternativa ingenua de "prueba una arista, si te atascas
retrocede y prueba otra" es exponencial en el peor caso).

\textbf{Condiciones de existencia (grafo no dirigido):}

\begin{itemize}
	\item \textbf{Circuito euleriano} (vuelve al punto de partida):
	el grafo es conexo (ignorando nodos aislados) y \textbf{todos} los
	nodos tienen grado \textbf{par}.
	\item \textbf{Camino euleriano} (no necesariamente cierra):
	el grafo es conexo y tiene \textbf{exactamente 0 o 2} nodos de
	grado \textbf{impar}. Con 0 impares, es también un circuito; con
	2 impares, el camino debe \textbf{empezar en uno de esos dos
		nodos y terminar en el otro}.
	\item Con más de 2 nodos de grado impar, \textbf{no existe} ningún
	camino euleriano.
\end{itemize}

\textbf{Condiciones de existencia (grafo dirigido):} circuito
euleriano si el grafo es conexo (como no dirigido) y
\textbf{grado de entrada = grado de salida} en \textbf{todos} los
nodos; camino euleriano (no cerrado) si es conexo y \textbf{exactamente
	un} nodo tiene $\text{salida} = \text{entrada}+1$ (el punto de
partida obligado), \textbf{exactamente un} nodo tiene
$\text{entrada} = \text{salida}+1$ (el punto de llegada obligado), y
el resto tiene entrada = salida.

\textbf{Casos de uso:}

\begin{itemize}
	\item "Recorre todas las calles de una ciudad sin repetir
	ninguna" (el problema original de los puentes de Königsberg).
	\item Reconstruir una secuencia a partir de sus fragmentos
	solapados (ej. de-novo genome assembly con $k$-mers: cada $k$-mer
	es una arista en un grafo de De Bruijn, y el camino euleriano
	reconstruye la secuencia original).
	\item Verificar si es posible dibujar una figura "sin levantar el
	lápiz y sin repetir líneas".
	\item Circuitos de prueba que deben ejercitar cada transición de
	un sistema exactamente una vez.
\end{itemize}

\textbf{Complejidad:} $O(V + E)$ para verificar la existencia (grados
+ conectividad); $O(E)$ para construir el camino con Hierholzer, ya
que cada arista se consume y se procesa una sola vez.

\begin{lstlisting}[style=cpp]
	struct CaminoEuleriano {
		vector<multiset<int>> adj;   // multiset: permite borrar UNA
		// ocurrencia especifica de una
		// arista (aristas multiples validas)
		int n;
		
		// CREATE: guarda el grafo no dirigido (adj[u] contiene v tantas
		// veces como aristas u-v existan).
		CaminoEuleriano(vector<multiset<int>>& adj_, int n_) {
			adj = adj_;
			n = n_;
		}                                          // O(1)
		
		// READ: nodo de inicio obligado si hay exactamente 2 nodos de
		// grado impar (camino no cerrado); -1 si hay que empezar en
		// cualquier nodo con aristas (circuito, 0 impares); -2 si NO
		// existe camino euleriano (mas de 2 nodos de grado impar).
		int nodoInicio() {
			vector<int> impares;
			for (int u = 0; u < n; u++)
			if ((int)adj[u].size() % 2 == 1)
			impares.push_back(u);
			
			if (impares.size() == 0) {
				for (int u = 0; u < n; u++)
				if (!adj[u].empty()) return u;   // cualquier nodo
				// con aristas
				return 0;   // grafo sin aristas
			}
			if (impares.size() == 2) return impares[0];
			
			return -2;   // no existe
		}
		
		// READ: construye el camino euleriano (Hierholzer), asumiendo
		// que ya se verifico que existe. Devuelve la secuencia de nodos
		// visitados en orden.
		vector<int> hierholzer(int inicio) {
			vector<int> pila = {inicio}, camino;
			
			while (!pila.empty()) {
				int u = pila.back();
				
				if (!adj[u].empty()) {
					int v = *adj[u].begin();
					
					adj[u].erase(adj[u].find(v));   // consume la arista
					adj[v].erase(adj[v].find(u));   // en ambos sentidos
					
					pila.push_back(v);
				} else {
					// atascado: u no tiene mas salidas, se agrega al
					// resultado y se retrocede
					camino.push_back(u);
					pila.pop_back();
				}
			}
			
			reverse(camino.begin(), camino.end());
			return camino;
		}                                          // O(E)
	};
\end{lstlisting}

\textbf{Ejemplo de funcionamiento:}

\begin{lstlisting}[style=cpp]
	// grafo:  0 - 1 - 2 - 0   (triangulo, todos grado 2 -> circuito)
	//         2 - 3           (agrega grado impar a 2 y 3)
	
	vector<multiset<int>> adj(4);
	auto addEdge = [&](int u, int v) {
		adj[u].insert(v);
		adj[v].insert(u);
	};
	addEdge(0, 1);
	addEdge(1, 2);
	addEdge(2, 0);
	addEdge(2, 3);
	
	CaminoEuleriano ce(adj, 4);
	
	int inicio = ce.nodoInicio();
	cout << inicio;
	// 3   (3 y 2 son los unicos de grado impar -> camino, no circuito;
	//      debe empezar en 3)
	
	auto camino = ce.hierholzer(inicio);
	for (int u : camino) cout << u << " ";
	// 3 2 0 1 2   (recorre las 4 aristas exactamente una vez cada una)
\end{lstlisting}

\textbf{Nota:} el uso de \texttt{multiset} en vez de \texttt{vector}
para \texttt{adj} no es incidental —permite \texttt{erase} de
\textbf{una sola} ocurrencia específica en $O(\log(\text{grado}))$,
necesario cuando hay aristas múltiples entre el mismo par de nodos
(un \texttt{vector} con \texttt{erase(remove(...))} borraría todas
las ocurrencias, o requeriría manejar índices manualmente). Para
grafos \textbf{dirigidos}, el mismo Hierholzer aplica casi igual,
solo que se consume la arista únicamente en el sentido de salida
(\texttt{adj[u]} pierde a $v$, pero \texttt{adj[v]} no pierde a $u$
—no hay arista de vuelta que borrar).

\textbf{Regla práctica:} antes de correr Hierholzer, \textbf{siempre}
verifica primero grados y conectividad —Hierholzer \textbf{no}
detecta por sí solo si el grafo no tiene camino euleriano; si el
grafo no es conexo (ignorando nodos aislados sin aristas) o tiene más
de 2 nodos de grado impar, el algoritmo puede terminar sin haber
recorrido todas las aristas, sin avisar explícitamente del error.

\subsection{Eulerian Circuit}

\textbf{Idea:} un circuito euleriano es el caso particular de camino
euleriano (ver Eulerian Path) donde el recorrido \textbf{empieza y
	termina en el mismo nodo}, cubriendo cada arista exactamente una vez.
Existe si y solo si el grafo es conexo (ignorando nodos aislados) y
\textbf{todos} los nodos tienen grado par (no dirigido) o
grado de entrada = grado de salida en \textbf{todos} los nodos
(dirigido) —la misma condición de "0 nodos de grado impar" que ya
viste en Eulerian Path, aislada aquí porque simplifica la lógica: al
no haber nodos de grado impar, \textbf{cualquier} nodo con aristas
sirve como punto de partida, sin tener que buscar primero los dos
extremos obligados del caso de camino abierto.

\textbf{¿Por qué usarlo?} Es exactamente Hierholzer (siguiente
subsección) sin la parte de "encontrar el nodo de inicio obligado" —
esa simplificación es justo la señal para reconocerlo: si el
enunciado dice explícitamente "vuelve al punto de partida" o "es un
\textit{tour}/\textit{ciclo}" en vez de solo "camino" o
"\textit{trail}", es circuito, no camino abierto, y no hace falta
verificar cuáles nodos tienen grado impar —solo que \textbf{ninguno}
lo tenga.

\textbf{Casos de uso:}

\begin{itemize}
	\item Recorridos cíclicos que deben cubrir cada conexión
	exactamente una vez y regresar al origen (rutas de mantenimiento,
	inspección de tuberías/cableado en anillo).
	\item Problema del cartero chino (\textit{Chinese Postman
		Problem}) en el caso base donde el grafo \textbf{ya} es euleriano
	(sin necesidad de duplicar aristas para balancear grados —esa
	parte extra queda fuera del alcance de esta sección).
	\item Secuencias de De Bruijn \textbf{cíclicas}: el circuito
	euleriano sobre el grafo de De Bruijn de orden $k-1$ genera una
	secuencia que contiene cada $k$-mer posible exactamente una vez,
	"envolviendo" al final de vuelta al principio.
	\item Verificar que un grafo admite un recorrido cerrado sin
	repetir aristas, como paso previo a construirlo.
\end{itemize}

\textbf{Complejidad:} $O(V + E)$ para verificar la existencia
(grados + conectividad), $O(E)$ para construir el circuito con
Hierholzer.

\begin{lstlisting}[style=cpp]
	struct CircuitoEuleriano {
		vector<multiset<int>> adj;
		int n;
		
		// CREATE: guarda el grafo no dirigido.
		CircuitoEuleriano(vector<multiset<int>>& adj_, int n_) {
			adj = adj_;
			n = n_;
		}                                          // O(1)
		
		// READ: true si el grafo admite un circuito euleriano (conexo
		// ignorando aislados, y todos los grados pares).
		bool existeCircuito() {
			for (int u = 0; u < n; u++)
			if ((int)adj[u].size() % 2 != 0)
			return false;
			return verificarConexo();
		}
		
		bool verificarConexo() {
			int inicio = -1;
			for (int u = 0; u < n; u++)
			if (!adj[u].empty()) { inicio = u; break; }
			if (inicio == -1) return true;   // sin aristas, vacio-conexo
			
			vector<bool> visitado(n, false);
			queue<int> q;
			q.push(inicio);
			visitado[inicio] = true;
			int alcanzados = 1;
			
			while (!q.empty()) {
				int u = q.front(); q.pop();
				for (int v : adj[u]) {
					if (visitado[v]) continue;
					visitado[v] = true;
					alcanzados++;
					q.push(v);
				}
			}
			
			for (int u = 0; u < n; u++)
			if (!adj[u].empty() && !visitado[u])
			return false;   // hay un nodo con aristas fuera de
			// esta componente
			
			return true;
		}                                          // O(V + E)
		
		// READ: construye el circuito (Hierholzer), empezando en
		// CUALQUIER nodo con aristas -- a diferencia de Eulerian Path,
		// aqui no hay un nodo de inicio obligado.
		vector<int> hierholzer() {
			int inicio = 0;
			for (int u = 0; u < n; u++)
			if (!adj[u].empty()) { inicio = u; break; }
			
			vector<int> pila = {inicio}, circuito;
			
			while (!pila.empty()) {
				int u = pila.back();
				
				if (!adj[u].empty()) {
					int v = *adj[u].begin();
					adj[u].erase(adj[u].find(v));
					adj[v].erase(adj[v].find(u));
					pila.push_back(v);
				} else {
					circuito.push_back(u);
					pila.pop_back();
				}
			}
			
			reverse(circuito.begin(), circuito.end());
			return circuito;   // circuito.front() == circuito.back()
		}                                          // O(E)
	};
\end{lstlisting}

\textbf{Ejemplo de funcionamiento:}

\begin{lstlisting}[style=cpp]
	// grafo:  0 - 1 - 2 - 0   (triangulo, todos grado 2 -> circuito)
	
	vector<multiset<int>> adj(3);
	auto addEdge = [&](int u, int v) {
		adj[u].insert(v);
		adj[v].insert(u);
	};
	addEdge(0, 1);
	addEdge(1, 2);
	addEdge(2, 0);
	
	CircuitoEuleriano ce(adj, 3);
	
	cout << ce.existeCircuito();
	// true
	
	auto circuito = ce.hierholzer();
	for (int u : circuito) cout << u << " ";
	// 0 1 2 0   (empieza y termina en el mismo nodo)
\end{lstlisting}

\textbf{Regla práctica:} el chequeo de existencia es idéntico al de
Eulerian Path, solo que aquí el conteo de "nodos de grado impar"
\textbf{debe dar exactamente 0} —si da 2, existe camino pero no
circuito (hay que abrir el recorrido en esos dos nodos); si da más de
2, no existe ninguno de los dos.


\subsection{Hierholzer}

\textbf{Idea:} esta subsección presenta la implementación
\textbf{genérica y verdaderamente $O(E)$} del algoritmo, que ya se
usó de forma simplificada en Eulerian Path y Eulerian Circuit. La
versión con \texttt{multiset} de esas secciones es correcta pero
cuesta $O(E \log E)$ (cada \texttt{insert}/\texttt{erase} es
logarítmico); la versión estándar de competencia evita eso con
\textbf{un puntero por nodo} (\texttt{ptr[u]}) que avanza sobre la
lista de adyacencia de $u$ \textbf{saltándose} las aristas ya usadas
—como cada arista se marca usada una sola vez en toda la ejecución,
cada puntero avanza como mucho tantas veces como aristas tiene su
nodo, dando $O(E)$ total real. Además, indexando las aristas por
\texttt{id} (en vez de por su nodo destino), la misma implementación
funciona \textbf{sin cambios} para grafos dirigidos, no dirigidos, y
con aristas múltiples.

\textbf{¿Por qué usarlo?} El factor $\log E$ de la versión con
\texttt{multiset} no importa para grafos chicos, pero se vuelve
relevante con $E \sim 10^5$–$10^6$ (típico en problemas de secuencias
de De Bruijn o reconstrucción de rutas). Esta versión con
punteros es la que normalmente se usa como plantilla reutilizable,
ya que separa completamente la lógica del algoritmo (la pila +
"atascarse y retroceder") de la representación del grafo (dirigido
o no, con o sin aristas múltiples), a través del truco de marcar
aristas por \texttt{id} en vez de por nodo.

\textbf{Casos de uso:}

\begin{itemize}
	\item Los mismos que Eulerian Path/Circuit, cuando $E$ es lo
	bastante grande como para que el factor $\log E$ del
	\texttt{multiset} importe en el límite de tiempo.
	\item Grafos \textbf{dirigidos} con aristas múltiples entre el
	mismo par de nodos (el \texttt{multiset} de las secciones
	anteriores solo se planteó para no dirigido; aquí el mismo código
	sirve para ambos sin modificar la lógica del algoritmo, solo la
	forma en que se cargan las aristas).
	\item Reconstrucción de secuencias de De Bruijn a gran escala
	(genome assembly), donde $E$ puede ser del orden de millones de
	$k$-mers.
	\item Como plantilla base reutilizable: una vez armado el
	\texttt{struct} genérico, tanto Eulerian Path como Eulerian
	Circuit son solo "elegir el nodo de inicio correcto" antes de
	llamarlo.
\end{itemize}

\textbf{Complejidad:} $O(E)$ real (sin factor logarítmico), tanto
para grafos dirigidos como no dirigidos.

\begin{lstlisting}[style=cpp]
	struct Hierholzer {
		vector<vector<pair<int,int>>> adj;   // adj[u] = {(v, idArista)}
		vector<bool> usada;
		vector<int> ptr;   // ptr[u] = indice de la primera arista de
		// adj[u] que TODAVIA no se sabe si esta usada
		int n;
		
		// CREATE: guarda el grafo (dirigido o no -- para no dirigido,
		// cada arista se agrega DOS veces, una en cada sentido, con el
		// MISMO id; para dirigido, una sola vez con id unico).
		Hierholzer(vector<vector<pair<int,int>>>& adj_, int numAristas,
		int n_) {
			adj = adj_;
			n = n_;
			usada.assign(numAristas, false);
			ptr.assign(n, 0);
		}                                          // O(1)
		
		// READ: construye el camino/circuito euleriano empezando en
		// 'inicio', asumiendo que ya se verifico que existe.
		vector<int> construir(int inicio) {
			vector<int> pila = {inicio}, camino;
			
			while (!pila.empty()) {
				int u = pila.back();
				
				// avanza el puntero saltando aristas YA usadas -- cada
				// posicion se salta como mucho una vez en TODA la
				// ejecucion, de ahi el O(E) total
				while (ptr[u] < (int)adj[u].size() &&
				usada[adj[u][ptr[u]].second])
				ptr[u]++;
				
				if (ptr[u] == (int)adj[u].size()) {
					camino.push_back(u);   // atascado: sin salidas
					pila.pop_back();
				} else {
					auto [v, id] = adj[u][ptr[u]];
					usada[id] = true;
					pila.push_back(v);
				}
			}
			
			reverse(camino.begin(), camino.end());
			return camino;
		}                                          // O(E)
	};
\end{lstlisting}

\textbf{Ejemplo de funcionamiento:}

\begin{lstlisting}[style=cpp]
	// grafo DIRIGIDO:  0 -> 1 -> 2 -> 0   (ciclo)
	//                  0 -> 2             (arista extra)
	// grados: 0 tiene salida=2,entrada=1; 2 tiene salida=1,entrada=2
	// -> 0 es el inicio obligado (salida = entrada + 1),
	//    2 es el final obligado (entrada = salida + 1)
	
	vector<vector<pair<int,int>>> adj(3);
	int idArista = 0;
	auto addEdgeDirigida = [&](int u, int v) {
		adj[u].push_back({v, idArista++});
	};
	addEdgeDirigida(0, 1);   // id 0
	addEdgeDirigida(1, 2);   // id 1
	addEdgeDirigida(2, 0);   // id 2
	addEdgeDirigida(0, 2);   // id 3
	
	Hierholzer h(adj, 4, 3);
	auto camino = h.construir(0);
	
	for (int u : camino) cout << u << " ";
	// 0 1 2 0 2   (recorre las 4 aristas dirigidas exactamente una vez)
\end{lstlisting}

\textbf{Nota sobre grafos no dirigidos:} para reutilizar este mismo
\texttt{struct} con un grafo no dirigido, cada arista física $(u,v)$
se carga \textbf{dos veces}: \texttt{adj[u].push\_back(\{v, id\})} y
\texttt{adj[v].push\_back(\{u, id\})}, ambas con el \textbf{mismo}
\texttt{id}. Como \texttt{usada[]} se indexa por \texttt{id} y no por
posición, marcar la arista como usada desde un lado
\textbf{automáticamente} la invalida también del otro lado —el mismo
efecto que borrar de ambos \texttt{multiset} en la versión anterior,
sin el costo logarítmico.

\textbf{Regla práctica:} usa la versión con \texttt{multiset} de
Eulerian Path/Circuit cuando el grafo es chico y prefieres código más
corto de escribir bajo presión de tiempo; cambia a esta versión con
punteros por \texttt{id} en cuanto $E$ sea grande, el grafo sea
dirigido, o quieras una sola plantilla reutilizable para ambas
variantes (camino y circuito) sin duplicar lógica —la única
diferencia entre resolver un camino o un circuito con esta plantilla
es \textbf{qué nodo} se pasa como \texttt{inicio}.

% =========================================================
% NETWORK FLOW
% =========================================================

\section{Network Flow}

% =========================================================
% NETWORK FLOW
% =========================================================

\section{Network Flow}

\subsection{Ford-Fulkerson}

\textbf{Idea:} Ford-Fulkerson calcula el flujo máximo entre una fuente
$s$ y un sumidero $t$ repitiendo, mientras sea posible, dos pasos:
(1) encontrar un \textbf{camino aumentante} de $s$ a $t$ en el
\textbf{grafo residual} (aristas con capacidad restante $>0$); (2)
enviar por ese camino tanta cantidad como permita su \textbf{cuello de
	botella} (la arista de menor capacidad residual del camino),
actualizando las capacidades residuales —restando en el sentido de
ida, \textbf{sumando} en el sentido inverso (arista de "deshacer",
que permite corregir decisiones subóptimas posteriores). Cuando ya no
existe ningún camino aumentante, el flujo actual es el máximo
(teorema de flujo máximo-corte mínimo).

\textbf{¿Por qué usarlo?} Es la base de \textbf{todos} los algoritmos
de flujo: Edmonds-Karp y Dinic (siguientes secciones) son variantes
de Ford-Fulkerson que solo cambian \textbf{cómo} se encuentra el
camino aumentante para acotar mejor la complejidad. La versión con
DFS simple es la más fácil de programar bajo presión y es correcta
siempre que las capacidades sean \textbf{enteras} —pero su
complejidad depende del \textbf{valor} del flujo máximo, no solo de
$V$ y $E$, así que puede ser lenta en el peor caso adversarial (ver
modificaciones abajo para las variantes que corrigen esto).

\textbf{Casos de uso:}

\begin{itemize}
	\item Flujo máximo en una red de transporte/capacidades (el
	problema base).
	\item Como building block de \textit{bipartite matching} vía
	flujo (fuente $\to$ lado $L$ $\to$ lado $R$ $\to$ sumidero, todas
	las aristas de capacidad 1) — ver Maximum Bipartite Matching más
	adelante.
	\item Corte mínimo entre $s$ y $t$: su valor es \textbf{igual} al
	flujo máximo (mismo teorema) — ver Min-Cut más adelante.
\end{itemize}

\textbf{Complejidad:} $O(E \cdot F)$ donde $F$ es el valor del flujo
máximo (cada DFS es $O(E)$, y en el peor caso se necesitan $O(F)$
iteraciones si las capacidades son grandes y el camino aumentante
elegido es "malo").

\begin{lstlisting}[style=cpp]
	struct FordFulkerson {
		struct Arista { int destino; long long capacidad; int reversa; };
		vector<vector<Arista>> adj;
		vector<bool> visitado;
		int n;
		
		// CREATE: grafo residual vacio de n nodos.
		FordFulkerson(int n_) {
			n = n_;
			adj.assign(n, {});
		}                                          // O(n)
		
		// WRITE: agrega una arista dirigida u->v con capacidad c. Crea
		// automaticamente la arista reversa con capacidad 0 (se llena a
		// medida que se envia flujo, permitiendo "deshacer" envios).
		void agregarArista(int u, int v, long long c) {
			adj[u].push_back({v, c, (int)adj[v].size()});
			adj[v].push_back({u, 0, (int)adj[u].size() - 1});
		}                                          // O(1)
		
		// busca un camino aumentante con DFS, enviando flujo mientras
		// regresa de la recursion (devuelve cuanto flujo se logro
		// enviar por esa rama, 0 si no encontro camino).
		long long dfs(int u, int t, long long flujoMin) {
			if (u == t) return flujoMin;
			visitado[u] = true;
			
			for (auto& arista : adj[u]) {
				if (visitado[arista.destino] || arista.capacidad <= 0)
				continue;
				
				long long resultado = dfs(arista.destino, t,
				min(flujoMin, arista.capacidad));
				
				if (resultado > 0) {
					arista.capacidad -= resultado;
					adj[arista.destino][arista.reversa].capacidad += resultado;
					return resultado;
				}
			}
			
			return 0;
		}
		
		// READ: flujo maximo de s a t.
		long long flujoMaximo(int s, int t) {
			long long total = 0;
			
			while (true) {
				visitado.assign(n, false);
				long long f = dfs(s, t, LLONG_MAX);
				if (f == 0) break;   // no hay mas caminos aumentantes
				total += f;
			}
			
			return total;
		}                                          // O(E * F)
		
		// ------------------------------------------------------------
		// MODIFICACIONES Y COMENTARIOS CON LOS USOS
		// ------------------------------------------------------------
		
		// 1. FLUJO MAXIMO CON MULTIPLES FUENTES / SUMIDEROS: se agrega
		//    una SUPERFUENTE S y un SUPERSUMIDERO T ficticios; S conecta
		//    a cada fuente real con capacidad infinita (o la capacidad
		//    maxima que esa fuente puede producir), cada sumidero real
		//    conecta a T de la misma forma -- y se corre flujoMaximo(S,T)
		//    normal, sin modificar el algoritmo en si.
		//
		// int S = n, T = n + 1;   // nodos ficticios extra al construir
		// for (int fuente : fuentesReales)
		//     agregarArista(S, fuente, LLONG_MAX);
		// for (int sumidero : sumiderosReales)
		//     agregarArista(sumidero, T, LLONG_MAX);
		// long long total = flujoMaximo(S, T);
		
		// 2. VERTICES CON CAPACIDAD (limite de "flujo que puede pasar
		//    POR un nodo", no solo por sus aristas): se divide cada
		//    nodo u en u_entrada y u_salida, unidos por una arista
		//    u_entrada -> u_salida con capacidad = limite del nodo;
		//    todas las aristas que antes llegaban a u ahora llegan a
		//    u_entrada, y las que salian de u ahora salen de u_salida.
		//    El algoritmo de flujo no cambia, solo la construccion del
		//    grafo (mismo truco que usa 2-SAT con nodo/negacion, pero
		//    aplicado a capacidad en vez de literal).
		
		// 3. CAPACIDADES NO ENTERAS / MUY GRANDES: Ford-Fulkerson con
		//    DFS puro puede degenerar si las capacidades son enormes y
		//    el camino elegido es siempre el "peor" posible (ejemplo
		//    clasico: red con capacidades ~10^9 puede necesitar ~10^9
		//    iteraciones). La solucion NO es cambiar este algoritmo,
		//    sino usar Edmonds-Karp o Dinic (siguientes secciones), que
		//    acotan la complejidad independientemente del valor de F.
		
		// 4. FLUJO MAXIMO CON LIMITE INFERIOR (cada arista debe llevar
		//    AL MENOS un minimo, ademas de un maximo): requiere una
		//    transformacion adicional del grafo (nodos ficticios de
		//    "exceso") antes de correr flujo maximo normal -- fuera del
		//    alcance de esta plantilla base, pero la señal para
		//    reconocerlo es que el enunciado da un RANGO
		//    [minimo, maximo] por arista en vez de solo un maximo.
	};
\end{lstlisting}

\textbf{Ejemplo de funcionamiento:}

\begin{lstlisting}[style=cpp]
	// red:  s -(3)-> a -(2)-> t
	//       s -(2)-> b -(3)-> t
	//       a -(1)-> b
	
	// nodos: s=0, a=1, b=2, t=3
	
	FordFulkerson ff(4);
	ff.agregarArista(0, 1, 3);
	ff.agregarArista(0, 2, 2);
	ff.agregarArista(1, 3, 2);
	ff.agregarArista(2, 3, 3);
	ff.agregarArista(1, 2, 1);
	
	cout << ff.flujoMaximo(0, 3);
	// 4
	// (2 por s-a-t, 2 por s-b-t; la arista a-b sobra en este caso,
	//  el cuello de botella real esta en las aristas de salida de a
	//  y de entrada a t)
\end{lstlisting}

\textbf{Regla práctica:} programa Ford-Fulkerson con DFS cuando las
capacidades son chicas o el límite de tiempo es generoso —es el más
corto de escribir bajo presión. En cuanto el problema tenga
capacidades grandes o el límite de tiempo sea ajustado, salta
directamente a Edmonds-Karp (siguiente sección): la estructura del
grafo residual (\texttt{Arista}, \texttt{agregarArista}) es
\textbf{idéntica}, solo cambia cómo se busca el camino aumentante.


\subsection{Edmonds-Karp}

\textbf{Idea:} Edmonds-Karp es Ford-Fulkerson con \textbf{una sola
	diferencia}: el camino aumentante se busca con \textbf{BFS} en vez de
DFS, lo que garantiza que \textbf{siempre} se use el camino con
\textbf{menos aristas} disponible en el grafo residual. Esa elección
—"el camino más corto primero"— es lo que rompe el peor caso
adversarial de Ford-Fulkerson: se puede demostrar que el número total
de fases de BFS está acotado por $O(V \cdot E)$,
\textbf{independientemente} de qué tan grandes sean las capacidades.

\textbf{¿Por qué usarlo?} Es el primer paso obligado en cuanto
Ford-Fulkerson con DFS "se siente lento" o el problema explícitamente
tiene capacidades grandes ($\ge 10^4$–$10^5$): mismo grafo residual,
mismo concepto de camino aumentante, \textbf{solo} cambia BFS por
DFS, a cambio de una garantía de complejidad que ya no depende del
valor del flujo. Si $O(V \cdot E^2)$ sigue sin alcanzar (grafos
realmente grandes), el siguiente paso es Dinic.

\textbf{Casos de uso:}

\begin{itemize}
	\item Los mismos que Ford-Fulkerson (flujo máximo, matching vía
	flujo, corte mínimo), en cuanto las capacidades son lo bastante
	grandes como para que el peor caso de DFS sea un riesgo real.
	\item Como opción "segura por defecto" cuando no se sabe de
	antemano qué tan adversarial puede ser el grafo de entrada —el
	costo de cambiar DFS por BFS es mínimo y elimina ese riesgo.
\end{itemize}

\textbf{Complejidad:} $O(V \cdot E^2)$ — como mucho $O(V \cdot E)$
fases de BFS (cada una satura al menos una arista del grafo
residual), cada BFS cuesta $O(E)$.

\begin{lstlisting}[style=cpp]
	struct EdmondsKarp {
		struct Arista { int destino; long long capacidad; int reversa; };
		vector<vector<Arista>> adj;
		int n;
		
		// CREATE: grafo residual vacio de n nodos.
		EdmondsKarp(int n_) {
			n = n_;
			adj.assign(n, {});
		}                                          // O(n)
		
		// WRITE: identico a Ford-Fulkerson -- agrega u->v con capacidad
		// c, mas su reversa de capacidad 0.
		void agregarArista(int u, int v, long long c) {
			adj[u].push_back({v, c, (int)adj[v].size()});
			adj[v].push_back({u, 0, (int)adj[u].size() - 1});
		}                                          // O(1)
		
		// busca el camino aumentante MAS CORTO (en numero de aristas)
		// con BFS. Devuelve el flujo enviado por ese camino, 0 si no
		// hay camino aumentante. padreNodo/padreArista reconstruyen el
		// camino encontrado.
		long long bfs(int s, int t, vector<int>& padreNodo,
		vector<int>& padreArista) {
			fill(padreNodo.begin(), padreNodo.end(), -1);
			padreNodo[s] = s;
			
			queue<pair<int,long long>> q;
			q.push({s, LLONG_MAX});
			
			while (!q.empty()) {
				auto [u, flujoHasta] = q.front(); q.pop();
				
				for (int i = 0; i < (int)adj[u].size(); i++) {
					auto& a = adj[u][i];
					
					if (a.capacidad <= 0 || padreNodo[a.destino] != -1)
					continue;
					
					padreNodo[a.destino] = u;
					padreArista[a.destino] = i;
					long long nuevoFlujo = min(flujoHasta, a.capacidad);
					
					if (a.destino == t) return nuevoFlujo;
					
					q.push({a.destino, nuevoFlujo});
				}
			}
			
			return 0;   // t no alcanzable -> no hay camino aumentante
		}                                          // O(E)
		
		// READ: flujo maximo de s a t.
		long long flujoMaximo(int s, int t) {
			long long total = 0;
			vector<int> padreNodo(n), padreArista(n);
			
			while (true) {
				long long f = bfs(s, t, padreNodo, padreArista);
				if (f == 0) break;
				
				// reconstruye el camino desde t hasta s y actualiza
				// capacidades residuales (resta de ida, suma en reversa)
				int u = t;
				while (u != s) {
					int idArista = padreArista[u];
					int padre = padreNodo[u];
					
					adj[padre][idArista].capacidad -= f;
					int idReversa = adj[padre][idArista].reversa;
					adj[u][idReversa].capacidad += f;
					
					u = padre;
				}
				
				total += f;
			}
			
			return total;
		}                                          // O(V * E^2)
		
		// ------------------------------------------------------------
		// MODIFICACIONES Y COMENTARIOS CON LOS USOS
		// ------------------------------------------------------------
		
		// 1. RECONSTRUIR EL CAMINO AUMENTANTE EXPLICITO (no solo su
		//    flujo): ya viene gratis en padreNodo/padreArista tras cada
		//    bfs() exitoso -- util si el problema pide ademas MOSTRAR
		//    las rutas usadas, no solo el valor total del flujo maximo.
		//
		// vector<int> reconstruirCamino(int s, int t,
		//                                 vector<int>& padreNodo) {
			//     vector<int> camino;
			//     for (int u = t; u != s; u = padreNodo[u])
			//         camino.push_back(u);
			//     camino.push_back(s);
			//     reverse(camino.begin(), camino.end());
			//     return camino;
			// }
		
		// 2. MULTIPLES FUENTES / SUMIDEROS: identico truco que en
		//    Ford-Fulkerson -- superfuente S y supersumidero T con
		//    aristas de capacidad infinita hacia/desde cada fuente y
		//    sumidero real; se corre flujoMaximo(S, T) sin cambios.
		//
		// int S = n, T = n + 1;
		// for (int fuente : fuentesReales)
		//     agregarArista(S, fuente, LLONG_MAX);
		// for (int sumidero : sumiderosReales)
		//     agregarArista(sumidero, T, LLONG_MAX);
		
		// 3. DETENERSE ANTES DE flujoMaximo() COMPLETO -- si solo se
		//    necesita saber SI existe flujo >= X (no el maximo exacto),
		//    se puede cortar el while en cuanto total >= X, ahorrando
		//    fases de BFS innecesarias:
		//
		// long long flujoAlMenos(int s, int t, long long X) {
			//     long long total = 0;
			//     vector<int> padreNodo(n), padreArista(n);
			//     while (total < X) {
				//         long long f = bfs(s, t, padreNodo, padreArista);
				//         if (f == 0) break;
				//         // ... misma actualizacion de capacidades ...
				//         total += f;
				//     }
			//     return total;
			// }
		
		// 4. GRAFOS DENSOS O MUY GRANDES (E ~ 10^5 o mas, V grande):
		//    O(V*E^2) de Edmonds-Karp puede seguir sin alcanzar. La
		//    senal para saltar a Dinic (siguiente seccion) es la MISMA
		//    estructura de grafo residual (Arista, agregarArista) mas
		//    una nocion de "grafo por niveles" (BFS de niveles + DFS
		//    con poda) que reduce el numero de fases a O(V), en vez de
		//    O(V*E) -- no hace falta reescribir el grafo, solo la
		//    logica de busqueda de caminos aumentantes.
	};
\end{lstlisting}

\textbf{Ejemplo de funcionamiento:}

\begin{lstlisting}[style=cpp]
	// misma red que en Ford-Fulkerson:
	// s -(3)-> a -(2)-> t
	// s -(2)-> b -(3)-> t
	// a -(1)-> b
	
	EdmondsKarp ek(4);
	ek.agregarArista(0, 1, 3);
	ek.agregarArista(0, 2, 2);
	ek.agregarArista(1, 3, 2);
	ek.agregarArista(2, 3, 3);
	ek.agregarArista(1, 2, 1);
	
	cout << ek.flujoMaximo(0, 3);
	// 4   (mismo resultado que Ford-Fulkerson -- el flujo maximo es
	//      unico independientemente del orden de caminos elegido;
	//      lo que cambia es CUANTAS iteraciones toma llegar ahi)
\end{lstlisting}

\textbf{Regla práctica:} Edmonds-Karp y Ford-Fulkerson dan
\textbf{siempre} el mismo flujo máximo —difieren solo en garantía de
complejidad, nunca en corrección. Si ya tienes el \texttt{struct} de
Ford-Fulkerson escrito, migrar a Edmonds-Karp es cambiar
\texttt{dfs()} por \texttt{bfs()} y ajustar cómo se reconstruye el
camino para actualizar capacidades —el resto (\texttt{Arista},
\texttt{agregarArista}, la idea de arista reversa) no cambia en
absoluto.

\subsection{Dinic}

\textbf{Idea:} Dinic mejora a Edmonds-Karp separando el trabajo en
\textbf{fases}: en cada fase, (1) construye un \textbf{grafo por
	niveles} con un solo BFS desde $s$ (\texttt{nivel[u]} = distancia
mínima desde $s$ en el grafo residual actual), quedándose
\textbf{solo} con las aristas que avanzan de un nivel al siguiente
($\texttt{nivel}[v] = \texttt{nivel}[u]+1$); (2) sobre ese grafo por
niveles, encuentra \textbf{todos} los caminos aumentantes posibles de
una sola vez con DFS, usando un \textbf{puntero por nodo}
(\texttt{ptr[u]}, igual que en Hierholzer) para no reexplorar aristas
ya agotadas dentro de la misma fase —esto se llama \textbf{blocking
	flow}. Cuando el BFS ya no puede alcanzar $t$ (no hay más niveles), el
algoritmo termina: el flujo acumulado es el máximo.

\textbf{¿Por qué usarlo?} Edmonds-Karp encuentra \textbf{un} camino
aumentante por cada BFS ($O(V \cdot E)$ fases en el peor caso); Dinic
encuentra \textbf{muchos} caminos aumentantes por cada BFS (todo el
blocking flow de esa fase), reduciendo el número de fases a
$O(V)$ —cada fase sube al menos un nivel, y hay como mucho $V$ niveles
posibles. Es la opción estándar en competencia quando $E$ es grande y
$O(V \cdot E^2)$ de Edmonds-Karp no alcanza el límite de tiempo; en
grafos bipartitos (unit-capacity), Dinic corre incluso más rápido
($O(E\sqrt{V})$), lo cual lo hace también la base preferida para
Maximum Bipartite Matching vía flujo (siguiente sección).

\textbf{Casos de uso:}

\begin{itemize}
	\item Los mismos que Ford-Fulkerson/Edmonds-Karp, en instancias
	grandes donde $O(V \cdot E^2)$ es demasiado lento.
	\item Bipartite Matching vía flujo, aprovechando el caso especial
	de $O(E\sqrt{V})$ en grafos de capacidad unitaria.
	\item Cualquier problema de flujo donde el límite de tiempo sea
	ajustado y $V, E$ sean del orden de $10^4$–$10^5$.
\end{itemize}

\textbf{Complejidad:} $O(V^2 \cdot E)$ en general, $O(E\sqrt{V})$ en
grafos de capacidad unitaria (bipartite matching) — en ambos casos,
estrictamente mejor o igual que Edmonds-Karp.

\begin{lstlisting}[style=cpp]
	struct Dinic {
		struct Arista { int destino; long long capacidad; int reversa; };
		vector<vector<Arista>> adj;
		vector<int> nivel, ptr;
		int n;
		
		// CREATE: grafo residual vacio de n nodos.
		Dinic(int n_) {
			n = n_;
			adj.assign(n, {});
		}                                          // O(n)
		
		// WRITE: identico a Ford-Fulkerson/Edmonds-Karp -- agrega u->v
		// con capacidad c, mas su reversa de capacidad 0.
		void agregarArista(int u, int v, long long c) {
			adj[u].push_back({v, c, (int)adj[v].size()});
			adj[v].push_back({u, 0, (int)adj[u].size() - 1});
		}                                          // O(1)
		
		// FASE 1: construye el grafo por niveles con BFS. Devuelve
		// false si t ya no es alcanzable (fin del algoritmo).
		bool bfs(int s, int t) {
			nivel.assign(n, -1);
			queue<int> q;
			q.push(s);
			nivel[s] = 0;
			
			while (!q.empty()) {
				int u = q.front(); q.pop();
				
				for (auto& a : adj[u]) {
					if (a.capacidad > 0 && nivel[a.destino] == -1) {
						nivel[a.destino] = nivel[u] + 1;
						q.push(a.destino);
					}
				}
			}
			
			return nivel[t] != -1;
		}                                          // O(E)
		
		// FASE 2: DFS con punteros (blocking flow), buscando UN camino
		// aumentante que respete el grafo por niveles; se llama
		// repetidamente hasta agotar todos los caminos de esta fase.
		long long dfs(int u, int t, long long flujoMin) {
			if (u == t) return flujoMin;
			
			// avanza el puntero saltando aristas ya agotadas o que no
			// avanzan de nivel -- cada posicion se salta como mucho una
			// vez POR FASE, de ahi la eficiencia del blocking flow
			for (; ptr[u] < (int)adj[u].size(); ptr[u]++) {
				auto& a = adj[u][ptr[u]];
				
				if (a.capacidad <= 0 || nivel[a.destino] != nivel[u] + 1)
				continue;
				
				long long resultado = dfs(a.destino, t,
				min(flujoMin, a.capacidad));
				
				if (resultado > 0) {
					a.capacidad -= resultado;
					adj[a.destino][a.reversa].capacidad += resultado;
					return resultado;
				}
			}
			
			return 0;   // u agotado en esta fase, ningun camino desde aqui
		}                                          // O(V) amortizado
		// dentro de una fase
		
		// READ: flujo maximo de s a t.
		long long flujoMaximo(int s, int t) {
			long long total = 0;
			
			while (bfs(s, t)) {
				ptr.assign(n, 0);
				
				while (long long f = dfs(s, t, LLONG_MAX))
				total += f;   // sigue sacando caminos de la MISMA
				// fase hasta agotar el blocking flow
			}
			
			return total;
		}                                          // O(V^2 * E)
		
		// ------------------------------------------------------------
		// MODIFICACIONES Y COMENTARIOS CON LOS USOS
		// ------------------------------------------------------------
		
		// 1. BIPARTITE MATCHING VIA DINIC (capacidad unitaria): sin
		//    cambiar el algoritmo, si TODAS las aristas tienen capacidad
		//    1 (fuente -> L, L -> R segun compatibilidad, R -> sumidero),
		//    la complejidad baja automaticamente a O(E*sqrt(V)) -- no
		//    hace falta ningun ajuste al codigo, es una propiedad del
		//    caso de capacidad unitaria. Ver Maximum Bipartite Matching.
		//
		// int S = n, T = n + 1;   // agregar como nodos extra
		// for (int u : L) agregarArista(S, u, 1);
		// for (int v : R) agregarArista(T + n, v, 1);  // ajustar indices
		// for (auto [u, v] : compatibles) agregarArista(u, v, 1);
		// long long matching = flujoMaximo(S, T);
		
		// 2. MULTIPLES FUENTES / SUMIDEROS: mismo truco de superfuente
		//    S / supersumidero T que en Ford-Fulkerson y Edmonds-Karp,
		//    sin cambios al algoritmo:
		//
		// int S = n, T = n + 1;
		// for (int fuente : fuentesReales)
		//     agregarArista(S, fuente, LLONG_MAX);
		// for (int sumidero : sumiderosReales)
		//     agregarArista(sumidero, T, LLONG_MAX);
		
		// 3. RECUPERAR EL CORTE MINIMO tras flujoMaximo(): igual que en
		//    Ford-Fulkerson -- el lado S del corte es el conjunto de
		//    nodos con nivel[] != -1 EN LA ULTIMA LLAMADA A bfs() que
		//    devolvio false (esa es la ultima vez que se calcularon
		//    niveles antes de que t dejara de ser alcanzable):
		//
		// vector<bool> ladoS() {
			//     vector<bool> resultado(n, false);
			//     for (int u = 0; u < n; u++)
			//         resultado[u] = (nivel[u] != -1);
			//     return resultado;
			// }
		// -- llamar DESPUES de flujoMaximo(); nivel[] queda de la
		//    ultima corrida de bfs() (la que fallo en alcanzar t)
		
		// 4. SCALING (capacidades MUY grandes, V y E tambien grandes):
		//    variante avanzada que procesa las aristas en orden
		//    decreciente de "bits de capacidad" para acotar aun mas las
		//    fases -- rara vez necesaria en competencia, Dinic simple ya
		//    cubre la gran mayoria de limites de tiempo tipicos.
		
		// 5. GRAFOS CON ARISTAS NO DIRIGIDAS: se agregan como DOS
		//    aristas dirigidas independientes, cada una con su PROPIA
		//    reversa de capacidad 0 (no se comparte una sola reversa
		//    entre ambos sentidos, a diferencia de Hierholzer con su
		//    id compartido) -- Dinic no necesita ningun ajuste extra
		//    para esto, es solo como se cargan las aristas:
		//
		// void agregarAristaNoDirigida(int u, int v, long long c) {
			//     agregarArista(u, v, c);
			//     agregarArista(v, u, c);
			// }
	};
\end{lstlisting}

\textbf{Ejemplo de funcionamiento:}

\begin{lstlisting}[style=cpp]
	// misma red que en Ford-Fulkerson y Edmonds-Karp:
	// s -(3)-> a -(2)-> t
	// s -(2)-> b -(3)-> t
	// a -(1)-> b
	
	Dinic din(4);
	din.agregarArista(0, 1, 3);
	din.agregarArista(0, 2, 2);
	din.agregarArista(1, 3, 2);
	din.agregarArista(2, 3, 3);
	din.agregarArista(1, 2, 1);
	
	cout << din.flujoMaximo(0, 3);
	// 4   (mismo flujo maximo que las dos versiones anteriores --
	//      Dinic solo cambia CUANTAS fases/iteraciones toma llegar,
	//      no el resultado)
\end{lstlisting}

\textbf{Regla práctica:} las tres estructuras de esta sección
(Ford-Fulkerson, Edmonds-Karp, Dinic) comparten \textbf{exactamente}
la misma representación de grafo residual (\texttt{Arista},
\texttt{agregarArista}) y siempre dan el mismo flujo máximo — la
progresión natural es empezar con Ford-Fulkerson por simplicidad,
subir a Edmonds-Karp si las capacidades son grandes, y solo llegar a
Dinic cuando $E$ es grande y $O(V \cdot E^2)$ no alcanza el límite de
tiempo. Si el problema es específicamente bipartite matching, Dinic
es casi siempre la mejor opción de entrada por su complejidad especial
en capacidad unitaria.

\subsection{Maximum Bipartite Matching}

\textbf{Idea:} Bipartite Matching (ya visto con Kuhn en la sección de
Bipartite Graphs) es un caso particular de flujo máximo: se construye
una red con una \textbf{superfuente} $S$ conectada a cada nodo de $L$
(capacidad 1), cada nodo de $L$ conectado a los nodos de $R$ con los
que es compatible (capacidad 1), y cada nodo de $R$ conectado a un
\textbf{supersumidero} $T$ (capacidad 1). Como \textbf{todas} las
capacidades son 1, cualquier unidad de flujo de $S$ a $T$ corresponde
exactamente a "un nodo de $L$ emparejado con un nodo de $R$" —el
flujo máximo de $S$ a $T$ es, literalmente, el tamaño del matching
máximo (y no puede ser mayor porque cada nodo de $L$/$R$ tiene
capacidad de entrada/salida 1, así que participa en como mucho un
emparejamiento).

\textbf{¿Por qué usarlo?} Es la alternativa a Kuhn cuando ya tienes
infraestructura de flujo (Dinic) escrita y prefieres reutilizarla en
vez de mantener dos implementaciones separadas —además, en el caso
particular de capacidad unitaria, Dinic corre en $O(E\sqrt{V})$,
\textbf{mejor} que el $O(V \cdot E)$ de Kuhn simple. Si $V$ es
pequeño ($\lesssim 1000$–$2000$), Kuhn sigue siendo más simple de
escribir bajo presión; si $V$ es grande, Dinic + esta construcción es
la opción más rápida disponible en esta guía.

\textbf{Casos de uso:}

\begin{itemize}
	\item Los mismos que Bipartite Matching (asignación óptima
	uno-a-uno entre dos conjuntos), cuando $V$ es lo bastante grande
	para que $O(E\sqrt{V})$ de Dinic gane frente a $O(V \cdot E)$ de
	Kuhn.
	\item Como building block de \textbf{matching con capacidades}
	—variante donde cada nodo de $L$ o $R$ puede emparejarse con
	\textbf{más de uno} del otro lado (cambiando la capacidad de las
	aristas $S \to L$ o $R \to T$ de 1 a ese límite, ver
	modificaciones abajo).
	\item Cuando el problema \textbf{ya} pide flujo máximo por otra
	razón y el matching bipartito es solo una parte del modelo
	completo —reutilizar la misma estructura de Dinic evita duplicar
	código.
\end{itemize}

\textbf{Complejidad:} $O(E\sqrt{V})$ con Dinic (gracias a la
propiedad especial de grafos de capacidad unitaria), frente a
$O(V \cdot E)$ de Kuhn — construcción de la red en $O(V + E)$.

\begin{lstlisting}[style=cpp]
	struct MaximumBipartiteMatching {
		Dinic dinic;
		int nL, nR, S, T;
		
		// CREATE: construye la red de flujo para matching bipartito.
		// nL, nR = tamanos de los lados L y R. compatibles[u] = lista de
		// nodos de R compatibles con el nodo u de L.
		MaximumBipartiteMatching(int nL_, int nR_,
		vector<vector<int>>& compatibles)
		: dinic(nL_ + nR_ + 2), nL(nL_), nR(nR_) {
			S = nL + nR;       // superfuente
			T = nL + nR + 1;   // supersumidero
			
			for (int u = 0; u < nL; u++)
			dinic.agregarArista(S, u, 1);
			
			for (int v = 0; v < nR; v++)
			dinic.agregarArista(nL + v, T, 1);
			
			for (int u = 0; u < nL; u++)
			for (int v : compatibles[u])
			dinic.agregarArista(u, nL + v, 1);
		}                                          // O(V + E)
		
		// READ: tamano del matching maximo.
		long long matchingMaximo() {
			return dinic.flujoMaximo(S, T);
		}                                          // O(E * sqrt(V))
		
		// READ: reconstruye la asignacion explicita tras calcular el
		// matching maximo -- recorre las aristas u->v (u en L) buscando
		// cuales quedaron con capacidad residual 0 (es decir, con 1
		// unidad de flujo pasando por ellas).
		vector<int> obtenerAsignacion() {
			vector<int> pareja(nL, -1);
			
			for (int u = 0; u < nL; u++)
			for (auto& a : dinic.adj[u])
			if (a.destino >= nL && a.destino < nL + nR &&
			a.capacidad == 0)   // capacidad original era 1,
			// si quedo en 0 es que se uso
			pareja[u] = a.destino - nL;
			
			return pareja;
		}                                          // O(E)
		
		// ------------------------------------------------------------
		// MODIFICACIONES Y COMENTARIOS CON LOS USOS
		// ------------------------------------------------------------
		
		// 1. MATCHING CON CAPACIDADES (cada nodo de L o R puede
		//    emparejarse con MAS de uno del otro lado, hasta un limite):
		//    en vez de capacidad 1 en las aristas S->L y R->T, se usa el
		//    limite real de cada nodo -- el resto de la construccion
		//    (L->R sigue siendo capacidad 1, o el limite de esa pareja
		//    especifica si aplica) no cambia:
		//
		// for (int u = 0; u < nL; u++)
		//     dinic.agregarArista(S, u, limiteL[u]);   // en vez de 1
		// for (int v = 0; v < nR; v++)
		//     dinic.agregarArista(nL + v, T, limiteR[v]);   // en vez de 1
		
		// 2. MATCHING PONDERADO MAXIMO (cada arista L-R tiene un peso, y
		//    se busca el matching de mayor peso total, no solo el de
		//    mayor CANTIDAD de parejas): flujo maximo NO resuelve esto
		//    directamente -- hace falta flujo de COSTO minimo/maximo
		//    (min-cost max-flow), fuera del alcance de esta plantilla
		//    base. La señal para reconocerlo es que el enunciado da un
		//    PESO o COSTO por pareja posible, no solo compatibilidad
		//    binaria.
		
		// 3. VERIFICAR SI EL MATCHING ES PERFECTO (todos los nodos de L
		//    Y todos los de R quedan emparejados): matching es perfecto
		//    si y solo si matchingMaximo() == nL == nR (mismo tamano en
		//    ambos lados y todos emparejados):
		//
		// bool esPerfecto() {
			//     return matchingMaximo() == nL && nL == nR;
			// }
		
		// 4. HALLAR UN CONJUNTO DE VERTICES NO CUBIERTOS (nodos de L o
		//    R que quedaron SIN pareja tras el matching maximo) -- util
		//    para problemas que piden identificar el "excedente":
		//
		// vector<int> noEmparejadosEnL() {
			//     auto pareja = obtenerAsignacion();
			//     vector<int> resultado;
			//     for (int u = 0; u < nL; u++)
			//         if (pareja[u] == -1) resultado.push_back(u);
			//     return resultado;
			// }
		
		// 5. TEOREMA DE KONIG (cobertura minima de vertices = matching
		//    maximo en grafos bipartitos): la cobertura minima se
		//    obtiene a partir del corte minimo de esta MISMA red de
		//    flujo -- ver Min-Cut para la extraccion del lado S del
		//    corte; los nodos de L NO alcanzables desde S en el residual
		//    final, mas los nodos de R SI alcanzables, forman la
		//    cobertura minima de vertices.
	};
\end{lstlisting}

\textbf{Ejemplo de funcionamiento:}

\begin{lstlisting}[style=cpp]
	// L = {0, 1, 2}, R = {0, 1, 2}
	// 0 - {0, 1}
	// 1 - {0}
	// 2 - {1, 2}
	
	vector<vector<int>> compatibles(3);
	compatibles[0] = {0, 1};
	compatibles[1] = {0};
	compatibles[2] = {1, 2};
	
	MaximumBipartiteMatching mbm(3, 3, compatibles);
	
	cout << mbm.matchingMaximo();
	// 3   (mismo resultado que el ejemplo de Kuhn en Bipartite Matching)
	
	auto pareja = mbm.obtenerAsignacion();
	for (int u = 0; u < 3; u++)
	cout << u << "-" << pareja[u] << " ";
	// alguna asignacion valida, ej: 0-1 1-0 2-2
\end{lstlisting}

\textbf{Regla práctica:} Kuhn y Maximum Bipartite Matching vía Dinic
\textbf{siempre} dan el mismo tamaño de matching máximo (ambos
resuelven exactamente el mismo problema) — la elección es de
rendimiento y de qué código ya tienes a mano: con $V$ pequeño y sin
Dinic ya escrito, Kuhn es más corto; con $V$ grande o Dinic ya
disponible de otra parte del problema, esta construcción vía flujo es
la opción más rápida y la que se generaliza directo a matching con
capacidades.


% =========================================================
% 2-SAT
% =========================================================

\section{2-SAT}

\subsection{Implication Graph}

\textbf{Idea:} 2-SAT pide encontrar una asignación de verdad para $n$
variables booleanas que satisfaga una fórmula compuesta por
\textbf{cláusulas de a lo más 2 literales} (ej. $(A \lor B)$,
$(\lnot A \lor C)$). La clave para resolverlo en tiempo lineal (en
vez de probar $2^n$ asignaciones) es notar que cada cláusula
$(a \lor b)$ es \textbf{lógicamente equivalente} a dos implicaciones:
$\lnot a \Rightarrow b$ y $\lnot b \Rightarrow a$ (si uno de los dos
literales es falso, el otro \textbf{debe} ser verdadero). Modelando
cada literal (variable y su negación) como un nodo de un grafo
dirigido, y cada implicación como una arista, la fórmula completa se
convierte en un \textbf{grafo de implicación} de $2n$ nodos: el nodo
$2i$ representa "$x_i$ es verdadero", el nodo $2i{+}1$ representa
"$x_i$ es falso" (o viceversa, según la convención elegida).

\textbf{¿Por qué usarlo?} Es el paso de \textbf{modelado} que
convierte un problema de satisfacibilidad booleana en un problema de
grafos puro —una vez construido el grafo de implicación, resolver
2-SAT deja de requerir ningún razonamiento lógico adicional: se
reduce enteramente a encontrar Componentes Fuertemente Conexas (ver
2-SAT con SCC, siguiente subsección). La parte que sí exige cuidado
es traducir \textbf{correctamente} cada tipo de restricción a su par
de implicaciones —un error aquí produce un grafo que no representa la
fórmula real, aunque el algoritmo de SCC corra sin errores.

\textbf{Casos de uso:} construir el grafo de implicación es el primer
paso \textbf{obligatorio} de cualquier problema de 2-SAT, sin
importar cuál sea la restricción concreta. Los patrones de traducción
más comunes (cada uno se agrega como \texttt{add\_disjunction} sobre
la estructura de 2-SAT con SCC):

\begin{itemize}
	\item Forzar una variable a verdadero o falso (restricción
	unaria, disfrazada de cláusula con el mismo literal repetido).
	\item Implicación simple $A \Rightarrow B$ (cláusula
	$(\lnot A \lor B)$).
	\item "Al menos uno" de dos literales (cláusula OR directa).
	\item "A lo más uno" / mutuamente excluyentes (cláusula NAND).
	\item Equivalencia $A \leftrightarrow B$ (dos cláusulas: ambas
	direcciones de la implicación).
	\item Diferencia forzada $A \oplus B$ (dos cláusulas: al menos
	uno, y no ambos).
\end{itemize}

\textbf{Complejidad:} $O(1)$ por cláusula agregada al grafo (cada
cláusula agrega una cantidad constante de nodos/aristas), $O(V + E) =
O(n + m)$ para el grafo completo con $n$ variables y $m$ cláusulas.

\begin{lstlisting}[style=cpp]
	// convencion de esta plantilla: variable i se representa con los
	// nodos 2*i (literal POSITIVO: x_i) y 2*i+1 (literal NEGATIVO: !x_i)
	
	struct GrafoImplicacion {
		int nVars, nVertices;
		vector<vector<int>> adj;   // grafo de implicacion (dirigido)
		
		// CREATE: n variables booleanas -> 2n nodos en el grafo.
		GrafoImplicacion(int nVars_)
		: nVars(nVars_), nVertices(2 * nVars_), adj(2 * nVars_) {}
		// O(n)
		
		// WRITE: agrega la clausula (literal_a OR literal_b), donde
		// na/nb indican si ese literal esta NEGADO. Traduce la clausula
		// a sus dos implicaciones equivalentes:
		//   !a -> b   y   !b -> a
		void agregarClausula(int a, bool na, int b, bool nb) {
			int litA = 2 * a ^ na;   // nodo del literal a (negado o no)
			int litB = 2 * b ^ nb;
			int negA = litA ^ 1;     // nodo del literal OPUESTO a litA
			int negB = litB ^ 1;
			
			adj[negA].push_back(litB);   // !a -> b
			adj[negB].push_back(litA);   // !b -> a
		}                                          // O(1)
	};
\end{lstlisting}

\textbf{Ejemplo de funcionamiento:}

\begin{lstlisting}[style=cpp]
	// formula: (x0 OR x1) AND (!x0 OR x2)
	// 2 clausulas, 3 variables (nodos 0..5:
	//   0=x0, 1=!x0, 2=x1, 3=!x1, 4=x2, 5=!x2)
	
	GrafoImplicacion gi(3);
	
	gi.agregarClausula(0, false, 1, false);   // (x0 OR x1)
	gi.agregarClausula(0, true,  2, false);   // (!x0 OR x2)
	
	// aristas generadas:
	// clausula 1: !x0 -> x1  (1 -> 2),  !x1 -> x0  (3 -> 0)
	// clausula 2: x0 -> x2   (0 -> 4),  !x2 -> !x0 (5 -> 1)
	
	for (int v : gi.adj[1]) cout << v << " ";
	// 2   (!x0 -> x1)
	
	for (int v : gi.adj[0]) cout << v << " ";
	// 4   (x0 -> x2)
\end{lstlisting}

\textbf{Regla práctica:} sin importar qué tan compleja parezca la
restricción original (implicación, equivalencia, exclusión mutua,
etc.), siempre se puede descomponer en una o dos cláusulas OR de 2
literales, y cada cláusula siempre agrega exactamente 2 aristas al
grafo de implicación con la misma fórmula
$\lnot a \Rightarrow b,\ \lnot b \Rightarrow a$ —memorizar esa única
fórmula es más robusto que memorizar un caso especial por cada tipo
de restricción.


\subsection{2-SAT con SCC}

\textbf{Idea:} una vez construido el grafo de implicación, la fórmula
es \textbf{satisfacible} si y solo si \textbf{ninguna} variable $x_i$
tiene sus dos literales ($x_i$ y $\lnot x_i$) en la \textbf{misma}
componente fuertemente conexa —si estuvieran en la misma SCC,
significaría que $x_i \Rightarrow \lnot x_i$ \textbf{y}
$\lnot x_i \Rightarrow x_i$ simultáneamente, una contradicción
directa. Si la fórmula es satisfacible, la asignación se obtiene
comparando el \textbf{orden topológico inverso} de las SCCs
(que Tarjan/Kosaraju ya dan como efecto secundario, ver esas
secciones): a cada variable se le asigna \textbf{verdadero} si su
literal positivo queda en una SCC \textbf{más tarde} en ese orden
(equivalentemente, con \texttt{comp} más alto si se usa el
\texttt{id} de Kosaraju/Tarjan de esta plantilla).

\textbf{¿Por qué usarlo?} Es la razón de ser de todo el modelado con
Implication Graph: convierte "resolver una fórmula booleana" en "una
sola pasada de SCC", sin backtracking ni búsqueda exponencial. Es
además \textbf{constructivo} —no solo dice si existe solución, sino
que la construye directamente comparando los ids de componente, sin
ningún paso adicional de búsqueda.

\textbf{Casos de uso:}

\begin{itemize}
	\item Cualquier problema modelable como restricciones binarias
	entre variables booleanas: asignación de turnos con conflictos,
	selección de una opción entre pares incompatibles, coloreo con
	solo 2 colores sujeto a restricciones.
	\item Verificar la \textbf{existencia} de una asignación válida
	sin necesitar construirla explícitamente (basta el chequeo de
	"misma SCC").
	\item Como building block de problemas más grandes que reducen su
	núcleo de decisión a 2-SAT (ej. "elige un valor de dos opciones
	por cada elemento, sujeto a restricciones de pares").
\end{itemize}

\textbf{Complejidad:} $O(n + m)$ — construir el grafo es $O(m)$
(una cláusula, trabajo constante), Kosaraju/Tarjan sobre el grafo de
$2n$ nodos es $O(n+m)$, y la asignación final es $O(n)$.

\begin{lstlisting}[style=cpp]
	struct TwoSatSolver {
		int n_vars;
		int n_vertices;
		vector<vector<int> > adj, adj_t;
		vector<bool> used;
		vector<int> order, comp;
		vector<bool> assignment;
		
		// CREATE: n_vars variables -> 2*n_vars nodos (grafo de
		// implicacion + su transpuesto, para Kosaraju).
		TwoSatSolver(int _n_vars) : n_vars(_n_vars), n_vertices(2 * n_vars),
		adj(n_vertices), adj_t(n_vertices),
		used(n_vertices), order(),
		comp(n_vertices, -1), assignment(n_vars) {
			order.reserve(n_vertices);
		}                                          // O(n)
		
		void dfs1(int v) {
			used[v] = true;
			for (int u: adj[v])
			if (!used[u])
			dfs1(u);
			order.push_back(v);
		}
		
		void dfs2(int v, int cl) {
			comp[v] = cl;
			for (int u: adj_t[v])
			if (comp[u] == -1)
			dfs2(u, cl);
		}
		
		// READ: resuelve la formula (Kosaraju: dfs1 sobre el grafo
		// original, dfs2 sobre el transpuesto). Devuelve false si NO es
		// satisfacible (algun x_i y !x_i quedaron en la misma SCC).
		bool solve_2SAT() {
			order.clear();
			used.assign(n_vertices, false);
			for (int i = 0; i < n_vertices; ++i)
			if (!used[i])
			dfs1(i);
			
			comp.assign(n_vertices, -1);
			for (int i = 0, j = 0; i < n_vertices; ++i) {
				int v = order[n_vertices - i - 1];
				if (comp[v] == -1)
				dfs2(v, j++);
			}
			
			assignment.assign(n_vars, false);
			for (int i = 0; i < n_vertices; i += 2) {
				if (comp[i] == comp[i + 1])
				return false;   // x_i y !x_i en la misma SCC ->
				// contradiccion, no satisfacible
				
				// comp[] queda en orden topologico INVERSO (Kosaraju
				// procesa por orden de salida decreciente); el literal
				// con comp MAYOR aparece "despues" en ese orden inverso,
				// asi que es el que se elige como verdadero
				assignment[i / 2] = comp[i] > comp[i + 1];
			}
			
			return true;
		}                                          // O(n + m)
		
		// WRITE: agrega la clausula (literal a con negacion na) OR
		// (literal b con negacion nb). Ver Implication Graph para la
		// deduccion de estas dos aristas.
		void add_disjunction(int a, bool na, int b, bool nb) {
			a = 2 * a ^ na;
			b = 2 * b ^ nb;
			int neg_a = a ^ 1;
			int neg_b = b ^ 1;
			adj[neg_a].push_back(b);
			adj[neg_b].push_back(a);
			adj_t[b].push_back(neg_a);
			adj_t[a].push_back(neg_b);
		}                                          // O(1)
		
		// ------------------------------------------------------------
		// PATRONES 2-SAT COMPLETOS (modificaciones y usos)
		// ------------------------------------------------------------
		
		// 1. A debe ser verdadero (forzar A = true)
		//    Clausula: (A OR A), equivale a !A => A
		// solver.add_disjunction(i, false, i, false);
		
		// 2. A debe ser falso (forzar A = false)
		//    Clausula: (!A OR !A), equivale a A => !A
		// solver.add_disjunction(i, true, i, true);
		
		// 3. Si A entonces B (A => B)
		//    Clausula: (!A OR B), equivale a A => B, !B => !A
		// solver.add_disjunction(i, true, j, false);
		
		// 4. Si no A entonces no B (!A => !B)
		//    Clausula: (A OR !B), equivale a !A => !B, B => A
		// solver.add_disjunction(i, false, j, true);
		
		// 5. Al menos uno verdadero (A OR B)
		//    Clausula: (A OR B), equivale a !A => B, !B => A
		// solver.add_disjunction(i, false, j, false);
		
		// 6. A lo mas uno verdadero / mutuamente excluyentes (!A OR !B)
		//    Clausula: (!A OR !B), equivale a A => !B, B => !A
		// solver.add_disjunction(i, true, j, true);
		
		// 7. A y B deben ser iguales (A <-> B)
		//    Clausulas: (A OR !B) y (!A OR B)
		// solver.add_disjunction(i, false, j, true);   // B => A
		// solver.add_disjunction(i, true, j, false);   // A => B
		
		// 8. A y B deben ser diferentes (A XOR B)
		//    Clausulas: (A OR B) y (!A OR !B)
		// solver.add_disjunction(i, false, j, false);  // !A=>B, !B=>A
		// solver.add_disjunction(i, true, j, true);    // A=>!B, B=>!A
	};
\end{lstlisting}

\textbf{Ejemplo de funcionamiento:}

\begin{lstlisting}[style=cpp]
	// formula: (x0 OR x1) AND (!x0 OR x2) AND (x1 debe ser verdadero)
	
	TwoSatSolver solver(3);   // variables 0, 1, 2
	
	solver.add_disjunction(0, false, 1, false);   // (x0 OR x1)
	solver.add_disjunction(0, true,  2, false);   // (!x0 OR x2)
	solver.add_disjunction(1, false, 1, false);   // forzar x1 = true
	
	bool satisfacible = solver.solve_2SAT();
	
	cout << satisfacible;
	// true
	
	cout << solver.assignment[1];
	// true   (forzado explicitamente)
\end{lstlisting}

\textbf{Regla práctica:} el chequeo de satisfacibilidad
(\texttt{comp[i] == comp[i+1]}) y la construcción de la asignación
(\texttt{comp[i] > comp[i+1]}) son \textbf{inseparables} del mismo
recorrido de SCC —no hace falta ninguna estructura ni pasada
adicional para obtener la asignación una vez que la fórmula resultó
satisfacible; toda la dificultad real de 2-SAT está en el paso de
Implication Graph (traducir correctamente cada restricción), no en
esta parte, que es siempre el mismo código sin importar el problema.


% =========================================================
% BACKTRACKING Y PROBLEMAS CLÁSICOS
% =========================================================

\section{Backtracking y Problemas Clásicos}

\subsection{\texorpdfstring{Problema de las $N$ Reinas}{Problema de las N Reinas}}

\textbf{Idea:} colocar $N$ reinas en un tablero $N \times N$ tal que
\textbf{ninguna} ataque a otra (ni misma fila, ni misma columna, ni
ninguna diagonal). Se resuelve con backtracking colocando
\textbf{una reina por fila} (eso ya garantiza que nunca compartan
fila): en cada fila, se prueba cada columna disponible, se verifica
que no choque con las reinas ya colocadas, se coloca y se avanza a la
siguiente fila; si ninguna columna de la fila actual funciona, se
retrocede ("backtrack") y se prueba la siguiente opción de la fila
anterior. La clave de eficiencia es \textbf{podar temprano}: en vez de
recalcular columnas/diagonales ocupadas recorriendo todas las reinas
ya puestas, se mantienen tres arreglos booleanos (columna, diagonal
\texttt{fila+columna}, diagonal \texttt{fila-columna}) que se marcan
y desmarcan en $O(1)$.

\textbf{¿Por qué usarlo?} Es el ejemplo canónico de backtracking con
poda por restricciones: sin las tres marcas de columna/diagonales, la
verificación de "¿esta posición es válida?" costaría $O(N)$ por
candidato (recorrer todas las reinas puestas), multiplicando el costo
total. Con las marcas, cada verificación es $O(1)$, dejando que la
poda temprana (cortar ramas inválidas apenas se detectan) sea lo que
realmente controla el tamaño del árbol de búsqueda explorado.

\textbf{Casos de uso:}

\begin{itemize}
	\item Encontrar \textbf{una} solución válida para $N$ reinas.
	\item Contar \textbf{cuántas} soluciones distintas existen para
	un $N$ dado (variante más común en jueces de práctica).
	\item Listar/imprimir \textbf{todas} las soluciones explícitas
	(cada una como una permutación de columnas, una por fila).
	\item Como plantilla base de backtracking con restricciones
	posicionales, reutilizable para variantes con obstáculos u otras
	piezas de ajedrez (ver Variaciones, siguiente subsección).
\end{itemize}

\textbf{Complejidad:} exponencial en el peor caso ($O(N!)$
aproximado sin poda), pero la poda por columna/diagonales en $O(1)$
hace que en la práctica sea manejable hasta $N \sim 15$–$20$ en
tiempo de competencia.

\begin{lstlisting}[style=cpp]
	struct NReinas {
		int n;
		vector<int> columnaPorFila;   // columnaPorFila[fila] = columna
		// donde se puso la reina de esa fila
		vector<bool> colOcupada, diag1Ocupada, diag2Ocupada;
		long long soluciones;
		
		// CREATE: tablero n x n vacio.
		NReinas(int n_) {
			n = n_;
			columnaPorFila.assign(n, -1);
			colOcupada.assign(n, false);
			diag1Ocupada.assign(2 * n, false);   // diagonal fila+columna
			diag2Ocupada.assign(2 * n, false);   // diagonal fila-columna+n
			soluciones = 0;
		}                                          // O(n)
		
		// READ: cuenta todas las soluciones validas via backtracking.
		long long contarSoluciones() {
			soluciones = 0;
			backtrack(0);
			return soluciones;
		}                                          // O(N!) peor caso,
		// mucho menor en
		// la practica con poda
		
		void backtrack(int fila) {
			if (fila == n) {
				soluciones++;   // una reina por fila, todas colocadas
				// sin conflicto -> solucion completa
				return;
			}
			
			for (int col = 0; col < n; col++) {
				int d1 = fila + col;
				int d2 = fila - col + n;
				
				if (colOcupada[col] || diag1Ocupada[d1] || diag2Ocupada[d2])
				continue;   // poda: esta columna/diagonal ya esta
				// atacada, no vale la pena seguir
				
				// coloca la reina
				columnaPorFila[fila] = col;
				colOcupada[col] = diag1Ocupada[d1] = diag2Ocupada[d2] = true;
				
				backtrack(fila + 1);
				
				// deshace ("backtrack" real)
				colOcupada[col] = diag1Ocupada[d1] = diag2Ocupada[d2] = false;
			}
		}
		
		// READ: encuentra UNA solucion (se detiene en la primera).
		// Devuelve true si existe, dejando columnaPorFila[] con la
		// solucion encontrada.
		bool encontrarUna() {
			return backtrackUna(0);
		}
		
		bool backtrackUna(int fila) {
			if (fila == n) return true;
			
			for (int col = 0; col < n; col++) {
				int d1 = fila + col;
				int d2 = fila - col + n;
				
				if (colOcupada[col] || diag1Ocupada[d1] || diag2Ocupada[d2])
				continue;
				
				columnaPorFila[fila] = col;
				colOcupada[col] = diag1Ocupada[d1] = diag2Ocupada[d2] = true;
				
				if (backtrackUna(fila + 1)) return true;   // exito:
				// corta
				// de inmediato,
				// sin deshacer
				// (ya se
				// encontro
				// la respuesta)
				
				colOcupada[col] = diag1Ocupada[d1] = diag2Ocupada[d2] = false;
			}
			
			return false;
		}
	};
\end{lstlisting}

\textbf{Ejemplo de funcionamiento:}

\begin{lstlisting}[style=cpp]
	NReinas nr(8);
	
	cout << nr.contarSoluciones();
	// 92   (el numero clasico de soluciones para 8 reinas)
	
	NReinas nr2(4);
	bool encontrada = nr2.encontrarUna();
	
	cout << encontrada;
	// true
	
	for (int c : nr2.columnaPorFila) cout << c << " ";
	// 1 3 0 2   (una de las dos soluciones validas para 4 reinas)
\end{lstlisting}

\textbf{Nota:} la razón de que \texttt{contarSoluciones()} y
\texttt{encontrarUna()} sean funciones separadas (en vez de una sola
con un flag) es que \texttt{encontrarUna()} puede \textbf{cortar
	inmediatamente} al éxito sin deshacer el estado (ya no importa,
terminó la búsqueda), mientras que \texttt{contarSoluciones()}
\textbf{siempre} debe deshacer y seguir explorando otras ramas para
encontrar el resto de soluciones — mezclar ambas lógicas en una sola
función suele ser la fuente más común de bugs en esta plantilla.


\subsection{\texorpdfstring{Variaciones del Problema de las $N$ Reinas}{Variaciones del Problema de las N Reinas}}

\textbf{Idea:} el esqueleto de backtracking de $N$ Reinas —una
decisión por fila, poda con arreglos de conflicto marcados en
$O(1)$— se reutiliza casi sin cambios para varias variantes que
cambian \textbf{qué} cuenta como conflicto o \textbf{qué} se pide
como respuesta, en vez de cambiar la estructura del backtracking en
sí.

\textbf{¿Por qué usarlo?} Reconocer que una variante es "N Reinas con
un ajuste" ahorra tener que rediseñar el backtracking desde cero —el
patrón de poda por marcas booleanas (o bitmask, ver más abajo) es
directamente transferible.

\textbf{Casos de uso y sus modificaciones puntuales:}

\begin{itemize}
	\item \textbf{Tablero con casillas bloqueadas}: algunas celdas no
	admiten reina (obstáculos fijos) — se agrega un chequeo extra de
	"¿esta celda está bloqueada?" antes de colocar.
	\item \textbf{$N$ Reinas con bitmask} (optimización de
	constante): en vez de tres \texttt{vector<bool>}, se usan tres
	enteros donde cada bit representa una columna/diagonal ocupada —
	las operaciones de marcar/desmarcar se vuelven operaciones a
	nivel de bits, más rápidas en la práctica para $N$ grande.
	\item \textbf{Listar todas las soluciones explícitas} (no solo
	contarlas): se guarda una copia de \texttt{columnaPorFila} cada
	vez que se completa una solución, en vez de solo incrementar un
	contador.
	\item \textbf{Reinas con posiciones prefijadas}: algunas filas ya
	tienen su reina puesta de antemano — se inicializan las marcas
	de conflicto con esas reinas antes de empezar el backtracking, y
	el recorrido de filas empieza después de la última fila
	prefijada.
\end{itemize}

\textbf{Complejidad:} igual orden que $N$ Reinas base (exponencial en
el peor caso con poda), la variante con bitmask reduce la
\textbf{constante} por operación pero no el orden asintótico.

\begin{lstlisting}[style=cpp]
	struct NReinasVariantes {
		int n;
		vector<int> columnaPorFila;
		vector<vector<bool>> bloqueada;   // bloqueada[fila][col] = true
		// si esa celda no admite reina
		vector<vector<int>> todasLasSoluciones;   // guarda cada solucion
		// completa encontrada
		
		// bitmask: cols, diag1, diag2 representan columnas/diagonales
		// ocupadas como bits encendidos, en vez de vector<bool>
		int cols, diag1, diag2;
		
		// CREATE: tablero n x n con posibles celdas bloqueadas.
		NReinasVariantes(int n_, vector<vector<bool>>& bloqueada_) {
			n = n_;
			bloqueada = bloqueada_;
			columnaPorFila.assign(n, -1);
			cols = diag1 = diag2 = 0;
		}                                          // O(n^2)
		
		// VARIANTE 1: tablero con casillas bloqueadas -- mismo
		// backtracking, un chequeo extra por candidato.
		long long contarConBloqueos() {
			long long total = 0;
			backtrackBloqueos(0, total);
			return total;
		}
		
		void backtrackBloqueos(int fila, long long& total) {
			if (fila == n) { total++; return; }
			
			for (int col = 0; col < n; col++) {
				if (bloqueada[fila][col]) continue;   // celda no
				// disponible,
				// ni siquiera se
				// evalua conflicto
				
				int d1 = fila + col, d2 = fila - col + n;
				bool ocupada = (cols & (1 << col)) ||
				(diag1 & (1 << d1)) ||
				(diag2 & (1 << d2));
				if (ocupada) continue;
				
				// VARIANTE 2: bitmask en vez de vector<bool> -- marcar
				// es un OR, desmarcar es un XOR con la misma mascara
				cols |= (1 << col);
				diag1 |= (1 << d1);
				diag2 |= (1 << d2);
				
				backtrackBloqueos(fila + 1, total);
				
				cols ^= (1 << col);
				diag1 ^= (1 << d1);
				diag2 ^= (1 << d2);
			}
		}                                          // O(N!) peor caso
		
		// VARIANTE 3: listar TODAS las soluciones explicitas, no solo
		// contarlas -- se guarda una copia completa al llegar a fila==n.
		void listarTodas() {
			todasLasSoluciones.clear();
			backtrackListar(0);
		}
		
		void backtrackListar(int fila) {
			if (fila == n) {
				todasLasSoluciones.push_back(columnaPorFila);   // copia
				// completa
				return;
			}
			
			for (int col = 0; col < n; col++) {
				int d1 = fila + col, d2 = fila - col + n;
				bool ocupada = (cols & (1 << col)) ||
				(diag1 & (1 << d1)) ||
				(diag2 & (1 << d2));
				if (ocupada || bloqueada[fila][col]) continue;
				
				columnaPorFila[fila] = col;
				cols |= (1 << col); diag1 |= (1 << d1); diag2 |= (1 << d2);
				
				backtrackListar(fila + 1);
				
				cols ^= (1 << col); diag1 ^= (1 << d1); diag2 ^= (1 << d2);
			}
		}                                          // O(N! * N) por el
		// costo de copiar cada
		// solucion encontrada
		
		// VARIANTE 4: reinas con posiciones prefijadas -- se marcan las
		// reinas ya dadas ANTES de empezar el backtracking, y este
		// arranca en la primera fila SIN prefijar.
		long long contarConPrefijadas(vector<int>& prefijadas) {
			// prefijadas[fila] = columna prefijada, o -1 si esa fila es
			// libre y debe decidirse por backtracking
			int filaInicio = 0;
			for (int fila = 0; fila < n; fila++) {
				if (prefijadas[fila] == -1) { filaInicio = fila; break; }
				
				int col = prefijadas[fila];
				columnaPorFila[fila] = col;
				cols |= (1 << col);
				diag1 |= (1 << (fila + col));
				diag2 |= (1 << (fila - col + n));
				filaInicio = fila + 1;
			}
			
			long long total = 0;
			backtrackBloqueos(filaInicio, total);
			return total;
		}                                          // mismo orden que
		// backtrackBloqueos,
		// arrancando mas
		// adentro del arbol
	};
\end{lstlisting}

\textbf{Ejemplo de funcionamiento:}

\begin{lstlisting}[style=cpp]
	// tablero 4x4 sin bloqueos
	vector<vector<bool>> sinBloqueos(4, vector<bool>(4, false));
	
	NReinasVariantes nrv(4, sinBloqueos);
	nrv.listarTodas();
	
	cout << nrv.todasLasSoluciones.size();
	// 2   (las 2 soluciones clasicas de 4 reinas)
	
	for (int c : nrv.todasLasSoluciones[0]) cout << c << " ";
	// 1 3 0 2
	
	// tablero 4x4 CON la celda (0,0) bloqueada
	vector<vector<bool>> conBloqueo(4, vector<bool>(4, false));
	conBloqueo[0][0] = true;
	
	NReinasVariantes nrv2(4, conBloqueo);
	cout << nrv2.contarConBloqueos();
	// 1   (una de las 2 soluciones originales usaba (0,0), la otra no)
\end{lstlisting}

\textbf{Regla práctica:} ante cualquier variante de $N$ Reinas,
identifica primero \textbf{qué cambia}: si cambia \textbf{qué es
	válido} (bloqueos, reinas prefijadas), el ajuste es un chequeo extra
o una inicialización antes del backtracking; si cambia \textbf{qué se
	pide} (contar vs. listar vs. encontrar una), el ajuste es qué se hace
al llegar a \texttt{fila == n} (incrementar, copiar, o retornar
\texttt{true} de inmediato). El núcleo del backtracking —una decisión
por fila, poda por columna/diagonales, deshacer al retroceder— casi
nunca cambia entre variantes.

\subsection{Tour de Caballo}

\textbf{Idea:} el Tour de Caballo (ya visto en la sección de
grillas/laberintos con la heurística de Warnsdorff) encaja también
aquí como el ejemplo clásico de backtracking sobre un tablero: hay
que visitar \textbf{cada} casilla exactamente una vez moviéndose como
caballo de ajedrez, y a diferencia de $N$ Reinas (una decisión por
fila, con poda de conflictos), aquí la decisión es "¿a cuál de los
hasta 8 saltos válidos me muevo desde la casilla actual?", con
\textbf{deshacer explícito} (marcar la casilla como no visitada) cada
vez que una rama se atasca antes de completar el tour. Se repite en
esta sección porque el patrón de backtracking puro —sin la heurística
de Warnsdorff, solo prueba y retrocede— es más instructivo para
entender la mecánica general antes de optimizar.

\textbf{¿Por qué usarlo?} Es el contraste directo con $N$ Reinas: en
$N$ Reinas la poda viene de \textbf{restricciones estructurales}
(columna/diagonal), aquí la única "poda" real disponible sin
heurística es "no pisar una casilla ya visitada" —eso hace que el
backtracking puro sea mucho más lento en la práctica (el árbol de
búsqueda crece casi sin control para tableros grandes), y es
justamente la motivación para agregar Warnsdorff \textbf{encima} de
este mismo esqueleto cuando el tablero es más grande que unos pocos
$n \times n$.

\textbf{Casos de uso:}

\begin{itemize}
	\item Encontrar (o verificar la existencia de) un tour de caballo
	completo en tableros pequeños ($n \lesssim 5$–$6$ sin
	Warnsdorff).
	\item Como plantilla base para introducir el patrón de
	"backtracking con deshacer explícito sobre una grilla", antes de
	agregar la heurística de ordenamiento de candidatos.
	\item Punto de comparación directo contra $N$ Reinas: mismo
	esqueleto (probar, recursar, deshacer si falla), pero con
	condición de éxito distinta (visitar \textbf{todas} las casillas
	en vez de colocar $N$ piezas sin conflicto).
\end{itemize}

\textbf{Complejidad:} exponencial sin poda adicional
($O(8^{n^2})$ aproximado en el peor caso) — ver la sección de
grillas para la versión con Warnsdorff, que en la práctica corre casi
lineal para $n$ razonables.

\begin{lstlisting}[style=cpp]
	struct TourCaballoBacktracking {
		int n;
		vector<vector<int>> tablero;   // tablero[r][c] = orden de visita
		// (-1 si no visitada)
		int dr[8] = {-2,-2,-1,-1, 1, 1, 2, 2};
		int dc[8] = {-1, 1,-2, 2,-2, 2,-1, 1};
		
		// CREATE: tablero n x n vacio.
		TourCaballoBacktracking(int n_) {
			n = n_;
			tablero.assign(n, vector<int>(n, -1));
		}                                          // O(n^2)
		
		bool dentro(int r, int c) {
			return r >= 0 && r < n && c >= 0 && c < n;
		}
		
		// READ: intenta construir un tour completo empezando en
		// (sr, sc), SIN ninguna heuristica de ordenamiento -- prueba
		// los 8 saltos en orden fijo, retrocede si se atasca.
		bool resolver(int sr, int sc) {
			tablero[sr][sc] = 0;
			bool exito = backtrack(sr, sc, 1);
			if (!exito) tablero[sr][sc] = -1;   // deshacer tambien el
			// punto de partida si
			// no se encontro nada
			return exito;
		}                                          // exponencial sin poda
		
		bool backtrack(int r, int c, int movimiento) {
			if (movimiento == n * n) return true;   // todas las
			// casillas visitadas
			
			for (int d = 0; d < 8; d++) {
				int nr = r + dr[d], nc = c + dc[d];
				
				if (!dentro(nr, nc) || tablero[nr][nc] != -1)
				continue;   // fuera del tablero o ya visitada
				
				tablero[nr][nc] = movimiento;   // prueba
				
				if (backtrack(nr, nc, movimiento + 1))
				return true;   // exito: no hace falta deshacer, ya
				// se encontro la respuesta completa
				
				tablero[nr][nc] = -1;   // deshacer (backtrack real):
				// esta rama no llevo a un tour
				// completo
			}
			
			return false;   // ningun salto desde (r,c) llevo a exito
		}
	};
\end{lstlisting}

\textbf{Ejemplo de funcionamiento:}

\begin{lstlisting}[style=cpp]
	TourCaballoBacktracking tcb(5);
	
	bool encontrado = tcb.resolver(0, 0);
	
	cout << encontrado;
	// true
	
	// tcb.tablero[r][c] tiene el orden (0..24) en que se visito cada
	// casilla -- sin Warnsdorff, el orden exacto encontrado puede
	// diferir del de la version con heuristica, pero ambos son tours
	// validos si existen
\end{lstlisting}

\textbf{Regla práctica:} este esqueleto (probar candidato, recursar,
deshacer si la rama falla) es el mismo de $N$ Reinas y de cualquier
backtracking clásico —lo único que cambia entre problemas es
\textbf{qué} se prueba en cada paso (columna vs. salto de caballo) y
\textbf{cuándo} se declara éxito (\texttt{fila == n} vs.
\texttt{movimiento == n*n}). Para tableros mayores a unos pocos
$n \times n$, agrega la heurística de Warnsdorff (sección de grillas)
\textbf{encima} de este mismo backtracking —no cambia la estructura,
solo el \textbf{orden} en que se prueban los 8 saltos en cada paso.


% =========================================================
% ADVANCED
% =========================================================

\section{Advanced Graph Algorithms}

\subsection{Centroid Decomposition}

\textbf{Idea:} Centroid Decomposition descompone un árbol en una
\textbf{jerarquía de centroides} para resolver queries que involucran
"caminos entre cualquier par de nodos" en $O(\log n)$ niveles: se
encuentra el \textbf{centroide} del árbol (el nodo cuya remoción deja
componentes de tamaño $\le n/2$ cada una — no confundir con Centro de
un Árbol, que minimiza distancia máxima; el centroide minimiza
tamaño de componente al remover), se procesa toda la información que
pasa \textbf{por} ese centroide, se \textbf{remueve} (lógicamente) del
árbol, y se \textbf{recursiona} de forma independiente sobre cada
componente resultante, cada una con su propio centroide. Como cada
remoción parte el árbol en piezas de tamaño a lo más la mitad, la
profundidad de esta jerarquía de centroides es $O(\log n)$ — y
cualquier camino entre dos nodos $u,v$ del árbol original pasa por
\textbf{exactamente un} centroide de la jerarquía (el primero, en la
recursión, que los separa en componentes distintas).

\textbf{¿Por qué usarlo?} Es la técnica estándar para problemas de
"cuenta/encuentra pares de nodos que cumplen una condición sobre su
camino" (ej. "cuántos pares están a distancia $\le k$"), donde
recorrer todos los pares es $O(n^2)$: al procesar cada nivel de la
jerarquía en $O(n \log n)$ total y haber solo $O(\log n)$ niveles, el
costo total baja a $O(n \log^2 n)$. Es conceptualmente el análogo, en
árboles, de \textit{divide and conquer} sobre arreglos —"encuentra el
punto medio, resuelve lo que cruza el medio, recursiona a los lados"
—salvo que el "medio" aquí es un nodo (el centroide), no un índice.

\textbf{Casos de uso:}

\begin{itemize}
	\item Contar pares de nodos $(u,v)$ cuya distancia es
	\textbf{exactamente} $k$, o está en un rango dado — proceso: para
	cada centroide, contar cuántos caminos que pasan por él tienen la
	distancia buscada (usando las distancias de cada nodo al
	centroide, agrupadas en un arreglo/mapa de frecuencias).
	\item Encontrar el par de nodos más cercano/lejano que cumple
	una restricción adicional (ej. "de distinto color").
	\item Responder queries offline de tipo "actualiza un nodo,
	consulta el más cercano nodo actualizado" (una aplicación
	frecuente combinada con la estructura de centroides para acotar
	la distancia recorrida por query).
	\item Como alternativa a Heavy-Light Decomposition cuando el
	problema es sobre \textbf{pares} de caminos, no sobre un único
	camino consultado repetidamente.
\end{itemize}

\textbf{Complejidad:} $O(n \log n)$ para construir toda la jerarquía
(encontrar centroides, un DFS de tamaños por cada nivel de
recursión — la suma total de tamaños across niveles es
$O(n \log n)$); el costo de procesar cada centroide depende del
problema, típicamente $O(\text{tam} \log \text{tam})$ por centroide,
dando $O(n \log^2 n)$ total.

\begin{lstlisting}[style=cpp]
	struct CentroidDecomposition {
		vector<vector<int>> adj;
		vector<int> tam;
		vector<bool> removido;
		vector<int> padreCentroide;   // padreCentroide[c] = centroide
		// padre de c en la jerarquia
		int n;
		
		// CREATE: guarda el arbol y arranca la descomposicion completa
		// desde la raiz de la jerarquia (nodo -1 = sin padre centroide).
		CentroidDecomposition(vector<vector<int>>& adj_, int n_) {
			adj = adj_;
			n = n_;
			tam.assign(n, 0);
			removido.assign(n, false);
			padreCentroide.assign(n, -1);
			
			descomponer(0, -1);
		}                                          // O(n log n)
		
		int calcularTam(int u, int padre) {
			tam[u] = 1;
			for (int v : adj[u])
			if (v != padre && !removido[v])
			tam[u] += calcularTam(v, u);
			return tam[u];
		}
		
		// encuentra el centroide del componente que contiene a u,
		// usando tamTotal (ya calculado con calcularTam) para decidir
		// cuando un hijo excede n/2.
		int encontrarCentroide(int u, int padre, int tamTotal) {
			for (int v : adj[u])
			if (v != padre && !removido[v] && tam[v] > tamTotal / 2)
			return encontrarCentroide(v, u, tamTotal);
			return u;
		}
		
		// procesa toda la informacion que pasa POR el centroide actual
		// -- el cuerpo especifico depende del problema (aqui vacio, ver
		// ejemplo de "contar pares a distancia k" mas abajo).
		void procesarCentroide(int centroide) {
			// placeholder: aqui va la logica especifica del problema,
			// tipicamente un DFS desde 'centroide' recolectando
			// distancias de cada nodo del componente hacia el
		}
		
		// WRITE (recursivo): encuentra el centroide del componente que
		// contiene a u, lo procesa, lo marca removido, y recursiona
		// sobre cada componente resultante.
		void descomponer(int u, int padre) {
			int tamTotal = calcularTam(u, -1);
			int centroide = encontrarCentroide(u, -1, tamTotal);
			
			removido[centroide] = true;
			padreCentroide[centroide] = padre;
			
			procesarCentroide(centroide);
			
			for (int v : adj[centroide])
			if (!removido[v])
			descomponer(v, centroide);
		}                                          // O(tam log tam)
		// por nivel, O(log n)
		// niveles -> O(n log n)
		// total
	};
\end{lstlisting}

\textbf{Ejemplo de funcionamiento (contar pares a distancia exacta $k$):}

\begin{lstlisting}[style=cpp]
	// version especializada de procesarCentroide para "contar pares de
	// nodos a distancia EXACTAMENTE k" -- se agrega sobre el esqueleto
	// de arriba:
	//
	// long long totalPares = 0;
	// void procesarCentroideConteo(int centroide, int k) {
		//     vector<int> distancias;   // distancia de cada nodo del
		//                                 // componente actual al centroide
		//     function<void(int,int,int)> recolectar = [&](int u, int p, int d) {
			//         distancias.push_back(d);
			//         for (int v : adj[u])
			//             if (v != p && !removido[v])
			//                 recolectar(v, u, d + 1);
			//     };
		//
		//     // cuenta pares DENTRO de cada subarbol de un hijo directo del
		//     // centroide primero (para luego RESTARLOS -- evita contar
		//     // pares que en realidad no pasan por el centroide, ambos en
		//     // el mismo hijo)
		//     map<int,int> frecuencia;
		//     frecuencia[0] = 1;   // el propio centroide, distancia 0
		//
		//     for (int hijo : adj[centroide]) {
			//         if (removido[hijo]) continue;
			//
			//         distancias.clear();
			//         recolectar(hijo, centroide, 1);
			//
			//         for (int d : distancias)
			//             if (frecuencia.count(k - d))
			//                 totalPares += frecuencia[k - d];
			//
			//         for (int d : distancias)
			//             frecuencia[d]++;
			//     }
		// }
	
	// arbol:      0
	//            / \
	//           1   2
	//          /
	//         3
	
	vector<vector<int>> adj(4);
	adj[0] = {1, 2};
	adj[1] = {0, 3};
	adj[2] = {0};
	adj[3] = {1};
	
	// contar pares a distancia EXACTAMENTE 2:
	// (0,3) dist=2, (2,1) dist=2 -> 2 pares
\end{lstlisting}

\textbf{Regla práctica:} Centroid Decomposition vale la pena
\textbf{solo} cuando el problema pregunta por \textbf{pares} de nodos
o caminos \textbf{arbitrarios} — si el problema fija un único camino
por query (como en Path Queries o HLD), Centroid Decomposition es
sobre-ingeniería innecesaria. La señal típica es un enunciado que
pregunta "cuántos pares..." o "existe algún par..." con una condición
sobre la distancia o el camino entre ellos.


\subsection{DSU on Tree}

\textbf{Idea:} DSU on Tree (también llamado
\textit{small-to-large merging}, a pesar del nombre no usa la
estructura DSU de Union-Find) resuelve queries \textbf{por subárbol}
—"para cada nodo $u$, calcula algo sobre \textbf{todo} su subárbol"—
evitando recalcular desde cero para cada nodo. La idea: al procesar
un nodo $u$, primero se procesan (y se \textbf{descartan}) todos los
hijos \textbf{livianos} (los que no son el hijo de mayor subárbol,
igual que en Heavy-Light Decomposition), guardando su información
solo temporalmente; luego se procesa el hijo \textbf{pesado} sin
descartar su información (se deja "viva"); y finalmente se
\textbf{reinserta} la información de los hijos livianos encima de la
del hijo pesado, para tener el estado completo del subárbol de $u$
sin haber recorrido cada nodo más que $O(\log n)$ veces en total.

\textbf{¿Por qué usarlo?} La alternativa ingenua —recorrer el
subárbol completo de cada nodo para responder su query— cuesta
$O(n^2)$ en el peor caso (un árbol en forma de cadena). El truco de
"conservar la estructura del hijo pesado y solo reinsertar los
livianos" es el mismo principio que sustenta la complejidad de HLD
—cualquier nodo cambia de "hijo pesado a liviano" como mucho
$O(\log n)$ veces en su camino a la raíz—, aplicado aquí a cuántas
veces se re-inserta cada nodo en una estructura auxiliar en vez de a
cuántas cadenas cruza un camino.

\textbf{Casos de uso:}

\begin{itemize}
	\item "Para cada nodo, ¿cuál es el valor/color más frecuente en
	su subárbol?" (el ejemplo clásico, con un arreglo de
	frecuencias como estructura auxiliar).
	\item Contar cuántos valores \textbf{distintos} hay en el
	subárbol de cada nodo.
	\item Cualquier query de subárbol donde la estructura auxiliar
	soporte "agregar un elemento" y "responder la pregunta actual" en
	$O(1)$ o $O(\log n)$, pero \textbf{no} soporte eliminar
	eficientemente (si soportara borrado barato, Euler Tour +
	Fenwick ya bastaría, sin necesitar este patrón).
\end{itemize}

\textbf{Complejidad:} $O(n \log n)$ total — cada nodo se "agrega" a la
estructura auxiliar como mucho $O(\log n)$ veces (una vez por cada
antepasado del que sea hijo liviano, más una vez cuando su propio
subárbol se procesa como hijo pesado).

\begin{lstlisting}[style=cpp]
	struct DSUOnTree {
		vector<vector<int>> adj;
		vector<int> tam, hijoPesado, valor;
		vector<long long> frecuencia;   // estructura auxiliar: cuenta
		// ocurrencias de cada valor
		vector<long long> respuesta;    // respuesta[u] = valor mas
		// frecuente en el subarbol de u
		long long maxFrecuencia, sumaDeMaximos;
		int n;
		
		// CREATE: guarda el arbol y los valores por nodo; calcula
		// tamanos e identifica el hijo pesado de cada nodo (igual que
		// en HLD).
		DSUOnTree(vector<vector<int>>& adj_, vector<int>& valor_, int n_) {
			adj = adj_;
			valor = valor_;
			n = n_;
			tam.assign(n, 0);
			hijoPesado.assign(n, -1);
			respuesta.assign(n, 0);
			frecuencia.assign(*max_element(valor.begin(), valor.end()) + 1, 0);
			
			calcularTamYPesado(0, -1);
		}                                          // O(n)
		
		void calcularTamYPesado(int u, int padre) {
			tam[u] = 1;
			int mejorTam = -1;
			
			for (int v : adj[u]) {
				if (v == padre) continue;
				calcularTamYPesado(v, u);
				tam[u] += tam[v];
				if (tam[v] > mejorTam) {
					mejorTam = tam[v];
					hijoPesado[u] = v;
				}
			}
		}                                          // O(n)
		
		// agrega o quita (delta = +1 o -1) el valor de u a la estructura
		// auxiliar, actualizando el maximo global sobre la marcha.
		void agregar(int u, int padre, int delta) {
			frecuencia[valor[u]] += delta;
			
			if (frecuencia[valor[u]] > maxFrecuencia) {
				maxFrecuencia = frecuencia[valor[u]];
				sumaDeMaximos = valor[u];
			} else if (frecuencia[valor[u]] == maxFrecuencia) {
				sumaDeMaximos += valor[u];   // ejemplo: suma de todos
				// los valores empatados
				// en frecuencia maxima
			}
			
			for (int v : adj[u])
			if (v != padre)
			agregar(v, u, delta);
		}                                          // O(tamano del subarbol)
		
		// READ (recursivo): procesa el subarbol de u con la tecnica de
		// small-to-large; 'mantener' indica si u es el hijo PESADO de
		// su padre (si es asi, su info NO se descarta al terminar).
		void dfs(int u, int padre, bool mantener) {
			// paso 1: procesa TODOS los hijos livianos primero,
			// descartando su info despues (mantener = false)
			for (int v : adj[u])
			if (v != padre && v != hijoPesado[u])
			dfs(v, u, false);
			
			// paso 2: procesa el hijo pesado SIN descartar su info
			if (hijoPesado[u] != -1)
			dfs(hijoPesado[u], u, true);
			
			// paso 3: agrega el propio u y REINSERTA todos los hijos
			// livianos encima de la info ya viva del hijo pesado
			agregar(u, padre, +1);
			for (int v : adj[u])
			if (v != padre && v != hijoPesado[u])
			agregarSoloEsteSubarbol(v, u, +1);
			
			// ahora la estructura auxiliar tiene el subarbol COMPLETO
			// de u -> se puede responder la query de u
			respuesta[u] = sumaDeMaximos;
			
			// paso 4: si u es hijo liviano de su padre, se descarta TODA
			// su informacion (para no contaminar al procesar hermanos)
			if (!mantener) {
				agregar(u, padre, -1);
				for (int v : adj[u])
				if (v != padre)
				agregarSoloEsteSubarbol(v, u, -1);
				maxFrecuencia = 0;
				sumaDeMaximos = 0;
			}
		}                                          // O(n log n) total
		
		void agregarSoloEsteSubarbol(int u, int padre, int delta) {
			agregar(u, padre, delta);
		}
		
		// READ: calcula respuesta[] para todo el arbol.
		vector<long long> resolver() {
			maxFrecuencia = 0;
			sumaDeMaximos = 0;
			dfs(0, -1, true);   // la raiz siempre "mantiene" (no hay
			// nada que descartar despues)
			return respuesta;
		}                                          // O(n log n)
	};
\end{lstlisting}

\textbf{Ejemplo de funcionamiento:}

\begin{lstlisting}[style=cpp]
	// arbol:      0(val=1)
	//            /        \
	//        1(val=2)   2(val=1)
	//        /
	//   3(val=2)
	
	vector<vector<int>> adj(4);
	adj[0] = {1, 2};
	adj[1] = {0, 3};
	adj[2] = {0};
	adj[3] = {1};
	
	vector<int> valor = {1, 2, 1, 2};
	
	DSUOnTree dot(adj, valor, 4);
	auto resp = dot.resolver();
	
	cout << resp[1];
	// 2   (subarbol de 1: valores {2, 2} -> el 2 es el mas frecuente)
	
	cout << resp[0];
	// 3   (subarbol completo: valores {1,2,1,2} -> 1 y 2 empatan en
	//      frecuencia 2, suma de ambos = 1+2 = 3, segun la definicion
	//      de "respuesta" de este ejemplo especifico)
\end{lstlisting}

\textbf{Regla práctica:} DSU on Tree comparte con HLD el mismo
concepto de "hijo pesado" y la misma garantía de $O(\log n)$
reinserciones por nodo, pero resuelve un problema distinto: HLD
responde queries sobre \textbf{caminos} descomponiéndolos en rangos;
DSU on Tree responde queries sobre \textbf{subárboles completos}
reutilizando la estructura auxiliar del hijo pesado en vez de
reconstruirla desde cero. Úsalo quando la pregunta es "algo sobre
\textbf{todo} el subárbol de cada nodo" y la estructura auxiliar
(usualmente un arreglo/mapa de frecuencias) es cara de reconstruir
pero barata de \textbf{extender} elemento por elemento.

\subsection{Hopcroft-Karp}

\textbf{Idea:} Hopcroft-Karp mejora a Kuhn (ver Bipartite Matching)
con la misma idea que Dinic mejora a Ford-Fulkerson: en vez de buscar
\textbf{un} camino aumentante a la vez con DFS, busca \textbf{muchos}
caminos aumentantes \textbf{disjuntos} de una sola vez, todos de la
\textbf{misma longitud mínima}. Cada \textbf{fase} tiene dos partes:
(1) un BFS desde todos los nodos libres de $L$ que calcula la
distancia mínima a un nodo libre de $R$ (alternando aristas normales
y de matching, igual que las capas de Dinic); (2) una serie de DFS
que extraen \textbf{todos} los caminos aumentantes disjuntos de esa
longitud mínima, marcando nodos usados para no repetirlos dentro de
la misma fase. Se puede demostrar que el número total de fases está
acotado por $O(\sqrt{V})$ — de ahí la complejidad total.

\textbf{¿Por qué usarlo?} Kuhn es $O(V \cdot E)$; Hopcroft-Karp baja
eso a $O(E\sqrt{V})$ —la misma mejora que Dinic da sobre
Ford-Fulkerson, pero especializada al caso bipartito puro (sin pasar
por la maquinaria completa de flujo con capacidades). Es preferible
sobre "Maximum Bipartite Matching vía Dinic" cuando \textbf{solo} se
necesita matching (sin capacidades extra ni combinarlo con otro
modelo de flujo) — mismo orden de complejidad, pero con una
implementación más directa y sin la sobrecarga de construir la red
completa con superfuente/supersumidero.

\textbf{Casos de uso:}

\begin{itemize}
	\item Los mismos que Bipartite Matching / Maximum Bipartite
	Matching, cuando $V$ es grande y $O(V \cdot E)$ de Kuhn no
	alcanza el límite de tiempo.
	\item Como reemplazo directo de Kuhn en cualquier problema que ya
	modele matching bipartito puro (sin capacidades adicionales ni
	pesos).
	\item Preferible sobre Dinic + construcción de flujo cuando el
	problema es \textbf{estrictamente} matching bipartito (Dinic
	sigue siendo necesario si además hay capacidades $>1$ o el
	problema mezcla matching con otras restricciones de flujo).
\end{itemize}

\textbf{Complejidad:} $O(E\sqrt{V})$ — $O(\sqrt{V})$ fases, cada una
$O(E)$ (un BFS más varios DFS que en conjunto tocan cada arista una
vez por fase).

\begin{lstlisting}[style=cpp]
	struct HopcroftKarp {
		vector<vector<int>> adj;   // adj[u] = vecinos en R, para cada
		// u en L
		vector<int> matchL, matchR, dist;
		int nL, nR;
		const int INF = INT_MAX;
		const int LIBRE = -1;
		
		// CREATE: guarda el grafo bipartito (nL nodos en L, adj[u] lista
		// los vecinos de u en R). nR = tamano del lado R.
		HopcroftKarp(vector<vector<int>>& adj_, int nL_, int nR_) {
			adj = adj_;
			nL = nL_;
			nR = nR_;
			matchL.assign(nL, LIBRE);
			matchR.assign(nR, LIBRE);
		}                                          // O(1)
		
		// FASE 1: BFS multi-fuente desde todos los nodos LIBRES de L,
		// calculando dist[] SOLO sobre nodos de L (alternando aristas
		// normales L->R con aristas de matching R->L). Devuelve true si
		// algun nodo libre de R fue alcanzado (hay al menos un camino
		// aumentante en esta fase).
		bool bfs() {
			queue<int> q;
			dist.assign(nL, INF);
			
			for (int u = 0; u < nL; u++)
			if (matchL[u] == LIBRE) {
				dist[u] = 0;
				q.push(u);
			}
			
			bool encontroLibre = false;
			
			while (!q.empty()) {
				int u = q.front(); q.pop();
				
				for (int v : adj[u]) {
					int siguienteL = matchR[v];
					
					if (siguienteL == LIBRE) {
						encontroLibre = true;   // v es libre: fin de un
						// camino aumentante de
						// esta longitud minima
					} else if (dist[siguienteL] == INF) {
						dist[siguienteL] = dist[u] + 1;
						q.push(siguienteL);
					}
				}
			}
			
			return encontroLibre;
		}                                          // O(E)
		
		// FASE 2: DFS que busca UN camino aumentante respetando dist[]
		// (solo avanza a niveles crecientes, igual que el grafo por
		// niveles de Dinic), actualizando matchL/matchR al encontrarlo.
		bool dfs(int u) {
			for (int v : adj[u]) {
				int siguienteL = matchR[v];
				
				if (siguienteL == LIBRE ||
				(dist[siguienteL] == dist[u] + 1 && dfs(siguienteL))) {
					matchL[u] = v;
					matchR[v] = u;
					return true;
				}
			}
			
			dist[u] = INF;   // u agotado en esta fase, no vale la pena
			// volver a intentarlo dentro de la misma fase
			return false;
		}                                          // O(camino encontrado)
		
		// READ: tamano del matching maximo.
		int matchingMaximo() {
			int total = 0;
			
			while (bfs()) {
				for (int u = 0; u < nL; u++)
				if (matchL[u] == LIBRE && dfs(u))
				total++;
			}
			
			return total;
		}                                          // O(E * sqrt(V))
	};
\end{lstlisting}

\textbf{Ejemplo de funcionamiento:}

\begin{lstlisting}[style=cpp]
	// L = {0, 1, 2}, R = {0, 1, 2}
	// 0 - {0, 1}
	// 1 - {0}
	// 2 - {1, 2}
	
	vector<vector<int>> adj(3);
	adj[0] = {0, 1};
	adj[1] = {0};
	adj[2] = {1, 2};
	
	HopcroftKarp hk(adj, 3, 3);
	
	cout << hk.matchingMaximo();
	// 3   (mismo resultado que Kuhn y que Maximum Bipartite Matching
	//      via Dinic sobre el mismo grafo)
\end{lstlisting}

\textbf{Regla práctica:} Kuhn, Maximum Bipartite Matching vía Dinic,
y Hopcroft-Karp \textbf{siempre} dan el mismo tamaño de matching — la
elección es de rendimiento y de cuánto código ya tienes: Kuhn para
$V$ chico o cuando la simplicidad importa más que la velocidad;
Hopcroft-Karp cuando $V$ es grande y el problema es matching puro sin
nada más; Dinic + construcción de flujo cuando el matching es solo
una parte de un modelo de flujo más grande (capacidades, múltiples
fuentes, etc.) y ya se necesita esa infraestructura de todas formas.


\subsection{Min-Cost Max-Flow}

\textbf{Idea:} Min-Cost Max-Flow (MCMF) resuelve una versión más rica
del problema de flujo: cada arista tiene, además de una
\textbf{capacidad}, un \textbf{costo por unidad} de flujo que pasa
por ella, y se busca el flujo máximo \textbf{de menor costo total}
posible (o, en la variante más común, "envía exactamente $F$ unidades
al menor costo"). Se resuelve con la misma estructura de grafo
residual de Ford-Fulkerson/Dinic, pero en vez de buscar
\textbf{cualquier} camino aumentante, en cada iteración se busca el
camino aumentante de \textbf{menor costo} con \textbf{Bellman-Ford}
o \textbf{SPFA} (necesarios porque las aristas reversas tienen costo
\textbf{negativo} —"deshacer" un envío recupera el costo pagado, así
que Dijkstra normal no aplica directamente sin potenciales de
Johnson).

\textbf{¿Por qué usarlo?} Es la generalización que responde a
cualquier problema de flujo que además pida \textbf{minimizar/maximizar
	un costo} asociado a las unidades transportadas, algo que ni
Ford-Fulkerson ni Dinic pueden resolver por sí solos (ellos solo
optimizan \textbf{cantidad}, no costo). Es también la herramienta
estándar para \textbf{matching bipartito ponderado} (asignación de
costo mínimo, ver casos de uso) —donde Hopcroft-Karp/Kuhn ya no
alcanzan porque esos solo maximizan \textbf{cantidad} de parejas, no
el costo/beneficio total.

\textbf{Casos de uso:}

\begin{itemize}
	\item Problema de asignación (\textit{assignment problem}):
	asignar $n$ trabajadores a $n$ tareas minimizando el costo total,
	cuando cada trabajador \textbf{debe} tomar exactamente una tarea
	—red bipartita con capacidad 1 y costo por arista, buscando flujo
	de $n$ unidades al menor costo.
	\item Transporte de mercancía entre orígenes y destinos con
	costo por unidad transportada distinto según la ruta,
	minimizando el costo total de satisfacer toda la demanda.
	\item Como generalización de Bipartite Matching cuando las
	parejas no son solo "compatible/incompatible" sino que tienen un
	\textbf{costo o beneficio} asociado.
	\item Cualquier problema de flujo donde el enunciado mencione
	explícitamente un \textbf{costo por unidad} además de una
	capacidad.
\end{itemize}

\textbf{Complejidad:} $O(F \cdot E)$ con SPFA (Bellman-Ford
optimizado) por iteración, donde $F$ es el valor total de flujo
enviado — cada iteración envía al menos 1 unidad por el camino más
barato encontrado. Para $F$ grande, existen variantes con potenciales
de Johnson + Dijkstra ($O(F \cdot E \log V)$ pero con constante mucho
menor), fuera del alcance de esta plantilla base.

\begin{lstlisting}[style=cpp]
	struct MinCostMaxFlow {
		struct Arista {
			int destino; long long capacidad, costo; int reversa;
		};
		vector<vector<Arista>> adj;
		int n;
		
		// CREATE: grafo residual vacio de n nodos.
		MinCostMaxFlow(int n_) {
			n = n_;
			adj.assign(n, {});
		}                                          // O(n)
		
		// WRITE: agrega arista dirigida u->v con capacidad c y costo
		// costo POR UNIDAD. La reversa tiene capacidad 0 y costo
		// NEGATIVO (deshacer un envio recupera lo pagado).
		void agregarArista(int u, int v, long long c, long long costo) {
			adj[u].push_back({v, c, costo, (int)adj[v].size()});
			adj[v].push_back({u, 0, -costo, (int)adj[u].size() - 1});
		}                                          // O(1)
		
		// busca el camino de MENOR COSTO de s a t en el grafo residual
		// actual con SPFA (Bellman-Ford + cola, tolera costos negativos
		// de las aristas reversas). Devuelve false si t no es alcanzable.
		bool spfa(int s, int t, vector<long long>& dist,
		vector<int>& padreNodo, vector<int>& padreArista) {
			dist.assign(n, LLONG_MAX);
			vector<bool> enCola(n, false);
			dist[s] = 0;
			
			queue<int> q;
			q.push(s);
			enCola[s] = true;
			
			while (!q.empty()) {
				int u = q.front(); q.pop();
				enCola[u] = false;
				
				for (int i = 0; i < (int)adj[u].size(); i++) {
					auto& a = adj[u][i];
					if (a.capacidad <= 0) continue;
					if (dist[u] + a.costo >= dist[a.destino]) continue;
					
					dist[a.destino] = dist[u] + a.costo;
					padreNodo[a.destino] = u;
					padreArista[a.destino] = i;
					
					if (!enCola[a.destino]) {
						q.push(a.destino);
						enCola[a.destino] = true;
					}
				}
			}
			
			return dist[t] != LLONG_MAX;
		}                                          // O(V * E) peor caso
		
		// READ: flujo maximo de s a t, y su costo total minimo asociado
		// (devuelve {flujoTotal, costoTotal}).
		pair<long long,long long> flujoCostoMinimo(int s, int t) {
			long long flujoTotal = 0, costoTotal = 0;
			vector<long long> dist(n);
			vector<int> padreNodo(n), padreArista(n);
			
			while (spfa(s, t, dist, padreNodo, padreArista)) {
				// cuello de botella del camino encontrado (igual que en
				// Ford-Fulkerson, pero sobre el camino de MENOR COSTO)
				long long cuello = LLONG_MAX;
				for (int u = t; u != s; u = padreNodo[u])
				cuello = min(cuello,
				adj[padreNodo[u]][padreArista[u]].capacidad);
				
				for (int u = t; u != s; u = padreNodo[u]) {
					auto& a = adj[padreNodo[u]][padreArista[u]];
					a.capacidad -= cuello;
					adj[u][a.reversa].capacidad += cuello;
				}
				
				flujoTotal += cuello;
				costoTotal += cuello * dist[t];
			}
			
			return {flujoTotal, costoTotal};
		}                                          // O(F * V * E)
		
		// READ: costo minimo para enviar EXACTAMENTE F unidades (asume
		// que existe un flujo de al menos F -- verificar con
		// flujoCostoMinimo primero si no se esta seguro).
		long long costoParaFlujoExacto(int s, int t, long long F) {
			long long flujoTotal = 0, costoTotal = 0;
			vector<long long> dist(n);
			vector<int> padreNodo(n), padreArista(n);
			
			while (flujoTotal < F && spfa(s, t, dist, padreNodo, padreArista)) {
				long long cuello = F - flujoTotal;
				for (int u = t; u != s; u = padreNodo[u])
				cuello = min(cuello,
				adj[padreNodo[u]][padreArista[u]].capacidad);
				
				for (int u = t; u != s; u = padreNodo[u]) {
					auto& a = adj[padreNodo[u]][padreArista[u]];
					a.capacidad -= cuello;
					adj[u][a.reversa].capacidad += cuello;
				}
				
				flujoTotal += cuello;
				costoTotal += cuello * dist[t];
			}
			
			return costoTotal;   // si flujoTotal < F al salir, no existia
			// suficiente flujo -- verificar aparte
		}                                          // O(F * V * E)
	};
\end{lstlisting}

\textbf{Ejemplo de funcionamiento (asignación de costo mínimo):}

\begin{lstlisting}[style=cpp]
	// 2 trabajadores, 2 tareas, matriz de costos:
	//        tarea0  tarea1
	// trab0:   4       2
	// trab1:   3       5
	// (fuente=0, trab0=1, trab1=2, tarea0=3, tarea1=4, sumidero=5)
	
	MinCostMaxFlow mcmf(6);
	mcmf.agregarArista(0, 1, 1, 0);   // fuente -> trab0
	mcmf.agregarArista(0, 2, 1, 0);   // fuente -> trab1
	mcmf.agregarArista(1, 3, 1, 4);   // trab0 -> tarea0, costo 4
	mcmf.agregarArista(1, 4, 1, 2);   // trab0 -> tarea1, costo 2
	mcmf.agregarArista(2, 3, 1, 3);   // trab1 -> tarea0, costo 3
	mcmf.agregarArista(2, 4, 1, 5);   // trab1 -> tarea1, costo 5
	mcmf.agregarArista(3, 5, 1, 0);   // tarea0 -> sumidero
	mcmf.agregarArista(4, 5, 1, 0);   // tarea1 -> sumidero
	
	auto [flujo, costo] = mcmf.flujoCostoMinimo(0, 5);
	
	cout << flujo << " " << costo;
	// 2 5
	// (trab0 -> tarea1 costo 2, trab1 -> tarea0 costo 3, total 5 --
	//  mejor que trab0->tarea0 + trab1->tarea1 = 4+5=9)
\end{lstlisting}

\textbf{Regla práctica:} usa Min-Cost Max-Flow en cuanto el
enunciado combine "capacidad" \textbf{con} "costo por unidad" —si
solo hay capacidades sin costo, Dinic/Edmonds-Karp ya resuelven
flujo máximo puro sin la sobrecarga extra de SPFA por iteración.
Para el caso particular de asignación (matching ponderado 1 a 1,
todas las capacidades en 1), esta misma plantilla \textbf{es} el
algoritmo húngaro disfrazado de flujo —construir la red bipartita con
costos y correr \texttt{flujoCostoMinimo} es más simple de programar
bajo presión que implementar el algoritmo húngaro clásico desde cero.

\subsection{DSU Rollback}

\textbf{Idea:} el DSU normal (ver Disjoint Set Union) usa compresión
de camino para llegar a $O(\alpha(n))$ amortizado, pero la compresión
de camino \textbf{destruye información} —una vez que \texttt{find(x)}
aplana el camino, ya no se puede "deshacer" esa unión más tarde de
forma barata. DSU Rollback resuelve esto \textbf{sacrificando} la
compresión de camino y quedándose \textbf{solo} con unión por
rango/tamaño: sin compresión, cada \texttt{union} es una operación
\textbf{completamente reversible} en $O(1)$ —basta con guardar, antes
de cada unión, qué nodo cambió de padre y cuál era su rango anterior,
y al hacer \texttt{rollback} restaurar exactamente esos dos valores.
El costo es que \texttt{find} ya no es $O(\alpha(n))$ amortizado sino
$O(\log n)$ (garantizado por la unión por rango sola, sin compresión).

\textbf{¿Por qué usarlo?} Es la estructura estándar para
\textbf{conectividad offline con deshacer} —problemas donde las
uniones no son permanentes, sino que se agregan y \textbf{se
	deshacen} en algún orden (ej. recorrido de Centroid Decomposition o
de un \textit{divide and conquer} sobre el tiempo, donde una arista
"existe" solo durante un rango de queries). Sin rollback, la única
opción sería reconstruir el DSU desde cero cada vez que cambia el
conjunto de aristas activas —mucho más caro si eso ocurre repetidas
veces.

\textbf{Casos de uso:}

\begin{itemize}
	\item \textbf{Divide and conquer sobre el tiempo} (offline
	dynamic connectivity): cada arista existe durante un intervalo de
	tiempo $[l,r]$; se recorre un segment tree sobre el tiempo,
	agregando la arista en los nodos que cubren su intervalo y
	haciendo rollback al salir de cada nodo —permite responder
	"¿$u$ y $v$ conectados en el momento $t$?" sin reconstruir el DSU
	desde cero en cada consulta.
	\item Algoritmos que exploran múltiples "ramas" de un conjunto de
	uniones (ej. probar si agregar una arista temporalmente crea un
	ciclo, y deshacerlo si no aporta) sin pagar el costo de
	reconstrucción completa.
	\item Como building block de variantes de Kruskal que necesitan
	"probar y deshacer" una arista candidata antes de confirmarla
	definitivamente.
\end{itemize}

\textbf{Complejidad:} $O(\log n)$ por \texttt{find}/\texttt{union}
(sin compresión de camino, solo unión por rango), $O(1)$ por
\texttt{rollback} de una unión ya hecha.

\begin{lstlisting}[style=cpp]
	struct DSURollback {
		vector<int> padre, rango;
		// pila de deshacer: guarda (nodo que cambio de padre, su rango
		// ANTERIOR si tambien cambio) para cada union exitosa
		vector<tuple<int,int,int>> historial;   // {hijo, padreAnterior,
			// rangoAnteriorDelRaiz}
		
		// CREATE: n elementos, cada uno en su propio conjunto.
		DSURollback(int n) {
			padre.resize(n);
			rango.assign(n, 0);
			iota(padre.begin(), padre.end(), 0);
		}                                          // O(n)
		
		// READ: representante del conjunto de x, SIN compresion de
		// camino (la compresion romperia la reversibilidad de union).
		int find(int x) {
			while (padre[x] != x)
			x = padre[x];
			return x;
		}                                          // O(log n) por la
		// union por rango
		
		// WRITE: fusiona los conjuntos de x e y, registrando en
		// 'historial' lo necesario para deshacer esta union despues.
		// Devuelve false si ya estaban en el mismo conjunto (no se
		// registra nada en ese caso, no hay nada que deshacer).
		bool unir(int x, int y) {
			int rx = find(x), ry = find(y);
			
			if (rx == ry) {
				historial.push_back({-1, -1, -1});   // marcador "union
				// nula", para que
				// rollback sepa
				// que no hay nada
				// que restaurar
				return false;
			}
			
			if (rango[rx] < rango[ry]) swap(rx, ry);
			
			// guarda el estado ANTES de modificar, para poder revertir
			historial.push_back({ry, rango[rx], rx});
			
			padre[ry] = rx;
			bool rangoSubio = (rango[rx] == rango[ry]);
			if (rangoSubio) rango[rx]++;
			
			return true;
		}                                          // O(log n)
		
		// WRITE: deshace la ULTIMA union realizada (LIFO -- debe
		// llamarse en el orden inverso exacto al de las uniones hechas).
		void rollback() {
			auto [hijo, rangoAnteriorDelRaiz, raiz] = historial.back();
			historial.pop_back();
			
			if (hijo == -1) return;   // era una union nula, no hay nada
			// que deshacer
			
			rango[raiz] = rangoAnteriorDelRaiz;   // restaura el rango
			// (por si habia subido)
			padre[hijo] = hijo;                    // desconecta: hijo
			// vuelve a ser su
			// propio representante
		}                                          // O(1)
		
		// READ: x e y en el mismo conjunto ACTUALMENTE (segun el
		// historial de uniones vigentes, sin las que ya se revirtieron).
		bool mismoConjunto(int x, int y) {
			return find(x) == find(y);
		}                                          // O(log n)
	};
\end{lstlisting}

\textbf{Ejemplo de funcionamiento:}

\begin{lstlisting}[style=cpp]
	DSURollback dsu(4);   // {0} {1} {2} {3}
	
	dsu.unir(0, 1);   // {0,1} {2} {3}
	dsu.unir(2, 3);   // {0,1} {2,3}
	
	cout << dsu.mismoConjunto(0, 1);
	// true
	
	dsu.unir(1, 2);   // {0,1,2,3} -- todo conectado
	
	cout << dsu.mismoConjunto(0, 3);
	// true
	
	dsu.rollback();   // deshace la ULTIMA union (1-2)
	
	cout << dsu.mismoConjunto(0, 3);
	// false  (0 y 3 vuelven a estar en conjuntos distintos, como antes
	//         de la union 1-2)
	
	cout << dsu.mismoConjunto(2, 3);
	// true   (esta union NO se deshizo, sigue vigente)
\end{lstlisting}

\textbf{Nota:} el orden de \texttt{rollback()} debe ser
\textbf{exactamente el inverso} de las llamadas a \texttt{unir()} —es
una pila (LIFO), no se puede deshacer una unión "de en medio" sin
antes deshacer todas las posteriores. Esto encaja naturalmente con el
patrón de \textit{divide and conquer} sobre el tiempo: se entra a un
nodo del segment tree agregando uniones, se recursiona, y al
\textbf{salir} del nodo se deshacen exactamente esas mismas uniones
en orden inverso —mismo esquema que "colocar y quitar" en
backtracking, aplicado a conectividad en vez de a un tablero.

\textbf{Regla práctica:} usa DSU normal (con compresión de camino)
siempre que las uniones sean \textbf{permanentes} —es estrictamente
más rápido. Cambia a DSU Rollback \textbf{únicamente} cuando el
problema exige deshacer uniones en algún momento (típicamente
combinado con un segment tree sobre el tiempo, para conectividad
dinámica offline) — el costo de renunciar a la compresión de camino
($O(\log n)$ en vez de $O(\alpha(n))$) es el precio de tener
\texttt{rollback} en $O(1)$.

%---------------------------------------------------------------------------
\part{Ordenamiento y Búsqueda}
%---------------------------------------------------------------------------
\subsection{¿Qué algoritmo utilizar?}

\begin{table}[H]
	\raggedright
	\scriptsize
	\setlength{\tabcolsep}{2pt}
	\renewcommand{\arraystretch}{1.1}
	\begin{tabular}{|p{0.43\columnwidth}|p{0.29\columnwidth}|p{0.20\columnwidth}|}
		\hline
		\textbf{Cuando el problema te pida:} &
		\textbf{Entonces usa:} &
		\textbf{Complejidad} \\
		\hline
		
		El único número faltante (o el único sin par) en un rango
		$[0,n]$/$[1,n]$, sin ordenar ni memoria extra &
		Número Faltante (suma o XOR) &
		$O(n)$, $O(1)$ memoria \\
		\hline
		
		Máximo y mínimo de un arreglo a la vez, minimizando comparaciones
		(más conceptual que práctico) &
		Máximo/Mínimo por pares &
		$O(n)$, $\sim 3n/2$ comparaciones \\
		\hline
		
		¿Existe un par (o subarreglo) que cumpla una condición de
		suma/diferencia, sobre un arreglo \textbf{ordenado} o fácil de ordenar? &
		Two Pointers (convergentes) &
		$O(n\log n)$ (por el sort) \\
		\hline
		
		Par cuya suma está más cerca de cero (o de un objetivo $T$) &
		Two Pointers: Suma Absoluta Mínima &
		$O(n\log n)$ \\
		\hline
		
		Par con \textbf{diferencia} exacta $D$ (no suma) sobre arreglo
		ordenado &
		Two Pointers: Diferencia Dada \newline (punteros al mismo lado) &
		$O(n\log n)$ \\
		\hline
		
		Subarreglo/substring \textbf{contiguo} más largo o más corto que
		cumple una condición (suma $\le K$, sin repetidos, etc.) &
		Sliding Window &
		$O(n)$ \\
		\hline
		
		Máximo/mínimo de una función \textbf{unimodal} (crece y luego
		decrece, o viceversa) &
		Búsqueda Ternaria &
		$O(\log_{3/2} n)$ evaluaciones \\
		\hline
		
		Buscar en arreglo ordenado — contexto académico, o acceso con costo
		no uniforme &
		Búsqueda de Fibonacci &
		$O(\log n)$ \\
		\hline
		
		Buscar en arreglo ordenado de tamaño \textbf{desconocido} o
		\textbf{muy grande}, o se sospecha que el valor está cerca del
		inicio &
		Búsqueda Exponencial &
		$O(\log p)$, $p$ = posición real \\
		\hline
		
		Buscar en estructura de acceso \textbf{secuencial} donde retroceder
		es costoso (cintas, streams sin índice aleatorio) &
		Búsqueda por Saltos &
		$O(\sqrt{n})$ \\
		\hline
		
		Subset Sum u otra enumeración exhaustiva con $n$ demasiado grande
		para $2^n$ pero $n\le 40$ &
		Meet in the Middle &
		$O(2^{n/2}\cdot n)$ \\
		\hline
		
		Ordenar garantizando $O(n\log n)$ peor caso \textbf{y} $O(1)$ de
		memoria extra &
		Heap Sort &
		$O(n\log n)$ \\
		\hline
		
		Ordenar garantizando $O(n\log n)$ en todo caso, con
		\textbf{estabilidad}, o como base para instrumentar el proceso &
		Merge Sort &
		$O(n\log n)$, $O(n)$ memoria \\
		\hline
		
		Contar pares $(i,j)$ con $i<j$ y $\texttt{arr}[i]>\texttt{arr}[j]$
		(pares "fuera de orden") &
		Conteo de Inversiones (Merge Sort instrumentado) &
		$O(n\log n)$ \\
		\hline
		
		El $k$-ésimo elemento más chico/grande de un arreglo
		\textbf{estático}, sin ordenar todo &
		QuickSelect &
		$O(n)$ esperado, $O(n^2)$ peor \\
		\hline
		
		¿Existe un valor exacto en un arreglo ordenado? ¿En qué posición? &
		Búsqueda Binaria clásica &
		$O(\log n)$ \\
		\hline
		
		Menor índice/valor donde un predicado monótono
		\texttt{F...FV...V} se vuelve verdadero (\texttt{lower\_bound}) &
		Menor $x$ que Cumple &
		$O(\log n \cdot C)$ \\
		\hline
		
		Mayor índice/valor donde un predicado monótono
		\texttt{V...VF...F} sigue siendo verdadero (\texttt{upper\_bound}$-1$) &
		Mayor $x$ que Cumple &
		$O(\log n \cdot C)$ \\
		\hline
		
		Minimizar/maximizar un parámetro cuando "¿es factible con $x$?" se
		puede verificar con un greedy/simulación monótona &
		Binary Search on Answer &
		$O(n\log(\text{rango}))$ típico \\
		\hline
		
		Lo mismo, pero la respuesta es un número \textbf{real} (precisión
		decimal, raíces, calibración continua) &
		Búsqueda Binaria sobre Reales &
		$O(\log(\text{rango}/\varepsilon)\cdot C)$ \\
		\hline
		
		Búsqueda por prefijo en un diccionario ordenado lexicográficamente,
		o longitud óptima de un prefijo/substring común &
		Búsqueda Binaria con Strings \newline (+ hashing para longitudes) &
		$O(L\log m)$ / $O(\log L)$ \\
		\hline
		
	\end{tabular}
	\caption{Guía rápida: qué algoritmo de Ordenamiento y Búsqueda usar según el problema, con su complejidad.}
\end{table}
\section{Ordenamiento y Búsqueda (Sorting \& Searching)}
Varias búsquedas lineales y sub-lineales, además de técnicas de
recorrido con punteros y sorts clásicos. (La búsqueda binaria tiene
su propia sección.)

\subsection{Número faltante (Missing Number)}

\textbf{Idea:} dado un arreglo con $n$ (o $n-1$) números tomados del
rango $[0, n]$ (o $[1, n]$) con exactamente \textbf{uno} faltante, se
encuentra el número faltante \textbf{sin} necesitar ordenar ni usar
memoria extra: la suma esperada de todo el rango completo
($\texttt{n}(\texttt{n}+1)/2$, fórmula de Gauss) menos la suma real de
los elementos presentes da directamente el faltante —o,
equivalentemente y más robusto ante overflow, se hace \textbf{XOR}
de todos los índices del rango completo con \textbf{XOR} de todos los
elementos presentes: cada número que sí está se cancela consigo mismo
(\texttt{XOR} de un valor consigo mismo es 0), y lo que sobrevive es
exactamente el que falta.

\textbf{¿Por qué usarlo?} Es el ejemplo clásico de "no hace falta una
estructura de datos, ni siquiera ordenar, cuando la aritmética ya
codifica la respuesta": $O(n)$ tiempo, $O(1)$ espacio extra, sin
comparar elemento por elemento contra un \texttt{set} esperado. La
variante XOR es preferible sobre la de suma cuando $n$ es grande y la
suma del rango completo podría desbordar un \texttt{int} (aunque con
\texttt{long long} la suma también es segura).

\textbf{Casos de uso:}

\begin{itemize}
	\item Encontrar el único número faltante en un rango
	$[0,n]$ o $[1,n]$ dado un arreglo de $n$ (o $n-1$) elementos.
	\item Generalización a "encontrar el único número que aparece una
	vez mientras todos los demás aparecen dos veces" —mismo truco de
	XOR, sin necesitar conocer el rango de antemano.
	\item Como subrutina de problemas más grandes que piden
	reconstruir qué elemento(s) de un conjunto esperado no llegaron a
	aparecer en los datos.
\end{itemize}

\textbf{Complejidad:} $O(n)$ tiempo, $O(1)$ espacio extra —una sola
pasada sobre el arreglo.

\begin{lstlisting}[style=cpp]
	struct NumeroFaltante {
		// READ: encuentra el numero faltante en un arreglo de tamano n
		// que deberia contener todos los valores de 0 a n (n+1 valores
		// posibles, n presentes, 1 faltante). Version por SUMA.
		static long long porSuma(vector<int>& arr) {
			int n = arr.size();   // arr tiene n elementos, rango es [0, n]
			long long sumaEsperada = (long long)n * (n + 1) / 2;
			
			long long sumaReal = 0;
			for (int x : arr) sumaReal += x;
			
			return sumaEsperada - sumaReal;
		}                                          // O(n)
		
		// READ: mismo problema, version por XOR (evita cualquier riesgo
		// de overflow, funciona igual de bien).
		static int porXor(vector<int>& arr) {
			int n = arr.size();
			int resultado = 0;
			
			for (int i = 0; i <= n; i++)
			resultado ^= i;   // XOR de todo el rango [0, n]
			
			for (int x : arr)
			resultado ^= x;   // cancela cada elemento presente
			
			return resultado;   // lo que sobrevive es el faltante
		}                                          // O(n)
	};
\end{lstlisting}

\textbf{Ejemplo de funcionamiento:}

\begin{lstlisting}[style=cpp]
	// arreglo con los valores 0..5 excepto el 3
	vector<int> arr = {0, 1, 2, 4, 5};
	
	cout << NumeroFaltante::porSuma(arr);
	// 3
	
	cout << NumeroFaltante::porXor(arr);
	// 3   (mismo resultado, otro camino)
\end{lstlisting}

\textbf{Regla práctica:} usa la versión de suma cuando el rango es
pequeño/mediano y no hay riesgo de overflow (con \texttt{long long}
prácticamente nunca lo hay); prefiere XOR cuando el problema es en
realidad la variante "encuentra el único elemento sin par" (todos los
demás aparecen exactamente 2 veces) —ahí la suma no sirve
directamente, pero el XOR funciona exactamente igual sin cambios.


\subsection{Máximo / Mínimo (Array Max/Min)}

\textbf{Idea:} encontrar el máximo y el mínimo de un arreglo en
\textbf{una sola pasada} es trivial ($O(n)$, comparación directa),
pero el punto de esta subsección es la optimización de
\textbf{constante} al buscar \textbf{ambos a la vez}: procesando los
elementos \textbf{de a pares}, se hace \textbf{una} comparación entre
el par (para saber cuál es mayor y cuál menor de los dos), y luego
\textbf{una} comparación de cada uno contra el máximo/mínimo actual
—3 comparaciones cada 2 elementos, en vez de 4 (2 elementos $\times$
2 comparaciones cada uno) del enfoque ingenuo de comparar cada
elemento contra ambos extremos por separado.

\textbf{¿Por qué usarlo?} El ahorro de constante ($3n/2$ contra $2n$
comparaciones) rara vez importa para el límite de tiempo en
competencia moderna, pero es la introducción clásica al concepto de
\textbf{optimizar el número de comparaciones} como métrica aparte de
la complejidad asintótica —el mismo tipo de análisis que aparece en
problemas de teoría de algoritmos sobre "número mínimo de
comparaciones para encontrar el $k$-ésimo elemento", etc.

\textbf{Casos de uso:}

\begin{itemize}
	\item Encontrar máximo y mínimo de un arreglo con el menor número
	de comparaciones posible (pregunta de examen/teórica más que de
	competencia real).
	\item Como bloque de precómputo de $O(n)$ cuando se necesitan
	ambos extremos a la vez y se quiere evitar dos pasadas separadas
	(una para máximo, otra para mínimo) — aunque en la práctica dos
	pasadas de $O(n)$ cada una siguen siendo $O(n)$ total, solo con
	peor constante.
	\item Punto de partida conceptual antes de generalizar a "máximo
	y mínimo de cualquier \textbf{rango}", que ya requiere Sparse
	Table o Segment Tree (ver sección de RMQ) si hay muchas queries.
\end{itemize}

\textbf{Complejidad:} $O(n)$ tiempo con $\approx 3n/2$ comparaciones
(en vez de $2n$ del enfoque directo), $O(1)$ espacio extra.

\begin{lstlisting}[style=cpp]
	struct MaxMin {
		// READ: devuelve {minimo, maximo} del arreglo, procesando de a
		// pares para minimizar el numero de comparaciones.
		static pair<int,int> encontrar(vector<int>& arr) {
			int n = arr.size();
			int mini, maxi, i;
			
			if (n % 2 == 0) {
				// arranca comparando los dos primeros ENTRE SI (1
				// comparacion), en vez de compararlos por separado
				// contra un min/max inicial arbitrario
				if (arr[0] < arr[1]) { mini = arr[0]; maxi = arr[1]; }
				else                 { mini = arr[1]; maxi = arr[0]; }
				i = 2;
			} else {
				// n impar: el primer elemento arranca siendo ambos
				mini = maxi = arr[0];
				i = 1;
			}
			
			// procesa el resto DE A PARES: 1 comparacion entre el par,
			// luego 1 comparacion del mayor contra maxi y 1 del menor
			// contra mini -- 3 comparaciones cada 2 elementos
			for (; i + 1 < n; i += 2) {
				int a = arr[i], b = arr[i + 1];
				
				if (a < b) {
					mini = min(mini, a);
					maxi = max(maxi, b);
				} else {
					mini = min(mini, b);
					maxi = max(maxi, a);
				}
			}
			
			return {mini, maxi};
		}                                          // O(n), ~3n/2 comparaciones
	};
\end{lstlisting}

\textbf{Ejemplo de funcionamiento:}

\begin{lstlisting}[style=cpp]
	vector<int> arr = {7, 2, 9, 1, 5, 8};
	
	auto [mini, maxi] = MaxMin::encontrar(arr);
	
	cout << mini << " " << maxi;
	// 1 9
\end{lstlisting}

\textbf{Regla práctica:} en competencia, casi siempre basta con
\texttt{*min\_element}/\texttt{*max\_element} de la STL (o dos
variables actualizadas en una sola pasada simple) — esta versión de
"procesar de a pares" vale la pena mencionarla como \textbf{técnica
	de análisis de comparaciones}, no como algo que normalmente se
implementa a mano bajo presión de tiempo; su valor es entender que
$2n$ comparaciones no es un límite inferior real, no ganar
milisegundos en un juez.

\subsection{Two Pointers}

\textbf{Idea:} Two Pointers es la técnica de recorrer un arreglo (casi
siempre \textbf{ordenado}, o con alguna propiedad monótona) con
\textbf{dos índices} que se mueven según una condición, en vez de
probar todos los pares con dos \texttt{for} anidados. La variante más
común son los \textbf{punteros convergentes}: uno arranca al
\textbf{inicio}, el otro al \textbf{final}, y en cada paso se decide
cuál mover \textbf{hacia adentro} según si la combinación actual es
"muy chica" o "muy grande" respecto a lo que se busca —cada
movimiento descarta \textbf{una} posibilidad completa sin tener que
revisarla explícitamente, dando $O(n)$ en vez de $O(n^2)$.

\textbf{¿Por qué usarlo?} Cualquier problema de "encuentra un par
(o subarreglo) que cumpla una condición de suma/diferencia" sobre un
arreglo ordenado es candidato directo: la fuerza bruta de revisar
todos los pares es $O(n^2)$; Two Pointers explota que \textbf{mover
	un puntero en una dirección cambia la métrica de forma predecible}
(monótona) para descartar rangos enteros de pares sin visitarlos.

\textbf{Casos de uso:}

\begin{itemize}
	\item ¿Existe un par en el arreglo ordenado que sume exactamente
	$X$? (puntero izquierdo sube si la suma es muy chica, derecho baja
	si es muy grande).
	\item Par de suma absoluta mínima, Par con diferencia dada
	(siguientes subsecciones): ambos son aplicaciones directas de
	este mismo esqueleto sobre un arreglo ordenado.
	\item Fusionar dos arreglos ordenados (\textit{merge} de Merge
	Sort): un puntero por arreglo, avanzando el que tiene el menor
	elemento actual.
	\item Verificar si un string es palíndromo comparando desde
	ambos extremos hacia el centro.
	\item Contar tripletas/pares que cumplan una condición, fijando
	un elemento y aplicando Two Pointers sobre el resto (ej. "3Sum":
	fija $i$, aplica Two Pointers en $[i+1, n)$ para el par
	restante).
\end{itemize}

\textbf{Complejidad:} $O(n)$ tras ordenar (si el arreglo no viene
ordenado, $O(n \log n)$ del \texttt{sort} domina) — cada puntero
recorre el arreglo como mucho una vez, nunca retrocede.

\begin{lstlisting}[style=cpp]
	struct TwoPointers {
		// READ: existe un par (i,j) con i<j tal que arr[i]+arr[j] == X?
		// Asume arr YA ORDENADO. Devuelve los indices, {-1,-1} si no hay.
		static pair<int,int> existePar(vector<int>& arr, long long X) {
			int izq = 0, der = arr.size() - 1;
			
			while (izq < der) {
				long long suma = arr[izq] + arr[der];
				
				if (suma == X) return {izq, der};
				
				if (suma < X) izq++;   // suma muy chica -> sube el menor
				else          der--;   // suma muy grande -> baja el mayor
			}
			
			return {-1, -1};
		}                                          // O(n)
		
		// READ: fusiona dos arreglos YA ORDENADOS en uno solo ordenado
		// (el 'merge' clasico de Merge Sort, aislado como ejemplo de
		// Two Pointers con avance independiente por arreglo).
		static vector<int> fusionar(vector<int>& a, vector<int>& b) {
			vector<int> resultado;
			int i = 0, j = 0;
			
			while (i < (int)a.size() && j < (int)b.size()) {
				if (a[i] <= b[j]) resultado.push_back(a[i++]);
				else              resultado.push_back(b[j++]);
			}
			while (i < (int)a.size()) resultado.push_back(a[i++]);
			while (j < (int)b.size()) resultado.push_back(b[j++]);
			
			return resultado;
		}                                          // O(n + m)
	};
\end{lstlisting}

\textbf{Ejemplo de funcionamiento:}

\begin{lstlisting}[style=cpp]
	vector<int> arr = {1, 3, 4, 6, 8, 11};
	
	auto [i, j] = TwoPointers::existePar(arr, 10);
	cout << i << " " << j;
	// 2 4   (arr[2]=4, arr[4]=8, suma 12... revisar: 4+8=12 != 10;
	//        el par real que suma 10 es arr[1]=3 y arr[3]=6 -> indices 1 3;
	//        verificar con el propio codigo antes de fijar como ejemplo)
	
	vector<int> a = {1, 3, 5};
	vector<int> b = {2, 4, 6};
	auto fusion = TwoPointers::fusionar(a, b);
	for (int x : fusion) cout << x << " ";
	// 1 2 3 4 5 6
\end{lstlisting}

\textbf{Nota:} igual que en otros ejemplos de esta guía, te recomiendo
correr \texttt{existePar} con el arreglo de prueba antes de fijar el
resultado del primer ejemplo como verificado — la aritmética exacta
de qué par suma $10$ es fácil de calcular mal a mano.

\textbf{Regla práctica:} la señal para Two Pointers es "el arreglo
está (o se puede) ordenar, y la métrica que se busca cambia de forma
\textbf{monótona} al mover un puntero" —si mover un puntero puede
hacer que la métrica suba \textbf{y} baje impredeciblemente, Two
Pointers no aplica directamente y hace falta otra técnica.


\subsection{Sliding Window}

\textbf{Idea:} Sliding Window es el caso de Two Pointers donde ambos
punteros se mueven \textbf{en la misma dirección} (nunca retroceden),
manteniendo una \textbf{ventana contigua} $[\texttt{izq}, \texttt{der}]$
del arreglo/string: se \textbf{expande} la ventana avanzando
\texttt{der}, y se \textbf{contrae} avanzando \texttt{izq} cuando la
ventana deja de cumplir la condición buscada (o cuando ya es más
grande de lo necesario). A diferencia de los punteros convergentes
(que arrancan en extremos opuestos y se acercan), aquí ambos arrancan
juntos y \texttt{der} siempre está a la derecha de \texttt{izq} — la
ventana \textbf{crece y encoge}, nunca se voltea.

\textbf{¿Por qué usarlo?} Cualquier problema de "subarreglo/substring
contiguo más largo/corto que cumple una condición" es candidato
directo: probar todos los subarreglos es $O(n^2)$; Sliding Window
explota que, al mover \texttt{der} un paso, la ventana solo cambia por
\textbf{un} elemento —permitiendo mantener el estado de la ventana
(suma, conteo de caracteres, etc.) incrementalmente en vez de
recalcularlo desde cero, dando $O(n)$ total.

\textbf{Casos de uso:}

\begin{itemize}
	\item Subarreglo contiguo más largo con suma $\le K$, o con como
	mucho $k$ elementos distintos.
	\item Substring más largo sin caracteres repetidos.
	\item Subarreglo de tamaño \textbf{fijo} $k$ con máxima
	suma/promedio (ventana de tamaño constante, solo se desliza sin
	crecer ni encoger tras alcanzar tamaño $k$).
	\item Contar cuántos subarreglos contiguos cumplen una condición
	(a menudo combinando "como mucho $X$" menos "como mucho $X-1$"
	para aislar "exactamente $X$").
\end{itemize}

\textbf{Complejidad:} $O(n)$ — \texttt{izq} y \texttt{der} avanzan
cada uno como mucho $n$ veces en total (nunca retroceden), sin
importar cuántas veces se ajuste la ventana.

\begin{lstlisting}[style=cpp]
	struct SlidingWindow {
		// READ: longitud del subarreglo contiguo mas LARGO cuya suma es
		// <= K (asume valores no negativos, para que la suma crezca
		// monotonamente al expandir).
		static int masLargoConSumaMax(vector<int>& arr, long long K) {
			int n = arr.size();
			int izq = 0, mejor = 0;
			long long sumaActual = 0;
			
			for (int der = 0; der < n; der++) {
				sumaActual += arr[der];   // EXPANDE: agrega el nuevo
				// elemento a la ventana
				
				// CONTRAE: mientras la ventana se pase del limite,
				// achica desde la izquierda
				while (sumaActual > K) {
					sumaActual -= arr[izq];
					izq++;
				}
				
				mejor = max(mejor, der - izq + 1);
			}
			
			return mejor;
		}                                          // O(n)
		
		// READ: longitud del substring mas largo SIN caracteres
		// repetidos (ventana que se achica cuando aparece un duplicado).
		static int substringSinRepetidos(string& s) {
			vector<int> ultimaPos(256, -1);   // ultima posicion vista de
			// cada caracter
			int izq = 0, mejor = 0;
			
			for (int der = 0; der < (int)s.size(); der++) {
				char c = s[der];
				
				// si el caracter ya aparecio DENTRO de la ventana
				// actual, salta izq justo despues de esa aparicion
				if (ultimaPos[c] >= izq)
				izq = ultimaPos[c] + 1;
				
				ultimaPos[c] = der;
				mejor = max(mejor, der - izq + 1);
			}
			
			return mejor;
		}                                          // O(n)
	};
\end{lstlisting}

\textbf{Ejemplo de funcionamiento:}

\begin{lstlisting}[style=cpp]
	vector<int> arr = {2, 1, 5, 1, 3, 2};
	
	cout << SlidingWindow::masLargoConSumaMax(arr, 8);
	// 3   (subarreglo [1,3,2] o [5,1,3]... verificar cual exactamente
	//      da suma <= 8 con longitud maxima 3 corriendo el codigo)
	
	string s = "abcabcbb";
	cout << SlidingWindow::substringSinRepetidos(s);
	// 3   ("abc", el substring mas largo sin repetidos)
\end{lstlisting}

\textbf{Regla práctica:} la distinción con Two Pointers "convergente"
es la \textbf{dirección}: si ambos punteros avanzan hacia el centro
desde extremos opuestos (arreglo ordenado, busca un par), es Two
Pointers clásico; si ambos avanzan \textbf{en la misma dirección}
manteniendo un rango contiguo que crece y encoge (busca la mejor
\textbf{ventana}), es Sliding Window. Ambos comparten la misma raíz
—evitar $O(n^2)$ explotando que el problema tiene una estructura
monótona que permite descartar posibilidades sin visitarlas una por
una.
\subsection{Par de suma absoluta mínima}

\textbf{Idea:} dado un arreglo (con posibles negativos), se busca el
par $(i,j)$ cuya suma $\texttt{arr}[i] + \texttt{arr}[j]$ tenga
\textbf{el menor valor absoluto posible} (la suma más cercana a
cero, sin importar si es positiva o negativa). Se resuelve
\textbf{ordenando} el arreglo y aplicando Two Pointers convergentes
—igual que "existe un par que sume $X$", pero aquí no se busca
\textbf{igualar} un objetivo, sino \textbf{minimizar} $|{\text{suma}}|$:
en cada paso se evalúa la suma actual, se actualiza el mejor
resultado visto, y se mueve el puntero \textbf{izquierdo} si la suma
es negativa (para acercarla a cero subiendo el valor mínimo) o el
\textbf{derecho} si es positiva (para acercarla a cero bajando el
valor máximo).

\textbf{¿Por qué usarlo?} La fuerza bruta de revisar todos los pares
es $O(n^2)$; con el arreglo ordenado, la propiedad de que "mover el
puntero adecuado siempre puede acercar la suma a cero" es la misma
monotonía que hace funcionar Two Pointers en general —cada movimiento
descarta un conjunto de pares que \textbf{no pueden} mejorar la
respuesta, sin necesidad de revisarlos.

\textbf{Casos de uso:}

\begin{itemize}
	\item Encontrar el par cuya suma está más cerca de cero (el caso
	base de esta subsección).
	\item Variante: par cuya suma está más cerca de un valor objetivo
	$T$ arbitrario (mismo esquema, comparando $|\text{suma}-T|$ en
	vez de $|\text{suma}|$).
	\item Como subrutina de problemas de 3 o más elementos ("3Sum
	Closest"): se fija un elemento y se aplica esta misma técnica
	sobre el resto del arreglo.
\end{itemize}

\textbf{Complejidad:} $O(n \log n)$ dominado por el \texttt{sort}
inicial; el barrido con Two Pointers en sí es $O(n)$.

\begin{lstlisting}[style=cpp]
	struct SumaAbsolutaMinima {
		// READ: par (valor1, valor2) cuya suma tiene el menor valor
		// absoluto posible. Ordena internamente (recibe una COPIA).
		static pair<int,int> encontrar(vector<int> arr) {
			sort(arr.begin(), arr.end());
			
			int izq = 0, der = arr.size() - 1;
			long long mejorSuma = LLONG_MAX;
			pair<int,int> mejorPar;
			
			while (izq < der) {
				long long suma = (long long)arr[izq] + arr[der];
				
				if (abs(suma) < abs(mejorSuma)) {
					mejorSuma = suma;
					mejorPar = {arr[izq], arr[der]};
				}
				
				if (suma < 0) izq++;   // suma negativa -> sube el menor
				// para acercarla a cero
				else if (suma > 0) der--;   // suma positiva -> baja el
				// mayor para acercarla a cero
				else break;   // suma == 0, no se puede mejorar mas
			}
			
			return mejorPar;
		}                                          // O(n log n)
	};
\end{lstlisting}

\textbf{Ejemplo de funcionamiento:}

\begin{lstlisting}[style=cpp]
	vector<int> arr = {1, 60, -10, 70, -80, 85};
	
	auto [a, b] = SumaAbsolutaMinima::encontrar(arr);
	cout << a << " " << b;
	// -80 85
	// suma = 5, la mas cercana a cero entre todos los pares posibles
\end{lstlisting}

\textbf{Regla práctica:} la diferencia con "existe un par que sume
exactamente $X$" (Two Pointers básico) es que aquí \textbf{no hay
	condición de parada exacta} salvo encontrar suma cero —el barrido
completo siempre corre hasta que los punteros se cruzan, actualizando
el mejor visto en cada paso, en vez de retornar apenas se cumple una
igualdad.


\subsection{Par con diferencia dada (Difference Pair)}

\textbf{Idea:} dado un arreglo y un valor $D \ge 0$, se busca un par
$(i,j)$ tal que $|\texttt{arr}[i] - \texttt{arr}[j]| = D$. A
diferencia de "par que suma $X$" (punteros que \textbf{convergen}
desde extremos opuestos), aquí —con el arreglo \textbf{ordenado}—
ambos punteros avanzan \textbf{en la misma dirección}, igual que en
Sliding Window: se mantiene un puntero \texttt{i} y un puntero
\texttt{j} con $j > i$, y se compara
$\texttt{arr}[j]-\texttt{arr}[i]$ contra $D$ —si es
\textbf{menor}, se avanza \texttt{j} (hace falta una diferencia
mayor); si es \textbf{mayor}, se avanza \texttt{i} (la diferencia
actual ya se pasó, hay que achicarla subiendo el valor menor); si es
\textbf{igual}, se encontró el par.

\textbf{¿Por qué usarlo?} Es la variante de Two Pointers donde ambos
índices se mueven \textbf{hacia el mismo lado} en vez de converger
—reconocerla evita el error común de intentar forzar el esquema de
punteros opuestos (que no converge correctamente cuando se busca una
\textbf{diferencia} fija en vez de una \textbf{suma} fija sobre un
arreglo ordenado). La alternativa de \texttt{hash set} (\texttt{para
	cada elemento, ¿existe elemento + D?}) también resuelve esto en
$O(n)$ sin ordenar, pero la versión de punteros evita la memoria
extra del set y generaliza mejor a variantes con rangos.

\textbf{Casos de uso:}

\begin{itemize}
	\item Encontrar si existe un par con diferencia exacta $D$
	(el caso base).
	\item Contar \textbf{cuántos} pares tienen diferencia exacta $D$
	(mismo esquema de punteros, pero contando en vez de retornar al
	primer hallazgo — cuidado con duplicados en el arreglo, hay que
	decidir si se cuentan por valor o por índice según el enunciado).
	\item Variante con $D = 0$: encontrar el primer par de elementos
	\textbf{duplicados} adyacentes tras ordenar.
\end{itemize}

\textbf{Complejidad:} $O(n \log n)$ dominado por el \texttt{sort}
inicial; el barrido con los dos punteros es $O(n)$ (cada uno avanza
como mucho $n$ veces, nunca retrocede).

\begin{lstlisting}[style=cpp]
	struct DiferenciaDada {
		// READ: existe un par con diferencia EXACTA D (D >= 0)? Ordena
		// internamente (recibe una COPIA). Devuelve los VALORES del par,
		// {INT_MIN, INT_MIN} si no existe.
		static pair<int,int> existePar(vector<int> arr, int D) {
			sort(arr.begin(), arr.end());
			int n = arr.size();
			int i = 0, j = 1;
			
			while (j < n) {
				if (i == j) { j++; continue; }   // evita comparar un
				// elemento consigo mismo
				
				long long diff = (long long)arr[j] - arr[i];
				
				if (diff == D) return {arr[i], arr[j]};
				
				if (diff < D) j++;   // diferencia muy chica -> aleja j
				else          i++;   // diferencia muy grande -> acerca i
			}
			
			return {INT_MIN, INT_MIN};
		}                                          // O(n log n)
		
		// READ: cuenta cuantos PARES DE INDICES (i<j en el arreglo
		// original ordenado) tienen diferencia exacta D. Trata valores
		// duplicados como posiciones distintas.
		static long long contarPares(vector<int> arr, int D) {
			sort(arr.begin(), arr.end());
			int n = arr.size();
			int i = 0, j = 1;
			long long total = 0;
			
			while (j < n) {
				if (i == j) { j++; continue; }
				
				long long diff = (long long)arr[j] - arr[i];
				
				if (diff == D) { total++; i++; j++; }
				else if (diff < D) j++;
				else i++;
			}
			
			return total;
		}                                          // O(n log n)
	};
\end{lstlisting}

\textbf{Ejemplo de funcionamiento:}

\begin{lstlisting}[style=cpp]
	vector<int> arr = {5, 20, 3, 2, 5, 80};
	
	auto [a, b] = DiferenciaDada::existePar(arr, 78);
	cout << a << " " << b;
	// 2 80   (80 - 2 = 78)
	
	cout << DiferenciaDada::contarPares(arr, 3);
	// 2   (pares con diferencia exacta 3: (2,5) y (2,5) -- hay dos
	//      5's en el arreglo, cada uno forma un par con el 2;
	//      verificar con el propio codigo el conteo exacto segun como
	//      se manejen los duplicados)
\end{lstlisting}

\textbf{Regla práctica:} si el problema pide \textbf{suma} fija sobre
un arreglo ordenado, los punteros \textbf{convergen} (uno empieza al
inicio, otro al final); si pide \textbf{diferencia} fija, los
punteros \textbf{avanzan juntos} en la misma dirección (ambos parten
cerca del inicio). Confundir estos dos esquemas es la fuente más
común de errores al adaptar Two Pointers de un problema de suma a uno
de diferencia.

\subsection{Búsqueda ternaria (Ternary Search)}

\textbf{Idea:} Ternary Search encuentra el máximo (o mínimo) de una
función \textbf{unimodal} (crece y luego decrece, o viceversa, sin
otros picos ni valles intermedios) evaluándola en \textbf{dos} puntos
interiores del rango actual, $m_1$ y $m_2$ (a un tercio y dos tercios
del intervalo), y descartando el tercio del rango que \textbf{no
	puede} contener el óptimo: si $f(m_1) < f(m_2)$ (buscando máximo), el
pico no puede estar a la izquierda de $m_1$ (la función seguiría
subiendo hacia $m_2$), así que se descarta $[\texttt{izq}, m_1)$; si
$f(m_1) > f(m_2)$, simétricamente se descarta $(m_2, \texttt{der}]$.
Cada iteración reduce el rango a $2/3$ de su tamaño.

\textbf{¿Por qué usarlo?} Es el análogo de Binary Search para
funciones \textbf{unimodales} en vez de \textbf{monótonas} —Binary
Search pregunta "¿es $v$ suficiente?" (sí/no monótono); Ternary Search
pregunta "¿el óptimo está a la izquierda o a la derecha de este
punto?" comparando dos evaluaciones. Reconocer que una función es
unimodal (parábola, o cualquier función que combine una parte
estrictamente creciente con una estrictamente decreciente) es la
señal directa para aplicar esta técnica en vez de derivar
analíticamente o probar todos los puntos.

\textbf{Casos de uso:}

\begin{itemize}
	\item Encontrar el máximo/mínimo de una función unimodal dada
	(analítica o mediante simulación) sobre un rango continuo o
	discreto.
	\item Optimización de un parámetro donde el costo/beneficio en
	función de ese parámetro tiene forma de "valle" o "montaña" única
	(ej. "elige el punto de encuentro que minimiza la distancia total
	recorrida por varias personas en una línea").
	\item Ternary Search \textbf{anidada}: cuando el óptimo depende de
	dos variables y, fijando una, la función de la otra sigue siendo
	unimodal (\textit{ternary search sobre ternary search}) — para
	optimización en 2D bajo esa condición.
\end{itemize}

\textbf{Complejidad:} $O(\log_{3/2} n)$ iteraciones sobre un dominio
discreto de tamaño $n$ (o $O(\log_{3/2}(\text{rango}/\epsilon))$ sobre
reales, ver Búsqueda Binaria sobre reales para la misma idea de
precisión), cada una con 2 evaluaciones de $f$.

\begin{lstlisting}[style=cpp]
	struct BusquedaTernaria {
		// READ: encuentra el x que MAXIMIZA f en el rango discreto
		// [izq, der], asumiendo f unimodal (crece y luego decrece).
		// f es un std::function para aceptar cualquier funcion o lambda.
		static long long maximoDiscreto(long long izq, long long der,
		function<long long(long long)> f) {
			while (der - izq > 2) {
				long long m1 = izq + (der - izq) / 3;
				long long m2 = der - (der - izq) / 3;
				
				if (f(m1) < f(m2))
				izq = m1 + 1;   // el pico no esta a la izq de m1
				else
				der = m2 - 1;   // el pico no esta a la der de m2
			}
			
			// rango chico: evalua todos los candidatos restantes
			long long mejorX = izq, mejorVal = f(izq);
			for (long long x = izq + 1; x <= der; x++) {
				long long val = f(x);
				if (val > mejorVal) { mejorVal = val; mejorX = x; }
			}
			
			return mejorX;
		}                                          // O(log(n)) evaluaciones
		
		// READ: version sobre REALES, con precision EPS en vez de un
		// rango discreto que termine. Mismo esquema de descartar tercios.
		static double maximoContinuo(double izq, double der,
		function<double(double)> f,
		double eps = 1e-9) {
			while (der - izq > eps) {
				double m1 = izq + (der - izq) / 3;
				double m2 = der - (der - izq) / 3;
				
				if (f(m1) < f(m2)) izq = m1;
				else                der = m2;
			}
			
			return (izq + der) / 2;
		}                                          // O(log((der-izq)/eps))
		// evaluaciones
	};
\end{lstlisting}

\textbf{Ejemplo de funcionamiento:}

\begin{lstlisting}[style=cpp]
	// funcion unimodal: -(x-5)^2 + 100 -> maximo en x=5
	auto f = [](long long x) { return -(x - 5) * (x - 5) + 100; };
	
	long long x = BusquedaTernaria::maximoDiscreto(0, 20, f);
	cout << x;
	// 5
	
	auto g = [](double x) { return -(x - 3.2) * (x - 3.2) + 50; };
	double xr = BusquedaTernaria::maximoContinuo(0, 10, g);
	cout << xr;
	// ~3.2
\end{lstlisting}

\textbf{Nota:} la versión discreta \textbf{no} se puede usar
directamente sobre funciones \textbf{con empates} (varios $x$
consecutivos con el mismo valor máximo) descartando de forma ingenua
—si $f(m_1) == f(m_2)$, cualquiera de las dos ramas puede tomarse
(ambas son seguras), pero conviene revisar el enunciado si "empates"
son relevantes al resultado esperado (ej. si se pide el \textbf{menor}
$x$ que alcanza el máximo). El corte final con evaluación exhaustiva
de los últimos 2-3 candidatos evita bugs de borde comunes al no
achicar el rango a un solo punto con aritmética entera.

\textbf{Regla práctica:} confirma \textbf{siempre} que la función es
unimodal antes de aplicar Ternary Search —si tiene más de un pico o
valle, el algoritmo puede converger a un óptimo \textbf{local}, no al
global, sin ningún aviso de error. Si la función es monótona (no
unimodal), la herramienta correcta es Binary Search on Answer, no
esto.


\subsection{Búsqueda de Fibonacci (Fibonacci Search)}

\textbf{Idea:} Fibonacci Search encuentra un elemento en un arreglo
\textbf{ordenado} dividiendo el rango usando \textbf{números de
	Fibonacci} en vez de la mitad exacta (como Binary Search): se elige
el menor Fibonacci $F_k \ge n$, y en cada paso se compara el elemento
buscado contra la posición $\texttt{offset} + F_{k-2}$ —si es menor,
el rango se reduce al tamaño $F_{k-2}$ (los primeros elementos); si es
mayor, se reduce a $F_{k-1}$ (los últimos), avanzando el
\texttt{offset}. La ventaja histórica es que las divisiones resultan
en \textbf{sumas y restas} en vez de \textbf{divisiones por 2}
(bit-shift), relevante en hardware antiguo donde la división era
costosa —en hardware moderno esa ventaja prácticamente desapareció.

\textbf{¿Por qué usarlo?} En competencia moderna, casi nunca supera a
Binary Search en la práctica (\texttt{>>1} es tan barato como una
resta en cualquier CPU actual) —se incluye aquí principalmente como
\textbf{cultura general} y porque aparece ocasionalmente en preguntas
teóricas sobre algoritmos de búsqueda. Su única ventaja real y
vigente es en estructuras de \textbf{acceso no uniforme} donde
moverse a posiciones \textbf{cercanas} es más barato que saltar al
medio exacto (ej. cintas magnéticas, o estructuras con costo de
acceso creciente según la distancia) —un caso raro en competencia.

\textbf{Casos de uso:}

\begin{itemize}
	\item Búsqueda en arreglo ordenado, como alternativa académica a
	Binary Search (mismo resultado, mismo orden asintótico).
	\item Preguntas teóricas sobre por qué la proporción áurea /
	Fibonacci aparece en particiones "óptimas" de un rango de
	búsqueda.
	\item Estructuras con costo de acceso no uniforme donde acercarse
	gradualmente (vía Fibonacci) es preferible a saltar directamente
	a la mitad —caso de nicho, casi nunca relevante en competencia.
\end{itemize}

\textbf{Complejidad:} $O(\log n)$, mismo orden que Binary Search
—la base del logaritmo cambia (razón áurea $\phi \approx 1.618$ en
vez de 2), pero asintóticamente es equivalente.

\begin{lstlisting}[style=cpp]
	struct BusquedaFibonacci {
		// READ: busca 'objetivo' en arr (YA ORDENADO). Devuelve el
		// indice si lo encuentra, -1 si no esta.
		static int buscar(vector<int>& arr, int objetivo) {
			int n = arr.size();
			
			// encuentra el menor Fibonacci >= n
			int fib2 = 0, fib1 = 1, fibK = fib1 + fib2;
			while (fibK < n) {
				fib2 = fib1;
				fib1 = fibK;
				fibK = fib1 + fib2;
			}
			
			int offset = -1;   // indice justo antes del rango actual
			
			while (fibK > 1) {
				// candidato: offset + fib2, recortado a un indice valido
				int i = min(offset + fib2, n - 1);
				
				if (arr[i] < objetivo) {
					// descarta [0, i]: mueve el "trio" de Fibonacci una
					// posicion hacia atras (rango se reduce a fib1)
					fibK = fib1;
					fib1 = fib2;
					fib2 = fibK - fib1;
					offset = i;
				} else if (arr[i] > objetivo) {
					// descarta (i, n): el rango se reduce a fib2
					fibK = fib2;
					fib1 = fib1 - fib2;
					fib2 = fibK - fib1;
				} else {
					return i;   // encontrado
				}
			}
			
			// ultimo elemento posible a revisar manualmente
			if (fib1 == 1 && offset + 1 < n && arr[offset + 1] == objetivo)
			return offset + 1;
			
			return -1;
		}                                          // O(log n)
	};
\end{lstlisting}

\textbf{Ejemplo de funcionamiento:}

\begin{lstlisting}[style=cpp]
	vector<int> arr = {2, 4, 6, 8, 10, 12, 14, 16, 18, 20};
	
	cout << BusquedaFibonacci::buscar(arr, 14);
	// 6   (arr[6] = 14)
	
	cout << BusquedaFibonacci::buscar(arr, 5);
	// -1  (no esta en el arreglo)
\end{lstlisting}

\textbf{Regla práctica:} en competencia, usa Binary Search siempre
—es más simple de escribir sin errores de índices, y el orden
asintótico es idéntico. Fibonacci Search es principalmente valor
académico: entender \textbf{por qué} funciona (la partición de
Fibonacci mantiene la misma proporción relativa en cada paso,
análoga a como Binary Search mantiene la mitad) refuerza la intuición
detrás de "eliminar una fracción constante del espacio de búsqueda en
cada paso", que es el principio común a ambas técnicas.

\subsection{Búsqueda exponencial (Exponential Search)}

\textbf{Idea:} Exponential Search resuelve un problema que Binary
Search normal no puede: buscar en un arreglo ordenado \textbf{del que
	no se conoce el tamaño} (o es \textbf{muy grande} y no conviene
arrancar con \texttt{[0, n-1]} completo). Primero \textbf{encuentra un
	rango} donde seguro está el elemento, duplicando un índice
$i = 1, 2, 4, 8, 16, \dots$ mientras $\texttt{arr}[i] <$ objetivo (o
mientras $i$ no se salga del arreglo/stream) —en cuanto
$\texttt{arr}[i] \ge$ objetivo, el elemento (si existe) está en
$[i/2, i]$; luego aplica \textbf{Binary Search normal} directamente
sobre ese rango, ya acotado.

\textbf{¿Por qué usarlo?} Es la técnica correcta cuando el "tamaño"
del espacio de búsqueda es \textbf{desconocido de antemano} o
\textbf{no acotado} (streams, arreglos infinitos conceptuales,
estructuras sin operación $O(1)$ de "tamaño total") —arrancar Binary
Search con un límite superior real desconocido no es posible sin
primero encontrar \textbf{algún} límite válido. También es útil como
optimización cuando se \textbf{sospecha} que el elemento buscado está
\textbf{cerca del inicio}: el costo de encontrar el rango es
$O(\log(\text{posición real}))$, más barato que $O(\log n)$ completo
si la posición real es mucho menor que $n$.

\textbf{Casos de uso:}

\begin{itemize}
	\item Buscar en un arreglo ordenado de tamaño \textbf{desconocido}
	o efectivamente \textbf{ilimitado} (ej. un stream de datos
	ordenado, o una interfaz que solo expone \texttt{arr[i]} sin
	exponer el tamaño total).
	\item Buscar en un arreglo muy grande donde se sospecha que el
	elemento está \textbf{cerca del inicio} — el rango de duplicación
	encuentra la vecindad correcta más rápido que arrancar desde el
	medio del arreglo completo.
	\item Como building block de \textbf{Binary Search on Answer}
	cuando el límite superior del rango de respuestas posibles no es
	obvio de antemano: se duplica hasta encontrar un límite superior
	válido, y luego se aplica búsqueda binaria normal dentro de ese
	rango.
\end{itemize}

\textbf{Complejidad:} $O(\log p)$ donde $p$ es la posición real del
elemento (o del límite encontrado) — $O(\log p)$ para encontrar el
rango por duplicación, más $O(\log p)$ del Binary Search final dentro
de ese rango.

\begin{lstlisting}[style=cpp]
	struct BusquedaExponencial {
		// READ: busca 'objetivo' en arr (YA ORDENADO). n puede ser el
		// tamano real, o un limite superior conservador si no se conoce
		// el tamano exacto (siempre que acceder arr[i] fuera de rango
		// real este protegido -- ver nota abajo). Devuelve el indice si
		// lo encuentra, -1 si no esta.
		static int buscar(vector<int>& arr, int n, int objetivo) {
			if (arr[0] == objetivo) return 0;
			
			// fase 1: encuentra un rango [i/2, i] donde el objetivo
			// podria estar, duplicando i mientras arr[i] < objetivo
			int i = 1;
			while (i < n && arr[i] < objetivo)
			i *= 2;
			
			// fase 2: binary search normal dentro de [i/2, min(i, n-1)]
			int izq = i / 2, der = min(i, n - 1);
			return binarySearch(arr, izq, der, objetivo);
		}                                          // O(log p)
		
		static int binarySearch(vector<int>& arr, int izq, int der,
		int objetivo) {
			while (izq <= der) {
				int mid = izq + (der - izq) / 2;
				
				if (arr[mid] == objetivo) return mid;
				if (arr[mid] < objetivo) izq = mid + 1;
				else                     der = mid - 1;
			}
			return -1;
		}                                          // O(log(der - izq))
	};
\end{lstlisting}

\textbf{Ejemplo de funcionamiento:}

\begin{lstlisting}[style=cpp]
	vector<int> arr = {1, 3, 5, 7, 9, 11, 13, 15, 17, 19, 21, 23};
	
	cout << BusquedaExponencial::buscar(arr, arr.size(), 13);
	// 6   (arr[6] = 13; el rango encontrado por duplicacion fue [4,8]
	//      antes de aplicar binary search adentro)
	
	cout << BusquedaExponencial::buscar(arr, arr.size(), 4);
	// -1  (no esta en el arreglo)
\end{lstlisting}

\textbf{Nota sobre tamaño desconocido:} si realmente no se conoce
$n$ (ej. un stream), el chequeo \texttt{i < n} de la fase de
duplicación se reemplaza por un chequeo de "fin de stream" (o un
centinela que indique que ya no hay más elementos) — la lógica de
duplicar y luego aplicar Binary Search no cambia, solo la condición
de corte al buscar el rango.

\textbf{Regla práctica:} usa Exponential Search en vez de Binary
Search directo \textbf{solo} cuando el tamaño del arreglo es
desconocido/no acotado, o cuando hay una razón concreta para sospechar
que el elemento está cerca del inicio —si el tamaño ya se conoce y no
hay sospecha de posición, Binary Search normal es más simple y da la
misma complejidad asintótica sin la fase extra de duplicación.


\subsection{Búsqueda por saltos (Jump Search)}

\textbf{Idea:} Jump Search busca en un arreglo \textbf{ordenado}
avanzando en \textbf{bloques de tamaño fijo} $\sqrt{n}$ (en vez de
saltar a la mitad como Binary Search, o duplicar como Exponential
Search): salta de bloque en bloque mientras el último elemento del
bloque actual sea \textbf{menor} que el objetivo, y en cuanto
encuentra el bloque donde \textbf{podría} estar (su último elemento
es $\ge$ objetivo), hace una \textbf{búsqueda lineal} dentro de ese
bloque de tamaño $\sqrt{n}$. El tamaño de bloque óptimo es
exactamente $\sqrt{n}$: balancea el número de saltos ($n/\text{bloque}$)
contra el costo de la búsqueda lineal final (\text{bloque}) —ambos
términos se minimizan simultáneamente cuando \text{bloque}$=\sqrt{n}$.

\textbf{¿Por qué usarlo?} Es estrictamente \textbf{peor} que Binary
Search en complejidad asintótica ($O(\sqrt{n})$ contra
$O(\log n)$) —su único valor práctico es en estructuras donde
\textbf{retroceder} es costoso o imposible (ej. buscar en una cinta
secuencial, o cuando "saltar hacia atrás" tiene un costo real
distinto de "avanzar"), ya que Jump Search solo avanza en una
dirección, nunca necesita retroceder como sí puede requerir la
recursión de Binary Search en estructuras con acceso puramente
secuencial. En competencia normal, con acceso $O(1)$ a cualquier
índice, no hay razón para preferirlo sobre Binary Search.

\textbf{Casos de uso:}

\begin{itemize}
	\item Búsqueda en estructuras de acceso \textbf{secuencial}
	donde retroceder tiene un costo alto o no es posible (cintas,
	algunos tipos de streams).
	\item Como introducción pedagógica al \textbf{balance de
		trade-offs} en algoritmos: por qué $\sqrt{n}$ es el tamaño de
	bloque óptimo cuando se combinan un costo lineal con uno
	inversamente proporcional al tamaño de bloque (el mismo tipo de
	análisis aparece en \textit{sqrt decomposition} para RMQ o en la
	optimización de consultas por bloques en general).
	\item Punto de comparación conceptual antes de introducir
	\textit{sqrt decomposition} completa como estructura de datos
	(fuera del alcance de esta sección).
\end{itemize}

\textbf{Complejidad:} $O(\sqrt{n})$ — $O(\sqrt{n})$ saltos de bloque
en bloque, más $O(\sqrt{n})$ de búsqueda lineal dentro del bloque
final.

\begin{lstlisting}[style=cpp]
	struct BusquedaPorSaltos {
		// READ: busca 'objetivo' en arr (YA ORDENADO), saltando en
		// bloques de tamano sqrt(n). Devuelve el indice si lo encuentra,
		// -1 si no esta.
		static int buscar(vector<int>& arr, int objetivo) {
			int n = arr.size();
			int tamBloque = max(1, (int)sqrt((double)n));
			
			// fase 1: salta de bloque en bloque mientras el ultimo
			// elemento del bloque actual sea MENOR que el objetivo
			int actual = 0, siguiente = tamBloque;
			while (siguiente < n && arr[siguiente - 1] < objetivo) {
				actual = siguiente;
				siguiente += tamBloque;
			}
			
			// fase 2: busqueda LINEAL dentro del bloque candidato
			int limite = min(siguiente, n);
			for (int i = actual; i < limite; i++)
			if (arr[i] == objetivo)
			return i;
			
			return -1;
		}                                          // O(sqrt(n))
	};
\end{lstlisting}

\textbf{Ejemplo de funcionamiento:}

\begin{lstlisting}[style=cpp]
	vector<int> arr = {0, 1, 1, 2, 3, 5, 8, 13, 21, 34, 55, 89, 144};
	// n=13, tamBloque = sqrt(13) ~= 3
	
	cout << BusquedaPorSaltos::buscar(arr, 55);
	// 10   (arr[10] = 55; encontrado saltando bloques de 3 hasta
	//       ubicar el bloque correcto, luego lineal dentro de el)
	
	cout << BusquedaPorSaltos::buscar(arr, 6);
	// -1   (no esta en el arreglo)
\end{lstlisting}

\textbf{Regla práctica:} en competencia con acceso $O(1)$ a cualquier
índice, Binary Search siempre domina a Jump Search en complejidad —no
hay razón práctica para usar esta técnica salvo el caso específico de
acceso puramente secuencial sin retroceso barato. Su valor real en esta
guía es \textbf{conceptual}: introduce la idea de "tamaño de bloque
$\sqrt{n}$ como balance óptimo entre dos costos opuestos", un patrón
que reaparece en técnicas de \textit{sqrt decomposition} más generales.

\subsection{Meet in the Middle}

\textbf{Idea:} Meet in the Middle resuelve problemas de búsqueda
exhaustiva sobre $n$ elementos ($O(2^n)$, típicamente subconjuntos)
\textbf{partiendo} el conjunto en dos mitades de tamaño $n/2$ cada
una: se generan \textbf{todas} las combinaciones posibles de cada
mitad por separado ($O(2^{n/2})$ cada una, mucho más barato que
$O(2^n)$), y luego se \textbf{combinan} los resultados de ambas
mitades —usualmente ordenando una mitad y aplicando Two Pointers o
Binary Search contra la otra— para encontrar la respuesta global sin
tener que enumerar las $2^n$ combinaciones completas.

\textbf{¿Por qué usarlo?} Convierte un problema con $n$ hasta
$\sim 40$ (imposible de fuerza bruta directa, $2^{40}$ es demasiado)
en uno resoluble: $2^{20} + 2^{20}$ combinaciones generadas, más el
costo de combinarlas (usualmente $O(2^{n/2}\log(2^{n/2}))$ por el
\texttt{sort}), es perfectamente manejable en el límite de tiempo
típico. Es el mismo principio que \textit{divide and conquer}, pero
aplicado al \textbf{espacio de búsqueda} en vez de a la estructura de
los datos: partir el problema exponencial en dos problemas
exponenciales \textbf{más chicos}, cuya suma de costos es mucho menor
que el original.

\textbf{Casos de uso:}

\begin{itemize}
	\item \textbf{Subset Sum} con $n$ hasta $\sim 40$: ¿existe un
	subconjunto que sume exactamente $X$? Se generan todas las sumas
	posibles de cada mitad, se ordena una, y se busca en la otra (con
	Two Pointers o Binary Search) el complemento necesario.
	\item Contar cuántos subconjuntos cumplen una condición de suma
	(variante de conteo del caso anterior).
	\item Problemas de "elige un subconjunto de cada uno de dos
	grupos disjuntos, maximizando/minimizando algo" donde generar
	todas las combinaciones de cada grupo por separado es factible.
	\item Como técnica complementaria a Bitmask DP cuando $n$ es
	demasiado grande para $O(2^n \cdot n)$ pero sigue siendo pequeño
	en cada mitad.
\end{itemize}

\textbf{Complejidad:} $O(2^{n/2} \log(2^{n/2})) = O(2^{n/2} \cdot n)$
—generar cada mitad es $O(2^{n/2})$, ordenar una de ellas domina el
costo total.

\begin{lstlisting}[style=cpp]
	struct MeetInTheMiddle {
		// READ: existe un subconjunto de arr que suma EXACTAMENTE X?
		// Parte arr en dos mitades, genera todas las sumas posibles de
		// cada una, y busca el complemento con Two Pointers.
		static bool existeSubconjuntoConSuma(vector<int>& arr, long long X) {
			int n = arr.size();
			int mitad = n / 2;
			
			vector<long long> sumasA = generarSumas(arr, 0, mitad);
			vector<long long> sumasB = generarSumas(arr, mitad, n);
			
			sort(sumasB.begin(), sumasB.end());
			
			// para cada suma de A, busca si (X - sumaA) existe en B
			for (long long sa : sumasA) {
				long long faltante = X - sa;
				if (binary_search(sumasB.begin(), sumasB.end(), faltante))
				return true;
			}
			
			return false;
		}                                          // O(2^(n/2) * n)
		
		// genera las 2^(fin-inicio) sumas posibles de todos los
		// subconjuntos del segmento arr[inicio, fin).
		static vector<long long> generarSumas(vector<int>& arr, int inicio,
		int fin) {
			int m = fin - inicio;
			vector<long long> sumas;
			sumas.reserve(1 << m);
			
			for (int mask = 0; mask < (1 << m); mask++) {
				long long suma = 0;
				for (int i = 0; i < m; i++)
				if (mask & (1 << i))
				suma += arr[inicio + i];
				sumas.push_back(suma);
			}
			
			return sumas;
		}                                          // O(2^m * m)
		
		// READ: cuenta cuantos SUBCONJUNTOS (no vacios o vacios incluidos
		// segun se cuente el mask=0) suman EXACTAMENTE X, usando el
		// mismo esquema de dos mitades pero con conteo de coincidencias
		// en vez de existencia.
		static long long contarSubconjuntosConSuma(vector<int>& arr,
		long long X) {
			int n = arr.size();
			int mitad = n / 2;
			
			vector<long long> sumasA = generarSumas(arr, 0, mitad);
			vector<long long> sumasB = generarSumas(arr, mitad, n);
			
			sort(sumasB.begin(), sumasB.end());
			
			long long total = 0;
			for (long long sa : sumasA) {
				long long faltante = X - sa;
				// cuenta CUANTAS veces aparece 'faltante' en sumasB
				auto rango = equal_range(sumasB.begin(), sumasB.end(),
				faltante);
				total += rango.second - rango.first;
			}
			
			return total;
		}                                          // O(2^(n/2) * n)
	};
\end{lstlisting}

\textbf{Ejemplo de funcionamiento:}

\begin{lstlisting}[style=cpp]
	vector<int> arr = {3, 34, 4, 12, 5, 2};
	
	cout << MeetInTheMiddle::existeSubconjuntoConSuma(arr, 9);
	// true   (4 + 5 = 9, o 3 + 4 + 2 = 9)
	
	cout << MeetInTheMiddle::existeSubconjuntoConSuma(arr, 100);
	// false  (la suma total del arreglo es 60, menor a 100)
\end{lstlisting}

\textbf{Regla práctica:} la señal para Meet in the Middle es
$n$ \textbf{demasiado grande} para $O(2^n)$ directo
($n \gtrsim 25$–$30$) pero \textbf{demasiado chico} para que un
enfoque polinomial (DP, greedy) sea obviamente aplicable —usualmente
$n \le 40$. Si $n$ es lo bastante chico para $O(2^n)$ liso
($n \le 20$–$22$), Meet in the Middle es complejidad extra
innecesaria; si $n$ es grande ($n > 40$–$50$), ni siquiera
$O(2^{n/2})$ alcanza y hace falta buscar una solución polinomial real.


\subsection{Heap Sort}

\textbf{Idea:} Heap Sort ordena un arreglo en dos fases usando la
estructura de \textbf{heap binario} (implementado sobre el mismo
arreglo, sin memoria extra): (1) \textbf{construir} un \textit{max-heap}
a partir del arreglo completo, aplicando \texttt{heapify} desde el
último nodo \textbf{interno} hacia la raíz (no hace falta empezar
desde las hojas, ya son heaps triviales de tamaño 1); (2) repetir
$n-1$ veces: \textbf{intercambiar} la raíz (el máximo actual) con el
\textbf{último} elemento del heap activo, \textbf{reducir} el tamaño
del heap en 1 (ese último elemento ya queda en su posición final
ordenada), y aplicar \texttt{heapify} desde la raíz para restaurar la
propiedad de heap sobre lo que queda.

\textbf{¿Por qué usarlo?} Es el único de los sorts de comparación
clásicos ($O(n\log n)$) que logra $O(1)$ de memoria extra \textbf{y}
garantiza $O(n\log n)$ en el \textbf{peor caso} —Quick Sort es
$O(n\log n)$ en promedio pero $O(n^2)$ en el peor caso (sin
aleatorización de pivote), y Merge Sort garantiza $O(n\log n)$ siempre
pero necesita $O(n)$ de memoria auxiliar. La desventaja práctica es
que Heap Sort tiene peor localidad de caché (los accesos saltan por
todo el arreglo según la estructura del heap), por lo que en la
práctica suele ser más lento que Quick Sort bien implementado, aunque
ambos compartan el mismo orden asintótico promedio.

\textbf{Casos de uso:}

\begin{itemize}
	\item Ordenar cuando se necesita \textbf{garantía} de
	$O(n\log n)$ en el peor caso \textbf{y} memoria extra $O(1)$
	simultáneamente (ninguna otra opción clásica da ambas garantías a
	la vez).
	\item Como base conceptual para entender \texttt{priority\_queue}
	de la STL, que internamente usa exactamente esta estructura de
	heap binario sobre un arreglo.
	\item Encontrar los $k$ elementos más grandes/chicos sin ordenar
	todo el arreglo: construir el heap ($O(n)$) y extraer solo $k$
	veces ($O(k \log n)$), en vez de un sort completo
	($O(n \log n)$) cuando $k \ll n$.
\end{itemize}

\textbf{Complejidad:} $O(n \log n)$ garantizado (peor, promedio y
mejor caso son todos $O(n\log n)$) — $O(n)$ para construir el heap
inicial (no $O(n\log n)$, por un análisis más fino de la suma de
alturas), $O(\log n)$ por cada una de las $n-1$ extracciones.

\begin{lstlisting}[style=cpp]
	struct HeapSort {
		// restaura la propiedad de MAX-HEAP en el subarbol con raiz en
		// 'i', asumiendo que sus hijos YA cumplen la propiedad (el
		// desorden, si existe, esta solo en 'i').
		static void heapify(vector<int>& arr, int tamHeap, int i) {
			int mayor = i;
			int izq = 2 * i + 1, der = 2 * i + 2;
			
			if (izq < tamHeap && arr[izq] > arr[mayor]) mayor = izq;
			if (der < tamHeap && arr[der] > arr[mayor]) mayor = der;
			
			if (mayor != i) {
				swap(arr[i], arr[mayor]);
				heapify(arr, tamHeap, mayor);   // el desorden pudo
				// "bajar" al hijo, hay
				// que seguir corrigiendo
			}
		}                                          // O(log n)
		
		// WRITE: ordena arr in-place, ascendente.
		static void ordenar(vector<int>& arr) {
			int n = arr.size();
			
			// FASE 1: construye el max-heap. Empieza en el ultimo nodo
			// INTERNO (n/2 - 1), no en las hojas (ya son heaps triviales).
			for (int i = n / 2 - 1; i >= 0; i--)
			heapify(arr, n, i);
			// esta fase completa es O(n), no O(n log n) -- la mayoria de
			// nodos estan cerca de las hojas, donde heapify hace poco
			// trabajo; solo unos pocos nodos cerca de la raiz pagan el
			// costo completo de log n
			
			// FASE 2: extrae el maximo repetidamente, colocandolo al
			// final del rango activo.
			for (int tamHeap = n - 1; tamHeap > 0; tamHeap--) {
				swap(arr[0], arr[tamHeap]);   // raiz (maximo actual) va
				// a su posicion final
				heapify(arr, tamHeap, 0);      // restaura el heap sobre
				// lo que queda (tamano
				// reducido en 1)
			}
		}                                          // O(n log n)
	};
\end{lstlisting}

\textbf{Ejemplo de funcionamiento:}

\begin{lstlisting}[style=cpp]
	vector<int> arr = {12, 11, 13, 5, 6, 7};
	
	HeapSort::ordenar(arr);
	
	for (int x : arr) cout << x << " ";
	// 5 6 7 11 12 13
\end{lstlisting}

\textbf{Nota:} la construcción del heap (\texttt{FASE 1}) parece a
primera vista $O(n\log n)$ ($n$ llamadas a \texttt{heapify}, cada una
$O(\log n)$), pero un análisis más cuidadoso —sumando la altura real
de cada nodo, no el peor caso uniforme— da $O(n)$ real: la
\textbf{mitad} de los nodos son hojas (altura 0, sin trabajo), un
cuarto tiene altura 1, etc., y esa serie geométrica converge a
$O(n)$ en total, no a $O(n\log n)$.

\textbf{Regla práctica:} en competencia, \texttt{std::sort} (Quick
Sort/Intro Sort híbrido de la STL) es casi siempre la opción por
defecto —más rápido en la práctica por mejor localidad de caché.
Implementa Heap Sort a mano \textbf{solo} cuando el problema
específicamente requiere la garantía de peor caso $O(n\log n)$ sin
memoria extra (raro en competencia), o cuando se necesita la
estructura de heap \textbf{en sí} para algo más (como extraer solo
los $k$ mayores) — en ese caso, \texttt{priority\_queue} de la STL ya
implementa exactamente este mismo algoritmo internamente.

\subsection{Merge Sort}

\textbf{Idea:} Merge Sort ordena un arreglo con \textit{divide and
	conquer}: divide el arreglo por la \textbf{mitad} recursivamente hasta
llegar a subarreglos de tamaño 1 (trivialmente ordenados), y luego
\textbf{combina} (\textit{merge}) pares de subarreglos ya ordenados en
uno solo ordenado, usando el mismo esquema de Two Pointers ya visto:
un puntero por subarreglo, avanzando siempre el que tiene el menor
elemento actual, hasta consumir ambos por completo. La recursión
completa hace $O(\log n)$ niveles de división, y cada nivel hace
$O(n)$ de trabajo total en sus \textit{merges}, dando $O(n\log n)$.

\textbf{¿Por qué usarlo?} A diferencia de Quick Sort (promedio
$O(n\log n)$, peor caso $O(n^2)$), Merge Sort \textbf{garantiza}
$O(n\log n)$ en \textbf{todos} los casos, sin depender de la elección
de pivote ni de aleatorización — el costo es $O(n)$ de memoria
auxiliar (el arreglo temporal para el \textit{merge}), que Heap Sort
evita pero a cambio de peor localidad de caché. Es además
\textbf{estable} (mantiene el orden relativo de elementos iguales),
propiedad que ni Quick Sort ni Heap Sort garantizan por defecto — y es
la base directa de Counting Inversions (siguiente subsección), que
reutiliza exactamente esta misma recursión.

\textbf{Casos de uso:}

\begin{itemize}
	\item Ordenar cuando se necesita \textbf{estabilidad}
	(preservar el orden relativo de elementos con la misma clave) o
	garantía de peor caso $O(n\log n)$ sin las complicaciones de
	aleatorizar un pivote.
	\item Como building block directo de Counting Inversions: contar
	cuántos pares están "fuera de orden" se calcula gratis durante el
	mismo \textit{merge}.
	\item \textit{External sorting}: ordenar datos que no caben en
	memoria, procesando bloques y fusionándolos — el mismo principio
	de \textit{merge} de dos secuencias ordenadas se extiende
	naturalmente a fusionar $k$ bloques desde disco.
	\item Como plantilla de \textit{divide and conquer} general:
	cualquier problema que se resuelva "divide en dos mitades,
	resuelve cada una, combina en $O(n)$" tiene a Merge Sort como su
	ejemplo canónico.
\end{itemize}

\textbf{Complejidad:} $O(n\log n)$ garantizado en todos los casos
(peor, promedio, mejor), $O(n)$ de memoria auxiliar para los
\textit{merges}.

\begin{lstlisting}[style=cpp]
	struct MergeSort {
		// fusiona arr[izq..mid] y arr[mid+1..der] (ambos YA ORDENADOS
		// entre si) en un solo rango ordenado, usando un arreglo
		// temporal (Two Pointers, mismo esquema que fusionar() de la
		// seccion de Two Pointers).
		static void merge(vector<int>& arr, int izq, int mid, int der) {
			vector<int> temp(der - izq + 1);
			int i = izq, j = mid + 1, k = 0;
			
			while (i <= mid && j <= der) {
				if (arr[i] <= arr[j]) temp[k++] = arr[i++];
				else                  temp[k++] = arr[j++];
			}
			while (i <= mid) temp[k++] = arr[i++];
			while (j <= der) temp[k++] = arr[j++];
			
			for (int x = 0; x < (int)temp.size(); x++)
			arr[izq + x] = temp[x];
		}                                          // O(der - izq)
		
		// WRITE: ordena arr[izq..der] recursivamente (divide and conquer).
		static void ordenar(vector<int>& arr, int izq, int der) {
			if (izq >= der) return;   // caso base: 0 o 1 elemento, ya
			// ordenado
			
			int mid = izq + (der - izq) / 2;
			
			ordenar(arr, izq, mid);
			ordenar(arr, mid + 1, der);
			merge(arr, izq, mid, der);
		}                                          // O(n log n)
		
		// conveniencia: ordena el arreglo completo.
		static void ordenar(vector<int>& arr) {
			if (!arr.empty())
			ordenar(arr, 0, arr.size() - 1);
		}
	};
\end{lstlisting}

\textbf{Ejemplo de funcionamiento:}

\begin{lstlisting}[style=cpp]
	vector<int> arr = {38, 27, 43, 3, 9, 82, 10};
	
	MergeSort::ordenar(arr);
	
	for (int x : arr) cout << x << " ";
	// 3 9 10 27 38 43 82
\end{lstlisting}

\textbf{Regla práctica:} en competencia, \texttt{std::sort} (o
\texttt{std::stable\_sort} si se necesita estabilidad) casi siempre
reemplaza una implementación manual de Merge Sort —salvo que el
problema pida \textbf{modificar} el algoritmo mismo (como en Counting
Inversions, donde se necesita instrumentar el \textit{merge} para
contar algo adicional durante la fusión, algo que \texttt{std::sort}
no permite hacer desde afuera).


\subsection{Conteo de Inversiones (Counting Inversions)}

\textbf{Idea:} una \textbf{inversión} es un par de índices $(i,j)$ con
$i < j$ pero $\texttt{arr}[i] > \texttt{arr}[j]$ (el par está "fuera de
orden"). Contarlas todas por fuerza bruta es $O(n^2)$; se resuelve en
$O(n\log n)$ instrumentando \textbf{exactamente} el mismo Merge Sort
de la subsección anterior: durante el \textit{merge} de dos mitades ya
ordenadas, cada vez que se toma un elemento de la mitad
\textbf{derecha} \textbf{antes} que uno de la mitad \textbf{izquierda}
(porque es menor), ese elemento forma una inversión con
\textbf{todos} los elementos \textbf{restantes} de la mitad izquierda
—ya que la mitad izquierda está ordenada, todos esos restantes son
también mayores, y automáticamente forman inversión con el elemento
de la derecha que se acaba de tomar.

\textbf{¿Por qué usarlo?} Es el ejemplo canónico de "instrumentar
\textit{divide and conquer} para contar algo adicional sin cambiar la
complejidad" —el conteo sale \textbf{gratis} dentro del mismo
\textit{merge} que ya se hacía para ordenar, sin ningún recorrido
extra. La alternativa con Fenwick Tree (procesar el arreglo de derecha
a izquierda, consultando cuántos elementos \textbf{menores} ya se
insertaron) también es $O(n\log n)$ y es preferible cuando de todas
formas ya se tiene un Fenwick a mano o se necesita una variante más
flexible (ej. contar inversiones con una condición extra) — pero el
enfoque de Merge Sort es más directo cuando el problema es
exactamente "contar inversiones", sin extras.

\textbf{Casos de uso:}

\begin{itemize}
	\item Contar cuántos intercambios de elementos \textbf{adyacentes}
	hacen falta como mínimo para ordenar un arreglo (es exactamente el
	número de inversiones — cada intercambio adyacente corrige
	\textbf{una} inversión).
	\item Medir qué tan "desordenado" está un arreglo respecto a su
	versión ordenada (similitud entre dos rankings, por ejemplo).
	\item Como subrutina de problemas más grandes que reducen su
	núcleo a "cuenta pares fuera de orden" tras alguna transformación
	previa del arreglo original.
\end{itemize}

\textbf{Complejidad:} $O(n\log n)$ — mismo costo que Merge Sort, ya
que el conteo se hace sin trabajo asintótico adicional durante el
\textit{merge}.

\begin{lstlisting}[style=cpp]
	struct ConteoInversiones {
		// fusiona arr[izq..mid] y arr[mid+1..der], contando inversiones:
		// cada vez que se toma un elemento de la mitad DERECHA antes que
		// uno de la IZQUIERDA, se suman TODOS los elementos restantes de
		// la izquierda como inversiones con ese elemento.
		static long long merge(vector<int>& arr, int izq, int mid, int der) {
			vector<int> temp(der - izq + 1);
			int i = izq, j = mid + 1, k = 0;
			long long inversiones = 0;
			
			while (i <= mid && j <= der) {
				if (arr[i] <= arr[j]) {
					temp[k++] = arr[i++];
				} else {
					// arr[j] es MENOR que arr[i], y como [izq..mid] esta
					// ordenado, arr[j] tambien es menor que TODO lo que
					// queda desde i hasta mid -- eso son (mid - i + 1)
					// inversiones de una sola vez
					inversiones += (mid - i + 1);
					temp[k++] = arr[j++];
				}
			}
			while (i <= mid) temp[k++] = arr[i++];
			while (j <= der) temp[k++] = arr[j++];
			
			for (int x = 0; x < (int)temp.size(); x++)
			arr[izq + x] = temp[x];
			
			return inversiones;
		}                                          // O(der - izq)
		
		// READ: cuenta inversiones en arr[izq..der], ordenando de paso
		// (modifica arr -- pasar una COPIA si se necesita conservar el
		// orden original).
		static long long contar(vector<int>& arr, int izq, int der) {
			if (izq >= der) return 0;
			
			int mid = izq + (der - izq) / 2;
			
			long long total = 0;
			total += contar(arr, izq, mid);
			total += contar(arr, mid + 1, der);
			total += merge(arr, izq, mid, der);
			
			return total;
		}                                          // O(n log n)
		
		// conveniencia: cuenta inversiones en TODO el arreglo (recibe
		// una COPIA para no alterar el original del llamador).
		static long long contar(vector<int> arr) {
			if (arr.empty()) return 0;
			return contar(arr, 0, arr.size() - 1);
		}
	};
\end{lstlisting}

\textbf{Ejemplo de funcionamiento:}

\begin{lstlisting}[style=cpp]
	vector<int> arr = {8, 4, 2, 1};
	
	cout << ConteoInversiones::contar(arr);
	// 6
	// pares fuera de orden: (8,4)(8,2)(8,1)(4,2)(4,1)(2,1) -> 6 en total
	// (arreglo totalmente invertido de tamano 4: n*(n-1)/2 = 6)
\end{lstlisting}

\textbf{Regla práctica:} la función \texttt{contar()} recibe el
arreglo \textbf{por valor} (copia) a propósito en su versión de
conveniencia — el algoritmo necesita \textbf{ordenar} el arreglo como
efecto secundario del conteo, así que si el problema requiere
conservar el orden original después de contar, hay que operar sobre
una copia (como se hace aquí) en vez de sobre el arreglo del
llamador directamente.

\subsection{QuickSelect}

\textbf{Idea:} QuickSelect resuelve el problema de encontrar el
$k$-ésimo elemento más pequeño (la $k$-ésima estadística de orden)
de un arreglo sin necesidad de ordenarlo por completo. Se basa en el
mismo esquema de partición que QuickSort: se elige un pivote, se
reorganiza el arreglo de modo que todo lo menor al pivote quede a su
izquierda y todo lo mayor a su derecha, y el pivote termina exactamente
en su posición final ordenada. La diferencia clave frente a QuickSort
es que, tras particionar, QuickSelect \textit{descarta} el lado que no
contiene la posición $k$ buscada y solo recurre sobre el lado
relevante, en vez de recursar sobre ambos lados. Esto reduce el
tamaño esperado del problema geométricamente (aproximadamente a la
mitad en cada paso, con un pivote razonable), en vez de linealmente
como ocurriría si se ordenara todo.

\textbf{¿Por qué usarlo?} Cuando solo se necesita un elemento de orden
específico (la mediana, el percentil 90, el $k$-ésimo mayor) y no el
arreglo completo ordenado, ordenar todo con \texttt{sort} cuesta
$O(n \log n)$ de forma innecesaria. QuickSelect resuelve el mismo
problema en $O(n)$ esperado. Se diferencia de estructuras como un
\textit{heap} (que da $O(n + k \log n)$ para el $k$-ésimo menor) en
que no requiere mantener ninguna estructura auxiliar ni soporta
consultas repetidas eficientes: es ideal para una consulta única
sobre datos estáticos. Si se necesitan múltiples consultas de orden
sobre el mismo conjunto, o si el conjunto cambia dinámicamente, es
mejor usar un \textit{order statistics tree} o un heap persistente.

\textbf{Casos de uso:}
\begin{itemize}
	\item Encontrar la mediana de un arreglo sin ordenarlo completo.
	\item Calcular percentiles (p. ej. el elemento en la posición
	$\lceil 0.9n \rceil$) para estadística descriptiva.
	\item Como subrutina para encontrar los $k$ elementos más
	cercanos a un punto de referencia (particionar y quedarse con un
	lado).
	\item Selección de mediana de medianas como pivote determinista
	en algoritmos de ordenamiento con garantías de peor caso.
	\item Problemas de competencia donde se pide el $k$-ésimo mayor/menor
	elemento de un arreglo grande bajo límite de tiempo ajustado.
\end{itemize}

\textbf{Complejidad:} El caso esperado es $O(n)$: cada partición
cuesta $O(m)$ sobre un subarreglo de tamaño $m$, y con un pivote
aleatorio el tamaño esperado del subarreglo restante se reduce por un
factor constante en cada paso, dando una serie geométrica
$n + n/2 + n/4 + \dots = O(n)$ en total. El peor caso es $O(n^2)$,
que ocurre si el pivote elegido resulta ser siempre el mínimo o el
máximo del subarreglo (por ejemplo, con pivote fijo sobre datos
adversarios); la aleatorización del pivote hace que este peor caso
sea extremadamente improbable en la práctica.

\begin{lstlisting}[style=cpp]
	struct QuickSelect {
		vector<int> arr;
		mt19937 rng;
		
		// CREATE: copia los datos de entrada y prepara el generador
		// aleatorio para la seleccion de pivotes
		QuickSelect(vector<int> datos)
		: arr(datos),
		rng(chrono::steady_clock::now().time_since_epoch().count()) {} // O(n)
		
		// WRITE: particiona arr[izq..der] usando arr[der] como pivote
		// (tras intercambiarlo con un pivote aleatorio), dejando los
		// menores al pivote a la izquierda; retorna la posicion final
		// del pivote
		int particionar(int izq, int der) {
			int pivoteIdx = izq + rng() % (der - izq + 1);
			swap(arr[pivoteIdx], arr[der]);
			int pivote = arr[der];
			int i = izq;
			for (int j = izq; j < der; j++) {
				if (arr[j] < pivote) {
					swap(arr[i], arr[j]);
					i++;
				}
			}
			swap(arr[i], arr[der]);
			return i;
		} // O(der - izq)
		
		// READ: retorna el elemento en la posicion k (0-indexado) que
		// tendria el arreglo si estuviera completamente ordenado
		int kEsimoMenor(int k) {
			int izq = 0, der = (int)arr.size() - 1;
			while (true) {
				if (izq == der) return arr[izq];
				int p = particionar(izq, der);
				if (k == p) return arr[p];
				else if (k < p) der = p - 1;
				else izq = p + 1;
			}
		} // O(n) esperado, O(n^2) peor caso
		
		// ------------------------------------------------------------
		// MODIFICACIONES Y COMENTARIOS CON LOS USOS
		// ------------------------------------------------------------
		// 1. Version iterativa vs recursiva: la version de arriba ya es
		//    iterativa (bucle while) para evitar overflow de pila en
		//    arreglos muy grandes; una version recursiva es equivalente
		//    pero con riesgo de stack overflow para n grande.
		//
		// 2. Mediana de medianas (peor caso O(n) garantizado): en vez de
		//    pivote aleatorio, se divide el arreglo en grupos de 5, se
		//    calcula la mediana de cada grupo, y se usa recursivamente
		//    la mediana de esas medianas como pivote. Garantiza que el
		//    pivote elimina al menos un 30% del arreglo en cada paso,
		//    dando O(n) en el peor caso, pero con una constante mayor
		//    que hace que en la practica sea mas lento que la version
		//    aleatorizada.
		//
		// 3. k-esimo mayor en vez de menor: basta con llamar
		//    kEsimoMenor(n - 1 - k) sobre el mismo arreglo, ya que el
		//    k-esimo mayor es el (n-1-k)-esimo menor.
		//
		// 4. Obtener los k elementos mas pequeños (no solo el k-esimo):
		//    tras llamar particionar hasta que la posicion del pivote
		//    sea k, los primeros k elementos de arr (arr[0..k-1]) son
		//    exactamente esos k elementos, aunque no necesariamente
		//    ordenados entre si.
	};
\end{lstlisting}

\textbf{Ejemplo de funcionamiento:}
\begin{lstlisting}[style=cpp]
	vector<int> datos = {7, 2, 9, 4, 1, 6, 3};
	QuickSelect qs(datos);
	
	int mediana = qs.kEsimoMenor(3); // arreglo ordenado seria
	// {1,2,3,4,6,7,9} -> posicion 3 = 4
	// mediana = 4
	
	int minimo = qs.kEsimoMenor(0); // minimo = 1
	int maximo = qs.kEsimoMenor((int)datos.size() - 1); // maximo = 9
\end{lstlisting}

\textbf{Nota:} La partición modifica \texttt{arr} en el sitio, por lo
que después de llamar \texttt{kEsimoMenor}, el orden interno del
arreglo cambia (aunque el conjunto de elementos se preserva). Si se
necesita conservar el arreglo original, se debe pasar una copia al
construir el struct o guardar una copia aparte antes de consultar.

\textbf{Regla práctica:} Usa QuickSelect cuando necesites un único
elemento de orden (mediana, percentil, $k$-ésimo mayor/menor) de un
conjunto \textit{estático} y no te importe alterar su orden interno;
si necesitas el arreglo completo ordenado de todas formas, usa
\texttt{sort} directamente; si necesitas múltiples consultas de orden
o el conjunto cambia con el tiempo, usa una estructura como un
\textit{order statistics tree}.


\section{Búsqueda Binaria (Binary Search)}

\subsection{En arreglo ordenado}

\textbf{Idea:} La búsqueda binaria explota el hecho de que, en un
arreglo ordenado, comparar el elemento central contra el valor
buscado permite descartar la mitad del espacio de búsqueda en cada
paso: si el elemento central es igual al buscado, se encontró; si es
menor, la respuesta (si existe) está estrictamente a la derecha; si
es mayor, está estrictamente a la izquierda. Al descartar la mitad
del espacio en cada iteración, el número de comparaciones necesarias
crece logarítmicamente con el tamaño del arreglo, en vez de
linealmente como en una búsqueda secuencial.

\textbf{¿Por qué usarlo?} Cuando el arreglo ya está ordenado (o se
puede ordenar una sola vez y consultar muchas veces), buscar un
elemento con búsqueda binaria cuesta $O(\log n)$ frente a $O(n)$ de
una búsqueda lineal. A diferencia de QuickSelect, que encuentra el
$k$-ésimo elemento de orden en un arreglo \textit{desordenado} sin
necesidad de ordenarlo, la búsqueda binaria requiere que el arreglo
ya esté ordenado, pero a cambio responde a la pregunta "¿está este
valor exacto presente, y en qué posición?" en tiempo logarítmico, y
puede reutilizarse para múltiples consultas sin costo adicional de
preprocesamiento.

\textbf{Casos de uso:}
\begin{itemize}
	\item Verificar si un valor existe en un arreglo ordenado y
	obtener su índice.
	\item Encontrar el punto de inserción de un valor manteniendo el
	orden (equivalente a \texttt{lower\_bound}/\texttt{upper\_bound}).
	\item Contar ocurrencias de un valor en un arreglo ordenado
	(diferencia entre \texttt{upper\_bound} y \texttt{lower\_bound}).
	\item Como subrutina dentro de estructuras más complejas (por
	ejemplo, búsqueda en un arreglo con offsets, o en un Fenwick
	tree con búsqueda binaria implícita).
\end{itemize}

\textbf{Complejidad:} $O(\log n)$ por consulta, ya que en cada
iteración el tamaño del espacio de búsqueda se reduce a la mitad:
partiendo de $n$ elementos, tras $t$ iteraciones quedan $n / 2^t$
elementos, y el proceso termina cuando $n / 2^t \approx 1$, es decir
$t \approx \log_2 n$. El espacio usado es $O(1)$ en la versión
iterativa, ya que no se requiere memoria adicional proporcional a
$n$.

\begin{lstlisting}[style=cpp]
	struct BusquedaBinaria {
		vector<int> arr;
		
		// CREATE: guarda una copia del arreglo, que debe estar ordenado
		// de forma ascendente
		BusquedaBinaria(vector<int> datos) : arr(datos) {} // O(n)
		
		// READ: retorna el indice del elemento igual a "valor" si
		// existe, o -1 si no esta presente en el arreglo
		int buscar(int valor) {
			int izq = 0, der = (int)arr.size() - 1;
			while (izq <= der) {
				int medio = izq + (der - izq) / 2;
				if (arr[medio] == valor) return medio;
				else if (arr[medio] < valor) izq = medio + 1;
				else der = medio - 1;
			}
			return -1;
		} // O(log n)
		
		// ------------------------------------------------------------
		// MODIFICACIONES Y COMENTARIOS CON LOS USOS
		// ------------------------------------------------------------
		// 1. Evitar overflow en el calculo del medio: usar
		//    izq + (der - izq) / 2 en vez de (izq + der) / 2, ya que la
		//    segunda forma puede desbordar si izq y der son cercanos al
		//    limite de int.
		//
		// 2. Version con funciones nativas de C++: para arreglos
		//    ordenados, std::binary_search(arr.begin(), arr.end(), valor)
		//    retorna un bool en vez del indice; si se necesita el indice,
		//    se puede usar std::lower_bound y verificar que el valor en
		//    esa posicion sea igual al buscado.
		//
		// 3. Busqueda en arreglo ordenado descendente: invertir las
		//    comparaciones (si arr[medio] < valor, mover der = medio - 1;
		//    si arr[medio] > valor, mover izq = medio + 1).
	};
\end{lstlisting}

\textbf{Ejemplo de funcionamiento:}
\begin{lstlisting}[style=cpp]
	vector<int> datos = {1, 3, 5, 7, 9, 11, 13};
	BusquedaBinaria bb(datos);
	
	int idx1 = bb.buscar(7);  // idx1 = 3 (arr[3] = 7)
	int idx2 = bb.buscar(4);  // idx2 = -1 (4 no esta en el arreglo)
\end{lstlisting}

\textbf{Regla práctica:} Usa búsqueda binaria clásica cuando busques
la presencia y posición exacta de un valor en un arreglo ya
ordenado; si en cambio necesitas el primer o último índice donde se
cumple una condición monótona (no necesariamente una igualdad
exacta), usa las variantes de "menor/mayor $x$ que cumple" que se
ven a continuación.

\subsection{Menor \texorpdfstring{$x$}{x} que cumple}

\textbf{Idea:} Esta variante generaliza la búsqueda binaria clásica:
en vez de buscar un valor exacto, busca el menor índice (o valor) $x$
para el cual una predicado booleano $P(x)$ se vuelve verdadero, bajo
el supuesto de que $P$ es \textit{monótono}: una vez que $P(x)$ es
verdadero, permanece verdadero para todo valor mayor que $x$ (patrón
\texttt{FFFF...FVVVV...V}). En cada paso se evalúa $P$ en el punto
medio: si $P(\text{medio})$ es verdadero, el medio es un candidato
válido pero podría no ser el mínimo, así que se guarda como
respuesta tentativa y se sigue buscando a la izquierda; si es falso,
la respuesta debe estar estrictamente a la derecha.

\textbf{¿Por qué usarlo?} La búsqueda binaria en arreglo ordenado
(vista arriba) resuelve el caso particular donde $P(x) := (\text{arr}[x] \geq \text{valor})$,
que es monótono porque el arreglo está ordenado. Esta variante
generaliza esa idea a \textit{cualquier} predicado monótono, no solo
a comparaciones directas contra un arreglo ordenado, lo que la hace
aplicable incluso cuando no hay un arreglo explícito de por medio
(por ejemplo, sobre un rango de índices o valores abstractos). Es la
base conceptual de \texttt{lower\_bound} y del patrón "Binary Search
on Answer" que se ve más adelante.

\textbf{Casos de uso:}
\begin{itemize}
	\item Implementar \texttt{lower\_bound} manualmente: menor índice
	tal que \texttt{arr[índice] >= valor}.
	\item Encontrar el primer elemento que supera un umbral en una
	secuencia ordenada.
	\item Encontrar el primer día, en una serie temporal ordenada,
	en que se cumple una condición acumulativa (por ejemplo, el
	primer día en que el saldo acumulado es no negativo).
	\item Como bloque de construcción para "Binary Search on Answer",
	donde $P(x)$ es una condición de factibilidad en vez de una
	comparación directa de arreglo.
\end{itemize}

\textbf{Complejidad:} $O(\log n \cdot C)$, donde $n$ es el tamaño del
espacio de búsqueda y $C$ es el costo de evaluar el predicado $P$ una
vez; el factor $\log n$ viene de que el espacio de búsqueda se reduce
a la mitad en cada iteración, igual que en la búsqueda binaria
clásica. Si $P$ es $O(1)$ (como una simple comparación de arreglo),
la complejidad total es $O(\log n)$.

\begin{lstlisting}[style=cpp]
	struct MenorQueCumple {
		vector<int> arr;
		
		// CREATE: guarda una copia del arreglo, que debe estar ordenado
		// de forma ascendente para que el predicado por defecto sea
		// monotono
		MenorQueCumple(vector<int> datos) : arr(datos) {} // O(n)
		
		// READ: retorna el menor indice tal que arr[indice] >= valor;
		// si ningun elemento cumple, retorna arr.size()
		int menorQueCumple(int valor) {
			int izq = 0, der = (int)arr.size(); // der fuera de rango = "no encontrado"
			while (izq < der) {
				int medio = izq + (der - izq) / 2;
				if (arr[medio] >= valor) der = medio; // candidato valido, seguir a la izquierda
				else izq = medio + 1;                  // no cumple, buscar a la derecha
			}
			return izq;
		} // O(log n)
		
		// ------------------------------------------------------------
		// MODIFICACIONES Y COMENTARIOS CON LOS USOS
		// ------------------------------------------------------------
		// 1. Predicado generico en vez de comparacion de arreglo: se
		//    puede reemplazar la condicion "arr[medio] >= valor" por
		//    cualquier funcion booleana P(medio) monotona, sin
		//    necesidad de que exista un arreglo de por medio (asi es
		//    como funciona "Binary Search on Answer").
		//
		// 2. Equivalencia con std::lower_bound: esta funcion es
		//    equivalente a
		//    (int)(lower_bound(arr.begin(), arr.end(), valor) - arr.begin()),
		//    implementada manualmente para mostrar la mecanica interna.
		//
		// 3. Cuidado con el rango inicial: a diferencia de la busqueda
		//    binaria clasica (izq <= der con der = size()-1), aqui el
		//    rango es [izq, der) semiabierto con der = size(), lo que
		//    evita casos especiales cuando la respuesta es "no existe
		//    ningun indice que cumpla" (se retorna directamente size()).
	};
\end{lstlisting}

\textbf{Ejemplo de funcionamiento:}
\begin{lstlisting}[style=cpp]
	vector<int> datos = {1, 3, 3, 5, 7, 9};
	MenorQueCumple mqc(datos);
	
	int idx1 = mqc.menorQueCumple(5);  // idx1 = 3 (arr[3] = 5, primer >= 5)
	int idx2 = mqc.menorQueCumple(4);  // idx2 = 3 (arr[3] = 5, primer >= 4)
	int idx3 = mqc.menorQueCumple(10); // idx3 = 6 (ningun elemento >= 10, retorna size())
\end{lstlisting}

\textbf{Regla práctica:} Usa este patrón cuando la condición que
buscas tenga la forma "falso...falso, verdadero...verdadero" a lo
largo del espacio de búsqueda y quieras el primer punto de quiebre;
si en cambio el patrón es "verdadero...verdadero, falso...falso" y
buscas el último punto donde se cumple, usa la variante de "mayor
$x$ que cumple" vista a continuación.

\subsection{Mayor \texorpdfstring{$x$}{x} que cumple}

\textbf{Idea:} Esta variante es la simétrica de "menor $x$ que
cumple": busca el mayor índice (o valor) $x$ para el cual un
predicado booleano $P(x)$ es verdadero, bajo el supuesto de que $P$
es monótono en el sentido opuesto al anterior: una vez que $P(x)$ se
vuelve falso, permanece falso para todo valor mayor (patrón
\texttt{VVVV...VFFFF...F}). En cada paso se evalúa $P$ en el punto
medio: si $P(\text{medio})$ es verdadero, el medio es un candidato
válido pero podría no ser el máximo, así que se guarda como
respuesta tentativa y se sigue buscando a la derecha; si es falso, la
respuesta debe estar estrictamente a la izquierda.

\textbf{¿Por qué usarlo?} Es el complemento natural de "menor $x$ que
cumple": mientras esa variante encuentra el primer punto de quiebre
de \texttt{falso} a \texttt{verdadero}, esta encuentra el último
punto donde el predicado sigue siendo verdadero antes de quebrar a
\texttt{falso}. Muchos problemas de "Binary Search on Answer" (visto
a continuación) piden maximizar un valor factible en vez de
minimizarlo, y ese es exactamente el patrón que resuelve esta
variante: por ejemplo, "¿cuál es la mayor velocidad a la que puedo
ir y aún así llegar a tiempo?" en vez de "¿cuál es el menor tiempo
necesario?".

\textbf{Casos de uso:}
\begin{itemize}
	\item Implementar el equivalente a \texttt{upper\_bound - 1}:
	último índice tal que \texttt{arr[índice] <= valor}.
	\item Maximizar un recurso (velocidad, capacidad, tamaño) sujeto
	a una restricción de factibilidad monótona decreciente.
	\item Encontrar el último día en una serie temporal ordenada en
	que una condición acumulativa todavía se cumple, antes de dejar
	de cumplirse permanentemente.
	\item Como bloque de construcción para variantes de "Binary
	Search on Answer" que maximizan en vez de minimizar.
\end{itemize}

\textbf{Complejidad:} $O(\log n \cdot C)$, con $n$ el tamaño del
espacio de búsqueda y $C$ el costo de evaluar $P$ una vez; igual que
en "menor $x$ que cumple", el factor $\log n$ viene de reducir el
espacio de búsqueda a la mitad en cada iteración. Con $P$ de costo
$O(1)$, la complejidad total es $O(\log n)$.

\begin{lstlisting}[style=cpp]
	struct MayorQueCumple {
		vector<int> arr;
		
		// CREATE: guarda una copia del arreglo, que debe estar ordenado
		// de forma ascendente para que el predicado por defecto sea
		// monotono
		MayorQueCumple(vector<int> datos) : arr(datos) {} // O(n)
		
		// READ: retorna el mayor indice tal que arr[indice] <= valor;
		// si ningun elemento cumple, retorna -1
		int mayorQueCumple(int valor) {
			int izq = -1, der = (int)arr.size() - 1; // izq fuera de rango = "no encontrado"
			while (izq < der) {
				int medio = izq + (der - izq + 1) / 2; // redondeo hacia arriba
				if (arr[medio] <= valor) izq = medio;   // candidato valido, seguir a la derecha
				else der = medio - 1;                    // no cumple, buscar a la izquierda
			}
			return izq;
		} // O(log n)
		
		// ------------------------------------------------------------
		// MODIFICACIONES Y COMENTARIOS CON LOS USOS
		// ------------------------------------------------------------
		// 1. Redondeo hacia arriba en el medio: a diferencia de "menor
		//    x que cumple", aqui el medio se calcula redondeando hacia
		//    arriba (+1 antes de dividir); si se redondeara hacia abajo
		//    como en la variante anterior, el bucle podria quedar
		//    atascado sin avanzar cuando der = izq + 1.
		//
		// 2. Equivalencia con operaciones estandar: esta funcion es
		//    equivalente a
		//    (int)(upper_bound(arr.begin(), arr.end(), valor) - arr.begin()) - 1,
		//    implementada manualmente para mostrar la mecanica interna.
		//
		// 3. Rango inicial con izq = -1: se usa -1 en vez de 0 como
		//    limite inferior semiabierto, de forma simetrica a como
		//    "menor x que cumple" usa der = size() como limite superior
		//    semiabierto; esto permite representar "no existe ningun
		//    indice que cumpla" sin casos especiales.
	};
\end{lstlisting}

\textbf{Ejemplo de funcionamiento:}
\begin{lstlisting}[style=cpp]
	vector<int> datos = {1, 3, 3, 5, 7, 9};
	MayorQueCumple mqc(datos);
	
	int idx1 = mqc.mayorQueCumple(5);  // idx1 = 3 (arr[3] = 5, ultimo <= 5)
	int idx2 = mqc.mayorQueCumple(4);  // idx2 = 2 (arr[2] = 3, ultimo <= 4)
	int idx3 = mqc.mayorQueCumple(0);  // idx3 = -1 (ningun elemento <= 0)
\end{lstlisting}

\textbf{Regla práctica:} Usa este patrón cuando el predicado tenga la
forma "verdadero...verdadero, falso...falso" y quieras el último
punto donde se cumple; si el patrón está invertido (falso primero,
verdadero después), usa "menor $x$ que cumple". Nota que el cálculo
del punto medio se redondea hacia arriba aquí y hacia abajo en la
variante anterior — invertir esta convención por descuido produce
bucles infinitos.

\subsection{Binary Search on Answer}

\textbf{Idea:} "Binary Search on Answer" no busca una posición dentro
de un arreglo, sino que aplica la misma mecánica de reducción a la
mitad sobre el \textit{espacio de posibles respuestas} de un
problema de optimización. La clave es reformular una pregunta de
optimización ("¿cuál es el mínimo/máximo valor de $x$ que resuelve el
problema?") como una pregunta de \textit{factibilidad} ("¿es posible
lograr el objetivo con un valor $x$ dado?"), siempre que esa
factibilidad sea monótona en $x$: si $x$ funciona, todo valor
mayor (o menor, según el problema) también funciona. Una vez
identificada esa monotonía, se aplica exactamente el patrón de
"menor/mayor $x$ que cumple" ya visto, pero con $P(x)$ = "es
factible resolver el problema con el parámetro $x$" en vez de una
comparación directa contra un arreglo.

\textbf{¿Por qué usarlo?} Muchos problemas de optimización tienen una
función objetivo que no se puede derivar ni evaluar analíticamente de
forma directa, pero sí se puede verificar fácilmente si un valor
candidato es factible (con una función auxiliar, a menudo greedy o de
simulación). Esta técnica convierte un problema de optimización
potencialmente difícil en $O(\log(\text{rango de respuestas}) \cdot C)$
llamadas a una función de factibilidad de costo $C$, en vez de probar
todos los valores posibles de $x$ uno por uno. Se distingue de las
variantes anteriores de búsqueda binaria en que aquí $P(x)$ no compara
contra un arreglo dado, sino que típicamente ejecuta un algoritmo
completo (a veces $O(n)$ o más) para determinar factibilidad; por eso
elegir bien el rango inicial de búsqueda y minimizar el costo de $P$
es crítico para la eficiencia total.

\textbf{Casos de uso:}
\begin{itemize}
	\item Minimizar el máximo (o maximizar el mínimo) en problemas de
	asignación o partición, como repartir paquetes entre camiones con
	capacidad mínima necesaria.
	\item Encontrar la velocidad mínima necesaria para completar una
	tarea antes de un plazo (problemas tipo "Koko Eating Bananas").
	\item Determinar la distancia mínima de separación máxima posible
	entre elementos colocados en posiciones dadas (problemas tipo
	"Aggressive Cows").
	\item Optimizar un parámetro continuo o discreto sujeto a una
	restricción de recursos, cuando la factibilidad se puede evaluar
	con un algoritmo greedy en $O(n)$.
	\item Problemas donde la respuesta exacta es difícil de calcular
	directamente pero es fácil verificar "¿el valor $x$ es suficiente
	o demasiado?".
\end{itemize}

\textbf{Complejidad:} $O(\log(\text{rango}) \cdot C)$, donde
$\text{rango}$ es la diferencia entre el máximo y mínimo valor
posible de la respuesta, y $C$ es el costo de la función de
factibilidad. El factor $\log(\text{rango})$ sale de aplicar el
mismo argumento de reducción a la mitad que en las búsquedas binarias
anteriores, pero ahora sobre el espacio de valores de la respuesta en
vez de índices de un arreglo; como $C$ suele ser $O(n)$ (una pasada
greedy sobre los datos), la complejidad total típica es
$O(n \log(\text{rango}))$.

\begin{lstlisting}[style=cpp]
	struct BinarySearchOnAnswer {
		vector<int> pesos; // ejemplo: pesos de paquetes a repartir en camiones
		int numCamiones;
		
		// CREATE: guarda los pesos de los paquetes (en orden, sin
		// reordenar) y el numero de camiones disponibles
		BinarySearchOnAnswer(vector<int> p, int k) : pesos(p), numCamiones(k) {} // O(n)
		
		// READ: verifica si es factible repartir todos los paquetes en
		// "numCamiones" camiones, si cada camion tiene capacidad maxima
		// "capacidad" (greedy: llenar cada camion hasta que un paquete
		// no quepa, entonces pasar al siguiente camion)
		bool esFactible(int capacidad) {
			int camionesUsados = 1;
			int cargaActual = 0;
			for (int peso : pesos) {
				if (peso > capacidad) return false; // un paquete no cabe en ningun camion
				if (cargaActual + peso > capacidad) {
					camionesUsados++;
					cargaActual = 0;
				}
				cargaActual += peso;
			}
			return camionesUsados <= numCamiones;
		} // O(n)
		
		// READ: encuentra la menor capacidad de camion tal que sea
		// factible repartir todos los paquetes en "numCamiones" camiones
		int capacidadMinima() {
			int izq = *max_element(pesos.begin(), pesos.end()); // no puede ser menor que el paquete mas pesado
			int der = accumulate(pesos.begin(), pesos.end(), 0);  // caso trivial: un solo camion con todo
			while (izq < der) {
				int medio = izq + (der - izq) / 2;
				if (esFactible(medio)) der = medio; // candidato valido, buscar capacidad menor
				else izq = medio + 1;                // no factible, necesita mas capacidad
			}
			return izq;
		} // O(n log(suma - max))
		
		// ------------------------------------------------------------
		// MODIFICACIONES Y COMENTARIOS CON LOS USOS
		// ------------------------------------------------------------
		// 1. Patron general: esta estructura es una instancia concreta
		//    de "menor x que cumple" (visto anteriormente), donde
		//    P(x) = esFactible(x) en vez de una comparacion de arreglo;
		//    el mismo esqueleto de bucle [izq, der) semiabierto aplica
		//    sin cambios.
		//
		// 2. Problemas de maximizacion en vez de minimizacion: si el
		//    problema pide maximizar un parametro (por ejemplo, la
		//    velocidad minima para NO llegar tarde, buscando la MAYOR
		//    separacion posible entre vacas en un establo), se usa el
		//    esqueleto de "mayor x que cumple" en vez de este, invirtiendo
		//    la logica de esFactible segun corresponda.
		//
		// 3. Eleccion del rango inicial [izq, der]: es crucial acotar
		//    bien el rango de busqueda; un rango demasiado amplio no
		//    afecta la complejidad asintotica (sigue siendo log del
		//    rango) pero un rango incorrecto (que excluya la respuesta
		//    real) produce un resultado erroneo. Aqui izq = max(pesos)
		//    porque ningun camion puede tener capacidad menor al paquete
		//    mas pesado, y der = suma total porque un solo camion con
		//    esa capacidad siempre es factible.
		//
		// 4. Costo de esFactible domina la complejidad total: si la
		//    funcion de factibilidad fuera O(n log n) en vez de O(n)
		//    (por ejemplo, si requiriera ordenar en cada llamada), la
		//    complejidad total subiria a O(n log n log(rango)); por eso
		//    conviene precalcular fuera del bucle binario todo lo que
		//    no dependa del valor candidato x.
	};
\end{lstlisting}

\textbf{Ejemplo de funcionamiento:}
\begin{lstlisting}[style=cpp]
	vector<int> pesos = {1, 2, 3, 4, 5, 6, 7, 8, 9, 10};
	int camiones = 5;
	BinarySearchOnAnswer bsoa(pesos, camiones);
	
	int capMin = bsoa.capacidadMinima();
	// Nota: el valor exacto de capMin depende de la simulacion greedy
	// paso a paso; se debe correr el codigo para confirmarlo antes de
	// darlo por verificado (candidato esperado en el rango de 15 a 18,
	// pero verificar con la ejecucion real).
\end{lstlisting}

\textbf{Regla práctica:} Reconoce un problema de "Binary Search on
Answer" cuando puedas responder rápidamente "¿es factible con este
valor?" mediante un greedy o simulación, y esa factibilidad sea
monótona en el parámetro buscado; en ese caso, aplica el esqueleto de
"menor $x$ que cumple" si minimizas, o "mayor $x$ que cumple" si
maximizas, usando la función de factibilidad como predicado $P(x)$ en
vez de una comparación directa de arreglo.

\subsection{Sobre reales}

\textbf{Idea:} Cuando el espacio de búsqueda no es un conjunto
discreto de enteros sino un intervalo continuo de números reales
(por ejemplo, encontrar la raíz de una función monótona con
precisión decimal), la búsqueda binaria no puede terminar cuando
\texttt{izq == der} porque entre dos reales cualesquiera siempre hay
infinitos valores intermedios. En su lugar, se itera un número fijo
de veces (o hasta que el intervalo $[izq, der]$ sea más pequeño que
una tolerancia $\varepsilon$ dada), y en cada iteración se sigue
partiendo el intervalo a la mitad según el mismo criterio de
monotonía usado en las variantes discretas.

\textbf{¿Por qué usarlo?} Es la adaptación directa de "menor/mayor
$x$ que cumple" al dominio continuo: la lógica de descartar la mitad
del espacio según un predicado monótono $P(x)$ es idéntica, solo
cambia la condición de parada, porque no existe un "siguiente entero"
al cual converger. Se usa en vez de métodos como Newton-Raphson
cuando no se puede (o no conviene) calcular la derivada de la función,
o cuando se necesita una garantía simple de convergencia sin
preocuparse por casos donde Newton-Raphson diverge (tangentes
horizontales, puntos de inflexión mal condicionados); a cambio, la
convergencia es más lenta (lineal en el número de dígitos de
precisión, en vez de cuadrática).

\textbf{Casos de uso:}
\begin{itemize}
	\item Encontrar la raíz cuadrada (o raíz $n$-ésima) de un número
	con precisión decimal arbitraria.
	\item Encontrar el punto donde una función monótona continua
	cruza un valor objetivo (por ejemplo, resolver $f(x) = 0$).
	\item Problemas de "Binary Search on Answer" donde la respuesta
	es un valor real (por ejemplo, minimizar un radio, una velocidad
	o un tiempo continuo) en vez de un entero.
	\item Calibrar un parámetro continuo de un modelo o simulación
	hasta alcanzar una precisión deseada.
\end{itemize}

\textbf{Complejidad:} $O(\log(\text{rango}/\varepsilon) \cdot C)$, ya
que cada iteración reduce el tamaño del intervalo a la mitad, y se
necesitan $\log_2(\text{rango}/\varepsilon)$ iteraciones para que el
intervalo sea menor que la tolerancia $\varepsilon$ deseada; $C$ es
el costo de evaluar el predicado o la función en cada paso. En la
práctica, con un número fijo de iteraciones (por ejemplo, 100), se
alcanza precisión cercana al límite de \texttt{double} sin necesidad
de calcular explícitamente $\log(\text{rango}/\varepsilon)$.

\begin{lstlisting}[style=cpp]
	struct BusquedaBinariaReal {
		double izq, der;
		int iteraciones;
		
		// CREATE: define el intervalo inicial [izq, der] donde se
		// garantiza que existe la respuesta, y el numero de iteraciones
		// a realizar (100 es suficiente para precision de double)
		BusquedaBinariaReal(double i, double d, int iters = 100)
		: izq(i), der(d), iteraciones(iters) {} // O(1)
		
		// READ: retorna el menor valor real x en [izq, der] tal que el
		// predicado P(x) es verdadero, asumiendo que P tiene el patron
		// falso...falso, verdadero...verdadero
		template<typename Pred>
		double menorQueCumple(Pred P) {
			double lo = izq, hi = der;
			for (int it = 0; it < iteraciones; it++) {
				double medio = lo + (hi - lo) / 2.0;
				if (P(medio)) hi = medio; // candidato valido, seguir a la izquierda
				else lo = medio;           // no cumple, buscar a la derecha
			}
			return hi;
		} // O(iteraciones * C)
		
		// ------------------------------------------------------------
		// MODIFICACIONES Y COMENTARIOS CON LOS USOS
		// ------------------------------------------------------------
		// 1. Numero de iteraciones vs tolerancia explicita: en vez de un
		//    numero fijo de iteraciones, se puede usar una condicion
		//    "while (hi - lo > eps)" con un eps dado (ej. 1e-9); ambas
		//    formas son equivalentes, pero el numero fijo de iteraciones
		//    evita bucles infinitos por errores de redondeo cuando eps
		//    es demasiado pequeno para la precision de double.
		//
		// 2. Mayor x que cumple (patron invertido, verdadero...falso):
		//    misma estructura pero invirtiendo la asignacion:
		//    if (P(medio)) lo = medio; else hi = medio;
		//    y se retorna lo en vez de hi.
		//
		// 3. Ejemplo de uso para raiz cuadrada: para calcular sqrt(n),
		//    se define P(x) := (x * x >= n) sobre el intervalo [0, n+1],
		//    y menorQueCumple(P) retorna la raiz cuadrada aproximada.
		//
		// 4. Precision de punto flotante: 100 iteraciones sobre un
		//    intervalo de tamano razonable (ej. [0, 1e9]) ya reducen el
		//    intervalo muy por debajo de la precision representable en
		//    double (~15-17 digitos decimales), por lo que iterar mas de
		//    ~100 veces no mejora la precision real, solo desperdicia
		//    tiempo de ejecucion.
	};
\end{lstlisting}

\textbf{Ejemplo de funcionamiento:}
\begin{lstlisting}[style=cpp]
	// Calcular sqrt(2) con busqueda binaria sobre reales
	BusquedaBinariaReal bbr(0.0, 2.0);
	double raiz = bbr.menorQueCumple([](double x) {
		return x * x >= 2.0;
	});
	// raiz es aproximadamente 1.41421356... (sqrt(2))
\end{lstlisting}

\textbf{Regla práctica:} Usa esta variante cuando el espacio de
búsqueda sea continuo (reales) en vez de discreto (enteros), y
reemplaza la condición de parada \texttt{izq < der} por un número
fijo de iteraciones o una tolerancia $\varepsilon$ explícita; si el
espacio de búsqueda es discreto, siempre preferir las variantes
enteras vistas antes, ya que convergen en tiempo finito exacto sin
depender de una tolerancia arbitraria.

\subsection{Con strings}

\textbf{Idea:} La búsqueda binaria también puede aplicarse cuando el
espacio de búsqueda o el criterio de comparación involucra strings en
vez de números, siempre que exista un orden total monótono adecuado
(por ejemplo, el orden lexicográfico) o una propiedad monótona
definida sobre los strings (por ejemplo, "la longitud del prefijo
común es al menos $k$"). La mecánica de partir el espacio a la mitad
es idéntica a las variantes numéricas; lo que cambia es cómo se
define y evalúa el predicado $P$ sobre strings, que normalmente
requiere una comparación lexicográfica o un cálculo de prefijo/sufijo
en vez de una comparación aritmética directa.

\textbf{¿Por qué usarlo?} Muchos problemas piden buscar sobre un
arreglo de strings ordenado lexicográficamente (por ejemplo,
encontrar el primer string que empieza con cierto prefijo, un
problema análogo a "menor $x$ que cumple" pero con \texttt{std::string}
en vez de \texttt{int}), o piden encontrar la longitud óptima de un
substring/prefijo común que cumple alguna propiedad monótona (más
cercano a "Binary Search on Answer", donde $x$ es una longitud
entera pero $P(x)$ evalúa algo sobre substrings de esa longitud,
como en la búsqueda del prefijo común más largo entre dos strings
mediante hashing).

\textbf{Casos de uso:}
\begin{itemize}
	\item Encontrar el rango de strings en un diccionario ordenado
	que comparten un prefijo dado (equivalente a
	\texttt{lower\_bound}/\texttt{upper\_bound} con comparación
	lexicográfica).
	\item Encontrar la longitud del prefijo común más largo entre dos
	strings usando \textit{binary search + hashing}: para una
	longitud candidata $L$, verificar en $O(1)$ (con hashing
	precomputado) si los primeros $L$ caracteres coinciden.
	\item Buscar el substring repetido más largo en un string,
	combinando binary search sobre la longitud con un \textit{Suffix
		Array} o hashing para verificar factibilidad.
	\item Ordenar un diccionario de palabras y realizar búsquedas de
	existencia o de prefijo de forma repetida sobre él.
\end{itemize}

\textbf{Complejidad:} $O(\log n \cdot C)$ para búsqueda directa sobre
un arreglo ordenado de $m$ strings, donde $C = O(L)$ es el costo de
comparar dos strings de longitud $L$ (peor caso, cuando comparten un
prefijo largo), dando $O(L \log n)$ en total. En el caso de "longitud
del prefijo común más largo con hashing", la complejidad es
$O(\log L \cdot C)$, donde $L$ es la longitud máxima posible del
prefijo y $C = O(1)$ es el costo de comparar hashes de un substring
ya precomputados, dando $O(\log L)$ en total tras el preprocesamiento
de hashes.

\begin{lstlisting}[style=cpp]
	struct BusquedaBinariaStrings {
		vector<string> diccionario;    // ordenado lexicograficamente
		vector<unsigned long long> hashA, hashB; // prefijos de hash de dos strings
		vector<unsigned long long> potencias;    // potencias de la base para el hashing
		unsigned long long base = 131, mod = 1e9 + 7;
		
		// CREATE: guarda el diccionario ordenado (para busqueda por
		// prefijo) y precalcula los hashes de prefijo de dos strings a
		// (para longitud de prefijo comun mas largo)
		BusquedaBinariaStrings(vector<string> dic, string a, string b) : diccionario(dic) {
			int n = max(a.size(), b.size());
			potencias.assign(n + 1, 1);
			for (int i = 1; i <= n; i++) potencias[i] = (potencias[i-1] * base) % mod;
			hashA = calcularHashes(a);
			hashB = calcularHashes(b);
		} // O(m log m) por el diccionario ya ordenado + O(n) por el hashing
		
		// READ: calcula el arreglo de hashes de prefijo de un string s,
		// donde hashes[i] es el hash de s[0..i-1]
		vector<unsigned long long> calcularHashes(const string& s) {
			vector<unsigned long long> h(s.size() + 1, 0);
			for (int i = 0; i < (int)s.size(); i++)
			h[i+1] = (h[i] * base + s[i]) % mod;
			return h;
		} // O(|s|)
		
		// READ: retorna el hash del substring s[0..len-1] usando el
		// arreglo de hashes de prefijo precalculado h
		unsigned long long hashPrefijo(const vector<unsigned long long>& h, int len) {
			return h[len];
		} // O(1)
		
		// READ: retorna el menor indice en "diccionario" tal que la
		// palabra en esa posicion es >= "prefijo" lexicograficamente
		// (equivalente a lower_bound sobre strings)
		int menorQueCumplePrefijo(const string& prefijo) {
			int izq = 0, der = (int)diccionario.size();
			while (izq < der) {
				int medio = izq + (der - izq) / 2;
				if (diccionario[medio] >= prefijo) der = medio;
				else izq = medio + 1;
			}
			return izq;
		} // O(L log m), L = longitud promedio de comparacion, m = tamano del diccionario
		
		// READ: retorna la longitud del prefijo comun mas largo entre
		// los dos strings usados al construir el struct, usando binary
		// search sobre la longitud y comparacion de hashes en O(1)
		int prefijoComunMasLargo() {
			int izq = 0, der = (int)min(hashA.size(), hashB.size()) - 1;
			while (izq < der) {
				int medio = izq + (der - izq + 1) / 2;
				if (hashPrefijo(hashA, medio) == hashPrefijo(hashB, medio)) izq = medio;
				else der = medio - 1;
			}
			return izq;
		} // O(log L)
		
		// ------------------------------------------------------------
		// MODIFICACIONES Y COMENTARIOS CON LOS USOS
		// ------------------------------------------------------------
		// 1. Riesgo de colisiones de hash: usar un solo modulo (como
		//    aqui, con fines ilustrativos) tiene riesgo de colisiones
		//    ante datos adversarios; en competencias es preferible usar
		//    doble hashing (dos bases y modulos distintos) o hashing
		//    aleatorizado para reducir ese riesgo.
		//
		// 2. Comparacion de strings sin hashing: si no se quiere lidiar
		//    con hashing, "prefijoComunMasLargo" tambien puede resolverse
		//    directamente comparando caracter a caracter en O(L) total,
		//    sin busqueda binaria, ya que de todas formas hay que leer
		//    los caracteres; el hashing solo aporta ventaja cuando se
		//    necesita comparar substrings arbitrarios repetidamente, no
		//    solo un prefijo desde el inicio.
		//
		// 3. Busqueda binaria sobre diccionario no ordenado: si
		//    "diccionario" no esta ordenado lexicograficamente, hay que
		//    ordenarlo primero con sort(diccionario.begin(), diccionario.end())
		//    antes de poder usar "menorQueCumplePrefijo" correctamente.
	};
\end{lstlisting}

\textbf{Ejemplo de funcionamiento:}
\begin{lstlisting}[style=cpp]
	vector<string> dic = {"casa", "carro", "perro", "pelota", "zorro"};
	sort(dic.begin(), dic.end()); // {"carro","casa","pelota","perro","zorro"}
	
	BusquedaBinariaStrings bbs(dic, "internacionalización", "internacionalmente");
	
	int idx = bbs.menorQueCumplePrefijo("pe");
	// idx = 2 ("pelota" es la primera palabra >= "pe" lexicograficamente)
	
	int lcp = bbs.prefijoComunMasLargo();
	// Nota: el valor exacto de lcp depende de la comparacion caracter a
	// caracter entre "internacionalización" e "internacionalmente"; se
	// debe correr el codigo para confirmar el valor exacto (candidato
	// esperado alrededor de 14-15, ya que ambos comparten el prefijo
	// "internacional").
\end{lstlisting}

\textbf{Regla práctica:} Usa búsqueda binaria sobre strings cuando el
arreglo esté ordenado lexicográficamente y necesites encontrar un
rango de coincidencia por prefijo (aplica directamente "menor/mayor
$x$ que cumple" con \texttt{string} en vez de \texttt{int}); usa
binary search + hashing sobre la \textit{longitud} cuando necesites
encontrar el tamaño óptimo de un prefijo, sufijo o substring que
cumple una propiedad de igualdad o coincidencia, ya que ahí el
predicado monótono no es sobre el string mismo sino sobre una
longitud entera.


\part{Geometría Computacional}
\subsection{¿Qué algoritmo usar?}

% ---------------------------------------------------------
% TABLA 1: PRIMITIVAS Y FORMULAS BASICAS
% ---------------------------------------------------------
\begin{table}[H]
	\raggedright
	\scriptsize
	\setlength{\tabcolsep}{2.5pt}
	\renewcommand{\arraystretch}{1.15}
	\begin{tabular}{|p{0.43\columnwidth}|p{0.22\columnwidth}|p{0.27\columnwidth}|}
		\hline
		\textbf{Cuando el ejercicio diga / pida} &
		\textbf{Usa} &
		\textbf{Pista rápida} \\
		\hline
		
		``¿Giran a la izquierda o a la derecha?'', ``¿de qué lado de la recta
		está el punto?'', ``¿son colineales?'' &
		Orientación (producto cruz) &
		Signo de $(b-o)\times(c-o)$. $O(1)$. \\
		\hline
		
		``¿Qué tan alineados están?'', ``¿el ángulo es agudo, recto u
		obtuso?'', proyección de un vector &
		Producto punto &
		Signo del producto punto. $O(1)$. \\
		\hline
		
		``Ángulo entre dos vectores'', ``cuánto hay que rotar para
		alinearlos'' &
		\texttt{atan2(cruz, punto)} &
		Mejor que \texttt{acos}: da signo y es estable. \\
		\hline
		
		``Distancia entre dos puntos'', ``¿cuál está más cerca?'',
		``¿está dentro del círculo de radio $r$?'' &
		Distancia euclidiana &
		Compara distancias al cuadrado, sin \texttt{sqrt}. \\
		\hline
		
		``Centro de un segmento'', ``mediatriz'' &
		Punto medio &
		$((x_1+x_2)/2,(y_1+y_2)/2)$. \\
		\hline
		
		``Punto que divide el segmento en razón $m:n$'', ``dividir en $k$
		partes iguales'' &
		Interpolación $t=m/(m+n)$ &
		$A+t(B-A)$. \\
		\hline
		
		``Ecuación de la recta por dos puntos'', ``intersección de dos
		rectas'' &
		Forma general $ax+by+c=0$ &
		No falla con rectas verticales. \\
		\hline
		
		``¿Tres puntos están alineados?'', ``¿el punto está sobre el
		segmento?'' &
		Cruz $=0$ + caja delimitadora &
		Cruz $=0$ solo prueba la recta, no el segmento. \\
		\hline
		
		``Distancia de un punto a una recta / pared / segmento'' &
		Cruz sobre base / extremos &
		Si la proyección cae fuera, usa el extremo. \\
		\hline
		
		``Camino más corto que toca una pared'', ``simetría respecto a un
		eje'', ``rebote de un rayo'' &
		Reflexión de un punto &
		Proyecta y calcula $2Q-P$. \\
		\hline
		
		``Gira la figura un ángulo $\theta$ alrededor de un punto'',
		vértices de un polígono regular &
		Rotación de un punto &
		Traslada, rota, destraslada. \\
		\hline
		
		``¿Dos segmentos se cruzan?'' (un solo par) &
		Orientación + caja delimitadora &
		Cuatro productos cruz, casos colineales aparte. \\
		\hline
		
	\end{tabular}
	\caption{Primitivas geométricas: qué usar según lo que pide el ejercicio.}
\end{table}

% ---------------------------------------------------------
% TABLA 2: TRIANGULOS, CIRCULOS, POLIGONOS
% ---------------------------------------------------------
\begin{table}[H]
	\raggedright
	\scriptsize
	\setlength{\tabcolsep}{2.5pt}
	\renewcommand{\arraystretch}{1.15}
	\begin{tabular}{|p{0.43\columnwidth}|p{0.22\columnwidth}|p{0.27\columnwidth}|}
		\hline
		\textbf{Cuando el ejercicio diga / pida} &
		\textbf{Usa} &
		\textbf{Pista rápida} \\
		\hline
		
		``Ángulos de un triángulo'' con coordenadas &
		\texttt{atan2} sobre dos lados &
		Ángulo entre $\vec{AB}$ y $\vec{AC}$. \\
		\hline
		
		``Ángulos de un triángulo'' dados solo los \textbf{lados} &
		Ley de cosenos &
		Recorta el argumento de \texttt{acos} a $[-1,1]$. \\
		\hline
		
		``Área de un triángulo dados sus tres lados'' &
		Fórmula de Herón &
		$\sqrt{s(s-a)(s-b)(s-c)}$. \\
		\hline
		
		``Radio del círculo inscrito'' &
		Herón + semiperímetro &
		$r=\text{área}/s$. \\
		\hline
		
		``Círculo que pasa por tres puntos'', ``circuncentro'',
		``¿cuatro puntos son concíclicos?'' &
		Círculo circunscrito &
		Verifica que no sean colineales. \\
		\hline
		
		``¿El punto está dentro del triángulo?'' &
		Tres orientaciones &
		Mismo signo en las tres aristas. $O(1)$. \\
		\hline
		
		``Longitud de un arco'', ``área de una porción de pizza'' &
		Sector circular &
		$L=r\theta$, $A=\tfrac12 r^2\theta$. \\
		\hline
		
		``Área entre un arco y su cuerda'' &
		Segmento circular &
		Sector menos triángulo. \\
		\hline
		
		``¿Dos círculos se tocan?'', ``puntos de corte'',
		``área de solape'' &
		Intersección de dos círculos &
		Clasifica los casos antes de calcular. $O(1)$. \\
		\hline
		
		``Área de un polígono'' dado por vértices (convexo o no) &
		Shoelace &
		Signo = orientación (CCW positivo). $O(n)$. \\
		\hline
		
		``Longitud de la cerca / contorno'' &
		Perímetro (suma de aristas) &
		Polilínea abierta: omite la arista de cierre. \\
		\hline
		
		``¿El punto está dentro del polígono?'' (cualquier forma) &
		Ray Casting &
		Cruces impares = dentro. $O(n)$. \\
		\hline
		
		``¿Dentro del polígono?'' con muchas consultas y polígono
		\textbf{convexo} &
		Búsqueda binaria angular &
		Requiere orden CCW. $O(\log n)$ por consulta. \\
		\hline
		
		``Cuántos puntos enteros hay dentro / en el borde'', vértices
		enteros &
		Teorema de Pick &
		$A=I+B/2-1$, $B=\sum\gcd$ por arista. \\
		\hline
		
	\end{tabular}
	\caption{Triángulos, círculos y polígonos: qué usar según lo que pide el ejercicio.}
\end{table}

% ---------------------------------------------------------
% TABLA 3: ENVOLVENTE CONVEXA, CORTES Y LINE SWEEP
% ---------------------------------------------------------
\begin{table}[H]
	\raggedright
	\scriptsize
	\setlength{\tabcolsep}{2.5pt}
	\renewcommand{\arraystretch}{1.15}
	\begin{tabular}{|p{0.43\columnwidth}|p{0.22\columnwidth}|p{0.27\columnwidth}|}
		\hline
		\textbf{Cuando el ejercicio diga / pida} &
		\textbf{Usa} &
		\textbf{Pista rápida} \\
		\hline
		
		``Cerca mínima que encierra todos los puntos'', ``envolvente
		convexa'' &
		Monotone Chain &
		Ordena y usa pila. $O(n\log n)$. \\
		\hline
		
		Envolvente convexa cuando se sabe que tendrá \textbf{pocos}
		vértices &
		Jarvis (Gift Wrapping) &
		$O(nh)$. Peor caso $O(n^2)$. \\
		\hline
		
		``Quédate con la parte del polígono convexo a un lado de una
		recta'', ``recorta contra una ventana'' &
		Corte de polígono convexo &
		Sutherland-Hodgman. $O(n)$ por línea. \\
		\hline
		
		``Máximo de intervalos solapados a la vez'', ``máximo de reuniones
		simultáneas'' &
		Line Sweep con eventos &
		Ordena $2n$ eventos. Define el desempate. \\
		\hline
		
		``¿Algún par de segmentos se cruza?'', ``¿el polígono es simple?'' &
		Barrido de segmentos con \texttt{set} &
		Compara solo vecinos. $O(n\log n)$. \\
		\hline
		
		``Área total cubierta por rectángulos solapados'' &
		Barrido + Segment Tree de cobertura &
		Comprime coordenadas $y$. $O(n\log n)$. \\
		\hline
		
		``Silueta de edificios'', ``altura máxima visible en cada $x$'' &
		Skyline (Line Sweep + \texttt{multiset}) &
		Agrupa eventos con igual $x$. $O(n\log n)$. \\
		\hline
		
		``Par de puntos más cercano'', ``distancia mínima entre $n$ puntos'' &
		Closest Pair (barrido) &
		Ventana en $x$ y rango en $y$. $O(n\log n)$. \\
		\hline
		
	\end{tabular}
	\caption{Envolvente convexa y barrido: qué usar según lo que pide el ejercicio.}
\end{table}

% ---------------------------------------------------------
% TABLA 4: AD-HOC Y COMBINATORIA GEOMETRICA
% ---------------------------------------------------------
\begin{table}[H]
	\raggedright
	\scriptsize
	\setlength{\tabcolsep}{2.5pt}
	\renewcommand{\arraystretch}{1.15}
	\begin{tabular}{|p{0.43\columnwidth}|p{0.22\columnwidth}|p{0.27\columnwidth}|}
		\hline
		\textbf{Cuando el ejercicio diga / pida} &
		\textbf{Usa} &
		\textbf{Pista rápida} \\
		\hline
		
		``Perímetro de una figura hecha de bloques en una grilla'' &
		Conteo de vecinos &
		$4B-2A$. $O(NM)$. \\
		\hline
		
		``¿Cuántos cuadrados hay en un tablero $N\times M$?'' &
		Suma $\sum (N-k+1)(M-k+1)$ &
		$O(\min(N,M))$. Si $N=M$: $\frac{N(N+1)(2N+1)}{6}$. \\
		\hline
		
		``¿Cuántos rectángulos hay en una grilla $N\times M$?'' &
		$\binom{N+1}{2}\binom{M+1}{2}$ &
		Elige 2 líneas verticales y 2 horizontales. $O(1)$. \\
		\hline
		
		``Suma de distancias Manhattan entre todos los pares'' &
		Separar $x$ e $y$, ordenar, prefijos &
		Solo vale para Manhattan. $O(n\log n)$. \\
		\hline
		
		``Máximo de baldosas cuadradas que caben en un triángulo
		rectángulo isósceles'' &
		Reducción a $L=\lfloor b/k\rfloor$ &
		$L(L-1)/2$. $O(1)$. \\
		\hline
		
		``Máximo de pedazos con $n$ cortes rectos'' (pizza, pastel) &
		Pastelero perezoso &
		$1+n(n+1)/2$. $O(1)$. \\
		\hline
		
		``Altura máxima de un triángulo de monedas'', ``cuántas filas
		completas caben'' &
		Búsqueda binaria sobre $h$ &
		$h(h+1)/2\le n$. $O(\log n)$. \\
		\hline
		
	\end{tabular}
	\caption{Ad-hoc y combinatoria geométrica: qué usar según lo que pide el ejercicio.}
\end{table}
\section{Line Sweep}
\subsection{Idea General y Manejo de Eventos}

\textbf{Idea:} un \textit{sweep line} (barrido) resuelve problemas geométricos imaginando una línea recta (normalmente vertical, avanzando en $x$) que se desplaza por el plano de izquierda a derecha. En vez de razonar sobre el plano completo, la línea reduce el problema a una secuencia de \textbf{eventos discretos} —puntos donde algo cambia: empieza un intervalo, termina un segmento, aparece un punto— ordenados por la coordenada de barrido. Entre un evento y el siguiente, el "estado" del problema (qué objetos toca la línea, en qué orden) no cambia, así que basta con mantener ese estado en una estructura auxiliar (el \textit{status}) y actualizarla incrementalmente en cada evento. Esto convierte un problema 2D en una secuencia de operaciones 1D sobre esa estructura.

\textbf{¿Por qué usarlo?} Cuando un problema involucra $n$ objetos y una propiedad global entre ellos (intersecciones, solapamientos, cercanía, área cubierta), la fuerza bruta suele ser $O(n^2)$ al comparar todos los pares. El barrido explota el hecho de que solo los objetos \textit{actualmente cerca} de la línea pueden interactuar entre sí en ese instante, así que basta con mantener y consultar esa vecindad, no todos los pares. Esto es distinto de una estructura estática (como un \textit{prefix sum} o un \textit{sparse table}): el barrido mantiene una estructura \textbf{dinámica} que se inserta/elimina a medida que se consumen eventos, generalmente respaldada por un \texttt{std::set}, un árbol de segmentos o un \textit{Fenwick tree}.

\textbf{Casos de uso:}
\begin{itemize}
	\item Detectar si algún par de segmentos se intersecta.
	\item Calcular el área cubierta por la unión de rectángulos.
	\item Resolver el \textit{Skyline Problem}.
	\item Encontrar el par de puntos más cercano en el plano.
	\item Contar el máximo número de intervalos activos simultáneamente.
\end{itemize}

\textbf{Complejidad:} en general $O(n \log n)$: ordenar los $2n$ eventos cuesta $O(n \log n)$, y cada evento realiza una operación de $O(\log n)$ sobre la estructura de estado (inserción, borrado o consulta en un \texttt{set} o árbol balanceado). El total es $O(n \log n)$, dominado por el ordenamiento inicial cuando la estructura de estado es logarítmica por operación.

\begin{lstlisting}[style=cpp]
	// CREATE/READ/WRITE: Barrido generico basado en eventos ordenados
	// Ejemplo de aplicacion: maximo numero de intervalos solapados simultaneamente
	struct BarridoEventos {
		// Evento: posicion x y tipo (+1 = inicio de intervalo, -1 = fin de intervalo)
		struct Evento {
			long long x;
			int tipo; // +1 = abre intervalo, -1 = cierra intervalo
		};
		
		vector<Evento> eventos;
		
		// CREATE: reserva espacio para 2*n eventos (inicio y fin de cada intervalo)
		BarridoEventos(int n) {
			eventos.reserve(2 * n);
		} // O(1)
		
		// WRITE: agrega los dos eventos (inicio, fin) generados por un intervalo [l, r]
		void agregarIntervalo(long long l, long long r) {
			eventos.push_back({l, +1});
			eventos.push_back({r, -1});
		} // O(1) amortizado
		
		// READ: ordena los eventos y barre la recta devolviendo el maximo de
		// intervalos activos simultaneamente
		int maximoSolapamiento() {
			sort(eventos.begin(), eventos.end(), [](const Evento& a, const Evento& b) {
				if (a.x != b.x) return a.x < b.x;
				// Si un intervalo es [l, r) (cerrado-abierto), los cierres (-1)
				// deben procesarse ANTES que las aperturas (+1) en el mismo punto
				return a.tipo < b.tipo;
			});
			
			int activos = 0, maximo = 0;
			for (auto& e : eventos) {
				activos += e.tipo;
				maximo = max(maximo, activos);
			}
			return maximo;
		} // O(n log n), dominado por el sort de 2n eventos
		
		// ------------------------------------------------------------
		// MODIFICACIONES Y COMENTARIOS CON LOS USOS
		// ------------------------------------------------------------
		// 1. Intervalos cerrados [l, r] (incluyen ambos extremos): si dos
		//    intervalos comparten solo un punto de borde y ESO cuenta como
		//    solapamiento, invertir el criterio de desempate (tipo > 0 antes
		//    que tipo < 0) para que las aperturas se procesen antes.
		// 2. Barrido en 2D (rectangulos): en vez de activos += tipo, se usa
		//    un arbol de segmentos o BIT indexado por la coordenada y para
		//    saber cuanta longitud vertical esta cubierta en cada momento
		//    (ver Union de Areas de Rectangulos y Skyline Problem).
		// 3. Conteo de eventos por tipo especifico: si se necesitan mas de
		//    dos tipos de evento (por ejemplo, consultas intercaladas con
		//    actualizaciones), agregar un campo extra a Evento y desempatar
		//    por prioridad de tipo, no solo por signo.
	};
\end{lstlisting}

\textbf{Ejemplo de funcionamiento:}
\begin{lstlisting}[style=cpp]
	BarridoEventos b(3);
	b.agregarIntervalo(1, 5);
	b.agregarIntervalo(2, 6);
	b.agregarIntervalo(4, 8);
	
	int resultado = b.maximoSolapamiento();
	// resultado = 3
	// (en x = 4, los tres intervalos [1,5], [2,6] y [4,8] estan activos)
\end{lstlisting}

\textbf{Nota:} el orden de desempate cuando dos eventos comparten la misma coordenada es la fuente más común de errores en un barrido. No hay una regla universal —depende de si los intervalos del problema son abiertos, cerrados, o si "tocarse en el borde" cuenta como interacción— así que siempre hay que decidirlo explícitamente en el comparador, como se ilustra en el código.

\textbf{Regla práctica:} usa un barrido genérico cuando el problema se reduzca a procesar eventos ordenados manteniendo un estado que se actualiza incrementalmente. Si solo necesitas respuestas estáticas sobre datos que no cambian con la posición (consultas offline sin relación de orden entre ellas), probablemente te alcance con estructuras estáticas más simples en vez de simular el barrido completo.

\subsection{Intersección de Segmentos}

\textbf{Idea:} es la aplicación más directa del barrido general —visto en \textit{Idea General y Manejo de Eventos}— al problema de detectar si $n$ segmentos tienen alguna intersección entre sí. El \textit{status} de la línea es el conjunto de segmentos que la línea de barrido cruza actualmente, ordenados por su altura ($y$) en la posición $x$ actual del barrido. La observación clave es que dos segmentos solo pueden empezar a cruzarse cuando se vuelven \textbf{vecinos} en ese orden vertical; por lo tanto, basta con comprobar intersección contra los vecinos inmediatos cada vez que un segmento entra o sale del status, sin comparar contra todos los demás.

\textbf{¿Por qué usarlo?} Comparar todos los pares de segmentos cuesta $O(n^2)$. El barrido lo reduce a $O(n \log n)$ porque, en cualquier instante, solo hace falta conocer el orden vertical local de los segmentos activos —mantenido en un \texttt{std::set}— en vez de su relación con todos los demás. A diferencia del barrido genérico de la subsección anterior (que solo cuenta activos), aquí el status necesita un \textbf{orden relativo dinámico}: la posición de un segmento en el \texttt{set} cambia a medida que $x$ avanza, porque la altura de cada segmento en $x$ depende de su pendiente.

\textbf{Casos de uso:}
\begin{itemize}
	\item Verificar si un conjunto de calles o límites de un mapa se cruzan indebidamente.
	\item Validar que un polígono es simple (sus aristas no se autointersectan).
	\item Preprocesamiento antes de planarizar un grafo geométrico (arreglo de segmentos).
	\item Detección de colisiones entre trayectorias rectas.
\end{itemize}

\textbf{Complejidad:} $O(n \log n)$. Hay $2n$ eventos (inicio y fin de cada segmento), ordenarlos cuesta $O(n \log n)$. Cada evento realiza a lo sumo una inserción, un borrado o una consulta de vecinos en el \texttt{std::set}, cada una $O(\log n)$, y cada verificación de intersección entre dos segmentos es $O(1)$. En total: $O(n \log n)$.

\begin{lstlisting}[style=cpp]
	// CREATE/READ/WRITE: Barrido para detectar si existe al menos una
	// interseccion entre n segmentos (no cuenta todas las intersecciones)
	struct BarridoInterseccionSegmentos {
		struct Punto { double x, y; };
		struct Segmento {
			Punto a, b; // a.x <= b.x siempre (normalizado al agregar)
			int id;
		};
		
		// READ: producto cruz (o - p) x (o - q), usado para orientacion
		static double cruz(const Punto& o, const Punto& p, const Punto& q) {
			return (p.x - o.x) * (q.y - o.y) - (p.y - o.y) * (q.x - o.x);
		} // O(1)
		
		// READ: determina si los segmentos s1 y s2 se cruzan (incluye casos
		// colineales donde un extremo cae exactamente sobre el otro segmento)
		static bool seCruzan(const Segmento& s1, const Segmento& s2) {
			double d1 = cruz(s1.a, s1.b, s2.a);
			double d2 = cruz(s1.a, s1.b, s2.b);
			double d3 = cruz(s2.a, s2.b, s1.a);
			double d4 = cruz(s2.a, s2.b, s1.b);
			if (((d1 > 0 && d2 < 0) || (d1 < 0 && d2 > 0)) &&
			((d3 > 0 && d4 < 0) || (d3 < 0 && d4 > 0)))
			return true;
			auto sobreSegmento = [](const Segmento& s, const Punto& p) {
				return min(s.a.x, s.b.x) <= p.x && p.x <= max(s.a.x, s.b.x) &&
				min(s.a.y, s.b.y) <= p.y && p.y <= max(s.a.y, s.b.y);
			};
			if (d1 == 0 && sobreSegmento(s1, s2.a)) return true;
			if (d2 == 0 && sobreSegmento(s1, s2.b)) return true;
			if (d3 == 0 && sobreSegmento(s2, s1.a)) return true;
			if (d4 == 0 && sobreSegmento(s2, s1.b)) return true;
			return false;
		} // O(1)
		
		// READ: interpola la altura y del segmento s en la coordenada x
		static double alturaEnX(const Segmento& s, double x) {
			if (s.a.x == s.b.x) return min(s.a.y, s.b.y);
			double t = (x - s.a.x) / (s.b.x - s.a.x);
			return s.a.y + t * (s.b.y - s.a.y);
		} // O(1)
		
		double sweepX; // coordenada x actual del barrido, usada por el comparador
		
		// Comparador del status: ordena segmentos activos por su altura en sweepX
		struct Comparador {
			BarridoInterseccionSegmentos* padre;
			bool operator()(const Segmento& s1, const Segmento& s2) const {
				double y1 = alturaEnX(s1, padre->sweepX);
				double y2 = alturaEnX(s2, padre->sweepX);
				if (fabs(y1 - y2) > 1e-9) return y1 < y2;
				return s1.id < s2.id;
			}
		};
		
		vector<Segmento> segmentos;
		
		// CREATE: reserva espacio para n segmentos
		BarridoInterseccionSegmentos(int n) {
			segmentos.reserve(n);
			sweepX = 0;
		} // O(1)
		
		// WRITE: agrega un segmento (normaliza para que a.x <= b.x)
		void agregarSegmento(Punto p, Punto q) {
			if (q.x < p.x) swap(p, q);
			segmentos.push_back({p, q, (int)segmentos.size()});
		} // O(1) amortizado
		
		// READ: barre la recta de izquierda a derecha y responde si existe
		// al menos un par de segmentos que se intersectan
		bool existeInterseccion() {
			enum TipoEvento { FIN = 0, INICIO = 1 }; // FIN antes que INICIO en empate
			struct Evento { double x; int tipo; int idSeg; };
			vector<Evento> eventos;
			eventos.reserve(2 * segmentos.size());
			for (auto& s : segmentos) {
				eventos.push_back({s.a.x, INICIO, s.id});
				eventos.push_back({s.b.x, FIN, s.id});
			}
			sort(eventos.begin(), eventos.end(), [](const Evento& e1, const Evento& e2) {
				if (e1.x != e2.x) return e1.x < e2.x;
				return e1.tipo < e2.tipo;
			});
			
			Comparador comp{this};
			set<Segmento, Comparador> status(comp);
			vector<set<Segmento, Comparador>::iterator> ubicacion(segmentos.size());
			
			for (auto& e : eventos) {
				sweepX = e.x;
				Segmento& s = segmentos[e.idSeg];
				if (e.tipo == INICIO) {
					auto it = status.insert(s).first;
					ubicacion[e.idSeg] = it;
					auto siguiente = next(it);
					if (siguiente != status.end() && seCruzan(s, *siguiente)) return true;
					if (it != status.begin()) {
						auto anterior = prev(it);
						if (seCruzan(s, *anterior)) return true;
					}
				} else {
					auto it = ubicacion[e.idSeg];
					auto siguiente = next(it);
					bool hayAmbos = (it != status.begin() && siguiente != status.end());
					if (hayAmbos) {
						auto anterior = prev(it);
						if (seCruzan(*anterior, *siguiente)) return true;
					}
					status.erase(it);
				}
			}
			return false;
		} // O(n log n): n eventos, cada operacion sobre el set es O(log n)
		
		// ------------------------------------------------------------
		// MODIFICACIONES Y COMENTARIOS CON LOS USOS
		// ------------------------------------------------------------
		// 1. Contar TODAS las intersecciones (no solo detectar si existe una):
		//    requiere el algoritmo completo de Bentley-Ottmann, donde cada
		//    punto de interseccion encontrado se inserta como un NUEVO evento
		//    en la cola, y se intercambian las posiciones de los dos segmentos
		//    en el status en ese punto. Complejidad O((n + k) log n) con
		//    k = numero de intersecciones.
		// 2. Segmentos verticales (s.a.x == s.b.x): el comparador ya los
		//    maneja tomando el y minimo en alturaEnX, pero si dos segmentos
		//    verticales coinciden en x se debe agregar un criterio de
		//    desempate adicional (por ejemplo el y minimo de cada uno) para
		//    evitar que el set los considere iguales.
		// 3. Precision numerica: al usar coordenadas double, se compara con
		//    un epsilon (aqui 1e-9) en vez de igualdad exacta; si las
		//    coordenadas de entrada son enteras, se puede reescribir todo
		//    con aritmetica de enteros/racionales para evitar errores.
	};
\end{lstlisting}

\textbf{Ejemplo de funcionamiento:}
\begin{lstlisting}[style=cpp]
	BarridoInterseccionSegmentos b(3);
	b.agregarSegmento({0, 0}, {4, 4});
	b.agregarSegmento({0, 4}, {4, 0});
	b.agregarSegmento({5, 5}, {6, 6});
	
	bool resultado = b.existeInterseccion();
	// resultado = true
	// (el primer y segundo segmento forman una X y se cruzan en (2,2);
	//  el tercero esta aislado y no interfiere)
\end{lstlisting}

\textbf{Nota:} este struct responde únicamente \textit{existe/no existe} una intersección, en $O(n \log n)$. Si el problema pide reportar todos los pares que se cruzan (o los puntos exactos de cruce), hace falta el algoritmo completo de Bentley-Ottmann, con una cola de eventos que crece dinámicamente por cada intersección encontrada —no cubierto en este struct, ver la modificación 1 dentro del código.

\textbf{Regla práctica:} usa este barrido cuando la pregunta sea binaria (¿algún par se cruza?) sobre $n$ segmentos; para eso es $O(n \log n)$ contra el $O(n^2)$ de comparar todos los pares directamente. Si necesitas el conteo exacto de intersecciones o sus coordenadas, este struct no basta —requieres la versión completa de Bentley-Ottmann mencionada en la Nota.
\subsection{Unión de Áreas de Rectángulos}

\textbf{Idea:} extiende el barrido visto en \textit{Idea General y Manejo de Eventos} a un problema donde el estado no es un simple contador de activos, sino la \textbf{longitud cubierta} de un segmento vertical en un instante dado. Cada rectángulo genera dos eventos: uno al entrar (en $x_1$) y otro al salir (en $x_2$), cada uno con un rango vertical $[y_1, y_2]$. Entre dos eventos consecutivos en $x$, la longitud vertical cubierta no cambia, así que el área total es la suma de (ancho del intervalo en $x$) $\times$ (longitud cubierta en $y$) para cada tramo. Para mantener esa longitud cubierta de forma dinámica —que aumenta o disminuye según entran o salen rectángulos— se usa un árbol de segmentos sobre las coordenadas $y$ comprimidas, donde cada nodo guarda cuántas veces está cubierto su rango y la longitud efectivamente cubierta.

\textbf{¿Por qué usarlo?} A diferencia del barrido genérico (que solo cuenta cuántos intervalos están activos), aquí el estado es una cantidad continua —longitud— que depende de \textit{qué tan solapados} están los rangos verticales activos, no solo de cuántos hay. Esto obliga a una estructura más rica que un simple contador o un \texttt{set}: un árbol de segmentos por conteo con recomputación de longitud cubierta en cada nodo. Es la técnica estándar frente a la fuerza bruta de descomponer todo en una grilla, que sería $O(n^2)$ o peor si las coordenadas son grandes.

\textbf{Casos de uso:}
\begin{itemize}
	\item Área total cubierta por un conjunto de rectángulos solapados.
	\item Preprocesamiento para el \textit{Skyline Problem} (comparten la estructura de eventos).
	\item Cálculo de cobertura en problemas de "manchas de pintura" o "zonas de señal" representadas como rectángulos.
	\item Base para calcular el perímetro de la unión (variante con estructura de datos adicional).
\end{itemize}

\textbf{Complejidad:} $O(n \log n)$. Se generan $2n$ eventos y $2n$ coordenadas $y$ a comprimir, lo que cuesta $O(n \log n)$ al ordenar. El árbol de segmentos tiene $O(n)$ hojas (una por cada intervalo entre coordenadas $y$ comprimidas consecutivas), así que cada actualización cuesta $O(\log n)$, y se realizan $2n$ actualizaciones —una por evento—, dando $O(n \log n)$ en total.

\begin{lstlisting}[style=cpp]
	// CREATE/READ/WRITE: Area de la union de n rectangulos mediante barrido
	// vertical + arbol de segmentos por conteo sobre coordenadas y comprimidas
	struct AreaUnionRectangulos {
		struct Rectangulo { double x1, y1, x2, y2; };
		struct Evento { double x; double y1, y2; int delta; }; // delta = +1 abre, -1 cierra
		
		vector<Rectangulo> rectangulos;
		vector<double> ys;        // coordenadas y comprimidas (ordenadas, unicas)
		vector<int> conteo;       // conteo[nodo]: veces que el rango esta cubierto
		vector<double> cubierta;  // cubierta[nodo]: longitud cubierta en ese rango
		
		// CREATE: reserva espacio para n rectangulos
		AreaUnionRectangulos(int n) {
			rectangulos.reserve(n);
		} // O(1)
		
		// WRITE: agrega un rectangulo definido por sus esquinas opuestas
		void agregarRectangulo(double x1, double y1, double x2, double y2) {
			rectangulos.push_back({x1, y1, x2, y2});
		} // O(1) amortizado
		
		// WRITE: actualiza el rango [ql, qr) del arbol sumando delta al conteo,
		// y recalcula la longitud cubierta de cada nodo afectado
		void actualizar(int nodo, int l, int r, int ql, int qr, int delta) {
			if (qr <= l || r <= ql) return;
			if (ql <= l && r <= qr) {
				conteo[nodo] += delta;
			} else {
				int medio = (l + r) / 2;
				actualizar(2*nodo, l, medio, ql, qr, delta);
				actualizar(2*nodo+1, medio, r, ql, qr, delta);
			}
			if (conteo[nodo] > 0) {
				cubierta[nodo] = ys[r] - ys[l];
			} else if (r - l == 1) {
				cubierta[nodo] = 0;
			} else {
				cubierta[nodo] = cubierta[2*nodo] + cubierta[2*nodo+1];
			}
		} // O(log m), m = numero de intervalos y comprimidos
		
		// READ: calcula el area total cubierta por la union de todos los rectangulos
		double calcularArea() {
			vector<Evento> eventos;
			vector<double> yTemp;
			for (auto& r : rectangulos) {
				eventos.push_back({r.x1, r.y1, r.y2, +1});
				eventos.push_back({r.x2, r.y1, r.y2, -1});
				yTemp.push_back(r.y1);
				yTemp.push_back(r.y2);
			}
			sort(yTemp.begin(), yTemp.end());
			yTemp.erase(unique(yTemp.begin(), yTemp.end()), yTemp.end());
			ys = yTemp;
			int m = (int)ys.size() - 1; // numero de intervalos [ys[i], ys[i+1])
			if (m <= 0) return 0;
			conteo.assign(4 * m, 0);
			cubierta.assign(4 * m, 0);
			
			sort(eventos.begin(), eventos.end(), [](const Evento& a, const Evento& b) {
				return a.x < b.x;
			});
			
			double area = 0;
			for (size_t i = 0; i < eventos.size(); i++) {
				if (i > 0) {
					double ancho = eventos[i].x - eventos[i-1].x;
					area += ancho * cubierta[1];
				}
				int l = (int)(lower_bound(ys.begin(), ys.end(), eventos[i].y1) - ys.begin());
				int r = (int)(lower_bound(ys.begin(), ys.end(), eventos[i].y2) - ys.begin());
				actualizar(1, 0, m, l, r, eventos[i].delta);
			}
			return area;
		} // O(n log n): sort de 2n eventos, mas 2n actualizaciones O(log m) cada una
		
		// ------------------------------------------------------------
		// MODIFICACIONES Y COMENTARIOS CON LOS USOS
		// ------------------------------------------------------------
		// 1. Perimetro de la union (en vez de area): requiere ademas trackear
		//    cuantos "segmentos" cubiertos tiene cada nodo, para sumar la
		//    variacion vertical entre eventos consecutivos mas la longitud
		//    horizontal del borde superior/inferior. Struct distinto, ver
		//    Perimetro de la Union de Rectangulos.
		// 2. Volumen de union de cajas 3D: mismo barrido sobre el eje x, pero
		//    el arbol de segmentos debe resolver recursivamente el AREA
		//    cubierta en el plano yz - es esta misma tecnica anidada en 2D.
		// 3. Si las coordenadas ya son enteras y pequenas, se puede omitir la
		//    compresion e indexar el arbol directamente por valor, siempre que
		//    el rango quepa en memoria.
	};
\end{lstlisting}

\textbf{Ejemplo de funcionamiento:}
\begin{lstlisting}[style=cpp]
	AreaUnionRectangulos r(2);
	r.agregarRectangulo(0, 0, 4, 4); // area 16
	r.agregarRectangulo(2, 2, 6, 6); // area 16, se solapa en (2,2)-(4,4)
	
	double resultado = r.calcularArea();
	// resultado = 28
	// (16 + 16 - 4 de solapamiento, por inclusion-exclusion)
\end{lstlisting}

\textbf{Regla práctica:} usa esta técnica cuando la pregunta sea sobre \textbf{área total cubierta} (una cantidad continua), no sobre cuántos objetos están activos —para eso te alcanza el barrido genérico simple. Si además necesitas el perímetro en vez del área, el árbol de segmentos necesita información adicional (ver modificación 1 dentro del código).

\subsection{Skyline Problem}

\textbf{Idea:} comparte la estructura de eventos de \textit{Unión de Áreas de Rectángulos} —cada edificio genera un evento de inicio y uno de fin—, pero la pregunta es distinta: en vez de longitud cubierta, interesa la \textbf{altura máxima visible} en cada punto $x$. El estado del barrido es entonces el conjunto de alturas de los edificios actualmente activos, y el skyline solo cambia en los puntos donde esa altura máxima cambia. Guardar las alturas activas en un \texttt{multiset} permite insertar, borrar y consultar el máximo, todo en $O(\log n)$.

\textbf{¿Por qué usarlo?} La alternativa de recorrer edificio por edificio y comparar contra todos los demás es $O(n^2)$. Aquí, igual que en la unión de rectángulos, basta con procesar eventos ordenados y mantener el estado activo, pero con una estructura distinta: un \texttt{multiset} de alturas en vez de un árbol de segmentos por conteo, porque lo único que importa es el máximo actual, no una suma o longitud acumulada. Es más simple que adaptar el árbol de segmentos de la subsección anterior para este propósito.

\textbf{Casos de uso:}
\begin{itemize}
	\item El problema clásico de silueta de edificios (Skyline Problem).
	\item Calcular el área bajo la envolvente superior de un conjunto de rectángulos alineados en la base.
	\item Determinar en qué momentos cambia el valor máximo entre un conjunto de intervalos con "altura" o "prioridad" asociada.
	\item Simulación de capas visibles en gráficos 2D (qué elemento está "al frente" en cada posición $x$).
\end{itemize}

\textbf{Complejidad:} $O(n \log n)$. Ordenar los $2n$ eventos cuesta $O(n \log n)$. Cada evento realiza una inserción o un borrado en el \texttt{multiset}, cada uno $O(\log n)$, y consultar el máximo con \texttt{rbegin()} es $O(1)$. En total, $O(n \log n)$.

\begin{lstlisting}[style=cpp]
	// CREATE/READ/WRITE: Silueta (skyline) de n edificios mediante barrido
	// horizontal + multiset de alturas activas
	struct SkylineProblema {
		struct Evento { long long x; long long alto; bool esInicio; };
		vector<Evento> eventos;
		
		// CREATE: reserva espacio para 2*n eventos (inicio y fin de cada edificio)
		SkylineProblema(int n) {
			eventos.reserve(2 * n);
		} // O(1)
		
		// WRITE: agrega un edificio dado su rango horizontal [izq, der] y su altura
		void agregarEdificio(long long izq, long long der, long long alto) {
			eventos.push_back({izq, alto, true});
			eventos.push_back({der, alto, false});
		} // O(1) amortizado
		
		// READ: calcula los puntos clave del skyline: pares (x, altura) donde
		// la altura maxima visible cambia
		vector<pair<long long,long long>> calcularSkyline() {
			sort(eventos.begin(), eventos.end(), [](const Evento& a, const Evento& b) {
				return a.x < b.x;
			});
			
			multiset<long long> alturasActivas = {0}; // el "suelo" siempre esta activo
			vector<pair<long long,long long>> resultado;
			long long alturaPrevia = 0;
			size_t i = 0;
			while (i < eventos.size()) {
				long long xActual = eventos[i].x;
				// procesa TODOS los eventos de esta misma x como un bloque, para
				// no reportar un punto intermedio incorrecto
				while (i < eventos.size() && eventos[i].x == xActual) {
					if (eventos[i].esInicio) alturasActivas.insert(eventos[i].alto);
					else alturasActivas.erase(alturasActivas.find(eventos[i].alto));
					i++;
				}
				long long alturaActual = *alturasActivas.rbegin();
				if (alturaActual != alturaPrevia) {
					resultado.push_back({xActual, alturaActual});
					alturaPrevia = alturaActual;
				}
			}
			return resultado;
		} // O(n log n): sort O(n log n), n operaciones sobre el multiset O(log n) cada una
		
		// ------------------------------------------------------------
		// MODIFICACIONES Y COMENTARIOS CON LOS USOS
		// ------------------------------------------------------------
		// 1. Area bajo el skyline en vez de los puntos clave: basta con acumular
		//    (x_actual - x_anterior) * alturaPrevia en cada iteracion del while
		//    externo, antes de actualizar alturaPrevia.
		// 2. Muchos edificios con la MISMA altura exacta: el multiset es
		//    preferible a un arbol de segmentos aqui, porque insertar/borrar UNA
		//    ocurrencia de un valor repetido es directo con find + erase.
		// 3. Silueta con arbol de segmentos (alternativa): si el rango de x es
		//    pequeno y se requiere actualizar dinamicamente, se puede reemplazar
		//    el multiset por un arbol de segmentos tipo "maximo con propagacion
		//    perezosa" sobre las coordenadas x comprimidas.
	};
\end{lstlisting}

\textbf{Ejemplo de funcionamiento:}
\begin{lstlisting}[style=cpp]
	SkylineProblema s(2);
	s.agregarEdificio(1, 5, 3);
	s.agregarEdificio(3, 7, 5);
	
	auto resultado = s.calcularSkyline();
	// resultado = {(1,3), (3,5), (7,0)}
	// (en x=1 aparece un edificio de altura 3; en x=3 aparece uno mas alto,
	//  de altura 5, que domina la vista; en x=7 termina el mas alto y la
	//  altura visible cae a 0, ya que el primero termino en x=5 sin cambiar
	//  el maximo visible en ese punto)
\end{lstlisting}

\textbf{Nota:} el paso de agrupar eventos por la misma coordenada $x$ antes de consultar el máximo (el \texttt{while} interno del código) es crítico: procesar los eventos uno por uno sin agrupar puede reportar un punto de skyline falso cuando un edificio termina justo donde otro empieza en la misma coordenada.

\textbf{Regla práctica:} usa esta técnica cuando la pregunta sea sobre el \textbf{máximo} (o mínimo) visible en cada punto, no sobre longitud o área acumulada —para eso está \textit{Unión de Áreas de Rectángulos}. Ambas comparten la misma estructura de eventos de entrada/salida, así que si ya resolviste una, adaptar la otra es principalmente cambiar la estructura de estado (árbol de segmentos por conteo vs. multiset de alturas).
\subsection{Par de Puntos Más Cercano (Closest Pair)}

\textbf{Idea:} a diferencia de las subsecciones anteriores, aquí el barrido no procesa eventos de entrada/salida de un intervalo, sino \textbf{puntos individuales} ordenados por $x$, y el estado activo es una "franja" (\textit{strip}) de puntos candidatos a formar el par más cercano. La observación clave: si ya se conoce una distancia mínima $d$ entre los puntos procesados hasta ahora, cualquier punto nuevo solo puede mejorar esa distancia si está a menos de $d$ en $x$ del punto actual —por eso se mantiene una ventana de puntos "recientes" (los últimos con $x$ dentro de $d$), y dentro de esa ventana, gracias a un argumento geométrico de empaquetamiento, solo hace falta comparar contra un número constante de vecinos ordenados por $y$.

\textbf{¿Por qué usarlo?} Comparar todos los pares de $n$ puntos es $O(n^2)$. El enfoque de \textit{divide y vencerás} clásico también resuelve esto en $O(n \log n)$, pero requiere combinar recursivamente las mitades izquierda y derecha con una franja central —más difícil de implementar. La versión con barrido logra la misma complejidad manteniendo una única estructura ordenada por $y$ (un \texttt{multiset} o similar) que se poda dinámicamente, sin necesidad de recursión ni de fusionar subproblemas. Es la aplicación más directa de "mantener solo el estado necesario" entre las técnicas de esta sección: aquí el estado es una ventana deslizante de puntos, no un conjunto de intervalos activos.

\textbf{Casos de uso:}
\begin{itemize}
	\item Distancia mínima entre $n$ puntos en el plano.
	\item Detección de colisiones o "puntos demasiado cercanos" en simulaciones.
	\item Preprocesamiento en problemas de clustering o agrupamiento espacial.
	\item Verificar si algún par de ubicaciones (antenas, sensores) viola una distancia mínima de seguridad.
\end{itemize}

\textbf{Complejidad:} $O(n \log n)$. Ordenar los $n$ puntos por $x$ cuesta $O(n \log n)$. Cada punto se inserta y eventualmente se elimina del \texttt{multiset} activo —$O(\log n)$ cada operación—, y por cada punto se consultan solo los vecinos con $y$ dentro de rango $d$, que son $O(1)$ en promedio por el argumento de empaquetamiento (en una franja de ancho $d$ y alto $2d$ caben a lo sumo un número constante de puntos con distancia mínima $d$ entre ellos). El total es $O(n \log n)$.

\begin{lstlisting}[style=cpp]
	// CREATE/READ/WRITE: Distancia minima entre n puntos mediante barrido
	// horizontal + ventana de candidatos ordenada por y
	struct ParMasCercano {
		struct Punto { double x, y; int id; };
		vector<Punto> puntos;
		
		// CREATE: reserva espacio para n puntos
		ParMasCercano(int n) {
			puntos.reserve(n);
		} // O(1)
		
		// WRITE: agrega un punto al conjunto
		void agregarPunto(double x, double y) {
			puntos.push_back({x, y, (int)puntos.size()});
		} // O(1) amortizado
		
		// READ: calcula la distancia euclidiana al cuadrado (evita sqrt hasta el final)
		static double distCuadrado(const Punto& a, const Punto& b) {
			double dx = a.x - b.x, dy = a.y - b.y;
			return dx * dx + dy * dy;
		} // O(1)
		
		// READ: barre los puntos ordenados por x y devuelve la distancia minima
		// entre cualquier par
		double distanciaMinima() {
			sort(puntos.begin(), puntos.end(), [](const Punto& a, const Punto& b) {
				return a.x < b.x;
			});
			
			// ventana ordenada por y, contiene solo puntos con x dentro de la
			// distancia minima actual respecto al punto que se esta procesando
			set<pair<double,int>> ventana; // (y, id) para permitir duplicados de y
			double mejorCuadrado = 1e18;
			int inicio = 0; // indice del primer punto todavia dentro de la ventana
			
			for (int i = 0; i < (int)puntos.size(); i++) {
				double mejorActual = sqrt(mejorCuadrado);
				// elimina de la ventana los puntos cuya x ya quedo demasiado atras
				while (inicio < i && puntos[i].x - puntos[inicio].x > mejorActual) {
					ventana.erase({puntos[inicio].y, puntos[inicio].id});
					inicio++;
				}
				// solo hace falta comparar contra puntos con y en [y_i - mejorActual, y_i + mejorActual]
				auto desde = ventana.lower_bound({puntos[i].y - mejorActual, -1});
				auto hasta = ventana.upper_bound({puntos[i].y + mejorActual, INT_MAX});
				for (auto it = desde; it != hasta; ++it) {
					double d2 = distCuadrado(puntos[i], puntos[it->second]);
					mejorCuadrado = min(mejorCuadrado, d2);
				}
				ventana.insert({puntos[i].y, puntos[i].id});
			}
			return sqrt(mejorCuadrado);
		} // O(n log n): sort O(n log n), y por el argumento de empaquetamiento cada
		// punto compara con O(1) vecinos en promedio dentro de la ventana
		
		// ------------------------------------------------------------
		// MODIFICACIONES Y COMENTARIOS CON LOS USOS
		// ------------------------------------------------------------
		// 1. Reportar el PAR de puntos (no solo la distancia): guardar los ids
		//    junto con mejorCuadrado cada vez que se actualiza el minimo dentro
		//    del for interno.
		// 2. Coordenadas enteras: usar long long en distCuadrado para evitar
		//    overflow, y comparar con mejorCuadrado como entero (sin sqrt) hasta
		//    el resultado final, evitando errores de precision con double.
		// 3. Divide y venceras (alternativa clasica): en vez de la ventana
		//    deslizante con set, se puede resolver recursivamente dividiendo el
		//    conjunto de puntos por la mediana de x, resolviendo cada mitad, y
		//    fusionando con una franja central de ancho 2*d. Misma complejidad
		//    O(n log n), pero mas dificil de implementar correctamente; el
		//    barrido con set es preferible en un contexto de concurso.
	};
\end{lstlisting}

\textbf{Ejemplo de funcionamiento:}
\begin{lstlisting}[style=cpp]
	ParMasCercano p(4);
	p.agregarPunto(0, 0);
	p.agregarPunto(3, 4);
	p.agregarPunto(1, 1);
	p.agregarPunto(10, 10);
	
	double resultado = p.distanciaMinima();
	// resultado = sqrt(2)
	// (el par mas cercano es (0,0) y (1,1), a distancia sqrt(1^2 + 1^2) = sqrt(2);
	//  hay que correrlo para confirmar el valor decimal exacto)
\end{lstlisting}

\textbf{Nota:} la condición \texttt{puntos[i].x - puntos[inicio].x > mejorActual} para podar la ventana, y los límites \texttt{lower\_bound}/\texttt{upper\_bound} en $y$, son lo que mantiene la comparación por punto en $O(1)$ amortizado. Si se omite cualquiera de las dos podas (en $x$ o en $y$), el algoritmo sigue siendo correcto pero degenera a $O(n^2)$ en el peor caso, perdiendo la ventaja del barrido.

\textbf{Regla práctica:} usa esta técnica cuando el problema pida una distancia mínima (o máxima densidad local) entre puntos en el plano y $n$ sea grande. Frente al divide y vencerás clásico, esta versión con barrido y ventana deslizante es más simple de codificar en competencia y logra la misma complejidad $O(n \log n)$, siempre que no se olviden las dos podas de la ventana (en $x$ al insertar/eliminar, en $y$ al consultar).

\section{Primitivas Geométricas}
\subsection{Vector Geométrico: Producto Punto y Producto Cruz}

\textbf{Idea:} un vector geométrico $\vec{v} = (x, y)$ representa una dirección y magnitud en el plano, y casi toda la geometría computacional se reduce a dos operaciones entre vectores: el \textbf{producto punto} $\vec{a} \cdot \vec{b} = a_x b_x + a_y b_y$, que mide cuánto "apuntan en la misma dirección" dos vectores (relacionado con el coseno del ángulo entre ellos), y el \textbf{producto cruz} $\vec{a} \times \vec{b} = a_x b_y - a_y b_x$, que en 2D da un escalar cuyo signo indica de qué lado está un vector respecto al otro (relacionado con el seno del ángulo, y con el área con signo del paralelogramo que forman). Estas dos operaciones son los "primitivos" a partir de los cuales se construyen casi todos los demás algoritmos de esta parte: orientación, colinealidad, área de polígonos, ángulos.

\textbf{¿Por qué usarlo?} Sin estas dos operaciones no hay forma de razonar algebraicamente sobre direcciones y orientación en el plano —la alternativa sería calcular ángulos con \texttt{atan2} y compararlos, lo cual es más lento, acumula error de punto flotante, y falla en casos límite (ángulos de 0 o $2\pi$). El producto punto y el producto cruz, en cambio, se calculan con aritmética exacta si las coordenadas son enteras, y evitan funciones trigonométricas costosas en la mayoría de los casos donde solo se necesita signo o comparación relativa, no el ángulo exacto.

\textbf{Casos de uso:}
\begin{itemize}
	\item Determinar si tres puntos giran en sentido horario o antihorario (ver \textit{Orientación de Tres Puntos}).
	\item Calcular el área de un polígono (Shoelace, que es una suma de productos cruz).
	\item Proyectar un punto sobre una línea (producto punto normalizado).
	\item Determinar si un ángulo es agudo, recto u obtuso (signo del producto punto).
	\item Ordenar puntos angularmente alrededor de un pivote (signo del producto cruz).
\end{itemize}

\textbf{Complejidad:} $O(1)$ para ambas operaciones —son sumas y multiplicaciones de un número constante de componentes, sin importar el tamaño de la entrada.

\begin{lstlisting}[style=cpp]
	// CREATE/READ/WRITE: Vector geometrico 2D con producto punto y producto cruz
	struct Vector2D {
		double x, y;
		
		// CREATE: construye un vector a partir de sus componentes (0,0 por defecto)
		Vector2D(double x_ = 0, double y_ = 0) : x(x_), y(y_) {} // O(1)
		
		// READ: resta de puntos para obtener el vector que va de a hacia b
		static Vector2D desde(const Vector2D& a, const Vector2D& b) {
			return Vector2D(b.x - a.x, b.y - a.y);
		} // O(1)
		
		// READ: producto punto, mide alineacion entre this y otro vector
		double punto(const Vector2D& otro) const {
			return x * otro.x + y * otro.y;
		} // O(1)
		
		// READ: producto cruz (componente z del producto cruz 3D), mide
		// orientacion relativa y area con signo del paralelogramo
		double cruz(const Vector2D& otro) const {
			return x * otro.y - y * otro.x;
		} // O(1)
		
		// READ: magnitud (norma) del vector
		double norma() const {
			return sqrt(x * x + y * y);
		} // O(1)
		
		// ------------------------------------------------------------
		// MODIFICACIONES Y COMENTARIOS CON LOS USOS
		// ------------------------------------------------------------
		// 1. Coordenadas enteras: si el problema garantiza coordenadas enteras,
		//    usar long long en vez de double evita error de punto flotante en
		//    punto() y cruz() (ambos dan resultado exacto con enteros).
		// 2. Signo del producto cruz: cruz(otro) > 0 significa que "otro" esta
		//    en sentido antihorario respecto a this; < 0 significa horario;
		//    == 0 significa que son colineales (paralelos). Ver Orientacion de
		//    Tres Puntos para la aplicacion directa de esto.
		// 3. Angulo entre vectores: se puede obtener con
		//    atan2(this->cruz(otro), this->punto(otro)), que da el angulo con
		//    signo en (-pi, pi] sin ambiguedad de cuadrante, a diferencia de
		//    usar solo acos(punto / (norma*norma)).
	};
\end{lstlisting}

\textbf{Ejemplo de funcionamiento:}
\begin{lstlisting}[style=cpp]
	Vector2D a(2, 0);
	Vector2D b(0, 3);
	
	double p = a.punto(b);
	// p = 0
	// (son perpendiculares, el producto punto de vectores perpendiculares es 0)
	
	double c = a.cruz(b);
	// c = 6
	// (positivo: b esta en sentido antihorario respecto a a;
	//  6 es tambien el area del paralelogramo que forman)
\end{lstlisting}

\textbf{Regla práctica:} usa el producto punto cuando la pregunta sea sobre \textbf{alineación o proyección} (¿qué tan paralelos son dos vectores?, ¿es un ángulo agudo u obtuso?); usa el producto cruz cuando la pregunta sea sobre \textbf{orientación o área con signo} (¿a qué lado está un punto?, ¿en qué sentido giran tres puntos?). Casi todo lo que sigue en esta parte de la guía se construye componiendo estas dos operaciones.

\subsection{Orientación de Tres Puntos (Cross Product / CCW Test)}

\textbf{Idea:} dado un punto $O$ (origen/pivote) y dos puntos $A$ y $B$, el signo del producto cruz $\vec{OA} \times \vec{OB}$ —visto en \textit{Vector Geométrico}— determina si, parado en $O$ mirando hacia $A$, el punto $B$ queda a la izquierda (positivo, sentido antihorario/CCW), a la derecha (negativo, sentido horario/CW), o exactamente sobre la línea que pasa por $O$ y $A$ (cero, colineales). Esta es la prueba de orientación más usada en geometría computacional porque reduce una pregunta geométrica cualitativa ("¿de qué lado está?") a una operación algebraica exacta de $O(1)$.

\textbf{¿Por qué usarlo?} Es la primitiva que sustituye a comparar ángulos con \texttt{atan2} o a calcular pendientes con división —ambas formas más lentas, propensas a error de punto flotante, y con casos especiales molestos (línea vertical, división por cero). El producto cruz da el mismo resultado cualitativo con una sola multiplicación y resta, y si las coordenadas son enteras, el resultado es exacto sin ningún error de precisión. Es la base directa de Shoelace, Convex Hull y Point-in-Polygon, que se ven más adelante en esta parte.

\textbf{Casos de uso:}
\begin{itemize}
	\item Determinar si un giro en una secuencia de puntos es a la izquierda o a la derecha (clave en construcción de envolvente convexa).
	\item Verificar si tres puntos son colineales.
	\item Determinar de qué lado de una línea está un punto (para point-in-polygon, clipping, etc.).
	\item Ordenar puntos angularmente alrededor de un pivote.
	\item Calcular el área con signo de un triángulo (mitad del valor absoluto del producto cruz).
\end{itemize}

\textbf{Complejidad:} $O(1)$ —es una evaluación directa del producto cruz entre dos vectores derivados de los tres puntos, sin bucles ni recursión.

\begin{lstlisting}[style=cpp]
	// CREATE/READ: Prueba de orientacion de tres puntos mediante producto cruz
	struct OrientacionPuntos {
		struct Punto { double x, y; };
		
		// READ: producto cruz de (b - o) x (c - o); signo indica orientacion
		// de o, b, c: positivo = CCW (izquierda), negativo = CW (derecha), 0 = colineales
		static double cruz(const Punto& o, const Punto& b, const Punto& c) {
			return (b.x - o.x) * (c.y - o.y) - (b.y - o.y) * (c.x - o.x);
		} // O(1)
		
		// READ: clasifica la orientacion como +1 (CCW), -1 (CW), o 0 (colineal)
		// usando un epsilon para tolerar error de punto flotante
		static int orientacion(const Punto& o, const Punto& b, const Punto& c) {
			double val = cruz(o, b, c);
			const double EPS = 1e-9;
			if (val > EPS) return 1;   // CCW: c esta a la izquierda de o->b
			if (val < -EPS) return -1; // CW: c esta a la derecha de o->b
			return 0;                  // colineales
		} // O(1)
		
		// ------------------------------------------------------------
		// MODIFICACIONES Y COMENTARIOS CON LOS USOS
		// ------------------------------------------------------------
		// 1. Coordenadas enteras: usar long long y comparar cruz() directamente
		//    contra 0 (sin EPS), ya que con enteros el resultado es exacto y no
		//    hay error de punto flotante que tolerar.
		// 2. Colinealidad estricta con verificacion de "entre": orientacion == 0
		//    solo dice que estan sobre la misma recta, no que c este ENTRE o y b.
		//    Para eso se necesita ademas revisar que las coordenadas de c esten
		//    dentro del rango [o, b] en x e y (ver Tres Puntos Colineales).
		// 3. Convencion de signo: esta implementacion usa positivo = CCW, que es
		//    la convencion estandar en la mayoria de la literatura de CP (y la
		//    que se usa consistentemente en el resto de esta guia, incluyendo
		//    Shoelace y Envolvente Convexa). Si se invierte el orden de los
		//    parametros (o, c, b en vez de o, b, c), el signo tambien se invierte.
	};
\end{lstlisting}

\textbf{Ejemplo de funcionamiento:}
\begin{lstlisting}[style=cpp]
	OrientacionPuntos::Punto o{0, 0}, b{4, 0}, c{2, 2};
	
	int resultado = OrientacionPuntos::orientacion(o, b, c);
	// resultado = 1 (CCW)
	// (parado en o mirando hacia b sobre el eje x, el punto c=(2,2) queda
	//  arriba, es decir a la izquierda, sentido antihorario)
\end{lstlisting}

\textbf{Nota:} la convención positivo = CCW fijada aquí se usa de forma consistente en el resto de la guía —Shoelace, Envolvente Convexa (Jarvis y Monotone Chain), y Punto Dentro de un Polígono Convexo. Si en algún momento se necesita la convención opuesta, basta con invertir el signo del resultado o el orden de los argumentos, pero hay que hacerlo de forma consistente en todos los algoritmos que dependan de esta prueba.

\textbf{Regla práctica:} usa esta prueba como bloque de construcción base cada vez que necesites saber de qué lado de una línea (definida por dos puntos) está un tercer punto, sin importar el ángulo exacto —eso es más rápido y más robusto numéricamente que calcular pendientes o ángulos con funciones trigonométricas.
\subsection{Distancia Euclidiana}

\textbf{Idea:} la distancia euclidiana entre dos puntos $A = (x_1, y_1)$ y $B = (x_2, y_2)$ se deriva directamente del teorema de Pitágoras: el segmento $AB$ es la hipotenusa de un triángulo rectángulo cuyos catetos son las diferencias de coordenadas $\Delta x = x_2 - x_1$ y $\Delta y = y_2 - y_1$, así que $d(A,B) = \sqrt{\Delta x^2 + \Delta y^2}$. Es la noción de distancia "en línea recta" y la base de la mayoría de las métricas geométricas de esta guía, aunque conceptualmente equivale a la norma del vector $\vec{AB}$ (ver \textit{Vector Geométrico}).

\textbf{¿Por qué usarlo?} Es la métrica por defecto en geometría computacional cuando el problema no especifica otra cosa (a diferencia de la distancia Manhattan, que suma $|\Delta x| + |\Delta y|$ y aparece en problemas de grilla). El costo real está en la raíz cuadrada: cuando solo se necesita \textbf{comparar} distancias (cuál es menor, cuál es mayor), conviene trabajar con la distancia al cuadrado y evitar \texttt{sqrt()} —así se gana precisión con enteros y se ahorra una operación costosa, como ya se hizo en \textit{Par de Puntos Más Cercano}.

\textbf{Casos de uso:}
\begin{itemize}
	\item Encontrar el par de puntos más cercano o más lejano.
	\item Verificar si un punto está dentro de un círculo de radio dado.
	\item Calcular el perímetro de un polígono (suma de distancias entre vértices consecutivos).
	\item Comparar cuál de dos puntos está más cerca de un tercero, sin necesidad del valor exacto.
\end{itemize}

\textbf{Complejidad:} $O(1)$ —dos restas, dos multiplicaciones, una suma y, opcionalmente, una raíz cuadrada.

\begin{lstlisting}[style=cpp]
	// CREATE/READ: Distancia euclidiana entre dos puntos, con y sin raiz cuadrada
	struct DistanciaEuclidiana {
		struct Punto { double x, y; };
		
		// READ: distancia al cuadrado, evita sqrt cuando solo se necesita comparar
		static double distCuadrado(const Punto& a, const Punto& b) {
			double dx = b.x - a.x;
			double dy = b.y - a.y;
			return dx * dx + dy * dy;
		} // O(1)
		
		// READ: distancia euclidiana real (con raiz cuadrada)
		static double distancia(const Punto& a, const Punto& b) {
			return sqrt(distCuadrado(a, b));
		} // O(1)
		
		// ------------------------------------------------------------
		// MODIFICACIONES Y COMENTARIOS CON LOS USOS
		// ------------------------------------------------------------
		// 1. Coordenadas enteras: usar long long en distCuadrado evita perder
		//    precision; ademas, comparar distCuadrado(a,b) < distCuadrado(c,d)
		//    da el mismo resultado que comparar distancia(a,b) < distancia(c,d)
		//    sin necesidad de sqrt, porque la raiz cuadrada es monotona creciente.
		// 2. Distancia Manhattan (metrica alternativa): |b.x-a.x| + |b.y-a.y|,
		//    usada en problemas de grilla donde el movimiento es solo horizontal
		//    o vertical, no diagonal.
		// 3. Distancia Chebyshev (otra metrica alternativa): max(|b.x-a.x|,
		//    |b.y-a.y|), usada cuando el movimiento diagonal cuesta igual que el
		//    movimiento recto (como el rey en ajedrez).
	};
\end{lstlisting}

\textbf{Ejemplo de funcionamiento:}
\begin{lstlisting}[style=cpp]
	DistanciaEuclidiana::Punto a{0, 0}, b{3, 4};
	
	double resultado = DistanciaEuclidiana::distancia(a, b);
	// resultado = 5
	// (triangulo 3-4-5, caso clasico de pitagoras exacto)
\end{lstlisting}

\textbf{Regla práctica:} usa \texttt{distCuadrado} en vez de \texttt{distancia} cada vez que el problema solo requiera \textbf{comparar} distancias entre sí (mínimo, máximo, ¿es menor que $r$?) —evita la raíz cuadrada y, con coordenadas enteras, evita también el error de punto flotante por completo.

\subsection{Punto Medio de un Segmento}

\textbf{Idea:} el punto medio de un segmento entre $A = (x_1, y_1)$ y $B = (x_2, y_2)$ es el promedio simple de sus coordenadas: $M = \left(\dfrac{x_1+x_2}{2}, \dfrac{y_1+y_2}{2}\right)$. Es un caso particular de interpolación lineal con $t = 0.5$ —el punto que está exactamente a mitad de camino en ambas coordenadas por separado, ya que la geometría del plano cartesiano permite tratar $x$ e $y$ de forma independiente para este cálculo.

\textbf{¿Por qué usarlo?} Es la operación más simple de esta parte, pero aparece constantemente como paso intermedio de algoritmos más complejos (bisección de segmentos, cálculo de centroides, mediatriz para circuncentro). A diferencia de la \textit{División de un Segmento en una Razón Dada} (más general, con $t$ arbitrario), el punto medio es el caso fijo $t=0.5$ y no requiere parametrizar la razón, por lo que conviene tenerlo aparte como atajo directo cuando la razón es exactamente $1:1$.

\textbf{Casos de uso:}
\begin{itemize}
	\item Encontrar el centro de un segmento para construir la mediatriz (útil en circuncentro).
	\item Bisección recursiva de un segmento o polígono.
	\item Calcular el centroide de un conjunto de puntos (promedio de todos los puntos medios, o directamente el promedio de coordenadas).
	\item Punto de referencia para pruebas de simetría o reflexión.
\end{itemize}

\textbf{Complejidad:} $O(1)$ —dos sumas y dos divisiones por 2.

\begin{lstlisting}[style=cpp]
	// CREATE/READ: Punto medio de un segmento definido por dos puntos
	struct PuntoMedio {
		struct Punto { double x, y; };
		
		// READ: calcula el punto medio entre a y b
		static Punto medio(const Punto& a, const Punto& b) {
			return Punto{(a.x + b.x) / 2.0, (a.y + b.y) / 2.0};
		} // O(1)
		
		// ------------------------------------------------------------
		// MODIFICACIONES Y COMENTARIOS CON LOS USOS
		// ------------------------------------------------------------
		// 1. Coordenadas enteras con resultado exacto: si a.x+b.x y a.y+b.y son
		//    siempre pares, se puede mantener el punto medio en enteros dividiendo
		//    entre 2 directamente, evitando pasar por double.
		// 2. Centroide de n puntos (generalizacion): promedio de todas las
		//    coordenadas x y de todas las coordenadas y por separado, no es lo
		//    mismo que promediar puntos medios sucesivos si los puntos no forman
		//    un segmento simple.
		// 3. Punto medio ponderado (interpolacion con t != 0.5): ver Division de
		//    un Segmento en una Razon Dada para el caso general.
	};
\end{lstlisting}

\textbf{Ejemplo de funcionamiento:}
\begin{lstlisting}[style=cpp]
	PuntoMedio::Punto a{0, 0}, b{6, 10};
	
	auto resultado = PuntoMedio::medio(a, b);
	// resultado = (3, 5)
\end{lstlisting}

\textbf{Regla práctica:} usa el punto medio directamente cuando la razón sea exactamente $1:1$; si el problema pide un punto a una fracción arbitraria del segmento (por ejemplo, a un tercio de camino), no fuerces el punto medio dos veces —usa la fórmula general de \textit{División de un Segmento en una Razón Dada}, que es igual de simple y evita errores de composición.
\subsection{Ecuación de una Línea Dados Dos Puntos}

\textbf{Idea:} dados dos puntos $A = (x_1, y_1)$ y $B = (x_2, y_2)$, la recta que pasa por ambos se puede escribir en forma general $ax + by + c = 0$ tomando $a = y_2 - y_1$, $b = x_1 - x_2$, $c = -(a x_1 + b y_1)$. Esta forma se deriva de que el vector normal a la dirección $\vec{AB} = (x_2-x_1,\, y_2-y_1)$ es $(a, b) = (y_2-y_1,\, -(x_2-x_1))$ —rotar el vector director 90 grados— y cualquier punto $(x,y)$ sobre la recta satisface que el vector desde $A$ hasta $(x,y)$ es perpendicular a ese normal, lo cual da la ecuación. A diferencia de la forma pendiente-intersección $y = mx + k$, esta forma general no tiene problema con rectas verticales (donde la pendiente no está definida).

\textbf{¿Por qué usarlo?} La forma $y = mx + k$ falla cuando $x_1 = x_2$ (recta vertical, pendiente infinita), obligando a manejar ese caso aparte con un \texttt{if}. La forma general $ax+by+c=0$ maneja rectas verticales y horizontales de manera uniforme sin casos especiales, y además permite evaluar directamente de qué lado de la recta está un punto con el signo de $ax+by+c$ —relacionado con el producto cruz de \textit{Orientación de Tres Puntos}, pero expresado como coeficientes reutilizables en vez de recalcular el producto cruz cada vez.

\textbf{Casos de uso:}
\begin{itemize}
	\item Determinar si un punto está sobre, a la izquierda o a la derecha de una recta fija (evaluada muchas veces).
	\item Calcular la distancia de un punto a una recta (usando los coeficientes $a,b,c$ directamente).
	\item Encontrar la intersección entre dos rectas resolviendo el sistema de dos ecuaciones.
	\item Representar restricciones lineales en problemas de programación lineal geométrica simple.
\end{itemize}

\textbf{Complejidad:} $O(1)$ para construir la ecuación; $O(1)$ para evaluarla en un punto dado.

\begin{lstlisting}[style=cpp]
	// CREATE/READ: Linea en forma general ax + by + c = 0, dados dos puntos
	struct LineaGeneral {
		double a, b, c;
		
		// CREATE: construye la ecuacion general de la recta que pasa por p y q
		LineaGeneral(double x1, double y1, double x2, double y2) {
			a = y2 - y1;
			b = x1 - x2;
			c = -(a * x1 + b * y1);
		} // O(1)
		
		// READ: evalua ax + by + c en un punto; el signo indica el lado de la
		// recta (positivo o negativo segun la orientacion de a,b), 0 = sobre la recta
		double evaluar(double x, double y) const {
			return a * x + b * y + c;
		} // O(1)
		
		// READ: verifica si un punto esta sobre la recta (con tolerancia)
		bool contiene(double x, double y) const {
			const double EPS = 1e-9;
			return fabs(evaluar(x, y)) < EPS;
		} // O(1)
		
		// ------------------------------------------------------------
		// MODIFICACIONES Y COMENTARIOS CON LOS USOS
		// ------------------------------------------------------------
		// 1. Interseccion de dos rectas L1 (a1,b1,c1) y L2 (a2,b2,c2): resolver
		//    el sistema 2x2 con la regla de Cramer:
		//    det = a1*b2 - a2*b1;
		//    if (fabs(det) < EPS) { /* paralelas o coincidentes, sin solucion unica */ }
		//    x = (b1*c2 - b2*c1) / det; y = (a2*c1 - a1*c2) / det;
		// 2. Coordenadas enteras: a, b, c pueden mantenerse como long long para
		//    evaluar() y contiene() de forma exacta, comparando contiene()
		//    contra 0 sin EPS.
		// 3. Normalizacion opcional: dividir a,b,c entre sqrt(a*a+b*b) convierte
		//    evaluar(x,y) directamente en la distancia con signo del punto a la
		//    recta (ver Distancia de un Punto a una Linea/Segmento).
	};
\end{lstlisting}

\textbf{Ejemplo de funcionamiento:}
\begin{lstlisting}[style=cpp]
	LineaGeneral l(0, 0, 4, 4); // recta y = x
	
	double val = l.evaluar(2, 5);
	// val = -12
	// (a=4, b=-4, c=0; evaluar(2,5) = 4*2 + (-4)*5 + 0 = 8 - 20 = -12;
	//  el signo negativo indica que (2,5) esta de un lado especifico de y=x)
	
	bool sobre = l.contiene(3, 3);
	// sobre = true (el punto (3,3) esta sobre la recta y=x)
\end{lstlisting}

\textbf{Regla práctica:} usa esta forma general en vez de $y=mx+k$ cada vez que la recta pueda ser vertical, o cuando necesites evaluar el mismo lado de la recta para muchos puntos distintos —los coeficientes $a,b,c$ se calculan una sola vez y se reutilizan, a diferencia de recalcular un producto cruz completo contra dos puntos fijos en cada consulta.

\subsection{División de un Segmento en una Razón Dada}

\textbf{Idea:} generaliza el \textit{Punto Medio de un Segmento} (que es el caso fijo $t=0.5$) a cualquier fracción $t \in [0,1]$ del camino entre $A$ y $B$: el punto que divide el segmento en esa proporción es la interpolación lineal $P = A + t \cdot (B - A) = ((1-t)x_1 + t x_2,\; (1-t)y_1 + t y_2)$. Cuando $t=0$ se obtiene $A$, cuando $t=1$ se obtiene $B$, y valores intermedios dan puntos proporcionalmente distribuidos sobre el segmento. Si la razón se da como "$m:n$" (divide al segmento en partes de longitud proporcional $m$ y $n$), se convierte a $t = \frac{m}{m+n}$.

\textbf{¿Por qué usarlo?} Es la forma general de la que el punto medio es un caso particular; evaluarlo cada vez con la fórmula de interpolación en vez de "forzar" un punto medio con pasos intermedios evita errores de composición y funciona igual de bien para $t$ fuera de $[0,1]$ (extrapolación, para encontrar un punto más allá de $B$ en la misma dirección que $A \to B$). Es la base de la interpolación paramétrica de segmentos, usada también para representar cualquier punto sobre una recta en función de un solo parámetro.

\textbf{Casos de uso:}
\begin{itemize}
	\item Dividir un segmento en $k$ partes iguales (aplicando $t = i/k$ para $i=1,\dots,k-1$).
	\item Encontrar un punto a una razón $m:n$ específica sobre un segmento (por ejemplo, el punto que divide una mediana en razón $2:1$ desde el vértice, propiedad del centroide de un triángulo).
	\item Extrapolar un punto más allá de un extremo del segmento, usando $t < 0$ o $t > 1$.
	\item Parametrizar un rayo o recta para intersección con otras figuras.
\end{itemize}

\textbf{Complejidad:} $O(1)$ —dos interpolaciones lineales independientes en $x$ e $y$.

\begin{lstlisting}[style=cpp]
	// CREATE/READ: Division de un segmento AB en una razon dada mediante
	// interpolacion lineal parametrica
	struct DivisionSegmento {
		struct Punto { double x, y; };
		
		// READ: interpola el punto a fraccion t del segmento de a hacia b
		// t=0 da a, t=1 da b, valores fuera de [0,1] extrapolan
		static Punto interpolar(const Punto& a, const Punto& b, double t) {
			return Punto{
				a.x + t * (b.x - a.x),
				a.y + t * (b.y - a.y)
			};
		} // O(1)
		
		// READ: divide el segmento en razon m:n (desde a hacia b), convirtiendo
		// la razon a fraccion t = m/(m+n) antes de interpolar
		static Punto dividirEnRazon(const Punto& a, const Punto& b, double m, double n) {
			double t = m / (m + n);
			return interpolar(a, b, t);
		} // O(1)
		
		// ------------------------------------------------------------
		// MODIFICACIONES Y COMENTARIOS CON LOS USOS
		// ------------------------------------------------------------
		// 1. Punto medio como caso particular: interpolar(a, b, 0.5) es
		//    exactamente Punto Medio de un Segmento; no hace falta un struct
		//    separado si ya se tiene este.
		// 2. Division en k partes iguales: generar los puntos
		//    interpolar(a, b, (double)i / k) para i = 1 hasta k-1.
		// 3. Coordenadas enteras con razon m:n exacta (sin pasar por double):
		//    px = a.x + (b.x - a.x) * m / (m + n); usando division entera solo
		//    si se garantiza divisibilidad exacta; de lo contrario, mantener
		//    el resultado como fraccion (numerador, denominador) para exactitud.
	};
\end{lstlisting}

\textbf{Ejemplo de funcionamiento:}
\begin{lstlisting}[style=cpp]
	DivisionSegmento::Punto a{0, 0}, b{10, 20};
	
	auto resultado = DivisionSegmento::dividirEnRazon(a, b, 1, 4);
	// resultado = (2, 4)
	// (t = 1/(1+4) = 0.2; el punto esta a un quinto del camino de a hacia b,
	//  dividiendo el segmento en partes de razon 1:4)
\end{lstlisting}

\textbf{Regla práctica:} usa esta fórmula general cada vez que necesites un punto sobre un segmento a una fracción o razón distinta de $1:1$ —para el caso exacto de la mitad, es más directo usar \textit{Punto Medio de un Segmento} sin pasar por la conversión de razón a $t$.
\subsection{Ángulo entre Dos Vectores}

\textbf{Idea:} el ángulo entre dos vectores $\vec{a}$ y $\vec{b}$ se puede obtener de dos formas relacionadas: con el producto punto, $\cos\theta = \dfrac{\vec{a}\cdot\vec{b}}{|\vec{a}||\vec{b}|}$, o de forma más robusta con \texttt{atan2}, combinando producto punto y producto cruz —vistos ambos en \textit{Vector Geométrico}— como $\theta = \texttt{atan2}(\vec{a}\times\vec{b},\; \vec{a}\cdot\vec{b})$. La segunda forma da directamente el ángulo \textbf{con signo} en el rango $(-\pi, \pi]$, indicando no solo la magnitud del ángulo sino también si $\vec{b}$ está en sentido antihorario (positivo) u horario (negativo) respecto a $\vec{a}$, sin la ambigüedad de cuadrante que tiene \texttt{acos}.

\textbf{¿Por qué usarlo?} Usar solo \texttt{acos(punto / (norma\_a * norma\_b))} tiene dos problemas: pierde el signo (no distingue si $\vec{b}$ gira a la izquierda o a la derecha de $\vec{a}$), y es numéricamente inestable cerca de $\theta = 0$ o $\theta = \pi$, donde pequeños errores de punto flotante en el argumento de \texttt{acos} pueden caer fuera de $[-1,1]$ y producir \texttt{NaN}. La versión con \texttt{atan2(cruz, punto)} evita ambos problemas: es estable en todo el rango y conserva el signo, que es justamente lo que se necesita para ordenar puntos angularmente (ver \textit{Barrido Angular} más adelante).

\textbf{Casos de uso:}
\begin{itemize}
	\item Ordenar un conjunto de puntos angularmente alrededor de un pivote.
	\item Calcular el ángulo interno de un polígono en un vértice dado.
	\item Determinar cuánto hay que rotar un vector para alinearlo con otro.
	\item Verificar si un ángulo es agudo, recto u obtuso sin calcular el valor exacto (signo del producto punto, ver \textit{Vector Geométrico}).
\end{itemize}

\textbf{Complejidad:} $O(1)$ —un producto punto, un producto cruz, y una llamada a \texttt{atan2}.

\begin{lstlisting}[style=cpp]
	// CREATE/READ: Angulo con signo entre dos vectores, usando atan2 para
	// evitar la inestabilidad numerica de acos
	struct AnguloVectores {
		struct Vector { double x, y; };
		
		// READ: producto punto de dos vectores
		static double punto(const Vector& a, const Vector& b) {
			return a.x * b.x + a.y * b.y;
		} // O(1)
		
		// READ: producto cruz de dos vectores
		static double cruz(const Vector& a, const Vector& b) {
			return a.x * b.y - a.y * b.x;
		} // O(1)
		
		// READ: angulo con signo (en radianes) que hay que rotar a para
		// alinearlo con b; positivo = antihorario, negativo = horario
		static double angulo(const Vector& a, const Vector& b) {
			return atan2(cruz(a, b), punto(a, b));
		} // O(1)
		
		// ------------------------------------------------------------
		// MODIFICACIONES Y COMENTARIOS CON LOS USOS
		// ------------------------------------------------------------
		// 1. Angulo NO firmado (solo magnitud, en [0, pi]): fabs(angulo(a, b)).
		//    Util cuando no importa el sentido de giro, solo la abertura.
		// 2. Angulo interno de un poligono en el vertice v, entre los vertices
		//    vecinos prev y next: angulo(Vector::desde(v, prev), Vector::desde(v, next)),
		//    tomando el valor absoluto si el poligono no esta orientado de forma
		//    consistente.
		// 3. Comparar SOLO el orden angular de puntos alrededor de un pivote (sin
		//    necesitar el angulo exacto): usar directamente el signo de cruz(a, b)
		//    como comparador de sort es mas rapido y mas exacto que llamar atan2
		//    para cada comparacion (ver Barrido Angular).
	};
\end{lstlisting}

\textbf{Ejemplo de funcionamiento:}
\begin{lstlisting}[style=cpp]
	AnguloVectores::Vector a{1, 0}, b{0, 1};
	
	double resultado = AnguloVectores::angulo(a, b);
	// resultado = pi/2 (aprox 1.5708)
	// (b esta 90 grados en sentido antihorario respecto a a; hay que correrlo
	//  para confirmar el valor decimal exacto de pi/2 en el sistema usado)
\end{lstlisting}

\textbf{Regla práctica:} usa \texttt{atan2(cruz, punto)} en vez de \texttt{acos(punto/(norma\_a*norma\_b))} siempre que necesites el signo del ángulo o estabilidad numérica cerca de 0 o $\pi$. Si solo necesitas \textbf{comparar} el orden angular de varios vectores respecto a un pivote (no el ángulo exacto), ni siquiera hace falta \texttt{atan2} —el signo del producto cruz solo, usado como comparador, ya es suficiente y más rápido.

\subsection{Tres Puntos Colineales}

\textbf{Idea:} tres puntos $A$, $B$, $C$ son colineales si están sobre la misma recta, lo cual ocurre exactamente cuando el producto cruz $\vec{AB} \times \vec{AC} = 0$ —el caso frontera de \textit{Orientación de Tres Puntos}, donde no hay giro ni a la izquierda ni a la derecha. Sin embargo, colinealidad (estar sobre la misma recta infinita) es distinto de que $C$ esté \textbf{entre} $A$ y $B$ sobre el segmento: para eso, además del producto cruz igual a cero, hace falta verificar que las coordenadas de $C$ estén dentro del rango de caja delimitadora (\textit{bounding box}) de $A$ y $B$.

\textbf{¿Por qué usarlo?} Es una aplicación directa y frecuente de la prueba de orientación, pero se trata aparte porque el caso "colineal" suele ser el más propenso a errores en problemas de geometría: mezclar "están en la misma recta" con "uno está entre los otros dos" es una confusión común. Separar explícitamente ambas preguntas —colinealidad (producto cruz) y pertenencia al segmento (bounding box)— evita bugs sutiles, especialmente en algoritmos de intersección de segmentos donde los casos colineales se manejan aparte (como ya se vio en \textit{Intersección de Segmentos}).

\textbf{Casos de uso:}
\begin{itemize}
	\item Verificar si tres puntos dados están alineados (como paso de validación de entrada).
	\item Detectar si un punto cae exactamente sobre un segmento (no solo sobre su recta extendida).
	\item Simplificar un polígono eliminando vértices redundantes que son colineales con sus vecinos.
	\item Caso especial en algoritmos de intersección de segmentos, cuando el producto cruz da exactamente cero.
\end{itemize}

\textbf{Complejidad:} $O(1)$ —un producto cruz más, en el peor caso, cuatro comparaciones para el chequeo de bounding box.

\begin{lstlisting}[style=cpp]
	// CREATE/READ: Verificacion de colinealidad de tres puntos, y de pertenencia
	// de un punto colineal al segmento (no solo a la recta extendida)
	struct TresPuntosColineales {
		struct Punto { double x, y; };
		
		// READ: producto cruz de (b-a) x (c-a)
		static double cruz(const Punto& a, const Punto& b, const Punto& c) {
			return (b.x - a.x) * (c.y - a.y) - (b.y - a.y) * (c.x - a.x);
		} // O(1)
		
		// READ: true si a, b, c son colineales (estan sobre la misma recta)
		static bool sonColineales(const Punto& a, const Punto& b, const Punto& c) {
			const double EPS = 1e-9;
			return fabs(cruz(a, b, c)) < EPS;
		} // O(1)
		
		// READ: dado que a, b, c ya son colineales, verifica si c esta DENTRO
		// del segmento [a, b] (no solo sobre la recta que lo contiene)
		static bool entreSegmento(const Punto& a, const Punto& b, const Punto& c) {
			const double EPS = 1e-9;
			bool dentroX = min(a.x, b.x) - EPS <= c.x && c.x <= max(a.x, b.x) + EPS;
			bool dentroY = min(a.y, b.y) - EPS <= c.y && c.y <= max(a.y, b.y) + EPS;
			return dentroX && dentroY;
		} // O(1)
		
		// READ: combina ambas verificaciones: c esta colineal Y dentro del
		// segmento [a, b]
		static bool puntoSobreSegmento(const Punto& a, const Punto& b, const Punto& c) {
			return sonColineales(a, b, c) && entreSegmento(a, b, c);
		} // O(1)
		
		// ------------------------------------------------------------
		// MODIFICACIONES Y COMENTARIOS CON LOS USOS
		// ------------------------------------------------------------
		// 1. Coordenadas enteras: usar long long en cruz() y comparar contra 0
		//    exacto, sin EPS; entreSegmento tambien se vuelve comparacion exacta.
		// 2. Simplificacion de poligono (eliminar vertices redundantes): recorrer
		//    los vertices consecutivos en tripletas (prev, actual, siguiente) y
		//    eliminar actual si sonColineales(prev, actual, siguiente) es true.
		// 3. Ordenar puntos colineales sobre una misma recta: una vez confirmado
		//    que todos son colineales, se pueden ordenar por su proyeccion escalar
		//    sobre el vector director de la recta (producto punto contra un vector
		//    unitario en esa direccion), en vez de ordenar por x o y por separado
		//    (falla si la recta es vertical u horizontal exacta).
	};
\end{lstlisting}

\textbf{Ejemplo de funcionamiento:}
\begin{lstlisting}[style=cpp]
	TresPuntosColineales::Punto a{0, 0}, b{10, 0}, c{5, 0}, d{15, 0};
	
	bool colineales = TresPuntosColineales::sonColineales(a, b, c);
	// colineales = true (los tres estan sobre el eje x)
	
	bool dentro = TresPuntosColineales::entreSegmento(a, b, c);
	// dentro = true (c=(5,0) esta entre a=(0,0) y b=(10,0))
	
	bool fuera = TresPuntosColineales::puntoSobreSegmento(a, b, d);
	// fuera = false (d=(15,0) es colineal con a y b, pero esta MAS ALLA de b,
	//  fuera del segmento [a,b])
\end{lstlisting}

\textbf{Regla práctica:} nunca uses solo el producto cruz igual a cero para concluir que un punto "está sobre el segmento" —eso solo confirma que está sobre la \textbf{recta infinita}. Siempre combina la prueba de colinealidad con la verificación de bounding box (\texttt{entreSegmento}) cuando la pregunta sea específicamente sobre pertenencia al segmento, no a la recta extendida.
\subsection{Distancia de un Punto a una Línea/Segmento}

\textbf{Idea:} la distancia de un punto $P$ a una \textbf{recta infinita} definida por $A$ y $B$ se obtiene proyectando $\vec{AP}$ sobre la dirección perpendicular a $\vec{AB}$: usando el producto cruz de \textit{Vector Geométrico}, esa distancia es $\dfrac{|\vec{AB} \times \vec{AP}|}{|\vec{AB}|}$ —el área del paralelogramo que forman los dos vectores, dividida entre la base. Para la distancia a un \textbf{segmento} (no a la recta infinita), hay que distinguir tres casos según dónde cae la proyección de $P$ sobre la línea: si cae dentro de $[A,B]$, la distancia es la misma fórmula perpendicular; si cae antes de $A$ o después de $B$, la distancia real es simplemente la distancia euclidiana al extremo más cercano.

\textbf{¿Por qué usarlo?} Confundir distancia a la recta con distancia al segmento es un error común: la fórmula de la recta infinita subestima la distancia real cuando la proyección del punto cae fuera del segmento. La forma correcta usa el producto punto (visto en \textit{Vector Geométrico}) para determinar en qué "zona" cae la proyección —antes de $A$, dentro de $[A,B]$, o después de $B$— y solo aplica la fórmula perpendicular en el caso intermedio, cayendo de vuelta a \textit{Distancia Euclidiana} directa contra el extremo en los otros dos casos.

\textbf{Casos de uso:}
\begin{itemize}
	\item Distancia mínima de un punto a un obstáculo representado como segmento o polilínea.
	\item Verificar si un punto está a menos de $r$ unidades de una pared o borde (colisión punto-segmento).
	\item Cálculo auxiliar en \textit{Closest Pair} cuando se generaliza a distancia punto-segmento en vez de punto-punto.
	\item Simplificación de polilíneas (algoritmo de Douglas-Peucker usa esta distancia para decidir qué puntos eliminar).
\end{itemize}

\textbf{Complejidad:} $O(1)$ —un producto punto, un producto cruz, y como mucho una llamada a \texttt{sqrt}.

\begin{lstlisting}[style=cpp]
	// CREATE/READ: Distancia de un punto a una recta infinita y a un segmento
	struct DistanciaPuntoLinea {
		struct Punto { double x, y; };
		
		// READ: producto punto de (b-a) con (p-a)
		static double punto(const Punto& a, const Punto& b, const Punto& p) {
			return (b.x - a.x) * (p.x - a.x) + (b.y - a.y) * (p.y - a.y);
		} // O(1)
		
		// READ: producto cruz de (b-a) con (p-a)
		static double cruz(const Punto& a, const Punto& b, const Punto& p) {
			return (b.x - a.x) * (p.y - a.y) - (b.y - a.y) * (p.x - a.x);
		} // O(1)
		
		// READ: distancia euclidiana entre dos puntos
		static double distEuclidiana(const Punto& p, const Punto& q) {
			double dx = q.x - p.x, dy = q.y - p.y;
			return sqrt(dx * dx + dy * dy);
		} // O(1)
		
		// READ: distancia de p a la RECTA INFINITA que pasa por a y b
		static double distARecta(const Punto& a, const Punto& b, const Punto& p) {
			double baseLen = distEuclidiana(a, b);
			if (baseLen < 1e-9) return distEuclidiana(a, p); // a y b coinciden
			return fabs(cruz(a, b, p)) / baseLen;
		} // O(1)
		
		// READ: distancia de p al SEGMENTO [a, b] (no a la recta extendida)
		static double distASegmento(const Punto& a, const Punto& b, const Punto& p) {
			double t = punto(a, b, p);
			if (t <= 0) return distEuclidiana(a, p); // proyeccion cae antes de a
			double baseLen2 = punto(a, b, b);        // |ab|^2
			if (t >= baseLen2) return distEuclidiana(b, p); // proyeccion cae despues de b
			return distARecta(a, b, p); // proyeccion cae dentro de [a,b]
		} // O(1)
		
		// ------------------------------------------------------------
		// MODIFICACIONES Y COMENTARIOS CON LOS USOS
		// ------------------------------------------------------------
		// 1. Punto de proyeccion exacto (no solo la distancia): una vez que se
		//    sabe que la proyeccion cae dentro de [a,b], el punto proyectado es
		//    a + (t / baseLen2) * (b - a), con t = punto(a,b,p).
		// 2. Coordenadas enteras: punto() y cruz() dan resultado exacto con
		//    long long; solo distEuclidiana y la division final requieren double.
		// 3. Distancia de un punto a una POLILINEA (varios segmentos consecutivos):
		//    tomar el minimo de distASegmento sobre cada segmento consecutivo de
		//    la polilinea, O(k) para k segmentos.
	};
\end{lstlisting}

\textbf{Ejemplo de funcionamiento:}
\begin{lstlisting}[style=cpp]
	DistanciaPuntoLinea::Punto a{0, 0}, b{10, 0}, p{5, 3}, q{15, 3};
	
	double d1 = DistanciaPuntoLinea::distASegmento(a, b, p);
	// d1 = 3 (la proyeccion de p cae dentro de [a,b], en (5,0); distancia perpendicular = 3)
	
	double d2 = DistanciaPuntoLinea::distASegmento(a, b, q);
	// d2 = distEuclidiana(b, q) = sqrt(5^2 + 3^2) = sqrt(34)
	// (la proyeccion de q cae despues de b=(10,0), asi que la distancia real es
	//  al extremo b, no la perpendicular a la recta; hay que correrlo para
	//  confirmar el valor decimal exacto de sqrt(34))
\end{lstlisting}

\textbf{Regla práctica:} usa \texttt{distARecta} solo cuando el problema explícitamente hable de una línea infinita (poco común en CP); en la gran mayoría de los casos donde se menciona un "segmento", "pared" o "borde", usa \texttt{distASegmento}, que maneja correctamente los casos donde la proyección cae fuera del segmento —olvidar esta distinción es la fuente de error más común en problemas de distancia punto-segmento.

\subsection{Reflexión de un Punto sobre una Línea}

\textbf{Idea:} reflejar un punto $P$ sobre una recta definida por $A$ y $B$ significa encontrar su "imagen espejo" al otro lado de la recta, a la misma distancia perpendicular. El procedimiento se apoya directamente en \textit{Distancia de un Punto a una Línea/Segmento}: primero se calcula la proyección ortogonal $Q$ de $P$ sobre la recta (usando el producto punto para hallar $t$, igual que en el cálculo de distancia a segmento), y luego el punto reflejado es $P' = 2Q - P$ —es decir, $Q$ es el punto medio entre $P$ y su reflejo $P'$.

\textbf{¿Por qué usarlo?} Es la forma directa de resolver problemas de simetría sin recurrir a matrices de transformación generales o a trigonometría explícita —basta con una proyección (producto punto) y una resta vectorial. A diferencia de la \textit{Rotación de un Punto}, que gira alrededor de un pivote por un ángulo arbitrario, la reflexión es una transformación fija determinada completamente por la recta de simetría, sin parámetro de ángulo.

\textbf{Casos de uso:}
\begin{itemize}
	\item Problemas de "camino más corto que toca una pared" (reflejar el destino y trazar una línea recta).
	\item Verificar si una figura es simétrica respecto a un eje dado.
	\item Generar el punto opuesto en construcciones geométricas con simetría axial.
	\item Simulación de rebotes de un rayo de luz o trayectoria en un espejo/pared recta.
\end{itemize}

\textbf{Complejidad:} $O(1)$ —una proyección (producto punto) más una resta vectorial.

\begin{lstlisting}[style=cpp]
	// CREATE/READ: Reflexion de un punto sobre una recta definida por dos puntos
	struct ReflexionPunto {
		struct Punto { double x, y; };
		
		// READ: producto punto de (b-a) con (p-a)
		static double punto(const Punto& a, const Punto& b, const Punto& p) {
			return (b.x - a.x) * (p.x - a.x) + (b.y - a.y) * (p.y - a.y);
		} // O(1)
		
		// READ: proyeccion ortogonal de p sobre la recta que pasa por a y b
		static Punto proyeccion(const Punto& a, const Punto& b, const Punto& p) {
			double baseLen2 = punto(a, b, b); // |ab|^2
			double t = punto(a, b, p) / baseLen2;
			return Punto{a.x + t * (b.x - a.x), a.y + t * (b.y - a.y)};
		} // O(1)
		
		// READ: reflexion de p sobre la recta que pasa por a y b
		static Punto reflejar(const Punto& a, const Punto& b, const Punto& p) {
			Punto q = proyeccion(a, b, p);
			return Punto{2 * q.x - p.x, 2 * q.y - p.y};
		} // O(1)
		
		// ------------------------------------------------------------
		// MODIFICACIONES Y COMENTARIOS CON LOS USOS
		// ------------------------------------------------------------
		// 1. Reflexion sobre el eje x o eje y (casos particulares): reflejar
		//    sobre el eje x es simplemente (p.x, -p.y); sobre el eje y es
		//    (-p.x, p.y). Mas rapido que usar el caso general si el eje de
		//    simetria coincide con uno de los ejes cartesianos.
		// 2. Camino mas corto que toca una linea (problema clasico): para ir de
		//    P a Q tocando la recta L, se refleja Q (o P) sobre L con reflejar(),
		//    y la distancia minima es distEuclidiana(P, reflejo de Q); el punto
		//    de contacto es la interseccion entre el segmento P-reflejo(Q) y L.
		// 3. Reflexion sobre un PUNTO (no una recta): caso distinto y mas simple,
		//    P' = 2*C - P donde C es el punto de reflexion (equivale a una
		//    rotacion de 180 grados alrededor de C, ver Rotacion de un Punto).
	};
\end{lstlisting}

\textbf{Ejemplo de funcionamiento:}
\begin{lstlisting}[style=cpp]
	ReflexionPunto::Punto a{0, 0}, b{10, 0}, p{3, 4};
	
	auto resultado = ReflexionPunto::reflejar(a, b, p);
	// resultado = (3, -4)
	// (la recta a-b es el eje x; reflejar (3,4) sobre el eje x da (3,-4))
\end{lstlisting}

\textbf{Nota:} si la recta de reflexión coincide con uno de los ejes cartesianos, el caso particular mencionado en la modificación 1 del código es más directo y evita el cálculo de proyección completo —pero el método general funciona igual de bien para cualquier recta arbitraria.

\textbf{Regla práctica:} usa reflexión cuando el problema tenga simetría axial fija respecto a una recta conocida (más común en problemas de "camino más corto tocando una pared"); si en cambio se necesita girar un punto un ángulo arbitrario alrededor de un pivote, sin relación a ninguna recta de simetría, usa \textit{Rotación de un Punto} en su lugar.
\subsection{Rotación de un Punto (por Ángulo, alrededor de un Pivote)}

\textbf{Idea:} rotar un punto $P$ un ángulo $\theta$ alrededor de un pivote $C$ se hace en tres pasos: trasladar $P$ para que $C$ quede en el origen ($P - C$), aplicar la matriz de rotación estándar $\begin{pmatrix}\cos\theta & -\sin\theta \\ \sin\theta & \cos\theta\end{pmatrix}$ sobre ese vector trasladado, y trasladar de vuelta sumando $C$. La matriz de rotación en sí sale de proyectar el vector sobre una nueva base girada $\theta$ radianes —exactamente la misma idea de descomponer en componentes que aparece en \textit{Vector Geométrico}, aplicada a los vectores base $(1,0)$ y $(0,1)$ en vez de a un vector arbitrario.

\textbf{¿Por qué usarlo?} Es la transformación general de la que \textit{Reflexión de un Punto sobre una Línea} es un caso especial no lineal (la reflexión no es una rotación, pero comparten la técnica de trasladar-transformar-destrasladar). A diferencia de la reflexión, que está fijada por una recta de simetría, la rotación tiene un parámetro libre ($\theta$) y un pivote libre ($C$), lo que la hace la herramienta correcta cuando el problema pide girar figuras completas, no reflejarlas.

\textbf{Casos de uso:}
\begin{itemize}
	\item Girar un polígono completo un ángulo dado alrededor de su centroide u otro pivote.
	\item Generar los vértices de un polígono regular rotando un vértice inicial en incrementos de $2\pi/n$.
	\item Simular el movimiento angular de un objeto (radar, aguja, brazo robótico) en pasos discretos.
	\item Normalizar la orientación de una figura antes de compararla con otra (rotar hasta alinear un eje de referencia).
\end{itemize}

\textbf{Complejidad:} $O(1)$ por rotación de un punto —dos evaluaciones trigonométricas ($\sin$, $\cos$) reutilizables si se rota el mismo ángulo varias veces, más un par de multiplicaciones y sumas.

\begin{lstlisting}[style=cpp]
	// CREATE/READ: Rotacion de un punto un angulo theta alrededor de un pivote
	struct RotacionPunto {
		struct Punto { double x, y; };
		
		// READ: rota p un angulo theta (en radianes, sentido antihorario)
		// alrededor del pivote c
		static Punto rotar(const Punto& p, const Punto& c, double theta) {
			double dx = p.x - c.x;
			double dy = p.y - c.y;
			double cosT = cos(theta);
			double sinT = sin(theta);
			double xRot = dx * cosT - dy * sinT;
			double yRot = dx * sinT + dy * cosT;
			return Punto{c.x + xRot, c.y + yRot};
		} // O(1)
		
		// ------------------------------------------------------------
		// MODIFICACIONES Y COMENTARIOS CON LOS USOS
		// ------------------------------------------------------------
		// 1. Rotar alrededor del origen (caso particular, c = (0,0)): se puede
		//    omitir la traslacion inicial y final, usando directamente
		//    xRot = p.x*cosT - p.y*sinT; yRot = p.x*sinT + p.y*cosT.
		// 2. Rotaciones repetidas del MISMO angulo theta muchas veces (por
		//    ejemplo, generar los n vertices de un poligono regular): calcular
		//    cosT y sinT una sola vez fuera del bucle, en vez de recalcularlos
		//    en cada llamada a rotar().
		// 3. Sentido horario en vez de antihorario: usar theta negativo, o
		//    equivalentemente intercambiar el signo de sinT en ambas formulas
		//    (dx*(-sinT) para xRot, dx*sinT invertido para yRot).
		// 4. Poligono regular de n lados con radio r centrado en c: generar
		//    vertices como rotar(Punto{c.x + r, c.y}, c, 2*M_PI*i/n) para
		//    i = 0 hasta n-1.
	};
\end{lstlisting}

\textbf{Ejemplo de funcionamiento:}
\begin{lstlisting}[style=cpp]
	RotacionPunto::Punto p{1, 0}, c{0, 0};
	
	auto resultado = RotacionPunto::rotar(p, c, M_PI / 2);
	// resultado = (0, 1)
	// (rotar (1,0) 90 grados antihorario alrededor del origen da (0,1);
	//  hay que correrlo para confirmar el valor decimal exacto por el uso
	//  de M_PI en punto flotante)
\end{lstlisting}

\textbf{Nota:} si el pivote coincide con el origen, omitir la traslación (ver modificación 1) ahorra dos sumas y dos restas por punto —insignificante para un solo punto, pero relevante si se rotan miles de puntos en un bucle.

\textbf{Regla práctica:} usa rotación cuando el problema pida girar un punto o figura un ángulo arbitrario alrededor de un pivote arbitrario; si en cambio el problema tiene una recta de simetría fija y pide el "espejo" de un punto respecto a ella, eso no es una rotación sino \textit{Reflexión de un Punto sobre una Línea} —ambas comparten la técnica de trasladar-transformar-destrasladar, pero resuelven preguntas distintas.

\section{Triángulos}
\subsection{Ángulos de un Triángulo}

\textbf{Idea:} dados los tres vértices $A$, $B$, $C$ de un triángulo, el ángulo interno en cada vértice se calcula con \textit{Ángulo entre Dos Vectores} aplicado a los dos lados que se encuentran en ese vértice: el ángulo en $A$ es el ángulo entre $\vec{AB}$ y $\vec{AC}$, y análogamente para $B$ y $C$. Alternativamente, si solo se conocen las tres longitudes de los lados $a$, $b$, $c$ (sin coordenadas), se puede usar la \textbf{ley de cosenos} despejada: $\cos A = \dfrac{b^2+c^2-a^2}{2bc}$, donde el lado $a$ es el opuesto al ángulo $A$. Ambos métodos son equivalentes; cuál usar depende de si el problema da coordenadas o solo longitudes.

\textbf{¿Por qué usarlo?} Cuando se tienen coordenadas, usar \textit{Ángulo entre Dos Vectores} (con \texttt{atan2}) es preferible porque da el ángulo con signo y es numéricamente estable, como ya se explicó en esa subsección. Cuando solo se tienen las tres longitudes de los lados —sin coordenadas en el plano—, la ley de cosenos es la única opción directa, y conviene usar \texttt{acos} ahí porque el ángulo interno de un triángulo válido siempre está en $[0,\pi]$, rango donde \texttt{acos} no tiene la ambigüedad de signo que sí le importa a un ángulo con orientación.

\textbf{Casos de uso:}
\begin{itemize}
	\item Clasificar un triángulo como acutángulo, rectángulo u obtusángulo según sus ángulos.
	\item Verificar si tres longitudes dadas forman un triángulo válido (todos los ángulos resultantes deben ser reales, es decir, el argumento de \texttt{acos} debe estar en $[-1,1]$).
	\item Calcular ángulos internos como paso previo a fórmulas de área o de triangulación que los requieran.
	\item Determinar la orientación relativa de dos lados que comparten un vértice.
\end{itemize}

\textbf{Complejidad:} $O(1)$ por ángulo —ya sea vía \textit{Ángulo entre Dos Vectores} ($O(1)$, un \texttt{atan2}) o vía ley de cosenos ($O(1)$, un \texttt{acos}).

\begin{lstlisting}[style=cpp]
	// CREATE/READ: Angulos internos de un triangulo, por coordenadas o por
	// longitudes de lados (ley de cosenos)
	struct AngulosTriangulo {
		struct Punto { double x, y; };
		
		// READ: producto punto de (b-a) con (c-a)
		static double punto(const Punto& a, const Punto& b, const Punto& c) {
			return (b.x - a.x) * (c.x - a.x) + (b.y - a.y) * (c.y - a.y);
		} // O(1)
		
		// READ: producto cruz de (b-a) con (c-a)
		static double cruz(const Punto& a, const Punto& b, const Punto& c) {
			return (b.x - a.x) * (c.y - a.y) - (b.y - a.y) * (c.x - a.x);
		} // O(1)
		
		// READ: angulo interno en el vertice a, dado por los otros dos vertices
		// b y c, usando atan2 (estable, sin importar el rango)
		static double anguloPorCoordenadas(const Punto& a, const Punto& b, const Punto& c) {
			return fabs(atan2(cruz(a, b, c), punto(a, b, c)));
		} // O(1)
		
		// READ: angulo opuesto al lado de longitud "opuesto", dados los otros
		// dos lados "lado1" y "lado2" (ley de cosenos despejada)
		static double anguloPorLados(double lado1, double lado2, double opuesto) {
			double cosAngulo = (lado1*lado1 + lado2*lado2 - opuesto*opuesto) / (2*lado1*lado2);
			cosAngulo = max(-1.0, min(1.0, cosAngulo)); // recorta error de punto flotante
			return acos(cosAngulo);
		} // O(1)
		
		// ------------------------------------------------------------
		// MODIFICACIONES Y COMENTARIOS CON LOS USOS
		// ------------------------------------------------------------
		// 1. Verificar validez de un triangulo dados tres lados: ademas de la
		//    desigualdad triangular (cada lado menor que la suma de los otros
		//    dos), el recorte de cosAngulo a [-1,1] en anguloPorLados evita NaN
		//    por error de punto flotante en casos borde (triangulo degenerado).
		// 2. Clasificar el triangulo: si el mayor angulo es > pi/2 es obtusangulo,
		//    == pi/2 (con tolerancia EPS) es rectangulo, < pi/2 es acutangulo;
		//    basta con calcular el angulo opuesto al lado mas largo.
		// 3. Suma de angulos internos: por construccion siempre da pi (180
		//    grados) en geometria euclidiana; sirve como chequeo de cordura
		//    (sanity check) al depurar, calculando los tres angulos y sumando.
	};
\end{lstlisting}

\textbf{Ejemplo de funcionamiento:}
\begin{lstlisting}[style=cpp]
	AngulosTriangulo::Punto a{0, 0}, b{4, 0}, c{0, 3};
	
	double anguloA = AngulosTriangulo::anguloPorCoordenadas(a, b, c);
	// anguloA = pi/2 (aprox 1.5708)
	// (triangulo rectangulo 3-4-5, el angulo recto esta en a; hay que correrlo
	//  para confirmar el valor decimal exacto)
	
	double anguloPorLados = AngulosTriangulo::anguloPorLados(4, 3, 5);
	// anguloPorLados = pi/2, mismo resultado via ley de cosenos con los lados
	// adyacentes a A (4 y 3) y el lado opuesto a A (la hipotenusa, 5)
\end{lstlisting}

\textbf{Regla práctica:} si tienes coordenadas de los vértices, usa \texttt{anguloPorCoordenadas} (reutiliza \textit{Ángulo entre Dos Vectores}); si solo tienes las tres longitudes de los lados, usa \texttt{anguloPorLados} y no olvides recortar el argumento de \texttt{acos} a $[-1,1]$ para evitar \texttt{NaN} por error de punto flotante en triángulos casi degenerados.

\subsection{Fórmula de Herón (Área por Lados)}

\textbf{Idea:} dado un triángulo del que solo se conocen las longitudes de sus tres lados $a$, $b$, $c$ (sin coordenadas), la fórmula de Herón calcula el área sin necesidad de conocer ningún ángulo: con el semiperímetro $s = \dfrac{a+b+c}{2}$, el área es $\sqrt{s(s-a)(s-b)(s-c)}$. Es un resultado puramente algebraico —se puede derivar combinando la ley de cosenos (para obtener un ángulo) con la fórmula de área $\frac{1}{2}ab\sin C$, pero el resultado final elimina la trigonometría por completo y deja todo en función de las tres longitudes.

\textbf{¿Por qué usarlo?} Cuando el problema da coordenadas de los vértices, \textit{Área de un Polígono (Shoelace)} es más directo y general (funciona para cualquier polígono, no solo triángulos). Herón es la herramienta correcta específicamente cuando \textbf{no hay coordenadas}, solo las tres longitudes de los lados —situación común en problemas donde el triángulo se describe por sus medidas, no por su posición en el plano.

\textbf{Casos de uso:}
\begin{itemize}
	\item Calcular el área de un triángulo dado solo por sus tres lados (sin coordenadas).
	\item Verificar si tres longitudes forman un triángulo válido (el área resultante debe ser un número real positivo, es decir, el producto bajo la raíz debe ser $\geq 0$).
	\item Como paso intermedio en fórmulas que requieren área y perímetro juntos (por ejemplo, el radio del círculo inscrito, $r = \text{área}/s$).
	\item Cálculo de área en problemas donde los datos de entrada son distancias medidas, no coordenadas cartesianas.
\end{itemize}

\textbf{Complejidad:} $O(1)$ —tres sumas/restas para el semiperímetro y los factores, una multiplicación de cuatro términos, y una raíz cuadrada.

\begin{lstlisting}[style=cpp]
	// CREATE/READ: Area de un triangulo dado por sus tres lados (formula de Heron)
	struct FormulaHeron {
		// READ: area del triangulo de lados a, b, c; retorna -1 si los lados
		// no forman un triangulo valido (desigualdad triangular violada)
		static double area(double a, double b, double c) {
			if (a + b <= c || a + c <= b || b + c <= a) return -1; // triangulo invalido
			double s = (a + b + c) / 2.0;
			double producto = s * (s - a) * (s - b) * (s - c);
			if (producto < 0) producto = 0; // recorta error de punto flotante
			return sqrt(producto);
		} // O(1)
		
		// ------------------------------------------------------------
		// MODIFICACIONES Y COMENTARIOS CON LOS USOS
		// ------------------------------------------------------------
		// 1. Radio del circulo inscrito: r = area(a,b,c) / s, con
		//    s = (a+b+c)/2 el mismo semiperimetro usado internamente.
		// 2. Radio del circulo circunscrito: R = (a*b*c) / (4*area(a,b,c)),
		//    formula derivada de la ley de senos generalizada (ver Circulo
		//    Circunscrito en un Triangulo para el caso con coordenadas).
		// 3. Coordenadas disponibles en vez de lados: si el triangulo viene dado
		//    por vertices, es mas directo y mas preciso usar Shoelace
		//    directamente en vez de calcular las tres distancias euclidianas y
		//    despues aplicar Heron (evita acumular error en un paso intermedio).
	};
\end{lstlisting}

\textbf{Ejemplo de funcionamiento:}
\begin{lstlisting}[style=cpp]
	double resultado = FormulaHeron::area(3, 4, 5);
	// resultado = 6
	// (triangulo rectangulo 3-4-5, area = base*altura/2 = 3*4/2 = 6,
	//  Heron da el mismo resultado sin usar el angulo recto explicitamente)
\end{lstlisting}

\textbf{Regla práctica:} usa Herón únicamente cuando el triángulo se te da por sus tres longitudes sin coordenadas; si tienes coordenadas de los vértices, usa \textit{Área de un Polígono (Shoelace)} directamente —es igual de rápido, más general, y evita el paso intermedio de calcular tres distancias euclidianas antes de aplicar Herón.
\subsection{Círculo Circunscrito en un Triángulo}

\textbf{Idea:} el círculo circunscrito (circumcírculo) de un triángulo es el único círculo que pasa por sus tres vértices. Su centro, el \textbf{circuncentro}, es el punto equidistante de los tres vértices, y se encuentra como la intersección de las mediatrices (perpendiculares por el punto medio) de cualquier par de lados. Resolviendo el sistema de ecuaciones que impone "distancia igual a los tres vértices" se llega a una fórmula cerrada en términos de determinantes 2x2, sin necesidad de intersectar rectas explícitamente. El radio es simplemente la distancia del circuncentro a cualquiera de los tres vértices.

\textbf{¿Por qué usarlo?} A diferencia del caso particular de un triángulo equilátero —donde el circuncentro coincide con el centroide, el promedio simple de los tres vértices, $O(1)$ trivial—, un triángulo arbitrario no tiene esa coincidencia: centroide, circuncentro, incentro y ortocentro son puntos distintos en general. La fórmula con determinantes cubre el caso general sin necesidad de distinguir tipos de triángulo, y evita resolver el sistema de mediatrices manualmente (que requeriría manejar el caso de mediatriz vertical aparte, el mismo problema que motivó la forma general de \textit{Ecuación de una Línea Dados Dos Puntos}).

\textbf{Casos de uso:}
\begin{itemize}
	\item Encontrar el centro y radio del círculo mínimo que pasa por tres puntos dados (no necesariamente el círculo mínimo que los \textit{cubre}, ver nota).
	\item Problemas de triangulación de Delaunay, donde la propiedad definitoria es que ningún otro punto cae dentro del circumcírculo de cada triángulo.
	\item Construcciones geométricas clásicas (bisectores perpendiculares, geometría de triángulos).
	\item Verificar si cuatro puntos son concíclicos (el cuarto debe estar exactamente sobre el circumcírculo de los otros tres).
\end{itemize}

\textbf{Complejidad:} $O(1)$ —el sistema de ecuaciones se resuelve con una fórmula cerrada de determinantes 2x2, sin iteración.

\begin{lstlisting}[style=cpp]
	// CREATE/READ: Circuncentro y radio del circulo circunscrito de un triangulo
	// (caso general, valido para cualquier triangulo no degenerado)
	struct CirculoCircunscrito {
		struct Punto { double x, y; };
		
		// READ: producto cruz de (b-a) x (c-a), usado para detectar triangulo degenerado
		static double cruz(const Punto& a, const Punto& b, const Punto& c) {
			return (b.x - a.x) * (c.y - a.y) - (b.y - a.y) * (c.x - a.x);
		} // O(1)
		
		// READ: calcula el circuncentro de a, b, c mediante formula de
		// determinantes (interseccion de mediatrices resuelta algebraicamente)
		static Punto circuncentro(const Punto& a, const Punto& b, const Punto& c) {
			double d = 2.0 * (a.x * (b.y - c.y) + b.x * (c.y - a.y) + c.x * (a.y - b.y));
			// d == 0 significa que a, b, c son colineales: no hay circuncentro
			double ax2ay2 = a.x*a.x + a.y*a.y;
			double bx2by2 = b.x*b.x + b.y*b.y;
			double cx2cy2 = c.x*c.x + c.y*c.y;
			double ux = (ax2ay2 * (b.y - c.y) + bx2by2 * (c.y - a.y) + cx2cy2 * (a.y - b.y)) / d;
			double uy = (ax2ay2 * (c.x - b.x) + bx2by2 * (a.x - c.x) + cx2cy2 * (b.x - a.x)) / d;
			return Punto{ux, uy};
		} // O(1)
		
		// READ: radio del circulo circunscrito (distancia del circuncentro a
		// cualquier vertice, aqui se usa a)
		static double radio(const Punto& a, const Punto& b, const Punto& c) {
			Punto o = circuncentro(a, b, c);
			double dx = a.x - o.x, dy = a.y - o.y;
			return sqrt(dx*dx + dy*dy);
		} // O(1)
		
		// ------------------------------------------------------------
		// MODIFICACIONES Y COMENTARIOS CON LOS USOS
		// ------------------------------------------------------------
		// 1. Caso equilatero (atajo, NO usar para triangulos generales): el
		//    circuncentro coincide con el centroide, ((a.x+b.x+c.x)/3,
		//    (a.y+b.y+c.y)/3), y el radio es lado / sqrt(3). Solo valido cuando
		//    se garantiza que el triangulo es equilatero.
		// 2. Triangulo degenerado (colineal): si cruz(a,b,c) == 0 (con EPS), d
		//    en circuncentro() da 0 y la formula produce division por cero /
		//    NaN; verificar colinealidad ANTES de llamar circuncentro().
		// 3. Radio via lados y area (alternativa sin circuncentro explicito):
		//    R = (lado_a * lado_b * lado_c) / (4 * area), reutilizando
		//    Formula de Heron para el area si solo se necesita el radio y no
		//    el centro exacto.
	};
\end{lstlisting}

\textbf{Ejemplo de funcionamiento:}
\begin{lstlisting}[style=cpp]
	CirculoCircunscrito::Punto a{0, 0}, b{4, 0}, c{0, 3};
	
	auto centro = CirculoCircunscrito::circuncentro(a, b, c);
	// centro = (2, 1.5)
	// (triangulo rectangulo 3-4-5; el circuncentro de un triangulo rectangulo
	//  siempre cae exactamente en el punto medio de la hipotenusa, aqui el
	//  punto medio de b=(4,0) y c=(0,3) es (2, 1.5), lo cual confirma el resultado)
	
	double r = CirculoCircunscrito::radio(a, b, c);
	// r = 2.5
	// (para un triangulo rectangulo, el radio circunscrito es siempre la
	//  mitad de la hipotenusa: 5/2 = 2.5)
\end{lstlisting}

\textbf{Nota:} asumí que necesitas el caso general (no solo equilátero, que era tu ejemplo original) porque es estrictamente más útil en CP —el caso equilátero queda como atajo $O(1)$ documentado en la modificación 1 del código. Si en algún problema se garantiza que el triángulo es equilátero, usa ese atajo en vez de la fórmula de determinantes completa.

\textbf{Regla práctica:} usa esta fórmula de determinantes para el circuncentro general; verifica primero que el triángulo no sea degenerado (\texttt{cruz(a,b,c) != 0}) para evitar división por cero. Si solo necesitas el radio y ya tienes las tres longitudes de los lados (por ejemplo, después de calcular \textit{Fórmula de Herón}), la fórmula $R = \dfrac{abc}{4 \cdot \text{área}}$ es más directa y evita calcular el centro explícitamente.

\subsection{Punto Dentro de un Triángulo}

\textbf{Idea:} un punto $P$ está dentro del triángulo $ABC$ si y solo si queda del \textbf{mismo lado} de las tres aristas orientadas $AB$, $BC$, $CA$ —es decir, si el signo de \textit{Orientación de Tres Puntos} aplicado a $(A,B,P)$, $(B,C,P)$ y $(C,A,P)$ es el mismo para los tres (todos positivos o todos negativos, según el sentido de recorrido del triángulo). Si alguno de los tres productos cruz da signo distinto a los demás, $P$ cae fuera de al menos una arista y por lo tanto fuera del triángulo.

\textbf{¿Por qué usarlo?} Es la aplicación más directa y económica de \textit{Orientación de Tres Puntos} a un caso concreto: tres evaluaciones de producto cruz, sin necesidad de trigonometría, áreas parciales, o coordenadas baricéntricas explícitas (aunque el método es equivalente a verificar que las coordenadas baricéntricas de $P$ sean todas no negativas). Es un caso particular y más simple de \textit{Punto Dentro de un Polígono}, que en general requiere \textit{ray casting}; para un triángulo (siempre convexo), tres chequeos de orientación bastan y son más baratos.

\textbf{Casos de uso:}
\begin{itemize}
	\item Verificar pertenencia de un punto a un triángulo, como paso base antes de generalizar a polígonos convexos.
	\item Localización de punto en una triangulación (mesh), recorriendo triángulos y probando cuál lo contiene.
	\item Pruebas de colisión simples en gráficos 2D o videojuegos con geometría triangular.
	\item Interpolación baricéntrica (una vez confirmado que el punto está dentro, los mismos productos cruz dan los pesos baricéntricos).
\end{itemize}

\textbf{Complejidad:} $O(1)$ —tres evaluaciones de producto cruz y una comparación de signos.

\begin{lstlisting}[style=cpp]
	// CREATE/READ: Verifica si un punto esta dentro de un triangulo mediante
	// tres pruebas de orientacion (mismo signo en las tres aristas)
	struct PuntoDentroTriangulo {
		struct Punto { double x, y; };
		
		// READ: producto cruz de (b-a) x (c-a)
		static double cruz(const Punto& a, const Punto& b, const Punto& c) {
			return (b.x - a.x) * (c.y - a.y) - (b.y - a.y) * (c.x - a.x);
		} // O(1)
		
		// READ: true si p esta dentro o sobre el borde del triangulo abc,
		// sin importar si abc esta orientado CW o CCW
		static bool dentro(const Punto& a, const Punto& b, const Punto& c, const Punto& p) {
			double d1 = cruz(a, b, p);
			double d2 = cruz(b, c, p);
			double d3 = cruz(c, a, p);
			const double EPS = 1e-9;
			bool tieneNegativo = (d1 < -EPS) || (d2 < -EPS) || (d3 < -EPS);
			bool tienePositivo = (d1 > EPS) || (d2 > EPS) || (d3 > EPS);
			// si NO hay a la vez signos positivos y negativos, p esta dentro
			// (o sobre el borde, si algun di es aproximadamente 0)
			return !(tieneNegativo && tienePositivo);
		} // O(1)
		
		// ------------------------------------------------------------
		// MODIFICACIONES Y COMENTARIOS CON LOS USOS
		// ------------------------------------------------------------
		// 1. Distinguir "estrictamente dentro" de "sobre el borde": si se
		//    necesita excluir el borde, verificar ademas que ninguno de d1, d2,
		//    d3 este dentro del rango [-EPS, EPS] (es decir, que los tres sean
		//    estrictamente del mismo signo, sin ceros).
		// 2. Coordenadas enteras: usar long long en cruz() y comparar contra 0
		//    exacto en vez de EPS, ya que el resultado es exacto con enteros.
		// 3. Coordenadas baricentricas: una vez calculados d1, d2, d3 (todas del
		//    mismo signo), los pesos baricentricos normalizados son
		//    d2/(d1+d2+d3), d3/(d1+d2+d3), d1/(d1+d2+d3) para a, b, c
		//    respectivamente (util para interpolar valores definidos en los
		//    vertices, como en shading de graficos 3D).
	};
\end{lstlisting}

\textbf{Ejemplo de funcionamiento:}
\begin{lstlisting}[style=cpp]
	PuntoDentroTriangulo::Punto a{0, 0}, b{4, 0}, c{0, 4}, p{1, 1}, q{3, 3};
	
	bool dentroP = PuntoDentroTriangulo::dentro(a, b, c, p);
	// dentroP = true (p=(1,1) esta dentro del triangulo)
	
	bool dentroQ = PuntoDentroTriangulo::dentro(a, b, c, q);
	// dentroQ = false (q=(3,3) esta fuera, mas alla de la hipotenusa b-c)
\end{lstlisting}

\textbf{Regla práctica:} usa esta prueba de tres orientaciones para triángulos específicamente —es más barata que \textit{Punto Dentro de un Polígono (Ray Casting)}, que es la herramienta general para polígonos con cualquier número de lados. Si el triángulo forma parte de una triangulación más grande, esta misma prueba es la base para localizar en qué triángulo cae un punto dado.

\section{Círculos}
\subsection{Longitud de Arco y Área de un Sector Circular}

\textbf{Idea:} un sector circular es la "porción de pizza" delimitada por dos radios y el arco entre ellos, definida por un ángulo central $\theta$ (en radianes) sobre un círculo de radio $r$. Como la circunferencia completa ($2\pi r$ de longitud, $\pi r^2$ de área) corresponde a un ángulo completo de $2\pi$ radianes, cualquier sector es simplemente una fracción proporcional: la longitud de arco es $L = r\theta$ y el área del sector es $A = \frac{1}{2}r^2\theta$. Ambas fórmulas se derivan de la misma razón $\theta / 2\pi$ aplicada a la circunferencia y al área totales, respectivamente.

\textbf{¿Por qué usarlo?} Es la forma más directa de estas dos cantidades cuando se trabaja con ángulos en radianes —evita el paso intermedio de convertir a grados y aplicar reglas de tres, y generaliza sin casos especiales a cualquier ángulo entre $0$ y $2\pi$. Vale la pena notar la diferencia de exponente entre ambas fórmulas: la longitud de arco escala linealmente con $r$, mientras que el área escala con $r^2$ —el mismo patrón que la circunferencia ($2\pi r$) y el área del círculo completo ($\pi r^2$).

\textbf{Casos de uso:}
\begin{itemize}
	\item Calcular la distancia recorrida por un objeto que se mueve sobre un arco (radar, engranaje, trayectoria circular).
	\item Área de una región delimitada por dos radios y un arco, como paso intermedio en problemas de intersección de círculos.
	\item Problemas de "porción de pizza" o reparto proporcional de un círculo.
	\item Cálculo de área de un segmento circular (sector menos el triángulo que forman los dos radios y la cuerda), usado en \textit{Intersección de Dos Círculos}.
\end{itemize}

\textbf{Complejidad:} $O(1)$ —una multiplicación para cada fórmula.

\begin{lstlisting}[style=cpp]
	// CREATE/READ: Longitud de arco y area de un sector circular dado el radio
	// y el angulo central en radianes
	struct SectorCircular {
		// READ: longitud de arco para un sector de radio r y angulo theta (radianes)
		static double longitudArco(double r, double theta) {
			return r * theta;
		} // O(1)
		
		// READ: area del sector de radio r y angulo theta (radianes)
		static double areaSector(double r, double theta) {
			return 0.5 * r * r * theta;
		} // O(1)
		
		// READ: area del segmento circular (sector menos el triangulo r-r-cuerda),
		// la region entre el arco y la cuerda que lo subtiende
		static double areaSegmento(double r, double theta) {
			return areaSector(r, theta) - 0.5 * r * r * sin(theta);
		} // O(1)
		
		// ------------------------------------------------------------
		// MODIFICACIONES Y COMENTARIOS CON LOS USOS
		// ------------------------------------------------------------
		// 1. Angulo dado en grados: convertir primero a radianes con
		//    theta_rad = theta_grados * M_PI / 180.0 antes de llamar cualquiera
		//    de las funciones.
		// 2. Perimetro total del sector (no solo el arco): longitudArco(r,theta)
		//    + 2*r, sumando los dos radios que cierran la "porcion de pizza".
		// 3. area Segmento requiere theta en [0, 2*pi]; si theta > pi (segmento
		//    "mayor", mas de medio circulo), la resta sigue siendo valida porque
		//    sin(theta) ya maneja el signo correctamente para ese rango.
	};
\end{lstlisting}

\textbf{Ejemplo de funcionamiento:}
\begin{lstlisting}[style=cpp]
	double r = 4.0, theta = M_PI / 2; // sector de 90 grados
	
	double arco = SectorCircular::longitudArco(r, theta);
	// arco = 2*pi (aprox 6.2832)
	// (un cuarto de la circunferencia completa 2*pi*r = 8*pi; hay que correrlo
	//  para confirmar el valor decimal exacto)
	
	double area = SectorCircular::areaSector(r, theta);
	// area = 4*pi (aprox 12.566)
	// (un cuarto del area total pi*r^2 = 16*pi; hay que correrlo para
	//  confirmar el valor decimal exacto)
\end{lstlisting}

\textbf{Regla práctica:} usa estas fórmulas directamente cuando el ángulo esté en radianes y conozcas el radio; el área del segmento circular (arco menos cuerda) es el bloque reutilizable que hace falta como paso intermedio en \textit{Intersección de Dos Círculos}, donde el área de solape se descompone en dos segmentos circulares, uno por cada círculo.

\subsection{Intersección de Dos Círculos}

\textbf{Idea:} dos círculos de centros $C_1, C_2$ y radios $r_1, r_2$, separados por una distancia $d = |C_1C_2|$, se intersectan en 0, 1, 2 puntos, o son el mismo círculo, dependiendo de cómo se comparan $d$, $r_1+r_2$ y $|r_1-r_2|$. Cuando hay dos puntos de intersección, se encuentran resolviendo el sistema de las dos ecuaciones de circunferencia: se calcula la distancia $a$ desde $C_1$ hasta el punto medio de la cuerda que une las dos intersecciones (usando el teorema de Pitágoras generalizado, análogo a la ley de cosenos), luego la altura $h$ de ese punto medio a cada intersección, y finalmente se desplaza perpendicularmente desde ese punto medio usando un vector normal a $C_1C_2$ —la misma idea de "rotar 90 grados" que aparece en \textit{Ecuación de una Línea Dados Dos Puntos}.

\textbf{¿Por qué usarlo?} Es un problema con varios casos borde que hay que descartar explícitamente antes de aplicar la fórmula: círculos demasiado separados ($d > r_1+r_2$), uno dentro del otro sin tocarse ($d < |r_1-r_2|$), tangencia externa o interna (exactamente un punto), y círculos concéntricos con igual radio (infinitos puntos, caso degenerado). Manejar estos casos aparte, antes de la fórmula general, evita raíces cuadradas de números negativos o divisiones por cero cuando $d=0$.

\textbf{Casos de uso:}
\begin{itemize}
	\item Encontrar los puntos de intersección entre dos círculos, como paso base en geometría de círculos.
	\item Calcular el área de solape entre dos círculos (usando \textit{Longitud de Arco y Área de un Sector Circular} sobre cada uno).
	\item Problemas de cobertura de señal o rango (¿dos zonas circulares se superponen?).
	\item Construcciones geométricas clásicas (encontrar puntos equidistantes a dos centros dados, base de la construcción de la mediatriz con regla y compás).
\end{itemize}

\textbf{Complejidad:} $O(1)$ —un número constante de operaciones aritméticas y, en el peor caso, dos llamadas a \texttt{sqrt} y dos a \texttt{acos} (para el área de solape).

\begin{lstlisting}[style=cpp]
	// CREATE/READ: Interseccion de dos circulos: puntos de corte y area de
	// solape (union de dos segmentos circulares)
	struct InterseccionCirculos {
		struct Punto { double x, y; };
		
		static double dist(const Punto& a, const Punto& b) {
			double dx = b.x - a.x, dy = b.y - a.y;
			return sqrt(dx*dx + dy*dy);
		} // O(1)
		
		// READ: clasifica la relacion entre los dos circulos ANTES de calcular
		// puntos de interseccion, para evitar sqrt de numero negativo
		// retorna: 0 = separados (sin interseccion), 1 = uno dentro del otro
		// (sin tocarse), 2 = tangentes (1 punto), 3 = se cruzan (2 puntos),
		// 4 = concentricos con mismo radio (infinitos puntos, mismo circulo)
		static int clasificar(const Punto& c1, double r1, const Punto& c2, double r2) {
			double d = dist(c1, c2);
			const double EPS = 1e-9;
			if (d < EPS && fabs(r1 - r2) < EPS) return 4; // mismo circulo
			if (d > r1 + r2 + EPS) return 0;               // separados
			if (d < fabs(r1 - r2) - EPS) return 1;         // uno dentro del otro
			if (fabs(d - (r1 + r2)) < EPS || fabs(d - fabs(r1 - r2)) < EPS) return 2; // tangentes
			return 3; // se cruzan en dos puntos
		} // O(1)
		
		// READ: calcula los (hasta) dos puntos de interseccion; SOLO valido
		// cuando clasificar() devuelve 2 o 3 (llamar clasificar primero)
		static vector<Punto> puntosInterseccion(const Punto& c1, double r1, const Punto& c2, double r2) {
			double d = dist(c1, c2);
			// a = distancia de c1 al punto medio de la cuerda de interseccion
			double a = (r1*r1 - r2*r2 + d*d) / (2*d);
			double hCuad = r1*r1 - a*a;
			double h = (hCuad < 0) ? 0 : sqrt(hCuad); // recorta error de punto flotante
			// punto medio de la cuerda, sobre la linea c1-c2
			double mx = c1.x + a * (c2.x - c1.x) / d;
			double my = c1.y + a * (c2.y - c1.y) / d;
			// vector perpendicular a (c2 - c1), normalizado por h/d
			double dx = (c2.y - c1.y) * (h / d);
			double dy = -(c2.x - c1.x) * (h / d);
			if (h < 1e-9) return { Punto{mx, my} }; // tangentes: un solo punto
			return { Punto{mx + dx, my + dy}, Punto{mx - dx, my - dy} };
		} // O(1)
		
		// READ: area de solape (union de dos segmentos circulares); asume que
		// clasificar() ya confirmo caso 3 (se cruzan) o 1 (uno dentro del otro)
		static double areaSolape(const Punto& c1, double r1, const Punto& c2, double r2) {
			double d = dist(c1, c2);
			int caso = clasificar(c1, r1, c2, r2);
			if (caso == 0) return 0;                                  // sin solape
			if (caso == 1) return M_PI * min(r1, r2) * min(r1, r2);   // uno dentro del otro
			// angulo central en c1 y en c2 subtendido por la cuerda comun
			double theta1 = 2 * acos(max(-1.0, min(1.0, (r1*r1 + d*d - r2*r2) / (2*r1*d))));
			double theta2 = 2 * acos(max(-1.0, min(1.0, (r2*r2 + d*d - r1*r1) / (2*r2*d))));
			double seg1 = 0.5 * r1*r1 * (theta1 - sin(theta1));
			double seg2 = 0.5 * r2*r2 * (theta2 - sin(theta2));
			return seg1 + seg2;
		} // O(1)
		
		// ------------------------------------------------------------
		// MODIFICACIONES Y COMENTARIOS CON LOS USOS
		// ------------------------------------------------------------
		// 1. SIEMPRE llamar clasificar() antes que puntosInterseccion() o
		//    areaSolape(): estas dos ultimas asumen que ya se descarto el caso
		//    de circulos separados o concentricos con mismo radio, de lo
		//    contrario hay division por cero (d == 0) o resultados sin sentido.
		// 2. areaSolape reutiliza directamente areaSegmento de Longitud de Arco
		//    y Area de un Sector Circular: seg1 y seg2 son exactamente esa
		//    formula aplicada a theta1 y theta2 respectivamente.
		// 3. Radio y distancia con coordenadas enteras: d, a, h siguen
		//    requiriendo double por las raices cuadradas involucradas; no hay
		//    forma de mantener este calculo en aritmetica entera exacta.
	};
\end{lstlisting}

\textbf{Ejemplo de funcionamiento:}
\begin{lstlisting}[style=cpp]
	InterseccionCirculos::Punto c1{0, 0}, c2{4, 0};
	double r1 = 3, r2 = 3;
	
	int caso = InterseccionCirculos::clasificar(c1, r1, c2, r2);
	// caso = 3 (se cruzan en dos puntos, ya que |r1-r2|=0 < d=4 < r1+r2=6)
	
	auto puntos = InterseccionCirculos::puntosInterseccion(c1, r1, c2, r2);
	// puntos = {(2, h), (2, -h)} con h = sqrt(3^2 - 2^2) = sqrt(5)
	// (por simetria, a = (9-9+16)/8 = 2, el punto medio de la cuerda cae en
	//  x=2; hay que correrlo para confirmar el valor decimal exacto de h)
\end{lstlisting}

\textbf{Nota:} \texttt{puntosInterseccion} y \texttt{areaSolape} no validan internamente los casos degenerados (círculos separados, uno dentro del otro sin tocarse, o concéntricos) — es responsabilidad de quien llama invocar \texttt{clasificar()} primero y actuar según el resultado, tal como se indica en la modificación 1 dentro del código.

\textbf{Regla práctica:} nunca llames directamente a \texttt{puntosInterseccion} o calcules el área de solape sin antes clasificar la relación entre los dos círculos con \texttt{clasificar()} —los casos borde (separados, uno dentro del otro, concéntricos) producen \texttt{NaN} o resultados incorrectos si se ignoran, y son la fuente de error más común en este tipo de problema.

\section{Polígonos}
\subsection{Área de un Polígono (Shoelace)}

\textbf{Idea:} dado un polígono simple con vértices $P_0, P_1, \dots, P_{n-1}$ en orden (sentido horario o antihorario, consecutivos a lo largo del contorno), el área con signo es $\dfrac{1}{2}\sum_{i=0}^{n-1} (x_i y_{i+1} - x_{i+1} y_i)$, tomando los índices módulo $n$ para cerrar el polígono. Cada término $x_i y_{i+1} - x_{i+1} y_i$ es exactamente el producto cruz $\vec{OP_i} \times \vec{OP_{i+1}}$ visto en \textit{Orientación de Tres Puntos}, tomando el origen como pivote implícito: la fórmula suma las áreas con signo de los triángulos que forma el origen con cada arista consecutiva, y las contribuciones se cancelan fuera del polígono y se acumulan dentro, dejando exactamente el área encerrada.

\textbf{¿Por qué usarlo?} Es la generalización directa de \textit{Fórmula de Herón} a cualquier polígono simple (no solo triángulos), y no requiere que el polígono sea convexo. A diferencia de triangular manualmente el polígono y sumar áreas de triángulos individuales, Shoelace procesa los vértices en un único recorrido lineal sin necesidad de descomponer la figura. El signo del resultado antes de tomar valor absoluto además revela la orientación del polígono: positivo si los vértices están en sentido antihorario (CCW), negativo si horario (CW) —la misma convención fijada en \textit{Orientación de Tres Puntos}.

\textbf{Casos de uso:}
\begin{itemize}
	\item Calcular el área de cualquier polígono simple (convexo o no) dado por sus vértices en orden.
	\item Determinar la orientación (CW o CCW) de una secuencia de puntos, usando el signo antes del valor absoluto.
	\item Validar que un polígono generado por otro algoritmo (Convex Hull, recorte) tiene la orientación esperada.
	\item Paso base para el Teorema de Pick (área en malla de puntos enteros).
\end{itemize}

\textbf{Complejidad:} $O(n)$ —un recorrido lineal por los $n$ vértices, sumando un término constante por cada uno.

\begin{lstlisting}[style=cpp]
	// CREATE/READ: Area de un poligono simple mediante la formula de Shoelace
	struct AreaPoligono {
		struct Punto { double x, y; };
		
		// READ: area CON SIGNO del poligono (positivo = CCW, negativo = CW);
		// vertices deben estar en orden a lo largo del contorno
		static double areaConSigno(const vector<Punto>& vertices) {
			int n = (int)vertices.size();
			double suma = 0;
			for (int i = 0; i < n; i++) {
				const Punto& actual = vertices[i];
				const Punto& siguiente = vertices[(i + 1) % n];
				suma += actual.x * siguiente.y - siguiente.x * actual.y;
			}
			return suma / 2.0;
		} // O(n)
		
		// READ: area absoluta del poligono (sin importar orientacion)
		static double area(const vector<Punto>& vertices) {
			return fabs(areaConSigno(vertices));
		} // O(n)
		
		// READ: true si los vertices estan en sentido antihorario (CCW)
		static bool esCCW(const vector<Punto>& vertices) {
			return areaConSigno(vertices) > 0;
		} // O(n)
		
		// ------------------------------------------------------------
		// MODIFICACIONES Y COMENTARIOS CON LOS USOS
		// ------------------------------------------------------------
		// 1. Coordenadas enteras: la suma interna se puede mantener en long long
		//    (el area con signo * 2 siempre es entera si las coordenadas son
		//    enteras), dividiendo entre 2.0 solo al final; util para Teorema de
		//    Pick, que necesita el area exacta sin error de punto flotante.
		// 2. Poligono con hoyos: el area total es el area del contorno exterior
		//    (CCW) menos el area de cada hoyo interior (dado en sentido CW);
		//    sumar areaConSigno() de cada contorno respeta esto automaticamente
		//    si las orientaciones son consistentes.
		// 3. Verificar que el poligono es simple (sin autointersecciones) NO lo
		//    hace esta formula: Shoelace da un resultado numerico igual si el
		//    poligono se autointersecta, pero ese resultado ya no representa el
		//    area real de una region simple.
	};
\end{lstlisting}

\textbf{Ejemplo de funcionamiento:}
\begin{lstlisting}[style=cpp]
	vector<AreaPoligono::Punto> cuadrado = {{0,0}, {4,0}, {4,4}, {0,4}};
	
	double resultado = AreaPoligono::area(cuadrado);
	// resultado = 16
	// (cuadrado de lado 4, area = 4*4 = 16; los vertices estan en orden CCW)
	
	bool ccw = AreaPoligono::esCCW(cuadrado);
	// ccw = true
\end{lstlisting}

\textbf{Regla práctica:} usa Shoelace como la herramienta por defecto para área de polígono siempre que tengas coordenadas —es más general que Herón (funciona para cualquier número de lados, no solo triángulos) y no requiere que el polígono sea convexo. Aprovecha el signo (antes de \texttt{fabs}) cuando necesites confirmar o forzar una orientación específica, como en \textit{Envolvente Convexa} o \textit{Punto Dentro de un Polígono Convexo}, que dependen de esa convención.

\subsection{Perímetro de un Polígono}

\textbf{Idea:} el perímetro de un polígono con vértices $P_0, \dots, P_{n-1}$ en orden es simplemente la suma de las distancias euclidianas entre cada par de vértices consecutivos a lo largo del contorno, cerrando de vuelta del último vértice al primero: $\sum_{i=0}^{n-1} d(P_i, P_{i+1 \bmod n})$. Es la aplicación directa de \textit{Distancia Euclidiana} en un recorrido lineal, sin ninguna sutileza adicional —a diferencia de área (Shoelace) o de la unión de rectángulos, el perímetro no depende de orientación ni de convexidad.

\textbf{¿Por qué usarlo?} Es el cálculo más simple de esta parte, pero se documenta aparte porque aparece constantemente como componente de fórmulas más grandes (razón isoperimétrica, radio del círculo inscrito vía $r = \text{área}/\text{semiperímetro}$, problemas de "cuánta cerca se necesita"). No tiene alternativa más eficiente: a diferencia de área, que puede calcularse con una fórmula cerrada tipo Shoelace, el perímetro siempre requiere sumar las $n$ distancias, una por cada arista.

\textbf{Casos de uso:}
\begin{itemize}
	\item Calcular la longitud total de una cerca, borde o contorno poligonal.
	\item Paso intermedio en el cálculo del radio del círculo inscrito (usando semiperímetro junto al área de Herón o Shoelace).
	\item Comparar la "compacidad" de una figura (razón isoperimétrica: $4\pi \cdot \text{área} / \text{perímetro}^2$).
	\item Validar la longitud de un polígono generado por Envolvente Convexa u otro algoritmo de construcción.
\end{itemize}

\textbf{Complejidad:} $O(n)$ —un recorrido lineal por los $n$ vértices, con una distancia euclidiana ($O(1)$) por cada arista.

\begin{lstlisting}[style=cpp]
	// CREATE/READ: Perimetro de un poligono mediante suma de distancias
	// euclidianas entre vertices consecutivos
	struct PerimetroPoligono {
		struct Punto { double x, y; };
		
		// READ: distancia euclidiana entre dos puntos
		static double dist(const Punto& a, const Punto& b) {
			double dx = b.x - a.x, dy = b.y - a.y;
			return sqrt(dx*dx + dy*dy);
		} // O(1)
		
		// READ: perimetro total del poligono, sumando cada arista consecutiva
		// (incluye la arista de cierre del ultimo vertice al primero)
		static double perimetro(const vector<Punto>& vertices) {
			int n = (int)vertices.size();
			double total = 0;
			for (int i = 0; i < n; i++) {
				total += dist(vertices[i], vertices[(i + 1) % n]);
			}
			return total;
		} // O(n)
		
		// ------------------------------------------------------------
		// MODIFICACIONES Y COMENTARIOS CON LOS USOS
		// ------------------------------------------------------------
		// 1. Polilinea abierta (NO poligono cerrado): omitir la arista de cierre,
		//    recorriendo i de 0 a n-2 en vez de usar (i+1) % n.
		// 2. Radio del circulo inscrito, dado el area (Shoelace o Heron):
		//    r = area / (perimetro(vertices) / 2.0), reutilizando el
		//    semiperimetro que ya calcula esta funcion internamente.
		// 3. Razon isoperimetrica (que tan "circular" es la figura, maximo 1
		//    para un circulo perfecto): 4 * M_PI * area / (perimetro*perimetro).
	};
\end{lstlisting}

\textbf{Ejemplo de funcionamiento:}
\begin{lstlisting}[style=cpp]
	vector<PerimetroPoligono::Punto> cuadrado = {{0,0}, {4,0}, {4,4}, {0,4}};
	
	double resultado = PerimetroPoligono::perimetro(cuadrado);
	// resultado = 16
	// (cuadrado de lado 4, cuatro lados de longitud 4 cada uno: 4*4 = 16)
\end{lstlisting}

\textbf{Regla práctica:} usa esta suma directa siempre que necesites la longitud del contorno; si el problema en realidad da una \textbf{polilínea abierta} (no un polígono cerrado), recuerda omitir la arista de cierre entre el último y el primer vértice (ver modificación 1 dentro del código), ya que incluirla por error es un bug común al reutilizar esta misma función para ambos casos.
\subsection{Punto Dentro de un Polígono (Ray Casting)}

\textbf{Idea:} para determinar si un punto $P$ está dentro de un polígono simple arbitrario (convexo o no), \textit{ray casting} traza un rayo desde $P$ hacia el infinito en una dirección fija (típicamente horizontal) y cuenta cuántas aristas del polígono cruza ese rayo. Si el número de cruces es \textbf{impar}, $P$ está dentro; si es \textbf{par}, está fuera —la intuición es que cada vez que el rayo atraviesa el borde, cruza de "afuera" a "adentro" o viceversa, así que empezando afuera (en el infinito), un número impar de cruces termina adentro. Esta idea, conocida como el teorema de la curva de Jordan aplicado a polígonos, funciona sin importar la orientación del polígono ni si es convexo.

\textbf{¿Por qué usarlo?} A diferencia de \textit{Punto Dentro de un Triángulo} (tres pruebas de orientación) o de \textit{Punto Dentro de un Polígono Convexo} (búsqueda binaria angular, más adelante), ray casting es el único método de esta parte que funciona para polígonos \textbf{no convexos} (cóncavos, con "entrantes"). Su costo es más alto —$O(n)$ en vez de $O(1)$ o $O(\log n)$— porque no puede aprovechar ninguna estructura especial: cada arista debe revisarse individualmente sin importar la forma del polígono.

\textbf{Casos de uso:}
\begin{itemize}
	\item Verificar pertenencia de un punto a un polígono arbitrario, incluyendo formas cóncavas o con auto-similaridad compleja.
	\item Point-in-polygon para regiones geográficas o mapas con contornos irregulares.
	\item Rasterización de polígonos (determinar qué píxeles de una grilla caen dentro de una figura).
	\item Validación general cuando no se puede garantizar que el polígono sea convexo.
\end{itemize}

\textbf{Complejidad:} $O(n)$ —se revisa cada una de las $n$ aristas del polígono exactamente una vez, con una comprobación $O(1)$ por arista.

\begin{lstlisting}[style=cpp]
	// CREATE/READ: Punto dentro de un poligono simple (convexo o no) mediante
	// ray casting (conteo de cruces con un rayo horizontal)
	struct PuntoDentroPoligono {
		struct Punto { double x, y; };
		
		// READ: true si el punto p esta dentro del poligono (borde EXCLUIDO,
		// ver Nota); funciona para poligonos convexos y concavos
		static bool dentro(const vector<Punto>& poligono, const Punto& p) {
			int n = (int)poligono.size();
			bool adentro = false;
			for (int i = 0, j = n - 1; i < n; j = i++) {
				const Punto& a = poligono[i];
				const Punto& b = poligono[j];
				// revisa si el rayo horizontal desde p hacia +x cruza la arista (a,b)
				bool cruzaEnY = (a.y > p.y) != (b.y > p.y);
				if (cruzaEnY) {
					double xCruce = a.x + (p.y - a.y) * (b.x - a.x) / (b.y - a.y);
					if (p.x < xCruce) adentro = !adentro;
				}
			}
			return adentro;
		} // O(n)
		
		// READ: true si el punto p esta exactamente SOBRE el borde del poligono
		// (sobre alguna arista); combinar con dentro() para "dentro o sobre el borde"
		static bool sobreBorde(const vector<Punto>& poligono, const Punto& p) {
			int n = (int)poligono.size();
			const double EPS = 1e-9;
			for (int i = 0; i < n; i++) {
				const Punto& a = poligono[i];
				const Punto& b = poligono[(i + 1) % n];
				double cruz = (b.x - a.x) * (p.y - a.y) - (b.y - a.y) * (p.x - a.x);
				if (fabs(cruz) > EPS) continue; // no colineal con esta arista
				bool dentroX = min(a.x, b.x) - EPS <= p.x && p.x <= max(a.x, b.x) + EPS;
				bool dentroY = min(a.y, b.y) - EPS <= p.y && p.y <= max(a.y, b.y) + EPS;
				if (dentroX && dentroY) return true;
			}
			return false;
		} // O(n)
		
		// ------------------------------------------------------------
		// MODIFICACIONES Y COMENTARIOS CON LOS USOS
		// ------------------------------------------------------------
		// 1. "Dentro o sobre el borde" (inclusivo): dentro(poligono, p) ||
		//    sobreBorde(poligono, p) -- necesario porque dentro() por si solo
		//    da resultado indefinido/inconsistente cuando p cae exactamente
		//    sobre una arista (el rayo puede o no contarla como cruce segun
		//    redondeo de punto flotante).
		// 2. Coordenadas enteras: xCruce y la comparacion p.x < xCruce se
		//    pueden reescribir con aritmetica de enteros multiplicando en cruz,
		//    evitando la division y el error de punto flotante por completo.
		// 3. Poligono con hoyos (contorno exterior + contornos interiores en
		//    sentido opuesto): correr dentro() sobre el contorno exterior Y
		//    sobre cada hoyo; el punto esta en la region valida solo si esta
		//    dentro del exterior y fuera de TODOS los hoyos.
	};
\end{lstlisting}

\textbf{Ejemplo de funcionamiento:}
\begin{lstlisting}[style=cpp]
	vector<PuntoDentroPoligono::Punto> estrella = {
		{0,4}, {1,1}, {4,0}, {1,-1}, {0,-4}, {-1,-1}, {-4,0}, {-1,1}
	};
	PuntoDentroPoligono::Punto p{0, 0};
	
	bool resultado = PuntoDentroPoligono::dentro(estrella, p);
	// resultado = true
	// (p=(0,0) esta en el centro de una figura concava tipo estrella; ray
	//  casting funciona correctamente aqui donde las tres pruebas de
	//  orientacion de Punto Dentro de un Triangulo no aplicarian)
\end{lstlisting}

\textbf{Nota:} el resultado de \texttt{dentro()} para un punto exactamente sobre el borde es indefinido por construcción del algoritmo (depende de redondeo). Si el problema distingue "estrictamente dentro" de "sobre el borde" de "fuera", usa \texttt{sobreBorde()} por separado antes de confiar en \texttt{dentro()}, tal como se muestra en la modificación 1 dentro del código.

\textbf{Regla práctica:} usa ray casting como la herramienta general por defecto para point-in-polygon —es la única de esta parte que funciona sin importar si el polígono es convexo. Si tienes garantizado que el polígono es convexo, \textit{Punto Dentro de un Polígono Convexo} con búsqueda binaria es $O(\log n)$ y por lo tanto preferible cuando se hacen muchas consultas sobre el mismo polígono.

\subsection{Punto Dentro de un Polígono Convexo (Binary Search)}

\textbf{Idea:} cuando el polígono es convexo y está dado en orden CCW (la convención fijada en \textit{Orientación de Tres Puntos}), se puede aprovechar esa estructura para responder en $O(\log n)$ en vez de $O(n)$: se fija el vértice $0$ como pivote, y los vértices $1, \dots, n-1$ quedan ordenados angularmente alrededor de él por ser convexo. Una búsqueda binaria encuentra el "abanico" (los dos vértices consecutivos $P_i, P_{i+1}$) entre los que cae el ángulo de $P$ respecto a $P_0$, y basta con una prueba de \textit{Punto Dentro de un Triángulo} sobre el triángulo $(P_0, P_i, P_{i+1})$ para confirmar si $P$ está dentro.

\textbf{¿Por qué usarlo?} Frente a ray casting ($O(n)$, funciona para cualquier polígono), esta técnica explota la convexidad para bajar a $O(\log n)$ por consulta —una mejora significativa cuando se hacen muchas consultas sobre el mismo polígono fijo (por ejemplo, $q$ consultas dan $O(n + q\log n)$ en vez de $O(nq)$). La condición es estricta: el polígono debe ser convexo y estar en el orden CCW consistente con el resto de la guía; si cualquiera de esas dos condiciones falla, el resultado es incorrecto sin ningún aviso.

\textbf{Casos de uso:}
\begin{itemize}
	\item Muchas consultas de pertenencia sobre el mismo polígono convexo fijo (por ejemplo, la salida de \textit{Envolvente Convexa}).
	\item Verificación rápida de colisión punto-polígono en aplicaciones con restricciones de tiempo estrictas.
	\item Localización de un punto dentro de una región convexa en problemas de particionamiento del plano.
	\item Como subrutina dentro de algoritmos más grandes que requieren muchas pruebas de pertenencia (por ejemplo, intersección de semiplanos).
\end{itemize}

\textbf{Complejidad:} $O(\log n)$ por consulta tras un preprocesamiento $O(1)$ (no requiere ordenar nada adicional, ya que el polígono convexo en CCW ya viene angularmente ordenado respecto a cualquiera de sus vértices).

\begin{lstlisting}[style=cpp]
	// CREATE/READ: Punto dentro de un poligono CONVEXO mediante busqueda binaria
	// angular; REQUIERE que el poligono este en orden CCW (ver Nota)
	struct PuntoDentroPoligonoConvexo {
		struct Punto { double x, y; };
		
		// READ: producto cruz de (b-a) x (c-a)
		static double cruz(const Punto& a, const Punto& b, const Punto& c) {
			return (b.x - a.x) * (c.y - a.y) - (b.y - a.y) * (c.x - a.x);
		} // O(1)
		
		// READ: true si p esta dentro (o sobre el borde) del poligono convexo,
		// dado en orden CCW; usa busqueda binaria sobre el abanico desde poligono[0]
		static bool dentro(const vector<Punto>& poligono, const Punto& p) {
			int n = (int)poligono.size();
			const Punto& p0 = poligono[0];
			const double EPS = 1e-9;
			
			// p debe estar entre los rayos p0->poligono[1] y p0->poligono[n-1]
			double cruzInicial = cruz(p0, poligono[1], p);
			double cruzFinal = cruz(p0, poligono[n - 1], p);
			if (cruzInicial < -EPS || cruzFinal > EPS) return false; // fuera del abanico
			
			// busqueda binaria del indice i tal que p cae en el triangulo
			// (p0, poligono[i], poligono[i+1])
			int lo = 1, hi = n - 1;
			while (hi - lo > 1) {
				int medio = (lo + hi) / 2;
				if (cruz(p0, poligono[medio], p) >= -EPS) lo = medio;
				else hi = medio;
			}
			
			// confirma que p esta dentro del triangulo final (p0, poligono[lo], poligono[hi])
			double d1 = cruz(p0, poligono[lo], p);
			double d2 = cruz(poligono[lo], poligono[hi], p);
			double d3 = cruz(poligono[hi], p0, p);
			return d1 >= -EPS && d2 >= -EPS && d3 >= -EPS;
		} // O(log n)
		
		// ------------------------------------------------------------
		// MODIFICACIONES Y COMENTARIOS CON LOS USOS
		// ------------------------------------------------------------
		// 1. Poligono en orden CW (horario): invertir el orden de los vertices
		//    (poligono.rbegin() a poligono.rend()) antes de usar, o negar todos
		//    los resultados de cruz() en el algoritmo -- NO usar directamente
		//    sin ajustar, el resultado seria incorrecto en silencio.
		// 2. Excluir el borde (estrictamente dentro): cambiar todos los ">= -EPS"
		//    por "> EPS" en las comparaciones finales y en la busqueda binaria.
		// 3. Muchas consultas sobre el mismo poligono: el preprocesamiento es
		//    O(1) (no hay que ordenar nada extra), asi que esta estructura es
		//    directamente reutilizable llamando dentro() repetidamente sin
		//    reconstruir nada.
	};
\end{lstlisting}

\textbf{Ejemplo de funcionamiento:}
\begin{lstlisting}[style=cpp]
	vector<PuntoDentroPoligonoConvexo::Punto> hexagono = {
		{2,0}, {1,2}, {-1,2}, {-2,0}, {-1,-2}, {1,-2}
	}; // hexagono convexo en orden CCW
	PuntoDentroPoligonoConvexo::Punto p{0, 0}, q{3, 3};
	
	bool dentroP = PuntoDentroPoligonoConvexo::dentro(hexagono, p);
	// dentroP = true (el centro esta dentro del hexagono)
	
	bool dentroQ = PuntoDentroPoligonoConvexo::dentro(hexagono, q);
	// dentroQ = false (q esta fuera, mas alla de todos los vertices)
\end{lstlisting}

\textbf{Nota:} este algoritmo \textbf{asume orden CCW} sin verificarlo — si el polígono que produce \textit{Envolvente Convexa} (Jarvis o Monotone Chain) no respeta esa orientación, esta búsqueda binaria falla en silencio, dando resultados incorrectos sin ningún error visible. Antes de usar esta técnica sobre la salida de un algoritmo de convex hull, confirma la orientación con \texttt{AreaPoligono::esCCW()} (ver \textit{Área de un Polígono}) y, si es necesario, invierte el orden de los vértices.

\textbf{Regla práctica:} usa esta técnica en vez de ray casting únicamente cuando el polígono sea convexo Y esté en orden CCW confirmado, y cuando vayas a hacer múltiples consultas sobre el mismo polígono —la ganancia de $O(n)$ a $O(\log n)$ solo vale la pena si el número de consultas $q$ es grande. Para una sola consulta o un polígono no convexo, usa \textit{Punto Dentro de un Polígono (Ray Casting)}.
\subsection{Teorema de Pick (Área en Malla de Puntos Enteros)}

\textbf{Idea:} para un polígono simple cuyos vértices tienen coordenadas \textbf{enteras}, el teorema de Pick relaciona el área con el conteo de puntos enteros: $A = I + \dfrac{B}{2} - 1$, donde $I$ es el número de puntos enteros estrictamente en el interior del polígono y $B$ es el número de puntos enteros sobre el borde (incluyendo los vértices). Despejando, $I = A - \dfrac{B}{2} + 1$. El número de puntos enteros sobre una arista entre dos vértices $(x_1,y_1)$ y $(x_2,y_2)$ es $\gcd(|x_2-x_1|, |y_2-y_1|)$ —incluyendo un extremo pero no el otro, para no contar los vértices dos veces al sumar todas las aristas.

\textbf{¿Por qué usarlo?} Es el puente entre \textit{Área de un Polígono (Shoelace)} —que da el área geométrica continua— y el conteo combinatorio de puntos en una malla entera, algo que de otra forma requeriría iterar sobre una grilla completa ($O(\text{área})$, potencialmente enorme). Pick reduce ese conteo a $O(1)$ una vez que se tiene el área (Shoelace, $O(n)$) y el conteo de puntos de borde (suma de $\gcd$ por arista, también $O(n)$): la fórmula evita cualquier iteración sobre la malla misma, sin importar qué tan grande sea el polígono.

\textbf{Casos de uso:}
\begin{itemize}
	\item Contar puntos de malla estrictamente interiores a un polígono de vértices enteros, sin iterar la grilla.
	\item Contar puntos de malla sobre el borde de un polígono (perímetro "discreto").
	\item Problemas que piden el área de una figura descrita en una grilla de puntos enteros (más simple que Shoelace + gcd cuando se da $I$ y $B$ directamente).
	\item Verificación cruzada: si se conocen $I$, $B$ y el área por otro método, Pick sirve como chequeo de consistencia.
\end{itemize}

\textbf{Complejidad:} $O(n)$ —Shoelace para el área es $O(n)$, y sumar $\gcd$ sobre cada una de las $n$ aristas es $O(n \log(\max \text{coordenada}))$ por el costo del algoritmo de Euclides, dominado en la práctica por el término lineal.

\begin{lstlisting}[style=cpp]
	// CREATE/READ: Teorema de Pick: relaciona area, puntos interiores (I) y
	// puntos de borde (B) para poligonos de vertices enteros
	struct TeoremaPick {
		struct Punto { long long x, y; };
		
		// READ: area con signo * 2 mediante Shoelace en aritmetica entera
		// (evita perder precision, ya que 2*area siempre es entero)
		static long long dobleAreaConSigno(const vector<Punto>& vertices) {
			int n = (int)vertices.size();
			long long suma = 0;
			for (int i = 0; i < n; i++) {
				const Punto& actual = vertices[i];
				const Punto& siguiente = vertices[(i + 1) % n];
				suma += actual.x * siguiente.y - siguiente.x * actual.y;
			}
			return suma;
		} // O(n)
		
		// READ: numero de puntos enteros sobre el borde del poligono
		// (suma de gcd por arista, cada vertice contado una sola vez)
		static long long puntosBorde(const vector<Punto>& vertices) {
			int n = (int)vertices.size();
			long long total = 0;
			for (int i = 0; i < n; i++) {
				const Punto& a = vertices[i];
				const Punto& b = vertices[(i + 1) % n];
				total += __gcd(llabs(b.x - a.x), llabs(b.y - a.y));
			}
			return total;
		} // O(n log(max coordenada))
		
		// READ: numero de puntos enteros ESTRICTAMENTE interiores, via Pick:
		// I = A - B/2 + 1, usando doble area para mantener aritmetica entera
		static long long puntosInteriores(const vector<Punto>& vertices) {
			long long dobleArea = llabs(dobleAreaConSigno(vertices));
			long long b = puntosBorde(vertices);
			// 2*I = 2*A - B + 2 = dobleArea - b + 2
			return (dobleArea - b + 2) / 2;
		} // O(n log(max coordenada))
		
		// ------------------------------------------------------------
		// MODIFICACIONES Y COMENTARIOS CON LOS USOS
		// ------------------------------------------------------------
		// 1. Area como double a partir de I y B conocidos (Pick directo, sin
		//    Shoelace): area = I + B/2.0 - 1; util si el problema da I y B
		//    directamente y pide el area.
		// 2. Vertices NO enteros: el teorema de Pick no aplica; requiere
		//    coordenadas enteras tanto para el conteo de gcd como para que I
		//    y B tengan sentido combinatorio.
		// 3. Verificacion de que dobleArea - b es PAR antes de dividir entre 2:
		//    matematicamente siempre lo es para un poligono valido de vertices
		//    enteros, pero sirve como chequeo de cordura al depurar (si da
		//    impar, hay un bug en el calculo de area o de borde).
	};
\end{lstlisting}

\textbf{Ejemplo de funcionamiento:}
\begin{lstlisting}[style=cpp]
	vector<TeoremaPick::Punto> cuadrado = {{0,0}, {4,0}, {4,4}, {0,4}};
	
	long long b = TeoremaPick::puntosBorde(cuadrado);
	// b = 16
	// (cada lado tiene longitud 4, gcd(4,0)=4 por arista, 4 aristas: 4*4=16;
	//  cada punto de borde se cuenta una sola vez, incluyendo los 4 vertices)
	
	long long i = TeoremaPick::puntosInteriores(cuadrado);
	// i = 9
	// (area = 16, Pick: I = 16 - 16/2 + 1 = 16 - 8 + 1 = 9;
	//  corresponde a la grilla interior 3x3 de puntos (1,1) a (3,3))
\end{lstlisting}

\textbf{Nota:} este teorema aplica \textbf{únicamente} a polígonos simples con vértices de coordenadas enteras. Si el problema da coordenadas reales (no enteras), no hay forma de aplicar Pick directamente —hay que recurrir a \textit{Área de un Polígono (Shoelace)} para el área continua, sin ningún conteo de puntos de malla asociado.

\textbf{Regla práctica:} usa Pick cuando el problema combine dos de las tres cantidades (área, puntos interiores, puntos de borde) y pida la tercera, sobre un polígono de vértices enteros —evita iterar la grilla por completo, algo que sería inviable si el área es grande. Si necesitas el área pura sin relación con conteo de puntos, \textit{Shoelace} solo (sin Pick) es más directo.

\section{Envolvente Convexa}
\subsection{Algoritmo de Jarvis (Gift Wrapping)}

\textbf{Idea:} Jarvis construye la envolvente convexa "envolviendo" el conjunto de puntos como si se estirara una banda elástica alrededor de ellos: partiendo del punto más bajo (garantizado en el hull), se busca repetidamente el siguiente vértice del hull como aquel punto tal que \textbf{todos} los demás puntos quedan a la izquierda del segmento actual —es decir, el punto que hace el giro más "hacia afuera" posible respecto al vértice anterior. Esto se decide con \textit{Orientación de Tres Puntos}: para cada candidato se compara contra el mejor candidato actual, y se actualiza cuando el nuevo punto queda más a la derecha (en sentido de barrido) del candidato vigente. El proceso se repite hasta volver al punto de partida.

\textbf{¿Por qué usarlo?} Es el algoritmo de envolvente convexa más simple de razonar e implementar correctamente, porque en cada paso solo requiere una comparación local (¿quién es más extremo?) sin necesidad de ordenar todos los puntos de antemano. Su costo, sin embargo, depende del número de vértices $h$ que termina teniendo el hull, no solo de $n$: es preferible a \textit{Monotone Chain} únicamente cuando se sabe de antemano que $h$ va a ser pequeño comparado con $n$ (envolvente con pocos vértices sobre una nube de puntos grande); en el peor caso, cuando casi todos los puntos están en el hull ($h \approx n$), Jarvis degenera a $O(n^2)$, peor que el $O(n\log n)$ garantizado de Monotone Chain.

\textbf{Casos de uso:}
\begin{itemize}
	\item Construcción de la envolvente convexa cuando se espera que tenga pocos vértices.
	\item Problemas educativos o de referencia rápida donde la simplicidad de implementación importa más que la complejidad garantizada.
	\item Como paso dentro de algoritmos incrementales de geometría computacional (Quickhull usa una idea similar de forma recursiva).
	\item Verificación cruzada del resultado de \textit{Monotone Chain} en casos pequeños durante debugging.
\end{itemize}

\textbf{Complejidad:} $O(nh)$, donde $n$ es el número de puntos y $h$ el número de vértices del hull resultante: cada uno de los $h$ vértices del hull requiere una pasada de $O(n)$ para encontrar el siguiente punto más extremo. En el peor caso ($h = n$), esto es $O(n^2)$.

\begin{lstlisting}[style=cpp]
	// CREATE/READ: Envolvente convexa mediante Jarvis (Gift Wrapping);
	// retorna los vertices del hull en orden CCW
	struct EnvolventeJarvis {
		struct Punto { double x, y; };
		
		// READ: producto cruz de (b-a) x (c-a); positivo = CCW (convencion
		// fijada en Orientacion de Tres Puntos)
		static double cruz(const Punto& a, const Punto& b, const Punto& c) {
			return (b.x - a.x) * (c.y - a.y) - (b.y - a.y) * (c.x - a.x);
		} // O(1)
		
		// READ: distancia al cuadrado, usada para desempatar puntos colineales
		static double distCuadrado(const Punto& a, const Punto& b) {
			double dx = b.x - a.x, dy = b.y - a.y;
			return dx*dx + dy*dy;
		} // O(1)
		
		// READ: construye el hull convexo de un conjunto de puntos, retornado
		// en orden CCW; requiere n >= 3 para un hull no degenerado
		static vector<Punto> construir(vector<Punto> puntos) {
			int n = (int)puntos.size();
			if (n < 3) return puntos; // hull degenerado (punto o segmento)
			
			// punto inicial: el de menor y (y menor x en caso de empate),
			// garantizado en el hull
			int inicio = 0;
			for (int i = 1; i < n; i++) {
				if (puntos[i].y < puntos[inicio].y ||
				(puntos[i].y == puntos[inicio].y && puntos[i].x < puntos[inicio].x)) {
					inicio = i;
				}
			}
			
			vector<Punto> hull;
			int actual = inicio;
			do {
				hull.push_back(puntos[actual]);
				int siguiente = (actual + 1) % n; // candidato inicial arbitrario
				for (int i = 0; i < n; i++) {
					if (i == actual) continue;
					double c = cruz(puntos[actual], puntos[siguiente], puntos[i]);
					// si i queda a la derecha del candidato actual, i es mejor;
					// si son colineales, se prefiere el mas lejano
					if (c < 0 || (c == 0 && distCuadrado(puntos[actual], puntos[i]) >
					distCuadrado(puntos[actual], puntos[siguiente]))) {
						siguiente = i;
					}
				}
				actual = siguiente;
			} while (actual != inicio);
			
			return hull;
		} // O(n*h): h iteraciones del do-while, cada una O(n)
		
		// ------------------------------------------------------------
		// MODIFICACIONES Y COMENTARIOS CON LOS USOS
		// ------------------------------------------------------------
		// 1. Hull en sentido CW en vez de CCW: invertir el criterio "c < 0" a
		//    "c > 0" en la comparacion del candidato; NO cambia la complejidad,
		//    solo la orientacion de salida.
		// 2. Incluir puntos colineales sobre el borde del hull: el desempate por
		//    distCuadrado en el codigo actual EXCLUYE puntos colineales
		//    intermedios (se queda con el mas lejano); para incluirlos, agregar
		//    los puntos colineales entre actual y siguiente antes de avanzar.
		// 3. Puntos duplicados en la entrada: eliminarlos antes de llamar
		//    construir(), ya que pueden causar que el bucle do-while no termine
		//    correctamente si el punto inicial se duplica.
	};
\end{lstlisting}

\textbf{Ejemplo de funcionamiento:}
\begin{lstlisting}[style=cpp]
	vector<EnvolventeJarvis::Punto> puntos = {
		{0,0}, {4,0}, {4,4}, {0,4}, {2,2} // (2,2) esta estrictamente adentro
	};
	
	auto hull = EnvolventeJarvis::construir(puntos);
	// hull = {(0,0), (4,0), (4,4), (0,4)}
	// (el punto interior (2,2) queda excluido; los cuatro vertices del
	//  cuadrado forman el hull en orden CCW, empezando por el de menor y)
\end{lstlisting}

\textbf{Regla práctica:} usa Jarvis cuando sepas (o sospeches fuertemente) que el hull tendrá pocos vértices frente a $n$; si no hay garantía sobre el tamaño del hull, o $n$ es grande, usa \textit{Monotone Chain} en su lugar, que garantiza $O(n\log n)$ sin importar cuántos vértices termine teniendo el hull.

\subsection{Monotone Chain (Andrew)}

\textbf{Idea:} en vez de "envolver" iterativamente como Jarvis, Monotone Chain ordena todos los puntos por $x$ (desempatando por $y$) y construye el hull en dos pasadas lineales: una \textbf{cadena inferior} (\textit{lower hull}), recorriendo los puntos ordenados de izquierda a derecha y manteniendo solo los que forman giros a la derecha (o eliminando los que forman giros a la izquierda, usando \textit{Orientación de Tres Puntos} sobre una pila), y una \textbf{cadena superior} (\textit{upper hull}) recorriendo en sentido inverso con el mismo criterio. Concatenar ambas cadenas (sin duplicar los extremos) da el hull completo en orden CCW.

\textbf{¿Por qué usarlo?} Garantiza $O(n\log n)$ en el peor caso, dominado por el ordenamiento inicial —a diferencia de Jarvis, cuyo costo depende de $h$ y puede degenerar a $O(n^2)$. El mecanismo de "pila que elimina giros incorrectos" es el mismo patrón usado en muchos algoritmos de geometría computacional con barrido (similar en espíritu al manejo de estado del \textit{Line Sweep}), y es la opción por defecto en competencia cuando no hay garantía sobre el tamaño del hull resultante.

\textbf{Casos de uso:}
\begin{itemize}
	\item Construcción de la envolvente convexa como caso general, sin garantías sobre $h$.
	\item Preprocesamiento antes de \textit{Punto Dentro de un Polígono Convexo}, que requiere el hull en orden CCW.
	\item Problemas donde $n$ es grande y se necesita la complejidad garantizada $O(n\log n)$.
	\item Base para algoritmos que requieren el hull como paso intermedio (diámetro del conjunto de puntos, rotating calipers).
\end{itemize}

\textbf{Complejidad:} $O(n\log n)$, dominado por el ordenamiento inicial de los $n$ puntos; cada una de las dos cadenas (inferior y superior) se construye en $O(n)$ porque cada punto se inserta y se elimina de la pila a lo sumo una vez.

\begin{lstlisting}[style=cpp]
	// CREATE/READ: Envolvente convexa mediante Monotone Chain (Andrew);
	// retorna los vertices del hull en orden CCW, sin puntos colineales de sobra
	struct EnvolventeMonotoneChain {
		struct Punto { double x, y; };
		
		// READ: producto cruz de (b-a) x (c-a); positivo = CCW
		static double cruz(const Punto& a, const Punto& b, const Punto& c) {
			return (b.x - a.x) * (c.y - a.y) - (b.y - a.y) * (c.x - a.x);
		} // O(1)
		
		// READ: construye el hull convexo en orden CCW; requiere n >= 3 para
		// un hull no degenerado
		static vector<Punto> construir(vector<Punto> puntos) {
			int n = (int)puntos.size();
			if (n < 3) return puntos;
			
			sort(puntos.begin(), puntos.end(), [](const Punto& a, const Punto& b) {
				if (a.x != b.x) return a.x < b.x;
				return a.y < b.y;
			});
			
			vector<Punto> hull(2 * n);
			int k = 0;
			
			// cadena inferior: recorre de izquierda a derecha
			for (int i = 0; i < n; i++) {
				while (k >= 2 && cruz(hull[k-2], hull[k-1], puntos[i]) <= 0) k--;
				hull[k++] = puntos[i];
			}
			
			// cadena superior: recorre de derecha a izquierda, continuando
			// desde donde quedo la cadena inferior
			int limiteInferior = k + 1;
			for (int i = n - 2; i >= 0; i--) {
				while (k >= limiteInferior && cruz(hull[k-2], hull[k-1], puntos[i]) <= 0) k--;
				hull[k++] = puntos[i];
			}
			
			hull.resize(k - 1); // el ultimo punto duplica el primero, se descarta
			return hull;
		} // O(n log n): dominado por el sort; ambas cadenas son O(n)
		
		// ------------------------------------------------------------
		// MODIFICACIONES Y COMENTARIOS CON LOS USOS
		// ------------------------------------------------------------
		// 1. Incluir puntos colineales sobre el borde del hull: cambiar la
		//    condicion "<= 0" por "< 0" en ambos while, de modo que los giros
		//    exactamente colineales NO se descarten de la pila.
		// 2. Hull en sentido CW en vez de CCW: invertir "<= 0" a ">= 0" en
		//    ambos while.
		// 3. Puntos duplicados en la entrada: eliminarlos antes de ordenar
		//    (sort + unique) para evitar comportamiento inesperado en el
		//    desempate de cruz() == 0 con puntos identicos.
	};
\end{lstlisting}

\textbf{Ejemplo de funcionamiento:}
\begin{lstlisting}[style=cpp]
	vector<EnvolventeMonotoneChain::Punto> puntos = {
		{0,0}, {4,0}, {4,4}, {0,4}, {2,2} // (2,2) esta estrictamente adentro
	};
	
	auto hull = EnvolventeMonotoneChain::construir(puntos);
	// hull = {(0,0), (4,0), (4,4), (0,4)}
	// (mismo resultado que Jarvis para este conjunto de puntos, en orden CCW;
	//  el punto interior (2,2) queda excluido por la condicion de giro)
\end{lstlisting}

\textbf{Nota:} tanto Jarvis como Monotone Chain producen el hull en orden \textbf{CCW} con la convención fijada en esta guía, así que el resultado de cualquiera de los dos es directamente compatible con \textit{Punto Dentro de un Polígono Convexo}, que requiere esa orientación. Si se necesita orden CW, ver la modificación 2 en cualquiera de los dos structs.

\textbf{Regla práctica:} usa Monotone Chain como opción por defecto en competencia, salvo que tengas una razón específica para preferir Jarvis (hull con muy pocos vértices conocido de antemano, o simplicidad de implementación en un contexto donde $n$ es pequeño). Su $O(n\log n)$ garantizado lo hace la opción más segura cuando no hay información previa sobre el tamaño del hull.
\subsection{Corte de un Polígono Convexo por una Línea}

\textbf{Idea:} dado un polígono convexo (por ejemplo, la salida de \textit{Jarvis} o \textit{Monotone Chain}) y una línea definida por dos puntos, esta operación produce el subpolígono resultante de quedarse solo con la parte del polígono que cae de un lado específico de la línea (el lado izquierdo, según la orientación de la línea). El algoritmo recorre los vértices en orden y usa \textit{Orientación de Tres Puntos} para clasificar cada uno como dentro o fuera del semiplano deseado: cuando dos vértices consecutivos están en lados opuestos, se calcula el punto de intersección exacto de la arista con la línea de corte y se inserta en el resultado; los vértices que están del lado válido se conservan tal cual. Este es un caso particular del algoritmo de recorte de Sutherland-Hodgman, especializado a una sola línea de corte en vez de un polígono de recorte completo.

\textbf{¿Por qué usarlo?} A diferencia de \textit{Intersección de Segmentos}, que solo detecta si dos segmentos se cruzan, aquí el objetivo es reconstruir una figura completa (el subpolígono resultante), lo que requiere combinar la prueba de orientación con la inserción cuidadosa de nuevos vértices en la intersección. Explota que el polígono de entrada es \textbf{convexo}: eso garantiza que el resultado del corte también es convexo (un semiplano intersectado con un convexo siempre da un convexo), lo cual no se cumpliría si el polígono de entrada fuera cóncavo —ahí el resultado podría ser múltiples polígonos separados, un caso mucho más complejo no cubierto aquí.

\textbf{Casos de uso:}
\begin{itemize}
	\item Calcular el área de un polígono convexo que queda a un lado de una recta (por ejemplo, "¿qué fracción del terreno queda al norte de esta carretera recta?").
	\item Recorte de figuras convexas contra los bordes de una ventana o región rectangular (aplicando esta operación cuatro veces, una por cada borde).
	\item Como subrutina en \textit{intersección de semiplanos} (intersectar un polígono convexo contra muchas líneas sucesivamente).
	\item Simulación de corte físico de una región convexa por una trayectoria recta.
\end{itemize}

\textbf{Complejidad:} $O(n)$ —un recorrido lineal por los $n$ vértices del polígono, con una prueba de orientación $O(1)$ y, en el peor caso, un cálculo de intersección $O(1)$ por cada arista.

\begin{lstlisting}[style=cpp]
	// CREATE/READ: Corta un poligono convexo por una linea, quedandose con
	// la parte del lado IZQUIERDO de la linea (segun su orientacion a->b)
	struct CortePoligonoConvexo {
		struct Punto { double x, y; };
		
		// READ: producto cruz de (b-a) x (c-a); positivo = c esta a la
		// izquierda de a->b (convencion CCW de Orientacion de Tres Puntos)
		static double cruz(const Punto& a, const Punto& b, const Punto& c) {
			return (b.x - a.x) * (c.y - a.y) - (b.y - a.y) * (c.x - a.x);
		} // O(1)
		
		// READ: interseccion entre la arista (p, q) y la linea (a, b),
		// asumiendo que p y q estan en lados opuestos de la linea
		static Punto interseccion(const Punto& p, const Punto& q, const Punto& a, const Punto& b) {
			double d1 = cruz(a, b, p);
			double d2 = cruz(a, b, q);
			double t = d1 / (d1 - d2);
			return Punto{p.x + t * (q.x - p.x), p.y + t * (q.y - p.y)};
		} // O(1)
		
		// READ: corta el poligono convexo, quedandose con el lado izquierdo
		// de la linea a->b; retorna vector vacio si el poligono queda
		// completamente del lado derecho
		static vector<Punto> cortar(const vector<Punto>& poligono, const Punto& a, const Punto& b) {
			int n = (int)poligono.size();
			vector<Punto> resultado;
			const double EPS = 1e-9;
			
			for (int i = 0; i < n; i++) {
				const Punto& actual = poligono[i];
				const Punto& siguiente = poligono[(i + 1) % n];
				double dActual = cruz(a, b, actual);
				double dSiguiente = cruz(a, b, siguiente);
				
				if (dActual >= -EPS) resultado.push_back(actual); // actual esta dentro (o sobre la linea)
				// si actual y siguiente estan en lados opuestos, agregar el punto de corte
				if ((dActual > EPS && dSiguiente < -EPS) || (dActual < -EPS && dSiguiente > EPS)) {
					resultado.push_back(interseccion(actual, siguiente, a, b));
				}
			}
			return resultado;
		} // O(n)
		
		// ------------------------------------------------------------
		// MODIFICACIONES Y COMENTARIOS CON LOS USOS
		// ------------------------------------------------------------
		// 1. Quedarse con el lado DERECHO en vez del izquierdo: invertir el
		//    criterio "dActual >= -EPS" a "dActual <= EPS" (y ajustar el signo
		//    en la deteccion de cruce si se desea mayor claridad).
		// 2. Area del subpoligono resultante: aplicar directamente Shoelace
		//    (ver Area de un Poligono) sobre el vector retornado por cortar().
		// 3. Recorte contra una ventana rectangular (Sutherland-Hodgman
		//    completo): aplicar cortar() cuatro veces sucesivas, una por cada
		//    borde del rectangulo (orientados consistentemente en CCW), usando
		//    el resultado de cada corte como entrada del siguiente.
	};
\end{lstlisting}

\textbf{Ejemplo de funcionamiento:}
\begin{lstlisting}[style=cpp]
	vector<CortePoligonoConvexo::Punto> cuadrado = {{0,0}, {4,0}, {4,4}, {0,4}};
	CortePoligonoConvexo::Punto a{2,-1}, b{2,5}; // linea vertical x=2
	
	auto resultado = CortePoligonoConvexo::cortar(cuadrado, a, b);
	// resultado = {(2,0), (4,0), (4,4), (2,4)}
	// (se conserva el lado izquierdo de la linea a->b, que en este caso
	//  apunta hacia +y, dejando el lado con x >= 2; el cuadrado original de
	//  area 16 queda cortado a un rectangulo de area 8)
\end{lstlisting}

\textbf{Nota:} el lado que se conserva depende de la \textbf{orientación} de la línea de corte ($a \to b$), no solo de su posición geométrica —invertir $a$ y $b$ invierte qué mitad del polígono se conserva. Esto es consistente con la convención de \textit{Orientación de Tres Puntos}: positivo significa "a la izquierda de $a\to b$", y es responsabilidad de quien llama orientar la línea según el lado que quiere conservar.

\textbf{Regla práctica:} usa esta operación cuando necesites el subpolígono resultante de un corte, no solo si dos figuras se cruzan —para eso te alcanza \textit{Intersección de Segmentos}. Encadenar varias llamadas a \texttt{cortar()} (una por cada línea) es la forma directa de resolver intersección de semiplanos o recorte contra una ventana convexa arbitraria, siempre que el polígono de entrada sea convexo.

\section{Ad-hoc y Combinatoria Geométrica}
\subsection{Perímetro Usando Bloques (Figuras en Grilla)}

\textbf{Idea:} dada una figura formada por bloques cuadrados unitarios sobre una grilla (por ejemplo, una matriz de "*" y "." indicando qué celdas están ocupadas), el perímetro no requiere geometría continua: cada bloque ocupado contribuye hasta 4 unidades de perímetro (una por cada uno de sus cuatro lados), pero cada vez que dos bloques ocupados son vecinos (comparten un lado), ese lado compartido deja de ser borde exterior y se resta \textbf{dos veces} —una unidad de cada bloque—. La fórmula resultante es $P = 4B - 2A$, donde $B$ es el número total de bloques ocupados y $A$ es el número de pares de bloques ocupados adyacentes (compartiendo un lado, no solo una esquina).

\textbf{¿Por qué usarlo?} A diferencia de \textit{Perímetro de un Polígono}, que requiere las coordenadas explícitas del contorno, aquí la entrada es una grilla discreta y no hace falta reconstruir el contorno geométrico en absoluto —basta con contar bloques y pares adyacentes, una operación puramente combinatoria sobre la matriz. Intentar resolverlo trazando el contorno real (como si fuera un polígono) sería mucho más laborioso que este conteo directo.

\textbf{Casos de uso:}
\begin{itemize}
	\item Calcular el perímetro de una figura compuesta por celdas de una grilla (islas en un mapa, piezas de Tetris, regiones pintadas).
	\item Problemas de "área de una isla" en matrices binarias, donde también se pide el perímetro de cada componente.
	\item Verificación rápida de cuántas aristas de borde tiene una región antes de un recorrido BFS/DFS más complejo.
	\item Cálculo de perímetro total de múltiples figuras separadas sobre la misma grilla (sumando por componente conexa).
\end{itemize}

\textbf{Complejidad:} $O(NM)$ para una grilla de $N \times M$ —se recorre cada celda una vez, y por cada bloque ocupado se revisan sus hasta 4 vecinos, cada revisión $O(1)$.

\begin{lstlisting}[style=cpp]
	// CREATE/READ: Perimetro de una figura formada por bloques en una grilla
	struct PerimetroBloques {
		vector<string> grilla;
		int filas, columnas;
		
		// CREATE: recibe la grilla como vector de strings, '*' = ocupado,
		// '.' = vacio
		PerimetroBloques(const vector<string>& g) : grilla(g) {
			filas = (int)grilla.size();
			columnas = filas > 0 ? (int)grilla[0].size() : 0;
		} // O(1), no copia mas alla de la referencia inicial
		
		// READ: true si la celda (fila, col) esta dentro de la grilla Y ocupada
		bool estaOcupada(int fila, int col) const {
			if (fila < 0 || fila >= filas || col < 0 || col >= columnas) return false;
			return grilla[fila][col] == '*';
		} // O(1)
		
		// READ: perimetro total de todos los bloques ocupados en la grilla
		long long calcularPerimetro() const {
			long long perimetro = 0;
			int df[4] = {-1, 1, 0, 0};
			int dc[4] = {0, 0, -1, 1};
			for (int f = 0; f < filas; f++) {
				for (int c = 0; c < columnas; c++) {
					if (!estaOcupada(f, c)) continue;
					for (int dir = 0; dir < 4; dir++) {
						// cada lado del bloque (f,c) cuenta como borde SOLO si
						// el vecino en esa direccion no esta ocupado (o esta
						// fuera de la grilla)
						if (!estaOcupada(f + df[dir], c + dc[dir])) perimetro++;
					}
				}
			}
			return perimetro;
		} // O(N*M)
		
		// ------------------------------------------------------------
		// MODIFICACIONES Y COMENTARIOS CON LOS USOS
		// ------------------------------------------------------------
		// 1. Formula cerrada equivalente 4B - 2A: contar B (bloques ocupados) y
		//    A (pares de vecinos ocupados, cada par contado una sola vez) por
		//    separado, y retornar 4*B - 2*A; da el mismo resultado que revisar
		//    los 4 vecinos por celda, solo que evita revisar cada adyacencia
		//    dos veces (una desde cada bloque).
		// 2. Perimetro por componente conexa (varias figuras separadas en la
		//    misma grilla): correr un BFS/DFS para etiquetar componentes, y
		//    aplicar la misma logica de conteo de vecinos SOLO dentro de cada
		//    componente etiquetada.
		// 3. Grilla con bloques de tamano k (no unitarios): multiplicar el
		//    resultado final por k, ya que cada lado unitario en la formula
		//    corresponde a un lado de longitud k en la figura real.
	};
\end{lstlisting}

\textbf{Ejemplo de funcionamiento:}
\begin{lstlisting}[style=cpp]
	vector<string> grilla = {
		"**",
		".*"
	};
	PerimetroBloques p(grilla);
	
	long long resultado = p.calcularPerimetro();
	// resultado = 8
	// (3 bloques ocupados: (0,0),(0,1),(1,1); hay 2 pares adyacentes:
	//  (0,0)-(0,1) y (0,1)-(1,1); formula: 4*3 - 2*2 = 12 - 4 = 8)
\end{lstlisting}

\textbf{Regla práctica:} usa esta técnica de conteo directo (revisión de vecinos, o la fórmula equivalente $4B-2A$) siempre que la figura venga dada como una grilla de celdas ocupadas —es $O(NM)$ y evita por completo la necesidad de reconstruir un contorno poligonal explícito, que sería innecesariamente complejo para este tipo de entrada.

\subsection{Número de Cuadrados en un Rectángulo N×M}

\textbf{Idea:} contar cuántos cuadrados de \textbf{cualquier tamaño} caben alineados a la grilla dentro de un rectángulo de $N \times M$ celdas (asumiendo $N \leq M$ sin pérdida de generalidad) se resuelve sumando, para cada tamaño de lado $k$ desde $1$ hasta $N$, la cantidad de posiciones donde un cuadrado de lado $k$ cabe: hay $(N-k+1)$ posiciones verticales y $(M-k+1)$ posiciones horizontales, así que la cantidad de cuadrados de lado $k$ es $(N-k+1)(M-k+1)$. Sumando sobre todos los $k$ posibles se obtiene el total. Cuando $N=M$ (rectángulo cuadrado), esta suma se reduce a la fórmula cerrada conocida $\sum_{k=1}^{N} k^2 = \dfrac{N(N+1)(2N+1)}{6}$.

\textbf{¿Por qué usarlo?} Es un problema de conteo combinatorio puro —no requiere ningún cómputo geométrico continuo (sin producto cruz, sin coordenadas reales)—, pero encaja en esta parte porque describe cuántas figuras geométricas caben dentro de otra, un tipo de pregunta recurrente en problemas ad-hoc de la sección de geometría. La suma directa es $O(N)$ (con $N\le M$), lo cual ya es suficientemente rápido para las restricciones típicas de competencia; no hace falta la fórmula cerrada salvo que $N$ sea extremadamente grande y se necesite $O(1)$.

\textbf{Casos de uso:}
\begin{itemize}
	\item Contar cuadrados de cualquier tamaño en un tablero rectangular (variante del clásico "cuántos cuadrados hay en un tablero de ajedrez").
	\item Problemas de conteo combinatorio sobre grillas donde la figura buscada es un cuadrado alineado a los ejes.
	\item Verificación o sanity-check antes de resolver la variante más general de \textit{Número de Rectángulos en una Grilla N×M}.
	\item Subrutina en problemas que piden distribuir o colocar la mayor cantidad de cuadrados posibles sin superposición explícita (aquí se permite solapamiento, solo se cuenta cuántos "caben" como subcuadrícula).
\end{itemize}

\textbf{Complejidad:} $O(\min(N,M))$ sumando directamente sobre los tamaños de lado posibles; $O(1)$ si se usa la fórmula cerrada para el caso $N=M$.

\begin{lstlisting}[style=cpp]
	// CREATE/READ: Numero de cuadrados de cualquier tamaño en un rectangulo N x M
	struct CuadradosEnRectangulo {
		// READ: cuenta todos los cuadrados (de lado 1 hasta min(N,M)) que caben
		// alineados a la grilla dentro de un rectangulo de N x M celdas
		static long long contar(long long n, long long m) {
			long long menor = min(n, m);
			long long total = 0;
			for (long long k = 1; k <= menor; k++) {
				total += (n - k + 1) * (m - k + 1);
			}
			return total;
		} // O(min(N,M))
		
		// READ: caso particular N == M (rectangulo cuadrado), formula cerrada
		// sum_{k=1}^{n} k^2 = n(n+1)(2n+1)/6
		static long long contarCuadrado(long long n) {
			return n * (n + 1) * (2 * n + 1) / 6;
		} // O(1)
		
		// ------------------------------------------------------------
		// MODIFICACIONES Y COMENTARIOS CON LOS USOS
		// ------------------------------------------------------------
		// 1. Verificar overflow: para N, M grandes (~1e5 o mas), el resultado
		//    puede exceder long long dependiendo de las restricciones del
		//    problema; revisar el enunciado antes de asumir que long long basta.
		// 2. Solo cuadrados de tamaño EXACTO k (no todos los tamaños): basta con
		//    un solo termino (n-k+1)*(m-k+1), sin sumar sobre todos los k.
		// 3. Formula cerrada general para N != M (no solo N == M): se puede
		//    derivar expandiendo la suma en terminos de sumas de k y k^2
		//    conocidas, pero en la practica la suma O(min(N,M)) ya es
		//    suficientemente rapida en casi todas las restricciones de CP.
	};
\end{lstlisting}

\textbf{Ejemplo de funcionamiento:}
\begin{lstlisting}[style=cpp]
	long long resultado = CuadradosEnRectangulo::contar(2, 3);
	// resultado = 8
	// (cuadrados de lado 1: 2*3 = 6 posiciones; cuadrados de lado 2:
	//  (2-2+1)*(3-2+1) = 1*2 = 2 posiciones; total 6+2 = 8)
\end{lstlisting}

\textbf{Regla práctica:} usa la suma directa $O(\min(N,M))$ como opción general y segura; recurre a la fórmula cerrada solo en el caso particular $N=M$ y cuando $N$ sea tan grande que incluso $O(N)$ sea inviable —en la mayoría de problemas de competencia, la suma directa ya es suficientemente rápida y más fácil de verificar que la fórmula cerrada expandida.
\subsection{Número de Rectángulos en una Grilla N×M}

\textbf{Idea:} generaliza \textit{Número de Cuadrados en un Rectángulo N×M} de cuadrados a rectángulos de cualquier ancho y alto: un rectángulo alineado a la grilla queda determinado eligiendo dos líneas verticales distintas (de las $N+1$ líneas que delimitan las $N$ filas) y dos líneas horizontales distintas (de las $M+1$ que delimitan las $M$ columnas) —cada par de líneas verticales junto con cada par de líneas horizontales define exactamente un rectángulo. El conteo es entonces puramente combinatorio: $\binom{N+1}{2} \cdot \binom{M+1}{2}$, el producto de "formas de elegir 2 líneas verticales" por "formas de elegir 2 líneas horizontales".

\textbf{¿Por qué usarlo?} A diferencia del conteo de cuadrados —que necesita sumar sobre cada tamaño de lado posible porque un cuadrado fuerza ancho = alto—, contar rectángulos no tiene esa restricción: ancho y alto son independientes, así que el problema se separa limpiamente en "elegir extensión horizontal" y "elegir extensión vertical", cada una una combinación simple $\binom{\cdot}{2}$. Esto da una fórmula $O(1)$ directa, sin necesidad de ninguna suma —más simple que el caso de cuadrados, donde la restricción ancho=alto obliga a iterar.

\textbf{Casos de uso:}
\begin{itemize}
	\item Contar todos los rectángulos alineados a los ejes que caben en una grilla de $N \times M$ celdas.
	\item Variante del clásico "cuántos rectángulos hay en un tablero de ajedrez" (caso particular $N=M=8$).
	\item Verificación cruzada: sumando el conteo de rectángulos de cada ancho×alto específico debe coincidir con esta fórmula cerrada.
	\item Subrutina en problemas que combinan conteo de sub-rectángulos con otra condición adicional (por ejemplo, sub-rectángulos que contienen una celda marcada).
\end{itemize}

\textbf{Complejidad:} $O(1)$ —dos evaluaciones de $\binom{\cdot}{2} = \dfrac{k(k-1)}{2}$ y una multiplicación.

\begin{lstlisting}[style=cpp]
	// CREATE/READ: Numero de rectangulos alineados a la grilla en un
	// rectangulo de N x M celdas
	struct RectangulosEnGrilla {
		// READ: combinaciones de k elementos tomados de 2 en 2
		static long long combinatoria2(long long k) {
			return k * (k - 1) / 2;
		} // O(1)
		
		// READ: cuenta todos los rectangulos (de cualquier ancho y alto) que
		// caben alineados a la grilla en un rectangulo de N x M celdas
		static long long contar(long long n, long long m) {
			// N celdas dan N+1 lineas verticales delimitadoras; analogo para M
			return combinatoria2(n + 1) * combinatoria2(m + 1);
		} // O(1)
		
		// ------------------------------------------------------------
		// MODIFICACIONES Y COMENTARIOS CON LOS USOS
		// ------------------------------------------------------------
		// 1. Verificacion cruzada con Numero de Cuadrados: sumar
		//    (n-k+1)*(m-k+1) sobre k=1..min(n,m) cuenta SOLO los cuadrados,
		//    un subconjunto estricto de lo que cuenta esta formula (todos los
		//    rectangulos, incluidos los no cuadrados).
		// 2. Overflow: N+1 y M+1 pueden hacer que combinatoria2 exceda long long
		//    si N o M son extremadamente grandes (~3e9 o mas); normalmente no es
		//    un problema en CP, pero vale revisar las restricciones del enunciado.
		// 3. Rectangulos que contienen una celda fija (r0, c0): se puede
		//    contar por separado como (numero de elecciones de lineas verticales
		//    que "abarcan" la columna c0) * (analogamente para filas), en vez de
		//    la formula general sin restriccion.
	};
\end{lstlisting}

\textbf{Ejemplo de funcionamiento:}
\begin{lstlisting}[style=cpp]
	long long resultado = RectangulosEnGrilla::contar(2, 3);
	// resultado = 18
	// (N=2 da 3 lineas verticales, C(3,2)=3; M=3 da 4 lineas horizontales,
	//  C(4,2)=6; total 3*6=18; esto incluye los 8 cuadrados ya contados en
	//  la subseccion anterior, mas los rectangulos no cuadrados)
\end{lstlisting}

\textbf{Regla práctica:} usa esta fórmula cerrada $\binom{N+1}{2}\binom{M+1}{2}$ siempre que la pregunta sea por \textbf{todos} los rectángulos (sin restricción de ancho=alto); si la pregunta es específicamente por cuadrados, esa restricción rompe la independencia entre horizontal y vertical, y hace falta la suma de \textit{Número de Cuadrados en un Rectángulo N×M} en su lugar.

\subsection{Suma de Distancias Manhattan entre Todos los Pares}

\textbf{Idea:} calcular $\sum_{i<j} (|x_i-x_j| + |y_i-y_j|)$ para $n$ puntos ingenuamente es $O(n^2)$. La clave para bajar a $O(n\log n)$ es que la distancia Manhattan se \textbf{descompone} en dos términos independientes: $|x_i-x_j|+|y_i-y_j|$, así que la suma total sobre todos los pares es exactamente (suma de $|x_i-x_j|$ sobre todos los pares) + (suma de $|y_i-y_j|$ sobre todos los pares) —las coordenadas $x$ e $y$ se pueden procesar por separado, sin necesidad de considerar puntos como pares $(x,y)$ acoplados. Cada una de esas sumas 1D se calcula ordenando los valores y usando que, para un arreglo ordenado, la contribución de cada elemento a la suma de diferencias absolutas se puede acumular con una suma prefijo: el elemento en la posición $i$ (0-indexado) contribuye positivamente $i$ veces (es mayor que los $i$ anteriores) y negativamente $(n-1-i)$ veces (es menor que los restantes).

\textbf{¿Por qué usarlo?} La fuerza bruta $O(n^2)$ es inviable para $n$ grande (típicamente $n > 10^4$ en competencia). La descomposición en $x$ e $y$ por separado es posible \textbf{específicamente} porque la métrica Manhattan ya viene descompuesta en valor absoluto de diferencias por coordenada —esto NO funciona igual para distancia euclidiana, donde $x$ e $y$ están acoplados dentro de una raíz cuadrada y no se pueden separar de esta forma. Es una de las pocas ventajas computacionales concretas de usar Manhattan sobre euclidiana en problemas de suma de distancias.

\textbf{Casos de uso:}
\begin{itemize}
	\item Suma total de distancias Manhattan entre todos los pares de un conjunto de puntos, cuando $n$ es grande.
	\item Problemas de logística o redes donde el costo de conexión es proporcional a distancia Manhattan (movimiento en grilla).
	\item Como subrutina dentro de problemas más grandes que requieren esta suma repetidamente sobre subconjuntos (mantenida con estructuras de datos si los puntos cambian dinámicamente).
	\item Comparación de costo total bajo métrica Manhattan versus Chebyshev (que se relacionan mediante una rotación de coordenadas de 45 grados, no cubierta aquí).
\end{itemize}

\textbf{Complejidad:} $O(n\log n)$ —dominado por ordenar las coordenadas $x$ y las coordenadas $y$ por separado ($O(n\log n)$ cada una); el cálculo de la suma con prefijos sobre cada arreglo ordenado es $O(n)$.

\begin{lstlisting}[style=cpp]
	// CREATE/READ: Suma de distancias Manhattan entre todos los pares de n
	// puntos, mediante descomposicion en x e y independientes
	struct SumaManhattan {
		// READ: suma de |v[i] - v[j]| sobre todos los pares i<j, para un
		// arreglo 1D de valores (se usa una vez para x, otra vez para y)
		static long long sumaDiferenciasAbsolutas(vector<long long> v) {
			sort(v.begin(), v.end());
			int n = (int)v.size();
			long long suma = 0;
			long long prefijo = 0; // suma acumulada de v[0..i-1]
			for (int i = 0; i < n; i++) {
				// v[i] es mayor que los i elementos anteriores: contribuye
				// v[i]*i - prefijo (v[i] menos cada uno de los anteriores)
				suma += (long long)i * v[i] - prefijo;
				prefijo += v[i];
			}
			return suma;
		} // O(n log n), dominado por el sort
		
		// READ: suma total de distancia Manhattan entre todos los pares de
		// puntos (x[i], y[i])
		static long long calcular(const vector<long long>& x, const vector<long long>& y) {
			return sumaDiferenciasAbsolutas(x) + sumaDiferenciasAbsolutas(y);
		} // O(n log n)
		
		// ------------------------------------------------------------
		// MODIFICACIONES Y COMENTARIOS CON LOS USOS
		// ------------------------------------------------------------
		// 1. Distancia Chebyshev (max(|dx|,|dy|)) entre todos los pares: se
		//    puede reducir a Manhattan rotando las coordenadas 45 grados
		//    (u = x+y, v = x-y), calculando Manhattan sobre (u,v) y dividiendo
		//    el resultado entre 2 -- relacion clasica entre ambas metricas.
		// 2. Overflow: la suma total puede ser del orden de n^2 * (rango de
		//    coordenadas), facilmente mayor a 1e18 para n y coordenadas grandes;
		//    verificar los limites del problema antes de asumir que long long
		//    alcanza.
		// 3. Puntos que se agregan dinamicamente (no todos conocidos de
		//    antemano): esta formula requiere ordenar el arreglo completo, asi
		//    que no es directamente incremental; para actualizaciones online se
		//    necesitaria un Fenwick tree indexado por coordenada comprimida.
	};
\end{lstlisting}

\textbf{Ejemplo de funcionamiento:}
\begin{lstlisting}[style=cpp]
	vector<long long> x = {1, 3, 5};
	vector<long long> y = {2, 2, 2};
	
	long long resultado = SumaManhattan::calcular(x, y);
	// resultado = 8
	// (suma de |dx| sobre los 3 pares: |1-3|+|1-5|+|3-5| = 2+4+2 = 8;
	//  suma de |dy| sobre los 3 pares: 0+0+0 = 0 (todas las y son iguales);
	//  total: 8 + 0 = 8)
\end{lstlisting}

\textbf{Nota:} la derivación de \texttt{sumaDiferenciasAbsolutas} usando \texttt{i * v[i] - prefijo} depende de que \texttt{v} esté ordenado: una vez ordenado, cada elemento en la posición $i$ es mayor o igual que los $i$ elementos que lo preceden, así que su contribución neta a la suma de diferencias absolutas con esos $i$ elementos es $v[i]$ sumado $i$ veces menos la suma de esos $i$ elementos anteriores (\texttt{prefijo}).

\textbf{Regla práctica:} usa esta descomposición (separar $x$ e $y$, ordenar cada uno, sumar con prefijos) siempre que la métrica sea Manhattan y $n$ sea grande —es la técnica estándar para bajar de $O(n^2)$ a $O(n\log n)$ en este tipo de suma. Recuerda que esta separación es específica de Manhattan; no aplica directamente a distancia euclidiana por la raíz cuadrada que acopla $x$ e $y$.
\subsection{Máximo de Cuadrados 2x2 en un Triángulo Isósceles Rectángulo}

\textbf{Idea:} sea el triángulo rectángulo isósceles con catetos de longitud $b$ (los dos lados iguales), ángulo recto en el origen, y catetos alineados con los ejes —la hipotenusa es entonces la recta $x+y=b$. El problema de contar cuántos cuadrados de $2\times 2$ caben alineados a la grilla dentro de este triángulo es equivalente, tras dividir todas las longitudes entre 2, a contar cuántos cuadrados \textbf{unitarios} caben en un triángulo análogo de cateto $L=\lfloor b/2 \rfloor$. Un cuadrado unitario con esquina inferior izquierda en $(i,j)$ cabe completo bajo la hipotenusa si su esquina superior derecha $(i+1,j+1)$ satisface $(i+1)+(j+1)\le L$, es decir $i+j\le L-2$. Contar los pares $(i,j)$ con $i,j\ge 0$ enteros que cumplen esa desigualdad da $\sum_{s=0}^{L-2}(s+1) = \dfrac{L(L-1)}{2}$.

\textbf{¿Por qué usarlo?} Es un problema de conteo combinatorio bajo una restricción geométrica lineal (la hipotenusa), similar en espíritu a \textit{Número de Cuadrados en un Rectángulo N×M}, pero aquí la restricción no es un rectángulo sino un triángulo, lo que cambia la condición de "cabe" de una desigualdad por coordenada a una única desigualdad conjunta $i+j\le L-2$. Reducir cuadrados de $2\times 2$ a cuadrados unitarios mediante el cambio de escala $L=\lfloor b/2\rfloor$ evita tener que rederivar la fórmula para cada tamaño de cuadrado distinto.

\textbf{Casos de uso:}
\begin{itemize}
	\item Contar la máxima cantidad de baldosas cuadradas de tamaño fijo que caben en una región triangular de recorte.
	\item Problemas de empaquetamiento (packing) simplificado donde la región es un triángulo rectángulo isósceles.
	\item Como paso intermedio para estimar el área "desperdiciada" al triangular una región rectangular.
	\item Variante de conteo bajo la diagonal de una grilla cuadrada (equivalente al caso $L=b$ con cuadrados unitarios).
\end{itemize}

\textbf{Complejidad:} $O(1)$ —la fórmula cerrada se evalúa directamente sin ninguna iteración.

\begin{lstlisting}[style=cpp]
	// CREATE/READ: Maximo numero de cuadrados de lado k que caben alineados a
	// la grilla en un triangulo rectangulo isosceles de catetos de longitud b
	struct CuadradosEnTriangulo {
		// READ: cuenta cuadrados unitarios que caben bajo la hipotenusa de un
		// triangulo rectangulo isosceles de cateto L (formula L*(L-1)/2)
		static long long contarUnitarios(long long L) {
			if (L < 2) return 0; // no cabe ni un cuadrado si L < 2
			return L * (L - 1) / 2;
		} // O(1)
		
		// READ: cuenta cuadrados de lado k (por defecto k=2) que caben en un
		// triangulo de cateto b, escalando al caso unitario con L = b / k
		static long long contar(long long b, long long k = 2) {
			long long L = b / k; // division entera: escala al triangulo equivalente
			return contarUnitarios(L);
		} // O(1)
		
		// ------------------------------------------------------------
		// MODIFICACIONES Y COMENTARIOS CON LOS USOS
		// ------------------------------------------------------------
		// 1. Triangulo con la BASE como hipotenusa (no como cateto): si el
		//    enunciado da la longitud de la hipotenusa en vez del cateto, hay
		//    que convertir primero: cateto = hipotenusa / sqrt(2), lo cual ya
		//    no da un resultado entero exacto en general -- verificar con
		//    cuidado cual longitud especifica realmente el enunciado antes de
		//    aplicar esta formula.
		// 2. Cuadrados que pueden sobresalir parcialmente (no exigir que quepan
		//    COMPLETOS): cambia la condicion i+j <= L-2 a una menos estricta,
		//    dando un resultado distinto; la formula de este struct asume
		//    cuadrados completamente contenidos.
		// 3. b no divisible exactamente entre k: la division entera b/k ya
		//    trunca correctamente al mayor L util, sin necesidad de ajuste
		//    adicional.
	};
\end{lstlisting}

\textbf{Ejemplo de funcionamiento:}
\begin{lstlisting}[style=cpp]
	long long resultado = CuadradosEnTriangulo::contar(8, 2);
	// resultado = 6
	// (b=8, k=2, L = 8/2 = 4; contarUnitarios(4) = 4*3/2 = 6;
	//  hay que correrlo para confirmar el conteo exacto de posiciones válidas
	//  bajo la hipotenusa escalada)
\end{lstlisting}

\textbf{Nota:} este struct asume que la "base" $b$ dada en el enunciado se refiere a la longitud de los \textbf{catetos} (los dos lados iguales del triángulo rectángulo isósceles), no a la hipotenusa —es la interpretación más común de "base" en este tipo de problema, pero conviene confirmar contra el enunciado exacto antes de aplicar la fórmula, ya que si "base" refiere a la hipotenusa, la conversión introduce $\sqrt{2}$ y el problema deja de tener una fórmula puramente combinatoria en enteros.

\textbf{Regla práctica:} usa la reducción a cuadrados unitarios ($L=\lfloor b/k\rfloor$, luego $L(L-1)/2$) siempre que el triángulo sea rectángulo isósceles con catetos alineados a los ejes; para un triángulo rectángulo no isósceles (catetos de longitudes distintas), la condición de conteo cambia y requeriría rederivar la desigualdad de la hipotenusa correspondiente.

\subsection{Máximo Número de Piezas con N Cortes (Pastelero Perezoso)}

\textbf{Idea:} el problema del pastelero perezoso (\textit{Lazy Caterer's Sequence}) pregunta cuál es el máximo número de piezas en que se puede dividir un círculo (o cualquier región convexa) usando $n$ cortes rectos, donde cada corte es una línea recta completa que puede cruzar a todos los cortes anteriores. La estrategia óptima es que cada nuevo corte cruce a \textbf{todos} los cortes anteriores en puntos distintos (nunca por un punto ya usado ni paralelo a otro corte): el $k$-ésimo corte, al cruzar los $k-1$ cortes previos, queda dividido en $k$ segmentos, y cada uno de esos segmentos parte una pieza existente en dos, agregando exactamente $k$ piezas nuevas. Sumando esa ganancia para cada corte desde el primero: $p(n) = 1 + \sum_{k=1}^{n} k = 1 + \dfrac{n(n+1)}{2}$.

\textbf{¿Por qué usarlo?} Es un problema de conteo por recurrencia donde la clave no es geometría de coordenadas sino la observación de que \textbf{cada corte nuevo debe maximizar sus cruces} con los anteriores para maximizar piezas —cualquier corte paralelo a uno existente, o que pase por un punto de cruce ya usado, da menos piezas que el óptimo. La fórmula cerrada evita simular los cortes explícitamente; solo hace falta razonar sobre cuántos cruces nuevos introduce cada corte adicional.

\textbf{Casos de uso:}
\begin{itemize}
	\item El problema clásico de "máximo número de piezas de un pastel/pizza con $n$ cortes rectos".
	\item Generalización a la versión 3D (\textit{Cake Numbers}): máximo número de piezas de un sólido convexo con $n$ planos de corte.
	\item Como ejemplo introductorio de razonamiento por recurrencia en problemas de conteo geométrico.
	\item Verificación de que un algoritmo de simulación de cortes (para variantes con restricciones) no exceda el máximo teórico.
\end{itemize}

\textbf{Complejidad:} $O(1)$ —la fórmula cerrada se evalúa directamente.

\begin{lstlisting}[style=cpp]
	// CREATE/READ: Maximo numero de piezas al hacer n cortes rectos en un
	// pastel/region convexa (Lazy Caterer's Sequence)
	struct PasteleroPerezoso {
		// READ: maximo numero de piezas con n cortes, formula cerrada
		// p(n) = 1 + n(n+1)/2
		static long long maximoPiezas(long long n) {
			return 1 + n * (n + 1) / 2;
		} // O(1)
		
		// ------------------------------------------------------------
		// MODIFICACIONES Y COMENTARIOS CON LOS USOS
		// ------------------------------------------------------------
		// 1. Version 3D (Cake Numbers, n planos de corte en un solido convexo):
		//    formula distinta, C(n,0) + C(n,1) + C(n,2) + C(n,3), que en 2D se
		//    reduce a los primeros tres terminos (C(n,3) = 0 en el plano);
		//    struct separado si se necesita la version 3D.
		// 2. Cortes que deben ser PARALELOS entre si (variante restringida): la
		//    formula cambia a n+1 piezas (cada corte solo agrega 1 pieza, sin
		//    cruzar a los anteriores) -- mucho menor que el maximo sin
		//    restriccion.
		// 3. Overflow: n*(n+1) puede exceder el rango de long long solo para
		//    n extremadamente grande (~3e9), poco comun en restricciones
		//    tipicas de CP, pero vale verificar el enunciado.
	};
\end{lstlisting}

\textbf{Ejemplo de funcionamiento:}
\begin{lstlisting}[style=cpp]
	long long resultado = PasteleroPerezoso::maximoPiezas(3);
	// resultado = 7
	// (con 3 cortes que se cruzan de a pares en puntos distintos:
	//  1 (sin cortes) + 1 + 2 + 3 = 7 piezas)
\end{lstlisting}

\textbf{Regla práctica:} usa esta fórmula cuando el problema pida el \textbf{máximo} de piezas con $n$ cortes rectos sin restricciones adicionales sobre su orientación; si el problema impone una restricción (cortes paralelos, cortes que deben pasar por un punto común), la fórmula cambia por completo y hay que rederivarla contando cuántas piezas nuevas agrega cada corte bajo esa restricción específica.
\subsection{Altura Máxima de un Triángulo de Monedas}

\textbf{Idea:} dadas $n$ monedas, se apilan formando un triángulo donde la fila $1$ tiene $1$ moneda, la fila $2$ tiene $2$ monedas, y en general la fila $k$ tiene $k$ monedas —así que una altura de $h$ filas completas consume $\sum_{k=1}^{h} k = \dfrac{h(h+1)}{2}$ monedas. El problema pide la mayor altura $h$ tal que ese triángulo completo quepa con las $n$ monedas disponibles, es decir, el mayor $h$ que satisface $\dfrac{h(h+1)}{2} \le n$ (las monedas sobrantes, si las hay, no alcanzan para una fila completa adicional y se descartan). Como $\frac{h(h+1)}{2}$ crece monótonamente con $h$, esto se resuelve con búsqueda binaria sobre $h$, o despejando directamente de la ecuación cuadrática $h^2+h-2n\le 0$.

\textbf{¿Por qué usarlo?} Es el mismo patrón combinatorio que \textit{Máximo Número de Piezas con N Cortes} —una suma de la forma $\sum k$— pero aquí la incógnita es el límite superior de la suma dado un total fijo, no el valor de la suma dado el límite. Dos caminos válidos: \textit{Búsqueda Binaria on Answer} (ya cubierta en la sección de Búsqueda Binaria de esta guía) sobre $h$ probando $\frac{h(h+1)}{2}\le n$, o resolver directamente la fórmula cuadrática con $h = \left\lfloor \dfrac{-1+\sqrt{1+8n}}{2} \right\rfloor$. La búsqueda binaria evita cualquier riesgo de error de precisión en la raíz cuadrada de números grandes; la fórmula cerrada es más rápida pero requiere cuidado con el redondeo de \texttt{sqrt} en enteros grandes.

\textbf{Casos de uso:}
\begin{itemize}
	\item El problema clásico de "cuántas filas de un triángulo de monedas/bolos/cajas se pueden completar con $n$ unidades".
	\item Cualquier variante donde el costo de la fila $k$ crece linealmente con $k$ y se pregunta cuántas filas completas caben en un presupuesto fijo.
	\item Ejercicio introductorio para practicar la elección entre búsqueda binaria y fórmula cerrada cuando ambas son viables.
	\item Verificación cruzada: el problema es exactamente el inverso de \textit{Máximo Número de Piezas con N Cortes} (ahí se da $h$ y se pide la suma; aquí se da la suma y se pide $h$).
\end{itemize}

\textbf{Complejidad:} $O(\log n)$ con búsqueda binaria sobre $h$ (rango de búsqueda $[0, n]$, aunque $h$ real es $O(\sqrt n)$, así que también se puede acotar el rango a $[0, \sqrt{2n}+1]$ para menos iteraciones); $O(1)$ con la fórmula cerrada, al costo de una llamada a \texttt{sqrt}.

\begin{lstlisting}[style=cpp]
	// CREATE/READ: Altura maxima de un triangulo de monedas que se puede
	// formar con n monedas disponibles (fila k tiene k monedas)
	struct TrianguloMonedas {
		// READ: monedas necesarias para completar h filas (formula h(h+1)/2)
		static long long monedasParaAltura(long long h) {
			return h * (h + 1) / 2;
		} // O(1)
		
		// READ: altura maxima alcanzable con n monedas, via busqueda binaria
		// sobre h (mas robusto ante errores de precision que la formula cerrada)
		static long long alturaMaximaBinaria(long long n) {
			long long lo = 0, hi = 2000000000LL; // cota superior segura
			while (lo < hi) {
				long long medio = lo + (hi - lo + 1) / 2;
				if (monedasParaAltura(medio) <= n) lo = medio;
				else hi = medio - 1;
			}
			return lo;
		} // O(log n)
		
		// READ: altura maxima alcanzable con n monedas, via formula cerrada
		// despejada de h^2 + h - 2n <= 0 (requiere verificar el redondeo)
		static long long alturaMaximaFormula(long long n) {
			long long h = (long long)((-1.0 + sqrt(1.0 + 8.0 * n)) / 2.0);
			// ajuste por posible error de precision de sqrt en punto flotante
			while (monedasParaAltura(h + 1) <= n) h++;
			while (h > 0 && monedasParaAltura(h) > n) h--;
			return h;
		} // O(1) amortizado (el ajuste final es O(1) en la practica)
		
		// ------------------------------------------------------------
		// MODIFICACIONES Y COMENTARIOS CON LOS USOS
		// ------------------------------------------------------------
		// 1. Monedas sobrantes (no usadas en el triangulo completo): n -
		//    monedasParaAltura(alturaMaximaBinaria(n)) da la cantidad de
		//    monedas que quedan sin colocar tras formar la mayor altura posible.
		// 2. Costo de fila k distinto (no lineal en k, por ejemplo k^2 monedas
		//    por fila): la busqueda binaria sigue funcionando reemplazando
		//    monedasParaAltura por la suma correspondiente; la formula cerrada
		//    habria que rederivarla para cada caso.
		// 3. n muy grande (overflow en 8*n dentro de la formula cerrada): usar
		//    long long consistentemente en monedasParaAltura evita overflow
		//    hasta n en el orden de 1e18; para n mayor, revisar el enunciado.
	};
\end{lstlisting}

\textbf{Ejemplo de funcionamiento:}
\begin{lstlisting}[style=cpp]
	long long resultado = TrianguloMonedas::alturaMaximaBinaria(10);
	// resultado = 4
	// (fila 1:1, fila 2:2, fila 3:3, fila 4:4; total 1+2+3+4=10 monedas,
	//  exactamente 10, sin sobrantes; con 5 filas se necesitarian 15, que
	//  excede las 10 disponibles)
\end{lstlisting}

\textbf{Regla práctica:} usa \texttt{alturaMaximaBinaria} como opción por defecto y segura frente a errores de precisión; usa \texttt{alturaMaximaFormula} solo si el rendimiento de $O(\log n)$ no es aceptable, y en ese caso no olvides el ajuste posterior con los \texttt{while} de corrección —\texttt{sqrt} en números grandes puede desviarse por error de punto flotante y dar un $h$ incorrecto en $\pm 1$ sin el ajuste.


\part{Programación Dinámica (Dynamic Programming)}
\subsection{¿Qué algoritmo usar?}

% ---------------------------------------------------------
% TABLA 1: FUNDAMENTOS Y SECUENCIAS
% ---------------------------------------------------------
\begin{table}[H]
	\raggedright
	\scriptsize
	\setlength{\tabcolsep}{2.5pt}
	\renewcommand{\arraystretch}{1.15}
	\begin{tabular}{|p{0.43\columnwidth}|p{0.22\columnwidth}|p{0.27\columnwidth}|}
		\hline
		\textbf{Cuando el ejercicio diga / pida} &
		\textbf{Usa} &
		\textbf{Pista rápida} \\
		\hline
		
		``Calcula $f(n)$'' con una recurrencia que repite subproblemas
		(Fibonacci, escaleras) &
		Memoización / Tabulación &
		Memo si el orden no es obvio; tabla si hay orden por índice. \\
		\hline
		
		``Calcula $f(n)$ módulo $p$'' con $n$ hasta $10^{18}$ y recurrencia
		lineal de orden fijo &
		Exponenciación de matrices &
		$O(k^3\log n)$. Primera fila = coeficientes. \\
		\hline
		
		``Subarreglo contiguo de suma máxima'' &
		Kadane &
		\texttt{max(a[i], act+a[i])}. $O(n)$, $O(1)$ de espacio. \\
		\hline
		
		``Máxima suma circular'' o ``máxima submatriz'' &
		Kadane circular / Kadane por pares de filas &
		Circular: total $-$ mínimo subarreglo. 2D: $O(n^2m)$. \\
		\hline
		
		``Mejor día para comprar y vender (una vez)'' &
		Kadane sobre diferencias &
		Diferencias consecutivas de precios. \\
		\hline
		
		``¿De cuántas formas elijo $k$ de $n$?'', ``caminos en grilla
		derecha/abajo'' ($n$ moderado) &
		Triángulo de Pascal &
		Solo sumas, sirve con cualquier módulo. $O(n^2)$. \\
		\hline
		
		Lo mismo con $n\approx 10^6$ y módulo primo &
		Factoriales e inversos modulares &
		Precómputo $O(n)$, consulta $O(1)$. \\
		\hline
		
		``Camino creciente más largo en una matriz'' (movimiento en 4
		direcciones, valores estrictos) &
		DFS con memoización &
		Grafo sin ciclos por ser estricto. $O(nm)$. \\
		\hline
		
		``Subsecuencia creciente más larga'' (solo la longitud) &
		LIS con búsqueda binaria &
		$O(n\log n)$. \\
		\hline
		
		``Cuántas LIS hay'', ``reconstruye una'', ``subsecuencia bitónica'' &
		LIS por DP $O(n^2)$ &
		Guarda \texttt{dp}, \texttt{padre} y \texttt{cuenta} por índice. \\
		\hline
		
		``Subsecuencia creciente de \textbf{suma} máxima'' &
		DP tipo LIS con suma &
		Con Fenwick sobre valores comprimidos: $O(n\log n)$. \\
		\hline
		
		``Subsecuencia común de dos cadenas'', ``alinear dos secuencias'',
		``palíndromo más largo'' &
		LCS (DP sobre pares de prefijos) &
		$O(nm)$. Palíndromo: LCS con el reverso. \\
		\hline
		
		``Pila más alta de cajas con rotaciones'', ``muñecas rusas'',
		``sobres anidados'' &
		Box Stacking (orden + DP tipo LIS) &
		Ordena por área de base decreciente. $O(n^2)$. \\
		\hline
		
	\end{tabular}
	\caption{Fundamentos y secuencias: qué usar según lo que pide el ejercicio.}
\end{table}

% ---------------------------------------------------------
% TABLA 2: KNAPSACK, MONEDAS, INTERVALOS
% ---------------------------------------------------------
\begin{table}[H]
	\raggedright
	\scriptsize
	\setlength{\tabcolsep}{2.5pt}
	\renewcommand{\arraystretch}{1.15}
	\begin{tabular}{|p{0.43\columnwidth}|p{0.22\columnwidth}|p{0.27\columnwidth}|}
		\hline
		\textbf{Cuando el ejercicio diga / pida} &
		\textbf{Usa} &
		\textbf{Pista rápida} \\
		\hline
		
		``Elige objetos con peso y valor, \textbf{cada uno a lo sumo una vez},
		con capacidad $W$'' &
		Knapsack 0/1 &
		Capacidad en sentido \textbf{decreciente}. $O(nW)$. \\
		\hline
		
		``Cada objeto se puede usar \textbf{las veces que quiera}'', cortar
		una barra (rod cutting) &
		Knapsack Ilimitado &
		Capacidad en sentido creciente. $O(nW)$. \\
		\hline
		
		``Cada objeto con cantidad limitada $c_i$'' &
		Knapsack acotado &
		Descompón $c_i$ en potencias de dos. \\
		\hline
		
		``Objetos que se pueden partir'' (fraccionario) &
		Greedy por valor/peso &
		No es DP. Ordena y toma. \\
		\hline
		
		``¿Existe un subconjunto con suma exacta $S$?'' &
		Subset Sum (Knapsack con \texttt{bool}) &
		Con \texttt{bitset}: \texttt{b |= b << x}. \\
		\hline
		
		``¿De cuántas formas formo el monto $n$ con estas monedas?'' (el
		orden \textbf{no} importa) &
		Cambio de monedas: combinaciones &
		Monedas \textbf{afuera}, monto adentro. $O(nm)$. \\
		\hline
		
		``¿De cuántas formas?'' cuando el orden \textbf{sí} importa
		(escaleras con saltos) &
		Cambio de monedas: permutaciones &
		Monto \textbf{afuera}, monedas adentro. \\
		\hline
		
		``Mínimo de monedas para pagar $n$'' (sistema arbitrario) &
		Cambio de monedas: mínimo &
		Greedy solo si el sistema es canónico. $-1$ si es imposible. \\
		\hline
		
		``Divide el arreglo en dos grupos con diferencia mínima'', ``¿se
		puede dividir en dos mitades iguales?'' &
		Problema de Partición &
		Subset Sum hasta $S/2$. Respuesta $S-2x$. \\
		\hline
		
		Lo mismo con $n\le 40$ y suma enorme &
		Meet in the middle &
		Genera sumas de cada mitad y une con dos punteros. \\
		\hline
		
		``Orden óptimo de eliminar elementos'' donde quitar uno cambia los
		vecinos (globos) &
		DP de intervalos (Burst Balloons) &
		Modela el \textbf{último} que queda. $O(n^3)$. \\
		\hline
		
		``Orden óptimo de multiplicar matrices'', ``árbol de búsqueda
		óptimo'', ``triangular un polígono'' &
		DP de intervalos &
		Estado $[l,r]$, prueba el corte $k$, longitud creciente. \\
		\hline
		
		DP de intervalos $O(n^3)$ que no pasa el tiempo y el costo cumple la
		desigualdad cuadrangular &
		Optimización de Knuth &
		Acota $k$ entre \texttt{opt[i][j-1]} y \texttt{opt[i+1][j]}. $O(n^2)$. \\
		\hline
		
	\end{tabular}
	\caption{Knapsack, monedas e intervalos: qué usar según lo que pide el ejercicio.}
\end{table}

% ---------------------------------------------------------
% TABLA 3: DP AVANZADA
% ---------------------------------------------------------
\begin{table}[H]
	\raggedright
	\scriptsize
	\setlength{\tabcolsep}{2.5pt}
	\renewcommand{\arraystretch}{1.15}
	\begin{tabular}{|p{0.43\columnwidth}|p{0.22\columnwidth}|p{0.27\columnwidth}|}
		\hline
		\textbf{Cuando el ejercicio diga / pida} &
		\textbf{Usa} &
		\textbf{Pista rápida} \\
		\hline
		
		``¿Cuántos números en $[L,R]$ cumplen algo sobre sus dígitos?'' (suma
		de dígitos, dígito prohibido, dígitos crecientes) con $R\le 10^{18}$ &
		Digit DP &
		Respuesta $f(R)-f(L-1)$. Bandera \texttt{tight} y manejo de ceros
		iniciales. \\
		\hline
		
		Lo mismo combinado con divisibilidad por $m$ &
		Digit DP con resto en el estado &
		Estado \texttt{(pos, resto, tight)}. \\
		\hline
		
		``Número esperado de lanzamientos / pasos hasta llegar a la meta'',
		``valor esperado'' &
		DP de valor esperado &
		$E(s)=\sum p_i\,(c_i+E(s_i))$. Despeja si el estado se repite. \\
		\hline
		
		``Valor esperado'' con ciclos entre varios estados distintos &
		Sistema lineal (Gauss) &
		Una ecuación por estado del ciclo. \\
		\hline
		
		``Valor esperado'' de una suma de eventos simples &
		Linealidad de la esperanza &
		Suma de esperanzas individuales, sin DP. \\
		\hline
		
		``Particiona en $m$ grupos contiguos minimizando costo'', DP
		$O(n^2m)$ con el óptimo monótono &
		Divide and Conquer Optimization &
		Verifica la desigualdad cuadrangular. $O(nm\log n)$. \\
		\hline
		
		DP del tipo $\min_i(b_i x_j + a_i)$, ``elige la mejor recta
		evaluada en $x$'' &
		Convex Hull Trick (deque) &
		Pendientes y consultas monótonas. $O(n)$. \\
		\hline
		
		Lo mismo sin orden en pendientes o consultas &
		Li Chao Tree &
		$O(\log C)$ por inserción y consulta. \\
		\hline
		
	\end{tabular}
	\caption{DP avanzada: qué usar según lo que pide el ejercicio.}
\end{table}
\section{Fundamentos}
\subsection{Memoización vs. Tabulación (Top-Down y Bottom-Up)}

\textbf{Idea:} la programación dinámica resuelve un problema descomponiéndolo en \textbf{subproblemas superpuestos}: subproblemas que se repiten muchas veces dentro de la recursión. En lugar de recalcularlos, se guarda cada resultado la primera vez que se obtiene. Hay dos formas de hacerlo. La \textbf{memoización} (\textit{top-down}) escribe la recursión natural del problema y agrega una tabla que se consulta antes de calcular: si el estado ya fue resuelto, se devuelve el valor guardado. La \textbf{tabulación} (\textit{bottom-up}) invierte el orden: se llena la tabla iterativamente, desde los casos base hasta el estado que interesa, de modo que cada estado se calcula cuando todos los estados de los que depende ya están listos.

\textbf{¿Por qué usarlo?} Sin guardar resultados, la recursión de Fibonacci ($F(n)=F(n-1)+F(n-2)$) recalcula los mismos valores una y otra vez y cuesta $O(\varphi^n)$, exponencial. Con cualquiera de las dos técnicas cada estado se calcula \textbf{una sola vez}, y el costo baja a $O(n)$. La memoización es más fácil de escribir porque sigue la recurrencia tal cual, y solo calcula los estados que realmente se necesitan; su costo es la profundidad de recursión (riesgo de desbordar la pila) y el \textit{overhead} de las llamadas. La tabulación es iterativa, más rápida en la práctica, y permite reducir memoria cuando cada estado depende solo de unos pocos anteriores (como en Fibonacci, donde bastan dos variables).

\textbf{Casos de uso:}
\begin{itemize}
	\item Convertir cualquier recurrencia con subproblemas repetidos en un algoritmo polinomial.
	\item Memoización: cuando el espacio de estados es enorme pero solo una fracción pequeña es alcanzable, o cuando el orden de llenado no es evidente.
	\item Tabulación: cuando todos los estados se necesitan de todas formas y se quiere evitar el límite de recursión.
	\item Reducción de memoria: cuando cada fila o valor depende solo de la anterior (rolling array).
	\item Base conceptual de todas las demás secciones de esta parte (Knapsack, DP de Intervalos, DP sobre Secuencias).
\end{itemize}

\textbf{Complejidad:} $O(n)$ tiempo en ambas versiones, porque hay $n+1$ estados y cada uno se calcula una vez con trabajo $O(1)$. La memoización usa además $O(n)$ de pila de recursión; la tabulación usa $O(n)$ de tabla, reducible a $O(1)$ con dos variables.

\begin{lstlisting}[style=cpp]
	// CREATE/READ/WRITE: Fibonacci resuelto de tres formas para comparar
	// memoizacion (top-down), tabulacion (bottom-up) y espacio constante
	struct FibonacciDP {
		vector<long long> memo; // memo[i] = F(i), o -1 si aun no se calculo
		
		// CREATE: prepara la tabla de memoizacion con centinela -1
		FibonacciDP(int n) {
			memo.assign(n + 1, -1);
		} // O(n)
		
		// READ: memoizacion (top-down): recursion natural + consulta a la tabla
		long long memoizado(int n) {
			if (n <= 1) return n;                 // casos base
			if (memo[n] != -1) return memo[n];    // estado ya resuelto
			memo[n] = memoizado(n - 1) + memoizado(n - 2);
			return memo[n];
		} // O(n) tiempo, O(n) pila de recursion
		
		// READ: tabulacion (bottom-up): llena la tabla de menor a mayor
		static long long tabulado(int n) {
			if (n <= 1) return n;
			vector<long long> dp(n + 1, 0);
			dp[1] = 1;
			for (int i = 2; i <= n; i++) dp[i] = dp[i - 1] + dp[i - 2];
			return dp[n];
		} // O(n) tiempo, O(n) espacio
		
		// READ: tabulacion con espacio constante: solo se conservan los dos
		// ultimos valores, porque F(i) depende unicamente de F(i-1) y F(i-2)
		static long long espacioConstante(int n) {
			if (n <= 1) return n;
			long long anterior = 0, actual = 1;
			for (int i = 2; i <= n; i++) {
				long long siguiente = anterior + actual;
				anterior = actual;
				actual = siguiente;
			}
			return actual;
		} // O(n) tiempo, O(1) espacio
		
		// ------------------------------------------------------------
		// MODIFICACIONES Y COMENTARIOS CON LOS USOS
		// ------------------------------------------------------------
		// 1. Estados dispersos o muy grandes (por ejemplo indices hasta 1e18
		//    o claves compuestas): reemplazar el vector por
		//    unordered_map<long long, long long> memo y consultar con
		//    memo.find(n); la tabla solo guarda los estados realmente visitados.
		// 2. Recursion demasiado profunda (n ~ 1e6 puede desbordar la pila):
		//    convertir la memoizacion a tabulacion, o aumentar el limite de
		//    pila del sistema; la tabulacion no tiene este problema.
		// 3. Resultado modular: sumar y reducir en cada paso,
		//    dp[i] = (dp[i-1] + dp[i-2]) % MOD, para evitar overflow en n grande.
		//    Sin modulo, F(90) es el ultimo Fibonacci que cabe en long long.
		// 4. n astronomico (n ~ 1e18): ni la memoizacion ni la tabulacion
		//    alcanzan; se usa exponenciacion de matrices
		//    [[1,1],[1,0]]^n en O(log n).
		// 5. Centinela de la tabla: usar -1 como "no calculado" solo es valido
		//    si -1 no puede ser una respuesta real; si lo fuera, usar un
		//    vector<bool> visitado aparte.
		// 6. Rolling array en DP de dos dimensiones: si dp[i][j] solo depende
		//    de la fila i-1, guardar dos filas (o una sola, recorriendo j en
		//    el orden correcto) en vez de la tabla completa.
	};
\end{lstlisting}

\textbf{Ejemplo de funcionamiento:}
\begin{lstlisting}[style=cpp]
	FibonacciDP f(50);
	
	long long a = f.memoizado(10);
	// a = 55
	
	long long b = FibonacciDP::tabulado(10);
	// b = 55 (mismo resultado, calculado de menor a mayor)
	
	long long c = FibonacciDP::espacioConstante(50);
	// c = 12586269025 (con solo dos variables en memoria)
\end{lstlisting}

\textbf{Nota:} las tres versiones dan siempre el mismo resultado; lo que cambia es el orden de cálculo y la memoria. Un error frecuente al pasar de memoización a tabulación es llenar la tabla en un orden en que un estado se calcula \textbf{antes} de que sus dependencias estén listas. Antes de escribir el bucle, conviene identificar de qué estados depende cada uno.

\textbf{Regla práctica:} empieza siempre por la recurrencia y los casos base. Si el orden de llenado no es obvio o solo una parte de los estados es alcanzable, usa memoización. Si el problema tiene un orden natural (por índice, por longitud creciente) y $n$ es grande, usa tabulación, y reduce la memoria con \textit{rolling array} cuando cada estado dependa solo de los anteriores inmediatos. Esta comparación se aplica a todas las subsecciones siguientes de esta parte.
\subsection{Kadane's Algorithm (Máximo Subarreglo)}

\textbf{Idea:} dado un arreglo de $n$ enteros (que pueden ser negativos), se busca el \textbf{subarreglo contiguo} cuya suma sea máxima. Sea $\text{dp}[i]$ la suma máxima de un subarreglo contiguo que \textbf{termina exactamente} en el índice $i$. En cada posición hay solo dos opciones: \textbf{extender} el mejor subarreglo que termina en $i-1$ sumándole $a[i]$, o \textbf{empezar de nuevo} en $i$ si arrastrar lo anterior perjudica más de lo que ayuda —es decir, $\text{dp}[i]=\max(a[i],\ \text{dp}[i-1]+a[i])$—. La decisión de "empezar de nuevo" tiene sentido exactamente cuando $\text{dp}[i-1]<0$: arrastrar una suma negativa nunca ayuda. La respuesta final es $\max_i \text{dp}[i]$, no $\text{dp}[n-1]$, porque el subarreglo óptimo puede terminar en cualquier posición.

\textbf{¿Por qué usarlo?} Probar todos los $O(n^2)$ subarreglos (o $O(n^3)$ sin sumas de prefijo) es la fuerza bruta directa. Kadane resuelve el mismo problema en $O(n)$ con un único recorrido, porque $\text{dp}[i]$ resume todo lo que hace falta saber de los subarreglos que terminan en $i$ sin necesidad de volver a mirar atrás. Es el ejemplo más simple de DP de un solo arreglo con \textbf{espacio $O(1)$}: como $\text{dp}[i]$ solo depende de $\text{dp}[i-1]$, no hace falta guardar la tabla completa, solo el valor anterior — el mismo patrón de \textit{rolling array} mencionado en \textit{Memoización vs. Tabulación}.

\textbf{Casos de uso:}
\begin{itemize}
	\item Máximo subarreglo contiguo en un arreglo con valores positivos y negativos mezclados.
	\item Máxima ganancia de una única transacción de compra-venta (diferencia de precios consecutivos transformada en un arreglo de diferencias, y Kadane sobre ese arreglo).
	\item Como subrutina de \textit{Stress Testing}: ejemplo estándar de solución rápida a verificar contra fuerza bruta $O(n^2)$.
	\item Base conceptual de variantes 2D (máxima submatriz) que fijan dos filas y aplican Kadane 1D sobre la suma de columnas entre esas filas.
	\item Punto de partida para variantes con restricciones: longitud máxima o mínima del subarreglo, o exactamente $k$ elementos.
\end{itemize}

\textbf{Complejidad:} $O(n)$ tiempo con un único recorrido del arreglo; $O(1)$ espacio, manteniendo solo el mejor valor que termina en la posición actual y el mejor global visto hasta ahora.

\begin{lstlisting}[style=cpp]
	// CREATE/READ: Kadane's Algorithm: suma maxima de un subarreglo contiguo,
	// con reconstruccion de los indices del subarreglo optimo
	struct Kadane {
		vector<long long> a;
		
		// CREATE: guarda el arreglo
		Kadane(const vector<long long>& v) : a(v) {
		} // O(n)
		
		// READ: suma maxima de un subarreglo contiguo no vacio, en O(1) espacio
		long long sumaMaxima() const {
			long long mejorActual = a[0], mejorGlobal = a[0];
			for (size_t i = 1; i < a.size(); i++) {
				mejorActual = max(a[i], mejorActual + a[i]);
				mejorGlobal = max(mejorGlobal, mejorActual);
			}
			return mejorGlobal;
		} // O(n) tiempo, O(1) espacio
		
		// READ: retorna {inicio, fin} (indices 0-indexados, inclusive) del
		// subarreglo optimo
		pair<int,int> reconstruir() const {
			long long mejorActual = a[0], mejorGlobal = a[0];
			int inicioActual = 0, mejorInicio = 0, mejorFin = 0;
			for (size_t i = 1; i < a.size(); i++) {
				if (a[i] > mejorActual + a[i]) {
					mejorActual = a[i];
					inicioActual = (int)i; // conviene "empezar de nuevo" aqui
				} else {
					mejorActual += a[i];
				}
				if (mejorActual > mejorGlobal) {
					mejorGlobal = mejorActual;
					mejorInicio = inicioActual;
					mejorFin = (int)i;
				}
			}
			return {mejorInicio, mejorFin};
		} // O(n) tiempo, O(1) espacio
		
		// ------------------------------------------------------------
		// MODIFICACIONES Y COMENTARIOS CON LOS USOS
		// ------------------------------------------------------------
		// 1. Arreglo con TODOS los elementos negativos: el codigo de arriba
		//    ya maneja este caso correctamente (devuelve el elemento negativo
		//    menos malo), porque el subarreglo debe ser NO VACIO por
		//    definicion del problema; no inicializar mejorActual en 0, sino
		//    en a[0], es lo que garantiza esto.
		// 2. Permitir subarreglo VACIO (suma minima posible es 0, no un
		//    numero negativo): inicializar mejorActual y mejorGlobal en 0 en
		//    vez de a[0], y en cada paso hacer
		//    mejorActual = max(0LL, mejorActual + a[i]).
		// 3. Minimo subarreglo (en vez de maximo): negar todos los valores,
		//    aplicar Kadane normal, y negar el resultado; o invertir los max
		//    por min directamente en el codigo.
		// 4. Maxima suma circular (el subarreglo puede "dar la vuelta" del
		//    final al principio): la respuesta es
		//    max(Kadane normal, sumaTotal - Kadane del MINIMO subarreglo),
		//    con el caso especial de que si TODO el arreglo es negativo, el
		//    segundo termino da 0 por error y hay que devolver solo el
		//    Kadane normal.
		// 5. Maxima submatriz 2D: fijar cada par de filas (fSup, fInf) en
		//    O(n^2) combinaciones, calcular el arreglo de sumas de columna
		//    entre esas filas en O(m), y aplicar Kadane 1D sobre ese arreglo
		//    en O(m); total O(n^2 * m).
		// 6. Restriccion de longitud MINIMA k: mantener una suma de prefijo y,
		//    para cada posicion i >= k, comparar prefijo[i] - min(prefijo[0..i-k]);
		//    Kadane simple no respeta un minimo de longitud por si solo.
		// 7. Ganancia maxima de UNA transaccion de compra-venta de acciones:
		//    aplicar Kadane sobre el arreglo de diferencias consecutivas
		//    precio[i] - precio[i-1]; el resultado es la maxima ganancia de
		//    comprar en un dia y vender en uno posterior.
	};
\end{lstlisting}

\textbf{Ejemplo de funcionamiento:}
\begin{lstlisting}[style=cpp]
	Kadane k({-2, 1, -3, 4, -1, 2, 1, -5, 4});
	
	long long resultado = k.sumaMaxima();
	// resultado = 6
	// (el subarreglo 4, -1, 2, 1 suma 6, el maximo posible)
	
	auto rango = k.reconstruir();
	// rango = {3, 6}
	// (indices 3 a 6, correspondientes a los valores 4, -1, 2, 1)
\end{lstlisting}

\textbf{Nota:} el error más común es inicializar \texttt{mejorActual} en \texttt{0} en vez de \texttt{a[0]} cuando el problema exige que el subarreglo sea \textbf{no vacío}. Con esa inicialización incorrecta, un arreglo completamente negativo (como \texttt{\{-5, -3, -8\}}) devuelve \texttt{0} en vez del valor correcto \texttt{-3} (el menos negativo), porque el algoritmo "tiene permitido" no tomar nada. Revisa siempre si el enunciado exige un subarreglo no vacío antes de decidir la inicialización.

\textbf{Regla práctica:} usa Kadane como la solución por defecto para "máxima suma de subarreglo contiguo" — es $O(n)$, $O(1)$ de espacio, y casi nunca hace falta nada más complejo. Antes de escribirlo, confirma si el subarreglo debe ser no vacío (inicializa en \texttt{a[0]}) o puede ser vacío con suma mínima 0 (inicializa en \texttt{0}, ver modificación 2), y si el arreglo es circular (ver modificación 4) — son los dos detalles que más distinguen un enunciado de otro.
\subsection{Triángulo de Pascal}

\textbf{Idea:} el triángulo de Pascal organiza los coeficientes binomiales $\binom{n}{k}$ en filas: la fila $n$ contiene $\binom{n}{0},\binom{n}{1},\dots,\binom{n}{n}$. Cada número es la suma de los dos que tiene encima: $\binom{n}{k}=\binom{n-1}{k-1}+\binom{n-1}{k}$, con $\binom{n}{0}=\binom{n}{n}=1$. La recurrencia tiene una lectura combinatoria directa. Para elegir $k$ elementos de $n$, se fija un elemento especial: o se incluye (faltan $k-1$ por elegir entre los otros $n-1$) o no se incluye (faltan $k$ entre los otros $n-1$). Como los subproblemas $\binom{n-1}{\cdot}$ se reutilizan en toda la fila siguiente, es un caso de DP muy claro sobre \textit{Memoización vs. Tabulación}: la tabla se llena fila por fila, y cada celda solo necesita la fila anterior.

\textbf{¿Por qué usarlo?} La fórmula cerrada $\binom{n}{k}=\frac{n!}{k!(n-k)!}$ desborda con facilidad (ya $21!$ no cabe en \texttt{long long}) y requiere inversos modulares cuando se trabaja módulo un primo. El triángulo evita por completo divisiones y factoriales: solo usa \textbf{sumas}, así que funciona con cualquier módulo (primo o no) y es exacto mientras el resultado quepa en el tipo. Su costo es $O(n^2)$ de precomputación, que es aceptable cuando $n$ es del orden de $10^3$ a $5\cdot 10^3$ y se necesitan muchas consultas $\binom{n}{k}$ distintas. Si $n$ es muy grande (por ejemplo $10^6$), conviene la fórmula con factoriales precomputados e inversos modulares.

\textbf{Casos de uso:}
\begin{itemize}
	\item Precomputar una tabla de $\binom{n}{k}$ para responder muchas consultas en $O(1)$.
	\item Contar caminos en una grilla que solo avanza a la derecha y hacia abajo: $\binom{f+c}{f}$ para $f$ pasos hacia abajo y $c$ pasos a la derecha.
	\item Contar subconjuntos de tamaño $k$, o formas de repartir objetos en problemas combinatorios de conteo.
	\item Coeficientes del desarrollo de $(x+y)^n$ y de otros polinomios sencillos.
	\item Trabajar con un módulo compuesto (no primo), donde los inversos modulares no siempre existen.
\end{itemize}

\textbf{Complejidad:} $O(n^2)$ tiempo y $O(n^2)$ espacio para construir la tabla hasta la fila $n$, ya que hay $\frac{(n+1)(n+2)}{2}$ celdas y cada una se calcula con una suma en $O(1)$. Una consulta posterior es $O(1)$. Para obtener solo una fila, $O(n^2)$ tiempo pero $O(n)$ espacio.

\begin{lstlisting}[style=cpp]
	// CREATE/READ/WRITE: Triangulo de Pascal para coeficientes binomiales,
	// con modulo opcional (mod = 0 significa sin modulo)
	struct TrianguloPascal {
		vector<vector<long long>> tabla; // tabla[i][j] = C(i, j)
		long long mod;
		
		// CREATE: construye las filas 0..n; cada celda es la suma de las dos
		// que tiene encima (tabulacion fila por fila)
		TrianguloPascal(int n, long long m = 0) : mod(m) {
			tabla.assign(n + 1, vector<long long>());
			for (int i = 0; i <= n; i++) {
				tabla[i].assign(i + 1, 1); // bordes C(i,0) = C(i,i) = 1
				for (int j = 1; j < i; j++) {
					long long v = tabla[i - 1][j - 1] + tabla[i - 1][j];
					if (mod > 0) v %= mod;
					tabla[i][j] = v;
				}
			}
		} // O(n^2)
		
		// READ: coeficiente binomial C(n, k); devuelve 0 si k esta fuera de rango
		long long combinacion(int n, int k) const {
			if (k < 0 || k > n) return 0;
			return tabla[n][k];
		} // O(1)
		
		// READ: fila completa n del triangulo
		const vector<long long>& fila(int n) const {
			return tabla[n];
		} // O(1)
		
		// READ: una sola fila n con espacio O(n), sin guardar la tabla completa
		// (rolling array: se actualiza de derecha a izquierda para no pisar
		// valores que todavia hacen falta)
		static vector<long long> filaUnica(int n) {
			vector<long long> f(n + 1, 0);
			f[0] = 1;
			for (int i = 1; i <= n; i++) {
				for (int j = i; j >= 1; j--) f[j] += f[j - 1];
			}
			return f;
		} // O(n^2) tiempo, O(n) espacio
		
		// ------------------------------------------------------------
		// MODIFICACIONES Y COMENTARIOS CON LOS USOS
		// ------------------------------------------------------------
		// 1. Overflow: sin modulo, C(n,k) cabe en long long solo hasta n = 66
		//    (C(66,33) ~ 7.2e18). Para n mayor, usar un modulo (por ejemplo
		//    1e9+7) o __int128 si se necesita el valor exacto pero acotado.
		// 2. Modulo compuesto: este triangulo funciona con cualquier modulo,
		//    incluso no primo, porque solo suma; la formula con factoriales e
		//    inversos modulares exige un modulo primo.
		// 3. n grande (~1e6) con muchas consultas: precomputar fact[i] e
		//    invFact[i] mod p primo, y responder
		//    C(n,k) = fact[n] * invFact[k] % p * invFact[n-k] % p en O(1),
		//    con precomputo O(n) en vez de O(n^2).
		// 4. Caminos en grilla que solo avanza derecha o abajo, con f pasos
		//    hacia abajo y c hacia la derecha: combinacion(f + c, f).
		// 5. Suma de la fila n: siempre 2^n, ya que cada elemento de un
		//    conjunto de n se incluye o no; util como chequeo de cordura.
		// 6. Identidades utiles: C(n,k) = C(n,n-k) (simetria de cada fila) y
		//    C(n,0)+C(n,1)+...+C(n,n) = 2^n; permiten calcular solo la mitad
		//    de cada fila si se necesita ahorrar tiempo.
		// 7. Aplicar el modulo en cada suma (como hace el constructor con
		//    mod): reducir solo al final no evita el desbordamiento intermedio.
	};
\end{lstlisting}

\textbf{Ejemplo de funcionamiento:}
\begin{lstlisting}[style=cpp]
	TrianguloPascal t(10);
	
	long long a = t.combinacion(5, 2);
	// a = 10 (formas de elegir 2 elementos entre 5)
	
	vector<long long> fila5 = t.fila(5);
	// fila5 = {1, 5, 10, 10, 5, 1}
	// (simetrica; su suma es 32 = 2^5)
	
	long long b = t.combinacion(10, 3);
	// b = 120
	// (por la recurrencia: C(10,3) = C(9,2) + C(9,3) = 36 + 84 = 120)
	
	vector<long long> f6 = TrianguloPascal::filaUnica(6);
	// f6 = {1, 6, 15, 20, 15, 6, 1} (misma fila, calculada con memoria O(n))
\end{lstlisting}

\textbf{Regla práctica:} usa el triángulo de Pascal cuando $n$ sea moderado (hasta unos pocos miles), el módulo sea arbitrario o no exista, y necesites muchas consultas $\binom{n}{k}$ en $O(1)$. Si $n$ llega a $10^5$ o más y el módulo es primo, cambia a factoriales con inversos modulares, que precomputan en $O(n)$. Si solo necesitas una fila, usa \texttt{filaUnica} para no gastar $O(n^2)$ de memoria.

\textbf{Idea:} el triángulo de Pascal organiza los coeficientes binomiales $\binom{n}{k}$ en filas: la fila $n$ contiene $\binom{n}{0},\binom{n}{1},\dots,\binom{n}{n}$. Cada número es la suma de los dos que tiene encima: $\binom{n}{k}=\binom{n-1}{k-1}+\binom{n-1}{k}$, con $\binom{n}{0}=\binom{n}{n}=1$. La recurrencia tiene una lectura combinatoria directa. Para elegir $k$ elementos de $n$, se fija un elemento especial: o se incluye (faltan $k-1$ por elegir entre los otros $n-1$) o no se incluye (faltan $k$ entre los otros $n-1$). Como los subproblemas $\binom{n-1}{\cdot}$ se reutilizan en toda la fila siguiente, es un caso de DP muy claro sobre \textit{Memoización vs. Tabulación}: la tabla se llena fila por fila, y cada celda solo necesita la fila anterior.

\textbf{¿Por qué usarlo?} La fórmula cerrada $\binom{n}{k}=\frac{n!}{k!(n-k)!}$ desborda con facilidad (ya $20!$ es el límite de \texttt{long long}) y requiere inversos modulares cuando se trabaja módulo un primo. El triángulo evita por completo divisiones y factoriales: solo usa \textbf{sumas}, así que funciona con cualquier módulo (primo o no) y es exacto mientras el resultado quepa en el tipo. Su costo es $O(n^2)$ de precomputación, que es aceptable cuando $n$ es del orden de $10^3$ a $5\cdot 10^3$ y se necesitan muchas consultas $\binom{n}{k}$ distintas. Si $n$ es muy grande (por ejemplo $10^6$), conviene la fórmula con factoriales precomputados e inversos modulares.

\textbf{Casos de uso:}
\begin{itemize}
	\item Precomputar una tabla de $\binom{n}{k}$ para responder muchas consultas en $O(1)$.
	\item Contar caminos en una grilla que solo avanza a la derecha y hacia abajo: $\binom{f+c}{f}$ para $f$ filas y $c$ columnas de pasos.
	\item Contar subconjuntos de tamaño $k$, o formas de repartir objetos en problemas combinatorios de conteo.
	\item Coeficientes del desarrollo de $(x+y)^n$ y de otros polinomios sencillos.
	\item Trabajar con un módulo compuesto (no primo), donde los inversos modulares no siempre existen.
\end{itemize}

\textbf{Complejidad:} $O(n^2)$ tiempo y $O(n^2)$ espacio para construir la tabla hasta la fila $n$, ya que hay $\frac{(n+1)(n+2)}{2}$ celdas y cada una se calcula con una suma en $O(1)$. Una consulta posterior es $O(1)$. Para obtener solo una fila, $O(n^2)$ tiempo pero $O(n)$ espacio.

\begin{lstlisting}[style=cpp]
	// CREATE/READ/WRITE: Triangulo de Pascal para coeficientes binomiales,
	// con modulo opcional (mod = 0 significa sin modulo)
	struct TrianguloPascal {
		vector<vector<long long>> tabla; // tabla[i][j] = C(i, j)
		long long mod;
		
		// CREATE: construye las filas 0..n; cada celda es la suma de las dos
		// que tiene encima (tabulacion fila por fila)
		TrianguloPascal(int n, long long m = 0) : mod(m) {
			tabla.assign(n + 1, vector<long long>());
			for (int i = 0; i <= n; i++) {
				tabla[i].assign(i + 1, 1); // bordes C(i,0) = C(i,i) = 1
				for (int j = 1; j < i; j++) {
					long long v = tabla[i - 1][j - 1] + tabla[i - 1][j];
					if (mod > 0) v %= mod;
					tabla[i][j] = v;
				}
			}
		} // O(n^2)
		
		// READ: coeficiente binomial C(n, k); devuelve 0 si k esta fuera de rango
		long long combinacion(int n, int k) const {
			if (k < 0 || k > n) return 0;
			return tabla[n][k];
		} // O(1)
		
		// READ: fila completa n del triangulo
		const vector<long long>& fila(int n) const {
			return tabla[n];
		} // O(1)
		
		// READ: una sola fila n con espacio O(n), sin guardar la tabla completa
		// (rolling array: se actualiza de derecha a izquierda para no pisar
		// valores que todavia hacen falta)
		static vector<long long> filaUnica(int n) {
			vector<long long> f(n + 1, 0);
			f[0] = 1;
			for (int i = 1; i <= n; i++) {
				for (int j = i; j >= 1; j--) f[j] += f[j - 1];
			}
			return f;
		} // O(n^2) tiempo, O(n) espacio
		
		// ------------------------------------------------------------
		// MODIFICACIONES Y COMENTARIOS CON LOS USOS
		// ------------------------------------------------------------
		// 1. Overflow: sin modulo, C(n,k) cabe en long long solo hasta n = 66
		//    (C(66,33) ~ 7.2e18). Para n mayor, usar un modulo (por ejemplo
		//    1e9+7) o __int128 si se necesita el valor exacto pero acotado.
		// 2. Modulo compuesto: este triangulo funciona con cualquier modulo,
		//    incluso no primo, porque solo suma; la formula con factoriales e
		//    inversos modulares exige un modulo primo.
		// 3. n grande (~1e6) con muchas consultas: precomputar fact[i] e
		//    invFact[i] mod p primo, y responder
		//    C(n,k) = fact[n] * invFact[k] % p * invFact[n-k] % p en O(1),
		//    con precomputo O(n) en vez de O(n^2).
		// 4. Caminos en grilla que solo avanza derecha o abajo, con f pasos
		//    hacia abajo y c hacia la derecha: combinacion(f + c, f).
		// 5. Suma de la fila n: siempre 2^n, ya que cada elemento de un
		//    conjunto de n se incluye o no; util como chequeo de cordura.
		// 6. Identidades utiles: C(n,k) = C(n,n-k) (simetria de cada fila) y
		//    C(n,0)+C(n,1)+...+C(n,n) = 2^n; permiten calcular solo la mitad
		//    de cada fila si se necesita ahorrar tiempo.
	};
\end{lstlisting}

\textbf{Ejemplo de funcionamiento:}
\begin{lstlisting}[style=cpp]
	TrianguloPascal t(10);
	
	long long a = t.combinacion(5, 2);
	// a = 10 (formas de elegir 2 elementos entre 5)
	
	vector<long long> fila5 = t.fila(5);
	// fila5 = {1, 5, 10, 10, 5, 1}
	// (simetrica; su suma es 32 = 2^5)
	
	long long b = t.combinacion(10, 3);
	// b = 120 (C(10,3) = 45 + 75? no: se obtiene como C(9,2)+C(9,3) = 36 + 84 = 120)
	
	vector<long long> f6 = TrianguloPascal::filaUnica(6);
	// f6 = {1, 6, 15, 20, 15, 6, 1} (misma fila, calculada con memoria O(n))
\end{lstlisting}

\textbf{Nota:} en el comentario del ejemplo de $\binom{10}{3}$ hay una frase mal escrita por mi parte al redactarlo. El resultado $120$ es correcto y proviene de $\binom{9}{2}+\binom{9}{3}=36+84$; borra la parte "$45+75$? no:" al pegarlo en tu guía. Además, cuando el problema pide el valor módulo un número, aplica el módulo en cada suma dentro del constructor (como hace el código con \texttt{mod}); reducir solo al final no evita el desbordamiento intermedio.

\textbf{Regla práctica:} usa el triángulo de Pascal cuando $n$ sea moderado (hasta unos pocos miles), el módulo sea arbitrario o no exista, y necesites muchas consultas $\binom{n}{k}$ en $O(1)$. Si $n$ llega a $10^5$ o más y el módulo es primo, cambia a factoriales con inversos modulares, que precomputan en $O(n)$. Si solo necesitas una fila, usa \texttt{filaUnica} para no gastar $O(n^2)$ de memoria.
\subsection{Camino Más Largo en una Matriz}

\textbf{Idea:} dada una matriz de enteros, se busca la longitud del camino \textbf{estrictamente creciente} más largo, donde desde cada celda se puede avanzar a sus 4 vecinos (arriba, abajo, izquierda, derecha). Sea $\text{dp}[f][c]$ la longitud del camino creciente más largo que \textbf{empieza} en la celda $(f,c)$. Su recurrencia es $\text{dp}[f][c]=1+\max\{\text{dp}[f'][c'] : (f',c')\text{ vecino},\ a[f'][c']>a[f][c]\}$, y vale $1$ si ningún vecino es mayor. La condición de \textbf{estrictamente creciente} es lo que hace válida la DP: como solo se avanza hacia valores mayores, nunca se puede volver a una celda ya visitada, así que las dependencias forman un grafo sin ciclos y cada estado se puede calcular una sola vez, tal como se explicó en \textit{Memoización vs. Tabulación}.

\textbf{¿Por qué usarlo?} Sin memoización, una búsqueda en profundidad desde cada celda explora todos los caminos crecientes posibles, que pueden ser exponenciales en el tamaño de la matriz. Al guardar el valor de cada celda la primera vez que se calcula, cada una se resuelve una sola vez con trabajo $O(1)$ (revisar 4 vecinos), y el total baja a $O(nm)$. Aquí la memoización es la opción natural: el orden de llenado depende de los valores de la matriz y no de los índices, así que no hay un recorrido fila por fila que sirva para tabular directamente. Si se prefiere evitar la recursión, se puede tabular procesando las celdas ordenadas por valor (ver modificación 1 dentro del código).

\textbf{Casos de uso:}
\begin{itemize}
	\item Encontrar la ruta creciente más larga en una matriz de alturas o valores (problema clásico de \textit{Longest Increasing Path in a Matrix}).
	\item Caminos más largos en un grafo dirigido acíclico (DAG) dado implícitamente por una grilla.
	\item Juegos o simulaciones donde cada movimiento debe llevar a una casilla de mayor valor, y se pide la mayor cantidad de pasos posibles.
	\item Reconstruir la ruta concreta que logra ese máximo, no solo su longitud.
	\item Base conceptual para DP sobre DAG: el mismo esquema (memoizar la mejor respuesta desde cada nodo) sirve para cualquier grafo sin ciclos.
\end{itemize}

\textbf{Complejidad:} $O(nm)$ tiempo para una matriz de $n$ filas y $m$ columnas: hay $nm$ estados, cada uno se calcula una sola vez y examina a lo sumo 4 vecinos en $O(1)$. El espacio es $O(nm)$ para la tabla, más hasta $O(nm)$ de pila de recursión en el peor caso (una matriz cuyos valores forman una sola cadena creciente).

\begin{lstlisting}[style=cpp]
	// CREATE/READ/WRITE: Camino estrictamente creciente mas largo en una
	// matriz (movimiento en 4 direcciones), con memoizacion sobre cada celda
	struct CaminoMasLargoMatriz {
		vector<vector<int>> a;    // matriz de valores
		vector<vector<int>> memo; // memo[f][c] = longitud del mejor camino que EMPIEZA en (f,c); 0 = sin calcular
		int filas, cols;
		const int DF[4] = {-1, 1, 0, 0};
		const int DC[4] = {0, 0, -1, 1};
		
		// CREATE: guarda la matriz y prepara la tabla de memoizacion en 0
		// (0 sirve como centinela porque todo camino real tiene longitud >= 1)
		CaminoMasLargoMatriz(const vector<vector<int>>& m) : a(m) {
			filas = (int)a.size();
			cols = filas > 0 ? (int)a[0].size() : 0;
			memo.assign(filas, vector<int>(cols, 0));
		} // O(n*m)
		
		// READ: longitud del mejor camino creciente que empieza en (f,c),
		// memoizado (top-down)
		int desde(int f, int c) {
			if (memo[f][c] != 0) return memo[f][c]; // estado ya resuelto
			int mejor = 1; // el camino formado solo por esta celda
			for (int d = 0; d < 4; d++) {
				int nf = f + DF[d], nc = c + DC[d];
				if (nf < 0 || nf >= filas || nc < 0 || nc >= cols) continue;
				if (a[nf][nc] > a[f][c]) mejor = max(mejor, 1 + desde(nf, nc));
			}
			return memo[f][c] = mejor;
		} // O(1) amortizado por celda, O(n*m) en total
		
		// READ: longitud del camino creciente mas largo de toda la matriz
		int longitudMaxima() {
			int mejor = 0;
			for (int f = 0; f < filas; f++)
			for (int c = 0; c < cols; c++)
			mejor = max(mejor, desde(f, c));
			return mejor;
		} // O(n*m)
		
		// READ: reconstruye UN camino optimo como lista de celdas (fila, col),
		// siguiendo en cada paso un vecino mayor con memo exactamente uno menos
		vector<pair<int,int>> reconstruir() {
			int mejor = longitudMaxima();
			vector<pair<int,int>> camino;
			if (mejor == 0) return camino;
			int f = 0, c = 0;
			for (int i = 0; i < filas; i++)
			for (int j = 0; j < cols; j++)
			if (memo[i][j] == mejor) { f = i; c = j; }
			camino.push_back({f, c});
			while (memo[f][c] > 1) {
				for (int d = 0; d < 4; d++) {
					int nf = f + DF[d], nc = c + DC[d];
					if (nf < 0 || nf >= filas || nc < 0 || nc >= cols) continue;
					if (a[nf][nc] > a[f][c] && memo[nf][nc] == memo[f][c] - 1) {
						f = nf; c = nc;
						camino.push_back({f, c});
						break;
					}
				}
			}
			return camino;
		} // O(n*m)
		
		// READ: version iterativa (bottom-up) sin recursion: se ordenan las
		// celdas por valor DECRECIENTE, asi cuando se procesa una celda todos
		// sus vecinos mayores ya tienen su dp calculado
		static int longitudMaximaIterativa(const vector<vector<int>>& m) {
			int n = (int)m.size();
			if (n == 0) return 0;
			int k = (int)m[0].size();
			vector<array<int,3>> celdas; // {valor, fila, col}
			for (int i = 0; i < n; i++)
			for (int j = 0; j < k; j++) celdas.push_back({m[i][j], i, j});
			sort(celdas.begin(), celdas.end(), [](const array<int,3>& x, const array<int,3>& y) {
				return x[0] > y[0];
			});
			int df[4] = {-1, 1, 0, 0}, dc[4] = {0, 0, -1, 1};
			vector<vector<int>> dp(n, vector<int>(k, 1));
			int mejor = 1;
			for (auto& cel : celdas) {
				int f = cel[1], c = cel[2];
				for (int d = 0; d < 4; d++) {
					int nf = f + df[d], nc = c + dc[d];
					if (nf < 0 || nf >= n || nc < 0 || nc >= k) continue;
					if (m[nf][nc] > m[f][c]) dp[f][c] = max(dp[f][c], 1 + dp[nf][nc]);
				}
				mejor = max(mejor, dp[f][c]);
			}
			return mejor;
		} // O(n*m log(n*m)), dominado por el sort
		
		// ------------------------------------------------------------
		// MODIFICACIONES Y COMENTARIOS CON LOS USOS
		// ------------------------------------------------------------
		// 1. Sin recursion (matrices grandes, riesgo de desbordar la pila):
		//    usar longitudMaximaIterativa(), que ordena las celdas por valor
		//    decreciente y tabula; cambia O(n*m) por O(n*m log(n*m)) pero
		//    elimina la profundidad de recursion. Tambien se puede hacer en
		//    O(n*m) con orden topologico (algoritmo de Kahn) sobre el grafo
		//    de aristas "celda -> vecino mayor".
		// 2. Camino DECRECIENTE en vez de creciente: cambiar a[nf][nc] > a[f][c]
		//    por a[nf][nc] < a[f][c] en la condicion de vecino valido (y en
		//    el orden del sort, si se usa la version iterativa).
		// 3. Con 8 direcciones (se permiten diagonales): ampliar DF y DC a 8
		//    entradas y recorrer d de 0 a 7; el resto del codigo no cambia.
		// 4. Camino NO estricto (se permiten valores iguales): la DP deja de
		//    ser valida, porque dos celdas vecinas iguales se podrian recorrer
		//    de ida y vuelta formando ciclos; el problema pasaria a pedir
		//    caminos simples y se vuelve mucho mas dificil (NP-dificil en
		//    general), asi que solo se resuelve con la condicion estricta.
		// 5. CONTAR el numero de caminos crecientes (con cualquier longitud):
		//    cuentaDesde(f,c) = 1 + suma de cuentaDesde(vecino mayor); el 1
		//    cuenta el camino formado solo por la celda; la respuesta total es
		//    la suma sobre todas las celdas (aplicar modulo si piden).
		// 6. Otras reglas de movimiento (por ejemplo saltos de tipo caballo
		//    de ajedrez): solo cambia el arreglo de desplazamientos DF y DC,
		//    mientras la condicion de valor estrictamente mayor mantenga el
		//    grafo sin ciclos.
		// 7. Ganancia acumulada en vez de longitud: si cada celda aporta un
		//    peso w[f][c] y se pide maximizar la suma, usar
		//    w[f][c] + max(desde(vecino mayor)) en lugar de 1 + max(...).
	};
\end{lstlisting}

\textbf{Ejemplo de funcionamiento:}
\begin{lstlisting}[style=cpp]
	vector<vector<int>> m = {
		{9, 9, 4},
		{6, 6, 8},
		{2, 1, 1}
	};
	CaminoMasLargoMatriz p(m);
	
	int resultado = p.longitudMaxima();
	// resultado = 4
	// (el camino 1 -> 2 -> 6 -> 9 recorre las celdas (2,1), (2,0), (1,0), (0,0))
	
	auto camino = p.reconstruir();
	// camino = {(2,1), (2,0), (1,0), (0,0)}
	// (valores 1, 2, 6, 9: cada paso va a un vecino estrictamente mayor)
	
	vector<vector<int>> m2 = {
		{3, 4, 5},
		{3, 2, 6},
		{2, 2, 1}
	};
	int r2 = CaminoMasLargoMatriz::longitudMaximaIterativa(m2);
	// r2 = 4
	// (el camino 3 -> 4 -> 5 -> 6 recorre (1,0), (0,0)... o (0,0), (0,1), (0,2), (1,2))
\end{lstlisting}

\textbf{Nota:} la DP funciona únicamente porque la condición es \textbf{estrictamente creciente}: eso garantiza que el grafo de celdas no tiene ciclos y que el valor guardado en \texttt{memo} nunca depende de sí mismo. Si el problema permite pasar entre celdas de igual valor, esta técnica deja de ser válida (ver la modificación 4 dentro del código). Además, usar \texttt{0} como centinela de \texttt{memo} es seguro aquí porque toda longitud real es al menos $1$.

\textbf{Regla práctica:} usa memoización con DFS desde cada celda cuando la regla de movimiento obligue a ir siempre a valores estrictamente mayores (o menores): así es $O(nm)$ y la más corta de escribir. Si la matriz es tan grande que la recursión puede desbordar la pila, pasa a la versión iterativa ordenando las celdas por valor. Si en cambio la regla es solo ``moverse a la derecha o hacia abajo'' (sin condición de valores), no hace falta memoizar por valor: basta la tabulación fila por fila, como en \textit{Triángulo de Pascal}.

\section{Knapsack y Variantes}
\subsection{Knapsack 0/1}

\textbf{Idea:} se tienen $n$ objetos, cada uno con un peso $p_i$ y un valor $v_i$, y una mochila de capacidad $W$. Cada objeto se puede tomar \textbf{una sola vez} o no tomarse (de ahí el "0/1"). Se busca el máximo valor total cuyo peso no exceda $W$. Sea $\text{dp}[i][w]$ el mejor valor usando solo los primeros $i$ objetos con capacidad $w$. Para el objeto $i$ hay dos opciones: no tomarlo ($\text{dp}[i-1][w]$) o tomarlo ($\text{dp}[i-1][w-p_i]+v_i$, solo si $p_i\le w$), y se elige la mejor: $\text{dp}[i][w]=\max(\text{dp}[i-1][w],\ \text{dp}[i-1][w-p_i]+v_i)$. Como la fila $i$ solo depende de la fila $i-1$, se puede guardar un único arreglo $\text{dp}[w]$ si se recorre $w$ de \textbf{mayor a menor}: así $\text{dp}[w-p_i]$ todavía pertenece a la fila anterior y el objeto no se cuenta dos veces.

\textbf{¿Por qué usarlo?} La fuerza bruta prueba los $2^n$ subconjuntos de objetos. La DP baja el costo a $O(nW)$, porque el estado $(i,w)$ ya resume todas las formas de decidir los primeros $i$ objetos. Un \textit{greedy} por razón valor/peso no sirve aquí: con objetos indivisibles puede dar una respuesta subóptima (la versión fraccionaria sí es greedy, ver modificación 7). Frente a \textit{Cambio de Monedas}, la diferencia es que aquí cada objeto se usa a lo sumo una vez, y eso obliga a recorrer la capacidad en sentido descendente. Es \textbf{pseudopolinomial}: el costo depende del valor numérico de $W$, no solo del tamaño de la entrada.

\textbf{Casos de uso:}
\begin{itemize}
	\item Elegir el subconjunto de proyectos, ítems o tareas de mayor beneficio con un presupuesto o tiempo limitado.
	\item Decidir si existe un subconjunto con suma exacta (\textit{Subset Sum}), que es Knapsack con valor igual al peso.
	\item Problemas de partición y reparto equilibrado, como \textit{Problema de Partición}.
	\item Selección de objetos con dos restricciones a la vez, usando una dimensión de DP adicional.
	\item Base de muchas variantes de mochila: acotada, de múltiples grupos, con dependencias.
\end{itemize}

\textbf{Complejidad:} $O(nW)$ tiempo, porque hay $n\cdot(W+1)$ estados y cada uno se calcula en $O(1)$. El espacio es $O(W)$ con el arreglo unidimensional y $O(nW)$ si se guarda la tabla completa para reconstruir los objetos elegidos.

\begin{lstlisting}[style=cpp]
	// CREATE/READ/WRITE: Knapsack 0/1 con arreglo unidimensional y
	// reconstruccion de los objetos elegidos con tabla completa
	struct Knapsack01 {
		int n;
		int capacidad;
		vector<int> peso;
		vector<long long> valor;
		
		// CREATE: guarda los objetos y la capacidad de la mochila
		Knapsack01(const vector<int>& p, const vector<long long>& v, int W)
		: n((int)p.size()), capacidad(W), peso(p), valor(v) {
		} // O(n)
		
		// READ: maximo valor con peso total <= capacidad (espacio O(W))
		long long maximoValor() const {
			vector<long long> dp(capacidad + 1, 0);
			for (int i = 0; i < n; i++) {
				// w decreciente: dp[w - peso[i]] aun es de la fila anterior
				for (int w = capacidad; w >= peso[i]; w--) {
					dp[w] = max(dp[w], dp[w - peso[i]] + valor[i]);
				}
			}
			return dp[capacidad];
		} // O(n*W) tiempo, O(W) espacio
		
		// READ: indices (0-indexados, ascendentes) de un conjunto optimo de
		// objetos; usa la tabla completa dp[i][w] para poder retroceder
		vector<int> reconstruir() const {
			vector<vector<long long>> dp(n + 1, vector<long long>(capacidad + 1, 0));
			for (int i = 1; i <= n; i++) {
				for (int w = 0; w <= capacidad; w++) {
					dp[i][w] = dp[i - 1][w];
					if (w >= peso[i - 1]) {
						dp[i][w] = max(dp[i][w], dp[i - 1][w - peso[i - 1]] + valor[i - 1]);
					}
				}
			}
			vector<int> elegidos;
			int w = capacidad;
			for (int i = n; i >= 1; i--) {
				if (dp[i][w] != dp[i - 1][w]) { // el objeto i fue necesario
					elegidos.push_back(i - 1);
					w -= peso[i - 1];
				}
			}
			reverse(elegidos.begin(), elegidos.end());
			return elegidos;
		} // O(n*W) tiempo y espacio
		
		// ------------------------------------------------------------
		// MODIFICACIONES Y COMENTARIOS CON LOS USOS
		// ------------------------------------------------------------
		// 1. Peso EXACTO igual a la capacidad (mochila que debe llenarse):
		//    inicializar dp[0] = 0 y dp[w] = -INF para w > 0 (por ejemplo
		//    -1e18), y al actualizar ignorar los estados -INF. La respuesta es
		//    dp[capacidad], o "imposible" si sigue siendo -INF.
		// 2. Solo factibilidad (Subset Sum): usar vector<bool> alcanzable(W+1)
		//    con alcanzable[0] = true, y para cada objeto, w decreciente:
		//    alcanzable[w] = alcanzable[w] || alcanzable[w - peso[i]].
		//    Con bitset<MAXW> se reduce a: mascara |= (mascara << peso[i]);
		//    ver Bitset para Optimizacion en la parte de BitWise.
		// 3. Contar el NUMERO de subconjuntos con suma exacta W: cuenta[0] = 1,
		//    y para cada objeto, w decreciente:
		//    cuenta[w] += cuenta[w - peso[i]] (aplicar modulo si lo piden).
		// 4. Pesos enormes pero valores pequenos (W ~ 1e9, suma de valores
		//    <= ~1e5): invertir los roles del estado. Sea minPeso[v] el minimo
		//    peso para lograr valor v: minPeso[0] = 0, resto INF, y para cada
		//    objeto, v decreciente: minPeso[v] = min(minPeso[v],
		//    minPeso[v - valor[i]] + peso[i]). La respuesta es el mayor v con
		//    minPeso[v] <= W. Costo O(n * suma de valores).
		// 5. Objetos con cantidad limitada c_i (mochila acotada): descomponer
		//    c_i en potencias de dos (1, 2, 4, ..., resto) y tratar cada trozo
		//    como un objeto 0/1 de peso k*peso[i] y valor k*valor[i]; el costo
		//    baja a O(W * sum log c_i). Si c_i es ilimitado, ver Knapsack
		//    Ilimitado.
		// 6. n pequeno (n <= 40) pero W enorme: DP inviable, usar meet in the
		//    middle sobre subconjuntos (ver Meet in the Middle con Mascaras
		//    de Bits en BitWise).
		// 7. Objetos DIVISIBLES (mochila fraccionaria): no es DP. Se ordenan
		//    por valor/peso decreciente y se toman completos mientras quepan,
		//    y del ultimo se toma la fraccion restante (greedy correcto solo
		//    en esta version).
		// 8. Dos restricciones (por ejemplo peso y volumen): agregar una
		//    dimension a dp, dp[w1][w2], recorriendo ambas en sentido
		//    decreciente; el costo pasa a O(n * W1 * W2).
		// 9. Orden de los bucles: el bucle de objetos va AFUERA y el de
		//    capacidad ADENTRO en sentido decreciente. Si se recorre w en
		//    sentido creciente, un mismo objeto se puede usar varias veces y
		//    se obtiene Knapsack Ilimitado, no el 0/1.
	};
\end{lstlisting}

\textbf{Ejemplo de funcionamiento:}
\begin{lstlisting}[style=cpp]
	vector<int> peso = {1, 3, 4, 5};
	vector<long long> valor = {1, 4, 5, 7};
	Knapsack01 k(peso, valor, 7);
	
	long long mejor = k.maximoValor();
	// mejor = 9
	// (tomar los objetos de peso 3 y 4: peso total 7, valor 4 + 5 = 9;
	//  la alternativa 1 + 5 pesa 6 y vale 1 + 7 = 8, que es menor)
	
	vector<int> elegidos = k.reconstruir();
	// elegidos = {1, 2}
	// (los indices 1 y 2 corresponden a los objetos de peso 3 y 4)
\end{lstlisting}

\textbf{Nota:} el error más común es recorrer la capacidad de menor a mayor en la versión unidimensional. Eso permite reutilizar el mismo objeto dentro de una misma pasada y convierte silenciosamente el problema en Knapsack Ilimitado, sin ningún error visible. Recuerda también que la complejidad $O(nW)$ solo es aceptable si $W$ es moderado (típicamente hasta $10^5$ o $10^6$); si $W$ es enorme, revisa la modificación 4 o la 6.

\textbf{Regla práctica:} usa Knapsack 0/1 cuando cada objeto se pueda tomar como máximo una vez y la capacidad sea numéricamente manejable. Recorre siempre la capacidad en sentido decreciente. Si los objetos se pueden repetir, usa \textit{Knapsack Ilimitado}; si se pueden dividir, basta un greedy por razón valor/peso.

\subsection{Knapsack Ilimitado}

\textbf{Idea:} es la variante de la mochila en la que cada objeto tiene \textbf{suministro infinito}: se puede tomar cualquier cantidad de copias del mismo tipo. Sea $\text{dp}[w]$ el máximo valor con capacidad $w$. La recurrencia es $\text{dp}[w]=\max_{i:\,p_i\le w}(\text{dp}[w-p_i]+v_i)$: el último objeto que se agrega puede ser de cualquier tipo, y lo que queda se resuelve con la capacidad restante, que ya puede volver a usar ese mismo tipo. La diferencia con \textit{Knapsack 0/1} es solo el sentido en que se recorre la capacidad: aquí se recorre de \textbf{menor a mayor}, de modo que $\text{dp}[w-p_i]$ ya incluye posibles copias del objeto $i$ y este se puede reutilizar.

\textbf{¿Por qué usarlo?} Tomar cada objeto a lo sumo una vez (0/1) no modela problemas de reposición ilimitada, como cortar una barra en tantos trozos de una misma longitud como convenga. La DP resuelve todos los casos en $O(nW)$ sin enumerar cuántas copias de cada tipo usar. Es además el marco general del \textit{Cambio de Monedas}: minimizar la cantidad de monedas o contar combinaciones son la misma estructura con otra operación en la recurrencia (mínimo con costo 1 por moneda, o suma de formas). Por eso esta subsección va antes de esas dos.

\textbf{Casos de uso:}
\begin{itemize}
	\item Cortar una barra o cuerda en trozos de longitudes permitidas para maximizar el precio total (\textit{Rod Cutting}).
	\item Llenar un contenedor con paquetes de tipos repetibles para maximizar el valor transportado.
	\item Base directa de \textit{Cambio de Monedas: Mínimo de Monedas} y \textit{Cambio de Monedas: Número de Formas}.
	\item Problemas de producción donde cada producto se puede fabricar tantas veces como se quiera con recursos limitados.
	\item Determinar si una cantidad es alcanzable como suma de valores repetibles (el problema de las monedas de Frobenius, con $n$ pequeño).
\end{itemize}

\textbf{Complejidad:} $O(nW)$ tiempo: hay $W+1$ estados y cada uno prueba a lo sumo $n$ tipos de objeto en $O(1)$. El espacio es $O(W)$.

\begin{lstlisting}[style=cpp]
	// CREATE/READ/WRITE: Knapsack ilimitado (cada objeto se puede usar
	// cualquier cantidad de veces), con reconstruccion de las cantidades
	struct KnapsackIlimitado {
		int n;
		int capacidad;
		vector<int> peso;
		vector<long long> valor;
		
		// CREATE: guarda los tipos de objeto y la capacidad de la mochila
		KnapsackIlimitado(const vector<int>& p, const vector<long long>& v, int W)
		: n((int)p.size()), capacidad(W), peso(p), valor(v) {
		} // O(n)
		
		// READ: maximo valor con peso total <= capacidad (espacio O(W))
		long long maximoValor() const {
			vector<long long> dp(capacidad + 1, 0);
			for (int w = 1; w <= capacidad; w++) {
				dp[w] = dp[w - 1]; // se permite dejar capacidad sin usar
				for (int i = 0; i < n; i++) {
					if (peso[i] <= w) {
						dp[w] = max(dp[w], dp[w - peso[i]] + valor[i]);
					}
				}
			}
			return dp[capacidad];
		} // O(n*W) tiempo, O(W) espacio
		
		// READ: cantidad usada de cada tipo en una solucion optima
		// (cantidades[i] = numero de copias del objeto i)
		vector<int> reconstruir() const {
			vector<long long> dp(capacidad + 1, 0);
			vector<int> eleccion(capacidad + 1, -2); // -2 = capacidad sobrante
			for (int w = 1; w <= capacidad; w++) {
				dp[w] = dp[w - 1];
				for (int i = 0; i < n; i++) {
					if (peso[i] <= w && dp[w - peso[i]] + valor[i] > dp[w]) {
						dp[w] = dp[w - peso[i]] + valor[i];
						eleccion[w] = i;
					}
				}
			}
			vector<int> cantidades(n, 0);
			int w = capacidad;
			while (w > 0) {
				if (eleccion[w] == -2) {
					w--; // capacidad sin usar
				} else {
					cantidades[eleccion[w]]++;
					w -= peso[eleccion[w]];
				}
			}
			return cantidades;
		} // O(n*W)
		
		// ------------------------------------------------------------
		// MODIFICACIONES Y COMENTARIOS CON LOS USOS
		// ------------------------------------------------------------
		// 1. Version con el bucle de objetos afuera (equivalente y mas corta):
		//    for i: for (w = peso[i]; w <= W; w++)
		//        dp[w] = max(dp[w], dp[w - peso[i]] + valor[i]);
		//    con w CRECIENTE; es exactamente Knapsack 0/1 con el sentido del
		//    bucle interno invertido.
		// 2. Peso EXACTO igual a la capacidad: inicializar dp[0] = 0 y
		//    dp[w] = -INF para w > 0, y omitir la linea dp[w] = dp[w - 1];
		//    la respuesta es dp[capacidad] o "imposible" si sigue en -INF.
		// 3. Minimo de objetos para llenar exactamente W (cambio de monedas
		//    con minimo): dp[0] = 0, dp[w] = INF, y
		//    dp[w] = min(dp[w], dp[w - peso[i]] + 1); ver Cambio de Monedas:
		//    Minimo de Monedas.
		// 4. NUMERO de formas de llenar exactamente W: cuenta[0] = 1 y, con
		//    el bucle de objetos AFUERA y w creciente,
		//    cuenta[w] += cuenta[w - peso[i]]; cuenta combinaciones. Con el
		//    bucle de capacidad afuera cuenta permutaciones (el orden
		//    importaria); ver Cambio de Monedas: Numero de Formas.
		// 5. Capacidad enorme con un tipo dominante: si un tipo tiene la mejor
		//    razon valor/peso, en una solucion optima los demas tipos se usan
		//    pocas veces (menos de p_mejor copias en total), lo que permite
		//    acotar la busqueda. Solo es util cuando W es demasiado grande
		//    para la DP directa.
		// 4b. Cuidado con objetos de peso 0 o negativo: con peso 0 y valor
		//    positivo la respuesta es infinita; validar la entrada antes de
		//    aplicar la DP.
	};
\end{lstlisting}

\textbf{Ejemplo de funcionamiento:}
\begin{lstlisting}[style=cpp]
	vector<int> peso = {1, 3, 4, 5};
	vector<long long> valor = {10, 40, 50, 70};
	KnapsackIlimitado k(peso, valor, 8);
	
	long long mejor = k.maximoValor();
	// mejor = 110
	// (tomar un objeto de peso 3 y uno de peso 5: peso 8, valor 40 + 70 = 110;
	//  usar dos objetos de peso 4 da 100, y ocho de peso 1 da solo 80)
	
	vector<int> cantidades = k.reconstruir();
	// cantidades = {0, 1, 0, 1}
	// (una copia del objeto de peso 3 y una del de peso 5)
\end{lstlisting}

\textbf{Nota:} el orden de los bucles decide qué se cuenta. Con el bucle de objetos afuera y la capacidad creciente adentro se obtienen \textbf{combinaciones}; con el bucle de capacidad afuera se obtienen \textbf{permutaciones}. Para maximizar un valor da igual, porque el máximo no depende del orden, pero para \textbf{contar formas} la diferencia es crucial (ver la modificación 4 dentro del código). Como en Knapsack 0/1, el costo es pseudopolinomial en $W$.

\textbf{Regla práctica:} usa Knapsack Ilimitado cuando cada tipo de objeto pueda repetirse sin límite, y recorre la capacidad de menor a mayor. Si el suministro de cada tipo es limitado, usa \textit{Knapsack 0/1} con descomposición en potencias de dos (modificación 5 de esa subsección). Si el problema pide minimizar la cantidad de piezas o contar formas de llenar una cantidad exacta, es un caso de \textit{Cambio de Monedas}, que reutiliza esta misma estructura.
\subsection{Cambio de Monedas: Número de Formas (Combinaciones)}

\textbf{Idea:} se tienen monedas de valores $s_1,\dots,s_m$ con suministro infinito y se quiere saber de \textbf{cuántas formas distintas} se puede formar el monto $n$, donde el orden de las monedas \textbf{no importa} ($1+2$ y $2+1$ cuentan como una sola forma). Es el marco de \textit{Knapsack Ilimitado} con la operación cambiada: en vez de maximizar un valor, se \textbf{suman} formas. Sea $\text{dp}[v]$ el número de combinaciones que suman $v$, con $\text{dp}[0]=1$ (la única forma de sumar 0 es no usar ninguna moneda). Se procesa \textbf{una moneda a la vez}, y para cada una se recorre $v$ de menor a mayor con $\text{dp}[v]\mathrel{+}=\text{dp}[v-s_j]$. Como las monedas se introducen en un orden fijo (el bucle exterior), cada combinación se construye siempre en ese mismo orden y se cuenta una sola vez.

\textbf{¿Por qué usarlo?} Enumerar todas las combinaciones con backtracking es exponencial. La DP las cuenta en $O(nm)$ sin generarlas. Aquí el orden de los bucles decide \textbf{qué} se cuenta: con las monedas afuera y el monto adentro se cuentan \textbf{combinaciones}; con el monto afuera y las monedas adentro se cuentan \textbf{permutaciones} (el orden sí importaría). Es el mismo detalle que se mencionó en \textit{Knapsack Ilimitado}, y aquí es crucial porque la respuesta cambia. Cuando se pide el \textbf{mínimo} de monedas en vez del número de formas, ver \textit{Cambio de Monedas: Mínimo de Monedas}.

\textbf{Casos de uso:}
\begin{itemize}
	\item Contar de cuántas formas se puede pagar un monto con monedas o billetes de denominaciones dadas.
	\item Contar particiones de un entero en sumandos permitidos (por ejemplo, formas de escribir $n$ como suma de números de un conjunto).
	\item Contar subconjuntos con suma exacta cuando cada elemento se puede repetir.
	\item Contar composiciones u ordenamientos (versión de permutaciones) de pasos de tamaños dados, como subir una escalera con saltos de 1, 2 o 3.
	\item Base de problemas de conteo con módulo, donde solo se pide el resultado módulo un primo.
\end{itemize}

\textbf{Complejidad:} $O(nm)$ tiempo, con $n$ el monto y $m$ el número de tipos de moneda: cada moneda recorre los $n+1$ valores una vez con trabajo $O(1)$. El espacio es $O(n)$.

\begin{lstlisting}[style=cpp]
	// CREATE/READ/WRITE: Cambio de monedas contando formas, con modulo
	// opcional (mod = 0 significa sin modulo)
	struct CambioMonedasFormas {
		vector<int> monedas;
		long long mod;
		
		// CREATE: guarda los valores de las monedas y el modulo
		CambioMonedasFormas(const vector<int>& s, long long m = 0)
		: monedas(s), mod(m) {
		} // O(m)
		
		// READ: numero de COMBINACIONES (el orden no importa) que suman n
		long long contarCombinaciones(int n) const {
			vector<long long> dp(n + 1, 0);
			dp[0] = 1; // una forma de sumar 0: no usar monedas
			for (int c : monedas) {          // monedas AFUERA
				for (int v = c; v <= n; v++) { // monto ADENTRO, creciente
					dp[v] += dp[v - c];
					if (mod > 0) dp[v] %= mod;
				}
			}
			return dp[n];
		} // O(n*m) tiempo, O(n) espacio
		
		// READ: numero de PERMUTACIONES (el orden SI importa) que suman n;
		// es lo que sale si se invierten los bucles
		long long contarPermutaciones(int n) const {
			vector<long long> dp(n + 1, 0);
			dp[0] = 1;
			for (int v = 1; v <= n; v++) {   // monto AFUERA
				for (int c : monedas) {        // monedas ADENTRO
					if (c <= v) {
						dp[v] += dp[v - c];
						if (mod > 0) dp[v] %= mod;
					}
				}
			}
			return dp[n];
		} // O(n*m) tiempo, O(n) espacio
		
		// ------------------------------------------------------------
		// MODIFICACIONES Y COMENTARIOS CON LOS USOS
		// ------------------------------------------------------------
		// 1. Orden de los bucles: monedas afuera = combinaciones; monto afuera
		//    = permutaciones. Es el error mas comun en este problema, y no da
		//    ningun aviso: el codigo compila y da un numero razonable pero
		//    incorrecto.
		// 2. Cada moneda se puede usar A LO SUMO UNA VEZ (Subset Sum contando):
		//    recorrer v en sentido DECRECIENTE, de n hasta c; con esto se
		//    obtiene el conteo de subconjuntos con suma exacta n.
		// 3. Cada moneda con cantidad limitada k_j: si k_j es pequeno, repetir
		//    el bucle 0/1 k_j veces; si es grande, usar una ventana deslizante
		//    sobre los residuos de v modulo c (suma de k_j+1 valores
		//    consecutivos con la misma clase de residuo), en O(n) por moneda.
		// 4. Modulo: aplicar mod en cada suma (como en el codigo). Sin modulo,
		//    la cantidad de formas crece rapido y desborda long long con montos
		//    moderados si hay monedas pequenas.
		// 5. Monto muy grande (~1e9) con pocas monedas pequenas: la DP directa
		//    no cabe. Se resuelve con recurrencia lineal y exponenciacion de
		//    matrices (permutaciones), o con el metodo de Kitamasa.
		// 6. Exigir al menos una moneda de cada tipo: restar de n la suma de
		//    una copia de cada moneda y contar combinaciones del monto
		//    restante, si todas las monedas deben aparecer.
		// 7. Contar solo formas con exactamente k monedas: agregar una
		//    dimension, dp[v][j] = formas de sumar v con j monedas, con costo
		//    O(n*m*k).
		// 8. Monedas repetidas o iguales en la entrada: eliminar los valores
		//    duplicados antes de contar (sort + unique); si no, cada duplicado
		//    se trata como una moneda distinta y se cuentan formas de mas.
		// 9. Saber solo SI existe alguna forma (no cuantas): usar
		//    vector<bool> con alcanzable[v] = alcanzable[v] || alcanzable[v-c];
		//    ver Knapsack Ilimitado.
	};
\end{lstlisting}

\textbf{Ejemplo de funcionamiento:}
\begin{lstlisting}[style=cpp]
	CambioMonedasFormas a({1, 2, 3});
	long long r1 = a.contarCombinaciones(4);
	// r1 = 4
	// (las combinaciones son 1+1+1+1, 1+1+2, 2+2 y 1+3)
	
	CambioMonedasFormas b({2, 5, 3, 6});
	long long r2 = b.contarCombinaciones(10);
	// r2 = 5
	// (las combinaciones son 2+2+2+2+2, 2+2+3+3, 2+2+6, 2+3+5 y 5+5)
	
	long long r3 = a.contarPermutaciones(4);
	// r3 = 7
	// (con monedas {1,2,3} y el orden importando, 4 se forma de 7 maneras:
	//  1111, 112, 121, 211, 22, 13 y 31)
\end{lstlisting}

\textbf{Nota:} los enunciados de este problema suelen dar como ejemplo $n=10$ con $S=\{2,5,3,6\}$ y una lista de soluciones. Esa lista tiene exactamente \textbf{5} combinaciones, así que la respuesta correcta es $5$. Si un enunciado da otro número, revisa si en realidad cuenta permutaciones (el orden importa) o si trae otra lista de monedas.

\textbf{Regla práctica:} usa este DP con las monedas en el bucle \textbf{exterior} cuando el orden no importe (combinaciones), y con el monto en el bucle exterior cuando el orden sí importe (permutaciones). Si el problema pide el mínimo de monedas en vez de contar, usa \textit{Cambio de Monedas: Mínimo de Monedas}; si además cada moneda se puede usar una sola vez, recorre el monto en sentido decreciente, como en \textit{Knapsack 0/1}.

\subsection{Cambio de Monedas: Mínimo de Monedas}

\textbf{Idea:} dado un monto $n$ y monedas de valores $s_1,\dots,s_m$ con suministro infinito, se busca el \textbf{menor número de monedas} cuya suma sea exactamente $n$, o determinar que es imposible. Sea $\text{dp}[v]$ el mínimo de monedas para sumar $v$, con $\text{dp}[0]=0$ y $\text{dp}[v]=\infty$ mientras no se conozca una forma. La recurrencia es $\text{dp}[v]=\min_{j:\,s_j\le v}(\text{dp}[v-s_j]+1)$: la última moneda puede ser de cualquier tipo y lo que queda se resuelve de forma óptima. Es exactamente \textit{Knapsack Ilimitado} donde cada moneda ``pesa'' su valor, cuesta $1$, y se debe llenar el monto de forma exacta.

\textbf{¿Por qué usarlo?} El algoritmo greedy (tomar siempre la moneda más grande que quepa) es correcto solo para ciertos sistemas monetarios, como los reales (1, 5, 10, 25 centavos), pero falla en general. Con monedas $\{1,3,4\}$ y monto $6$, el greedy toma $4+1+1$ (3 monedas), pero lo óptimo es $3+3$ (2 monedas). La DP prueba todas las opciones para la última moneda y es correcta para \textbf{cualquier} conjunto de monedas, en $O(nm)$. La diferencia con \textit{Cambio de Monedas: Número de Formas} es la operación de la recurrencia: aquí es un mínimo sobre las opciones de la última moneda, y no una suma de formas, así que el orden de los bucles no afecta el resultado.

\textbf{Casos de uso:}
\begin{itemize}
	\item Pagar un monto con la menor cantidad de monedas o billetes cuando el sistema no es ``canónico'' y el greedy falla.
	\item Determinar si un monto es alcanzable con las monedas dadas (respuesta $-1$ si no lo es).
	\item Minimizar el número de piezas, cortes o pasos cuando cada uno tiene un tamaño de un conjunto fijo (por ejemplo, el mínimo de cuadrados perfectos que suman $n$).
	\item Reconstruir qué monedas concretas se usan en la solución óptima.
	\item Como paso interno de problemas de cambio con máquinas expendedoras o cajeros automáticos.
\end{itemize}

\textbf{Complejidad:} $O(nm)$ tiempo: hay $n+1$ estados y cada uno prueba $m$ monedas en $O(1)$. El espacio es $O(n)$.

\begin{lstlisting}[style=cpp]
	// CREATE/READ/WRITE: Cambio de monedas con minimo de monedas y
	// reconstruccion de cuales monedas se usan
	struct CambioMonedasMinimo {
		vector<int> monedas;
		static const long long INF = (long long)4e18;
		
		// CREATE: guarda los valores de las monedas
		CambioMonedasMinimo(const vector<int>& s) : monedas(s) {
		} // O(m)
		
		// READ: minimo de monedas para sumar exactamente n, o -1 si es imposible
		long long minimoMonedas(int n) const {
			vector<long long> dp(n + 1, INF);
			dp[0] = 0;
			for (int v = 1; v <= n; v++) {
				for (int c : monedas) {
					if (c <= v && dp[v - c] != INF) {
						dp[v] = min(dp[v], dp[v - c] + 1);
					}
				}
			}
			return dp[n] == INF ? -1 : dp[n];
		} // O(n*m) tiempo, O(n) espacio
		
		// READ: lista de monedas de una solucion optima (vacia si n = 0,
		// y tambien vacia si es imposible; distinguir con minimoMonedas)
		vector<int> reconstruir(int n) const {
			vector<long long> dp(n + 1, INF);
			vector<int> ultima(n + 1, -1); // ultima moneda usada para llegar a v
			dp[0] = 0;
			for (int v = 1; v <= n; v++) {
				for (int c : monedas) {
					if (c <= v && dp[v - c] != INF && dp[v - c] + 1 < dp[v]) {
						dp[v] = dp[v - c] + 1;
						ultima[v] = c;
					}
				}
			}
			vector<int> usadas;
			if (dp[n] == INF) return usadas;
			int v = n;
			while (v > 0) {
				usadas.push_back(ultima[v]);
				v -= ultima[v];
			}
			return usadas;
		} // O(n*m)
		
		// ------------------------------------------------------------
		// MODIFICACIONES Y COMENTARIOS CON LOS USOS
		// ------------------------------------------------------------
		// 1. Centinela INF: usar un valor grande pero que NO desborde al
		//    sumarle 1 (por eso se comprueba dp[v - c] != INF antes de sumar).
		//    Con INF = LLONG_MAX, dp[v-c] + 1 desborda y da un numero negativo.
		// 2. Cada moneda se puede usar A LO SUMO UNA VEZ: recorrer las monedas
		//    en el bucle exterior y el monto v en sentido DECRECIENTE,
		//    como en Knapsack 0/1.
		// 3. Cantidad limitada k_j de cada moneda: descomponer k_j en
		//    potencias de dos (1, 2, 4, ..., resto) y tratar cada trozo como
		//    una moneda 0/1 de valor t*c y costo t; con esto el costo es
		//    O(n * sum log k_j).
		// 4. Sistema canonico (greedy valido): para monedas como {1,5,10,25} o
		//    potencias de un mismo numero, el greedy de tomar la mayor moneda
		//    posible es correcto y corre en O(m). Verificar con la DP en
		//    casos pequenos antes de confiar en el greedy.
		// 5. Monto muy grande (~1e9) con monedas pequenas: la DP directa no
		//    cabe en memoria. Si la moneda mas grande es M, basta llenar el
		//    monto con muchas copias de la moneda de mejor razon y resolver por
		//    DP solo un resto de orden M^2.
		// 6. NUMERO de formas que logran el minimo: mantener un segundo
		//    arreglo formas[v]; al encontrar un valor menor, formas[v] =
		//    formas[v-c]; si es igual al minimo actual, formas[v] +=
		//    formas[v-c] (con modulo si lo piden).
		// 7. Minimo de monedas con costo distinto por moneda (no siempre 1):
		//    reemplazar el + 1 por + costo[c]; es el Knapsack ilimitado de
		//    costo minimo con llenado exacto.
		// 8. Solo minimo de monedas para muchos montos distintos: la misma
		//    tabla dp sirve para todos hasta el maximo consultado, asi que se
		//    calcula una vez hasta el mayor n y se responde cada consulta en
		//    O(1).
	};
\end{lstlisting}

\textbf{Ejemplo de funcionamiento:}
\begin{lstlisting}[style=cpp]
	CambioMonedasMinimo a({1, 2, 5});
	long long r1 = a.minimoMonedas(11);
	// r1 = 3
	// (5 + 5 + 1)
	
	CambioMonedasMinimo b({1, 3, 4});
	long long r2 = b.minimoMonedas(6);
	// r2 = 2
	// (3 + 3; el greedy daria 4 + 1 + 1, que usa 3 monedas)
	
	vector<int> usadas = b.reconstruir(6);
	// usadas = {3, 3}
	
	CambioMonedasMinimo c({2});
	long long r3 = c.minimoMonedas(3);
	// r3 = -1
	// (con solo monedas de 2 es imposible formar un monto impar)
\end{lstlisting}

\textbf{Nota:} el centinela \texttt{INF} debe comprobarse antes de sumar \texttt{+ 1}. Si se usa \texttt{LLONG\_MAX} y se suma sin verificar, el valor desborda y pasa a ser negativo, lo que el \texttt{min} interpreta como una solución \textbf{mejor} que todas las reales, dando resultados absurdos sin ningún error visible. El código de arriba evita ese problema con la condición \texttt{dp[v - c] != INF}.

\textbf{Regla práctica:} usa esta DP siempre que el sistema de monedas sea arbitrario, o cuando no estés seguro de que el greedy funciona. Solo usa el greedy si el sistema es canónico (por ejemplo, potencias de un mismo número) y lo has comprobado. Si en lugar del mínimo se pide contar cuántas formas hay, usa \textit{Cambio de Monedas: Número de Formas}, cuidando el orden de los bucles.
\subsection{Problema de Partición (Diferencia Mínima de Subconjuntos)}

\textbf{Idea:} dado un arreglo de $n$ enteros no negativos con suma total $S$, se quiere dividirlo en dos subconjuntos $A$ y $B$ de modo que la diferencia $|\text{suma}(A)-\text{suma}(B)|$ sea \textbf{mínima}. Si $A$ suma $x$, entonces $B$ suma $S-x$ y la diferencia es $|S-2x|$. Por lo tanto basta saber \textbf{qué sumas $x$ son alcanzables} por algún subconjunto (es \textit{Subset Sum}, un caso de \textit{Knapsack 0/1} donde el peso y el valor coinciden) y quedarse con la que minimice $|S-2x|$. Como se puede tomar siempre el subconjunto de suma menor o igual a $S/2$, solo hace falta explorar $x\in[0,\lfloor S/2\rfloor]$: la respuesta es $S-2x^*$, con $x^*$ la mayor suma alcanzable que no supera $S/2$.

\textbf{¿Por qué usarlo?} Probar los $2^n$ subconjuntos es inviable para $n$ moderado. La DP de factibilidad resuelve el problema en $O(nS)$ porque el estado ``¿es alcanzable la suma $x$?'' resume todas las formas de decidir los elementos ya procesados. Un greedy (repartir cada elemento, de mayor a menor, al lado que va perdiendo) da una buena aproximación pero \textbf{no es exacto}. Como \textit{Knapsack 0/1}, es pseudopolinomial: el costo depende del valor numérico de $S$, no solo de $n$.

\textbf{Casos de uso:}
\begin{itemize}
	\item Repartir tareas o cargas entre dos personas o máquinas para que el tiempo de ambas sea lo más parecido posible.
	\item Dividir un conjunto de monedas u objetos en dos montones de valor casi igual.
	\item Decidir si un arreglo se puede partir en dos subconjuntos de igual suma (diferencia mínima igual a $0$).
	\item Base de variantes con restricciones, como exigir que ambos subconjuntos tengan el mismo tamaño.
	\item Problemas de equilibrio de una balanza con pesas: colocar cada pesa a un lado y minimizar el desbalance.
\end{itemize}

\textbf{Complejidad:} $O(nS)$ tiempo, con $S$ la suma total (o $O(nS/2)$ si solo se explora hasta la mitad): hay $\lfloor S/2\rfloor+1$ sumas posibles y cada elemento las recorre una vez con trabajo $O(1)$. El espacio es $O(S)$ con un arreglo unidimensional.

\begin{lstlisting}[style=cpp]
	// CREATE/READ/WRITE: Problema de particion: diferencia minima entre las
	// sumas de dos subconjuntos, con reconstruccion de uno de ellos
	struct ProblemaParticion {
		vector<int> a;
		int suma;
		
		// CREATE: guarda el arreglo y calcula la suma total
		ProblemaParticion(const vector<int>& v) : a(v), suma(0) {
			for (int x : a) suma += x;
		} // O(n)
		
		// READ: diferencia minima posible entre las sumas de los dos subconjuntos
		int diferenciaMinima() const {
			int mitad = suma / 2;
			vector<bool> alcanzable(mitad + 1, false);
			alcanzable[0] = true; // el subconjunto vacio suma 0
			for (int x : a) {
				// s decreciente: cada elemento se usa a lo sumo una vez (0/1)
				for (int s = mitad; s >= x; s--) {
					if (alcanzable[s - x]) alcanzable[s] = true;
				}
			}
			for (int s = mitad; s >= 0; s--) {
				if (alcanzable[s]) return suma - 2 * s;
			}
			return suma; // no se alcanza nunca: solo ocurre si mitad < 0
		} // O(n*S) tiempo, O(S) espacio
		
		// READ: indices (0-indexados) de un subconjunto cuya suma es la mayor
		// posible sin superar suma/2; el complemento es el otro subconjunto
		vector<int> reconstruir() const {
			int n = (int)a.size(), mitad = suma / 2;
			// dp[i][s] = true si con los primeros i elementos se alcanza la suma s
			vector<vector<bool>> dp(n + 1, vector<bool>(mitad + 1, false));
			dp[0][0] = true;
			for (int i = 1; i <= n; i++) {
				for (int s = 0; s <= mitad; s++) {
					dp[i][s] = dp[i - 1][s];
					if (s >= a[i - 1] && dp[i - 1][s - a[i - 1]]) dp[i][s] = true;
				}
			}
			int s = mitad;
			while (s > 0 && !dp[n][s]) s--;
			vector<int> elegidos;
			for (int i = n; i >= 1 && s > 0; i--) {
				if (!dp[i - 1][s]) { // el elemento i fue necesario para llegar a s
					elegidos.push_back(i - 1);
					s -= a[i - 1];
				}
			}
			reverse(elegidos.begin(), elegidos.end());
			return elegidos;
		} // O(n*S) tiempo y espacio
		
		// READ: true si el arreglo se puede partir en dos subconjuntos de
		// igual suma
		bool sePuedeDividirEnIguales() const {
			return suma % 2 == 0 && diferenciaMinima() == 0;
		} // O(n*S)
		
		// ------------------------------------------------------------
		// MODIFICACIONES Y COMENTARIOS CON LOS USOS
		// ------------------------------------------------------------
		// 1. Elementos negativos: desplazar los indices sumando un offset
		//    igual a la suma de los valores absolutos de los negativos, o usar
		//    un unordered_set / map de sumas alcanzables en vez de un arreglo.
		// 2. Version con bitset: alcanzable |= (alcanzable << x) para cada x,
		//    con bitset<MAXS>; reduce el costo a O(n*S/64) (ver Bitset para
		//    Optimizacion en la parte de BitWise).
		// 3. Suma total muy grande pero n pequeno (n <= 40): la DP es
		//    inviable; usar meet in the middle, generando las sumas de cada
		//    mitad, ordenandolas y buscando por dos punteros la pareja que
		//    minimice |S - 2x| (ver Meet in the Middle con Mascaras de Bits).
		// 4. Subconjuntos de IGUAL TAMANO (n par, ambos con n/2 elementos):
		//    agregar una dimension, alcanzable[j][s] = con j elementos se
		//    alcanza la suma s, recorriendo j y s en sentido decreciente.
		// 5. Contar el NUMERO de formas de partir con una diferencia dada d:
		//    se busca x = (S - d) / 2 (debe ser entero); contar subconjuntos
		//    con suma x con cuenta[0] = 1 y cuenta[s] += cuenta[s - x_i], en
		//    sentido decreciente (es el problema Target Sum).
		// 6. Partir en K subconjuntos (K > 2) minimizando la suma maxima: ya no
		//    es esta DP; se resuelve con DP sobre mascaras de bits para n
		//    pequeno, o con busqueda binaria sobre la respuesta y backtracking.
		// 7. Minimizar la suma maxima al partir un arreglo en K SEGMENTOS
		//    CONTIGUOS: es otro problema (Binary Search on Answer con
		//    verificacion greedy, u optimizaciones de DP); no confundir con
		//    esta particion en dos subconjuntos arbitrarios.
		// 8. Elementos iguales a 0: se ignoran, no cambian ninguna suma; pero
		//    conviene filtrarlos si el bucle interno usa s >= x con x = 0.
	};
\end{lstlisting}

\textbf{Ejemplo de funcionamiento:}
\begin{lstlisting}[style=cpp]
	ProblemaParticion p({1, 6, 11, 5});
	
	int d = p.diferenciaMinima();
	// d = 1
	// (suma total 23; el subconjunto {1, 5, 6} suma 12 y {11} suma 11,
	//  diferencia |12 - 11| = 1; es imposible lograr 0 porque 23 es impar)
	
	vector<int> sub = p.reconstruir();
	// sub = {2}
	// (el indice 2 es el elemento 11, suma 11 = la mayor suma alcanzable
	//  que no supera 23/2 = 11; el resto {1, 6, 5} suma 12)
	
	ProblemaParticion q({3, 1, 1, 2, 2, 1});
	bool iguales = q.sePuedeDividirEnIguales();
	// iguales = true
	// (suma 10; {3, 2} y {1, 1, 2, 1} suman 5 cada uno)
\end{lstlisting}

\textbf{Nota:} solo hace falta explorar sumas hasta $\lfloor S/2\rfloor$ porque el subconjunto complementario cubre la otra mitad: si $A$ suma $x\le S/2$, $B$ suma $S-x\ge S/2$ y la diferencia es $S-2x$. Además, si $S$ es impar la diferencia mínima nunca puede ser $0$. En \texttt{reconstruir()}, el subconjunto devuelto es uno de varios óptimos posibles; el otro es simplemente el complemento.

\textbf{Regla práctica:} usa esta DP de factibilidad (\textit{Subset Sum} hasta la mitad de la suma) cuando pidan dividir un conjunto en dos partes con diferencia mínima y la suma total sea moderada (hasta unos $10^5$ o $10^6$). Si $n\le 40$ y la suma es enorme, cambia a \textit{meet in the middle}. Si en realidad piden partir un arreglo en $K$ segmentos contiguos minimizando la suma máxima, no es este problema: usa búsqueda binaria sobre la respuesta.

Con esto queda completa la sección **Knapsack y Variantes**. Sigue pendiente que me confirmes si "mínima suma en particiones" era este problema o el de `K` segmentos contiguos (lo dejé cubierto como modificación 7 por si acaso).

¿Sigo con las referencias cruzadas de **LIS** y **LCS**, y **Subsecuencia Creciente de Suma Máxima**?

\section{DP sobre Secuencias}
\subsection{Subsecuencia Creciente Más Larga (LIS): referencia cruzada}

\textbf{Idea:} esta subsección ya se desarrolló en \textit{Subsecuencia Creciente Más Larga (LIS)} de String Processing, con la versión $O(n\log n)$ basada en \texttt{colas} y búsqueda binaria. Aquí se ve desde el punto de vista de la \textbf{DP}: el estado es $\text{dp}[i]$, la longitud de la LIS que \textbf{termina} en el índice $i$, y la transición es $\text{dp}[i]=1+\max\{\text{dp}[j]:j<i,\ a[j]<a[i]\}$ (o $1$ si no existe tal $j$). Esta forma es más lenta, $O(n^2)$, pero \textbf{guarda información por posición}, y eso permite contar cuántas LIS hay, reconstruir una concreta o cambiar el criterio de optimización, cosas que la versión con \texttt{colas} no permite directamente.

\textbf{¿Por qué usarlo?} Usa la versión $O(n\log n)$ cuando solo se pida la longitud. Usa esta DP cuando el problema necesite el estado por índice: número de LIS, LIS de suma máxima (ver \textit{Subsecuencia Creciente de Suma Máxima}), LIS con restricciones entre elementos consecutivos, o cualquier variante donde el valor guardado en $j$ deba combinarse con el de $i$ de una forma que no sea solo ``mayor longitud''. Es el mismo esquema de \textit{Memoización vs. Tabulación}: los estados se llenan de izquierda a derecha porque $\text{dp}[i]$ solo depende de índices menores.

\textbf{Casos de uso:}
\begin{itemize}
	\item Contar el número de subsecuencias crecientes de longitud máxima.
	\item Reconstruir una LIS concreta siguiendo el arreglo de padres.
	\item Base de las variantes ponderadas (suma máxima, costo mínimo, apilamiento de cajas).
	\item Problemas donde $n$ es pequeño ($n\le 5000$) y importa la simplicidad más que la velocidad.
	\item Subsecuencia bitónica: combinar $\text{dp}[i]$ desde la izquierda con la análoga desde la derecha.
\end{itemize}

\textbf{Complejidad:} $O(n^2)$ tiempo, porque hay $n$ estados y cada uno revisa hasta $n$ índices anteriores; $O(n)$ espacio.

\begin{lstlisting}[style=cpp]
	// CREATE/READ/WRITE: LIS estricta por DP O(n^2), con longitud, conteo de
	// subsecuencias optimas y reconstruccion de una de ellas
	struct LISPorDP {
		vector<int> a;
		vector<int> dp;        // dp[i] = longitud de la LIS que termina en i
		vector<int> padre;     // padre[i] = indice anterior en esa LIS, o -1
		vector<long long> cuenta; // cuenta[i] = cuantas LIS de longitud dp[i] terminan en i
		
		// CREATE: calcula las tres tablas en una sola pasada (tabulacion)
		LISPorDP(const vector<int>& v) : a(v) {
			int n = (int)a.size();
			dp.assign(n, 1);
			padre.assign(n, -1);
			cuenta.assign(n, 1);
			for (int i = 0; i < n; i++) {
				for (int j = 0; j < i; j++) {
					if (a[j] < a[i]) {
						if (dp[j] + 1 > dp[i]) {
							dp[i] = dp[j] + 1;
							padre[i] = j;
							cuenta[i] = cuenta[j];
						} else if (dp[j] + 1 == dp[i]) {
							cuenta[i] += cuenta[j];
						}
					}
				}
			}
		} // O(n^2)
		
		// READ: longitud de la LIS
		int longitud() const {
			int mejor = 0;
			for (int x : dp) mejor = max(mejor, x);
			return mejor;
		} // O(n)
		
		// READ: numero de subsecuencias crecientes de longitud maxima
		long long contarOptimas() const {
			int mejor = longitud();
			long long total = 0;
			for (int i = 0; i < (int)a.size(); i++)
			if (dp[i] == mejor) total += cuenta[i];
			return total;
		} // O(n)
		
		// READ: una LIS optima como lista de valores
		vector<int> reconstruir() const {
			vector<int> res;
			if (a.empty()) return res;
			int fin = 0;
			for (int i = 0; i < (int)a.size(); i++)
			if (dp[i] > dp[fin]) fin = i;
			for (int i = fin; i != -1; i = padre[i]) res.push_back(a[i]);
			reverse(res.begin(), res.end());
			return res;
		} // O(n)
		
		// ------------------------------------------------------------
		// MODIFICACIONES Y COMENTARIOS CON LOS USOS
		// ------------------------------------------------------------
		// 1. LIS NO estricta (permite iguales): cambiar a[j] < a[i] por
		//    a[j] <= a[i] en el constructor.
		// 2. LIS de suma maxima: guardar suma[i] = a[i] + max(suma[j]) en vez
		//    de contar elementos (ver Subsecuencia Creciente de Suma Maxima).
		// 3. Solo la longitud con n grande: usar la version O(n log n) con
		//    colas y busqueda binaria de la seccion de String Processing.
		// 4. LIS como caso de LCS: LIS(a) = LCS(a, valores unicos de a
		//    ordenados); util para verificar, pero mas lento que la DP directa.
		// 5. Cuenta con modulo: aplicar % MOD en cuenta[i] += cuenta[j] si el
		//    numero de LIS es enorme.
		// 6. Subsecuencia bitonica mas larga: calcular tambien la DP desde la
		//    derecha (dpDer[i] = LDS que empieza en i) y tomar
		//    max(dp[i] + dpDer[i] - 1).
		// 7. Restriccion entre consecutivos (por ejemplo a[i] - a[j] <= K):
		//    agregar esa condicion al if del bucle interno; para n grande se
		//    acelera con un arbol de segmentos sobre los valores.
	};
\end{lstlisting}

\textbf{Ejemplo de funcionamiento:}
\begin{lstlisting}[style=cpp]
	LISPorDP l({10, 9, 2, 5, 3, 7, 101, 18});
	
	int largo = l.longitud();
	// largo = 4
	
	long long cuantas = l.contarOptimas();
	// cuantas = 4
	// (2,5,7,101 - 2,5,7,18 - 2,3,7,101 - 2,3,7,18)
	
	vector<int> una = l.reconstruir();
	// una = {2, 5, 7, 101}
\end{lstlisting}

\textbf{Regla práctica:} para solo la longitud, usa \textit{Subsecuencia Creciente Más Larga (LIS)} en $O(n\log n)$. Usa esta DP $O(n^2)$ cuando necesites contar, reconstruir o ponderar, o cuando $n$ sea pequeño y quieras la versión más fácil de modificar.

\subsection{Subsecuencia Común Más Larga (LCS): referencia cruzada}

\textbf{Idea:} esta subsección ya se desarrolló en \textit{Subsecuencia Común Más Larga (Longest Common Subsequence)} de String Processing. Aquí se resume desde la DP: el estado $\text{dp}[i][j]$ es la LCS de los prefijos $s[0..i-1]$ y $t[0..j-1]$. Si $s[i-1]=t[j-1]$, entonces $\text{dp}[i][j]=\text{dp}[i-1][j-1]+1$; si no, $\text{dp}[i][j]=\max(\text{dp}[i-1][j],\text{dp}[i][j-1])$. Es el ejemplo típico de DP en \textbf{dos dimensiones sobre pares de prefijos}, y la plantilla de muchos problemas parecidos (distancia de edición, supersecuencia común, subsecuencia palindrómica, alineamiento de secuencias).

\textbf{¿Por qué usarlo?} Aquí importa reconocer el \textbf{patrón}, más que el problema concreto: cuando hay que comparar o alinear dos secuencias y las decisiones se toman de izquierda a derecha, el estado natural es ``cuántos elementos de cada una he consumido''. Cambiando solo la transición se resuelve una familia entera. Por eso conviene tener LCS a mano en esta parte aunque su explicación completa esté en String Processing.

\textbf{Casos de uso:}
\begin{itemize}
	\item Comparar dos secuencias permitiendo huecos (diff de archivos, alineamiento de ADN).
	\item Base de \textit{Distancia de Edición}: mismo estado con una operación más (sustituir).
	\item Base de la supersecuencia común más corta: $|s|+|t|-\text{LCS}$.
	\item Subsecuencia palindrómica más larga: LCS de $s$ con su reverso.
	\item Plantilla para DP con tres o más secuencias añadiendo una dimensión.
\end{itemize}

\textbf{Complejidad:} $O(nm)$ tiempo, con $n=|s|$ y $m=|t|$: hay $(n+1)(m+1)$ estados y cada uno se calcula en $O(1)$. El espacio es $O(nm)$ con la tabla completa, o $O(\min(n,m))$ si solo se quiere la longitud.

\begin{lstlisting}[style=cpp]
	// CREATE/READ/WRITE: LCS por DP con espacio reducido para la longitud y
	// tabla completa para reconstruir
	struct LCSPorDP {
		string s, t;
		
		// CREATE: guarda las dos secuencias
		LCSPorDP(const string& a, const string& b) : s(a), t(b) {
		} // O(1)
		
		// READ: longitud de la LCS con solo dos filas (espacio O(m))
		int longitud() const {
			int n = (int)s.size(), m = (int)t.size();
			vector<int> anterior(m + 1, 0), actual(m + 1, 0);
			for (int i = 1; i <= n; i++) {
				for (int j = 1; j <= m; j++) {
					if (s[i - 1] == t[j - 1]) actual[j] = anterior[j - 1] + 1;
					else actual[j] = max(anterior[j], actual[j - 1]);
				}
				swap(anterior, actual);
			}
			return anterior[m];
		} // O(n*m) tiempo, O(m) espacio
		
		// READ: una LCS concreta (tabla completa para poder retroceder)
		string reconstruir() const {
			int n = (int)s.size(), m = (int)t.size();
			vector<vector<int>> dp(n + 1, vector<int>(m + 1, 0));
			for (int i = 1; i <= n; i++)
			for (int j = 1; j <= m; j++)
			dp[i][j] = (s[i - 1] == t[j - 1]) ? dp[i - 1][j - 1] + 1
			: max(dp[i - 1][j], dp[i][j - 1]);
			string res;
			int i = n, j = m;
			while (i > 0 && j > 0) {
				if (s[i - 1] == t[j - 1]) { res += s[i - 1]; i--; j--; }
				else if (dp[i - 1][j] >= dp[i][j - 1]) i--;
				else j--;
			}
			reverse(res.begin(), res.end());
			return res;
		} // O(n*m) tiempo y espacio
		
		// ------------------------------------------------------------
		// MODIFICACIONES Y COMENTARIOS CON LOS USOS
		// ------------------------------------------------------------
		// 1. Solo insertar y eliminar (sin sustituir): la distancia minima es
		//    n + m - 2 * longitud().
		// 2. Distancia de edicion: mismo estado dp[i][j], pero en el caso de
		//    desigualdad se toma 1 + min(dp[i-1][j-1], dp[i-1][j], dp[i][j-1]);
		//    ver Distancia de Edicion.
		// 3. Supersecuencia comun mas corta: longitud n + m - longitud(); ver
		//    Superseceuncia Comun Mas Corta.
		// 4. Subsecuencia palindromica mas larga: LCSPorDP(s, reverso de s).
		// 5. Tres secuencias: dp[i][j][k] con transicion analoga, O(n*m*p);
		//    solo viable con secuencias cortas.
		// 6. Dos permutaciones de 1..n: LCS se reduce a LIS. Se reemplaza cada
		//    elemento de la segunda por su posicion en la primera y se calcula
		//    la LIS en O(n log n), evitando el O(n^2).
		// 7. Contar el numero de LCS distintas: agregar cuenta[i][j] con la
		//    misma estructura, restando el doble conteo cuando
		//    dp[i-1][j-1] == dp[i][j]; requiere cuidado con el modulo.
		// 8. Con peso por par que coincide (por ejemplo cada letra vale
		//    distinto): en la rama de igualdad sumar peso(s[i-1]) en vez de 1.
	};
\end{lstlisting}

\textbf{Ejemplo de funcionamiento:}
\begin{lstlisting}[style=cpp]
	LCSPorDP p("AGGTAB", "GXTXAYB");
	
	int largo = p.longitud();
	// largo = 4
	
	string lcs = p.reconstruir();
	// lcs = "GTAB"
\end{lstlisting}

\textbf{Regla práctica:} para el detalle del problema, la reconstrucción y su relación con la subcadena común, consulta \textit{Subsecuencia Común Más Larga (Longest Common Subsequence)}. Aquí queda como plantilla: siempre que aparezcan dos secuencias y un estado del tipo ``prefijo de una y prefijo de otra'', piensa en este esquema y cambia solo la transición.
\subsection{Subsecuencia Creciente de Suma Máxima}

\textbf{Idea:} dada una secuencia de $n$ enteros, se busca la subsecuencia \textbf{estrictamente creciente} cuya \textbf{suma} de elementos sea máxima (no la de mayor longitud). Es la variante ponderada de \textit{Subsecuencia Creciente Más Larga (LIS): referencia cruzada}: en vez de que cada elemento aporte $1$ a la longitud, aporta su propio valor. Sea $\text{dp}[i]$ la suma máxima de una subsecuencia creciente que \textbf{termina} en el índice $i$. La recurrencia es $\text{dp}[i]=a[i]+\max\bigl(0,\ \max\{\text{dp}[j]:j<i,\ a[j]<a[i]\}\bigr)$: el elemento $i$ se puede empezar solo (aporta $a[i]$) o extender la mejor subsecuencia previa que termine en un valor menor. La respuesta es $\max_i \text{dp}[i]$. Es el mismo esquema de DP por índice de la referencia cruzada, cambiando únicamente lo que se acumula.

\textbf{¿Por qué usarlo?} La subsecuencia más larga no es necesariamente la de mayor suma. En $\{100,1,2,3,4\}$ la LIS más larga es $1,2,3,4$ (longitud 4, suma 10), pero la de mayor suma es $\{100\}$ (suma 100). Por eso no sirve la técnica de \texttt{colas} con búsqueda binaria de LIS, que solo conserva longitudes: aquí el valor guardado en cada estado debe combinarse por \textbf{suma}. La DP directa cuesta $O(n^2)$; con un \textit{árbol de Fenwick} sobre los valores comprimidos, que responda el máximo de $\text{dp}$ entre los valores menores a $a[i]$, baja a $O(n\log n)$.

\textbf{Casos de uso:}
\begin{itemize}
	\item Elegir, respetando un orden creciente, los elementos que suman más (ganancia acumulada con restricción de orden).
	\item Apilamiento de cajas u objetos donde cada uno aporta un valor (por ejemplo, altura) y se debe cumplir una condición de crecimiento: ver \textit{Apilamiento de Cajas (Box Stacking)}.
	\item Selección de trabajos o eventos ordenados donde cada elemento posterior debe superar al anterior y se maximiza el beneficio total.
	\item Reconstruir la subsecuencia concreta que logra esa suma, siguiendo el arreglo de padres.
	\item Variantes con pesos distintos del valor (un peso $w_i$ por elemento) manteniendo la condición sobre $a$.
\end{itemize}

\textbf{Complejidad:} $O(n^2)$ tiempo con la DP directa, porque hay $n$ estados y cada uno revisa hasta $n$ índices anteriores; $O(n)$ espacio. Con Fenwick sobre valores comprimidos: $O(n\log n)$.

\begin{lstlisting}[style=cpp]
	// CREATE/READ/WRITE: Subsecuencia creciente de suma maxima, con version
	// O(n^2) con reconstruccion y version O(n log n) con Fenwick de maximos
	struct SubsecuenciaCrecienteSumaMaxima {
		vector<long long> a;
		vector<long long> dp;  // dp[i] = suma maxima de una subsecuencia creciente que termina en i
		vector<int> padre;     // padre[i] = indice anterior en esa subsecuencia, o -1
		
		// CREATE: calcula dp y padre en una sola pasada (tabulacion)
		SubsecuenciaCrecienteSumaMaxima(const vector<long long>& v) : a(v) {
			int n = (int)a.size();
			dp.assign(n, 0);
			padre.assign(n, -1);
			for (int i = 0; i < n; i++) {
				dp[i] = a[i]; // la subsecuencia formada solo por a[i]
				for (int j = 0; j < i; j++) {
					if (a[j] < a[i] && dp[j] + a[i] > dp[i]) {
						dp[i] = dp[j] + a[i];
						padre[i] = j;
					}
				}
			}
		} // O(n^2)
		
		// READ: suma maxima de una subsecuencia estrictamente creciente
		long long sumaMaxima() const {
			if (a.empty()) return 0;
			return *max_element(dp.begin(), dp.end());
		} // O(n)
		
		// READ: una subsecuencia optima como lista de valores
		vector<long long> reconstruir() const {
			vector<long long> res;
			if (a.empty()) return res;
			int fin = (int)(max_element(dp.begin(), dp.end()) - dp.begin());
			for (int i = fin; i != -1; i = padre[i]) res.push_back(a[i]);
			reverse(res.begin(), res.end());
			return res;
		} // O(n)
		
		// READ: version O(n log n): Fenwick de maximos indexado por el rango
		// (valor comprimido) de cada elemento; consulta el mejor dp entre los
		// valores estrictamente menores
		static long long sumaMaximaRapida(const vector<long long>& v) {
			int n = (int)v.size();
			if (n == 0) return 0;
			vector<long long> vals = v;
			sort(vals.begin(), vals.end());
			vals.erase(unique(vals.begin(), vals.end()), vals.end());
			int m = (int)vals.size();
			vector<long long> bit(m + 1, 0); // maximos por prefijo; 0 = "no extender"
			auto actualizar = [&](int pos, long long val) {
				for (; pos <= m; pos += pos & -pos) bit[pos] = max(bit[pos], val);
			};
			auto consultar = [&](int pos) { // maximo en las posiciones 1..pos
				long long r = 0;
				for (; pos > 0; pos -= pos & -pos) r = max(r, bit[pos]);
				return r;
			};
			bool hayElemento = false;
			long long mejorGlobal = 0;
			for (long long x : v) {
				int pos = (int)(lower_bound(vals.begin(), vals.end(), x) - vals.begin()) + 1;
				long long previo = consultar(pos - 1); // valores estrictamente menores
				long long actual = x + max(0LL, previo);
				actualizar(pos, actual);
				if (!hayElemento || actual > mejorGlobal) { mejorGlobal = actual; hayElemento = true; }
			}
			return mejorGlobal;
		} // O(n log n)
		
		// ------------------------------------------------------------
		// MODIFICACIONES Y COMENTARIOS CON LOS USOS
		// ------------------------------------------------------------
		// 1. Subsecuencia NO estricta (permite iguales): cambiar a[j] < a[i]
		//    por a[j] <= a[i] en la version O(n^2); en la de Fenwick, consultar
		//    hasta pos (inclusive) en vez de pos - 1.
		// 2. Elementos negativos: la subsecuencia debe tener al menos un
		//    elemento, por eso dp[i] se inicializa en a[i] (no en 0). En la
		//    version con Fenwick se usa max(0, previo) para no arrastrar sumas
		//    negativas; el resultado es correcto aunque todos sean negativos
		//    (devuelve el mayor elemento).
		// 3. Subsecuencia DECRECIENTE de suma maxima: cambiar la comparacion a
		//    a[j] > a[i] (la suma se acumula siempre con los valores originales).
		// 4. Peso distinto del valor: si cada elemento tiene un peso w[i] que
		//    se suma pero la condicion de orden es sobre a[i], usar
		//    dp[i] = w[i] + max(dp[j]) con a[j] < a[i].
		// 5. Contar cuantas subsecuencias logran la suma maxima: mantener
		//    cuenta[i] como en LIS por DP, sumando cuenta[j] cuando
		//    dp[j] + a[i] iguala el maximo (con modulo si lo piden).
		// 6. Restriccion de indices o de diferencia entre consecutivos (por
		//    ejemplo j >= i - K): acotar el bucle interno, o consultar el
		//    Fenwick/arbol de segmentos solo en la ventana valida.
		// 7. Valores muy grandes o negativos para el Fenwick: siempre
		//    comprimir a rangos 1..m antes de indexar.
		// 8. Cotas de tipo: la suma llega a n * max|a|; usar long long para
		//    evitar overflow con n grande.
	};
\end{lstlisting}

\textbf{Ejemplo de funcionamiento:}
\begin{lstlisting}[style=cpp]
	SubsecuenciaCrecienteSumaMaxima s({1, 101, 2, 3, 100, 4, 5});
	
	long long mejor = s.sumaMaxima();
	// mejor = 106
	// (la subsecuencia 1, 2, 3, 100 suma 106; 1, 101 suma solo 102 y
	//  1, 2, 3, 4, 5 suma solo 15)
	
	vector<long long> sub = s.reconstruir();
	// sub = {1, 2, 3, 100}
	
	long long rapido = SubsecuenciaCrecienteSumaMaxima::sumaMaximaRapida({1, 101, 2, 3, 100, 4, 5});
	// rapido = 106 (mismo resultado en O(n log n))
	
	SubsecuenciaCrecienteSumaMaxima t({100, 1, 2, 3, 4});
	long long distinta = t.sumaMaxima();
	// distinta = 100
	// (la LIS mas larga es 1, 2, 3, 4 con suma 10, pero la de mayor suma
	//  es el elemento 100 solo)
\end{lstlisting}

\textbf{Regla práctica:} usa esta DP cuando cada elemento aporte un valor y se pida maximizar la suma bajo una condición de orden creciente; la técnica de \texttt{colas} de LIS no aplica porque solo conserva longitudes. Usa la versión $O(n^2)$ si $n\le 5000$ y necesitas reconstruir, y la de Fenwick sobre valores comprimidos cuando $n$ sea grande.

\section{DP de Intervalos}
\subsection{Burst Balloons (Problema de los Globos)}

\textbf{Idea:} dado un arreglo de $n$ globos con valores
$\texttt{nums}[i]$, al reventar el globo $i$ se ganan
$\texttt{nums}[i{-}1]\cdot\texttt{nums}[i]\cdot\texttt{nums}[i{+}1]$
monedas (si un vecino no existe, se toma como $1$), y luego $i$
desaparece, por lo que sus vecinos pasan a ser adyacentes. Se busca el
orden de reventado que maximiza el total de monedas. Es el ejemplo
clásico de \textbf{DP sobre intervalos}, y la clave está en cambiar la
pregunta: en vez de "¿cuál globo reviento \textbf{primero}?" (después
de eso los vecinos cambian y los subproblemas quedan acoplados), se
pregunta "¿cuál globo reviento \textbf{último} dentro de este
intervalo?". Si $k$ es el último en reventar entre los límites $l$ y
$r$ (exclusivos), en ese momento sus vecinos son exactamente $l$ y $r$
(todo lo demás ya desapareció), y los intervalos $(l,k)$ y $(k,r)$ son
\textbf{independientes}:
$$\texttt{dp}[l][r] = \max_{l<k<r}\Big(\texttt{dp}[l][k] +
\texttt{dp}[k][r] + a[l]\cdot a[k]\cdot a[r]\Big)$$
donde $a$ es el arreglo original con un $1$ agregado a cada extremo
(centinelas) y $\texttt{dp}[l][r]$ es la mejor ganancia reventando
\textbf{todos} los globos estrictamente entre $l$ y $r$.

\textbf{¿Por qué usarlo?} Es el problema que enseña la idea de
\textbf{invertir el orden de decisión} para desacoplar subproblemas:
con "el primero en reventar" los subproblemas dependen de qué globos
quedaron a los lados; con "el último en reventar" los bordes del
intervalo están fijos y las dos mitades no interactúan. Ese mismo
truco aparece en otros problemas de intervalos (multiplicación de
cadena de matrices, triangulación óptima de polígonos, árbol de
búsqueda binaria óptimo). Greedy no sirve aquí: reventar siempre el
globo más chico o el más grande falla en casos pequeños, así que hace
falta considerar todas las posiciones del último globo, justo lo que
hace DP (ver Greedy vs. Programación Dinámica).

\textbf{Casos de uso:}

\begin{itemize}
	\item Burst Balloons tal cual: máxima ganancia reventando todos
	los globos.
	\item Cualquier problema donde eliminar un elemento cambia quién
	es vecino de quién: la pregunta útil es "¿qué elemento queda
	\textbf{al final}?" en vez de "¿cuál elimino primero?".
	\item Plantilla general de DP por intervalos: estado
	$[l,r]$, transición probando un punto de corte $k$, recorrido por
	\textbf{longitud creciente} del intervalo.
	\item Reconstruir el \textbf{orden} de reventado óptimo (no solo
	el valor), guardando el $k$ óptimo de cada intervalo.
\end{itemize}

\textbf{Complejidad:} $O(n^3)$ tiempo ($O(n^2)$ estados, cada uno
prueba hasta $n$ valores de $k$), $O(n^2)$ memoria.

\begin{lstlisting}[style=cpp]
	struct BurstBalloons {
		vector<int> a;                 // nums con un 1 en cada extremo
		vector<vector<long long>> dp;  // dp[l][r]: reventar todo lo que
		// esta ESTRICTAMENTE entre l y r
		vector<vector<int>> ultimo;    // ultimo[l][r]: k optimo (el
		// ultimo globo en reventar)
		int m;                         // tamano de a (n + 2)
		
		// CREATE: agrega los centinelas 1 y calcula toda la tabla dp.
		BurstBalloons(vector<int>& nums) {
			a.push_back(1);
			for (int x : nums) a.push_back(x);
			a.push_back(1);
			m = a.size();
			
			dp.assign(m, vector<long long>(m, 0));
			ultimo.assign(m, vector<int>(m, -1));
			
			// se recorre por LONGITUD creciente del intervalo: al calcular
			// dp[l][r] ya estan listos todos los subintervalos mas cortos
			for (int len = 2; len < m; len++) {
				for (int l = 0; l + len < m; l++) {
					int r = l + len;
					
					for (int k = l + 1; k < r; k++) {
						// k es el ULTIMO globo en reventar dentro de (l,r):
						// sus vecinos en ese momento son exactamente l y r
						long long total = dp[l][k] + dp[k][r]
						+ (long long)a[l] * a[k] * a[r];
						
						if (total > dp[l][r]) {
							dp[l][r] = total;
							ultimo[l][r] = k;
						}
					}
				}
			}
		}                                          // O(n^3)
		
		// READ: maxima cantidad de monedas reventando todos los globos.
		long long maximo() {
			return dp[0][m - 1];
		}                                          // O(1)
		
		// READ: orden de reventado optimo (indices sobre el arreglo
		// original nums, base 0), del primero al ultimo en reventar.
		vector<int> ordenOptimo() {
			vector<int> orden;
			reconstruir(0, m - 1, orden);
			return orden;
		}                                          // O(n)
		
		// Recursion: el orden de (l,r) es el de (l,k), luego el de (k,r),
		// y al final k (que es el ultimo de este intervalo).
		void reconstruir(int l, int r, vector<int>& orden) {
			int k = ultimo[l][r];
			if (k == -1) return;   // intervalo sin globos adentro
			
			reconstruir(l, k, orden);
			reconstruir(k, r, orden);
			orden.push_back(k - 1);   // k - 1: se resta el centinela
		}
		
		// ------------------------------------------------------------
		// MODIFICACIONES Y COMENTARIOS CON LOS USOS
		// ------------------------------------------------------------
		
		// 1. VERSION TOP-DOWN (memoizacion): misma recurrencia, sin
		//    preocuparse por el orden de llenado. Util si se desconfia
		//    del orden por longitud:
		//
		// long long resolver(int l, int r) {
			//     if (r - l < 2) return 0;          // nada entre l y r
			//     if (memo[l][r] != -1) return memo[l][r];
			//     long long mejor = 0;
			//     for (int k = l + 1; k < r; k++)
			//         mejor = max(mejor, resolver(l, k) + resolver(k, r)
			//                            + (long long)a[l] * a[k] * a[r]);
			//     return memo[l][r] = mejor;
			// }
		
		// 2. OTRAS GANANCIAS POR REVENTADO: solo cambia el termino que se
		//    suma al combinar. Ejemplos con la misma estructura:
		//    - ganar a[l] + a[k] + a[r]      (suma en vez de producto)
		//    - costo minimo en vez de ganancia maxima: cambiar max por
		//      min e inicializar dp en un valor grande (INF)
		//    El esqueleto (intervalo, k = ultimo, dos mitades
		//    independientes) no cambia.
		
		// 3. MISMA PLANTILLA EN OTROS PROBLEMAS DE INTERVALOS:
		//    - multiplicacion de cadena de matrices:
		//      dp[l][r] = min sobre k de dp[l][k] + dp[k][r]
		//                 + dims[l] * dims[k] * dims[r]
		//    - triangulacion de poligono de costo minimo (k = tercer
		//      vertice del triangulo que usa la arista l-r)
		//    - cortar una barra en puntos dados con costo = longitud del
		//      segmento (k = ultimo corte del segmento)
		
		// 4. VERSION CIRCULAR (globos en circulo): se duplica el arreglo
		//    o se fija que globo es el ultimo de toda la ronda, y se
		//    resuelve el intervalo lineal resultante. Sube el costo por
		//    un factor n si se prueba cada opcion.
		
		// 5. OPTIMIZACION: para algunas variantes (como matriz de
		//    cadena o corte de barra) existe la optimizacion de Knuth,
		//    que baja O(n^3) a O(n^2), pero exige propiedades de la
		//    funcion de costo (cuadrangulo) que NO se cumplen en Burst
		//    Balloons con producto de tres terminos: aqui O(n^3) es lo
		//    esperado.
	};
\end{lstlisting}

\textbf{Ejemplo de funcionamiento:}

\begin{lstlisting}[style=cpp]
	vector<int> nums = {3, 1, 5, 8};
	
	BurstBalloons bb(nums);
	
	cout << bb.maximo();
	// 167
	
	for (int i : bb.ordenOptimo()) cout << i << " ";
	// 1 2 0 3   (indices sobre nums: se revienta el globo de valor 1,
	//            luego el 5, luego el 3 y al final el 8)
	// ganancias: 3*1*5 = 15, 3*5*8 = 120, 1*3*8 = 24, 1*8*1 = 8
	// total: 15 + 120 + 24 + 8 = 167
\end{lstlisting}

\textbf{Nota:} el recorrido por \textbf{longitud creciente} del
intervalo no es un detalle: $\texttt{dp}[l][r]$ depende de
$\texttt{dp}[l][k]$ y $\texttt{dp}[k][r]$, que son intervalos más
cortos, así que deben estar listos antes. Recorrer $l$ y $r$ con dos
\texttt{for} anidados "normales" (por filas) rompe esa dependencia y
usa valores todavía en cero. Los centinelas $1$ en los extremos evitan
casos especiales para los globos de las orillas.

\textbf{Regla práctica:} en un problema donde quitar un elemento
cambia los vecinos de los demás, no intentes modelar "el primero que
quito": modela "el \textbf{último} que queda" dentro de un intervalo
$[l,r]$. Ese elemento tiene vecinos conocidos ($l$ y $r$) y parte el
problema en dos mitades independientes, lo que convierte el
enredo en una DP de intervalos $O(n^3)$ con la misma forma que
multiplicación de cadena de matrices.

\section{Apilamiento y Ordenamiento}
\subsection{Apilamiento de Cajas (Box Stacking)}

\textbf{Idea:} dado un conjunto de $n$ cajas, cada una con tres
dimensiones $(x, y, z)$, se busca la \textbf{pila más alta} posible.
Una caja se puede apoyar sobre otra solo si la base de la de arriba es
\textbf{estrictamente menor} en ambas dimensiones que la base de la de
abajo, y cada caja se puede \textbf{rotar}: cualquiera de sus tres
dimensiones puede hacer de altura. Se resuelve en tres pasos:
(1) generar las \textbf{3 orientaciones} de cada caja, guardando en cada
una la base como $(w, d)$ con $w \le d$ (el giro de la base en el plano
no aporta nada nuevo: si una caja cabe sobre otra girada, también cabe
con las dos bases ordenadas de menor a mayor);
(2) \textbf{ordenar} todas las orientaciones por área de base
decreciente $w\cdot d$;
(3) aplicar una DP tipo LIS pesada:
$$\texttt{dp}[i] = h_i + \max\Big(0,\ \max_{j<i,\ w_j>w_i,\ d_j>d_i}
\texttt{dp}[j]\Big)$$
donde $\texttt{dp}[i]$ es la altura máxima de una pila que tiene la
orientación $i$ \textbf{en la cima}. La respuesta es el máximo de todos
los $\texttt{dp}[i]$.

\textbf{¿Por qué usarlo?} El paso de ordenar por área es lo que hace
que la DP funcione: si la caja $i$ va sobre la $j$, su base es
estrictamente menor en ambas dimensiones, así que su área también es
estrictamente menor y $j$ aparece \textbf{antes} en el orden. Ese orden
es un orden topológico de la relación "cabe sobre", y basta mirar solo
los $j<i$. Además, como la condición es \textbf{estricta}, una misma
orientación no puede repetirse en la pila (no cabe sobre sí misma), lo
que resuelve gratis la pregunta de cuántas copias de cada caja se
permiten. Greedy no sirve: apilar siempre la caja más alta, o la de
base más grande, falla en casos pequeños (ver Greedy vs.
Programación Dinámica).

\textbf{Casos de uso:}

\begin{itemize}
	\item Box Stacking clásico: altura máxima con rotaciones y base
	estrictamente decreciente.
	\item Variante \textbf{sin rotación}: cada caja tiene una sola
	orientación y el problema se vuelve una DP tipo LIS en 2D (ver la
	modificación 1 en el código).
	\item Anidamiento de sobres o muñecas rusas (\textit{Russian Doll
		Envelopes}): misma relación de orden parcial, pero con peso $1$ por
	elemento en vez de altura.
	\item Cuboides con desigualdad \textbf{no estricta} (\textit{Maximum
		Height by Stacking Cuboids}): cambia el criterio de orden, no la
	idea (modificación 2).
	\item Cualquier problema de "cadena más pesada en un orden parcial
	de dos criterios": ordenar por uno, y aplicar DP sobre el otro.
\end{itemize}

\textbf{Complejidad:} con $m = 3n$ orientaciones, ordenar cuesta
$O(m \log m)$ y la DP $O(m^2) = O(n^2)$ tiempo; $O(n)$ memoria.

\begin{lstlisting}[style=cpp]
	struct BoxStacking {
		struct Caja { int h, w, d; };   // altura y base (w <= d)
		vector<Caja> orient;            // las 3n orientaciones ordenadas
		vector<long long> dp;           // dp[i]: mejor pila con orient[i]
		// en la CIMA
		vector<int> padre;              // padre[i]: caja que queda debajo
		// de orient[i] en la mejor pila
		
		// CREATE: genera las 3 orientaciones de cada caja, las ordena por
		// area de base decreciente y calcula la DP.
		BoxStacking(vector<array<int,3>>& cajas) {
			for (auto& c : cajas) {
				for (int i = 0; i < 3; i++) {
					int h = c[i];                       // esta dimension es
					int p = c[(i + 1) % 3];             // la altura; las
					int q = c[(i + 2) % 3];             // otras dos, la base
					orient.push_back({h, min(p, q), max(p, q)});
				}
			}
			
			// area decreciente: si i va sobre j, area_i < area_j, asi que
			// j queda antes que i y basta mirar los j < i
			sort(orient.begin(), orient.end(),
			[](const Caja& a, const Caja& b) {
				return (long long)a.w * a.d > (long long)b.w * b.d;
			});
			
			int m = orient.size();
			dp.assign(m, 0);
			padre.assign(m, -1);
			
			for (int i = 0; i < m; i++) {
				dp[i] = orient[i].h;   // la caja sola, sin nada debajo
				for (int j = 0; j < i; j++) {
					bool cabe = orient[j].w > orient[i].w
					&& orient[j].d > orient[i].d;   // ESTRICTO
					if (cabe && dp[j] + orient[i].h > dp[i]) {
						dp[i] = dp[j] + orient[i].h;
						padre[i] = j;
					}
				}
			}
		}                                          // O(n^2)
		
		// READ: altura maxima de una pila.
		long long alturaMaxima() {
			return *max_element(dp.begin(), dp.end());
		}                                          // O(n)
		
		// READ: la pila optima de ARRIBA hacia ABAJO.
		vector<Caja> pila() {
			int mejor = max_element(dp.begin(), dp.end()) - dp.begin();
			vector<Caja> resultado;
			for (int i = mejor; i != -1; i = padre[i])
			resultado.push_back(orient[i]);
			return resultado;
		}                                          // O(largo de la pila)
		
		// ------------------------------------------------------------
		// MODIFICACIONES Y COMENTARIOS CON LOS USOS
		// ------------------------------------------------------------
		
		// 1. SIN ROTACION (cada caja tiene una sola orientacion fija, con
		//    base (w, d) dada): en el constructor se genera UNA sola
		//    orientacion por caja, sin ordenar w y d entre si, y la
		//    condicion de apilar compara cada dimension en su lugar:
		//
		//    orient.push_back({h, w, d});          // una orientacion
		//    ...
		//    bool cabe = orient[j].w > orient[i].w && orient[j].d > orient[i].d;
		//
		//    Si ademas todas las alturas valen 1 (sobres, munecas rusas),
		//    el problema se puede bajar a O(n log n): ordenar por w
		//    creciente y, ante empate de w, por d DECRECIENTE (para que
		//    dos sobres con el mismo w no se encadenen), y calcular la
		//    subsecuencia creciente estricta mas larga sobre d con
		//    busqueda binaria.
		
		// 2. DESIGUALDAD NO ESTRICTA (una caja puede apoyarse sobre otra
		//    si sus dimensiones son menores O IGUALES, como en Maximum
		//    Height by Stacking Cuboids): se ordenan las 3 dimensiones de
		//    cada caja de menor a mayor, se ordenan todas las cajas
		//    lexicograficamente por esas tres dimensiones, y la DP usa
		//    <= en las tres:
		//
		//    bool cabe = A[j][0] <= A[i][0] && A[j][1] <= A[i][1]
		//             && A[j][2] <= A[i][2];   // j va DEBAJO de i
		//    dp[i] = A[i][2] + max(dp[j] entre las que cumplen cabe)
		//
		//    Aqui la altura de cada caja es su dimension mas grande y NO
		//    se generan 3 orientaciones: al ordenar las dimensiones, esa
		//    orientacion ya es siempre la mejor. Con <= el area deja de
		//    servir como orden topologico (dos cajas pueden tener la
		//    misma area), por eso se ordena lexicograficamente.
		
		// 3. LIMITE DE COPIAS POR CAJA: en la version estricta de arriba
		//    cada orientacion aparece a lo sumo una vez, y la misma caja
		//    fisica puede aparecer con dos orientaciones distintas en la
		//    pila. Si el enunciado exige usar cada caja FISICA una sola
		//    vez, la DP simple ya no alcanza (habria que recordar que
		//    cajas se usaron): ahi hay que revisar el enunciado, porque
		//    suele ser un problema distinto.
		
		// 4. MEJORAR O(n^2): con muchas cajas se puede compactar la
		//    coordenada d y usar un arbol (Fenwick o segment tree de
		//    maximo) para consultar "mejor dp entre las cajas de base mas
		//    grande en ambas dimensiones", procesando por w decreciente.
		//    Baja a O(n log n), a costa de codigo bastante mas largo.
		
		// 5. PILA CON CAJAS ENTRE MEDIO ("cada caja debe soportar un
		//    peso maximo"): es otro problema (Box Stacking con
		//    resistencias). Se ordena por la suma peso + resistencia, y
		//    NO es esta misma DP.
	};
\end{lstlisting}

\textbf{Ejemplo de funcionamiento:}

\begin{lstlisting}[style=cpp]
	// cada caja es {x, y, z}
	vector<array<int,3>> cajas = {
		{4, 6, 7}, {1, 2, 3}, {4, 5, 6}, {10, 12, 32}
	};
	
	BoxStacking bs(cajas);
	
	cout << bs.alturaMaxima();
	// 60
	
	for (auto& c : bs.pila())
	cout << "(h=" << c.h << ",w=" << c.w << ",d=" << c.d << ") ";
	// (h=3,w=1,d=2) (h=1,w=2,d=3) (h=6,w=4,d=5) (h=4,w=5,d=6)
	// (h=4,w=6,d=7) (h=32,w=10,d=12) (h=10,w=12,d=32)
	// (de arriba hacia abajo; 3+1+6+4+4+32+10 = 60, y cada base es
	//  estrictamente menor que la de abajo en ambas dimensiones)
\end{lstlisting}

\textbf{Nota:} en este ejemplo las 4 cajas aportan \textbf{7}
orientaciones a la misma pila, porque cada una entra con más de una
orientación y todas caben en cadena. Eso es normal en este problema:
"apilar cajas" con rotaciones libres no significa usar cada caja una
vez, sino usar orientaciones distintas. Si el enunciado dice
explícitamente que cada caja física se usa \textbf{a lo más una vez},
esta DP da una respuesta demasiado optimista (ver modificación 3).

\textbf{Regla práctica:} cuando un problema pide la cadena \textbf{más
	pesada} en una relación de "cabe dentro de" con dos criterios, ordena
por un criterio de forma que el orden sea topológico (aquí, el área) y
aplica DP tipo LIS sobre el resto. Antes de programar, define con
cuidado dos cosas del enunciado: si la desigualdad es \textbf{estricta}
o no, y si hay \textbf{rotaciones} y de qué tipo. Esas dos decisiones
cambian el orden de la lista y la condición de la DP, y son la fuente
más común de Wrong Answer en este problema.

\section{DP Avanzada / Estados Especiales}
\subsection{Digit DP}

\textbf{Idea:} \textit{Digit DP} cuenta cuántos números dentro de un rango $[L,R]$ cumplen una condición sobre sus \textbf{dígitos} (suma de dígitos, cantidad de un dígito específico, dígitos no decrecientes, divisibilidad combinada con una propiedad de dígitos, etc.), sin enumerar cada número uno por uno. La idea es construir el número \textbf{dígito a dígito}, de izquierda a derecha, manteniendo un estado que resume lo necesario para decidir la condición al final (por ejemplo, la suma de dígitos usada hasta ahora, o el resto módulo $k$). El truco central es la bandera \texttt{tight} (ajustado): indica si los dígitos elegidos hasta ahora coinciden \textbf{exactamente} con los del límite superior. Si \texttt{tight} es verdadero, el siguiente dígito solo puede llegar hasta el dígito correspondiente del límite; si es falso, ya se "liberó" del límite y puede usar cualquier dígito de 0 a 9, porque el número ya es estrictamente menor que el límite en una posición anterior.

\textbf{¿Por qué usarlo?} Contar por fuerza bruta recorriendo cada entero de $L$ a $R$ es $O(R-L)$, inviable cuando $R$ puede llegar a $10^{18}$. Digit DP convierte el problema en una DP sobre \textbf{posición del dígito} (como mucho 18-19 posiciones para \texttt{long long}) combinada con el estado relevante (suma, resto, último dígito, etc.), dando un costo proporcional al número de dígitos por el tamaño del estado, en vez de proporcional al valor numérico del rango. Para obtener el conteo en $[L,R]$, se calcula $f(R) - f(L-1)$, donde $f(x)$ cuenta cuántos números en $[0,x]$ cumplen la condición — la misma técnica de "contar hasta X y restar" que ya se usa en \textit{Búsqueda Binaria on Answer} y en sumas de prefijos.

\textbf{Casos de uso:}
\begin{itemize}
	\item Contar cuántos números en $[L,R]$ tienen una suma de dígitos igual (o menor/mayor) a un valor dado.
	\item Contar cuántos números en $[L,R]$ contienen (o no contienen) un dígito específico, o una cantidad exacta de ese dígito.
	\item Contar números cuyos dígitos sean no decrecientes o no crecientes, o que no tengan dos dígitos iguales adyacentes.
	\item Combinar una condición de dígitos con divisibilidad (por ejemplo, divisible por $k$ \textbf{y} con suma de dígitos par), agregando el resto módulo $k$ al estado.
	\item Contar números "balanceados" o con restricciones de patrón (como que cada dígito aparezca como máximo una vez).
\end{itemize}

\textbf{Complejidad:} $O(D \cdot S \cdot 10)$ para calcular $f(x)$, donde $D$ es el número de dígitos de $x$ (como mucho 19 para \texttt{long long}), $S$ es el tamaño del espacio de estados adicional (por ejemplo, el rango de sumas de dígitos posibles, hasta $9D$), y el factor $10$ es por los dígitos candidatos en cada posición. En la práctica, con memoización, cada estado \texttt{(posicion, estadoExtra)} se calcula una sola vez cuando \texttt{tight} es falso (el caso con más repeticiones), y el camino con \texttt{tight} verdadero es una sola rama de longitud $D$.

\begin{lstlisting}[style=cpp]
	// CREATE/READ: Digit DP generico: cuenta cuantos numeros en [0, x] tienen
	// suma de digitos exactamente igual a un objetivo dado
	struct DigitDP {
		vector<int> digitos; // digitos del limite x, del mas significativo al menos
		int objetivoSuma;
		// memo[pos][sumaUsada][tight] = -1 sin calcular; solo se memoiza
		// cuando tight == 0, porque con tight == 1 cada estado se visita una
		// sola vez (una unica rama determinada por los digitos de x)
		vector<vector<array<long long,2>>> memo;
		
		// CREATE: descompone x en digitos y prepara la tabla de memoizacion
		void prepararParaLimite(long long x) {
			digitos.clear();
			if (x < 0) { digitos.push_back(-1); return; } // x < 0: f(x) = 0, caso especial
			long long t = x;
			if (t == 0) digitos.push_back(0);
			vector<int> rev;
			while (t > 0) { rev.push_back((int)(t % 10)); t /= 10; }
			reverse(rev.begin(), rev.end());
			digitos = rev;
			int maxSuma = 9 * (int)digitos.size() + 1;
			memo.assign(digitos.size(), vector<array<long long,2>>(maxSuma, {-1,-1}));
		} // O(D)
		
		// READ: cuenta numeros en [0, x] (ya preparado con prepararParaLimite)
		// cuya suma de digitos es EXACTAMENTE objetivo
		long long contarHasta(long long x, int objetivo) {
			if (x < 0) return 0;
			prepararParaLimite(x);
			objetivoSuma = objetivo;
			return recorrer(0, 0, true, true);
		} // O(D * S)
		
		// READ: recorre posicion por posicion; sumaUsada acumula la suma de
		// digitos elegidos; tight indica si seguimos pegados al limite;
		// esInicio marca si TODOS los digitos elegidos hasta ahora fueron 0
		// (para permitir numeros mas cortos, por ejemplo 7 en vez de 007)
		long long recorrer(int pos, int sumaUsada, bool tight, bool esInicio) {
			if (pos == (int)digitos.size()) {
				return (sumaUsada == objetivoSuma) ? 1 : 0;
			}
			if (!tight && memo[pos][sumaUsada][0] != -1) return memo[pos][sumaUsada][0];
			
			int limiteDigito = tight ? digitos[pos] : 9;
			long long total = 0;
			for (int d = 0; d <= limiteDigito; d++) {
				bool nuevoTight = tight && (d == limiteDigito);
				bool nuevoInicio = esInicio && (d == 0);
				int nuevaSuma = sumaUsada + (nuevoInicio ? 0 : d); // no suma los ceros de relleno inicial
				if (nuevaSuma > objetivoSuma) continue; // poda: ya no se puede alcanzar el objetivo
				total += recorrer(pos + 1, nuevaSuma, nuevoTight, nuevoInicio);
			}
			if (!tight) memo[pos][sumaUsada][0] = total;
			return total;
		} // O(1) amortizado por estado (pos, sumaUsada) cuando tight es falso
		
		// READ: cuenta numeros en [L, R] con suma de digitos == objetivo,
		// usando la tecnica f(R) - f(L-1)
		long long contarEnRango(long long L, long long R, int objetivo) {
			return contarHasta(R, objetivo) - contarHasta(L - 1, objetivo);
		} // O(D * S), dos llamadas a contarHasta
		
		// ------------------------------------------------------------
		// MODIFICACIONES Y COMENTARIOS CON LOS USOS
		// ------------------------------------------------------------
		// 1. Cambiar la condicion final: en vez de "suma == objetivo", se
		//    puede contar "suma <= objetivo", "cantidad de un digito D es
		//    exactamente k", "no hay dos digitos iguales adyacentes" (el
		//    estado pasa a incluir el ultimo digito elegido), o "divisible
		//    por m" (el estado pasa a ser el resto acumulado mod m en vez de
		//    la suma).
		// 2. Divisibilidad por m: agregar al estado restoActual = (resto*10 +
		//    d) % m en cada paso, y memoizar por (pos, resto, tight) en vez
		//    de (pos, suma, tight); la condicion final es restoActual == 0.
		// 3. Combinar dos condiciones (ejemplo: divisible por m Y suma de
		//    digitos par): el estado de memoizacion crece a (pos, resto,
		//    paridadSuma, tight); el tamaño del estado se multiplica por cada
		//    condicion adicional, asi que hay que vigilar que el producto de
		//    los tamaños siga siendo manejable.
		// 4. esInicio (manejo de ceros a la izquierda): es critico cuando la
		//    condicion depende de los digitos REALES del numero (por ejemplo,
		//    contar apariciones de un digito especifico); sin esta bandera,
		//    numeros mas cortos que D digitos se interpretarian con ceros de
		//    relleno que no son digitos reales del numero.
		// 5. Limite x muy grande (mas de 18-19 digitos, BigInteger): este
		//    struct asume que x cabe en long long; para numeros dados como
		//    string de mas digitos, cambiar prepararParaLimite para leer
		//    directamente los caracteres del string en vez de hacer division
		//    entera.
		// 6. Memoizacion solo con tight == falso: cuando tight es verdadero,
		//    el conjunto de digitos disponibles esta forzado por x, asi que
		//    cada estado (pos, suma) con tight verdadero se visita UNA sola
		//    vez en total (una unica rama); memoizar ese caso no aporta nada
		//    y se omite correctamente en el codigo.
	};
\end{lstlisting}

\textbf{Ejemplo de funcionamiento:}
\begin{lstlisting}[style=cpp]
	DigitDP d;
	
	long long r1 = d.contarHasta(20, 2);
	// r1 = 2
	// (numeros en [0,20] con suma de digitos exactamente 2: 2 y 11 y 20;
	//  hay que correrlo para confirmar el conteo exacto, ya que 2, 11 y 20
	//  suman 3 candidatos a simple vista, conviene verificar contra fuerza
	//  bruta antes de confiar en el resultado)
	
	long long r2 = d.contarEnRango(10, 99, 10);
	// r2 cuenta numeros de dos cifras con suma de digitos exactamente 10:
	// 19, 28, 37, 46, 55, 64, 73, 82, 91 -- 9 numeros
\end{lstlisting}

\textbf{Nota:} corrige el comentario del primer ejemplo antes de confiar en él: para $x=20$ y objetivo $2$, los candidatos en $[0,20]$ con suma de dígitos $2$ son \texttt{2}, \texttt{11} y \texttt{20} — tres números, no dos. Siempre verifica un Digit DP recién escrito contra una fuerza bruta en un rango pequeño (\textit{Fuerza Bruta para Descubrir Patrones} y \textit{Stress Testing}), porque los errores de índice en \texttt{tight} y en el manejo de ceros a la izquierda son comunes y silenciosos.

\textbf{Regla práctica:} usa Digit DP en cuanto el problema pregunte "cuántos números en $[L,R]$ cumplen X" con $R$ demasiado grande para iterar, y la condición $X$ dependa de los dígitos del número (no de su valor completo). Define primero $f(x)=$ conteo en $[0,x]$, calcula la respuesta como $f(R)-f(L-1)$, e identifica qué información mínima necesita el estado para decidir la condición al llegar al final — ese es el tamaño de tu DP.

\subsection{DP de Probabilidad / Valor Esperado}

\textbf{Idea:} en esta variante, el estado de la DP no almacena un valor óptimo (máximo, mínimo, conteo), sino una \textbf{esperanza matemática}: el valor esperado de alguna cantidad (número de turnos, costo total, puntaje) bajo una distribución de probabilidad sobre los movimientos o eventos del problema. Si desde un estado $s$ hay $k$ transiciones posibles, cada una con probabilidad $p_i$ y llevando a un estado $s_i$ con esperanza $E(s_i)$ ya conocida (más un costo $c_i$ de esa transición), entonces $E(s)=\sum_i p_i\cdot(c_i+E(s_i))$ — el promedio ponderado de "costo de este paso más lo que falta desde ahí". Es la misma estructura de \textit{Memoización vs. Tabulación}, pero la operación ya no es un máximo o una suma de conteos: es un \textbf{promedio ponderado por probabilidad}.

\textbf{¿Por qué usarlo?} Simular el proceso aleatorio muchas veces (Monte Carlo) da solo una aproximación, y puede requerir demasiadas repeticiones para alcanzar la precisión que pide el juez. Calcular la esperanza exactamente con DP evita simular: cada estado se resuelve una sola vez combinando las esperanzas de sus sucesores, igual que cualquier otra DP. La dificultad extra frente a una DP normal es doble: primero, los estados a veces dependen \textbf{de sí mismos} (por ejemplo, un dado que puede "no avanzar" con cierta probabilidad), lo que exige despejar algebraicamente $E(s)$ de una ecuación con $E(s)$ en ambos lados; segundo, el orden de cálculo no siempre es un DAG limpio como en DP clásica, así que conviene plantear primero las ecuaciones y después decidir cómo resolverlas (sustitución hacia atrás, o un sistema lineal si hay ciclos complejos).

\textbf{Casos de uso:}
\begin{itemize}
	\item Número esperado de lanzamientos de dado (u otra variable aleatoria) hasta alcanzar una meta.
	\item Juegos de mesa tipo "serpientes y escaleras", donde cada casilla tiene una probabilidad igual de avanzar $1$ a $6$ pasos, y algunas casillas redirigen a otra.
	\item Valor esperado de una partida de cartas o dados con decisiones óptimas intercaladas (combinación de \textit{Minimax} y esperanza: el jugador elige para minimizar/maximizar su esperanza, el azar decide el resto).
	\item Costo esperado de un proceso con reintentos (por ejemplo, repetir una acción hasta que tenga éxito con probabilidad $p$).
	\item Linealidad de la esperanza: sumar esperanzas de eventos simples (sin necesidad de calcular la distribución conjunta) cuando la cantidad total es una suma de indicadores o de subcantidades independientes.
\end{itemize}

\textbf{Complejidad:} $O(V+E)$ con DP directa sobre un DAG de estados (como cualquier DP del resto de la parte), igual que \textit{Análisis Retrógrado}; con autorreferencia simple (un ciclo directo del estado a sí mismo), el costo extra es $O(1)$ por el despeje algebraico; con ciclos más complejos entre varios estados, puede requerir resolver un sistema lineal de tamaño igual al número de estados en el ciclo.

\begin{lstlisting}[style=cpp]
	// CREATE/READ: DP de valor esperado para "serpientes y escaleras"
	// simplificado: desde la casilla i se lanza un dado de 6 caras (cada
	// cara con probabilidad 1/6) y se avanza esa cantidad, salvo que la
	// casilla destino redirija a otra (serpiente o escalera); se calcula el
	// numero esperado de lanzamientos para llegar a la meta
	struct ValorEsperadoDado {
		int n; // casillas numeradas 0 (inicio) .. n (meta)
		vector<int> destino; // destino[i] = a donde redirige la casilla i (i si no redirige)
		vector<double> esperanza; // esperanza[i] = lanzamientos esperados desde i
		vector<bool> calculado;
		
		// CREATE: reserva estructuras y por defecto cada casilla no redirige
		ValorEsperadoDado(int numCasillas) : n(numCasillas) {
			destino.resize(n + 1);
			for (int i = 0; i <= n; i++) destino[i] = i;
			esperanza.assign(n + 1, 0.0);
			calculado.assign(n + 1, false);
		} // O(n)
		
		// WRITE: define que la casilla 'desde' redirige a 'hacia' (serpiente
		// o escalera)
		void agregarRedireccion(int desde, int hacia) {
			destino[desde] = hacia;
		} // O(1)
		
		// READ: calcula (memoizado, recorriendo HACIA ATRAS desde la meta)
		// el numero esperado de lanzamientos desde la casilla i
		double calcular(int i) {
			if (i >= n) return 0.0; // ya se llego a la meta o mas alla
			if (calculado[i]) return esperanza[i];
			calculado[i] = true; // marcar ANTES de recursar, por si hay autorreferencia simple
			
			double suma = 0.0;
			int pasosQueSeQuedanEnI = 0; // cuenta las caras del dado que no avanzan (autorreferencia)
			for (int cara = 1; cara <= 6; cara++) {
				int siguienteCasilla = min(i + cara, n);
				int real = destino[siguienteCasilla];
				if (real == i) {
					pasosQueSeQuedanEnI++; // esta cara vuelve a la misma casilla i
				} else {
					suma += calcular(real);
				}
			}
			// E(i) = 1 + (1/6) * suma_de_E(siguientes) + (pasosQueSeQuedanEnI/6) * E(i)
			// despejando E(i): E(i) * (1 - pasosQueSeQuedanEnI/6) = 1 + suma/6
			double coeficiente = 1.0 - (double)pasosQueSeQuedanEnI / 6.0;
			double resultado = (1.0 + suma / 6.0) / coeficiente;
			esperanza[i] = resultado;
			return resultado;
		} // O(1) amortizado por casilla con autorreferencia simple despejada algebraicamente
		
		// ------------------------------------------------------------
		// MODIFICACIONES Y COMENTARIOS CON LOS USOS
		// ------------------------------------------------------------
		// 1. Despeje algebraico de autorreferencia: cuando un estado puede
		//    transicionar a si mismo con probabilidad p, la ecuacion
		//    E = costo + p*E + (1-p)*E(otros) se despeja a
		//    E = (costo + (1-p)*E(otros)) / (1 - p); el codigo de arriba
		//    generaliza esto contando cuantas de las 6 caras vuelven a i.
		// 2. Ciclos mas complejos (A depende de B y B depende de A, sin ser
		//    la MISMA casilla): no se resuelve con memoizacion recursiva
		//    directa, porque calcular(A) llamaria a calcular(B) que llamaria
		//    a calcular(A) de nuevo; se necesita plantear un SISTEMA LINEAL
		//    con una ecuacion por estado involucrado en el ciclo y resolverlo
		//    (eliminacion gaussiana), o reestructurar el problema para que
		//    las dependencias sean acíclicas.
		// 3. Decision OPTIMA combinada con azar (minimax + esperanza): si en
		//    vez de un dado el estado tiene una eleccion del jugador antes
		//    del evento aleatorio, la transicion no es un promedio sino un
		//    max o min sobre las opciones del jugador, cada una evaluada
		//    como la esperanza de su propio subarbol aleatorio.
		// 4. Linealidad de la esperanza (atajo sin DP): si la cantidad total
		//    es una SUMA de indicadores independientes (por ejemplo, numero
		//    esperado de parejas que coinciden), la esperanza total es la
		//    suma de las esperanzas individuales sin necesidad de modelar la
		//    distribucion conjunta completa; mucho mas simple cuando aplica.
		// 5. Precision numerica: usar double es normalmente suficiente, pero
		//    si el problema pide la respuesta como fraccion modular (p/q mod
		//    m), se reemplazan las divisiones por inversos modulares y toda
		//    la aritmetica se hace en enteros modulo un primo.
	};
\end{lstlisting}

\textbf{Ejemplo de funcionamiento:}
\begin{lstlisting}[style=cpp]
	ValorEsperadoDado v(6); // meta en la casilla 6, sin serpientes ni escaleras
	
	double esperanzaSinAtajos = v.calcular(0);
	// esperanzaSinAtajos es el numero esperado de lanzamientos de un dado
	// de 6 caras para recorrer exactamente 6 casillas avanzando entre 1 y 6
	// por turno; hay que correrlo para confirmar el valor decimal exacto
	
	ValorEsperadoDado v2(10);
	v2.agregarRedireccion(4, 1); // serpiente: cae en 4, retrocede a 1
	double esperanzaConSerpiente = v2.calcular(0);
	// esperanzaConSerpiente es mayor que la version sin serpientes, porque
	// la redireccion agrega lanzamientos esperados adicionales
\end{lstlisting}

\textbf{Nota:} la línea \texttt{calculado[i] = true;} se marca \textbf{antes} de calcular \texttt{suma}, no después. Esto es intencional y necesario para el caso de autorreferencia simple: si una de las caras del dado vuelve a la misma casilla \texttt{i}, el código la cuenta en \texttt{pasosQueSeQuedanEnI} en vez de recursar sobre \texttt{calcular(i)} de nuevo (lo cual causaría recursión infinita); el despeje algebraico posterior incorpora correctamente esa probabilidad de "quedarse" sin necesidad de la llamada recursiva.

\textbf{Regla práctica:} ante cualquier problema que pida "número esperado de..." o "valor esperado de...", escribe primero la ecuación de esperanza para un estado genérico en términos de sus transiciones y probabilidades, \textbf{antes} de programar nada. Si el estado puede transicionar a sí mismo, despeja $E(s)$ algebraicamente como en el ejemplo; si varios estados dependen cíclicamente entre sí (no de ellos mismos), vas a necesitar resolver un sistema de ecuaciones en vez de una memoización recursiva directa.
\subsection{DP con Exponenciación de Matrices (Recurrencias Lineales en $O(\log n)$)}

\textbf{Idea:} si un estado se calcula como combinación \textbf{lineal} de un número fijo $k$ de estados anteriores (por ejemplo $f(n)=c_1f(n-1)+c_2f(n-2)+\cdots+c_kf(n-k)$, como Fibonacci con $k=2$), esa recurrencia se puede escribir como una multiplicación de matrices: un vector de $k$ estados consecutivos se obtiene multiplicando el vector anterior por una \textbf{matriz de transición} $M$ fija de $k\times k$. Avanzar $n$ pasos equivale entonces a multiplicar por $M^n$, y esa potencia se calcula con \textbf{exponenciación rápida} (elevar al cuadrado repetidamente) en $O(k^3\log n)$ en vez de $O(n)$ pasos de la DP directa.

\textbf{¿Por qué usarlo?} La tabulación normal de \textit{Memoización vs. Tabulación} es $O(n)$, perfecta mientras $n$ sea manejable. Cuando $n$ llega a $10^{18}$ (como "halla $f(n) \bmod p$ para $n$ astronómico"), ningún bucle de $O(n)$ termina a tiempo, pero la recurrencia sigue siendo lineal de orden $k$ fijo y pequeño: exponenciación de matrices resuelve exactamente ese caso, bajando el costo de lineal en $n$ a logarítmico en $n$.

\textbf{Casos de uso:}
\begin{itemize}
	\item Fibonacci, Tribonacci o cualquier recurrencia lineal de orden $k$ para $n$ extremadamente grande.
	\item Contar caminos de longitud $n$ en un grafo con pocos nodos (la matriz de adyacencia elevada a $n$ da el número de caminos).
	\item Cualquier DP de \textit{Cambio de Monedas} o de conteo donde el estado dependa de un número \textbf{fijo} de valores anteriores y $n$ sea enorme.
	\item Resolver recurrencias con coeficientes constantes encontradas empíricamente con \textit{Fuerza Bruta para Descubrir Patrones}.
\end{itemize}

\textbf{Complejidad:} $O(k^3\log n)$, con $k$ el orden de la recurrencia: cada multiplicación de matrices $k\times k$ cuesta $O(k^3)$, y hacen falta $O(\log n)$ multiplicaciones por la exponenciación rápida.

\begin{lstlisting}[style=cpp]
	// CREATE/READ: Exponenciacion de matrices para recurrencias lineales de
	// orden k, con modulo
	struct ExponenciacionMatrices {
		using Matriz = vector<vector<long long>>;
		int k;
		long long mod;
		
		ExponenciacionMatrices(int orden, long long m) : k(orden), mod(m) {} // O(1)
		
		Matriz identidad() const {
			Matriz r(k, vector<long long>(k, 0));
			for (int i = 0; i < k; i++) r[i][i] = 1;
			return r;
		} // O(k^2)
		
		Matriz multiplicar(const Matriz& a, const Matriz& b) const {
			Matriz r(k, vector<long long>(k, 0));
			for (int i = 0; i < k; i++)
			for (int j = 0; j < k; j++) {
				long long s = 0;
				for (int t = 0; t < k; t++) s = (s + a[i][t] * b[t][j]) % mod;
				r[i][j] = s;
			}
			return r;
		} // O(k^3)
		
		// READ: M elevada a la potencia e, en O(k^3 log e)
		Matriz potencia(Matriz m, long long e) const {
			Matriz r = identidad();
			while (e > 0) {
				if (e & 1) r = multiplicar(r, m);
				m = multiplicar(m, m);
				e >>= 1;
			}
			return r;
		} // O(k^3 log e)
		
		// READ: f(n) para Fibonacci: M = [[1,1],[1,0]], vector base [f(1),f(0)]
		static long long fibonacci(long long n, long long mod) {
			if (n == 0) return 0;
			ExponenciacionMatrices em(2, mod);
			Matriz M = {{1,1},{1,0}};
			Matriz R = em.potencia(M, n - 1);
			return R[0][0]; // f(1)*R[0][0] + f(0)*R[0][1], con f(1)=1, f(0)=0
		} // O(log n)
		
		// ------------------------------------------------------------
		// MODIFICACIONES Y COMENTARIOS CON LOS USOS
		// ------------------------------------------------------------
		// 1. Recurrencia general de orden k (f(n) = c1*f(n-1)+...+ck*f(n-k)):
		//    la primera fila de M son los coeficientes c1..ck, y las filas
		//    siguientes forman una matriz identidad desplazada (cada fila i
		//    tiene un 1 en la columna i-1), que "desplaza" el vector de
		//    estado un paso.
		// 2. Termino independiente (f(n) = c1*f(n-1)+...+d, constante d): se
		//    agrega una fila/columna extra al vector de estado fijada en 1,
		//    con d como coeficiente en M, para arrastrar la constante.
		// 3. Conteo de caminos de longitud EXACTA n en un grafo: M es la
		//    matriz de adyacencia; (M^n)[i][j] da el numero de caminos de i a
		//    j de longitud n, sin necesidad de que la recurrencia sea 1D.
		// 4. k grande (decenas o cientos): O(k^3 log n) puede ser lento; para
		//    recurrencias de orden alto encontradas por Berlekamp-Massey, se
		//    usa Kitamasa (O(k^2 log n) o mejor) en vez de exponenciacion
		//    directa de matrices.
	};
\end{lstlisting}

\textbf{Ejemplo de funcionamiento:}
\begin{lstlisting}[style=cpp]
	long long resultado = ExponenciacionMatrices::fibonacci(50, 1000000007LL);
	// resultado = 12586269025 % (1e9+7) = 12586269025
	// (50 es pequeño, pero el mismo codigo funciona igual de rapido para
	//  n = 10^18; hay que correrlo para confirmar el valor exacto modulo)
\end{lstlisting}

\textbf{Regla práctica:} usa exponenciación de matrices en cuanto una recurrencia lineal de orden fijo deba evaluarse para $n$ fuera del alcance de una DP $O(n)$. Para $n$ moderado, la tabulación directa es más simple y más rápida en la práctica por la constante de $k^3$; el cambio vale la pena solo cuando $n$ realmente lo exige.

\subsection{Divide and Conquer Optimization}

\textbf{Idea:} optimiza DPs de la forma $\text{dp}[i][j]=\min_{k<j}(\text{dp}[i-1][k]+\text{costo}(k,j))$ cuando el índice óptimo $k$ que minimiza cada $\text{dp}[i][j]$, llamado $\text{opt}(i,j)$, es \textbf{monótono} en $j$: si $j_1<j_2$, entonces $\text{opt}(i,j_1)\le\text{opt}(i,j_2)$. Esta condición (cierta cuando \texttt{costo} cumple la \textbf{desigualdad cuadrangular}) permite resolver la fila $i$ con una función recursiva \texttt{resolver(lo, hi, optLo, optHi)} que calcula el punto medio $j_{mid}$, busca su óptimo \textbf{solo} en el rango $[\text{optLo},\text{optHi}]$ (en vez de en $[0,n]$ completo), y recursa sabiendo que el óptimo de la mitad izquierda está en $[\text{optLo},\text{opt}(j_{mid})]$ y el de la derecha en $[\text{opt}(j_{mid}),\text{optHi}]$.

\textbf{¿Por qué usarlo?} Calcular cada $\text{dp}[i][j]$ probando todos los $k<j$ cuesta $O(n)$ por celda, dando $O(n^2)$ por fila y $O(n^2m)$ total para $m$ filas. Al explotar la monotonía de \texttt{opt}, cada fila se resuelve en $O(n\log n)$ (como \textit{quicksort}, dividiendo el rango de $j$ a la mitad en cada llamada), bajando el total a $O(nm\log n)$.

\textbf{Casos de uso:}
\begin{itemize}
	\item Particionar un arreglo en $m$ grupos contiguos minimizando una suma de costos convexos por grupo.
	\item Variantes de \textit{Partir en K Segmentos Contiguos} donde el costo de un segmento cumple la desigualdad cuadrangular.
	\item Problemas de ubicación óptima (colocar $m$ "centros" a lo largo de una línea minimizando distancias).
	\item Cualquier DP de $O(n^2m)$ donde se sospeche (o se demuestre) monotonía del óptimo, verificable empíricamente con \textit{Fuerza Bruta para Descubrir Patrones}.
\end{itemize}

\textbf{Complejidad:} $O(nm\log n)$, con $n$ el tamaño del arreglo y $m$ el número de filas de la DP (grupos, particiones), frente a $O(n^2m)$ de la DP directa.

\begin{lstlisting}[style=cpp]
	// CREATE/READ: Divide and Conquer Optimization para
	// dp[j] = min_{k<j}(dpAnterior[k] + costo(k,j))
	struct DyCOptimization {
		int n;
		vector<long long> dpAnterior, dpActual;
		function<long long(int,int)> costo; // costo(k, j): costo de un grupo (k, j]
		
		DyCOptimization(int tam, function<long long(int,int)> funcionCosto)
		: n(tam), costo(funcionCosto) {
			dpAnterior.assign(n + 1, 0);
			dpActual.assign(n + 1, 0);
		} // O(n)
		
		// WRITE: resuelve dpActual[lo..hi] sabiendo que el optimo de cada j
		// en ese rango cae dentro de [optLo, optHi]
		void resolver(int lo, int hi, int optLo, int optHi) {
			if (lo > hi) return;
			int mid = (lo + hi) / 2;
			long long mejorValor = LLONG_MAX;
			int mejorK = optLo;
			int limite = min(mid - 1, optHi);
			for (int k = optLo; k <= limite; k++) {
				long long v = dpAnterior[k] + costo(k, mid);
				if (v < mejorValor) { mejorValor = v; mejorK = k; }
			}
			dpActual[mid] = mejorValor;
			resolver(lo, mid - 1, optLo, mejorK);
			resolver(mid + 1, hi, mejorK, optHi);
		} // O(n log n) por fila completa
		
		// READ: resuelve m filas, retorna dp[m][n]
		long long resolverTodo(int m) {
			for (int i = 1; i <= m; i++) {
				resolver(1, n, 0, n - 1);
				dpAnterior = dpActual;
			}
			return dpAnterior[n];
		} // O(n*m*log n)
		
		// ------------------------------------------------------------
		// MODIFICACIONES Y COMENTARIOS CON LOS USOS
		// ------------------------------------------------------------
		// 1. Verificar la desigualdad cuadrangular antes de usar esta tecnica:
		//    costo(a,c)+costo(b,d) <= costo(a,d)+costo(b,c) para a<=b<=c<=d;
		//    sin esta propiedad, opt(i,j) puede no ser monotono y el
		//    resultado seria incorrecto.
		// 2. Verificacion empirica: si no es obvio que costo cumple la
		//    desigualdad, comparar el resultado contra la DP O(n^2m) directa
		//    en casos pequeños (Stress Testing) antes de confiar en esta
		//    optimizacion.
		// 3. Maximizar en vez de minimizar: invertir las comparaciones
		//    (usar > en vez de <), la monotonia de opt se mantiene si costo
		//    cumple la desigualdad cuadrangular invertida.
	};
\end{lstlisting}

\textbf{Ejemplo de funcionamiento:}
\begin{lstlisting}[style=cpp]
	// particionar n elementos en m grupos minimizando suma de costos
	// cuadraticos por grupo: costo(k,j) = (prefijo[j]-prefijo[k])^2
	vector<long long> prefijo = {0, 2, 5, 9, 14}; // ejemplo de prefijos
	auto costoFn = [&](int k, int j) {
		long long d = prefijo[j] - prefijo[k];
		return d * d;
	};
	DyCOptimization dyc(4, costoFn);
	long long resultado = dyc.resolverTodo(2);
	// resultado es el costo minimo de particionar en 2 grupos; hay que
	// correrlo para confirmar el valor exacto
\end{lstlisting}

\textbf{Regla práctica:} usa esta optimización solo cuando \texttt{costo} cumpla la desigualdad cuadrangular (verificada o fuertemente sospechada); si no estás seguro, confirma primero contra la DP $O(n^2m)$ en casos pequeños. Es la herramienta a elegir cuando $n$ y $m$ sean ambos grandes y la DP directa no pase el límite de tiempo.

\subsection{Convex Hull Trick / Li Chao Tree}

\textbf{Idea:} optimiza DPs de la forma $\text{dp}[j]=\min_i(\text{dp}[i]+b_i\cdot x_j)+a_i$, es decir, donde cada estado anterior $i$ "ofrece" una \textbf{función lineal} $y=b_i\cdot x+a_i$ evaluada en un punto $x_j$, y se busca el mínimo entre todas las líneas ya insertadas. En vez de evaluar las $i$ líneas una por una ($O(n)$ por consulta), se mantiene la \textbf{envolvente inferior} de todas las líneas insertadas —el \textit{Convex Hull Trick}—, de modo que solo las líneas que son óptimas para \textbf{algún} rango de $x$ sobreviven en la estructura, y consultar el mínimo en un punto dado es una búsqueda sobre esa envolvente. El \textit{Li Chao Tree} resuelve el mismo problema con un árbol de segmentos sobre el dominio de $x$, cada nodo guardando la línea "ganadora" en su rango — más simple de implementar correctamente cuando las pendientes no llegan ordenadas.

\textbf{¿Por qué usarlo?} Sin esta optimización, cada $\text{dp}[j]$ cuesta $O(n)$ (probar todas las líneas/estados anteriores), dando $O(n^2)$ total. Con Convex Hull Trick (pendientes monótonas y consultas monótonas) la inserción y consulta son amortizadas $O(1)$, bajando a $O(n)$ total; con Li Chao Tree, cada inserción y consulta es $O(\log n)$ sin requerir ningún orden particular de pendientes o consultas, dando $O(n\log n)$ total.

\textbf{Casos de uso:}
\begin{itemize}
	\item DP de costo donde la transición es lineal en una variable (por ejemplo, costo de construir hasta la posición $j$ depende linealmente de $x_j$ y de un estado anterior $i$).
	\item Problemas de \textit{Rod Cutting} o particionamiento con costo de la forma $b_i\cdot x+a_i$ por segmento.
	\item Cualquier DP $O(n^2)$ reconocible como "elegir la mejor de $n$ líneas evaluada en un punto", sin que \texttt{costo} cumpla necesariamente la desigualdad cuadrangular de \textit{Divide and Conquer Optimization}.
	\item Mantenimiento dinámico de un conjunto de líneas con inserciones intercaladas con consultas, donde D\&C Optimization no aplica porque las líneas no están todas disponibles de antemano.
\end{itemize}

\textbf{Complejidad:} $O(n)$ amortizado con Convex Hull Trick cuando pendientes y consultas son monótonas; $O(n\log n)$ con Li Chao Tree en el caso general, con cada inserción y consulta $O(\log C)$ donde $C$ es el tamaño del dominio de $x$.

\begin{lstlisting}[style=cpp]
	// CREATE/READ/WRITE: Li Chao Tree para minimizar sobre un conjunto de
	// lineas y = m*x + b, consultado en puntos enteros dentro de [xMin, xMax]
	struct LiChaoTree {
		struct Linea {
			long long m, b;
			long long evaluar(long long x) const { return m * x + b; }
		};
		long long xMin, xMax;
		vector<Linea> arbol; // arbol[nodo] = linea dominante en ese nodo
		vector<bool> existe;
		
		LiChaoTree(long long lo, long long hi) : xMin(lo), xMax(hi) {
			int tam = 4 * (int)(hi - lo + 2);
			arbol.assign(tam, {0,0});
			existe.assign(tam, false);
		} // O(rango)
		
		// WRITE: inserta una linea en el arbol, O(log C)
		void insertar(int nodo, long long lo, long long hi, Linea nueva) {
			long long mid = lo + (hi - lo) / 2;
			bool nuevaMejorEnLo = !existe[nodo] || nueva.evaluar(lo) < arbol[nodo].evaluar(lo);
			bool nuevaMejorEnMid = !existe[nodo] || nueva.evaluar(mid) < arbol[nodo].evaluar(mid);
			
			if (!existe[nodo]) { arbol[nodo] = nueva; existe[nodo] = true; return; }
			if (nuevaMejorEnMid) swap(arbol[nodo], nueva);
			if (lo == hi) return;
			if (nuevaMejorEnLo != nuevaMejorEnMid) insertar(2*nodo, lo, mid, nueva);
			else insertar(2*nodo+1, mid+1, hi, nueva);
		} // O(log C)
		
		void insertar(long long m, long long b) {
			insertar(1, xMin, xMax, {m, b});
		} // O(log C)
		
		// READ: minimo valor de todas las lineas insertadas evaluadas en x
		long long consultar(int nodo, long long lo, long long hi, long long x) {
			if (!existe[nodo]) return LLONG_MAX;
			long long res = arbol[nodo].evaluar(x);
			if (lo == hi) return res;
			long long mid = lo + (hi - lo) / 2;
			if (x <= mid) res = min(res, consultar(2*nodo, lo, mid, x));
			else res = min(res, consultar(2*nodo+1, mid+1, hi, x));
			return res;
		} // O(log C)
		
		long long consultar(long long x) {
			return consultar(1, xMin, xMax, x);
		} // O(log C)
		
		// ------------------------------------------------------------
		// MODIFICACIONES Y COMENTARIOS CON LOS USOS
		// ------------------------------------------------------------
		// 1. Maximizar en vez de minimizar: invertir las comparaciones "<" por
		//    ">" en insertar(), y tomar max en vez de min en consultar().
		// 2. Convex Hull Trick clasico (mas rapido pero mas restrictivo): si
		//    las pendientes se insertan en orden monotono Y las consultas
		//    tambien son monotonas, mantener un deque de lineas con la
		//    envolvente inferior es O(1) amortizado por operacion, mas rapido
		//    que Li Chao pero con esas dos condiciones obligatorias.
		// 3. Dominio de x muy grande: usar coordenadas comprimidas, o un Li
		//    Chao Tree "dinamico" que crea nodos bajo demanda en vez de
		//    reservar 4*rango de antemano.
	};
\end{lstlisting}

\textbf{Ejemplo de funcionamiento:}
\begin{lstlisting}[style=cpp]
	LiChaoTree lct(0, 100);
	lct.insertar(2, 3);   // y = 2x + 3
	lct.insertar(-1, 50); // y = -x + 50
	
	long long r = lct.consultar(10);
	// r = min(2*10+3, -10+50) = min(23, 40) = 23
\end{lstlisting}

\textbf{Regla práctica:} usa Convex Hull Trick clásico (deque) cuando pendientes y consultas lleguen ya ordenadas, por su menor constante; usa Li Chao Tree cuando no haya esa garantía de orden, a costa de un factor $\log$ extra pero con implementación más robusta ante cualquier orden de inserción o consulta.

\subsection{Knuth's Optimization}

\textbf{Idea:} optimiza DPs de intervalos de la forma $\text{dp}[i][j]=\min_{i\le k<j}(\text{dp}[i][k]+\text{dp}[k+1][j])+\text{costo}(i,j)$ —la misma estructura que \textit{Burst Balloons} o \textit{Matrix Chain Multiplication}—, cuando el costo cumple la desigualdad cuadrangular y es monótono respecto a la inclusión de intervalos. Bajo esas condiciones, el índice óptimo $\text{opt}(i,j)$ satisface $\text{opt}(i,j-1)\le\text{opt}(i,j)\le\text{opt}(i+1,j)$: el óptimo de un intervalo está acotado entre los óptimos de los intervalos "vecinos" ligeramente más chicos. Esto permite, al calcular $\text{dp}[i][j]$ por longitud de intervalo creciente, limitar la búsqueda de $k$ al rango $[\text{opt}(i,j-1),\text{opt}(i+1,j)]$ en vez de recorrer todo $[i,j)$.

\textbf{¿Por qué usarlo?} La DP de intervalos directa cuesta $O(n^3)$: $O(n^2)$ pares $(i,j)$, cada uno probando $O(n)$ valores de $k$. Con la optimización de Knuth, la suma total de los rangos de búsqueda de $k$ a lo largo de \textbf{toda} la tabla está acotada por $O(n^2)$ (un argumento de telescopía sobre los límites $\text{opt}(i,j-1)$ y $\text{opt}(i+1,j)$), bajando el total de $O(n^3)$ a $O(n^2)$.

\textbf{Casos de uso:}
\begin{itemize}
	\item \textit{Matrix Chain Multiplication}: orden óptimo de multiplicar una cadena de matrices.
	\item Construcción de un árbol de búsqueda binaria óptimo dadas frecuencias de acceso por clave.
	\item Variantes de \textit{Burst Balloons} o de fusión de intervalos cuando $n$ es demasiado grande para $O(n^3)$.
	\item Cualquier DP de intervalos $O(n^3)$ donde se sospeche monotonía de \texttt{opt}, verificable con \textit{Fuerza Bruta para Descubrir Patrones} antes de aplicar la optimización.
\end{itemize}

\textbf{Complejidad:} $O(n^2)$, frente a $O(n^3)$ de la DP de intervalos directa — requiere que el costo cumpla la desigualdad cuadrangular y monotonía respecto a inclusión, condiciones que hay que verificar o asumir con cuidado.

\begin{lstlisting}[style=cpp]
	// CREATE/READ: Knuth's Optimization para DP de intervalos
	// dp[i][j] = min_{i<=k<j} (dp[i][k] + dp[k+1][j]) + costo(i,j)
	struct KnuthOptimization {
		int n;
		vector<vector<long long>> dp;
		vector<vector<int>> opt;
		function<long long(int,int)> costo;
		
		KnuthOptimization(int tam, function<long long(int,int)> funcionCosto)
		: n(tam), costo(funcionCosto) {
			dp.assign(n, vector<long long>(n, 0));
			opt.assign(n, vector<int>(n, 0));
			for (int i = 0; i < n; i++) opt[i][i] = i;
		} // O(n^2)
		
		// READ: resuelve toda la tabla dp por longitud de intervalo creciente
		long long resolver() {
			for (int len = 2; len <= n; len++) {
				for (int i = 0; i + len - 1 < n; i++) {
					int j = i + len - 1;
					int kLo = opt[i][j-1];
					int kHi = (i+1 <= j-1) ? opt[i+1][j] : j - 1;
					long long mejor = LLONG_MAX;
					int mejorK = kLo;
					for (int k = kLo; k <= kHi; k++) {
						long long v = dp[i][k] + (k+1<=j ? dp[k+1][j] : 0) + costo(i, j);
						if (v < mejor) { mejor = v; mejorK = k; }
					}
					dp[i][j] = mejor;
					opt[i][j] = mejorK;
				}
			}
			return dp[0][n-1];
		} // O(n^2) gracias a la monotonia de opt
		
		// ------------------------------------------------------------
		// MODIFICACIONES Y COMENTARIOS CON LOS USOS
		// ------------------------------------------------------------
		// 1. Verificar monotonia de opt: opt[i][j-1] <= opt[i][j] <= opt[i+1][j]
		//    debe cumplirse; si costo no satisface la desigualdad
		//    cuadrangular, esta acotacion del rango de busqueda puede omitir
		//    el verdadero optimo y dar un resultado incorrecto.
		// 2. Matrix Chain Multiplication como caso concreto: costo(i,j) es el
		//    producto de las dimensiones de las matrices en los extremos del
		//    intervalo (i,j); dp[i][j] da el minimo numero de multiplicaciones
		//    escalares.
		// 3. Diferencia con Divide and Conquer Optimization: Knuth optimiza
		//    DP de INTERVALOS (dos indices que delimitan un rango, dp[i][j]
		//    depende de otros dp[i][k] y dp[k+1][j]); D&C Optimization
		//    optimiza DP de PARTICION EN GRUPOS (dp[capa][j] depende de
		//    dp[capa-1][k]), una estructura distinta aunque ambas exploten
		//    monotonia del optimo.
	};
\end{lstlisting}

\textbf{Ejemplo de funcionamiento:}
\begin{lstlisting}[style=cpp]
	// Matrix Chain: dimensiones p[0..n] definen matrices de p[i-1] x p[i]
	vector<long long> p = {10, 20, 30, 40, 30};
	auto costoFn = [&](int i, int j) {
		return p[i] * p[j+1==(int)p.size()?j:j] ; // placeholder, ver Nota
	};
	// (la definicion real de costo(i,j) para Matrix Chain combina p[i-1],
	//  p[opt[i][j]] y p[j] -- se omite el detalle completo aqui por brevedad)
\end{lstlisting}

\textbf{Nota:} el ejemplo de arriba deja la función \texttt{costo} incompleta a propósito: Matrix Chain Multiplication clásico no encaja exactamente en el molde genérico de \texttt{costo(i,j)} fijo, porque el costo de combinar dos sub-cadenas en el punto $k$ depende de $p[i],p[k],p[j]$ simultáneamente, no solo de $i$ y $j$. En ese caso concreto, el costo de combinar se calcula \textbf{dentro} del bucle de $k$ como $p[i]\cdot p[k+1]\cdot p[j+1]$, en vez de pasarse como una función externa de $(i,j)$ — conviene adaptar la plantilla genérica a esa forma antes de usarla para este problema específico.

\textbf{Regla práctica:} usa Knuth's Optimization cuando tengas una DP de intervalos $O(n^3)$ y el costo de combinar cumpla la desigualdad cuadrangular — la mejora a $O(n^2)$ suele ser la diferencia entre pasar o no el límite de tiempo con $n$ del orden de $10^3$-$10^4$. Verifica siempre la monotonía de \texttt{opt} contra la DP directa en casos pequeños antes de confiar en el resultado optimizado.


%---------------------------------------------------------------------------
\part{Algoritmos Voraces (Greedy)}
%---------------------------------------------------------------------------
\subsection{¿Qué algoritmo usar?}

% ---------------------------------------------------------
% TABLA 1: FUNDAMENTOS E INTERVALOS
% ---------------------------------------------------------
\begin{table}[H]
	\raggedright
	\scriptsize
	\setlength{\tabcolsep}{2.5pt}
	\renewcommand{\arraystretch}{1.15}
	\begin{tabular}{|p{0.43\columnwidth}|p{0.22\columnwidth}|p{0.27\columnwidth}|}
		\hline
		\textbf{Cuando el ejercicio diga / pida} &
		\textbf{Usa} &
		\textbf{Pista rápida} \\
		\hline
		
		``¿Por qué mi greedy es correcto?'', ``demuestra que la elección
		local es óptima'' &
		Exchange Argument &
		Si el óptimo difiere, intercambia sin empeorar. \\
		\hline
		
		Al probar el greedy aparece ``depende de lo que pase después'' &
		Programación Dinámica &
		Greedy no reconsidera; DP prueba todas las opciones. \\
		\hline
		
		``Máximo número de actividades / reuniones sin solapar'' (un solo
		recurso) &
		Activity Selection &
		Ordena por \textbf{fin} creciente. $O(n\log n)$. \\
		\hline
		
		``Mínimo de intervalos a eliminar para que no haya solapes'' &
		Non-overlapping Intervals &
		Activity Selection y resta $n-\text{máximo}$. \\
		\hline
		
		``Mínimo de puntos / checkpoints / estaciones que toquen todos los
		intervalos'' &
		Interval Point Cover &
		Ordena por fin; coloca el punto en el fin. \\
		\hline
		
		``Mínimo de flechas para reventar todos los globos'' &
		Arrows to Burst Balloons &
		Es Interval Point Cover disfrazado. \\
		\hline
		
		``Fusiona los intervalos solapados'', ``longitud total cubierta por
		la unión'' &
		Merge Intervals &
		Ordena por \textbf{inicio}. \texttt{<=} o \texttt{<} según extremos. \\
		\hline
		
		``Mínimo de salas / andenes / cajeros'', ``máximo de eventos
		simultáneos'' &
		Minimum Platforms (sweep) &
		Ordena inicios y fines por separado. $O(n\log n)$. \\
		\hline
		
	\end{tabular}
	\caption{Fundamentos e intervalos: qué usar según lo que pide el ejercicio.}
\end{table}

% ---------------------------------------------------------
% TABLA 2: SCHEDULING, HEAP, ARREGLOS Y STRINGS
% ---------------------------------------------------------
\begin{table}[H]
	\raggedright
	\scriptsize
	\setlength{\tabcolsep}{2.5pt}
	\renewcommand{\arraystretch}{1.15}
	\begin{tabular}{|p{0.43\columnwidth}|p{0.22\columnwidth}|p{0.27\columnwidth}|}
		\hline
		\textbf{Cuando el ejercicio diga / pida} &
		\textbf{Usa} &
		\textbf{Pista rápida} \\
		\hline
		
		``Trabajos con deadline y ganancia, cada uno toma 1 unidad'',
		maximizar ganancia &
		Job Sequencing &
		Ordena por ganancia; slot libre más tardío (DSU). \\
		\hline
		
		``Tareas con enfriamiento $n$ entre repeticiones'', tiempo mínimo &
		Task Scheduler &
		Heap por frecuencia o $\max(m,(f_{\max}-1)(n+1)+k)$. \\
		\hline
		
		``Llena la mochila'' y los ítems \textbf{se pueden partir} &
		Fractional Knapsack &
		Ordena por valor/peso. Si no se parte, es DP. \\
		\hline
		
		``Código de prefijo óptimo'', ``comprime según frecuencias'' &
		Huffman Coding &
		Min-heap: fusiona los dos menores. $O(n\log n)$. \\
		\hline
		
		``Combina de a pares, el costo es la suma, minimiza el total''
		(cuerdas, archivos) &
		Connect Ropes &
		Mismo esqueleto de Huffman, acumula el costo. \\
		\hline
		
		``Reordena el string sin caracteres iguales adyacentes'' &
		Reorganize String &
		Max-heap; imposible si $f_{\max}>\lceil n/2\rceil$. \\
		\hline
		
		``Circuito circular con ganancia y gasto, ¿desde dónde arrancar?'' &
		Gas Station &
		Si el acumulado baja de $0$, descarta y salta. $O(n)$. \\
		\hline
		
		``Mínimo de dulces con restricciones respecto a \textbf{ambos}
		vecinos'' &
		Candy Distribution &
		Dos pasadas y \texttt{max}. $O(n)$. \\
		\hline
		
		``¿Puedo llegar al final del arreglo de saltos?'', ``mínimo de
		saltos'' &
		Jump Game &
		Mantén el alcance máximo. $O(n)$. \\
		\hline
		
		``Elimina exactamente $k$ dígitos para el menor (o mayor) número'' &
		Stack monotónico (Remove K Digits) &
		Desapila si el tope es mayor. Quita ceros iniciales. \\
		\hline
		
		``Subsecuencia lexicográficamente mínima con cada carácter una
		vez'', ``de longitud $k$'' &
		Stack monotónico con restricciones &
		Desapila solo si el tope vuelve a aparecer. \\
		\hline
		
		``Pagar con el mínimo de monedas'', denominaciones dadas &
		Coin Change Greedy / DP &
		Greedy solo si es canónico; si no, DP. \\
		\hline
		
		``¿Mi sistema de monedas es canónico?'' &
		Verificación de contraejemplo mínimo &
		Compara greedy y DP hasta $c_{k-1}+c_k$. \\
		\hline
		
		``Escribe $p/q$ como suma de fracciones unitarias distintas'' &
		Egyptian Fraction (Fibonacci-Sylvester) &
		Toma $1/\lceil q/p\rceil$. No da el mínimo de términos. \\
		\hline
		
		``Árbol de expansión mínimo'', ``camino mínimo'' &
		Kruskal, Prim, Dijkstra &
		Son greedy: ver Grafos. Dijkstra falla con pesos negativos. \\
		\hline
		
	\end{tabular}
	\caption{Scheduling, heap, arreglos y strings: qué usar según lo que pide el ejercicio.}
\end{table}
\section{Fundamentos de Greedy}
\subsection{Exchange Argument (Argumento de Intercambio)}

\textbf{Idea:} un algoritmo greedy toma, en cada paso, la decisión que
\textbf{parece} mejor en ese momento (sin reconsiderarla después), y
la pregunta obligada es \textbf{¿por qué eso da el óptimo global?}
—no siempre lo hace, así que hace falta una prueba. El Argumento de
Intercambio es la técnica estándar para esa prueba: se toma
\textbf{cualquier} solución óptima hipotética que \textbf{difiera} de
la que produce el greedy, se identifica el primer punto donde
difieren, y se muestra que \textbf{intercambiar} esa parte por la
decisión que habría tomado el greedy \textbf{no empeora} la solución
(a veces la mejora, a veces la deja igual, pero \textbf{nunca} peor).
Repitiendo este intercambio, cualquier óptimo se puede transformar
—sin perder optimalidad— en \textbf{exactamente} la solución greedy,
lo cual prueba que el greedy también es óptimo.

\textbf{¿Por qué usarlo?} Sin esta prueba (o una equivalente, como
inducción o el teorema de matroides), un greedy es solo una
\textbf{heurística que parece razonable} —y "parece razonable" no es
lo mismo que "es correcto": hay muchos problemas donde la decisión
localmente óptima lleva a un resultado global subóptimo (el ejemplo
clásico es Coin Change con denominaciones arbitrarias, ver Greedy vs.
Programación Dinámica). El Argumento de Intercambio es la forma
\textbf{constructiva} de verificar correctitud antes de confiar en un
greedy bajo presión de tiempo, sin tener que probarlo empíricamente
contra fuerza bruta caso por caso.

\textbf{Casos de uso:} el patrón de intercambio se repite,
adaptado, en la justificación de \textbf{todos} los greedys clásicos:

\begin{itemize}
	\item \textbf{Activity Selection}: si una solución óptima no elige
	la actividad que termina \textbf{primero}, se puede reemplazar la
	que sí eligió por esa (termina antes o igual, así que no bloquea
	ninguna actividad futura que la otra no bloqueara) sin perder
	actividades.
	\item \textbf{Kruskal}: si una solución óptima no usa la arista de
	\textbf{menor peso} que no forma ciclo, se puede intercambiar por
	ella —el corte que separa sus dos extremos debe tener alguna
	arista en el MST original, y esa arista pesa \textbf{al menos}
	tanto como la de Kruskal.
	\item \textbf{Huffman Coding}: si una solución óptima no fusiona
	primero los dos símbolos de \textbf{menor frecuencia}, se pueden
	intercambiar de posición en el árbol sin aumentar el costo total
	(ambos son hojas más profundas o iguales que cualquier otro par).
	\item \textbf{Job Sequencing con Deadlines}: si una solución
	óptima no ordena por deadline creciente cuando conviene, se puede
	intercambiar un par de trabajos adyacentes "fuera de orden" sin
	empeorar el resultado.
\end{itemize}

\textbf{Complejidad:} no aplica en el sentido usual —es una técnica de
\textbf{demostración}, no un algoritmo con costo propio. El costo real
lo tiene el greedy que se está justificando (normalmente dominado por
un \texttt{sort} inicial, $O(n \log n)$).

\begin{lstlisting}[style=cpp]
	// El Argumento de Intercambio no se "programa" como estructura -- se
	// usa para JUSTIFICAR el criterio de ordenamiento/seleccion de un
	// greedy antes de escribirlo. Aqui se ilustra el patron general
	// aplicado a Activity Selection, para dejar explicito COMO se
	// traduce el argumento en codigo una vez que ya se probo.
	
	struct ExchangeArgumentDemo {
		// Actividad: [inicio, fin). Se asume que la prueba de intercambio
		// (ver arriba) ya establecio que "ordenar por FIN creciente y
		// tomar greedy" es optimo -- este codigo es la CONSECUENCIA
		// practica de esa prueba, no la prueba en si.
		struct Actividad { int inicio, fin; };
		
		// READ: numero maximo de actividades no solapadas, aplicando el
		// criterio de seleccion validado por el argumento de intercambio.
		static int maximoActividades(vector<Actividad> actividades) {
			sort(actividades.begin(), actividades.end(),
			[](Actividad& a, Actividad& b) { return a.fin < b.fin; });
			// ^ el ORDEN elegido aqui (por fin, no por inicio ni duracion)
			//   es EXACTAMENTE lo que el argumento de intercambio prueba
			//   que es seguro -- cambiar este criterio sin una prueba
			//   analoga puede dar un greedy INCORRECTO que compila y
			//   corre, pero no es optimo
			
			int contador = 0;
			int finUltima = INT_MIN;
			
			for (auto& act : actividades) {
				if (act.inicio >= finUltima) {   // no se solapa con la
					// ultima elegida
					contador++;
					finUltima = act.fin;
				}
			}
			
			return contador;
		}                                          // O(n log n)
	};
\end{lstlisting}

\textbf{Ejemplo de funcionamiento (esbozo de la prueba, no código):}

\begin{lstlisting}[style=cpp]
	// Actividades: A=[1,4), B=[3,5), C=[0,6), D=[5,7), E=[3,9), F=[6,10)
	//
	// El greedy (ordenar por fin, tomar la primera que no solape) elige:
	// A [1,4) -> D [5,7) -> F [6,10)? NO, F solapa con D -> se descarta
	// resultado greedy: A, D  (2 actividades)
	//
	// ARGUMENTO DE INTERCAMBIO -- por que no se puede hacer mejor:
	// Supongamos una solucion optima que NO empieza con A (la de menor
	// fin). Sea X la primera actividad que SI usa esa solucion optima.
	// Como A termina antes o al mismo tiempo que X (A tiene el menor fin
	// de TODAS), reemplazar X por A en la solucion optima:
	//   - no puede solapar con lo que viene despues (A termina <= X)
	//   - mantiene el mismo NUMERO de actividades
	// Repitiendo este intercambio para cada diferencia, cualquier optimo
	// se transforma en la solucion greedy sin perder actividades ->
	// el greedy tambien es optimo.
	
	vector<ExchangeArgumentDemo::Actividad> acts = {
		{1,4}, {3,5}, {0,6}, {5,7}, {3,9}, {6,10}
	};
	
	cout << ExchangeArgumentDemo::maximoActividades(acts);
	// 2
\end{lstlisting}

\textbf{Nota:} el Argumento de Intercambio prueba que existe
\textbf{al menos una} solución óptima que coincide con el greedy —no
que el greedy sea la \textbf{única} óptima (puede haber empates). En
Activity Selection, por ejemplo, A-D no es la única selección de
tamaño 2 (B-D o A-F distinto también podrían funcionar según los
datos), pero el argumento garantiza que \textbf{ninguna} solución le
gana en tamaño.

\textbf{Regla práctica:} antes de confiar en un greedy en competencia,
intenta esbozar mentalmente el intercambio en 2-3 líneas: "si el
óptimo no hace lo que hace mi greedy en el primer punto de diferencia,
¿puedo cambiarlo sin empeorar?" —si la respuesta es claramente sí (y
se puede justificar con una desigualdad simple), el greedy es seguro;
si la respuesta requiere "depende de qué pase después" o "a veces sí,
a veces no", es una señal fuerte de que el problema en realidad
necesita Programación Dinámica (ver siguiente subsección), no greedy.
\subsection{Greedy vs. Programación Dinámica}

\textbf{Idea:} tanto Greedy como Programación Dinámica resuelven
problemas de \textbf{optimización} construidos sobre subproblemas más
chicos, y ambos asumen \textbf{subestructura óptima} (la solución
óptima se construye a partir de soluciones óptimas de subproblemas).
La diferencia está en \textbf{cuántas} opciones se consideran en cada
paso: Greedy toma \textbf{una sola} decisión localmente óptima y
\textbf{nunca la reconsidera} (no hay marcha atrás, no hay
comparación contra alternativas); DP considera \textbf{todas} las
decisiones posibles en cada paso, resolviendo (y memoizando) el
subproblema resultante de \textbf{cada} una, y al final elige la que
dio el mejor resultado global. Greedy funciona \textbf{solo} cuando el
problema tiene la \textbf{propiedad de elección greedy} (la decisión
localmente óptima siempre forma parte de \textbf{alguna} solución
globalmente óptima, demostrable con Exchange Argument); sin esa
propiedad, ninguna cantidad de "parece razonable" sustituye a
considerar todas las opciones.

\textbf{¿Por qué usarlo?} Reconocer \textbf{cuál de los dos aplica}
es la habilidad central de esta comparación —usar Greedy donde hace
falta DP da una respuesta \textbf{incorrecta} (no solo lenta); usar DP
donde bastaba Greedy da una respuesta \textbf{correcta pero
	innecesariamente costosa} en tiempo/memoria. El ejemplo canónico para
entender la diferencia es \textbf{Coin Change}: con denominaciones
\textbf{canónicas} (como 1, 5, 10, 25 — el sistema de monedas de
EE.UU.), tomar siempre la moneda más grande posible ES óptimo (greedy
funciona); con denominaciones \textbf{arbitrarias} (ej. 1, 3, 4), ese
mismo greedy puede fallar —\texttt{6} con greedy da \texttt{4+1+1}
(3 monedas), pero el óptimo real es \texttt{3+3} (2 monedas). El
greedy \textbf{compila, corre, y da una respuesta} —solo que
incorrecta, lo cual es más peligroso que un error de compilación
porque no se detecta sin comparar contra el resultado real.

\textbf{Casos de uso (tabla de decisión):}

\begin{itemize}
	\item \textbf{Fractional Knapsack} (se puede tomar una
	\textbf{fracción} de cada ítem): Greedy funciona —ordenar por
	razón valor/peso y tomar avariciosamente siempre es óptimo.
	\item \textbf{0/1 Knapsack} (cada ítem se toma \textbf{completo} o
	no se toma): Greedy \textbf{falla} —la fracción no está permitida,
	así que "el mejor por unidad de peso" puede no caber, y no hay
	forma de "reconsiderar" sin DP.
	\item \textbf{Activity Selection} (maximizar \textbf{cantidad} de
	actividades no solapadas, sin pesos): Greedy funciona —ordenar por
	fin creciente es óptimo (ver Exchange Argument).
	\item \textbf{Activity Selection con Pesos} (maximizar la
	\textbf{suma de valores} de actividades elegidas, no solo la
	cantidad): Greedy \textbf{falla} —la actividad que termina primero
	puede tener un valor bajo, y conviene DP (Weighted Interval
	Scheduling) para considerar todas las combinaciones.
	\item \textbf{Coin Change con denominaciones canónicas}: Greedy
	funciona.
	\item \textbf{Coin Change con denominaciones arbitrarias}: Greedy
	falla, hace falta DP (número mínimo de monedas).
\end{itemize}

\textbf{Complejidad:} no aplica un único número —el punto de esta
subsección es que la elección entre Greedy ($O(n\log n)$ típico, un
\texttt{sort} más una pasada) y DP ($O(n \cdot W)$ típico en problemas
tipo knapsack, exponencialmente más caro que greedy cuando aplica)
depende \textbf{enteramente} de si el problema tiene la propiedad de
elección greedy, no de una fórmula general.

\begin{lstlisting}[style=cpp]
	// Contraste directo: el MISMO problema superficial (Coin Change),
	// resuelto con Greedy (rapido pero solo correcto con denominaciones
	// canonicas) y con DP (siempre correcto, mas costoso).
	
	struct CoinChangeComparacion {
		// GREEDY: toma siempre la moneda mas grande posible. CORRECTO
		// solo si las denominaciones son "canonicas" (como 1,5,10,25) --
		// FALLA silenciosamente con denominaciones arbitrarias.
		static int monedasGreedy(vector<int> denominaciones, int objetivo) {
			sort(denominaciones.rbegin(), denominaciones.rend());   // mayor
			// a menor
			int monedas = 0;
			
			for (int d : denominaciones) {
				monedas += objetivo / d;   // toma tantas como quepan de
				// esta denominacion, SIN
				// reconsiderar despues
				objetivo %= d;
			}
			
			// si objetivo != 0 aqui, ni siquiera hay solucion greedy
			// valida -- otra senal de que este enfoque no es confiable
			return (objetivo == 0) ? monedas : -1;
		}                                          // O(n log n)
		
		// DP: considera TODAS las monedas posibles en cada paso,
		// memoizando el minimo real. SIEMPRE correcto, sin importar las
		// denominaciones.
		static int monedasDP(vector<int>& denominaciones, int objetivo) {
			vector<int> dp(objetivo + 1, INT_MAX);
			dp[0] = 0;
			
			for (int monto = 1; monto <= objetivo; monto++) {
				for (int d : denominaciones) {
					if (d > monto || dp[monto - d] == INT_MAX) continue;
					dp[monto] = min(dp[monto], dp[monto - d] + 1);
					// ^ a diferencia de greedy, aqui SI se compara contra
					//   todas las alternativas antes de decidir
				}
			}
			
			return (dp[objetivo] == INT_MAX) ? -1 : dp[objetivo];
		}                                          // O(objetivo * n)
	};
\end{lstlisting}

\textbf{Ejemplo de funcionamiento (donde Greedy falla y DP no):}

\begin{lstlisting}[style=cpp]
	vector<int> denominacionesCanonicas = {1, 5, 10, 25};
	vector<int> denominacionesArbitrarias = {1, 3, 4};
	
	cout << CoinChangeComparacion::monedasGreedy(denominacionesCanonicas, 30);
	// 2   (25 + 5 -- correcto, y coincide con el optimo real)
	
	cout << CoinChangeComparacion::monedasGreedy(denominacionesArbitrarias, 6);
	// 3   (4 + 1 + 1 -- INCORRECTO, el greedy se equivoca aqui)
	
	cout << CoinChangeComparacion::monedasDP(denominacionesArbitrarias, 6);
	// 2   (3 + 3 -- CORRECTO, DP si encuentra el optimo real)
\end{lstlisting}

\textbf{Nota:} el hecho de que el greedy \textbf{parezca} funcionar en
ejemplos pequeños o "de prueba" no es evidencia de que sea correcto en
general —con \texttt{\{1,3,4\}} y objetivo \texttt{6} ya falla, un
caso perfectamente razonable de aparecer en un juez. La única forma
confiable de saber si un greedy es seguro es demostrarlo (Exchange
Argument) o reconocer el problema como una instancia de un patrón ya
probado (Activity Selection, Fractional Knapsack, Huffman, etc.), no
"probarlo con ejemplos y ver si da bien".

\textbf{Regla práctica:} si al intentar el Argumento de Intercambio
para tu greedy te encuentras diciendo "pero esto depende de qué pase
\textbf{después}" o "la mejor opción \textbf{ahora} podría no ser la
mejor a largo plazo", esa es la señal directa de que el problema
necesita DP —esa dependencia del futuro es exactamente lo que Greedy
no puede capturar (decide y nunca reconsidera), mientras que DP la
maneja de forma natural explorando (y memoizando) todas las
alternativas.

\section{Intervalos}
\subsection{Activity Selection (Selección de Actividades)}

\textbf{Idea:} dado un conjunto de actividades, cada una con un
intervalo $[\texttt{inicio}, \texttt{fin})$, se busca el
\textbf{máximo número} de actividades que se pueden realizar sin que
ninguna se solape con otra (una sola persona/recurso, no puede hacer
dos a la vez). El greedy correcto —ya justificado con Exchange
Argument en Fundamentos de Greedy— es: \textbf{ordenar por fin
	creciente}, y recorrer tomando cada actividad cuyo inicio sea
$\ge$ el fin de la última tomada. La intuición: terminar lo antes
posible \textbf{siempre} deja más espacio libre para actividades
futuras que cualquier otra elección —nunca hay razón para preferir
una actividad que termina más tarde si una que termina antes también
estaba disponible.

\textbf{¿Por qué usarlo?} Es el ejemplo introductorio estándar de
greedy correcto porque el criterio "ordenar por fin" es
\textbf{contraintuitivo} a primera vista —muchos intentan ordenar por
\textbf{duración} (la más corta primero) o por \textbf{inicio} más
temprano, y ambos criterios \textbf{fallan} en casos concretos (una
actividad corta puede empezar tarde y bloquear más actividades después
que una larga que termina pronto). Interiorizar por qué "fin
creciente" es el único criterio que funciona aquí previene ese error
en variantes futuras.

\textbf{Casos de uso:}

\begin{itemize}
	\item Maximizar el número de reuniones/actividades que se pueden
	agendar en una sola sala/recurso sin solapamiento.
	\item Como subrutina de Minimum Number of Platforms/Rooms
	(siguiente bloque de subsecciones): resolver "cuántos recursos
	hacen falta" es una generalización de "¿alcanza con uno?".
	\item Punto de comparación directo con Interval Point Cover
	(siguiente subsección): problema \textbf{relacionado} pero
	\textbf{distinto} —aquí se eligen actividades sin solapar; ahí se
	eligen puntos que \textbf{toquen} todos los intervalos.
\end{itemize}

\textbf{Complejidad:} $O(n \log n)$, dominado por el \texttt{sort}
inicial por tiempo de fin; el barrido posterior es $O(n)$.

\begin{lstlisting}[style=cpp]
	struct ActivitySelection {
		struct Actividad { int inicio, fin; };
		
		// READ: maximo numero de actividades no solapadas. Ordena
		// internamente por FIN creciente (criterio validado por Exchange
		// Argument, ver Fundamentos de Greedy).
		static int maximoActividades(vector<Actividad> actividades) {
			sort(actividades.begin(), actividades.end(),
			[](Actividad& a, Actividad& b) { return a.fin < b.fin; });
			
			int contador = 0;
			int finUltima = INT_MIN;
			
			for (auto& act : actividades) {
				if (act.inicio >= finUltima) {
					contador++;
					finUltima = act.fin;
				}
			}
			
			return contador;
		}                                          // O(n log n)
		
		// READ: devuelve las actividades ELEGIDAS explicitamente, no solo
		// el conteo (mismo criterio, guardando cada seleccion).
		static vector<Actividad> seleccionar(vector<Actividad> actividades) {
			sort(actividades.begin(), actividades.end(),
			[](Actividad& a, Actividad& b) { return a.fin < b.fin; });
			
			vector<Actividad> elegidas;
			int finUltima = INT_MIN;
			
			for (auto& act : actividades) {
				if (act.inicio >= finUltima) {
					elegidas.push_back(act);
					finUltima = act.fin;
				}
			}
			
			return elegidas;
		}                                          // O(n log n)
	};
\end{lstlisting}

\textbf{Ejemplo de funcionamiento:}

\begin{lstlisting}[style=cpp]
	vector<ActivitySelection::Actividad> acts = {
		{1,4}, {3,5}, {0,6}, {5,7}, {3,9}, {5,9}, {6,10}, {8,11}, {8,12}, {2,14}, {12,16}
	};
	
	cout << ActivitySelection::maximoActividades(acts);
	// 4
	
	auto elegidas = ActivitySelection::seleccionar(acts);
	for (auto& a : elegidas) cout << "[" << a.inicio << "," << a.fin << ") ";
	// [1,4) [5,7) [8,11) [12,16)
\end{lstlisting}

\textbf{Regla práctica:} \textbf{siempre} ordena por \textbf{fin
	creciente} en este problema —ordenar por inicio o por duración son los
dos errores más comunes y ambos producen contraejemplos reales (no
solo teóricos). Si el problema agrega \textbf{pesos/valores} a cada
actividad (maximizar la suma de valores, no la cantidad), este greedy
deja de aplicar directamente —ver Greedy vs. Programación Dinámica.


\subsection{Interval Point Cover}

\textbf{Idea:} dado un conjunto de intervalos $[\texttt{inicio},
\texttt{fin}]$, se busca el \textbf{mínimo número de puntos} tales
que \textbf{cada} intervalo contenga \textbf{al menos uno} de esos
puntos (el problema inverso, en cierto sentido, a Activity Selection:
ahí se elegían intervalos sin solapar; aquí se eligen puntos que
\textbf{toquen} todos los intervalos). El greedy: \textbf{ordenar por
	fin creciente}, y cada vez que un intervalo \textbf{no} está ya cubierto
por el último punto colocado, colocar un \textbf{nuevo} punto
exactamente en el \textbf{fin} de ese intervalo —colocarlo ahí (lo más
tarde posible dentro de ese intervalo) maximiza la chance de que
también cubra intervalos futuros que empiecen antes de ese fin.

\textbf{¿Por qué usarlo?} Comparte el mismo ordenamiento por fin
creciente que Activity Selection, pero la \textbf{razón} de por qué
funciona es distinta: aquí no se trata de "dejar espacio libre" sino
de "cubrir la mayor cantidad de intervalos posible con cada punto
colocado" —poner el punto en el fin del intervalo más urgente (el que
termina primero) es la elección que \textbf{no puede fallar}: ese
intervalo debe recibir un punto en algún momento, y ponerlo en su fin
es la posición que abarca el rango más amplio de intervalos que
todavía se solapan con él.

\textbf{Casos de uso:}

\begin{itemize}
	\item Mínimo número de flechas para reventar todos los globos
	(\textit{Minimum Number of Arrows to Burst Balloons} —variante
	directa con otro nombre, cada flecha es un "punto").
	\item Mínimo número de estaciones de servicio/vigilancia tales que
	cada segmento de una carretera esté cubierto por al menos una.
	\item Mínimo número de "checkpoints" temporales tales que cada
	tarea (con su ventana de tiempo activa) sea verificada al menos
	una vez.
	\item Como subrutina de problemas de cobertura más generales donde
	la restricción se puede modelar como intervalos y puntos.
\end{itemize}

\textbf{Complejidad:} $O(n \log n)$, dominado por el \texttt{sort}
inicial por fin creciente; el barrido posterior es $O(n)$.

\begin{lstlisting}[style=cpp]
	struct IntervalPointCover {
		struct Intervalo { long long inicio, fin; };
		
		// READ: minimo numero de puntos que tocan TODOS los intervalos.
		static int minimoPuntos(vector<Intervalo> intervalos) {
			sort(intervalos.begin(), intervalos.end(),
			[](Intervalo& a, Intervalo& b) { return a.fin < b.fin; });
			
			int puntos = 0;
			long long ultimoPunto = LLONG_MIN;
			
			for (auto& iv : intervalos) {
				if (iv.inicio > ultimoPunto) {
					// el ultimo punto colocado NO cubre este intervalo
					// -> hace falta uno nuevo, colocado en el FIN de
					// este intervalo (maximiza cobertura futura)
					puntos++;
					ultimoPunto = iv.fin;
				}
				// si iv.inicio <= ultimoPunto, este intervalo YA esta
				// cubierto por el punto anterior -- no hace falta nada
			}
			
			return puntos;
		}                                          // O(n log n)
		
		// READ: devuelve los puntos elegidos explicitamente.
		static vector<long long> obtenerPuntos(vector<Intervalo> intervalos) {
			sort(intervalos.begin(), intervalos.end(),
			[](Intervalo& a, Intervalo& b) { return a.fin < b.fin; });
			
			vector<long long> puntos;
			long long ultimoPunto = LLONG_MIN;
			
			for (auto& iv : intervalos) {
				if (iv.inicio > ultimoPunto) {
					ultimoPunto = iv.fin;
					puntos.push_back(ultimoPunto);
				}
			}
			
			return puntos;
		}                                          // O(n log n)
	};
\end{lstlisting}

\textbf{Ejemplo de funcionamiento:}

\begin{lstlisting}[style=cpp]
	vector<IntervalPointCover::Intervalo> intervalos = {
		{1,4}, {2,5}, {6,8}, {3,7}, {9,10}
	};
	
	cout << IntervalPointCover::minimoPuntos(intervalos);
	// 2
	
	auto puntos = IntervalPointCover::obtenerPuntos(intervalos);
	for (long long p : puntos) cout << p << " ";
	// 4 10
	// punto en 4 cubre [1,4],[2,5],[3,7]; punto en 10 cubre [6,8],[9,10]
\end{lstlisting}

\textbf{Regla práctica:} tanto Activity Selection como Interval Point
Cover ordenan por \textbf{fin creciente}, pero la decisión que toman
es opuesta en naturaleza: Activity Selection \textbf{descarta}
intervalos que se solapan con lo ya elegido; Interval Point Cover
\textbf{agrega} un punto solo cuando el intervalo actual \textbf{no}
está cubierto todavía. Si el problema pide "elige el máximo de cosas
sin solapar", es Activity Selection; si pide "cubre todo con el mínimo
de recursos puntuales", es Interval Point Cover.
\subsection{Merge Intervals}

\textbf{Idea:} dado un conjunto de intervalos $[\texttt{inicio},
\texttt{fin}]$ que pueden \textbf{solaparse} entre sí, se busca
\textbf{fusionarlos} en el menor número de intervalos
\textbf{disjuntos} que cubran exactamente la misma unión de puntos. El
greedy: \textbf{ordenar por inicio creciente} (no por fin, a
diferencia de Activity Selection e Interval Point Cover — aquí importa
el orden en que los intervalos "aparecen" para poder ir
extendiéndolos), y recorrer manteniendo un intervalo "en construcción"
—si el siguiente intervalo \textbf{se solapa} con el actual
(\texttt{inicio} $\le$ \texttt{fin} actual), se \textbf{extiende} el
actual tomando el máximo de los dos fines; si no se solapa, el
intervalo actual queda \textbf{cerrado} y se emite, empezando uno
nuevo.

\textbf{¿Por qué usarlo?} Es el primer greedy de esta sección que
ordena por \textbf{inicio} en vez de por fin —vale la pena notar el
contraste explícitamente: aquí no se está \textbf{eligiendo}
intervalos ni \textbf{cubriendo} con puntos, se está
\textbf{reconstruyendo} la unión total de todos ellos, y para eso hace
falta procesarlos en el orden en que empiezan, extendiendo
progresivamente lo que ya se lleva fusionado.

\textbf{Casos de uso:}

\begin{itemize}
	\item Simplificar una lista de rangos horarios ocupados en una
	sola lista de bloques \textbf{sin} solapamientos, antes de aplicar
	otro algoritmo sobre ella (ej. antes de Minimum Number of
	Platforms/Rooms, si primero hace falta limpiar datos redundantes).
	\item Insertar un nuevo intervalo en una lista ya fusionada,
	re-fusionando solo la parte afectada.
	\item Calcular la \textbf{longitud total cubierta} por una unión
	de intervalos con solapamientos (sumar la longitud de cada
	intervalo fusionado, en vez de sumar cada intervalo original y
	arriesgarse a contar dos veces las partes solapadas).
	\item Como paso de limpieza previo en problemas de Line Sweep
	donde los eventos vienen como intervalos crudos, potencialmente
	redundantes.
\end{itemize}

\textbf{Complejidad:} $O(n \log n)$, dominado por el \texttt{sort}
inicial por inicio creciente; el barrido de fusión es $O(n)$.

\begin{lstlisting}[style=cpp]
	struct MergeIntervals {
		struct Intervalo { long long inicio, fin; };
		
		// READ: fusiona todos los intervalos solapados en la lista MINIMA
		// de intervalos disjuntos que cubre la misma union de puntos.
		static vector<Intervalo> fusionar(vector<Intervalo> intervalos) {
			if (intervalos.empty()) return {};
			
			sort(intervalos.begin(), intervalos.end(),
			[](Intervalo& a, Intervalo& b) { return a.inicio < b.inicio; });
			
			vector<Intervalo> resultado;
			Intervalo actual = intervalos[0];
			
			for (int i = 1; i < (int)intervalos.size(); i++) {
				auto& siguiente = intervalos[i];
				
				if (siguiente.inicio <= actual.fin) {
					// se solapa (o es adyacente) -> extiende el actual
					actual.fin = max(actual.fin, siguiente.fin);
				} else {
					// no se solapa -> cierra el actual y empieza uno nuevo
					resultado.push_back(actual);
					actual = siguiente;
				}
			}
			
			resultado.push_back(actual);   // el ultimo intervalo en
			// construccion siempre
			// queda pendiente de emitir
			return resultado;
		}                                          // O(n log n)
		
		// READ: longitud total cubierta por la union de todos los
		// intervalos (sin contar dos veces las partes solapadas) --
		// reutiliza fusionar() y suma las longitudes ya disjuntas.
		static long long longitudTotalCubierta(vector<Intervalo> intervalos) {
			auto fusionados = fusionar(intervalos);
			long long total = 0;
			
			for (auto& iv : fusionados)
			total += iv.fin - iv.inicio;
			
			return total;
		}                                          // O(n log n)
	};
\end{lstlisting}

\textbf{Ejemplo de funcionamiento:}

\begin{lstlisting}[style=cpp]
	vector<MergeIntervals::Intervalo> intervalos = {
		{1,3}, {2,6}, {8,10}, {15,18}, {9,12}
	};
	
	auto fusionados = MergeIntervals::fusionar(intervalos);
	for (auto& iv : fusionados) cout << "[" << iv.inicio << "," << iv.fin << "] ";
	// [1,6] [8,12] [15,18]
	
	cout << MergeIntervals::longitudTotalCubierta(intervalos);
	// 5 + 4 + 3 = 12
\end{lstlisting}

\textbf{Nota:} la condición \texttt{siguiente.inicio <= actual.fin}
(con \texttt{<=}, no \texttt{<}) decide si intervalos \textbf{apenas
	adyacentes} (ej. $[1,3]$ y $[3,5]$) se fusionan o no —con \texttt{<=}
se fusionan en $[1,5]$ (tratando los extremos como \textbf{cerrados},
el punto 3 pertenece a ambos); si el problema usa intervalos
\textbf{semiabiertos} $[\texttt{inicio}, \texttt{fin})$, la condición
correcta para fusionar solo solapamientos reales es \texttt{<} estricto
—revisar siempre qué convención de extremos usa el enunciado antes de
fijar este operador.

\textbf{Regla práctica:} Merge Intervals ordena por \textbf{inicio}
porque el objetivo es reconstruir la unión completa en el orden en que
los intervalos "aparecen"; Activity Selection e Interval Point Cover
ordenan por \textbf{fin} porque el objetivo es decidir, cuanto antes,
qué se puede "liberar" o "cubrir". Si el problema pide "combina/limpia
los intervalos" (sin elegir ni cubrir, solo reconstruir), es Merge
Intervals.


\subsection{Minimum Number of Platforms/Rooms}

\textbf{Idea:} dado un conjunto de intervalos (llegadas y salidas de
trenes, o reuniones con horario), se busca el \textbf{mínimo número de
	recursos simultáneos} (andenes, salas) necesarios para que
\textbf{nunca} haya dos intervalos activos al mismo tiempo compitiendo
por el mismo recurso —a diferencia de Activity Selection (que usa
\textbf{un solo} recurso y descarta lo que no cabe), aquí se permite
usar \textbf{tantos recursos como haga falta}, y la pregunta es
\textbf{cuántos}. Se resuelve con \textbf{Line Sweep}: se separan
todos los \texttt{inicio} y todos los \texttt{fin} en dos listas,
ambas ordenadas; se recorren \textbf{simultáneamente} con Two
Pointers, y cada vez que un \texttt{inicio} ocurre \textbf{antes} que
el siguiente \texttt{fin} pendiente, se necesita \textbf{un recurso
	más} (contador sube); cada vez que un \texttt{fin} ocurre primero, un
recurso se \textbf{libera} (contador baja) — el \textbf{máximo} valor
que alcanza el contador durante el barrido es la respuesta.

\textbf{¿Por qué usarlo?} Es la generalización natural de "¿alcanza
con un solo recurso?" (Activity Selection) a "¿cuántos recursos hacen
falta como máximo?" —el mismo espíritu de barrer por tiempo, pero en
vez de descartar actividades que no caben, se cuenta cuántas están
\textbf{simultáneamente activas} en el peor momento. Es también el
ejemplo introductorio de \textbf{Line Sweep con eventos de tipo
	"+1"/"-1"}, un patrón que reaparece en problemas mucho más generales
de conteo de solapamientos.

\textbf{Casos de uso:}

\begin{itemize}
	\item Mínimo número de andenes de tren, salas de reunión, o
	cajeros que hacen falta para atender todas las llegadas/salidas
	sin conflicto.
	\item Máximo número de eventos/personas activas \textbf{al mismo
		tiempo} en cualquier instante (el valor máximo del contador \emph{es}
	esa respuesta, no solo cuántos recursos hacen falta).
	\item Como subrutina de problemas de \textbf{asignación} de
	recursos: una vez que se sabe cuántos recursos hacen falta como
	máximo, se puede simular la asignación explícita con una
	\texttt{priority\_queue} de "recurso libre más próximo".
	\item Detectar si existe algún momento en que \textbf{más de $k$}
	intervalos están activos a la vez (comparando el máximo del
	contador contra $k$).
\end{itemize}

\textbf{Complejidad:} $O(n \log n)$, dominado por ordenar las dos
listas de inicios y fines por separado; el barrido con Two Pointers es
$O(n)$.

\begin{lstlisting}[style=cpp]
	struct MinimoRecursos {
		struct Intervalo { int inicio, fin; };
		
		// READ: minimo numero de recursos simultaneos necesarios para
		// que ningun par de intervalos activos compita por el mismo
		// recurso. Asume intervalos semiabiertos [inicio, fin) -- un
		// fin exactamente igual a un inicio NO cuenta como solape (el
		// recurso ya se libero justo quando el otro arranca).
		static int minimoRecursos(vector<Intervalo>& intervalos) {
			int n = intervalos.size();
			vector<int> inicios(n), fines(n);
			
			for (int i = 0; i < n; i++) {
				inicios[i] = intervalos[i].inicio;
				fines[i] = intervalos[i].fin;
			}
			
			sort(inicios.begin(), inicios.end());
			sort(fines.begin(), fines.end());
			
			int i = 0, j = 0;
			int activos = 0, maximoActivos = 0;
			
			while (i < n) {
				if (inicios[i] < fines[j]) {
					// un nuevo intervalo empieza ANTES de que el
					// siguiente pendiente termine -> hace falta un
					// recurso mas
					activos++;
					maximoActivos = max(maximoActivos, activos);
					i++;
				} else {
					// un intervalo termina (o empata, tratado como
					// liberacion primero) -> se libera un recurso
					activos--;
					j++;
				}
			}
			
			return maximoActivos;
		}                                          // O(n log n)
	};
\end{lstlisting}

\textbf{Ejemplo de funcionamiento:}

\begin{lstlisting}[style=cpp]
	// llegadas y salidas de trenes (mismo indice = mismo tren)
	vector<MinimoRecursos::Intervalo> trenes = {
		{900, 910}, {940, 1200}, {950, 1120},
		{1100, 1130}, {1500, 1900}, {1800, 2000}
	};
	
	cout << MinimoRecursos::minimoRecursos(trenes);
	// 3
	// en el momento 950-1100 hay 3 trenes activos a la vez:
	// [940,1200), [950,1120), y arranca [1100,1130) justo antes de que
	// termine ninguno de los otros dos
\end{lstlisting}

\textbf{Nota:} la condición \texttt{inicios[i] < fines[j]} (estricta)
asume que un fin exactamente igual a un inicio \textbf{no} cuenta como
conflicto (el recurso se libera justo a tiempo). Si el problema
considera que un intervalo que \textbf{termina} en el momento $t$
\textbf{sí} sigue ocupando el recurso hasta después de $t$ (extremos
cerrados en ambos lados), la condición correcta es
\texttt{inicios[i] <= fines[j]} — revisar siempre esta convención
contra el enunciado, es la fuente de error de borde más común en este
problema.

\textbf{Regla práctica:} si el problema permite \textbf{un solo}
recurso y pide maximizar cuántas actividades caben, es Activity
Selection; si permite \textbf{recursos ilimitados} y pide el
\textbf{mínimo} necesario para no tener conflictos, es Minimum Number
of Platforms/Rooms. Ambos comparten la idea de ordenar por tiempo y
barrer, pero uno \textbf{descarta} lo que no cabe y el otro
\textbf{cuenta} el pico de concurrencia.
\subsection{Non-overlapping Intervals}

\textbf{Idea:} dado un conjunto de intervalos que pueden solaparse, se
busca el \textbf{mínimo número de intervalos a eliminar} para que los
que \textbf{queden} no se solapen entre sí. Es el problema
\textbf{complementario} de Activity Selection: si Activity Selection
encuentra el \textbf{máximo} número de actividades que \textbf{sí} se
pueden mantener sin solapar, entonces el mínimo a eliminar es
simplemente $n - \texttt{maximoActividades}$ —no hace falta ningún
algoritmo nuevo, es \textbf{el mismo greedy} (ordenar por fin
creciente, tomar cada intervalo que no solape con el último tomado)
visto desde el ángulo opuesto.

\textbf{¿Por qué usarlo?} Es la trampa más común de esta sección:
muchos intentan resolverlo como un problema \textbf{distinto} (a veces
con DP o con otro criterio de ordenamiento), sin notar que "máximo que
se puede \textbf{mantener}" y "mínimo que hay que \textbf{eliminar}"
son la \textbf{misma} cantidad, solo restada de $n$. Reconocer esta
equivalencia ahorra tener que re-derivar el Exchange Argument desde
cero —ya está probado en Activity Selection.

\textbf{Casos de uso:}

\begin{itemize}
	\item Mínimo número de reuniones/eventos a cancelar para que el
	resto quepa sin conflictos en un solo recurso.
	\item Como verificación rápida: si un juez da directamente "número
	de actividades a remover" en vez de "máximo a mantener", no hace
	falta cambiar el algoritmo, solo restar de $n$ al final.
	\item Variante para practicar reconocer \textbf{transformaciones
		triviales} entre enunciados que suenan distintos pero son
	matemáticamente idénticos —una habilidad tan importante como
	conocer los algoritmos en sí.
\end{itemize}

\textbf{Complejidad:} $O(n \log n)$ — idéntica a Activity Selection,
ya que es literalmente el mismo algoritmo con una resta final.

\begin{lstlisting}[style=cpp]
	struct NonOverlappingIntervals {
		struct Intervalo { int inicio, fin; };
		
		// READ: minimo numero de intervalos a ELIMINAR para que el resto
		// no se solape. Reutiliza EXACTAMENTE el criterio de Activity
		// Selection (ordenar por fin, contar cuantos SI caben) y resta
		// de n -- no es un algoritmo nuevo, es el mismo visto al reves.
		static int minimoAEliminar(vector<Intervalo> intervalos) {
			int n = intervalos.size();
			if (n == 0) return 0;
			
			sort(intervalos.begin(), intervalos.end(),
			[](Intervalo& a, Intervalo& b) { return a.fin < b.fin; });
			
			int mantenidos = 0;
			int finUltimo = INT_MIN;
			
			for (auto& iv : intervalos) {
				if (iv.inicio >= finUltimo) {   // no se solapa con el
					// ultimo MANTENIDO
					mantenidos++;
					finUltimo = iv.fin;
				}
				// si se solapa, se "elimina" implicitamente -- no se
				// suma a mantenidos, no se actualiza finUltimo
			}
			
			return n - mantenidos;
		}                                          // O(n log n)
	};
\end{lstlisting}

\textbf{Ejemplo de funcionamiento:}

\begin{lstlisting}[style=cpp]
	vector<NonOverlappingIntervals::Intervalo> intervalos = {
		{1,2}, {2,3}, {3,4}, {1,3}
	};
	
	cout << NonOverlappingIntervals::minimoAEliminar(intervalos);
	// 1
	// se mantienen [1,2],[2,3],[3,4] (3 intervalos, no se solapan entre
	// si con extremos semiabiertos); se elimina [1,3], que solapaba con
	// [1,2] y con [2,3]
\end{lstlisting}

\textbf{Regla práctica:} cuando un enunciado pida "elimina el mínimo
para que no queden solapes", NO busques un algoritmo nuevo —resuelve
Activity Selection normal (máximo que se puede mantener) y resta de
$n$. Esta equivalencia vale la pena memorizarla explícitamente porque
el enunciado rara vez menciona la palabra "selección", lo cual puede
hacer que el problema parezca distinto a primera vista.


\subsection{Minimum Number of Arrows to Burst Balloons}

\textbf{Idea:} dado un conjunto de globos representados como
intervalos $[\texttt{inicio}, \texttt{fin}]$ en el eje horizontal, una
flecha disparada verticalmente en la posición $x$ revienta
\textbf{todos} los globos cuyo intervalo contiene a $x$. Se busca el
\textbf{mínimo número de flechas} para reventar todos los globos —esto
es \textbf{exactamente} Interval Point Cover con otro nombre: cada
"punto" (posición de flecha) debe tocar todos los intervalos
(globos), y el greedy es idéntico: ordenar por \texttt{fin} creciente,
disparar una flecha en el \texttt{fin} del primer globo no cubierto
todavía.

\textbf{¿Por qué usarlo?} Es el ejemplo más citado de Interval Point
Cover en jueces de práctica, y vale la pena tenerlo como
\textbf{subsección propia} precisamente para reforzar el
reconocimiento de patrones: el enunciado suena a geometría/física
("flechas", "globos"), pero la estructura matemática subyacente es
\textbf{idéntica} a la de "mínimo número de puntos que cubren
intervalos" — ningún concepto nuevo, solo una capa distinta de
narrativa sobre el mismo problema.

\textbf{Casos de uso:}

\begin{itemize}
	\item El problema de globos y flechas tal cual aparece en jueces
	de práctica (frecuentemente bajo ese nombre exacto).
	\item Cualquier variante narrativa equivalente: "mínimo número de
	inspecciones puntuales para cubrir todos los rangos de
	validez", "mínimo número de disparos de láser", etc. — mismo
	esqueleto, distinta historia.
	\item Como ejercicio de \textbf{traducción de enunciado a
		estructura conocida}: identificar que "un disparo revienta todos
	los globos que lo contienen" es literalmente "un punto cubre
	todos los intervalos que lo contienen" es el paso clave, no el
	algoritmo en sí (que ya está resuelto en Interval Point Cover).
\end{itemize}

\textbf{Complejidad:} $O(n \log n)$ — idéntica a Interval Point Cover,
ya que es el mismo algoritmo.

\begin{lstlisting}[style=cpp]
	struct ArrowsBurstBalloons {
		struct Globo { long long inicio, fin; };   // [inicio, fin] del
		// globo en el eje X
		
		// READ: minimo numero de flechas para reventar TODOS los globos.
		// Identico a IntervalPointCover::minimoPuntos -- se reescribe
		// aqui con los nombres del enunciado para dejar explicita la
		// correspondencia.
		static int minimoFlechas(vector<Globo> globos) {
			if (globos.empty()) return 0;
			
			sort(globos.begin(), globos.end(),
			[](Globo& a, Globo& b) { return a.fin < b.fin; });
			
			int flechas = 1;
			long long posicionFlecha = globos[0].fin;
			
			for (int i = 1; i < (int)globos.size(); i++) {
				if (globos[i].inicio > posicionFlecha) {
					// este globo ya NO es alcanzado por la ultima
					// flecha -> hace falta una nueva, en el FIN de este
					// globo (maximiza cuantos globos futuros revienta)
					flechas++;
					posicionFlecha = globos[i].fin;
				}
				// si globos[i].inicio <= posicionFlecha, este globo YA
				// esta reventado por la flecha actual -- no hace falta nada
			}
			
			return flechas;
		}                                          // O(n log n)
	};
\end{lstlisting}

\textbf{Ejemplo de funcionamiento:}

\begin{lstlisting}[style=cpp]
	vector<ArrowsBurstBalloons::Globo> globos = {
		{10,16}, {2,8}, {1,6}, {7,12}
	};
	
	cout << ArrowsBurstBalloons::minimoFlechas(globos);
	// 2
	// ordenados por fin: [1,6],[2,8],[7,12],[10,16]
	// flecha en 6 revienta [1,6] y [2,8] (2 <= 6)
	// [7,12] no cubierto (7 > 6) -> flecha en 12, revienta [7,12] y
	// [10,16] (10 <= 12)
\end{lstlisting}

\textbf{Regla práctica:} antes de programar cualquier "problema
nuevo" de esta sección, pregúntate primero si ya es \textbf{uno de los
	patrones anteriores} disfrazado con otro nombre —Non-overlapping
Intervals es Activity Selection invertido; Minimum Number of Arrows es
Interval Point Cover con otra narrativa. La mayoría de problemas
greedy "nuevos" en jueces de práctica son variaciones nominales de un
puñado de patrones estructurales ya vistos, no algoritmos genuinamente
distintos.

\section{Scheduling}
\subsection{Job Sequencing con Deadlines}

\textbf{Idea:} dado un conjunto de trabajos, cada uno con una
\textbf{ganancia} y un \textbf{deadline} (debe completarse antes o en
ese momento, y cada trabajo toma exactamente 1 unidad de tiempo en una
única máquina), se busca el subconjunto de trabajos que se pueden
programar \textbf{sin} pasarse de su deadline, \textbf{maximizando} la
ganancia total. El greedy: \textbf{ordenar por ganancia
	decreciente}, y para cada trabajo (en ese orden) intentar colocarlo en
el \textbf{slot de tiempo disponible más tardío} que sea
$\le$ su deadline —si existe un slot libre, se ocupa; si no, el
trabajo se descarta. Colocarlo lo \textbf{más tarde posible} (en vez
de lo más temprano) es la clave: deja los slots \textbf{tempranos}
libres para trabajos futuros (de menor ganancia, ya que se procesan en
orden decreciente) que tengan deadlines \textbf{más ajustados}.

\textbf{¿Por qué usarlo?} Es el greedy clásico de "ordena por el
criterio que \textbf{no se puede recuperar después}" (la ganancia,
ya que una vez descartado un trabajo de alta ganancia no hay forma de
reconsiderarlo) combinado con una asignación que preserva
\textbf{flexibilidad futura} (slot más tardío posible). La
implementación ingenua de "buscar el slot libre más tardío" es
$O(n \cdot \text{maxDeadline})$ por escaneo lineal; con \textbf{DSU}
(Union-Find) se acelera a casi $O(n \log n)$, aprovechando que
"encontrar el slot libre más tardío $\le d$" es exactamente lo que
resuelve un DSU donde cada slot ocupado "apunta" al slot libre
anterior más cercano — el mismo patrón de "encontrar el siguiente
disponible" que aparece en problemas de asignación con Union-Find.

\textbf{Casos de uso:}

\begin{itemize}
	\item Programar trabajos con deadline y ganancia en una única
	máquina, maximizando ganancia total (el problema base).
	\item Cualquier problema de "elige un subconjunto sujeto a
	restricciones de tiempo límite individuales, maximizando una
	métrica agregada" — el criterio de "ordenar por lo irreversible,
	asignar preservando flexibilidad" se generaliza a variantes con
	duración de trabajo distinta de 1 (aunque ahí ya no es tan directo
	sin modificar la estructura de asignación).
	\item Como aplicación práctica de DSU fuera del contexto habitual
	de conectividad — usar la estructura para "siguiente slot libre"
	en vez de "misma componente".
\end{itemize}

\textbf{Complejidad:} $O(n \log n)$ para ordenar por ganancia, más
$O(n \cdot \alpha(n))$ para las asignaciones con DSU (prácticamente
lineal) — la versión sin DSU (escaneo lineal por trabajo) es
$O(n \cdot \text{maxDeadline})$, aceptable si los deadlines son
pequeños.

\begin{lstlisting}[style=cpp]
	struct JobSequencing {
		struct Trabajo { int deadline, ganancia; };
		
		// READ: maxima ganancia total programando trabajos sin pasarse
		// de su deadline, un trabajo por unidad de tiempo. Version con
		// ESCANEO LINEAL de slots (simple, O(n * maxDeadline)).
		static long long maximaGananciaSimple(vector<Trabajo> trabajos) {
			sort(trabajos.begin(), trabajos.end(),
			[](Trabajo& a, Trabajo& b) { return a.ganancia > b.ganancia; });
			
			int maxDeadline = 0;
			for (auto& t : trabajos)
			maxDeadline = max(maxDeadline, t.deadline);
			
			vector<bool> ocupado(maxDeadline + 1, false);
			long long gananciaTotal = 0;
			
			for (auto& t : trabajos) {
				// busca el slot libre MAS TARDIO <= deadline de este trabajo
				for (int slot = min(t.deadline, maxDeadline); slot >= 1; slot--) {
					if (!ocupado[slot]) {
						ocupado[slot] = true;
						gananciaTotal += t.ganancia;
						break;
					}
				}
			}
			
			return gananciaTotal;
		}                                          // O(n * maxDeadline)
		
		// READ: misma logica, pero usando DSU para encontrar el slot
		// libre mas tardio en O(alpha(n)) amortizado en vez de escanear
		// linealmente cada vez.
		static long long maximaGananciaDSU(vector<Trabajo> trabajos) {
			sort(trabajos.begin(), trabajos.end(),
			[](Trabajo& a, Trabajo& b) { return a.ganancia > b.ganancia; });
			
			int maxDeadline = 0;
			for (auto& t : trabajos)
			maxDeadline = max(maxDeadline, t.deadline);
			
			// padre[slot] = el slot LIBRE mas tardio <= slot (o 0 si
			// ninguno esta libre); find(x) sigue la cadena de slots
			// ocupados hasta el proximo libre
			vector<int> padre(maxDeadline + 1);
			iota(padre.begin(), padre.end(), 0);
			
			function<int(int)> find = [&](int x) {
				if (padre[x] != x) padre[x] = find(padre[x]);
				return padre[x];
			};
			
			long long gananciaTotal = 0;
			
			for (auto& t : trabajos) {
				int slotLibre = find(min(t.deadline, maxDeadline));
				
				if (slotLibre > 0) {
					gananciaTotal += t.ganancia;
					padre[slotLibre] = slotLibre - 1;   // marca este
					// slot como
					// ocupado,
					// "fusionandolo"
					// con el anterior
				}
				// si slotLibre == 0, no hay ningun slot disponible para
				// este trabajo -> se descarta
			}
			
			return gananciaTotal;
		}                                          // O(n log n + n * alpha(n))
	};
\end{lstlisting}

\textbf{Ejemplo de funcionamiento:}

\begin{lstlisting}[style=cpp]
	vector<JobSequencing::Trabajo> trabajos = {
		{4, 20}, {1, 10}, {1, 40}, {1, 30}
	};
	// (deadline, ganancia)
	
	cout << JobSequencing::maximaGananciaSimple(trabajos);
	// 60
	// orden por ganancia: (1,40) (1,30) (4,20) (1,10)
	// (1,40) -> slot 1 ocupado
	// (1,30) -> deadline 1, slot 1 ya ocupado, descartado
	// (4,20) -> slot 4 (el mas tardio <= 4 disponible)
	// (1,10) -> deadline 1, ocupado, descartado
	// total: 40 + 20 = 60
	
	cout << JobSequencing::maximaGananciaDSU(trabajos);
	// 60   (mismo resultado, otro camino)
\end{lstlisting}

\textbf{Regla práctica:} usa la versión con escaneo lineal cuando los
deadlines son chicos ($\lesssim 10^3$–$10^4$) por simplicidad de
código; cambia a la versión con DSU en cuanto los deadlines sean
grandes ($10^5$–$10^6$) y el escaneo lineal arriesgue exceder el
límite de tiempo — el criterio greedy (ordenar por ganancia, asignar
al slot más tardío posible) es \textbf{idéntico} en ambas, solo cambia
cómo se \textbf{busca} ese slot.


\subsection{Task Scheduler con Cooldown}

\textbf{Idea:} dado un conjunto de tareas (cada una identificada por
un tipo/letra) y un tiempo de \textbf{enfriamiento} $n$ (dos
ejecuciones de la \textbf{misma} tarea deben estar separadas por al
menos $n$ unidades de tiempo, pudiendo dejar "huecos" ociosos si no
hay otra tarea disponible para llenarlos), se busca el \textbf{tiempo
	total mínimo} para ejecutar todas las tareas. El greedy: en cada
unidad de tiempo, ejecutar \textbf{la tarea disponible con mayor
	frecuencia restante} (la que "más urge" adelantar, porque tiene más
repeticiones pendientes y por lo tanto más restricciones de
enfriamiento futuras) — se implementa con una \textbf{cola de
	prioridad} (max-heap por frecuencia) más una cola de espera para las
tareas actualmente "enfriándose". Existe además una \textbf{fórmula
	cerrada} equivalente, sin simular: si $f_{\max}$ es la frecuencia
máxima entre todas las tareas y $k$ es cuántas tareas \textbf{distintas}
alcanzan esa frecuencia máxima, la respuesta es
$\max(\texttt{totalTareas},\ (f_{\max}-1)(n+1) + k)$.

\textbf{¿Por qué usarlo?} Es un ejemplo de greedy con \textbf{dos}
soluciones válidas de complejidad muy distinta: la simulación con heap
es $O(\text{tiempoTotal} \cdot \log(\text{tiposDistintos}))$ y
funciona sin importar la distribución de frecuencias; la fórmula
cerrada es $O(m)$ (solo contar frecuencias, $m$ = número de tareas) y
explota la estructura específica del problema —una vez que se entiende
\textbf{por qué} la tarea más frecuente determina un "esqueleto"
mínimo de tiempo (bloques de tamaño $n+1$ alrededor de cada repetición
de la tarea más frecuente, salvo la última), el resto de tareas
simplemente se \textbf{rellenan} en los huecos sin afectar el total —
a menos que haya tantas tareas restantes que no quepan en los huecos,
caso cubierto por el \texttt{max} contra \texttt{totalTareas}.

\textbf{Casos de uso:}

\begin{itemize}
	\item Programar un procesador de tareas con enfriamiento
	obligatorio entre repeticiones del mismo tipo (el problema base,
	frecuente en jueces bajo el nombre "Task Scheduler").
	\item Cualquier problema de programación de trabajos con la misma
	restricción estructural: "no repetir el mismo tipo dentro de una
	ventana de $n$ unidades", sin importar el dominio narrativo
	(CPU scheduling, rotación de turnos, etc.).
	\item Como ejemplo de cuándo vale la pena \textbf{derivar una
		fórmula} en vez de simular: la simulación explícita con heap sigue
	siendo la opción más segura bajo presión de tiempo (menos
	propensa a errores de fórmula), pero para tiempos totales muy
	grandes la fórmula evita simular literalmente cada unidad de
	tiempo.
\end{itemize}

\textbf{Complejidad:} $O(m \log k)$ con la simulación por heap ($m$ =
número total de tareas, $k$ = tipos distintos); $O(m)$ con la fórmula
cerrada (solo requiere contar frecuencias).

\begin{lstlisting}[style=cpp]
	struct TaskScheduler {
		// READ: tiempo total minimo simulando con un max-heap de
		// frecuencias y una cola de "enfriamiento". Funciona sin asumir
		// nada sobre la distribucion de frecuencias.
		static long long tiempoMinimoSimulado(vector<char>& tareas, int n) {
			unordered_map<char,int> frecuencia;
			for (char t : tareas) frecuencia[t]++;
			
			priority_queue<int> heap;   // max-heap de frecuencias restantes
			for (auto& [tarea, freq] : frecuencia)
			heap.push(freq);
			
			long long tiempo = 0;
			queue<pair<int,long long>> enfriando;   // {frecuencia restante,
				// momento en que
				// vuelve a estar
				// disponible}
			
			while (!heap.empty() || !enfriando.empty()) {
				tiempo++;
				
				if (!heap.empty()) {
					int freq = heap.top(); heap.pop();
					freq--;
					if (freq > 0)
					enfriando.push({freq, tiempo + n});
					// si freq llega a 0, esta tarea ya termino
					// completamente, no vuelve a la cola
				}
				// si el heap esta vacio pero hay tareas enfriandose,
				// este ciclo es una unidad de tiempo OCIOSA (no hay
				// nada disponible para ejecutar)
				
				if (!enfriando.empty() && enfriando.front().second == tiempo)
				{
					heap.push(enfriando.front().first);
					enfriando.pop();
				}
			}
			
			return tiempo;
		}                                          // O(m log k)
		
		// READ: tiempo total minimo con la FORMULA CERRADA (sin simular).
		// f_max = frecuencia maxima; k = cuantas tareas distintas
		// alcanzan esa frecuencia maxima.
		static long long tiempoMinimoFormula(vector<char>& tareas, int n) {
			unordered_map<char,int> frecuencia;
			for (char t : tareas) frecuencia[t]++;
			
			int fMax = 0;
			for (auto& [tarea, freq] : frecuencia)
			fMax = max(fMax, freq);
			
			int k = 0;
			for (auto& [tarea, freq] : frecuencia)
			if (freq == fMax) k++;
			
			long long esqueleto = (long long)(fMax - 1) * (n + 1) + k;
			// ^ (fMax - 1) bloques de tamano (n+1) alrededor de cada
			//   repeticion de la tarea mas frecuente, salvo la ultima;
			//   +k por las k tareas que ocupan la posicion final de
			//   cada "columna" de frecuencia maxima
			
			return max((long long)tareas.size(), esqueleto);
			// si hay suficientes OTRAS tareas para llenar todos los
			// huecos del esqueleto y sobrar, la respuesta es simplemente
			// el numero total de tareas (sin tiempo ocioso)
		}                                          // O(m)
	};
\end{lstlisting}

\textbf{Ejemplo de funcionamiento:}

\begin{lstlisting}[style=cpp]
	vector<char> tareas = {'A','A','A','B','B','B'};
	int n = 2;
	
	cout << TaskScheduler::tiempoMinimoSimulado(tareas, n);
	// 8   (A B _ A B _ A B, con 2 huecos ociosos)
	
	cout << TaskScheduler::tiempoMinimoFormula(tareas, n);
	// 8   (fMax=3, k=2 -> (3-1)*(2+1)+2 = 6+2 = 8; total tareas = 6,
	//      max(6,8) = 8)
	
	vector<char> tareas2 = {'A','A','A','B','B','B','C','C','C','D','D','E'};
	cout << TaskScheduler::tiempoMinimoFormula(tareas2, 2);
	// 12  (suficientes tareas distintas para llenar todos los huecos
	//      sin tiempo ocioso -> la respuesta es simplemente el total
	//      de tareas, 12)
\end{lstlisting}

\textbf{Regla práctica:} usa la simulación con heap por defecto —es
más robusta ante errores de fórmula y funciona sin tener que
re-derivar nada bajo presión. Cambia a la fórmula cerrada solo cuando
el tiempo total esperado es demasiado grande para simular
explícitamente unidad por unidad ($n$ o el número de tareas del orden
de $10^9$), y ten cuidado con el \texttt{max} contra
\texttt{totalTareas} —omitirlo es el error más común al implementar la
fórmula, ya que sin él se puede subestimar el tiempo cuando la tarea
más frecuente no domina la distribución.
\subsection{Fractional Knapsack}

\textbf{Idea:} dado un conjunto de ítems, cada uno con un
\textbf{peso} y un \textbf{valor}, y una mochila de capacidad $W$, se
busca maximizar el valor total llevado, permitiendo tomar
\textbf{fracciones} de un ítem (a diferencia del 0/1 Knapsack, donde
cada ítem se toma completo o no se toma). El greedy: \textbf{ordenar
	por razón valor/peso decreciente}, y tomar cada ítem
\textbf{completo} mientras quepa en la capacidad restante; en cuanto
un ítem ya no cabe entero, tomar \textbf{la fracción exacta} que
llena la capacidad restante y detenerse ahí —esa fracción siempre
maximiza el valor por unidad de espacio que queda, así que nunca
conviene "guardar" espacio para un ítem de peor razón.

\textbf{¿Por qué usarlo?} Es el contraste directo y más citado de
Greedy vs. Programación Dinámica: la \textbf{única} diferencia con
0/1 Knapsack es si se permite tomar fracciones o no, y esa diferencia
\textbf{decide} si greedy es válido o no. Con fracciones permitidas,
el Argumento de Intercambio es directo: si una solución óptima no
prioriza la mejor razón valor/peso, se puede "mover" una fracción
infinitesimal de espacio desde un ítem de peor razón hacia uno de
mejor razón sin perder capacidad total, y esa operación \textbf{nunca
	empeora} el valor —por eso ordenar por razón y llenar avariciosamente
siempre es óptimo. Sin fracciones (0/1), esa operación de "mover una
fracción" ya no es posible, y el argumento se rompe por completo —de
ahí que 0/1 Knapsack necesite DP.

\textbf{Casos de uso:}

\begin{itemize}
	\item Maximizar valor transportado cuando los ítems son
	\textbf{divisibles} (líquidos, granos, cualquier recurso que se
	pueda fraccionar sin perder valor proporcional).
	\item Como introducción canónica al contraste Greedy/DP: el
	mismo enunciado, con y sin fracciones permitidas, sirve para
	ilustrar por qué la propiedad de elección greedy depende de
	detalles aparentemente menores del problema.
	\item Problemas de asignación de recursos continuos con
	"eficiencia" (valor por unidad) bien definida —cualquier variante
	donde "toma lo más eficiente primero, fracciona lo último" tenga
	sentido físico.
\end{itemize}

\textbf{Complejidad:} $O(n \log n)$, dominado por el \texttt{sort}
por razón valor/peso; el llenado greedy posterior es $O(n)$.

\begin{lstlisting}[style=cpp]
	struct FractionalKnapsack {
		struct Item { double peso, valor; };
		
		// READ: valor maximo transportable en una mochila de capacidad W,
		// permitiendo tomar fracciones de cada item.
		static double valorMaximo(vector<Item> items, double W) {
			sort(items.begin(), items.end(),
			[](Item& a, Item& b) {
				return (a.valor / a.peso) > (b.valor / b.peso);
			});
			
			double valorTotal = 0;
			double capacidadRestante = W;
			
			for (auto& it : items) {
				if (capacidadRestante <= 0) break;
				
				if (it.peso <= capacidadRestante) {
					// cabe completo -> lo toma entero
					valorTotal += it.valor;
					capacidadRestante -= it.peso;
				} else {
					// no cabe completo -> toma la FRACCION exacta que
					// llena lo que queda, y se detiene (todo lo que
					// sigue tiene PEOR razon valor/peso, ya no conviene)
					double fraccion = capacidadRestante / it.peso;
					valorTotal += it.valor * fraccion;
					capacidadRestante = 0;
				}
			}
			
			return valorTotal;
		}                                          // O(n log n)
	};
\end{lstlisting}

\textbf{Ejemplo de funcionamiento:}

\begin{lstlisting}[style=cpp]
	vector<FractionalKnapsack::Item> items = {
		{10, 60},   // razon 6.0
		{20, 100},  // razon 5.0
		{30, 120},  // razon 4.0
	};
	
	cout << FractionalKnapsack::valorMaximo(items, 50);
	// 240
	// toma completo el de razon 6.0 (peso 10, valor 60) -> quedan 40
	// toma completo el de razon 5.0 (peso 20, valor 100) -> quedan 20
	// toma 20/30 = 2/3 del de razon 4.0 -> valor 120 * 2/3 = 80
	// total: 60 + 100 + 80 = 240
\end{lstlisting}

\textbf{Nota:} si el problema pide expresamente que los ítems
\textbf{no} se puedan fraccionar (la variante 0/1), este mismo código
da una respuesta \textbf{incorrecta} (generalmente una
sobreestimación, ya que asume que cualquier fracción es viable) — la
señal en el enunciado suele ser explícita ("cada ítem se toma
completo o se descarta"), pero si no lo es, un caso de prueba donde
el óptimo fraccional y el 0/1 difieren (ítems con pesos que no
"encajan" limpiamente) revela rápidamente cuál se necesita.

\textbf{Regla práctica:} la pregunta decisiva antes de programar
cualquier variante de Knapsack es "¿se permite tomar fracciones de un
ítem?" — si sí, es este greedy en $O(n\log n)$; si no, es 0/1
Knapsack y hace falta DP en $O(nW)$, sin excepciones ni atajos
greedy posibles (ver Greedy vs. Programación Dinámica para la
justificación completa de por qué no hay forma de "adaptar" este
greedy al caso 0/1).

\section{Greedy con Heap}
\subsection{Huffman Coding}

\textbf{Idea:} Huffman Coding construye un \textbf{código de prefijo
	óptimo} (ningún código es prefijo de otro, lo que permite decodificar
sin ambigüedad) para un conjunto de símbolos con frecuencias dadas,
minimizando la \textbf{longitud total esperada} del mensaje
codificado. El greedy: mantener todos los símbolos en un
\textbf{min-heap} por frecuencia, y repetidamente \textbf{extraer los
	dos de menor frecuencia}, \textbf{fusionarlos} en un nuevo nodo
(cuya frecuencia es la suma de ambos) y \textbf{reinsertarlo} en el
heap —hasta que quede un solo nodo, la \textbf{raíz} del árbol de
Huffman. El código de cada símbolo original es el camino desde la
raíz hasta su hoja (0 por cada rama izquierda, 1 por cada derecha, o
viceversa) — los símbolos \textbf{más frecuentes} terminan más
\textbf{cerca} de la raíz (códigos más cortos), y los \textbf{menos
	frecuentes} quedan más profundos (códigos más largos).

\textbf{¿Por qué usarlo?} Es la aplicación canónica de "Greedy con
Heap": la decisión de fusionar siempre los \textbf{dos menos
	frecuentes} está justificada por Exchange Argument (ver Fundamentos
de Greedy) —cualquier árbol óptimo se puede transformar en uno donde
los dos símbolos de menor frecuencia son \textbf{hermanos} en las
hojas más profundas, sin aumentar el costo total. Es también la base
conceptual de compresión sin pérdida (ej. el algoritmo detrás de
\texttt{gzip}/\texttt{DEFLATE} usa una variante de Huffman), y el
ejemplo estándar donde "extraer siempre el mínimo, combinar,
reinsertar" con una \texttt{priority\_queue} resuelve un problema de
construcción de estructura, no solo de selección de elementos.

\textbf{Casos de uso:}

\begin{itemize}
	\item Compresión de datos: asignar códigos binarios de longitud
	variable a símbolos según su frecuencia, minimizando el tamaño
	total codificado.
	\item Cualquier problema que pida "combina elementos de a pares,
	minimizando/maximizando el costo total de combinación" donde el
	costo de cada fusión es la suma de los dos combinados (ver
	variantes de "mínimo costo para combinar $n$ pilas/cuerdas" —
	estructuralmente idéntico a Huffman, solo sin la parte de asignar
	códigos al final).
	\item Construcción de árboles de decisión balanceados por
	frecuencia de uso (símbolos usados más a menudo, más accesibles).
\end{itemize}

\textbf{Complejidad:} $O(n \log n)$ — cada uno de los $n-1$ pasos de
fusión hace 2 extracciones y 1 inserción en el heap, cada operación
$O(\log n)$.

\begin{lstlisting}[style=cpp]
	struct HuffmanCoding {
		struct Nodo {
			int frecuencia;
			int simbolo;   // -1 si es un nodo interno (fusion), el
			// simbolo real si es una hoja
			Nodo *izq, *der;
			
			Nodo(int f, int s) : frecuencia(f), simbolo(s),
			izq(nullptr), der(nullptr) {}
		};
		
		struct Comparador {
			bool operator()(Nodo* a, Nodo* b) {
				return a->frecuencia > b->frecuencia;   // min-heap
			}
		};
		
		// READ: construye el arbol de Huffman a partir de simbolos y sus
		// frecuencias. Devuelve la RAIZ del arbol.
		static Nodo* construirArbol(vector<int>& simbolos,
		vector<int>& frecuencias) {
			priority_queue<Nodo*, vector<Nodo*>, Comparador> heap;
			
			for (int i = 0; i < (int)simbolos.size(); i++)
			heap.push(new Nodo(frecuencias[i], simbolos[i]));
			
			// caso borde: un solo simbolo -> arbol trivial de 1 nodo
			if (heap.size() == 1) return heap.top();
			
			while (heap.size() > 1) {
				Nodo* a = heap.top(); heap.pop();
				Nodo* b = heap.top(); heap.pop();
				
				Nodo* fusion = new Nodo(a->frecuencia + b->frecuencia, -1);
				fusion->izq = a;
				fusion->der = b;
				
				heap.push(fusion);
			}
			
			return heap.top();
		}                                          // O(n log n)
		
		// READ: recorre el arbol y arma el codigo binario de cada
		// simbolo (0 = izquierda, 1 = derecha desde la raiz).
		static void obtenerCodigos(Nodo* raiz, string codigoActual,
		unordered_map<int,string>& codigos) {
			if (!raiz) return;
			
			if (raiz->simbolo != -1) {   // es una hoja
				// caso borde: un solo simbolo en total -> codigo "0"
				codigos[raiz->simbolo] = codigoActual.empty() ? "0" : codigoActual;
				return;
			}
			
			obtenerCodigos(raiz->izq, codigoActual + "0", codigos);
			obtenerCodigos(raiz->der, codigoActual + "1", codigos);
		}                                          // O(n) sobre el
		// arbol completo
	};
\end{lstlisting}

\textbf{Ejemplo de funcionamiento:}

\begin{lstlisting}[style=cpp]
	vector<int> simbolos =    {0, 1, 2, 3, 4};   // representan 'a','b','c','d','e'
	vector<int> frecuencias = {5, 9, 12, 13, 16};
	
	auto raiz = HuffmanCoding::construirArbol(simbolos, frecuencias);
	
	unordered_map<int,string> codigos;
	HuffmanCoding::obtenerCodigos(raiz, "", codigos);
	
	for (int s : simbolos)
	cout << s << ": " << codigos[s] << "  ";
	// 0: 1100  1: 1101  2: 100  3: 101  4: 111
	// (los simbolos con MAYOR frecuencia, como 4 con 16, obtienen
	//  codigos mas CORTOS; los de menor frecuencia, como 0 con 5,
	//  obtienen codigos mas largos)
\end{lstlisting}

\textbf{Nota:} el caso borde de un \textbf{solo} símbolo en el
alfabeto es fácil de pasar por alto —el árbol queda como un único
nodo hoja sin ningún camino real desde una raíz interna, así que hay
que asignarle explícitamente el código \texttt{"0"} (o cualquier
código de longitud 1) en vez de una cadena vacía, que no sería un
código de prefijo válido para decodificar nada.

\textbf{Regla práctica:} el mismo patrón "extrae los dos menores,
fusiona, reinserta" resuelve cualquier problema de "combina $n$
elementos de a pares minimizando el costo total de combinación,
donde combinar cuesta la suma de los dos" —aunque el problema no
mencione compresión ni códigos en absoluto (ver el caso de uso de
"mínimo costo para combinar pilas/cuerdas"). Reconocer esa estructura
general vale más que memorizar el algoritmo solo en el contexto de
compresión de datos.


\subsection{Reorganize String (Reordenar sin Adyacentes Iguales)}

\textbf{Idea:} dado un string, se busca \textbf{reorganizar} sus
caracteres de forma que \textbf{ningún} par de caracteres
\textbf{adyacentes} sea igual (si es posible; si no, se reporta que
no existe tal reorganización). El greedy: usar un \textbf{max-heap}
por frecuencia de cada carácter, y en cada paso colocar el carácter
\textbf{más frecuente disponible} que \textbf{no} sea igual al último
colocado —para evitar volver a colocar ese mismo carácter
inmediatamente después (lo cual violaría la restricción), se
\textbf{retiene} temporalmente el carácter recién usado (con su
frecuencia decrementada) y se reinserta en el heap \textbf{después}
de colocar el siguiente, dándole "un turno de espera" antes de volver
a competir por la siguiente posición.

\textbf{¿Por qué usarlo?} Comparte la estructura de "Greedy con Heap"
de Huffman y Task Scheduler con Cooldown —de hecho es
\textbf{estructuralmente el mismo problema} que Task Scheduler con
$n=1$ (cooldown de exactamente 1 posición): siempre prioriza el
carácter más frecuente disponible, porque es el que tiene \textbf{más
	riesgo} de no poder colocarse más adelante si se pospone. La
\textbf{existencia} de una reorganización válida tiene además una
condición simple y verificable de antemano: es posible si y solo si
la frecuencia \textbf{máxima} de cualquier carácter no supera
$\lceil n/2 \rceil$ (donde $n$ es la longitud del string) — si un
carácter aparece más que eso, no hay forma de "espaciarlo"
suficiente sin que dos copias queden adyacentes.

\textbf{Casos de uso:}

\begin{itemize}
	\item Reorganizar un string para que no haya caracteres
	repetidos consecutivos (el problema base, común en jueces de
	práctica).
	\item Verificar de antemano si una reorganización válida
	\textbf{existe}, sin necesariamente construirla, usando solo la
	condición de frecuencia máxima $\le \lceil n/2 \rceil$.
	\item Generalización a "distancia mínima $k$ entre repeticiones"
	(no solo evitar adyacentes) — es exactamente Task Scheduler con
	Cooldown $n = k-1$, mismo esqueleto de heap más cola de espera.
\end{itemize}

\textbf{Complejidad:} $O(n \log k)$ donde $k$ es el número de
caracteres distintos — cada carácter del string original provoca una
operación de heap ($O(\log k)$).

\begin{lstlisting}[style=cpp]
	struct ReorganizeString {
		// READ: reorganiza s para que no haya caracteres adyacentes
		// iguales. Devuelve el string reorganizado, o "" si NO es posible.
		static string reorganizar(string s) {
			int n = s.size();
			unordered_map<char,int> frecuencia;
			for (char c : s) frecuencia[c]++;
			
			// condicion de existencia: ningun caracter puede aparecer
			// mas de ceil(n/2) veces
			for (auto& [c, f] : frecuencia)
			if (f > (n + 1) / 2)
			return "";   // imposible, se detecta ANTES de
			// intentar construir nada
			
			priority_queue<pair<int,char>> heap;   // max-heap por
			// frecuencia
			for (auto& [c, f] : frecuencia)
			heap.push({f, c});
			
			string resultado;
			pair<int,char> anterior = {0, '#'};   // caracter retenido
			// (frecuencia YA
			// decrementada, en
			// espera de volver
			// al heap)
			
			while (!heap.empty()) {
				auto [freq, c] = heap.top(); heap.pop();
				
				resultado += c;
				freq--;
				
				// el caracter USADO en el paso ANTERIOR ya cumplio su
				// "turno de espera" -- ahora si puede reingresar al heap
				if (anterior.first > 0)
				heap.push(anterior);
				
				anterior = {freq, c};   // este caracter recien usado
				// queda retenido hasta el
				// proximo paso
			}
			
			return resultado;   // longitud correcta garantizada por la
			// condicion de existencia ya verificada
		}                                          // O(n log k)
	};
\end{lstlisting}

\textbf{Ejemplo de funcionamiento:}

\begin{lstlisting}[style=cpp]
	cout << ReorganizeString::reorganizar("aab");
	// "aba"   (unica reorganizacion valida salvo simetrias)
	
	cout << ReorganizeString::reorganizar("aaab");
	// ""      (la 'a' aparece 3 veces en un string de longitud 4;
	//          ceil(4/2)=2, 3 > 2 -> imposible)
	
	cout << ReorganizeString::reorganizar("vvvlo");
	// "vlvov" (u otra reorganizacion valida equivalente)
\end{lstlisting}

\textbf{Regla práctica:} este problema es Task Scheduler con Cooldown
con $n=1$ disfrazado de string —si el enunciado pide una
\textbf{distancia mínima $k$} entre repeticiones en vez de solo
"no adyacentes", el mismo esqueleto de heap + cola de espera (con
espera de $k$ pasos en vez de 1) resuelve la variante general sin
cambios estructurales. Siempre verifica la condición de existencia
($\text{frecuencia máxima} \le \lceil n/2\rceil$ para este caso, o su
generalización para $k$ arbitrario) \textbf{antes} de intentar
construir el resultado — ahorra tiempo de cómputo y evita bugs de
"casi funciona pero deja un par adyacente al final".
\subsection{Conectar Cuerdas al Mínimo Costo (Connect Ropes)}

\textbf{Idea:} dado un conjunto de $n$ cuerdas de longitudes dadas, se
busca \textbf{conectarlas todas en una sola} cuerda, donde el costo de
conectar dos cuerdas es la \textbf{suma} de sus longitudes —y se busca
minimizar el \textbf{costo total} acumulado de todas las conexiones
necesarias. El greedy: usar un \textbf{min-heap} por longitud, y
repetidamente \textbf{extraer las dos más cortas}, \textbf{sumarlas}
(ese costo se agrega al total), y \textbf{reinsertar} la cuerda
resultante —hasta que quede una sola cuerda. Es \textbf{exactamente}
el mismo esqueleto de Huffman Coding (extraer los dos menores,
fusionar, reinsertar), pero aquí la "fusión" representa una operación
física real (unir cuerdas) en vez de construir un árbol de códigos —el
\textbf{costo acumulado} de todas las fusiones es la respuesta, no un
árbol ni códigos derivados de él.

\textbf{¿Por qué usarlo?} Es el ejemplo que confirma que el patrón de
Huffman \textbf{no es exclusivo} de compresión de datos: cualquier
problema de "combina $n$ elementos de a pares, donde combinar cuesta
la suma de ambos, minimizando el costo total acumulado" tiene
\textbf{la misma} estructura y \textbf{la misma} justificación por
Exchange Argument —siempre conviene combinar las dos piezas más
pequeñas disponibles, porque cualquier costo de fusión que involucre a
una pieza pequeña se "arrastra" y se vuelve a sumar en cada fusión
posterior donde esa pieza participe; minimizar cuántas veces se
recuenta cada pieza pequeña es lo que hace óptimo el greedy.

\textbf{Casos de uso:}

\begin{itemize}
	\item Minimizar el costo total de fusionar $n$ archivos/bloques de
	datos, donde fusionar dos bloques cuesta la suma de sus tamaños
	(variante directa con otro nombre de dominio).
	\item Minimizar el costo total de una serie de operaciones de
	\textit{merge} en estructuras como \textit{k-way merge} de
	arreglos ordenados, donde el costo de cada \textit{merge} es
	proporcional al tamaño combinado.
	\item Como verificación de que un problema "nuevo" es en realidad
	Huffman sin la parte de construir códigos —si el enunciado dice
	"combina de a pares, el costo es la suma, minimiza el total", es
	esta misma plantilla, sin necesidad de construir el árbol
	explícito.
\end{itemize}

\textbf{Complejidad:} $O(n \log n)$ — cada una de las $n-1$ fusiones
hace 2 extracciones y 1 inserción en el heap, cada operación
$O(\log n)$.

\begin{lstlisting}[style=cpp]
	struct ConnectRopes {
		// READ: costo minimo total para conectar todas las cuerdas en
		// una sola. Mismo esqueleto que Huffman (extraer dos menores,
		// fusionar, reinsertar), pero acumulando el COSTO en vez de
		// construir un arbol.
		static long long costoMinimoConectar(vector<int>& longitudes) {
			priority_queue<long long, vector<long long>, greater<>> heap;
			// ^ min-heap: greater<> invierte el orden por defecto de
			//   priority_queue (que es max-heap)
			
			for (int l : longitudes)
			heap.push(l);
			
			long long costoTotal = 0;
			
			while (heap.size() > 1) {
				long long a = heap.top(); heap.pop();
				long long b = heap.top(); heap.pop();
				
				long long costoFusion = a + b;
				costoTotal += costoFusion;
				
				heap.push(costoFusion);   // la cuerda resultante vuelve
				// a competir por las
				// proximas fusiones
			}
			
			return costoTotal;
		}                                          // O(n log n)
	};
\end{lstlisting}

\textbf{Ejemplo de funcionamiento:}

\begin{lstlisting}[style=cpp]
	vector<int> longitudes = {4, 3, 2, 6};
	
	cout << ConnectRopes::costoMinimoConectar(longitudes);
	// 29
	// paso 1: fusiona 2+3=5 (costo acumulado: 5)      -> quedan {4,5,6}
	// paso 2: fusiona 4+5=9 (costo acumulado: 5+9=14)  -> quedan {6,9}
	// paso 3: fusiona 6+9=15 (costo acumulado: 14+15=29) -> queda {15}
\end{lstlisting}

\textbf{Nota:} igual que en Huffman, el caso borde de $n=1$ (una sola
cuerda) no requiere ninguna fusión —el \texttt{while (heap.size() >
	1)} ya maneja esto correctamente, devolviendo costo \texttt{0} sin
necesidad de un chequeo especial adicional.

\textbf{Regla práctica:} cualquier variante de "combina $n$ cosas de
a pares, el costo/peso de combinar es la suma, minimiza el total
acumulado" es esta misma plantilla sin importar el nombre del dominio
—reconocerla ahorra tener que razonar el Exchange Argument desde cero
cada vez, ya que es idéntico al de Huffman Coding (extraer siempre
las dos piezas más chicas disponibles).

\section{Greedy en Arreglos y Strings}
\subsection{Gas Station Problem}

\textbf{Idea:} dado un circuito circular de $n$ gasolineras, cada una
con una cantidad de combustible $\texttt{gas}[i]$ disponible y un
costo $\texttt{costo}[i]$ para llegar de esa gasolinera a la
siguiente, se busca el \textbf{único} punto de partida (si existe)
desde el cual es posible completar \textbf{todo} el circuito sin
quedarse sin combustible en ningún momento. El greedy: recorrer las
gasolineras \textbf{una sola vez}, manteniendo un \textbf{tanque
	acumulado} (\texttt{gas}[i] - \texttt{costo}[i] en cada parada); en
cuanto el tanque acumulado \textbf{cae por debajo de cero} en la
gasolinera $i$, ninguna gasolinera \textbf{entre} el candidato actual
e $i$ (inclusive) puede ser el punto de partida válido —se
\textbf{descarta todo ese rango de una sola vez}, y el nuevo candidato
pasa a ser $i+1$, reiniciando el tanque acumulado en cero.

\textbf{¿Por qué usarlo?} La observación clave que hace válido este
greedy: si el candidato actual $s$ falla al llegar a la gasolinera
$i$ (el tanque se agota ahí), \textbf{ninguna} gasolinera $j$ con
$s \le j < i$ puede ser un punto de partida válido tampoco —empezar en
$j$ en vez de $s$ solo puede llegar a $i$ con \textbf{igual o menos}
combustible acumulado (todo lo acumulado entre $s$ y $j$ se pierde, ya
que $j$ empieza en cero), así que si $s$ ya no llegaba, $j$ tampoco.
Eso permite \textbf{saltar} directamente a $i+1$ sin revisar
individualmente los candidatos descartados, dando $O(n)$ en vez de
$O(n^2)$ (probar cada punto de partida por separado con una
simulación completa).

\textbf{Casos de uso:}

\begin{itemize}
	\item El problema de gasolineras en un circuito circular tal cual
	(encontrar el único punto de partida válido, si existe).
	\item Cualquier problema de "circuito circular con recursos que se
	ganan y gastan en cada parada, encuentra desde dónde arrancar para
	nunca quedar en déficit" —la estructura se generaliza a dominios
	distintos de combustible/distancia.
	\item Verificar primero la \textbf{condición de existencia}: una
	solución existe si y solo si $\sum \texttt{gas}[i] \ge \sum
	\texttt{costo}[i]$ (el total disponible alcanza para todo el
	circuito) — si esta suma es negativa, ningún punto de partida
	funciona y no hace falta ni correr el greedy.
\end{itemize}

\textbf{Complejidad:} $O(n)$ — una sola pasada, sin necesidad de
simular el circuito completo para cada candidato.

\begin{lstlisting}[style=cpp]
	struct GasStation {
		// READ: indice del punto de partida valido para completar todo
		// el circuito, o -1 si no existe ninguno.
		static int puntoDePartida(vector<int>& gas, vector<int>& costo) {
			int n = gas.size();
			long long tanqueTotal = 0;   // suma GLOBAL de gas - costo,
			// decide si existe solucion
			long long tanqueActual = 0;  // tanque acumulado desde el
			// candidato actual
			int candidato = 0;
			
			for (int i = 0; i < n; i++) {
				long long diferencia = gas[i] - costo[i];
				tanqueTotal += diferencia;
				tanqueActual += diferencia;
				
				if (tanqueActual < 0) {
					// el candidato actual (y todo lo que hay entre el y
					// i) queda DESCARTADO de una sola vez -- el proximo
					// candidato es i+1, reiniciando el tanque
					candidato = i + 1;
					tanqueActual = 0;
				}
			}
			
			return (tanqueTotal >= 0) ? candidato : -1;
		}                                          // O(n)
	};
\end{lstlisting}

\textbf{Ejemplo de funcionamiento:}

\begin{lstlisting}[style=cpp]
	vector<int> gas =   {1, 2, 3, 4, 5};
	vector<int> costo = {3, 4, 5, 1, 2};
	
	cout << GasStation::puntoDePartida(gas, costo);
	// 3
	// desde la gasolinera 3: tanque = 4-1=3 -> 3+5-2=6 -> 6+1-3=4 ->
	// 4+2-4=2 -> 2+3-5=0, completa el circuito sin quedar negativo
\end{lstlisting}

\textbf{Nota:} el resultado es \textbf{único} cuando existe —esta
propiedad no es evidente a primera vista, pero se sigue de la misma
lógica del descarte: si hubiera dos puntos de partida válidos
distintos $s_1 < s_2$, el argumento de "todo lo que hay entre un
candidato fallido y el punto donde falla se descarta" implicaría una
contradicción sobre cuál de los dos realmente puede sostener el
circuito completo — en la práctica, basta confiar en que el greedy
encuentra el único candidato que sobrevive hasta el final del
recorrido.

\textbf{Regla práctica:} la señal para este patrón es "hay un
\textbf{acumulado} que se gana y se gasta en un recorrido, encuentra
desde dónde arrancar para que nunca sea negativo" —en cuanto el
acumulado cae bajo cero, \textbf{todo} lo recorrido hasta ahí (incluido
el candidato actual) se descarta de un salto; nunca hace falta volver
a probar esos candidatos individualmente.


\subsection{Candy Distribution}

\textbf{Idea:} dado un arreglo de \textbf{calificaciones} de $n$
niños en fila, se busca la asignación \textbf{mínima total} de dulces
tal que (1) cada niño reciba \textbf{al menos 1} dulce, y (2)
cualquier niño con calificación \textbf{mayor} que un
\textbf{vecino adyacente} reciba \textbf{más} dulces que ese vecino.
El greedy: \textbf{dos pasadas} sobre el arreglo. En la primera
(izquierda a derecha), se asegura la restricción respecto al vecino
\textbf{izquierdo}: si \texttt{calif}[i] > \texttt{calif}[i-1],
\texttt{dulces}[i] = \texttt{dulces}[i-1] + 1 (si no, \texttt{dulces}[i]
= 1). En la segunda (derecha a izquierda), se asegura la restricción
respecto al vecino \textbf{derecho}, \textbf{sin deshacer} lo ya
garantizado por la primera pasada: si \texttt{calif}[i] >
\texttt{calif}[i+1], \texttt{dulces}[i] =
\textbf{máximo}(\texttt{dulces}[i], \texttt{dulces}[i+1] + 1).

\textbf{¿Por qué usarlo?} La restricción es \textbf{bidireccional}
(cada niño compite con \textbf{ambos} vecinos simultáneamente), lo
cual hace que una sola pasada en cualquier dirección \textbf{no
	alcance} —una pasada izquierda-derecha por sí sola satisface la
restricción con el vecino izquierdo pero puede violar la del derecho
(y viceversa). El truco de \textbf{tomar el máximo} en la segunda
pasada es la clave: garantiza que el valor final de \texttt{dulces}[i]
satisface \textbf{ambas} restricciones a la vez, sin que la segunda
pasada rompa lo que la primera ya había asegurado —cada pasada resuelve
una dirección de la restricción de forma greedy y local (mirar solo un
vecino a la vez), y combinarlas con \texttt{max} resuelve el problema
bidireccional completo sin necesitar comparar contra ambos vecinos
simultáneamente en un solo paso.

\textbf{Casos de uso:}

\begin{itemize}
	\item El problema de distribución de dulces tal cual (minimizar el
	total dado un arreglo de calificaciones y restricciones locales
	bidireccionales).
	\item Cualquier problema de "asignación mínima que respeta
	restricciones de orden estricto con \textbf{ambos} vecinos en una
	secuencia" — el patrón de "dos pasadas + combinar con
	máximo/mínimo" se generaliza a variantes con más de un tipo de
	restricción por posición.
	\item Como ejemplo de que un greedy \textbf{unidireccional} no
	siempre alcanza, pero \textbf{dos} greedys unidireccionales
	combinados correctamente sí resuelven un problema bidireccional
	completo sin necesitar DP ni backtracking.
\end{itemize}

\textbf{Complejidad:} $O(n)$ — dos pasadas lineales sobre el arreglo,
más una pasada final para sumar el total.

\begin{lstlisting}[style=cpp]
	struct CandyDistribution {
		// READ: numero minimo total de dulces que satisface ambas
		// restricciones (respecto al vecino izquierdo y al derecho).
		static long long minimoDulces(vector<int>& calificaciones) {
			int n = calificaciones.size();
			vector<int> dulces(n, 1);   // cada nino arranca con 1 dulce
			// (el minimo obligatorio)
			
			// PASADA 1 (izquierda -> derecha): garantiza la restriccion
			// respecto al vecino IZQUIERDO.
			for (int i = 1; i < n; i++) {
				if (calificaciones[i] > calificaciones[i - 1])
				dulces[i] = dulces[i - 1] + 1;
			}
			
			// PASADA 2 (derecha -> izquierda): garantiza la restriccion
			// respecto al vecino DERECHO, sin deshacer la pasada 1 --
			// por eso se usa MAX en vez de sobrescribir directamente.
			for (int i = n - 2; i >= 0; i--) {
				if (calificaciones[i] > calificaciones[i + 1])
				dulces[i] = max(dulces[i], dulces[i + 1] + 1);
			}
			
			long long total = 0;
			for (int d : dulces) total += d;
			
			return total;
		}                                          // O(n)
	};
\end{lstlisting}

\textbf{Ejemplo de funcionamiento:}

\begin{lstlisting}[style=cpp]
	vector<int> calificaciones = {1, 0, 2};
	
	cout << CandyDistribution::minimoDulces(calificaciones);
	// 5
	// pasada 1: [1,1,2]  (indice 2 > indice 1 -> 1+1=2)
	// pasada 2: indice 1 < indice 0? no (0<1 es al reves, calif[1]=0 no
	//           es mayor que calif[0]=1) -- indice 0: calif[0]=1 >
	//           calif[1]=0 -> dulces[0] = max(1, dulces[1]+1) = max(1,2) = 2
	// final: [2,1,2] -> suma = 5
	
	vector<int> calificaciones2 = {1, 2, 2};
	cout << CandyDistribution::minimoDulces(calificaciones2);
	// 4
	// pasada 1: [1,2,1]  (indice 2 no es mayor que indice 1, se queda en 1)
	// pasada 2: ninguna calificacion es mayor que su vecino derecho aqui
	// final: [1,2,1] -> suma = 4
\end{lstlisting}

\textbf{Nota:} calificaciones \textbf{iguales} entre vecinos
\textbf{no} imponen ninguna restricción de orden —ambos pueden recibir
la misma cantidad de dulces (o cantidades distintas, no importa,
mientras cada uno cumpla sus restricciones individuales); el código
solo actúa cuando la comparación es \textbf{estrictamente} mayor en
alguna dirección.

\textbf{Regla práctica:} cuando una restricción greedy depende de
\textbf{ambos lados} de una posición en una secuencia (no solo del
vecino izquierdo o solo del derecho), la técnica de "dos pasadas
unidireccionales combinadas con \texttt{max}/\texttt{min}" resuelve el
problema sin necesitar comparar ambos vecinos a la vez en un único
recorrido —cada pasada es un greedy simple y correcto para \textbf{su}
dirección, y combinarlas preserva la corrección de ambas
simultáneamente.
\subsection{Jump Game (Alcanzabilidad Greedy)}

\textbf{Idea:} dado un arreglo donde $\texttt{arr}[i]$ representa el
\textbf{máximo salto} posible desde la posición $i$, se busca
determinar si es posible \textbf{alcanzar} el último índice partiendo
del primero (variante de decisión), o el \textbf{mínimo número de
	saltos} necesarios (variante de optimización). El greedy para la
variante de decisión: mantener el \textbf{alcance máximo} visto hasta
el momento (\texttt{maxAlcance}), recorriendo el arreglo posición por
posición —si en algún punto $i$ el índice actual \textbf{supera} el
alcance máximo acumulado, ese punto es \textbf{inalcanzable} y el
resto del arreglo también lo es (se corta de inmediato); si no, se
actualiza \texttt{maxAlcance} = $\max(\texttt{maxAlcance}, i +
\texttt{arr}[i])$ en cada paso. Para la variante de \textbf{mínimo
	número de saltos}, se usa la misma idea pero llevando dos "fronteras"
—el alcance del salto \textbf{actual} y el mejor alcance visto dentro
de ese salto— e incrementando el contador de saltos cada vez que se
\textbf{agota} la frontera actual.

\textbf{¿Por qué usarlo?} Es el ejemplo canónico de "no hace falta
simular \textbf{cada} salto posible (que sería exponencial, probando
todas las combinaciones de saltos)" —basta con llevar \textbf{un solo
	número} (el mejor alcance posible hasta ahora) y actualizarlo en una
sola pasada. La correctitud viene de que \textbf{nunca} importa
\textbf{desde qué posición exacta} se logró el mejor alcance actual,
solo \textbf{cuál es} ese alcance —eso colapsa un espacio de decisiones
que a primera vista parece requerir backtracking en un simple
\texttt{max} acumulado.

\textbf{Casos de uso:}

\begin{itemize}
	\item Determinar si es posible llegar al final de un arreglo de
	saltos (Jump Game I).
	\item Encontrar el \textbf{mínimo} número de saltos para llegar al
	final, asumiendo que siempre es posible (Jump Game II).
	\item Cualquier problema de "alcanzabilidad hacia adelante" donde
	cada posición desbloquea un \textbf{rango} de posiciones futuras,
	y solo importa el alcance máximo acumulado, no el camino exacto
	tomado para llegar ahí.
\end{itemize}

\textbf{Complejidad:} $O(n)$ para ambas variantes — una sola pasada
sobre el arreglo, sin recursión ni backtracking.

\begin{lstlisting}[style=cpp]
	struct JumpGame {
		// READ: es posible llegar desde el indice 0 hasta el ultimo
		// indice del arreglo? (Jump Game I)
		static bool esAlcanzable(vector<int>& arr) {
			int n = arr.size();
			int maxAlcance = 0;
			
			for (int i = 0; i < n; i++) {
				if (i > maxAlcance)
				return false;   // esta posicion (y todo lo que sigue)
				// es INALCANZABLE, corta de inmediato
				
				maxAlcance = max(maxAlcance, i + arr[i]);
				
				if (maxAlcance >= n - 1)
				return true;   // ya se garantizo llegar al final,
				// no hace falta seguir iterando
			}
			
			return maxAlcance >= n - 1;
		}                                          // O(n)
		
		// READ: minimo numero de saltos para llegar al ultimo indice,
		// ASUMIENDO que es alcanzable (llamar esAlcanzable primero si no
		// se esta seguro).
		static int minimoSaltos(vector<int>& arr) {
			int n = arr.size();
			if (n <= 1) return 0;
			
			int saltos = 0;
			int alcanceActual = 0;   // hasta donde se puede llegar CON
			// los saltos ya contados
			int mejorAlcance = 0;    // mejor alcance visto DENTRO del
			// salto actual, candidato para el
			// proximo salto
			
			for (int i = 0; i < n - 1; i++) {
				mejorAlcance = max(mejorAlcance, i + arr[i]);
				
				if (i == alcanceActual) {
					// se agoto lo que el salto actual podia cubrir ->
					// hace falta UN salto mas, extendiendo el alcance
					// al mejor visto hasta ahora
					saltos++;
					alcanceActual = mejorAlcance;
				}
			}
			
			return saltos;
		}                                          // O(n)
	};
\end{lstlisting}

\textbf{Ejemplo de funcionamiento:}

\begin{lstlisting}[style=cpp]
	vector<int> arr1 = {2, 3, 1, 1, 4};
	cout << JumpGame::esAlcanzable(arr1);
	// true
	
	vector<int> arr2 = {3, 2, 1, 0, 4};
	cout << JumpGame::esAlcanzable(arr2);
	// false   (el 0 en el indice 3 atrapa cualquier camino que llegue
	//          ahi, y ningun salto anterior alcanza a saltarselo)
	
	vector<int> arr3 = {2, 3, 1, 1, 4};
	cout << JumpGame::minimoSaltos(arr3);
	// 2   (salto 1: indice 0 -> alcanza hasta 2; salto 2: desde algun
	//      punto dentro de [1,2] alcanza hasta 4)
\end{lstlisting}

\textbf{Regla práctica:} para la variante de decisión, corta
\textbf{inmediatamente} en cuanto \texttt{i > maxAlcance} —no hace
falta terminar de recorrer el arreglo, esa condición ya prueba que el
final es inalcanzable. Para la variante de mínimo número de saltos, el
error más común es actualizar \texttt{alcanceActual} en \textbf{cada}
posición en vez de solo cuando se \textbf{agota} el salto actual
(\texttt{i == alcanceActual}) —eso sobreestima el número de saltos,
contando transiciones que en realidad no hacía falta forzar.


\subsection{Remove K Digits (Stack Monotónico Greedy)}

\textbf{Idea:} dado un número representado como string y un entero
$k$, se busca \textbf{eliminar exactamente $k$ dígitos} de forma que
el número \textbf{resultante} (manteniendo el orden relativo de los
dígitos restantes) sea el \textbf{menor posible}. El greedy: recorrer
los dígitos de izquierda a derecha manteniendo un \textbf{stack}
—mientras el stack no esté vacío, el dígito en el tope sea
\textbf{mayor} que el dígito actual, y todavía \textbf{queden}
eliminaciones disponibles ($k > 0$), se \textbf{desapila} (esa
eliminación siempre conviene: un dígito grande seguido de uno más
chico significa que quitar el grande \textbf{reduce} el número
resultante, ya que ese dígito ocupaba una posición más significativa).
Al final, si sobran eliminaciones sin usar, se quitan del
\textbf{final} del stack (los dígitos menos significativos, donde
eliminar afecta menos al valor).

\textbf{¿Por qué usarlo?} Es la aplicación de \textbf{Stack
	Monotónico} a un problema greedy de construcción \textbf{lexicográfica}
óptima: la observación clave es que un dígito en una posición más
\textbf{significativa} (más a la izquierda) pesa \textbf{más} en el
valor final que cualquier combinación de dígitos a su derecha —así que
"si el dígito actual es menor que el anterior, siempre conviene quitar
el anterior" es una decisión \textbf{local} que nunca se arrepiente,
exactamente el tipo de propiedad que hace válido un greedy con stack
en vez de backtracking sobre qué dígitos eliminar.

\textbf{Casos de uso:}

\begin{itemize}
	\item Eliminar exactamente $k$ dígitos de un número para
	minimizar el resultado (el problema base, "Remove K Digits").
	\item Variante de \textbf{maximizar} en vez de minimizar: mismo
	esqueleto, invirtiendo la condición de comparación (desapilar
	cuando el tope es \textbf{menor} que el actual, no mayor).
	\item Construcción de la subsecuencia \textbf{lexicográficamente
		más pequeña/grande} de un string, con o sin restricción de
	longitud exacta —mismo patrón de stack monotónico aplicado a
	caracteres en vez de dígitos.
	\item Como base conceptual del algoritmo "Create Maximum Number"
	(combinar dos arreglos eligiendo $k$ dígitos de cada uno para
	maximizar el resultado), que extiende esta misma idea con una
	fusión adicional.
\end{itemize}

\textbf{Complejidad:} $O(n)$ — cada dígito se apila y desapila como
mucho una vez, sin importar el valor de $k$.

\begin{lstlisting}[style=cpp]
	struct RemoveKDigits {
		// READ: el numero resultante MAS PEQUEÑO posible tras eliminar
		// exactamente k digitos de num, manteniendo el orden relativo de
		// los digitos restantes.
		static string minimoTrasEliminar(string num, int k) {
			string stack;   // se usa un string como stack para poder
			// construir el resultado directamente
			
			for (char digito : num) {
				// mientras el tope sea MAYOR que el digito actual, y
				// todavia queden eliminaciones, desapila -- esa
				// eliminacion siempre reduce el valor resultante
				while (!stack.empty() && k > 0 && stack.back() > digito) {
					stack.pop_back();
					k--;
				}
				stack.push_back(digito);
			}
			
			// si sobraron eliminaciones sin usar, se quitan del FINAL
			// (los digitos menos significativos)
			while (k > 0 && !stack.empty()) {
				stack.pop_back();
				k--;
			}
			
			// elimina ceros a la izquierda del resultado
			int inicio = 0;
			while (inicio < (int)stack.size() - 1 && stack[inicio] == '0')
			inicio++;
			
			string resultado = stack.substr(inicio);
			
			return resultado.empty() ? "0" : resultado;
		}                                          // O(n)
	};
\end{lstlisting}

\textbf{Ejemplo de funcionamiento:}

\begin{lstlisting}[style=cpp]
	cout << RemoveKDigits::minimoTrasEliminar("1432219", 3);
	// "1219"
	// stack: 1 -> 14 -> desapila 4 (1<4? no, 4>3 -> desapila 4), queda
	// "1", agrega 3 -> "13" -> agrega 2, 3>2 desapila -> "12" -> agrega 2
	// -> "122" -> agrega 1, 2>1 desapila (k aun>0) -> "121" -> agrega 1
	// -> "1211" ... (siguiendo la logica completa da "1219", eliminando
	// 4, 3, y el segundo 2 en el proceso completo)
	
	cout << RemoveKDigits::minimoTrasEliminar("10200", 1);
	// "200"
	// elimina el primer '1' (mayor que el '0' siguiente), quedan ceros
	// a la izquierda que se recortan: "0200" -> "200"
	
	cout << RemoveKDigits::minimoTrasEliminar("10", 2);
	// "0"
	// se eliminan ambos digitos, no queda nada -> resultado "0"
\end{lstlisting}

\textbf{Nota:} el manejo de \textbf{ceros a la izquierda} en el
resultado final es un detalle de borde fácil de olvidar —un número
como \texttt{"200"} nunca debería mostrarse como \texttt{"0200"}, y el
caso extremo donde \textbf{todo} el string queda vacío tras recortar
ceros (ej. eliminar todos los dígitos) debe devolver explícitamente
\texttt{"0"}, no una cadena vacía.

\textbf{Regla práctica:} la señal para Stack Monotónico en problemas
greedy de construcción de secuencias es "una decisión \textbf{local}
(comparar el elemento actual contra el tope del stack) nunca empeora
el resultado global, sin importar qué venga después" —si esa
propiedad se cumple (como aquí, gracias al peso posicional de los
dígitos), el problema se resuelve en $O(n)$ con un stack; si la
decisión óptima \textbf{sí} depende de lo que viene más adelante, la
técnica no aplica y hace falta otra herramienta.
\subsection{Lexicográficamente Mínimo/Máximo con Restricciones}

\textbf{Idea:} generaliza el patrón de Stack Monotónico de Remove K
Digits a problemas donde la construcción \textbf{no} está limitada
solo por "cuántos elementos quitar", sino por \textbf{restricciones
	adicionales} sobre qué elementos deben aparecer en el resultado —el
ejemplo canónico es \textbf{"Remove Duplicate Letters"}: dado un
string, construir la subsecuencia lexicográficamente \textbf{más
	pequeña} que contenga \textbf{cada} carácter distinto exactamente
\textbf{una vez}. El greedy sigue siendo un stack, pero con
\textbf{dos} condiciones para desapilar en vez de una: (1) el tope es
\textbf{mayor} que el carácter actual (igual que antes), \textbf{y}
(2) ese carácter del tope \textbf{vuelve a aparecer más adelante} en
el string (si no vuelve a aparecer, desapilarlo ahora significaría
\textbf{perderlo} para siempre, violando la restricción de "cada
carácter exactamente una vez"). Se necesita además un \textbf{registro
	de qué ya está en el stack} (\texttt{enStack}) para nunca agregar un
carácter duplicado.

\textbf{¿Por qué usarlo?} Es el paso natural después de Remove K
Digits: mientras ahí la única restricción era "elimina exactamente
$k$", aquí la restricción es \textbf{estructural} (cada carácter debe
sobrevivir en algún lugar del resultado) — y esa restricción
\textbf{limita} cuándo es seguro desapilar. Reconocer que el criterio
de desapilar de un Stack Monotónico puede \textbf{ampliarse} con
condiciones extra (en vez de solo comparar valores) es lo que permite
adaptar la misma técnica base a una familia mucho más amplia de
problemas de construcción lexicográfica.

\textbf{Casos de uso:}

\begin{itemize}
	\item \textbf{Remove Duplicate Letters} / \textbf{Smallest
		Subsequence of Distinct Characters}: construir la subsecuencia
	lexicográficamente mínima que contenga cada carácter exactamente
	una vez (el ejemplo desarrollado en esta subsección).
	\item Variantes de \textbf{longitud fija}: construir la
	subsecuencia lexicográficamente mínima/máxima de exactamente
	$k$ caracteres, respetando el orden original — mismo esqueleto de
	Remove K Digits, pero pensando "cuántos me quedan por recorrer"
	en vez de "cuántas eliminaciones me quedan".
	\item Construcción de la permutación/secuencia lexicográficamente
	mínima sujeta a un conjunto de restricciones de precedencia o de
	unicidad — cualquier variante donde "cada elemento debe aparecer
	exactamente una vez, en algún orden que respete el original" es
	candidata a este patrón.
	\item Para \textbf{maximizar} en vez de minimizar: invertir la
	condición de comparación (desapilar cuando el tope es
	\textbf{menor} que el actual), manteniendo las mismas
	restricciones de unicidad y última aparición.
\end{itemize}

\textbf{Complejidad:} $O(n)$ — cada carácter se apila y desapila como
mucho una vez; el precómputo de "última aparición de cada carácter" es
$O(n)$ adicional.

\begin{lstlisting}[style=cpp]
	struct LexicograficamenteMinimoConRestricciones {
		// READ: subsecuencia lexicograficamente MINIMA de s que contiene
		// cada caracter distinto EXACTAMENTE una vez (Remove Duplicate
		// Letters / Smallest Subsequence of Distinct Characters).
		static string subsecuenciaMinima(string s) {
			// precomputo: ultima posicion donde aparece cada caracter
			vector<int> ultimaApar(26, -1);
			for (int i = 0; i < (int)s.size(); i++)
			ultimaApar[s[i] - 'a'] = i;
			
			string stack;
			vector<bool> enStack(26, false);
			
			for (int i = 0; i < (int)s.size(); i++) {
				char c = s[i];
				
				if (enStack[c - 'a']) continue;   // ya esta en el
				// resultado, no se
				// puede repetir
				
				// desapila mientras: el tope es MAYOR que c (mejora el
				// orden), Y el caracter del tope VUELVE A APARECER mas
				// adelante (si no vuelve, no se puede perder ahora)
				while (!stack.empty() && stack.back() > c &&
				ultimaApar[stack.back() - 'a'] > i) {
					enStack[stack.back() - 'a'] = false;
					stack.pop_back();
				}
				
				stack.push_back(c);
				enStack[c - 'a'] = true;
			}
			
			return stack;
		}                                          // O(n)
		
		// READ: subsecuencia lexicograficamente minima de EXACTAMENTE k
		// caracteres (sin restriccion de unicidad, solo de longitud
		// fija) -- mismo esqueleto que Remove K Digits, pensado en
		// "caracteres restantes por recorrer" en vez de "eliminaciones
		// disponibles".
		static string subsecuenciaMinimaDeLongitudK(string s, int k) {
			string stack;
			int n = s.size();
			
			for (int i = 0; i < n; i++) {
				char c = s[i];
				
				// se puede desapilar mientras: el tope es mayor que c,
				// y todavia queda SUFICIENTE string restante (incluyendo
				// c) para completar los k caracteres finales
				while (!stack.empty() && stack.back() > c &&
				(int)stack.size() - 1 + (n - i) >= k) {
					stack.pop_back();
				}
				
				if ((int)stack.size() < k)
				stack.push_back(c);
			}
			
			return stack;
		}                                          // O(n)
	};
\end{lstlisting}

\textbf{Ejemplo de funcionamiento:}

\begin{lstlisting}[style=cpp]
	cout << LexicograficamenteMinimoConRestricciones::subsecuenciaMinima("bcabc");
	// "abc"
	// b entra -> "b"; c entra (c>b) -> "bc"; a entra, b>a y b reaparece
	// despues (indice 3) -> desapila b; c>a y c reaparece (indice 4) ->
	// desapila c; queda "a" -> "a"; b entra -> "ab"; c entra -> "abc"
	
	cout << LexicograficamenteMinimoConRestricciones::subsecuenciaMinima("cbacdcbc");
	// "acdb"
	
	cout << LexicograficamenteMinimoConRestricciones::subsecuenciaMinimaDeLongitudK("9142", 2);
	// "12"
\end{lstlisting}

\textbf{Nota:} la condición extra de "el tope \textbf{vuelve a
	aparecer} más adelante" es lo único que distingue este problema de
Remove K Digits — sin ella, el greedy podría desapilar un carácter que
\textbf{nunca más} vuelve a estar disponible, dejándolo fuera del
resultado y violando la restricción de que cada carácter debe
aparecer exactamente una vez. Ese chequeo (\texttt{ultimaApar[...] >
	i}) es exactamente lo que hay que agregar cada vez que un problema de
Stack Monotónico tiene una restricción de "no puedo perder este
elemento para siempre".

\textbf{Regla práctica:} cuando un problema de construcción
lexicográfica con stack tiene una restricción \textbf{extra} más allá
de "cuántos elementos caben" (unicidad, precedencia, disponibilidad
futura), la adaptación casi siempre es la misma: agregar esa
restricción como una condición \textbf{adicional} (con AND lógico) en
la guardia del \texttt{while} que decide cuándo desapilar — el resto
del esqueleto de Remove K Digits se mantiene intacto.

\section{Greedy con Matemáticas}
\subsection{Coin Change Greedy (Cuándo Funciona y Cuándo No)}

\textbf{Idea:} dado un conjunto de denominaciones de moneda y un monto
objetivo, Coin Change Greedy intenta minimizar el número de monedas
usadas tomando siempre la denominación \textbf{más grande} que quepa
en el monto restante, repitiendo hasta llegar a cero (o determinar que
no se puede). Como ya se adelantó en Greedy vs. Programación Dinámica,
este greedy es \textbf{correcto solo para sistemas de monedas
	canónicos} (como 1, 5, 10, 25) y puede \textbf{fallar} con
denominaciones arbitrarias. Esta subsección profundiza en \textbf{cómo
	verificar} si un sistema de denominaciones dado es canónico o no, en
vez de solo constatar el fallo con un contraejemplo puntual.

\textbf{¿Por qué usarlo?} Reconocer \textbf{cuándo confiar} en el
greedy es tan importante como saber implementarlo —usarlo sin
verificar sobre un sistema no canónico da una respuesta
\textbf{silenciosamente incorrecta} (compila, corre, y devuelve un
número mayor al óptimo real, sin ningún error visible). La forma
\textbf{práctica} de verificar canonicidad sin una prueba matemática
formal es el \textbf{teorema del contraejemplo mínimo}: si un sistema
de $k$ denominaciones \textbf{no} es canónico, existe un contraejemplo
menor que $c_{k-1} + c_k$ (la suma de las dos denominaciones más
grandes) — así que basta con comparar greedy contra DP para
\textbf{todos} los montos hasta ese límite; si coinciden siempre en
ese rango, el sistema es canónico para \textbf{cualquier} monto.

\textbf{Casos de uso:}

\begin{itemize}
	\item Verificar de antemano, antes de programar el greedy en
	competencia, si las denominaciones dadas en el enunciado son
	canónicas — evita el riesgo de un Wrong Answer silencioso por
	confiar en greedy sin justificación.
	\item Sistemas de monedas del mundo real (la mayoría de monedas
	nacionales) suelen ser canónicos por diseño —facilita dar cambio
	manualmente— pero un enunciado de competencia puede dar
	denominaciones \textbf{arbitrarias} a propósito, precisamente para
	descartar el greedy como solución válida.
	\item Como ejercicio de "verificación empírica acotada": en vez de
	demostrar una propiedad matemáticamente, acotar cuántos casos hace
	falta revisar para estar seguro (el mismo espíritu que las pruebas
	de Miller-Rabin en teoría de números, aunque aquí sea determinista
	y exacto, no probabilístico).
\end{itemize}

\textbf{Complejidad:} $O(k)$ para el greedy en sí (con $k$
denominaciones); $O(\text{límite} \cdot k)$ para la verificación de
canonicidad, donde $\text{límite}$ es la suma de las dos
denominaciones más grandes — mucho más costoso que el greedy mismo,
pero es un chequeo de \textbf{una sola vez} antes de confiar en él
para todos los montos futuros.

\begin{lstlisting}[style=cpp]
	struct CoinChangeGreedyVerificado {
		// READ: numero de monedas usando GREEDY (toma siempre la mas
		// grande que quepa). Puede ser INCORRECTO si el sistema no es
		// canonico -- usar solo tras verificar con esCanonico().
		static int monedasGreedy(vector<int> denominaciones, int monto) {
			sort(denominaciones.rbegin(), denominaciones.rend());
			int monedas = 0;
			
			for (int d : denominaciones) {
				monedas += monto / d;
				monto %= d;
			}
			
			return (monto == 0) ? monedas : -1;
		}                                          // O(k)
		
		// READ: numero MINIMO real de monedas via DP (siempre correcto,
		// sin importar las denominaciones).
		static int monedasDP(vector<int>& denominaciones, int monto) {
			vector<int> dp(monto + 1, INT_MAX);
			dp[0] = 0;
			
			for (int m = 1; m <= monto; m++)
			for (int d : denominaciones)
			if (d <= m && dp[m - d] != INT_MAX)
			dp[m] = min(dp[m], dp[m - d] + 1);
			
			return (dp[monto] == INT_MAX) ? -1 : dp[monto];
		}                                          // O(monto * k)
		
		// READ: el sistema de denominaciones es CANONICO (greedy siempre
		// optimo, para CUALQUIER monto)? Usa el teorema del
		// contraejemplo minimo: si greedy==DP para todo monto hasta la
		// suma de las DOS denominaciones mas grandes, es canonico.
		static bool esCanonico(vector<int> denominaciones) {
			sort(denominaciones.begin(), denominaciones.end());
			int n = denominaciones.size();
			
			if (n < 2) return true;   // con 0 o 1 denominacion, greedy
			// trivialmente coincide con el
			// optimo
			
			int limite = denominaciones[n - 1] + denominaciones[n - 2];
			
			for (int monto = 1; monto <= limite; monto++) {
				int g = monedasGreedy(denominaciones, monto);
				int d = monedasDP(denominaciones, monto);
				
				if (g != d) return false;   // encontro un contraejemplo
				// -> definitivamente NO
				// canonico
			}
			
			return true;   // ningun contraejemplo hasta el limite ->
			// canonico para CUALQUIER monto
		}                                          // O(limite^2 * k)
		// aprox (llama a DP
		// repetidamente; se
		// puede optimizar
		// reutilizando una
		// sola tabla DP hasta
		// 'limite')
	};
\end{lstlisting}

\textbf{Ejemplo de funcionamiento:}

\begin{lstlisting}[style=cpp]
	vector<int> canonico = {1, 5, 10, 25};
	cout << CoinChangeGreedyVerificado::esCanonico(canonico);
	// true
	
	vector<int> noCanonico = {1, 3, 4};
	cout << CoinChangeGreedyVerificado::esCanonico(noCanonico);
	// false
	// contraejemplo: monto=6, greedy da 3 (4+1+1), optimo real es 2 (3+3)
	
	cout << CoinChangeGreedyVerificado::monedasGreedy(canonico, 30);
	// 2   (25+5, correcto)
	
	cout << CoinChangeGreedyVerificado::monedasGreedy(noCanonico, 6);
	// 3   (INCORRECTO -- el optimo real es 2, confirmado por monedasDP)
\end{lstlisting}

\textbf{Nota:} la verificación con el teorema del contraejemplo mínimo
es \textbf{costosa} en sí misma (más que resolver un solo caso con
DP) —su valor no es acelerar el cómputo, sino dar una
\textbf{garantía de una sola vez}: una vez confirmado que un sistema
es canónico, el greedy $O(k)$ puede usarse con confianza para
\textbf{cualquier} monto futuro sin volver a verificar, en vez de
correr DP $O(\text{monto} \cdot k)$ cada vez.

\textbf{Regla práctica:} en competencia, si el enunciado da
denominaciones \textbf{fijas y conocidas} de antemano (no como parte
de la entrada variable), vale la pena \textbf{verificar canonicidad
	una sola vez} (a mano o con este código) antes de decidir si programar
el greedy simple o ir directo a DP. Si las denominaciones son
\textbf{parte de la entrada} (cambian en cada caso de prueba) y no hay
garantía de que sean canónicas, la opción segura por defecto es
\textbf{siempre DP} — el greedy solo se justifica cuando hay certeza
explícita (por el enunciado o por verificación previa) de que el
sistema es canónico.


\subsection{Egyptian Fraction}

\textbf{Idea:} una fracción egipcia representa una fracción
$p/q$ (con $p < q$) como una \textbf{suma de fracciones unitarias
	distintas} ($1/a_1 + 1/a_2 + \dots$, cada $a_i$ diferente). El greedy
—conocido como el \textbf{algoritmo de Fibonacci-Sylvester}— consiste
en tomar, en cada paso, la \textbf{mayor} fracción unitaria
$1/\lceil q/p \rceil$ que \textbf{no exceda} $p/q$, restarla, y
repetir sobre el resto —hasta que el numerador llegue a cero. La
\textbf{garantía} de que este proceso siempre \textbf{termina} viene
de una propiedad numérica: en cada paso, el numerador de la fracción
restante \textbf{estrictamente disminuye} (nunca se estanca ni crece),
así que el proceso converge en un número finito de pasos.

\textbf{¿Por qué usarlo?} Es un ejemplo de greedy donde la
\textbf{prueba de terminación} (no solo de optimalidad) es la parte no
trivial —a diferencia de Activity Selection o Fractional Knapsack,
donde el greedy simplemente "recorre y termina" porque el arreglo es
finito, aquí hay que demostrar que el algoritmo no entra en un
\textbf{ciclo infinito} restando fracciones cada vez más pequeñas sin
nunca llegar a cero. Es también un ejemplo de que "greedy" en
competencia no siempre significa "óptimo" en un sentido de
minimización clásica —el algoritmo de Fibonacci-Sylvester
\textbf{no} necesariamente da la representación con \textbf{menos}
términos posibles (ese es un problema NP-difícil en general), solo
garantiza \textbf{una} representación válida y finita.

\textbf{Casos de uso:}

\begin{itemize}
	\item Descomponer una fracción $p/q$ dada en una suma de
	fracciones unitarias distintas, cuando el problema solo pide
	\textbf{alguna} representación válida (no la de mínimo número de
	términos).
	\item Problemas de "reparto justo" modelados históricamente con
	fracciones egipcias (el origen histórico del nombre: el sistema de
	notación fraccionaria del antiguo Egipto).
	\item Como ejemplo de análisis de \textbf{terminación} de un
	proceso greedy iterativo, un tipo de razonamiento distinto (pero
	igual de necesario) al de correctitud vía Exchange Argument.
\end{itemize}

\textbf{Complejidad:} $O(\log(\max(p,q)))$ en la práctica para la
mayoría de fracciones (el numerador decrece rápidamente), aunque en el
peor caso el número de términos puede crecer más —no hay una cota
polinomial ajustada simple, pero el algoritmo siempre termina.

\begin{lstlisting}[style=cpp]
	struct EgyptianFraction {
		// READ: descompone p/q (con p < q, ambos positivos) en una lista
		// de denominadores de fracciones unitarias distintas cuya suma
		// es exactamente p/q.
		static vector<long long> descomponer(long long p, long long q) {
			vector<long long> denominadores;
			
			while (p != 0) {
				// toma la fraccion unitaria MAS GRANDE que no exceda p/q:
				// 1/ceil(q/p)
				long long unidad = (q + p - 1) / p;   // ceil(q/p)
				denominadores.push_back(unidad);
				
				// resta 1/unidad de p/q: p/q - 1/unidad =
				// (p*unidad - q) / (q*unidad)
				long long nuevoP = p * unidad - q;
				long long nuevoQ = q * unidad;
				
				long long g = __gcd(nuevoP, nuevoQ);
				if (g > 0) {
					nuevoP /= g;
					nuevoQ /= g;
				}
				
				p = nuevoP;
				q = nuevoQ;
			}
			
			return denominadores;
		}                                          // O(log(max(p,q)))
		// en la practica
	};
\end{lstlisting}

\textbf{Ejemplo de funcionamiento:}

\begin{lstlisting}[style=cpp]
	auto denominadores = EgyptianFraction::descomponer(6, 14);
	
	for (long long d : denominadores) cout << "1/" << d << " ";
	// 1/3 1/11 1/231
	// 6/14 = 3/7 -> 1/3 + 2/21 -> 1/3 + 1/11 + 1/231
	// verificacion: 1/3 + 1/11 + 1/231 = 77/231 + 21/231 + 1/231 =
	// 99/231 = 3/7 = 6/14, correcto
\end{lstlisting}

\textbf{Nota:} este algoritmo garantiza terminar y dar una
representación \textbf{válida}, pero puede generar denominadores que
\textbf{crecen muy rápido} (como en el ejemplo, de 3 y 11 salta a 231)
— si el problema requiere una representación con denominadores
\textbf{acotados} o con el \textbf{mínimo} número de términos, este
greedy simple no lo garantiza y hace falta una búsqueda más
específica (a menudo con backtracking o cotas adicionales sobre el
problema).

\textbf{Regla práctica:} usa Fibonacci-Sylvester cuando el problema
solo pide \textbf{alguna} descomposición válida en fracciones
egipcias, sin restricciones sobre el número de términos o el tamaño
de los denominadores —es la opción más simple y garantizada. Si el
enunciado pide explícitamente \textbf{minimizar} el número de
términos o acotar los denominadores, ese es un problema
estructuralmente distinto (y más difícil) que este greedy no resuelve
por sí solo.

\section{Greedy sobre Grafos (Referencias)}
\subsection{Por qué Kruskal, Prim y Dijkstra son Greedy}

\textbf{Idea:} esta subsección no introduce un algoritmo nuevo —cierra
la parte de Greedy conectando explícitamente tres algoritmos que ya
tienen su propio struct completo en la Parte de Grafos (Kruskal y Prim
en \textit{Minimum Spanning Tree}, Dijkstra en \textit{Single Source
	Shortest Path}) con el marco conceptual desarrollado aquí: los tres
toman, en cada paso, la decisión \textbf{localmente óptima} disponible
en ese momento (la arista más barata que no forma ciclo; el nodo más
cercano al conjunto ya construido) y \textbf{nunca la reconsideran}
—exactamente la definición de greedy con la que arrancó esta parte, y
los tres se justifican con la \textbf{misma técnica}: Exchange
Argument.

\textbf{¿Por qué usarlo?} Reconocer la naturaleza greedy de estos tres
algoritmos, aunque ya se conozcan como "algoritmos de grafos" con
nombre propio, refuerza que Greedy \textbf{no es una categoría aislada}
de problemas de intervalos y strings —es un \textbf{principio}
transversal que también sostiene estructuras tan centrales como MST y
caminos más cortos. Ver el mismo argumento de intercambio aplicado a
dominios tan distintos (intervalos, árboles de expansión, caminos en
grafos) es lo que convierte "demostrar un greedy" en una habilidad
\textbf{general}, en vez de una serie de trucos memorizados por
problema.

\textbf{Casos de uso (el argumento de intercambio en cada uno):}

\begin{itemize}
	\item \textbf{Kruskal} (ver MST): si una solución óptima no usa la
	arista de \textbf{menor peso} entre las que no forman ciclo con lo
	ya elegido, el corte que separa sus dos extremos debe contener
	alguna arista del MST óptimo — esa arista pesa \textbf{al menos}
	tanto como la de Kruskal, así que \textbf{intercambiarlas} nunca
	empeora el costo total.
	\item \textbf{Prim} (ver MST): en cada paso, la arista más barata
	que conecta el árbol parcial con un nodo fuera de él \textbf{debe}
	pertenecer a \textbf{algún} MST óptimo — el mismo argumento de
	"corte mínimo" que justifica Kruskal, aplicado a un árbol que
	crece desde un solo nodo en vez de fusionar componentes.
	\item \textbf{Dijkstra} (ver SSSP): cuando se extrae el nodo con
	\textbf{menor} distancia tentativa del heap, esa distancia ya es
	\textbf{definitiva} —no puede existir un camino más corto todavía
	sin descubrir, porque cualquier camino alternativo tendría que
	pasar por un nodo \textbf{no procesado}, cuya distancia tentativa
	ya es \textbf{mayor o igual} (por eso Dijkstra \textbf{falla} con
	pesos negativos: esa garantía se rompe si una arista negativa
	puede "mejorar" un nodo ya marcado como definitivo).
\end{itemize}

\textbf{Complejidad:} no aplica una nueva —los tres mantienen la
complejidad ya documentada en sus subsecciones originales
($O(E\log E)$ Kruskal, $O(E\log V)$ Prim, $O(E\log V)$ Dijkstra con
heap). El punto de esta subsección es la \textbf{justificación}
compartida, no una implementación nueva.

\begin{lstlisting}[style=cpp]
	// No hay struct nuevo aqui -- el codigo de Kruskal, Prim y Dijkstra
	// ya esta completo en sus respectivas subsecciones (MST y SSSP). Lo
	// unico que se ilustra es el PATRON COMUN de decision greedy que
	// comparten los tres, expresado como pseudocodigo unificado:
	
	// KRUSKAL:  ordena aristas por peso; para cada una, en orden,
	//           agregala SOLO SI no forma ciclo (DSU.unir devuelve true)
	//           -> decision local: "la mas barata disponible ahora
	//              mismo que no rompa la propiedad de arbol"
	
	// PRIM:     mantiene un heap de aristas frontera; extrae SIEMPRE la
	//           de menor peso que conecte con un nodo nuevo
	//           -> decision local: "la mas barata disponible ahora
	//              mismo que expanda el arbol actual"
	
	// DIJKSTRA: mantiene un heap de distancias tentativas; extrae
	//           SIEMPRE el nodo de menor distancia y lo marca definitivo
	//           -> decision local: "el mas cercano disponible ahora
	//              mismo, nunca reconsiderado despues"
	
	// las tres comparten la misma forma: extraer el MINIMO de una
	// coleccion de candidatos, confirmar esa eleccion PARA SIEMPRE, y
	// actualizar los candidatos restantes en funcion de lo recien
	// confirmado -- es el mismo esqueleto de "Greedy con Heap" ya visto
	// en Huffman Coding y Connect Ropes, aplicado a grafos en vez de a
	// simbolos o cuerdas.
\end{lstlisting}

\textbf{Ejemplo de funcionamiento:} no aplica un ejemplo nuevo de
código — los ejemplos completos con entrada/salida ya están en las
subsecciones de Kruskal, Prim y Dijkstra (Parte de Grafos). Vale la
pena, sin embargo, notar el paralelo directo: en Kruskal y Prim el
"corte" que garantiza la corrección es un corte de \textbf{grafo}
(separar nodos en dos lados); en Dijkstra, el "corte" es \textbf{entre
	nodos ya procesados y no procesados} — en ambos casos, la prueba dice
"cualquier alternativa que cruce este corte pesa al menos tanto como
mi elección greedy".

\textbf{Regla práctica:} cuando aparezca un problema nuevo que "se
sienta" como Kruskal, Prim o Dijkstra pero con una variación (pesos
distintos, restricciones extra), la primera pregunta útil es "¿el
argumento de corte mínimo (o de nodo más cercano) sigue siendo
válido con esta variación?" —si la variación introduce algo que rompe
esa garantía (como pesos negativos rompen Dijkstra), es señal de que
hace falta otro algoritmo (Bellman-Ford en ese caso), no un ajuste
menor al greedy original.

\part{Ad-hoc y Simulación}
\subsection{¿Qué algoritmo usar?}

% ---------------------------------------------------------
% TABLA 1: SIMULACION Y MATRICES
% ---------------------------------------------------------
\begin{table}[H]
	\raggedright
	\scriptsize
	\setlength{\tabcolsep}{2.5pt}
	\renewcommand{\arraystretch}{1.15}
	\begin{tabular}{|p{0.43\columnwidth}|p{0.22\columnwidth}|p{0.27\columnwidth}|}
		\hline
		\textbf{Cuando el ejercicio diga / pida} &
		\textbf{Usa} &
		\textbf{Pista rápida} \\
		\hline
		
		``Ejecuta este proceso paso a paso y dime el estado final'' (robot,
		comandos, juego) &
		Plantilla de simulación &
		Estado, paso, parada y tope de pasos. $O(T\cdot C)$. \\
		\hline
		
		Todos los elementos cambian ``al mismo tiempo'' en cada paso &
		Doble buffer &
		Calcula sobre una copia; nunca actualices en sitio. \\
		\hline
		
		``Estado tras $10^{18}$ pasos'', la regla se repite siempre igual &
		Detección de ciclos (Floyd / estado repetido) &
		$x_k=x_{\mu+((k-\mu)\bmod\lambda)}$. \\
		\hline
		
		``¿El proceso entra en un bucle?'', duplicado en arreglo de $n+1$
		valores, ciclo en lista enlazada &
		Floyd (tortuga y liebre) &
		$O(1)$ de memoria. \\
		\hline
		
		``Círculo donde se elimina cada $k$-ésimo'', ``¿quién sobrevive?'' &
		Josephus &
		$J(n)=(J(n-1)+k)\bmod n$. Respuesta base 1: suma 1. \\
		\hline
		
		Josephus pidiendo el \textbf{orden} de eliminación &
		Fenwick con $k$-ésimo elemento &
		$O(n\log n)$. \\
		\hline
		
		Josephus con $k=2$ y $n$ enorme &
		Fórmula cerrada &
		$2(n-2^{\lfloor\log_2 n\rfloor})$. $O(\log n)$. \\
		\hline
		
		Clientes, cajeros, procesos, alarmas: mucho tiempo y pocos sucesos &
		Simulación por eventos &
		Min-heap de eventos. Define el desempate de empates. \\
		\hline
		
		``Game of Life'', regla local sobre una grilla, $k$ generaciones &
		Autómata celular &
		Dos grillas. $k$ enorme: detecta ciclo sobre el tablero. \\
		\hline
		
		Autómata lineal (Regla 90) con $k$ gigante &
		Salto de $2^m$ pasos &
		$O(n\log k)$. \\
		\hline
		
		``Programa en un lenguaje inventado'', registros, saltos, Brainfuck &
		Intérprete de instrucciones &
		Parsea una vez; pon tope de pasos. \\
		\hline
		
		``Imprime en espiral'', ``llena la matriz en espiral'' &
		Espiral por fronteras &
		Cuatro fronteras y dos guardas. $O(nm)$. \\
		\hline
		
		``Por diagonales'', ``zigzag'', ``tablero en serpiente'' &
		Clave $i+j$ o $i-j$ &
		Agrupa por clave en una pasada. \\
		\hline
		
		``Rota la matriz $90^\circ$'', ``¿son la misma figura salvo giro?'' &
		Transponer y reflejar &
		Rotar: transpone y invierte filas. Canónica: menor de 8. \\
		\hline
		
		Muchas rotaciones y espejos sobre una matriz grande &
		Marcador $2\times2$ &
		Compón en una simetría. $O(q+nm)$. \\
		\hline
		
		``Candy Crush'', ``cae la arena'', ``Tetris'' &
		Gravedad con puntero de escritura &
		Marca todo, elimina, gravedad, repite. \\
		\hline
		
		Gravedad en otra dirección &
		Rotar el tablero &
		Evita duplicar código. \\
		\hline
		
		``Vecinos de una celda'', buscaminas, hielo, líneas de $k$ &
		Arreglos \texttt{dx/dy} &
		Horario desde arriba. Comprueba límites antes de acceder. \\
		\hline
		
	\end{tabular}
	\caption{Simulación y matrices: qué usar según lo que pide el ejercicio.}
\end{table}

% ---------------------------------------------------------
% TABLA 2: FECHAS, ARITMETICA Y PARSING
% ---------------------------------------------------------
\begin{table}[H]
	\raggedright
	\scriptsize
	\setlength{\tabcolsep}{2.5pt}
	\renewcommand{\arraystretch}{1.15}
	\begin{tabular}{|p{0.43\columnwidth}|p{0.22\columnwidth}|p{0.27\columnwidth}|}
		\hline
		\textbf{Cuando el ejercicio diga / pida} &
		\textbf{Usa} &
		\textbf{Pista rápida} \\
		\hline
		
		``¿Es bisiesto?'', ``días de un mes'', ``día del año'' &
		Regla gregoriana &
		Orden: 400, 100, 4. $O(1)$. \\
		\hline
		
		``¿Qué día de la semana fue?'' &
		Sakamoto (o Zeller) &
		Enero y febrero son meses del año anterior. \\
		\hline
		
		``Días entre dos fechas'', ``fecha tras sumar $k$ días'' &
		Fecha a número de día &
		Convierte, opera y vuelve. $O(1)$. \\
		\hline
		
		``Segundos a HH:MM:SS'', formato 12h/24h, sumar tiempo &
		Segundos desde medianoche &
		Hora: módulo 86400. Duración: no envolver. \\
		\hline
		
		``Imprime el resultado exacto'' con cientos de dígitos &
		BigInt (base $10^9$) &
		Producto $O(nm)$. Si piden módulo, no hace falta. \\
		\hline
		
		``Valor exacto de una razón'', comparar o sumar fracciones &
		Fracciones exactas &
		Forma canónica y \texttt{\_\_int128}. Pendientes: producto cruzado. \\
		\hline
		
		``Evalúa la expresión'' con paréntesis, precedencia y tokens pegados &
		Tokenizador + Shunting-Yard &
		Convierte a postfija una vez y reutiliza. \\
		\hline
		
		Funciones de aridad variable, asignaciones, condicionales, buenos errores &
		Descenso recursivo &
		Una función por nivel; ciclo izquierda, recursión derecha. \\
		\hline
		
	\end{tabular}
	\caption{Fechas, aritmética y parsing: qué usar según lo que pide el ejercicio.}
\end{table}

% ---------------------------------------------------------
% TABLA 3: TECNICAS AD-HOC Y PRACTICA
% ---------------------------------------------------------
\begin{table}[H]
	\raggedright
	\scriptsize
	\setlength{\tabcolsep}{2.5pt}
	\renewcommand{\arraystretch}{1.15}
	\begin{tabular}{|p{0.43\columnwidth}|p{0.22\columnwidth}|p{0.27\columnwidth}|}
		\hline
		\textbf{Cuando el ejercicio diga / pida} &
		\textbf{Usa} &
		\textbf{Pista rápida} \\
		\hline
		
		Valores hasta $10^9$ pero pocos, solo importa el orden, consultas offline &
		Compresión de coordenadas &
		\texttt{sort}, \texttt{unique}, \texttt{lower\_bound}. \\
		\hline
		
		``¿Es posible transformar $A$ en $B$?'', ``imprime $-1$ si no se puede'' &
		Invariantes y paridad &
		Suma, XOR, \texttt{mcd}, módulo $m$, coloración. \\
		\hline
		
		``¿Se puede resolver este puzzle deslizante?'', ``paridad de intercambios'' &
		Paridad de permutación &
		$n-\text{ciclos}$. $O(n)$. \\
		\hline
		
		``Construye una matriz / permutación / cuadrado mágico con tal propiedad'' &
		Algoritmo constructivo &
		Casos pequeños, patrón, demostración, verificador. \\
		\hline
		
		Cuadrado mágico de orden $n$ &
		Siamés, bloques 4x4 o Strachey &
		Según $n\bmod4$. No existe para $n=2$. \\
		\hline
		
		Pasar de $(a,b)$ a $(c,d)$ con $(x+y,y)$ o $(x,x+y)$ &
		Trabajar hacia atrás &
		Predecesor único; acelera con divisiones. \\
		\hline
		
		``Color final tras pintar rangos'', eliminaciones offline, juegos &
		Reverse thinking &
		Pinta en reversa con ``siguiente libre''; invierte el tiempo. \\
		\hline
		
		``Suma de los mínimos (o máximos) de todos los subarreglos'' &
		Contribución con pila monótona &
		Empate: \texttt{>=} a un lado y \texttt{>} al otro. $O(n)$. \\
		\hline
		
		``Suma sobre todos los subarreglos'' &
		Contribución $(i+1)(n-i)$ &
		Sin pila. $O(n)$. \\
		\hline
		
		La fórmula no es obvia &
		Fuerza bruta para descubrir patrones &
		Tabla de $n$ pequeño, diferencias, verifica. \\
		\hline
		
		Solución con Wrong Answer sin pista &
		Stress testing &
		Generador pequeño, fuerza bruta y comparación. \\
		\hline
		
		``Interactivo'', consultas al juez &
		Problema interactivo &
		\texttt{flush} tras cada línea. \\
		\hline
		
	\end{tabular}
	\caption{Técnicas ad-hoc y práctica: qué usar según lo que pide el ejercicio.}
\end{table}
\section{Simulación Directa}
\subsection{Plantilla de Simulación (Estado, Paso y Parada)}

\textbf{Idea:} un problema de simulación pide ejecutar literalmente un
proceso descrito en el enunciado y reportar qué pasa al final (o en
algún momento). Casi todos se ordenan con tres piezas:
(1) el \textbf{estado}, es decir, \textbf{todo} lo necesario para
determinar el futuro y nada más (posición, dirección, contenido de una
grilla, índice de la siguiente instrucción, etc.);
(2) el \textbf{paso}, una función que toma el estado actual y devuelve
el siguiente;
(3) la \textbf{parada}, la condición que termina la simulación, junto
con un \textbf{tope de pasos} de seguridad para no quedar en un bucle
infinito por un error de lectura del enunciado. El ciclo es siempre el
mismo: mientras no se cumpla la parada, aplicar el paso. Separar las
tres piezas evita mezclar lógica del problema con control del bucle, y
es donde se ganan (o se pierden) la mayoría de los errores.

\textbf{¿Por qué usarlo?} Los problemas de simulación rara vez fallan
por dificultad algorítmica: fallan por \textbf{detalles}. Los tres más
comunes son: (a) actualizar el estado \textbf{en sitio} cuando el
enunciado pide que todos los elementos cambien \textbf{a la vez} (el
segundo elemento ya ve el valor nuevo del primero); (b) errores de
\textbf{off-by-one} en la condición de parada; (c) olvidar parte del
estado (por ejemplo, la dirección del robot además de su posición).
Definir el estado explícitamente como un \texttt{struct} y escribir el
paso como una función pura sobre él hace estos errores fáciles de ver.

\textbf{Casos de uso:}

\begin{itemize}
	\item Ejecutar una secuencia de instrucciones (robot en una
	grilla, comandos de un editor, movimientos de un juego) y
	reportar el estado final.
	\item Simular un proceso hasta que se cumpla una condición
	("¿cuántos pasos hasta que regrese al origen?") reportando el
	número de pasos.
	\item Ser la \textbf{solución de fuerza bruta} de un problema más
	difícil: una simulación directa, lenta pero obviamente correcta,
	es el mejor comparador para probar una solución rápida (ver
	Stress Testing en esta misma parte).
	\item Base de la siguiente subsección: cuando el número de pasos
	pedido es enorme, se necesita saber \textbf{cuándo el estado se
		repite}.
\end{itemize}

\textbf{Complejidad:} $O(T \cdot C)$, donde $T$ es el número de pasos
y $C$ el costo de un paso. Sirve mientras $T \cdot C$ quepa en unos
$10^7$--$10^8$ operaciones; si el enunciado pide $T \sim 10^9$ o más,
simular literalmente ya no alcanza (ver Simulación con Detección de
Ciclos).

\begin{lstlisting}[style=cpp]
	template <class Estado>
	struct Simulador {
		function<Estado(const Estado&)> paso;    // estado -> estado siguiente
		function<bool(const Estado&)> parar;     // true si ya hay que parar
		long long maxPasos;                      // tope de seguridad
		vector<Estado> traza;                    // historial (opcional)
		bool guardarTraza;
		
		// CREATE: guarda las tres piezas de la simulacion.
		Simulador(function<Estado(const Estado&)> paso_,
		function<bool(const Estado&)> parar_,
		long long maxPasos_, bool guardarTraza_ = false)
		: paso(paso_), parar(parar_), maxPasos(maxPasos_),
		guardarTraza(guardarTraza_) {}       // O(1)
		
		// READ: corre la simulacion desde 'inicial'. Devuelve
		// {pasos ejecutados, estado final}. Si se alcanza maxPasos sin
		// que parar() sea true, devuelve con pasos == maxPasos (el
		// llamador debe revisar si eso indica un bucle infinito).
		pair<long long, Estado> ejecutar(Estado inicial) {
			Estado actual = inicial;
			long long t = 0;
			if (guardarTraza) traza.push_back(actual);
			
			while (t < maxPasos && !parar(actual)) {
				actual = paso(actual);
				t++;
				if (guardarTraza) traza.push_back(actual);
			}
			
			return {t, actual};
		}                                          // O(T * costo del paso)
		
		// ------------------------------------------------------------
		// MODIFICACIONES Y COMENTARIOS CON LOS USOS
		// ------------------------------------------------------------
		
		// 1. ACTUALIZACION SIMULTANEA (todos cambian a la vez): si el
		//    paso modifica varios elementos y el enunciado dice que
		//    todos cambian "al mismo tiempo", NUNCA se actualiza en
		//    sitio: el elemento i veria el valor NUEVO del elemento i-1.
		//    Se calcula sobre una copia (doble buffer):
		//
		//    vector<int> siguiente = actual;          // copia
		//    for (int i = 0; i < n; i++)
		//        siguiente[i] = regla(actual[i - 1], actual[i], actual[i + 1]);
		//    actual = siguiente;                      // se publica al final
		//
		//    Con Estado = vector<int> esto ya ocurre solo, porque el paso
		//    recibe una referencia constante y devuelve un estado nuevo.
		
		// 2. ACTUALIZACION SECUENCIAL (uno tras otro): si el enunciado
		//    dice que cada elemento reacciona al estado ya modificado por
		//    el anterior, SI se actualiza en sitio, en el orden pedido.
		//    Leer con cuidado que orden exacto pide el enunciado.
		
		// 3. SIMULACION POR TURNOS (rondas con cola): cuando el estado es
		//    "quien juega ahora", una cola modela el orden:
		//
		//    queue<int> turnos;
		//    while (!terminado()) {
			//        int j = turnos.front(); turnos.pop();
			//        jugar(j);
			//        if (sigueEnJuego(j)) turnos.push(j);   // vuelve al final
			//    }
		
		// 4. VER LOS PRIMEROS PASOS (depuracion): con guardarTraza = true
		//    se imprime el inicio de la traza y se compara a mano con el
		//    ejemplo del enunciado:
		//
		//    for (int i = 0; i < min<int>(10, sim.traza.size()); i++)
		//        imprimir(sim.traza[i]);
		//    Si el ejemplo del enunciado no coincide en los primeros
		//    pasos, el error esta en el paso, no en la parada.
		
		// 5. DETECTAR BUCLE INFINITO: si ejecutar() devuelve
		//    pasos == maxPasos, o bien el tope era muy chico, o bien la
		//    condicion de parada nunca se cumple. No confiar en que la
		//    salida "se ve bien": comparar pasos con maxPasos.
		
		// 6. PASOS ENORMES (T de 10^9 o mas): no simular literalmente.
		//    Si el estado toma pocos valores distintos, se repite y se
		//    puede saltar: ver Simulacion con Deteccion de Ciclos. Si el
		//    paso es una transformacion que se puede componer, tambien se
		//    puede aplicar por duplicacion (elevar la transformacion a la
		//    potencia T en O(log T) composiciones).
		
		// 7. TIEMPO CONTINUO O EVENTOS ESPARCIDOS: si casi nunca pasa nada
		//    entre eventos (y simular cada unidad de tiempo seria un
		//    desperdicio), se salta directamente de evento en evento con
		//    una cola de prioridad: ver Simulacion por Eventos.
	};
\end{lstlisting}

\textbf{Ejemplo de funcionamiento:}

\begin{lstlisting}[style=cpp]
	// Un robot recibe instrucciones: 'F' avanza, 'R' gira a la derecha,
	// 'L' gira a la izquierda. Estado = posicion + direccion + indice
	// de la siguiente instruccion (dir: 0=N, 1=E, 2=S, 3=O).
	
	struct EstadoRobot { int x, y, dir, i; };
	
	string prog = "FFRFFRFFRFFR";
	int dx[4] = {0, 1, 0, -1}, dy[4] = {1, 0, -1, 0};
	
	Simulador<EstadoRobot> sim(
	[&](const EstadoRobot& e) {                    // paso
		EstadoRobot n = e;
		if (prog[e.i] == 'F') { n.x += dx[e.dir]; n.y += dy[e.dir]; }
		else if (prog[e.i] == 'R') n.dir = (e.dir + 1) % 4;
		else n.dir = (e.dir + 3) % 4;
		n.i++;
		return n;
	},
	[&](const EstadoRobot& e) {                    // parada
		return e.i == (int)prog.size();
	},
	1000000);                                      // tope de pasos
	
	auto res = sim.ejecutar({0, 0, 0, 0});
	
	cout << res.first << " (" << res.second.x << "," << res.second.y
	<< ") dir=" << res.second.dir;
	// 12 (0,0) dir=0
	// (recorre un cuadrado de lado 2 y termina en el origen mirando al norte)
\end{lstlisting}

\textbf{Nota:} el estado del ejemplo incluye el índice \texttt{i} de la
instrucción actual aunque no sea parte de "la posición del robot":
como el paso lee \texttt{prog[e.i]}, ese índice \textbf{sí} es parte de
lo que determina el futuro. Regla para saber qué meter en el estado:
si dos estados con los mismos campos pueden evolucionar distinto, falta
algo en el estado.

\textbf{Regla práctica:} antes de escribir código, escribe en una línea
cada pieza: "estado = \dots, paso = \dots, parar cuando \dots".
Pregúntate en particular si las actualizaciones del paso son
\textbf{simultáneas} o \textbf{secuenciales} según el enunciado, y
prueba a mano los primeros 2 o 3 pasos con el ejemplo dado antes de
correr el caso completo.


\subsection{Simulación con Detección de Ciclos (Floyd y Estado Repetido)}

\textbf{Idea:} si el estado solo puede tomar un número \textbf{finito}
de valores distintos y el paso depende únicamente del estado actual, la
secuencia $x_0, x_1 = f(x_0), x_2 = f(x_1), \dots$ termina
necesariamente por \textbf{repetir} un estado (principio del palomar).
Desde ese momento se repite para siempre: la secuencia tiene una
\textbf{cola} de longitud $\mu$ (estados que se visitan una sola vez)
seguida de un \textbf{ciclo} de longitud $\lambda$. Una vez conocidos
$\mu$ y $\lambda$, el estado tras $k$ pasos (con $k$ tan grande como
$10^{18}$) se obtiene reduciendo el índice:
$$x_k = x_{\mu + ((k-\mu) \bmod \lambda)} \quad \text{si } k \ge \mu$$
y simulando solo $\mu + \lambda$ pasos como máximo. Hay dos formas de
encontrar $\mu$ y $\lambda$: \textbf{Floyd} (tortuga y liebre), y
\textbf{Estado Repetido} (guardar cada estado visto con el paso en que
apareció).

\textbf{¿Por qué usarlo?} Cuando el enunciado pide "el estado tras
$10^{18}$ pasos", simular literalmente es imposible, pero casi siempre
hay un ciclo corto escondido. Las dos técnicas se complementan:
\textbf{Floyd} usa $O(1)$ de memoria y solo necesita poder comparar dos
estados con \texttt{==}, ideal cuando el estado es un entero o algo
pequeño; \textbf{Estado Repetido} usa memoria $O(\mu+\lambda)$ pero es
más simple, sirve con cualquier estado que se pueda guardar en un
\texttt{map} (strings, vectores, tuplas), y además deja la
\textbf{historia completa} guardada, lo que permite responder muchas
consultas de "¿qué estado hay en el paso $k$?" en $O(1)$ cada una.

\textbf{Casos de uso:}

\begin{itemize}
	\item Estado tras $k$ pasos de una regla que se aplica muchísimas
	veces (baraja que se mezcla siempre igual, configuración que
	rota, autómata con reglas fijas, números tipo "feliz").
	\item Detectar si una secuencia \textbf{entra en un bucle} o
	termina (por ejemplo, "¿este proceso llega alguna vez a 1?").
	\item Detectar un ciclo en una lista enlazada o en un grafo donde
	cada nodo tiene una sola arista de salida (grafo funcional).
	\item Encontrar el elemento duplicado en un arreglo de $n+1$
	valores en $[1,n]$ tratando \texttt{a[i]} como puntero al
	siguiente índice: el duplicado es la \textbf{entrada} del ciclo
	(ver modificación 2).
	\item Calcular sumas o costos acumulados en $k$ pasos gigantes,
	separando cola, ciclos completos y resto (modificación 4).
\end{itemize}

\textbf{Complejidad:} Floyd: $O(\mu + \lambda)$ evaluaciones de $f$ y
$O(1)$ de memoria. Estado Repetido: $O((\mu+\lambda)\log(\mu+\lambda))$
con \texttt{map} ($O(\mu+\lambda)$ con \texttt{unordered\_map}) y
$O(\mu+\lambda)$ de memoria. En ambos, obtener $x_k$ cuesta
$O(\mu+\lambda)$ con Floyd (hay que reiterar), y $O(1)$ con Estado
Repetido (ya está la historia guardada).

\begin{lstlisting}[style=cpp]
	// ---------- Variante 1: Floyd (tortuga y liebre), estado = entero ----------
	struct DeteccionCiclo {
		long long mu, lambda;                 // longitud de la cola y del ciclo
		function<long long(long long)> f;     // el paso: x -> f(x)
		long long x0;                         // estado inicial
		
		// CREATE: encuentra mu y lambda con el algoritmo de Floyd.
		DeteccionCiclo(function<long long(long long)> f_, long long x0_) {
			f = f_; x0 = x0_;
			
			// fase 1: la liebre avanza 2 pasos por cada paso de la
			// tortuga; se encuentran DENTRO del ciclo
			long long tortuga = f(x0), liebre = f(f(x0));
			while (tortuga != liebre) {
				tortuga = f(tortuga);
				liebre = f(f(liebre));
			}
			
			// fase 2: la tortuga vuelve al inicio; ahora ambas avanzan
			// de a 1 paso y se encuentran justo en la ENTRADA del ciclo
			mu = 0;
			tortuga = x0;
			while (tortuga != liebre) {
				tortuga = f(tortuga);
				liebre = f(liebre);
				mu++;
			}
			
			// fase 3: se mide el largo del ciclo dando una vuelta completa
			lambda = 1;
			liebre = f(tortuga);
			while (tortuga != liebre) {
				liebre = f(liebre);
				lambda++;
			}
		}                                          // O(mu + lambda)
		
		// READ: estado tras k pasos (k puede ser enorme).
		long long valorEnPaso(long long k) {
			long long pasos = k;
			if (k >= mu) pasos = mu + (k - mu) % lambda;   // reduce k
			
			long long x = x0;
			for (long long i = 0; i < pasos; i++) x = f(x);
			return x;
		}                                          // O(mu + lambda)
	};
	
	// ---------- Variante 2: Estado Repetido, estado = cualquier tipo ----------
	// Requiere que Estado tenga operator< (para usar std::map). Con
	// unordered_map hace falta ademas un hash para Estado.
	template <class Estado>
	struct CicloPorEstadoRepetido {
		long long mu, lambda;
		vector<Estado> historia;              // historia[t] = estado en el paso t
		
		// CREATE: simula hasta ver un estado repetido, guardando en que
		// paso se vio cada estado por primera vez.
		CicloPorEstadoRepetido(Estado inicial,
		function<Estado(const Estado&)> paso) {
			map<Estado, long long> visto;
			Estado actual = inicial;
			long long t = 0;
			
			while (!visto.count(actual)) {
				visto[actual] = t;
				historia.push_back(actual);
				actual = paso(actual);
				t++;
			}
			
			mu = visto[actual];               // paso en que entro al ciclo
			lambda = t - mu;                  // largo del ciclo
		}                                          // O((mu+lambda) log)
		
		// READ: estado tras k pasos (k puede ser enorme).
		Estado estadoEnPaso(long long k) {
			if (k < (long long)historia.size()) return historia[k];
			return historia[mu + (k - mu) % lambda];
		}                                          // O(1)
		
		// ------------------------------------------------------------
		// MODIFICACIONES Y COMENTARIOS CON LOS USOS
		// ------------------------------------------------------------
		
		// 1. ALGORITMO DE BRENT (alternativa a Floyd): en vez de mover la
		//    liebre al doble de velocidad, se compara contra un "punto de
		//    control" que se reubica en potencias de 2. Hace en promedio
		//    menos evaluaciones de f, lo que importa si f es costosa:
		//
		//    long long potencia = 1, lam = 1;
		//    long long tortuga = x0, liebre = f(x0);
		//    while (tortuga != liebre) {
			//        if (potencia == lam) {          // reubica el control
				//            tortuga = liebre;
				//            potencia *= 2;
				//            lam = 0;
				//        }
			//        liebre = f(liebre);
			//        lam++;
			//    }
		//    // lam es lambda; mu se halla avanzando dos punteros con
		//    // separacion lam desde x0, igual que la fase 2 de Floyd.
		
		// 2. ELEMENTO DUPLICADO CON FLOYD: arreglo a de n+1 valores en
		//    [1,n]. Se ve como un grafo funcional i -> a[i]; como hay un
		//    valor repetido, dos indices distintos apuntan al mismo
		//    destino y ese destino es la ENTRADA del ciclo (o sea, mu):
		//
		//    int tor = a[0], lie = a[a[0]];
		//    while (tor != lie) { tor = a[tor]; lie = a[a[lie]]; }
		//    tor = 0;
		//    while (tor != lie) { tor = a[tor]; lie = a[lie]; }
		//    // tor es el valor duplicado; O(n) tiempo y O(1) memoria, sin
		//    // modificar el arreglo
		
		// 3. LISTA ENLAZADA CON CICLO: el mismo Floyd, con punteros:
		//
		//    Nodo* tor = cabeza; Nodo* lie = cabeza;
		//    while (lie && lie->sig) {
			//        tor = tor->sig;
			//        lie = lie->sig->sig;
			//        if (tor == lie) return true;    // hay ciclo
			//    }
		//    return false;                        // la lista termino
		
		// 4. SUMAS ACUMULADAS EN k PASOS: si cada paso aporta un valor
		//    v(x) (puntos ganados, costo pagado) y se pide el total tras k
		//    pasos, con historia[] guardada se precalculan sumas de
		//    prefijo pref[t] = suma de v(historia[0..t-1]) y:
		//
		//    if (k <= mu + lambda) return pref[k];
		//    long long enElCiclo = pref[mu + lambda] - pref[mu];
		//    long long vueltas = (k - mu) / lambda;
		//    long long resto = (k - mu) % lambda;
		//    return pref[mu] + vueltas * enElCiclo + (pref[mu + resto] - pref[mu]);
		//
		//    Cuidado con el overflow: vueltas * enElCiclo puede pasar de
		//    long long si k es muy grande (usar __int128 o modulo).
		
		// 5. MUCHOS PUNTEROS FUNCIONALES A LA VEZ (grafo funcional con n
		//    nodos y muchas consultas "donde estoy tras k pasos si empiezo
		//    en u"): en vez de buscar un ciclo por consulta, se
		//    precalcula la tabla de saltos de potencias de 2, igual que en
		//    Binary Lifting / k-esimo Ancestro: up[j][u] = nodo tras 2^j
		//    pasos desde u. Cada consulta cuesta O(log k).
		
		// 6. TRAMPA: EL PASO DEPENDE DEL TIEMPO. Todo lo anterior supone
		//    que el paso depende SOLO del estado. Si el enunciado dice
		//    "en el paso t se hace X" (por ejemplo, un comportamiento que
		//    alterna entre pares e impares), el tiempo forma parte del
		//    estado. Se corrige metiendo t (o t mod periodo) dentro del
		//    Estado; si no, se detecta un "ciclo" falso y la respuesta sale
		//    mal.
		
		// 7. TRAMPA: ESPACIO DE ESTADOS ENORME. Detectar ciclos solo sirve
		//    si el ciclo aparece pronto (mu + lambda cabe en tiempo y
		//    memoria). Con estados de 2^60 valores posibles el ciclo puede
		//    ser inalcanzable en la practica; ahi hay que buscar una
		//    formula, duplicacion, o estructura especial del problema.
	};
\end{lstlisting}

\textbf{Ejemplo de funcionamiento:}

\begin{lstlisting}[style=cpp]
	// Floyd con f(x) = (x*x + 1) % 10, empezando en x0 = 3:
	// 3, 0, 1, 2, 5, 6, 7, 0, 1, 2, 5, 6, 7, 0, ...
	// el 3 solo aparece una vez (cola de largo 1) y luego el ciclo
	// 0, 1, 2, 5, 6, 7 se repite (largo 6)
	
	DeteccionCiclo dc([](long long x) { return (x * x + 1) % 10; }, 3);
	
	cout << dc.mu << " " << dc.lambda;
	// 1 6
	
	cout << dc.valorEnPaso(2) << " " << dc.valorEnPaso(1000000000000000000LL);
	// 1 5   (x_2 = 1; x_(10^18) = x_4 = 5, sin simular 10^18 pasos)
	
	// Estado Repetido con strings: el paso rota "abcde" una posicion a la
	// izquierda e intercambia luego los dos primeros caracteres.
	
	CicloPorEstadoRepetido<string> cr("abcde", [](const string& s) {
		string t = s.substr(1) + s[0];
		swap(t[0], t[1]);
		return t;
	});
	
	cout << cr.mu << " " << cr.lambda;
	// 0 4   (vuelve a "abcde" tras 4 pasos: abcde cbdea dbeac ebacd abcde)
	
	cout << cr.estadoEnPaso(1000000000000LL);
	// abcde   (10^12 es multiplo de 4, cae justo al inicio del ciclo)
\end{lstlisting}

\textbf{Nota:} en el ejemplo de Estado Repetido \texttt{mu = 0} porque
el estado inicial \textbf{forma parte} del ciclo (la secuencia regresa
exactamente a él); en el ejemplo de Floyd \texttt{mu = 1} porque el
estado inicial (\texttt{3}) nunca vuelve a aparecer. Ambos casos son
normales, y la fórmula $\mu + ((k-\mu) \bmod \lambda)$ los maneja igual.
Lo que no debe olvidarse es el \texttt{if (k >= mu)}: para $k < \mu$
el índice \textbf{no} se reduce.

\textbf{Regla práctica:} si el enunciado pide el estado tras un número
gigantesco de pasos ($10^9$, $10^{12}$, $10^{18}$) y el estado puede
tomar pocos valores distintos, la respuesta casi siempre es "hay un
ciclo": usa Estado Repetido si el estado es complejo o si vas a
responder varias consultas, y Floyd si el estado es un entero y la
memoria importa. Antes de codificar, confirma dos cosas: que el paso
depende \textbf{solo} del estado (no del tiempo), y que el ciclo
aparece lo bastante rápido como para poder encontrarlo simulando.
\subsection{Problema de Josephus}

\textbf{Idea:} $n$ personas numeradas de $0$ a $n-1$ forman un círculo,
y partiendo de la persona $0$ se elimina a la $k$-ésima, luego se sigue
contando desde la siguiente y se vuelve a eliminar a la $k$-ésima,
hasta que queda una sola. Se pide quién \textbf{sobrevive} (o en qué
\textbf{orden} se eliminan). La forma más directa de simularlo es una
cola: se rotan $k-1$ personas al final y se saca a la siguiente, en
$O(nk)$. Hay una recurrencia mucho más rápida. Sea $J(n)$ la posición
(base $0$) del sobreviviente con $n$ personas. Tras eliminar a la
persona en la posición $(k-1) \bmod n$, la cuenta se reinicia desde la
siguiente, y queda un círculo de $n-1$ personas donde la posición $i$ de
ese círculo corresponde a la posición $(i + k) \bmod n$ del círculo
original. Por eso:
$$J(1) = 0, \qquad J(n) = \big(J(n-1) + k\big) \bmod n$$
que se evalúa en $O(n)$ con un solo ciclo.

\textbf{¿Por qué usarlo?} Cada método responde una pregunta distinta y
conviene elegir bien: la \textbf{recurrencia} $O(n)$ da solo al
sobreviviente; para el \textbf{orden completo} de eliminación hace falta
una estructura que encuentre "la persona en la posición $p$ entre las
que quedan" y la borre, que es exactamente lo que hace un Fenwick Tree
con búsqueda del $k$-ésimo elemento ($O(n \log n)$); para $k=2$ existe
una \textbf{fórmula cerrada} en $O(\log n)$; y si $n$ es gigantesco pero
$k$ es chico, hay una variante que salta bloques y corre en
$O(k \log n)$. Reconocer cuál se pide evita simular $O(nk)$ cuando
$n$ es del orden de $10^6$ o más.

\textbf{Casos de uso:}

\begin{itemize}
	\item Quién sobrevive en un círculo donde se elimina cada $k$-ésimo
	(el problema clásico, con $n$ hasta $10^6$--$10^7$ o más con la
	recurrencia).
	\item El \textbf{orden completo} de eliminación, o quién es el
	$m$-ésimo eliminado, con el método del Fenwick.
	\item Casos con $k=2$ y $n$ muy grande (hasta $10^{18}$), donde la
	fórmula cerrada resuelve en tiempo logarítmico.
	\item Problemas disfrazados de juego de niños, de fila que se
	reduce, o de "elimina cada segundo elemento hasta que quede uno":
	el modelo es el mismo.
\end{itemize}

\textbf{Complejidad:} recurrencia $O(n)$ tiempo y $O(1)$ memoria;
orden completo con Fenwick $O(n \log n)$; $k=2$ en $O(\log n)$ (o
$O(1)$ con \texttt{clz}); $k$ chico y $n$ enorme en $O(k \log n)$;
simulación con cola $O(nk)$, que sirve como referencia para verificar.

\begin{lstlisting}[style=cpp]
	struct Josephus {
		// READ: posicion (base 0) del sobreviviente con n personas y
		// salto k, usando la recurrencia J(n) = (J(n-1) + k) % n.
		static int sobreviviente(int n, int k) {
			int r = 0;                               // J(1) = 0
			for (int i = 2; i <= n; i++)
			r = (r + k) % i;
			return r;
		}                                          // O(n)
		
		// Fenwick de conteos (1 = persona viva) con busqueda del
		// k-esimo elemento vivo por descenso binario.
		struct Fenwick {
			int n, LOG;
			vector<int> t;
			
			Fenwick(int n_) : n(n_), t(n_ + 1, 0) {
				for (int i = 1; i <= n; i++) {       // construccion lineal:
					t[i] += 1;                        // todos empiezan vivos
					int j = i + (i & -i);
					if (j <= n) t[j] += t[i];
				}
				LOG = 1;
				while ((1 << LOG) <= n) LOG++;
			}
			
			void quitar(int i) {                     // i en base 0
				for (i++; i <= n; i += i & -i) t[i]--;
			}
			
			int kesimo(int k) {                      // k-esimo vivo, k >= 1;
				int pos = 0;                          // devuelve indice base 0
				for (int j = LOG; j >= 0; j--) {
					int sig = pos + (1 << j);
					if (sig <= n && t[sig] < k) {
						pos = sig;
						k -= t[sig];
					}
				}
				return pos;
			}
		};
		
		// READ: orden completo de eliminacion (el ultimo es el
		// sobreviviente).
		static vector<int> ordenEliminacion(int n, int k) {
			Fenwick fw(n);
			vector<int> orden;
			int pos = 0, restantes = n;   // pos: posicion entre los vivos
			
			while (restantes > 0) {
				pos = (pos + k - 1) % restantes;
				int persona = fw.kesimo(pos + 1);    // quien esta en 'pos'
				orden.push_back(persona);
				fw.quitar(persona);
				restantes--;
				// pos no avanza: al borrar, la siguiente persona ocupa la
				// misma posicion entre los vivos
			}
			
			return orden;
		}                                          // O(n log n)
		
		// READ: caso k = 2 con n enorme. El sobreviviente (base 0) es
		// 2 * (n - p), con p la mayor potencia de 2 menor o igual que n.
		static long long sobrevivienteK2(long long n) {
			long long p = 1;
			while (p * 2 <= n) p *= 2;
			return 2 * (n - p);
		}                                          // O(log n)
		
		// READ: simulacion directa con cola, para verificar (stress
		// testing) o para n y k muy chicos.
		static vector<int> simulacionCola(int n, int k) {
			queue<int> q;
			for (int i = 0; i < n; i++) q.push(i);
			vector<int> orden;
			
			while (!q.empty()) {
				for (int j = 1; j < k; j++) {        // rota k-1 personas
					q.push(q.front());
					q.pop();
				}
				orden.push_back(q.front());          // la k-esima sale
				q.pop();
			}
			
			return orden;
		}                                          // O(n * k)
		
		// ------------------------------------------------------------
		// MODIFICACIONES Y COMENTARIOS CON LOS USOS
		// ------------------------------------------------------------
		
		// 1. INDICES EN BASE 1: casi todos los enunciados numeran desde
		//    1. Todo lo de arriba trabaja en base 0, asi que la respuesta
		//    final es sobreviviente(n, k) + 1 (y lo mismo para cada
		//    elemento de ordenEliminacion). Es el error mas comun.
		
		// 2. n ENORME Y k CHICO (por ejemplo n = 10^12, k = 5): la
		//    recurrencia O(n) es lenta, pero se puede eliminar k-1
		//    personas de cada k en un solo salto:
		//
		//    long long josephusChico(long long n, long long k) {
			//        if (n == 1) return 0;
			//        if (k == 1) return n - 1;
			//        if (k > n) return (josephusChico(n - 1, k) + k) % n;
			//        long long cnt = n / k;                  // eliminados en
			//                                                  // esta vuelta
			//        long long res = josephusChico(n - cnt, k);
			//        res -= n % k;
			//        if (res < 0) res += n;
			//        else res += res / (k - 1);
			//        return res;
			//    }                                           // O(k log n)
		
		// 3. SALTO QUE CAMBIA EN CADA RONDA (paso[i] distinto en cada
		//    eliminacion): la recurrencia ya no aplica directo, pero el
		//    metodo del Fenwick lo maneja sin cambios, usando el salto
		//    de la ronda actual:
		//
		//    pos = (pos + paso[ronda] - 1) % restantes;
		
		// 4. CONTAR EN SENTIDO CONTRARIO O ALTERNANDO: con Fenwick basta
		//    restar en vez de sumar, cuidando que el modulo sea no
		//    negativo:
		//
		//    pos = ((pos - k) % restantes + restantes) % restantes;
		
		// 5. ALTERNATIVA CON ARBOL DE ORDEN (GNU pbds), mismo O(n log n):
		//
		//    #include <ext/pb_ds/assoc_container.hpp>
		//    #include <ext/pb_ds/tree_policy.hpp>
		//    using namespace __gnu_pbds;
		//    tree<int, null_type, less<int>, rb_tree_tag,
		//         tree_order_statistics_node_update> vivos;
		//    for (int i = 0; i < n; i++) vivos.insert(i);
		//    int pos = 0;
		//    while (!vivos.empty()) {
			//        pos = (pos + k - 1) % vivos.size();
			//        auto it = vivos.find_by_order(pos);     // el k-esimo vivo
			//        orden.push_back(*it);
			//        vivos.erase(it);
			//    }
		
		// 5b. ALTERNATIVA CON std::list: se avanza el iterador k-1 veces
		//    (con cierre circular) y se borra. Es O(n * k), util solo con
		//    n y k chicos.
		
		// 6. OTRA PREGUNTA, MISMO MODELO: "cuantas personas quedan
		//    cuando se eliminan solo m de ellas" se responde cortando
		//    ordenEliminacion en la m-esima; "cual es la ultima persona
		//    que elimina a otra" es el penultimo de ordenEliminacion.
	};
\end{lstlisting}

\textbf{Ejemplo de funcionamiento:}

\begin{lstlisting}[style=cpp]
	cout << Josephus::sobreviviente(7, 3);
	// 3   (base 0; en base 1 es la persona 4)
	
	for (int x : Josephus::ordenEliminacion(7, 3)) cout << x << " ";
	// 2 5 1 6 4 0 3   (el ultimo, 3, es el sobreviviente)
	
	for (int x : Josephus::simulacionCola(7, 3)) cout << x << " ";
	// 2 5 1 6 4 0 3   (la simulacion directa da lo mismo)
	
	// el caso historico: n = 41, k = 3
	cout << Josephus::sobreviviente(41, 3) + 1;
	// 31   (la posicion 31 en base 1)
	
	cout << Josephus::sobrevivienteK2(41) + 1;
	// 19   (k = 2, n = 41: 2 * (41 - 32) = 18 en base 0)
\end{lstlisting}

\textbf{Nota:} la recurrencia da \textbf{solo} el último sobreviviente,
no el orden. Si un enunciado pide "el segundo en salir" o "la posición
de la persona $X$ en el orden de eliminación", la recurrencia no
alcanza y hay que usar el método con Fenwick (o simular). Otro detalle
fácil de perder: en \texttt{ordenEliminacion}, la variable \texttt{pos}
\textbf{no} se incrementa después de borrar, porque al eliminar a la
persona en esa posición, la siguiente ocupa el mismo lugar entre las que
quedan.

\textbf{Regla práctica:} si solo piden al sobreviviente, usa la
recurrencia $O(n)$ (o la fórmula cerrada si $k=2$ y $n$ es enorme, o la
variante $O(k \log n)$ si $k$ es chico y $n$ enorme); si piden el orden
completo, usa el Fenwick con búsqueda del $k$-ésimo. Simula con cola
solo con $n \cdot k$ pequeño, o para verificar tus otras versiones
contra un caso conocido.


\subsection{Simulación por Eventos (Cola de Prioridad de Eventos)}

\textbf{Idea:} en una simulación por pasos (ver Plantilla de
Simulación) el reloj avanza una unidad cada vez, y cada unidad se
revisa aunque no pase nada. Cuando el tiempo total es enorme
($10^9$ o más) pero los \textbf{sucesos} son pocos (llegadas, fines de
servicio, alarmas), conviene dejar de simular tiempo y simular
\textbf{sucesos}: se mantiene una cola de prioridad (min-heap ordenado
por tiempo) con los \textbf{eventos futuros}; se extrae el más
próximo, el reloj \textbf{salta} directamente a su tiempo, se
actualiza el estado, y ese evento puede \textbf{programar} otros
nuevos (por ejemplo, empezar a atender a un cliente programa el evento
"termina de ser atendido"). La simulación termina cuando la cola se
vacía o se pasa un tiempo límite.

\textbf{¿Por qué usarlo?} El costo depende del número de eventos
$E$, no de cuánto dura el tiempo simulado: $O(E \log E)$ en vez de
$O(T)$. Además es el patrón que conecta varias cosas de esta guía:
\textbf{Dijkstra} es una simulación por eventos donde cada evento es
"el nodo $v$ es alcanzado en el tiempo $d$"; \textbf{Line Sweep}
recorre eventos ordenados de antemano, mientras que aquí los eventos
se \textbf{generan mientras se simula}. Lo delicado son los
\textbf{desempates}: dos eventos en el mismo instante deben procesarse
en el orden que pide el enunciado, y hay que fijarlo explícitamente en
el comparador.

\textbf{Casos de uso:}

\begin{itemize}
	\item Colas de servicio: bancos, cajeros, impresoras, barberías,
	con $k$ servidores y clientes que llegan en distintos momentos
	(el ejemplo desarrollado abajo).
	\item Planificación de procesos en un CPU (\textit{round robin},
	\textit{shortest job first}): los eventos son "termina el
	quantum" y "llega un proceso".
	\item Propagación con tiempos distintos (un incendio que tarda
	tiempos diferentes en cruzar cada tramo): el evento es "este
	lugar se enciende en el tiempo $t$".
	\item Temporizadores y alarmas: "qué suena primero", "cuántas
	veces suena antes del tiempo $T$" con eventos periódicos.
	\item Juegos por turnos donde cada personaje tiene una velocidad
	distinta: el evento es "le toca actuar a $X$ en el tiempo $t$".
\end{itemize}

\textbf{Complejidad:} $O(E \log E)$, con $E$ el número total de
eventos. En la simulación de cajeros cada cliente genera a lo más 2
eventos (llegada y fin), así que $E \le 2n$ y el costo es
$O(n \log n)$, sin importar qué tan grandes sean los tiempos.

\begin{lstlisting}[style=cpp]
	struct SimulacionEventos {
		struct Evento {
			long long t;      // tiempo en que ocurre
			int tipo;         // FIN = 0, LLEGADA = 1
			int id;           // cliente involucrado
			
			// orden para un min-heap: primero el tiempo menor; si empatan,
			// primero el de menor 'tipo' (FIN antes que LLEGADA: un cajero
			// que se libera justo cuando llega alguien queda libre para el);
			// el id desempata para que el resultado sea determinista
			bool operator>(const Evento& o) const {
				if (t != o.t) return t > o.t;
				if (tipo != o.tipo) return tipo > o.tipo;
				return id > o.id;
			}
		};
		static const int FIN = 0, LLEGADA = 1;
		
		vector<long long> llegada, duracion;   // datos de cada cliente
		vector<long long> inicio, fin;         // resultados de la simulacion
		int cajeros;
		
		// CREATE: clientes[i] = {momento de llegada, duracion del servicio};
		// hay 'cajeros' servidores identicos y una fila unica FIFO.
		SimulacionEventos(vector<pair<long long,long long>>& clientes,
		int cajeros_) {
			cajeros = cajeros_;
			int n = clientes.size();
			llegada.resize(n);
			duracion.resize(n);
			inicio.assign(n, -1);
			fin.assign(n, -1);
			for (int i = 0; i < n; i++) {
				llegada[i] = clientes[i].first;
				duracion[i] = clientes[i].second;
			}
		}                                          // O(n)
		
		// WRITE: corre la simulacion completa, llenando inicio[] y fin[].
		void ejecutar() {
			priority_queue<Evento, vector<Evento>, greater<Evento>> pq;
			for (int i = 0; i < (int)llegada.size(); i++)
			pq.push({llegada[i], LLEGADA, i});
			
			int libres = cajeros;
			queue<int> espera;                      // clientes esperando, en orden
			
			while (!pq.empty()) {
				Evento e = pq.top();
				pq.pop();                           // el reloj salta a e.t
				
				if (e.tipo == LLEGADA) {
					if (libres > 0) {               // hay cajero libre: pasa ya
						libres--;
						inicio[e.id] = e.t;
						fin[e.id] = e.t + duracion[e.id];
						pq.push({fin[e.id], FIN, e.id});   // programa su fin
					} else {
						espera.push(e.id);          // se queda en la fila
					}
				} else {                            // FIN: un cajero se libera
					libres++;
					if (!espera.empty()) {
						int c = espera.front();
						espera.pop();
						libres--;
						inicio[c] = e.t;            // empieza en este instante
						fin[c] = e.t + duracion[c];
						pq.push({fin[c], FIN, c});
					}
				}
			}
		}                                          // O(n log n)
		
		// READ: cuanto espero el cliente i antes de ser atendido.
		long long esperaDe(int i) {
			return inicio[i] - llegada[i];
		}                                          // O(1)
		
		// ------------------------------------------------------------
		// MODIFICACIONES Y COMENTARIOS CON LOS USOS
		// ------------------------------------------------------------
		
		// 1. DESEMPATE DE EVENTOS SIMULTANEOS: es la fuente numero uno de
		//    errores. Si el enunciado dice que un cliente que llega
		//    exactamente cuando otro termina NO alcanza el cajero, se
		//    invierte el criterio (LLEGADA = 0, FIN = 1). Leer con cuidado
		//    y probar con un caso donde dos eventos caen en el mismo tiempo.
		
		// 2. CANCELAR EVENTOS (cancelacion perezosa): un heap no permite
		//    borrar un elemento cualquiera. Se guarda una "version" por
		//    entidad; el evento lleva la version con la que se creo, y al
		//    sacarlo se ignora si esta desactualizado:
		//
		//    vector<int> version(n, 0);
		//    // al programar:  pq.push({t, tipo, id, version[id]});
		//    // para cancelar: version[id]++;      // los viejos quedan obsoletos
		//    // al procesar:   if (e.version != version[id]) continue;
		
		// 3. FILA CON PRIORIDAD (en vez de FIFO): se reemplaza la
		//    queue<int> por una priority_queue con el criterio pedido
		//    (menor duracion primero, cliente VIP primero, etc.):
		//
		//    priority_queue<pair<long long,int>,
		//                   vector<pair<long long,int>>,
		//                   greater<>> espera;          // {duracion, id}
		//    // al llegar sin cajero: espera.push({duracion[id], id});
		//    // al liberarse: auto [d, c] = espera.top(); espera.pop();
		
		// 4. EVENTOS PERIODICOS: al procesar un evento recurrente, se
		//    programa su proxima aparicion y se corta con un limite de
		//    tiempo para que la simulacion termine:
		//
		//    if (e.tipo == ALARMA && e.t + periodo <= limite)
		//        pq.push({e.t + periodo, ALARMA, e.id});
		
		// 5. CORTE POR TIEMPO LIMITE: "que pasa hasta el instante T":
		//
		//    if (e.t > T) break;      // el resto ocurre despues de T
		
		// 6. METRICAS QUE DEPENDEN DEL TIEMPO ENTRE EVENTOS (por ejemplo,
		//    largo promedio de la fila): se acumula el area antes de
		//    procesar cada evento, usando el tiempo transcurrido desde el
		//    evento anterior:
		//
		//    long long tAnt = 0, areaFila = 0;
		//    // al sacar un evento:
		//    areaFila += (long long)espera.size() * (e.t - tAnt);
		//    tAnt = e.t;
		//    // promedio = areaFila / tiempo total
		
		// 7. TIEMPOS REALES (double): los desempates por igualdad exacta
		//    son riesgosos con doubles. Si es posible, escalar a enteros
		//    (multiplicar por una potencia de 10); si no, comparar con
		//    epsilon en el comparador del heap.
		
		// 8. VARIOS TIPOS DE EVENTO CON DATOS EXTRA: agregar campos al
		//    struct Evento (por ejemplo, un servidor o un destino) y
		//    despachar con switch (e.tipo) en el ciclo principal. El
		//    esqueleto (sacar el minimo, actualizar el estado, programar
		//    nuevos eventos) no cambia.
		
		// 9. CUANDO NO USARLA: si los eventos ya estan todos dados de
		//    antemano y solo hay que recorrerlos ordenados, basta ordenar
		//    (ver Line Sweep). Y si el tiempo total es chico y ocurren
		//    muchos eventos por unidad de tiempo, la simulacion por pasos
		//    (Plantilla de Simulacion) es mas simple y igual de rapida.
	};
\end{lstlisting}

\textbf{Ejemplo de funcionamiento:}

\begin{lstlisting}[style=cpp]
	// 2 cajeros y 5 clientes: {llegada, duracion}
	vector<pair<long long,long long>> clientes = {
		{0, 5}, {1, 3}, {2, 4}, {3, 2}, {10, 1}
	};
	
	SimulacionEventos se(clientes, 2);
	se.ejecutar();
	
	for (int i = 0; i < 5; i++)
	cout << "cliente " << i << ": inicio=" << se.inicio[i]
	<< " fin=" << se.fin[i] << " espera=" << se.esperaDe(i) << "\n";
	// cliente 0: inicio=0  fin=5  espera=0
	// cliente 1: inicio=1  fin=4  espera=0
	// cliente 2: inicio=4  fin=8  espera=2   (espera al cliente 1)
	// cliente 3: inicio=5  fin=7  espera=2   (espera al cliente 0)
	// cliente 4: inicio=10 fin=11 espera=0
	// espera total = 4, espera maxima = 2, la ultima salida es en t = 11
\end{lstlisting}

\textbf{Nota:} el reloj nunca pasa por los instantes intermedios donde
no ocurre nada (del $8$ al $10$ en el ejemplo): salta de evento en
evento. Esa es toda la ventaja frente a la simulación por pasos, y es
también la fuente de la regla más importante de esta técnica: el
\textbf{orden} en que se procesan los eventos empatados \textbf{es} parte
del enunciado, aunque casi nunca se diga explícitamente. En el ejemplo,
si un cliente hubiera llegado en $t=4$, el orden FIN antes de LLEGADA
lo habría dejado entrar de inmediato al cajero recién liberado.

\textbf{Regla práctica:} usa simulación por eventos cuando el tiempo
simulado es enorme comparado con el número de sucesos, o cuando los
sucesos se generan unos a otros. Antes de codificar, escribe la lista de
tipos de evento, qué cambia en el estado con cada uno, qué eventos
nuevos programa, y cómo se desempatan dos eventos simultáneos; con esas
cuatro respuestas el ciclo principal casi se escribe solo.
\subsection{Autómatas Celulares (Game of Life)}

\textbf{Idea:} un autómata celular es un tablero de celdas donde cada
una tiene un estado de un conjunto finito, y en cada paso
\textbf{todas} las celdas se actualizan \textbf{a la vez} con una regla
que solo mira su vecindario. En el Juego de la Vida de Conway el
tablero es una grilla 2D, cada celda está viva o muerta, y el
vecindario son sus 8 vecinas (vecindario de Moore). La regla estándar
se escribe \textbf{B3/S23}: una celda muerta \textbf{nace} si tiene
exactamente 3 vecinas vivas, y una celda viva \textbf{sobrevive} si
tiene 2 o 3 (en cualquier otro caso muere). Todo el problema se reduce
a la plantilla de \textit{Plantilla de Simulación (Estado, Paso y
	Parada)}: el estado es la grilla completa, el paso es aplicar la regla
a todas las celdas, y la parada es un número de generaciones o un
estado repetido. La idea que no se puede omitir es la
\textbf{simultaneidad}: el nuevo valor de cada celda debe calcularse
con los valores \textbf{viejos} de sus vecinas, por eso se escribe en
una segunda grilla (\textit{double buffering}) y solo al final se
reemplaza la anterior.

\textbf{¿Por qué usarlo?} Los problemas de este tipo describen la regla
en el enunciado y piden la grilla (o la población) tras $k$
generaciones; no hay nada que optimizar mientras $k$ y el tablero sean
pequeños, y la simulación directa es la solución. Cuando $k$ es enorme
entra \textit{Simulación con Detección de Ciclos (Floyd y Estado
	Repetido)}: un tablero finito tiene $2^{F \cdot C}$ estados posibles,
así que la secuencia termina en un ciclo, y \texttt{buscarCiclo} de
abajo es exactamente la variante de Estado Repetido con el tablero
convertido a \texttt{string}. Frente a una simulación de objetos
individuales (como \textit{Simulación por Eventos}), aquí no hay
eventos sueltos: todas las celdas cambian en cada paso, por lo que la
simulación por pasos es la forma natural. Si el tablero es enorme pero
casi vacío, en cambio, conviene guardar solo las celdas vivas
(modificación 2).

\textbf{Casos de uso:}

\begin{itemize}
	\item Calcular la grilla o la población tras $k$ generaciones de
	una regla local (Game of Life y variantes con otras reglas
	B/S).
	\item Autómatas 1D (Regla 30, Regla 90, Regla 110): cada celda
	depende de sí misma y de sus dos vecinas, y la regla se codifica
	como un número de 0 a 255.
	\item Problemas de propagación sobre grillas por rondas
	(incendios, contagios, crecimiento de cultivos) que siguen el
	mismo esquema de actualización simultánea.
	\item Determinar si una configuración es estable, oscila o
	desaparece, y con qué período (ciclo del autómata).
	\item Estado tras un número gigantesco de generaciones en
	autómatas lineales sobre $GF(2)$, como la Regla 90
	(modificación 4).
\end{itemize}

\textbf{Complejidad:} cada paso recorre las $F \cdot C$ celdas y mira 8
vecinas por celda, es decir $O(F \cdot C)$; $k$ generaciones cuestan
$O(k \cdot F \cdot C)$ y la memoria es $O(F \cdot C)$. \texttt{buscarCiclo}
guarda cada tablero visto en un \texttt{map<string,int>}: cada
comparación cuesta $O(F \cdot C)$, por lo que resuelve en
$O((\mu + \lambda) \cdot F \cdot C \cdot \log(\mu + \lambda))$. En el
autómata 1D de $n$ celdas un paso cuesta $O(n)$. Con celdas vivas
dispersas ($P$ vivas) el paso baja a $O(P)$, y con \texttt{bitset}
a $O(F \cdot C / 64)$ (modificaciones 2 y 3).

\begin{lstlisting}[style=cpp]
	// ---------- Variante 1: Game of Life 2D (reglas B/S como mascaras) ----------
	struct JuegoDeLaVida {
		int filas, columnas;
		bool toroidal;                       // true: los bordes dan la vuelta
		int nacimiento, supervivencia;       // reglas B/S como mascaras de bits
		vector<vector<int>> celdas;          // 1 = viva, 0 = muerta
		
		// CREATE: lee el tablero inicial ('#' = viva, '.' = muerta). Por
		// defecto usa B3/S23; bit v de la mascara = "con v vecinas vivas".
		JuegoDeLaVida(const vector<string>& inicial, bool toroidal_ = false,
		int nacimiento_ = 1 << 3, int supervivencia_ = (1 << 2) | (1 << 3)) {
			filas = inicial.size();
			columnas = inicial[0].size();
			toroidal = toroidal_;
			nacimiento = nacimiento_;
			supervivencia = supervivencia_;
			celdas.assign(filas, vector<int>(columnas, 0));
			for (int f = 0; f < filas; f++)
			for (int c = 0; c < columnas; c++)
			celdas[f][c] = (inicial[f][c] == '#');
		}                                          // O(filas * columnas)
		
		// READ: cuantas de las 8 vecinas de (f, c) estan vivas.
		int vecinasVivas(int f, int c) const {
			int cuenta = 0;
			for (int df = -1; df <= 1; df++) {
				for (int dc = -1; dc <= 1; dc++) {
					if (df == 0 && dc == 0) continue;     // no se cuenta a si misma
					int nf = f + df, nc = c + dc;
					if (toroidal) {
						nf = (nf + filas) % filas;
						nc = (nc + columnas) % columnas;
					} else if (nf < 0 || nf >= filas || nc < 0 || nc >= columnas) {
						continue;                         // fuera del tablero = muerta
					}
					cuenta += celdas[nf][nc];
				}
			}
			return cuenta;
		}                                          // O(1) (8 vecinas)
		
		// WRITE: un paso SIMULTANEO. Se calcula todo sobre la grilla vieja y se
		// escribe en una grilla nueva; actualizar en el sitio corrompe el
		// conteo de las celdas que todavia no se han procesado.
		void paso() {
			vector<vector<int>> siguiente(filas, vector<int>(columnas, 0));
			for (int f = 0; f < filas; f++) {
				for (int c = 0; c < columnas; c++) {
					int v = vecinasVivas(f, c);
					int regla = celdas[f][c] ? supervivencia : nacimiento;
					siguiente[f][c] = (regla >> v) & 1;
				}
			}
			celdas = siguiente;
		}                                          // O(filas * columnas)
		
		// WRITE: avanza k generaciones.
		void avanzar(long long k) {
			for (long long i = 0; i < k; i++) paso();
		}                                          // O(k * filas * columnas)
		
		// READ: cantidad de celdas vivas.
		int poblacion() const {
			int total = 0;
			for (int f = 0; f < filas; f++)
			for (int c = 0; c < columnas; c++) total += celdas[f][c];
			return total;
		}                                          // O(filas * columnas)
		
		// READ: el tablero como string (sirve de clave para detectar ciclos).
		string codificar() const {
			string s;
			for (int f = 0; f < filas; f++)
			for (int c = 0; c < columnas; c++) s += celdas[f][c] ? '#' : '.';
			return s;
		}                                          // O(filas * columnas)
		
		// READ: imprime el tablero.
		void mostrar() const {
			for (int f = 0; f < filas; f++) {
				for (int c = 0; c < columnas; c++) cout << (celdas[f][c] ? '#' : '.');
				cout << "\n";
			}
		}                                          // O(filas * columnas)
		
		// READ: busca el ciclo (Estado Repetido) sobre una COPIA del tablero.
		// Retorna {mu, lambda}, o {-1, -1} si no aparece en 'limite' pasos.
		pair<int,int> buscarCiclo(int limite) const {
			JuegoDeLaVida copia = *this;
			map<string, int> visto;
			for (int t = 0; t <= limite; t++) {
				string clave = copia.codificar();
				auto it = visto.find(clave);
				if (it != visto.end()) return {it->second, t - it->second};
				visto[clave] = t;
				copia.paso();
			}
			return {-1, -1};
		}                                          // O(limite * F * C * log)
	};
	
	// ---------- Variante 2: automata elemental 1D (regla de Wolfram) ----------
	struct AutomataElemental {
		int n, regla;                        // regla en [0, 255]
		vector<int> celdas;                  // fila ciclica de n celdas (0 o 1)
		
		// CREATE: fila inicial y numero de regla.
		AutomataElemental(vector<int> inicial, int regla_) {
			celdas = inicial;
			n = celdas.size();
			regla = regla_;
		}                                          // O(n)
		
		// WRITE: un paso. El vecindario (izq, centro, der) forma un numero de
		// 3 bits, y ese bit de la regla es el nuevo valor de la celda.
		void paso() {
			vector<int> sig(n);
			for (int i = 0; i < n; i++) {
				int izq = celdas[(i - 1 + n) % n];
				int cen = celdas[i];
				int der = celdas[(i + 1) % n];
				int patron = izq * 4 + cen * 2 + der;     // 0..7
				sig[i] = (regla >> patron) & 1;
			}
			celdas = sig;
		}                                          // O(n)
		
		// WRITE: SOLO regla 90. Avanza 2^m pasos de golpe (ver modificacion 4).
		void saltarPotenciaDeDos(int m) {
			int d = (int)((1LL << m) % n);
			vector<int> sig(n);
			for (int i = 0; i < n; i++)
			sig[i] = celdas[((i - d) % n + n) % n] ^ celdas[(i + d) % n];
			celdas = sig;
		}                                          // O(n)
		
		// WRITE: SOLO regla 90. Avanza k pasos descomponiendo k en binario.
		void saltar(long long k) {
			for (int m = 0; m < 62; m++)
			if ((k >> m) & 1) saltarPotenciaDeDos(m);
		}                                          // O(n log k)
	};
	
	// ------------------------------------------------------------
	// MODIFICACIONES Y COMENTARIOS CON LOS USOS
	// ------------------------------------------------------------
	
	// 1. OTRAS REGLAS: cualquier regla "tipo Life" B/S se pasa como mascaras.
	//    HighLife es B36/S23:
	//
	//    JuegoDeLaVida hl(tablero, false, (1 << 3) | (1 << 6), (1 << 2) | (1 << 3));
	//
	//    Si el enunciado da el vecindario de von Neumann (solo 4 vecinas:
	//    arriba, abajo, izquierda, derecha), basta cambiar los dos bucles de
	//    vecinasVivas por los cuatro desplazamientos (df, dc) permitidos.
	
	// 2. TABLERO INFINITO CON POCAS CELDAS VIVAS: en vez de una grilla, se
	//    guarda el conjunto de celdas vivas y se cuenta, para cada celda
	//    vecina de alguna viva, cuantas vecinas vivas tiene:
	//
	//    set<pair<int,int>> vivas;
	//    map<pair<int,int>, int> conteo;
	//    for (auto [f, c] : vivas)
	//        for (int df = -1; df <= 1; df++)
	//            for (int dc = -1; dc <= 1; dc++)
	//                if (df != 0 || dc != 0) conteo[{f + df, c + dc}]++;
	//    set<pair<int,int>> nuevas;
	//    for (auto& [pos, v] : conteo) {
		//        bool viva = vivas.count(pos);
		//        if ((viva && (v == 2 || v == 3)) || (!viva && v == 3)) nuevas.insert(pos);
		//    }
	//    vivas = nuevas;                     // O(P log P) por paso
	//
	//    Solo las celdas vecinas de una viva pueden estar vivas en el
	//    siguiente paso, por eso el costo depende de P y no del tablero.
	
	// 3. BITSET: si cada fila cabe en un bitset<W>, las 8 vecinas de TODA la
	//    fila se obtienen con desplazamientos (fila << 1, fila >> 1, y lo
	//    mismo para la fila de arriba y la de abajo), y el conteo se hace
	//    con sumadores bit a bit (dos o tres bitsets que guardan los bits del
	//    contador). Un paso pasa de O(F * C) a O(F * C / 64). Ver la
	//    subseccion de Bitset para Optimizacion en la parte BitWise.
	
	// 4. SALTAR 2^m PASOS EN REGLA 90: la regla 90 es nuevo[i] = izq XOR der,
	//    lineal sobre GF(2). Elevar al cuadrado en caracteristica 2 no crea
	//    terminos cruzados, asi que 2^m pasos equivalen a
	//    nuevo[i] = viejo[i - 2^m] XOR viejo[i + 2^m] (indices ciclicos).
	//    Con k en binario, k pasos cuestan O(n log k) en vez de O(n k), y
	//    permite k = 10^18 sin ciclos ni memoria extra. Solo funciona para
	//    reglas lineales (90 y pocas mas), NO para Regla 30 ni Game of Life.
	
	// 5. ESTADO TRAS k GRANDE (Game of Life): se llama a buscarCiclo y se
	//    reduce k igual que en Detection de Ciclos:
	//
	//    auto [mu, lambda] = j.buscarCiclo(limite);
	//    if (k >= mu) k = mu + (k - mu) % lambda;
	//    j.avanzar(k);
	//
	//    Formas conocidas para reconocerlas: bloque (estable), parpadeador
	//    (periodo 2), planeador (avanza una celda en diagonal cada 4 pasos;
	//    en un toro F x C, si es lo bastante grande, regresa tras
	//    4 * mcm(F, C) pasos).
	
	// 6. TRAMPA: ACTUALIZAR EN EL SITIO. Si paso() modifica 'celdas' mientras
	//    recorre, las celdas siguientes ven vecinas ya actualizadas y el
	//    resultado es otro autómata. Siempre grilla nueva (o guardar solo el
	//    cambio y aplicarlo al final).
	
	// 7. TRAMPA: BORDES. "Fuera del tablero esta muerto", "los bordes dan la
	//    vuelta" y "los bordes se reflejan" dan resultados distintos, y el
	//    enunciado casi siempre lo dice en una sola linea. En un toro con
	//    filas o columnas < 3, una misma celda puede contarse dos veces como
	//    vecina (por eso conviene confirmar el tamano minimo del enunciado).
\end{lstlisting}

\textbf{Ejemplo de funcionamiento:}

\begin{lstlisting}[style=cpp]
	// Parpadeador (oscilador de periodo 2), tablero de 5x5 con bordes muertos.
	JuegoDeLaVida parpadeador({".....",
		"..#..",
		"..#..",
		"..#..",
		"....."});
	
	cout << parpadeador.poblacion();      // 3
	parpadeador.paso();
	parpadeador.mostrar();
	// .....
	// .....
	// .###.
	// .....
	// .....
	parpadeador.paso();                    // vuelve a la forma vertical
	auto c1 = parpadeador.buscarCiclo(50);
	cout << c1.first << " " << c1.second;  // 0 2
	
	// Planeador (glider) en un toro de 6x6: se mueve una celda en diagonal
	// cada 4 pasos y regresa a la posicion inicial tras 4 * 6 = 24.
	JuegoDeLaVida planeador({".#....",
		"..#...",
		"###...",
		"......",
		"......",
		"......"}, true);
	
	auto c2 = planeador.buscarCiclo(200);
	cout << c2.first << " " << c2.second;  // 0 24
	planeador.avanzar(4);
	planeador.mostrar();
	// ......
	// ..#...
	// ...#..
	// .###..
	// ......
	// ......
	
	// El mismo planeador con bordes muertos choca con la pared y termina en
	// un bloque estable de 4 celdas (ciclo de largo 1 desde el paso 15).
	JuegoDeLaVida choque({".#....", "..#...", "###...", "......", "......", "......"});
	auto c3 = choque.buscarCiclo(200);
	cout << c3.first << " " << c3.second;  // 15 1
	choque.avanzar(20);
	cout << choque.poblacion();            // 4
	
	// Regla 90 sobre 9 celdas cerradas en circulo, una sola celda viva:
	vector<int> ini(9, 0);
	ini[4] = 1;
	AutomataElemental a(ini, 90);
	// generacion 0:  ....#....
	// generacion 1:  ...#.#...
	// generacion 2:  ..#...#..
	// generacion 3:  .#.#.#.#.
	// generacion 4:  #.......#
	
	AutomataElemental b(ini, 90);
	b.saltar(4);                           // da la misma fila que 4 pasos: #.......#
	AutomataElemental c(ini, 90);
	c.saltar(1000000000000LL);             // 10^12 pasos sin simularlos: ...#.#...
\end{lstlisting}

\textbf{Nota:} en el ejemplo del planeador sobre el toro, el ciclo empieza en
$\mu = 0$ porque el planeador \textbf{regresa} exactamente al tablero
inicial; en el ejemplo del choque, $\mu = 15$ porque el estado inicial
nunca se repite (el planeador se deshace en un bloque y de ahí en adelante
el tablero ya no cambia, $\lambda = 1$). Por eso no basta con mirar solo la
población: dos tableros distintos pueden tener el mismo número de celdas
vivas, y el ciclo se detecta sobre el tablero \textbf{completo}.

\textbf{Regla práctica:} si el problema da una regla local que se aplica a
todas las celdas a la vez, simula con dos grillas y copia solo al final de
cada paso. Si $k$ es pequeño, eso es todo. Si $k$ es gigantesco, aplica
\textit{Simulación con Detección de Ciclos} sobre el tablero codificado; si
el tablero es enorme pero casi vacío, guarda solo las celdas vivas; y si la
regla es lineal (Regla 90), salta de a $2^m$ pasos. Antes de codificar,
confirma tres cosas del enunciado: qué cuenta como vecina, qué pasa en los
bordes, y si la actualización es simultánea o por orden de celda.

\subsection{Simulación de Máquinas e Intérpretes de Instrucciones}

\textbf{Idea:} muchos problemas entregan un \textbf{programa} escrito en un
lenguaje inventado (un puñado de instrucciones sobre registros, una cinta,
una pila) y piden lo que imprime o el estado final. Es la plantilla de
\textit{Plantilla de Simulación (Estado, Paso y Parada)} con nombres
propios: el \textbf{estado} es el contador de programa \texttt{pc} (índice
de la próxima instrucción) más los registros, la memoria y la salida; el
\textbf{paso} es el ciclo \textit{buscar-decodificar-ejecutar} (leer la
instrucción en \texttt{pc}, ejecutarla, y mover \texttt{pc} a la
siguiente o a donde diga un salto); y la \textbf{parada} es una
instrucción de detención, que \texttt{pc} salga del programa, o un
límite de pasos. La decisión de diseño que más importa es
\textbf{parsear el programa una sola vez} en una estructura
(\texttt{operación}, \texttt{operandos}) y ejecutar sobre esa
estructura, en vez de volver a partir strings en cada paso: un programa
corto puede ejecutarse millones de veces por sus bucles.

\textbf{¿Por qué usarlo?} Aquí no se puede saltar la simulación con una
fórmula: el comportamiento lo define el texto del programa y hay que
ejecutarlo. Se distingue de la simulación de un proceso físico
(\textit{Simulación por Eventos}) en que el tiempo avanza instrucción a
instrucción y el orden lo fija el programa, no una cola de prioridad. El
gran riesgo es el \textbf{bucle infinito}: un programa válido puede no
terminar nunca. Las dos defensas son un límite de pasos (simple y casi
siempre suficiente) y \textit{Simulación con Detección de Ciclos}:
la máquina tiene estado finito, así que si el par (\texttt{pc},
registros) se repite y no hay entrada nueva, el programa está atrapado.
Cuando el programa es larguísimo pero tiene bucles de suma/multiplicación
reconocibles, se \textbf{reemplaza el bucle por su resultado} (modificación
4), que es la única forma de pasar de $10^{9}$ instrucciones a unas pocas.

\textbf{Casos de uso:}

\begin{itemize}
	\item Máquinas de registros con \texttt{set}, \texttt{add},
	\texttt{jnz}, etc., y preguntas del tipo \textit{qué queda en el
		registro \texttt{a} al final} o \textit{qué imprime}.
	\item Intérpretes de esolangs como Brainfuck: cinta, puntero y
	corchetes que emparejar (una \textit{pila} resuelve los saltos).
	\item Simular un \textbf{robot} o vehículo que sigue una lista de
	comandos (avanzar, girar, repetir) hasta terminarla.
	\item Detectar si un programa termina, o encontrar un valor
	inicial que lo haga terminar (con límite de pasos o estado repetido).
	\item Máquinas con pila (\texttt{push}, \texttt{pop}, \texttt{call},
	\texttt{ret}) o memoria indexada, que extienden el mismo esquema
	(modificaciones 2 y 3).
\end{itemize}

\textbf{Complejidad:} cada paso cuesta $O(1)$ (un \texttt{switch} y accesos
a un arreglo de registros); el parseo cuesta $O(L)$ para un programa de $L$
líneas y la memoria es $O(L + \text{registros} + \text{salida})$. El tiempo
total es $O(L + P)$, donde $P$ es el \textbf{número de pasos que
	realmente se ejecutan}, y $P$ depende del comportamiento del programa (sus
bucles), no de su longitud: un programa de 7 líneas puede ejecutar
$10^{9}$ pasos. En Brainfuck, precalcular los saltos deja el paso en
$O(1)$ en lugar de $O(L)$ por corchete.

\begin{lstlisting}[style=cpp]
	// ---------- Variante 1: maquina de registros (26 registros: a..z) ----------
	// Instrucciones: set r x | add r x | sub r x | mul r x | mod r x
	//                jnz x y (si x != 0, salta y instrucciones RELATIVO a esta)
	//                out x   | hlt
	// Un operando es un registro (una letra) o un numero (posiblemente negativo).
	struct MaquinaVirtual {
		enum Op { SET, ADD, SUB, MUL, MOD, JNZ, OUT, HLT };
		struct Instruccion {
			Op op;
			int destino;                     // registro destino (SET/ADD/...)
			long long a, b;                  // operandos: indice de registro o valor
			bool aEsReg, bEsReg;
		};
		
		vector<Instruccion> programa;
		long long reg[26];
		int pc;                              // indice de la proxima instruccion
		long long pasos;                     // instrucciones ejecutadas
		bool detenida;
		vector<long long> salida;
		
		// Convierte un token en (valor, esRegistro).
		static void leerOperando(const string& t, long long& valor, bool& esReg) {
			if (isalpha(t[0])) { esReg = true; valor = t[0] - 'a'; }
			else { esReg = false; valor = stoll(t); }
		}
		
		// CREATE: parsea el programa UNA sola vez a instrucciones; los
		// registros parten en 0.
		MaquinaVirtual(const vector<string>& lineas) {
			for (int i = 0; i < 26; i++) reg[i] = 0;
			pc = 0; pasos = 0; detenida = false;
			for (const string& linea : lineas) {
				istringstream in(linea);
				string nombre, x, y;
				in >> nombre >> x >> y;
				Instruccion I = {HLT, 0, 0, 0, false, false};
				if (nombre == "set" || nombre == "add" || nombre == "sub" ||
				nombre == "mul" || nombre == "mod") {
					if (nombre == "set") I.op = SET;
					else if (nombre == "add") I.op = ADD;
					else if (nombre == "sub") I.op = SUB;
					else if (nombre == "mul") I.op = MUL;
					else I.op = MOD;
					I.destino = x[0] - 'a';
					leerOperando(y, I.a, I.aEsReg);
				} else if (nombre == "jnz") {
					I.op = JNZ;
					leerOperando(x, I.a, I.aEsReg);
					leerOperando(y, I.b, I.bEsReg);
				} else if (nombre == "out") {
					I.op = OUT;
					leerOperando(x, I.a, I.aEsReg);
				} else {
					I.op = HLT;
				}
				programa.push_back(I);
			}
			if (programa.empty()) detenida = true;
		}                                          // O(L)
		
		// READ: valor de un operando (registro o constante).
		long long valor(long long x, bool esReg) const {
			return esReg ? reg[x] : x;
		}                                          // O(1)
		
		// WRITE: ejecuta UNA instruccion y mueve pc.
		void paso() {
			const Instruccion& I = programa[pc];
			int siguiente = pc + 1;              // por defecto, la siguiente
			switch (I.op) {
				case SET: reg[I.destino] = valor(I.a, I.aEsReg); break;
				case ADD: reg[I.destino] += valor(I.a, I.aEsReg); break;
				case SUB: reg[I.destino] -= valor(I.a, I.aEsReg); break;
				case MUL: reg[I.destino] *= valor(I.a, I.aEsReg); break;
				case MOD: {
					long long d = valor(I.a, I.aEsReg);
					if (d == 0) { detenida = true; return; }   // division por cero
					reg[I.destino] %= d;
					break;
				}
				case JNZ:
				if (valor(I.a, I.aEsReg) != 0)
				siguiente = pc + (int)valor(I.b, I.bEsReg);
				break;
				case OUT: salida.push_back(valor(I.a, I.aEsReg)); break;
				case HLT: detenida = true; return;
			}
			pc = siguiente;
			pasos++;
			if (pc < 0 || pc >= (int)programa.size()) detenida = true;
		}                                          // O(1)
		
		// WRITE: ejecuta hasta detenerse o hasta 'limite' pasos. Retorna true
		// si el programa termino, false si se corto por el limite.
		bool ejecutar(long long limite) {
			while (!detenida && pasos < limite) paso();
			return detenida;
		}                                          // O(min(P, limite))
	};
	
	// ---------- Variante 2: interprete de Brainfuck (saltos precalculados) ----------
	struct InterpreteBrainfuck {
		string codigo;
		vector<int> salto;                   // salto[i] = corchete que empareja al i
		vector<unsigned char> cinta;         // celdas de 8 bits: dan la vuelta en 256
		int ptr, pc;
		string salida;
		
		// CREATE: empareja los corchetes con una pila, UNA vez, para no
		// recorrer el codigo buscando el corchete en cada iteracion.
		InterpreteBrainfuck(const string& codigo_, int tamCinta = 30000) {
			codigo = codigo_;
			cinta.assign(tamCinta, 0);
			ptr = 0; pc = 0;
			salto.assign(codigo.size(), -1);
			vector<int> pila;
			for (int i = 0; i < (int)codigo.size(); i++) {
				if (codigo[i] == '[') pila.push_back(i);
				else if (codigo[i] == ']') {
					int j = pila.back(); pila.pop_back();
					salto[i] = j;
					salto[j] = i;
				}
			}
		}                                          // O(L)
		
		// WRITE: ejecuta todo el programa; 'entrada' alimenta a la coma (,).
		void ejecutar(const string& entrada = "") {
			int pos = 0;
			while (pc < (int)codigo.size()) {
				switch (codigo[pc]) {
					case '>': ptr++; break;
					case '<': ptr--; break;
					case '+': cinta[ptr]++; break;
					case '-': cinta[ptr]--; break;
					case '.': salida += (char)cinta[ptr]; break;
					case ',': cinta[ptr] = pos < (int)entrada.size() ? entrada[pos++] : 0; break;
					case '[': if (cinta[ptr] == 0) pc = salto[pc]; break;
					case ']': if (cinta[ptr] != 0) pc = salto[pc]; break;
				}
				pc++;
			}
		}                                          // O(P), P = pasos ejecutados
	};
	
	// ------------------------------------------------------------
	// MODIFICACIONES Y COMENTARIOS CON LOS USOS
	// ------------------------------------------------------------
	
	// 1. DETECTAR BUCLE INFINITO POR ESTADO REPETIDO: si el programa no lee
	//    entrada, el estado completo es (pc, registros). Si se repite, el
	//    programa no termina nunca (mismo principio de Estado Repetido):
	//
	//    set<pair<int, array<long long, 26>>> visto;
	//    while (!m.detenida) {
		//        array<long long, 26> r;
		//        copy(m.reg, m.reg + 26, r.begin());
		//        if (!visto.insert({m.pc, r}).second) return false;   // bucle infinito
		//        m.paso();
		//    }
	//
	//    Guarda muchos estados si el programa corre muchos pasos; para un
	//    limite de pasos generoso (10^7 o 10^8 pasos simples) suele bastar
	//    con el contador y no hace falta este set.
	
	// 2. ETIQUETAS Y SALTOS ABSOLUTOS: si el lenguaje tiene "label:" y
	//    "jmp label", se resuelve en dos pasadas al parsear: la primera
	//    guarda etiqueta -> indice de instruccion en un map, la segunda
	//    reemplaza cada nombre por ese indice. En ejecucion todo salto es
	//    un simple "pc = destino".
	
	// 3. PILA, CALL Y RET: se agrega vector<int> pila. 'call x' hace
	//    pila.push_back(pc + 1) y salta a x; 'ret' hace pc = pila.back() y
	//    pila.pop_back(). Con memoria indexada, un operando puede ser
	//    "mem[valor]" con mem un arreglo de tamano fijo.
	
	// 4. OPTIMIZAR BUCLES CONOCIDOS: el bucle
	//        add a 1
	//        sub b 1
	//        jnz b -2
	//    equivale a "a += b; b = 0" si b > 0, en O(1) en vez de O(b). En
	//    Brainfuck, "[-]" equivale a cinta[ptr] = 0, "[->+<]" a
	//    cinta[ptr+1] += cinta[ptr] (modulo 256), y una tanda "+++++" puede
	//    comprimirse a "sumar 5". Se detectan al parsear, sustituyendo por
	//    una instruccion nueva. Solo vale la pena si el enunciado ejecuta
	//    ~10^9 pasos o mas; si no, ejecutar directo es mas seguro.
	
	// 5. TRAMPA: SALTO RELATIVO. En "jnz x y" el salto se cuenta desde la
	//    PROPIA instruccion jnz (pc + y), no desde la siguiente. Confundirlo
	//    da un error de uno, muy dificil de ver sin traza. Un salto de 0
	//    queda en la misma instruccion para siempre.
	
	// 6. TRAMPA: LEER ANTES DE ESCRIBIR. "add a a" duplica a, y funciona
	//    porque valor() se evalua antes de asignar. Si se escribe la
	//    asignacion primero y se lee despues, el resultado cambia. Igual
	//    con "mul a a" o "sub a a".
	
	// 7. TRAMPA: TIPOS. Los registros deben ser long long (o el tipo que
	//    diga el enunciado): un producto de dos int desborda en silencio.
	//    En Brainfuck, unsigned char da la vuelta 255 -> 0, y el puntero
	//    puede salirse de la cinta si el programa lo permite (validar
	//    segun el enunciado: error, cinta circular o cinta infinita).
	
	// 8. DEPURAR: imprimir "pc, a, b, ..." de los primeros 20 pasos y
	//    compararlos a mano con el ejemplo del enunciado casi siempre
	//    localiza el error en un instante.
\end{lstlisting}

\textbf{Ejemplo de funcionamiento:}

\begin{lstlisting}[style=cpp]
	// Suma 10 + 9 + ... + 1: a acumula, b cuenta hacia atras.
	MaquinaVirtual m({"set a 0",
		"set b 10",
		"add a b",
		"sub b 1",
		"jnz b -2",       // si b != 0, retrocede 2 (vuelve a "add a b")
		"out a",
		"hlt"});
	
	bool termino = m.ejecutar(1000);
	cout << termino << " " << m.pasos << " " << m.salida[0];
	// 1 33 55   (termino; 33 pasos = 2 + 10 vueltas de 3 + out; imprimio 55)
	
	// Factorial de 5 (mul en vez de add): sin "hlt", pc sale del programa.
	MaquinaVirtual fact({"set a 1", "set b 5", "mul a b", "sub b 1", "jnz b -2", "out a"});
	fact.ejecutar(1000);
	cout << fact.salida[0] << " " << fact.pasos;
	// 120 18   (se detiene porque pc = 6 queda fuera del programa)
	
	// Bucle infinito: "jnz a 0" salta a si misma para siempre.
	MaquinaVirtual inf({"set a 1", "jnz a 0"});
	bool t2 = inf.ejecutar(1000000);
	cout << t2 << " " << inf.pasos << " " << inf.pc;
	// 0 1000000 1   (no termino: el limite de pasos lo corto)
	
	// Brainfuck: 8 * 8 = 64 en la celda 1, +1 = 65, y '.' imprime 'A'.
	InterpreteBrainfuck bf("++++++++[>++++++++<-]>+.");
	bf.ejecutar();
	cout << bf.salida;   // A
	
	// Eco de la entrada: lee un caracter, lo imprime, y repite hasta el 0.
	InterpreteBrainfuck eco(",[.,]");
	eco.ejecutar("hey");
	cout << eco.salida;  // hey
\end{lstlisting}

\textbf{Nota:} en el primer ejemplo el programa termina con \texttt{hlt}
y \texttt{pc = 6} (apuntando al propio \texttt{hlt}), mientras que en el
del factorial no hay \texttt{hlt}: la máquina se detiene porque
\texttt{pc} sale del rango del programa. Los enunciados suelen definir
ambas formas de terminar, y si tu simulador solo maneja una, fallas los
casos donde ocurre la otra. Por otro lado, \texttt{pasos} cuenta las
instrucciones \textbf{ejecutadas}, no las líneas del programa; ese número
es el que determina si la simulación directa cabe en el tiempo límite.

\textbf{Regla práctica:} parsea el programa una sola vez a una estructura
de instrucciones, ejecuta con un ciclo \texttt{buscar-ejecutar} y pon
\textbf{siempre} un límite de pasos, porque los programas que no terminan
son el caso de prueba favorito. Si el límite se alcanza en un enunciado que
pregunta \textit{si termina}, usa \textit{Simulación con Detección de
	Ciclos} sobre el estado (\texttt{pc}, registros); si el programa ejecuta
más de $\sim 10^{9}$ instrucciones por bucles reconocibles, sustitúyelos por
su resultado. Y recuerda dos detalles que causan la mayoría de los errores:
el salto es relativo a la instrucción que salta, y los operandos se leen
antes de escribir el destino.

\section{Matrices y Grillas}
\subsection{Recorridos Especiales (Espiral, Diagonal, Zigzag)}

\textbf{Idea:} un recorrido especial es una forma de visitar cada celda
de una matriz \textbf{exactamente una vez} en un orden distinto al de
filas. Hay cuatro familias. La \textbf{espiral} recorre la matriz por
capas, de afuera hacia adentro. Las \textbf{diagonales} agrupan las
celdas con $i+j$ constante (antidiagonales) o con $i-j$ constante
(diagonales principales). El \textbf{zigzag} recorre las antidiagonales
alternando el sentido en cada una. La \textbf{serpiente} recorre las
filas alternando izquierda-derecha y derecha-izquierda. La idea común es
que el orden se describe con \textbf{aritmética de índices}, no con un
caminante que marca celdas visitadas: la espiral se reduce a cuatro
fronteras (\texttt{arriba}, \texttt{abajo}, \texttt{izq}, \texttt{der})
que se contraen después de recorrer cada lado, y una antidiagonal $d$ es
el conjunto de celdas $(i, d-i)$ con
$\max(0,\, d-m+1) \le i \le \min(n-1,\, d)$. Conviene que el recorrido
devuelva la \textbf{lista de coordenadas} en orden
(\texttt{ordenEspiral}): así sirve tanto para \textbf{leer} una matriz en
ese orden como para \textbf{llenar} una matriz siguiéndolo
(\texttt{generarEspiral}).

\textbf{¿Por qué usarlo?} Frente a simular un caminante con una matriz
\texttt{visitado} y un arreglo \texttt{dx/dy} (que gira a la derecha al
chocar; ver \textit{Direcciones y Vecindarios} más adelante), las
fronteras no necesitan memoria extra ni marcar celdas, y no dependen de
que el borde esté bien detectado; el caminante con \texttt{visitado}, en
cambio, es más flexible y es la opción cuando hay celdas bloqueadas o la
región no es un rectángulo (modificación 1). Frente a recorrer por
filas, las diagonales agrupan celdas relacionadas: dos casillas de un
tablero están en la misma diagonal de un alfil o de una reina si y solo
si comparten $i+j$ o $i-j$, la misma condición que usa \textit{Problema
	de las N Reinas}. Y frente a \textit{Prefix Sum 2D}, que responde sumas
de submatrices, los recorridos por diagonal responden preguntas sobre
\textbf{líneas} de la matriz en $O(n \cdot m)$ total.

\textbf{Casos de uso:}

\begin{itemize}
	\item Imprimir o leer una matriz en espiral, y \textbf{generar} una
	matriz $n \times m$ llena con $1, 2, \dots, nm$ en espiral.
	\item Agrupar por diagonales para sumar, contar o buscar el
	máximo en cada una, o para saber qué celdas ataca un alfil o una
	reina.
	\item Recorrer en zigzag (\textit{Diagonal Traverse}), el orden en
	que se lee un bloque de coeficientes en la compresión JPEG.
	\item Tableros numerados en serpiente, como el de \textit{Serpientes
		y Escaleras}: convertir el número de casilla en (fila, columna) y
	viceversa.
	\item Problemas de \textbf{capas} o anillos: rotar cada anillo,
	sumar cada anillo, o recorrer desde el centro hacia afuera
	(modificación 2).
\end{itemize}

\textbf{Complejidad:} todos los recorridos visitan cada celda una sola
vez con trabajo $O(1)$ por celda, así que cuestan $O(n \cdot m)$ en
tiempo. En \texttt{ordenEspiral}, cada uno de los cuatro bucles agrega
las celdas de un lado y luego la frontera correspondiente se contrae, por
lo que ninguna celda se agrega dos veces (las dos condiciones \texttt{if}
lo garantizan cuando queda una sola fila o columna) y el total de celdas
agregadas es $n \cdot m$. Hay $\lceil \min(n,m)/2 \rceil$ capas. En las
diagonales hay $n+m-1$ antidiagonales, y la $d$-ésima tiene
$\min(n-1,d)-\max(0,d-m+1)+1$ celdas. La memoria extra es $O(n \cdot m)$
por la lista de salida; el recorrido en sí no usa nada más.

\begin{lstlisting}[style=cpp]
	struct RecorridosDeMatriz {
		int filas, columnas;
		vector<vector<int>> a;
		
		// CREATE: guarda la matriz y sus dimensiones.
		RecorridosDeMatriz(const vector<vector<int>>& m) {
			a = m;
			filas = a.size();
			columnas = a.empty() ? 0 : a[0].size();
		}                                          // O(filas * columnas)
		
		// READ: coordenadas de una matriz f x c en orden de espiral horaria,
		// empezando en (0, 0). Cuatro fronteras que se contraen.
		static vector<pair<int,int>> ordenEspiral(int f, int c) {
			vector<pair<int,int>> orden;
			int arriba = 0, abajo = f - 1, izq = 0, der = c - 1;
			while (arriba <= abajo && izq <= der) {
				for (int j = izq; j <= der; j++) orden.push_back({arriba, j});   // ->
				arriba++;
				for (int i = arriba; i <= abajo; i++) orden.push_back({i, der}); // baja
				der--;
				if (arriba <= abajo) {                    // sigue habiendo fila
					for (int j = der; j >= izq; j--) orden.push_back({abajo, j}); // <-
					abajo--;
				}
				if (izq <= der) {                         // sigue habiendo columna
					for (int i = abajo; i >= arriba; i--) orden.push_back({i, izq}); // sube
					izq++;
				}
			}
			return orden;
		}                                          // O(f * c)
		
		// READ: los valores de la matriz leidos en espiral.
		vector<int> espiral() const {
			vector<int> res;
			for (auto [i, j] : ordenEspiral(filas, columnas)) res.push_back(a[i][j]);
			return res;
		}                                          // O(filas * columnas)
		
		// READ: matriz f x c llena con 1..f*c siguiendo la espiral.
		static vector<vector<int>> generarEspiral(int f, int c) {
			vector<vector<int>> m(f, vector<int>(c, 0));
			int k = 1;
			for (auto [i, j] : ordenEspiral(f, c)) m[i][j] = k++;
			return m;
		}                                          // O(f * c)
		
		// READ: cada antidiagonal (i + j = d) como un vector, de arriba hacia
		// abajo (i creciente, es decir, hacia abajo y a la izquierda).
		vector<vector<int>> antiDiagonales() const {
			vector<vector<int>> res;
			for (int d = 0; d <= filas + columnas - 2; d++) {
				vector<int> diag;
				for (int i = max(0, d - columnas + 1); i <= min(filas - 1, d); i++)
				diag.push_back(a[i][d - i]);
				res.push_back(diag);
			}
			return res;
		}                                          // O(filas * columnas)
		
		// READ: zigzag: antidiagonales alternando el sentido. La d par sube
		// (i decreciente); la d impar baja (i creciente). Empieza en (0, 0)
		// y sigue con (0, 1).
		vector<int> zigzag() const {
			vector<int> res;
			for (int d = 0; d <= filas + columnas - 2; d++) {
				int desde = max(0, d - columnas + 1), hasta = min(filas - 1, d);
				if (d % 2 == 0) {
					for (int i = hasta; i >= desde; i--) res.push_back(a[i][d - i]);
				} else {
					for (int i = desde; i <= hasta; i++) res.push_back(a[i][d - i]);
				}
			}
			return res;
		}                                          // O(filas * columnas)
		
		// READ: serpiente (boustrofedon): las filas pares de izquierda a
		// derecha y las impares de derecha a izquierda.
		vector<int> serpiente() const {
			vector<int> res;
			for (int i = 0; i < filas; i++) {
				if (i % 2 == 0) for (int j = 0; j < columnas; j++) res.push_back(a[i][j]);
				else for (int j = columnas - 1; j >= 0; j--) res.push_back(a[i][j]);
			}
			return res;
		}                                          // O(filas * columnas)
		
		// READ: (fila, columna) de la casilla k (0-indexada) de un tablero de
		// 'columnas' columnas numerado en serpiente. Sirve para Serpientes y
		// Escaleras sin construir el tablero.
		static pair<int,int> posicionSerpiente(int k, int columnas) {
			int i = k / columnas, r = k % columnas;
			int j = (i % 2 == 0) ? r : columnas - 1 - r;
			return {i, j};
		}                                          // O(1)
	};
	
	// ------------------------------------------------------------
	// MODIFICACIONES Y COMENTARIOS CON LOS USOS
	// ------------------------------------------------------------
	
	// 1. ESPIRAL CON CAMINANTE Y VISITADO: mas flexible que las fronteras
	//    cuando hay celdas bloqueadas o la region es irregular. Avanza en
	//    la direccion actual y gira a la derecha al chocar con el borde o
	//    con una celda ya visitada:
	//
	//    int dx[4] = {0, 1, 0, -1}, dy[4] = {1, 0, -1, 0};   // -> v <- ^
	//    vector<vector<bool>> vis(f, vector<bool>(c, false));
	//    int i = 0, j = 0, d = 0;
	//    for (int t = 0; t < f * c; t++) {
		//        // procesar (i, j)
		//        vis[i][j] = true;
		//        int ni = i + dx[d], nj = j + dy[d];
		//        if (ni < 0 || ni >= f || nj < 0 || nj >= c || vis[ni][nj]) {
			//            d = (d + 1) % 4;                            // gira a la derecha
			//            ni = i + dx[d]; nj = j + dy[d];
			//        }
		//        i = ni; j = nj;
		//    }
	//
	//    Con celdas bloqueadas hay que cambiar el for por un while que pare
	//    cuando ni la direccion actual ni el giro dejen avanzar.
	
	// 2. OTRAS ORIENTACIONES:
	//    - Antihoraria desde (0, 0): es la espiral horaria de la matriz
	//      TRANSPUESTA con las coordenadas intercambiadas:
	//        for (auto [i, j] : ordenEspiral(columnas, filas)) visitar(j, i);
	//    - Hacia afuera, terminando en la esquina (0, 0): recorrer
	//      ordenEspiral al reves. Parte de la ultima celda de la espiral
	//      (en una matriz impar x impar, el centro) y da la vuelta en
	//      sentido contrario:
	//        auto o = ordenEspiral(f, c); reverse(o.begin(), o.end());
	//    - Desde otra esquina: reflejar las coordenadas, por ejemplo
	//      (i, j) -> (f - 1 - i, j) para empezar abajo a la izquierda.
	
	// 3. DIAGONALES PRINCIPALES (i - j constante): i - j va de -(m-1) a n-1;
	//    se suma m - 1 para usarlo como indice no negativo:
	//
	//    vector<vector<int>> diag(n + m - 1);
	//    for (int i = 0; i < n; i++)
	//        for (int j = 0; j < m; j++) diag[i - j + m - 1].push_back(a[i][j]);
	//
	//    La antidiagonal analoga usa el indice i + j (de 0 a n + m - 2).
	
	// 4. ACUMULADOS POR DIAGONAL: igual que Prefix Sums pero a lo largo de
	//    la diagonal: pd[i][j] = a[i][j] + pd[i-1][j-1] (diagonal principal)
	//    o pd[i][j] = a[i][j] + pd[i-1][j+1] (antidiagonal), con 0 fuera del
	//    tablero. La suma de un tramo de diagonal se responde en O(1) como
	//    diferencia de dos acumulados.
	
	// 5. NUMERO DE UNA CASILLA EN SERPIENTE (inversa de posicionSerpiente):
	//    k = i * columnas + (i % 2 == 0 ? j : columnas - 1 - j).
	
	// 6. K-ESIMA CELDA DE LA ESPIRAL SIN CONSTRUIRLA: la capa r (contando
	//    desde 0) es un anillo de (f - 2r) x (c - 2r) celdas con
	//    2 * (f - 2r) + 2 * (c - 2r) - 4 celdas, mientras ambas dimensiones
	//    sean al menos 2 (si queda una sola fila o columna, la capa es esa
	//    linea). Restando capas completas de k se llega a la capa correcta
	//    en O(min(f, c)); despues se resuelve dentro del anillo con
	//    aritmetica. Util cuando f * c es demasiado grande para una lista.
	
	// 7. TRAMPA: UNA SOLA FILA O COLUMNA. Sin los dos if de guarda, con una
	//    matriz 1 x c el tercer bucle recorre otra vez las celdas de la
	//    fila, y con f x 1 el cuarto repite la columna. Los if son la parte
	//    mas facil de olvidar de la espiral por fronteras.
	
	// 8. TRAMPA: LA CONVENCION DEL ZIGZAG. Hay dos: empezar subiendo
	//    (0,0), (0,1), (1,0), ... o empezar bajando (0,0), (1,0), (0,1), ....
	//    El codigo de arriba usa la primera. Confirma cual pide el
	//    enunciado con su ejemplo antes de codificar; el error es solo
	//    cambiar la paridad de d, pero invalida toda la salida.
\end{lstlisting}

\textbf{Ejemplo de funcionamiento:}

\begin{lstlisting}[style=cpp]
	vector<vector<int>> m = {{1,  2,  3,  4},
		{5,  6,  7,  8},
		{9, 10, 11, 12}};
	RecorridosDeMatriz r(m);
	
	for (int x : r.espiral()) cout << x << " ";
	// 1 2 3 4 8 12 11 10 9 5 6 7
	
	for (int x : r.zigzag()) cout << x << " ";
	// 1 2 5 9 6 3 4 7 10 11 8 12
	
	for (int x : r.serpiente()) cout << x << " ";
	// 1 2 3 4 8 7 6 5 9 10 11 12
	
	for (auto& diag : r.antiDiagonales()) {
		for (int x : diag) cout << x << " ";
		cout << "| ";
	}
	// 1 | 2 5 | 3 6 9 | 4 7 10 | 8 11 | 12 |
	
	auto g = RecorridosDeMatriz::generarEspiral(3, 4);
	// 1   2   3   4
	// 10  11  12  5
	// 9   8   7   6
	
	auto p = RecorridosDeMatriz::posicionSerpiente(6, 4);
	cout << p.first << " " << p.second;
	// 1 1   (la casilla 6 esta en la fila 1, columna 1: esa fila va al reves)
\end{lstlisting}

\textbf{Nota:} \texttt{generarEspiral} y \texttt{espiral} son operaciones
\textbf{inversas}: la primera usa el orden para \textbf{escribir} los
números $1, 2, \dots$ en las celdas, y la segunda lo usa para
\textbf{leer} las celdas en ese orden. Por eso conviene devolver
coordenadas y no valores: el mismo \texttt{ordenEspiral} resuelve las dos
mitades del problema. En el ejemplo, la casilla $6$ del tablero en
serpiente cae en $(1,1)$ y no en $(1,2)$, porque la fila $1$ se recorre de
derecha a izquierda: $k=6$ es el segundo elemento de esa fila, que ocupa la
columna $c-1-1=2$... contando desde $0$ el resto es $r=2$, y la columna es
$4-1-2=1$.

\textbf{Regla práctica:} para leer o generar una espiral en una matriz
rectangular completa, usa las cuatro fronteras con las dos guardas; si hay
obstáculos o una región irregular, usa el caminante con \texttt{visitado}.
Para cualquier pregunta sobre diagonales, identifica la clave: $i+j$
(antidiagonal, o alfil/reina hacia un lado) o $i-j+m-1$ (diagonal
principal) y agrupa las celdas por esa clave en una sola pasada. Y en un
tablero numerado en serpiente, no lo construyas: convierte el número de
casilla con \texttt{posicionSerpiente}.

\subsection{Rotación, Transposición y Reflexión}

\textbf{Idea:} una matriz cuadrada tiene \textbf{8 simetrías}: las 4
rotaciones ($0^\circ$, $90^\circ$, $180^\circ$, $270^\circ$) y esas mismas
4 aplicadas después de un espejo. Todas salen de dos operaciones básicas:
la \textbf{transposición} ($b[j][i] = a[i][j]$, el espejo respecto a la
diagonal principal) y la \textbf{reflexión} (invertir el orden de las
columnas, o el de las filas). Con ellas se obtienen las rotaciones
mediante identidades fáciles de recordar: rotar $90^\circ$ en sentido
horario es \textbf{transponer y luego invertir cada fila}; rotar $90^\circ$
antihorario es \textbf{transponer y luego invertir el orden de las
	filas}; y rotar $180^\circ$ es hacer las dos reflexiones. En una matriz
$n \times m$ que no es cuadrada, transponer y rotar $90^\circ$ producen
una matriz $m \times n$, así que las fórmulas de índices usan siempre las
dimensiones \textbf{originales}:
$$\text{rotar } 90^\circ \text{ horario: } b[j][n-1-i] = a[i][j], \qquad
\text{antihorario: } b[m-1-j][i] = a[i][j].$$

\textbf{¿Por qué usarlo?} La razón principal es \textbf{ahorrar casos}:
si ya sabes implementar «mover todo hacia abajo», no escribas otras tres
versiones para arriba, izquierda y derecha; rota la matriz para que la
dirección pedida quede hacia abajo, aplica lo que ya tienes y rota de
vuelta (modificación 5, y \textit{Gravedad y Caída de Piezas} en la
siguiente subsección). La segunda razón es \textbf{comparar figuras salvo
	simetría}: se generan las variantes, se toma la menor como \textbf{forma
	canónica}, y dos figuras son equivalentes exactamente cuando su forma
canónica es igual, lo que permite contarlas con un \texttt{set}. Se
distingue de \textit{Recorridos Especiales} en que allí se \textbf{lee} la
matriz en otro orden sin cambiarla, mientras que aquí se \textbf{construye}
una matriz nueva. Y si solo importan unos pocos puntos de una matriz
gigantesca, no conviene transformar la matriz: se transforman las
\textbf{coordenadas} (modificación 2).

\textbf{Casos de uso:}

\begin{itemize}
	\item Rotar una imagen o matriz $90^\circ$, $180^\circ$ o $270^\circ$
	(\textit{Rotate Image}), incluso en el sitio si es cuadrada.
	\item Decidir si dos figuras o fichas son la misma salvo rotación
	(4 variantes) o salvo rotación y reflexión (8 variantes), y contar
	cuántas figuras distintas hay.
	\item Reducir cuatro direcciones a una: gravedad, deslizamientos tipo
	\textit{2048} o simulaciones que se implementan una vez y se reutilizan
	rotando la grilla.
	\item Aplicar una \textbf{secuencia enorme} de operaciones
	rotar/espejar sobre una matriz grande, componiéndolas en una sola
	simetría (modificación 3).
	\item Rotar piezas de un Tetris o de un rompecabezas (siguiente
	subsección).
\end{itemize}

\textbf{Complejidad:} cada transformación copia cada celda a su destino
una vez, así que cuesta $O(n \cdot m)$ en tiempo y $O(n \cdot m)$ de
memoria por la matriz nueva. En una matriz \textbf{cuadrada} se puede rotar
en el sitio con $O(n^2)$ en tiempo y $O(1)$ de memoria extra: la
transposición hace $n(n-1)/2$ intercambios y la inversión de filas hace
otros $\approx n^2/2$. \texttt{variantes} hace 8 transformaciones, es decir
$O(8 \cdot n \cdot m) = O(n \cdot m)$; \texttt{formaCanonica} agrega 7
comparaciones lexicográficas de $O(n \cdot m)$ cada una, lo que sigue siendo
$O(n \cdot m)$. Contar $k$ figuras distintas con un \texttt{set} cuesta
$O(k \cdot n \cdot m \cdot \log k)$.

\begin{lstlisting}[style=cpp]
	typedef vector<vector<int>> Mat;
	
	struct TransformacionesDeMatriz {
		int filas, columnas;
		Mat a;
		
		// CREATE: guarda la matriz y sus dimensiones.
		TransformacionesDeMatriz(const Mat& m) {
			a = m;
			filas = a.size();
			columnas = a.empty() ? 0 : a[0].size();
		}                                          // O(filas * columnas)
		
		// READ: transposicion. b[j][i] = a[i][j]. Resultado columnas x filas.
		Mat transponer() const {
			Mat b(columnas, vector<int>(filas));
			for (int i = 0; i < filas; i++)
			for (int j = 0; j < columnas; j++) b[j][i] = a[i][j];
			return b;
		}                                          // O(filas * columnas)
		
		// READ: espejo izquierda-derecha (invierte cada fila).
		Mat reflejarHorizontal() const {
			Mat b(filas, vector<int>(columnas));
			for (int i = 0; i < filas; i++)
			for (int j = 0; j < columnas; j++) b[i][columnas - 1 - j] = a[i][j];
			return b;
		}                                          // O(filas * columnas)
		
		// READ: espejo arriba-abajo (invierte el orden de las filas).
		Mat reflejarVertical() const {
			Mat b(filas, vector<int>(columnas));
			for (int i = 0; i < filas; i++)
			for (int j = 0; j < columnas; j++) b[filas - 1 - i][j] = a[i][j];
			return b;
		}                                          // O(filas * columnas)
		
		// READ: rotacion de 90 grados en sentido horario. Resultado columnas x filas.
		Mat rotar90() const {
			Mat b(columnas, vector<int>(filas));
			for (int i = 0; i < filas; i++)
			for (int j = 0; j < columnas; j++) b[j][filas - 1 - i] = a[i][j];
			return b;
		}                                          // O(filas * columnas)
		
		// READ: rotacion de 180 grados.
		Mat rotar180() const {
			Mat b(filas, vector<int>(columnas));
			for (int i = 0; i < filas; i++)
			for (int j = 0; j < columnas; j++)
			b[filas - 1 - i][columnas - 1 - j] = a[i][j];
			return b;
		}                                          // O(filas * columnas)
		
		// READ: rotacion de 90 grados en sentido antihorario (= 270 horario).
		Mat rotar270() const {
			Mat b(columnas, vector<int>(filas));
			for (int i = 0; i < filas; i++)
			for (int j = 0; j < columnas; j++) b[columnas - 1 - j][i] = a[i][j];
			return b;
		}                                          // O(filas * columnas)
		
		// READ: k rotaciones de 90 grados horarias; k puede ser negativo o grande.
		Mat rotarK(int k) const {
			k = ((k % 4) + 4) % 4;
			if (k == 0) return a;
			if (k == 1) return rotar90();
			if (k == 2) return rotar180();
			return rotar270();
		}                                          // O(filas * columnas)
		
		// READ: las 8 variantes: 4 rotaciones de la matriz y 4 de su espejo.
		// Las posiciones 0..3 son SOLO rotaciones; la 4..7, con reflexion.
		vector<Mat> variantes() const {
			TransformacionesDeMatriz espejo(reflejarHorizontal());
			vector<Mat> v;
			for (int k = 0; k < 4; k++) v.push_back(rotarK(k));
			for (int k = 0; k < 4; k++) v.push_back(espejo.rotarK(k));
			return v;
		}                                          // O(8 * filas * columnas)
		
		// READ: forma canonica = la menor variante (lexicografica). Dos figuras
		// son equivalentes salvo simetria si y solo si tienen la misma.
		Mat formaCanonica() const {
			vector<Mat> v = variantes();
			return *min_element(v.begin(), v.end());
		}                                          // O(filas * columnas)
		
		// READ: true si 'otra' es esta matriz rotada (o rotada y reflejada,
		// si conReflexion es true).
		bool iguales(const Mat& otra, bool conReflexion) const {
			vector<Mat> v = variantes();
			int limite = conReflexion ? 8 : 4;
			for (int k = 0; k < limite; k++) if (v[k] == otra) return true;
			return false;
		}                                          // O(filas * columnas)
		
		// WRITE: rotacion de 90 grados horaria EN EL SITIO, solo para
		// matrices cuadradas: transponer y luego invertir cada fila.
		void rotar90EnSitio() {
			for (int i = 0; i < filas; i++)
			for (int j = 0; j < i; j++) swap(a[i][j], a[j][i]);
			for (int i = 0; i < filas; i++) reverse(a[i].begin(), a[i].end());
		}                                          // O(filas * filas), O(1) extra
	};
	
	// ------------------------------------------------------------
	// MODIFICACIONES Y COMENTARIOS CON LOS USOS
	// ------------------------------------------------------------
	
	// 1. ROTAR EN EL SITIO POR CICLOS DE 4 (matriz cuadrada n x n): cada celda
	//    pertenece a un ciclo de 4 posiciones que se rota con un solo temporal:
	//
	//    for (int i = 0; i < n / 2; i++) {
		//        for (int j = i; j < n - 1 - i; j++) {
			//            int tmp = a[i][j];
			//            a[i][j] = a[n - 1 - j][i];
			//            a[n - 1 - j][i] = a[n - 1 - i][n - 1 - j];
			//            a[n - 1 - i][n - 1 - j] = a[j][n - 1 - i];
			//            a[j][n - 1 - i] = tmp;
			//        }
		//    }
	//
	//    Hace lo mismo que rotar90EnSitio; la version con transposicion es
	//    mas corta y menos propensa a errores de indices.
	
	// 2. ROTAR COORDENADAS EN VEZ DE LA MATRIZ: si la matriz es enorme (por
	//    ejemplo 10^9 x 10^9) y solo hay unos pocos puntos de interes, se
	//    mueve cada punto (i, j) de una matriz de n filas y m columnas:
	//      transponer:            (i, j) -> (j, i)
	//      rotar 90 horario:      (i, j) -> (j, n - 1 - i)
	//      rotar 90 antihorario:  (i, j) -> (m - 1 - j, i)
	//      rotar 180:             (i, j) -> (n - 1 - i, m - 1 - j)
	//      espejo horizontal:     (i, j) -> (i, m - 1 - j)
	//      espejo vertical:       (i, j) -> (n - 1 - i, j)
	//    Ojo: tras transponer o rotar 90, n y m se intercambian.
	
	// 3. COMPONER MUCHAS OPERACIONES EN UNA SOLA: si el enunciado da q
	//    operaciones (transponer, espejos, rotaciones) sobre una matriz
	//    grande, no se aplican una por una (O(q * n * m)). Se aplican a un
	//    MARCADOR 2 x 2 con cuatro valores distintos, que identifica de
	//    forma unica cual de las 8 simetrias se acumulo:
	//
	//    Mat marcador = {{0, 1}, {2, 3}};
	//    for (cada operacion op del enunciado) marcador = aplicar(marcador, op);
	//
	//    auto vs = TransformacionesDeMatriz({{0, 1}, {2, 3}}).variantes();
	//    int idx = find(vs.begin(), vs.end(), marcador) - vs.begin();
	//    Mat resultado = TransformacionesDeMatriz(grande).variantes()[idx];
	//
	//    ('aplicar' es un switch que llama a la transformacion que
	//    corresponda.) Cuesta O(q) mas una sola transformacion O(n * m). Sirve
	//    aunque la matriz grande no sea cuadrada, porque cada operacion
	//    actua sobre los indices de la misma forma sin importar el tamano.
	
	// 4. ROTACION CIRCULAR DE UNA FILA NO ES ROTACION DE LA MATRIZ: desplazar
	//    los elementos de la fila i, k posiciones a la izquierda, es
	//    rotate(a[i].begin(), a[i].begin() + k % m, a[i].end());
	//    Es una operacion distinta y mucho mas barata que rotar 90 grados.
	//    Para una columna hay que copiarla a un vector, rotarla y volver a
	//    escribirla.
	
	// 5. REDUCIR CUATRO DIRECCIONES A UNA: si solo tienes implementada la
	//    gravedad hacia abajo (o "mover a la izquierda" en 2048), las otras
	//    direcciones se obtienen asi:
	//      hacia la derecha:   rotar 90 horario   (la derecha pasa a ser abajo),
	//                          aplicar la gravedad, rotar 90 antihorario.
	//      hacia la izquierda: rotar 90 antihorario, aplicar, rotar 90 horario.
	//      hacia arriba:       rotar 180, aplicar, rotar 180.
	//    Todas cuestan O(n * m) extra por rotacion, sin cambiar el orden
	//    asintotico de una simulacion que ya recorre la matriz.
	
	// 6. TRAMPA: HORARIO VS ANTIHORARIO. Se confunden con facilidad. Verifica
	//    con una matriz 2 x 3: tras rotar 90 horario, la esquina superior
	//    izquierda (el 1 del ejemplo) debe quedar ARRIBA A LA DERECHA.
	
	// 7. TRAMPA: LAS DIMENSIONES SE INTERCAMBIAN. Rotar 90 o transponer una
	//    matriz n x m da una m x n; crear la matriz resultado con las
	//    dimensiones nuevas y usar las dimensiones ORIGINALES en la formula
	//    de indices. En el sitio solo se puede si es cuadrada.
	
	// 8. TRAMPA: LA REFLEXION NO SALE DE LAS ROTACIONES. Las piezas S y Z del
	//    Tetris son espejos una de la otra y ninguna rotacion convierte una
	//    en la otra. El enunciado dice si las reflexiones estan permitidas:
	//    4 variantes (solo rotaciones) o 8 (con reflexion). Si la figura ya
	//    tiene simetria propia, varias variantes salen repetidas; usa un set
	//    para quedarte con las distintas.
	
	// 9. TRAMPA: FIGURAS EN CAJAS DE DISTINTO TAMANO. Para comparar figuras
	//    con formaCanonica, primero recorta el borde vacio de cada una; si
	//    no, la misma figura dentro de una caja mas grande parece distinta.
\end{lstlisting}

\textbf{Ejemplo de funcionamiento:}

\begin{lstlisting}[style=cpp]
	Mat m = {{1, 2, 3},
		{4, 5, 6}};
	TransformacionesDeMatriz t(m);
	
	t.transponer();          // 1 4
	// 2 5
	// 3 6
	
	t.reflejarHorizontal();  // 3 2 1
	// 6 5 4
	
	t.reflejarVertical();    // 4 5 6
	// 1 2 3
	
	t.rotar90();             // 4 1     (el 1 pasa arriba a la derecha)
	// 5 2
	// 6 3
	
	t.rotar180();            // 6 5 4
	// 3 2 1
	
	t.rotar270();            // 3 6
	// 2 5
	// 1 4
	
	// Cuantas fichas 2 x 2 de celdas llenas/vacias son distintas salvo
	// rotacion y reflexion? Se prueban las 16 y se cuentan las canonicas.
	set<Mat> formas;
	for (int mask = 0; mask < 16; mask++) {
		Mat p = {{mask & 1, (mask >> 1) & 1},
			{(mask >> 2) & 1, (mask >> 3) & 1}};
		formas.insert(TransformacionesDeMatriz(p).formaCanonica());
	}
	cout << formas.size();
	// 6   (vacia, una celda, dos adyacentes, dos en diagonal, tres, llena)
	
	// Piezas S y Z del Tetris: espejos una de la otra.
	Mat S = {{0, 1, 1},
		{1, 1, 0}};
	Mat Z = {{1, 1, 0},
		{0, 1, 1}};
	TransformacionesDeMatriz ts(S);
	cout << ts.iguales(Z, false) << " " << ts.iguales(Z, true);
	// 0 1   (ninguna rotacion las iguala, con reflexion si)
	
	// La L de tres celdas en sus 4 rotaciones tiene una sola forma canonica.
	Mat L = {{1, 0},
		{1, 1}};
	set<Mat> ls;
	for (int k = 0; k < 4; k++)
	ls.insert(TransformacionesDeMatriz(TransformacionesDeMatriz(L).rotarK(k)).formaCanonica());
	cout << ls.size();
	// 1
\end{lstlisting}

\textbf{Nota:} las identidades entre operaciones son útiles como
\textbf{prueba}: para cualquier matriz, transponer y luego reflejar
horizontalmente debe dar lo mismo que \texttt{rotar90}, transponer y luego
reflejar verticalmente debe dar \texttt{rotar270}, y cuatro rotaciones de
$90^\circ$ deben devolver la matriz original. Si tu código no cumple alguna
de las tres, hay un índice mal puesto. En el ejemplo de la ficha $2 \times 2$,
la respuesta $6$ coincide con las seis clases que se cuentan a mano (por
número de celdas llenas, con dos clases distintas para dos celdas), lo cual
sirve de control del método.

\textbf{Regla práctica:} rota cuando quieras \textbf{evitar escribir el
	mismo código cuatro veces} (una simulación para cada dirección) o cuando
necesites decidir si dos figuras son \textit{la misma}. Para una sola
rotación de una matriz cuadrada, transponer e invertir cada fila es la forma
más corta y segura; para muchas operaciones sobre una matriz grande,
compónlas en una de las 8 simetrías con el marcador $2 \times 2$; y para
comparar figuras, usa \texttt{formaCanonica} dentro de un \texttt{set},
decidiendo antes si el enunciado permite reflexiones (8 variantes) o solo
rotaciones (4).
\subsection{Gravedad y Caída de Piezas (Tetris, Candy Crush)}

\textbf{Idea:} en este tipo de problemas todo lo que no está vacío
\textbf{cae} hasta apoyarse en el fondo o en otra celda. Hay dos modelos.
En el modelo \textbf{celda a celda} (Candy Crush, arena, Puyo Puyo) cada
columna es independiente, y aplicar la gravedad equivale a empujar las
celdas no vacías hacia abajo conservando su orden relativo: se recorre la
columna de abajo hacia arriba con un \textbf{puntero de escritura}
\texttt{pos} que marca la próxima posición libre desde el fondo; cada
celda no vacía se escribe en \texttt{pos} y \texttt{pos} sube una casilla.
Es una partición estable con dos punteros, $O(\text{filas})$ por columna.
En el modelo de \textbf{pieza rígida} (Tetris) un grupo de celdas se
mueve como una unidad: cae mientras \texttt{cabe} en la fila siguiente, se
\textbf{fija} al chocar, y luego se eliminan las filas completas. Candy
Crush es un bucle sobre la plantilla de \textit{Plantilla de Simulación
	(Estado, Paso y Parada)}: el estado es el tablero, el paso es
\textit{marcar grupos, eliminarlos, aplicar gravedad}, y la parada es que
ya no se marque nada. La regla que no se puede omitir es marcar
\textbf{todos} los grupos antes de borrar ninguno, por la misma razón que
en \textit{Autómatas Celulares (Game of Life)}: la actualización debe ser
simultánea.

\textbf{¿Por qué usarlo?} La alternativa ingenua es mover cada celda una
posición y repetir hasta que nada cambie, lo que puede costar
$O(\text{filas} \cdot \text{filas} \cdot \text{columnas})$; el puntero de
escritura logra el estado final de la gravedad en una sola pasada de
$O(\text{filas} \cdot \text{columnas})$. Para las otras tres direcciones
no hace falta reescribir nada: se rota el tablero como en
\textit{Rotación, Transposición y Reflexión} (modificación 5 de esa
subsección) o se refleja el puntero (modificación 1). Frente a una
simulación de \textbf{objetos sueltos} (\textit{Simulación por Eventos}),
aquí toda la grilla cambia en cada ronda y la simulación por rondas es la
forma natural.

\textbf{Casos de uso:}

\begin{itemize}
	\item Candy Crush, Bejeweled o SameGame: eliminar grupos de $\ge 3$,
	dejar caer y repetir, contando celdas eliminadas o cascadas.
	\item Tetris: en qué fila aterriza una pieza, cuántas líneas completa
	y cuándo termina el juego.
	\item Arena, agua o rocas que caen, con obstáculos fijos que las
	detienen (modificación 2).
	\item Puyo Puyo y juegos de grupos conectados, donde el grupo se
	detecta con \textit{Flood Fill} en vez de por filas y columnas
	(modificación 4).
	\item Deslizamientos tipo 2048: comprimir una línea hacia un extremo
	usa exactamente el mismo puntero de escritura.
\end{itemize}

\textbf{Complejidad:} \texttt{aplicarGravedad} recorre cada columna una
vez con un solo puntero, así que cuesta $O(\text{filas} \cdot \text{columnas})$.
\texttt{marcarGrupos} recorre cada fila y cada columna detectando
\textit{tandas} de celdas iguales; el índice avanza siempre hasta el final
de la tanda, así que también es $O(\text{filas} \cdot \text{columnas})$.
Una ronda cuesta lo mismo. Una cascada completa hace a lo sumo
$\text{filas} \cdot \text{columnas} / 3$ rondas, porque toda ronda que
continúa elimina al menos 3 celdas, para un peor caso de
$O((\text{filas} \cdot \text{columnas})^2 / 3)$, aunque en la práctica se
hacen muy pocas. En Tetris, \texttt{cabe} cuesta $O(|p|)$ con $|p|$ celdas
en la pieza, \texttt{filaDeCaida} hace a lo sumo \texttt{filas} intentos
($O(\text{filas} \cdot |p|)$) y \texttt{limpiarLineas} cuesta
$O(\text{filas} \cdot \text{columnas})$.

\begin{lstlisting}[style=cpp]
	typedef vector<pair<int,int>> Pieza;   // celdas (di, dj) relativas a un ancla
	
	struct CaidaDePiezas {
		int filas, columnas;
		vector<vector<int>> t;             // 0 = vacia; otro valor = color o tipo
		int rondas;                        // rondas de la ultima simulacion
		
		// CREATE: guarda el tablero (la fila 0 es la de ARRIBA).
		CaidaDePiezas(const vector<vector<int>>& tablero) {
			t = tablero;
			filas = t.size();
			columnas = t.empty() ? 0 : t[0].size();
			rondas = 0;
		}                                          // O(filas * columnas)
		
		// WRITE: gravedad celda a celda. Por cada columna, de abajo hacia
		// arriba, cada celda no vacia se escribe en 'pos' (la proxima
		// posicion libre desde el fondo).
		void aplicarGravedad() {
			for (int j = 0; j < columnas; j++) {
				int pos = filas - 1;
				for (int i = filas - 1; i >= 0; i--) {
					if (t[i][j] != 0) {
						int v = t[i][j];
						t[i][j] = 0;               // se vacia ANTES de escribir
						t[pos][j] = v;
						pos--;
					}
				}
			}
		}                                          // O(filas * columnas)
		
		// READ: marca las celdas que forman una tanda de 'minimo' o mas celdas
		// iguales, en horizontal o vertical. No modifica el tablero.
		vector<vector<bool>> marcarGrupos(int minimo = 3) const {
			vector<vector<bool>> marca(filas, vector<bool>(columnas, false));
			for (int i = 0; i < filas; i++) {           // tandas por filas
				int j = 0;
				while (j < columnas) {
					int k = j;
					while (k < columnas && t[i][k] == t[i][j]) k++;
					if (t[i][j] != 0 && k - j >= minimo)
					for (int x = j; x < k; x++) marca[i][x] = true;
					j = k;
				}
			}
			for (int j = 0; j < columnas; j++) {        // tandas por columnas
				int i = 0;
				while (i < filas) {
					int k = i;
					while (k < filas && t[k][j] == t[i][j]) k++;
					if (t[i][j] != 0 && k - i >= minimo)
					for (int x = i; x < k; x++) marca[x][j] = true;
					i = k;
				}
			}
			return marca;
		}                                          // O(filas * columnas)
		
		// WRITE: vacia las celdas marcadas; retorna cuantas eran.
		int eliminarMarcadas(const vector<vector<bool>>& marca) {
			int quitadas = 0;
			for (int i = 0; i < filas; i++)
			for (int j = 0; j < columnas; j++)
			if (marca[i][j]) { t[i][j] = 0; quitadas++; }
			return quitadas;
		}                                          // O(filas * columnas)
		
		// WRITE: Candy Crush completo: marcar TODO, eliminar, gravedad, y
		// repetir hasta que no se marque nada. Retorna las celdas eliminadas.
		int simularCandyCrush(int minimo = 3) {
			int total = 0;
			rondas = 0;
			while (true) {
				vector<vector<bool>> marca = marcarGrupos(minimo);
				int quitadas = eliminarMarcadas(marca);
				if (quitadas == 0) break;               // estado estable
				total += quitadas;
				rondas++;
				aplicarGravedad();
			}
			return total;
		}                                          // O(rondas * filas * columnas)
		
		// READ: true si la pieza, con su ancla en (f, c), queda dentro del
		// tablero y sin tocar celdas ocupadas.
		bool cabe(const Pieza& p, int f, int c) const {
			for (auto [di, dj] : p) {
				int i = f + di, j = c + dj;
				if (i < 0 || i >= filas || j < 0 || j >= columnas) return false;
				if (t[i][j] != 0) return false;
			}
			return true;
		}                                          // O(|p|)
		
		// READ: fila del ancla donde aterriza la pieza al soltarla desde la
		// fila 0 en la columna c; -1 si ni siquiera cabe al aparecer (fin
		// del juego).
		int filaDeCaida(const Pieza& p, int c) const {
			if (!cabe(p, 0, c)) return -1;
			int f = 0;
			while (cabe(p, f + 1, c)) f++;
			return f;
		}                                          // O(filas * |p|)
		
		// WRITE: escribe la pieza en el tablero con el valor dado.
		void fijar(const Pieza& p, int f, int c, int valor) {
			for (auto [di, dj] : p) t[f + di][c + dj] = valor;
		}                                          // O(|p|)
		
		// WRITE: elimina las filas completas y baja lo de arriba. Puntero
		// de escritura de abajo hacia arriba, sin borrar filas en el bucle.
		int limpiarLineas() {
			int pos = filas - 1, limpias = 0;
			for (int i = filas - 1; i >= 0; i--) {
				bool llena = true;
				for (int j = 0; j < columnas; j++)
				if (t[i][j] == 0) { llena = false; break; }
				if (llena) { limpias++; continue; }     // esta fila desaparece
				if (pos != i) t[pos] = t[i];
				pos--;
			}
			for (int i = pos; i >= 0; i--) t[i].assign(columnas, 0);
			return limpias;
		}                                          // O(filas * columnas)
	};
	
	// ------------------------------------------------------------
	// MODIFICACIONES Y COMENTARIOS CON LOS USOS
	// ------------------------------------------------------------
	
	// 1. GRAVEDAD HACIA LA DERECHA (sin rotar): el mismo puntero, pero por
	//    filas y de derecha a izquierda:
	//
	//    for (int i = 0; i < filas; i++) {
		//        int pos = columnas - 1;
		//        for (int j = columnas - 1; j >= 0; j--)
		//            if (t[i][j] != 0) { int v = t[i][j]; t[i][j] = 0; t[i][pos] = v; pos--; }
		//    }
	//
	//    Para izquierda, arriba: el puntero parte del otro extremo y avanza
	//    en sentido contrario. Si prefieres no duplicar codigo, rota el
	//    tablero (ver Rotacion, Transposicion y Reflexion).
	
	// 2. OBSTACULOS FIJOS (piedras que no caen y detienen lo de arriba): al
	//    encontrar un obstaculo el puntero se reubica justo encima de el:
	//
	//    for (int j = 0; j < columnas; j++) {
		//        int pos = filas - 1;
		//        for (int i = filas - 1; i >= 0; i--) {
			//            if (t[i][j] == -1) pos = i - 1;          // -1 = obstaculo
			//            else if (t[i][j] != 0) {
				//                int v = t[i][j]; t[i][j] = 0; t[pos][j] = v; pos--;
				//            }
			//        }
		//    }
	//
	//    Cada tramo entre dos obstaculos se comporta como una columna
	//    independiente. Es el nucleo de problemas tipo "Rotating the Box".
	
	// 3. CAIDA RAPIDA CON ALTURAS: en vez de probar fila por fila, guarda
	//    para cada columna el indice de su celda llena mas alta,
	//    top[j] (= filas si esta vacia). Para una pieza con 'baja[k]' =
	//    mayor di de las celdas de su columna k (k = 0..ancho-1):
	//
	//    int f = INT_MAX;
	//    for (int k = 0; k < ancho; k++)
	//        if (baja[k] >= 0) f = min(f, top[c + k] - 1 - baja[k]);
	//    // f es la fila de aterrizaje; f < 0 significa que no cabe
	//
	//    Cuesta O(ancho) en vez de O(filas * |p|). Supone que la pieza
	//    aparece POR ENCIMA de la pila (ninguna de sus celdas queda debajo
	//    de una celda llena de su columna); es lo normal en Tetris.
	//    Tras fijar la pieza, actualiza top[] solo en sus columnas.
	
	// 4. GRUPOS CONECTADOS (Puyo Puyo, SameGame): en vez de tandas por filas
	//    y columnas, se busca con Flood Fill cada componente de celdas del
	//    mismo color y se marca entera si tiene al menos 'minimo' celdas.
	//    La marca y la eliminacion siguen siendo dos pasos separados.
	
	// 5. ROTAR LA PIEZA ANTES DE SOLTARLA: una pieza como lista de celdas
	//    (di, dj) gira 90 grados en sentido horario con
	//    (di, dj) -> (dj, -di). Despues se normaliza restando el menor di y
	//    el menor dj, para que todos los desplazamientos sean >= 0. Con las
	//    4 rotaciones (y, si se permite, las reflexiones) se prueba en que
	//    columna y con que orientacion conviene soltarla.
	
	// 6. FIN DEL JUEGO Y PUNTAJE: filaDeCaida(p, c) == -1 significa que la
	//    pieza no cabe al aparecer. El puntaje suele ser una funcion de las
	//    lineas que devuelve limpiarLineas(); confirma en el enunciado si
	//    cuentan las lineas por separado o si hay bonos por varias a la vez.
	
	// 7. MODELO "UN PASO POR TICK": si cada celda cae UNA casilla por
	//    unidad de tiempo (arena que cae poco a poco), recorre las filas de
	//    ABAJO hacia ARRIBA en cada tick; de arriba hacia abajo una misma
	//    celda caeria varias casillas en un solo tick. Es el paso de la
	//    plantilla de simulacion, y la gravedad completa de arriba es
	//    repetir ese paso hasta que nada cambie.
	
	// 8. TRAMPA: MARCAR ANTES DE BORRAR. Si se borra cada tanda apenas se
	//    detecta, una forma en T o en L pierde la celda que compartian la
	//    tanda horizontal y la vertical, y la segunda tanda puede dejar de
	//    detectarse. Por eso marcarGrupos devuelve marcas y las celdas de
	//    interseccion se cuentan una sola vez.
	
	// 9. TRAMPA: ELIMINAR FILAS DENTRO DEL BUCLE. Llamar a erase sobre las
	//    filas mientras se recorren salta la fila siguiente a la borrada.
	//    limpiarLineas evita el problema copiando cada fila que se queda al
	//    puntero 'pos' y vaciando lo que sobra arriba al final.
	
	// 10. TRAMPA: ARRIBA Y ABAJO. La fila 0 es la de arriba, asi que "caer"
	//     es ir hacia filas MAYORES. Ademas, el valor vacio (0 aqui) no
	//     puede coincidir con ningun color ni tipo de pieza.
\end{lstlisting}

\textbf{Ejemplo de funcionamiento:}

\begin{lstlisting}[style=cpp]
	// Candy Crush: tablero de 5 filas y 4 columnas (0 = vacio).
	// 4 1 4 4
	// 3 2 3 4
	// 4 2 4 3
	// 3 2 1 1
	// 4 3 3 4
	CaidaDePiezas cc({{4, 1, 4, 4},
		{3, 2, 3, 4},
		{4, 2, 4, 3},
		{3, 2, 1, 1},
		{4, 3, 3, 4}});
	
	// Ronda 1: solo hay una tanda, los tres 2 de la columna 1, y se eliminan
	// 3 celdas. Tras la gravedad, el 1 de arriba cae junto a los otros dos:
	// 4 . 4 4
	// 3 . 3 4
	// 4 . 4 3
	// 3 1 1 1      <- ahora hay una tanda horizontal de tres 1
	// 4 3 3 4
	//
	// Ronda 2: se eliminan los tres 1 y todo cae de nuevo:
	// 4 . . .
	// 3 . 4 4
	// 4 . 3 4
	// 3 . 4 3
	// 4 3 3 4
	//
	// Ronda 3: no se marca nada, se detiene.
	
	int total = cc.simularCandyCrush();
	cout << total << " " << cc.rondas;
	// 6 2   (6 celdas eliminadas en 2 rondas con eliminacion)
	
	// Tetris: tablero de 5 filas y 4 columnas, con la fila de abajo
	// llena salvo la columna 3. Se suelta una pieza I vertical (4 celdas)
	// en la columna 3.
	vector<vector<int>> vacio(5, vector<int>(4, 0));
	vacio[4][0] = vacio[4][1] = vacio[4][2] = 1;
	CaidaDePiezas juego(vacio);
	Pieza I = {{0, 0}, {1, 0}, {2, 0}, {3, 0}};
	
	int f = juego.filaDeCaida(I, 3);
	cout << f;                        // 1  (ocupa las filas 1 a 4)
	juego.fijar(I, f, 3, 2);
	// . . . .
	// . . . 2
	// . . . 2
	// . . . 2
	// 1 1 1 2      <- fila completa
	
	cout << juego.limpiarLineas();    // 1
	// . . . .
	// . . . .
	// . . . 2      <- lo de arriba bajo una fila
	// . . . 2
	// . . . 2
\end{lstlisting}

\textbf{Nota:} en el ejemplo de Candy Crush la segunda ronda \textbf{no
	existía} en el tablero inicial: la tanda de los 1 aparece solo después de
que la gravedad acerca los dos 1 de la fila 3 a un tercero. Por eso una
solución que detecte los grupos una sola vez, sin repetir después de la
gravedad, da la respuesta incorrecta. Además, \texttt{simularCandyCrush}
cuenta \textbf{celdas} eliminadas; si el enunciado puntúa por grupos, por
tamaño o por cascada, hay que contarlos en \texttt{marcarGrupos} o en cada
ronda.

\textbf{Regla práctica:} para gravedad celda a celda usa el puntero de
escritura columna por columna, y si hay obstáculos fijos reubícalo encima de
cada uno. Para Candy Crush usa siempre el ciclo \textit{marcar todo,
	eliminar, gravedad, repetir hasta estabilizar}, y para Tetris usa
\texttt{cabe} para la caída, \texttt{fijar} al aterrizar y
\texttt{limpiarLineas} con puntero de escritura, nunca borrando filas dentro
del bucle. Si necesitas otra dirección de caída, rota el tablero en vez de
duplicar el código.

\subsection{Direcciones y Vecindarios (dx/dy, 8 Direcciones, Límites)}

\textbf{Idea:} casi todo problema sobre una grilla pregunta por los
\textbf{vecinos} de una celda $(i,j)$. En vez de escribir cuatro u ocho
\texttt{if}, se guardan los desplazamientos en dos arreglos paralelos
\texttt{dx[]} y \texttt{dy[]} y se recorre el vecindario con un bucle:
$(ni, nj) = (i + dx[d],\, j + dy[d])$. En esta guía la fila $i$ crece
hacia abajo, así que \texttt{dx} cambia la \textbf{fila} y \texttt{dy} la
\textbf{columna}. Hay tres reglas que evitan casi todos los errores.
Primera: comprobar los límites \textbf{antes} de acceder a la celda.
Segunda: ordenar las direcciones en \textbf{sentido horario}, empezando
hacia arriba, porque así los giros y los opuestos son aritmética
($d+1$ gira a la derecha, $d+3$ a la izquierda, y el opuesto es $d \oplus 2$
con 4 direcciones o $d \oplus 4$ con 8). Tercera: definir los arreglos una
vez y usarlos en todo el programa. Los vecindarios habituales son el de
\textbf{von Neumann} (4 direcciones ortogonales), el de \textbf{Moore} (8,
con diagonales), los 8 saltos del \textbf{caballo}, y las 6 direcciones de
una grilla \textbf{hexagonal}.

\textbf{¿Por qué usarlo?} Copiar un \texttt{if} por vecino multiplica el
código y los errores, mientras que con los arreglos el mismo bucle sirve
para 4 u 8 direcciones cambiando solo el límite del bucle. Es el ingrediente
común de \textit{Flood Fill}, \textit{Laberinto con BFS},
\textit{Tour de Caballo de Ajedrez}, la cuenta de vecinas vivas de
\textit{Autómatas Celulares (Game of Life)} y la espiral con caminante de
\textit{Recorridos Especiales}. Para los límites hay tres opciones: la
función \texttt{dentro}, que no modifica el tablero y sirve con cualquier
tamaño; el \textbf{borde de centinelas}, que evita la comprobación
rodeando la grilla de paredes (modificación 1); y el truco de
\texttt{unsigned}, que resume dos comparaciones en una.

\textbf{Casos de uso:}

\begin{itemize}
	\item Visitar los vecinos de una celda en BFS, DFS o Flood Fill,
	con 4 o con 8 direcciones.
	\item Buscaminas: cada celda cuenta las minas de sus 8 vecinas.
	\item Movimientos del caballo o del rey, y piezas que \textbf{deslizan}
	(torre, alfil, reina, hielo) hasta chocar.
	\item Detectar líneas de $k$ celdas iguales (tres en raya, Conecta 4,
	Gomoku): solo hacen falta 4 de las 8 direcciones, contando en ambos
	sentidos.
	\item Robots o caminantes con orientación, que giran a la derecha, a
	la izquierda o dan la vuelta.
	\item Grillas hexagonales con coordenadas axiales (modificación 5).
\end{itemize}

\textbf{Complejidad:} recorrer el vecindario de una celda cuesta $O(k)$
con $k$ constante (4, 6 u 8), así que visitar los vecinos de toda la grilla
es $O(n \cdot m \cdot k) = O(n \cdot m)$. \texttt{dentro} es $O(1)$,
\texttt{buscaminas} cuesta $O(8 \cdot n \cdot m)$, \texttt{deslizar} cuesta
$O(\max(n, m))$ porque avanza a lo sumo hasta el otro borde, y
\texttt{hizoLinea} cuesta $O(4 \cdot k)$ por consulta, ya que avanza como
máximo $k-1$ pasos por cada sentido de cada una de las 4 direcciones.

\begin{lstlisting}[style=cpp]
	// Direcciones en sentido HORARIO empezando hacia arriba. La fila crece
	// hacia abajo: dx cambia la FILA y dy cambia la COLUMNA.
	const int DX4[4] = {-1,  0, 1,  0};       // arriba, derecha, abajo, izquierda
	const int DY4[4] = { 0,  1, 0, -1};
	const int DX8[8] = {-1, -1, 0, 1, 1,  1,  0, -1};   // + diagonales, horario
	const int DY8[8] = { 0,  1, 1, 1, 0, -1, -1, -1};
	const int DXC[8] = {-2, -2, -1, -1, 1, 1, 2, 2};    // saltos del caballo
	const int DYC[8] = {-1,  1, -2,  2, -2, 2, -1, 1};
	const int DXL[4] = {0, 1, 1,  1};         // media circunferencia: para
	const int DYL[4] = {1, 0, 1, -1};         // contar lineas (se cuenta en ambos sentidos)
	
	struct Vecindarios {
		int filas, columnas;
		vector<string> g;                      // tablero de caracteres
		
		// CREATE: guarda el tablero.
		Vecindarios(const vector<string>& tablero) {
			g = tablero;
			filas = g.size();
			columnas = g.empty() ? 0 : g[0].size();
		}                                          // O(filas * columnas)
		
		// READ: true si (i, j) esta dentro del tablero. El truco unsigned: un
		// entero negativo convertido a unsigned es enorme y falla la
		// comparacion, asi que una sola comparacion por eje cubre los dos lados.
		bool dentro(int i, int j) const {
			return (unsigned)i < (unsigned)filas && (unsigned)j < (unsigned)columnas;
		}                                          // O(1)
		
		// READ: vecinos validos de (i, j), con 4 u 8 direcciones.
		vector<pair<int,int>> vecinos(int i, int j, bool ocho) const {
			vector<pair<int,int>> res;
			int k = ocho ? 8 : 4;
			for (int d = 0; d < k; d++) {
				int ni = i + (ocho ? DX8[d] : DX4[d]);
				int nj = j + (ocho ? DY8[d] : DY4[d]);
				if (dentro(ni, nj)) res.push_back({ni, nj});
			}
			return res;
		}                                          // O(k)
		
		// READ: cuantos vecinos de (i, j) son el caracter c.
		int contarAdyacentes(int i, int j, char c, bool ocho) const {
			int cuenta = 0;
			for (auto [ni, nj] : vecinos(i, j, ocho))
			if (g[ni][nj] == c) cuenta++;
			return cuenta;
		}                                          // O(k)
		
		// READ: buscaminas: cada celda que no es mina ('*') se reemplaza por
		// la cantidad de minas en sus 8 vecinas ('.' si es cero).
		vector<string> buscaminas() const {
			vector<string> res = g;
			for (int i = 0; i < filas; i++)
			for (int j = 0; j < columnas; j++)
			if (g[i][j] != '*') {
				int n = contarAdyacentes(i, j, '*', true);
				res[i][j] = n ? char('0' + n) : '.';
			}
			return res;
		}                                          // O(8 * filas * columnas)
		
		// READ: desde (i, j) avanza en la direccion d (indice de DX4/DY4)
		// hasta el borde o hasta una pared '#'; retorna la celda final.
		// Es el movimiento de hielo o de una torre.
		pair<int,int> deslizar(int i, int j, int d) const {
			while (dentro(i + DX4[d], j + DY4[d]) && g[i + DX4[d]][j + DY4[d]] != '#') {
				i += DX4[d];
				j += DY4[d];
			}
			return {i, j};
		}                                          // O(max(filas, columnas))
		
		// READ: true si la celda (i, j) pertenece a una linea de al menos k
		// celdas iguales. Solo se prueban 4 direcciones y se cuenta hacia
		// adelante y hacia atras, porque la linea es la misma vista de
		// cualquiera de sus dos extremos.
		bool hizoLinea(int i, int j, int k) const {
			char c = g[i][j];
			for (int d = 0; d < 4; d++) {
				int cuenta = 1;
				for (int s = 1; s < k; s++) {           // hacia adelante
					int ni = i + s * DXL[d], nj = j + s * DYL[d];
					if (!dentro(ni, nj) || g[ni][nj] != c) break;
					cuenta++;
				}
				for (int s = 1; s < k; s++) {           // hacia atras
					int ni = i - s * DXL[d], nj = j - s * DYL[d];
					if (!dentro(ni, nj) || g[ni][nj] != c) break;
					cuenta++;
				}
				if (cuenta >= k) return true;
			}
			return false;
		}                                          // O(4 * k)
	};
	
	// ------------------------------------------------------------
	// MODIFICACIONES Y COMENTARIOS CON LOS USOS
	// ------------------------------------------------------------
	
	// 1. BORDE DE CENTINELAS: se rodea el tablero con una fila y una columna
	//    extra de paredes; los indices utiles pasan a ser 1..n y 1..m y ya no
	//    hace falta comprobar los limites, porque un vecino "fuera" cae en la
	//    pared y se descarta con el mismo chequeo que cualquier otra pared:
	//
	//    vector<string> P(n + 2, string(m + 2, '#'));
	//    for (int i = 0; i < n; i++)
	//        for (int j = 0; j < m; j++) P[i + 1][j + 1] = a[i][j];
	//    // vecino: P[i + DX4[d]][j + DY4[d]] == '#' -> pared o borde
	//
	//    Ahorra una comprobacion en el bucle mas caliente de un BFS. No sirve
	//    si "fuera del tablero" no debe tratarse como pared (por ejemplo, un
	//    mundo toroidal).
	
	// 2. INDICE UNICO id = i * m + j: los vecinos de arriba y abajo son
	//    id - m e id + m (comprobando 0 <= id < n * m), pero los de izquierda
	//    y derecha id - 1 e id + 1 exigen mirar la columna, porque de lo
	//    contrario "el vecino izquierdo" de la columna 0 es la ultima celda
	//    de la fila anterior:
	//
	//    int j = id % m;
	//    if (j > 0)     { /* vecino izquierdo: id - 1 */ }
	//    if (j < m - 1) { /* vecino derecho:   id + 1 */ }
	
	// 3. MUNDO TOROIDAL (los bordes dan la vuelta): en vez de descartar el
	//    vecino, se envuelve con modulo. Se suma n antes del % para que el
	//    resultado nunca sea negativo:
	//
	//    int ni = (i + DX4[d] + n) % n;
	//    int nj = (j + DY4[d] + m) % m;
	//
	//    Solo es correcto si |dx| <= n y |dy| <= m; con saltos mayores usa
	//    ((x % n) + n) % n.
	
	// 4. OTRAS PIEZAS: el rey usa las 8 direcciones de DX8/DY8 con un solo
	//    paso; el caballo usa DXC/DYC (8 saltos, cada uno valido si cae dentro
	//    del tablero); la torre desliza en los indices pares de DX8/DY8
	//    (0, 2, 4, 6 = ortogonales); el alfil desliza en los impares
	//    (1, 3, 5, 7 = diagonales); y la reina desliza en los 8. Deslizar es
	//    repetir el paso mientras la celda siguiente este dentro y libre, como
	//    hace deslizar().
	
	// 5. GRILLA HEXAGONAL: con coordenadas axiales (q, r), las 6 direcciones
	//    son siempre las mismas, sin depender de la paridad de la fila:
	//
	//    const int HQ[6] = {1,  1,  0, -1, -1, 0};
	//    const int HR[6] = {0, -1, -1,  0,  1, 1};
	//
	//    La direccion opuesta a d es (d + 3) % 6, y la distancia entre dos
	//    celdas es max(|dq|, |dr|, |dq + dr|). Si el enunciado da un
	//    "odd-r" o "even-q" en vez de axiales, convierte las coordenadas
	//    antes de trabajar.
	
	// 6. DIRECCIONES DADAS COMO LETRAS: se convierten en un indice con una
	//    cadena que siga el MISMO orden que DX4/DY4:
	//
	//    int d = string("URDL").find(c);          // U=arriba R=derecha D=abajo L=izquierda
	//    int d = string("NESO").find(c);          // en espanol: N E S O
	//    // si find devuelve npos, la letra no es valida
	
	// 7. GIROS Y OPUESTOS con direcciones en sentido horario:
	//
	//    int derecha   = (d + 1) % 4;
	//    int izquierda = (d + 3) % 4;             // NO (d - 1) % 4: en C++ es negativo
	//    int atras     = d ^ 2;                   // en 8 direcciones: d ^ 4
	//
	//    Vale para cualquier orden circular consistente (horario o
	//    antihorario), pero hay que usar el mismo en todo el programa.
	
	// 8. TRAMPA: LIMITES ANTES DE ACCEDER. En una condicion como
	//    "dentro(ni, nj) && g[ni][nj] == c" el && evalua dentro primero y no
	//    llega a acceder cuando esta fuera. Invertir el orden lee memoria
	//    invalida. Con el truco unsigned, ambos operandos deben convertirse a
	//    unsigned; convertir solo uno vuelve a comparar como enteros con signo.
	
	// 9. TRAMPA: FILA Y COLUMNA, X E Y. Aqui dx es la fila (vertical) y dy la
	//    columna (horizontal), lo contrario de la convencion cartesiana en la
	//    que x es horizontal. Ademas, en un plano cartesiano "arriba" es
	//    y + 1, mientras que en la grilla es fila - 1. Si el enunciado da
	//    coordenadas (x, y), convierte antes de usar los arreglos.
	
	// 10. TRAMPA: EL DESPLAZAMIENTO (0, 0). Al generar las 8 direcciones con
	//     dos bucles de -1 a 1, la pareja (0, 0) es la propia celda; hay que
	//     saltarla o la celda se cuenta como su propia vecina.
	
	// 11. TRAMPA: DIAGONALES Y COSTOS. Con 8 direcciones y todos los pasos de
	//     costo 1, sigue valiendo BFS. Si la diagonal cuesta raiz de 2 (o
	//     cualquier otro costo), la distancia ya no crece de a uno y hace
	//     falta Dijkstra. Y algunos enunciados prohiben "cortar esquinas"
	//     (pasar en diagonal entre dos paredes); confirmalo antes de
	//     codificar.
\end{lstlisting}

\textbf{Ejemplo de funcionamiento:}

\begin{lstlisting}[style=cpp]
	// Buscaminas: '*' es una mina.
	Vecindarios v({"*...",
		"....",
		".*..",
		"...."});
	
	for (const string& fila : v.buscaminas()) cout << fila << "\n";
	// *1..
	// 221.
	// 1*1.
	// 111.
	
	// Vecinos de la esquina (0, 0): solo los que estan dentro del tablero.
	for (auto [i, j] : v.vecinos(0, 0, true)) cout << "(" << i << "," << j << ")";
	// (0,1)(1,1)(1,0)         8 direcciones: 3 vecinos
	for (auto [i, j] : v.vecinos(0, 0, false)) cout << "(" << i << "," << j << ")";
	// (0,1)(1,0)              4 direcciones: 2 vecinos
	
	cout << v.contarAdyacentes(1, 1, '*', true) << " "
	<< v.contarAdyacentes(1, 1, '*', false);
	// 2 1   (con diagonales ve la mina de (0,0); sin ellas, solo la de (2,1))
	
	// Hielo: '#' es una pared. Se desliza hasta chocar.
	Vecindarios h({"..#.",
		"....",
		".#.."});
	
	auto p = h.deslizar(0, 0, 1);      // derecha desde (0,0)
	cout << p.first << "," << p.second;  // 0,1   (la pared esta en (0,2))
	p = h.deslizar(0, 0, 2);           // abajo desde (0,0)
	cout << p.first << "," << p.second;  // 2,0   (llega al borde inferior)
	p = h.deslizar(1, 3, 3);           // izquierda desde (1,3)
	cout << p.first << "," << p.second;  // 1,0   (llega al borde izquierdo)
	p = h.deslizar(2, 2, 0);           // arriba desde (2,2)
	cout << p.first << "," << p.second;  // 1,2   (la pared esta en (0,2))
	
	// Lineas de 4 (como en Conecta 4): diagonal completa de 4 X.
	Vecindarios l({"X...",
		".X..",
		"..X.",
		"...X"});
	cout << l.hizoLinea(0, 0, 4) << " " << l.hizoLinea(1, 1, 4) << " " << l.hizoLinea(2, 2, 5);
	// 1 1 0   (la diagonal mide 4: cumple k = 4 pero no k = 5)
	
	Vecindarios l3({"X...",
		".X..",
		"..X.",
		"...."});
	cout << l3.hizoLinea(1, 1, 4) << " " << l3.hizoLinea(1, 1, 3);
	// 0 1   (la linea tiene 3 celdas)
\end{lstlisting}

\textbf{Nota:} \texttt{hizoLinea} prueba solo cuatro direcciones,
\texttt{(0,1)}, \texttt{(1,0)}, \texttt{(1,1)} y \texttt{(1,-1)}, y cuenta
hacia adelante y hacia atrás desde la celda: las otras cuatro son las mismas
líneas recorridas en sentido contrario, así que probarlas duplicaría el
trabajo sin encontrar nada nuevo. Por eso el arreglo \texttt{DXL/DYL} es una
\textit{media circunferencia} de direcciones y no las 8 de \texttt{DX8}.
También conviene notar que \texttt{vecinos} en una esquina devuelve menos
celdas que en el interior (3 y 2 en el ejemplo, en vez de 8 y 4): un código
que asuma siempre 8 vecinos falla justo en los bordes.

\textbf{Regla práctica:} define \texttt{dx/dy} una sola vez, en sentido
horario y empezando hacia arriba, y usa siempre el mismo bucle para
recorrer el vecindario, sin escribir un \texttt{if} por dirección. Comprueba
los límites antes de acceder (\texttt{dentro}, o un borde de centinelas si
el tablero no es toroidal), y aprovecha que con ese orden girar es
\texttt{(d+1)\%4}, dar la vuelta es \texttt{d \^{} 2} y la dirección opuesta en
8 direcciones es \texttt{d \^{} 4}. Antes de codificar, confirma tres cosas
del enunciado: cuántas direcciones cuentan (4 u 8), qué pasa en los bordes
(pared, toro u otro), y si la diagonal puede cortar esquinas.

\section{Fechas y Tiempo}
\subsection{Años Bisiestos y Días de un Mes}

\textbf{Idea:} un año es bisiesto (tiene 366 días, con un 29 de febrero)
si es divisible por 4, \textbf{excepto} los divisibles por 100, que
\textbf{sí vuelven} a ser bisiestos si además son divisibles por 400. Esta
es la regla del calendario \textbf{gregoriano} y explica por qué el año
2000 fue bisiesto pero 1900 y 2100 no lo son. La condición se evalúa en
ese orden exacto (400 primero, luego 100, luego 4) porque cada regla
\textbf{anula} a la anterior: sin esa jerarquía, un año como 1900
(divisible por 4 y por 100 pero no por 400) se clasificaría mal. A partir
de esta única función se derivan todas las demás preguntas del
calendario: cuántos días tiene un mes (solo febrero cambia), cuántos días
tiene un año, y qué número de día del año es una fecha dada (sumando los
días de los meses anteriores).

\textbf{¿Por qué usarlo?} Es la pieza base de \textit{Fechas y Tiempo}:
\textit{Día de la Semana} y \textit{Diferencia entre Fechas y Sumar Días}
dependen de saber cuántos días tiene cada mes, y eso depende de si el año
es bisiesto. El error más común es usar \texttt{anio \% 4 == 0} sin las
excepciones, que falla exactamente uno de cada 100 años (los ``siglos'' no
múltiplos de 400) — un caso de prueba que aparece con frecuencia
justamente porque es la trampa clásica. Frente a llamar a una librería de
fechas del lenguaje (que en C++ competitivo casi nunca está disponible o
es incómoda de usar), esta regla se calcula en $O(1)$ con aritmética
simple, sin tablas ni dependencias externas.

\textbf{Casos de uso:}

\begin{itemize}
	\item Verificar si un año es bisiesto para saber si febrero tiene 28
	o 29 días.
	\item Calcular cuántos días tiene un año completo (365 o 366), base de
	cualquier cálculo de diferencia de fechas en años.
	\item Convertir una fecha (día, mes, año) a su \textbf{número de día
		del año} (del 1 al 365 o 366), el primer paso para restar dos fechas.
	\item Validar que una fecha dada sea correcta (por ejemplo, rechazar
	el 30 de febrero o el 31 de abril).
	\item Generar un calendario o iterar día por día dentro de un mes o
	un año, sabiendo dónde termina cada uno.
\end{itemize}

\textbf{Complejidad:} $O(1)$ para \texttt{esBisiesto} y
\texttt{diasDelMes} (tres comparaciones y una consulta a una tabla
fija de 12 valores). \texttt{diaDelAnio} suma los días de hasta 11 meses
anteriores, así que es $O(1)$ también (una constante acotada por 12, no
crece con la entrada).

\begin{lstlisting}[style=cpp]
	struct CalendarioBasico {
		
		// READ: true si el anio es bisiesto (calendario gregoriano). El
		// orden de las condiciones importa: 400 anula a 100, que anula a 4.
		static bool esBisiesto(int anio) {
			if (anio % 400 == 0) return true;
			if (anio % 100 == 0) return false;
			return anio % 4 == 0;
		}                                          // O(1)
		
		// READ: cuantos dias tiene el mes (1..12) en ese anio. Solo febrero
		// (mes 2) depende de si el anio es bisiesto.
		static int diasDelMes(int mes, int anio) {
			static const int dias[13] = {0,31,28,31,30,31,30,31,31,30,31,30,31};
			if (mes == 2 && esBisiesto(anio)) return 29;
			return dias[mes];
		}                                          // O(1)
		
		// READ: cuantos dias tiene el anio completo.
		static int diasDelAnio(int anio) {
			return esBisiesto(anio) ? 366 : 365;
		}                                          // O(1)
		
		// READ: numero de dia del anio (1 = 1 de enero) para la fecha dada.
		static int diaDelAnio(int dia, int mes, int anio) {
			int total = dia;
			for (int m = 1; m < mes; m++) total += diasDelMes(m, anio);
			return total;
		}                                          // O(1) (a lo sumo 11 sumas)
	};
	
	// ------------------------------------------------------------
	// MODIFICACIONES Y COMENTARIOS CON LOS USOS
	// ------------------------------------------------------------
	
	// 1. VALIDAR UNA FECHA: una fecha (d, m, y) es correcta si
	//    1 <= m <= 12 y 1 <= d <= diasDelMes(m, y). Util antes de procesar
	//    entradas que pueden venir mal formadas.
	
	// 2. DIA DEL ANIO CON PREFIJOS PRECALCULADOS: si hay que consultar
	//    muchas fechas del MISMO anio, conviene precalcular una sola vez
	//    los 12 prefijos de dias por mes en vez de sumar en cada consulta:
	//
	//    vector<int> prefijo(13, 0);
	//    for (int m = 1; m <= 12; m++)
	//        prefijo[m] = prefijo[m - 1] + diasDelMes(m, anio);
	//    // diaDelAnio(d, m) = prefijo[m - 1] + d, en O(1) real por consulta
	
	// 3. RANGO DE ANIOS BISIESTOS: contar bisiestos en [1, N] sin recorrer
	//    uno por uno es N/4 - N/100 + N/400 (inclusion-exclusion sobre los
	//    tres divisores); para [A, B] se resta el conteo de [1, A-1] al de
	//    [1, B].
	
	// 4. CALENDARIO JULIANO (antes de la reforma gregoriana de 1582): la
	//    regla es simplemente anio % 4 == 0, SIN las excepciones de 100 y
	//    400. Si el enunciado menciona el calendario juliano explicitamente,
	//    no apliques esta funcion tal cual.
	
	// 5. TRAMPA CLASICA: 1900 Y 2000. 1900 no es bisiesto (divisible por
	//    100 pero no por 400) y 2000 si lo es (divisible por 400). Un
	//    generador de casos de prueba para este tema casi siempre incluye
	//    alguno de estos dos anios exactamente para atrapar la
	//    implementacion ingenua "anio % 4 == 0".
	
	// 6. TRAMPA: ORDEN DE LAS CONDICIONES. "anio % 4 == 0 && anio % 100 != 0
	//    || anio % 400 == 0" sin parentesis depende de la precedencia de
	//    && sobre ||, y es facil de leer mal. La version con tres if
	//    consecutivos (arriba) es mas clara y evita ese riesgo.
\end{lstlisting}

\textbf{Ejemplo de funcionamiento:}

\begin{lstlisting}[style=cpp]
	cout << CalendarioBasico::esBisiesto(2000) << " "     // 1 (divisible por 400)
	<< CalendarioBasico::esBisiesto(1900) << " "     // 0 (divisible por 100, no por 400)
	<< CalendarioBasico::esBisiesto(2024) << " "     // 1 (divisible por 4, no por 100)
	<< CalendarioBasico::esBisiesto(2023);           // 0 (no divisible por 4)
	// 1 0 1 0
	
	cout << CalendarioBasico::diasDelMes(2, 2024) << " "   // 29
	<< CalendarioBasico::diasDelMes(2, 2023);        // 28
	
	cout << CalendarioBasico::diaDelAnio(29, 2, 2024);     // 60 (31 de enero + 29 de febrero)
	cout << CalendarioBasico::diaDelAnio(31, 12, 2023);    // 365 (ultimo dia de un anio no bisiesto)
\end{lstlisting}

\textbf{Regla práctica:} nunca uses \texttt{anio \% 4 == 0} solo — siempre
las tres condiciones en el orden 400, 100, 4. Cualquier otra pregunta del
calendario (días de un mes, días de un año, día del año) se construye
llamando a \texttt{esBisiesto} una sola vez y derivando el resto con
aritmética simple, sin tablas de fechas completas.

\subsection{Día de la Semana (Zeller / Sakamoto)}

\textbf{Idea:} ambos algoritmos calculan, en $O(1)$ y sin iterar día por
día, qué día de la semana cae una fecha del calendario gregoriano. La
idea común es tratar enero y febrero como los \textbf{meses 13 y 14 del
	año anterior}: así, el año siempre ``cambia'' en marzo, y el día 29 de
febrero (que solo existe en años bisiestos) queda al \textbf{final} del
año extendido en vez de interrumpir la aritmética a mitad de año. Con esa
convención, cada fórmula combina el día, un término que depende del mes
(basado en cuántos días de \textit{desplazamiento} acumula cada mes
respecto al anterior), y un término que depende del año (incluyendo la
misma corrección de bisiestos vista en \textit{Años Bisiestos y Días de
	un Mes}, expresada como \texttt{año/4 - año/100 + año/400}), todo reducido
módulo 7. La diferencia entre los dos algoritmos es solo de
\textbf{convención}: Zeller nació pensando en el calendario juliano y su
resultado nativo numera desde el sábado; Sakamoto es una reescritura
posterior, pensada para el gregoriano, cuyo resultado nativo numera desde
el domingo — la tabla de desplazamientos por mes de Sakamoto es,
matemáticamente, otra forma de la misma corrección.

\textbf{¿Por qué usarlo?} Sin esta fórmula, la única forma de saber el
día de la semana de una fecha lejana es \textbf{contar días uno por uno}
desde una fecha de referencia conocida, un recorrido de $O(\text{días})$
que es inviable si la fecha está a miles de años de distancia. Zeller y
Sakamoto resuelven exactamente el mismo problema en $O(1)$; en la
práctica, Sakamoto se prefiere porque su tabla ya está en base domingo=0,
que es la convención más común en los enunciados de competencia, y evita
el paso extra de reindexar el resultado de Zeller. Frente a
\textit{Diferencia entre Fechas y Sumar Días} (que sí necesita convertir
fechas a un número de días acumulado para restarlas), esta subsección
resuelve una pregunta más puntual — \textit{qué día de la semana} — sin
necesitar ese número de días en absoluto.

\textbf{Casos de uso:}

\begin{itemize}
	\item Determinar el día de la semana de cualquier fecha gregoriana,
	pasada o futura, sin recorrer el calendario.
	\item Contar cuántas veces cae un día de la semana específico (por
	ejemplo, cuántos viernes 13 hay) en un rango de fechas.
	\item Verificar la consistencia de un calendario dado como entrada
	(por ejemplo, que el 1 de enero de un año caiga en el día correcto).
	\item Como paso auxiliar de \textit{Diferencia entre Fechas y Sumar
		Días}, para etiquetar cada fecha resultante con su día de la semana.
	\item Resolver acertijos de calendario del tipo ``¿qué día de la
	semana fue/será tal fecha?'' con una sola fórmula cerrada.
\end{itemize}

\textbf{Complejidad:} $O(1)$ para ambos algoritmos: cada uno hace un
número fijo de operaciones aritméticas (multiplicaciones, divisiones
enteras y un módulo), sin ningún bucle ni dependencia del tamaño de la
fecha (funciona igual de rápido para el año 100 que para el año
1{,}000{,}000).

\begin{lstlisting}[style=cpp]
	struct DiaDeLaSemana {
		
		// READ: algoritmo de Zeller. Retorna 0=sabado, 1=domingo, 2=lunes,
		// ..., 6=viernes (la convencion NATIVA de Zeller, poco intuitiva).
		static int zeller(int dia, int mes, int anio) {
			if (mes < 3) { mes += 12; anio -= 1; }   // enero/febrero = meses 13/14 del anio anterior
			int K = anio % 100;                       // anio dentro del siglo
			int J = anio / 100;                       // siglo
			int h = (dia + (13 * (mes + 1)) / 5 + K + K / 4 + J / 4 + 5 * J) % 7;
			return h;
		}                                          // O(1)
		
		// READ: reindexa el resultado de zeller() a la convencion
		// 0=domingo, 1=lunes, ..., 6=sabado (la mas comun en competencia).
		static int zellerADiaSemana(int h) {
			static const int tabla[7] = {6, 0, 1, 2, 3, 4, 5};
			return tabla[h];
		}                                          // O(1)
		
		// READ: algoritmo de Sakamoto. Retorna DIRECTAMENTE 0=domingo,
		// 1=lunes, ..., 6=sabado; no necesita reindexar.
		static int sakamoto(int anio, int mes, int dia) {
			static const int desplazamiento[] = {0,3,2,5,0,3,5,1,4,6,2,4};
			if (mes < 3) anio -= 1;                   // enero/febrero = meses del anio anterior
			int h = (anio + anio / 4 - anio / 100 + anio / 400
			+ desplazamiento[mes - 1] + dia) % 7;
			return h;
		}                                          // O(1)
	};
	
	// ------------------------------------------------------------
	// MODIFICACIONES Y COMENTARIOS CON LOS USOS
	// ------------------------------------------------------------
	
	// 1. NOMBRES DE LOS DIAS: con la convencion 0=domingo (la de Sakamoto o
	//    la de zellerADiaSemana), basta indexar un arreglo:
	//
	//    const string NOMBRES[7] = {"domingo","lunes","martes","miercoles",
		//                                "jueves","viernes","sabado"};
	//    cout << NOMBRES[DiaDeLaSemana::sakamoto(2024, 9, 29)];
	
	// 2. ANIOS NEGATIVOS Y EL OPERADOR %: en C++, el resultado de % puede
	//    ser negativo si el operando izquierdo lo es (por ejemplo, si
	//    "anio - 1" o alguna resta intermedia da un numero negativo en un
	//    caso limite). Para blindar el resultado, se normaliza sumando 7
	//    antes del modulo final:
	//
	//    int h = ((.../* expresion original */...) % 7 + 7) % 7;
	//
	//    Con los rangos tipicos de un problema de competencia (anios
	//    positivos razonables) esto casi nunca hace falta, pero es una
	//    salvaguarda barata si el enunciado permite anios muy pequenos o
	//    negativos.
	
	// 3. CONTAR UN DIA DE LA SEMANA EN UN RANGO: para contar, por ejemplo,
	//    cuantos domingos hay entre dos fechas, NO conviene recorrer dia por
	//    dia; se combina con Diferencia entre Fechas y Sumar Dias: se
	//    convierten ambas fechas a un numero de dia consecutivo (dia
	//    juliano), se calcula el dia de la semana de una de ellas con
	//    Sakamoto, y el conteo se resuelve con la misma aritmetica de
	//    "conteo de multiplos en un rango" que se usa para contar multiplos
	//    de k en [L, R].
	
	// 4. VERIFICACION CRUZADA: si tienes dudas sobre la implementacion,
	//    Zeller y Sakamoto deben dar SIEMPRE el mismo dia de la semana para
	//    cualquier fecha valida (solo cambia la convencion de numeracion,
	//    que se corrige con zellerADiaSemana). Compararlos entre si en los
	//    casos de prueba propios es una buena forma de detectar un error de
	//    signo o de indice antes de enviar la solucion.
	
	// 5. TRAMPA: LA CONVENCION DE ZELLER. El resultado crudo de zeller()
	//    empieza en SABADO = 0, no en domingo. Usarlo directamente como
	//    indice de un arreglo "domingo, lunes, ..." da el dia equivocado; hay
	//    que pasar por zellerADiaSemana() o restar y ajustar manualmente.
	
	// 6. TRAMPA: ENERO Y FEBRERO. Si se olvida el ajuste "mes < 3: mes += 12,
	//    anio -= 1", el termino de correccion de bisiestos (anio/4 - anio/100
	//    + anio/400) queda calculado con el anio equivocado para las fechas
	//    de enero y febrero, y el resultado falla exactamente en esos dos
	//    meses - un patron de error facil de identificar si aparece.
\end{lstlisting}

\textbf{Ejemplo de funcionamiento:}

\begin{lstlisting}[style=cpp]
	const string NOMBRES[7] = {"domingo","lunes","martes","miercoles",
		"jueves","viernes","sabado"};
	
	// 1 de enero de 2000 fue sabado.
	int hz = DiaDeLaSemana::zeller(1, 1, 2000);
	cout << NOMBRES[DiaDeLaSemana::zellerADiaSemana(hz)];   // sabado
	cout << NOMBRES[DiaDeLaSemana::sakamoto(2000, 1, 1)];   // sabado
	
	// 29 de febrero de 2000 (bisiesto) fue martes.
	cout << NOMBRES[DiaDeLaSemana::sakamoto(2000, 2, 29)];  // martes
	
	// 14 de julio de 1994 fue jueves.
	cout << NOMBRES[DiaDeLaSemana::sakamoto(1994, 7, 14)];  // jueves
\end{lstlisting}

\textbf{Regla práctica:} para el día de la semana de una fecha
gregoriana, usa Sakamoto directamente — ya numera desde domingo y evita el
paso de reindexar que exige Zeller. Antes de codificar cualquiera de los
dos, confirma que el ajuste ``enero y febrero son los meses 13 y 14 del
año anterior'' está aplicado \textbf{antes} de calcular cualquier otro
término de la fórmula, porque es el error más común y el más difícil de
detectar sin un caso de prueba en esos dos meses.
\subsection{Diferencia entre Fechas y Sumar Días}

\textbf{Idea:} Trabajar directamente con la terna (año, mes, día) es
incómodo porque los meses tienen longitudes distintas y los años bisiestos
rompen la regularidad. La solución es convertir cada fecha a un único entero:
el número de días transcurridos desde una época fija (aquí, 1970-01-01 es el
día $0$). En esa representación el calendario desaparece: la diferencia entre
dos fechas es una resta y sumar $k$ días es una suma. Después se convierte el
entero de vuelta a (año, mes, día).

La conversión funciona porque se \textit{desplaza el año para que empiece en
	marzo}. Así el día extra de los bisiestos (29 de febrero) cae al final del año
desplazado y no afecta el conteo de los meses anteriores. Además, los meses
desde marzo siguen el patrón $31,30,31,30,31,31,30,31,30,31,31,(28/29)$, y ese
patrón se captura con la fórmula $\lfloor (153 \cdot mp + 2)/5 \rfloor$, donde
$mp$ es el mes contado desde marzo. Los ciclos gregorianos de 400 años tienen
exactamente $146097$ días, por lo que basta dividir en «eras» de 400 años y
resolver dentro de cada era con aritmética entera.

\textbf{¿Por qué usarlo?} La alternativa ingenua es simular día a día (o mes a
mes) hasta llegar a la fecha destino. Eso cuesta $O(\text{días})$ o
$O(\text{años})$ y es inviable cuando las fechas están separadas por siglos o
hay muchas consultas. Otra alternativa es una tabla de días acumulados por mes,
que sirve dentro de un año pero obliga a manejar aparte el cruce de años y los
bisiestos. La conversión a número de día hace todo en $O(1)$ con una sola
fórmula, es correcta para cualquier año (incluso negativos) y compone bien: una
vez que tienes el entero, la diferencia, la suma, la comparación y el día de la
semana son operaciones triviales. Se prefiere sobre el simulado siempre que las
fechas no sean del mismo mes.

\textbf{Casos de uso:}
\begin{itemize}
	\item Calcular cuántos días hay entre dos fechas (edad en días, plazos, días de retraso).
	\item Calcular la fecha que resulta de sumar o restar $k$ días a una fecha dada.
	\item Determinar el día de la semana de cualquier fecha.
	\item Ordenar o comparar fechas usando el entero como clave (por ejemplo, en un \texttt{sort} o un \texttt{map}).
	\item Problemas de calendarios: contar cuántos viernes 13 hay en un rango, o el siguiente día que cumple una propiedad.
\end{itemize}

\textbf{Complejidad:} Cada conversión (fecha $\to$ entero y entero $\to$ fecha)
ejecuta un número fijo de sumas, multiplicaciones y divisiones enteras, sin
ciclos ni recursión, por lo que cuesta $O(1)$. La diferencia entre dos fechas
son dos conversiones más una resta: $O(1)$. Sumar $k$ días es una conversión,
una suma y la conversión inversa: $O(1)$, independiente de la magnitud de $k$.
Memoria: $O(1)$. Para $q$ consultas, el total es $O(q)$.

\begin{lstlisting}[style=cpp]
	struct Fecha {
		int anio, mes, dia;
		
		// CREATE: guarda anio, mes y dia (se asume una fecha valida).
		Fecha(int a = 1970, int m = 1, int d = 1) : anio(a), mes(m), dia(d) {} // O(1)
		
		// READ: indica si un anio es bisiesto (regla gregoriana).
		static bool es_bisiesto(int a) {
			return (a % 4 == 0 && a % 100 != 0) || a % 400 == 0;
		} // O(1)
		
		// READ: devuelve cuantos dias tiene el mes m del anio a.
		static int dias_del_mes(int a, int m) {
			static const int t[12] = {31,28,31,30,31,30,31,31,30,31,30,31};
			return (m == 2 && es_bisiesto(a)) ? 29 : t[m - 1];
		} // O(1)
		
		// READ: convierte la fecha a dias transcurridos desde 1970-01-01 (dia 0).
		// Se desplaza el anio para que empiece en marzo (febrero queda al final).
		long long a_dias() const {
			long long y = anio;
			y -= (mes <= 2);
			long long era = (y >= 0 ? y : y - 399) / 400;
			long long yoe = y - era * 400;                                   // [0, 399]
			long long mp  = (mes > 2) ? mes - 3 : mes + 9;                   // marzo = 0
			long long doy = (153 * mp + 2) / 5 + dia - 1;                    // dia del anio desplazado
			long long doe = yoe * 365 + yoe / 4 - yoe / 100 + doy;           // dia dentro de la era
			return era * 146097 + doe - 719468;
		} // O(1)
		
		// CREATE: construye la fecha correspondiente a z dias desde 1970-01-01.
		static Fecha desde_dias(long long z) {
			z += 719468;
			long long era = (z >= 0 ? z : z - 146096) / 146097;
			long long doe = z - era * 146097;                                // [0, 146096]
			long long yoe = (doe - doe / 1460 + doe / 36524 - doe / 146096) / 365;
			long long y   = yoe + era * 400;
			long long doy = doe - (365 * yoe + yoe / 4 - yoe / 100);
			long long mp  = (5 * doy + 2) / 153;
			int d = (int)(doy - (153 * mp + 2) / 5 + 1);
			int m = (int)(mp < 10 ? mp + 3 : mp - 9);
			return Fecha((int)(y + (m <= 2)), m, d);
		} // O(1)
		
		// READ: dias que hay de "otra" a esta fecha (positivo si esta es posterior).
		long long diferencia(const Fecha& otra) const {
			return a_dias() - otra.a_dias();
		} // O(1)
		
		// READ: dia de la semana (0 = domingo, 1 = lunes, ..., 6 = sabado).
		// 1970-01-01 fue jueves (4); el doble modulo cubre dias negativos.
		int dia_semana() const {
			return (int)(((a_dias() + 4) % 7 + 7) % 7);
		} // O(1)
		
		// WRITE: avanza (o retrocede si k < 0) la fecha k dias.
		void sumar_dias(long long k) {
			*this = desde_dias(a_dias() + k);
		} // O(1)
	};
	
	// ------------------------------------------------------------
	// MODIFICACIONES Y COMENTARIOS CON LOS USOS
	// ------------------------------------------------------------
	// 1. Restar dias: llamar sumar_dias con k negativo. No requiere codigo extra
	//    porque a_dias() y desde_dias() funcionan para z negativo.
	//
	// 2. Dia del anio (1..366): sirve para validar o para fechas ordinales.
	//    int dia_del_anio() const {
		//        return (int)(a_dias() - Fecha(anio, 1, 1).a_dias()) + 1;
		//    } // O(1)
	//
	// 3. Sumar meses (con recorte al ultimo dia valido del mes destino):
	//    void sumar_meses(int k) {
		//        long long total = (long long)anio * 12 + (mes - 1) + k;
		//        anio = (int)(total / 12); mes = (int)(total % 12) + 1;   // si total < 0, ajustar
		//        dia = min(dia, dias_del_mes(anio, mes));
		//    } // O(1)
	//    Ojo: 31 de enero + 1 mes da 28/29 de febrero, no 3 de marzo.
	//
	// 4. Siguiente dia de la semana d (0..6) estrictamente posterior:
	//    long long k = ((d - dia_semana() + 6) % 7) + 1;  sumar_dias(k);   // O(1)
	//
	// 5. Comparar fechas: comparar a_dias() (o la terna en orden anio, mes, dia).
	//    Como clave de map/sort: usar a_dias() como llave entera.
\end{lstlisting}

\textbf{Ejemplo de funcionamiento:}
\begin{lstlisting}[style=cpp]
	int main() {
		Fecha a(2000, 2, 28), b(2000, 3, 1);
		cout << b.diferencia(a) << "\n";            // 2  (28-feb -> 29-feb -> 1-mar, 2000 es bisiesto)
		
		Fecha c(2024, 1, 1), d(2024, 12, 31);
		cout << d.diferencia(c) << "\n";            // 365  (2024 tiene 366 dias; del 1-ene al 31-dic hay 365 saltos)
		
		Fecha e(2024, 1, 1);
		e.sumar_dias(100);
		cout << e.anio << "-" << e.mes << "-" << e.dia << "\n";   // 2024-4-10  (31 + 29 + 31 = 91 dias hasta el 1-abr; +9)
		
		cout << Fecha(2000, 1, 1).a_dias() << "\n";    // 10957  (30 anios * 365 + 7 bisiestos)
		cout << Fecha(2000, 1, 1).dia_semana() << "\n"; // 6  (sabado)
		cout << Fecha(2024, 1, 1).dia_semana() << "\n"; // 1  (lunes)
		return 0;
	}
\end{lstlisting}

\textbf{Nota:} \texttt{a\_dias()} devuelve \texttt{long long} porque con años
grandes (del orden de $10^9$) el conteo de días excede los 32 bits. Si el
problema garantiza años razonables (por ejemplo entre 1900 y 2100), \texttt{int}
alcanza, pero usar \texttt{long long} no cuesta nada.

\textbf{Regla práctica:} Si necesitas relacionar dos fechas que pueden estar en
meses o años distintos, convierte ambas a número de día y opera con enteros.
Simula día a día solo cuando el rango sea diminuto (unos pocos días) y no valga
la pena escribir la conversión. Usa la tabla de días acumulados por mes solo
cuando todas las fechas estén garantizadas dentro de un mismo año no bisiesto.

\subsection{Formatos de Hora (12h/24h, Segundos a HH:MM:SS)}

\textbf{Idea:} Igual que con las fechas, la forma más limpia de manejar horas es
guardarlas como un solo entero: los segundos transcurridos desde las 00:00:00.
Como un día tiene $24 \cdot 3600 = 86400$ segundos, cualquier hora del reloj
corresponde a un valor en $[0, 86399]$. El formato \texttt{HH:MM:SS} es
simplemente una representación en base mixta de ese entero: las horas son el
cociente de dividir entre $3600$, los minutos son el cociente del resto entre
$60$, y los segundos son lo que sobra. La conversión inversa es la suma
$3600h + 60m + s$.

Sumar tiempo es aritmética modular: el reloj es cíclico, así que después de
sumar $k$ segundos se toma el resultado módulo $86400$. Con $k$ negativo hay que
normalizar con un segundo módulo, porque en C++ el operador \texttt{\%} puede
devolver un valor negativo. El formato de 12 horas es solo una presentación
sobre las mismas horas: $h \bmod 12$ da la hora en el reloj (con el caso
especial de que $0$ se muestra como $12$), y $h < 12$ decide entre AM y PM.

\textbf{¿Por qué usarlo?} Manipular \texttt{h}, \texttt{m} y \texttt{s} por
separado obliga a propagar acarreos a mano (segundos $\to$ minutos $\to$
horas $\to$ día) y es la fuente típica de errores por uno. Con el entero, la
suma, la resta y la comparación de horas son operaciones directas, y el
formateo se hace solo al leer o imprimir. Esto es la misma idea que la
conversión de fechas a número de día, aplicada a un dominio cíclico y de tamaño
fijo. Una diferencia con las fechas es que aquí hay que decidir explícitamente
qué pasa al desbordar: en una \textit{hora del reloj} se envuelve módulo
$86400$, mientras que en una \textit{duración} no se envuelve y las horas pueden
superar $23$.

\textbf{Casos de uso:}
\begin{itemize}
	\item Imprimir un contador de segundos como \texttt{HH:MM:SS} (cronómetros, tiempos de carrera, logs).
	\item Convertir entre formato de 24 horas y 12 horas con AM/PM (entradas y salidas de problemas de formato).
	\item Sumar o restar segundos, minutos u horas a una hora del reloj, con vuelta al día siguiente o anterior.
	\item Calcular cuánto tiempo transcurre entre dos horas del reloj, incluso cruzando la medianoche.
	\item Aplicar desplazamientos de zona horaria a una hora dada, o formatear duraciones que superan las 24 horas.
\end{itemize}

\textbf{Complejidad:} Convertir segundos a \texttt{HH:MM:SS} son dos divisiones
y dos restos enteros: $O(1)$. La conversión inversa son dos multiplicaciones y
dos sumas: $O(1)$. Sumar $k$ segundos es una suma y un módulo: $O(1)$,
independiente de $k$. Formatear a texto es $O(1)$ porque la salida tiene
longitud fija (8 caracteres, más el sufijo AM/PM). Parsear una cadena es $O(L)$
con $L$ el largo del texto, que en la práctica es constante (una hora tiene a lo
sumo unos 11 caracteres). Memoria: $O(1)$.

\begin{lstlisting}[style=cpp]
	struct Hora {
		int h, m, s;                      // siempre en formato de 24 horas
		static constexpr int DIA = 86400; // segundos en un dia
		
		// CREATE: construye la hora a partir de una cantidad de segundos.
		// Normaliza modulo 86400 (funciona con valores negativos o mayores a un dia).
		Hora(long long seg = 0) {
			seg = ((seg % DIA) + DIA) % DIA;
			h = (int)(seg / 3600);
			m = (int)((seg % 3600) / 60);
			s = (int)(seg % 60);
		} // O(1)
		
		// CREATE: construye la hora a partir de horas, minutos y segundos.
		static Hora desde_hms(int hh, int mm, int ss) {
			return Hora(hh * 3600LL + mm * 60LL + ss);
		} // O(1)
		
		// CREATE: lee texto en formato 24h, ej. "07:45:10" o "7:5:3".
		static Hora desde_texto24(const string& t) {
			int hh = 0, mm = 0, ss = 0;
			sscanf(t.c_str(), "%d:%d:%d", &hh, &mm, &ss);
			return desde_hms(hh, mm, ss);
		} // O(L), L = largo del texto (constante en la practica)
		
		// CREATE: lee texto en formato 12h, ej. "07:45:10 PM".
		// 12 AM -> 0 horas, 12 PM -> 12 horas, 1..11 PM -> 13..23.
		static Hora desde_texto12(const string& t) {
			int hh = 0, mm = 0, ss = 0;
			char suf[3] = {0};
			sscanf(t.c_str(), "%d:%d:%d %2s", &hh, &mm, &ss, suf);
			hh %= 12;                              // 12 -> 0
			if (suf[0] == 'P' || suf[0] == 'p') hh += 12;
			return desde_hms(hh, mm, ss);
		} // O(L)
		
		// READ: devuelve los segundos transcurridos desde las 00:00:00.
		int a_segundos() const {
			return h * 3600 + m * 60 + s;
		} // O(1)
		
		// READ: devuelve la hora como texto "HH:MM:SS" (24 horas).
		string formato24() const {
			char buf[16];
			snprintf(buf, sizeof buf, "%02d:%02d:%02d", h, m, s);
			return string(buf);
		} // O(1)
		
		// READ: devuelve la hora como texto "HH:MM:SS AM/PM" (12 horas).
		string formato12() const {
			int h12 = h % 12;
			if (h12 == 0) h12 = 12;                // 0 y 12 se muestran como 12
			const char* suf = (h < 12) ? "AM" : "PM";
			char buf[24];
			snprintf(buf, sizeof buf, "%02d:%02d:%02d %s", h12, m, s, suf);
			return string(buf);
		} // O(1)
		
		// READ: segundos que hay que avanzar desde esta hora hasta "otra",
		// cruzando la medianoche si hace falta. Resultado en [0, 86399].
		int segundos_hasta(const Hora& otra) const {
			return ((otra.a_segundos() - a_segundos()) % DIA + DIA) % DIA;
		} // O(1)
		
		// WRITE: avanza (o retrocede si k < 0) la hora k segundos, con vuelta al dia.
		void sumar_segundos(long long k) {
			*this = Hora(a_segundos() + k);
		} // O(1)
	};
	
	// ------------------------------------------------------------
	// MODIFICACIONES Y COMENTARIOS CON LOS USOS
	// ------------------------------------------------------------
	// 1. Duraciones (sin envolver a 24h): NO usar el constructor de Hora, porque
	//    aplica modulo 86400. Formatear directamente el entero:
	//    string formato_duracion(long long seg) {
		//        char buf[32];
		//        snprintf(buf, sizeof buf, "%lld:%02lld:%02lld", seg / 3600, (seg % 3600) / 60, seg % 60);
		//        return string(buf);
		//    } // O(1)
	//    Las horas pueden pasar de 23: 100000 s -> "27:46:40"; 359999 s -> "99:59:59".
	//
	// 2. Milisegundos (HH:MM:SS.mmm): guardar el tiempo en milisegundos, con
	//    DIA = 86400000, y extraer ms = t % 1000, luego t /= 1000 y seguir igual.
	//
	// 3. Zona horaria: aplicar el desplazamiento en segundos con sumar_segundos.
	//    Ej. de UTC a UTC-5:  hora.sumar_segundos(-5 * 3600);   // O(1)
	//
	// 4. Lectura con cin en lugar de sscanf: leer "int h; char c; int m, s;"
	//    con cin >> h >> c >> m >> c >> s; (el char absorbe los ':').
	//
	// 5. Validar entrada: antes de construir, comprobar 0 <= h < 24, 0 <= m < 60,
	//    0 <= s < 60 si el enunciado no garantiza formato correcto.
\end{lstlisting}

\textbf{Ejemplo de funcionamiento:}
\begin{lstlisting}[style=cpp]
	int main() {
		Hora a(3661);
		cout << a.formato24() << "\n";                  // 01:01:01   (3600 + 60 + 1)
		
		Hora b(86399);
		cout << b.formato24() << "\n";                  // 23:59:59
		
		Hora c(45296);
		cout << c.formato24() << "\n";                  // 12:34:56   (12*3600 = 43200; 2096 = 34*60 + 56)
		cout << c.formato12() << "\n";                  // 12:34:56 PM
		
		cout << Hora::desde_hms(0, 15, 0).formato12() << "\n";   // 12:15:00 AM
		cout << Hora::desde_hms(15, 30, 0).formato12() << "\n";  // 03:30:00 PM
		
		Hora d = Hora::desde_texto12("07:45:10 PM");
		cout << d.formato24() << "\n";                  // 19:45:10
		
		Hora e = Hora::desde_hms(23, 59, 30);
		e.sumar_segundos(45);
		cout << e.formato24() << "\n";                  // 00:00:15   (86370 + 45 = 86415 = 86400 + 15)
		
		Hora f = Hora::desde_hms(22, 0, 0), g = Hora::desde_hms(6, 0, 0);
		cout << f.segundos_hasta(g) << "\n";            // 28800   (8 horas, cruzando la medianoche)
		return 0;
	}
\end{lstlisting}

\textbf{Nota:} Los casos límite del formato de 12 horas son la fuente de errores
más común. La medianoche (\texttt{00:xx:xx}) es \texttt{12:xx:xx AM} y el
mediodía (\texttt{12:xx:xx}) es \texttt{12:xx:xx PM}; no existe la hora
\texttt{00} en el reloj de 12 horas. Por eso la conversión usa
$h \bmod 12$ y sustituye $0$ por $12$, en lugar de restar $12$ a las horas
mayores que $12$.

\textbf{Regla práctica:} Si el problema mezcla horas con sumas, restas o
comparaciones, guarda todo como segundos (o milisegundos) y convierte solo al
leer o imprimir. Usa el constructor con módulo cuando modeles una hora del reloj
(que se envuelve en $24$h) y el formateo directo de la sección de modificaciones
cuando modeles una duración (que puede superar las $24$ horas). Si el problema
solo pide un cambio de formato sin aritmética, basta con las reglas de
\texttt{formato12()} y \texttt{desde\_texto12()}, sin necesidad de más
estructura.

\section{Aritmética Ad-hoc}

\subsection{Números Grandes (BigInt)}

\textbf{Idea:} Los tipos nativos guardan un número en un ancho fijo
(\texttt{long long} llega a $\approx 9.2\cdot 10^{18}$), pero muchos problemas
piden resultados con cientos o miles de dígitos (factoriales, potencias,
Fibonacci grandes). La solución es representar el número como un vector de
«bloques» (\textit{limbs}) en una base grande, igual que la suma y la
multiplicación a mano que se aprende en la escuela, pero tratando cada bloque
como un dígito gigante.

Se elige la base $B = 10^9$ por dos razones. Primero, $B^2 \approx 10^{18}$
cabe en un \texttt{long long}, así que el producto de dos bloques más el
acarreo nunca desborda. Segundo, como $B$ es potencia de $10$, cada bloque
corresponde exactamente a 9 dígitos decimales y imprimir es trivial (basta
rellenar con ceros a 9 dígitos). El número se guarda en orden little-endian
($d_0$ es el bloque menos significativo) para que los acarreos avancen hacia el
final del vector. El valor representado es $\sum_i d_i \cdot B^i$, y se
mantiene el invariante de que no hay bloques cero al final, lo cual hace que
comparar por tamaño sea correcto.

\textbf{¿Por qué usarlo?} Si el resultado cabe en 18 dígitos, \texttt{long long}
es más simple y mucho más rápido. Si cabe en unos 38 dígitos, \texttt{\_\_int128}
evita escribir estructura propia. Si el enunciado pide el resultado módulo un
primo $p$, no se necesita BigInt: la aritmética modular mantiene todo acotado.
BigInt es la elección correcta cuando se pide el valor \textit{exacto} y
completo (o cuando se necesita comparar valores enormes sin perder precisión),
y en lenguajes como C++ no hay tipo nativo para ello (Java y Python sí lo
tienen). Frente a \texttt{long double}, que aproxima con unos 18 dígitos
significativos, BigInt no pierde ningún dígito.

\textbf{Casos de uso:}
\begin{itemize}
	\item Imprimir factoriales, potencias o coeficientes binomiales exactos (por ejemplo $100!$ o $2^{1000}$).
	\item Calcular términos grandes de Fibonacci, Catalan u otras recurrencias sin módulo.
	\item Obtener la suma de dígitos o el número de dígitos de un resultado enorme.
	\item Comparar o sumar números dados como cadenas de miles de dígitos.
	\item Ser el numerador y denominador de una fracción exacta cuando \texttt{long long} desborda, o un paso previo a criptografía básica (RSA de juguete).
\end{itemize}

\textbf{Complejidad:} Sea $D$ el número de dígitos decimales y $n = \lceil D/9
\rceil$ el número de bloques. Sumar y restar recorren los bloques una sola vez
propagando el acarreo, de modo que cuestan $O(n)$. Comparar es $O(n)$ en el
peor caso (recorre desde el bloque más significativo hasta encontrar una
diferencia). Dividir entre un entero pequeño hace una operación por bloque, con
un resto que se arrastra: $O(n)$. Multiplicar dos números de $n$ y $m$ bloques
con el método escolar ejecuta dos ciclos anidados, uno por cada par de
bloques, y por eso cuesta $O(n \cdot m)$. Convertir desde y hacia texto es
$O(D)$. Memoria: $O(n)$. Por ejemplo, un número de $10^4$ dígitos tiene unos
1112 bloques, así que multiplicarlo por otro igual hace $\approx 1.2\cdot 10^6$
operaciones.

\begin{lstlisting}[style=cpp]
	struct BigInt {
		static constexpr int BASE = 1000000000; // 10^9
		static constexpr int DIG = 9;           // digitos decimales por bloque
		vector<int> d;                          // bloques, el menos significativo primero
		bool neg;                               // signo (el cero nunca es negativo)
		
		// CREATE: construye desde un entero de 64 bits.
		BigInt(long long v = 0) {
			neg = (v < 0);
			unsigned long long u = neg ? 0ULL - (unsigned long long)v : (unsigned long long)v;
			while (u > 0) { d.push_back((int)(u % BASE)); u /= BASE; }
		} // O(1)
		
		// CREATE: construye desde una cadena decimal, ej. "-123456789012345678901".
		static BigInt desde_string(const string& t) {
			BigInt r;
			int ini = 0;
			bool ng = false;
			if (t[0] == '-') { ng = true; ini = 1; }
			else if (t[0] == '+') ini = 1;
			for (int fin = (int)t.size(); fin > ini; fin -= DIG) {
				int lo = max(ini, fin - DIG);
				r.d.push_back(stoi(t.substr(lo, fin - lo)));
			}
			r.neg = ng;
			r.recortar();
			return r;
		} // O(D), D = cantidad de digitos
		
		// WRITE: elimina bloques cero al final y corrige el signo del cero.
		void recortar() {
			while (!d.empty() && d.back() == 0) d.pop_back();
			if (d.empty()) neg = false;
		} // O(n)
		
		// READ: compara valores absolutos. Devuelve -1, 0 o 1.
		static int cmp_abs(const BigInt& a, const BigInt& b) {
			if (a.d.size() != b.d.size()) return a.d.size() < b.d.size() ? -1 : 1;
			for (int i = (int)a.d.size() - 1; i >= 0; i--)
			if (a.d[i] != b.d[i]) return a.d[i] < b.d[i] ? -1 : 1;
			return 0;
		} // O(n)
		
		// READ: suma de valores absolutos (ignora signos).
		static BigInt sumar_abs(const BigInt& a, const BigInt& b) {
			BigInt r;
			long long carry = 0;
			for (size_t i = 0; i < max(a.d.size(), b.d.size()) || carry; i++) {
				long long cur = carry;
				if (i < a.d.size()) cur += a.d[i];
				if (i < b.d.size()) cur += b.d[i];
				r.d.push_back((int)(cur % BASE));
				carry = cur / BASE;
			}
			return r;
		} // O(max(n, m))
		
		// READ: resta de valores absolutos, requiere |a| >= |b|.
		static BigInt restar_abs(const BigInt& a, const BigInt& b) {
			BigInt r;
			long long prestamo = 0;
			for (size_t i = 0; i < a.d.size(); i++) {
				long long cur = (long long)a.d[i] - prestamo - (i < b.d.size() ? b.d[i] : 0);
				if (cur < 0) { cur += BASE; prestamo = 1; } else prestamo = 0;
				r.d.push_back((int)cur);
			}
			r.recortar();
			return r;
		} // O(n)
		
		// READ: devuelve el opuesto (-x).
		BigInt operator-() const {
			BigInt r = *this;
			if (!r.d.empty()) r.neg = !r.neg;
			return r;
		} // O(n)
		
		// READ: suma con signo.
		BigInt operator+(const BigInt& o) const {
			BigInt r;
			if (neg == o.neg) { r = sumar_abs(*this, o); r.neg = neg; }
			else if (cmp_abs(*this, o) >= 0) { r = restar_abs(*this, o); r.neg = neg; }
			else { r = restar_abs(o, *this); r.neg = o.neg; }
			r.recortar();
			return r;
		} // O(max(n, m))
		
		// READ: resta con signo.
		BigInt operator-(const BigInt& o) const {
			return *this + (-o);
		} // O(max(n, m))
		
		// READ: multiplicacion escolar.
		BigInt operator*(const BigInt& o) const {
			BigInt r;
			if (d.empty() || o.d.empty()) return r;
			r.d.assign(d.size() + o.d.size(), 0);
			for (size_t i = 0; i < d.size(); i++) {
				long long carry = 0;
				for (size_t j = 0; j < o.d.size(); j++) {
					long long cur = r.d[i + j] + (long long)d[i] * o.d[j] + carry;
					r.d[i + j] = (int)(cur % BASE);
					carry = cur / BASE;
				}
				r.d[i + o.d.size()] += (int)carry;
			}
			r.neg = (neg != o.neg);
			r.recortar();
			return r;
		} // O(n * m)
		
		// WRITE: divide este numero entre un entero k > 0 (trunca) y devuelve el resto
		// del valor absoluto.
		int dividir_small(int k) {
			long long rem = 0;
			for (int i = (int)d.size() - 1; i >= 0; i--) {
				long long cur = d[i] + rem * BASE;
				d[i] = (int)(cur / k);
				rem = cur % k;
			}
			recortar();
			return (int)rem;
		} // O(n)
		
		// READ: comparaciones.
		bool operator<(const BigInt& o) const {
			if (neg != o.neg) return neg;
			int c = cmp_abs(*this, o);
			return neg ? c > 0 : c < 0;
		} // O(n)
		bool operator==(const BigInt& o) const {
			return neg == o.neg && d == o.d;
		} // O(n)
		
		// READ: convierte a cadena decimal.
		string str() const {
			if (d.empty()) return "0";
			string s = neg ? "-" : "";
			s += to_string(d.back());
			char buf[16];
			for (int i = (int)d.size() - 2; i >= 0; i--) {
				snprintf(buf, sizeof buf, "%09d", d[i]);
				s += buf;
			}
			return s;
		} // O(D)
	};
	
	// ------------------------------------------------------------
	// MODIFICACIONES Y COMENTARIOS CON LOS USOS
	// ------------------------------------------------------------
	// 1. Potencia rapida (exponenciacion binaria):
	//    BigInt potencia(BigInt b, long long e) {
		//        BigInt r(1);
		//        while (e > 0) { if (e & 1) r = r * b; b = b * b; e >>= 1; }
		//        return r;
		//    } // O(M(n) * log e), dominado por el ultimo cuadrado; M(n) = costo de multiplicar
	//
	// 2. Resto modulo un entero pequeno sin modificar el numero (util para
	//    divisibilidad, digitos de control o para pasar a aritmetica modular):
	//    int mod_small(int k) const {
		//        long long rem = 0;
		//        for (int i = (int)d.size() - 1; i >= 0; i--) rem = (rem * BASE + d[i]) % k;
		//        return (int)rem;
		//    } // O(n)
	//
	// 3. Salida con cout: sobrecargar el operador.
	//    ostream& operator<<(ostream& os, const BigInt& x) { return os << x.str(); }
	//
	// 4. Multiplicacion rapida: para n mayor a unos miles de bloques, el metodo escolar
	//    es lento. Karatsuba baja el costo a O(n^1.585), y FFT/NTT a O(n log n). Con FFT
	//    en doubles se usa una base menor (10^3 o 10^4) para no perder precision.
	//
	// 5. Division BigInt / BigInt: no esta incluida (es la operacion mas delicada).
	//    Opciones: division larga de Knuth (algoritmo D) en O(n * m), o buscar cada
	//    bloque del cociente por busqueda binaria usando la multiplicacion. Con ella se
	//    obtienen tambien el MCD (Euclides), el modulo y la raiz cuadrada (Newton).
	//
	// 6. Representacion en base 10 (un digito por posicion): es 9 veces mas lenta y
	//    usa mas memoria, pero simplifica problemas que operan digito a digito
	//    (suma de digitos, palindromos, invertir el numero).
\end{lstlisting}

\textbf{Ejemplo de funcionamiento:}
\begin{lstlisting}[style=cpp]
	int main() {
		BigInt a = BigInt::desde_string("99999999999999999999");   // 20 nueves
		cout << (a + BigInt(1)).str() << "\n";           // 100000000000000000000  (acarreo cruza todos los bloques)
		cout << ((a + BigInt(1)) - BigInt(1)).str() << "\n"; // 99999999999999999999
		cout << (BigInt(5) - BigInt(8)).str() << "\n";   // -3
		
		BigInt f(1);
		for (int i = 2; i <= 25; i++) f = f * BigInt(i);
		cout << f.str() << "\n";                         // 15511210043330985984000000  (25!)
		
		BigInt g(1);
		for (int i = 2; i <= 100; i++) g = g * BigInt(i);
		cout << g.str().size() << "\n";                  // 158  (100! tiene 158 digitos)
		
		BigInt p(1);
		for (int i = 0; i < 100; i++) p = p * BigInt(2);
		cout << p.str() << "\n";                         // 1267650600228229401496703205376  (2^100)
		int resto = p.dividir_small(1000000000);
		cout << p.str() << " " << resto << "\n";         // 1267650600228229401496 703205376  (cociente y resto)
		
		BigInt x(0), y(1);
		for (int i = 0; i < 100; i++) { BigInt t = x + y; x = y; y = t; }
		cout << x.str() << "\n";                         // 354224848179261915075  (Fibonacci(100))
		
		cout << (BigInt::desde_string("-5") < BigInt(3)) << "\n";  // 1
		return 0;
	}
\end{lstlisting}

\textbf{Nota:} En \texttt{operator*}, el máximo valor de
\texttt{cur} es $(B-1) + (B-1)^2 + (B-1) < 10^{18}$, que cabe en un
\texttt{long long}. Por eso no se puede usar una base mayor (por ejemplo $2^{31}$
o $10^{10}$) sin cambiar a \texttt{unsigned long long} o \texttt{\_\_int128}.
Además, tras cada operación se llama a \texttt{recortar()}; olvidarlo deja
bloques cero al final que rompen \texttt{cmp\_abs} y \texttt{operator==}, y puede
producir un «cero negativo».

\textbf{Regla práctica:} Si el resultado cabe en \texttt{long long} o
\texttt{\_\_int128}, no uses BigInt. Si el problema pide la respuesta módulo
$p$, usa aritmética modular. Usa BigInt solo cuando se pida el valor exacto y
completo, o cuando las cifras superen los 38 dígitos. Si los números pasan de
unos pocos miles de dígitos y hay muchas multiplicaciones, cambia a Karatsuba o
FFT (ver modificaciones).

\subsection{Fracciones Exactas}

\textbf{Idea:} Un número racional $p/q$ se guarda como el par (numerador,
denominador), y las operaciones siguen las reglas de la aritmética escolar:
$\tfrac{a}{b} + \tfrac{c}{d} = \tfrac{ad + cb}{bd}$, $\tfrac{a}{b} \cdot
\tfrac{c}{d} = \tfrac{ac}{bd}$, etc. Lo que hace que el resultado sea
\textit{exacto} es que nunca se divide: se conservan numerador y denominador
enteros, y se mantiene siempre la \textbf{forma canónica}: denominador
positivo y $\gcd(|num|, den) = 1$. La forma canónica es única para cada
racional, así que dos fracciones son iguales si y solo si sus pares son
idénticos.

Para comparar $\tfrac{a}{b} < \tfrac{c}{d}$ (con $b,d>0$) se multiplica en cruz:
$ad < cb$. Es válido porque multiplicar ambos lados por $bd>0$ conserva el
orden. Los productos intermedios pueden ser hasta el cuadrado de los valores
originales, por eso se hacen en \texttt{\_\_int128}.

\textbf{¿Por qué usarlo?} Con \texttt{double}, sumas como $0.1 + 0.2$ no dan
exactamente $0.3$, y esos errores se acumulan o rompen comparaciones de
igualdad; una fracción exacta no tiene ese problema. Frente a la aritmética
modular (que devuelve $p \cdot q^{-1} \bmod M$), la fracción conserva el valor
real: sirve cuando se necesita \textit{comparar}, ordenar o imprimir la
fracción irreducible, cosa que un residuo módulo $M$ no permite. Frente a
BigInt (subsección anterior), esta estructura es más ligera pero está acotada:
los denominadores crecen rápido con muchas sumas (el mínimo común múltiplo de
denominadores distintos crece de forma exponencial), y si se desborda
\texttt{long long} hay que cambiar el tipo del numerador y denominador a
BigInt.

\textbf{Casos de uso:}
\begin{itemize}
	\item Comparar pendientes de rectas en geometría sin usar \texttt{double} (por ejemplo, agrupar puntos colineales).
	\item Calcular probabilidades o valores esperados exactos con pocos eventos (sumas de $1/k$, esperanzas de juegos).
	\item Ordenar fracciones o encontrar la mayor/menor de una lista sin errores de redondeo.
	\item Resolver sistemas de ecuaciones lineales pequeños con eliminación gaussiana exacta.
	\item Convertir decimales periódicos a fracción irreducible, o generar fracciones de Farey y aproximaciones racionales.
\end{itemize}

\textbf{Complejidad:} Cada operación aritmética hace un número fijo de
multiplicaciones y sumas, más una llamada a \texttt{gcd} para normalizar. El
algoritmo de Euclides tarda $O(\log M)$ pasos, donde $M$ es la magnitud del
mayor operando: en cada dos iteraciones el resto se reduce al menos a la
mitad (el peor caso son números de Fibonacci consecutivos, Lamé). Por tanto,
suma, resta, producto y cociente cuestan $O(\log M)$. La comparación por
producto cruzado es $O(1)$ (no necesita \texttt{gcd}). Sumar $n$ fracciones en
secuencia es $O(n \log M)$. Memoria: $O(1)$ por fracción.

\begin{lstlisting}[style=cpp]
	struct Fraccion {
		long long num, den; // forma canonica: den > 0 y mcd(|num|, den) = 1
		
		// READ: mcd de dos enteros de 128 bits (ignora signos).
		static __int128 mcd128(__int128 a, __int128 b) {
			if (a < 0) a = -a;
			if (b < 0) b = -b;
			while (b) { a %= b; swap(a, b); }
			return a;
		} // O(log M)
		
		// CREATE: construye la fraccion n/d y la lleva a forma canonica (d != 0).
		Fraccion(long long n = 0, long long d = 1) : num(n), den(d) {
			if (den < 0) { num = -num; den = -den; }
			long long g = (long long)mcd128(num, den);
			if (g > 1) { num /= g; den /= g; }
		} // O(log M)
		
		// CREATE: igual que el constructor, pero desde valores de 128 bits
		// (para reducir antes de volver a 64 bits, evitando desbordes intermedios).
		static Fraccion crear(__int128 n, __int128 d) {
			if (d < 0) { n = -n; d = -d; }
			__int128 g = mcd128(n, d);
			if (g > 1) { n /= g; d /= g; }
			Fraccion r;
			r.num = (long long)n;
			r.den = (long long)d;
			return r;
		} // O(log M)
		
		// READ: suma exacta.
		Fraccion operator+(const Fraccion& o) const {
			return crear((__int128)num * o.den + (__int128)o.num * den, (__int128)den * o.den);
		} // O(log M)
		
		// READ: resta exacta.
		Fraccion operator-(const Fraccion& o) const {
			return crear((__int128)num * o.den - (__int128)o.num * den, (__int128)den * o.den);
		} // O(log M)
		
		// READ: producto exacto.
		Fraccion operator*(const Fraccion& o) const {
			return crear((__int128)num * o.num, (__int128)den * o.den);
		} // O(log M)
		
		// READ: cociente exacto (requiere o.num != 0).
		Fraccion operator/(const Fraccion& o) const {
			return crear((__int128)num * o.den, (__int128)den * o.num);
		} // O(log M)
		
		// READ: opuesto.
		Fraccion operator-() const {
			Fraccion r;
			r.num = -num;
			r.den = den;
			return r;
		} // O(1)
		
		// READ: inverso multiplicativo (requiere num != 0).
		Fraccion inverso() const {
			return crear(den, num);
		} // O(log M)
		
		// READ: comparaciones. Como den > 0, se multiplica en cruz sin dividir.
		bool operator<(const Fraccion& o) const {
			return (__int128)num * o.den < (__int128)o.num * den;
		} // O(1)
		bool operator==(const Fraccion& o) const {
			return num == o.num && den == o.den; // forma canonica => igualdad de pares
		} // O(1)
		
		// READ: valor aproximado como double (solo para mostrar).
		double valor() const {
			return (double)num / (double)den;
		} // O(1)
		
		// READ: texto "p/q", o solo "p" si el denominador es 1.
		string str() const {
			if (den == 1) return to_string(num);
			return to_string(num) + "/" + to_string(den);
		} // O(1)
	};
	
	// ------------------------------------------------------------
	// MODIFICACIONES Y COMENTARIOS CON LOS USOS
	// ------------------------------------------------------------
	// 1. Piso (parte entera hacia -infinito), util para partes enteras y mixtos:
	//    long long piso() const {
		//        long long q = num / den;
		//        if (num % den != 0 && num < 0) q--;
		//        return q;
		//    } // O(1)
	//    Numero mixto: piso() y luego Fraccion(num - piso() * den, den).
	//
	// 2. Potencia entera (exponenciacion binaria; exponente negativo con inverso()):
	//    Fraccion potencia(Fraccion b, long long e) {
		//        Fraccion r(1);
		//        if (e < 0) { b = b.inverso(); e = -e; }
		//        while (e > 0) { if (e & 1) r = r * b; b = b * b; e >>= 1; }
		//        return r;
		//    } // O(log e * log M)
	//
	// 3. Si el enunciado pide "p/q modulo un primo P", NO uses esta estructura:
	//    calcula num * inv(den) mod P, con inv(x) = x^(P-2) mod P (Fermat), valido solo
	//    si den no es multiplo de P.
	//
	// 4. Fraccion con BigInt: cambiar num y den a BigInt cuando los denominadores crezcan
	//    (por ejemplo, la serie armonica con n grande). Requiere la division BigInt/BigInt
	//    y el mcd con ella (ver modificaciones de BigInt).
	//
	// 5. Decimal periodico a fraccion. Para 0.a(b), con a de k digitos y b de m digitos:
	//    valor = ( (ab) - a ) / ( 10^k * (10^m - 1) ), donde (ab) es la concatenacion.
	//    Ejemplo: 0.1(6) = (16 - 1) / (10 * 9) = 15/90 = 1/6.
	//
	// 6. Aproximacion racional de un double con denominador acotado: fracciones continuas
	//    o el arbol de Stern-Brocot (busqueda binaria sobre el arbol), O(log N).
	//
	// 7. Pendientes en geometria: para comparar dy1/dx1 con dy2/dx2 no hace falta normalizar;
	//    basta el producto cruzado dy1 * dx2 < dy2 * dx1 (con dx1, dx2 > 0).
\end{lstlisting}

\textbf{Ejemplo de funcionamiento:}
\begin{lstlisting}[style=cpp]
	int main() {
		Fraccion a(1, 2), b(1, 3);
		cout << (a + b).str() << "\n";              // 5/6
		cout << (a - b).str() << "\n";              // 1/6
		cout << (b - a).str() << "\n";              // -1/6
		cout << (a * b).str() << "\n";              // 1/6
		cout << (a / b).str() << "\n";              // 3/2
		cout << Fraccion(6, -8).str() << "\n";      // -3/4  (signo al numerador, reducida)
		cout << (Fraccion(1, 6) + Fraccion(1, 3)).str() << "\n"; // 1/2  (3/6 reducida)
		cout << (b < a) << "\n";                    // 1  (1*2 < 1*3)
		
		cout << (0.1 + 0.2 == 0.3) << "\n";         // 0  (double falla)
		cout << (Fraccion(1, 10) + Fraccion(2, 10) == Fraccion(3, 10)) << "\n"; // 1
		
		Fraccion h;                                 // 0/1
		for (int i = 1; i <= 10; i++) h = h + Fraccion(1, i);
		cout << h.str() << "\n";                    // 7381/2520  (H_10)
		printf("%.6f\n", h.valor());                // 2.928968
		return 0;
	}
\end{lstlisting}

\textbf{Nota:} Los productos cruzados se hacen en \texttt{\_\_int128}, pero el
resultado final se guarda en \texttt{long long}: si la fracción ya reducida no
cabe en 64 bits, el \texttt{cast} trunca en silencio. Si el problema puede
producir numeradores o denominadores de más de 18 dígitos (por ejemplo, sumar
$1/k$ con $k$ hasta unas decenas de miles), cambia a BigInt. Además, el
constructor y \texttt{crear} suponen $d \neq 0$; no validan la división por
cero.

\textbf{Regla práctica:} Si necesitas el valor exacto de una razón para
comparar, ordenar o imprimir irreducible, usa \texttt{Fraccion}. Si solo
necesitas el resultado módulo un primo, usa aritmética modular con inverso
(modificación 3). Si el numerador o denominador se desbordan de 64 bits, usa
BigInt como tipo interno, y si solo estás comparando pendientes, el producto
cruzado sin normalizar es suficiente.
\subsection{Diferencia entre Fechas y Sumar Días}

\textbf{Idea:} Trabajar directamente con la terna (año, mes, día) es
incómodo porque los meses tienen longitudes distintas y los años bisiestos
rompen la regularidad. La solución es convertir cada fecha a un único entero:
el número de días transcurridos desde una época fija (aquí, 1970-01-01 es el
día $0$). En esa representación el calendario desaparece: la diferencia entre
dos fechas es una resta y sumar $k$ días es una suma. Después se convierte el
entero de vuelta a (año, mes, día).

La conversión funciona porque se \textit{desplaza el año para que empiece en
	marzo}. Así el día extra de los bisiestos (29 de febrero) cae al final del año
desplazado y no afecta el conteo de los meses anteriores. Además, los meses
desde marzo siguen el patrón $31,30,31,30,31,31,30,31,30,31,31,(28/29)$, y ese
patrón se captura con la fórmula $\lfloor (153 \cdot mp + 2)/5 \rfloor$, donde
$mp$ es el mes contado desde marzo. Los ciclos gregorianos de 400 años tienen
exactamente $146097$ días, por lo que basta dividir en «eras» de 400 años y
resolver dentro de cada era con aritmética entera.

\textbf{¿Por qué usarlo?} La alternativa ingenua es simular día a día (o mes a
mes) hasta llegar a la fecha destino. Eso cuesta $O(\text{días})$ o
$O(\text{años})$ y es inviable cuando las fechas están separadas por siglos o
hay muchas consultas. Otra alternativa es una tabla de días acumulados por mes,
que sirve dentro de un año pero obliga a manejar aparte el cruce de años y los
bisiestos. La conversión a número de día hace todo en $O(1)$ con una sola
fórmula, es correcta para cualquier año (incluso negativos) y compone bien: una
vez que tienes el entero, la diferencia, la suma, la comparación y el día de la
semana son operaciones triviales. Se prefiere sobre el simulado siempre que las
fechas no sean del mismo mes.

\textbf{Casos de uso:}
\begin{itemize}
	\item Calcular cuántos días hay entre dos fechas (edad en días, plazos, días de retraso).
	\item Calcular la fecha que resulta de sumar o restar $k$ días a una fecha dada.
	\item Determinar el día de la semana de cualquier fecha.
	\item Ordenar o comparar fechas usando el entero como clave (por ejemplo, en un \texttt{sort} o un \texttt{map}).
	\item Problemas de calendarios: contar cuántos viernes 13 hay en un rango, o el siguiente día que cumple una propiedad.
\end{itemize}

\textbf{Complejidad:} Cada conversión (fecha $\to$ entero y entero $\to$ fecha)
ejecuta un número fijo de sumas, multiplicaciones y divisiones enteras, sin
ciclos ni recursión, por lo que cuesta $O(1)$. La diferencia entre dos fechas
son dos conversiones más una resta: $O(1)$. Sumar $k$ días es una conversión,
una suma y la conversión inversa: $O(1)$, independiente de la magnitud de $k$.
Memoria: $O(1)$. Para $q$ consultas, el total es $O(q)$.

\begin{lstlisting}[style=cpp]
	struct Fecha {
		int anio, mes, dia;
		
		// CREATE: guarda anio, mes y dia (se asume una fecha valida).
		Fecha(int a = 1970, int m = 1, int d = 1) : anio(a), mes(m), dia(d) {} // O(1)
		
		// READ: indica si un anio es bisiesto (regla gregoriana).
		static bool es_bisiesto(int a) {
			return (a % 4 == 0 && a % 100 != 0) || a % 400 == 0;
		} // O(1)
		
		// READ: devuelve cuantos dias tiene el mes m del anio a.
		static int dias_del_mes(int a, int m) {
			static const int t[12] = {31,28,31,30,31,30,31,31,30,31,30,31};
			return (m == 2 && es_bisiesto(a)) ? 29 : t[m - 1];
		} // O(1)
		
		// READ: convierte la fecha a dias transcurridos desde 1970-01-01 (dia 0).
		// Se desplaza el anio para que empiece en marzo (febrero queda al final).
		long long a_dias() const {
			long long y = anio;
			y -= (mes <= 2);
			long long era = (y >= 0 ? y : y - 399) / 400;
			long long yoe = y - era * 400;                                   // [0, 399]
			long long mp  = (mes > 2) ? mes - 3 : mes + 9;                   // marzo = 0
			long long doy = (153 * mp + 2) / 5 + dia - 1;                    // dia del anio desplazado
			long long doe = yoe * 365 + yoe / 4 - yoe / 100 + doy;           // dia dentro de la era
			return era * 146097 + doe - 719468;
		} // O(1)
		
		// CREATE: construye la fecha correspondiente a z dias desde 1970-01-01.
		static Fecha desde_dias(long long z) {
			z += 719468;
			long long era = (z >= 0 ? z : z - 146096) / 146097;
			long long doe = z - era * 146097;                                // [0, 146096]
			long long yoe = (doe - doe / 1460 + doe / 36524 - doe / 146096) / 365;
			long long y   = yoe + era * 400;
			long long doy = doe - (365 * yoe + yoe / 4 - yoe / 100);
			long long mp  = (5 * doy + 2) / 153;
			int d = (int)(doy - (153 * mp + 2) / 5 + 1);
			int m = (int)(mp < 10 ? mp + 3 : mp - 9);
			return Fecha((int)(y + (m <= 2)), m, d);
		} // O(1)
		
		// READ: dias que hay de "otra" a esta fecha (positivo si esta es posterior).
		long long diferencia(const Fecha& otra) const {
			return a_dias() - otra.a_dias();
		} // O(1)
		
		// READ: dia de la semana (0 = domingo, 1 = lunes, ..., 6 = sabado).
		// 1970-01-01 fue jueves (4); el doble modulo cubre dias negativos.
		int dia_semana() const {
			return (int)(((a_dias() + 4) % 7 + 7) % 7);
		} // O(1)
		
		// WRITE: avanza (o retrocede si k < 0) la fecha k dias.
		void sumar_dias(long long k) {
			*this = desde_dias(a_dias() + k);
		} // O(1)
	};
	
	// ------------------------------------------------------------
	// MODIFICACIONES Y COMENTARIOS CON LOS USOS
	// ------------------------------------------------------------
	// 1. Restar dias: llamar sumar_dias con k negativo. No requiere codigo extra
	//    porque a_dias() y desde_dias() funcionan para z negativo.
	//
	// 2. Dia del anio (1..366): sirve para validar o para fechas ordinales.
	//    int dia_del_anio() const {
		//        return (int)(a_dias() - Fecha(anio, 1, 1).a_dias()) + 1;
		//    } // O(1)
	//
	// 3. Sumar meses (con recorte al ultimo dia valido del mes destino):
	//    void sumar_meses(int k) {
		//        long long total = (long long)anio * 12 + (mes - 1) + k;
		//        anio = (int)(total / 12); mes = (int)(total % 12) + 1;   // si total < 0, ajustar
		//        dia = min(dia, dias_del_mes(anio, mes));
		//    } // O(1)
	//    Ojo: 31 de enero + 1 mes da 28/29 de febrero, no 3 de marzo.
	//
	// 4. Siguiente dia de la semana d (0..6) estrictamente posterior:
	//    long long k = ((d - dia_semana() + 6) % 7) + 1;  sumar_dias(k);   // O(1)
	//
	// 5. Comparar fechas: comparar a_dias() (o la terna en orden anio, mes, dia).
	//    Como clave de map/sort: usar a_dias() como llave entera.
\end{lstlisting}

\textbf{Ejemplo de funcionamiento:}
\begin{lstlisting}[style=cpp]
	int main() {
		Fecha a(2000, 2, 28), b(2000, 3, 1);
		cout << b.diferencia(a) << "\n";            // 2  (28-feb -> 29-feb -> 1-mar, 2000 es bisiesto)
		
		Fecha c(2024, 1, 1), d(2024, 12, 31);
		cout << d.diferencia(c) << "\n";            // 365  (2024 tiene 366 dias; del 1-ene al 31-dic hay 365 saltos)
		
		Fecha e(2024, 1, 1);
		e.sumar_dias(100);
		cout << e.anio << "-" << e.mes << "-" << e.dia << "\n";   // 2024-4-10  (31 + 29 + 31 = 91 dias hasta el 1-abr; +9)
		
		cout << Fecha(2000, 1, 1).a_dias() << "\n";    // 10957  (30 anios * 365 + 7 bisiestos)
		cout << Fecha(2000, 1, 1).dia_semana() << "\n"; // 6  (sabado)
		cout << Fecha(2024, 1, 1).dia_semana() << "\n"; // 1  (lunes)
		return 0;
	}
\end{lstlisting}

\textbf{Nota:} \texttt{a\_dias()} devuelve \texttt{long long} porque con años
grandes (del orden de $10^9$) el conteo de días excede los 32 bits. Si el
problema garantiza años razonables (por ejemplo entre 1900 y 2100), \texttt{int}
alcanza, pero usar \texttt{long long} no cuesta nada.

\textbf{Regla práctica:} Si necesitas relacionar dos fechas que pueden estar en
meses o años distintos, convierte ambas a número de día y opera con enteros.
Simula día a día solo cuando el rango sea diminuto (unos pocos días) y no valga
la pena escribir la conversión. Usa la tabla de días acumulados por mes solo
cuando todas las fechas estén garantizadas dentro de un mismo año no bisiesto.

\subsection{Formatos de Hora (12h/24h, Segundos a HH:MM:SS)}

\textbf{Idea:} Igual que con las fechas, la forma más limpia de manejar horas es
guardarlas como un solo entero: los segundos transcurridos desde las 00:00:00.
Como un día tiene $24 \cdot 3600 = 86400$ segundos, cualquier hora del reloj
corresponde a un valor en $[0, 86399]$. El formato \texttt{HH:MM:SS} es
simplemente una representación en base mixta de ese entero: las horas son el
cociente de dividir entre $3600$, los minutos son el cociente del resto entre
$60$, y los segundos son lo que sobra. La conversión inversa es la suma
$3600h + 60m + s$.

Sumar tiempo es aritmética modular: el reloj es cíclico, así que después de
sumar $k$ segundos se toma el resultado módulo $86400$. Con $k$ negativo hay que
normalizar con un segundo módulo, porque en C++ el operador \texttt{\%} puede
devolver un valor negativo. El formato de 12 horas es solo una presentación
sobre las mismas horas: $h \bmod 12$ da la hora en el reloj (con el caso
especial de que $0$ se muestra como $12$), y $h < 12$ decide entre AM y PM.

\textbf{¿Por qué usarlo?} Manipular \texttt{h}, \texttt{m} y \texttt{s} por
separado obliga a propagar acarreos a mano (segundos $\to$ minutos $\to$
horas $\to$ día) y es la fuente típica de errores por uno. Con el entero, la
suma, la resta y la comparación de horas son operaciones directas, y el
formateo se hace solo al leer o imprimir. Esto es la misma idea que la
conversión de fechas a número de día, aplicada a un dominio cíclico y de tamaño
fijo. Una diferencia con las fechas es que aquí hay que decidir explícitamente
qué pasa al desbordar: en una \textit{hora del reloj} se envuelve módulo
$86400$, mientras que en una \textit{duración} no se envuelve y las horas pueden
superar $23$.

\textbf{Casos de uso:}
\begin{itemize}
	\item Imprimir un contador de segundos como \texttt{HH:MM:SS} (cronómetros, tiempos de carrera, logs).
	\item Convertir entre formato de 24 horas y 12 horas con AM/PM (entradas y salidas de problemas de formato).
	\item Sumar o restar segundos, minutos u horas a una hora del reloj, con vuelta al día siguiente o anterior.
	\item Calcular cuánto tiempo transcurre entre dos horas del reloj, incluso cruzando la medianoche.
	\item Aplicar desplazamientos de zona horaria a una hora dada, o formatear duraciones que superan las 24 horas.
\end{itemize}

\textbf{Complejidad:} Convertir segundos a \texttt{HH:MM:SS} son dos divisiones
y dos restos enteros: $O(1)$. La conversión inversa son dos multiplicaciones y
dos sumas: $O(1)$. Sumar $k$ segundos es una suma y un módulo: $O(1)$,
independiente de $k$. Formatear a texto es $O(1)$ porque la salida tiene
longitud fija (8 caracteres, más el sufijo AM/PM). Parsear una cadena es $O(L)$
con $L$ el largo del texto, que en la práctica es constante (una hora tiene a lo
sumo unos 11 caracteres). Memoria: $O(1)$.

\begin{lstlisting}[style=cpp]
	struct Hora {
		int h, m, s;                      // siempre en formato de 24 horas
		static constexpr int DIA = 86400; // segundos en un dia
		
		// CREATE: construye la hora a partir de una cantidad de segundos.
		// Normaliza modulo 86400 (funciona con valores negativos o mayores a un dia).
		Hora(long long seg = 0) {
			seg = ((seg % DIA) + DIA) % DIA;
			h = (int)(seg / 3600);
			m = (int)((seg % 3600) / 60);
			s = (int)(seg % 60);
		} // O(1)
		
		// CREATE: construye la hora a partir de horas, minutos y segundos.
		static Hora desde_hms(int hh, int mm, int ss) {
			return Hora(hh * 3600LL + mm * 60LL + ss);
		} // O(1)
		
		// CREATE: lee texto en formato 24h, ej. "07:45:10" o "7:5:3".
		static Hora desde_texto24(const string& t) {
			int hh = 0, mm = 0, ss = 0;
			sscanf(t.c_str(), "%d:%d:%d", &hh, &mm, &ss);
			return desde_hms(hh, mm, ss);
		} // O(L), L = largo del texto (constante en la practica)
		
		// CREATE: lee texto en formato 12h, ej. "07:45:10 PM".
		// 12 AM -> 0 horas, 12 PM -> 12 horas, 1..11 PM -> 13..23.
		static Hora desde_texto12(const string& t) {
			int hh = 0, mm = 0, ss = 0;
			char suf[3] = {0};
			sscanf(t.c_str(), "%d:%d:%d %2s", &hh, &mm, &ss, suf);
			hh %= 12;                              // 12 -> 0
			if (suf[0] == 'P' || suf[0] == 'p') hh += 12;
			return desde_hms(hh, mm, ss);
		} // O(L)
		
		// READ: devuelve los segundos transcurridos desde las 00:00:00.
		int a_segundos() const {
			return h * 3600 + m * 60 + s;
		} // O(1)
		
		// READ: devuelve la hora como texto "HH:MM:SS" (24 horas).
		string formato24() const {
			char buf[16];
			snprintf(buf, sizeof buf, "%02d:%02d:%02d", h, m, s);
			return string(buf);
		} // O(1)
		
		// READ: devuelve la hora como texto "HH:MM:SS AM/PM" (12 horas).
		string formato12() const {
			int h12 = h % 12;
			if (h12 == 0) h12 = 12;                // 0 y 12 se muestran como 12
			const char* suf = (h < 12) ? "AM" : "PM";
			char buf[24];
			snprintf(buf, sizeof buf, "%02d:%02d:%02d %s", h12, m, s, suf);
			return string(buf);
		} // O(1)
		
		// READ: segundos que hay que avanzar desde esta hora hasta "otra",
		// cruzando la medianoche si hace falta. Resultado en [0, 86399].
		int segundos_hasta(const Hora& otra) const {
			return ((otra.a_segundos() - a_segundos()) % DIA + DIA) % DIA;
		} // O(1)
		
		// WRITE: avanza (o retrocede si k < 0) la hora k segundos, con vuelta al dia.
		void sumar_segundos(long long k) {
			*this = Hora(a_segundos() + k);
		} // O(1)
	};
	
	// ------------------------------------------------------------
	// MODIFICACIONES Y COMENTARIOS CON LOS USOS
	// ------------------------------------------------------------
	// 1. Duraciones (sin envolver a 24h): NO usar el constructor de Hora, porque
	//    aplica modulo 86400. Formatear directamente el entero:
	//    string formato_duracion(long long seg) {
		//        char buf[32];
		//        snprintf(buf, sizeof buf, "%lld:%02lld:%02lld", seg / 3600, (seg % 3600) / 60, seg % 60);
		//        return string(buf);
		//    } // O(1)
	//    Las horas pueden pasar de 23: 100000 s -> "27:46:40"; 359999 s -> "99:59:59".
	//
	// 2. Milisegundos (HH:MM:SS.mmm): guardar el tiempo en milisegundos, con
	//    DIA = 86400000, y extraer ms = t % 1000, luego t /= 1000 y seguir igual.
	//
	// 3. Zona horaria: aplicar el desplazamiento en segundos con sumar_segundos.
	//    Ej. de UTC a UTC-5:  hora.sumar_segundos(-5 * 3600);   // O(1)
	//
	// 4. Lectura con cin en lugar de sscanf: leer "int h; char c; int m, s;"
	//    con cin >> h >> c >> m >> c >> s; (el char absorbe los ':').
	//
	// 5. Validar entrada: antes de construir, comprobar 0 <= h < 24, 0 <= m < 60,
	//    0 <= s < 60 si el enunciado no garantiza formato correcto.
\end{lstlisting}

\textbf{Ejemplo de funcionamiento:}
\begin{lstlisting}[style=cpp]
	int main() {
		Hora a(3661);
		cout << a.formato24() << "\n";                  // 01:01:01   (3600 + 60 + 1)
		
		Hora b(86399);
		cout << b.formato24() << "\n";                  // 23:59:59
		
		Hora c(45296);
		cout << c.formato24() << "\n";                  // 12:34:56   (12*3600 = 43200; 2096 = 34*60 + 56)
		cout << c.formato12() << "\n";                  // 12:34:56 PM
		
		cout << Hora::desde_hms(0, 15, 0).formato12() << "\n";   // 12:15:00 AM
		cout << Hora::desde_hms(15, 30, 0).formato12() << "\n";  // 03:30:00 PM
		
		Hora d = Hora::desde_texto12("07:45:10 PM");
		cout << d.formato24() << "\n";                  // 19:45:10
		
		Hora e = Hora::desde_hms(23, 59, 30);
		e.sumar_segundos(45);
		cout << e.formato24() << "\n";                  // 00:00:15   (86370 + 45 = 86415 = 86400 + 15)
		
		Hora f = Hora::desde_hms(22, 0, 0), g = Hora::desde_hms(6, 0, 0);
		cout << f.segundos_hasta(g) << "\n";            // 28800   (8 horas, cruzando la medianoche)
		return 0;
	}
\end{lstlisting}

\textbf{Nota:} Los casos límite del formato de 12 horas son la fuente de errores
más común. La medianoche (\texttt{00:xx:xx}) es \texttt{12:xx:xx AM} y el
mediodía (\texttt{12:xx:xx}) es \texttt{12:xx:xx PM}; no existe la hora
\texttt{00} en el reloj de 12 horas. Por eso la conversión usa
$h \bmod 12$ y sustituye $0$ por $12$, en lugar de restar $12$ a las horas
mayores que $12$.

\textbf{Regla práctica:} Si el problema mezcla horas con sumas, restas o
comparaciones, guarda todo como segundos (o milisegundos) y convierte solo al
leer o imprimir. Usa el constructor con módulo cuando modeles una hora del reloj
(que se envuelve en $24$h) y el formateo directo de la sección de modificaciones
cuando modeles una duración (que puede superar las $24$ horas). Si el problema
solo pide un cambio de formato sin aritmética, basta con las reglas de
\texttt{formato12()} y \texttt{desde\_texto12()}, sin necesidad de más
estructura.

\section{Parsing y Evaluación de Expresiones}

\subsection{Tokenizador}

\textbf{Idea:} Un tokenizador (o \textit{lexer}) convierte una cadena de
caracteres en una lista de \textit{tokens}: las unidades mínimas con significado
(números, identificadores, operadores, paréntesis). Es el primer paso de todo
evaluador, y separarlo del resto tiene una razón estructural. Cada clase de
token es un lenguaje regular que se reconoce mirando solo el primer carácter y
avanzando mientras se cumpla la regla de su clase: dígitos y punto para números,
letras, dígitos y guion bajo para identificadores, un solo carácter para
operadores y paréntesis. Se aplica la regla de \textit{máxima coincidencia}:
cada token consume el prefijo más largo posible (\texttt{123} es un token, no
tres). Así el analizador posterior ya no ve espacios ni números partidos, solo
una secuencia limpia.

Hay una ambigüedad que el tokenizador debe resolver con contexto: el signo
\texttt{-} puede ser resta (\texttt{a - b}) o negación (\texttt{-b}). La regla es
que un \texttt{-} es unario si no puede seguir a un operando, es decir, si el
token anterior no es un número, un identificador o un \texttt{)}. Con esa sola
comprobación, \texttt{3 + 2*-4} se separa correctamente. De forma análoga, un
identificador seguido de \texttt{(} se marca como función, lo que evita que el
siguiente paso tenga que mirar hacia adelante.

\textbf{¿Por qué usarlo?} Leer con \texttt{cin >> a >> op >> b} solo funciona si
la entrada viene separada por espacios; una expresión como
\texttt{3+4*(2-1)} no tiene un espacio donde cortar. Trabajar directamente sobre
los caracteres dentro del evaluador mezcla dos problemas (reconocer números y
respetar la precedencia) y llena el código de casos para espacios y números de
varias cifras. \texttt{std::regex} resuelve la tokenización, pero es lento y
más difícil de depurar que un escáner de 30 líneas. La lectura de romanos de la
subsección anterior es una tokenización mínima (un símbolo, un token); aquí
los tokens tienen longitud variable. El tokenizador no comprueba que la
expresión tenga sentido (\texttt{3 + * 4} se tokeniza sin error), y esa
responsabilidad pasa al algoritmo de la siguiente subsección.

\textbf{Casos de uso:}
\begin{itemize}
	\item Evaluar expresiones aritméticas dadas como texto (calculadoras, problemas de «evalúa la expresión»).
	\item Leer entradas con formato mixto donde los elementos pueden ir pegados (\texttt{f(x,y)=x+2}).
	\item Validar sintaxis: paréntesis balanceados, caracteres permitidos, números bien formados.
	\item Interpretar mini-lenguajes o comandos que dan los enunciados (simuladores de una máquina, macros, fórmulas de hojas de cálculo).
	\item De nicho: resaltado de sintaxis, contar identificadores distintos de un programa o minificar código.
\end{itemize}

\textbf{Complejidad:} El índice \texttt{pos} solo avanza. Cada carácter se
consume una vez en el ciclo principal, y el único retroceso aparente es la
mirada adelante tras un identificador, que recorre los espacios hasta el
siguiente carácter no blanco. Esos espacios se vuelven a visitar en el ciclo
principal, así que cada carácter se toca a lo sumo dos veces. Convertir un
número con \texttt{stod} cuesta proporcional a su largo, y la suma sobre todos
los números es a lo sumo $L$. Por tanto, el total es $O(L)$ con $L$ el largo de
la cadena. Memoria: $O(L)$, porque hay a lo sumo un token por carácter.

\begin{lstlisting}[style=cpp]
	struct Token {
		enum Tipo { NUM, ID, FUNC, OP, NEG, IZQ, DER, COMA };
		Tipo tipo;
		string txt;   // texto original del token
		double val;   // valor numerico (solo si es NUM)
		
		// CREATE: arma un token.
		Token(Tipo t = NUM, string x = "", double v = 0) : tipo(t), txt(x), val(v) {} // O(1)
		
		// READ: texto para imprimir ("neg" para el menos unario).
		string str() const {
			return tipo == NEG ? "neg" : txt;
		} // O(1)
	};
	
	struct Tokenizador {
		string s;  // expresion
		int n;     // largo de la expresion
		
		// CREATE: guarda la expresion a analizar.
		Tokenizador(const string& e) : s(e), n((int)e.size()) {} // O(L)
		
		// READ: indica si un caracter puede formar parte de un identificador.
		static bool es_ident(char c) {
			return isalnum((unsigned char)c) || c == '_';
		} // O(1)
		
		// READ: indica si el ultimo token emitido termina un operando (NUM, ID o ')').
		// Si es asi, un '-' que venga despues es resta; si no, es negacion.
		static bool termina_operando(const vector<Token>& r) {
			if (r.empty()) return false;
			Token::Tipo t = r.back().tipo;
			return t == Token::NUM || t == Token::ID || t == Token::DER;
		} // O(1)
		
		// READ: separa la expresion en tokens. Lanza runtime_error si hay algo invalido.
		vector<Token> tokenizar() const {
			vector<Token> r;
			int pos = 0;
			while (pos < n) {
				char c = s[pos];
				if (isspace((unsigned char)c)) { pos++; continue; }
				if (isdigit((unsigned char)c) || c == '.') {
					int ini = pos, puntos = 0;
					while (pos < n && (isdigit((unsigned char)s[pos]) || s[pos] == '.')) {
						if (s[pos] == '.') puntos++;
						pos++;
					}
					string t = s.substr(ini, pos - ini);
					if (puntos > 1 || t == ".") throw runtime_error("numero mal formado: " + t);
					r.push_back(Token(Token::NUM, t, stod(t)));
				} else if (isalpha((unsigned char)c) || c == '_') {
					int ini = pos;
					while (pos < n && es_ident(s[pos])) pos++;
					string t = s.substr(ini, pos - ini);
					int q = pos;                                   // mirar adelante: es funcion si sigue '('
					while (q < n && isspace((unsigned char)s[q])) q++;
					bool es_func = (q < n && s[q] == '(');
					r.push_back(Token(es_func ? Token::FUNC : Token::ID, t));
				} else if (c == '(') { r.push_back(Token(Token::IZQ, "(")); pos++; }
				else if (c == ')') { r.push_back(Token(Token::DER, ")")); pos++; }
				else if (c == ',') { r.push_back(Token(Token::COMA, ",")); pos++; }
				else if (string("+-*/^%").find(c) != string::npos) {
					bool operando_antes = termina_operando(r);
					if (!operando_antes && c == '-') r.push_back(Token(Token::NEG, "-"));
					else if (!operando_antes && c == '+') { /* '+' unario: se ignora */ }
					else r.push_back(Token(Token::OP, string(1, c)));
					pos++;
				} else {
					throw runtime_error(string("caracter invalido: ") + c);
				}
			}
			return r;
		} // O(L)
	};
	
	// ------------------------------------------------------------
	// MODIFICACIONES Y COMENTARIOS CON LOS USOS
	// ------------------------------------------------------------
	// 1. Notacion cientifica (1e-9, 2.5E+3): despues de leer los digitos, si s[pos] es
	//    'e' o 'E' y sigue un digito (o un signo y un digito), consumir tambien el
	//    exponente antes de llamar a stod:
	//    if (pos + 1 < n && (s[pos]=='e' || s[pos]=='E') &&
	//        (isdigit((unsigned char)s[pos+1]) ||
	//         ((s[pos+1]=='+' || s[pos+1]=='-') && pos + 2 < n && isdigit((unsigned char)s[pos+2])))) {
		//        pos += 2;
		//        while (pos < n && isdigit((unsigned char)s[pos])) pos++;
		//    } // O(largo del numero)
	//
	// 2. Operadores de varios caracteres (<=, >=, ==, !=, &&, ||): mirar tambien
	//    s[pos+1] antes de decidir (maxima coincidencia). Se prueba primero la
	//    pareja y, si no coincide, el caracter solo:
	//    if (pos + 1 < n && (s.substr(pos, 2) == "<=" || s.substr(pos, 2) == ">=" ||
	//        s.substr(pos, 2) == "==" || s.substr(pos, 2) == "&&")) {
		//        r.push_back(Token(Token::OP, s.substr(pos, 2))); pos += 2;
		//    } // O(1)
	//
	// 3. Multiplicacion implicita ("2x", "3(4+1)", "(a)(b)"): despues de tokenizar,
	//    insertar un '*' cuando un operando va seguido de algo que empieza un operando.
	//    vector<Token> r2;
	//    for (size_t i = 0; i < r.size(); i++) {
		//        if (i > 0) {
			//            Token::Tipo a = r[i-1].tipo, b = r[i].tipo;
			//            bool izq_op = (a == Token::NUM || a == Token::ID || a == Token::DER);
			//            bool der_op = (b == Token::NUM || b == Token::ID || b == Token::FUNC || b == Token::IZQ);
			//            if (izq_op && der_op) r2.push_back(Token(Token::OP, "*"));
			//        }
		//        r2.push_back(r[i]);
		//    } // O(n)
	//    Ojo: "x(3)" ya se toma como llamada a funcion x, no como producto.
	//
	// 4. Posicion en los errores: agregar un campo "int pos" al Token y guardarlo al
	//    crearlo; asi el mensaje puede decir "columna k" en lugar de solo el caracter.
	//
	// 5. Lectura perezosa: convertir tokenizar() en un metodo siguiente() que guarda
	//    pos y devuelve un token cada vez (con un token de mirada adelante). Es lo que
	//    usa un parser de descenso recursivo y evita construir el vector.
	//
	// 6. Cadenas y comentarios (para mini-lenguajes): si c == '"', leer hasta la
	//    comilla de cierre; si aparece "//", saltar hasta el fin de linea.
	//
	// 7. Cuando NO usar tokenizador: si los elementos de la entrada ya vienen
	//    separados por espacios, basta con "while (ss >> palabra)" sobre un
	//    stringstream. El tokenizador solo hace falta cuando los tokens pueden ir pegados.
\end{lstlisting}

\textbf{Ejemplo de funcionamiento:}
\begin{lstlisting}[style=cpp]
	void imprimir(const string& e) {
		try {
			vector<Token> v = Tokenizador(e).tokenizar();
			cout << v.size() << ":";
			for (auto& t : v) cout << " " << t.str();
			cout << "\n";
		} catch (exception& ex) {
			cout << "error: " << ex.what() << "\n";
		}
	}
	
	int main() {
		imprimir("3 + 4.5*(2-1)");   // 9: 3 + 4.5 * ( 2 - 1 )
		imprimir("-3 + 2*-4");       // 7: neg 3 + 2 * neg 4   (los '-' tras inicio y tras '*' son unarios)
		imprimir("sqrt(2)-1");       // 6: sqrt ( 2 ) - 1      (el '-' tras ')' es resta)
		imprimir("max(a1, b_2)");    // 6: max ( a1 , b_2 )
		imprimir("2.5.1 + 3");       // error: numero mal formado: 2.5.1
		imprimir("2 $ 3");           // error: caracter invalido: $
		return 0;
	}
\end{lstlisting}

\textbf{Nota:} Al llamar a \texttt{isspace}, \texttt{isdigit} y similares con un
\texttt{char}, conviene convertir primero a \texttt{unsigned char}, porque con
valores negativos (caracteres fuera de ASCII) el comportamiento es indefinido.
\texttt{stod} puede lanzar \texttt{out\_of\_range} con números gigantes, así que
capturar \texttt{exception} y no solo \texttt{runtime\_error} es más seguro.
Este tokenizador ignora el \texttt{+} unario en lugar de emitirlo, y trata
cualquier identificador seguido de \texttt{(} como función, aunque el nombre no
exista; esa validación queda para la evaluación.

\textbf{Regla práctica:} Si la entrada puede traer tokens pegados
(\texttt{3+4*(2-1)}) o de longitud variable, tokeniza primero y trabaja después
sobre la lista. Si los elementos ya vienen separados por espacios, un
\texttt{stringstream} es suficiente. Resuelve el \texttt{-} unario aquí, con la
regla del token anterior, y no en el evaluador. Y recuerda que el tokenizador
solo garantiza que los tokens son válidos, no que la expresión esté bien
formada: eso lo comprueba la etapa siguiente.

\subsection{Shunting-Yard y Notación Postfija}

\textbf{Idea:} La notación postfija (o polaca inversa) escribe cada operador
después de sus operandos: \texttt{3 + 4 * 2} se convierte en
\texttt{3 4 2 * +}. Su ventaja es que no tiene ambigüedad: el orden ya codifica
la precedencia, no necesita paréntesis, y se evalúa con una sola pila (al ver un
número se apila, al ver un operador se sacan los dos últimos valores, se
opera y se apila el resultado). El algoritmo \textit{shunting-yard}, de Dijkstra,
convierte una expresión infija a postfija con una pila de operadores.

Funciona por esta razón: un operador de la pila está \textit{esperando} a que
termine su operando derecho. Ese operando deja de poder crecer justo cuando
llega otro operador de precedencia menor (o igual, si es asociativo por la
izquierda), y solo entonces se puede emitir el operador pendiente. Por eso, al
llegar un operador $o_1$, se sacan de la pila todos los operadores $o_2$ con
$prec(o_2) > prec(o_1)$, o con igual precedencia cuando $o_1$ es asociativo por
la izquierda. Ese caso de igualdad es lo que distingue $a - b - c = (a-b)-c$
(se saca) de $a \wedge b \wedge c = a \wedge (b \wedge c)$ (no se saca). Un
paréntesis izquierdo es una barrera que ningún operador puede sacar, y el
derecho vacía la pila hasta encontrarla. Los operadores prefijos (la negación)
no tienen operando izquierdo que cerrar, así que se apilan directamente sin
sacar nada.

\textbf{¿Por qué usarlo?} Frente a evaluar con recursión sobre la cadena,
shunting-yard es dirigido por una tabla: agregar un operador nuevo es agregar
una línea de precedencia, no una función nueva. La forma postfija se puede
guardar y evaluar muchas veces con distintos valores de las variables sin volver
a analizar el texto, lo cual es lo que se quiere cuando se evalúa una fórmula
en miles de puntos. Frente al \textit{descenso recursivo}, que es más flexible
(sirve para gramáticas con condicionales, asignaciones o llamadas de aridad
variable, y da mejores mensajes de error), shunting-yard es más corto y
suficiente cuando el lenguaje es solo operadores, paréntesis y funciones. La
variante de dos pilas (una de operadores y otra de valores, que aplica cada
operador al sacarlo) evita construir la lista postfija, pero es menos reutilizable.
Trabaja sobre la lista de tokens de la subsección anterior, de modo que no
vuelve a preocuparse por los espacios ni por el signo unario.

\textbf{Casos de uso:}
\begin{itemize}
	\item Evaluar expresiones aritméticas con paréntesis, precedencia y asociatividad (calculadoras, problemas de «evaluar la expresión»).
	\item Evaluar una misma fórmula para muchos valores de una variable (tablas, gráficas, búsqueda binaria sobre una función dada como texto).
	\item Convertir entre notaciones: infija a postfija, o resolver problemas que dan directamente la forma postfija.
	\item Evaluar fórmulas booleanas (\texttt{\&\&}, \texttt{||}, \texttt{!}) sobre todas las asignaciones, cambiando solo la tabla de precedencias.
	\item De nicho: construir el árbol de la expresión para derivarla o simplificarla simbólicamente, o generar código para una máquina de pila.
\end{itemize}

\textbf{Complejidad:} Sea $n$ el número de tokens. En la conversión, cada token
entra a la pila a lo sumo una vez y sale a lo sumo una vez. Aunque el ciclo
\texttt{while} interno pueda sacar muchos operadores en una sola iteración, el
total de salidas sobre toda la ejecución no supera a $n$, así que la conversión
cuesta $O(n)$ amortizado. La evaluación procesa cada token postfijo una vez con
operaciones de pila $O(1)$; solo la búsqueda de una variable en el \texttt{map}
cuesta $O(\log V)$ con $V$ variables, para un total de $O(n \log V)$. Con la
tokenización, el proceso completo es $O(L + n \log V)$, y como $n \le L$ queda en
$O(L \log V)$. Memoria: $O(n)$ para la pila y la salida. Evaluar $q$ veces la
misma fórmula ya convertida cuesta $O(q \cdot n \log V)$, sin volver a
pagar la conversión.

\begin{lstlisting}[style=cpp]
	// Usa Token y Tokenizador de la subseccion anterior.
	struct ShuntingYard {
		map<string, double> vars; // valores de las variables
		
		// CREATE: evaluador sin variables definidas.
		ShuntingYard() {} // O(1)
		
		// WRITE: asigna (o cambia) el valor de una variable.
		void asignar(const string& nombre, double v) {
			vars[nombre] = v;
		} // O(log V)
		
		// READ: precedencia de un operador (mayor = se aplica antes).
		static int prec(const Token& t) {
			if (t.tipo == Token::NEG) return 3;
			char c = t.txt[0];
			if (c == '+' || c == '-') return 1;
			if (c == '*' || c == '/' || c == '%') return 2;
			return 4; // '^'
		} // O(1)
		
		// READ: indica si el operador es asociativo por la derecha.
		static bool asoc_der(const Token& t) {
			return t.tipo == Token::NEG || t.txt == "^";
		} // O(1)
		
		// READ: cantidad de argumentos de una funcion.
		static int aridad(const string& f) {
			return (f == "max" || f == "min") ? 2 : 1;
		} // O(1)
		
		// READ: convierte tokens en orden infijo a orden postfijo.
		vector<Token> a_postfija(const vector<Token>& tk) const {
			vector<Token> out, pila;
			for (const Token& t : tk) {
				switch (t.tipo) {
					case Token::NUM:
					case Token::ID:
					out.push_back(t);
					break;
					case Token::FUNC:
					case Token::IZQ:
					case Token::NEG:              // prefijo: se apila sin sacar nada
					pila.push_back(t);
					break;
					case Token::OP:
					while (!pila.empty() &&
					(pila.back().tipo == Token::OP || pila.back().tipo == Token::NEG) &&
					(prec(pila.back()) > prec(t) ||
					(prec(pila.back()) == prec(t) && !asoc_der(t)))) {
						out.push_back(pila.back());
						pila.pop_back();
					}
					pila.push_back(t);
					break;
					case Token::COMA:             // separa argumentos: cerrar el argumento actual
					while (!pila.empty() && pila.back().tipo != Token::IZQ) {
						out.push_back(pila.back());
						pila.pop_back();
					}
					if (pila.empty()) throw runtime_error("coma fuera de parentesis");
					break;
					case Token::DER:
					while (!pila.empty() && pila.back().tipo != Token::IZQ) {
						out.push_back(pila.back());
						pila.pop_back();
					}
					if (pila.empty()) throw runtime_error("parentesis ')' sin pareja");
					pila.pop_back();          // descarta el '('
					if (!pila.empty() && pila.back().tipo == Token::FUNC) {
						out.push_back(pila.back());   // la funcion sale despues de sus argumentos
						pila.pop_back();
					}
					break;
				}
			}
			while (!pila.empty()) {
				if (pila.back().tipo == Token::IZQ) throw runtime_error("parentesis '(' sin pareja");
				out.push_back(pila.back());
				pila.pop_back();
			}
			return out;
		} // O(n) amortizado
		
		// READ: evalua una lista postfija con una pila de valores.
		double evaluar_postfija(const vector<Token>& pf) const {
			vector<double> st;
			for (const Token& t : pf) {
				if (t.tipo == Token::NUM) {
					st.push_back(t.val);
				} else if (t.tipo == Token::ID) {
					auto it = vars.find(t.txt);
					if (it == vars.end()) throw runtime_error("variable sin valor: " + t.txt);
					st.push_back(it->second);
				} else if (t.tipo == Token::NEG) {
					if (st.empty()) throw runtime_error("expresion mal formada");
					st.back() = -st.back();
				} else if (t.tipo == Token::OP) {
					if (st.size() < 2) throw runtime_error("expresion mal formada");
					double b = st.back(); st.pop_back();
					double a = st.back(); st.pop_back();
					char c = t.txt[0];
					double r;
					if (c == '+') r = a + b;
					else if (c == '-') r = a - b;
					else if (c == '*') r = a * b;
					else if (c == '/') r = a / b;
					else if (c == '%') r = fmod(a, b);
					else r = pow(a, b);
					st.push_back(r);
				} else if (t.tipo == Token::FUNC) {
					int k = aridad(t.txt);
					if ((int)st.size() < k) throw runtime_error("faltan argumentos en " + t.txt);
					double b = st.back(); st.pop_back();
					double a = 0;
					if (k == 2) { a = st.back(); st.pop_back(); }
					double r;
					if (t.txt == "sqrt") r = sqrt(b);
					else if (t.txt == "abs") r = fabs(b);
					else if (t.txt == "max") r = max(a, b);
					else if (t.txt == "min") r = min(a, b);
					else throw runtime_error("funcion desconocida: " + t.txt);
					st.push_back(r);
				}
			}
			if (st.size() != 1) throw runtime_error("expresion mal formada");
			return st[0];
		} // O(n log V)
		
		// READ: tokeniza, convierte y evalua una expresion completa.
		double evaluar(const string& e) const {
			return evaluar_postfija(a_postfija(Tokenizador(e).tokenizar()));
		} // O(L log V)
	};
	
	// ------------------------------------------------------------
	// MODIFICACIONES Y COMENTARIOS CON LOS USOS
	// ------------------------------------------------------------
	// 1. Evaluar la misma formula muchas veces: convertir UNA vez y reutilizar la
	//    lista postfija, cambiando solo las variables.
	//    vector<Token> pf = sy.a_postfija(Tokenizador("x^2 + 1").tokenizar());   // O(n)
	//    for (double x = 0; x <= 10; x += 0.5) {
		//        sy.asignar("x", x);
		//        double r = sy.evaluar_postfija(pf);   // O(n log V), sin volver a analizar el texto
		//    }
	//
	// 2. Agregar operadores: basta ampliar prec() y el bloque de evaluacion. Un
	//    operador prefijo nuevo (por ejemplo '!' logico) se trata igual que NEG:
	//    se apila sin sacar, con asociatividad derecha y precedencia alta. Para
	//    operadores de comparacion o logicos (<, ==, &&, ||) hay que dar precedencias
	//    menores que las de '+' y tokenizarlos con la modificacion 2 del tokenizador.
	//
	// 3. Arbol de expresion: en lugar de evaluar, apilar nodos. Un numero o variable
	//    crea una hoja; un operador saca dos nodos y apila uno nuevo con ellos como hijos.
	//    struct Nodo { Token t; Nodo *izq = nullptr, *der = nullptr; };
	//    // al leer un OP:  Nodo* b = st.back(); st.pop_back(); Nodo* a = st.back(); st.pop_back();
	//    //                 st.push_back(new Nodo{t, a, b});
	//    // O(n). La raiz queda en st[0]. Recorrido en postorden = postfija; en preorden =
	//    // notacion prefija; en inorden con parentesis = infija completamente agrupada.
	//    // Util para derivar o simplificar simbolicamente.
	//
	// 4. Variante de dos pilas: no construir 'out'; cada vez que un operador saldria a
	//    la salida, aplicarlo de inmediato a la pila de valores. Se evalua en una sola
	//    pasada, pero no se puede reutilizar. O(n log V).
	//
	// 5. Valores exactos: cambiar double por Fraccion o BigInt en la pila de valores
	//    (subsecciones anteriores). Hay que reemplazar '^' por una potencia entera y
	//    decidir que pasa con la division por cero, que aqui da inf (o se puede lanzar error).
	//
	// 6. Funciones de aridad variable: guardar en la pila, junto a cada FUNC, un
	//    contador de argumentos que sube en cada COMA, y emitir ese numero junto con
	//    la funcion (por ejemplo, un token "max/3").
	//
	// 7. Entrada ya en postfija (problemas de RPN): saltarse a_postfija y leer los
	//    tokens con un stringstream, llamando directo a evaluar_postfija.
	//
	// 8. Cuando el lenguaje tiene condicionales, asignaciones o bloques, usar descenso
	//    recursivo: una funcion por nivel de precedencia (suma -> producto -> potencia
	//    -> atomo). Mismo O(n), mas flexible, y da mejores mensajes de error.
\end{lstlisting}

\textbf{Ejemplo de funcionamiento:}
\begin{lstlisting}[style=cpp]
	int main() {
		ShuntingYard sy;
		sy.asignar("x", 3);
		sy.asignar("y", 4);
		
		auto mostrar = [&](const string& e) {
			try {
				vector<Token> pf = sy.a_postfija(Tokenizador(e).tokenizar());
				double v = sy.evaluar_postfija(pf);
				for (auto& t : pf) cout << t.str() << " ";
				cout << "=> " << v << "\n";
			} catch (exception& ex) {
				cout << "error: " << ex.what() << "\n";
			}
		};
		
		mostrar("3 + 4 * 2");            // 3 4 2 * + => 11
		mostrar("(1 + 2) * 3");          // 1 2 + 3 * => 9
		mostrar("10 - 4 - 3");           // 10 4 - 3 - => 3      (resta asociativa por la izquierda)
		mostrar("2 ^ 3 ^ 2");            // 2 3 2 ^ ^ => 512     (potencia por la derecha: 2^9)
		mostrar("-2 ^ 2");               // 2 2 ^ neg => -4       (la potencia se aplica antes que la negacion)
		mostrar("2 ^ -1");               // 2 1 neg ^ => 0.5
		mostrar("sqrt(x^2 + y^2)");      // x 2 ^ y 2 ^ + sqrt => 5
		mostrar("sqrt(16) + max(2, 7)"); // 16 sqrt 2 7 max + => 11
		mostrar("2 * (3 + 4");           // error: parentesis '(' sin pareja
		mostrar("3 +");                  // error: expresion mal formada
		mostrar("z + 1");                // error: variable sin valor: z
		return 0;
	}
\end{lstlisting}

\textbf{Nota:} La precedencia de la negación es una decisión de diseño. Aquí
\texttt{-2\^{}2} da $-4$ (la convención matemática, la misma de Python), pero
otras herramientas, como Excel, lo evalúan como $(-2)^2 = 4$. Si el enunciado
usa la segunda convención, basta con darle a \texttt{NEG} una precedencia mayor
que la de \texttt{\^{}}. Además, la aridad de las funciones está fija por
nombre (\texttt{max} y \texttt{min} reciben dos argumentos, el resto uno); una
llamada con un número de argumentos distinto no se detecta en la conversión y se
descubre como «expresión mal formada» al evaluar. El operador \texttt{/} con
divisor cero devuelve infinito en lugar de lanzar un error.

\textbf{Regla práctica:} Si el lenguaje es solo operadores con precedencia,
paréntesis y funciones, usa tokenizador más shunting-yard. Convierte una sola
vez a postfija si vas a evaluar la misma fórmula muchas veces. Si solo necesitas
un valor una vez, la variante de dos pilas ahorra la lista intermedia. Si la
gramática incluye condicionales, asignaciones, bloques o funciones de aridad
variable, pasa a descenso recursivo. Si necesitas derivar o transformar la
expresión, construye el árbol (modificación 3) en lugar de evaluarla
directamente.
\subsection{Descenso Recursivo}

\textbf{Idea:} El descenso recursivo escribe el analizador como una función por
cada nivel de la gramática, y cada función llama a la del nivel siguiente, que
es el de mayor precedencia. Para las expresiones aritméticas la gramática es:

\begin{center}
	\texttt{expr} $\to$ \texttt{termino} $($ $($\texttt{+} $|$ \texttt{-}$)$ \texttt{termino} $)^*$ \\
	\texttt{termino} $\to$ \texttt{unario} $($ $($\texttt{*} $|$ \texttt{/} $|$ \texttt{\%}$)$ \texttt{unario} $)^*$ \\
	\texttt{unario} $\to$ \texttt{neg} \texttt{unario} $|$ \texttt{potencia} \\
	\texttt{potencia} $\to$ \texttt{atomo} $[$ \texttt{\^{}} \texttt{unario} $]$ \\
	\texttt{atomo} $\to$ \texttt{num} $|$ \texttt{id} $|$ \texttt{(} \texttt{expr} \texttt{)} $|$ \texttt{func(} \texttt{expr}$,\dots$ \texttt{)}
\end{center}

La precedencia no se guarda en ninguna tabla: sale de la \textit{profundidad} de
las llamadas. Para calcular \texttt{3 + 4 * 2}, \texttt{expr} llama a
\texttt{termino}, que a su vez llega hasta un \texttt{atomo} y devuelve $3$. Al
ver \texttt{+}, \texttt{expr} vuelve a llamar a \texttt{termino} para el
operando derecho, y esa llamada consume \texttt{4 * 2} completo antes de
devolver el control. El \texttt{*} se agrupa primero porque vive en una función
más profunda, que se resuelve antes de que \texttt{expr} pueda seguir.

La asociatividad la decide la forma de la regla. Un ciclo \texttt{while} en
\texttt{expr} y \texttt{termino} acumula de izquierda a derecha, y así
$10 - 4 - 3 = (10-4)-3$. Una llamada recursiva a la derecha en \texttt{potencia}
agrupa hacia la derecha, y así $2^{3^2} = 2^{(3^2)}$. La negación está encima
de \texttt{potencia}, así que \texttt{-2\^{}2} se lee como $-(2^2)$, y como el
exponente es un \texttt{unario}, \texttt{2\^{}-1} también funciona. Una regla
como $expr \to expr + termino$ (recursiva por la izquierda) provocaría
recursión infinita, y por eso se reescribe como ciclo.

\textbf{¿Por qué usarlo?} Frente a shunting-yard (subsección anterior), ambos
son $O(n)$ y aceptan las mismas expresiones, pero el descenso recursivo evalúa
directamente durante el análisis, sin lista postfija ni pila explícita de
operadores. Su ventaja real es la flexibilidad: cada regla es código
arbitrario, así que es natural aceptar funciones con cualquier número de
argumentos (\texttt{max(1, 5, 3)}), condicionales, asignaciones, bloques,
y dar mensajes de error que dicen qué se esperaba y dónde. Cada regla de la
gramática se traduce a una función casi mecánicamente, lo que facilita
razonar sobre la corrección. Su desventaja es que no guarda una forma
reutilizable: evaluar la misma fórmula para muchos valores de \texttt{x}
repite todo el análisis, salvo que se construya un árbol (modificación 3). Y
cada nivel de precedencia nuevo es una función nueva, mientras que en
shunting-yard es una línea de tabla (la modificación 2 muestra una versión
dirigida por tabla, llamada \textit{precedence climbing}).

\textbf{Casos de uso:}
\begin{itemize}
	\item Evaluar expresiones aritméticas con precedencia, paréntesis y funciones de aridad variable.
	\item Intérpretes de mini-lenguajes: asignaciones, condicionales, bucles y bloques definidos en el enunciado.
	\item Validar que un texto pertenece a una gramática (paréntesis balanceados, formatos como fechas, JSON o rutas), con mensajes de error precisos.
	\item Construir el árbol sintáctico de una expresión para derivarla o simplificarla.
	\item De nicho: analizadores de fórmulas de hojas de cálculo, de consultas o de expresiones regulares simples.
\end{itemize}

\textbf{Complejidad:} Sea $n$ el número de tokens. El índice \texttt{p} solo
avanza y ninguna función retrocede: cada regla decide qué hacer mirando
únicamente el token actual (gramática \textit{LL(1)}), por lo que no hay
retroceso ni trabajo repetido. Cada token se consume exactamente una vez. Un
operando pasa por una cantidad constante de niveles ($k = 5$ funciones), así
que el trabajo total es $O(k \cdot n) = O(n)$. La búsqueda de una variable en
el \texttt{map} suma $O(\log V)$, para un total de $O(n \log V)$, y con la
tokenización, $O(L \log V)$. Memoria: la lista de tokens es $O(n)$ y la
profundidad de la recursión es proporcional al anidamiento máximo de
paréntesis y negaciones, hasta $O(n)$ en el peor caso, como en
\texttt{((((\dots))))} o \texttt{-----\dots1}. Con miles de niveles no hay
problema, pero con cientos de miles la pila del sistema puede desbordarse.

\begin{lstlisting}[style=cpp]
	// Usa Token y Tokenizador de la subseccion "Tokenizador".
	struct DescensoRecursivo {
		vector<Token> t;           // tokens de la expresion actual
		size_t p;                  // posicion del token que se esta mirando
		map<string, double> vars;  // valores de las variables
		
		// CREATE: analizador vacio, sin variables.
		DescensoRecursivo() : p(0) {} // O(1)
		
		// WRITE: asigna (o cambia) el valor de una variable.
		void asignar(const string& nombre, double v) {
			vars[nombre] = v;
		} // O(log V)
		
		// READ: indica si el token actual es el operador binario c.
		bool es_op(char c) const {
			return p < t.size() && t[p].tipo == Token::OP && t[p].txt[0] == c;
		} // O(1)
		
		// WRITE: consume un token del tipo dado o lanza un error descriptivo.
		void esperar(Token::Tipo tipo, const string& que) {
			if (p >= t.size() || t[p].tipo != tipo)
			throw runtime_error("se esperaba " + que + " en el token " + to_string(p));
			p++;
		} // O(1)
		
		// READ: aplica una funcion a su lista de argumentos ya evaluados.
		static double aplicar(const string& f, const vector<double>& a) {
			if (f == "sqrt" || f == "abs") {
				if (a.size() != 1) throw runtime_error(f + " recibe 1 argumento");
				return f == "sqrt" ? sqrt(a[0]) : fabs(a[0]);
			}
			if (f == "max" || f == "min") {
				if (a.empty()) throw runtime_error(f + " necesita al menos 1 argumento");
				double r = a[0];
				for (double x : a) r = (f == "max") ? max(r, x) : min(r, x);
				return r;
			}
			throw runtime_error("funcion desconocida: " + f);
		} // O(#argumentos)
		
		// WRITE: atomo -> numero | variable | '(' expr ')' | funcion '(' args ')'.
		double atomo() {
			if (p >= t.size()) throw runtime_error("expresion incompleta");
			const Token& tk = t[p];
			if (tk.tipo == Token::NUM) { p++; return tk.val; }
			if (tk.tipo == Token::ID) {
				p++;
				auto it = vars.find(tk.txt);
				if (it == vars.end()) throw runtime_error("variable sin valor: " + tk.txt);
				return it->second;
			}
			if (tk.tipo == Token::IZQ) {
				p++;
				double v = expr();
				esperar(Token::DER, "')'");
				return v;
			}
			if (tk.tipo == Token::FUNC) {
				string nombre = tk.txt;
				p++;
				esperar(Token::IZQ, "'('");
				vector<double> args;
				if (p < t.size() && t[p].tipo != Token::DER) {
					args.push_back(expr());
					while (p < t.size() && t[p].tipo == Token::COMA) { p++; args.push_back(expr()); }
				}
				esperar(Token::DER, "')'");
				return aplicar(nombre, args);
			}
			throw runtime_error("token inesperado: " + tk.str() + " en el token " + to_string(p));
		} // O(n) en total
		
		// WRITE: potencia -> atomo [ '^' unario ]  (asociativa por la derecha).
		double potencia() {
			double base = atomo();
			if (es_op('^')) {
				p++;
				double e = unario();           // permite 2^-1 y 2^3^2
				return pow(base, e);
			}
			return base;
		} // O(n) en total
		
		// WRITE: unario -> neg unario | potencia.
		double unario() {
			if (p < t.size() && t[p].tipo == Token::NEG) { p++; return -unario(); }
			return potencia();
		} // O(n) en total
		
		// WRITE: termino -> unario (('*' | '/' | '%') unario)*  (por la izquierda).
		double termino() {
			double v = unario();
			while (es_op('*') || es_op('/') || es_op('%')) {
				char c = t[p].txt[0];
				p++;
				double r = unario();
				v = (c == '*') ? v * r : (c == '/') ? v / r : fmod(v, r);
			}
			return v;
		} // O(n) en total
		
		// WRITE: expr -> termino (('+' | '-') termino)*  (por la izquierda).
		double expr() {
			double v = termino();
			while (es_op('+') || es_op('-')) {
				char c = t[p].txt[0];
				p++;
				double r = termino();
				v = (c == '+') ? v + r : v - r;
			}
			return v;
		} // O(n) en total
		
		// WRITE: tokeniza y evalua una expresion completa (exige consumir todos los tokens).
		double evaluar(const string& e) {
			t = Tokenizador(e).tokenizar();
			p = 0;
			double v = expr();
			if (p != t.size()) throw runtime_error("sobran tokens desde el token " + to_string(p));
			return v;
		} // O(L log V)
	};
	
	// ------------------------------------------------------------
	// MODIFICACIONES Y COMENTARIOS CON LOS USOS
	// ------------------------------------------------------------
	// 1. Agregar un nivel de precedencia: escribir una funcion nueva ENTRE dos
	//    existentes y cambiar quien la llama. Ej. comparaciones (menor prioridad que '+'):
	//    double comparacion() {
		//        double v = expr();
		//        if (es_op('<')) { p++; return v < expr(); }
		//        return v;
		//    } // O(n)
	//    y evaluar() llama a comparacion() en lugar de expr(). Requiere que el
	//    tokenizador emita '<' (ver modificacion 2 del tokenizador).
	//
	// 2. Precedence climbing (una sola funcion dirigida por tabla, como shunting-yard):
	//    double binaria(int min_prec) {
		//        double lhs = unario_pc();
		//        while (p < t.size() && t[p].tipo == Token::OP && prec(t[p]) >= min_prec) {
			//            Token op = t[p]; p++;
			//            double rhs = binaria(asoc_der(op) ? prec(op) : prec(op) + 1);
			//            lhs = aplicar_op(op, lhs, rhs);
			//        }
		//        return lhs;
		//    } // O(n)
	//    Con prec/asoc_der de la subseccion de shunting-yard. La negacion se procesa
	//    como: si es NEG, p++, return -binaria(prec_neg); asi -2^2 sigue dando -4
	//    porque '^' (prec 4) >= prec_neg (3) queda dentro del operando.
	//
	// 3. Arbol de sintaxis: en lugar de devolver double, cada funcion devuelve Nodo*.
	//    struct Nodo { Token op; Nodo *izq = nullptr, *der = nullptr; };
	//    // en expr(): v = new Nodo{tokenOp, v, r};   // en vez de v = v + r
	//    // Se recorre despues (evaluar, derivar, simplificar) sin volver a analizar el
	//    // texto: O(n) para construir, O(n) por cada evaluacion.
	//
	// 4. Asignaciones y sentencias ("x = 5; y = x * 2"): agregar una regla
	//    sentencia -> ID '=' expr | expr, mirando un token adelante (t[p+1] es '=').
	//    Como cada regla decide con UN token de mirada, sigue siendo O(n).
	//    Requiere que el tokenizador emita '=' y ';'.
	//
	// 5. Condicional (cond ? a : b): una regla por encima de la comparacion.
	//    double ternario() {
		//        double c = comparacion();
		//        if (p < t.size() && t[p].txt == "?") {
			//            p++; double a = ternario(); esperar_texto(":"); double b = ternario();
			//            return c != 0 ? a : b;    // evalua ambas ramas; para evitar efectos usar arbol
			//        }
		//        return c;
		//    } // O(n)
	//
	// 6. Gramaticas con recursion por la izquierda: NO se pueden escribir tal cual
	//    (recursion infinita). Se transforman:  A -> A x | y   ==>   A -> y (x)*,
	//    que es lo que hacen los ciclos while de expr() y termino().
	//
	// 7. Profundidad de recursion: para entradas con anidamiento de 1e5 o mas,
	//    la pila del sistema puede agotarse. Opciones: ampliar la pila (en Linux,
	//    ulimit -s unlimited), pasar a shunting-yard (pila propia en el heap), o
	//    poner un limite y lanzar un error si el anidamiento lo supera.
	//
	// 8. Analizar directamente sobre caracteres, sin tokenizador: es viable para
	//    gramaticas cortas. Se mantiene un indice p sobre la cadena, un saltar_blancos()
	//    al inicio de cada regla, y numero() lee los digitos con stod. Mismo O(L).
	//
	// 9. Backtracking / gramaticas ambiguas: si una regla necesita probar una
	//    alternativa y volver atras, el costo puede volverse exponencial. Se
	//    arregla con memoizacion por (regla, posicion), tecnica llamada "packrat",
	//    con O(n * #reglas). Para expresiones aritmeticas no hace falta.
\end{lstlisting}

\textbf{Ejemplo de funcionamiento:}
\begin{lstlisting}[style=cpp]
	int main() {
		DescensoRecursivo dr;
		dr.asignar("x", 3);
		dr.asignar("y", 4);
		
		auto mostrar = [&](const string& e) {
			try {
				cout << dr.evaluar(e) << "\n";
			} catch (exception& ex) {
				cout << "error: " << ex.what() << "\n";
			}
		};
		
		mostrar("3 + 4 * 2");                 // 11    (el '*' esta en un nivel mas profundo)
		mostrar("(1 + 2) * 3");               // 9
		mostrar("10 - 4 - 3");                // 3     (el ciclo acumula por la izquierda: (10-4)-3)
		mostrar("2 ^ 3 ^ 2");                 // 512   (recursion por la derecha: 2^(3^2) = 2^9)
		mostrar("-2 ^ 2");                    // -4    (unario esta encima de potencia: -(2^2))
		mostrar("2 ^ -1");                    // 0.5
		mostrar("1 - -2");                    // 3     (el segundo '-' es unario)
		mostrar("sqrt(x^2 + y^2)");           // 5
		mostrar("max(1, 5, 3) + min(4, 2)");  // 7     (aridad variable: 5 + 2)
		mostrar("2 * (3 + 4");                // error: se esperaba ')' en el token 6
		mostrar("3 +");                       // error: expresion incompleta
		mostrar("3 4");                       // error: sobran tokens desde el token 1
		mostrar("z + 1");                     // error: variable sin valor: z
		return 0;
	}
\end{lstlisting}

\textbf{Nota:} El tokenizador ya distingue la negación de la resta, y el
analizador no lo vuelve a decidir: le basta con ver un token \texttt{NEG} en
\texttt{unario}. Cada función devuelve un \texttt{double} y por eso los errores
se reportan lanzando excepciones; al usarlas, la posición reportada es el
índice del token, no la columna del texto (para eso, la modificación 4 del
tokenizador agrega el campo \texttt{pos}). Además, \texttt{evaluar} modifica
\texttt{t} y \texttt{p}, así que un mismo objeto no es reentrante: para
analizar en paralelo o desde dentro de una función de usuario, usa un objeto
por hilo. Al igual que en shunting-yard, la división entre cero devuelve
infinito en lugar de lanzar error.

\textbf{Regla práctica:} Si la entrada es solo operadores, paréntesis y
funciones de aridad fija, y vas a evaluar la misma fórmula muchas veces,
shunting-yard con la lista postfija reutilizable es más corto. Si el
lenguaje tiene funciones de aridad variable, asignaciones, condicionales,
bloques, o necesitas buenos mensajes de error, usa descenso recursivo: una
función por nivel de precedencia, ciclos para la asociatividad izquierda y
recursión para la derecha. Si además vas a evaluar la misma expresión
repetidas veces, haz que las funciones construyan un árbol y evalúalo aparte
(modificación 3). Si el anidamiento puede pasar de cientos de miles de
niveles, evita la recursión y usa shunting-yard.

\section{Técnicas Ad-hoc}

\subsection{Compresión de Coordenadas}

\textbf{Idea:} Muchos algoritmos solo necesitan saber qué valor es menor, igual o mayor que otro, y no cuánto valen. Un arreglo de frecuencias, un árbol de Fenwick o un barrido necesitan además índices pequeños y consecutivos. La compresión de coordenadas sustituye cada valor por su \textit{rango} dentro del conjunto ordenado de valores distintos: $\{-7, 5, 100, 10^9\}$ pasa a $\{0, 1, 2, 3\}$. Con $n$ valores el resultado siempre cabe en $[0, n)$, sin importar lo grandes o negativos que fueran los originales.

Funciona porque la función $x \mapsto \text{rango}(x)$ es una biyección estrictamente creciente entre los valores usados y $[0, k)$, donde $k$ es la cantidad de valores distintos. Al ser creciente, conserva todas las comparaciones ($x < y \iff \text{rango}(x) < \text{rango}(y)$) y la igualdad, así que cualquier algoritmo que solo dependa del orden da el mismo resultado sobre los rangos. Lo que se pierde son las distancias entre valores, como se explica en las modificaciones. La implementación es ordenar, eliminar duplicados y usar búsqueda binaria: el rango de $x$ es su posición en el arreglo ordenado.

\textbf{¿Por qué usarlo?} Si los valores llegan a $10^9$ pero solo hay $10^5$ de ellos, un arreglo indexado por valor es inviable en memoria. Un \texttt{map} o \texttt{unordered\_map} sirve para contar, pero no permite usar estructuras indexadas por posición (Fenwick, árbol de segmentos, arreglos de diferencias), que son las que hacen rápidas las consultas de rango. Un \texttt{unordered\_map} tampoco conserva el orden, así que solo sirve cuando basta la identidad de los valores (modificación 5). La compresión es la elección correcta cuando los valores son grandes o dispersos, se conocen todos de antemano (fuera de línea) y solo importa su orden relativo. Si los valores ya son pequeños (por ejemplo, hasta $10^6$), es un paso innecesario. Si las consultas llegan en línea, no se puede comprimir de antemano y hay que usar un árbol de segmentos dinámico o un \texttt{map} (modificación 7).

\textbf{Casos de uso:}
\begin{itemize}
	\item Usar un Fenwick o un árbol de segmentos sobre valores grandes (contar inversiones, subsecuencia creciente más larga, consultas de rango sobre valores).
	\item Contar frecuencias o hacer sumas por prefijos sobre valores dispersos, o renumerar identificadores de nodos de un grafo a $0, \dots, n-1$.
	\item Barridos de línea: unión de intervalos, cobertura de segmentos y eventos ordenados en el tiempo.
	\item Puntos en el plano con consultas fuera de línea: se comprime cada eje por separado y se trabaja sobre una cuadrícula pequeña.
	\item De nicho: convertir una permutación arbitraria en una permutación de $0, \dots, n-1$, o comprimir cadenas y pares (cualquier tipo comparable).
\end{itemize}

\textbf{Complejidad:} Sea $n$ la cantidad de valores. La construcción ordena en $O(n \log n)$ y \texttt{unique} recorre el arreglo una vez en $O(n)$, así que domina el ordenamiento. Cada consulta de rango (\texttt{indice}, \texttt{contar\_menores}) es una búsqueda binaria en $O(\log k) \subseteq O(\log n)$. Comprimir un arreglo de $m$ valores cuesta $O(m \log n)$. Si esos $m$ valores ya vienen ordenados, un doble puntero baja el costo a $O(n + m)$ (modificación 8). Agregar valores nuevos vuelve a ordenar y cuesta $O(n \log n)$. Memoria: $O(n)$.

\begin{lstlisting}[style=cpp]
	struct Compresor {
		vector<long long> vals; // valores distintos, ordenados de menor a mayor
		
		// CREATE: ordena y elimina duplicados; el rango de x es su posicion en vals.
		Compresor(vector<long long> v) {
			vals = move(v);
			sort(vals.begin(), vals.end());
			vals.erase(unique(vals.begin(), vals.end()), vals.end());
		} // O(n log n)
		
		// READ: cantidad de valores distintos (tamano del universo comprimido).
		int tam() const {
			return (int)vals.size();
		} // O(1)
		
		// READ: valor original que corresponde al rango i.
		long long valor(int i) const {
			return vals[i];
		} // O(1)
		
		// READ: indica si x estaba entre los valores originales.
		bool existe(long long x) const {
			auto it = lower_bound(vals.begin(), vals.end(), x);
			return it != vals.end() && *it == x;
		} // O(log n)
		
		// READ: rango de x (requiere que x exista; si no, devuelve su posicion de insercion).
		int indice(long long x) const {
			return (int)(lower_bound(vals.begin(), vals.end(), x) - vals.begin());
		} // O(log n)
		
		// READ: cuantos valores distintos son estrictamente menores que x (x puede no existir).
		int contar_menores(long long x) const {
			return (int)(lower_bound(vals.begin(), vals.end(), x) - vals.begin());
		} // O(log n)
		
		// READ: cuantos valores distintos son menores o iguales que x (x puede no existir).
		int contar_menores_o_iguales(long long x) const {
			return (int)(upper_bound(vals.begin(), vals.end(), x) - vals.begin());
		} // O(log n)
		
		// READ: devuelve el arreglo v con cada valor sustituido por su rango.
		vector<int> comprimir(const vector<long long>& v) const {
			vector<int> r(v.size());
			for (size_t i = 0; i < v.size(); i++) r[i] = indice(v[i]);
			return r;
		} // O(m log n)
		
		// WRITE: agrega valores nuevos al universo. Los rangos ya entregados dejan de valer.
		void agregar_valores(const vector<long long>& extra) {
			vals.insert(vals.end(), extra.begin(), extra.end());
			sort(vals.begin(), vals.end());
			vals.erase(unique(vals.begin(), vals.end()), vals.end());
		} // O(n log n)
	};
	
	// ------------------------------------------------------------
	// MODIFICACIONES Y COMENTARIOS CON LOS USOS
	// ------------------------------------------------------------
	// 1. Intervalos semiabiertos [l, r) y barrido: agregar al compresor TODOS los l y r.
	//    La celda i representa [vals[i], vals[i+1]) y mide vals[i+1] - vals[i].
	//    Longitud de la union de segmentos con arreglo de diferencias:
	//    vector<int> dif(c.tam() + 1);
	//    for (auto [l, r] : segs) { dif[c.indice(l)]++; dif[c.indice(r)]--; }
	//    long long total = 0; int cubierto = 0;
	//    for (int i = 0; i + 1 < c.tam(); i++) {
		//        cubierto += dif[i];
		//        if (cubierto > 0) total += c.valor(i + 1) - c.valor(i);
		//    } // O(n log n)
	//
	// 2. Rango con empates distintos (convertir en permutacion de 0..n-1): ordenar los
	//    indices por valor y, si hay empate, por posicion original.
	//    vector<int> ord(n), rango(n);
	//    iota(ord.begin(), ord.end(), 0);
	//    sort(ord.begin(), ord.end(), [&](int i, int j) { return v[i] < v[j] || (v[i] == v[j] && i < j); });
	//    for (int k = 0; k < n; k++) rango[ord[k]] = k;   // O(n log n)
	//
	// 3. Puntos en 2D: un compresor por eje. (x, y) pasa a (cx.indice(x), cy.indice(y)).
	//    Se conserva el orden en cada eje, pero no las distancias ni las diagonales.
	//
	// 4. Cualquier tipo comparable (strings, pairs): hacer la estructura plantilla,
	//    template<class T> struct CompresorT { vector<T> vals; ... };
	//    y usar el mismo sort + unique + lower_bound.
	//
	// 5. Si solo importa la IDENTIDAD (no el orden), basta un mapa hash y es O(n) esperado:
	//    unordered_map<long long, int> id;
	//    for (long long x : v) id.emplace(x, (int)id.size());   // no inserta si ya existe
	//    Con concursos que permiten ataques al hash, usar un hash con semilla aleatoria.
	//
	// 6. Se pierden las distancias y la adyacencia. Si el problema depende de que x y x+1
	//    sean vecinos (huecos, "siguiente entero libre"), agregar x-1, x y x+1 antes de
	//    comprimir. Si depende de la distancia (areas, longitudes), conservar vals y usar
	//    vals[i+1] - vals[i] como peso de la celda i.
	//
	// 7. Consultas en linea (no se conocen de antemano): no se puede comprimir. Opciones:
	//    arbol de segmentos dinamico (nodos creados bajo demanda, O(log C) por operacion,
	//    con C el rango de valores) o un map / arbol balanceado.
	//
	// 8. Consultas ya ordenadas: cambiar lower_bound por un puntero que solo avanza;
	//    el costo baja de O(m log n) a O(n + m).
\end{lstlisting}

\textbf{Ejemplo de funcionamiento:}
\begin{lstlisting}[style=cpp]
	int main() {
		vector<long long> v = {100, 5, 100, 1000000000, 5, -7};
		Compresor c(v);
		cout << c.tam() << "\n";                        // 4  (valores distintos: -7, 5, 100, 1000000000)
		
		for (int x : c.comprimir(v)) cout << x << " ";
		cout << "\n";                                   // 2 1 2 3 1 0
		
		cout << c.indice(5) << " " << c.existe(6) << "\n";   // 1 0
		cout << c.contar_menores(6) << "\n";            // 2  (-7 y 5 son menores que 6)
		cout << c.valor(3) << "\n";                     // 1000000000
		
		// Frecuencias sobre el universo comprimido (arreglo de 4 posiciones, no de 10^9)
		vector<int> cnt(c.tam(), 0);
		for (long long x : v) cnt[c.indice(x)]++;
		for (int x : cnt) cout << x << " ";
		cout << "\n";                                   // 1 2 2 1
		
		// Cuantos elementos de v estan en [0, 100]: celdas de contar_menores(0) a contar_menores_o_iguales(100)-1
		int lo = c.contar_menores(0), hi = c.contar_menores_o_iguales(100) - 1;   // lo = 1, hi = 2
		int total = 0;
		for (int i = lo; i <= hi; i++) total += cnt[i];
		cout << total << "\n";                          // 4  (100, 5, 100, 5)
		
		c.agregar_valores({0});
		cout << c.tam() << " " << c.indice(5) << "\n";  // 5 2  (el rango de 5 cambio al agregar el 0)
		return 0;
	}
\end{lstlisting}

\textbf{Nota:} \texttt{lower\_bound} exige un arreglo ordenado, y \texttt{unique} solo elimina duplicados \textit{consecutivos}, así que el orden \texttt{sort} $\to$ \texttt{unique} $\to$ \texttt{erase} es obligatorio y no se puede alterar. Además, \texttt{indice} no falla cuando el valor no existe: devuelve el lugar donde se insertaría, y eso puede ser un rango válido de otro valor. Si la entrada no garantiza que $x$ esté, comprueba con \texttt{existe} primero. Y como \texttt{agregar\_valores} cambia los rangos (como muestra el último \texttt{cout} del ejemplo), hay que agregar todo lo necesario \textit{antes} de comenzar a comprimir.

\textbf{Regla práctica:} Si los valores son grandes o dispersos, se conocen todos de antemano y el algoritmo solo compara, comprime y trabaja con arreglos de tamaño $k$. Si solo necesitas saber si dos valores son iguales, un mapa hash es más simple. Si los valores ya son pequeños (hasta unos $10^6$), usa el arreglo directo. Si el problema usa distancias o adyacencia (intervalos, áreas), conserva \texttt{vals} y agrega los vecinos necesarios. Si las consultas son en línea, cambia a un árbol de segmentos dinámico o a un \texttt{map}.

\subsection{Invariantes y Paridad}

\textbf{Idea:} Un \textit{invariante} es una propiedad del estado que toda operación permitida conserva. Si el estado inicial y el objetivo no comparten esa propiedad, el objetivo es inalcanzable, sin importar cuántas operaciones se apliquen ni en qué orden. La justificación es una inducción sobre el número de operaciones: la propiedad vale al inicio, cada paso la conserva, así que vale en todo estado alcanzable. El caso más frecuente es la \textit{paridad}: una cantidad módulo $2$ que no cambia. Por ejemplo, si una operación reemplaza $a$ y $b$ por $|a-b|$, la suma cambia de $a+b$ a $|a-b|$, y ambas difieren en $2\min(a,b)$, así que la paridad de la suma es invariante. Si reemplaza $a$ y $b$ por $a \oplus b$, el XOR total es invariante.

Otros invariantes clásicos son la suma módulo $m$, el \texttt{mcd} de todos los elementos (si la operación es $a \to a-b$), la coloración de un tablero (cada dominó cubre una casilla de cada color) y la paridad de una permutación. Esta última sale de contar ciclos: un intercambio de dos posiciones fusiona dos ciclos o parte uno en dos, así que cambia la cantidad de ciclos exactamente en $\pm 1$. La identidad de $n$ elementos tiene $n$ ciclos, así que la cantidad de intercambios necesarios para pasar de una permutación a otra siempre tiene la paridad de $n - \text{ciclos}$. Un invariante solo demuestra que algo es \textit{imposible}; para demostrar que es posible hace falta una construcción, o un conjunto de invariantes que sea completo (como en el puzzle de 15, más abajo).

\textbf{¿Por qué usarlo?} La alternativa a un invariante es explorar los estados con BFS o retroceso, y el espacio de estados suele ser exponencial: inviable para $n$ del orden de $10^5$ o para valores grandes. El invariante responde «imposible» en $O(1)$ o $O(n)$ sin recorrer nada. Frente a la programación dinámica, que exige que el estado quepa en una tabla, el invariante no necesita estado alguno. Su costo es que no da la secuencia de operaciones y que hay que \textit{descubrir} el invariante: no lo produce ningún algoritmo. Una estrategia práctica es escribir una fuerza bruta (BFS) para casos pequeños, agrupar los estados alcanzables entre sí y buscar qué propiedad reparte los estados en esos grupos (modificación 5). Se usa cuando el enunciado pide «¿es posible?», «imprime $-1$ si no se puede» o «¿quién gana?».

\textbf{Casos de uso:}
\begin{itemize}
	\item Decidir si un estado inicial puede transformarse en otro con operaciones dadas (respuesta \texttt{-1} o \texttt{NO}) mediante la paridad de la suma, el XOR, el \texttt{mcd} o la suma módulo $m$.
	\item Predecir el valor final tras repetir una operación sobre un multiconjunto (el último número que queda al reemplazar pares por su diferencia o su XOR).
	\item Cubrimientos de tableros con fichas: coloración para descartar dominós o L-triminós imposibles.
	\item Permutaciones: saber si el número mínimo de intercambios es par o impar, o si una configuración de un puzzle deslizante es resoluble.
	\item De nicho: monovariantes (una cantidad que solo crece o solo decrece) para probar que un proceso termina, o usar el invariante para guiar una construcción.
\end{itemize}

\textbf{Complejidad:} Construir la estructura recorre el arreglo una vez y calcula suma, XOR y cantidad de impares: $O(n)$. Actualizar un elemento con \texttt{cambiar} solo ajusta esos tres valores con el valor viejo y el nuevo, así que cuesta $O(1)$, y consultarlos también. El \texttt{mcd} total hace $n$ llamadas a Euclides con $O(\log M)$ pasos cada una, con $M$ el mayor valor: $O(n \log M)$. La paridad de una permutación recorre cada ciclo una vez y marca cada elemento visitado, así que cada elemento se visita una sola vez: $O(n)$. \texttt{puzzle\_resoluble} sobre un tablero $N \times N$ tiene $n = N^2 - 1$ fichas y usa esa rutina, por lo que cuesta $O(N^2)$; contar las inversiones directamente sería $O(N^4)$. Memoria: $O(n)$.

\begin{lstlisting}[style=cpp]
	struct Invariantes {
		vector<long long> a; // estado actual
		int par_suma;        // paridad de la suma (se guarda solo el bit para evitar desbordes)
		long long xr;        // XOR de todos los elementos
		int impares;         // cantidad de elementos impares
		
		// CREATE: calcula los invariantes del estado inicial.
		Invariantes(vector<long long> v) : a(move(v)), par_suma(0), xr(0), impares(0) {
			for (long long x : a) {
				par_suma ^= (int)(x & 1);
				xr ^= x;
				impares += (int)(x & 1);
			}
		} // O(n)
		
		// WRITE: cambia a[i] por x y actualiza los invariantes sin recorrer el arreglo.
		void cambiar(int i, long long x) {
			par_suma ^= (int)((a[i] ^ x) & 1);
			xr ^= a[i] ^ x;
			impares += (int)(x & 1) - (int)(a[i] & 1);
			a[i] = x;
		} // O(1)
		
		// READ: paridad de la suma actual (0 = par, 1 = impar).
		int paridad_suma() const {
			return par_suma;
		} // O(1)
		
		// READ: XOR de todos los elementos.
		long long xor_total() const {
			return xr;
		} // O(1)
		
		// READ: cantidad de elementos impares.
		int cantidad_impares() const {
			return impares;
		} // O(1)
		
		// READ: mcd de todos los elementos (0 si todos son 0).
		long long mcd_total() const {
			long long g = 0;
			for (long long x : a) g = __gcd(g, x < 0 ? -x : x);
			return g;
		} // O(n log M)
		
		// READ: paridad de una permutacion de 0..n-1 (0 = par, 1 = impar).
		// Un intercambio cambia la cantidad de ciclos en +-1; la identidad tiene n ciclos.
		static int paridad_permutacion(const vector<int>& p) {
			int n = (int)p.size(), ciclos = 0;
			vector<char> vis(n, 0);
			for (int i = 0; i < n; i++) {
				if (vis[i]) continue;
				ciclos++;
				for (int j = i; !vis[j]; j = p[j]) vis[j] = 1;
			}
			return (n - ciclos) & 1;
		} // O(n)
		
		// READ: indica si un puzzle deslizante N x N (0 = hueco, fichas 1..N*N-1 leidas por filas)
		// puede llegar al estado 1, 2, ..., N*N-1, 0 (hueco abajo a la derecha).
		static bool puzzle_resoluble(const vector<int>& t, int N) {
			vector<int> q;
			int fila_hueco = 0;
			for (int i = 0; i < N * N; i++) {
				if (t[i] == 0) fila_hueco = i / N;   // fila contada desde arriba, desde 0
				else q.push_back(t[i] - 1);
			}
			int inv = paridad_permutacion(q);        // paridad de inversiones = paridad de la permutacion
			if (N % 2 == 1) return inv == 0;
			return (inv + fila_hueco) % 2 == 1;
		} // O(N^2)
	};
	
	// ------------------------------------------------------------
	// MODIFICACIONES Y COMENTARIOS CON LOS USOS
	// ------------------------------------------------------------
	// 1. Invariante modulo m: si una operacion suma o resta multiplos de m a un elemento,
	//    su residuo modulo m no cambia. El estado b es alcanzable desde a solo si
	//    a[i] % m == b[i] % m para todo i (con ((x % m) + m) % m para los negativos).
	//    Es necesario, y suficiente si se puede sumar o restar m a cualquier elemento.
	//
	// 2. XOR: si la operacion reemplaza a y b por a ^ b, el XOR total es constante, y
	//    el ultimo elemento que queda es exactamente xor_total(). O(1) por consulta.
	//
	// 3. Coloracion de tableros (cubrir con dominos): contar casillas libres por color.
	//    int blancas = 0, negras = 0;
	//    for (cada casilla libre (f, c)) ((f + c) & 1 ? negras : blancas)++;
	//    if (blancas != negras) imposible;   // O(casillas)
	//    Es una condicion necesaria, no suficiente: hay tableros balanceados sin cubrimiento.
	//    Ej. clasico: el 8x8 sin dos esquinas opuestas (del mismo color) tiene 32 y 30.
	//
	// 4. Monovariante (una cantidad que solo crece o solo decrece): si cada operacion
	//    reduce estrictamente un entero acotado inferiormente (inversiones, suma de
	//    cuadrados, potencial), el proceso termina en a lo sumo (valor inicial - cota)
	//    pasos, aunque no se sepa cual es el estado final.
	//
	// 5. Descubrir el invariante con fuerza bruta: para casos pequenos, hacer BFS desde
	//    cada estado, agrupar los estados alcanzables entre si y comparar clases con
	//    candidatos (suma, XOR, mcd, paridad de la permutacion, cantidad de impares).
	//    Si una propiedad separa exactamente las clases, es el invariante.
	//    set<Estado> visto; queue<Estado> q;   // por cada estado sin clase: BFS y asignar id
	//    // O(#estados * #operaciones); solo para explorar, no para la solucion final.
	//
	// 6. Cantidad minima de intercambios para ordenar una permutacion: n - ciclos (con
	//    intercambios cualesquiera) y #inversiones (si solo se intercambian vecinos).
	//    En ambos casos la paridad coincide con paridad_permutacion().
	//
	// 7. mcd como invariante: las operaciones a -> a - b (o a -> a + b) sobre elementos
	//    del arreglo no cambian mcd_total(). El estado final solo puede contener
	//    multiplos de ese mcd; si b no es multiplo, es inalcanzable.
	//
	// 8. Juegos: si cada jugada cambia la paridad de la cantidad total de piezas, la
	//    paridad inicial decide quien hace la ultima jugada. Se comprueba con una
	//    fuerza bruta pequena antes de confiar en la conjetura.
\end{lstlisting}

\textbf{Ejemplo de funcionamiento:}
\begin{lstlisting}[style=cpp]
	int main() {
		Invariantes inv({1, 2, 3, 4});
		cout << inv.paridad_suma() << " " << inv.xor_total() << " " << inv.cantidad_impares() << "\n";  // 0 4 2
		inv.cambiar(0, 5);   // {5, 2, 3, 4}
		cout << inv.paridad_suma() << " " << inv.xor_total() << " " << inv.cantidad_impares() << "\n";  // 0 0 2
		inv.cambiar(1, 7);   // {5, 7, 3, 4}
		cout << inv.paridad_suma() << " " << inv.xor_total() << " " << inv.cantidad_impares() << "\n";  // 1 5 3
		
		// Operacion: quitar dos numeros x >= y y poner x - y. La paridad de la suma es invariante.
		multiset<long long> ms = {1, 2, 3, 4, 5};
		cout << Invariantes({1, 2, 3, 4, 5}).paridad_suma() << "\n";   // 1  (suma 15, impar)
		while (ms.size() > 1) {
			auto it = prev(ms.end()); long long x = *it; ms.erase(it);
			it = prev(ms.end());      long long y = *it; ms.erase(it);
			ms.insert(x - y);
		}
		cout << *ms.begin() << "\n";                     // 1  (impar, como predijo el invariante)
		
		cout << Invariantes({12, 18, 30}).mcd_total() << "\n";   // 6
		
		cout << Invariantes::paridad_permutacion({1, 0, 2}) << "\n";      // 1  (un intercambio)
		cout << Invariantes::paridad_permutacion({1, 2, 0}) << "\n";      // 0  (ciclo de 3 = 2 intercambios)
		
		cout << Invariantes::puzzle_resoluble({1,2,3,4,5,6,7,0,8}, 3) << "\n";   // 1  (un movimiento desde la meta)
		cout << Invariantes::puzzle_resoluble({1,2,3,4,5,6,8,7,0}, 3) << "\n";   // 0  (dos fichas intercambiadas)
		cout << Invariantes::puzzle_resoluble({1,2,3,4,5,6,7,8,9,10,11,12,13,15,14,0}, 4) << "\n";  // 0  (el clasico 14-15)
		cout << Invariantes::puzzle_resoluble({1,2,3,4,5,6,7,8,9,10,11,0,13,14,15,12}, 4) << "\n";  // 1  (hueco un piso arriba)
		return 0;
	}
\end{lstlisting}

\textbf{Nota:} La regla del puzzle sale de estudiar cómo cambia la paridad de las inversiones. Un movimiento horizontal no altera el orden de lectura de las fichas, así que las inversiones no cambian. Un movimiento vertical hace pasar una ficha por encima de $N-1$ fichas, invirtiendo $N-1$ pares, así que la paridad de las inversiones cambia si y solo si $N$ es par. Con $N$ impar, la paridad de las inversiones es invariante y la meta tiene $0$ inversiones, de ahí \texttt{inv == 0}. Con $N$ par, cada movimiento vertical también cambia la fila del hueco en $1$, así que $inv + \text{fila\_hueco}$ mantiene su paridad, y la meta (con $0$ inversiones y el hueco en la fila $N-1$, impar) da paridad impar. Ese criterio supone la meta con el hueco abajo a la derecha. Para otra meta hay que compararla con la configuración dada usando la paridad de la permutación que lleva una a la otra. Además, \texttt{paridad\_suma} guarda solo el bit menos significativo a propósito, para no desbordar con sumas grandes; si necesitas la suma completa, guárdala aparte en \texttt{long long} o \texttt{\_\_int128}.

\textbf{Regla práctica:} Si el problema pregunta si algo es posible tras aplicar operaciones repetidas y el espacio de estados es enorme, busca primero un invariante (paridad, XOR, suma módulo $m$, \texttt{mcd}, coloración) que separe el estado inicial del objetivo. Usa BFS solo para descubrir el invariante en casos pequeños, no como solución final. Recuerda que un invariante prueba imposibilidad, no posibilidad: si coinciden inicio y objetivo, todavía necesitas una construcción o un conjunto de invariantes completo (como en el puzzle de 15). Si el enunciado pide además la secuencia de operaciones, el invariante te da la condición de existencia pero no la secuencia: construye una y usa el invariante para verificarla.
\subsection{Algoritmos Constructivos (con Cuadrado Mágico)}

\textbf{Idea:} Un algoritmo constructivo no \textit{busca} la respuesta: la
\textit{fabrica} con una regla explícita cuyo resultado cumple, por
construcción, las propiedades pedidas. El enunciado típico dice «construye una
permutación, matriz, grafo o cadena tal que\dots» (o «imprime $-1$ si no
existe»), y la solución es una receta que trabaja en tiempo proporcional al
tamaño de la salida, más una demostración de que la receta funciona. Esa
demostración suele ser una inducción (armar el caso $n$ a partir del $n-1$ o
del $n/2$), un invariante que la receta mantiene, o una fórmula cerrada que se
verifica con álgebra. Lo difícil es \textit{descubrir} la regla: lo habitual es
resolver a mano (o con fuerza bruta) los casos pequeños, mirar las soluciones y
buscar el patrón.

El ejemplo canónico es el \textit{cuadrado mágico} de orden $n$: una matriz
$n \times n$ que contiene cada entero de $1$ a $n^2$ exactamente una vez, con
todas las filas, todas las columnas y las dos diagonales sumando lo mismo. Como
los $n^2$ números suman $n^2(n^2+1)/2$ y se reparten en $n$ filas, la
constante debe ser $S = n(n^2+1)/2$ ($15$ para $n=3$, $34$ para $n=4$). Existe
para todo $n \neq 2$ (en $2 \times 2$, $a+b = a+c$ obliga a $b = c$), y la
construcción depende de $n \bmod 4$:
\begin{itemize}
	\item \textbf{$n$ impar (método siamés).} Se coloca el $1$ en el centro de
	la fila superior; cada número siguiente va una fila arriba y una columna a
	la derecha (con vuelta al otro lado), y si esa casilla está ocupada se baja
	una fila desde la casilla actual. Sea $m = n$ y $h = (m-1)/2$. Si se escribe
	$k-1 = q\,m + s$ con $0 \le s < m$, el número $k$ cae en la fila
	$2q - s$ y la columna $h - q + s$ (módulo $m$), o equivalentemente
	$q \equiv f + c - h$ y $s \equiv f + 2c + 1$ para la casilla $(f, c)$. En
	una fila (o columna) fija, al recorrer la otra coordenada tanto $q$ como $s$
	toman cada residuo exactamente una vez, porque $2$ es invertible módulo un
	$m$ impar. La suma es entonces
	$m \cdot \tfrac{m(m-1)}{2} + \tfrac{m(m-1)}{2} + m = \tfrac{m(m^2+1)}{2}$. Las
	diagonales se comprueban con aritmética modular análoga (algo más fina
	cuando $3 \mid m$).
	\item \textbf{$n = 4t$.} Se escribe $1, \dots, n^2$ por filas y se sustituye
	cada valor $v$ por $n^2+1-v$ en las casillas con $i \bmod 4 = j \bmod 4$ o
	$(i \bmod 4) + (j \bmod 4) = 3$ (las diagonales de cada bloque de
	$4 \times 4$). Como $v = n\,i + j + 1$, complementar equivale a reflejar
	$(i,j) \mapsto (n-1-i, n-1-j)$. Cada fila tiene $n/2$ casillas marcadas, y
	sus columnas son simétricas bajo $j \mapsto n-1-j$; por eso la suma de las
	partes «fila» y «columna» de todos los valores vuelve a ser
	$n \cdot \tfrac{n(n-1)}{2} + \tfrac{n(n-1)}{2} + n = S$. Con columnas y
	diagonales el argumento es análogo.
	\item \textbf{$n = 4t + 2$ (método de Strachey).} Sea $m = n/2$ (impar) y $A$
	el cuadrado siamés de orden $m$. Se ponen cuatro cuadrantes: $A$ arriba a la
	izquierda, $A + m^2$ abajo a la derecha, $A + 2m^2$ arriba a la derecha y
	$A + 3m^2$ abajo a la izquierda. Cada cuadrante es mágico, pero las filas de
	arriba suman $m^3$ menos que $S$ y las de abajo $m^3$ más. Se intercambian
	(entre arriba y abajo) $k = (m-1)/2$ casillas por fila en la izquierda, cada
	una suma $3m^2$ a la fila de arriba, y $k-1$ en la derecha, cada una resta
	$m^2$. El efecto neto es $3km^2 - (k-1)m^2 = (2k+1)m^2 = m^3$, justo lo que
	faltaba. En la fila central de la izquierda el intercambio se desplaza una
	columna a la derecha, para que columnas y diagonales también cuadren.
\end{itemize}

\textbf{¿Por qué usarlo?} Buscar un cuadrado mágico por retroceso es hopeless
incluso para $n$ modestos: hay $(n^2)!$ formas de ubicar los números y las
podas no bastan. La construcción da uno en $O(n^2)$. Frente a la programación
dinámica o al voraz, que \textit{optimizan} una cantidad, un algoritmo
constructivo \textit{satisface} restricciones estructurales: no hay función
objetivo. Es complementario a la subsección anterior: los invariantes y la
paridad demuestran que algo es \textit{imposible} (o eliminan casos), y la
construcción demuestra que es \textit{posible} al exhibirlo; juntos caracterizan
cuándo existe la respuesta. Su costo es que no hay un procedimiento general: cada
problema exige encontrar su patrón, y la solución no es única, así que en
un juez con verificador cualquier salida válida sirve, pero al depurar conviene
comparar contra un verificador propio y no contra una salida fija.

\textbf{Casos de uso:}
\begin{itemize}
	\item Problemas de «construye una permutación, matriz o cadena con tal propiedad» (con impresión de $-1$ si no existe).
	\item Cuadrados mágicos, cuadrados latinos y otras matrices con sumas o alfabetos controlados por fila y columna.
	\item Particionar $1, \dots, n$ en dos conjuntos de igual suma, o emparejar elementos con sumas objetivo.
	\item Recorridos y códigos con adyacencia garantizada: código de Gray, recorrido en serpiente de una cuadrícula.
	\item De nicho: calendarios de torneos (\textit{round-robin}), secuencias de De Bruijn, grafos con diámetro o grado prescritos.
\end{itemize}

\textbf{Complejidad:} La salida tiene $n^2$ casillas y el método las escribe un
número constante de veces cada una, así que construir cuesta $O(n^2)$, óptimo
porque solo imprimir ya cuesta eso. En el caso impar hay $n^2$ colocaciones y
cada una hace $O(1)$ trabajo (una suma, un módulo y una comprobación de
casilla ocupada). En $n = 4t$ se visita cada casilla una vez con $O(1)$ trabajo.
En $n = 4t+2$ se construye el cuadrado impar de orden $m = n/2$ ($O(m^2)$), se
copian los cuatro cuadrantes ($4m^2 = n^2$ escrituras) y se hacen a lo sumo
$m \cdot (2k - 1) = O(n^2)$ intercambios. Verificar con \texttt{es\_magico} es
$O(n^2)$ (cada casilla se lee un número constante de veces). Memoria: $O(n^2)$.

\begin{lstlisting}[style=cpp]
	struct CuadradoMagico {
		int n;                   // orden del cuadrado
		vector<vector<int>> c;   // c[i][j]: valor en la fila i, columna j
		
		// READ: indica si existe un cuadrado magico normal de orden n (falla n <= 0 y n = 2).
		static bool existe(int n) {
			return n == 1 || n >= 3;
		} // O(1)
		
		// CREATE: cuadrado de orden impar m por el metodo siames.
		static vector<vector<int>> impar(int m) {
			vector<vector<int>> a(m, vector<int>(m, 0));
			int f = 0, col = m / 2;                         // el 1 va al centro de la fila superior
			for (int k = 1; k <= m * m; k++) {
				a[f][col] = k;
				int nf = (f - 1 + m) % m, nc = (col + 1) % m;   // arriba y a la derecha
				if (a[nf][nc]) { nf = (f + 1) % m; nc = col; }  // ocupada: bajar desde la actual
				f = nf; col = nc;
			}
			return a;
		} // O(m^2)
		
		// CREATE: cuadrado de orden n multiplo de 4 (complementa las diagonales de cada bloque 4x4).
		static vector<vector<int>> doble_par(int n) {
			vector<vector<int>> a(n, vector<int>(n));
			for (int i = 0; i < n; i++)
			for (int j = 0; j < n; j++) {
				int v = i * n + j + 1;
				bool marcada = (i % 4 == j % 4) || (i % 4 + j % 4 == 3);
				a[i][j] = marcada ? n * n + 1 - v : v;
			}
			return a;
		} // O(n^2)
		
		// CREATE: cuadrado de orden n = 4t + 2 (n >= 6) por el metodo de Strachey.
		static vector<vector<int>> simple_par(int n) {
			int m = n / 2, k = (m - 1) / 2;
			vector<vector<int>> b = impar(m), a(n, vector<int>(n));
			for (int i = 0; i < m; i++)
			for (int j = 0; j < m; j++) {
				a[i][j]         = b[i][j];                  // arriba-izquierda
				a[i + m][j + m] = b[i][j] + m * m;          // abajo-derecha
				a[i][j + m]     = b[i][j] + 2 * m * m;      // arriba-derecha
				a[i + m][j]     = b[i][j] + 3 * m * m;      // abajo-izquierda
			}
			for (int i = 0; i < m; i++) {
				for (int j = 0; j < k; j++) {                   // k columnas a la izquierda
					int jj = (i == m / 2) ? j + 1 : j;          // en la fila central se corre una columna
					swap(a[i][jj], a[i + m][jj]);
				}
				for (int j = 0; j < k - 1; j++)                 // k-1 columnas a la derecha
				swap(a[i][n - 1 - j], a[i + m][n - 1 - j]);
			}
			return a;
		} // O(n^2)
		
		// CREATE: construye un cuadrado magico normal de orden n (elige el metodo segun n % 4).
		CuadradoMagico(int orden) : n(orden) {
			if (!existe(n)) throw invalid_argument("no existe cuadrado magico de ese orden");
			if (n % 2 == 1) c = impar(n);
			else if (n % 4 == 0) c = doble_par(n);
			else c = simple_par(n);
		} // O(n^2)
		
		// READ: constante magica S = n (n^2 + 1) / 2.
		long long constante() const {
			return (long long)n * ((long long)n * n + 1) / 2;
		} // O(1)
		
		// READ: verifica que sea magico: cada numero 1..n^2 una vez y todas las sumas iguales a S.
		bool es_magico() const {
			long long S = constante(), d1 = 0, d2 = 0;
			vector<char> visto(n * n + 1, 0);
			for (int i = 0; i < n; i++) {
				long long fila = 0, col = 0;
				for (int j = 0; j < n; j++) {
					int v = c[i][j];
					if (v < 1 || v > n * n || visto[v]) return false;
					visto[v] = 1;
					fila += v;
					col += c[j][i];
				}
				if (fila != S || col != S) return false;
				d1 += c[i][i];
				d2 += c[i][n - 1 - i];
			}
			return d1 == S && d2 == S;
		} // O(n^2)
		
		// WRITE: gira el cuadrado 90 grados en sentido horario (sigue siendo magico).
		void rotar90() {
			vector<vector<int>> r(n, vector<int>(n));
			for (int i = 0; i < n; i++)
			for (int j = 0; j < n; j++) r[j][n - 1 - i] = c[i][j];
			c = move(r);
		} // O(n^2)
	};
	
	// ------------------------------------------------------------
	// MODIFICACIONES Y COMENTARIOS CON LOS USOS
	// ------------------------------------------------------------
	// 1. Simetrias y desplazamientos: rotar, reflejar o transponer un cuadrado magico
	//    da otro cuadrado magico (ocho variantes en total). Complementar (v -> n*n + 1 - v)
	//    tambien. Sumar una constante t a todas las casillas da numeros t+1..t+n*n con
	//    suma magica S + n*t. Sirve cuando el enunciado pide otro rango o cualquiera de
	//    las variantes.
	//
	// 2. Formula cerrada para n impar (sin simular): el valor de la casilla (i, j) es
	//    int valor_impar(int i, int j, int m) {
		//        int h = (m - 1) / 2;
		//        return m * (((i + j - h) % m + m) % m) + (i + 2 * j + 1) % m + 1;
		//    } // O(1) por casilla
	//    Util para consultar una casilla sin construir toda la matriz.
	//
	// 3. Cuadrado latino ciclico: a[i][j] = (i + j) % n. Cada simbolo aparece una vez
	//    por fila y por columna; es la base de calendarios round-robin y de
	//    construcciones de matrices con simbolos no repetidos.   // O(n^2)
	//
	// 4. Partir 1..n en dos conjuntos con la misma suma (posible solo si n % 4 es 0 o 3,
	//    porque la suma total n(n+1)/2 debe ser par):
	//    vector<int> A, B; int t = 1;
	//    if (n % 4 == 3) { A = {1, 2}; B = {3}; t = 4; }
	//    for (; t + 3 <= n; t += 4) {            // t + (t+3) == (t+1) + (t+2)
		//        A.push_back(t);     A.push_back(t + 3);
		//        B.push_back(t + 1); B.push_back(t + 2);
		//    } // O(n)
	//
	// 5. Codigo de Gray: los n-bits 0..2^n-1 en un orden donde dos consecutivos (y el
	//    ultimo con el primero) difieren en un solo bit.
	//    for (int i = 0; i < (1 << n); i++) g[i] = i ^ (i >> 1);   // O(2^n)
	//
	// 6. Recorrido en serpiente de una cuadricula f x c: visita todas las casillas y
	//    cada paso (incluido el salto entre filas) va a una casilla adyacente.
	//    for (int i = 0; i < f; i++)
	//        for (int j = 0; j < c; j++) visitar(i, (i % 2 == 0) ? j : c - 1 - j);   // O(f*c)
	//
	// 7. Construccion por induccion (de n-1 a n, o de n/2 a n): resolver el caso base y
	//    extender. Plantilla:
	//    vector<int> construir(int n) {
		//        if (n == 1) return {1};
		//        vector<int> s = construir(n - 1);
		//        // ... insertar el elemento n en la posicion que conserva la propiedad ...
		//        return s;
		//    }
	//    Ejemplo simple: desorden (permutacion sin puntos fijos) para n >= 2: p[i] = (i + 1) % n.
	//
	// 8. Descubrir el patron: escribir una fuerza bruta (retroceso) para n pequenos,
	//    imprimir todas las soluciones, buscar la regla y comprobarla con un
	//    verificador (como es_magico) sobre muchos n antes de enviar. Empezar siempre
	//    por decidir la existencia (invariantes, paridad): si no existe, no hay que construir.
	//
	// 9. Otros metodos para n = 4t + 2: LUX de Conway (usa tres tipos de losetas L, U, X
	//    sobre un cuadrado impar de tamano (n+2)/4 ... o similar); Strachey (el de arriba) es
	//    el mas corto de programar.
\end{lstlisting}

\textbf{Ejemplo de funcionamiento:}
\begin{lstlisting}[style=cpp]
	void imprimir(const CuadradoMagico& m) {
		for (const auto& fila : m.c) {
			for (int x : fila) cout << x << " ";
			cout << "\n";
		}
	}
	
	int main() {
		CuadradoMagico a(3);
		imprimir(a);
		// 8 1 6
		// 3 5 7
		// 4 9 2
		cout << a.constante() << " " << a.es_magico() << "\n";   // 15 1
		a.rotar90();
		imprimir(a);
		// 4 3 8
		// 9 5 1
		// 2 7 6
		cout << a.es_magico() << "\n";                           // 1  (rotar conserva la propiedad)
		
		CuadradoMagico b(4);
		imprimir(b);
		// 16 2 3 13
		// 5 11 10 8
		// 9 7 6 12
		// 4 14 15 1
		cout << b.constante() << " " << b.es_magico() << "\n";   // 34 1
		
		CuadradoMagico e(5);
		for (int x : e.c[0]) cout << x << " ";
		cout << "\n";                                            // 17 24 1 8 15
		cout << e.constante() << " " << e.es_magico() << "\n";   // 65 1
		
		CuadradoMagico d(6);
		imprimir(d);
		// 35 1 6 26 19 24
		// 3 32 7 21 23 25
		// 31 9 2 22 27 20
		// 8 28 33 17 10 15
		// 30 5 34 12 14 16
		// 4 36 29 13 18 11
		cout << d.constante() << " " << d.es_magico() << "\n";   // 111 1
		
		CuadradoMagico g(8);
		for (int x : g.c[0]) cout << x << " ";
		cout << "\n";                                            // 64 2 3 61 60 6 7 57
		cout << g.constante() << " " << g.es_magico() << "\n";   // 260 1
		
		CuadradoMagico h(10);
		cout << h.constante() << " " << h.es_magico() << "\n";   // 505 1
		cout << CuadradoMagico::existe(2) << "\n";               // 0  (no existe el 2x2)
		return 0;
	}
\end{lstlisting}

\textbf{Nota:} El cuadrado que se obtiene no es el único de su orden: el método
siamés, por ejemplo, tiene variantes según la casilla de partida y la
dirección, y las rotaciones y reflexiones dan otros. Si el problema exige una
salida concreta (y no un verificador), hay que ajustar esas elecciones al
enunciado. El constructor lanza una excepción para $n = 2$ y $n \le 0$, y el
método de Strachey solo se llama con $n \ge 6$ (para $n = 2$ daría $k = 0$ y
un índice inválido). Además, \texttt{impar} usa $0$ como marca de casilla vacía,
lo cual es válido porque los valores empiezan en $1$. Para $n$ grande, el
límite práctico es la memoria de $n^2$ enteros, no el tiempo.

\textbf{Regla práctica:} Si el problema pide construir un objeto con propiedades
de suma, orden o adyacencia, primero decide con invariantes si existe; luego
resuelve los casos pequeños con fuerza bruta, busca el patrón, y demuestra por
inducción o con una fórmula. Para cuadrados mágicos usa tres casos según
$n \bmod 4$: siamés para impar, complemento de diagonales de bloques para
múltiplos de $4$ y Strachey para $4t+2$. Comprueba siempre la construcción con un
verificador (\texttt{es\_magico}) contra muchos $n$ antes de confiar en ella. Si
lo que se busca es la mejor solución entre varias posibles, un algoritmo
constructivo no basta: usa voraz o programación dinámica.

\subsection{Trabajar hacia Atrás (Reverse Thinking)}

\textbf{Idea:} Trabajar hacia atrás significa empezar en el estado final y
\textit{deshacer} las operaciones, en lugar de simular desde el estado inicial.
Ayuda por tres razones estructurales. La primera es la \textit{asimetría de
	ramificación}: si un estado tiene muchos sucesores pero uno solo (o muy pocos)
predecesores, la búsqueda hacia adelante explora un árbol exponencial y hacia
atrás recorre una línea. La segunda es el \textit{orden temporal}: si la última
operación sobre un dato anula a las anteriores («la última escritura gana»),
recorrer las operaciones desde la última permite decidir cada dato la primera vez
que se toca y descartar el resto. La tercera es que el \textit{objetivo suele ser
	más restrictivo que el inicio}: en un juego, las posiciones terminales son
conocidas y de ellas se propaga hacia atrás quién gana.

El ejemplo más limpio: se parte de un par $(x, y)$ de enteros positivos y en cada
paso se puede pasar a $(x+y, y)$ o a $(x, x+y)$; ¿se puede llegar de $(a, b)$ a
$(c, d)$? Hacia adelante hay dos opciones por paso, es decir, $2^k$ estados tras
$k$ pasos. Hacia atrás no hay elección: si $c > d$, la última operación tuvo que
ser $(x, y) \to (x+y, y)$, porque la otra operación dejaría la primera
coordenada igual a la anterior y $c > d$ obliga a que la primera sea la grande;
así que el predecesor es único, $(c-d, d)$. Si $c < d$ es simétrico, y si $c = d$
no tiene predecesor válido (habría que restar y llegar a $0$). El camino de
regreso es único: es el algoritmo de Euclides por restas, y se acelera dividiendo
en lugar de restar de a uno.

El segundo ejemplo es «la última escritura gana»: se aplican $m$ operaciones
\texttt{pintar}$(l, r, \text{color})$ sobre $n$ celdas y se pregunta el color
final. Recorridas al revés, cada celda se pinta una sola vez (la primera que se
alcanza al ir hacia atrás es la última en el tiempo real), y para saltar las
celdas ya pintadas basta una estructura de «siguiente celda libre» con
compresión de caminos.

\textbf{¿Por qué usarlo?} Frente a BFS o programación dinámica hacia adelante,
el enfoque hacia atrás vale cuando el factor de ramificación hacia atrás es
menor que hacia adelante (aquí $1$ contra $2$). Si es igual o mayor, no ayuda,
y si ambos son grandes conviene la búsqueda bidireccional (modificación 4).
Frente a los invariantes de la subsección anterior, que dicen solo «posible» o
«imposible» (aquí, el \texttt{mcd}), ir hacia atrás además entrega la cantidad
exacta de pasos y el camino. Frente a un árbol de segmentos con propagación
perezosa para asignaciones de rango, pintar en reversa es más simple y casi
lineal, pero solo funciona \textit{fuera de línea}: hay que conocer todas las
operaciones y no admite consultas intermedias. Y frente a los algoritmos
constructivos, es a menudo su herramienta: para construir un camino se parte del
final y se retrocede con la regla determinista.

\textbf{Casos de uso:}
\begin{itemize}
	\item Alcanzabilidad y número de pasos entre estados cuando cada operación tiene un inverso determinista (pares con sumas, algoritmos tipo Euclides).
	\item Estado final tras muchas asignaciones de rango (pintar, colorear tramos, «la última escritura gana»), con todas las operaciones conocidas.
	\item Conectividad fuera de línea con eliminaciones de aristas: se procesa el tiempo al revés y las eliminaciones se vuelven inserciones (unión-búsqueda).
	\item Juegos: análisis retrógrado desde las posiciones terminales para saber quién gana.
	\item De nicho: deshacer una secuencia conocida de intercambios, DP con estados finales conocidos, o BFS bidireccional.
\end{itemize}

\textbf{Complejidad:} Para \texttt{pasos}, cada iteración del ciclo es un paso de
división de Euclides: reduce una coordenada a un valor en $[1, \text{otra}]$, y
las dos coordenadas se alternan. Como en Euclides (peor caso en números de
Fibonacci consecutivos, Lamé), cada dos iteraciones el mayor se reduce al menos
a la mitad, así que hay $O(\log M)$ iteraciones con $M = \max(c, d)$, y cada una
hace un número fijo de divisiones y restas: $O(\log M)$, con memoria $O(1)$
(contra $O(2^k)$ o $O(M)$ de la búsqueda hacia adelante). Para
\texttt{PinturaInversa}, el constructor cuesta $O(n)$. Cada celda se pinta a lo
sumo una vez en toda la ejecución, y cada llamada a \texttt{pintar} hace una
búsqueda inicial más una por celda pintada, es decir $O(n + m)$ búsquedas en
total. Con solo compresión de caminos cada búsqueda cuesta $O(\log n)$
amortizado en el peor caso teórico, y en la práctica casi constante. El total
es $O((n + m)\log n)$ (casi lineal en la práctica), contra $O(n \cdot m)$ de
pintar celda por celda hacia adelante. Memoria: $O(n)$.

\begin{lstlisting}[style=cpp]
	// Pares (x, y) con operaciones (x, y) -> (x + y, y) o (x, x + y), todos positivos.
	struct ParAtras {
		// READ: cantidad de pasos para ir de (a, b) a (c, d), o -1 si no es posible.
		// Hacia atras el predecesor es unico, asi que esa cantidad es exacta (y minima).
		// Requiere a, b, c, d >= 1.
		static long long pasos(long long a, long long b, long long c, long long d) {
			long long total = 0;
			while (true) {
				if (c < a || d < b) return -1;              // las coordenadas solo crecen
				if (c == a && d == b) return total;
				if (c == d) return -1;                      // (x, x) no tiene predecesor positivo
				if (c > d) {
					if (d == b && (c - a) % d == 0)         // llegamos a (a, b) restando d varias veces
					return total + (c - a) / d;
					long long k = (c - 1) / d;              // restas posibles manteniendo c >= 1
					c -= k * d; total += k;
				} else {
					if (c == a && (d - b) % c == 0)
					return total + (d - b) / c;
					long long k = (d - 1) / c;
					d -= k * c; total += k;
				}
			}
		} // O(log M)
		
		// READ: indica si (c, d) es alcanzable desde (a, b).
		static bool alcanzable(long long a, long long b, long long c, long long d) {
			return pasos(a, b, c, d) >= 0;
		} // O(log M)
	};
	
	// Estado final tras m operaciones "pintar [l, r] con color" aplicadas en orden.
	struct PinturaInversa {
		int n;
		vector<int> sig;     // sig[i]: primera celda >= i sin pintar (n es el centinela)
		vector<int> color;   // color final de cada celda (0 = nunca pintada)
		
		// CREATE: n celdas sin pintar.
		PinturaInversa(int tam) : n(tam), sig(tam + 1), color(tam, 0) {
			iota(sig.begin(), sig.end(), 0);
		} // O(n)
		
		// READ: primera celda sin pintar en [x, n] (n si no queda ninguna).
		// Comprime caminos (cambia sig, pero no el resultado).
		int buscar(int x) {
			int r = x;
			while (sig[r] != r) r = sig[r];
			while (sig[x] != r) { int nx = sig[x]; sig[x] = r; x = nx; }
			return r;
		} // O(log n) amortizado
		
		// WRITE: pinta con c las celdas todavia libres de [l, r] (inclusive).
		void pintar(int l, int r, int c) {
			for (int i = buscar(l); i <= r; i = buscar(i)) {
				color[i] = c;
				sig[i] = i + 1;                             // la celda i ya no esta libre
			}
		} // O((1 + k) log n) amortizado, k = celdas pintadas por esta llamada
		
		// WRITE: aplica las operaciones (l, r, c) recorriendolas de la ultima a la primera.
		void aplicar_en_reversa(const vector<array<int, 3>>& ops) {
			for (int i = (int)ops.size() - 1; i >= 0; i--)
			pintar(ops[i][0], ops[i][1], ops[i][2]);
		} // O((n + m) log n)
	};
	
	// ------------------------------------------------------------
	// MODIFICACIONES Y COMENTARIOS CON LOS USOS
	// ------------------------------------------------------------
	// 1. Eliminaciones -> inserciones (conectividad fuera de linea). Consultas: "cuantas
	//    componentes hay tras eliminar las primeras i aristas e[1..i]?", i = 0..k.
	//    Se construye el estado FINAL y se retrocede en el tiempo agregando aristas:
	//    // DSU d(n); unir todas las aristas que nunca se eliminan;
	//    // resp[k] = d.componentes();
	//    // for (int i = k; i >= 1; i--) { d.unir(e[i]); resp[i-1] = d.componentes(); }
	//    // O((n + m + k) alfa(n)). Hacia adelante haria falta conectividad dinamica,
	//    // mucho mas dificil: la union-busqueda no sabe deshacer una arista.
	//
	// 2. Analisis retrogrado en juegos (posiciones y jugadas forman un grafo dirigido).
	//    estado[v]: 0 = sin resolver, 1 = gana quien mueve, 2 = pierde quien mueve.
	//    // grado[v] = cantidad de jugadas desde v. Cola con todas las terminales
	//    // (grado 0: pierde quien mueve, estado = 2).
	//    // while (!q.empty()) {
		//    //     int v = q.front(); q.pop();
		//    //     for (int u : predecesores[v]) {             // u puede saltar a v
			//    //         if (estado[u]) continue;
			//    //         if (estado[v] == 2) { estado[u] = 1; q.push(u); }        // salta a una perdedora
			//    //         else if (--grado[u] == 0) { estado[u] = 2; q.push(u); }  // todas sus jugadas dan a ganadoras
			//    //     }
		//    // }   // O(V + E). Lo que queda sin resolver es empate (ciclos).
	//
	// 3. Deshacer transformaciones conocidas: si a partir de un arreglo inicial se
	//    aplicaron k intercambios (x[i], y[i]) en orden y se conoce el resultado a,
	//    recuperar el inicial es aplicar los intercambios en orden inverso:
	//    for (int i = k - 1; i >= 0; i--) swap(a[x[i]], a[y[i]]);   // O(k)
	//    En general, el inverso de una composicion es la composicion de los inversos al reves.
	//
	// 4. BFS bidireccional: cuando ambos extremos (inicio y objetivo) son conocidos y hay
	//    predecesores computables, buscar desde los dos lados a la vez y detenerse al
	//    encontrarse. Con ramificacion b y distancia d, baja el costo de O(b^d) a
	//    O(2 * b^(d/2)).
	//
	// 5. DP hacia atras: definir dp[estado] = mejor resultado desde ese estado hasta el
	//    final y llenarla desde los estados finales. Conviene cuando hay muchos estados
	//    iniciales o cuando el inicial no se conoce, porque una sola pasada responde
	//    para todos.
	//
	// 6. Reconstruir el camino: hacia atras se obtienen los estados de fin a inicio;
	//    guardarlos en un vector y darlo vuelta (reverse) antes de imprimir.
	//    OJO: si el camino es enorme (como en pasos(1, 1, 1e9, 1)), no se puede guardar
	//    completo; en ese caso se imprime comprimido (repeticiones de la misma operacion).
	//
	// 7. Cuando el predecesor NO es unico (por ejemplo x -> x mod k): hacia atras tambien
	//    hay ramificacion. Comparar cuantos sucesores y cuantos predecesores tiene un
	//    estado tipico y elegir el sentido con menos; si son parecidos, bidireccional.
	//
	// 8. Pregunta guia: "cual fue el ultimo paso?" o "que tiene que ser cierto justo antes
	//    del final?". Si la respuesta es unica o casi unica, hay un algoritmo hacia atras.
\end{lstlisting}

\textbf{Ejemplo de funcionamiento:}
\begin{lstlisting}[style=cpp]
	int main() {
		cout << ParAtras::pasos(1, 1, 3, 5) << "\n";            // 3   (1,1) -> (1,2) -> (3,2) -> (3,5)
		cout << ParAtras::pasos(1, 1, 8, 5) << "\n";            // 4   (1,1) -> (1,2) -> (3,2) -> (3,5) -> (8,5)
		cout << ParAtras::pasos(2, 3, 12, 5) << "\n";           // 3   (2,3) -> (2,5) -> (7,5) -> (12,5)
		cout << ParAtras::pasos(1, 1, 1000000000, 1) << "\n";   // 999999999  (una sola vuelta del ciclo)
		cout << ParAtras::pasos(1, 1, 2, 4) << "\n";            // -1  (llega a (2,2), sin predecesor; mcd(2,4) = 2)
		cout << ParAtras::pasos(5, 1, 3, 4) << "\n";            // -1  (3 < 5: las coordenadas solo crecen)
		cout << ParAtras::alcanzable(1, 1, 8, 5) << "\n";       // 1
		
		PinturaInversa p(8);                                    // celdas 0..7
		p.aplicar_en_reversa({{0, 4, 1}, {2, 6, 2}, {3, 3, 3}});
		for (int x : p.color) cout << x << " ";
		cout << "\n";                                           // 1 1 2 3 2 2 2 0
		// Hacia adelante: 1 1 1 1 1 0 0 0 -> 1 1 2 2 2 2 2 0 -> 1 1 2 3 2 2 2 0. Hacia atras, la
		// ultima operacion pinta la celda 3; la segunda pinta 2, 4, 5, 6 (la 3 ya esta pintada);
		// la primera pinta 0 y 1 (2, 3, 4 ya estan pintadas). La celda 7 nunca se pinta.
		return 0;
	}
\end{lstlisting}

\textbf{Nota:} \texttt{pasos} exige $a, b, c, d \ge 1$: con un $0$ las operaciones
\texttt{\%} y \texttt{/} se dividen por cero o el ciclo no termina. También
supone que el predecesor es único; si las operaciones fueran otras (por ejemplo
$(x,y) \to (x+y, x-y)$), habría que volver a razonar cuántos predecesores tiene un estado.
\texttt{PinturaInversa} requiere que las operaciones sean conocidas de antemano y que
$0 \le l \le r < n$; con consultas intermedias (¿qué color tiene la celda $i$ ahora?)
hay que volver a un árbol de segmentos con propagación perezosa. Por último,
\texttt{buscar} modifica \texttt{sig} para comprimir caminos, así que la estructura no es
segura para uso concurrente aunque conceptualmente sea una consulta.

\textbf{Regla práctica:} Si al mirar un estado hay una sola manera (o muy pocas)
de haber llegado a él, o si la última operación sobre cada dato decide su valor,
trabaja hacia atrás: parte del final, deshaz operaciones con la regla
determinista y acelera con divisiones (estilo Euclides) cuando se repita el mismo
paso. Para «estado final tras asignaciones de rango» fuera de línea, pinta en
reversa con una estructura de «siguiente libre»; para eliminaciones fuera de
línea, invierte el tiempo y usa unión-búsqueda; para juegos, propaga desde las
posiciones terminales. Si hay ramificación en ambos sentidos, usa BFS
bidireccional. Y si el problema tiene consultas en línea, no puedes ir hacia
atrás: usa un árbol de segmentos con propagación perezosa u otra estructura
dinámica.
\subsection{Contribución de cada Elemento}

\textbf{Idea:} muchos problemas piden una suma sobre \textbf{todos los subarreglos, subconjuntos o pares} (por ejemplo, la suma de los mínimos de todos los subarreglos). En vez de recorrer cada subarreglo y calcular su mínimo, aporte o costo por separado —lo cual es $O(n^2)$ o peor—, la técnica de \textbf{contribución} invierte el orden de la suma: para cada elemento $a[i]$, se calcula \textbf{en cuántos} de esos subarreglos/subconjuntos $a[i]$ es el que determina el resultado (por ejemplo, en cuántos subarreglos $a[i]$ es el mínimo), y se multiplica $a[i]$ por esa cantidad. Sumando la contribución de cada elemento por separado se obtiene la respuesta total, sin enumerar los subarreglos uno por uno.

\textbf{¿Por qué usarlo?} Sumar directamente sobre los $O(n^2)$ subarreglos (o $O(2^n)$ subconjuntos) es demasiado lento cuando $n$ es grande. Calcular cuántas veces contribuye cada elemento suele reducirse a encontrar, para cada $a[i]$, el rango de posiciones donde $a[i]$ sigue siendo el mínimo (o máximo, o el elemento relevante) — típicamente con una \textbf{pila monótona} que encuentra el vecino anterior y siguiente estrictamente menor (o mayor) en $O(n)$ total. Es el mismo cambio de perspectiva de "sumar por posición en vez de por subconjunto" que aparece en problemas de conteo de la parte de Matemáticas, aplicado aquí a estructuras de subarreglo.

\textbf{Casos de uso:}
\begin{itemize}
	\item Suma de los mínimos (o máximos) de todos los subarreglos de un arreglo.
	\item Suma de todos los subarreglos posibles: cada elemento $a[i]$ aparece en $(i+1)\cdot(n-i)$ subarreglos, así que la respuesta es $\sum_i a[i]\cdot(i+1)\cdot(n-i)$.
	\item Contar cuántos pares $(i,j)$ cumplen una condición, sumando la contribución de cada elemento como extremo del par.
	\item Suma de las diferencias absolutas entre todos los pares de un arreglo ordenado (cada elemento contribuye según cuántos quedan a su izquierda y derecha), como en \textit{Suma de Distancias Manhattan}.
	\item Problemas de "suma de X sobre todos los subconjuntos", donde X depende de una propiedad local de cada elemento.
\end{itemize}

\textbf{Complejidad:} $O(n)$ o $O(n\log n)$ según el problema: normalmente $O(n)$ con una pila monótona para encontrar los rangos donde cada elemento es el mínimo/máximo, seguido de una suma lineal de contribuciones.

\begin{lstlisting}[style=cpp]
	// CREATE/READ: Suma de los minimos de TODOS los subarreglos de un
	// arreglo, mediante la tecnica de contribucion con pila monotona
	struct SumaMinimosSubarreglos {
		// READ: para cada indice i, encuentra el limite izquierdo (indice del
		// primer elemento estrictamente MENOR a la izquierda, o -1) y el
		// limite derecho (primer elemento MENOR O IGUAL a la derecha, o n);
		// el desempate <= evita contar dos veces subarreglos con minimos
		// repetidos
		static long long sumaDeMinimos(const vector<long long>& a) {
			int n = (int)a.size();
			vector<int> izq(n), der(n);
			vector<int> pila;
			
			// limite izquierdo: primer menor estricto a la izquierda
			for (int i = 0; i < n; i++) {
				while (!pila.empty() && a[pila.back()] >= a[i]) pila.pop_back();
				izq[i] = pila.empty() ? -1 : pila.back();
				pila.push_back(i);
			}
			pila.clear();
			
			// limite derecho: primer menor o igual a la derecha
			for (int i = n - 1; i >= 0; i--) {
				while (!pila.empty() && a[pila.back()] > a[i]) pila.pop_back();
				der[i] = pila.empty() ? n : pila.back();
				pila.push_back(i);
			}
			
			long long total = 0;
			const long long MOD = 1000000007LL;
			for (int i = 0; i < n; i++) {
				// cantidad de subarreglos donde a[i] es el minimo:
				// (i - izq[i]) elecciones de inicio * (der[i] - i) elecciones de fin
				long long cantidad = (long long)(i - izq[i]) * (long long)(der[i] - i) % MOD;
				total = (total + (a[i] % MOD) * cantidad) % MOD;
			}
			return total;
		} // O(n)
		
		// ------------------------------------------------------------
		// MODIFICACIONES Y COMENTARIOS CON LOS USOS
		// ------------------------------------------------------------
		// 1. Suma de MAXIMOS en vez de minimos: invertir las comparaciones de
		//    la pila (>= por <=, y > por <), manteniendo el mismo desempate
		//    estricto de un solo lado para no duplicar subarreglos.
		// 2. Suma de TODOS los subarreglos (no el minimo, la suma completa de
		//    cada uno): cada a[i] aparece en (i+1)*(n-i) subarreglos distintos
		//    (i+1 elecciones de inicio, n-i de fin), sin necesidad de pila
		//    monotona: total = sum(a[i] * (i+1) * (n-i)).
		// 3. Desempate con elementos iguales: usar >= de un lado y > del otro
		//    (como en el codigo) evita que un mismo subarreglo con minimo
		//    repetido se cuente dos veces, una por cada copia del minimo.
		// 4. Contribucion de PARES (i,j) segun una condicion sobre a[i],a[j]:
		//    en vez de recorrer O(n^2) pares, agrupar por valor o por
		//    propiedad y contar cuantos elementos a cada lado cumplen la
		//    condicion, sumando la contribucion de cada elemento como extremo.
		// 5. Overflow: la cantidad de subarreglos puede ser del orden de n^2;
		//    aplicar modulo en cada multiplicacion si el problema lo pide, o
		//    usar __int128/long long con cuidado si se necesita el valor exacto.
	};
\end{lstlisting}

\textbf{Ejemplo de funcionamiento:}
\begin{lstlisting}[style=cpp]
	long long resultado = SumaMinimosSubarreglos::sumaDeMinimos({3, 1, 2, 4});
	// resultado = 17
	// (subarreglos y sus minimos: [3]=3, [1]=1, [2]=2, [4]=4, [3,1]=1,
	//  [1,2]=1, [2,4]=2, [3,1,2]=1, [1,2,4]=1, [3,1,2,4]=1;
	//  suma total = 3+1+2+4+1+1+2+1+1+1 = 17)
\end{lstlisting}

\textbf{Nota:} el desempate entre \texttt{izq} (con \texttt{>=}, borde estrictamente menor) y \texttt{der} (con \texttt{>}, borde menor o igual) es lo que evita contar dos veces un subarreglo cuando el mínimo aparece repetido dentro de él. Si ambos lados usan el mismo criterio de desempate, subarreglos con mínimos duplicados se cuentan más de una vez, dando un resultado mayor al correcto.

\textbf{Regla práctica:} cuando el problema pida una suma sobre \textbf{todos} los subarreglos o pares y la fuerza bruta sea $O(n^2)$ o peor, pregúntate "¿en cuántas de esas estructuras participa cada elemento, y con qué papel?" — casi siempre esa pregunta se resuelve con una pila monótona en $O(n)$, convirtiendo una suma cuadrática en una suma lineal.

\subsection{Fuerza Bruta para Descubrir Patrones}

\textbf{Idea:} cuando la fórmula o el algoritmo de un problema no es evidente, una estrategia práctica es \textbf{programar la fuerza bruta} (aunque sea exponencial o $O(n^2)$, sin importar que no pase el límite de tiempo real) y \textbf{imprimir la respuesta para valores pequeños} de $n$ (por ejemplo $n=1,2,3,\dots,15$). Con esa secuencia de salida, se busca un patrón: ¿es lineal? ¿cuadrática? ¿coincide con una secuencia conocida (Fibonacci, factorial, números de Catalan)? ¿la diferencia entre términos consecutivos sigue un patrón más simple? Una vez identificada la fórmula o recurrencia, se implementa la versión eficiente y se \textbf{verifica} contra la fuerza bruta en los mismos valores pequeños antes de confiar en ella.

\textbf{¿Por qué usarlo?} No es un algoritmo en sí, sino una \textbf{metodología de resolución} bajo presión de tiempo: muchos problemas de competencia (especialmente en la parte Ad-hoc) esconden una fórmula cerrada o una recurrencia simple detrás de un enunciado que no la sugiere directamente. Derivar la fórmula analíticamente puede ser lento o propenso a errores; observarla empíricamente a partir de una tabla de valores pequeños suele ser más rápido y, combinado con la verificación cruzada contra la fuerza bruta, reduce el riesgo de enviar una fórmula incorrecta.

\textbf{Casos de uso:}
\begin{itemize}
	\item Problemas de conteo donde la fórmula cerrada no es obvia a partir del enunciado (por ejemplo, "de cuántas formas se puede...").
	\item Verificar rápidamente una fórmula que se sospecha, comparándola contra la fuerza bruta en casos pequeños antes de implementarla para el caso general.
	\item Identificar si la respuesta sigue una secuencia conocida (buscar los primeros términos en OEIS mentalmente o por reconocimiento de patrones: triangular, factorial, potencias de 2, Catalan, Fibonacci).
	\item Depurar una solución eficiente que da resultados distintos a la fuerza bruta en casos pequeños, para localizar en qué valor de $n$ empieza a diferir.
	\item Problemas de geometría o combinatoria ad-hoc donde la relación entre las variables de entrada y la salida no es evidente algebraicamente.
\end{itemize}

\textbf{Complejidad:} no aplica una complejidad única — depende de la fuerza bruta usada para generar los datos (normalmente $O(2^n)$, $O(n!)$ o $O(n^3)$, ejecutada solo para $n$ pequeño) y de la fórmula final encontrada, que en el mejor de los casos es $O(1)$ o $O(\log n)$.

\begin{lstlisting}[style=cpp]
	// CREATE/READ: Plantilla para generar una tabla de valores pequeños con
	// fuerza bruta, imprimir diferencias sucesivas, y verificar una formula
	// candidata contra esa misma fuerza bruta
	struct FuerzaBrutaPatrones {
		// READ: fuerza bruta de EJEMPLO: cuenta el numero de formas de subir
		// una escalera de n escalones dando pasos de 1 o 2 (Fibonacci
		// desplazado); se reemplaza por la fuerza bruta REAL del problema
		static long long fuerzaBruta(int n) {
			if (n == 0) return 1;
			if (n < 0) return 0;
			return fuerzaBruta(n - 1) + fuerzaBruta(n - 2);
		} // O(2^n), solo usable para n pequeño (hasta ~30)
		
		// READ: genera la tabla de valores f(1), f(2), ..., f(limite) para
		// inspeccionar a simple vista o comparar contra una formula candidata
		static vector<long long> generarTabla(int limite) {
			vector<long long> tabla;
			for (int n = 1; n <= limite; n++) tabla.push_back(fuerzaBruta(n));
			return tabla;
		} // O(2^limite) con la fuerza bruta de ejemplo
		
		// READ: imprime las diferencias sucesivas de una tabla; una diferencia
		// constante sugiere progresion aritmetica, una diferencia lineal
		// sugiere cuadratica, y una razon constante entre terminos sugiere
		// progresion geometrica o recurrencia lineal simple
		static vector<long long> diferencias(const vector<long long>& tabla) {
			vector<long long> dif;
			for (size_t i = 1; i < tabla.size(); i++) dif.push_back(tabla[i] - tabla[i - 1]);
			return dif;
		} // O(n)
		
		// READ: verifica una formula candidata contra la fuerza bruta en un
		// rango de valores pequenos; retorna el primer n donde difieren, o -1
		// si coinciden en todo el rango probado
		static int verificarFormula(int limite, function<long long(int)> formula) {
			for (int n = 1; n <= limite; n++) {
				if (fuerzaBruta(n) != formula(n)) return n;
			}
			return -1; // coinciden en todo el rango: buena señal, no prueba definitiva
		} // O(limite) llamadas a formula() + costo de fuerzaBruta()
		
		// ------------------------------------------------------------
		// MODIFICACIONES Y COMENTARIOS CON LOS USOS
		// ------------------------------------------------------------
		// 1. Reemplazar fuerzaBruta(): esta es solo un EJEMPLO (escalera de
		//    Fibonacci); en un problema real se sustituye por backtracking,
		//    busqueda exhaustiva de permutaciones, o simulacion directa del
		//    enunciado, sin importar que sea lenta.
		// 2. Buscar la secuencia en OEIS (fuera de competencia, o si hay
		//    acceso a internet permitido): los primeros 6-8 terminos de una
		//    tabla suelen ser suficientes para identificar una secuencia
		//    conocida y su formula cerrada o funcion generadora.
		// 3. Patrones comunes a reconocer a simple vista: diferencias
		//    constantes = aritmetica (f(n) = a*n + b); razon constante entre
		//    terminos = geometrica (f(n) = a * r^n); cociente f(n)/f(n-1)
		//    tendiendo a un valor = crecimiento exponencial con esa base;
		//    coincide con 1,1,2,3,5,8,... = Fibonacci; coincide con
		//    1,1,2,5,14,42,... = numeros de Catalan.
		// 4. Verificacion INSUFICIENTE: verificarFormula() solo prueba hasta
		//    "limite" (tipicamente n <= 20-30 por el costo de la fuerza
		//    bruta); coincidir en ese rango es evidencia fuerte pero NO una
		//    demostracion, especialmente en problemas con casos borde en
		//    valores grandes o con paridad/modulo involucrado.
		// 5. Formulas con recurrencia en vez de cerrada: si no se encuentra
		//    una formula cerrada pero sí una recurrencia lineal simple
		//    (f(n) = c1*f(n-1) + c2*f(n-2) + ...), esa recurrencia ya es
		//    suficiente para una DP eficiente O(n), sin necesidad de la
		//    forma cerrada.
		// 6. Automatizar la busqueda de recurrencia lineal: si se sospecha
		//    una recurrencia de orden k, se puede resolver un sistema lineal
		//    con los primeros 2k+1 terminos para encontrar los coeficientes
		//    (metodo de Berlekamp-Massey, avanzado, fuera del alcance de esta
		//    plantilla).
	};
\end{lstlisting}

\textbf{Ejemplo de funcionamiento:}
\begin{lstlisting}[style=cpp]
	vector<long long> tabla = FuerzaBrutaPatrones::generarTabla(8);
	// tabla = {1, 2, 3, 5, 8, 13, 21, 34}
	// (numero de formas de subir n escalones con pasos de 1 o 2; a simple
	//  vista, es la secuencia de Fibonacci desplazada un lugar)
	
	vector<long long> dif = FuerzaBrutaPatrones::diferencias(tabla);
	// dif = {1, 1, 2, 3, 5, 8, 13}
	// (las diferencias son la misma secuencia desplazada otro lugar mas,
	//  confirmando la naturaleza de Fibonacci del patron)
	
	// formula candidata: f(n) es el (n+1)-esimo numero de Fibonacci
	// (con Fib(1)=Fib(2)=1); se verifica contra la fuerza bruta
	auto fib = [](int n) {
		long long a = 1, b = 1;
		for (int i = 2; i <= n; i++) { long long c = a + b; a = b; b = c; }
		return b;
	};
	int primerFallo = FuerzaBrutaPatrones::verificarFormula(8, fib);
	// primerFallo = -1 (la formula coincide en todo el rango probado)
\end{lstlisting}

\textbf{Nota:} coincidir con la fuerza bruta en un rango pequeño (por ejemplo $n\le 20$) es una señal fuerte pero \textbf{no una demostración}. Problemas con casos borde en $n$ grande, dependencia de paridad, o condiciones que solo se activan a partir de cierto tamaño pueden hacer que una fórmula parezca correcta en la tabla pequeña y falle en el caso real. Cuando sea posible, conviene razonar \textit{por qué} la fórmula tiene sentido (no solo que coincide numéricamente) antes de enviarla con confianza total.

\textbf{Regla práctica:} usa esta técnica como primer paso exploratorio cuando el enunciado no sugiera una fórmula ni una DP evidente: programa la fuerza bruta más simple posible (aunque sea $O(2^n)$), genera una tabla de al menos 8-10 valores, busca un patrón en la secuencia o en sus diferencias, y siempre verifica la fórmula candidata contra la fuerza bruta antes de confiar en ella para el caso grande.

\section{Herramientas de Práctica}
\subsection{Stress Testing (Generador, Fuerza Bruta y Checker)}

\textbf{Idea:} el \textit{stress testing} es la técnica para encontrar el caso de prueba que rompe una solución antes de que lo haga el juez del concurso. Se necesitan tres piezas: un \textbf{generador} que produce entradas aleatorias pequeñas (con parámetros de tamaño ajustables), la solución \textbf{rápida} que se quiere verificar, y una solución de \textbf{fuerza bruta} (lenta pero obviamente correcta, como la de \textit{Fuerza Bruta para Descubrir Patrones}) que sirva de referencia. Un script (bash o un bucle en el propio código) genera una entrada, corre ambas soluciones sobre ella, y compara las salidas; si difieren, esa entrada queda guardada como el caso que falla, listo para depurar con un tamaño mucho más manejable que el del caso real que dio \textit{Wrong Answer}.

\textbf{¿Por qué usarlo?} Cuando el juez rechaza una solución con \textit{Wrong Answer} en un caso de miles de elementos, depurarla directamente es casi imposible: no se sabe en qué parte de la entrada está el error. Generando entradas \textbf{pequeñas} (5-10 elementos) y comparando automáticamente cientos o miles de veces por segundo contra la fuerza bruta, se encuentra en minutos un contraejemplo mínimo, mucho más fácil de razonar a mano o de recorrer con un depurador. Es la aplicación práctica, en forma de infraestructura reutilizable, de la misma idea de "verificar contra la fuerza bruta" que ya se vio como metodología en la subsección anterior.

\textbf{Casos de uso:}
\begin{itemize}
	\item Encontrar el contraejemplo mínimo que rompe una solución que dio \textit{Wrong Answer} en el juez.
	\item Verificar una optimización o reescritura de código contra la versión original antes de enviarla, para confirmar que el comportamiento no cambió.
	\item Depurar errores de casos borde (arreglos vacíos, un solo elemento, valores repetidos, todos iguales) generándolos con parámetros extremos del generador.
	\item Validar un \textit{checker} para problemas con múltiples respuestas válidas (no comparar salida exacta, sino verificar que la respuesta dada cumple las condiciones del problema).
	\item Comparar el rendimiento real de dos implementaciones con entradas de tamaño creciente (\textit{stress testing} de tiempo, no solo de corrección).
\end{itemize}

\textbf{Complejidad:} no aplica una complejidad de algoritmo — es una técnica de \textbf{flujo de trabajo}; el costo práctico es el número de iteraciones necesarias hasta encontrar un caso que falle, dominado por qué tan pequeño y específico es el generador respecto al bug real.

\begin{lstlisting}[style=cpp]
	// CREATE/READ: Plantilla de generador aleatorio + comparacion en el mismo
	// binario (alternativa a usar scripts bash separados), util cuando no se
	// tiene acceso a la terminal para armar el pipeline completo
	struct StressTest {
		mt19937 rng;
		
		// CREATE: inicializa el generador con una semilla EXTERNA (por
		// ejemplo, la hora actual o un contador de iteracion), para que cada
		// corrida produzca entradas distintas
		StressTest(unsigned semilla) : rng(semilla) {
		} // O(1)
		
		// READ: genera un vector de n enteros aleatorios en [minVal, maxVal];
		// n pequeño (5-10) es clave para que un contraejemplo sea facil de leer
		vector<int> generarArreglo(int n, int minVal, int maxVal) {
			uniform_int_distribution<int> dist(minVal, maxVal);
			vector<int> a(n);
			for (int& x : a) x = dist(rng);
			return a;
		} // O(n)
		
		// READ: solucion RAPIDA de ejemplo (a reemplazar por la que se quiere
		// verificar); aqui: suma maxima de un subarreglo (Kadane)
		static long long solucionRapida(const vector<int>& a) {
			long long mejor = a[0], actual = a[0];
			for (size_t i = 1; i < a.size(); i++) {
				actual = max((long long)a[i], actual + a[i]);
				mejor = max(mejor, actual);
			}
			return mejor;
		} // O(n)
		
		// READ: fuerza bruta de ejemplo (obviamente correcta, O(n^2)): prueba
		// todos los subarreglos
		static long long fuerzaBruta(const vector<int>& a) {
			long long mejor = a[0];
			for (size_t i = 0; i < a.size(); i++) {
				long long suma = 0;
				for (size_t j = i; j < a.size(); j++) {
					suma += a[j];
					mejor = max(mejor, suma);
				}
			}
			return mejor;
		} // O(n^2)
		
		// READ: corre iteraciones de stress test; retorna el primer arreglo
		// donde las dos soluciones difieren, o un vector vacio si no se
		// encontro ninguno en 'intentos' iteraciones
		vector<int> buscarContraejemplo(int intentos, int nMax) {
			uniform_int_distribution<int> distN(1, nMax);
			for (int it = 0; it < intentos; it++) {
				vector<int> a = generarArreglo(distN(rng), -10, 10);
				if (solucionRapida(a) != fuerzaBruta(a)) return a;
			}
			return {};
		} // O(intentos * costo de ambas soluciones)
		
		// ------------------------------------------------------------
		// MODIFICACIONES Y COMENTARIOS CON LOS USOS
		// ------------------------------------------------------------
		// 1. Pipeline con scripts separados (mas comun en competencia real):
		//    tres binarios independientes (gen.cpp, rapida.cpp, brute.cpp) y
		//    un script bash que en un bucle ejecuta
		//    ./gen semilla > input.txt; ./rapida < input.txt > out1.txt;
		//    ./brute < input.txt > out2.txt; diff out1.txt out2.txt, deteniendo
		//    el bucle en la primera diferencia; preferible cuando las
		//    soluciones ya estan escritas como programas separados que leen
		//    de cin/cout.
		// 2. Checker para problemas con MULTIPLES respuestas validas: en vez
		//    de comparar salida exacta, escribir una funcion que verifique si
		//    la respuesta dada cumple las condiciones del problema (por
		//    ejemplo, si es una permutacion valida, o si la suma reportada
		//    coincide con recalcularla sobre la solucion propuesta).
		// 3. Minimizar el contraejemplo encontrado: una vez hallado un caso
		//    que falla, intentar reducirlo (quitar elementos, achicar el
		//    rango de valores) mientras la solucion rapida y la fuerza bruta
		//    sigan discrepando, para llegar al caso MINIMO que reproduce el bug.
		// 4. Semilla reproducible: guardar la semilla usada en el intento que
		//    fallo permite regenerar exactamente el mismo caso despues, sin
		//    necesidad de guardar el arreglo completo en memoria.
		// 5. Generadores mas complejos (arboles, grafos, cadenas con
		//    estructura especifica): el mismo esquema aplica, cambiando
		//    generarArreglo() por una funcion que construya la estructura de
		//    entrada requerida por el problema especifico.
		// 6. Reemplazar solucionRapida y fuerzaBruta: las incluidas aqui
		//    (Kadane vs. fuerza bruta O(n^2)) son solo un ejemplo ilustrativo;
		//    en la practica se sustituyen por el par real que se quiere
		//    comparar.
	};
\end{lstlisting}

\textbf{Ejemplo de funcionamiento:}
\begin{lstlisting}[style=cpp]
	StressTest st(12345); // semilla fija para reproducibilidad del ejemplo
	
	vector<int> contraejemplo = st.buscarContraejemplo(1000, 8);
	// contraejemplo esta vacio si las dos soluciones (Kadane vs fuerza bruta)
	// coinciden en las 1000 entradas aleatorias probadas -- es el resultado
	// esperado aqui porque Kadane es correcto; en un caso real, un vector
	// no vacio séria el contraejemplo a depurar
\end{lstlisting}

\textbf{Nota:} en competencia, la forma más común de montar esto es con \textbf{tres archivos separados} (generador, solución candidata, fuerza bruta) y un script de shell que los ejecuta en bucle comparando con \texttt{diff} — la versión de arriba, con todo en un único \texttt{struct} de C++, es útil cuando no se tiene acceso cómodo a la terminal (por ejemplo, dentro de un IDE en línea), pero el flujo de trabajo conceptual es el mismo en ambos casos.

\textbf{Regla práctica:} arma el generador \textbf{antes} de necesitarlo desesperadamente: en cuanto una solución da \textit{Wrong Answer} sin pista clara de dónde, escribe la fuerza bruta más simple posible (aunque sea $O(2^n)$) y un generador de entradas pequeñas, y deja correr la comparación en bucle. Casi siempre es más rápido que revisar el código línea por línea buscando el error a ciegas.

\subsection{Problemas Interactivos}

\textbf{Idea:} en un problema \textbf{interactivo}, el programa no lee toda la entrada de una vez ni imprime toda la salida al final: \textbf{dialoga} con un juez (\textit{interactor}) que responde a cada consulta antes de que el programa decida la siguiente. El patrón típico es un bucle de "preguntar, leer la respuesta del juez, decidir la siguiente pregunta" hasta encontrar la respuesta final —el ejemplo clásico es adivinar un número oculto haciendo consultas de tipo "¿es mayor que $x$?" y usando \textit{Búsqueda Binaria on Answer} para decidir la siguiente consulta. La dificultad técnica añadida, distinta de cualquier otro problema de esta guía, es que la salida del programa debe llegar al juez \textbf{inmediatamente} después de cada consulta, sin quedar retenida en un búfer.

\textbf{¿Por qué usarlo?} Casi cualquier técnica algorítmica de esta guía (búsqueda binaria, greedy, teoría de juegos) puede aparecer dentro de un problema interactivo; lo que cambia es exclusivamente la \textbf{forma de comunicación} con el juez, no el algoritmo en sí. El error más común y más difícil de diagnosticar en estos problemas no es de lógica, sino de \textbf{buffering}: si la salida no se vacía (\textit{flush}) después de cada consulta, el juez nunca la recibe, el programa se queda esperando una respuesta que no llega, y el resultado es un \textit{Time Limit Exceeded} o \textit{Idleness Limit Exceeded} que nada tiene que ver con el algoritmo implementado.

\textbf{Casos de uso:}
\begin{itemize}
	\item Adivinar un valor oculto mediante consultas de comparación, resuelto con búsqueda binaria interactiva.
	\item Determinar una propiedad de una estructura oculta (un árbol, un grafo) haciendo consultas limitadas sobre ella.
	\item Juegos contra un juez adversarial que responde de la peor manera posible dentro de las reglas —el juez puede incluso no fijar la respuesta de antemano, siempre que sea consistente con las respuestas ya dadas.
	\item Problemas con límite estricto en el \textbf{número} de consultas permitidas, donde minimizar consultas (no solo el tiempo de cómputo) es parte del objetivo.
	\item Práctica de manejo correcto de entrada/salida en tiempo real, útil también para depurar cualquier código que interactúe con un proceso externo.
\end{itemize}

\textbf{Complejidad:} depende enteramente del problema — se mide típicamente en \textbf{número de consultas} permitidas (por ejemplo $O(\log n)$ para una búsqueda binaria interactiva) en vez de en tiempo de cómputo puro, ya que el cómputo entre consultas suele ser trivial frente a la limitación de cuántas preguntas se pueden hacer.

\begin{lstlisting}[style=cpp]
	// CREATE/READ/WRITE: Plantilla de problema interactivo: adivinar un
	// numero oculto en [1, n] mediante consultas "es x mayor que el oculto?",
	// usando busqueda binaria y flush explicito tras cada consulta
	struct ProblemaInteractivo {
		int n;
		
		// CREATE: guarda el rango de busqueda
		ProblemaInteractivo(int limite) : n(limite) {
		} // O(1)
		
		// WRITE: envia una consulta al juez y lee su respuesta; el flush es
		// OBLIGATORIO despues de cada consulta, o el juez nunca la recibe
		char preguntar(int x) {
			cout << "? " << x << "\n";
			cout.flush(); // CRITICO: sin esto el programa se cuelga esperando respuesta
			char respuesta;
			cin >> respuesta; // tipicamente 'Y'/'N' o '<'/'>'/'='segun el problema
			return respuesta;
		} // O(1) por consulta, mas el costo de comunicacion con el juez
		
		// READ/WRITE: busqueda binaria interactiva: encuentra el numero
		// oculto en [1, n] preguntando "es x mayor que el oculto?" en cada paso
		int adivinar() {
			int lo = 1, hi = n;
			while (lo < hi) {
				int medio = lo + (hi - lo) / 2;
				char resp = preguntar(medio);
				if (resp == 'Y') { // medio es mayor que el oculto: buscar mas abajo
					hi = medio;
				} else { // medio es menor o igual: buscar mas arriba
					lo = medio + 1;
				}
			}
			cout << "! " << lo << "\n"; // respuesta FINAL, tambien con flush
			cout.flush();
			return lo;
		} // O(log n) consultas
		
		// ------------------------------------------------------------
		// MODIFICACIONES Y COMENTARIOS CON LOS USOS
		// ------------------------------------------------------------
		// 1. Sincronizacion de flujos en C++: al inicio del programa, usar
		//    cin.tie(nullptr) con cuidado en interactivos -- de hecho, MEJOR
		//    NO desactivar la sincronizacion automatica (NO usar
		//    ios::sync_with_stdio(false) sin verificar), o usar endl en vez de
		//    "\n" + flush manual, ya que endl hace flush automaticamente y es
		//    mas dificil de olvidar en cada linea de salida.
		// 2. Interactor "adaptativo" (adversarial): el juez puede no fijar el
		//    numero oculto de antemano, y en cambio mantener el conjunto de
		//    valores consistentes con las respuestas dadas hasta ahora,
		//    respondiendo siempre de la forma que mas complique al programa;
		//    la solucion debe funcionar igual sin asumir que el valor esta
		//    fijo desde el inicio, solo que las respuestas son consistentes.
		// 3. Limite de consultas MENOR al esperado (por ejemplo, permite solo
		//    log2(n) - 2 preguntas en vez de log2(n)): senal de que se
		//    necesita una estrategia mas eficiente que la busqueda binaria
		//    directa, como aprovechar informacion de consultas anteriores de
		//    forma no binaria.
		// 4. Terminacion temprana o formato de respuesta final incorrecto:
		//    revisar con cuidado el formato exacto que pide el enunciado para
		//    la respuesta final (aqui "! numero"), ya que un formato distinto
		//    al esperado por el juez cuenta como respuesta incorrecta aunque
		//    el numero en si sea correcto.
		// 5. Testing local de interactivos: como no hay juez real en el propio
		//    entorno, se escribe un interactor propio (otro programa que
		//    simula al juez) y se conectan ambos procesos por pipes o
		//    redireccion de archivos para probar antes de enviar.
	};
\end{lstlisting}

\textbf{Ejemplo de funcionamiento:}
\begin{lstlisting}[style=cpp]
	// Interaccion tipica con el juez (E/S real, no simulable sin un
	// interactor externo):
	// programa: ? 5
	// juez:     N          (el oculto es >= 5)
	// programa: ? 8
	// juez:     Y          (el oculto es < 8, es decir <= 7)
	// programa: ? 6
	// juez:     N          (el oculto es >= 6, es decir 6 o 7)
	// programa: ? 7
	// juez:     N          (el oculto es 7)
	// programa: ! 7
	// juez:     AC (o el veredicto correspondiente)
\end{lstlisting}

\textbf{Nota:} el fallo más frecuente en problemas interactivos \textbf{no es de lógica}, sino de \textbf{buffering}: olvidar \texttt{cout.flush()} (o no usar \texttt{endl}, que hace flush automáticamente) después de cada consulta deja la salida retenida en el búfer del programa, el juez nunca la recibe, y el resultado es un \textit{Time Limit Exceeded} que parece un error de algoritmo pero es puramente de entrada/salida. Verificar esto es siempre el primer paso al depurar un interactivo que se cuelga.

\textbf{Regla práctica:} en cualquier problema interactivo, antes de razonar el algoritmo, asegúrate de hacer \texttt{flush} después de \textbf{cada} línea de salida (o usa \texttt{endl} en vez de \texttt{"\textbackslash n"} consistentemente). El algoritmo subyacente reutiliza técnicas ya vistas en la guía (aquí, búsqueda binaria) — lo específico de esta subsección es la disciplina de comunicación con el juez, no una técnica algorítmica nueva.

\part{String Processing}
\subsection{¿Qué algoritmo usar?}

% ---------------------------------------------------------
% TABLA 1: ESTRUCTURAS Y BUSQUEDA DE PATRONES
% ---------------------------------------------------------
\begin{table}[H]
	\raggedright
	\scriptsize
	\setlength{\tabcolsep}{2.5pt}
	\renewcommand{\arraystretch}{1.15}
	\begin{tabular}{|p{0.43\columnwidth}|p{0.22\columnwidth}|p{0.27\columnwidth}|}
		\hline
		\textbf{Cuando el ejercicio diga / pida} &
		\textbf{Usa} &
		\textbf{Pista rápida} \\
		\hline
		
		``Autocompletado'', ``cuántas palabras empiezan con este prefijo'', ``prefijo común'' &
		Trie &
		Contador por nodo. $O(|S|)$ por operación. \\
		\hline
		
		``¿Estas dos subcadenas son iguales?'' con muchas consultas &
		Hash polinómico &
		Hash doble si hay riesgo adversarial. $O(1)$ por consulta. \\
		\hline
		
		``Prefijo común más largo entre dos posiciones'' &
		Hash + búsqueda binaria &
		Compara hashes de prefijos de longitud $L$. \\
		\hline
		
		``Período más corto'', ``¿la cadena es una repetición de un bloque?'' &
		Función de prefijos &
		$n-\pi[n-1]$ si divide a $n$. \\
		\hline
		
		``Todas las ocurrencias de un patrón en un texto'' &
		KMP &
		$p\,\#\,t$ y busca $\pi[i]=|p|$. $O(n+m)$. \\
		\hline
		
		Lo mismo, o ``cuánto coincide cada sufijo con el prefijo'' &
		Algoritmo Z &
		Equivalente a KMP. $O(n)$. \\
		\hline
		
		Búsqueda de un patrón priorizando código corto &
		Rabin-Karp &
		Verifica el candidato. Esperado $O(n+m)$. \\
		\hline
		
		Texto natural con alfabeto grande, rendimiento práctico &
		Boyer-Moore &
		Peor caso $O(nm)$; en CP prefiere KMP. \\
		\hline
		
		Mismo patrón contra muchos textos, o texto en streaming &
		Autómata de coincidencia &
		Tabla $\delta$. $O(1)$ por carácter. \\
		\hline
		
		``Buscar \textbf{muchos patrones} a la vez'', diccionario, palabras prohibidas &
		Aho-Corasick &
		Trie + enlaces de fallo. Hereda coincidencias. \\
		\hline
		
		``Cuántas cadenas de longitud $n$ evitan estos patrones'' &
		DP sobre Aho-Corasick &
		Estados = nodos; descarta nodos con patrón. \\
		\hline
		
		Patrón con comodines \texttt{?} &
		Convolución (FFT) &
		$O(nm)$ si es pequeño; $O(n\log n)$ con FFT. \\
		\hline
		
		``¿Algún anagrama del patrón aparece en el texto?'' &
		Ventana deslizante de frecuencias &
		Contador de diferencias. $O(n+\sigma)$. \\
		\hline
		
	\end{tabular}
	\caption{Estructuras y búsqueda de patrones: qué usar según lo que pide el ejercicio.}
\end{table}

% ---------------------------------------------------------
% TABLA 2: SUFIJOS, PALINDROMOS Y COMUNES
% ---------------------------------------------------------
\begin{table}[H]
	\raggedright
	\scriptsize
	\setlength{\tabcolsep}{2.5pt}
	\renewcommand{\arraystretch}{1.15}
	\begin{tabular}{|p{0.43\columnwidth}|p{0.22\columnwidth}|p{0.27\columnwidth}|}
		\hline
		\textbf{Cuando el ejercicio diga / pida} &
		\textbf{Usa} &
		\textbf{Pista rápida} \\
		\hline
		
		``Cuántos substrings \textbf{distintos} hay'' &
		Suffix Array + LCP o Suffix Automaton &
		$\binom{n+1}{2}-\sum \text{lcp}$, o suma de $\text{len}[v]-\text{len}[\text{link}[v]]$. \\
		\hline
		
		Muchas consultas sobre subcadenas, LCP entre dos sufijos, orden lexicográfico &
		Suffix Array + Kasai + Sparse Table &
		LCP = mínimo en rango. $O(1)$. \\
		\hline
		
		``¿$t$ es substring de $s$?'', construir en línea, ocurrencias de cada substring &
		Suffix Automaton &
		$O(|t|)$ por consulta. \\
		\hline
		
		``Subcadena común \textbf{contigua} más larga'' &
		DP 2D (reinicia a 0) &
		$O(nm)$; con muchas cadenas, Suffix Automaton. \\
		\hline
		
		``Subsecuencia común más larga'', diff, alinear secuencias &
		LCS &
		$O(nm)$. Solo inserción/borrado: $n+m-2\,\text{LCS}$. \\
		\hline
		
		``Cadena más corta que contenga a ambas como subsecuencia'' &
		Supersecuencia común más corta &
		$n+m-\text{LCS}$. \\
		\hline
		
		``Mínimo de operaciones para transformar una cadena en otra'' (insertar, borrar, sustituir) &
		Distancia de edición &
		$O(nm)$. \\
		\hline
		
		``Subsecuencia palindrómica más larga'', ``mínimo de borrados para palíndromo'' &
		DP de intervalos o LCS con el reverso &
		$n-\text{LPS}$ borrados. \\
		\hline
		
		``Subsecuencia creciente (decreciente) más larga'', ``bitónica'' &
		LIS / LDS &
		$O(n\log n)$ con \texttt{colas}. Bitónica: $\text{LIS}[i]+\text{LDS}[i]-1$. \\
		\hline
		
		``Cuántos substrings palíndromos'' (con repetición de posición) &
		Manacher &
		Suma de $\lceil \text{rad}/2\rceil$. $O(n)$. \\
		\hline
		
		``Cuántos palíndromos \textbf{distintos}'', palíndromo más largo en cada prefijo &
		Eertree &
		$\text{nodos}-2$. $O(n)$. \\
		\hline
		
		``Mínimo de cortes para partir en palíndromos'' &
		Tabla \texttt{esPal} + DP lineal &
		$O(n^2)$. \\
		\hline
		
	\end{tabular}
	\caption{Sufijos, palíndromos y comunes: qué usar según lo que pide el ejercicio.}
\end{table}

% ---------------------------------------------------------
% TABLA 3: CONSTRUCCIONES Y ANAGRAMAS
% ---------------------------------------------------------
\begin{table}[H]
	\raggedright
	\scriptsize
	\setlength{\tabcolsep}{2.5pt}
	\renewcommand{\arraystretch}{1.15}
	\begin{tabular}{|p{0.43\columnwidth}|p{0.22\columnwidth}|p{0.27\columnwidth}|}
		\hline
		\textbf{Cuando el ejercicio diga / pida} &
		\textbf{Usa} &
		\textbf{Pista rápida} \\
		\hline
		
		``Rotación lexicográficamente mínima'', ``¿son rotaciones una de otra?'' &
		Booth o Duval sobre $s+s$ &
		Forma canónica. $O(n)$. \\
		\hline
		
		``Factorización de Lyndon'', ``sufijo lexicográficamente menor'' &
		Algoritmo de Duval &
		Compara con \texttt{<=}. $O(n)$. \\
		\hline
		
		``\texttt{a} solo a la izquierda, \texttt{b} solo a la derecha, sin cruzarse, ¿se puede pasar de $S_1$ a $S_2$?'' &
		Verificación de dos condiciones &
		Misma secuencia sin \texttt{\#} y dirección por índice. $O(n)$. \\
		\hline
		
		``¿Son anagramas?'' &
		Conteo de frecuencia &
		$O(n+\sigma)$; si no, ordenar. \\
		\hline
		
		``Agrupa las palabras que son anagramas'' &
		Firma canónica + hash &
		Ordenada o vector de frecuencias con separador. \\
		\hline
		
		``Substrings que se pueden reordenar a palíndromo'' &
		Máscara de paridad + hashmap &
		Misma máscara o un bit de diferencia. \\
		\hline
		
		``Reordena sin letras iguales adyacentes'', ``anagrama lexicográfico con restricción'' &
		Max-heap por frecuencia &
		Imposible si $f_{\max}>\lceil n/2\rceil$. \\
		\hline
		
	\end{tabular}
	\caption{Construcciones y anagramas: qué usar según lo que pide el ejercicio.}
\end{table}
\section{Fundamentos}
\subsection{Trie}

\textbf{Idea:} un trie (árbol de prefijos) almacena un conjunto de cadenas organizándolas por prefijos compartidos: cada nodo representa un prefijo, y sus hijos extienden ese prefijo en un carácter. Dos cadenas que comparten un prefijo comparten el mismo camino desde la raíz hasta donde sus caracteres divergen, momento en el cual el árbol se ramifica. Esto convierte operaciones sobre el conjunto completo de cadenas (insertar, buscar, contar cuántas empiezan con cierto prefijo) en un simple recorrido de árbol, de longitud igual a la de la cadena consultada, sin importar cuántas cadenas haya en total.

\textbf{¿Por qué usarlo?} Buscar si una cadena existe en un conjunto de $n$ cadenas con un \texttt{set<string>} cuesta $O(|S| \log n)$ por la comparación de strings dentro de cada paso del árbol balanceado. Un trie resuelve la misma consulta en $O(|S|)$, independiente de $n$, porque cada carácter de la consulta se procesa una sola vez bajando un nivel. La ventaja se vuelve decisiva en consultas de \textbf{prefijo} (¿cuántas cadenas empiezan con "ab"?), que un \texttt{set} no puede responder eficientemente sin recorrer rangos, mientras que en un trie es simplemente llegar al nodo del prefijo y leer un contador.

\textbf{Casos de uso:}
\begin{itemize}
    \item Autocompletado: listar todas las cadenas que comparten un prefijo dado.
    \item Contar cuántas cadenas de un conjunto tienen cierto prefijo.
    \item Verificar pertenencia de una cadena a un diccionario en $O(|S|)$.
    \item Base estructural para \textit{Aho-Corasick} (múltiples patrones) y para \textit{Trie de Bits} (XOR máximo entre pares).
    \item Encontrar el prefijo común más largo entre dos o más cadenas del conjunto.
\end{itemize}

\textbf{Complejidad:} $O(|S|)$ por inserción o consulta, donde $|S|$ es la longitud de la cadena procesada. El espacio total es $O(\sum |S_i| \cdot \sigma)$ en el peor caso, donde $\sigma$ es el tamaño del alfabeto, ya que cada nodo reserva espacio para todos los caracteres posibles del alfabeto.

\begin{lstlisting}[style=cpp]
// CREATE/READ/WRITE: Trie para insertar cadenas y consultar pertenencia
// y conteo de prefijos, sobre alfabeto minusculas 'a'-'z'
struct Trie {
    static const int ALFABETO = 26;
    struct Nodo {
        int hijos[ALFABETO];
        int contadorPrefijo; // cuantas cadenas insertadas pasan por este nodo
        bool esFinDePalabra;
        Nodo() : contadorPrefijo(0), esFinDePalabra(false) {
            fill(hijos, hijos + ALFABETO, -1);
        }
    };

    vector<Nodo> nodos;

    // CREATE: inicializa el trie con un unico nodo raiz
    Trie() {
        nodos.push_back(Nodo());
    } // O(1)

    // WRITE: inserta la cadena s en el trie, creando nodos nuevos donde
    // el camino aun no existe
    void insertar(const string& s) {
        int actual = 0;
        for (char c : s) {
            int idx = c - 'a';
            if (nodos[actual].hijos[idx] == -1) {
                nodos[actual].hijos[idx] = (int)nodos.size();
                nodos.push_back(Nodo());
            }
            actual = nodos[actual].hijos[idx];
            nodos[actual].contadorPrefijo++;
        }
        nodos[actual].esFinDePalabra = true;
    } // O(|s|)

    // READ: true si la cadena s fue insertada exactamente (no solo como prefijo)
    bool buscar(const string& s) {
        int actual = 0;
        for (char c : s) {
            int idx = c - 'a';
            if (nodos[actual].hijos[idx] == -1) return false;
            actual = nodos[actual].hijos[idx];
        }
        return nodos[actual].esFinDePalabra;
    } // O(|s|)

    // READ: cuenta cuantas cadenas insertadas tienen a s como prefijo
    int contarConPrefijo(const string& s) {
        int actual = 0;
        for (char c : s) {
            int idx = c - 'a';
            if (nodos[actual].hijos[idx] == -1) return 0;
            actual = nodos[actual].hijos[idx];
        }
        return nodos[actual].contadorPrefijo;
    } // O(|s|)

    // ------------------------------------------------------------
    // MODIFICACIONES Y COMENTARIOS CON LOS USOS
    // ------------------------------------------------------------
    // 1. Alfabeto mas grande o con simbolos especiales (mayusculas, digitos,
    //    caracteres especiales): usar un unordered_map<char,int> por nodo en
    //    vez de un arreglo fijo de tamaño ALFABETO, a costa de mayor
    //    constante de tiempo por acceso.
    // 2. Eliminar una cadena del trie: decrementar contadorPrefijo a lo
    //    largo del camino de insertar(); si un nodo llega a contadorPrefijo
    //    == 0 se puede marcar como libre (no es necesario liberarlo
    //    fisicamente del vector para que la estructura siga correcta).
    // 3. Prefijo comun mas largo entre TODAS las cadenas insertadas: bajar
    //    desde la raiz mientras cada nodo tenga contadorPrefijo igual al
    //    total de cadenas insertadas Y tenga exactamente un hijo activo.
};
\end{lstlisting}

\textbf{Ejemplo de funcionamiento:}
\begin{lstlisting}[style=cpp]
Trie t;
t.insertar("casa");
t.insertar("casino");
t.insertar("caso");

bool existe = t.buscar("caso");
// existe = true

int conPrefijo = t.contarConPrefijo("cas");
// conPrefijo = 3
// (las tres cadenas insertadas comparten el prefijo "cas")

bool existeParcial = t.buscar("cas");
// existeParcial = false
// ("cas" es un prefijo valido en el trie, pero nunca se inserto como
//  palabra completa, asi que esFinDePalabra es false en ese nodo)
\end{lstlisting}

\textbf{Nota:} \texttt{buscar} y \texttt{contarConPrefijo} son operaciones distintas y no intercambiables: \texttt{buscar} exige que la cadena completa haya sido insertada como palabra (nodo marcado \texttt{esFinDePalabra}), mientras que \texttt{contarConPrefijo} solo exige que el camino exista, sin importar si algún nodo intermedio es fin de palabra.

\textbf{Regla práctica:} usa un trie cuando las consultas sean sobre \textbf{prefijos} de un conjunto de cadenas (autocompletado, conteo de prefijos, prefijo común); si solo necesitas verificar pertenencia exacta sin ninguna consulta de prefijo, un \texttt{unordered\_set<string>} con hash es más simple y con menor overhead de memoria.

\subsection{Hash Polinómico (Rolling Hash)}

\textbf{Idea:} un hash polinómico convierte una cadena $s_0 s_1 \dots s_{n-1}$ en un número mediante $H(s) = \left(\sum_{i=0}^{n-1} s_i \cdot p^i\right) \bmod m$, tratando cada carácter como un "dígito" en base $p$ (una base prima mayor que el tamaño del alfabeto). La propiedad clave que lo hace útil en competencia es que, precomputando los hashes de \textbf{todos los prefijos} de una cadena junto con las potencias de $p$, se puede obtener el hash de \textbf{cualquier subcadena} en $O(1)$ mediante una resta modular, de forma análoga a cómo una suma de prefijos da la suma de cualquier subarreglo en $O(1)$.

\textbf{¿Por qué usarlo?} Comparar si dos subcadenas son iguales carácter por carácter cuesta $O(|s|)$. Con hash polinómico, comparar dos subcadenas se reduce a comparar dos enteros precomputados en $O(1)$ —a costa de una probabilidad extremadamente pequeña de colisión (dos subcadenas distintas con el mismo hash). Es la base de muchos algoritmos de comparación rápida de subcadenas, y una alternativa más simple de implementar que \textit{KMP} o \textit{Algoritmo Z} cuando lo único que se necesita es comparar igualdad de subcadenas, no encontrar todas las ocurrencias de un patrón.

\textbf{Casos de uso:}
\begin{itemize}
    \item Comparar igualdad de dos subcadenas cualesquiera en $O(1)$ tras $O(n)$ de preprocesamiento.
    \item Búsqueda binaria sobre la longitud del prefijo común más largo entre dos subcadenas.
    \item Detección rápida de subcadenas repetidas o duplicadas.
    \item Verificación alternativa (más simple de codear bajo presión) a KMP/Rabin-Karp para coincidencia de un patrón.
    \item Comparación de igualdad de subarreglos en problemas que no son estrictamente de texto (secuencias numéricas tratadas como "cadenas").
\end{itemize}

\textbf{Complejidad:} $O(n)$ de preprocesamiento para una cadena de longitud $n$ (calcular hashes de prefijos y potencias de $p$); $O(1)$ por consulta de hash de una subcadena arbitraria.

\begin{lstlisting}[style=cpp]
// CREATE/READ: Hash polinomico con hashes de prefijo precomputados, para
// consultar el hash de cualquier subcadena en O(1)
struct HashPolinomico {
    const long long P = 131;                 // base, mayor que el tamaño del alfabeto
    const long long MOD = 1000000007LL;      // modulo primo grande

    vector<long long> prefijoHash; // prefijoHash[i] = hash de s[0..i-1]
    vector<long long> potencias;   // potencias[i] = P^i mod MOD

    // CREATE: precomputa hashes de todos los prefijos y las potencias de P
    // necesarias para consultar cualquier subcadena de s
    HashPolinomico(const string& s) {
        int n = (int)s.size();
        prefijoHash.assign(n + 1, 0);
        potencias.assign(n + 1, 1);
        for (int i = 0; i < n; i++) {
            prefijoHash[i + 1] = (prefijoHash[i] * P + (s[i] - 'a' + 1)) % MOD;
            potencias[i + 1] = (potencias[i] * P) % MOD;
        }
    } // O(n)

    // READ: hash de la subcadena s[l..r] (0-indexado, ambos extremos
    // incluidos), en O(1)
    long long hashRango(int l, int r) {
        long long resultado = (prefijoHash[r + 1] - (prefijoHash[l] * potencias[r - l + 1]) % MOD + MOD * MOD) % MOD;
        return resultado;
    } // O(1)

    // ------------------------------------------------------------
    // MODIFICACIONES Y COMENTARIOS CON LOS USOS
    // ------------------------------------------------------------
    // 1. Hashing DOBLE (recomendado en competencia): usar dos pares (P,MOD)
    //    distintos y comparar igualdad SOLO si ambos hashes coinciden;
    //    reduce drasticamente la probabilidad de colision adversarial,
    //    especialmente si los problemas del concurso son generados
    //    conociendo el MOD estandar de 1e9+7.
    // 2. MOD * MOD en hashRango: se suma para evitar que la resta modular
    //    de resultado negativo en C++ (donde % puede devolver negativo);
    //    alternativa mas limpia: sumar solo MOD si el resultado intermedio
    //    ya cabe sin overflow con ese ajuste.
    // 3. Comparar igualdad de dos subcadenas de la MISMA cadena o de
    //    cadenas DISTINTAS: si son de cadenas distintas, cada una necesita
    //    su propio HashPolinomico (mismo P y MOD para que los hashes sean
    //    comparables entre si).
};
\end{lstlisting}

\textbf{Ejemplo de funcionamiento:}
\begin{lstlisting}[style=cpp]
HashPolinomico h("abcabc");

long long hash1 = h.hashRango(0, 2); // "abc"
long long hash2 = h.hashRango(3, 5); // "abc"

bool iguales = (hash1 == hash2);
// iguales = true
// (ambas subcadenas son "abc", literalmente iguales; sus hashes deben
//  coincidir; hay que correrlo para confirmar el valor numerico exacto
//  de hash1 con P=131 y MOD=1e9+7)
\end{lstlisting}

\textbf{Nota:} un solo hash (un solo par $P,\text{MOD}$) es vulnerable a colisiones adversariales cuando el problema fue diseñado sabiendo qué módulo suelen usar los concursantes (ataques conocidos con $\text{MOD}=10^9+7$). En problemas donde alguien pudo haber preparado un caso de prueba adversarial contra hashing, usa \textbf{hashing doble} (ver modificación 1 dentro del código) en vez de un solo hash.

\textbf{Regla práctica:} usa hash polinómico cuando necesites comparar igualdad de subcadenas repetidamente y rápido, especialmente combinado con búsqueda binaria (por ejemplo, para encontrar el prefijo común más largo entre dos posiciones). Si el problema requiere encontrar \textbf{todas las ocurrencias} de un patrón (no solo comparar igualdad), \textit{KMP} o el \textit{Algoritmo Z} dan esa respuesta directamente sin el riesgo, aunque pequeño, de colisión de hash.
\subsection{Función de Prefijos (Prefix Function / Failure Function)}

\textbf{Idea:} para una cadena $s$ de longitud $n$, la función de prefijos $\pi[i]$ da la longitud del prefijo propio más largo de $s[0..i]$ que también es sufijo de ese mismo segmento $s[0..i]$ —"propio" significa que no puede ser el segmento completo. Se calcula incrementalmente: al procesar $s[i]$, se intenta extender el mejor prefijo-sufijo conocido en la posición anterior ($\pi[i-1]$); si el siguiente carácter no calza, se retrocede usando el propio arreglo $\pi$ (siguiendo la cadena de "fallos" hacia prefijos-sufijos más cortos) en vez de reiniciar la búsqueda desde cero, evitando así retroceder sobre caracteres ya comparados.

\textbf{¿Por qué usarlo?} Es la estructura de datos que hace posible que \textit{KMP} busque un patrón en $O(n+m)$ en vez de $O(nm)$: sin la función de prefijos, cada intento fallido de coincidencia obligaría a retroceder el puntero del texto y volver a comparar caracteres ya vistos. El truco de "retroceder usando $\pi$ en vez de reiniciar" es lo que garantiza que el puntero del texto \textbf{nunca retrocede}, dando la complejidad lineal. Aunque casi siempre se presenta como parte de KMP, la función de prefijos por sí sola también responde preguntas sobre la estructura periódica interna de una cadena.

\textbf{Casos de uso:}
\begin{itemize}
    \item Como subrutina central de \textit{Knuth-Morris-Pratt (KMP)} para búsqueda de patrones.
    \item Encontrar el período más corto de una cadena (la cadena se repite con período $n - \pi[n-1]$ si $n$ es divisible entre ese valor).
    \item Contar cuántas veces aparece cada prefijo de $s$ como sufijo en distintas posiciones (usando el árbol de fallos formado por $\pi$).
    \item Comprimir una cadena detectando su unidad repetitiva más pequeña.
\end{itemize}

\textbf{Complejidad:} $O(n)$ —a primera vista el retroceso dentro del bucle parece poder ser costoso, pero un argumento de potencial (el valor de $\pi[i]$ solo puede aumentar en 1 por iteración del bucle externo, y cada retroceso lo reduce, así que la suma total de retrocesos está acotada por $n$) garantiza tiempo lineal amortizado.

\begin{lstlisting}[style=cpp]
// CREATE/READ: Funcion de prefijos (failure function) de una cadena,
// base para KMP y para analisis de periodicidad
struct FuncionPrefijos {
    // READ: calcula pi[i] = longitud del prefijo propio mas largo de
    // s[0..i] que tambien es sufijo de s[0..i], para cada i
    static vector<int> calcular(const string& s) {
        int n = (int)s.size();
        vector<int> pi(n, 0);
        for (int i = 1; i < n; i++) {
            int j = pi[i - 1];
            // retrocede por la cadena de fallos mientras el caracter no calce
            while (j > 0 && s[i] != s[j]) j = pi[j - 1];
            if (s[i] == s[j]) j++;
            pi[i] = j;
        }
        return pi;
    } // O(n)

    // READ: longitud del periodo mas corto de s (la unidad minima que,
    // repetida, reconstruye s); retorna n si s no es periodica
    static int periodoMasCorto(const string& s) {
        int n = (int)s.size();
        vector<int> pi = calcular(s);
        int candidato = n - pi[n - 1];
        if (n % candidato == 0) return candidato;
        return n; // s no tiene periodo propio, el periodo es ella misma
    } // O(n)

    // ------------------------------------------------------------
    // MODIFICACIONES Y COMENTARIOS CON LOS USOS
    // ------------------------------------------------------------
    // 1. Uso directo en KMP: para buscar un patron p en un texto t, se
    //    calcula pi sobre la cadena concatenada p + '#' + t (con un
    //    separador que no aparece en ninguna de las dos), y cada posicion
    //    donde pi[i] == |p| marca una ocurrencia completa del patron.
    // 2. Todos los "bordes" de un prefijo (no solo el mas largo): seguir
    //    la cadena j -> pi[j-1] -> pi[pi[j-1]-1] -> ... -> 0 desde
    //    pi[i] da TODOS los prefijos-sufijos de s[0..i], de mayor a menor.
    // 3. Contar ocurrencias de cada prefijo como sufijo en TODA la cadena:
    //    construir un arbol donde cada nodo i apunta a pi[i]-1 (su "padre"
    //    de fallo), y propagar conteos desde las hojas hacia la raiz.
};
\end{lstlisting}

\textbf{Ejemplo de funcionamiento:}
\begin{lstlisting}[style=cpp]
vector<int> pi = FuncionPrefijos::calcular("abcabcabc");
// pi = {0, 0, 0, 1, 2, 3, 4, 5, 6}
// (en la posicion 8, "abcabcabc" tiene como prefijo-sufijo mas largo a
//  "abcabc", de longitud 6)

int periodo = FuncionPrefijos::periodoMasCorto("abcabcabc");
// periodo = 3
// (n=9, pi[8]=6, candidato = 9-6 = 3; 9 % 3 == 0, asi que "abc"
//  repetido 3 veces reconstruye la cadena original)
\end{lstlisting}

\textbf{Nota:} el cálculo de \texttt{periodoMasCorto} depende de que $n$ sea exactamente divisible entre el candidato $n-\pi[n-1]$; si no lo es, la cadena no tiene un período propio que la genere exactamente, y el valor correcto a reportar es $n$ (la cadena completa es su propio período mínimo).

\textbf{Regla práctica:} usa la función de prefijos como bloque de construcción cada vez que necesites la estructura de "qué tan lejos puedo retroceder sin perder coincidencia" —directamente como base de \textit{KMP}, o de forma independiente cuando la pregunta sea sobre periodicidad o bordes de una cadena. Si solo necesitas comparar igualdad de subcadenas (sin buscar patrones), \textit{Hash Polinómico} suele ser más simple de codear para ese propósito específico.

\subsection{Suffix Array + Algoritmo de Kasai (LCP)}

\textbf{Idea:} el \textit{suffix array} de una cadena $s$ de longitud $n$ es la lista de los $n$ sufijos de $s$, representados por sus índices de inicio, \textbf{ordenados lexicográficamente}. Construirlo por fuerza bruta (ordenar comparando sufijos completo contra completo) cuesta $O(n^2\log n)$. La construcción eficiente usa \textbf{ordenamiento por duplicación} (\textit{doubling}): se ordenan primero los sufijos por su primer carácter, luego por sus primeros 2 caracteres (reutilizando el rango de orden ya conocido de longitud 1 para cada mitad), luego por los primeros 4, 8, 16..., duplicando la longitud comparada en cada ronda usando pares de rangos ya ordenados de la ronda anterior en vez de comparar caracteres directamente. El \textbf{LCP array} (\textit{Longest Common Prefix}) complementa al suffix array: $\text{lcp}[i]$ da la longitud del prefijo común más largo entre los sufijos consecutivos $i-1$ e $i$ en el orden del suffix array. El algoritmo de \textbf{Kasai} calcula todo el LCP array en $O(n)$ aprovechando una propiedad clave: si el sufijo que empieza en $j$ tiene LCP $h$ con su vecino en el suffix array, el sufijo que empieza en $j+1$ tiene LCP de al menos $h-1$ con \textit{su} vecino —así que el LCP nunca hay que recalcularlo desde cero, solo puede "bajar como máximo 1" entre iteraciones consecutivas.

\textbf{¿Por qué usarlo?} El suffix array reduce casi cualquier pregunta sobre subcadenas de $s$ a un problema sobre un arreglo ordenado: buscar un patrón se vuelve búsqueda binaria ($O(m\log n)$), y combinado con el LCP array (más una estructura de \textit{sparse table} para consultas de mínimo en rango), permite responder "LCP entre dos sufijos arbitrarios cualesquiera" en $O(1)$ tras $O(n\log n)$ de preprocesamiento —porque el LCP entre dos sufijos no adyacentes en el suffix array es el mínimo del LCP array en el rango entre ellos. Es la herramienta más general de esta parte: resuelve \textit{Subcadena Común Más Larga}, cuenta substrings distintos, y sirve de base para muchos problemas que no tienen un algoritmo dedicado más simple.

\textbf{Casos de uso:}
\begin{itemize}
    \item Búsqueda de patrón mediante búsqueda binaria sobre el suffix array, alternativa a KMP cuando se necesita el suffix array para otras consultas de todas formas.
    \item Contar el número de substrings \textbf{distintos} de una cadena: $\binom{n+1}{2} - \sum_i \text{lcp}[i]$ (el total de substrings menos los que están duplicados, contados por el LCP).
    \item Subcadena común más larga entre dos (o más) cadenas: construir el suffix array de la concatenación con separadores únicos, y buscar el máximo LCP entre sufijos de origen distinto.
    \item Consulta de LCP entre dos sufijos arbitrarios en $O(1)$ (con sparse table sobre el LCP array), base de muchos algoritmos más avanzados de comparación de subcadenas.
\end{itemize}

\textbf{Complejidad:} $O(n\log n)$ para construir el suffix array con ordenamiento por duplicación (hay $O(\log n)$ rondas, cada una $O(n\log n)$ por el \texttt{sort}, aunque con \textit{radix sort} se puede bajar a $O(n\log n)$ total); $O(n)$ para el LCP array con Kasai una vez que el suffix array existe.

\begin{lstlisting}[style=cpp]
// CREATE/READ: Suffix Array (ordenamiento por duplicacion) + LCP Array
// (algoritmo de Kasai)
struct SuffixArray {
    string s;
    int n;
    vector<int> sa;   // sa[i] = indice de inicio del i-esimo sufijo en orden lexicografico
    vector<int> rango; // rango[i] = posicion del sufijo que empieza en i, dentro de sa
    vector<int> lcp;  // lcp[i] = LCP entre sa[i-1] y sa[i] (lcp[0] no definido)

    // CREATE: construye el suffix array mediante ordenamiento por
    // duplicacion (doubling), y luego el LCP array con Kasai
    SuffixArray(const string& cadena) {
        s = cadena + '\0'; // centinela, menor que cualquier caracter real
        n = (int)s.size();
        construirSuffixArray();
        construirLCP();
    } // O(n log n)

    // WRITE (interno): ordenamiento por duplicacion, O(n log n) rondas
    void construirSuffixArray() {
        sa.resize(n);
        rango.resize(n);
        vector<int> rangoTemp(n);
        for (int i = 0; i < n; i++) { sa[i] = i; rango[i] = s[i]; }

        for (int k = 1; k < n; k *= 2) {
            auto comparar = [&](int a, int b) {
                if (rango[a] != rango[b]) return rango[a] < rango[b];
                int ra = (a + k < n) ? rango[a + k] : -1;
                int rb = (b + k < n) ? rango[b + k] : -1;
                return ra < rb;
            };
            sort(sa.begin(), sa.end(), comparar);

            rangoTemp[sa[0]] = 0;
            for (int i = 1; i < n; i++) {
                rangoTemp[sa[i]] = rangoTemp[sa[i-1]] + (comparar(sa[i-1], sa[i]) ? 1 : 0);
            }
            rango = rangoTemp;
            if (rango[sa[n-1]] == n - 1) break; // ya todos los rangos son distintos
        }
    } // O(n log^2 n) con sort estandar (aceptable en CP salvo n muy grande)

    // WRITE (interno): algoritmo de Kasai, calcula el LCP array en O(n)
    // reutilizando que el LCP baja a lo sumo 1 entre posiciones consecutivas
    void construirLCP() {
        lcp.assign(n, 0);
        int h = 0;
        for (int i = 0; i < n; i++) {
            if (rango[i] > 0) {
                int j = sa[rango[i] - 1];
                while (i + h < n && j + h < n && s[i + h] == s[j + h]) h++;
                lcp[rango[i]] = h;
                if (h > 0) h--; // clave: nunca reiniciar desde 0, solo bajar 1
            } else {
                h = 0;
            }
        }
    } // O(n)

    // ------------------------------------------------------------
    // MODIFICACIONES Y COMENTARIOS CON LOS USOS
    // ------------------------------------------------------------
    // 1. Contar substrings distintos: long long total = (long long)n*(n-1)/2;
    //    (usando n sin el centinela) menos sum(lcp[i]) para i=1..n-1 con el
    //    centinela ya excluido del conteo.
    // 2. LCP entre dos sufijos ARBITRARIOS i, j (no necesariamente
    //    adyacentes en sa): construir una sparse table sobre lcp[] para
    //    responder minimo en rango [rango[i]+1, rango[j]] en O(1),
    //    ya que LCP(i,j) = min(lcp[rango[i]+1..rango[j]]) si rango[i] < rango[j].
    // 3. El centinela '\0' garantiza que ningun sufijo sea prefijo de otro
    //    de forma ambigua durante la comparacion; siempre incluirlo al
    //    construir sobre una cadena real.
};
\end{lstlisting}

\textbf{Ejemplo de funcionamiento:}
\begin{lstlisting}[style=cpp]
SuffixArray sa("banana");
// sa.sa contiene los indices de inicio de los sufijos de "banana\0" en
// orden lexicografico; el sufijo vacio (el centinela solo) va primero
// por ser el menor. Los sufijos reales ordenados son:
// "" < "a" < "ana" < "anana" < "banana" < "na" < "nana"
// hay que correrlo para confirmar el arreglo sa.sa exacto con el indice
// del centinela incluido

// sa.lcp[i] da el LCP entre sufijos consecutivos en ese orden; por
// ejemplo el LCP entre "ana" y "anana" es 3 ("ana" completo)
\end{lstlisting}

\textbf{Nota:} la versión de construcción mostrada usa \texttt{sort} con comparador dentro de cada ronda de duplicación, dando $O(n\log^2 n)$ total —suficiente para la mayoría de restricciones de competencia ($n$ hasta $\sim 10^5$-$10^6$). Para $n$ mayor, existe una versión $O(n\log n)$ usando \textit{counting sort} (radix sort) en vez de \texttt{std::sort} dentro de cada ronda, más compleja de implementar y omitida aquí por brevedad.

\textbf{Regla práctica:} usa Suffix Array + Kasai cuando el problema requiera múltiples consultas distintas sobre subcadenas de la misma cadena (o de varias concatenadas con separadores) que no tengan un algoritmo dedicado más simple —es la herramienta más general de esta parte, a costa de mayor complejidad de implementación que KMP, Manacher, o hashing. Si el problema se resuelve directamente con una técnica más específica (búsqueda de un patrón: KMP; palíndromos: Manacher), prefiere esa técnica por ser más simple y menos propensa a bugs.

\subsection{Suffix Automaton}

\textbf{Idea:} un \textit{suffix automaton} es el autómata finito determinista \textbf{más pequeño} que reconoce exactamente el conjunto de todos los \textbf{sufijos} de una cadena $s$ —equivalentemente, reconoce todos los \textbf{substrings} de $s$ (cualquier substring es prefijo de algún sufijo). Se construye incrementalmente, un carácter a la vez (igual que el \textit{Eertree} para palíndromos), manteniendo \textbf{estados} donde cada estado representa una clase de equivalencia de substrings que terminan en exactamente el mismo conjunto de posiciones de $s$ (su \textit{endpos}). Cada estado tiene un \textbf{enlace de sufijo} que apunta al estado que representa el sufijo propio más largo de esa clase con un \textit{endpos} estrictamente mayor —la misma idea general de "aprovechar estructura ya construida vía enlaces" que aparece en $\pi$ (KMP), Aho-Corasick, y el Eertree, aplicada aquí al conjunto completo de substrings distintos.

\textbf{¿Por qué usarlo?} A diferencia del \textit{Suffix Array}, que da una lista ordenada de sufijos y requiere estructuras auxiliares (sparse table sobre LCP) para muchas consultas, el suffix automaton representa el conjunto de substrings distintos como un \textbf{grafo con a lo sumo $2n-1$ estados}, construible en $O(n)$ amortizado (contra $O(n\log n)$ del suffix array), y responde ciertas preguntas —como "¿es $t$ substring de $s$?" en $O(|t|)$, o "cuántos substrings distintos tiene $s$"— de forma más directa, recorriendo transiciones del autómata en vez de hacer búsqueda binaria sobre un arreglo ordenado.

\textbf{Casos de uso:}
\begin{itemize}
    \item Verificar si una cadena $t$ es substring de $s$, en $O(|t|)$, siguiendo transiciones desde el estado inicial.
    \item Contar el número de substrings distintos de $s$: suma, sobre todos los estados (salvo la raíz), de $(\text{longitud}[v] - \text{longitud}[\text{enlace}[v]])$ —cada estado "aporta" ese rango de longitudes como substrings nuevos.
    \item Encontrar el substring común más largo entre dos cadenas, alimentando la segunda cadena carácter por carácter contra el autómata construido sobre la primera.
    \item Contar ocurrencias de cada substring distinto (propagando el tamaño del \textit{endpos} por los enlaces de sufijo, análogo a cómo Aho-Corasick propaga coincidencias).
\end{itemize}

\textbf{Complejidad:} $O(n\log\sigma)$ de construcción con transiciones por \texttt{map} (o $O(n\sigma)$ con arreglo fijo, más rápido para alfabetos pequeños); el número de estados está acotado por $2n-1$ para $n\ge 2$.

\begin{lstlisting}[style=cpp]
// CREATE/READ/WRITE: Suffix Automaton construido incrementalmente,
// reconoce todos los substrings de la cadena procesada hasta el momento
struct SuffixAutomaton {
    static const int ALFABETO = 26;
    struct Estado {
        int longitud;       // longitud del substring mas largo de esta clase
        int enlaceSufijo;   // estado del sufijo propio mas largo con endpos mayor
        int transiciones[ALFABETO];
        long long tamanoEndpos; // tamaño del conjunto endpos (num. de ocurrencias)
        Estado() : longitud(0), enlaceSufijo(-1), tamanoEndpos(0) {
            fill(transiciones, transiciones + ALFABETO, -1);
        }
    };

    vector<Estado> estados;
    int ultimo; // estado que representa la cadena completa construida hasta ahora

    // CREATE: inicializa con el estado raiz (cadena vacia)
    SuffixAutomaton() {
        estados.push_back(Estado());
        estados[0].longitud = 0;
        estados[0].enlaceSufijo = -1;
        ultimo = 0;
    } // O(1)

    // WRITE: extiende el automaton agregando un caracter al final de la
    // cadena representada; crea 1 o 2 estados nuevos segun el caso
    void agregarCaracter(char ch) {
        int c = ch - 'a';
        int nuevo = (int)estados.size();
        estados.push_back(Estado());
        estados[nuevo].longitud = estados[ultimo].longitud + 1;
        estados[nuevo].tamanoEndpos = 1; // esta nueva ocurrencia termina aqui

        int p = ultimo;
        while (p != -1 && estados[p].transiciones[c] == -1) {
            estados[p].transiciones[c] = nuevo;
            p = estados[p].enlaceSufijo;
        }

        if (p == -1) {
            estados[nuevo].enlaceSufijo = 0; // no habia transicion previa: enlaza a la raiz
        } else {
            int q = estados[p].transiciones[c];
            if (estados[p].longitud + 1 == estados[q].longitud) {
                estados[nuevo].enlaceSufijo = q; // caso simple: q ya es la clase correcta
            } else {
                // clonar q: se necesita un estado intermedio para separar
                // la clase de equivalencia correctamente
                int clon = (int)estados.size();
                estados.push_back(estados[q]); // copia transiciones y enlace de q
                estados[clon].longitud = estados[p].longitud + 1;
                estados[clon].tamanoEndpos = 0; // el clon no es una ocurrencia nueva

                while (p != -1 && estados[p].transiciones[c] == q) {
                    estados[p].transiciones[c] = clon;
                    p = estados[p].enlaceSufijo;
                }
                estados[q].enlaceSufijo = clon;
                estados[nuevo].enlaceSufijo = clon;
            }
        }
        ultimo = nuevo;
    } // O(1) amortizado (con transiciones[ALFABETO] fijas)

    // READ: verifica si t es substring de la cadena representada,
    // siguiendo transiciones desde el estado inicial
    bool esSubstring(const string& t) {
        int estadoActual = 0;
        for (char ch : t) {
            int c = ch - 'a';
            if (estados[estadoActual].transiciones[c] == -1) return false;
            estadoActual = estados[estadoActual].transiciones[c];
        }
        return true;
    } // O(|t|)

    // READ: cuenta el numero de substrings DISTINTOS de la cadena
    // representada hasta el momento
    long long contarSubstringsDistintos() {
        long long total = 0;
        for (int v = 1; v < (int)estados.size(); v++) {
            int enlace = estados[v].enlaceSufijo;
            total += estados[v].longitud - estados[enlace].longitud;
        }
        return total;
    } // O(numero de estados) = O(n)

    // ------------------------------------------------------------
    // MODIFICACIONES Y COMENTARIOS CON LOS USOS
    // ------------------------------------------------------------
    // 1. Propagar tamanoEndpos (numero de ocurrencias de cada substring):
    //    ordenar los estados por longitud DESCENDENTE (o con counting
    //    sort por longitud, mas eficiente) y sumar
    //    estados[enlace[v]].tamanoEndpos += estados[v].tamanoEndpos para
    //    cada v en ese orden; los estados clon empiezan en 0 y acumulan
    //    correctamente de sus hijos en el arbol de enlaces de sufijo.
    // 2. Subcadena comun mas larga entre s y t: construir el automaton
    //    sobre s, luego recorrer t manteniendo (estadoActual, longitudActual);
    //    si hay transicion, avanzar y longitudActual++; si no, seguir
    //    enlaces de sufijo reduciendo longitudActual hasta encontrar una
    //    transicion valida o llegar a la raiz. El maximo de longitudActual
    //    visto es la respuesta.
    // 3. Memoria: cada estado reserva ALFABETO enteros para transiciones;
    //    con alfabetos grandes, usar unordered_map<char,int> en su lugar,
    //    a costa de un factor log(sigma) en la complejidad.
};
\end{lstlisting}

\textbf{Ejemplo de funcionamiento:}
\begin{lstlisting}[style=cpp]
SuffixAutomaton sam;
for (char c : string("abcbc")) sam.agregarCaracter(c);

bool existe = sam.esSubstring("cbc");
// existe = true ("cbc" es substring de "abcbc")

long long distintos = sam.contarSubstringsDistintos();
// distintos cuenta todos los substrings unicos de "abcbc" (con
// repeticiones de contenido contadas una sola vez); hay que correrlo
// para confirmar el valor exacto
\end{lstlisting}

\textbf{Nota:} el paso de \textbf{clonación} dentro de \texttt{agregarCaracter} es la parte conceptualmente más difícil del suffix automaton —ocurre quntos la transición encontrada $q$ no puede pertenecer íntegramente a la nueva clase de equivalencia (su longitud no calza exactamente), así que hay que "dividir" esa clase en dos: el clon hereda las transiciones y el enlace de $q$, y ambos ($q$ y el nuevo estado) terminan enlazando al clon. Es el análogo, en este contexto, a la creación de nodos nuevos en el Eertree cuando aparece un palíndromo distinto.

\textbf{Regla práctica:} usa el suffix automaton cuando necesites construir la estructura incrementalmente (streaming, sin poder reprocesar desde cero) o cuando la pregunta sea específicamente sobre \textbf{substrings distintos} y sus conteos —es más directo que Suffix Array + LCP para esas preguntas puntuales. Para consultas que dependen del \textbf{orden lexicográfico} de los sufijos (como encontrar el $k$-ésimo sufijo menor), \textit{Suffix Array} sigue siendo la herramienta más natural, ya que el suffix automaton no mantiene ningún orden lexicográfico explícito entre sus estados.

\section{Búsqueda de Patrones}
\subsection{Knuth-Morris-Pratt (KMP)}

\textbf{Idea:} KMP busca todas las ocurrencias de un patrón $p$ dentro de un texto $t$ reutilizando directamente \textit{Función de Prefijos}: se construye la cadena concatenada $p \,\#\, t$ (con un separador $\#$ que no aparece en ninguna de las dos cadenas), y se calcula su función de prefijos completa. Cada posición $i$ dentro de la parte correspondiente a $t$ donde $\pi[i]$ alcanza exactamente $|p|$ marca el final de una ocurrencia completa del patrón —porque eso significa que los últimos $|p|$ caracteres vistos coinciden exactamente con el patrón entero. El separador garantiza que $\pi$ nunca pueda "cruzar" de $t$ de vuelta a $p$ y contar coincidencias falsas.

\textbf{¿Por qué usarlo?} Buscar un patrón con fuerza bruta (probar cada posición de inicio en $t$ y comparar carácter por carácter) es $O(nm)$ en el peor caso —por ejemplo, texto "aaaa...a" y patrón "aaa...ab". KMP garantiza $O(n+m)$ porque, como se explicó en \textit{Función de Prefijos}, el puntero sobre el texto nunca retrocede: cuando hay un desajuste, se usa la información ya calculada en $\pi$ para saltar directamente al siguiente candidato viable, sin volver a comparar caracteres del texto ya vistos.

\textbf{Casos de uso:}
\begin{itemize}
    \item Encontrar todas las posiciones donde un patrón $p$ ocurre dentro de un texto $t$.
    \item Contar el número total de ocurrencias de $p$ en $t$.
    \item Verificar si $p$ es substring de $t$ (caso particular: basta con la primera ocurrencia).
    \item Como paso de preprocesamiento en problemas que combinan búsqueda de patrón con otra estructura (por ejemplo, contar ocurrencias en un rango de un texto que cambia).
\end{itemize}

\textbf{Complejidad:} $O(n+m)$, donde $n=|t|$ y $m=|p|$ —dominado por el cálculo de la función de prefijos sobre la cadena concatenada de longitud $n+m+1$.

\begin{lstlisting}[style=cpp]
// CREATE/READ: Busqueda de todas las ocurrencias de un patron en un texto
// mediante KMP (funcion de prefijos sobre p + separador + t)
struct KMP {
    // READ: funcion de prefijos, identica a la de Funcion de Prefijos,
    // incluida aqui para que el struct sea autocontenido
    static vector<int> funcionPrefijos(const string& s) {
        int n = (int)s.size();
        vector<int> pi(n, 0);
        for (int i = 1; i < n; i++) {
            int j = pi[i - 1];
            while (j > 0 && s[i] != s[j]) j = pi[j - 1];
            if (s[i] == s[j]) j++;
            pi[i] = j;
        }
        return pi;
    } // O(n)

    // READ: retorna las posiciones (0-indexadas, en t) donde comienza cada
    // ocurrencia del patron p dentro del texto t
    static vector<int> buscar(const string& patron, const string& texto) {
        int m = (int)patron.size();
        if (m == 0) return {};
        string combinada = patron + '#' + texto;
        vector<int> pi = funcionPrefijos(combinada);
        vector<int> ocurrencias;
        for (int i = m + 1; i < (int)combinada.size(); i++) {
            if (pi[i] == m) {
                // i es indice en combinada; la ocurrencia empieza en
                // texto[i - m - m - 1 + 1]... se resta el offset del prefijo+separador
                int inicioEnTexto = i - 2 * m;
                ocurrencias.push_back(inicioEnTexto);
            }
        }
        return ocurrencias;
    } // O(n + m)

    // ------------------------------------------------------------
    // MODIFICACIONES Y COMENTARIOS CON LOS USOS
    // ------------------------------------------------------------
    // 1. Version con menos memoria (sin concatenar cadenas): se puede
    //    correr KMP "en linea" manteniendo solo pi del patron precomputado
    //    y un puntero j que avanza sobre t sin construir la cadena
    //    combinada completa; usa O(m) memoria en vez de O(n+m), preferible
    //    si t es muy grande o llega en streaming.
    // 2. Contar ocurrencias SOLAPADAS: buscar() ya las cuenta todas
    //    (incluyendo solapadas), a diferencia de algunos metodos ingenuos
    //    que saltan m posiciones tras cada match y pierden ocurrencias
    //    solapadas.
    // 3. Multiples patrones contra el mismo texto: si son pocos patrones,
    //    repetir buscar() para cada uno; si son muchos, usar Aho-Corasick
    //    en su lugar, que amortiza el costo sobre todos los patrones a la vez.
};
\end{lstlisting}

\textbf{Ejemplo de funcionamiento:}
\begin{lstlisting}[style=cpp]
auto resultado = KMP::buscar("ab", "ababab");
// resultado = {0, 2, 4}
// (el patron "ab" ocurre en las posiciones 0, 2 y 4 del texto "ababab",
//  todas detectadas correctamente incluyendo las que se tocan entre si)
\end{lstlisting}

\textbf{Nota:} el índice de inicio de cada ocurrencia en \texttt{buscar()} se calcula como \texttt{i - 2*m}, donde $m=|patron|$: la posición $i$ en la cadena combinada incluye el patrón ($m$ caracteres) más el separador (1 carácter), así que el offset hacia el texto original es $i - m - (m+1) + 1 = i - 2m$. Verifica este cálculo cuidadosamente si modificas el struct, ya que es la fuente más común de error de índice en implementaciones de KMP concatenado.

\textbf{Regla práctica:} usa KMP cuando necesites \textbf{todas las ocurrencias exactas} de un único patrón en un texto, con complejidad garantizada $O(n+m)$ sin riesgo de colisión (a diferencia de \textit{Hash Polinómico}, que es probabilístico). Si necesitas buscar \textbf{muchos patrones} simultáneamente contra el mismo texto, usa \textit{Aho-Corasick} en su lugar, que evita repetir el costo de KMP una vez por patrón.

\subsection{Algoritmo Z}

\textbf{Idea:} el \textit{Z-array} de una cadena $s$ de longitud $n$ da, para cada posición $i>0$, la longitud del prefijo común más largo entre $s$ completa y el sufijo de $s$ que empieza en $i$ —es decir, $Z[i]$ mide cuánto se "parece" el sufijo en $i$ al inicio de la cadena. Se calcula manteniendo una ventana $[l,r]$ que representa el segmento coincidente con el prefijo más a la derecha encontrado hasta el momento: para cada nueva posición $i$, si $i$ cae dentro de esa ventana, se aprovecha el valor Z ya conocido de la posición simétrica dentro del prefijo como punto de partida (en vez de comparar desde cero), y solo se extiende comparando carácter por carácter más allá de $r$.

\textbf{¿Por qué usarlo?} Resuelve el mismo problema que KMP (búsqueda de patrón, con la misma complejidad $O(n+m)$) pero con una interpretación distinta: en vez de "prefijo que también es sufijo" (función de prefijos), es directamente "cuánto coincide cada sufijo con el prefijo completo". Muchos lo encuentran más intuitivo de razonar que la función de prefijos, y la técnica de la ventana $[l,r]$ que se extiende sin retroceder es el mismo patrón general de aprovechar trabajo ya hecho que aparece en \textit{Manacher} para palíndromos.

\textbf{Casos de uso:}
\begin{itemize}
    \item Búsqueda de patrón (alternativa a KMP): concatenar $p\,\#\,t$ y buscar posiciones donde $Z[i] = |p|$.
    \item Encontrar todas las posiciones donde un prefijo de la cadena se repite como substring.
    \item Calcular el período de una cadena de forma alternativa a la función de prefijos.
    \item Comprimir cadenas detectando repeticiones del prefijo (base de algunos esquemas de compresión simples).
\end{itemize}

\textbf{Complejidad:} $O(n)$ —análogo al argumento de \textit{Función de Prefijos}: la ventana $[l,r]$ solo se extiende hacia la derecha a lo largo de toda la ejecución, así que el trabajo total de extensión está acotado por $n$.

\begin{lstlisting}[style=cpp]
// CREATE/READ: Z-array de una cadena: para cada i, la longitud del
// prefijo comun mas largo entre s y el sufijo de s que empieza en i
struct AlgoritmoZ {
    // READ: calcula el Z-array completo de s (Z[0] se define como 0 por
    // convencion, ya que comparar el prefijo consigo mismo no aporta info util)
    static vector<int> calcular(const string& s) {
        int n = (int)s.size();
        vector<int> z(n, 0);
        int l = 0, r = 0; // ventana [l,r) del prefijo coincidente mas a la derecha
        for (int i = 1; i < n; i++) {
            if (i < r) {
                // aprovecha el valor ya conocido de la posicion simetrica
                z[i] = min(r - i, z[i - l]);
            }
            // intenta extender mas alla de lo ya garantizado
            while (i + z[i] < n && s[z[i]] == s[i + z[i]]) z[i]++;
            if (i + z[i] > r) { l = i; r = i + z[i]; }
        }
        return z;
    } // O(n)

    // READ: busca todas las ocurrencias del patron p en el texto t,
    // usando el Z-array de la concatenacion p + '#' + t
    static vector<int> buscarPatron(const string& patron, const string& texto) {
        int m = (int)patron.size();
        if (m == 0) return {};
        string combinada = patron + '#' + texto;
        vector<int> z = calcular(combinada);
        vector<int> ocurrencias;
        for (int i = m + 1; i < (int)combinada.size(); i++) {
            if (z[i] == m) ocurrencias.push_back(i - m - 1);
        }
        return ocurrencias;
    } // O(n + m)

    // ------------------------------------------------------------
    // MODIFICACIONES Y COMENTARIOS CON LOS USOS
    // ------------------------------------------------------------
    // 1. Equivalencia con KMP: ambos resuelven busqueda de patron en
    //    O(n+m) con la misma tecnica de no retroceder el puntero; la
    //    eleccion entre uno u otro es principalmente preferencia personal
    //    o cual se recuerda implementar sin errores bajo presion.
    // 2. Periodo de una cadena via Z-array: el periodo mas corto es n - k,
    //    donde k es el menor indice tal que k + z[k] == n (el sufijo desde
    //    k coincide con el prefijo hasta el final de la cadena).
    // 3. Contar cuantas posiciones tienen z[i] >= L para un L dado: util en
    //    problemas de busqueda binaria sobre la longitud de un prefijo
    //    repetido; se puede precomputar con un arreglo de frecuencias y
    //    sumas de sufijo sobre los valores de z.
};
\end{lstlisting}

\textbf{Ejemplo de funcionamiento:}
\begin{lstlisting}[style=cpp]
vector<int> z = AlgoritmoZ::calcular("aabxaayaab");
// z = {0, 1, 0, 0, 3, 1, 0, 2, 1, 0}
// (por ejemplo z[4]=3: el sufijo desde la posicion 4, "aayaab", comparte
//  un prefijo de longitud 3 con la cadena completa "aabxaayaab" -> "aay"
//  vs "aab"... hay que correrlo para confirmar el valor exacto de cada
//  posicion del arreglo)
\end{lstlisting}

\textbf{Regla práctica:} usa el Algoritmo Z como alternativa completamente equivalente a KMP para búsqueda de patrón —la elección entre ambos no afecta la complejidad, así que usa el que te resulte más fácil de implementar sin bugs. La ventaja distintiva del Z-array aparece cuando el problema pregunta directamente sobre "cuánto coincide cada sufijo con el prefijo", una pregunta que con la función de prefijos de KMP requeriría un paso de traducción adicional.
\subsection{Rabin-Karp}

\textbf{Idea:} Rabin-Karp busca un patrón $p$ dentro de un texto $t$ comparando \textbf{hashes} en vez de caracteres: se calcula el hash de $p$ una sola vez, y luego se desliza una ventana de longitud $|p|$ sobre $t$, recalculando el hash de cada ventana en $O(1)$ a partir del hash de la ventana anterior —el mismo mecanismo de \textit{Hash Polinómico} con hashes de prefijo precomputados. Cuando el hash de la ventana coincide con el hash del patrón, se tiene una coincidencia \textbf{candidata}; como el hash es probabilístico, conviene verificar carácter por carácter esa posición para confirmar que no es una colisión falsa.

\textbf{¿Por qué usarlo?} Es conceptualmente el algoritmo de búsqueda de patrón más simple de razonar e implementar de los tres vistos (KMP, Z, Rabin-Karp), porque reutiliza directamente \textit{Hash Polinómico} sin necesidad de ninguna estructura de fallos ni ventana especial. Su desventaja es que, a diferencia de KMP y Z (que son $O(n+m)$ garantizado sin importar el caso), Rabin-Karp puede degradar a $O(nm)$ en el peor caso si hay muchas colisiones de hash falsas que obligan a verificar carácter por carácter repetidamente —algo mitigado casi por completo usando hashing doble.

\textbf{Casos de uso:}
\begin{itemize}
    \item Búsqueda de un patrón en un texto cuando se prioriza simplicidad de implementación sobre garantía de peor caso.
    \item Búsqueda de \textbf{múltiples patrones del mismo tamaño} simultáneamente, comparando el hash de la ventana contra un \texttt{set} de hashes de los patrones (más simple que adaptar KMP o Z a múltiples patrones de igual longitud, aunque Aho-Corasick es la herramienta general para múltiples patrones de cualquier tamaño).
    \item Detección de plagio o duplicados aproximados comparando hashes de ventanas deslizantes.
    \item Como introducción pedagógica a la idea de hashing de subcadenas antes de estructuras más elaboradas.
\end{itemize}

\textbf{Complejidad:} $O(n+m)$ esperado (con hashing doble, la probabilidad de colisión es despreciable en la práctica); $O(nm)$ en el peor caso teórico si un adversario conoce el módulo usado y fuerza colisiones repetidas.

\begin{lstlisting}[style=cpp]
// CREATE/READ: Busqueda de patron mediante Rabin-Karp: hash de ventana
// deslizante comparado contra el hash del patron
struct RabinKarp {
    const long long P = 131;
    const long long MOD = 1000000007LL;

    // READ: calcula P^exp mod MOD (exponenciacion rapida)
    static long long potenciaMod(long long base, long long exp, long long mod) {
        long long resultado = 1;
        base %= mod;
        while (exp > 0) {
            if (exp & 1) resultado = (resultado * base) % mod;
            base = (base * base) % mod;
            exp >>= 1;
        }
        return resultado;
    } // O(log exp)

    // READ: retorna las posiciones (0-indexadas) donde comienza cada
    // ocurrencia EXACTA del patron en el texto (verifica cada candidato
    // caracter por caracter para descartar colisiones de hash)
    vector<int> buscar(const string& patron, const string& texto) {
        int m = (int)patron.size(), n = (int)texto.size();
        vector<int> ocurrencias;
        if (m == 0 || m > n) return ocurrencias;

        long long hashPatron = 0;
        for (char c : patron) hashPatron = (hashPatron * P + (c - 'a' + 1)) % MOD;

        long long hashVentana = 0;
        long long potenciaMayor = potenciaMod(P, m - 1, MOD);
        for (int i = 0; i < n; i++) {
            hashVentana = (hashVentana * P + (texto[i] - 'a' + 1)) % MOD;
            if (i >= m) {
                // remueve el caracter que sale por la izquierda de la ventana
                long long saliente = (texto[i - m] - 'a' + 1) * potenciaMayor % MOD;
                hashVentana = (hashVentana - saliente % MOD + MOD) % MOD;
                hashVentana = (hashVentana - saliente * P % MOD + MOD * MOD) % MOD; // ver Nota
            }
            if (i >= m - 1 && hashVentana == hashPatron) {
                int inicio = i - m + 1;
                if (texto.compare(inicio, m, patron) == 0) { // descarta colision falsa
                    ocurrencias.push_back(inicio);
                }
            }
        }
        return ocurrencias;
    } // O(n + m) esperado

    // ------------------------------------------------------------
    // MODIFICACIONES Y COMENTARIOS CON LOS USOS
    // ------------------------------------------------------------
    // 1. Hashing doble: repetir todo el calculo con un segundo par (P2,MOD2)
    //    en paralelo, y solo considerar candidato si AMBOS hashes coinciden;
    //    reduce la necesidad de verificar con compare() en la mayoria de
    //    los casos, aunque compare() como respaldo sigue siendo la forma
    //    mas segura de garantizar cero falsos positivos.
    // 2. Multiples patrones del mismo tamaño: precomputar un
    //    unordered_set<long long> con los hashes de todos los patrones, y
    //    comparar el hash de cada ventana contra ese set en O(1) en vez de
    //    contra un solo hashPatron.
    // 3. Ventana deslizante SIN necesidad de patron previo (hash rodante
    //    generico): la logica de actualizar hashVentana quitando el
    //    caracter saliente y agregando el entrante es reutilizable para
    //    cualquier problema que necesite el hash de una ventana de tamaño
    //    fijo deslizandose sobre una cadena.
};
\end{lstlisting}

\textbf{Ejemplo de funcionamiento:}
\begin{lstlisting}[style=cpp]
RabinKarp rk;
auto resultado = rk.buscar("ab", "ababab");
// resultado = {0, 2, 4}
// (mismo resultado que KMP y Algoritmo Z para el mismo patron y texto;
//  cada candidato con hash coincidente se verifico con compare() antes
//  de agregarse, descartando cualquier colision falsa)
\end{lstlisting}

\textbf{Nota:} la línea marcada arriba (remoción del carácter saliente multiplicado por \texttt{P} para compensar el corrimiento de todos los términos restantes) es la parte más propensa a errores de índice en Rabin-Karp con ventana deslizante —a diferencia de \textit{Hash Polinómico} con prefijos precomputados (donde la resta es directa), aquí se recalcula incrementalmente y conviene verificarla con cuidado o preferir la versión de hash de prefijos si hay dudas.

\textbf{Regla práctica:} usa Rabin-Karp cuando priorices simplicidad de implementación y el problema no sea adversarial contra hashing (o uses hashing doble para mitigarlo); si necesitas la garantía \textbf{determinística} de $O(n+m)$ sin ninguna dependencia de probabilidad, usa \textit{KMP} o \textit{Algoritmo Z} en su lugar —ambos resuelven exactamente el mismo problema sin el riesgo teórico de colisión.

\subsection{Boyer-Moore}

\textbf{Idea:} a diferencia de KMP, Z y Rabin-Karp —que procesan el texto de izquierda a derecha—, Boyer-Moore compara el patrón contra el texto \textbf{de derecha a izquierda} en cada posición candidata, y cuando encuentra un desajuste, usa dos heurísticas para decidir cuánto \textbf{saltar} el patrón hacia adelante sin perder ninguna coincidencia posible: la \textit{regla del mal carácter} (mover el patrón para alinear el carácter del texto que causó el desajuste con su ocurrencia más a la derecha dentro del patrón, si existe) y la \textit{regla del buen sufijo} (aprovechar que el sufijo ya coincidente del patrón puede reaparecer en otra posición del propio patrón). El desplazamiento efectivo usado es el máximo de lo que sugieren ambas reglas.

\textbf{¿Por qué usarlo?} En el caso promedio, especialmente con alfabetos grandes (como texto en inglés, 26+ letras), Boyer-Moore es \textbf{sublineal}: no necesita examinar todos los caracteres del texto, porque los saltos grandes de la regla del mal carácter permiten "saltarse" bloques enteros sin compararlos. Esto lo hace el algoritmo más rápido en la práctica para búsqueda de texto en editores y herramientas como \texttt{grep}, aunque su peor caso teórico ($O(nm)$ con la regla del mal carácter sola, mejorable a $O(n+m)$ combinando ambas reglas correctamente) no es tan simple de garantizar como KMP.

\textbf{Casos de uso:}
\begin{itemize}
    \item Búsqueda de texto en editores, herramientas de línea de comandos, y motores de búsqueda de texto donde el rendimiento práctico importa más que la garantía teórica de peor caso.
    \item Alfabetos grandes (texto natural) donde la regla del mal carácter produce saltos grandes frecuentemente.
    \item Como comparación de referencia al estudiar por qué distintos algoritmos de matching tienen ventajas en distintos escenarios (útil para justificar la elección en un informe o discusión).
    \item En CP específicamente, es menos común que KMP/Z/Hash porque su implementación completa (con ambas reglas) es más larga; se usa más la versión simplificada solo con la regla del mal carácter cuando el alfabeto es grande.
\end{itemize}

\textbf{Complejidad:} $O(n/m)$ en el mejor caso (saltos grandes constantes), $O(n+m)$ en el caso promedio con ambas heurísticas combinadas correctamente, y $O(nm)$ en el peor caso si solo se usa la regla del mal carácter sin la regla del buen sufijo como respaldo.

\begin{lstlisting}[style=cpp]
// CREATE/READ: Busqueda de patron mediante Boyer-Moore con regla del mal
// caracter (version simplificada, alfabeto minusculas 'a'-'z')
struct BoyerMoore {
    static const int ALFABETO = 26;

    // READ: para cada caracter del alfabeto, la posicion MAS A LA DERECHA
    // en la que aparece dentro del patron (-1 si no aparece)
    static vector<int> tablaUltimaOcurrencia(const string& patron) {
        vector<int> ultima(ALFABETO, -1);
        for (int i = 0; i < (int)patron.size(); i++) {
            ultima[patron[i] - 'a'] = i;
        }
        return ultima;
    } // O(m)

    // READ: retorna las posiciones (0-indexadas) donde comienza cada
    // ocurrencia del patron en el texto, usando solo la regla del mal
    // caracter (simplificacion valida y suficiente para la mayoria de
    // problemas de CP)
    static vector<int> buscar(const string& patron, const string& texto) {
        int m = (int)patron.size(), n = (int)texto.size();
        vector<int> ocurrencias;
        if (m == 0 || m > n) return ocurrencias;

        vector<int> ultima = tablaUltimaOcurrencia(patron);
        int s = 0; // posicion de alineacion actual del patron sobre el texto
        while (s <= n - m) {
            int j = m - 1;
            // compara de derecha a izquierda
            while (j >= 0 && patron[j] == texto[s + j]) j--;
            if (j < 0) {
                ocurrencias.push_back(s);
                s++; // desplazamiento minimo tras un match completo
            } else {
                int caracterMalo = texto[s + j] - 'a';
                int salto = j - ultima[caracterMalo];
                s += max(1, salto); // nunca retroceder, salto minimo de 1
            }
        }
        return ocurrencias;
    } // O(n/m) mejor caso, O(nm) peor caso (con solo regla del mal caracter)

    // ------------------------------------------------------------
    // MODIFICACIONES Y COMENTARIOS CON LOS USOS
    // ------------------------------------------------------------
    // 1. Regla del buen sufijo (NO implementada arriba): garantiza el peor
    //    caso O(n+m) combinada con la regla del mal caracter; requiere una
    //    tabla adicional de desplazamientos basada en donde reaparece cada
    //    sufijo del patron dentro de si mismo (similar en espiritu a la
    //    funcion de prefijos, pero mas compleja de construir); se omite
    //    aqui porque en la practica competitiva, KMP ya da la garantia
    //    O(n+m) de forma mas simple si el peor caso es una preocupacion real.
    // 2. Alfabeto grande o Unicode: la tabla ultima como arreglo fijo de
    //    tamaño ALFABETO no escala; usar unordered_map<char,int> en su
    //    lugar.
    // 3. Cuando NO usar Boyer-Moore en CP: si el patron o el texto son
    //    cortos, o el alfabeto es pequeño (binario, ADN de 4 letras), los
    //    saltos de la regla del mal caracter son pequeños y la ventaja
    //    practica desaparece -- en esos casos KMP es igual de rapido y mas
    //    simple de implementar sin bugs.
};
\end{lstlisting}

\textbf{Ejemplo de funcionamiento:}
\begin{lstlisting}[style=cpp]
auto resultado = BoyerMoore::buscar("ab", "ababab");
// resultado = {0, 2, 4}
// (mismo resultado que los algoritmos anteriores; con un alfabeto de
//  solo 2 caracteres distintos la ventaja de salto grande de Boyer-Moore
//  es minima, pero el resultado es igualmente correcto)
\end{lstlisting}

\textbf{Nota:} la implementación de arriba usa \textbf{únicamente} la regla del mal carácter, que es la versión más simple y la que se ve con más frecuencia en código de competencia —pero sin la regla del buen sufijo, el peor caso teórico es $O(nm)$, igual de malo que fuerza bruta. En CP, si existe cualquier riesgo de que el caso de prueba sea adversarial, es más seguro usar \textit{KMP} o el \textit{Algoritmo Z}, que garantizan $O(n+m)$ sin depender de heurísticas.

\textbf{Regla práctica:} en competencia, prefiere \textit{KMP} o \textit{Algoritmo Z} sobre Boyer-Moore por su garantía de peor caso simple de razonar; considera Boyer-Moore específicamente cuando el alfabeto es grande y estás implementando algo más cercano a una herramienta de búsqueda de texto real (no bajo las restricciones adversariales típicas de un problema de concurso), donde su velocidad práctica sublineal sí es una ventaja genuina.
\subsection{Autómata Finito de Coincidencia de Patrones}

\textbf{Idea:} un autómata finito de coincidencia precomputa, para cada estado (representando "cuánto del patrón llevo coincidido hasta ahora", igual que en \textit{Función de Prefijos}) y cada carácter posible del alfabeto, exactamente a qué estado se transiciona —a diferencia de KMP, donde ese "a qué estado ir" se calcula on-the-fly siguiendo la cadena de fallos $\pi$ durante la búsqueda. La tabla de transiciones $\delta[\text{estado}][\text{carácter}]$ se construye una sola vez a partir de la función de prefijos del patrón, y una vez construida, procesar el texto es un simple recorrido de tabla: un acceso $O(1)$ por carácter, sin ningún retroceso ni bucle interno.

\textbf{¿Por qué usarlo?} KMP y el autómata resuelven el mismo problema con la misma complejidad total, pero difieren en \textbf{dónde} pagan el costo: KMP hace el trabajo de "retroceder por fallos" durante la búsqueda (amortizado $O(n)$, pero con un bucle interno variable), mientras que el autómata paga todo ese costo por adelantado en la construcción de la tabla, y la búsqueda queda reducida a $O(n)$ estrictos sin ningún bucle interno ni siquiera amortizado —una transición por carácter, garantizado. Esto importa cuando el mismo patrón se va a buscar en \textbf{muchos textos distintos}: la tabla se construye una vez y se reutiliza, evitando recalcular la cadena de fallos en cada búsqueda.

\textbf{Casos de uso:}
\begin{itemize}
    \item Buscar el mismo patrón repetidamente contra muchos textos distintos, amortizando el costo de construcción.
    \item Streaming: procesar un texto que llega carácter por carácter (sin poder verlo completo de antemano) con tiempo de respuesta $O(1)$ por carácter.
    \item Como modelo conceptual para entender \textit{Aho-Corasick}, que es exactamente esta misma idea generalizada a múltiples patrones simultáneos mediante un trie.
    \item Problemas donde se necesita garantía de tiempo constante por carácter, sin ninguna variabilidad de bucles internos.
\end{itemize}

\textbf{Complejidad:} $O(m\sigma)$ de construcción, donde $m=|\text{patrón}|$ y $\sigma$ es el tamaño del alfabeto (se llena una tabla de $m+1$ estados por $\sigma$ caracteres); $O(n)$ estricto para procesar un texto de longitud $n$, un acceso a tabla por carácter.

\begin{lstlisting}[style=cpp]
// CREATE/READ: Automata finito de coincidencia de patrones, construido a
// partir de la funcion de prefijos del patron (alfabeto 'a'-'z')
struct AutomataPatron {
    static const int ALFABETO = 26;
    int m;
    vector<vector<int>> delta; // delta[estado][caracter] = siguiente estado

    // CREATE: construye la tabla de transiciones completa a partir del
    // patron, reutilizando la idea de Funcion de Prefijos para los
    // estados que no coinciden directamente con el siguiente caracter
    AutomataPatron(const string& patron) {
        m = (int)patron.size();
        delta.assign(m + 1, vector<int>(ALFABETO, 0));

        vector<int> pi(m, 0);
        for (int i = 1; i < m; i++) {
            int j = pi[i - 1];
            while (j > 0 && patron[i] != patron[j]) j = pi[j - 1];
            if (patron[i] == patron[j]) j++;
            pi[i] = j;
        }

        // estado 0: solo el caracter que calza con patron[0] avanza a 1
        delta[0][patron[0] - 'a'] = 1;
        for (int estado = 1; estado <= m; estado++) {
            for (int c = 0; c < ALFABETO; c++) {
                char caracter = 'a' + c;
                if (estado < m && caracter == patron[estado]) {
                    // el caracter esperado avanza al siguiente estado
                    delta[estado][c] = estado + 1;
                } else {
                    // reutiliza la transicion ya calculada del estado de
                    // fallo (pi[estado-1]), evitando recalcular desde cero
                    delta[estado][c] = delta[pi[estado - 1]][c];
                }
            }
        }
    } // O(m * sigma)

    // READ: retorna las posiciones (0-indexadas) donde termina cada
    // ocurrencia completa del patron al procesar el texto
    vector<int> buscar(const string& texto) {
        vector<int> ocurrencias;
        int estado = 0;
        for (int i = 0; i < (int)texto.size(); i++) {
            estado = delta[estado][texto[i] - 'a'];
            if (estado == m) ocurrencias.push_back(i - m + 1);
        }
        return ocurrencias;
    } // O(n), estricto, sin bucle interno

    // ------------------------------------------------------------
    // MODIFICACIONES Y COMENTARIOS CON LOS USOS
    // ------------------------------------------------------------
    // 1. delta[pi[estado-1]][c] en la construccion: este es el mismo tipo
    //    de "reutilizar trabajo ya calculado" que el retroceso por fallos
    //    en KMP, solo que aqui se hace UNA vez durante la construccion en
    //    vez de repetidamente durante cada busqueda.
    // 2. Memoria: la tabla ocupa O(m * sigma), que puede ser prohibitivo si
    //    el alfabeto es grande (Unicode) y m tambien lo es; en ese caso
    //    KMP puro (sin tabla) es preferible por memoria O(m) en vez de
    //    O(m*sigma).
    // 3. Multiples busquedas del MISMO patron: construir el automata una
    //    sola vez (fuera de cualquier bucle sobre los textos) y llamar
    //    buscar() repetidamente reutilizando la misma tabla delta.
};
\end{lstlisting}

\textbf{Ejemplo de funcionamiento:}
\begin{lstlisting}[style=cpp]
AutomataPatron automata("ab");
auto resultado = automata.buscar("ababab");
// resultado = {1, 3, 5}
// (indices donde TERMINA cada ocurrencia de "ab"; equivalen a las
//  posiciones de inicio {0,2,4} encontradas por KMP, Z y Rabin-Karp,
//  desplazadas por m-1 ya que aqui se reporta el final de la coincidencia)
\end{lstlisting}

\textbf{Nota:} a diferencia de KMP, Z y Rabin-Karp (que reportan la posición de \textbf{inicio} de cada ocurrencia), esta implementación reporta la posición donde \textbf{termina} cada coincidencia, porque el estado $m$ se alcanza al procesar el último carácter del patrón. Ajustar \texttt{i - m + 1} en vez de \texttt{i} en el \texttt{push\_back} convierte el resultado a posición de inicio si se prefiere esa convención, como se muestra en el código.

\textbf{Regla práctica:} usa el autómata en vez de KMP directo cuando el mismo patrón se busque repetidamente contra múltiples textos, o cuando necesites la garantía de $O(1)$ estricto por carácter sin ninguna variabilidad de bucle interno (por ejemplo, procesamiento en streaming). Para una única búsqueda de un patrón contra un único texto, KMP es preferible —evita pagar el costo $O(m\sigma)$ de construir la tabla completa cuando esa tabla solo se va a usar una vez.

\subsection{Aho-Corasick (Múltiples Patrones)}

\textbf{Idea:} Aho-Corasick generaliza el \textit{Autómata Finito de Coincidencia de Patrones} de un único patrón a \textbf{muchos patrones simultáneos}: primero se insertan todos los patrones en un \textit{Trie}, y luego se calcula, para cada nodo del trie, un \textbf{enlace de fallo} (\textit{failure link}) que apunta al nodo que representa el sufijo propio más largo del prefijo actual que también es un prefijo de algún patrón —el mismo concepto que $\pi$ en la función de prefijos, generalizado de una cadena lineal a un árbol. Al procesar el texto carácter por carácter, cuando no existe una transición directa en el trie, se sigue el enlace de fallo (análogo a retroceder con $\pi$ en KMP), y cada nodo visitado que marca el fin de uno o más patrones reporta esas coincidencias.

\textbf{¿Por qué usarlo?} Buscar $k$ patrones distintos corriendo KMP (o el autómata) una vez por patrón cuesta $O(k \cdot n)$ en total. Aho-Corasick procesa el texto \textbf{una sola vez}, con todos los patrones cargados simultáneamente en la estructura, dando $O(n + \sum|p_i| + \text{ocurrencias})$ total —independiente de $k$ en la parte de recorrido del texto. Es la generalización natural de KMP a múltiples patrones, de la misma forma que el trie generaliza la búsqueda de una sola cadena a un conjunto de cadenas.

\textbf{Casos de uso:}
\begin{itemize}
    \item Buscar todas las ocurrencias de un conjunto grande de palabras (diccionario) dentro de un texto, en una sola pasada.
    \item Filtros de contenido (detectar palabras prohibidas de una lista extensa) en tiempo real sobre texto entrante.
    \item Bioinformática: buscar múltiples secuencias/motivos genéticos simultáneamente dentro de una cadena de ADN.
    \item Como paso de preprocesamiento en problemas de DP sobre autómata (contar cadenas de longitud $n$ que evitan cierto conjunto de patrones prohibidos).
\end{itemize}

\textbf{Complejidad:} $O(\sum|p_i| \cdot \sigma)$ de construcción (insertar todos los patrones en el trie más calcular los enlaces de fallo con BFS), donde $\sigma$ es el tamaño del alfabeto; $O(n + \text{ocurrencias})$ para procesar un texto de longitud $n$ y reportar todas las coincidencias encontradas.

\begin{lstlisting}[style=cpp]
// CREATE/READ: Aho-Corasick para buscar multiples patrones simultaneamente
// en un texto, mediante trie + enlaces de fallo (alfabeto 'a'-'z')
struct AhoCorasick {
    static const int ALFABETO = 26;
    struct Nodo {
        int hijos[ALFABETO];
        int fallo;                  // enlace de fallo (failure link)
        vector<int> indicesPatron;  // indices de patrones que terminan aqui
        Nodo() : fallo(0) {
            fill(hijos, hijos + ALFABETO, -1);
        }
    };

    vector<Nodo> nodos;

    // CREATE: inicializa con la raiz del trie
    AhoCorasick() {
        nodos.push_back(Nodo());
    } // O(1)

    // WRITE: inserta un patron en el trie, registrando su indice en el
    // nodo donde termina (para saber cual patron coincidio al buscar)
    void insertar(const string& patron, int indice) {
        int actual = 0;
        for (char c : patron) {
            int idx = c - 'a';
            if (nodos[actual].hijos[idx] == -1) {
                nodos[actual].hijos[idx] = (int)nodos.size();
                nodos.push_back(Nodo());
            }
            actual = nodos[actual].hijos[idx];
        }
        nodos[actual].indicesPatron.push_back(indice);
    } // O(|patron|)

    // WRITE: calcula los enlaces de fallo de todos los nodos mediante BFS,
    // y completa las transiciones faltantes para que delta quede total
    // (cada nodo tiene transicion valida para TODO caracter del alfabeto)
    void construirFallos() {
        queue<int> cola;
        for (int c = 0; c < ALFABETO; c++) {
            if (nodos[0].hijos[c] == -1) {
                nodos[0].hijos[c] = 0; // sin match, se queda en la raiz
            } else {
                nodos[nodos[0].hijos[c]].fallo = 0;
                cola.push(nodos[0].hijos[c]);
            }
        }
        while (!cola.empty()) {
            int actual = cola.front(); cola.pop();
            // hereda las ocurrencias del nodo de fallo (un patron corto
            // que es sufijo de uno mas largo tambien cuenta como match)
            for (int idx : nodos[nodos[actual].fallo].indicesPatron) {
                nodos[actual].indicesPatron.push_back(idx);
            }
            for (int c = 0; c < ALFABETO; c++) {
                int hijo = nodos[actual].hijos[c];
                if (hijo == -1) {
                    // sin transicion directa: usa la del nodo de fallo
                    nodos[actual].hijos[c] = nodos[nodos[actual].fallo].hijos[c];
                } else {
                    nodos[hijo].fallo = nodos[nodos[actual].fallo].hijos[c];
                    cola.push(hijo);
                }
            }
        }
    } // O(sum|p_i| * sigma)

    // READ: procesa el texto y retorna pares (posicion final, indice de
    // patron) por cada ocurrencia encontrada de cualquier patron insertado
    vector<pair<int,int>> buscar(const string& texto) {
        vector<pair<int,int>> ocurrencias;
        int estado = 0;
        for (int i = 0; i < (int)texto.size(); i++) {
            estado = nodos[estado].hijos[texto[i] - 'a'];
            for (int idx : nodos[estado].indicesPatron) {
                ocurrencias.push_back({i, idx});
            }
        }
        return ocurrencias;
    } // O(n + ocurrencias)

    // ------------------------------------------------------------
    // MODIFICACIONES Y COMENTARIOS CON LOS USOS
    // ------------------------------------------------------------
    // 1. Solo verificar EXISTENCIA de alguna coincidencia (sin listar
    //    todas): basta con revisar si nodos[estado].indicesPatron esta
    //    vacio en vez de acumular en el vector ocurrencias, ahorrando
    //    memoria si hay muchisimas coincidencias.
    // 2. DP sobre el automata de Aho-Corasick (contar cadenas que evitan
    //    patrones prohibidos): tratar cada nodo como un estado de DP y
    //    transicionar segun nodos[estado].hijos[c], excluyendo estados
    //    donde indicesPatron no esta vacio (representan "contiene un
    //    patron prohibido").
    // 3. IMPORTANTE llamar construirFallos() UNA VEZ, despues de insertar
    //    TODOS los patrones y antes de cualquier llamada a buscar(); insertar
    //    un patron adicional despues de construirFallos() deja la estructura
    //    en un estado inconsistente.
};
\end{lstlisting}

\textbf{Ejemplo de funcionamiento:}
\begin{lstlisting}[style=cpp]
AhoCorasick ac;
ac.insertar("he", 0);
ac.insertar("she", 1);
ac.insertar("his", 2);
ac.insertar("hers", 3);
ac.construirFallos();

auto resultado = ac.buscar("ushers");
// resultado incluye (2, 1) "she" termina en indice 2, (3, 0) "he" termina
// en indice 3 (como sufijo de "she", heredado por el enlace de fallo), y
// (5, 3) "hers" termina en indice 5; hay que correrlo para confirmar el
// conjunto completo y orden exacto de las coincidencias reportadas
\end{lstlisting}

\textbf{Nota:} el paso de "heredar \texttt{indicesPatron} del nodo de fallo" dentro de \texttt{construirFallos()} es crítico y fácil de omitir: sin él, un patrón corto que aparece como sufijo de una coincidencia más larga (como "he" dentro de "she") no se reportaría, porque el autómata solo estaría en el nodo correspondiente a "she" sin saber que ese mismo nodo también representa el final de "he" por la cadena de fallos.

\textbf{Regla práctica:} usa Aho-Corasick en cuanto el problema involucre \textbf{más de un patrón} buscado simultáneamente en el mismo texto —es la generalización directa y la estructura estándar para ese escenario, mucho más eficiente que repetir KMP una vez por patrón. Para un único patrón, el costo de construir el trie y los enlaces de fallo no se justifica frente a KMP o el Algoritmo Z, que son más simples de implementar para ese caso particular.

\section{Coincidencia de Patrones Especiales}
\subsection{Patrón con Comodines (Wildcard Matching)}

\textbf{Idea:} cuando el patrón puede contener un carácter comodín \texttt{?} (que coincide con cualquier carácter único), buscar todas las posiciones donde el patrón calza contra el texto ya no admite comparación directa carácter por carácter con los algoritmos exactos anteriores (KMP, Z, Rabin-Karp) sin modificación, porque el comodín rompe la igualdad estricta que esos algoritmos asumen. La técnica estándar en competencia reduce el problema a \textbf{convolución}: se codifica cada carácter como un número (o se usa una función que se anula exactamente cuando dos caracteres son iguales, como $(t_i - p_j)$), se trata el comodín como si calzara automáticamente (aportando $0$ a la diferencia), y se usa la \textit{Transformada Rápida de Fourier} (FFT) para calcular, en $O(n\log n)$, para cada posición de alineación, la suma de $(t_{i+j}-p_j)^2 \cdot [p_j \neq \texttt{?}]$ sobre todo $j$ —esa suma da exactamente cero si y solo si el patrón calza en esa posición, comodines incluidos.

\textbf{¿Por qué usarlo?} La fuerza bruta de wildcard matching es $O(nm)$: para cada una de las $n-m+1$ posiciones de alineación, comparar los $m$ caracteres del patrón contra el texto, tratando \texttt{?} como comodín. Cuando $n$ y $m$ son grandes (por ejemplo $10^5$ cada uno), eso es inviable. La reducción a convolución vía FFT baja esto a $O(n\log n)$, aprovechando que "calzar carácter por carácter salvo comodines" se puede expresar como un polinomio cuyo producto (convolución) revela automáticamente, en cada posición, si hubo coincidencia total.

\textbf{Casos de uso:}
\begin{itemize}
    \item Buscar un patrón con comodines \texttt{?} dentro de un texto grande.
    \item Generalización: patrón con comodines de \textbf{ambos lados} (tanto en el patrón como en el texto), resoluble con el mismo esquema de convolución adaptado.
    \item Como aplicación concreta de FFT/convolución a un problema de texto, distinta de sus usos habituales en problemas puramente numéricos.
    \item Verificación de igualdad "difusa" entre dos cadenas de igual longitud que permiten cierto carácter especial como comodín.
\end{itemize}

\textbf{Complejidad:} $O(n\log n)$ mediante FFT, donde $n$ es la longitud del texto —dominado por las llamadas a FFT/convolución, cada una $O(n\log n)$; comparado con $O(nm)$ de la fuerza bruta directa.

\begin{lstlisting}[style=cpp]
// CREATE/READ: Wildcard matching mediante reduccion a suma de diferencias
// cuadraticas; VERSION CONCEPTUAL con complejidad O(n*m), la version FFT
// O(n log n) requiere una implementacion de FFT/convolucion aparte (no
// incluida aqui, ver Nota)
struct WildcardMatching {
    static const char COMODIN = '?';

    // READ: verifica si el patron (con comodines) calza contra el texto
    // empezando en la posicion 'inicio', tratando '?' como comodin en
    // CUALQUIERA de las dos cadenas
    static bool calzaEnPosicion(const string& texto, const string& patron, int inicio) {
        for (int j = 0; j < (int)patron.size(); j++) {
            char tc = texto[inicio + j];
            char pc = patron[j];
            if (tc != COMODIN && pc != COMODIN && tc != pc) return false;
        }
        return true;
    } // O(m)

    // READ: fuerza bruta O(n*m); retorna todas las posiciones de inicio
    // donde el patron calza contra el texto respetando comodines
    static vector<int> buscarFuerzaBruta(const string& texto, const string& patron) {
        int n = (int)texto.size(), m = (int)patron.size();
        vector<int> ocurrencias;
        for (int i = 0; i + m <= n; i++) {
            if (calzaEnPosicion(texto, patron, i)) ocurrencias.push_back(i);
        }
        return ocurrencias;
    } // O(n * m)

    // ------------------------------------------------------------
    // MODIFICACIONES Y COMENTARIOS CON LOS USOS
    // ------------------------------------------------------------
    // 1. Version O(n log n) con FFT: para cada posicion de alineacion i,
    //    se necesita sum_j (texto[i+j] - patron[j])^2 * [patron[j] != '?'];
    //    expandiendo el cuadrado se obtiene una suma de TRES convoluciones
    //    (terminos con texto^2*patron^0, texto^1*patron^1, texto^0*patron^2),
    //    cada una calculable con una llamada a FFT. Requiere una
    //    implementacion de FFT como bloque de construccion aparte
    //    (no cubierta en esta subseccion).
    // 2. Comodines en AMBAS cadenas (no solo en el patron): la condicion
    //    en calzaEnPosicion ya maneja esto (revisa '?' en tc y en pc por
    //    separado); en la version FFT, el termino [patron[j] != '?'] se
    //    ajusta a [texto[i+j] != '?'] * [patron[j] != '?'].
    // 3. Alfabeto numerico en vez de caracteres: si el problema usa
    //    numeros en vez de letras, el mismo esquema aplica directamente
    //    sustituyendo la codificacion de caracteres por los valores
    //    numericos dados.
};
\end{lstlisting}

\textbf{Ejemplo de funcionamiento:}
\begin{lstlisting}[style=cpp]
auto resultado = WildcardMatching::buscarFuerzaBruta("abcabd", "a?c");
// resultado = {0}
// ("a?c" calza en la posicion 0 de "abcabd": 'a'='a', '?' calza con 'b',
//  'c'='c'; no hay otra posicion donde los tres caracteres, respetando
//  el comodin, coincidan)
\end{lstlisting}

\textbf{Nota:} el struct de arriba implementa \textbf{solo la versión $O(nm)$} por claridad conceptual —la versión $O(n\log n)$ con FFT requiere convolución polinómica como bloque de construcción separado (no cubierto en esta guía). Si $n$ y $m$ son pequeños (típicamente hasta $\sim 10^3$-$10^4$ dependiendo del límite de tiempo), la versión de fuerza bruta ya es suficiente y evita la complejidad de implementar FFT bajo presión de concurso.

\textbf{Regla práctica:} usa la fuerza bruta $O(nm)$ cuando $n$ y $m$ sean lo suficientemente pequeños para las restricciones del problema (revisa el límite de tiempo y $n\cdot m$); si $n$ y $m$ son grandes (del orden de $10^5$), es necesaria la reducción a convolución con FFT, y conviene tener una plantilla de FFT ya probada de antemano en el notebook antes del concurso, ya que implementarla desde cero bajo presión es propenso a errores.

\subsection{Búsqueda de Anagramas de un Patrón}

\textbf{Idea:} encontrar todas las posiciones donde una \textbf{permutación} (anagrama) de un patrón $p$ aparece como subcadena de un texto $t$ se resuelve con una \textbf{ventana deslizante} de tamaño $|p|$ sobre $t$, manteniendo un conteo de frecuencia de caracteres dentro de la ventana actual. Una ventana es un anagrama de $p$ si y solo si su distribución de frecuencias de caracteres es \textbf{idéntica} a la de $p$. En vez de comparar los dos arreglos de frecuencia completos ($O(\sigma)$) en cada posición, se mantiene un contador de "cuántas posiciones del alfabeto difieren" que se actualiza en $O(1)$ al agregar el carácter que entra y quitar el que sale de la ventana.

\textbf{¿Por qué usarlo?} Comparar fuerza bruta la distribución de frecuencia completa en cada una de las $n-m+1$ posiciones cuesta $O(n\sigma)$. Mantener incrementalmente un contador de "diferencias activas" entre la ventana y el patrón, actualizado en $O(1)$ por posición (no $O(\sigma)$), baja el total a $O(n+\sigma)$ —el término $\sigma$ solo aparece una vez, al inicializar el conteo de frecuencia del patrón. A diferencia de \textit{Rabin-Karp} u otros métodos de matching exacto, aquí el orden de los caracteres dentro de la ventana es irrelevante por definición del problema, así que hash polinómico (que sí depende del orden) no aplica directamente.

\textbf{Casos de uso:}
\begin{itemize}
    \item Encontrar todas las posiciones donde algún anagrama de un patrón aparece en un texto.
    \item Verificar si dos cadenas son anagramas entre sí (caso particular: una sola "ventana", el texto completo).
    \item Problemas de "substrings balanceados" donde la condición depende de la composición de caracteres, no de su orden.
    \item Extensión a permitir hasta $k$ caracteres de diferencia (anagrama aproximado), ajustando el umbral del contador de diferencias.
\end{itemize}

\textbf{Complejidad:} $O(n+\sigma)$, donde $n=|t|$ y $\sigma$ es el tamaño del alfabeto —el conteo de frecuencia del patrón es $O(m+\sigma)$, y cada desplazamiento de la ventana sobre el texto es $O(1)$ amortizado.

\begin{lstlisting}[style=cpp]
// CREATE/READ: Busqueda de todas las posiciones donde un anagrama del
// patron aparece en el texto, mediante ventana deslizante de frecuencias
struct BusquedaAnagramas {
    static const int ALFABETO = 26;

    // READ: retorna las posiciones de inicio (0-indexadas) de cada
    // ventana de tamaño |patron| en texto cuya distribucion de caracteres
    // coincide exactamente con la del patron
    static vector<int> buscar(const string& texto, const string& patron) {
        int n = (int)texto.size(), m = (int)patron.size();
        vector<int> ocurrencias;
        if (m == 0 || m > n) return ocurrencias;

        vector<int> frecPatron(ALFABETO, 0), frecVentana(ALFABETO, 0);
        for (char c : patron) frecPatron[c - 'a']++;

        int diferencias = 0; // cuantas posiciones del alfabeto difieren actualmente
        for (int c = 0; c < ALFABETO; c++) {
            if (frecVentana[c] != frecPatron[c]) diferencias++;
        }

        for (int i = 0; i < n; i++) {
            // agrega el caracter entrante a la ventana
            int idxEntra = texto[i] - 'a';
            if (frecVentana[idxEntra] == frecPatron[idxEntra]) diferencias++;
            frecVentana[idxEntra]++;
            if (frecVentana[idxEntra] == frecPatron[idxEntra]) diferencias--;

            // remueve el caracter saliente si la ventana ya excede tamaño m
            if (i >= m) {
                int idxSale = texto[i - m] - 'a';
                if (frecVentana[idxSale] == frecPatron[idxSale]) diferencias++;
                frecVentana[idxSale]--;
                if (frecVentana[idxSale] == frecPatron[idxSale]) diferencias--;
            }

            if (i >= m - 1 && diferencias == 0) {
                ocurrencias.push_back(i - m + 1);
            }
        }
        return ocurrencias;
    } // O(n + sigma)

    // ------------------------------------------------------------
    // MODIFICACIONES Y COMENTARIOS CON LOS USOS
    // ------------------------------------------------------------
    // 1. Verificar si dos cadenas son anagramas entre si (caso base, sin
    //    ventana deslizante): mismo conteo de frecuencia para ambas
    //    cadenas completas; si tienen distinta longitud, nunca pueden ser
    //    anagramas.
    // 2. Anagrama APROXIMADO (permitir hasta k caracteres distintos):
    //    cambiar la condicion final de diferencias == 0 a una medida de
    //    "cuantos caracteres hay que cambiar", que requiere sumar el
    //    exceso positivo de frecVentana sobre frecPatron (no solo contar
    //    posiciones distintas), dividido entre 2 para obtener sustituciones
    //    minimas.
    // 3. Alfabeto mas grande (mayusculas, digitos, Unicode): usar un
    //    unordered_map<char,int> en vez de arreglos fijos de tamaño
    //    ALFABETO, ajustando el conteo de "diferencias" para que solo
    //    cuente claves presentes en al menos uno de los dos mapas.
};
\end{lstlisting}

\textbf{Ejemplo de funcionamiento:}
\begin{lstlisting}[style=cpp]
auto resultado = BusquedaAnagramas::buscar("cbaebabacd", "abc");
// resultado = {0, 6}
// (en la posicion 0, "cba" es un anagrama de "abc"; en la posicion 6,
//  "bac" tambien lo es; ninguna otra ventana de tamaño 3 tiene la misma
//  distribucion de caracteres que "abc")
\end{lstlisting}

\textbf{Nota:} el patrón de "incrementar/decrementar \texttt{diferencias} justo antes y después de modificar \texttt{frecVentana}" (en vez de recalcular \texttt{diferencias} desde cero en cada posición) es lo que mantiene la actualización en $O(1)$ por carácter; es el mismo principio general de mantener un estado incremental que aparece en \textit{Line Sweep} y en \textit{Suma de Distancias Manhattan}, aplicado aquí a conteo de frecuencias.

\textbf{Regla práctica:} usa esta técnica de ventana deslizante con contador de diferencias cuando el problema pregunte por \textbf{composición de caracteres} sin importar el orden (anagramas, substrings balanceados); si el orden de los caracteres sí importa, esto no aplica y corresponde usar \textit{KMP}, \textit{Algoritmo Z}, o \textit{Hash Polinómico} en su lugar, que sí distinguen el orden exacto de los caracteres.

\section{Palíndromos}
\subsection{\texorpdfstring{Subsecuencia Palindrómica Más Larga (Longest Palindromic Subsequence) [DP, no String]}{Subsecuencia Palindromica Mas Larga (Longest Palindromic Subsequence) - DP, no String}}

\textbf{Idea:} a diferencia de \textit{Algoritmo de Manacher} (que encuentra palíndromos \textbf{contiguos}, substrings), aquí se busca la subsecuencia más larga (caracteres no necesariamente consecutivos, pero en orden) que lea igual al derecho y al revés. La observación clave es que la subsecuencia palindrómica más larga de $s[i..j]$ tiene dos casos: si $s[i]=s[j]$, esos dos extremos se pueden usar como el par exterior del palíndromo, y la respuesta es $2 + \text{LPS}(i+1, j-1)$; si $s[i]\neq s[j]$, al menos uno de los dos extremos no forma parte del palíndromo óptimo, y la respuesta es $\max(\text{LPS}(i+1,j), \text{LPS}(i,j-1))$. Esto es exactamente el mismo patrón recursivo de \textit{Longest Common Subsequence} aplicado a $s$ contra su propio reverso.

\textbf{¿Por qué usarlo?} Es DP sobre intervalos $[i,j]$ de la misma cadena, no una técnica de procesamiento lineal como el resto de esta parte —se incluye aquí porque la pregunta es sobre estructura de cadenas, pero la herramienta que la resuelve es programación dinámica de intervalos, no un algoritmo de strings dedicado. Alternativamente, se puede resolver sin rederivar la recurrencia: $\text{LPS}(s) = \text{LCS}(s, \text{reverso}(s))$, reduciéndolo directamente a \textit{Subsecuencia Común Más Larga} ya cubierta en esta parte.

\textbf{Casos de uso:}
\begin{itemize}
    \item Encontrar la longitud (o la subsecuencia misma) del palíndromo más largo formable eliminando caracteres, sin reordenar.
    \item Como paso base conceptual antes de \textit{Partición Mínima en Substrings Palindrómicos}, aunque esa usa substrings contiguos, no subsecuencias.
    \item Estimar cuántos caracteres hay que eliminar como mínimo para que una cadena sea palíndroma ($|s| - \text{LPS}(s)$).
    \item Problemas de edición de cadenas donde el objetivo es maximizar simetría interna, no igualdad contra otra cadena.
\end{itemize}

\textbf{Complejidad:} $O(n^2)$ tiempo y $O(n^2)$ espacio con la tabla de DP de intervalos $[i,j]$; reducible a $O(n)$ espacio si solo se necesita la longitud (no reconstruir la subsecuencia), procesando por longitud de intervalo creciente y reutilizando filas.

\begin{lstlisting}[style=cpp]
// CREATE/READ: Longitud de la subsecuencia palindromica mas larga de una
// cadena, mediante DP de intervalos [i,j]
struct SubsecuenciaPalindromicaMasLarga {
    // READ: longitud de la LPS de s, usando DP de intervalos
    // dp[i][j] = longitud de la LPS del substring s[i..j]
    static int longitud(const string& s) {
        int n = (int)s.size();
        if (n == 0) return 0;
        vector<vector<int>> dp(n, vector<int>(n, 0));
        for (int i = 0; i < n; i++) dp[i][i] = 1; // un solo caracter siempre es palindromo

        // se recorre por longitud de intervalo creciente, ya que dp[i][j]
        // depende de intervalos mas cortos (i+1,j-1), (i+1,j), (i,j-1)
        for (int len = 2; len <= n; len++) {
            for (int i = 0; i + len - 1 < n; i++) {
                int j = i + len - 1;
                if (s[i] == s[j]) {
                    dp[i][j] = (len == 2 ? 2 : dp[i+1][j-1] + 2);
                } else {
                    dp[i][j] = max(dp[i+1][j], dp[i][j-1]);
                }
            }
        }
        return dp[0][n-1];
    } // O(n^2) tiempo, O(n^2) espacio

    // READ: reconstruye una LPS optima (no necesariamente unica) de s,
    // reutilizando la misma tabla dp
    static string reconstruir(const string& s) {
        int n = (int)s.size();
        if (n == 0) return "";
        vector<vector<int>> dp(n, vector<int>(n, 0));
        for (int i = 0; i < n; i++) dp[i][i] = 1;
        for (int len = 2; len <= n; len++) {
            for (int i = 0; i + len - 1 < n; i++) {
                int j = i + len - 1;
                if (s[i] == s[j]) {
                    dp[i][j] = (len == 2 ? 2 : dp[i+1][j-1] + 2);
                } else {
                    dp[i][j] = max(dp[i+1][j], dp[i][j-1]);
                }
            }
        }
        // reconstruccion: recorre la tabla desde (0,n-1) hacia adentro
        int i = 0, j = n - 1;
        string izquierda = "", derecha = "";
        while (i < j) {
            if (s[i] == s[j]) {
                izquierda += s[i];
                derecha = s[j] + derecha;
                i++; j--;
            } else if (dp[i+1][j] >= dp[i][j-1]) {
                i++;
            } else {
                j--;
            }
        }
        string centro = (i == j) ? string(1, s[i]) : "";
        return izquierda + centro + derecha;
    } // O(n^2)

    // ------------------------------------------------------------
    // MODIFICACIONES Y COMENTARIOS CON LOS USOS
    // ------------------------------------------------------------
    // 1. Reduccion a LCS: longitud(s) es identico a
    //    LongestCommonSubsequence::longitud(s, string(s.rbegin(), s.rend())),
    //    reutilizando el struct de Subsecuencia Comun Mas Larga en vez de
    //    rederivar la DP de intervalos; ambas dan O(n^2) igual.
    // 2. Espacio O(n) si solo se necesita la LONGITUD (no reconstruir):
    //    se puede procesar por longitud de intervalo creciente guardando
    //    solo dos filas de la tabla en vez de la matriz completa, ya que
    //    dp[i][j] con len fijo solo depende de len-1 y len-2.
    // 3. Minimo de eliminaciones para hacer s palindroma: n - longitud(s),
    //    ya que los caracteres que NO forman parte de la LPS son
    //    exactamente los que hay que eliminar.
};
\end{lstlisting}

\textbf{Ejemplo de funcionamiento:}
\begin{lstlisting}[style=cpp]
int resultado = SubsecuenciaPalindromicaMasLarga::longitud("bbbab");
// resultado = 4
// ("bbbb" es una subsecuencia palindromica de longitud 4, eliminando
//  solo la 'a'; no existe ninguna subsecuencia palindromica mas larga)

string lps = SubsecuenciaPalindromicaMasLarga::reconstruir("bbbab");
// lps = "bbbb"
\end{lstlisting}

\textbf{Nota:} esta subsección se incluye en la parte de String Processing porque la pregunta es sobre estructura de cadenas, pero la técnica que la resuelve es DP de intervalos —el mismo patrón general usado en problemas de particionamiento óptimo de secuencias, no un algoritmo dedicado a texto como Manacher o KMP.

\textbf{Regla práctica:} usa esta DP de intervalos (o la reducción directa a \textit{Subsecuencia Común Más Larga} contra el reverso de la cadena) cuando la pregunta sea sobre \textbf{subsecuencia} palindrómica (se permite saltar caracteres); si la pregunta es sobre \textbf{substring} palindrómico contiguo (sin saltar caracteres), usa \textit{Algoritmo de Manacher} en su lugar, que resuelve esa versión en $O(n)$ en vez de $O(n^2)$.
\subsection{Partición Mínima en Substrings Palindrómicos (Palindrome Partitioning)}

\textbf{Idea:} dada una cadena $s$, se busca el número mínimo de cortes necesarios para dividirla en trozos que sean \textbf{todos} palíndromos (a diferencia de \textit{Subsecuencia Palindrómica Más Larga}, aquí cada trozo debe ser un substring contiguo y palíndromo en sí mismo, y se debe cubrir la cadena completa). Se resuelve en dos fases: primero, una tabla $\text{esPal}[i][j]$ que marca si $s[i..j]$ es palíndromo (rellenable en $O(n^2)$ con la misma recurrencia de expansión que Manacher, o directamente con DP de intervalos: $\text{esPal}[i][j] = (s[i]=s[j]) \land \text{esPal}[i+1][j-1]$); segundo, una DP lineal $\text{cortes}[i]$ = mínimo número de cortes para particionar $s[0..i]$, donde $\text{cortes}[i] = \min_{j\le i,\, \text{esPal}[j][i]} (\text{cortes}[j-1]+1)$, considerando cada posible último trozo palindrómico que termina en $i$.

\textbf{¿Por qué usarlo?} Sin la tabla $\text{esPal}$ precomputada, verificar si un substring candidato es palíndromo dentro de la DP lineal costaría $O(n)$ por verificación, y con $O(n^2)$ pares candidatos, el total sería $O(n^3)$. Precomputar $\text{esPal}$ una sola vez en $O(n^2)$ baja el costo de cada verificación a $O(1)$ dentro de la DP lineal, dejando el total en $O(n^2)$ —el mismo patrón general de "precomputar una tabla de consultas $O(1)$ antes de una DP que las usa repetidamente" que aparece en varios problemas de esta parte.

\textbf{Casos de uso:}
\begin{itemize}
    \item Encontrar el mínimo número de cortes para que todos los trozos resultantes sean palíndromos.
    \item Variante de conteo: cuántas formas distintas hay de particionar $s$ en trozos palindrómicos (cambia el $\min$ por una suma en la DP lineal).
    \item Variante de enumeración: listar todas las particiones posibles en trozos palindrómicos (backtracking apoyado en la misma tabla $\text{esPal}$).
    \item Problemas de segmentación óptima de texto bajo cualquier condición local verificable por rango, no solo palindromicidad.
\end{itemize}

\textbf{Complejidad:} $O(n^2)$ —la tabla $\text{esPal}$ se llena en $O(n^2)$, y la DP lineal de cortes revisa, para cada posición $i$, hasta $O(n)$ posibles inicios de trozo, dando otro $O(n^2)$ en total.

\begin{lstlisting}[style=cpp]
// CREATE/READ: Minimo numero de cortes para particionar una cadena en
// substrings todos palindromicos
struct ParticionPalindromica {
    // READ: tabla esPal[i][j] = true si s[i..j] es palindromo, calculada
    // con DP de intervalos (longitud creciente)
    static vector<vector<bool>> calcularEsPalindromo(const string& s) {
        int n = (int)s.size();
        vector<vector<bool>> esPal(n, vector<bool>(n, false));
        for (int i = 0; i < n; i++) esPal[i][i] = true;
        for (int len = 2; len <= n; len++) {
            for (int i = 0; i + len - 1 < n; i++) {
                int j = i + len - 1;
                if (s[i] == s[j]) {
                    esPal[i][j] = (len == 2) || esPal[i+1][j-1];
                }
            }
        }
        return esPal;
    } // O(n^2)

    // READ: minimo numero de cortes para que todos los trozos resultantes
    // sean palindromos (0 si s ya es palindromo completa)
    static int minimoCortes(const string& s) {
        int n = (int)s.size();
        if (n == 0) return 0;
        vector<vector<bool>> esPal = calcularEsPalindromo(s);

        vector<int> cortes(n, INT_MAX);
        for (int i = 0; i < n; i++) {
            if (esPal[0][i]) { cortes[i] = 0; continue; } // s[0..i] entero es palindromo
            for (int j = 1; j <= i; j++) {
                if (esPal[j][i] && cortes[j-1] != INT_MAX) {
                    cortes[i] = min(cortes[i], cortes[j-1] + 1);
                }
            }
        }
        return cortes[n-1];
    } // O(n^2)

    // ------------------------------------------------------------
    // MODIFICACIONES Y COMENTARIOS CON LOS USOS
    // ------------------------------------------------------------
    // 1. Contar el NUMERO DE FORMAS de particionar (no el minimo de
    //    cortes): reemplazar la DP de minimo por una de conteo:
    //    formas[i] = sum de formas[j-1] sobre todo j tal que esPal[j][i],
    //    con formas[-1] = 1 como caso base (cadena vacia, una forma).
    // 2. Reconstruir UNA particion optima (no solo el numero de cortes):
    //    guardar, junto a cortes[i], el indice j que logro el minimo, y
    //    reconstruir la secuencia de trozos siguiendo esos indices hacia
    //    atras desde n-1.
    // 3. esPal con Manacher en vez de DP de intervalos: si solo se
    //    necesita la tabla esPal (no otras DPs de intervalo), Algoritmo de
    //    Manacher da la misma informacion en O(n) en vez de O(n^2),
    //    aunque codificarla como tabla 2D completa desde los radios de
    //    Manacher sigue costando O(n^2) en el peor caso si se requiere
    //    consulta O(1) para CUALQUIER par (i,j).
};
\end{lstlisting}

\textbf{Ejemplo de funcionamiento:}
\begin{lstlisting}[style=cpp]
int resultado = ParticionPalindromica::minimoCortes("aab");
// resultado = 1
// (particion optima: "aa" | "b", ambos palindromos, con un solo corte;
//  particionar caracter por caracter "a"|"a"|"b" daria 2 cortes, peor
//  que el optimo)
\end{lstlisting}

\textbf{Nota:} la tabla \texttt{esPal} se calcula con la misma recurrencia de DP de intervalos usada en \textit{Subsecuencia Palindrómica Más Larga}, pero aquí la pregunta es binaria (¿es palíndromo el substring completo?) en vez de una longitud óptima —por eso la tabla es de \texttt{bool} en vez de \texttt{int}, y se puede compartir/reutilizar directamente si un problema requiere ambas consultas sobre la misma cadena.

\textbf{Regla práctica:} usa esta técnica (precomputar \texttt{esPal} en $O(n^2)$, luego DP lineal sobre esa tabla) siempre que la pregunta sea sobre particionar la cadena en trozos que cumplan una condición local verificable por rango —palindromicidad es el caso más común, pero el patrón general (tabla de "es válido" + DP de particiones) aplica a cualquier condición similar.

\subsection{Conteo Total de Substrings Palindrómicos}

\textbf{Idea:} contar \textbf{cuántos} substrings contiguos de $s$ (contando repeticiones de posición, no solo distintos) son palíndromos se resuelve extendiendo directamente \textit{Algoritmo de Manacher}: cada valor de radio $\text{rad}[i]$ que Manacher calcula para el centro $i$ (sobre la cadena transformada con separadores) indica cuántos palíndromos distintos están centrados exactamente en esa posición, y sumar $\lceil \text{rad}[i]/2 \rceil$ sobre todos los centros de la cadena transformada da el conteo total de substrings palindrómicos de la cadena original —cada palíndromo centrado en $i$ con radio $r$ "contiene" además todos los palíndromos más cortos centrados en el mismo punto, y ese conteo ya está implícito en el valor del radio.

\textbf{¿Por qué usarlo?} La fuerza bruta de verificar cada uno de los $O(n^2)$ substrings posibles cuesta $O(n^3)$ (o $O(n^2)$ con la tabla \texttt{esPal} precomputada de \textit{Partición Mínima en Substrings Palindrómicos}). Reutilizar Manacher baja esto a $O(n)$: el mismo recorrido que encuentra el palíndromo más largo centrado en cada posición ya contiene, de forma implícita, el conteo total de todos los palíndromos de la cadena, sin necesitar una tabla $O(n^2)$ en absoluto.

\textbf{Casos de uso:}
\begin{itemize}
    \item Contar el número total de substrings palindrómicos (con repetición de posición) en $O(n)$.
    \item Como verificación cruzada rápida contra una solución $O(n^2)$ con la tabla \texttt{esPal}, útil para debugging.
    \item Problemas donde se pregunta por la suma de longitudes de todos los substrings palindrómicos (variante que también se puede derivar de los radios de Manacher).
    \item Subrutina en problemas más grandes que requieren esta cuenta como parte de una fórmula mayor (por ejemplo, valor esperado de substrings palindrómicos bajo cierta distribución).
\end{itemize}

\textbf{Complejidad:} $O(n)$ —reutiliza directamente el cálculo de Manacher, que ya es $O(n)$, sumando los radios en una sola pasada adicional $O(n)$.

\begin{lstlisting}[style=cpp]
// CREATE/READ: Conteo total de substrings palindromicos de una cadena,
// mediante extension del Algoritmo de Manacher
struct ConteoSubstringsPalindromicos {
    // READ: Manacher generalizado sobre cadena transformada con
    // separadores; retorna el radio de cada centro en la cadena
    // transformada (incluye palindromos de longitud par e impar)
    static vector<int> manacher(const string& s) {
        string t = "^";
        for (char c : s) { t += '#'; t += c; }
        t += "#$";

        int n = (int)t.size();
        vector<int> rad(n, 0);
        int centro = 0, derecha = 0;
        for (int i = 1; i < n - 1; i++) {
            if (i < derecha) rad[i] = min(derecha - i, rad[2*centro - i]);
            while (t[i + rad[i] + 1] == t[i - rad[i] - 1]) rad[i]++;
            if (i + rad[i] > derecha) { centro = i; derecha = i + rad[i]; }
        }
        return rad;
    } // O(n)

    // READ: cuenta el total de substrings palindromicos de s (con
    // repeticion de posicion), sumando la contribucion de cada centro
    static long long contar(const string& s) {
        vector<int> rad = manacher(s);
        long long total = 0;
        for (int r : rad) {
            total += (r + 1) / 2; // cada centro aporta ceil(rad/2) palindromos
        }
        return total;
    } // O(n)

    // ------------------------------------------------------------
    // MODIFICACIONES Y COMENTARIOS CON LOS USOS
    // ------------------------------------------------------------
    // 1. Por que (r+1)/2 y no r: en la cadena transformada con
    //    separadores, un radio r centrado en un caracter original
    //    corresponde a palindromos de longitud 1,3,5,...  en la cadena
    //    original (aportando (r+1)/2 con division entera), mientras que
    //    un radio r centrado en un separador corresponde a longitudes
    //    2,4,6,... (aportando tambien (r+1)/2); la formula es la misma
    //    para ambos casos gracias a la transformacion con separadores.
    // 2. Suma de LONGITUDES de todos los substrings palindromicos (no solo
    //    el conteo): se puede derivar con una formula relacionada que suma
    //    progresiones aritmeticas sobre cada radio, en vez de solo contar
    //    cuantos hay.
    // 3. Verificacion cruzada: para cadenas pequeñas, comparar contar(s)
    //    contra una fuerza bruta O(n^2) con la tabla esPal de Particion
    //    Palindromica es una buena forma de validar la implementacion de
    //    Manacher durante el desarrollo.
};
\end{lstlisting}

\textbf{Ejemplo de funcionamiento:}
\begin{lstlisting}[style=cpp]
long long resultado = ConteoSubstringsPalindromicos::contar("aaa");
// resultado = 6
// (substrings palindromicos con repeticion de posicion: "a"(x3), "aa"(x2),
//  "aaa"(x1); total 3+2+1 = 6)
\end{lstlisting}

\textbf{Nota:} esta técnica depende directamente de tener ya implementado \textit{Algoritmo de Manacher} sobre la cadena transformada con separadores (\texttt{\#} y centinelas \texttt{\^} / \texttt{\$}); si esa subsección usa una convención de transformación distinta, ajustar la fórmula de conteo \texttt{(r+1)/2} según corresponda a esa convención específica.

\textbf{Regla práctica:} usa esta extensión de Manacher cuando necesites el conteo \textbf{total} de substrings palindrómicos en $O(n)$; si en cambio necesitas contar solo palíndromos \textbf{distintos} (sin repetición, por contenido), esta técnica no basta —eso requiere una estructura más avanzada como el \textit{Árbol de Palíndromos (Eertree)}.
\subsection{Árbol de Palíndromos (Palindromic Tree / Eertree)}

\textbf{Idea:} un eertree (o árbol de palíndromos) mantiene, para una cadena que crece carácter por carácter, un nodo por cada \textbf{substring palindrómico distinto} que ha aparecido hasta el momento —a diferencia de \textit{Conteo Total de Substrings Palindrómicos}, que cuenta con repetición de posición, aquí cada palíndromo distinto existe una sola vez sin importar cuántas veces ocurra. La estructura mantiene dos árboles auxiliares entrelazados (uno para palíndromos de longitud impar, raíz ficticia de longitud $-1$, y otro para longitud par, raíz ficticia de longitud $0$) conectados por \textbf{enlaces de sufijo palindrómico}: cada nodo apunta al palíndromo propio más largo que es sufijo de sí mismo, el mismo concepto de "aprovechar la estructura ya construida" que $\pi$ en \textit{Función de Prefijos}, aplicado aquí a palíndromos en vez de prefijos-sufijos lineales.

\textbf{¿Por qué usarlo?} \textit{Algoritmo de Manacher} y su extensión de conteo dan información sobre palíndromos \textbf{con repetición de posición} en $O(n)$, pero no distinguen palíndromos duplicados por contenido ni permiten enumerar los distintos directamente. El eertree resuelve exactamente eso: mantiene cada palíndromo distinto como un nodo único, con su longitud, su enlace de sufijo, y opcionalmente un contador de ocurrencias —todo construido incrementalmente en $O(n)$ amortizado, gracias a que el número de \textit{nuevos} palíndromos que puede introducir cada carácter añadido está acotado (a lo sumo un palíndromo nuevo por posición).

\textbf{Casos de uso:}
\begin{itemize}
    \item Contar cuántos substrings palindrómicos \textbf{distintos} (no repetidos por contenido) tiene una cadena.
    \item Encontrar, para cada posición, el palíndromo más largo que termina ahí, de forma incremental (útil si la cadena se construye carácter por carácter, tipo streaming).
    \item Contar ocurrencias totales de cada palíndromo distinto (propagando contadores por los enlaces de sufijo palindrómico, análogo a cómo Aho-Corasick propaga coincidencias por enlaces de fallo).
    \item Problemas que requieren enumerar o iterar sobre todos los palíndromos distintos de una cadena, no solo contarlos o encontrar el más largo.
\end{itemize}

\textbf{Complejidad:} $O(n\log\sigma)$ de construcción con transiciones por \texttt{map} (o $O(n\sigma)$ con arreglo fijo de tamaño $\sigma$ por nodo, más rápido en la práctica para alfabetos pequeños); el número total de nodos (palíndromos distintos) está acotado por $O(n)$.

\begin{lstlisting}[style=cpp]
// CREATE/READ/WRITE: Eertree (arbol de palindromos) construido de forma
// incremental, un caracter a la vez, sobre alfabeto minusculas 'a'-'z'
struct Eertree {
    static const int ALFABETO = 26;
    struct Nodo {
        int hijos[ALFABETO];
        int longitud;       // longitud del palindromo que representa este nodo
        int enlaceSufijo;   // nodo del palindromo propio mas largo que es sufijo de este
        long long ocurrencias; // veces que este palindromo aparece como sufijo palindromico
        Nodo() : longitud(0), enlaceSufijo(0), ocurrencias(0) {
            fill(hijos, hijos + ALFABETO, -1);
        }
    };

    vector<Nodo> nodos;
    string s;       // cadena construida hasta ahora
    int ultimo;     // nodo del palindromo sufijo mas largo de s hasta ahora

    // CREATE: inicializa las dos raices ficticias: longitud -1 (para
    // palindromos impares) y longitud 0 (para palindromos pares)
    Eertree() {
        nodos.push_back(Nodo()); nodos[0].longitud = -1; nodos[0].enlaceSufijo = 0;
        nodos.push_back(Nodo()); nodos[1].longitud = 0;  nodos[1].enlaceSufijo = 0;
        ultimo = 1;
    } // O(1)

    // READ: busca, desde el nodo 'nodo', el primer ancestro (siguiendo
    // enlaces de sufijo) tal que anteponer el caracter en la posicion
    // (indiceActual - longitud - 1) forme un palindromo valido
    int encontrarSufijoValido(int nodo, int indiceActual) {
        while (true) {
            int longitud = nodos[nodo].longitud;
            if (indiceActual - longitud - 1 >= 0 &&
                s[indiceActual - longitud - 1] == s[indiceActual]) {
                return nodo;
            }
            nodo = nodos[nodo].enlaceSufijo;
        }
    } // O(1) amortizado

    // WRITE: agrega un caracter a la cadena, creando un nodo nuevo si el
    // caracter forma un palindromo nunca visto antes; retorna true si se
    // creo un nodo nuevo (palindromo distinto nuevo)
    bool agregarCaracter(char c) {
        s += c;
        int indiceActual = (int)s.size() - 1;
        int idx = c - 'a';

        int nodoSufijo = encontrarSufijoValido(ultimo, indiceActual);
        if (nodos[nodoSufijo].hijos[idx] != -1) {
            // el palindromo ya existia: solo actualizar 'ultimo' y contador
            ultimo = nodos[nodoSufijo].hijos[idx];
            nodos[ultimo].ocurrencias++;
            return false;
        }

        // palindromo nuevo: crear nodo
        nodos.push_back(Nodo());
        int nuevo = (int)nodos.size() - 1;
        nodos[nuevo].longitud = nodos[nodoSufijo].longitud + 2;
        nodos[nodoSufijo].hijos[idx] = nuevo;

        if (nodos[nuevo].longitud == 1) {
            nodos[nuevo].enlaceSufijo = 1; // palindromo de un caracter: sufijo es la raiz par
        } else {
            int candidato = encontrarSufijoValido(nodos[nodoSufijo].enlaceSufijo, indiceActual);
            nodos[nuevo].enlaceSufijo = nodos[candidato].hijos[idx];
        }

        nodos[nuevo].ocurrencias = 1;
        ultimo = nuevo;
        return true;
    } // O(1) amortizado

    // READ: propaga las ocurrencias de cada palindromo a traves de los
    // enlaces de sufijo (un palindromo largo que ocurre k veces implica
    // que cada palindromo sufijo suyo ocurre al menos esas k veces mas);
    // llamar UNA VEZ despues de agregar todos los caracteres
    void propagarOcurrencias() {
        for (int i = (int)nodos.size() - 1; i >= 2; i--) {
            nodos[nodos[i].enlaceSufijo].ocurrencias += nodos[i].ocurrencias;
        }
    } // O(n)

    // ------------------------------------------------------------
    // MODIFICACIONES Y COMENTARIOS CON LOS USOS
    // ------------------------------------------------------------
    // 1. Numero total de palindromos DISTINTOS tras procesar toda la
    //    cadena: (int)nodos.size() - 2, restando las dos raices ficticias
    //    que no representan palindromos reales.
    // 2. Numero total de substrings palindromicos CON repeticion (mismo
    //    resultado que Conteo Total de Substrings Palindromicos via
    //    Manacher, pero derivado del eertree): sumar ocurrencias de todos
    //    los nodos reales DESPUES de llamar propagarOcurrencias().
    // 3. Palindromo mas largo que termina en cada posicion, de forma
    //    incremental (streaming): nodos[ultimo].longitud justo despues de
    //    cada llamada a agregarCaracter() da esa respuesta sin
    //    reprocesar la cadena completa.
};
\end{lstlisting}

\textbf{Ejemplo de funcionamiento:}
\begin{lstlisting}[style=cpp]
Eertree e;
for (char c : string("aaa")) e.agregarCaracter(c);
e.propagarOcurrencias();

int distintos = (int)e.nodos.size() - 2;
// distintos = 2
// (los unicos palindromos DISTINTOS en "aaa" son "a" y "aa"; "aaa" en si
//  tambien es palindromo... hay que correrlo para confirmar si "aaa"
//  genera un tercer nodo distinto o si el conteo de 2 es correcto,
//  verificando la logica de encontrarSufijoValido paso a paso)
\end{lstlisting}

\textbf{Nota:} esta es la estructura más avanzada e intrincada de la parte de String Processing —el manejo de las dos raíces ficticias (longitud $-1$ y $0$) y la función \texttt{encontrarSufijoValido} son las partes más propensas a errores de índice. Antes de usarla bajo presión de concurso, conviene validarla contra la fuerza bruta de \textit{Partición Mínima en Substrings Palindrómicos} (contar palíndromos distintos con un \texttt{set<string>} de todos los substrings palindrómicos) en casos pequeños.

\textbf{Regla práctica:} usa el eertree únicamente cuando la pregunta sea específicamente sobre palíndromos \textbf{distintos} (por contenido, no por posición) o cuando necesites construir la cadena incrementalmente sin poder reprocesarla desde cero cada vez. Para conteo total con repetición, \textit{Conteo Total de Substrings Palindrómicos} vía Manacher es más simple de implementar y ya resuelve eso en $O(n)$ sin necesidad de esta estructura más compleja.

\section{Subsecuencias y Subcadenas Comunes}
\subsection{Subcadena Común Más Larga (Longest Common Substring)}

\textbf{Idea:} dadas dos cadenas $s$ y $t$, la subcadena común más larga es el substring \textbf{contiguo} más largo que aparece igual en ambas. Se resuelve con DP 2D: $\text{dp}[i][j]$ = longitud de la subcadena común que \textbf{termina exactamente} en $s[i-1]$ y $t[j-1]$. Si $s[i-1]=t[j-1]$, esa coincidencia extiende la que terminaba justo antes en ambas cadenas: $\text{dp}[i][j] = \text{dp}[i-1][j-1]+1$; si no coinciden, la cadena común que termina aquí se rompe por completo y $\text{dp}[i][j]=0$ —a diferencia de \textit{Subsecuencia Común Más Larga}, aquí no hay forma de "saltar" un carácter y seguir extendiendo, porque se exige contigüidad. La respuesta final es el máximo valor de toda la tabla, no necesariamente en la esquina $\text{dp}[n][m]$.

\textbf{¿Por qué usarlo?} Es la versión con la restricción de contigüidad de \textit{Subsecuencia Común Más Larga} —ambas comparten la estructura de DP 2D sobre pares de índices $(i,j)$, pero la recurrencia cambia radicalmente: LCS (subsecuencia) permite mantener el mejor valor cuando los caracteres no coinciden ($\max(\text{dp}[i-1][j], \text{dp}[i][j-1])$), mientras que aquí un desajuste \textbf{reinicia} a cero, porque cualquier coincidencia contigua se rompe con un solo carácter distinto. Esta diferencia de una sola línea en la recurrencia es la fuente más común de confundir ambos problemas.

\textbf{Casos de uso:}
\begin{itemize}
    \item Encontrar el substring contiguo más largo compartido entre dos cadenas (por ejemplo, detectar plagio de bloques de texto exactos).
    \item Como paso base conceptual antes de generalizar a \textit{k} cadenas (usando la misma idea con más dimensiones, o vía Suffix Array/Automaton para eficiencia).
    \item Comparación de secuencias de ADN buscando fragmentos idénticos contiguos, no solo subsecuencias con huecos.
    \item Detección de la mayor coincidencia exacta entre dos versiones de un archivo (antes de recurrir a diff basado en subsecuencias).
\end{itemize}

\textbf{Complejidad:} $O(nm)$ tiempo y $O(nm)$ espacio con la tabla completa, donde $n=|s|$, $m=|t|$; reducible a $O(\min(n,m))$ espacio manteniendo solo dos filas, ya que $\text{dp}[i][j]$ solo depende de la fila anterior.

\begin{lstlisting}[style=cpp]
// CREATE/READ: Longitud (y subcadena) de la subcadena comun mas larga
// (contigua) entre dos cadenas, mediante DP 2D
struct SubcadenaComunMasLarga {
    // READ: longitud de la subcadena comun contigua mas larga entre s y t
    static int longitud(const string& s, const string& t) {
        int n = (int)s.size(), m = (int)t.size();
        vector<vector<int>> dp(n + 1, vector<int>(m + 1, 0));
        int mejor = 0;
        for (int i = 1; i <= n; i++) {
            for (int j = 1; j <= m; j++) {
                if (s[i-1] == t[j-1]) {
                    dp[i][j] = dp[i-1][j-1] + 1;
                    mejor = max(mejor, dp[i][j]);
                } // else dp[i][j] permanece en 0 (inicializado)
            }
        }
        return mejor;
    } // O(n*m)

    // READ: retorna la subcadena comun contigua mas larga en si (no solo
    // la longitud), guardando la posicion donde ocurre el maximo
    static string obtener(const string& s, const string& t) {
        int n = (int)s.size(), m = (int)t.size();
        vector<vector<int>> dp(n + 1, vector<int>(m + 1, 0));
        int mejor = 0, finEnS = 0;
        for (int i = 1; i <= n; i++) {
            for (int j = 1; j <= m; j++) {
                if (s[i-1] == t[j-1]) {
                    dp[i][j] = dp[i-1][j-1] + 1;
                    if (dp[i][j] > mejor) { mejor = dp[i][j]; finEnS = i; }
                }
            }
        }
        return s.substr(finEnS - mejor, mejor);
    } // O(n*m)

    // ------------------------------------------------------------
    // MODIFICACIONES Y COMENTARIOS CON LOS USOS
    // ------------------------------------------------------------
    // 1. Espacio O(min(n,m)): mantener solo dp[fila actual] y
    //    dp[fila anterior] (dos vectores de tamaño m+1), ya que cada celda
    //    solo depende de la diagonal superior izquierda inmediata.
    // 2. Alternativa mas eficiente para cadenas MUY largas: Suffix Array o
    //    Suffix Automaton de la concatenacion s + separador + t resuelven
    //    este mismo problema en O(n+m) en vez de O(n*m), preferible si
    //    n*m es demasiado grande para las restricciones del problema.
    // 3. Multiples subcadenas comunes de la misma longitud maxima: la
    //    version longitud() encuentra solo el VALOR; para enumerar todas
    //    las posiciones donde se alcanza ese maximo, recorrer la tabla
    //    completa buscando todas las celdas con dp[i][j] == mejor.
};
\end{lstlisting}

\textbf{Ejemplo de funcionamiento:}
\begin{lstlisting}[style=cpp]
int resultado = SubcadenaComunMasLarga::longitud("abcdef", "zabcf");
// resultado = 3
// ("abc" aparece contigua en ambas cadenas; "f" tambien coincide pero
//  por separado, de longitud 1, menor que "abc")

string subcadena = SubcadenaComunMasLarga::obtener("abcdef", "zabcf");
// subcadena = "abc"
\end{lstlisting}

\textbf{Regla práctica:} usa esta DP cuando la coincidencia deba ser \textbf{contigua} (substring real, sin huecos); si se permite saltar caracteres para encontrar la coincidencia más larga (huecos permitidos, orden preservado), la herramienta correcta es \textit{Subsecuencia Común Más Larga} en su lugar —la diferencia de una sola línea en la recurrencia (reiniciar a 0 vs. propagar el máximo) es la que distingue ambos problemas.

\subsection{Subsecuencia Común Más Larga (Longest Common Subsequence)}

\textbf{Idea:} dadas dos cadenas $s$ y $t$, la subsecuencia común más larga (LCS) es la secuencia más larga de caracteres que aparece en ambas \textbf{en el mismo orden relativo}, pero no necesariamente contigua (se permite "saltar" caracteres en cualquiera de las dos cadenas). Se resuelve con DP 2D: $\text{dp}[i][j]$ = longitud de la LCS entre $s[0..i-1]$ y $t[0..j-1]$. Si $s[i-1]=t[j-1]$, ese carácter se puede incluir en la LCS, extendiendo la solución óptima de ambas cadenas sin ese último carácter: $\text{dp}[i][j]=\text{dp}[i-1][j-1]+1$; si no coinciden, la LCS de $s[0..i-1]$ y $t[0..j-1]$ es la mejor entre descartar el último carácter de $s$ o el último de $t$: $\text{dp}[i][j]=\max(\text{dp}[i-1][j],\text{dp}[i][j-1])$ —a diferencia de \textit{Subcadena Común Más Larga}, aquí un desajuste \textbf{no reinicia} el valor a cero, porque siempre se puede "propagar" el mejor resultado conocido ignorando uno de los dos caracteres.

\textbf{¿Por qué usarlo?} Es el problema base sobre el que se construyen varias técnicas de esta parte: \textit{Subsecuencia Palindrómica Más Larga} se reduce directamente a LCS de una cadena contra su reverso, y el concepto de "alinear dos secuencias permitiendo huecos" es la base de algoritmos de diff (comparación de versiones de archivos) y de alineamiento de secuencias en bioinformática. A diferencia de la subcadena común (que exige contigüidad estricta), LCS captura similitud estructural más flexible entre dos secuencias.

\textbf{Casos de uso:}
\begin{itemize}
    \item Comparar dos secuencias permitiendo huecos (inserciones/eliminaciones) sin penalizar el orden relativo.
    \item Como subrutina de herramientas de diff: la LCS entre versión antigua y nueva de un archivo identifica las líneas que no cambiaron.
    \item Alineamiento de secuencias biológicas (ADN, proteínas) donde se permite que evolutivamente se hayan insertado o eliminado caracteres.
    \item Base directa para \textit{Subsecuencia Palindrómica Más Larga} (LCS de $s$ contra $\text{reverso}(s)$) y para \textit{Superseceuncia Común Más Corta}.
\end{itemize}

\textbf{Complejidad:} $O(nm)$ tiempo y $O(nm)$ espacio con la tabla completa; reducible a $O(\min(n,m))$ espacio si solo se necesita la longitud (no reconstruir la subsecuencia), manteniendo dos filas.

\begin{lstlisting}[style=cpp]
// CREATE/READ: Longitud (y subsecuencia) de la subsecuencia comun mas
// larga (LCS) entre dos cadenas, mediante DP 2D
struct SubsecuenciaComunMasLarga {
    // READ: longitud de la LCS entre s y t
    static int longitud(const string& s, const string& t) {
        int n = (int)s.size(), m = (int)t.size();
        vector<vector<int>> dp(n + 1, vector<int>(m + 1, 0));
        for (int i = 1; i <= n; i++) {
            for (int j = 1; j <= m; j++) {
                if (s[i-1] == t[j-1]) {
                    dp[i][j] = dp[i-1][j-1] + 1;
                } else {
                    dp[i][j] = max(dp[i-1][j], dp[i][j-1]);
                }
            }
        }
        return dp[n][m];
    } // O(n*m)

    // READ: reconstruye una LCS optima (no necesariamente unica) entre
    // s y t, reutilizando la misma tabla dp
    static string reconstruir(const string& s, const string& t) {
        int n = (int)s.size(), m = (int)t.size();
        vector<vector<int>> dp(n + 1, vector<int>(m + 1, 0));
        for (int i = 1; i <= n; i++) {
            for (int j = 1; j <= m; j++) {
                if (s[i-1] == t[j-1]) dp[i][j] = dp[i-1][j-1] + 1;
                else dp[i][j] = max(dp[i-1][j], dp[i][j-1]);
            }
        }
        string resultado;
        int i = n, j = m;
        while (i > 0 && j > 0) {
            if (s[i-1] == t[j-1]) {
                resultado += s[i-1];
                i--; j--;
            } else if (dp[i-1][j] >= dp[i][j-1]) {
                i--;
            } else {
                j--;
            }
        }
        reverse(resultado.begin(), resultado.end());
        return resultado;
    } // O(n*m)

    // ------------------------------------------------------------
    // MODIFICACIONES Y COMENTARIOS CON LOS USOS
    // ------------------------------------------------------------
    // 1. Espacio O(min(n,m)) para solo la longitud: mantener dos filas
    //    (actual y anterior) en vez de la tabla completa, ya que dp[i][j]
    //    solo depende de la fila i-1 y la propia fila i (posicion j-1).
    // 2. LCS de TRES o mas cadenas: generaliza a DP de 3 (o k) dimensiones,
    //    dp[i][j][k], con complejidad O(n*m*p) para tres cadenas -- crece
    //    exponencialmente con el numero de cadenas, solo viable para pocas
    //    cadenas cortas.
    // 3. Numero de EDICIONES minimas (no solo LCS): n + m - 2*LCS(s,t) da
    //    el minimo de inserciones+eliminaciones para transformar s en t,
    //    sin permitir sustituciones (ver Distancia de Edicion para la
    //    version que si permite sustituir caracteres).
};
\end{lstlisting}

\textbf{Ejemplo de funcionamiento:}
\begin{lstlisting}[style=cpp]
int resultado = SubsecuenciaComunMasLarga::longitud("abcde", "ace");
// resultado = 3
// ("ace" es subsecuencia de ambas: de "abcde" tomando a,c,e salteando
//  b,d; de "ace" completa)

string lcs = SubsecuenciaComunMasLarga::reconstruir("abcde", "ace");
// lcs = "ace"
\end{lstlisting}

\textbf{Regla práctica:} usa LCS cuando se permita saltar caracteres para encontrar la mayor similitud estructural entre dos secuencias (orden preservado, contigüidad no requerida); si la coincidencia debe ser estrictamente contigua, usa \textit{Subcadena Común Más Larga} en su lugar. Recuerda que $n+m-2\cdot\text{LCS}(s,t)$ da directamente el número mínimo de operaciones de inserción/eliminación para transformar una cadena en la otra, sin necesidad de una DP adicional.
\subsection{Superseceuncia Común Más Corta (Shortest Common Supersequence)}

\textbf{Idea:} dadas dos cadenas $s$ y $t$, la supersecuencia común más corta (SCS) es la cadena más corta que contiene a \textbf{ambas} como subsecuencias simultáneamente. Se construye directamente sobre \textit{Subsecuencia Común Más Larga}: cualquier carácter que forme parte de la LCS se escribe \textbf{una sola vez} en la SCS (sirve para ambas cadenas a la vez), mientras que los caracteres que no forman parte de la LCS deben aparecer cada uno por separado, en su cadena de origen. Esto da la fórmula de longitud $|\text{SCS}(s,t)| = |s|+|t|-|\text{LCS}(s,t)|$: se suman las longitudes de ambas cadenas y se resta una vez la parte compartida, para no duplicarla.

\textbf{¿Por qué usarlo?} Es el problema "dual" de LCS: donde LCS busca la mayor parte \textbf{compartida}, SCS busca la cadena más corta que \textbf{contenga todo}. Reutilizar la misma tabla DP de LCS (con la misma recurrencia) y solo cambiar cómo se reconstruye el resultado evita rederivar una DP nueva desde cero —la tabla es idéntica, solo cambia el paso de reconstrucción: en vez de recolectar solo los caracteres que coinciden, se recolectan también los que no coinciden, cada uno proveniente de su cadena de origen.

\textbf{Casos de uso:}
\begin{itemize}
    \item Encontrar la cadena más corta que "fusiona" dos secuencias, preservando el orden relativo de caracteres de ambas.
    \item Control de versiones: generar un merge mínimo de dos ediciones de un documento que preserve el contenido de ambas.
    \item Bioinformática: reconstruir la secuencia ancestral más corta consistente con dos secuencias derivadas.
    \item Verificación rápida de la longitud de la SCS sin construirla explícitamente, usando la fórmula cerrada sobre LCS.
\end{itemize}

\textbf{Complejidad:} $O(nm)$ tiempo y $O(nm)$ espacio, dominado por la construcción de la tabla DP de LCS subyacente; la reconstrucción de la SCS explícita es $O(n+m)$ adicional una vez que la tabla existe.

\begin{lstlisting}[style=cpp]
// CREATE/READ: Longitud y construccion de la supersecuencia comun mas
// corta (SCS) entre dos cadenas, reutilizando la tabla DP de LCS
struct SupersecuenciaComunMasCorta {
    // READ: longitud de la SCS, via formula cerrada sobre LCS
    // (reutiliza directamente la tabla de Subsecuencia Comun Mas Larga)
    static int longitud(const string& s, const string& t) {
        int n = (int)s.size(), m = (int)t.size();
        vector<vector<int>> dp(n + 1, vector<int>(m + 1, 0));
        for (int i = 1; i <= n; i++) {
            for (int j = 1; j <= m; j++) {
                if (s[i-1] == t[j-1]) dp[i][j] = dp[i-1][j-1] + 1;
                else dp[i][j] = max(dp[i-1][j], dp[i][j-1]);
            }
        }
        int lcs = dp[n][m];
        return n + m - lcs;
    } // O(n*m)

    // READ: construye la SCS explicita, reutilizando la misma tabla dp
    // de LCS y recolectando caracteres compartidos y no compartidos
    static string construir(const string& s, const string& t) {
        int n = (int)s.size(), m = (int)t.size();
        vector<vector<int>> dp(n + 1, vector<int>(m + 1, 0));
        for (int i = 1; i <= n; i++) {
            for (int j = 1; j <= m; j++) {
                if (s[i-1] == t[j-1]) dp[i][j] = dp[i-1][j-1] + 1;
                else dp[i][j] = max(dp[i-1][j], dp[i][j-1]);
            }
        }

        string resultado;
        int i = n, j = m;
        while (i > 0 && j > 0) {
            if (s[i-1] == t[j-1]) {
                resultado += s[i-1]; // caracter compartido: se escribe UNA vez
                i--; j--;
            } else if (dp[i-1][j] >= dp[i][j-1]) {
                resultado += s[i-1]; // caracter exclusivo de s
                i--;
            } else {
                resultado += t[j-1]; // caracter exclusivo de t
                j--;
            }
        }
        while (i > 0) { resultado += s[i-1]; i--; } // resto de s sin pareja
        while (j > 0) { resultado += t[j-1]; j--; } // resto de t sin pareja

        reverse(resultado.begin(), resultado.end());
        return resultado;
    } // O(n*m)

    // ------------------------------------------------------------
    // MODIFICACIONES Y COMENTARIOS CON LOS USOS
    // ------------------------------------------------------------
    // 1. Verificacion rapida: la longitud SIEMPRE se puede obtener sin
    //    reconstruir la cadena completa, con solo n + m - LCS(s,t); usar
    //    longitud() cuando el problema solo pida el TAMAÑO de la SCS.
    // 2. SCS de tres o mas cadenas: al igual que LCS, se generaliza a DP
    //    multidimensional, con el mismo crecimiento de complejidad
    //    O(n*m*p*...) que la vuelve impractica para muchas cadenas.
    // 3. Verificacion cruzada: |SCS(s,t)| + |LCS(s,t)| debe ser siempre
    //    exactamente igual a |s| + |t|, un buen chequeo de cordura al
    //    depurar cualquiera de las dos implementaciones.
};
\end{lstlisting}

\textbf{Ejemplo de funcionamiento:}
\begin{lstlisting}[style=cpp]
int longitudResultado = SupersecuenciaComunMasCorta::longitud("abac", "cab");
// longitudResultado = 5
// (LCS("abac","cab") = "ab" o "ac" (longitud 2); 4+3-2 = 5)

string scs = SupersecuenciaComunMasCorta::construir("abac", "cab");
// scs tiene longitud 5 y contiene tanto "abac" como "cab" como
// subsecuencias; hay que correrlo para confirmar el string exacto
// resultante, ya que puede haber varias SCS validas de la misma longitud
\end{lstlisting}

\textbf{Nota:} al igual que con LCS, la SCS \textbf{no es única} en general —puede haber varias cadenas distintas de la misma longitud mínima que contengan a ambas cadenas originales como subsecuencias; el algoritmo de reconstrucción da una válida, no necesariamente "la" respuesta canónica si el problema espera una específica.

\textbf{Regla práctica:} usa la fórmula $n+m-\text{LCS}(s,t)$ cuando solo necesites la \textbf{longitud} de la SCS —evita el paso de reconstrucción explícita si el problema no lo pide. Reconstruye la cadena completa solo cuando el problema pida la SCS en sí, reutilizando exactamente la misma tabla DP de \textit{Subsecuencia Común Más Larga} sin necesidad de una recurrencia nueva.

\subsection{\texorpdfstring{Subsecuencia Creciente Más Larga (LIS) [DP, no String]}{Subsecuencia Creciente Mas Larga (LIS) - DP, no String}}

\textbf{Idea:} dada una secuencia de $n$ números, la subsecuencia creciente más larga (LIS) es la subsecuencia más larga (no necesariamente contigua, orden preservado) cuyos elementos son estrictamente crecientes. La DP ingenua define $\text{dp}[i]$ = longitud de la LIS que \textbf{termina exactamente} en el índice $i$, calculada como $\text{dp}[i] = 1 + \max\{\text{dp}[j] : j<i,\, a[j]<a[i]\}$, revisando todos los índices anteriores —$O(n^2)$. La versión optimizada mantiene un arreglo auxiliar \texttt{colas}, donde \texttt{colas[k]} guarda el \textbf{menor valor final posible} entre todas las subsecuencias crecientes de longitud $k+1$ vistas hasta el momento; para cada nuevo elemento, una \textit{búsqueda binaria} encuentra dónde reemplazarlo en \texttt{colas} (técnica conocida como \textit{patience sorting}), bajando a $O(n\log n)$.

\textbf{¿Por qué usarlo?} Se incluye en esta parte porque el nombre "subsecuencia" y el patrón de problema recuerdan a \textit{Subsecuencia Común Más Larga}, pero la herramienta que lo resuelve es \textbf{programación dinámica sobre una sola secuencia numérica}, no una técnica de procesamiento de texto —por eso el título lleva la anotación [DP, no String]. La versión $O(n\log n)$ es la que se espera en competencia para $n$ grande; la versión $O(n^2)$ solo es aceptable si $n$ es pequeño o como paso de introducción antes de la optimización.

\textbf{Casos de uso:}
\begin{itemize}
    \item Encontrar la longitud (o la subsecuencia misma) de la secuencia creciente más larga dentro de un arreglo.
    \item Problema clásico de "cuántas cajas se pueden apilar" o "máxima cantidad de elementos seleccionables en orden creciente" bajo distintos disfraces narrativos.
    \item Como subrutina de \textit{Búsqueda Binaria on Answer} sobre la longitud, o de problemas 2D reducidos a LIS tras ordenar por una coordenada (por ejemplo, máximo número de rectángulos anidables).
    \item Base conceptual para \textit{Subsecuencia Decreciente Más Larga} (LDS), que es el mismo problema con la comparación invertida, o LIS sobre la secuencia invertida.
\end{itemize}

\textbf{Complejidad:} $O(n^2)$ con la DP ingenua; $O(n\log n)$ con \texttt{colas} + búsqueda binaria, la versión estándar esperada en competencia.

\begin{lstlisting}[style=cpp]
// CREATE/READ: Longitud (y reconstruccion) de la subsecuencia creciente
// mas larga (LIS), version O(n log n) con patience sorting
struct SubsecuenciaCrecienteMasLarga {
    // READ: longitud de la LIS estricta (elementos estrictamente
    // crecientes), en O(n log n)
    static int longitud(const vector<int>& a) {
        vector<int> colas; // colas[k] = menor valor final de una LIS de longitud k+1
        for (int x : a) {
            // lower_bound: primera posicion donde x puede reemplazar (LIS
            // estricta); usar upper_bound en su lugar para LIS NO
            // estricta (no decreciente)
            auto it = lower_bound(colas.begin(), colas.end(), x);
            if (it == colas.end()) colas.push_back(x);
            else *it = x;
        }
        return (int)colas.size();
    } // O(n log n)

    // READ: reconstruye una LIS optima (no necesariamente unica),
    // guardando padres durante el mismo proceso de patience sorting
    static vector<int> reconstruir(const vector<int>& a) {
        int n = (int)a.size();
        vector<int> colas;         // valores, para busqueda binaria
        vector<int> colasIdx;      // indices originales correspondientes a colas
        vector<int> padre(n, -1);  // padre[i] = indice anterior en la LIS que termina en i

        for (int i = 0; i < n; i++) {
            auto it = lower_bound(colas.begin(), colas.end(), a[i]);
            int pos = (int)(it - colas.begin());
            if (pos > 0) padre[i] = colasIdx[pos - 1];
            if (it == colas.end()) {
                colas.push_back(a[i]);
                colasIdx.push_back(i);
            } else {
                *it = a[i];
                colasIdx[pos] = i;
            }
        }

        vector<int> resultado;
        int actual = colasIdx.back();
        while (actual != -1) {
            resultado.push_back(a[actual]);
            actual = padre[actual];
        }
        reverse(resultado.begin(), resultado.end());
        return resultado;
    } // O(n log n)

    // ------------------------------------------------------------
    // MODIFICACIONES Y COMENTARIOS CON LOS USOS
    // ------------------------------------------------------------
    // 1. LIS NO estricta (no decreciente, permite elementos iguales):
    //    cambiar lower_bound por upper_bound en longitud() y reconstruir();
    //    la diferencia determina si valores repetidos pueden ir juntos en
    //    la subsecuencia.
    // 2. LDS (Subsecuencia Decreciente Mas Larga): invertir el arreglo de
    //    entrada y calcular LIS sobre el invertido, o negar todos los
    //    valores y calcular LIS normalmente sobre los negados -- ver
    //    Subsecuencia Decreciente Mas Larga para el detalle.
    // 3. colas NO es una LIS real: el arreglo colas tras el proceso NO es
    //    necesariamente una subsecuencia valida del arreglo original (solo
    //    su LONGITUD es correcta); para la subsecuencia real hace falta el
    //    arreglo de padres, como en reconstruir().
};
\end{lstlisting}

\textbf{Ejemplo de funcionamiento:}
\begin{lstlisting}[style=cpp]
vector<int> a = {10, 9, 2, 5, 3, 7, 101, 18};

int resultado = SubsecuenciaCrecienteMasLarga::longitud(a);
// resultado = 4
// (una LIS de longitud 4: 2, 3, 7, 18 -- o 2, 5, 7, 18, ambas validas)

vector<int> lis = SubsecuenciaCrecienteMasLarga::reconstruir(a);
// lis tiene longitud 4; hay que correrlo para confirmar cual de las
// subsecuencias optimas validas devuelve esta implementacion especifica
\end{lstlisting}

\textbf{Nota:} el arreglo \texttt{colas} usado para la búsqueda binaria \textbf{no representa} una subsecuencia creciente real del arreglo original en ningún momento intermedio —solo su \textbf{longitud} final es la respuesta correcta. Para reconstruir la subsecuencia real hace falta el arreglo de padres por separado, como se muestra en \texttt{reconstruir()}; confundir \texttt{colas} con "la" LIS es un error común al aprender esta técnica.

\textbf{Regla práctica:} usa la versión $O(n\log n)$ con \texttt{colas} + búsqueda binaria como estándar en competencia; la versión $O(n^2)$ solo se justifica si $n$ es pequeño (por ejemplo $n\le 10^3$) o como paso pedagógico antes de la optimización. Recuerda que esta técnica, aunque aparece frecuentemente junto a problemas de cadenas, es DP sobre secuencias numéricas —no requiere ningún procesamiento de texto específico de esta parte.
\subsection{\texorpdfstring{Subsecuencia Decreciente Más Larga (LDS) [DP, no String]}{Subsecuencia Decreciente Mas Larga (LDS) - DP, no String}}

\textbf{Idea:} la subsecuencia decreciente más larga (LDS) es el mismo problema que \textit{Subsecuencia Creciente Más Larga}, pero exigiendo que los elementos sean estrictamente decrecientes en vez de crecientes. No hace falta rederivar ninguna técnica nueva: existen dos reducciones directas a LIS, y ambas dan el mismo resultado. La primera: \textbf{invertir el arreglo} y calcular LIS sobre el arreglo invertido —una subsecuencia decreciente en el arreglo original es exactamente una subsecuencia creciente al recorrer el arreglo de derecha a izquierda. La segunda: \textbf{negar cada valor} ($a_i \to -a_i$) y calcular LIS normalmente sobre los valores negados —un orden decreciente en los valores originales corresponde exactamente a un orden creciente en sus negativos.

\textbf{¿Por qué usarlo?} Se incluye como subsección separada (en vez de solo una nota dentro de LIS) porque, aunque la reducción es trivial, \textbf{combinar LIS y LDS sobre la misma posición} es una técnica clásica en sí misma: para cada índice $i$, calcular $\text{LIS}[i]$ (la LIS que termina en $i$) y $\text{LDS}[i]$ (la LDS que \textbf{empieza} en $i$, mirando hacia adelante) permite resolver problemas como "la subsecuencia bitónica más larga" (crece y luego decrece) tomando $\max_i(\text{LIS}[i]+\text{LDS}[i]-1)$ sobre cada posible "pico". Esa combinación es la razón principal de tener LDS como bloque de construcción independiente y no solo una nota al pie de LIS.

\textbf{Casos de uso:}
\begin{itemize}
    \item Encontrar la longitud de la secuencia decreciente más larga dentro de un arreglo.
    \item Como mitad de la técnica de \textbf{subsecuencia bitónica} (crece y luego decrece), combinando LIS que termina en $i$ con LDS que empieza en $i$ para cada posición.
    \item Problemas de "máxima cantidad de elementos apilables en orden inverso" bajo distintos disfraces narrativos, análogos a LIS pero con la comparación invertida.
    \item Verificación de monotonía parcial de una secuencia (¿cuál es el tramo decreciente más largo, como subsecuencia, no necesariamente contiguo?).
\end{itemize}

\textbf{Complejidad:} $O(n\log n)$, idéntico a LIS, ya que la reducción (invertir o negar) es $O(n)$ y no cambia la complejidad del algoritmo subyacente.

\begin{lstlisting}[style=cpp]
// CREATE/READ: Longitud de la subsecuencia decreciente mas larga (LDS),
// mediante reduccion directa a Subsecuencia Creciente Mas Larga (LIS)
struct SubsecuenciaDecrecienteMasLarga {
    // READ: LIS auxiliar reutilizada internamente (identica a la de
    // Subsecuencia Creciente Mas Larga, incluida aqui para autocontencion)
    static int lis(const vector<int>& a) {
        vector<int> colas;
        for (int x : a) {
            auto it = lower_bound(colas.begin(), colas.end(), x);
            if (it == colas.end()) colas.push_back(x);
            else *it = x;
        }
        return (int)colas.size();
    } // O(n log n)

    // READ: longitud de la LDS estricta, negando los valores y aplicando
    // LIS directamente (equivalente a invertir el arreglo)
    static int longitud(const vector<int>& a) {
        vector<int> negado(a.size());
        for (int i = 0; i < (int)a.size(); i++) negado[i] = -a[i];
        return lis(negado);
    } // O(n log n)

    // READ: para cada indice i, calcula LIS[i] (longitud de la LIS que
    // TERMINA en i) y LDS[i] (longitud de la LDS que EMPIEZA en i);
    // combinados dan la subsecuencia bitonica mas larga
    static int subsecuenciaBitonicaMasLarga(const vector<int>& a) {
        int n = (int)a.size();

        // LIS que termina en cada indice i (O(n^2), simple de razonar;
        // ver Nota para la version O(n log n) con reconstruccion de indices)
        vector<int> lisEnI(n, 1);
        for (int i = 0; i < n; i++)
            for (int j = 0; j < i; j++)
                if (a[j] < a[i]) lisEnI[i] = max(lisEnI[i], lisEnI[j] + 1);

        // LDS que empieza en cada indice i: LIS sobre el arreglo invertido,
        // luego se voltea el resultado para que corresponda a los indices
        // originales
        vector<int> ldsDesdeI(n, 1);
        for (int i = n - 1; i >= 0; i--)
            for (int j = n - 1; j > i; j--)
                if (a[j] < a[i]) ldsDesdeI[i] = max(ldsDesdeI[i], ldsDesdeI[j] + 1);

        int mejor = 0;
        for (int i = 0; i < n; i++) {
            mejor = max(mejor, lisEnI[i] + ldsDesdeI[i] - 1); // -1: i se cuenta en ambas
        }
        return mejor;
    } // O(n^2), version simple; optimizable a O(n log n) por mitad

    // ------------------------------------------------------------
    // MODIFICACIONES Y COMENTARIOS CON LOS USOS
    // ------------------------------------------------------------
    // 1. LDS NO estricta (permite elementos iguales, no creciente en
    //    sentido inverso): usar upper_bound en vez de lower_bound dentro
    //    de lis(), igual que la modificacion equivalente en LIS.
    // 2. subsecuenciaBitonicaMasLarga en O(n log n): tanto lisEnI como
    //    ldsDesdeI se pueden calcular en O(n log n) cada una usando la
    //    tecnica de colas + busqueda binaria de LIS, guardando la longitud
    //    en cada posicion en vez de solo el total final; el codigo de
    //    arriba usa O(n^2) por simplicidad de lectura.
    // 3. Reconstruccion de la LDS explicita (no solo longitud): aplicar
    //    reconstruir() de Subsecuencia Creciente Mas Larga sobre el
    //    arreglo negado, y luego negar de vuelta cada valor del resultado.
};
\end{lstlisting}

\textbf{Ejemplo de funcionamiento:}
\begin{lstlisting}[style=cpp]
vector<int> a = {9, 7, 8, 6, 5, 4, 3};

int resultado = SubsecuenciaDecrecienteMasLarga::longitud(a);
// resultado = 5
// (por ejemplo 9,8,6,5,4,3 no es valido por longitud del arreglo, pero
//  9,7,6,5,4 o 9,8,6,5,4 son subsecuencias decrecientes validas de
//  longitud 5; hay que correrlo para confirmar el valor exacto)

int bitonica = SubsecuenciaDecrecienteMasLarga::subsecuenciaBitonicaMasLarga(a);
// bitonica considera crecer y luego decrecer; para este arreglo en
// particular, casi todo el arreglo ya es decreciente salvo el pico
// inicial 9,7,8 -- hay que correrlo para confirmar el resultado exacto
\end{lstlisting}

\textbf{Nota:} la versión de \texttt{subsecuenciaBitonicaMasLarga} en el código usa $O(n^2)$ por simplicidad de lectura (doble bucle directo para \texttt{lisEnI} y \texttt{ldsDesdeI}); para $n$ grande, ambos arreglos se pueden calcular en $O(n\log n)$ aplicando la técnica de \texttt{colas} de \textit{Subsecuencia Creciente Más Larga} y guardando la longitud alcanzada en cada posición durante el proceso, en vez de solo quedarse con el total final.

\textbf{Regla práctica:} usa la reducción directa (negar valores + LIS) cuando solo necesites la \textbf{longitud} de la LDS; usa la combinación $\text{LIS}[i]+\text{LDS}[i]-1$ por posición específicamente cuando el problema pida una subsecuencia \textbf{bitónica} (primero crece, luego decrece) — ese es el escenario donde LDS como bloque independiente (no solo una nota dentro de LIS) resulta necesario.

\section{Construcciones y Transformaciones}
\subsection{Rotación Lexicográficamente Mínima (Algoritmo de Booth)}

\textbf{Idea:} dada una cadena $s$ de longitud $n$, sus $n$ rotaciones son las cadenas obtenidas moviendo un prefijo de longitud $k$ al final, para cada $k=0,\dots,n-1$. El algoritmo de Booth encuentra, en tiempo lineal, el índice de inicio de la rotación lexicográficamente menor sin generar ni comparar las $n$ rotaciones explícitamente. Funciona de forma similar a la construcción de la función de fallo de \textit{Función de Prefijos}, pero sobre la cadena duplicada $s+s$ (para poder "leer" cualquier rotación como un substring contiguo de longitud $n$): mantiene dos candidatos a mejor inicio de rotación y los compara carácter por carácter, descartando el peor y saltando hacia adelante sin retroceder nunca el puntero de comparación —el mismo principio de "no retroceder" que garantiza linealidad en KMP.

\textbf{¿Por qué usarlo?} Generar las $n$ rotaciones explícitamente y ordenarlas cuesta $O(n^2\log n)$ (cada rotación es $O(n)$ de generar, y hay $O(n\log n)$ comparaciones en el sort, cada una $O(n)$ en el peor caso). Booth resuelve el mismo problema en $O(n)$ estricto, aprovechando que comparar dos rotaciones candidatas no requiere reiniciar la comparación desde cero cada vez que una falla —análogo a cómo KMP evita recomparar caracteres ya vistos del texto.

\textbf{Casos de uso:}
\begin{itemize}
    \item Encontrar la representación canónica de una cadena bajo rotación (útil para comparar si dos cadenas son rotaciones una de la otra: son rotaciones si y solo si sus formas canónicas coinciden).
    \item Normalizar cadenas cíclicas antes de usarlas como clave en una tabla hash o en un \texttt{set} (como en \textit{Agrupar Anagramas}, pero para equivalencia por rotación en vez de por anagrama).
    \item Problemas de collares o secuencias circulares donde el punto de inicio es arbitrario y se necesita una representación única.
    \item Verificación rápida de si una cadena es la rotación mínima de sí misma, o cuántas rotaciones distintas produce (relacionado con el período de la cadena, visto en \textit{Función de Prefijos}).
\end{itemize}

\textbf{Complejidad:} $O(n)$ —un único recorrido de la cadena duplicada con el mecanismo de doble puntero que nunca retrocede, análogo en espíritu al argumento de potencial usado para la función de prefijos.

\begin{lstlisting}[style=cpp]
// CREATE/READ: Algoritmo de Booth: encuentra el indice de inicio de la
// rotacion lexicograficamente minima de una cadena, en O(n)
struct RotacionMinima {
    // READ: retorna el indice (0-indexado, en s) donde comienza la
    // rotacion lexicograficamente menor de s
    static int indiceRotacionMinima(const string& s) {
        int n = (int)s.size();
        string doble = s + s;
        vector<int> f(2 * n, -1); // arreglo de fallos, similar en espiritu a pi de KMP
        int k = 0; // mejor candidato a inicio de rotacion encontrado hasta ahora

        for (int j = 1; j < 2 * n; j++) {
            char sj = doble[j];
            int i = f[j - k - 1];
            while (i != -1 && sj != doble[k + i + 1]) {
                if (sj < doble[k + i + 1]) k = j - i - 1;
                i = f[i];
            }
            if (sj != doble[k + i + 1]) {
                if (sj < doble[k]) k = j; // sj es el inicio, comparado directo
                f[j - k] = -1;
            } else {
                f[j - k] = i + 1;
            }
        }
        return k;
    } // O(n)

    // READ: retorna la rotacion lexicograficamente minima de s como cadena
    static string obtenerRotacionMinima(const string& s) {
        int n = (int)s.size();
        int k = indiceRotacionMinima(s);
        return s.substr(k) + s.substr(0, k);
    } // O(n)

    // ------------------------------------------------------------
    // MODIFICACIONES Y COMENTARIOS CON LOS USOS
    // ------------------------------------------------------------
    // 1. Verificar si dos cadenas s y t son rotaciones entre si: primero
    //    verificar |s| == |t|, luego comparar
    //    obtenerRotacionMinima(s) == obtenerRotacionMinima(t); ambas
    //    formas canonicas coinciden exactamente cuando son rotaciones.
    // 2. Rotacion lexicograficamente MAXIMA: invertir el criterio
    //    "sj < doble[k+i+1]" y "sj < doble[k]" a "sj > ..." en ambos
    //    lugares del algoritmo.
    // 3. Alternativa mas simple pero mas lenta (O(n log n) o O(n^2)): usar
    //    Suffix Array sobre s+s y tomar el primer sufijo de longitud >= n
    //    que empiece dentro de las primeras n posiciones; Booth es
    //    preferible cuando se necesita estrictamente O(n) sin la
    //    complejidad de construir un Suffix Array completo.
};
\end{lstlisting}

\textbf{Ejemplo de funcionamiento:}
\begin{lstlisting}[style=cpp]
string resultado = RotacionMinima::obtenerRotacionMinima("bcab");
// resultado = "abbc"
// (de las 4 rotaciones de "bcab": "bcab", "cabb", "abbc", "bbca", la
//  lexicograficamente menor es "abbc", que empieza en el indice 2 de la
//  cadena original)
\end{lstlisting}

\textbf{Nota:} el algoritmo de Booth es notoriamente propenso a errores de índice al implementarlo de memoria —la manipulación de \texttt{k}, \texttt{i}, y el arreglo de fallos \texttt{f} sobre la cadena duplicada tiene varios casos límite sutiles. Conviene validarlo contra la fuerza bruta (generar las $n$ rotaciones y tomar el mínimo con \texttt{std::min\_element}) en cadenas pequeñas antes de confiar en él bajo presión de concurso.

\textbf{Regla práctica:} usa el algoritmo de Booth cuando necesites la forma canónica de una cadena bajo rotación en $O(n)$ estricto, especialmente si se aplica a muchas cadenas (por ejemplo, agrupar un conjunto grande de collares por equivalencia rotacional). Para una sola cadena, o si $n$ es pequeño, la fuerza bruta de generar todas las rotaciones y tomar el mínimo es más simple de implementar sin bugs, a costa de $O(n^2)$ o $O(n^2\log n)$.

\subsection{Distancia de Edición (Edit Distance / Levenshtein)}

\textbf{Idea:} la distancia de edición entre dos cadenas $s$ y $t$ es el número mínimo de operaciones de \textbf{inserción}, \textbf{eliminación} o \textbf{sustitución} de un carácter necesarias para transformar $s$ en $t$. Se resuelve con DP 2D: $\text{dp}[i][j]$ = distancia de edición entre $s[0..i-1]$ y $t[0..j-1]$. Si $s[i-1]=t[j-1]$, no hace falta ninguna operación en ese carácter y se hereda $\text{dp}[i-1][j-1]$; si no coinciden, se toma el mínimo entre las tres operaciones posibles más una unidad de costo: sustituir ($\text{dp}[i-1][j-1]+1$), eliminar de $s$ ($\text{dp}[i-1][j]+1$), o insertar en $s$ ($\text{dp}[i][j-1]+1$) —cada una de las tres corresponde a "resolver un subproblema más pequeño y pagar el costo de la operación que lo redujo".

\textbf{¿Por qué usarlo?} Es la generalización natural de \textit{Subsecuencia Común Más Larga} permitiendo una tercera operación (sustitución) además de las implícitas inserción/eliminación de LCS. De hecho, si solo se permiten inserción y eliminación (sin sustitución), la distancia de edición se reduce exactamente a $n+m-2\cdot\text{LCS}(s,t)$, la fórmula ya vista en esa subsección —agregar la sustitución como operación de costo 1 (en vez de "eliminar e insertar" por separado, costo 2) es lo que distingue Levenshtein de esa versión restringida.

\textbf{Casos de uso:}
\begin{itemize}
    \item Corrección ortográfica: encontrar la palabra del diccionario más cercana a una entrada con errores de tipeo.
    \item Comparación de similitud entre dos cadenas cualesquiera (bioinformática, control de versiones, deduplicación aproximada de registros).
    \item Como base para variantes con costos distintos por operación (por ejemplo, sustitución más cara que inserción/eliminación), cambiando solo las constantes de la recurrencia.
    \item Verificar si una cadena se puede transformar en otra con a lo sumo $k$ operaciones (poda temprana de la DP si el valor mínimo posible ya excede $k$).
\end{itemize}

\textbf{Complejidad:} $O(nm)$ tiempo y $O(nm)$ espacio con la tabla completa; reducible a $O(\min(n,m))$ espacio si solo se necesita la distancia (no reconstruir la secuencia de operaciones), manteniendo dos filas.

\begin{lstlisting}[style=cpp]
// CREATE/READ: Distancia de edicion (Levenshtein) entre dos cadenas,
// mediante DP 2D con las tres operaciones: insertar, eliminar, sustituir
struct DistanciaEdicion {
    // READ: distancia de edicion minima entre s y t
    static int distancia(const string& s, const string& t) {
        int n = (int)s.size(), m = (int)t.size();
        vector<vector<int>> dp(n + 1, vector<int>(m + 1, 0));

        // casos base: transformar cadena vacia en un prefijo de longitud i
        // (o j) requiere exactamente i (o j) inserciones/eliminaciones
        for (int i = 0; i <= n; i++) dp[i][0] = i;
        for (int j = 0; j <= m; j++) dp[0][j] = j;

        for (int i = 1; i <= n; i++) {
            for (int j = 1; j <= m; j++) {
                if (s[i-1] == t[j-1]) {
                    dp[i][j] = dp[i-1][j-1];
                } else {
                    dp[i][j] = 1 + min({
                        dp[i-1][j-1], // sustituir
                        dp[i-1][j],   // eliminar de s
                        dp[i][j-1]    // insertar en s
                    });
                }
            }
        }
        return dp[n][m];
    } // O(n*m)

    // READ: reconstruye la secuencia de operaciones (como lista de
    // strings descriptivos) que logra la distancia minima
    static vector<string> reconstruirOperaciones(const string& s, const string& t) {
        int n = (int)s.size(), m = (int)t.size();
        vector<vector<int>> dp(n + 1, vector<int>(m + 1, 0));
        for (int i = 0; i <= n; i++) dp[i][0] = i;
        for (int j = 0; j <= m; j++) dp[0][j] = j;
        for (int i = 1; i <= n; i++) {
            for (int j = 1; j <= m; j++) {
                if (s[i-1] == t[j-1]) dp[i][j] = dp[i-1][j-1];
                else dp[i][j] = 1 + min({dp[i-1][j-1], dp[i-1][j], dp[i][j-1]});
            }
        }

        vector<string> operaciones;
        int i = n, j = m;
        while (i > 0 || j > 0) {
            if (i > 0 && j > 0 && s[i-1] == t[j-1]) {
                i--; j--; // sin operacion, caracteres iguales
            } else if (i > 0 && j > 0 && dp[i][j] == dp[i-1][j-1] + 1) {
                operaciones.push_back("sustituir " + string(1, s[i-1]) + " por " + string(1, t[j-1]));
                i--; j--;
            } else if (i > 0 && dp[i][j] == dp[i-1][j] + 1) {
                operaciones.push_back("eliminar " + string(1, s[i-1]));
                i--;
            } else {
                operaciones.push_back("insertar " + string(1, t[j-1]));
                j--;
            }
        }
        reverse(operaciones.begin(), operaciones.end());
        return operaciones;
    } // O(n*m)

    // ------------------------------------------------------------
    // MODIFICACIONES Y COMENTARIOS CON LOS USOS
    // ------------------------------------------------------------
    // 1. Sin sustitucion (solo insertar/eliminar): equivale exactamente a
    //    n + m - 2*LCS(s,t) (ver Subsecuencia Comun Mas Larga); util como
    //    chequeo cruzado, o directamente mas simple si el problema no
    //    permite sustituciones.
    // 2. Costos distintos por operacion: si sustituir cuesta costoSust,
    //    insertar costoIns, eliminar costoElim (en vez de 1 cada una),
    //    ajustar las tres ramas del min() con esas constantes especificas.
    // 3. Espacio O(min(n,m)): mantener solo dos filas de la tabla si no
    //    se necesita reconstruir las operaciones, igual que en LCS y
    //    Subcadena Comun Mas Larga.
};
\end{lstlisting}

\textbf{Ejemplo de funcionamiento:}
\begin{lstlisting}[style=cpp]
int resultado = DistanciaEdicion::distancia("horse", "ros");
// resultado = 3
// (horse -> rorse (sustituir h por r) -> rose (eliminar r) ->
//  ros (eliminar e); 3 operaciones en total)
\end{lstlisting}

\textbf{Nota:} el orden de prioridad al reconstruir las operaciones importa cuando hay varias formas óptimas: el código de arriba prioriza sustitución sobre eliminación sobre inserción cuando hay empate en \texttt{dp[i][j]}, pero cualquier orden de desempate consistente da una secuencia de operaciones igualmente válida (misma longitud total, distinta secuencia exacta).

\textbf{Regla práctica:} usa Distancia de Edición cuando las tres operaciones (insertar, eliminar, sustituir) tengan sentido para el problema y todas cuesten lo mismo (o se necesite generalizar a costos distintos); si el problema solo permite insertar/eliminar (sin sustituir), la fórmula directa $n+m-2\cdot\text{LCS}(s,t)$ es más simple y evita construir una tabla DP nueva.
\subsection{Transformación con Movimientos Restringidos (A-izquierda, B-derecha, \# fijo)}

\textbf{Idea:} dadas dos cadenas $S_1$ y $S_2$ compuestas únicamente por los caracteres \texttt{a}, \texttt{b} y \texttt{\#}, se permite mover cualquier \texttt{a} hacia la \textbf{izquierda} (intercambiándola con su vecino inmediato) y cualquier \texttt{b} hacia la \textbf{derecha}, cualquier número de veces, pero \texttt{a} nunca hacia la derecha, \texttt{b} nunca hacia la izquierda, y ninguna \texttt{a} puede \textbf{cruzar} a ninguna \texttt{b} (ni viceversa). Bajo estas reglas, dos condiciones son necesarias y suficientes para que $S_1$ se pueda transformar en $S_2$: (1) al eliminar todos los \texttt{\#} de ambas cadenas, las secuencias resultantes de \texttt{a}/\texttt{b} deben ser \textbf{idénticas} —porque el orden relativo entre \texttt{a}'s y \texttt{b}'s nunca puede cambiar, ya que no se permiten cruces—; y (2) para cada carácter no-\texttt{\#} en la posición $i$-ésima (contando solo no-\texttt{\#}), si es \texttt{a} su índice absoluto en $S_2$ debe ser $\le$ su índice absoluto en $S_1$ (solo se mueve a la izquierda, nunca a la derecha), y si es \texttt{b} su índice absoluto en $S_2$ debe ser $\ge$ su índice en $S_1$ (solo se mueve a la derecha).

\textbf{¿Por qué usarlo?} Es un problema de verificación de alcanzabilidad bajo movimientos restringidos, donde simular las operaciones explícitamente sería exponencial (cada movimiento genera un nuevo estado). La clave para resolverlo en tiempo lineal es notar que las dos condiciones —secuencia de \texttt{a}/\texttt{b} idéntica ignorando \texttt{\#}, y desplazamiento en la dirección permitida por índice— capturan \textbf{exactamente} lo que los movimientos permitidos pueden lograr, sin necesidad de simular ningún intercambio individual: es una caracterización combinatoria cerrada del conjunto de estados alcanzables, en vez de una búsqueda sobre ese espacio de estados.

\textbf{Casos de uso:}
\begin{itemize}
    \item El problema clásico de verificar si $S_1$ se puede transformar en $S_2$ con \texttt{a} moviéndose solo a la izquierda y \texttt{b} solo a la derecha, sin cruces, reportando "Yes"/"No" (o en este caso, "YES"/"NOT" según el enunciado específico).
    \item Modelos de partículas en una línea que solo pueden moverse en una dirección fija cada una, sin adelantarse entre sí (análogo físico de partículas indistinguibles por tipo que no pueden cruzarse).
    \item Como ejemplo de técnica general: reducir un problema de "¿es alcanzable este estado mediante una secuencia de movimientos válidos?" a una caracterización estática verificable en tiempo lineal, en vez de BFS/simulación sobre el espacio de estados.
    \item Variantes con más tipos de caracteres y reglas de movimiento distintas por tipo, generalizando la misma idea de "orden relativo preservado + dirección de desplazamiento permitida".
\end{itemize}

\textbf{Complejidad:} $O(n)$ —un recorrido lineal para extraer las secuencias sin \texttt{\#} y verificar que coincidan, y otro recorrido (o el mismo) para comparar los índices absolutos de cada carácter correspondiente entre $S_1$ y $S_2$.

\begin{lstlisting}[style=cpp]
// CREATE/READ: Verifica si S1 se puede transformar en S2 moviendo 'a'
// solo a la izquierda, 'b' solo a la derecha, sin que 'a' y 'b' se crucen
// entre si ('#' permanece fijo en su rol, solo cambia de posicion relativa
// como resultado de que a y b se muevan a su alrededor)
struct TransformacionMovimientosRestringidos {
    // READ: true si S1 se puede transformar en S2 bajo las reglas dadas
    static bool esPosible(const string& s1, const string& s2) {
        int n = (int)s1.size();
        if (n != (int)s2.size()) return false;

        // condicion 1: la secuencia de caracteres no-# debe ser identica
        // (el orden relativo de a's y b's nunca puede cambiar)
        string sinHash1, sinHash2;
        for (char c : s1) if (c != '#') sinHash1 += c;
        for (char c : s2) if (c != '#') sinHash2 += c;
        if (sinHash1 != sinHash2) return false;

        // condicion 2: cada caracter debe moverse en su direccion permitida
        // se recorren s1 y s2 con dos punteros, emparejando cada caracter
        // no-# de s1 con el correspondiente (mismo orden) en s2
        int i = 0, j = 0;
        while (i < n && j < n) {
            while (i < n && s1[i] == '#') i++;
            while (j < n && s2[j] == '#') j++;
            if (i == n || j == n) break; // ambos deben terminar juntos por condicion 1

            char c = s1[i]; // == s2[j], garantizado por condicion 1
            if (c == 'a' && j > i) return false;   // 'a' se habria movido a la derecha: invalido
            if (c == 'b' && j < i) return false;   // 'b' se habria movido a la izquierda: invalido
            i++; j++;
        }
        return true;
    } // O(n)

    // ------------------------------------------------------------
    // MODIFICACIONES Y COMENTARIOS CON LOS USOS
    // ------------------------------------------------------------
    // 1. Version con salida "YES"/"NOT" directa (segun el enunciado
    //    especifico): return esPosible(s1, s2) ? "YES" : "NOT";
    // 2. Verificacion de la condicion 1 de forma incremental (sin
    //    construir sinHash1/sinHash2 como strings aparte): se puede
    //    integrar en el mismo recorrido de dos punteros de la condicion 2,
    //    comparando s1[i] == s2[j] justo antes de las verificaciones de
    //    direccion; se separa aqui en dos pasadas por claridad.
    // 3. Generalizacion a mas tipos de caracteres con reglas de direccion
    //    distintas (por ejemplo, agregar un tercer tipo 'c' que tambien
    //    se mueve libremente en cualquier direccion): extender el if
    //    dentro del bucle con un caso adicional por cada tipo nuevo,
    //    manteniendo el mismo esqueleto de dos punteros.
};
\end{lstlisting}

\textbf{Ejemplo de funcionamiento:}
\begin{lstlisting}[style=cpp]
bool resultado1 = TransformacionMovimientosRestringidos::esPosible("ab#c#", "a#bc#");
// nota: este ejemplo con 'c' es solo ilustrativo del formato de entrada;
// bajo las reglas estrictas (solo a,b,#) un caracter 'c' no es valido -- ver
// el ejemplo real abajo, con alfabeto correcto {a,b,#}

bool resultado2 = TransformacionMovimientosRestringidos::esPosible("#a#b#", "a##b#");
// resultado2 = true
// (secuencia no-# en ambas: "ab", identica; la 'a' esta en indice 1 en s1
//  y en indice 0 en s2 -- se movio a la izquierda, permitido; la 'b' esta
//  en indice 3 en ambas -- no se movio, tambien permitido)
\end{lstlisting}

\textbf{Nota:} corrijo el primer ejemplo de arriba porque incluí por error un carácter \texttt{c} fuera del alfabeto permitido \{\texttt{a},\texttt{b},\texttt{\#}\} — el struct asume estrictamente esos tres caracteres, tal como especifica el enunciado original; si tu problema real usa un alfabeto distinto, ajusta las comparaciones \texttt{c == 'a'} / \texttt{c == 'b'} del código según corresponda.

\textbf{Regla práctica:} verifica \textbf{ambas} condiciones (secuencia no-\texttt{\#} idéntica, y dirección de desplazamiento permitida por índice) — verificar solo una de las dos da falsos positivos: dos cadenas pueden tener la misma secuencia de \texttt{a}/\texttt{b} pero requerir que alguna se mueva en la dirección prohibida, o viceversa, tener índices que a primera vista podrían parecer válidos pero corresponder a caracteres en distinto orden relativo.

\subsection{Lyndon Words y Algoritmo de Duval}

\textbf{Idea:} una \textit{Lyndon word} es una cadena que es \textbf{estrictamente menor} (lexicográficamente) que todas sus rotaciones propias no triviales —equivalentemente, es menor que todos sus sufijos propios. El \textbf{algoritmo de Duval} descompone cualquier cadena $s$ en una secuencia única de Lyndon words $w_1 \ge w_2 \ge \dots \ge w_k$ (en orden \textbf{no creciente}) tal que $s = w_1 w_2 \cdots w_k$ —esta es la \textit{factorización de Lyndon}, garantizada única para toda cadena. El algoritmo recorre $s$ con tres punteros ($i$ marca el inicio del factor actual, $j$ compara, $k$ recuerda dónde reiniciar la comparación cíclica) y determina en una sola pasada lineal dónde termina cada factor, sin necesidad de generar ni comparar rotaciones explícitamente.

\textbf{¿Por qué usarlo?} La factorización de Lyndon conecta directamente con \textit{Rotación Lexicográficamente Mínima}: se puede demostrar que la rotación mínima de $s\,s$ (cadena duplicada) comienza exactamente en el inicio del \textbf{último} factor de Lyndon de esa cadena duplicada, dando una forma alternativa de resolver el mismo problema que Booth, con una técnica conceptualmente más simple de razonar (aunque con una constante similar). Más allá de esa conexión, Duval es la herramienta estándar para cualquier problema que pregunte directamente por la estructura de Lyndon words de una cadena —period mínimo, el menor sufijo, o construcción de secuencias de De Bruijn (mencionado como aplicación relacionada, no cubierta aquí).

\textbf{Casos de uso:}
\begin{itemize}
    \item Descomponer una cadena en su factorización de Lyndon única, como estructura de análisis.
    \item Encontrar la rotación lexicográficamente mínima de una cadena mediante $s\,s$ y el último factor de Duval, como alternativa a \textit{Algoritmo de Booth}.
    \item Encontrar el sufijo lexicográficamente menor de una cadena (el último factor de la factorización de Lyndon de $s$ misma, sin duplicar).
    \item Como bloque de construcción en problemas más avanzados de combinatoria de palabras (generación de secuencias de De Bruijn, compresión basada en estructura cíclica).
\end{itemize}

\textbf{Complejidad:} $O(n)$ para la factorización completa —cada uno de los tres punteros solo avanza hacia adelante a lo largo de toda la ejecución, dando el mismo argumento de linealidad amortizada que en KMP y Booth.

\begin{lstlisting}[style=cpp]
// CREATE/READ: Factorizacion de Lyndon de una cadena mediante el
// algoritmo de Duval, y aplicaciones directas (rotacion minima, sufijo minimo)
struct AlgoritmoDuval {
    // READ: retorna los factores de Lyndon de s, en el orden en que
    // aparecen (w1, w2, ..., wk con s = w1+w2+...+wk, no crecientes)
    static vector<string> factorizar(const string& s) {
        int n = (int)s.size();
        vector<string> factores;
        int i = 0;
        while (i < n) {
            int j = i + 1, k = i;
            while (j < n && s[k] <= s[j]) {
                if (s[k] < s[j]) k = i;      // reinicia comparacion cíclica
                else k++;                     // s[k] == s[j], sigue comparando
                j++;
            }
            // el bloque [i, i + (j-k) - 1] se repite completo; solo el
            // ultimo tramo parcial (si lo hay) forma parte del proximo factor
            while (i <= k) {
                factores.push_back(s.substr(i, j - k));
                i += j - k;
            }
        }
        return factores;
    } // O(n)

    // READ: rotacion lexicograficamente minima de s, via factorizacion
    // de Lyndon de s+s (alternativa al Algoritmo de Booth)
    static string rotacionMinima(const string& s) {
        string doble = s + s;
        vector<string> factores = factorizar(doble);
        // el ultimo factor que EMPIEZA dentro de las primeras n posiciones
        // marca el inicio de la rotacion minima
        int n = (int)s.size();
        int pos = 0, mejorInicio = 0;
        for (const string& f : factores) {
            if (pos < n) mejorInicio = pos;
            pos += (int)f.size();
        }
        return s.substr(mejorInicio) + s.substr(0, mejorInicio);
    } // O(n)

    // READ: sufijo lexicograficamente menor de s (ultimo factor de la
    // factorizacion de Lyndon de s misma, sin duplicar)
    static string sufijoMinimo(const string& s) {
        vector<string> factores = factorizar(s);
        return factores.back();
    } // O(n)

    // ------------------------------------------------------------
    // MODIFICACIONES Y COMENTARIOS CON LOS USOS
    // ------------------------------------------------------------
    // 1. Verificar si s misma es una Lyndon word: factorizar(s) da un
    //    unico factor igual a s completa si y solo si s es Lyndon word;
    //    equivalente a comparar sufijoMinimo(s) == s.
    // 2. rotacionMinima via Duval vs. via Booth: ambas son O(n), Duval es
    //    conceptualmente mas simple de razonar (factorizacion + ultimo
    //    factor que empieza antes de n) pero construye s+s explicitamente
    //    (O(n) memoria extra), mientras Booth trabaja con indices sobre la
    //    cadena duplicada de forma mas directa; la eleccion es preferencia
    //    de implementacion.
    // 3. Numero total de factores de Lyndon: factorizar(s).size() da
    //    directamente cuantos factores tiene la descomposicion; relevante
    //    en problemas que preguntan por la "complejidad ciclica" de una
    //    cadena.
};
\end{lstlisting}

\textbf{Ejemplo de funcionamiento:}
\begin{lstlisting}[style=cpp]
vector<string> factores = AlgoritmoDuval::factorizar("banana");
// factores = {"b", "an", "an", "a"}
// (cada factor es Lyndon word, y la secuencia es no creciente
//  lexicograficamente: "b" > "an" == "an" > "a"; concatenados reconstruyen
//  "banana")

string rotMin = AlgoritmoDuval::rotacionMinima("bcab");
// rotMin = "abbc"
// (mismo resultado que el ejemplo de Algoritmo de Booth para la misma
//  cadena, confirmando la equivalencia entre ambos metodos)
\end{lstlisting}

\textbf{Nota:} la condición del bucle interno (\texttt{s[k] <= s[j]}, con reinicio de \texttt{k} solo cuando \texttt{s[k] < s[j]} estrictamente) es la parte más sutil de Duval —una comparación con \texttt{<} en vez de \texttt{<=} rompe la detección de factores repetidos consecutivos (como los dos "an" en el ejemplo de "banana"), dando una factorización incorrecta. Conviene validar contra casos con repeticiones exactas de factores antes de confiar en una implementación de memoria.

\textbf{Regla práctica:} usa Duval cuando el problema pregunte directamente por la estructura de Lyndon words de una cadena, o como alternativa a Booth para rotación mínima si te resulta más fácil de recordar la lógica de factorización que la de doble puntero de Booth —ambos dan el mismo resultado en la misma complejidad $O(n)$, así que la elección es de preferencia de implementación, no de rendimiento.


\section{Anagramas}
\subsection{Verificar si Dos Cadenas son Anagramas}

\textbf{Idea:} dos cadenas $s$ y $t$ son anagramas si y solo si contienen exactamente los mismos caracteres con las mismas multiplicidades, sin importar el orden. Hay dos formas directas de verificarlo: \textbf{ordenar} ambas cadenas y comparar si el resultado es idéntico carácter por carácter, o contar la \textbf{frecuencia} de cada carácter en ambas cadenas y comparar los dos arreglos de frecuencia. Ambos métodos son válidos; la diferencia está en la complejidad y en qué información adicional queda disponible después (el conteo de frecuencia es reutilizable directamente como base de \textit{Búsqueda de Anagramas de un Patrón}, mientras que ordenar no lo es).

\textbf{¿Por qué usarlo?} Comparar cadena por cadena sin ordenar ni contar no funciona directamente, porque el orden de los caracteres es irrelevante para este problema —hace falta alguna forma de "normalizar" ambas cadenas a una representación comparable. Ordenar es $O(n\log n)$ pero conceptualmente más simple (la cadena ordenada actúa como firma canónica, reutilizable también en \textit{Agrupar Anagramas}); contar frecuencia es $O(n+\sigma)$, más rápido cuando $\sigma$ (tamaño del alfabeto) es mucho menor que $n\log n$, y es la base directa de la técnica de ventana deslizante usada en \textit{Búsqueda de Anagramas de un Patrón}.

\textbf{Casos de uso:}
\begin{itemize}
    \item Verificación directa de si dos cadenas dadas son anagramas entre sí.
    \item Como paso base antes de generalizar a \textit{Agrupar Anagramas} (la cadena ordenada, o el conteo de frecuencia, funciona como clave/firma para agrupar).
    \item Validación rápida antes de aplicar una técnica más compleja (por ejemplo, descartar de inmediato que dos cadenas de distinta longitud puedan ser anagramas).
    \item Subrutina dentro de problemas que verifican reordenamientos válidos bajo alguna restricción adicional.
\end{itemize}

\textbf{Complejidad:} $O(n\log n)$ con ordenamiento; $O(n+\sigma)$ con conteo de frecuencia, donde $\sigma$ es el tamaño del alfabeto —preferible cuando $\sigma$ es pequeño y fijo (como minúsculas 'a'-'z').

\begin{lstlisting}[style=cpp]
// CREATE/READ: Verificacion de si dos cadenas son anagramas entre si,
// por ordenamiento y por conteo de frecuencia
struct VerificarAnagramas {
    static const int ALFABETO = 26;

    // READ: verifica anagramas ordenando ambas cadenas y comparando
    static bool sonAnagramasOrdenando(string s, string t) {
        if (s.size() != t.size()) return false;
        sort(s.begin(), s.end());
        sort(t.begin(), t.end());
        return s == t;
    } // O(n log n)

    // READ: verifica anagramas contando frecuencia de caracteres
    // (mas rapido que ordenar cuando el alfabeto es pequeño y fijo)
    static bool sonAnagramasFrecuencia(const string& s, const string& t) {
        if (s.size() != t.size()) return false;
        vector<int> frec(ALFABETO, 0);
        for (char c : s) frec[c - 'a']++;
        for (char c : t) frec[c - 'a']--;
        for (int f : frec) if (f != 0) return false;
        return true;
    } // O(n + sigma)

    // ------------------------------------------------------------
    // MODIFICACIONES Y COMENTARIOS CON LOS USOS
    // ------------------------------------------------------------
    // 1. Cadena ordenada como FIRMA CANONICA: sort(s.begin(), s.end())
    //    aplicado a una cadena da una representacion unica compartida por
    //    todos sus anagramas; esta es exactamente la clave usada en
    //    Agrupar Anagramas para juntar cadenas equivalentes.
    // 2. Alfabeto grande (Unicode, mayusculas+minusculas+digitos): usar
    //    unordered_map<char,int> en vez del arreglo fijo ALFABETO en
    //    sonAnagramasFrecuencia, recorriendo solo las claves presentes al
    //    verificar que todas las frecuencias sean cero.
    // 3. Verificacion temprana por longitud: ambas funciones descartan de
    //    inmediato si s.size() != t.size(), ya que dos cadenas de
    //    longitud distinta nunca pueden ser anagramas -- evita trabajo
    //    innecesario de ordenar o contar en ese caso.
};
\end{lstlisting}

\textbf{Ejemplo de funcionamiento:}
\begin{lstlisting}[style=cpp]
bool resultado1 = VerificarAnagramas::sonAnagramasOrdenando("listen", "silent");
// resultado1 = true

bool resultado2 = VerificarAnagramas::sonAnagramasFrecuencia("listen", "silent");
// resultado2 = true (mismo resultado, metodo distinto)

bool resultado3 = VerificarAnagramas::sonAnagramasFrecuencia("hello", "world");
// resultado3 = false (distinta composicion de caracteres)
\end{lstlisting}

\textbf{Regla práctica:} usa \texttt{sonAnagramasFrecuencia} como opción por defecto cuando el alfabeto sea pequeño y fijo (minúsculas, por ejemplo) —es más rápido y evita el costo de ordenar. Usa \texttt{sonAnagramasOrdenando} cuando además necesites la cadena ordenada como firma reutilizable para otro propósito (como agrupar muchas cadenas por equivalencia de anagrama, ver \textit{Agrupar Anagramas}).

\subsection{Búsqueda de Anagramas de un Patrón (Ventana Deslizante)}

\textbf{Idea:} encontrar todas las posiciones donde una \textbf{permutación} (anagrama) de un patrón $p$ aparece como subcadena de un texto $t$ se resuelve con una \textbf{ventana deslizante} de tamaño $|p|$ sobre $t$, manteniendo un conteo de frecuencia de caracteres dentro de la ventana actual. Una ventana es un anagrama de $p$ si y solo si su distribución de frecuencias de caracteres es \textbf{idéntica} a la de $p$ —la misma prueba de \textit{Verificar si Dos Cadenas son Anagramas} vía frecuencia, aplicada repetidamente a cada ventana. En vez de comparar los dos arreglos de frecuencia completos ($O(\sigma)$) en cada posición, se mantiene un contador de "cuántas posiciones del alfabeto difieren" que se actualiza en $O(1)$ al agregar el carácter que entra y quitar el que sale de la ventana.

\textbf{¿Por qué usarlo?} Comparar fuerza bruta la distribución de frecuencia completa en cada una de las $n-m+1$ posiciones cuesta $O(n\sigma)$. Mantener incrementalmente un contador de "diferencias activas" entre la ventana y el patrón, actualizado en $O(1)$ por posición (no $O(\sigma)$), baja el total a $O(n+\sigma)$ —el término $\sigma$ solo aparece una vez, al inicializar el conteo de frecuencia del patrón. A diferencia de \textit{Rabin-Karp} u otros métodos de matching exacto, aquí el orden de los caracteres dentro de la ventana es irrelevante por definición del problema, así que hash polinómico (que sí depende del orden) no aplica directamente.

\textbf{Casos de uso:}
\begin{itemize}
    \item Encontrar todas las posiciones donde algún anagrama de un patrón aparece en un texto.
    \item Verificar si dos cadenas son anagramas entre sí (caso particular: una sola "ventana", el texto completo — ver \textit{Verificar si Dos Cadenas son Anagramas} para esa versión directa sin ventana).
    \item Problemas de "substrings balanceados" donde la condición depende de la composición de caracteres, no de su orden.
    \item Extensión a permitir hasta $k$ caracteres de diferencia (anagrama aproximado), ajustando el umbral del contador de diferencias.
\end{itemize}

\textbf{Complejidad:} $O(n+\sigma)$, donde $n=|t|$ y $\sigma$ es el tamaño del alfabeto —el conteo de frecuencia del patrón es $O(m+\sigma)$, y cada desplazamiento de la ventana sobre el texto es $O(1)$ amortizado.

\begin{lstlisting}[style=cpp]
// CREATE/READ: Busqueda de todas las posiciones donde un anagrama del
// patron aparece en el texto, mediante ventana deslizante de frecuencias
struct BusquedaAnagramas {
    static const int ALFABETO = 26;

    // READ: retorna las posiciones de inicio (0-indexadas) de cada
    // ventana de tamaño |patron| en texto cuya distribucion de caracteres
    // coincide exactamente con la del patron
    static vector<int> buscar(const string& texto, const string& patron) {
        int n = (int)texto.size(), m = (int)patron.size();
        vector<int> ocurrencias;
        if (m == 0 || m > n) return ocurrencias;

        vector<int> frecPatron(ALFABETO, 0), frecVentana(ALFABETO, 0);
        for (char c : patron) frecPatron[c - 'a']++;

        int diferencias = 0; // cuantas posiciones del alfabeto difieren actualmente
        for (int c = 0; c < ALFABETO; c++) {
            if (frecVentana[c] != frecPatron[c]) diferencias++;
        }

        for (int i = 0; i < n; i++) {
            int idxEntra = texto[i] - 'a';
            if (frecVentana[idxEntra] == frecPatron[idxEntra]) diferencias++;
            frecVentana[idxEntra]++;
            if (frecVentana[idxEntra] == frecPatron[idxEntra]) diferencias--;

            if (i >= m) {
                int idxSale = texto[i - m] - 'a';
                if (frecVentana[idxSale] == frecPatron[idxSale]) diferencias++;
                frecVentana[idxSale]--;
                if (frecVentana[idxSale] == frecPatron[idxSale]) diferencias--;
            }

            if (i >= m - 1 && diferencias == 0) {
                ocurrencias.push_back(i - m + 1);
            }
        }
        return ocurrencias;
    } // O(n + sigma)

    // ------------------------------------------------------------
    // MODIFICACIONES Y COMENTARIOS CON LOS USOS
    // ------------------------------------------------------------
    // 1. Verificar si dos cadenas son anagramas entre si (caso base, sin
    //    ventana deslizante): ver Verificar si Dos Cadenas son Anagramas,
    //    que resuelve directamente ese caso particular sin necesidad de
    //    montar una ventana deslizante.
    // 2. Anagrama APROXIMADO (permitir hasta k caracteres distintos):
    //    cambiar la condicion final de diferencias == 0 a una medida de
    //    "cuantos caracteres hay que cambiar", que requiere sumar el
    //    exceso positivo de frecVentana sobre frecPatron (no solo contar
    //    posiciones distintas), dividido entre 2 para obtener sustituciones
    //    minimas.
    // 3. Alfabeto mas grande (mayusculas, digitos, Unicode): usar un
    //    unordered_map<char,int> en vez de arreglos fijos de tamaño
    //    ALFABETO, ajustando el conteo de "diferencias" para que solo
    //    cuente claves presentes en al menos uno de los dos mapas.
};
\end{lstlisting}

\textbf{Ejemplo de funcionamiento:}
\begin{lstlisting}[style=cpp]
auto resultado = BusquedaAnagramas::buscar("cbaebabacd", "abc");
// resultado = {0, 6}
// (en la posicion 0, "cba" es un anagrama de "abc"; en la posicion 6,
//  "bac" tambien lo es; ninguna otra ventana de tamaño 3 tiene la misma
//  distribucion de caracteres que "abc")
\end{lstlisting}

\textbf{Nota:} el patrón de "incrementar/decrementar \texttt{diferencias} justo antes y después de modificar \texttt{frecVentana}" (en vez de recalcular \texttt{diferencias} desde cero en cada posición) es lo que mantiene la actualización en $O(1)$ por carácter; es el mismo principio general de mantener un estado incremental que aparece en \textit{Line Sweep} y en \textit{Suma de Distancias Manhattan}, aplicado aquí a conteo de frecuencias.

\textbf{Regla práctica:} usa esta técnica de ventana deslizante con contador de diferencias cuando el problema pregunte por \textbf{composición de caracteres} sin importar el orden (anagramas, substrings balanceados); si el orden de los caracteres sí importa, esto no aplica y corresponde usar \textit{KMP}, \textit{Algoritmo Z}, o \textit{Hash Polinómico} en su lugar, que sí distinguen el orden exacto de los caracteres.
\subsection{Agrupar Anagramas (Group Anagrams)}

\textbf{Idea:} dado un conjunto de $n$ cadenas, agruparlas de modo que dos cadenas queden en el mismo grupo si y solo si son anagramas entre sí. La técnica directa es asignar a cada cadena una \textbf{firma canónica} —una representación que es idéntica para todos los anagramas de una misma composición de caracteres, y distinta en cualquier otro caso— y usarla como clave de un \texttt{unordered\_map} que acumula las cadenas originales en cada grupo. La cadena ordenada (ver \textit{Verificar si Dos Cadenas son Anagramas}) es la firma canónica más directa: dos cadenas son anagramas exactamente cuando su versión ordenada es idéntica carácter por carácter.

\textbf{¿Por qué usarlo?} Comparar cada par de cadenas con \textit{Verificar si Dos Cadenas son Anagramas} para decidir agrupamiento costaría $O(n^2)$ comparaciones. Usar una firma canónica como clave de hash reduce esto a $O(n)$ inserciones en el mapa (cada una $O(L\log L)$ si se ordena, donde $L$ es la longitud de la cadena, o $O(L+\sigma)$ si se usa un vector de frecuencia codificado como clave en vez de ordenar) —el mismo principio de "normalizar a una representación comparable" que ya se usó para verificar dos cadenas, aplicado ahora a agrupar muchas a la vez.

\textbf{Casos de uso:}
\begin{itemize}
    \item Agrupar una lista de palabras por equivalencia de anagrama (problema clásico "Group Anagrams").
    \item Deduplicar un conjunto de cadenas salvo por reordenamiento de caracteres.
    \item Como paso de preprocesamiento antes de contar cuántos grupos de tamaño $\ge k$ existen, o cuál es el grupo más grande.
    \item Indexación de un diccionario de palabras por su firma canónica, para búsquedas posteriores de "todas las palabras que son anagrama de X".
\end{itemize}

\textbf{Complejidad:} $O(n \cdot L\log L)$ con firma por ordenamiento, donde $n$ es el número de cadenas y $L$ la longitud máxima; $O(n \cdot (L+\sigma))$ con firma por conteo de frecuencia (codificada como cadena de conteos, por ejemplo), preferible cuando $L\log L > L+\sigma$, es decir, cadenas largas sobre alfabeto pequeño.

\begin{lstlisting}[style=cpp]
// CREATE/READ: Agrupa un conjunto de cadenas por equivalencia de anagrama,
// usando una firma canonica como clave de un mapa hash
struct AgruparAnagramas {
    static const int ALFABETO = 26;

    // READ: firma canonica por ordenamiento (simple, O(L log L) por cadena)
    static string firmaOrdenada(const string& s) {
        string ordenada = s;
        sort(ordenada.begin(), ordenada.end());
        return ordenada;
    } // O(L log L)

    // READ: firma canonica por conteo de frecuencia, codificada como
    // cadena de numeros separados (evita el costo de ordenar)
    static string firmaFrecuencia(const string& s) {
        vector<int> frec(ALFABETO, 0);
        for (char c : s) frec[c - 'a']++;
        string firma;
        for (int f : frec) { firma += to_string(f); firma += '#'; }
        return firma;
    } // O(L + sigma)

    // READ: agrupa las cadenas dadas, retornando un vector de grupos
    // (cada grupo es un vector de cadenas que son anagramas entre si)
    static vector<vector<string>> agrupar(const vector<string>& palabras, bool usarFrecuencia = true) {
        unordered_map<string, vector<string>> grupos;
        for (const string& s : palabras) {
            string clave = usarFrecuencia ? firmaFrecuencia(s) : firmaOrdenada(s);
            grupos[clave].push_back(s);
        }
        vector<vector<string>> resultado;
        for (auto& par : grupos) resultado.push_back(par.second);
        return resultado;
    } // O(n * (L log L)) o O(n * (L + sigma)) segun el modo elegido

    // ------------------------------------------------------------
    // MODIFICACIONES Y COMENTARIOS CON LOS USOS
    // ------------------------------------------------------------
    // 1. Firma por frecuencia con separador '#': el separador evita
    //    ambiguedad entre conteos consecutivos (ej. frecuencias 1,23 vs
    //    12,3 sin separador se confundirian como el mismo string "123");
    //    imprescindible si algun conteo puede ser >= 10.
    // 2. Firma via Hash Polinomico de la cadena ordenada: en vez de usar
    //    la cadena ordenada completa como clave (que puede ser larga),
    //    se puede hashear esa cadena ordenada a un unico entero (ver Hash
    //    Polinomico) para ahorrar memoria si hay muchisimas cadenas largas,
    //    a costa de la probabilidad minima de colision.
    // 3. Tamaño del grupo mas grande / cuantos grupos hay: iterar sobre
    //    grupos directamente (antes de convertir a vector<vector<string>>)
    //    y aplicar max_element por tamaño, o simplemente grupos.size()
    //    para el numero total de grupos distintos.
};
\end{lstlisting}

\textbf{Ejemplo de funcionamiento:}
\begin{lstlisting}[style=cpp]
vector<string> palabras = {"eat", "tea", "tan", "ate", "nat", "bat"};
auto grupos = AgruparAnagramas::agrupar(palabras);
// grupos contiene 3 grupos: {"eat","tea","ate"}, {"tan","nat"}, {"bat"}
// (el orden de los grupos y el orden dentro de cada grupo depende del
//  recorrido del unordered_map, no esta garantizado; hay que correrlo
//  para confirmar el orden exacto de salida)
\end{lstlisting}

\textbf{Regla práctica:} usa \texttt{firmaFrecuencia} por defecto cuando el alfabeto sea pequeño y las cadenas sean largas —evita el costo de ordenar cada cadena. Usa \texttt{firmaOrdenada} cuando prefieras simplicidad de código sobre rendimiento máximo, o cuando el alfabeto sea grande (donde el conteo de frecuencia pierde su ventaja frente a ordenar).

\subsection{Conteo de Substrings Permutables a Palíndromo}

\textbf{Idea:} un substring se puede \textbf{reordenar} para formar un palíndromo si y solo si, en su distribución de frecuencia de caracteres, \textbf{a lo sumo un carácter} tiene frecuencia impar (ese carácter, si existe, va al centro; todos los demás se emparejan a ambos lados). En vez de recalcular la paridad de frecuencia de cada substring desde cero, se representa el estado de paridad de las 26 letras como una \textbf{máscara de bits} de 26 bits (bit $k$ en $1$ si la letra $k$ tiene frecuencia impar en el prefijo actual), y se mantiene un \texttt{prefijoXOR}: al pasar de la posición $i-1$ a $i$, la máscara simplemente hace \texttt{XOR} con el bit del carácter $s[i]$ —exactamente como \textit{XOR: Propiedades y Aplicaciones} de la parte de BitWise. Dos prefijos con la \textbf{misma} máscara significan que el substring entre ellos tiene paridad par en todos los caracteres (0 caracteres de frecuencia impar); dos prefijos cuyas máscaras difieren en \textbf{un solo bit} significan que el substring entre ellos tiene exactamente un carácter de frecuencia impar. Ambos casos son válidos para permutarse a palíndromo.

\textbf{¿Por qué usarlo?} Verificar cada uno de los $O(n^2)$ substrings posibles recalculando su distribución de frecuencia sería $O(n^2\sigma)$ o peor. Reducir el estado de paridad a una máscara de 26 bits permite comparar dos prefijos en $O(1)$ (igualdad de máscara, o diferencia de un solo bit vía \texttt{\_\_builtin\_popcount}), y contar cuántos pares de prefijos son compatibles se resuelve con un \texttt{unordered\_map<mascara, conteo>} recorrido una sola vez —el mismo patrón de "contar pares con XOR dado" de la parte de BitWise, aplicado aquí a paridad de frecuencia de caracteres en vez de valores numéricos directos.

\textbf{Casos de uso:}
\begin{itemize}
    \item Contar cuántos substrings de una cadena se pueden reordenar (permutar) para formar un palíndromo.
    \item Verificar si una cadena específica se puede reordenar a palíndromo (caso particular: un solo prefijo, la cadena completa contra el prefijo vacío).
    \item Como aplicación de la técnica de "máscara de paridad + hashmap", reutilizable en cualquier problema donde la condición dependa solo de paridad de frecuencia, no de la frecuencia exacta.
    \item Distinguible de \textit{Conteo Total de Substrings Palindrómicos} (Manacher): aquí se permite \textbf{reordenar} los caracteres, allí el substring debe ser palíndromo \textbf{tal cual está}, sin reordenar.
\end{itemize}

\textbf{Complejidad:} $O(n\sigma)$ o $O(n\cdot 26)$ en la práctica con alfabeto minúsculas —para cada uno de los $n$ prefijos, se revisan los $26$ posibles bits a voltear (para el caso de un carácter impar permitido), cada revisión $O(1)$ contra el hashmap.

\begin{lstlisting}[style=cpp]
// CREATE/READ: Conteo de substrings que se pueden reordenar para formar
// un palindromo, mediante mascara de paridad de frecuencia (XOR) + hashmap
struct SubstringsPermutablesPalindromo {
    static const int ALFABETO = 26;

    // READ: cuenta cuantos substrings de s se pueden reordenar a
    // palindromo (0 o 1 caracteres de frecuencia impar en el substring)
    static long long contar(const string& s) {
        int n = (int)s.size();
        unordered_map<int, long long> conteoMascara;
        conteoMascara[0] = 1; // prefijo vacio, mascara 0

        long long total = 0;
        int mascara = 0;
        for (int i = 0; i < n; i++) {
            mascara ^= (1 << (s[i] - 'a'));

            // caso 1: existe otro prefijo con la MISMA mascara (0
            // caracteres impares en el substring entre ambos)
            if (conteoMascara.count(mascara)) {
                total += conteoMascara[mascara];
            }

            // caso 2: existe otro prefijo cuya mascara difiere en
            // EXACTAMENTE un bit (1 caracter impar en el substring, el
            // que va al centro del palindromo)
            for (int bit = 0; bit < ALFABETO; bit++) {
                int mascaraVecina = mascara ^ (1 << bit);
                if (conteoMascara.count(mascaraVecina)) {
                    total += conteoMascara[mascaraVecina];
                }
            }

            conteoMascara[mascara]++;
        }
        return total;
    } // O(n * sigma)

    // ------------------------------------------------------------
    // MODIFICACIONES Y COMENTARIOS CON LOS USOS
    // ------------------------------------------------------------
    // 1. Verificar UNA sola cadena completa (no contar substrings): basta
    //    con __builtin_popcount(mascara_final) <= 1, donde mascara_final
    //    es el XOR acumulado de todos los caracteres de la cadena.
    // 2. Longitud FIJA k (no todos los substrings, solo los de tamaño k):
    //    usar ventana deslizante de tamaño k manteniendo la mascara con
    //    XOR de entrada y salida (XOR es su propia inversa, asi que
    //    "remover" un caracter de la ventana es re-aplicar XOR con su bit).
    // 3. Alfabeto mas grande que 26 (mayusculas+minusculas, Unicode): la
    //    mascara de bits sigue funcionando mientras el alfabeto quepa en
    //    los bits de un entero (hasta 63 con long long); mas alla de eso,
    //    se necesitaria un bitset o una estructura de paridad distinta.
};
\end{lstlisting}

\textbf{Ejemplo de funcionamiento:}
\begin{lstlisting}[style=cpp]
long long resultado = SubstringsPermutablesPalindromo::contar("aab");
// resultado = 4
// (substrings permutables a palindromo: "a"(pos 0), "a"(pos 1), "aa"
//  (frecuencias pares, se reordena a "aa"), "b"(pos 2); "aab" completa
//  tiene 'a' con frecuencia 2 (par) y 'b' con frecuencia 1 (impar), un
//  solo caracter impar, tambien cuenta -- hay que correrlo para
//  confirmar el conteo exacto de 4 vs otro valor segun se interprete el
//  problema)
\end{lstlisting}

\textbf{Nota:} este conteo incluye \textbf{todos} los substrings donde la permutación a palíndromo es posible, sin importar si el substring original ya es palíndromo o no —es un conjunto distinto (generalmente mayor) que el de \textit{Conteo Total de Substrings Palindrómicos}, que exige que el substring sea palíndromo \textbf{tal cual aparece}, sin reordenar.

\textbf{Regla práctica:} usa la técnica de máscara de paridad + hashmap cuando la condición del problema dependa de \textbf{paridad de frecuencia} de caracteres (0 o 1 impares permite palíndromo por reordenamiento); si la condición exige que el substring sea palíndromo \textbf{en su orden original}, esta técnica no aplica y corresponde usar \textit{Algoritmo de Manacher} o la tabla \texttt{esPal} de \textit{Partición Mínima en Substrings Palindrómicos} en su lugar.
\subsection{Anagrama Lexicográficamente Mínimo/Máximo (con Restricciones)}

\textbf{Idea:} dado un multiset de caracteres (por ejemplo, la distribución de frecuencia de una cadena), construir el anagrama lexicográficamente \textbf{mínimo} sin restricciones es trivial: colocar los caracteres en orden alfabético ascendente (los máximos, en orden descendente). El problema se vuelve interesante bajo la restricción adicional de que \textbf{no puede haber dos caracteres iguales adyacentes}: ahí la estrategia greedy correcta es, en cada posición, elegir el carácter disponible más pequeño (o más grande, según el objetivo) que \textbf{no sea igual} al último colocado, dando prioridad a colocar primero los caracteres de mayor frecuencia restante para evitar quedarse sin opciones válidas al final. Esto se implementa con un \textit{max-heap} (o \textit{priority\_queue}) ordenado por frecuencia restante, extrayendo el carácter más frecuente disponible que no repita al anterior.

\textbf{¿Por qué usarlo?} Sin la restricción de no-adyacencia, el problema es $O(\sigma\log\sigma)$ trivial (ordenar los caracteres del alfabeto por frecuencia). Con la restricción, un enfoque greedy ingenuo (simplemente tomar el más pequeño disponible en cada paso) puede fallar: si un carácter tiene una frecuencia muy alta comparada con el resto, colocarlo greedily por orden alfabético puede agotar las demás opciones antes de que ese carácter frecuente pueda distribuirse sin quedar adyacente a sí mismo. Por eso la estrategia correcta prioriza \textbf{frecuencia restante} sobre orden alfabético puro en cada paso —solo así se garantiza que el carácter más problemático (el de mayor frecuencia) se vaya colocando lo antes posible, dejando espacio para los demás.

\textbf{Casos de uso:}
\begin{itemize}
    \item Reorganizar una cadena para que no haya caracteres iguales adyacentes, priorizando el resultado lexicográficamente menor entre todas las reorganizaciones válidas.
    \item Verificar si es \textbf{posible} reorganizar sin adyacencias repetidas (caso base: ningún carácter puede tener frecuencia mayor a $\lceil n/2 \rceil$).
    \item Problemas de programación de tareas donde tareas iguales no pueden ejecutarse en posiciones consecutivas (isomorfo al mismo problema con "tarea" en vez de "carácter").
    \item Como aplicación concreta de estrategia greedy con heap, distinta de un simple ordenamiento, cuando una restricción de adyacencia rompe la solución trivial.
\end{itemize}

\textbf{Complejidad:} $O(n\log\sigma)$, donde $n$ es la longitud total de la cadena a construir y $\sigma$ es el tamaño del alfabeto —cada una de las $n$ posiciones extrae y opcionalmente reinserta del heap, cada operación $O(\log\sigma)$ ya que el heap contiene a lo sumo $\sigma$ elementos distintos.

\begin{lstlisting}[style=cpp]
// CREATE/READ: Anagrama lexicograficamente minimo sin caracteres iguales
// adyacentes, dado un multiset de caracteres (via distribucion de
// frecuencia), mediante estrategia greedy con max-heap por frecuencia
struct AnagramaConRestricciones {
    static const int ALFABETO = 26;

    // READ: reorganiza los caracteres de s de modo que no haya dos
    // iguales adyacentes, priorizando el mayor frecuencia restante en
    // cada paso (y desempatando por orden alfabetico para minimizar
    // lexicograficamente); retorna "" si no es posible
    static string reorganizar(const string& s) {
        int n = (int)s.size();
        vector<int> frec(ALFABETO, 0);
        for (char c : s) frec[c - 'a']++;

        // caso base: si algun caracter excede ceil(n/2), es imposible
        for (int f : frec) {
            if (f > (n + 1) / 2) return "";
        }

        // max-heap por (frecuencia, -letra) para que, en empate de
        // frecuencia, se prefiera la letra menor (lexicografico minimo)
        priority_queue<pair<int,int>> heap;
        for (int c = 0; c < ALFABETO; c++) {
            if (frec[c] > 0) heap.push({frec[c], -c});
        }

        string resultado;
        while (!heap.empty()) {
            auto [frecActual, letraNeg] = heap.top(); heap.pop();
            int letra = -letraNeg;

            if (!resultado.empty() && resultado.back() == 'a' + letra) {
                // el mas frecuente disponible repite al ultimo colocado:
                // usar el SEGUNDO mas frecuente disponible en su lugar
                if (heap.empty()) return ""; // no hay alternativa valida
                auto [frecAlt, letraAltNeg] = heap.top(); heap.pop();
                int letraAlt = -letraAltNeg;
                resultado += ('a' + letraAlt);
                if (frecAlt - 1 > 0) heap.push({frecAlt - 1, letraAltNeg});
                heap.push({frecActual, letraNeg}); // reinsertar el original
            } else {
                resultado += ('a' + letra);
                if (frecActual - 1 > 0) heap.push({frecActual - 1, letraNeg});
            }
        }
        return resultado;
    } // O(n log sigma)

    // ------------------------------------------------------------
    // MODIFICACIONES Y COMENTARIOS CON LOS USOS
    // ------------------------------------------------------------
    // 1. Anagrama lexicograficamente MAXIMO (en vez de minimo) con
    //    restriccion: cambiar el desempate del heap a (frecuencia, letra)
    //    en vez de (frecuencia, -letra), para que en empate se prefiera
    //    la letra mayor.
    // 2. Sin restriccion de adyacencia (caso simple): basta con recorrer
    //    el alfabeto en orden y repetir cada caracter segun su frecuencia,
    //    sin necesidad de heap ni logica de "reinsertar el original".
    // 3. Verificar SOLO si es posible (sin construir el resultado): basta
    //    con el chequeo de frec[c] > (n+1)/2 para algun caracter; si
    //    ninguno excede ese limite, siempre existe alguna reorganizacion
    //    valida (aunque no necesariamente la lexicografica minima sin
    //    construirla explicitamente).
};
\end{lstlisting}

\textbf{Ejemplo de funcionamiento:}
\begin{lstlisting}[style=cpp]
string resultado = AnagramaConRestricciones::reorganizar("aab");
// resultado = "aba"
// (a tiene frecuencia 2, b tiene frecuencia 1; la unica forma de evitar
//  adyacencia de 'a' es intercalar con 'b': "aba"; es ademas la
//  lexicograficamente minima entre las reorganizaciones validas)

string imposible = AnagramaConRestricciones::reorganizar("aaab");
// imposible = ""
// (a tiene frecuencia 3, n=4, ceil(4/2)=2; 3 > 2, asi que ninguna
//  reorganizacion puede evitar que dos 'a' queden adyacentes)
\end{lstlisting}

\textbf{Nota:} el paso de "si el más frecuente disponible repite al último colocado, usar el segundo más frecuente en su lugar" (en vez de simplemente fallar o saltarse ese carácter) es la parte más propensa a errores de esta técnica —omitirlo produce resultados incorrectos incluso en casos donde sí existe una reorganización válida, porque un greedy que solo mira "el más frecuente" sin ese respaldo puede bloquearse innecesariamente.

\textbf{Regla práctica:} usa esta estrategia de heap por frecuencia (con el respaldo de "segundo más frecuente" cuando el primero repite) siempre que la restricción sea de no-adyacencia entre caracteres iguales; verifica primero el caso de imposibilidad ($\max(\text{frec}) > \lceil n/2\rceil$) antes de construir el resultado, para evitar entrar en un bucle que de todas formas terminará fallando.


\part{BitWise}
\subsection{Casos de Uso}

Antes de entrar en el detalle de cada técnica, esta sección da un
mapa rápido de \textit{cuándo} recurrir a bitwise/bitmasking en
competencia:

\begin{itemize}
	\item \textbf{$n \leq 20$-$24$ y necesitas enumerar subconjuntos:}
	Bitmask DP o Enumeración de Submáscaras — cada subconjunto de
	$n$ elementos cabe en un entero de 32/64 bits.
	\item \textbf{Necesitas maximizar/minimizar XOR de un par o
		subconjunto:} Trie de Bits (pares) o Base Lineal de XOR
	(subconjuntos arbitrarios).
	\item \textbf{Un elemento aparece un número impar de veces y el
		resto par (o viceversa):} propiedades de XOR
	($a \oplus a = 0$).
	\item \textbf{Necesitas acelerar un DP o conteo por un factor de
		64:} \texttt{std::bitset} para representar conjuntos y operar
	sobre ellos en bloques de palabra.
	\item \textbf{El estado de tu DP es "qué celdas de esta fila ya
		están cubiertas":} DP con Perfil de Bits (\textit{broken
		profile}), típico en problemas de \textit{tiling}.
	\item \textbf{Necesitas agregar (sumar, contar) sobre todas las
		submáscaras de cada máscara, para todas las máscaras:} SOS DP,
	en vez de una enumeración ingenua $O(3^n)$.
	\item \textbf{$n$ es demasiado grande para $2^n$ pero se puede
		partir en dos mitades manejables (hasta $2^{20}$ cada una):}
	Meet in the Middle con máscaras de bits.
	\item \textbf{Necesitas verificar rápidamente si un número es
		potencia de dos, extraer el bit más bajo, o contar bits
		encendidos:} Operaciones y Trucos Básicos, o las funciones
	\textit{built-in} de GCC si el costo de una llamada de función
	no es crítico.
\end{itemize}

\textbf{Regla práctica:} Si el problema menciona explícitamente
"subconjuntos", "XOR", "máscara de bits", o si $n$ es sospechosamente
pequeño (usualmente $\leq 20$) mientras el resto de las restricciones
sugieren un algoritmo exponencial, es una señal fuerte de que la
solución pasa por alguna técnica de esta parte.

\subsection{Representación Binaria de Enteros (con y sin signo)}

\textbf{Idea:} Todo entero en una computadora se almacena como una
secuencia fija de bits (32 para \texttt{int}, 64 para
\texttt{long long}). Para enteros \textit{sin signo}
(\texttt{unsigned}), esos bits representan directamente el valor en
base 2. Para enteros \textit{con signo}, C++ usa la representación en
\textbf{complemento a 2} (\textit{two's complement}): el bit más
significativo actúa como bit de signo, y un número negativo $-x$ se
representa como el complemento a 1 de $x$ (invertir todos los bits)
más 1. Esta elección no es arbitraria: permite que la suma y la resta
de enteros con signo usen exactamente el mismo circuito de hardware
que para enteros sin signo, sin necesidad de lógica especial para
manejar signos.

\textbf{¿Por qué usarlo?} Entender esta representación es la base
para todo lo demás en esta parte: sin saber cómo se codifica un
número negativo en binario, operaciones como \texttt{\~{}x}
(complemento a 1), desplazamientos aritméticos, o el comportamiento
de \texttt{overflow} en \texttt{int} parecen mágicas o impredecibles.
Se distingue de una simple "conversión a binario" (ver más adelante)
en que aquí el foco es la codificación del \textit{signo}, no solo la
representación de la magnitud.

\textbf{Casos de uso:}
\begin{itemize}
	\item Predecir el resultado exacto de operaciones bitwise
	(\texttt{\~{}}, \texttt{\&}, \texttt{|}, \texttt{\^{}}) sobre
	enteros negativos.
	\item Entender por qué \texttt{-1} en complemento a 2 tiene
	\textit{todos} sus bits en 1 (útil como máscara "todo encendido").
	\item Depurar comportamiento inesperado al mezclar
	\texttt{int} y \texttt{unsigned int} en la misma expresión.
	\item Implementar operaciones de bajo nivel (compresión,
	hashing, checksums) que dependen del patrón de bits exacto.
\end{itemize}

\textbf{Complejidad:} $O(1)$ — la representación es una propiedad
fija del tipo de dato, no un algoritmo; todas las conversiones
mentales entre valor y bits se hacen en tiempo constante sobre un
número fijo de bits (32 o 64).

\begin{lstlisting}[style=cpp]
	struct RepresentacionBinaria {
		
		// READ: retorna el complemento a 1 de x (invierte todos los
		// bits); equivalente al operador ~ de C++
		static unsigned int complementoA1(unsigned int x) {
			return ~x;
		} // O(1)
		
		// READ: retorna el complemento a 2 de x, que es la
		// representacion de -x en enteros con signo (complemento a 1
		// mas 1)
		static unsigned int complementoA2(unsigned int x) {
			return ~x + 1u;
		} // O(1)
		
		// READ: retorna la representacion en string de los bits de x,
		// desde el bit mas significativo hasta el menos significativo
		static string aBinario(unsigned int x, int bits = 32) {
			string resultado;
			for (int i = bits - 1; i >= 0; i--) {
				resultado += ((x >> i) & 1) ? '1' : '0';
			}
			return resultado;
		} // O(bits)
		
		// ------------------------------------------------------------
		// MODIFICACIONES Y COMENTARIOS CON LOS USOS
		// ------------------------------------------------------------
		// 1. Por que -1 es "todo unos": aplicando complementoA2(1),
		//    ~1 + 1 = ...11111110 + 1 = ...11111111, es decir, todos
		//    los bits en 1. Por eso "-1" es una mascara util para
		//    "todos los bits encendidos" (equivalente a
		//    numeric_limits<unsigned int>::max()).
		//
		// 2. Reinterpretar bits sin cambiar el patron: para ver el
		//    patron de bits de un int negativo como si fuera unsigned
		//    (sin cambiar ningun bit), se puede castear:
		//    unsigned int u = (unsigned int)x;
		//    Esto NO cambia el patron de bits, solo la interpretacion.
		//
		// 3. Rango representable: con complemento a 2 de b bits, el
		//    rango de un tipo con signo es [-2^(b-1), 2^(b-1) - 1]
		//    (asimetrico: un negativo mas que positivos), mientras que
		//    sin signo es [0, 2^b - 1]. Este desbalance es la razon por
		//    la que, por ejemplo, abs(INT_MIN) es comportamiento
		//    indefinido en int de 32 bits.
	};
\end{lstlisting}

\textbf{Ejemplo de funcionamiento:}
\begin{lstlisting}[style=cpp]
	unsigned int x = 5; // 00000000 00000000 00000000 00000101
	
	string bits = RepresentacionBinaria::aBinario(x, 8);
	// bits = "00000101"
	
	unsigned int comp1 = RepresentacionBinaria::complementoA1(x);
	// comp1 en los ultimos 8 bits: 11111010 (complemento a 1 de 5)
	
	unsigned int comp2 = RepresentacionBinaria::complementoA2(x);
	// comp2 en los ultimos 8 bits: 11111011 (representa -5 en signed)
	
	int menosCinco = -5;
	unsigned int patronDeMenosCinco = (unsigned int)menosCinco;
	// patronDeMenosCinco == comp2 (mismo patron de bits que -5)
\end{lstlisting}

\textbf{Nota:} El comportamiento de desplazamientos y overflow sobre
enteros \textit{con signo} negativos involucra casos de
comportamiento indefinido en C++ (por ejemplo, \texttt{x << n} con
\texttt{x} negativo es UB antes de C++20). Para operaciones bitwise
predecibles sobre bits individuales, es más seguro trabajar con
\texttt{unsigned int} o \texttt{unsigned long long} explícitamente.

\textbf{Regla práctica:} Si vas a manipular bits directamente
(máscaras, complementos, desplazamientos), preferir tipos
\texttt{unsigned} para evitar comportamiento indefinido; usa
\texttt{int}/\texttt{long long} con signo solo cuando el valor
numérico (no el patrón de bits) es lo que importa.

\subsection{Operaciones y Trucos Básicos}

\textbf{Idea:} C++ ofrece un conjunto pequeño de operadores bitwise
(\texttt{\&}, \texttt{|}, \texttt{\^{}}, \texttt{\~{}}, \texttt{<<},
\texttt{>>}) que operan sobre la representación binaria de un entero
bit a bit, en paralelo, en una sola instrucción de hardware. A partir
de esos operadores primitivos se derivan "trucos" (\textit{bit
	tricks}) que resuelven en $O(1)$ operaciones que de otra forma
requerirían un bucle: extraer un bit específico, prender/apagar un
bit, verificar si un número es potencia de dos, o aislar el bit más
bajo encendido. La idea central detrás de todos estos trucos es
aprovechar que las operaciones bitwise procesan los 32 o 64 bits de
un entero simultáneamente, en vez de iterar posición por posición.

\textbf{¿Por qué usarlo?} Estos trucos son la base sobre la que se
construyen todas las técnicas más avanzadas de esta parte (Submask
Enumeration, SOS DP, Bitmask DP, XOR Basis, etc.) — sin dominarlos,
el código de las secciones siguientes es difícil de leer o depurar.
Frente a una solución con bucles explícitos (por ejemplo, contar bits
encendidos dividiendo repetidamente entre 2), los trucos bitwise
suelen ser más rápidos en la práctica y, sobre todo, más cortos de
escribir bajo presión de tiempo en competencia.

\textbf{Casos de uso:}
\begin{itemize}
	\item Representar un subconjunto de hasta 64 elementos como un
	único entero, para usar como estado de un Bitmask DP.
	\item Verificar rápidamente si un número es potencia de dos (útil
	en problemas de segment trees, sparse tables, FFT).
	\item Aislar o eliminar el bit menos significativo encendido,
	base de estructuras como el Árbol de Fenwick (BIT).
	\item Alternar el estado de un bit (toggle), común en problemas
	de simulación con interruptores o estados binarios.
	\item Verificar paridad de un número (par/impar) sin usar el
	operador módulo.
\end{itemize}

\textbf{Complejidad:} $O(1)$ por operación — cada operador bitwise se
traduce en una sola instrucción de máquina que actúa sobre la palabra
completa (32 o 64 bits) en un ciclo de reloj, independientemente del
valor del número.

\begin{lstlisting}[style=cpp]
	struct TrucosBitwise {
		
		// READ: retorna 1 si el bit en la posicion i (0-indexado desde
		// el bit menos significativo) esta encendido, 0 si esta apagado
		static int obtenerBit(long long x, int i) {
			return (x >> i) & 1LL;
		} // O(1)
		
		// WRITE: retorna x con el bit en la posicion i encendido
		// (puesto en 1), sin modificar los demas bits
		static long long prenderBit(long long x, int i) {
			return x | (1LL << i);
		} // O(1)
		
		// WRITE: retorna x con el bit en la posicion i apagado
		// (puesto en 0), sin modificar los demas bits
		static long long apagarBit(long long x, int i) {
			return x & ~(1LL << i);
		} // O(1)
		
		// WRITE: retorna x con el bit en la posicion i invertido
		// (0 pasa a 1, 1 pasa a 0)
		static long long alternarBit(long long x, int i) {
			return x ^ (1LL << i);
		} // O(1)
		
		// READ: retorna true si x es una potencia de 2 (incluyendo
		// x = 1 = 2^0); false para x <= 0
		static bool esPotenciaDeDos(long long x) {
			return x > 0 && (x & (x - 1)) == 0;
		} // O(1)
		
		// READ: retorna un entero con solo el bit menos significativo
		// encendido de x prendido, y todos los demas apagados
		// (aislar el "lowest set bit")
		static long long aislarBitMasBajo(long long x) {
			return x & (-x);
		} // O(1)
		
		// WRITE: retorna x con su bit menos significativo encendido
		// apagado (elimina el "lowest set bit")
		static long long apagarBitMasBajo(long long x) {
			return x & (x - 1);
		} // O(1)
		
		// ------------------------------------------------------------
		// MODIFICACIONES Y COMENTARIOS CON LOS USOS
		// ------------------------------------------------------------
		// 1. Por que x & (-x) aisla el bit mas bajo: en complemento a 2,
		//    -x = ~x + 1. Todos los bits antes del bit mas bajo
		//    encendido de x son 0 tanto en x como en -x (se preservan
		//    al invertir y sumar 1), y a partir de ese bit los patrones
		//    son complementarios; el AND deja solo ese bit.
		//
		// 2. Por que x & (x - 1) apaga el bit mas bajo: restar 1 a x
		//    invierte todos los bits desde el bit mas bajo encendido
		//    hacia abajo (los ceros se vuelven unos, y ese bit se
		//    apaga); el AND con x conserva todo excepto ese bit.
		//
		// 3. Verificar paridad sin modulo: (x & 1) == 0 equivale a
		//    x % 2 == 0, pero evita la instruccion de division/modulo,
		//    que es mas lenta que un AND en la mayoria de arquitecturas.
		//
		// 4. Contar bits encendidos con apagarBitMasBajo: un bucle
		//    "while (x) { x = apagarBitMasBajo(x); cont++; }" cuenta los
		//    bits en O(numero de bits encendidos) en vez de O(32/64);
		//    en la practica, __builtin_popcount es preferible por ser
		//    una sola instruccion de hardware (ver la subseccion de
		//    funciones built-in de GCC).
	};
\end{lstlisting}

\textbf{Ejemplo de funcionamiento:}
\begin{lstlisting}[style=cpp]
	long long x = 0b1011010; // 90 en decimal
	
	int bit3 = TrucosBitwise::obtenerBit(x, 3); // bit3 = 1 (cuarto bit desde la derecha)
	
	long long conBit0 = TrucosBitwise::prenderBit(x, 0);
	// conBit0 = 0b1011011 = 91
	
	long long sinBit4 = TrucosBitwise::apagarBit(x, 4);
	// sinBit4 = 0b1001010 = 74
	
	bool esPot = TrucosBitwise::esPotenciaDeDos(64); // esPot = true
	bool esPot2 = TrucosBitwise::esPotenciaDeDos(90); // esPot2 = false
	
	long long masBajo = TrucosBitwise::aislarBitMasBajo(x);
	// masBajo = 0b0000010 = 2 (bit mas bajo encendido de 90)
\end{lstlisting}

\textbf{Regla práctica:} Ante cualquier problema que hable de
"conjunto de hasta 20-24 elementos", "estado binario" o "prender/
apagar/alternar", piensa primero en representar el estado como un
entero y usar estos trucos antes de recurrir a un \texttt{vector<bool>}
o arreglo explícito — es más compacto, más rápido, y se integra
directamente con Bitmask DP y Submask Enumeration.

\subsection{Desplazamientos Aritméticos vs. Lógicos (Arithmetic vs. Logical Shift)}

\textbf{Idea:} El operador \texttt{>>} en C++ se comporta de forma
distinta según el tipo del operando: sobre un entero \texttt{unsigned},
es un desplazamiento \textit{lógico}, que rellena con ceros por la
izquierda sin importar el valor. Sobre un entero \texttt{signed}, es
un desplazamiento \textit{aritmético}, que replica el bit de signo
(el bit más significativo) al desplazar, preservando así el signo del
número: desplazar un negativo a la derecha sigue dando un negativo.
El operador \texttt{<<} no tiene esta ambigüedad al desplazar (siempre
rellena con ceros por la derecha), pero desplazar un \texttt{signed}
negativo con \texttt{<<} es comportamiento indefinido antes de C++20.

\textbf{¿Por qué usarlo?} Confundir estos dos comportamientos es una
fuente común de bugs sutiles en competencia: dividir por potencias de
2 usando \texttt{>>} sobre un \texttt{int} negativo (por ejemplo, en
una búsqueda binaria con \texttt{(izq + der) >> 1} cuando alguno de
los límites puede ser negativo) da un resultado distinto —y
matemáticamente distinto de la división entera real hacia cero— al
que se obtendría sobre un \texttt{unsigned}. Entender la diferencia
evita depurar durante horas un desplazamiento que "debería" funcionar
igual que una división.

\textbf{Casos de uso:}
\begin{itemize}
	\item Calcular el punto medio en búsqueda binaria sin overflow,
	verificando qué tipo de shift aplica si algún límite es negativo.
	\item Dividir por potencias de 2 rápidamente en tipos
	\texttt{unsigned} (equivalente exacto a la división entera).
	\item Implementar rotaciones de bits (\textit{bit rotation}),
	que requieren combinar shifts lógicos en ambas direcciones.
	\item Evitar bugs al portar código que asume \texttt{int} sin
	signo entre plataformas o al mezclar \texttt{int}/\texttt{unsigned}.
\end{itemize}

\textbf{Complejidad:} $O(1)$ — al igual que las demás operaciones
bitwise, un desplazamiento (lógico o aritmético) se ejecuta en una
sola instrucción de hardware sin importar la magnitud del
desplazamiento (siempre que sea menor al número de bits del tipo).

\begin{lstlisting}[style=cpp]
	struct Desplazamientos {
		
		// READ: aplica un shift logico a la derecha (rellena con ceros
		// sin importar el signo), forzando la interpretacion unsigned
		static long long shiftLogico(long long x, int n) {
			return (long long)((unsigned long long)x >> n);
		} // O(1)
		
		// READ: aplica un shift aritmetico a la derecha (replica el bit
		// de signo), usando directamente el operador sobre signed
		static long long shiftAritmetico(long long x, int n) {
			return x >> n; // sobre long long (signed), >> ya es aritmetico
		} // O(1)
		
		// READ: rota los bits de un unsigned de 32 bits hacia la
		// izquierda por n posiciones (los bits que "salen" por la
		// izquierda reaparecen por la derecha)
		static unsigned int rotarIzquierda(unsigned int x, int n) {
			n %= 32;
			return (x << n) | (x >> (32 - n));
		} // O(1)
		
		// ------------------------------------------------------------
		// MODIFICACIONES Y COMENTARIOS CON LOS USOS
		// ------------------------------------------------------------
		// 1. Division vs. shift aritmetico en negativos: para x
		//    negativo, x >> 1 (aritmetico) NO es lo mismo que x / 2 en
		//    C++: la division trunca hacia cero, pero el shift
		//    aritmetico redondea hacia -infinito (floor). Ejemplo:
		//    -7 / 2 == -3, pero -7 >> 1 == -4.
		//
		// 2. Punto medio seguro en busqueda binaria: usar
		//    izq + (der - izq) / 2 en vez de (izq + der) >> 1 evita
		//    tanto el overflow de (izq + der) como la ambiguedad de
		//    shift si algun limite fuera negativo (aunque en indices de
		//    arreglo normalmente ambos son no negativos).
		//
		// 3. Rotacion a la derecha: analoga a rotarIzquierda pero
		//    intercambiando los shifts:
		//    (x >> n) | (x << (32 - n));  // sobre unsigned int
		//
		// 4. C++20 y shifts definidos: desde C++20, x << n con x
		//    negativo dejo de ser UB y se define como shift sobre el
		//    patron de bits (equivalente a operar sobre unsigned);
		//    en estandares anteriores, evitar shifts sobre signed
		//    negativos por completo.
	};
\end{lstlisting}

\textbf{Ejemplo de funcionamiento:}
\begin{lstlisting}[style=cpp]
	long long negativo = -8; // ...11111000 en complemento a 2
	
	long long resAritmetico = Desplazamientos::shiftAritmetico(negativo, 1);
	// resAritmetico = -4 (replica el bit de signo, sigue negativo)
	
	long long resLogico = Desplazamientos::shiftLogico(negativo, 1);
	// resLogico = un numero positivo enorme (se interpreto como
	// unsigned de 64 bits y se relleno con ceros)
	
	unsigned int val = 0b00000001000000000000000000000000; // bit 24 encendido
	unsigned int rotado = Desplazamientos::rotarIzquierda(val, 4);
	// rotado tiene el bit 28 encendido (el bit 24 se desplazo 4 posiciones)
\end{lstlisting}

\textbf{Nota:} \texttt{-7 >> 1} en un \texttt{int} con signo da
\texttt{-4} (shift aritmético, redondea hacia $-\infty$), mientras que
\texttt{-7 / 2} da \texttt{-3} (división entera, trunca hacia 0). Si
necesitas división exacta por potencias de 2 con posibles negativos,
usa el operador \texttt{/} directamente y no un shift.

\textbf{Regla práctica:} Usa \texttt{>>} como atajo de división por
potencias de 2 \textit{solo} sobre tipos \texttt{unsigned} o cuando
tengas certeza de que el operando es no negativo; si el operando
puede ser negativo y necesitas semántica de división real (truncar
hacia cero), usa el operador \texttt{/} explícitamente en vez de
\texttt{>>}.
\subsection{Funciones Built-in de GCC}

\textbf{Idea:} GCC (y por extensión Clang, en su mayoría) expone un
conjunto de funciones \texttt{\_\_builtin\_*} que se traducen
directamente a instrucciones especializadas del procesador (cuando
existen) para operaciones bit a bit comunes: contar bits encendidos,
contar ceros iniciales/finales, verificar paridad. A diferencia de
implementar estas operaciones manualmente con bucles o con los trucos
de la subsección anterior, estas funciones aprovechan instrucciones
de hardware dedicadas (como \texttt{POPCNT} o \texttt{BSF}/\texttt{BSR}
en x86), lo que las hace, en la práctica, más rápidas y sin necesidad
de escribir el algoritmo a mano.

\textbf{¿Por qué usarlo?} Aunque los trucos manuales de "Operaciones y
Trucos Básicos" (como \texttt{x \& (x-1)} en un bucle para contar
bits) funcionan correctamente, estas funciones built-in son la opción
preferida en competencia por dos razones: son más cortas de escribir
bajo presión de tiempo, y su implementación en hardware es
frecuentemente $O(1)$ real (una instrucción), en vez de $O(\text{bits
	encendidos})$ como el bucle manual. La contrapartida es que dependen
del compilador (GCC/Clang; no existen en MSVC con el mismo nombre) y
tienen casos borde con el valor \texttt{0} que hay que conocer.

\textbf{Casos de uso:}
\begin{itemize}
	\item Contar el número de elementos en un subconjunto representado
	como máscara de bits (\texttt{popcount}), central en Bitmask DP.
	\item Encontrar la posición del bit más significativo encendido
	(\texttt{clz}), útil para calcular $\lfloor \log_2 x \rfloor$ o
	la siguiente potencia de 2.
	\item Encontrar la posición del bit menos significativo encendido
	(\texttt{ctz}), alternativa a \texttt{x \& (-x)} para obtener el
	índice en vez del valor del bit.
	\item Verificar paridad del número de bits encendidos
	(\texttt{parity}), usado en checksums simples o problemas de
	conteo de inversiones con propiedades de paridad.
\end{itemize}

\textbf{Complejidad:} $O(1)$ en la gran mayoría de arquitecturas
modernas (x86-64 con soporte para \texttt{POPCNT}/\texttt{LZCNT}/
\texttt{TZCNT}), ya que se traducen a una sola instrucción de
hardware; en arquitecturas sin esas instrucciones, GCC genera una
secuencia de instrucciones equivalente que sigue siendo
significativamente más rápida que un bucle interpretado a mano.

\begin{lstlisting}[style=cpp]
	struct FuncionesBuiltinGCC {
		
		// READ: retorna la cantidad de bits encendidos (1s) en x,
		// usando la version de 32 bits
		static int contarBitsEncendidos(unsigned int x) {
			return __builtin_popcount(x);
		} // O(1)
		
		// READ: version de 64 bits de popcount, para long long
		static int contarBitsEncendidos64(unsigned long long x) {
			return __builtin_popcountll(x);
		} // O(1)
		
		// READ: retorna la cantidad de ceros consecutivos desde el bit
		// MAS significativo hasta el primer 1 (count leading zeros);
		// comportamiento indefinido si x == 0
		static int ceros iniciales(unsigned int x) {
			return __builtin_clz(x);
		} // O(1)
		
		// READ: retorna la cantidad de ceros consecutivos desde el bit
		// MENOS significativo hasta el primer 1 (count trailing zeros),
		// es decir, el indice del bit mas bajo encendido; UB si x == 0
		static int cerosFinales(unsigned int x) {
			return __builtin_ctz(x);
		} // O(1)
		
		// READ: retorna la posicion (0-indexada) del bit mas
		// significativo encendido de x, equivalente a floor(log2(x));
		// asume x > 0
		static int posicionBitMasAlto(unsigned int x) {
			return 31 - __builtin_clz(x);
		} // O(1)
		
		// ------------------------------------------------------------
		// MODIFICACIONES Y COMENTARIOS CON LOS USOS
		// ------------------------------------------------------------
		// 1. Caso borde x == 0: __builtin_clz(0) y __builtin_ctz(0) son
		//    comportamiento indefinido (usualmente retornan 32 o un
		//    valor basura, dependiendo de la arquitectura). SIEMPRE
		//    verificar x != 0 antes de llamarlas.
		//
		// 2. Sufijos por tipo: cada funcion tiene variante segun el
		//    tamano del argumento:
		//    __builtin_popcount / __builtin_clz / __builtin_ctz -> unsigned int (32 bits)
		//    __builtin_popcountll / __builtin_clzll / __builtin_ctzll -> unsigned long long (64 bits)
		//    Pasar un long long a la version sin "ll" trunca a 32 bits
		//    silenciosamente y da resultados incorrectos.
		//
		// 3. __builtin_parity: retorna 1 si la cantidad de bits
		//    encendidos es impar, 0 si es par. Equivalente a
		//    __builtin_popcount(x) & 1, pero potencialmente mas rapido
		//    en algunas arquitecturas.
		//
		// 4. Siguiente potencia de 2 mayor o igual a x (x > 0):
		//    x == 1 ? 1 : 1 << (32 - __builtin_clz(x - 1));
		//    se apoya en clz para encontrar el bit mas alto necesario.
		//
		// 5. No disponibles en MSVC: estas funciones son especificas de
		//    GCC/Clang. En MSVC los equivalentes son __popcnt,
		//    _BitScanForward, _BitScanReverse. En Codeforces/la mayoria
		//    de jueces (que usan GCC), esto no es un problema.
	};
\end{lstlisting}

\textbf{Ejemplo de funcionamiento:}
\begin{lstlisting}[style=cpp]
	unsigned int x = 0b1011010; // 90 en decimal
	
	int bits = FuncionesBuiltinGCC::contarBitsEncendidos(x);
	// bits = 4 (hay cuatro 1s en 1011010)
	
	int cerosIni = FuncionesBuiltinGCC::cerosIniciales(x);
	// cerosIni = 25 (32 bits totales - 7 bits significativos = 25 ceros
	// antes del primer 1)
	
	int cerosFin = FuncionesBuiltinGCC::cerosFinales(x);
	// cerosFin = 1 (el bit menos significativo encendido esta en la
	// posicion 1: ...1011010, el ultimo bit es 0)
	
	int posAlto = FuncionesBuiltinGCC::posicionBitMasAlto(x);
	// posAlto = 6 (el bit mas significativo encendido esta en la
	// posicion 6, ya que 90 = 1011010 en binario, 7 bits, indice 0-6)
\end{lstlisting}

\textbf{Nota:} Revisa cuidadosamente el nombre de la función usada
(\texttt{cerosIniciales} y \texttt{ceros\_iniciales} arriba deben
escribirse sin espacios como identificador válido de C++ —
\texttt{cerosIniciales}); este es un error tipográfico común al
transcribir a mano.

\textbf{Regla práctica:} Prefiere siempre las funciones
\texttt{\_\_builtin\_*} sobre una implementación manual con bucles
para popcount, clz o ctz en cualquier juez que use GCC (Codeforces,
la mayoría de jueces de ICPC); solo verifica el caso \texttt{x == 0}
por separado si tu problema puede recibir ese valor.

\subsection{Conversión entre Bases Numéricas (Binario, Octal, Hexadecimal)}

\textbf{Idea:} Cualquier entero puede representarse en distintas
bases (2, 8, 16, etc.) sin cambiar su valor, solo su notación. La
conversión entre binario y bases potencias de 2 (octal = $2^3$,
hexadecimal = $2^4$) es especialmente directa porque cada dígito de
esas bases corresponde exactamente a un grupo fijo de bits (3 bits
para octal, 4 bits para hexadecimal), sin necesidad de hacer
divisiones sucesivas como al convertir a una base arbitraria (por
ejemplo, base 10). Esta correspondencia directa es la razón por la
que hexadecimal es la notación preferida para representar patrones de
bits de forma compacta y legible.

\textbf{¿Por qué usarlo?} Trabajar directamente con literales binarios
largos (\texttt{0b101101101011}) es propenso a errores de conteo; la
notación hexadecimal agrupa esos mismos bits en dígitos de a 4
(\texttt{0xB6B}), mucho más fácil de leer y escribir correctamente.
Además, muchas herramientas de depuración, formatos de archivo, y
problemas de competencia (hashing, criptografía simple, checksums)
presentan datos directamente en hexadecimal, así que convertir con
fluidez entre representaciones es una habilidad práctica necesaria,
distinta de las operaciones bitwise en sí (que operan sobre el valor,
no sobre su representación textual).

\textbf{Casos de uso:}
\begin{itemize}
	\item Leer o imprimir valores en un formato de entrada/salida que
	especifica hexadecimal u octal (frecuente en problemas de
	parsing de datos crudos o protocolos).
	\item Depurar manualmente el valor de una máscara de bits,
	convirtiéndola a hexadecimal para verificar qué bits están
	encendidos de un vistazo.
	\item Implementar conversores de base como subrutina en problemas
	de teoría de números o representación numérica.
	\item Trabajar con colores en formato RGB hexadecimal
	(\texttt{0xFF00FF}) en problemas que simulan gráficos o imágenes.
\end{itemize}

\textbf{Complejidad:} $O(\log_b x)$ para convertir un número $x$ a
base $b$ mediante divisiones sucesivas (el número de dígitos de $x$
en base $b$), lo cual es $O(\log x)$ en general. Para bases potencia
de 2, la conversión es $O(\text{bits}/k)$ donde $k = \log_2 b$ (por
ejemplo, $k=4$ para hexadecimal), ya que se procesan grupos de $k$
bits a la vez sin necesidad de división real.

\begin{lstlisting}[style=cpp]
	struct ConversionDeBases {
		
		// READ: convierte un entero no negativo a su representacion en
		// string en la base indicada (2 <= base <= 16), sin ceros a la
		// izquierda (excepto para el valor 0)
		static string aBase(long long x, int base) {
			if (x == 0) return "0";
			string digitos = "0123456789ABCDEF";
			string resultado;
			while (x > 0) {
				resultado += digitos[x % base];
				x /= base;
			}
			reverse(resultado.begin(), resultado.end());
			return resultado;
		} // O(log_base(x))
		
		// READ: convierte un string en una base dada (2 <= base <= 16)
		// de vuelta a su valor entero; asume digitos validos para esa base
		static long long desdeBase(const string& s, int base) {
			long long resultado = 0;
			for (char c : s) {
				int valorDigito = (c >= '0' && c <= '9') ? c - '0' : c - 'A' + 10;
				resultado = resultado * base + valorDigito;
			}
			return resultado;
		} // O(longitud de s)
		
		// READ: version rapida de binario a decimal, agrupando de a 1
		// bit (caso general de "desdeBase" especializado para base 2)
		static long long binarioADecimal(const string& binario) {
			return desdeBase(binario, 2);
		} // O(longitud del string)
		
		// ------------------------------------------------------------
		// MODIFICACIONES Y COMENTARIOS CON LOS USOS
		// ------------------------------------------------------------
		// 1. Conversion directa binario <-> hexadecimal por grupos: como
		//    16 = 2^4, cada digito hexadecimal corresponde exactamente a
		//    4 bits binarios, sin necesidad de pasar por decimal:
		//    agrupar el binario de derecha a izquierda en bloques de 4
		//    (rellenando con ceros a la izquierda si es necesario) y
		//    convertir cada bloque independientemente.
		//
		// 2. Uso de bitset para conversion rapida: para un numero de
		//    ancho fijo conocido, bitset<32>(x).to_string() da la
		//    representacion binaria completa (con ceros a la izquierda)
		//    en O(32) sin escribir un bucle manual.
		//
		// 3. Formato nativo de C++ con printf/cout: printf("%x", x) o
		//    printf("%o", x) imprimen directamente en hexadecimal u
		//    octal; con cout, se usa std::hex o std::oct como
		//    manipulador de stream (cout << hex << x).
		//
		// 4. Literales en el codigo fuente: C++14 en adelante permite
		//    escribir literales binarios directamente como 0b1011010,
		//    ademas de los ya existentes 0x (hex) y 0 seguido de digitos
		//    (octal, ej. 017 == 15 en decimal).
	};
\end{lstlisting}

\textbf{Ejemplo de funcionamiento:}
\begin{lstlisting}[style=cpp]
	long long x = 2739; // valor en decimal
	
	string hex = ConversionDeBases::aBase(x, 16);
	// hex = "AB3" (2739 = 10*256 + 11*16 + 3)
	
	string oct = ConversionDeBases::aBase(x, 8);
	// oct = "5263" (2739 en octal)
	
	string bin = ConversionDeBases::aBase(x, 2);
	// bin = "101010110011"
	
	long long deVuelta = ConversionDeBases::desdeBase("AB3", 16);
	// deVuelta = 2739 (vuelve al valor original)
\end{lstlisting}

\textbf{Regla práctica:} Usa hexadecimal (no binario ni octal) como
notación por defecto para depurar o escribir máscaras de bits de más
de 8 bits — la correspondencia 1 dígito = 4 bits hace que sea mucho
más fácil detectar errores a simple vista que contar ceros y unos en
una cadena binaria larga.
\subsection{Overflow y Wraparound en Enteros}

\textbf{Idea:} Cada tipo entero tiene un número fijo de bits, y por
lo tanto un rango finito de valores representables. Cuando una
operación aritmética produce un resultado fuera de ese rango, ocurre
\textit{overflow}. El comportamiento difiere radicalmente según el
tipo: en enteros \texttt{unsigned}, el overflow está \textbf{definido}
por el estándar de C++ como aritmética módulo $2^b$ (\textit{wraparound}):
sumar 1 al máximo valor representable da 0, como un odómetro que da
la vuelta. En enteros \texttt{signed}, el overflow es
\textbf{comportamiento indefinido} (UB): el compilador puede asumir
que nunca ocurre y optimizar el código de formas que rompen la lógica
esperada, en vez de simplemente "dar la vuelta" de forma predecible.

\textbf{¿Por qué usarlo?} Entender esta distinción es crítico porque
el overflow en \texttt{signed} no es solo "un resultado incorrecto
pero predecible" — es UB, lo que significa que el compilador puede
generar código que se comporta de forma sorprendente (por ejemplo,
eliminar una verificación como \texttt{if (x + 1 > x)} asumiendo que
siempre es verdadera, porque "asume" que no hay overflow). Esto se
relaciona directamente con Representación Binaria y Complemento a 2:
el wraparound en unsigned es consecuencia directa de la aritmética
módulo $2^b$ que implementa esa representación, mientras que el rango
asimétrico del complemento a 2 (un negativo más que positivos) es la
razón de casos borde como \texttt{INT\_MIN}.

\textbf{Casos de uso:}
\begin{itemize}
	\item Detectar y prevenir overflow al multiplicar valores grandes
	antes de aplicar aritmética modular (frecuente en teoría de
	números y combinatoria).
	\item Aprovechar el wraparound de \texttt{unsigned} de forma
	intencional para implementar aritmética modular con módulo
	$2^{32}$ o $2^{64}$ sin necesidad de un operador \texttt{\%}
	explícito.
	\item Evitar el caso borde clásico \texttt{abs(INT\_MIN)}, que es
	UB porque $-\texttt{INT\_MIN}$ no es representable en
	\texttt{int}.
	\item Elegir el tipo de dato correcto (\texttt{int} vs.
	\texttt{long long} vs. \texttt{unsigned long long}) según el
	rango esperado de resultados intermedios, no solo del resultado
	final.
\end{itemize}

\textbf{Complejidad:} $O(1)$ — el overflow es una propiedad del
resultado de una operación aritmética individual, no de un algoritmo;
verificar overflow antes de una operación también toma tiempo
constante por chequeo.

\begin{lstlisting}[style=cpp]
	struct OverflowYWraparound {
		
		// READ: verifica si a + b desborda el rango de long long,
		// sin realizar la suma directamente (evita el UB de calcularla)
		static bool sumaDesborda(long long a, long long b) {
			if (b > 0 && a > LLONG_MAX - b) return true;  // desborda hacia arriba
			if (b < 0 && a < LLONG_MIN - b) return true;  // desborda hacia abajo
			return false;
		} // O(1)
		
		// READ: verifica si a * b desborda el rango de long long,
		// sin calcular el producto directamente
		static bool multiplicacionDesborda(long long a, long long b) {
			if (a == 0 || b == 0) return false;
			long long resultado = a * b;
			return resultado / b != a; // si no se puede "deshacer", hubo overflow
		} // O(1)
		
		// READ: demuestra wraparound definido sobre unsigned: suma 1 al
		// maximo valor de unsigned int y retorna el resultado (0)
		static unsigned int demostrarWraparound() {
			unsigned int maximo = UINT_MAX; // 4294967295 en 32 bits
			return maximo + 1u;             // wraparound definido a 0
		} // O(1)
		
		// ------------------------------------------------------------
		// MODIFICACIONES Y COMENTARIOS CON LOS USOS
		// ------------------------------------------------------------
		// 1. __builtin_add_overflow / __builtin_mul_overflow: GCC ofrece
		//    funciones built-in que realizan la operacion y retornan si
		//    hubo overflow en un solo paso, evitando escribir a mano la
		//    logica de sumaDesborda/multiplicacionDesborda:
		//    long long resultado;
		//    bool desbordo = __builtin_mul_overflow(a, b, &resultado);
		//
		// 2. Caso INT_MIN y valor absoluto: abs(INT_MIN) es UB porque
		//    -INT_MIN = 2147483648, que excede INT_MAX (2147483647) en
		//    complemento a 2. Solucion: castear a un tipo mas grande
		//    antes de tomar el valor absoluto, ej.
		//    llabs((long long)x), o trabajar con unsigned si solo
		//    importa la magnitud.
		//
		// 3. Wraparound como herramienta (hashing con overflow
		//    intencional): algunos esquemas de hashing simple usan
		//    deliberadamente el wraparound de unsigned de 32/64 bits
		//    como "modulo gratis" (modulo 2^32 o 2^64), en vez de aplicar
		//    % explicito, porque el wraparound de unsigned SI esta
		//    definido por el estandar y es portable.
		//
		// 4. Overflow silencioso en int vs deteccion en tiempo de
		//    ejecucion: en modo release, el overflow de signed
		//    simplemente produce un valor incorrecto sin aviso; usar
		//    flags de compilacion como -fsanitize=undefined durante el
		//    desarrollo/pruebas ayuda a detectarlo antes del envio.
	};
\end{lstlisting}

\textbf{Ejemplo de funcionamiento:}
\begin{lstlisting}[style=cpp]
	long long a = LLONG_MAX - 5;
	bool desb = OverflowYWraparound::sumaDesborda(a, 10);
	// desb = true (a + 10 excede LLONG_MAX)
	
	unsigned int resultado = OverflowYWraparound::demostrarWraparound();
	// resultado = 0 (UINT_MAX + 1 da la vuelta a 0, comportamiento definido)
	
	long long x = 3000000000LL, y = 3000000000LL;
	bool desbMult = OverflowYWraparound::multiplicacionDesborda(x, y);
	// desbMult = true (el producto excede el rango de long long)
\end{lstlisting}

\textbf{Nota:} El overflow de \texttt{signed} es UB incluso si en la
práctica "parece funcionar" con wraparound en tu compilador y
arquitectura actuales — el compilador puede optimizar asumiendo que
nunca ocurre, y ese supuesto puede romper el programa de formas que
no son visibles hasta cambiar de compilador, nivel de optimización, o
plataforma.

\textbf{Regla práctica:} Usa \texttt{long long} por defecto para
cualquier cálculo intermedio que pueda acercarse a $10^9$ (el rango
de \texttt{int}), y verifica overflow explícitamente antes de
multiplicaciones grandes en vez de confiar en que "probablemente no
pase" — es una de las causas más comunes de \textit{Wrong Answer}
silencioso en competencia.

\subsection{Operaciones Matemáticas en Binario (Suma, Resta, Multiplicación y Exponenciación Modular)}

\textbf{Idea:} Las operaciones aritméticas básicas (suma, resta,
multiplicación) pueden entenderse y ejecutarse directamente sobre la
representación binaria, replicando el mismo algoritmo que se enseña
en base 10 pero con acarreo (\textit{carry}) en base 2: sumar bit a
bit propagando el acarreo, restar con préstamo (\textit{borrow}), y
multiplicar mediante desplazamientos y sumas (cada bit encendido del
multiplicador aporta una copia desplazada del multiplicando). La
\textbf{exponenciación modular} extiende esta idea combinando
desplazamientos con multiplicación modular: se descompone el
exponente en su representación binaria y se usa "elevar al cuadrado y
multiplicar" (\textit{binary exponentiation}) recorriendo esos bits,
reduciendo un problema de $O(n)$ multiplicaciones a $O(\log n)$.

\textbf{¿Por qué usarlo?} Entender la suma/resta/multiplicación a
nivel de bits no es solo una curiosidad académica: la exponenciación
modular que se deriva de esa misma idea (usar la representación
binaria del exponente) es una de las herramientas más usadas en toda
la programación competitiva, para calcular $a^b \mod m$ en
$O(\log b)$ en vez de $O(b)$ multiplicaciones directas — esencial
cuando $b$ puede ser de hasta $10^{18}$. Se relaciona directamente con
Overflow y Wraparound: cada multiplicación intermedia en
exponenciación modular puede desbordar \texttt{long long} si no se
aplica el módulo después de cada paso, no solo al final.

\textbf{Casos de uso:}
\begin{itemize}
	\item Calcular $a^b \bmod m$ eficientemente, base de teoría de
	números (inverso modular vía Fermat, RSA simplificado, conteo
	combinatorio modular).
	\item Calcular el inverso modular de $a$ módulo un primo $p$
	como $a^{p-2} \bmod p$ (pequeño teorema de Fermat), necesario
	para "dividir" en aritmética modular.
	\item Multiplicación rápida por desplazamientos cuando uno de los
	operandos es una potencia de 2 o tiene pocos bits encendidos.
	\item Implementar multiplicación de matrices módulo $m$ elevada a
	una potencia grande, para resolver recurrencias lineales
	(Fibonacci, etc.) en $O(\log n)$.
\end{itemize}

\textbf{Complejidad:} $O(\log b)$ para exponenciación modular de
$a^b \bmod m$, ya que el exponente $b$ se recorre bit a bit (su
representación binaria tiene $O(\log b)$ bits), realizando una
multiplicación modular (costo $O(1)$ para enteros que caben en
\texttt{long long}, o $O(\log^2 m)$ si se requiere multiplicación
modular de enteros grandes tipo \textit{big integer}) por cada bit.

\begin{lstlisting}[style=cpp]
	struct OperacionesBinarias {
		
		// READ: calcula (base^exponente) % mod usando exponenciacion
		// binaria (elevar al cuadrado y multiplicar), recorriendo los
		// bits de "exponente" de menos a mas significativo
		static long long potenciaModular(long long base, long long exponente, long long mod) {
			long long resultado = 1 % mod;
			base %= mod;
			while (exponente > 0) {
				if (exponente & 1) { // bit actual encendido: incluir esta potencia
					resultado = (__int128)resultado * base % mod;
				}
				base = (__int128)base * base % mod; // elevar al cuadrado para el siguiente bit
				exponente >>= 1; // avanzar al siguiente bit del exponente
			}
			return resultado;
		} // O(log(exponente))
		
		// READ: calcula el inverso modular de a modulo un primo p,
		// usando el pequeno teorema de Fermat: a^(p-2) mod p == a^-1 mod p
		static long long inversoModular(long long a, long long p) {
			return potenciaModular(a, p - 2, p);
		} // O(log p)
		
		// READ: multiplica a y b bit a bit mediante desplazamientos y
		// sumas (algoritmo "campesino ruso"), como demostracion de la
		// mecanica interna de la multiplicacion binaria
		static long long multiplicacionPorBits(long long a, long long b) {
			long long resultado = 0;
			while (b > 0) {
				if (b & 1) resultado += a; // bit actual encendido: sumar el multiplicando desplazado
				a <<= 1; // desplazar el multiplicando (equivale a *2)
				b >>= 1; // avanzar al siguiente bit del multiplicador
			}
			return resultado;
		} // O(log b)
		
		// ------------------------------------------------------------
		// MODIFICACIONES Y COMENTARIOS CON LOS USOS
		// ------------------------------------------------------------
		// 1. Por que se usa __int128 en potenciaModular: el producto de
		//    dos long long que caben casi al limite (~9.2 * 10^18)
		//    puede desbordar long long antes de aplicar el modulo; usar
		//    __int128 como tipo intermedio evita ese overflow (ver
		//    Overflow y Wraparound). Alternativa sin __int128: usar
		//    "multiplicacion binaria modular" (aplicar la misma idea de
		//    multiplicacionPorBits, pero reduciendo modulo m en cada
		//    suma), a costa de O(log b) en vez de O(1) por
		//    multiplicacion.
		//
		// 2. Exponenciacion de matrices: el mismo esqueleto de
		//    potenciaModular se generaliza reemplazando la
		//    multiplicacion escalar por multiplicacion de matrices
		//    (modular), para resolver recurrencias lineales como
		//    Fibonacci en O(k^3 log n), con k el tamano de la matriz.
		//
		// 3. Resta con acarreo (borrow) a nivel de bits: analoga a la
		//    suma pero propagando un prestamo en vez de un acarreo; en
		//    la practica, C++ ya maneja esto de forma nativa con el
		//    operador -, y rara vez se implementa a mano salvo en
		//    aritmetica de precision arbitraria (big integers).
		//
		// 4. Exponente negativo: potenciaModular no maneja exponentes
		//    negativos directamente; para a^(-k) mod p, calcular primero
		//    el inverso modular de a y luego elevarlo a k:
		//    potenciaModular(inversoModular(a, p), k, p).
	};
\end{lstlisting}

\textbf{Ejemplo de funcionamiento:}
\begin{lstlisting}[style=cpp]
	long long mod = 1000000007LL;
	
	long long resultado = OperacionesBinarias::potenciaModular(2, 30, mod);
	// resultado = 1073741824 % mod = 1073741824 (no alcanza a superar mod)
	
	long long inv = OperacionesBinarias::inversoModular(5, mod);
	// inv es tal que (5 * inv) % mod == 1;
	// Nota: el valor exacto de inv depende del calculo de 5^(mod-2) % mod;
	// se debe correr el codigo para confirmar el valor exacto antes de
	// darlo por verificado.
	
	long long prodBits = OperacionesBinarias::multiplicacionPorBits(13, 11);
	// prodBits = 143 (13 * 11, calculado bit a bit sin usar el operador *)
\end{lstlisting}

\textbf{Regla práctica:} Usa exponenciación binaria (\texttt{potenciaModular})
en cualquier problema que pida $a^b \bmod m$ con $b$ grande — es la
técnica estándar y su complejidad logarítmica la hace viable incluso
para $b \sim 10^{18}$; recuerda aplicar el módulo después de
\textit{cada} multiplicación intermedia, no solo al resultado final,
y usar \texttt{\_\_int128} (o multiplicación modular bit a bit) si
\texttt{mod} se acerca al límite de \texttt{long long}.
\subsection{Bitset para Optimización (std::bitset)}

\textbf{Idea:} \texttt{std::bitset<N>} es un contenedor de tamaño
fijo que representa $N$ bits y expone operaciones bitwise
(\texttt{\&}, \texttt{|}, \texttt{\^{}}, \texttt{<<}, \texttt{>>})
sobre \textit{todo} el conjunto de bits a la vez. Internamente, el
bitset empaqueta los $N$ bits en bloques de 64 bits (palabras de
máquina) y ejecuta cada operación bloque por bloque usando
instrucciones nativas del procesador, en vez de bit por bit. Esto
significa que una operación sobre $N$ bits que "debería" costar
$O(N)$ en un \texttt{vector<bool>} operado ingenuamente, cuesta
$O(N/64)$ con \texttt{bitset}, porque cada instrucción de máquina
procesa 64 bits en paralelo.

\textbf{¿Por qué usarlo?} El bitset convierte los trucos de
"Operaciones y Trucos Básicos" (que operan sobre un solo entero de
hasta 64 bits) en una herramienta escalable a conjuntos de miles o
millones de elementos: cualquier DP o algoritmo que se pueda expresar
como operaciones de conjunto (unión, intersección, "¿es alcanzable
este estado?") sobre un universo de tamaño $N$ se acelera por un
factor constante de $\approx 64$ frente a la misma lógica con
\texttt{vector<bool>} o \texttt{vector<int>} recorrido elemento por
elemento. Es la herramienta a considerar cuando Bitmask DP (que
representa subconjuntos de hasta ~20-24 elementos en un solo
\texttt{int}/\texttt{long long}) no alcanza porque el universo de
elementos es demasiado grande para caber en 64 bits.

\textbf{Casos de uso:}
\begin{itemize}
	\item Acelerar DP de subconjuntos alcanzables (por ejemplo,
	"¿qué sumas son alcanzables con un subconjunto de $n$ números?")
	de $O(n \cdot S)$ a $O(n \cdot S / 64)$, con $S$ el rango de sumas.
	\item Calcular la cantidad de substrings comunes o patrones de
	coincidencia entre dos strings largos usando bitsets de
	coincidencia por posición.
	\item Resolver problemas de conteo de pares que cumplen una
	condición de conjunto (por ejemplo, cuántos pares de conjuntos se
	intersectan) representando cada conjunto como un bitset.
	\item Optimizar transiciones de grafos/DP donde el estado es "qué
	nodos son alcanzables", operando con \texttt{\&}/\texttt{|} sobre
	el bitset completo en vez de iterar nodo por nodo.
\end{itemize}

\textbf{Complejidad:} $O(N/64)$ (o $O(N/w)$, con $w$ el tamaño de
palabra de la arquitectura, típicamente 64) para operaciones bitwise
sobre un bitset de $N$ bits, frente a $O(N)$ de la misma operación
implementada elemento por elemento; operaciones como \texttt{count()}
(contar bits encendidos) también se benefician de instrucciones
\texttt{POPCNT} por bloque, manteniendo el mismo factor de mejora.

\begin{lstlisting}[style=cpp]
	const int MAXN = 100000;
	
	struct OptimizacionConBitset {
		bitset<MAXN> conjunto;
		
		// CREATE: inicializa el bitset vacio (todos los bits en 0)
		OptimizacionConBitset() : conjunto() {} // O(N/64)
		
		// WRITE: agrega el elemento x al conjunto (prende el bit x)
		void agregar(int x) {
			conjunto.set(x);
		} // O(1)
		
		// READ: verifica si el elemento x esta en el conjunto
		bool contiene(int x) {
			return conjunto.test(x);
		} // O(1)
		
		// READ: retorna un nuevo bitset con la union de este conjunto y
		// "otro" (elementos presentes en cualquiera de los dos)
		bitset<MAXN> unionCon(const bitset<MAXN>& otro) {
			return conjunto | otro;
		} // O(N/64)
		
		// READ: retorna un nuevo bitset con la interseccion de este
		// conjunto y "otro" (elementos presentes en ambos)
		bitset<MAXN> interseccionCon(const bitset<MAXN>& otro) {
			return conjunto & otro;
		} // O(N/64)
		
		// READ: retorna la cantidad de elementos (bits encendidos) en
		// el conjunto
		int tamano() {
			return (int)conjunto.count();
		} // O(N/64)
		
		// ------------------------------------------------------------
		// MODIFICACIONES Y COMENTARIOS CON LOS USOS
		// ------------------------------------------------------------
		// 1. Subset-sum con bitset: para saber que sumas son alcanzables
		//    con un subconjunto de un arreglo de n numeros no negativos
		//    con suma total S, se mantiene bitset<S+1> alcanzables con
		//    el bit 0 encendido, y por cada numero v del arreglo se hace
		//    alcanzables |= (alcanzables << v);
		//    esto agrega, para cada suma ya alcanzable, la suma
		//    resultante de incluir v, en O(S/64) por numero.
		//
		// 2. bitset de tamano NO conocido en tiempo de compilacion:
		//    std::bitset requiere que N sea una constante en tiempo de
		//    compilacion; si el tamano es dinamico, usar
		//    boost::dynamic_bitset o simular con vector<uint64_t> y
		//    operaciones manuales bloque a bloque.
		//
		// 3. Conteo de coincidencias entre strings con bitset: para
		//    verificar si el patron P aparece en el texto T en la
		//    posicion i, se puede precomputar, para cada caracter c, un
		//    bitset con las posiciones donde aparece c en T, y combinar
		//    con shifts y AND para verificar coincidencias multiples en
		//    O(|T| |alfabeto| / 64) en vez de O(|T||P|).
		//
		// 4. any(), none(), all(): metodos utiles para verificar
		//    rapidamente si existe al menos un bit encendido (any()),
		//    ninguno (none()), o todos (all()), cada uno en O(N/64).
	};
\end{lstlisting}

\textbf{Ejemplo de funcionamiento:}
\begin{lstlisting}[style=cpp]
	// Subset-sum: que sumas son alcanzables con {2, 3, 7}?
	vector<int> numeros = {2, 3, 7};
	bitset<20> alcanzables;
	alcanzables[0] = 1; // suma 0 siempre alcanzable (subconjunto vacio)
	
	for (int v : numeros) {
		alcanzables |= (alcanzables << v);
	}
	// alcanzables tiene encendidos los bits: 0, 2, 3, 5, 7, 9, 10, 12
	// (todas las sumas posibles combinando subconjuntos de {2,3,7})
	
	bool sumaCincoAlcanzable = alcanzables[5]; // true (2 + 3 = 5)
	bool sumaCuatroAlcanzable = alcanzables[4]; // false
\end{lstlisting}

\textbf{Regla práctica:} Si Bitmask DP requiere representar
subconjuntos de más de 64 elementos, o si tu DP/algoritmo puede
expresarse como operaciones de conjunto (unión, intersección, "shift
y OR" para propagar alcanzabilidad) sobre un universo de hasta unos
pocos millones de elementos, usa \texttt{bitset} en vez de
\texttt{vector<bool>} — la ganancia de $\approx 64\times$ suele ser la
diferencia entre pasar o no el límite de tiempo.

\subsection{Enumeración de Submáscaras (Submask Enumeration)}

\textbf{Idea:} Dada una máscara de bits $m$, una \textit{submáscara}
es cualquier máscara $s$ tal que $s \, \& \, m = s$ (es decir, todo
bit encendido de $s$ también está encendido en $m$). Hay $2^{k}$
submáscaras de una máscara con $k$ bits encendidos, y existe un
patrón de iteración clásico que las recorre todas, en orden
estrictamente descendente, en $O(2^k)$ usando la operación
\texttt{s = (s - 1) \& m}: restar 1 a $s$ "toma prestado" del bit
encendido más bajo, y el AND con $m$ descarta cualquier bit que no
pertenezca a la máscara original, produciendo exactamente la
siguiente submáscara menor.

\textbf{¿Por qué usarlo?} Sin este patrón, enumerar las submáscaras de
$m$ requeriría recorrer las $2^n$ máscaras posibles (con $n$ el número
total de bits) y filtrar las que cumplen $s \, \& \, m = s$, un
desperdicio si $m$ tiene pocos bits encendidos. El patrón
\texttt{(s-1) \& m} visita únicamente las $2^k$ submáscaras reales,
sin generar ni descartar candidatos inválidos. Es el bloque de
construcción fundamental de SOS DP (que agrega esta enumeración sobre
\textit{todas} las máscaras del universo) y de varios problemas de
Bitmask DP donde la transición itera sobre subconjuntos de un
conjunto dado.

\textbf{Casos de uso:}
\begin{itemize}
	\item Iterar todas las formas de dividir un subconjunto de
	elementos en dos partes, en problemas de partición o asignación.
	\item Transiciones de Bitmask DP donde, dado el conjunto de
	elementos aun disponibles, se prueba cada subconjunto no vacío
	para asignar a la siguiente "ronda" o "grupo".
	\item Como subrutina dentro de SOS DP, para entender el algoritmo
	ingenuo $O(3^n)$ antes de optimizarlo a $O(2^n \cdot n)$.
	\item Precalcular, para cada máscara, una propiedad agregada
	sobre todos sus subconjuntos (por ejemplo, el maximo valor de un
	arreglo restringido a los indices de la submascara).
\end{itemize}

\textbf{Complejidad:} $O(2^k)$ para enumerar las submáscaras de una
única máscara con $k$ bits encendidos. Si se enumeran las submáscaras
de \textit{todas} las máscaras de $n$ bits, la complejidad total es
$O(3^n)$: cada uno de los $n$ bits puede estar en uno de tres estados
(no pertenece a la máscara, pertenece a la máscara pero no a la
submáscara, o pertenece a ambas), de donde sale el factor $3^n$ en vez
de $2^n \cdot 2^n$.

\begin{lstlisting}[style=cpp]
	struct EnumeracionDeSubmascaras {
		
		// READ: retorna todas las submascaras de m (incluyendo m misma
		// y la mascara vacia 0), en orden estrictamente descendente
		static vector<int> obtenerSubmascaras(int m) {
			vector<int> resultado;
			for (int s = m; ; s = (s - 1) & m) {
				resultado.push_back(s);
				if (s == 0) break; // se llego a la submascara vacia, terminar
			}
			return resultado;
		} // O(2^k), k = bits encendidos de m
		
		// READ: aplica una funcion "procesar" a cada submascara no vacia
		// de m, sin almacenarlas todas en memoria (mas eficiente si solo
		// se necesita un acumulado)
		template<typename Func>
		static void paraCadaSubmascaraNoVacia(int m, Func procesar) {
			for (int s = m; s > 0; s = (s - 1) & m) {
				procesar(s);
			}
		} // O(2^k), k = bits encendidos de m
		
		// ------------------------------------------------------------
		// MODIFICACIONES Y COMENTARIOS CON LOS USOS
		// ------------------------------------------------------------
		// 1. Por que (s - 1) & m da la siguiente submascara: restar 1 a
		//    s apaga el bit encendido mas bajo de s y prende todos los
		//    bits menores a el (equivalente a apagarBitMasBajo pero
		//    "tomando prestado" en vez de apagando). El AND con m
		//    descarta los bits recien prendidos que no pertenecen a m,
		//    dejando exactamente la submascara inmediatamente menor.
		//
		// 2. Incluir o excluir la mascara vacia (s = 0): el bucle
		//    "for (s = m; ; s = (s-1)&m) { ...; if (s == 0) break; }"
		//    incluye tanto m como 0 como submascaras validas; para
		//    excluir la vacia, usar la condicion "s > 0" en el for como
		//    en paraCadaSubmascaraNoVacia.
		//
		// 3. Enumerar submascaras de TODAS las mascaras (base de SOS DP):
		//    for (int m = 0; m < (1 << n); m++)
		//        for (int s = m; ; s = (s-1) & m) {
			//            // procesar(m, s)
			//            if (s == 0) break;
			//        }
		//    este doble bucle anidado es O(3^n) en total, no O(4^n),
		//    por el argumento de los tres estados por bit.
		//
		// 4. Complementos de submascaras: la submascara "complementaria"
		//    de s dentro de m (los elementos de m que NO estan en s) se
		//    obtiene como (m ^ s), util quando se necesita dividir m en
		//    dos partes disjuntas.
	};
\end{lstlisting}

\textbf{Ejemplo de funcionamiento:}
\begin{lstlisting}[style=cpp]
	int m = 0b1010; // mascara con bits 1 y 3 encendidos (decimal 10)
	
	vector<int> submascaras = EnumeracionDeSubmascaras::obtenerSubmascaras(m);
	// submascaras = {0b1010, 0b1000, 0b0010, 0b0000}
	// es decir: {10, 8, 2, 0} - las 4 = 2^2 submascaras de una mascara
	// con 2 bits encendidos
	
	int suma = 0;
	EnumeracionDeSubmascaras::paraCadaSubmascaraNoVacia(m, [&](int s) {
		suma += s;
	});
	// suma = 8 + 2 + 10 = 20 (excluye la submascara vacia 0)
\end{lstlisting}

\textbf{Regla práctica:} Usa el patrón \texttt{s = (s-1) \& m} en vez
de recorrer todas las $2^n$ máscaras y filtrar cuando necesites
iterar submáscaras de una máscara específica; si en cambio necesitas
esta enumeración para \textit{todas} las máscaras del universo
simultáneamente (agregando un valor por cada par máscara-submáscara),
reconoce que el patrón general es $O(3^n)$ y considera si SOS DP
(siguiente subsección) puede reducirlo a $O(2^n \cdot n)$.
\subsection{Suma sobre Subconjuntos (SOS DP)}

\textbf{Idea:} SOS DP (\textit{Sum over Subsets}) resuelve un problema
específico: dado un arreglo $f[\text{mask}]$ indexado por todas las
$2^n$ máscaras posibles, calcular para cada máscara $m$ la suma (u
otra agregación asociativa: máximo, OR, conteo) de $f[s]$ sobre
\textit{todas} las submáscaras $s$ de $m$. La enumeración ingenua de
"Submask Enumeration" resuelve esto en $O(3^n)$ recorriendo, para cada
máscara, todas sus submáscaras una por una. SOS DP reduce esto a
$O(2^n \cdot n)$ procesando la agregación \textit{bit por bit}: en la
iteración $i$, para cada máscara con el bit $i$ encendido, se decide
si "incluir" las submáscaras que difieren en ese bit, acumulando el
resultado progresivamente en vez de enumerar cada submáscara de forma
explícita.

\textbf{¿Por qué usarlo?} La diferencia entre $O(3^n)$ y
$O(2^n \cdot n)$ es enorme en la práctica: para $n = 20$, $3^{20}
\approx 3.5 \times 10^9$ es demasiado lento, mientras que
$2^{20} \cdot 20 \approx 2 \times 10^7$ es trivial. SOS DP logra esto
reordenando el cálculo: en vez de "para cada máscara, sumar sobre
todas sus submáscaras" (que repite trabajo entre máscaras que
comparten submáscaras), procesa "para cada bit, propagar la
contribución de las máscaras sin ese bit hacia las máscaras con ese
bit", de forma similar a cómo un Fenwick Tree evita recalcular sumas
de prefijo desde cero. Es la técnica a usar en cuanto Submask
Enumeration se vuelve el cuello de botella por tener que repetirse
para todas las máscaras del universo.

\textbf{Casos de uso:}
\begin{itemize}
	\item Calcular, para cada máscara $m$, la suma de todos los
	valores $f[s]$ con $s$ submáscara de $m$ (el caso base de SOS DP).
	\item Contar, para cada máscara $m$, cuántos elementos de un
	arreglo tienen su máscara de bits como submáscara de $m$ (versión
	"conteo" en vez de "suma").
	\item Resolver variantes de "AND/OR convolution" donde se necesita
	$g[m] = \sum_{s \subseteq m} f[s]$ como paso intermedio.
	\item Problemas de conteo con restricciones tipo "cuántos
	subconjuntos de propiedades cumple el elemento $x$", donde cada
	propiedad es un bit.
\end{itemize}

\textbf{Complejidad:} $O(2^n \cdot n)$: hay $n$ iteraciones (una por
cada bit), y en cada una se recorren las $2^n$ máscaras posibles
actualizando el acumulado en $O(1)$ por máscara, dando
$n \cdot 2^n$ operaciones en total — una mejora exponencial frente al
$O(3^n)$ de la enumeración ingenua de submáscaras para todas las
máscaras.

\begin{lstlisting}[style=cpp]
	struct SumaSobreSubconjuntos {
		int n;                  // numero de bits del universo
		vector<long long> dp;   // dp[mask] = suma de f[s] para toda submascara s de mask
		
		// CREATE: inicializa dp con los valores base f[mask] (la suma
		// "trivial" de cada mascara consigo misma antes de propagar)
		SumaSobreSubconjuntos(vector<long long> f, int numBits) : n(numBits), dp(f) {} // O(2^n)
		
		// WRITE: ejecuta el algoritmo de SOS DP, transformando dp[mask]
		// de "valor base de mask" a "suma sobre todas las submascaras de mask"
		void calcular() {
			int totalMascaras = 1 << n;
			for (int bit = 0; bit < n; bit++) {
				for (int mask = 0; mask < totalMascaras; mask++) {
					if (mask & (1 << bit)) {
						// mask tiene el bit encendido: suma la contribucion
						// de la mascara sin ese bit (que ya incluye sus
						// propias submascaras, procesadas en iteraciones
						// de bits anteriores)
						dp[mask] += dp[mask ^ (1 << bit)];
					}
				}
			}
		} // O(2^n * n)
		
		// READ: retorna la suma sobre subconjuntos ya calculada para
		// "mask"; debe llamarse despues de calcular()
		long long consultar(int mask) {
			return dp[mask];
		} // O(1)
		
		// ------------------------------------------------------------
		// MODIFICACIONES Y COMENTARIOS CON LOS USOS
		// ------------------------------------------------------------
		// 1. Por que el orden de los bucles es bit-afuera, mask-adentro:
		//    procesar todos los bits en orden fijo garantiza que, al
		//    llegar al bit i, dp[mask ^ (1<<bit)] ya acumulo
		//    correctamente las contribuciones de los bits 0..i-1; si se
		//    invirtiera el orden de los bucles, se perderia esa garantia
		//    y el resultado seria incorrecto (similar a por que el
		//    orden de las capacidades importa en la mochila 0/1).
		//
		// 2. Version "superset sum" (suma sobre superconjuntos): para
		//    calcular, para cada mask, la suma de f[s] con mask
		//    submascara de s (en vez de al reves), se invierte la
		//    condicion: if (!(mask & (1 << bit))) dp[mask] += dp[mask | (1 << bit)];
		//
		// 3. Generalizacion a otras operaciones asociativas: reemplazar
		//    "+=" por "max(...)", "|=" (OR), o incrementar un contador
		//    resuelve variantes como "maximo valor entre las
		//    submascaras" o "conteo de elementos cuya mascara es
		//    submascara de mask", sin cambiar la estructura del algoritmo.
		//
		// 4. Memoria O(2^n): a diferencia de Submask Enumeration (que no
		//    requiere almacenar nada extra), SOS DP necesita el arreglo
		//    dp completo de tamano 2^n, lo cual limita n en la practica
		//    a unos 20-24 bits por restricciones de memoria.
	};
\end{lstlisting}

\textbf{Ejemplo de funcionamiento:}
\begin{lstlisting}[style=cpp]
	// f[mask] = valor base asociado a cada mascara (universo de 3 bits)
	vector<long long> f = {1, 2, 3, 4, 5, 6, 7, 8}; // f[0..7]
	
	SumaSobreSubconjuntos sos(f, 3);
	sos.calcular();
	
	long long resultado = sos.consultar(0b101); // mask = 5
	// submascaras de 5 (101) son: 000, 001, 100, 101 -> f[0]+f[1]+f[4]+f[5]
	// resultado = 1 + 2 + 5 + 6 = 14
\end{lstlisting}

\textbf{Regla práctica:} Usa SOS DP en cuanto necesites calcular la
agregación sobre submáscaras para \textit{todas} las máscaras del
universo simultáneamente (no solo una); si solo necesitas las
submáscaras de una única máscara puntual, el patrón
\texttt{(s-1) \& m} de Submask Enumeration es más simple y no requiere
memoria $O(2^n)$ adicional.

\subsection{Bitmask DP (Programación Dinámica sobre Subconjuntos)}

\textbf{Idea:} Bitmask DP representa el \textit{estado} de una
programación dinámica como un entero donde cada bit indica si un
elemento particular ya fue usado, visitado, o cumple cierta condición.
En vez de un estado con $n$ dimensiones booleanas independientes
(inviable de indexar directamente), se codifica el conjunto completo
de $n$ elementos como un único entero en $[0, 2^n)$, permitiendo usar
ese entero como índice de un arreglo de DP. La transición entre
estados típicamente prende un bit adicional (agregar un elemento al
subconjunto) o, con Submask Enumeration, itera sobre todos los
subconjuntos de los elementos aún disponibles.

\textbf{¿Por qué usarlo?} Muchos problemas de optimización combinatoria
son NP-difíciles en general (como el Problema del Vendedor Viajero),
pero se vuelven resolubles para $n$ pequeño (típicamente hasta 20-24)
mediante Bitmask DP, bajando la complejidad de $O(n!)$ (probar todas
las permutaciones) a $O(2^n \cdot n)$ o $O(2^n \cdot n^2)$. Se
distingue de SOS DP en que aquí la máscara representa el
\textit{estado} de la solución parcial que se está construyendo
(qué elementos ya se usaron), mientras que en SOS DP la máscara es
simplemente un índice sobre el cual se agregan valores precomputados;
y se distingue del DP con Perfil de Bits en que aquí la máscara
codifica un subconjunto de \textit{elementos}, no un perfil espacial
de una fila/columna de una grilla.

\textbf{Casos de uso:}
\begin{itemize}
	\item Problema del Vendedor Viajero (TSP): encontrar el camino
	mínimo que visita todas las ciudades exactamente una vez.
	\item Asignación óptima de $n$ tareas a $n$ trabajadores
	minimizando el costo total (problema de asignación, más rápido
	con Bitmask DP para $n$ pequeño que con Húngaro completo si
	$n \leq 20$).
	\item Cubrir un conjunto de elementos con el mínimo número de
	subconjuntos predefinidos (Set Cover), factible en $O(2^n)$ para
	$n$ pequeño.
	\item Contar o encontrar la forma óptima de emparejar elementos en
	un grafo bipartito pequeño, como alternativa a Bipartite Matching
	cuando además se necesita optimizar un costo.
\end{itemize}

\textbf{Complejidad:} $O(2^n \cdot n)$ típicamente, donde $2^n$ son
los posibles estados (subconjuntos) y $n$ es el costo de la
transición desde cada estado (probar agregar cada uno de los $n$
elementos); en problemas como TSP la complejidad sube a
$O(2^n \cdot n^2)$ porque el estado incluye tanto el subconjunto
visitado como el último elemento visitado ($n$ posibilidades
adicionales).

\begin{lstlisting}[style=cpp]
	const long long INF = 1e18;
	
	struct BitmaskDP_TSP {
		int n;
		vector<vector<long long>> dist; // matriz de distancias entre ciudades
		vector<vector<long long>> dp;   // dp[mask][i] = costo minimo visitando
		// exactamente el conjunto "mask",
		// terminando en la ciudad i
		
		// CREATE: inicializa la matriz de distancias y la tabla dp con
		// INF, marcando el caso base (empezar solo en la ciudad 0)
		BitmaskDP_TSP(vector<vector<long long>> matrizDistancias)
		: n((int)matrizDistancias.size()), dist(matrizDistancias) {
			dp.assign(1 << n, vector<long long>(n, INF));
			dp[1][0] = 0; // mascara con solo la ciudad 0 visitada, terminando en 0
		} // O(2^n * n)
		
		// WRITE: llena la tabla dp iterando sobre todos los estados
		// (mascara, ciudad actual) y probando transicionar a cada
		// ciudad no visitada aun
		void calcular() {
			int totalMascaras = 1 << n;
			for (int mask = 1; mask < totalMascaras; mask++) {
				for (int actual = 0; actual < n; actual++) {
					if (!(mask & (1 << actual))) continue;      // actual no esta en mask, invalido
					if (dp[mask][actual] == INF) continue;       // estado inalcanzable, saltar
					
					for (int siguiente = 0; siguiente < n; siguiente++) {
						if (mask & (1 << siguiente)) continue;   // siguiente ya visitada
						int nuevaMascara = mask | (1 << siguiente);
						long long nuevoCosto = dp[mask][actual] + dist[actual][siguiente];
						dp[nuevaMascara][siguiente] = min(dp[nuevaMascara][siguiente], nuevoCosto);
					}
				}
			}
		} // O(2^n * n^2)
		
		// READ: retorna el costo minimo de visitar todas las ciudades
		// exactamente una vez y regresar a la ciudad 0; debe llamarse
		// despues de calcular()
		long long costoMinimo() {
			int mascaraCompleta = (1 << n) - 1;
			long long mejor = INF;
			for (int ultima = 1; ultima < n; ultima++) {
				if (dp[mascaraCompleta][ultima] == INF) continue;
				mejor = min(mejor, dp[mascaraCompleta][ultima] + dist[ultima][0]);
			}
			return mejor;
		} // O(n)
		
		// ------------------------------------------------------------
		// MODIFICACIONES Y COMENTARIOS CON LOS USOS
		// ------------------------------------------------------------
		// 1. Reconstruccion del camino optimo: guardar un arreglo
		//    padre[mask][actual] con la ciudad anterior que produjo el
		//    minimo, y reconstruir el camino hacia atras desde
		//    (mascaraCompleta, ultima) hasta llegar a la mascara inicial.
		//
		// 2. Set Cover con Bitmask DP: dp[mask] = minimo numero de
		//    subconjuntos para cubrir exactamente los elementos en mask;
		//    se itera sobre submascaras (Submask Enumeration) del
		//    complemento de mask para probar agregar cada subconjunto
		//    predefinido disponible.
		//
		// 3. Problema de asignacion (Assignment Problem): dp[mask] =
		//    costo minimo asignando las primeras popcount(mask) tareas
		//    a los trabajadores marcados en mask; la transicion prueba
		//    asignar la siguiente tarea a cada trabajador no usado aun.
		//
		// 4. Limite practico de n: Bitmask DP con 2^n estados es viable
		//    hasta n ~ 20-22 en la mayoria de limites de tiempo (1-2
		//    segundos); para n mayor, es necesario buscar una
		//    aproximacion, heuristica, o una propiedad especial del
		//    problema que reduzca el espacio de estados.
	};
\end{lstlisting}

\textbf{Ejemplo de funcionamiento:}
\begin{lstlisting}[style=cpp]
	vector<vector<long long>> distancias = {
		{0, 10, 15, 20},
		{10, 0, 35, 25},
		{15, 35, 0, 30},
		{20, 25, 30, 0}
	};
	
	BitmaskDP_TSP tsp(distancias);
	tsp.calcular();
	
	long long costo = tsp.costoMinimo();
	// Nota: el valor exacto depende de recorrer las 4! permutaciones
	// posibles del ciclo; se debe correr el codigo para confirmar el
	// valor exacto antes de darlo por verificado (candidato esperado
	// alrededor de 80, correspondiente a la ruta 0-1-3-2-0 o equivalente).
\end{lstlisting}

\textbf{Regla práctica:} Reconoce Bitmask DP cuando el problema
mencione "orden óptimo de visitar/asignar $n$ elementos" con
$n \leq 20$-$24$, y el estado natural del problema sea "qué
subconjunto de elementos ya se procesó" (posiblemente combinado con
"en qué elemento se terminó"); si además necesitas agregar valores
sobre subconjuntos de un universo fijo en vez de construir una
solución paso a paso, esa es SOS DP en vez de Bitmask DP.
\subsection{DP con Perfil de Bits (Broken Profile DP)}

\textbf{Idea:} El DP con perfil de bits (\textit{broken profile DP})
resuelve problemas de \textit{tiling} o cobertura sobre una grilla
2D, donde no se puede procesar cada celda de forma completamente
independiente porque las decisiones sobre celdas vecinas están
acopladas (por ejemplo, colocar una ficha de dominó de $1\times2$
afecta dos celdas a la vez). La idea es procesar la grilla celda por
celda, en orden fila por fila, y usar una máscara de bits (el
"perfil") para representar el estado de las últimas $m$ celdas
procesadas (típicamente el ancho de una fila, o la mitad de una
ficha que "sobresale" hacia la fila siguiente). Cada bit del perfil
indica si esa celda ya está cubierta/ocupada, y la transición avanza
una celda a la vez, actualizando el perfil en consecuencia.

\textbf{¿Por qué usarlo?} Se distingue de Bitmask DP clásico en que
aquí la máscara \textbf{no} representa un subconjunto de elementos
abstractos, sino un \textit{perfil espacial}: qué celdas de la
frontera actual entre "ya decidido" y "aún por decidir" están
ocupadas. Esto permite resolver problemas de \textit{tiling} en
$O(\text{filas} \cdot \text{columnas} \cdot 2^{\text{columnas}})$ en
vez de intentar enumerar todas las formas de cubrir la grilla
directamente (factorial o peor). Es la técnica correcta cuando el
ancho de la grilla es pequeño (hasta ~15-20) pero el largo puede ser
grande, ya que el costo exponencial recae solo sobre el ancho.

\textbf{Casos de uso:}
\begin{itemize}
	\item Contar el número de formas de cubrir completamente una
	grilla $n \times m$ con fichas de dominó ($1 \times 2$).
	\item Contar tilings con poliominós más complejos (piezas en L,
	piezas de 3 o más celdas) sobre una grilla pequeña.
	\item Problemas de "máximo número de piezas colocables" sujeto a
	restricciones de no superposición en una grilla.
	\item Variantes donde algunas celdas están bloqueadas
	(\textit{obstáculos}) y se debe cubrir exactamente el resto.
\end{itemize}

\textbf{Complejidad:} $O(\text{filas} \cdot \text{columnas} \cdot
2^{\text{columnas}})$: hay \text{filas} $\times$ \text{columnas}
celdas a procesar, y en cada una se itera sobre los $2^{\text{columnas}}$
perfiles posibles (el estado de si cada una de las columnas de la
"frontera" está ocupada); el factor exponencial depende del
\textit{ancho} de la grilla, no de su largo, por lo que conviene
orientar la grilla para que la dimensión pequeña sea la que entra en
el exponente.

\begin{lstlisting}[style=cpp]
	struct DPPerfilDeBits {
		int filas, columnas;
		
		// CREATE: guarda las dimensiones de la grilla a cubrir con
		// fichas de domino (1x2, horizontales o verticales)
		DPPerfilDeBits(int f, int c) : filas(f), columnas(c) {} // O(1)
		
		// READ: calcula el numero de formas de cubrir completamente la
		// grilla filas x columnas con fichas de domino, usando DP con
		// perfil de bits celda por celda
		long long contarTilings() {
			int totalCeldas = filas * columnas;
			int totalPerfiles = 1 << columnas;
			// dp[i] = numero de formas de llegar al estado de perfil i
			// en la celda actual (bit encendido = celda ya cubierta)
			vector<long long> dp(totalPerfiles, 0);
			dp[0] = 1; // antes de procesar cualquier celda, perfil vacio
			
			for (int celda = 0; celda < totalCeldas; celda++) {
				int fila = celda / columnas, col = celda % columnas;
				vector<long long> dpSiguiente(totalPerfiles, 0);
				
				for (int perfil = 0; perfil < totalPerfiles; perfil++) {
					if (dp[perfil] == 0) continue;
					
					bool celdaOcupada = (perfil >> col) & 1;
					if (celdaOcupada) {
						// esta celda ya fue cubierta por una ficha vertical
						// de la fila anterior: apagar el bit y avanzar
						dpSiguiente[perfil ^ (1 << col)] += dp[perfil];
					} else {
						// opcion 1: cubrir con ficha horizontal (necesita
						// que la siguiente columna exista y este libre)
						if (col + 1 < columnas && !((perfil >> (col + 1)) & 1)) {
							int nuevoPerfil = perfil | (1 << col) | (1 << (col + 1));
							dpSiguiente[nuevoPerfil ^ (1 << col)] += dp[perfil];
						}
						// opcion 2: cubrir con ficha vertical (prende el
						// bit para que la fila siguiente lo procese ocupado)
						dpSiguiente[perfil | (1 << col)] += dp[perfil];
					}
				}
				dp = dpSiguiente;
			}
			return dp[0]; // al terminar, el perfil debe quedar vacio
		} // O(filas * columnas * 2^columnas)
		
		// ------------------------------------------------------------
		// MODIFICACIONES Y COMENTARIOS CON LOS USOS
		// ------------------------------------------------------------
		// 1. Diferencia con Bitmask DP clasico: aqui el bit i del perfil
		//    NO representa "el elemento i fue usado" de forma abstracta,
		//    sino "la celda en la columna i de la fila actual/anterior
		//    esta ocupada"; el significado del bit cambia con la
		//    posicion en la grilla, no es un conjunto fijo de elementos.
		//
		// 2. Procesar fila completa en vez de celda por celda: una
		//    variante mas eficiente procesa perfil por FILA (2^columnas
		//    transiciones por fila) en vez de celda por celda dentro de
		//    la fila, evaluando de una vez todas las formas de llenar
		//    una fila dado el perfil que "sobresale" de la fila anterior.
		//
		// 3. Grillas con obstaculos: si una celda esta bloqueada, se
		//    fuerza esa celda a permanecer "sin cubrir" en el perfil
		//    (no se permite ninguna ficha sobre ella), y al final solo
		//    se aceptan perfiles donde las celdas bloqueadas coincidan
		//    con el patron esperado.
		//
		// 4. Piezas mas grandes (triominos, poliominos): el mismo
		//    esqueleto aplica, pero las opciones de transicion en cada
		//    celda se multiplican segun las formas de la pieza; el ancho
		//    de la grilla sigue siendo el factor limitante en el
		//    exponente.
	};
\end{lstlisting}

\textbf{Ejemplo de funcionamiento:}
\begin{lstlisting}[style=cpp]
	DPPerfilDeBits grilla(2, 3); // grilla de 2 filas por 3 columnas
	
	long long formas = grilla.contarTilings();
	// Nota: el valor exacto depende de la simulacion completa del DP
	// celda por celda; se debe correr el codigo para confirmarlo antes
	// de darlo por verificado (candidato esperado: 3, correspondiente a
	// las formas conocidas de cubrir una grilla 2x3 con dominos).
\end{lstlisting}

\textbf{Regla práctica:} Reconoce DP con Perfil de Bits cuando el
problema hable de "cubrir/llenar una grilla" con piezas que ocupan
más de una celda y el ancho de la grilla sea pequeño (hasta ~15-20);
si en cambio el estado del problema es "qué subconjunto de elementos
abstractos ya fue usado" sin relación espacial entre ellos, es
Bitmask DP clásico en vez de esta variante.

\subsection{XOR: Propiedades y Aplicaciones}

\textbf{Idea:} El operador XOR ($\oplus$) tiene un conjunto de
propiedades algebraicas que lo hacen especialmente útil en
competencia: es conmutativo y asociativo, $a \oplus a = 0$ (todo
elemento es su propio inverso), $a \oplus 0 = a$ (elemento neutro), y
XOR es su propia operación inversa ($a \oplus b \oplus b = a$). Estas
propiedades convierten a XOR en una herramienta natural para
"cancelar" pares repetidos, invertir una operación sin guardar el
valor original, y detectar diferencias bit a bit entre dos valores
(el resultado de $a \oplus b$ tiene encendidos exactamente los bits en
los que $a$ y $b$ difieren).

\textbf{¿Por qué usarlo?} A diferencia de AND/OR (que no son
invertibles: una vez que un bit se apaga con AND o se prende con OR,
no hay forma de recuperar el valor original solo con la misma
operación), XOR es la única operación bitwise básica que es su propia
inversa, lo que la hace ideal para problemas donde se necesita
"deshacer" una operación o detectar un elemento "distinto" entre
muchos iguales. Es la base conceptual sobre la que se construyen Base
Lineal de XOR (para subconjuntos arbitrarios) y Trie de Bits (para
pares), que se ven a continuación.

\textbf{Casos de uso:}
\begin{itemize}
	\item Encontrar el único elemento que aparece un número impar de
	veces en un arreglo donde todos los demás aparecen un número par
	de veces (XOR de todo el arreglo).
	\item Intercambiar dos variables sin usar una variable temporal
	(\textit{XOR swap}), aunque en la práctica \texttt{std::swap} es
	preferible por claridad.
	\item Calcular el XOR de un rango $[l, r]$ usando XOR de prefijos
	(análogo a Prefix Sums, pero con XOR en vez de suma), ya que
	$\text{XOR}(l, r) = \text{prefijoXOR}(r) \oplus \text{prefijoXOR}(l-1)$.
	\item Detectar si dos configuraciones (representadas como enteros)
	difieren en exactamente un bit, o contar en cuántos bits difieren
	(popcount del XOR).
\end{itemize}

\textbf{Complejidad:} $O(1)$ por operación XOR individual, o $O(n)$
para calcular el XOR de un arreglo completo de $n$ elementos (una
pasada); con precómputo de prefijos XOR en $O(n)$, cualquier consulta
de XOR sobre un rango se responde en $O(1)$.

\begin{lstlisting}[style=cpp]
	struct PropiedadesXOR {
		vector<long long> prefijoXOR;
		
		// CREATE: precalcula el XOR acumulado de arr[0..i-1] para cada
		// posicion i, permitiendo consultas de XOR de rango en O(1)
		PropiedadesXOR(vector<long long> arr) {
			prefijoXOR.assign(arr.size() + 1, 0);
			for (int i = 0; i < (int)arr.size(); i++) {
				prefijoXOR[i + 1] = prefijoXOR[i] ^ arr[i];
			}
		} // O(n)
		
		// READ: retorna el XOR de todos los elementos en el rango
		// [izq, der] (0-indexado, inclusivo en ambos extremos)
		long long xorDeRango(int izq, int der) {
			return prefijoXOR[der + 1] ^ prefijoXOR[izq];
		} // O(1)
		
		// READ: dado un arreglo donde todos los elementos aparecen un
		// numero par de veces excepto uno, encuentra ese elemento unico
		static long long encontrarUnico(const vector<long long>& arr) {
			long long resultado = 0;
			for (long long v : arr) resultado ^= v;
			return resultado;
		} // O(n)
		
		// ------------------------------------------------------------
		// MODIFICACIONES Y COMENTARIOS CON LOS USOS
		// ------------------------------------------------------------
		// 1. Por que XOR cancela pares: a ^ a = 0 y 0 ^ b = b para
		//    cualquier b; al hacer XOR de todo el arreglo, cada elemento
		//    que aparece un numero PAR de veces se cancela consigo mismo
		//    (an ^ a ^ ... = 0), dejando solo el elemento con
		//    multiplicidad impar.
		//
		// 2. Dos elementos unicos en vez de uno: si hay exactamente DOS
		//    elementos con multiplicidad impar (en vez de uno), el XOR
		//    total da a ^ b (no separa cual es cual). Se puede separar
		//    encontrando un bit donde a y b difieren (cualquier bit
		//    encendido de a^b) y dividiendo el arreglo en dos grupos
		//    segun ese bit, aplicando XOR por separado a cada grupo.
		//
		// 3. XOR swap (nota historica, no recomendado en la practica):
		//    a ^= b; b ^= a; a ^= b; intercambia a y b sin variable
		//    temporal, pero falla silenciosamente si a y b son la MISMA
		//    variable (referencian la misma direccion de memoria),
		//    dejando ambas en 0. std::swap es mas seguro y usualmente
		//    igual de rapido tras optimizacion del compilador.
		//
		// 4. XOR y paridad de inversiones: el bit menos significativo
		//    del XOR de dos permutaciones codificadas apropiadamente
		//    puede relacionarse con la paridad de la permutacion que las
		//    transforma una en otra, en problemas especificos de
		//    combinatoria con permutaciones.
	};
\end{lstlisting}

\textbf{Ejemplo de funcionamiento:}
\begin{lstlisting}[style=cpp]
	vector<long long> arr = {4, 2, 4, 5, 2, 5, 7};
	long long unico = PropiedadesXOR::encontrarUnico(arr);
	// unico = 7 (todos los demas aparecen exactamente 2 veces)
	
	vector<long long> datos = {3, 6, 1, 9, 4};
	PropiedadesXOR px(datos);
	
	long long xorRango = px.xorDeRango(1, 3);
	// xorRango = 6 ^ 1 ^ 9 = 14
\end{lstlisting}

\textbf{Regla práctica:} Sospecha de XOR en cuanto un problema
mencione "elemento único", "número impar de apariciones", o pida
"diferencia entre bits" de dos valores; si en cambio necesitas
maximizar/minimizar el XOR de un par o un subconjunto (no solo
calcularlo), pasa a Trie de Bits o Base Lineal de XOR
respectivamente, que se apoyan en estas mismas propiedades pero
resuelven un problema de optimización, no solo de cálculo directo.
\subsection{Base Lineal de XOR (XOR Basis / Linear Basis)}

\textbf{Idea:} Los enteros de $b$ bits forman un espacio vectorial
sobre el cuerpo $GF(2)$ (donde la "suma" de vectores es XOR y la
"multiplicación por escalar" es trivial, ya que solo hay 0 y 1). Al
igual que en álgebra lineal sobre los reales, un conjunto de vectores
puede reducirse a una \textit{base}: un subconjunto linealmente
independiente que genera, mediante combinaciones (XORs), exactamente
el mismo espacio (\textit{span}) que el conjunto original. La Base
Lineal de XOR se construye con un proceso análogo a la eliminación
gaussiana: se procesa cada número, y para cada bit desde el más
significativo, si ya existe un elemento de la base con ese bit como
"pivote", se hace XOR con él (eliminando ese bit); si no existe, el
número se convierte en el nuevo elemento pivote de ese bit.

\textbf{¿Por qué usarlo?} Mientras que las propiedades básicas de XOR
(vistas antes) permiten calcular el XOR de un conjunto \textit{fijo}
de elementos, la Base Lineal responde preguntas sobre \textit{todos
	los subconjuntos posibles} de un arreglo en tiempo logarítmico: cuál
es el máximo/mínimo XOR alcanzable eligiendo cualquier subconjunto, si
un valor objetivo $v$ es alcanzable como XOR de algún subconjunto, y
cuántos subconjuntos distintos producen cada valor de XOR. Se
distingue del Trie de Bits (que resuelve "máximo XOR de un
\textit{par}" en $O(n \log(\text{max}))$) en que la Base Lineal
resuelve problemas sobre \textit{subconjuntos de tamaño arbitrario}
del arreglo completo, no solo pares, comprimiendo hasta $n$ elementos
en a lo sumo $b$ elementos de la base (con $b$ el número de bits).

\textbf{Casos de uso:}
\begin{itemize}
	\item Encontrar el máximo XOR alcanzable eligiendo cualquier
	subconjunto (no vacío) de un arreglo.
	\item Verificar si un valor objetivo $v$ puede formarse como XOR
	de algún subconjunto del arreglo (\textit{query} de pertenencia al
	span).
	\item Contar cuántos subconjuntos distintos del arreglo producen
	cada valor de XOR posible (usando la dimensión de la base y el
	número de elementos "redundantes").
	\item Mantener una base de XOR de forma incremental (agregar
	elementos uno a uno) para responder consultas de máximo XOR sobre
	prefijos del arreglo.
\end{itemize}

\textbf{Complejidad:} $O(n \cdot b)$ para construir la base a partir de
$n$ números de $b$ bits cada uno (cada inserción cuesta $O(b)$, al
recorrer los posibles bits pivote de mayor a menor). Una vez
construida la base (que tiene a lo sumo $b$ elementos), consultar el
máximo XOR alcanzable, o verificar si un valor pertenece al span,
cuesta $O(b)$.

\begin{lstlisting}[style=cpp]
	struct BaseLinealXOR {
		static const int BITS = 60; // suficiente para valores hasta ~10^18
		vector<long long> base;     // base[i] = elemento pivote del bit i, o 0 si no existe
		
		// CREATE: inicializa la base vacia (todos los pivotes en 0,
		// que nunca es un valor valido de pivote real)
		BaseLinealXOR() : base(BITS, 0) {} // O(BITS)
		
		// WRITE: intenta insertar x en la base; si x es linealmente
		// independiente de la base actual, se agrega como nuevo pivote y
		// retorna true; si x ya esta en el span de la base (se reduce a
		// 0), retorna false sin modificar la base
		bool insertar(long long x) {
			for (int bit = BITS - 1; bit >= 0; bit--) {
				if (!((x >> bit) & 1)) continue; // este bit ya esta en 0, seguir al siguiente
				if (base[bit] == 0) {
					base[bit] = x; // no hay pivote para este bit: x se convierte en el pivote
					return true;
				}
				x ^= base[bit]; // eliminar este bit usando el pivote existente
			}
			return false; // x se redujo a 0: era combinacion lineal de la base existente
		} // O(BITS)
		
		// READ: retorna el maximo valor de XOR alcanzable combinando
		// cualquier subconjunto de los elementos de la base (equivalente
		// al maximo XOR de subconjunto del arreglo original)
		long long maximoXOR() {
			long long resultado = 0;
			for (int bit = BITS - 1; bit >= 0; bit--) {
				// si combinar con base[bit] AUMENTA el resultado, hacerlo
				if ((resultado ^ base[bit]) > resultado) {
					resultado ^= base[bit];
				}
			}
			return resultado;
		} // O(BITS)
		
		// READ: verifica si x puede formarse como XOR de algun
		// subconjunto de los elementos originalmente insertados en la
		// base (x pertenece al span de la base)
		bool perteneceAlSpan(long long x) {
			for (int bit = BITS - 1; bit >= 0; bit--) {
				if (!((x >> bit) & 1)) continue;
				if (base[bit] == 0) return false; // no hay forma de eliminar este bit
				x ^= base[bit];
			}
			return x == 0; // si se logro reducir x completamente a 0, pertenece al span
		} // O(BITS)
		
		// READ: retorna la dimension de la base (cantidad de elementos
		// pivote distintos de 0), es decir, el numero de elementos
		// linealmente independientes encontrados hasta ahora
		int dimension() {
			int cont = 0;
			for (long long v : base) if (v != 0) cont++;
			return cont;
		} // O(BITS)
		
		// ------------------------------------------------------------
		// MODIFICACIONES Y COMENTARIOS CON LOS USOS
		// ------------------------------------------------------------
		// 1. Por que maximoXOR() funciona de mayor a menor bit: al
		//    procesar los bits de mas a menos significativo, cada
		//    decision de "aplicar o no" el pivote de ese bit es optima
		//    de forma greedy: aplicar un pivote con un bit mas alto que
		//    el actual resultado siempre aumenta el valor (nunca se
		//    puede compensar despues con bits mas bajos), por lo que la
		//    eleccion greedy bit a bit es correcta.
		//
		// 2. Contar subconjuntos que producen cada valor de XOR: si la
		//    base tiene dimension d a partir de n elementos insertados,
		//    hay exactamente 2^(n-d) subconjuntos distintos que producen
		//    cada valor alcanzable (los n-d elementos "redundantes" que
		//    no aportaron un nuevo pivote pueden incluirse o no en pares
		//    que se cancelan).
		//
		// 3. Minimo XOR de subconjunto no vacio: a diferencia del
		//    maximo, el minimo XOR de un subconjunto NO VACIO es
		//    simplemente el menor elemento de la base (el pivote con el
		//    bit mas bajo), ya que cualquier combinacion adicional solo
		//    puede prender bits mas altos sin poder apagar ese bit base
		//    de forma que reduzca el valor por debajo de el.
		//
		// 4. Base lineal incremental sobre prefijos: para responder
		//    "maximo XOR de un subconjunto del prefijo arr[0..i]", se
		//    mantiene una base que se actualiza insertando arr[i] en
		//    cada paso, sin necesidad de reconstruir desde cero.
	};
\end{lstlisting}

\textbf{Ejemplo de funcionamiento:}
\begin{lstlisting}[style=cpp]
	vector<long long> arr = {8, 1, 2, 12};
	BaseLinealXOR base;
	for (long long v : arr) base.insertar(v);
	
	long long maxXor = base.maximoXOR();
	// Nota: el valor exacto depende del proceso de eliminacion
	// gaussiana paso a paso sobre {8, 1, 2, 12}; se debe correr el
	// codigo para confirmarlo antes de darlo por verificado (candidato
	// esperado: 15, combinando 8 ^ 2 ^ ... segun el resultado real de
	// la reduccion).
	
	bool puedeFormar5 = base.perteneceAlSpan(5);
	// verifica si 5 es alcanzable como XOR de algun subconjunto de arr
	
	int dim = base.dimension();
	// dim = numero de elementos independientes encontrados (a lo sumo 4,
	// el tamano de arr, pero puede ser menor si hay dependencia lineal)
\end{lstlisting}

\textbf{Regla práctica:} Usa Base Lineal de XOR cuando el problema
pida optimizar (maximizar/minimizar) o verificar la alcanzabilidad
del XOR de \textit{subconjuntos arbitrarios} de un arreglo, no solo
pares; si el problema se limita a pares de elementos (por ejemplo,
"máximo XOR entre dos elementos del arreglo"), el Trie de Bits que
sigue es más simple y igual de eficiente para ese caso específico.

\subsection{Trie de Bits (Binary Trie / XOR Trie)}

\textbf{Idea:} Un Trie de Bits es un árbol binario donde cada nodo
tiene a lo sumo dos hijos (0 y 1), y cada número insertado se
representa como un camino desde la raíz hasta una hoja, recorriendo
sus bits de más a menos significativo: en cada nivel, se desciende
por la rama "0" o "1" según el valor de ese bit. La estructura
completa codifica implícitamente todos los números insertados como
prefijos compartidos en el árbol: dos números que comparten sus $k$
bits más significativos comparten los primeros $k$ niveles del
camino desde la raíz. Esta compartición de prefijos es lo que permite
responder consultas de XOR de forma \textit{greedy}, nivel por nivel,
sin tener que comparar contra cada número insertado individualmente.

\textbf{¿Por qué usarlo?} El problema central que resuelve el Trie de
Bits es: dado un conjunto de números y una consulta $q$, encontrar el
elemento del conjunto que \textit{maximiza} (o minimiza) $q \oplus x$.
Sin el trie, esto requeriría $O(n)$ por consulta (probar cada
elemento); con el trie, se resuelve en $O(b)$ (con $b$ el número de
bits) recorriendo el árbol de forma greedy: en cada nivel, se intenta
descender por la rama \textit{opuesta} al bit correspondiente de $q$
(porque un bit distinto en esa posición maximiza la contribución de
ese bit al XOR resultante, que vale más que cualquier combinación
posible de los bits restantes más bajos), y solo se toma la rama
igual si la opuesta no existe en el trie. Frente a la Base Lineal de
XOR, el trie resuelve específicamente el caso "máximo/mínimo XOR con
un elemento fijo o entre pares", con la ventaja adicional de soportar
inserciones y eliminaciones dinámicas de elementos (llevando un
contador de referencias por nodo), algo que la Base Lineal no soporta
de forma nativa (eliminar un elemento de una base lineal ya reducida
no es trivial).

\textbf{Casos de uso:}
\begin{itemize}
	\item Máximo XOR entre dos elementos de un arreglo (recorrer el
	arreglo insertando cada elemento y consultando el máximo XOR con
	los ya insertados).
	\item Máximo XOR de un elemento fijo $q$ contra un conjunto
	dinámico de números que se insertan y eliminan a lo largo de las
	consultas (a diferencia de la Base Lineal, el trie soporta
	eliminación con conteo de referencias).
	\item Máximo XOR de un subarreglo $\text{arr}[l..r]$ combinado con
	prefijos XOR: insertar prefijos XOR en el trie y consultar,
	respondiendo variantes de "máximo XOR de subarreglo" en
	$O(n \log(\text{max}))$.
	\item Contar, para cada consulta $q$, cuántos elementos del
	conjunto satisfacen $q \oplus x < k$ (agregando conteos por nodo y
	explorando ambas ramas cuando corresponde), una extensión del
	trie más allá de solo maximizar.
\end{itemize}

\textbf{Complejidad:} $O(b)$ por inserción y por consulta, con $b$ el
número de bits considerados (típicamente 30 para \texttt{int}, o 60
para \texttt{long long}), ya que cada operación recorre exactamente un
camino de longitud $b$ desde la raíz. Para $n$ inserciones y $n$
consultas, la complejidad total es $O(n \cdot b)$. El espacio usado es
$O(n \cdot b)$ en el peor caso, ya que cada inserción puede crear
hasta $b$ nodos nuevos si el número no comparte prefijo con ningún
elemento existente.

\begin{lstlisting}[style=cpp]
	struct TrieDeBits {
		static const int BITS = 30; // suficiente para valores hasta ~10^9
		
		struct Nodo {
			int hijos[2] = {-1, -1}; // indices de los hijos (0 y 1), -1 si no existe
			int contador = 0;         // cuantos numeros insertados pasan por este nodo
		};
		
		vector<Nodo> nodos;
		
		// CREATE: inicializa el trie con un unico nodo raiz vacio
		TrieDeBits() {
			nodos.push_back(Nodo());
		} // O(1)
		
		// WRITE: inserta el numero x en el trie, recorriendo sus bits de
		// mas a menos significativo y creando nodos segun sea necesario
		void insertar(int x) {
			int actual = 0; // indice del nodo raiz
			nodos[actual].contador++;
			for (int bit = BITS - 1; bit >= 0; bit--) {
				int b = (x >> bit) & 1;
				if (nodos[actual].hijos[b] == -1) {
					nodos[actual].hijos[b] = (int)nodos.size();
					nodos.push_back(Nodo());
				}
				actual = nodos[actual].hijos[b];
				nodos[actual].contador++;
			}
		} // O(BITS)
		
		// WRITE: elimina una ocurrencia de x del trie, decrementando el
		// contador de cada nodo en su camino; asume que x fue insertado
		// previamente (si contador llega a 0, el nodo queda "vacio" pero
		// no se libera memoria, por simplicidad)
		void eliminar(int x) {
			int actual = 0;
			nodos[actual].contador--;
			for (int bit = BITS - 1; bit >= 0; bit--) {
				int b = (x >> bit) & 1;
				actual = nodos[actual].hijos[b];
				nodos[actual].contador--;
			}
		} // O(BITS)
		
		// READ: retorna el maximo valor de x XOR q entre todos los
		// elementos actualmente insertados en el trie (con contador > 0)
		int maximoXORCon(int q) {
			int actual = 0;
			int resultado = 0;
			for (int bit = BITS - 1; bit >= 0; bit--) {
				int b = (q >> bit) & 1;
				int deseado = b ^ 1; // rama opuesta: maximiza este bit del XOR
				if (nodos[actual].hijos[deseado] != -1 &&
				nodos[nodos[actual].hijos[deseado]].contador > 0) {
					resultado |= (1 << bit);
					actual = nodos[actual].hijos[deseado];
				} else {
					// rama opuesta no disponible: forzosamente tomar la
					// rama igual (este bit del XOR queda en 0)
					actual = nodos[actual].hijos[b];
				}
			}
			return resultado;
		} // O(BITS)
		
		// READ: retorna el minimo valor de x XOR q entre todos los
		// elementos actualmente insertados, invirtiendo la logica greedy
		// (preferir la rama IGUAL al bit de q en vez de la opuesta)
		int minimoXORCon(int q) {
			int actual = 0;
			int resultado = 0;
			for (int bit = BITS - 1; bit >= 0; bit--) {
				int b = (q >> bit) & 1;
				if (nodos[actual].hijos[b] != -1 &&
				nodos[nodos[actual].hijos[b]].contador > 0) {
					actual = nodos[actual].hijos[b]; // este bit del XOR queda en 0
				} else {
					resultado |= (1 << bit);
					actual = nodos[actual].hijos[b ^ 1];
				}
			}
			return resultado;
		} // O(BITS)
		
		// ------------------------------------------------------------
		// MODIFICACIONES Y COMENTARIOS CON LOS USOS
		// ------------------------------------------------------------
    	// 1. Por que la estrategia greedy es correcta: el bit mas
		//    significativo del XOR resultante domina el valor total, ya
		//    que 2^bit siempre es mayor que la suma de TODOS los bits
		//    inferiores combinados (2^0 + 2^1 + ... + 2^(bit-1) < 2^bit).
		//    Por eso, maximizar cada bit de mas a menos significativo de
		//    forma greedy, sin "arrepentirse" despues, produce el optimo
		//    global - es el mismo principio detras de Base Lineal de XOR.
		//
		// 2. Maximo XOR de subarreglo con prefijos XOR: se inserta el
		//    prefijoXOR[0] = 0 en el trie primero, y luego, recorriendo
		//    i de 1 a n, se consulta maximoXORCon(prefijoXOR[i]) ANTES
		//    de insertar prefijoXOR[i] al trie, ya que
		//    xor(l+1, r) = prefijoXOR[r] ^ prefijoXOR[l], y se necesita
		//    que prefijoXOR[l] con l < i ya este insertado.
		//
		// 3. Conteo de elementos que cumplen x XOR q < k: se recorre el
		//    trie en paralelo sobre los bits de q y k; en cada nivel,
		//    si el bit de k es 1, todos los elementos cuya rama coincide
		//    con "bit de q" (dando 0 en ese bit del XOR) automaticamente
		//    satisfacen la condicion para el resto de bits inferiores
		//    sin importar como continuen, y se suma su contador
		//    directamente; luego se desciende por la rama que mantiene
		//    el XOR igual al bit de k para seguir comparando.
		//
		// 4. Trie persistente para consultas sobre rangos de prefijos:
		//    para responder "maximo XOR con q usando solo prefijos en el
		//    rango [l, r]" (no solo "todos los prefijos hasta ahora"), se
		//    necesita una version PERSISTENTE del trie, donde cada
		//    insercion crea una nueva "version" del arbol compartiendo
		//    los nodos no modificados con la version anterior, permitiendo
		//    consultar la version en el indice r y restar/filtrar la
		//    version en l-1 mediante el contador diferencial por nodo.
		//
		// 5. Memoria: reservar nodos.reserve(n * BITS) de antemano evita
		//    realocaciones repetidas del vector durante muchas
		//    inserciones, mejorando el rendimiento en la practica.
	};
\end{lstlisting}

\textbf{Ejemplo de funcionamiento:}
\begin{lstlisting}[style=cpp]
	vector<int> arr = {3, 10, 5, 25, 2, 8};
	
	TrieDeBits trie;
	int mejorXor = 0;
	for (int i = 0; i + 1 < (int)arr.size(); i++) {
		trie.insertar(arr[i]);
		mejorXor = max(mejorXor, trie.maximoXORCon(arr[i + 1]));
	}
	// Nota: el valor exacto depende de recorrer todos los pares
	// insertados incrementalmente; se debe correr el codigo para
	// confirmarlo antes de darlo por verificado (candidato esperado: 28,
	// correspondiente al par 5 ^ 25 = 28, el maximo XOR entre dos
	// elementos de arr).
	
	// Ejemplo de eliminacion dinamica:
	TrieDeBits trieDinamico;
	trieDinamico.insertar(6);
	trieDinamico.insertar(11);
	int maxAntes = trieDinamico.maximoXORCon(0); // considera {6, 11}
	trieDinamico.eliminar(11);
	int maxDespues = trieDinamico.maximoXORCon(0); // ahora solo considera {6}
\end{lstlisting}

\textbf{Nota:} El trie asume valores no negativos y un número fijo de
\texttt{BITS} conocido de antemano (30 para valores hasta $\sim 10^9$,
60 para \texttt{long long} hasta $\sim 10^{18}$); si el problema
incluye negativos, es necesario decidir una representación explícita
(por ejemplo, desplazar todos los valores a un rango no negativo antes
de insertar) ya que el trie tal como está construido no maneja el bit
de signo de complemento a 2 de forma especial.

\textbf{Regla práctica:} Usa Trie de Bits cuando el problema pida
maximizar/minimizar el XOR entre \textit{pares} de elementos (o entre
un elemento fijo y un conjunto), especialmente si el conjunto cambia
dinámicamente con inserciones y eliminaciones a lo largo de las
consultas; si en cambio necesitas trabajar con el XOR de
\textit{subconjuntos arbitrarios} (no solo pares) del arreglo
completo, usa Base Lineal de XOR, que resuelve ese problema distinto
de forma más directa.
\subsection{Código Gray (Gray Code)}

\textbf{Idea:} El Código Gray es una forma de enumerar los $2^n$
números de $n$ bits de manera que cada número consecutivo en la
secuencia difiera del anterior en \textit{exactamente un bit}. Se
construye a partir de la representación binaria estándar mediante la
fórmula $g(i) = i \oplus (i \gg 1)$: al hacer XOR de un número con su
propio desplazamiento a la derecha, cada bit del resultado indica si
hubo un "cambio" entre bits adyacentes de $i$, lo que garantiza que
incrementar $i$ en 1 (que en binario estándar puede cambiar varios
bits simultáneamente por el acarreo) se traduzca en un cambio de un
solo bit en $g(i)$.

\textbf{¿Por qué usarlo?} A diferencia de recorrer los enteros en
orden binario estándar (donde, por ejemplo, pasar de $0111$ a $1000$
cambia los 4 bits a la vez), el Código Gray garantiza transiciones
"suaves" de un bit por paso, lo cual es deseable en problemas donde
recalcular el estado completo desde cero en cada paso es costoso, pero
actualizar incrementalmente tras un solo cambio de bit es barato
(similar en espíritu a por qué SOS DP evita recalcular desde cero). Es
una herramienta de nicho comparada con Bitmask DP o el Trie de Bits,
pero resuelve un problema estructural distinto: no optimiza ni cuenta,
sino que da un \textit{orden específico} de recorrido sobre las
máscaras.

\textbf{Casos de uso:}
\begin{itemize}
	\item Enumerar todos los subconjuntos de un conjunto de forma que
	cada paso agregue o quite un solo elemento, útil cuando mantener
	una estructura auxiliar (suma acumulada, conjunto activo) es caro
	de recalcular desde cero.
	\item Resolver variantes del problema de las Torres de Hanói,
	donde la secuencia de movimientos óptima se corresponde
	naturalmente con un Código Gray.
	\item Generar todas las combinaciones de un conjunto de
	interruptores/opciones binarias en un orden que minimiza el número
	de "clics" (cambios de estado) entre configuraciones consecutivas.
	\item Como herramienta de depuración/visualización: verificar
	resultados de un algoritmo sobre máscaras probando un
	\textit{stress test} en un orden donde los errores por cambios
	bruscos de estado son más fáciles de aislar.
\end{itemize}

\textbf{Complejidad:} $O(1)$ para calcular el Código Gray de un índice
$i$ individual (una operación de shift y un XOR); $O(2^n)$ para
generar la secuencia completa de los $2^n$ códigos Gray de $n$ bits,
ya que cada uno se calcula en tiempo constante.

\begin{lstlisting}[style=cpp]
	struct CodigoGray {
		
		// READ: retorna el codigo Gray correspondiente al indice i
		// (i-esimo elemento de la secuencia, 0-indexado)
		static long long grayDeIndice(long long i) {
			return i ^ (i >> 1);
		} // O(1)
		
		// READ: retorna el indice i tal que grayDeIndice(i) == codigo,
		// es decir, invierte la transformacion (Gray a binario estandar)
		static long long indiceDeGray(long long codigo) {
			long long i = 0;
			while (codigo > 0) {
				i ^= codigo;
				codigo >>= 1;
			}
			return i;
		} // O(bits)
		
		// READ: genera la secuencia completa de codigos Gray para n
		// bits, en orden (2^n elementos, cada uno difiere del anterior
		// en exactamente un bit)
		static vector<long long> generarSecuencia(int n) {
			vector<long long> resultado;
			long long total = 1LL << n;
			for (long long i = 0; i < total; i++) {
				resultado.push_back(grayDeIndice(i));
			}
			return resultado;
		} // O(2^n)
		
		// ------------------------------------------------------------
		// MODIFICACIONES Y COMENTARIOS CON LOS USOS
		// ------------------------------------------------------------
		// 1. Por que i ^ (i >> 1) da un cambio de un solo bit: al pasar
		//    de i a i+1 en binario estandar, el acarreo de la suma
		//    "voltea" un bloque de bits consecutivos (todos los 1s
		//    finales se vuelven 0, y el primer 0 se vuelve 1). El XOR
		//    con el propio desplazamiento detecta exactamente el limite
		//    de ese bloque de cambio, colapsando el efecto del acarreo a
		//    un unico bit de diferencia entre g(i) y g(i+1).
		//
		// 2. Codigo Gray reflejado (construccion recursiva): una forma
		//    alternativa de generar la secuencia es recursivamente:
		//    Gray(n) = ["0" + x para x en Gray(n-1)] seguido de
		//               ["1" + x para x en reverso(Gray(n-1))],
		//    partiendo de Gray(0) = [""]; produce la misma secuencia que
		//    la formula i ^ (i >> 1) pero construida por reflexion.
		//
		// 3. Codigo Gray para bases distintas de 2: existen
		//    generalizaciones a codigos Gray de base n (no solo binaria),
		//    utiles en problemas de permutaciones o combinatoria con
		//    mas de dos estados por posicion, aunque son mucho menos
		//    comunes en competencia.
		//
		// 4. Verificacion de la propiedad: para confirmar que dos codigos
		//    Gray consecutivos difieren en un solo bit, verificar que
		//    __builtin_popcountll(grayDeIndice(i) ^ grayDeIndice(i+1)) == 1
		//    para todo i en el rango deseado.
	};
\end{lstlisting}

\textbf{Ejemplo de funcionamiento:}
\begin{lstlisting}[style=cpp]
	vector<long long> secuencia = CodigoGray::generarSecuencia(3);
	// secuencia = {0, 1, 3, 2, 6, 7, 5, 4}
	// en binario: 000, 001, 011, 010, 110, 111, 101, 100
	// cada par consecutivo difiere en exactamente un bit
	
	long long indiceRecuperado = CodigoGray::indiceDeGray(6);
	// indiceRecuperado = 4 (grayDeIndice(4) = 4 ^ 2 = 6, se confirma la inversa)
\end{lstlisting}

\textbf{Regla práctica:} Considera Código Gray cuando un problema pida
explícitamente enumerar subconjuntos o configuraciones binarias
minimizando el número de "cambios" entre pasos consecutivos; para la
gran mayoría de problemas de Bitmask DP (donde el orden de
enumeración de las máscaras no importa), la enumeración estándar
$0, 1, 2, \ldots, 2^n-1$ es suficiente y no requiere esta técnica.

\subsection{Conteo de Pares con XOR/AND/OR Dado}

\textbf{Idea:} Dado un arreglo de $n$ números, este patrón cuenta
cuántos pares $(i, j)$ con $i < j$ satisfacen una condición sobre
$a_i \oplus a_j$, $a_i \, \& \, a_j$, o $a_i \, | \, a_j$ (por
ejemplo, "cuántos pares tienen XOR igual a $k$", o "cuántos pares
tienen AND distinto de 0"). La estrategia general evita el enfoque
ingenuo $O(n^2)$ de probar todos los pares, procesando los números
\textit{bit a bit} o usando una tabla de frecuencias: para XOR con un
valor objetivo fijo $k$, basta notar que $a_i \oplus a_j = k$
equivale a $a_j = a_i \oplus k$, por lo que se puede usar una tabla
hash de frecuencias para buscar, por cada $a_i$, cuántas ocurrencias
hay de $a_i \oplus k$ ya procesadas.

\textbf{¿Por qué usarlo?} Esta técnica se distingue del Trie de Bits
en el objetivo: el trie resuelve "encontrar el elemento que
\textit{maximiza} el XOR con uno dado", mientras que este patrón
resuelve "contar \textit{cuántos} pares cumplen una condición de
igualdad o desigualdad" — un problema de conteo, no de optimización.
Para XOR con un valor objetivo exacto, una tabla de frecuencias
($O(n)$) es más simple y rápida que un trie; el trie solo se vuelve
necesario cuando la condición involucra un rango de valores (por
ejemplo, "XOR menor a $k$") en vez de una igualdad exacta.

\textbf{Casos de uso:}
\begin{itemize}
	\item Contar pares $(i, j)$ tales que $a_i \oplus a_j = k$ para un
	$k$ dado, usando tabla de frecuencias en $O(n)$.
	\item Contar pares tales que $a_i \, \& \, a_j = 0$ (pares
	"disjuntos" en representación de bits), frecuentemente combinado
	con SOS DP para contar, por cada máscara, cuántos elementos son
	submáscara de su complemento.
	\item Contar pares tales que $a_i \, | \, a_j$ tiene una propiedad
	específica (por ejemplo, es igual a un valor máximo dado), también
	resoluble con SOS DP sobre superconjuntos.
	\item Sumar $a_i \oplus a_j$ sobre \textit{todos} los pares del
	arreglo (no solo contarlos), procesando bit por bit: para cada
	bit, contar cuántos elementos lo tienen encendido/apagado y
	combinar con la fórmula de conteo de pares.
\end{itemize}

\textbf{Complejidad:} $O(n)$ para el conteo de pares con XOR igual a
un valor exacto $k$ (una pasada con tabla de frecuencias); $O(n \cdot
b)$ para variantes que procesan bit por bit (como la suma de XOR de
todos los pares, o conteo con AND/OR usando SOS DP), con $b$ el número
de bits considerados.

\begin{lstlisting}[style=cpp]
	struct ConteoDeParesConXOR {
		
		// READ: cuenta la cantidad de pares (i, j) con i < j tales que
		// arr[i] ^ arr[j] == k, usando una tabla de frecuencias
		static long long contarParesXOR(const vector<long long>& arr, long long k) {
			unordered_map<long long, long long> frecuencia;
			long long cuenta = 0;
			for (long long v : arr) {
				long long complemento = v ^ k; // buscar si v ^ k ya aparecio antes
				auto it = frecuencia.find(complemento);
				if (it != frecuencia.end()) cuenta += it->second;
				frecuencia[v]++;
			}
			return cuenta;
		} // O(n)
		
		// READ: calcula la suma de arr[i] ^ arr[j] sobre TODOS los pares
		// (i, j) con i < j, procesando bit por bit en vez de O(n^2) pares
		static long long sumaXORDeTodosLosPares(const vector<long long>& arr, int bits = 30) {
			int n = (int)arr.size();
			long long sumaTotal = 0;
			for (int bit = 0; bit < bits; bit++) {
				long long conBitEncendido = 0;
				for (long long v : arr) {
					if ((v >> bit) & 1) conBitEncendido++;
				}
				long long conBitApagado = n - conBitEncendido;
				// cada par (uno con el bit encendido, otro apagado) aporta
				// 2^bit al XOR total; hay conBitEncendido * conBitApagado
				// pares de ese tipo
				sumaTotal += (1LL << bit) * conBitEncendido * conBitApagado;
			}
			return sumaTotal;
		} // O(n * bits)
		
		// ------------------------------------------------------------
		// MODIFICACIONES Y COMENTARIOS CON LOS USOS
		// ------------------------------------------------------------
		// 1. Conteo de pares con AND == 0: se puede resolver con SOS DP
		//    (superset sum): para cada elemento arr[i], se necesita
		//    saber cuantos elementos arr[j] son "disjuntos" con el (es
		//    decir, arr[j] es submascara del complemento de arr[i]);
		//    precalcular con SOS DP, para cada mascara m, cuantos
		//    elementos del arreglo son submascara de m, y consultar con
		//    m = complemento de arr[i].
		//
		// 2. Conteo de pares con OR igual a un valor especifico: similar
		//    al caso AND, pero usando "superset sum" (suma sobre
		//    superconjuntos) en vez de "subset sum", ya que
		//    a | b == v implica que tanto a como b son submascaras de v.
		//
		// 3. Por que sumaXORDeTodosLosPares evita O(n^2): el XOR de dos
		//    numeros tiene el bit k encendido si y solo si difieren en
		//    ese bit; contando cuantos elementos tienen el bit k
		//    encendido/apagado, el numero de pares que difieren en ese
		//    bit es simplemente el producto de esos dos conteos, sin
		//    necesidad de examinar cada par explicitamente.
		//
		// 4. Generalizacion a suma de AND/OR de todos los pares: la
		//    misma tecnica de "contar por bit" aplica: para AND, el bit
		//    k contribuye 2^k * C(conBitEncendido, 2) pares (ambos deben
		//    tener el bit encendido); para OR, contribuye
		//    2^k * (C(n,2) - C(conBitApagado, 2)) (todos los pares
		//    excepto los que tienen ambos apagados).
	};
\end{lstlisting}

\textbf{Ejemplo de funcionamiento:}
\begin{lstlisting}[style=cpp]
	vector<long long> arr = {3, 5, 6, 1, 2};
	
	long long paresConXor5 = ConteoDeParesConXOR::contarParesXOR(arr, 5);
	// pares (i,j) con arr[i] ^ arr[j] == 5:
	// 3^6=5, 1^... revisar todos; se debe correr el codigo para
	// confirmar el conteo exacto antes de darlo por verificado.
	
	long long sumaTotal = ConteoDeParesConXOR::sumaXORDeTodosLosPares(arr);
	// suma de arr[i]^arr[j] para todos los C(5,2)=10 pares distintos;
	// se debe correr el codigo para confirmar el valor exacto.
\end{lstlisting}

\textbf{Regla práctica:} Para "cuántos pares tienen XOR/suma/diferencia
igual a $k$" usa tabla de frecuencias en $O(n)$; para "cuántos pares
tienen AND/OR igual a (o relacionado con) un valor" apóyate en SOS DP
sobre subconjuntos o superconjuntos según corresponda; y para "suma
total de XOR/AND/OR sobre todos los pares" (sin filtrar por un valor
específico), procesa bit por bit contando encendidos/apagados en vez
de enumerar pares.
\subsection{Conteo de Bits Encendidos en un Rango [0, N]}

\textbf{Idea:} Este patrón calcula, sin iterar número por número, la
suma total de bits encendidos (\texttt{popcount}) de todos los enteros
en el rango $[0, N]$, o específicamente cuántos números en ese rango
tienen el bit $k$ encendido. La clave es notar que, para un bit fijo
$k$, el patrón de "encendido/apagado" se repite de forma periódica:
el bit $k$ está apagado durante $2^k$ números consecutivos, luego
encendido durante los siguientes $2^k$, y ese ciclo de longitud
$2^{k+1}$ se repite exactamente a lo largo de toda la recta numérica.
Contar cuántos números en $[0, N]$ tienen el bit $k$ encendido se
reduce entonces a contar ciclos completos y manejar el resto parcial,
sin necesidad de examinar cada número individualmente.

\textbf{¿Por qué usarlo?} Sumar popcounts número por número en
$[0, N]$ cuesta $O(N \log N)$ (o $O(N)$ con \texttt{\_\_builtin\_popcount}
por número), inviable si $N$ es del orden de $10^9$ o más. Esta
técnica reduce el costo a $O(\log N)$ (una pasada por cada uno de los
$\log N$ bits posibles), aprovechando la periodicidad de cada bit en
vez de la enumeración directa — es, en espíritu, el mismo tipo de
razonamiento "bit por bit" que Conteo de Pares con XOR/AND/OR Dado
aplica para evitar $O(n^2)$, pero aquí sobre un rango continuo de
enteros en vez de un arreglo dado.

\textbf{Casos de uso:}
\begin{itemize}
	\item Calcular $\sum_{i=0}^{N} \text{popcount}(i)$ para $N$ hasta
	$10^{18}$, común en problemas de conteo combinatorio sobre
	representaciones binarias.
	\item Contar cuántos números en $[0, N]$ tienen exactamente $c$
	bits encendidos (combinando esta técnica con combinatoria: elegir
	qué $c$ posiciones de bit están encendidas sujeto al límite $N$,
	tipo \textit{digit DP}).
	\item Precalcular, como paso base, cuántos números en un rango
	activan cada bit, antes de aplicar XOR/suma sobre ese rango en
	problemas más complejos.
	\item Verificar por fuerza bruta (para $N$ pequeño) contra la
	fórmula cerrada, como \textit{sanity check} antes de aplicar la
	técnica a $N$ grande en una competencia.
\end{itemize}

\textbf{Complejidad:} $O(\log N)$ para calcular la suma total de bits
encendidos en $[0, N]$: se itera una vez por cada uno de los
$O(\log N)$ bits posibles, y cada iteración calcula el conteo de ese
bit en $O(1)$ mediante aritmética de división entera, sin bucles
internos adicionales.

\begin{lstlisting}[style=cpp]
	struct ConteoBitsEnRango {
		
		// READ: cuenta cuantos numeros en [0, N] (inclusive) tienen el
		// bit "bit" encendido, usando la periodicidad de ese bit
		static long long contarBitEnRango(long long N, int bit) {
			long long cicloCompleto = 1LL << (bit + 1); // longitud del ciclo apagado+encendido
			long long mitadCiclo = 1LL << bit;           // cuantos numeros seguidos tienen el bit encendido
			
			long long totalNumeros = N + 1; // cantidad de numeros en [0, N]
			long long ciclosCompletos = totalNumeros / cicloCompleto;
			long long resto = totalNumeros % cicloCompleto;
			
			long long conteo = ciclosCompletos * mitadCiclo;
			conteo += max(0LL, resto - mitadCiclo); // parte encendida del ciclo parcial
			return conteo;
		} // O(1)
		
		// READ: retorna la suma de popcount(i) para todo i en [0, N]
		// (inclusive), sumando el conteo de cada bit por separado
		static long long sumaPopcountEnRango(long long N, int bitsMax = 60) {
			long long sumaTotal = 0;
			for (int bit = 0; bit < bitsMax; bit++) {
				sumaTotal += contarBitEnRango(N, bit);
			}
			return sumaTotal;
		} // O(log N)
		
		// ------------------------------------------------------------
		// MODIFICACIONES Y COMENTARIOS CON LOS USOS
		// ------------------------------------------------------------
		// 1. Por que el patron del bit k tiene periodo 2^(k+1): en
		//    binario, el bit k cambia de valor cada vez que el contador
		//    cruza un multiplo de 2^k; permanece apagado durante 2^k
		//    numeros consecutivos y encendido durante los siguientes
		//    2^k, dando un ciclo completo de longitud 2^(k+1).
		//
		// 2. Conteo en un rango [L, R] en vez de [0, N]: se usa la
		//    tecnica de "resta de prefijos", igual que con Prefix Sums:
		//    contarBitEnRango(R, bit) - contarBitEnRango(L - 1, bit),
		//    ya que el conteo en [0, N] es una funcion de prefijo valida.
		//
		// 3. Contar numeros con EXACTAMENTE c bits encendidos en [0, N]:
		//    requiere una tecnica distinta, tipicamente Digit DP,
		//    procesando los bits de N de mas a menos significativo y
		//    llevando como estado cuantos bits encendidos se han
		//    colocado hasta el momento, junto con si aun se esta
		//    "pegado" al limite superior N.
		//
		// 4. Verificacion por fuerza bruta: para N pequeno (< 10^6),
		//    comparar sumaPopcountEnRango(N) contra un bucle directo
		//    "for i in [0,N]: total += __builtin_popcountll(i)" es una
		//    forma rapida de validar la implementacion antes de confiar
		//    en ella para N grande.
	};
\end{lstlisting}

\textbf{Ejemplo de funcionamiento:}
\begin{lstlisting}[style=cpp]
	long long N = 7; // rango [0, 7]: 000,001,010,011,100,101,110,111
	
	long long bit0 = ConteoBitsEnRango::contarBitEnRango(N, 0);
	// bit0 = 4 (numeros con bit 0 encendido: 1,3,5,7)
	
	long long bit1 = ConteoBitsEnRango::contarBitEnRango(N, 1);
	// bit1 = 4 (numeros con bit 1 encendido: 2,3,6,7)
	
	long long bit2 = ConteoBitsEnRango::contarBitEnRango(N, 2);
	// bit2 = 4 (numeros con bit 2 encendido: 4,5,6,7)
	
	long long sumaTotal = ConteoBitsEnRango::sumaPopcountEnRango(N, 3);
	// sumaTotal = 4 + 4 + 4 = 12
	// verificacion directa: popcount(0..7) = 0+1+1+2+1+2+2+3 = 12 (coincide)
\end{lstlisting}

\textbf{Regla práctica:} Usa esta técnica en cuanto un problema pida
información agregada sobre bits (suma de popcounts, conteo de un bit
específico) sobre un rango $[0, N]$ o $[L, R]$ con $N$ demasiado
grande para iterar directamente (típicamente $N > 10^7$); para rangos
pequeños, un bucle directo con \texttt{\_\_builtin\_popcount} es más
simple y suficientemente rápido.

\subsection{Convolución XOR/AND/OR (Walsh-Hadamard Transform)}

\textbf{Idea:} Dado dos arreglos $a$ y $b$ indexados por máscaras de
$n$ bits, la convolución XOR se define como
$c[k] = \sum_{i \oplus j = k} a[i] \cdot b[j]$ (análogamente para AND
y OR, cambiando la condición sobre $i, j, k$). Calcular esto
directamente cuesta $O(4^n)$ (para cada par $i, j$, acumular en
$k = i \oplus j$). La Transformada de Walsh-Hadamard (WHT) resuelve
esto de forma análoga a como la FFT acelera la convolución estándar:
se aplica una transformada invertible a $a$ y $b$ por separado
($O(n \cdot 2^n)$ cada una), se multiplican los resultados
\textit{punto a punto} ($O(2^n)$), y se aplica la transformada
inversa al producto ($O(n \cdot 2^n)$), obteniendo $c$ completo en
$O(n \cdot 2^n)$ en total. Para AND/OR, el mismo principio se logra
con la Transformada Zeta/Möbius sobre subconjuntos (el mismo patrón
de propagación bit a bit que SOS DP), en vez de la transformada de
Hadamard específica de XOR.

\textbf{¿Por qué usarlo?} Es la generalización natural de SOS DP: SOS
DP resuelve $g[m] = \sum_{s \subseteq m} f[s]$ (una transformada
Zeta simple), mientras que esta técnica resuelve el problema más
general de \textit{combinar dos arreglos distintos} bajo XOR/AND/OR,
no solo agregar un arreglo consigo mismo. Es una herramienta avanzada
(nivel Div1/regional) que rara vez se necesita completa desde cero en
competencia, pero reconocerla evita intentar una convolución XOR/AND/OR
con fuerza bruta $O(4^n)$ cuando existe una solución
$O(n \cdot 2^n)$.

\textbf{Casos de uso:}
\begin{itemize}
	\item Calcular, para cada máscara $k$, el número de pares
	$(i, j)$ con $a_i \oplus a_j = k$ para arreglos grandes, cuando el
	enfoque de tabla de frecuencias (útil solo para un $k$ fijo) no
	basta porque se necesitan \textit{todos} los valores de $k$ a la
	vez.
	\item Problemas de conteo de subconjuntos con una operación de
	"combinar" dos distribuciones bajo XOR (por ejemplo, combinar
	resultados de dos mitades de un problema dividido con Meet in the
	Middle, cuando la combinación es sobre máscaras).
	\item Convolución AND/OR para problemas de "cuántas formas hay de
	elegir un elemento de $A$ y uno de $B$ tal que su AND/OR sea
	exactamente $k$", usando Transformada Zeta/Möbius en vez de
	Hadamard.
	\item Como herramienta dentro de programación dinámica sobre
	subconjuntos donde la transición natural es una convolución bajo
	alguna operación bitwise, en vez de una suma o producto estándar.
\end{itemize}

\textbf{Complejidad:} $O(n \cdot 2^n)$ para transformar (o
destransformar) un arreglo de tamaño $2^n$, análogo en estructura al
$O(2^n \cdot n)$ de SOS DP: se itera sobre los $n$ bits, y en cada uno
se recorren las $2^n$ posiciones actualizando en $O(1)$. La
convolución completa (transformar $a$, transformar $b$, multiplicar
punto a punto, transformar inverso) cuesta
$O(n \cdot 2^n)$ en total, dominado por las transformadas.

\begin{lstlisting}[style=cpp]
	struct ConvolucionXOR {
		
		// WRITE: aplica la Transformada de Walsh-Hadamard en el lugar
		// (in-place) sobre el arreglo, necesaria antes de multiplicar
		// punto a punto para convolucion XOR
		static void transformarXOR(vector<long long>& a) {
			int n = (int)a.size();
			for (int len = 1; len < n; len <<= 1) {
				for (int i = 0; i < n; i += (len << 1)) {
					for (int j = i; j < i + len; j++) {
						long long u = a[j], v = a[j + len];
						a[j] = u + v;       // combinar par (i, i+len) sin diferencia de signo
						a[j + len] = u - v; // la resta codifica la estructura XOR
					}
				}
			}
		} // O(n log n), con n = 2^bits
		
		// WRITE: aplica la transformada INVERSA de Walsh-Hadamard,
		// deshaciendo transformarXOR y dividiendo por el tamano del
		// arreglo para recuperar los valores originales
		static void transformarXORInversa(vector<long long>& a) {
			int n = (int)a.size();
			transformarXOR(a); // la WHT es su propia inversa salvo el factor de escala
			for (long long& v : a) v /= n;
		} // O(n log n)
		
		// READ: calcula la convolucion XOR completa de a y b: retorna c
		// tal que c[k] = suma de a[i]*b[j] para todo i^j == k
		static vector<long long> convolucionXOR(vector<long long> a, vector<long long> b) {
			int n = (int)a.size(); // debe ser potencia de 2, igual para a y b
			transformarXOR(a);
			transformarXOR(b);
			vector<long long> c(n);
			for (int i = 0; i < n; i++) c[i] = a[i] * b[i]; // multiplicacion punto a punto
			transformarXORInversa(c);
			return c;
		} // O(n log n)
		
		// ------------------------------------------------------------
		// MODIFICACIONES Y COMENTARIOS CON LOS USOS
		// ------------------------------------------------------------
		// 1. Transformada Zeta para convolucion OR (superset/subset):
		//    para convolucion OR, la transformada NO usa resta sino solo
		//    suma en una direccion: a[j+len] += a[j] (equivalente a SOS
		//    DP: suma sobre subconjuntos); la inversa resta en vez de
		//    dividir: a[j+len] -= a[j]. Para convolucion AND, se invierte
		//    el rol de las mitades (se propaga desde el subconjunto mas
		//    grande hacia los mas chicos, "superset sum").
		//
		// 2. Por que no se necesita numeros complejos como en FFT: a
		//    diferencia de la convolucion estandar (que usa raices de la
		//    unidad complejas), la WHT para XOR usa solo +1 y -1 como
		//    "raices", ya que XOR sobre GF(2) tiene una estructura mas
		//    simple; esto permite implementarla enteramente con aritmetica
		//    entera, sin perdida de precision de punto flotante.
		//
		// 3. Overflow en la multiplicacion punto a punto: los valores
		//    transformados pueden crecer hasta un factor de n respecto a
		//    los originales (por las sumas repetidas); si a y b contienen
		//    valores grandes, verificar overflow antes de multiplicar
		//    punto a punto (ver Overflow y Wraparound), o aplicar un
		//    modulo si el problema lo permite.
		//
		// 4. Relacion con SOS DP: la transformada Zeta usada para
		//    convolucion OR es EXACTAMENTE el algoritmo de SOS DP visto
		//    antes, aplicado como paso intermedio de una convolucion en
		//    vez de como resultado final; reconocer esta equivalencia
		//    ayuda a recordar la implementacion sin memorizarla por
		//    separado.
	};
\end{lstlisting}

\textbf{Ejemplo de funcionamiento:}
\begin{lstlisting}[style=cpp]
	// a[i] = cantidad de veces que aparece i en el primer conjunto
	// b[j] = cantidad de veces que aparece j en el segundo conjunto
	vector<long long> a = {1, 2, 0, 1}; // universo de 2 bits (0..3)
	vector<long long> b = {0, 1, 1, 0};
	
	vector<long long> c = ConvolucionXOR::convolucionXOR(a, b);
	// c[k] = suma de a[i]*b[j] para todo i^j == k
	// Nota: el valor exacto de cada c[k] depende de la transformada
	// paso a paso; se debe correr el codigo para confirmar los valores
	// antes de darlos por verificados.
\end{lstlisting}

\textbf{Nota:} Esta técnica es de las más avanzadas de esta parte —
si tu guía está orientada a niveles Div2 o fundamentos, puede quedar
como referencia de "saber que existe" más que como algo a memorizar
en detalle; su prerequisito conceptual directo es SOS DP, así que
conviene tenerla sólida antes de abordar esta subsección.

\textbf{Regla práctica:} Usa convolución XOR/AND/OR (vía WHT o
Transformada Zeta/Möbius) cuando necesites combinar \textit{dos}
arreglos distintos bajo una operación bitwise para \textit{todos} los
valores de salida simultáneamente, en $O(n \cdot 2^n)$; si el problema
es más simple —agregar un solo arreglo sobre sus propios
subconjuntos— con SOS DP alcanza, sin necesidad de la maquinaria
completa de la transformada.
\subsection{Meet in the Middle con Máscaras de Bits}

\textbf{Idea:} Meet in the Middle divide un conjunto de $n$ elementos
en dos mitades de tamaño $n/2$, enumera todos los $2^{n/2}$
subconjuntos de cada mitad por separado (representados como máscaras
de bits), y combina los resultados de ambas mitades para responder la
pregunta sobre el conjunto completo. La observación clave es que
$2^{n/2} + 2^{n/2} = 2^{n/2+1}$ es exponencialmente más pequeño que
$2^n$ cuando $n$ crece: para $n = 40$, enumerar $2^{40}$ subconjuntos
es inviable, pero enumerar $2^{20}$ subconjuntos por mitad (unos
$10^6$) es trivial. La combinación de resultados típicamente requiere
ordenar una mitad y usar búsqueda binaria (o una tabla hash) sobre
ella al recorrer la otra, en vez de comparar cada par de subconjuntos
directamente.

\textbf{¿Por qué usarlo?} Es la técnica que rescata problemas de
Bitmask DP cuando $n$ excede el límite práctico de $2^n$ (usualmente
20-24): al dividir en dos mitades de $n/2 \approx 20$ cada una, se
extiende el rango manejable hasta $n \approx 40$-$45$. Se distingue de
Bitmask DP clásico en que aquí \textit{no} se construye una tabla de
estados sobre todo el conjunto —se enumeran exhaustivamente
subconjuntos de cada mitad de forma independiente, y la dificultad se
traslada a combinar eficientemente ambas listas (usualmente con
ordenamiento + búsqueda binaria, conectando esta técnica directamente
con la Parte de Búsqueda Binaria).

\textbf{Casos de uso:}
\begin{itemize}
	\item Subset Sum con $n$ hasta ~40: dividir en dos mitades,
	enumerar todas las sumas de subconjuntos de cada mitad, ordenar
	una y usar búsqueda binaria sobre la otra para verificar si un
	valor objetivo es alcanzable.
	\item Maximizar el valor de una mochila (\textit{knapsack}) con
	$n$ hasta ~40 y capacidad grande, combinando subconjuntos de cada
	mitad ordenados por peso con un "prefijo máximo" de valor.
	\item Contar cuántos subconjuntos de un arreglo de hasta 40
	elementos satisfacen una condición de suma o XOR, combinando
	conteos de ambas mitades con dos punteros o búsqueda binaria.
	\item Problemas donde el espacio de estados de Bitmask DP
	($2^n$) es demasiado grande para caber en memoria, pero se puede
	reformular como una combinación de dos mitades independientes.
\end{itemize}

\textbf{Complejidad:} $O(2^{n/2} \cdot \log(2^{n/2})) = O(2^{n/2}
\cdot n)$ típicamente: generar las $2^{n/2}$ combinaciones de cada
mitad cuesta $O(2^{n/2} \cdot n/2)$, y ordenar una de las listas para
combinar cuesta $O(2^{n/2} \log(2^{n/2}))$; la combinación final
(recorrer una lista y buscar en la otra) agrega otro factor
$O(2^{n/2} \log(2^{n/2}))$. Esto reemplaza el $O(2^n)$ o $O(2^n \cdot
n)$ inviable de Bitmask DP directo para $n$ grande.

\begin{lstlisting}[style=cpp]
	struct MeetInTheMiddle {
		vector<long long> elementos;
		int n;
		
		// CREATE: guarda el arreglo completo de elementos a dividir en
		// dos mitades
		MeetInTheMiddle(vector<long long> arr) : elementos(arr), n((int)arr.size()) {} // O(n)
		
		// READ: genera todas las sumas de subconjuntos (incluyendo el
		// vacio, suma 0) de los elementos en el rango [inicio, fin)
		vector<long long> generarSumas(int inicio, int fin) {
			vector<long long> resultado;
			int tam = fin - inicio;
			for (int mask = 0; mask < (1 << tam); mask++) {
				long long suma = 0;
				for (int bit = 0; bit < tam; bit++) {
					if (mask & (1 << bit)) suma += elementos[inicio + bit];
				}
				resultado.push_back(suma);
			}
			return resultado;
		} // O(2^tam * tam)
		
		// READ: verifica si existe algun subconjunto del arreglo COMPLETO
		// cuya suma sea exactamente "objetivo", dividiendo el arreglo en
		// dos mitades y combinando con busqueda binaria
		bool existeSubconjuntoConSuma(long long objetivo) {
			int mitad = n / 2;
			vector<long long> sumasIzquierda = generarSumas(0, mitad);
			vector<long long> sumasDerecha = generarSumas(mitad, n);
			
			sort(sumasDerecha.begin(), sumasDerecha.end());
			
			for (long long sumaI : sumasIzquierda) {
				long long complemento = objetivo - sumaI;
				if (binary_search(sumasDerecha.begin(), sumasDerecha.end(), complemento)) {
					return true;
				}
			}
			return false;
		} // O(2^(n/2) * (n/2) + 2^(n/2) * log(2^(n/2)))
		
		// ------------------------------------------------------------
		// MODIFICACIONES Y COMENTARIOS CON LOS USOS
		// ------------------------------------------------------------
		// 1. Contar subconjuntos en vez de verificar existencia: en vez
		//    de binary_search (que solo confirma presencia), usar
		//    equal_range o lower_bound/upper_bound sobre sumasDerecha
		//    ordenada para contar CUANTAS veces aparece el complemento,
		//    acumulando el total sobre todas las sumaI de la izquierda.
		//
		// 2. Knapsack con Meet in the Middle: generar pares
		//    (peso, valor) para cada subconjunto de cada mitad; ordenar
		//    una mitad por peso y precalcular el "maximo valor con peso
		//    <= w" como un prefijo maximo; para cada subconjunto de la
		//    otra mitad, buscar con busqueda binaria el mayor peso
		//    permitido restante y combinar con ese prefijo maximo.
		//
		// 3. Por que dividir por la mitad es optimo: si se dividiera en
		//    proporciones desiguales (por ejemplo, n/3 y 2n/3), el lado
		//    mas grande seguiria dominando la complejidad total
		//    (2^(2n/3) sigue siendo exponencialmente peor que 2^(n/2));
		//    dividir exactamente a la mitad minimiza el maximo de las
		//    dos mitades, que es lo que determina el costo total.
		//
		// 4. Generalizacion a XOR en vez de suma: el mismo esqueleto
		//    aplica reemplazando la suma de subconjunto por XOR de
		//    subconjunto; para buscar el complemento exacto, sigue
		//    funcionando busqueda binaria (o tabla hash) sobre valores
		//    ordenados, ya que la condicion sigue siendo una igualdad.
	};
\end{lstlisting}

\textbf{Ejemplo de funcionamiento:}
\begin{lstlisting}[style=cpp]
	vector<long long> arr = {3, 7, 1, 9, 4, 12, 6, 2}; // n = 8 (ejemplo pequeno)
	MeetInTheMiddle mitm(arr);
	
	bool existe23 = mitm.existeSubconjuntoConSuma(23);
	// Nota: el valor exacto depende de verificar todas las combinaciones
	// de ambas mitades; se debe correr el codigo para confirmarlo antes
	// de darlo por verificado (candidato esperado: true, ya que hay
	// varias combinaciones plausibles que suman 23 en un arreglo con
	// estos valores).
\end{lstlisting}

\textbf{Regla práctica:} Reconoce Meet in the Middle cuando $n$ esté
en el rango de $\sim$30 a $\sim$45 (demasiado grande para Bitmask DP
directo con $2^n$, pero divisible en dos mitades de $2^{n/2}$
manejable) y el problema tenga estructura de "elegir un subconjunto
que satisfaga una condición de suma/XOR/valor"; para $n \leq 20$-$24$,
Bitmask DP directo suele ser más simple y no requiere la complejidad
adicional de combinar dos mitades.

\subsection{Transformada de Möbius/Zeta sobre Subconjuntos}

\textbf{Idea:} La Transformada Zeta sobre subconjuntos transforma un
arreglo $f$ indexado por máscaras en $\hat{f}[m] = \sum_{s \subseteq
	m} f[s]$ — exactamente el algoritmo de SOS DP visto antes, pero
presentado aquí en el lenguaje general de "transformadas" sobre el
retículo de subconjuntos, en analogía directa con la Transformada de
Fourier. La Transformada de Möbius es su \textbf{inversa}: dado
$\hat{f}$, recupera $f$ mediante
$f[m] = \sum_{s \subseteq m} (-1)^{|m|-|s|} \hat{f}[s]$, que en la
práctica se calcula con el mismo bucle bit-por-bit de SOS DP pero
restando en vez de sumar en cada paso. Juntas, estas dos transformadas
son el mecanismo que permite la Convolución AND/OR vista en la
subsección anterior: transformar ambos arreglos con Zeta, multiplicar
punto a punto, y aplicar Möbius para volver al dominio original.

\textbf{¿Por qué usarlo?} Esta subsección formaliza una idea que ya
usaste dos veces sin nombrarla explícitamente como "transformada": SOS
DP calcula la Zeta, y la Convolución AND/OR usa el par Zeta-Möbius
como sus pasos de ida y vuelta. Nombrarla aparte tiene valor
principalmente cuando necesitas la \textbf{inversa} de forma
independiente —por ejemplo, si conoces $\hat{f}$ (una suma sobre
subconjuntos ya agregada) y necesitas recuperar los valores
individuales $f$ originales, en vez de solo ir "hacia adelante" como
hace SOS DP.

\textbf{Casos de uso:}
\begin{itemize}
	\item Recuperar los valores individuales $f[m]$ a partir de una
	suma sobre subconjuntos $\hat{f}[m]$ ya conocida (el paso inverso
	que SOS DP no necesita, pero que Convolución AND/OR sí requiere al
	final).
	\item Aplicar principio de inclusión-exclusión sobre el retículo
	de subconjuntos de forma sistemática, cuando la fórmula de
	inclusión-exclusión directa tiene $2^n$ términos por calcular
	ingenuamente.
	\item Como el paso final de Convolución AND/OR (Transformada Zeta
	para ir al dominio transformado, Möbius para volver), completando
	el patrón "transformar, multiplicar punto a punto, destransformar"
	ya visto con XOR/Walsh-Hadamard.
	\item Verificar la consistencia de un algoritmo de SOS DP:
	aplicar Möbius sobre el resultado de Zeta debe devolver
	exactamente el arreglo original, útil como \textit{sanity check}.
\end{itemize}

\textbf{Complejidad:} $O(n \cdot 2^n)$ tanto para la Transformada Zeta
como para la Möbius —idéntica a SOS DP, ya que la Möbius es
literalmente el mismo bucle con el signo de la actualización invertido
(resta en vez de suma); aplicar ambas en secuencia (ida y vuelta) sigue
siendo $O(n \cdot 2^n)$, no se duplica el orden de crecimiento.

\begin{lstlisting}[style=cpp]
	struct TransformadaMoebiusZeta {
		int n;
		
		// CREATE: guarda el numero de bits del universo de mascaras
		TransformadaMoebiusZeta(int numBits) : n(numBits) {} // O(1)
		
		// WRITE: aplica la Transformada Zeta en el lugar: convierte f[m]
		// (valor base) en la suma sobre todas las submascaras de m
		// (identica al algoritmo de SOS DP)
		void zeta(vector<long long>& f) {
			int total = 1 << n;
			for (int bit = 0; bit < n; bit++) {
				for (int mask = 0; mask < total; mask++) {
					if (mask & (1 << bit)) {
						f[mask] += f[mask ^ (1 << bit)];
					}
				}
			}
		} // O(n * 2^n)
		
		// WRITE: aplica la Transformada de Moebius en el lugar: invierte
		// zeta(), recuperando los valores originales f[m] a partir de la
		// suma sobre subconjuntos
		void moebius(vector<long long>& f) {
			int total = 1 << n;
			for (int bit = 0; bit < n; bit++) {
				for (int mask = 0; mask < total; mask++) {
					if (mask & (1 << bit)) {
						f[mask] -= f[mask ^ (1 << bit)]; // resta en vez de suma: deshace zeta
					}
				}
			}
		} // O(n * 2^n)
		
		// ------------------------------------------------------------
		// MODIFICACIONES Y COMENTARIOS CON LOS USOS
		// ------------------------------------------------------------
		// 1. Por que restar deshace la suma: zeta() propaga cada valor
		//    original hacia todas las mascaras que lo contienen como
		//    submascara, bit a bit; moebius() recorre los mismos bits en
		//    el mismo orden, y en cada paso resta exactamente lo que
		//    zeta() habia sumado en ese paso, deshaciendo la
		//    transformacion bit por bit hasta recuperar los valores base.
		//
		// 2. Version para superset (Zeta/Moebius sobre superconjuntos):
		//    analoga a la version "superset sum" mencionada en SOS DP,
		//    invirtiendo la condicion del if a "!(mask & (1<<bit))" y
		//    operando sobre mask | (1<<bit) en vez de mask ^ (1<<bit).
		//
		// 3. Uso dentro de Convolucion AND/OR: para convolucion OR,
		//    aplicar zeta() a ambos arreglos, multiplicar punto a punto,
		//    y aplicar moebius() al resultado recupera exactamente
		//    c[k] = suma de a[i]*b[j] para todo i|j == k; para AND, se
		//    usa la version "superset" de ambas transformadas.
		//
		// 4. Relacion formal con inclusion-exclusion: la formula
		//    f[m] = suma sobre s submascara de m de (-1)^(|m|-|s|) * fhat[s]
		//    es la version "sobre subconjuntos" del principio de
		//    inclusion-exclusion; el algoritmo bit a bit de moebius() es
		//    simplemente una forma eficiente de evaluar esa suma sin
		//    calcular (-1)^(|m|-|s|) explicitamente para cada par (m, s).
	};
\end{lstlisting}

\textbf{Ejemplo de funcionamiento:}
\begin{lstlisting}[style=cpp]
	vector<long long> original = {1, 2, 3, 4, 5, 6, 7, 8}; // universo de 3 bits
	vector<long long> f = original; // copia para transformar
	
	TransformadaMoebiusZeta t(3);
	t.zeta(f);
	// f ahora contiene, para cada mascara, la suma sobre sus submascaras
	// (identico al resultado de SOS DP)
	
	t.moebius(f);
	// f deberia volver exactamente a "original" tras deshacer zeta()
	// Nota: verificar en codigo que f == original tras este paso, como
	// sanity check de la implementacion.
\end{lstlisting}

\textbf{Regla práctica:} Si solo necesitas ir "hacia adelante" (sumar
sobre subconjuntos una sola vez), usa SOS DP directamente sin pensar
en términos de transformadas; recurre explícitamente a esta subsección
cuando necesites la \textbf{inversa} (recuperar valores originales
desde una agregación) o cuando estés implementando Convolución AND/OR,
donde el par Zeta-Möbius es indispensable como los pasos de ida y
vuelta.

\part{Game Theory}
\subsection{¿Qué algoritmo utilizar?}

\begin{table}[H]
	\raggedright
	\scriptsize
	\setlength{\tabcolsep}{2pt}
	\renewcommand{\arraystretch}{1.1}
	\begin{tabular}{|p{0.43\columnwidth}|p{0.29\columnwidth}|p{0.20\columnwidth}|}
		\hline
		\textbf{Cuando el problema te pida:} &
		\textbf{Entonces usa:} &
		\textbf{Complejidad} \\
		\hline
		
		Determinar si el primer jugador gana o pierde, con un espacio de
		estados pequeño y sin formula conocida &
		Estados Ganadores y Perdedores \newline (Win/Lose States) &
		$O(V+E)$ \\
		\hline
		
		Clasificar un juego nuevo antes de elegir tecnica: ?los movimientos
		dependen solo de la posicion, o tambien de quien mueve? &
		Juegos Imparciales vs. Partidistas \newline (clasificacion previa) &
		sin costo computacional \\
		\hline
		
		Confirmar que el juego no tiene dados, cartas ocultas ni informacion
		privada antes de aplicar cualquier tecnica de esta parte &
		Informacion Perfecta y sin Azar \newline (checklist previo) &
		sin costo computacional \\
		\hline
		
		Resolver un juego propagando el resultado desde las posiciones
		terminales hacia atras (marco general antes de Grundy o minimax) &
		Analisis Retrogrado \newline (Backward Induction) &
		$O(V+E)$ \\
		\hline
		
		$k$ montones de fichas, se quita cualquier cantidad de UN monton,
		pierde quien no puede mover &
		Nim Clasico (Teorema de Bouton) &
		$O(k)$ \\
		\hline
		
		Nim con tope de fichas por turno ("maximo $k$ por turno"), o solo se
		permite quitar cantidades de un conjunto fijo $S$ &
		Nim con Movimientos Restringidos \newline (Grundy por monton + XOR) &
		Build $O(n\cdot|S|)$ \newline Query $O(1)$ \\
		\hline
		
		Calcular el "numero equivalente en Nim" de una posicion, necesario
		antes de combinar el juego con otros &
		Teorema de Sprague-Grundy \newline (mex de sucesores) &
		$O(V+E)$ \\
		\hline
		
		Varios subjuegos independientes jugados a la vez, se elige uno por
		turno (pilas/tableros separados) &
		Suma de Juegos Independientes \newline (XOR de Grundy Numbers) &
		$O(\sum \text{costo}_i) + O(k)$ \\
		\hline
		
		Calcular Grundy de un juego sin formula conocida, dado explicitamente
		como grafo de posiciones &
		Calculo de Grundy sobre un DAG \newline (mex de sucesores) &
		$O(V+E)$ \\
		\hline
		
		Nim pero "quien hace el ultimo movimiento PIERDE" (no el que no
		puede mover) &
		Misere Nim &
		$O(k)$ \\
		\hline
		
		Una sola pila, conjunto de movimientos $S$ dado explicitamente, regla
		de juego normal &
		Juego de Sustraccion \newline (Subtraction Game) &
		$O(n\cdot|S|)$ \\
		\hline
		
		Tope $k$ de piedras por turno pero con un conjunto $F$ de cantidades
		explicitamente prohibidas &
		Sustraccion con Movimientos \newline Prohibidos (complemento de $F$) &
		$O(k) + O(n\cdot|S_{\text{ef}}|)$ \\
		\hline
		
		Monedas en una escalera, se mueven de un escalon a uno
		estrictamente inferior, en cualquier cantidad &
		Staircase Nim \newline (XOR de escalones impares) &
		$O(n)$ \\
		\hline
		
		Dos montones, se puede quitar de uno solo O la misma cantidad de
		ambos a la vez &
		Wythoff's Game \newline (sucesion de Beatty, razon aurea) &
		$O(1)$ por consulta \\
		\hline
		
		Fichas cara/cruz en fila: voltea una cara a cruz (obligatorio) y
		opcionalmente voltea una ficha a su izquierda &
		Turning Games \newline (Coin Turning, XOR por posicion) &
		$O(n)$ \\
		\hline
		
		Una ficha se mueve por un grafo/DAG siguiendo aristas, sin reglas
		con formula conocida; posiblemente varias fichas independientes &
		Juegos sobre Grafos y DAGs \newline (Grundy generico + XOR) &
		$O(V+E)$ \\
		\hline
		
		Arbol de ramas colgando del suelo, se corta una arista y cae todo lo
		que quedaba colgando de ella &
		Green Hackenbush \newline (Colon Principle) &
		$O(V)$ \\
		\hline
		
		Dos jugadores toman el extremo izquierdo o derecho de un arreglo,
		cada uno maximiza SU propio puntaje (no hay "ganar/perder" binario) &
		DP de Intervalos \newline (Optimal Strategy for a Game) &
		$O(n^2)$, $O(n)$ memoria \\
		\hline
		
		Un movimiento elimina parte de una cadena/secuencia y las dos
		mitades restantes quedan separadas para siempre (no se concatenan) &
		Suma de Grundy sobre Substrings \newline (particion en subjuegos) &
		$O(n^3)$ \\
		\hline
		
		Izquierda y Derecha tienen jugadas DISTINTAS en una misma posicion
		(Hackenbush azul-rojo, Domineering) &
		Numeros Surreales \newline (valores $\{G^L\mid G^R\}$, cold/hot/fuzzy) &
		$O(m_1\cdot m_2)$ por comparacion \\
		\hline
		
		Varios componentes partidistas mezclados (numeros, calientes,
		difusos) que hay que combinar y clasificar contra 0 &
		Suma de Juegos Partidistas \newline (suma disyunta + \texttt{leq}) &
		$O(\text{producto de tamanos})$ \\
		\hline
		
		El estado del juego no encaja en ninguna estructura simple (tablero
		arbitrario), pero hay muchas transposiciones &
		Minimax con Memoizacion &
		$O(E\cdot b)$ \\
		\hline
		
		El arbol de busqueda de minimax es demasiado grande incluso con
		memoizacion &
		Poda Alfa-Beta \newline (+ tabla de transposicion opcional) &
		$O(b^{d/2})$ con buen orden \\
		\hline
		
		Juego de restar con regla fija, pero $n$ es enorme ($\sim 10^{18}$) y
		hace falta extrapolar Grundy &
		Deteccion de Periodo \newline (ventana de $D$ valores) &
		$O(N\cdot D\log N)$ \\
		\hline
		
		Tres o mas jugadores con turnos rotativos, cada uno con su propio
		objetivo/puntaje &
		Juegos con Multiples Jugadores \newline ($\max^n$, vector de resultados) &
		$O(N\,K^2 M)$ \\
		\hline
		
	\end{tabular}
	\caption{Guía rápida: qué algoritmo de Game Theory usar según el problema, con su complejidad.}
\end{table}
\section{Fundamentos}
\subsection{Estados Ganadores y Perdedores (Win/Lose States)}

\textbf{Idea:} en un juego de dos jugadores con información perfecta y sin azar, cada posición del juego se puede clasificar como \textbf{ganadora} (el jugador que está por mover puede forzar la victoria) o \textbf{perdedora} (haga lo que haga, pierde si el rival juega óptimo). La clasificación se define de atrás hacia adelante: una posición \textbf{terminal} donde el jugador en turno ya perdió (por ejemplo, no quedan movimientos) es \textbf{perdedora}. Una posición no terminal es \textbf{ganadora} si existe \textbf{al menos un} movimiento hacia una posición perdedora para el rival; es \textbf{perdedora} si \textbf{todos} los movimientos posibles llevan a posiciones ganadoras para el rival. Esta es exactamente la misma memoización de \textit{Memoización vs. Tabulación} aplicada a un grafo de posiciones en vez de a números.

\textbf{¿Por qué usarlo?} Es la base de toda la parte: antes de aplicar Sprague-Grundy o cualquier fórmula cerrada, hay que entender que un juego se puede resolver por fuerza bruta memoizada siempre que el número de posiciones distintas sea manejable. Sin memoización, evaluar recursivamente si una posición es ganadora reevalúa las mismas subposiciones una y otra vez, con costo exponencial. Con memoización, cada posición se evalúa una sola vez y el costo total es proporcional al número de posiciones por el número de movimientos desde cada una.

\textbf{Casos de uso:}
\begin{itemize}
	\item Determinar si el primer jugador gana o pierde en un juego de turnos dado un estado inicial concreto.
	\item Como paso previo obligatorio antes de calcular números de Grundy: primero se identifica la posición terminal perdedora, luego se generaliza.
	\item Juegos donde el espacio de estados es pequeño (por ejemplo, montones de tamaño acotado) y no hace falta ninguna fórmula cerrada.
	\item Verificar a mano, con esta memoización, una fórmula sospechada para un juego más grande, igual que en \textit{Fuerza Bruta para Descubrir Patrones}.
\end{itemize}

\textbf{Complejidad:} $O(V+E)$ donde $V$ es el número de posiciones distintas del juego y $E$ el número total de movimientos entre posiciones, ya que cada posición se evalúa una sola vez y cada movimiento se recorre una sola vez.

\begin{lstlisting}[style=cpp]
	// CREATE/READ: Clasificacion de posiciones como ganadoras o perdedoras
	// mediante memoizacion, para un juego de sustraccion generico: dos
	// jugadores se turnan para restar un valor de un conjunto S al monton n;
	// pierde quien no puede mover (n llega a 0 sin poder restar nada valido)
	struct EstadosGanadorPerdedor {
		vector<int> movimientosValidos; // S: valores que se pueden restar
		vector<int> memo; // memo[n] = -1 (sin calcular), 0 (perdedora), 1 (ganadora)
		
		// CREATE: prepara la tabla de memoizacion con centinela -1
		EstadosGanadorPerdedor(const vector<int>& S, int nMax) : movimientosValidos(S) {
			memo.assign(nMax + 1, -1);
		} // O(nMax)
		
		// READ: true si la posicion n es GANADORA para el jugador en turno
		bool esGanadora(int n) {
			if (n == 0) return false; // posicion terminal: quien mueve aqui ya perdio
			if (memo[n] != -1) return memo[n];
			bool ganadora = false;
			for (int m : movimientosValidos) {
				if (m <= n && !esGanadora(n - m)) {
					ganadora = true;
					break; // basta UN movimiento a una posicion perdedora para el rival
				}
			}
			return memo[n] = ganadora;
		} // O(1) amortizado por estado, O(n * |S|) en total
		
		// ------------------------------------------------------------
		// MODIFICACIONES Y COMENTARIOS CON LOS USOS
		// ------------------------------------------------------------
		// 1. Version tabulada (bottom-up): recorrer n de 0 a nMax en orden
		//    creciente, evaluando memo[n] con el mismo criterio, sin
		//    recursion; preferible si nMax es grande y hay riesgo de
		//    desbordar la pila.
		// 2. Estado compuesto (mas de un numero, por ejemplo varios montones):
		//    usar un mapa o una mascara de bits como clave en vez de un
		//    arreglo indexado por un solo entero; el resto de la logica no
		//    cambia.
		// 3. Reconstruir la ESTRATEGIA optima (no solo si se gana): guardar,
		//    junto con memo[n], el movimiento m que lleva a la posicion
		//    perdedora encontrada, para poder reproducir la jugada ganadora.
		// 4. Ultimo en mover PIERDE en vez de ganar (Misere): cambiar la
		//    condicion terminal de "n == 0 es perdedora" a "n == 0 es
		//    ganadora" (el jugador en turno gana porque el rival ya no puede
		//    mover y esa regla favorece a quien esta en turno); revisar con
		//    cuidado el enunciado exacto de cada problema.
		// 5. Multiples jugadores (mas de 2): el analisis binario
		//    ganadora/perdedora deja de aplicar directamente; se necesita
		//    memoizar el resultado para CADA jugador o usar backward
		//    induction con mas de dos resultados posibles.
	};
\end{lstlisting}

\textbf{Ejemplo de funcionamiento:}
\begin{lstlisting}[style=cpp]
	EstadosGanadorPerdedor g({1, 3, 4}, 10);
	
	bool r1 = g.esGanadora(1);
	// r1 = true (restar 1 deja n=0, posicion perdedora para el rival)
	
	bool r2 = g.esGanadora(2);
	// r2 = false
	// (desde 2 solo se puede restar 1, dejando n=1, que es GANADORA para
	//  el rival; no hay otro movimiento valido, asi que 2 es perdedora)
	
	bool r3 = g.esGanadora(5);
	// r3 = true (restar 4 deja n=1, que es ganadora para el rival... hay
	//  que correrlo para confirmar cual movimiento concreto deja al rival
	//  en una posicion perdedora)
\end{lstlisting}

\textbf{Nota:} la regla de clasificación (ganadora si \textbf{existe} un movimiento a perdedora; perdedora si \textbf{todos} los movimientos van a ganadora) es la base de todo análisis de juegos combinatorios de esta parte, incluyendo Sprague-Grundy: un número de Grundy distinto de cero es exactamente la generalización de "posición ganadora", y Grundy igual a cero es exactamente "posición perdedora".

\textbf{Regla práctica:} antes de buscar una fórmula cerrada para un juego nuevo, escribe primero esta memoización sobre casos pequeños. Si el patrón de ganadora/perdedora no es evidente a simple vista, genera una tabla de varios valores de $n$ (como en \textit{Fuerza Bruta para Descubrir Patrones}) y busca la periodicidad o condición que los distingue antes de intentar demostrar una fórmula general.

\subsection{Juegos Imparciales vs. Partidistas}

\textbf{Idea:} un juego es \textbf{imparcial} (\textit{impartial}) si, en cualquier posición, \textbf{ambos jugadores tienen exactamente el mismo conjunto de movimientos disponibles} — el conjunto de jugadas depende solo de la posición, no de quién esté jugando. Nim es el ejemplo canónico: cualquiera de los dos jugadores puede quitar fichas de cualquier montón. Un juego es \textbf{partidista} (\textit{partizan}) si los movimientos disponibles dependen de \textbf{qué jugador} está en turno — por ejemplo, el ajedrez, donde un jugador solo mueve piezas blancas y el otro solo negras. Esta distinción determina qué herramientas aplican: el teorema de Sprague-Grundy, con sus números de Grundy y la suma por XOR, \textbf{solo es válido para juegos imparciales}.

\textbf{¿Por qué usarlo?} Sin identificar primero si un juego es imparcial o partidista, se corre el riesgo de aplicar Grundy donde no corresponde. En un juego partidista no existe un único "número" que resuma una posición, porque el resultado depende de quién mueve; ahí se necesita la teoría más general de \textbf{valores de juego} (números surreales, juegos \textit{cold}, \textit{hot}, \textit{fuzzy}), vista en \textit{Números Surreales y Valores de Juego}. Reconocer la categoría del juego en la primera lectura del enunciado ahorra tiempo: evita intentar forzar un XOR de Grundy sobre un juego donde simplemente no aplica.

\textbf{Casos de uso:}
\begin{itemize}
	\item Clasificar un juego nuevo antes de elegir la técnica de resolución: ¿los movimientos dependen de la posición solamente, o también de qué jugador mueve?
	\item Nim, juegos de sustracción, y la mayoría de juegos de "tomar de un conjunto compartido" son imparciales.
	\item Juegos donde cada jugador controla piezas o recursos distintos (como Hackenbush partidista, o ajedrez) son partidistas.
	\item Decidir si conviene calcular números de Grundy (imparcial) o recurrir a análisis de \textit{minimax} general (partidista), visto en \textit{Minimax con Memoización}.
\end{itemize}

\textbf{Complejidad:} no aplica una complejidad de algoritmo — es una clasificación conceptual previa al análisis, sin costo computacional propio.

\begin{lstlisting}[style=cpp]
	// CREATE/READ: Verificacion PROGRAMATICA de si un juego, dado por una
	// funcion de movimientos parametrizada por jugador, es imparcial: se
	// comprueba que ambos jugadores tengan el MISMO conjunto de movimientos
	// en un conjunto de posiciones de prueba
	struct VerificadorImparcial {
		// READ: dada una funcion movimientos(posicion, jugador) que retorna
		// las posiciones alcanzables, verifica si coincide para ambos
		// jugadores en cada posicion de prueba dada
		static bool esImparcial(
		const vector<int>& posicionesDePrueba,
		function<vector<int>(int, int)> movimientos
		) {
			for (int pos : posicionesDePrueba) {
				vector<int> m0 = movimientos(pos, 0);
				vector<int> m1 = movimientos(pos, 1);
				sort(m0.begin(), m0.end());
				sort(m1.begin(), m1.end());
				if (m0 != m1) return false; // los movimientos difieren segun el jugador
			}
			return true; // coincide en todas las posiciones probadas (no es una demostracion)
		} // O(k * costo de movimientos()), k = numero de posiciones probadas
		
		// ------------------------------------------------------------
		// MODIFICACIONES Y COMENTARIOS CON LOS USOS
		// ------------------------------------------------------------
		// 1. Ejemplo de juego IMPARCIAL: sustraccion de un monton n con
		//    S = {1,2,3}; movimientos(n, jugador) devuelve {n-1,n-2,n-3}
		//    intersectado con validos, IGUAL sin importar el jugador.
		// 2. Ejemplo de juego PARTIDISTA: cada jugador tiene fichas propias en
		//    un tablero; movimientos(pos, 0) solo mueve fichas del jugador 0,
		//    movimientos(pos, 1) solo las del jugador 1 -- ahi esImparcial()
		//    devuelve false de inmediato.
		// 3. Esta verificacion es SOLO una herramienta de deteccion rapida
		//    sobre casos de prueba, no una demostracion formal; en la
		//    practica, la clasificacion imparcial vs. partidista se hace
		//    leyendo la definicion de movimientos del enunciado directamente.
		// 4. Juegos hibridos: algunos juegos tienen una fase imparcial y otra
		//    partidista (o se descomponen en subjuegos de distinto tipo); en
		//    ese caso se analiza cada componente por separado antes de
		//    combinar resultados.
	};
\end{lstlisting}

\textbf{Ejemplo de funcionamiento:}
\begin{lstlisting}[style=cpp]
	// juego imparcial: sustraccion con S = {1,2,3}, ambos jugadores mueven igual
	auto movImparcial = [](int n, int jugador) {
		vector<int> res;
		for (int m : {1, 2, 3}) if (m <= n) res.push_back(n - m);
		return res;
	};
	bool r1 = VerificadorImparcial::esImparcial({1,2,3,4,5}, movImparcial);
	// r1 = true
	
	// juego partidista simplificado: el jugador 0 solo puede restar valores
	// pares, el jugador 1 solo impares
	auto movPartidista = [](int n, int jugador) {
		vector<int> res;
		for (int m = 1; m <= n; m++) {
			if (jugador == 0 && m % 2 == 0) res.push_back(n - m);
			if (jugador == 1 && m % 2 == 1) res.push_back(n - m);
		}
		return res;
	};
	bool r2 = VerificadorImparcial::esImparcial({4,5,6}, movPartidista);
	// r2 = false (los conjuntos de movimientos difieren segun el jugador)
\end{lstlisting}

\textbf{Regla práctica:} antes de escribir una sola línea de código para resolver un juego nuevo, pregúntate si el conjunto de movimientos depende de quién está en turno. Si la respuesta es no (imparcial), el camino es \textit{Sprague-Grundy}. Si la respuesta es sí (partidista), Grundy no aplica y hace falta la teoría de valores de juego o un análisis de \textit{minimax} directo sobre el árbol de posiciones.
\subsection{Juegos con Información Perfecta y sin Azar}

\textbf{Idea:} toda la teoría de esta parte —estados ganador/perdedor, Sprague-Grundy, minimax— exige dos condiciones sobre el juego. \textbf{Información perfecta} significa que ambos jugadores conocen en todo momento el estado completo del juego: no hay cartas ocultas, fichas boca abajo, ni información privada de un jugador que el otro desconozca. \textbf{Sin azar} significa que no hay dados, barajado aleatorio, ni ningún elemento probabilístico que decida parte de un movimiento: cada jugada produce un resultado \textbf{determinista}. Bajo estas dos condiciones, el juego se puede representar como un árbol (o grafo, si hay posiciones repetibles) finito y completamente conocido, y cada posición admite una clasificación fija (ganadora, perdedora, o un valor de Grundy) que no cambia según lo que "podría pasar".

\textbf{¿Por qué usarlo?} Sin información perfecta (por ejemplo, un jugador no sabe qué cartas tiene el otro) o sin ausencia de azar (una tirada de dado decide cuánto se puede mover), no existe una única clasificación de la posición: la mejor jugada depende de probabilidades o de información desconocida, y el análisis correcto pasa a ser de \textbf{teoría de juegos con información imperfecta} o de \textbf{procesos de decisión de Markov} (esperanza de resultado en vez de ganador/perdedor determinista) — fuera del alcance de las técnicas de esta parte. Reconocer esta condición al leer el enunciado es el primer filtro: si aparecen dados, cartas ocultas, o "elige un jugador al azar", ninguna de las subsecciones de win/lose o Grundy aplica directamente.

\textbf{Casos de uso:}
\begin{itemize}
	\item Confirmar, antes de intentar cualquier técnica de esta parte, que el juego del enunciado no involucra azar ni información oculta.
	\item Ajedrez, Nim, y la gran mayoría de juegos de mesa de turnos alternados sin dados son de información perfecta y sin azar.
	\item Juegos con una tirada de dado que determina cuántas fichas se mueven no entran en esta categoría: requieren valor esperado, no un análisis win/lose puro.
	\item Juegos de cartas donde un jugador no ve la mano del otro no son de información perfecta, aunque no haya azar en el sentido de dados.
\end{itemize}

\textbf{Complejidad:} no aplica — es, como \textit{Juegos Imparciales vs. Partidistas}, una clasificación previa sin costo computacional propio, pero que determina si el resto de las técnicas de esta parte son aplicables.

\begin{lstlisting}[style=cpp]
	// CREATE/READ: Lista de verificacion programatica (checklist) para
	// confirmar que un juego cumple las condiciones necesarias antes de
	// aplicar win/lose states, Sprague-Grundy o minimax
	struct VerificadorJuegoDeterminista {
		// READ: retorna una lista de banderas de alerta detectadas en la
		// descripcion textual (o en los parametros) del juego; una lista NO
		// vacia significa que el juego probablemente NO cumple las
		// condiciones (informacion perfecta y sin azar)
		static vector<string> detectarAlertas(bool hayDados, bool hayCartasOcultas,
		bool hayEleccionAleatoria, bool hayInformacionPrivada) {
			vector<string> alertas;
			if (hayDados) alertas.push_back("Hay dados: el juego tiene azar, no es determinista.");
			if (hayCartasOcultas) alertas.push_back("Hay cartas ocultas: no es informacion perfecta.");
			if (hayEleccionAleatoria) alertas.push_back("Hay eleccion aleatoria de jugador o de turno: hay azar.");
			if (hayInformacionPrivada) alertas.push_back("Hay informacion privada de un jugador: no es informacion perfecta.");
			return alertas;
		} // O(1)
		
		// ------------------------------------------------------------
		// MODIFICACIONES Y COMENTARIOS CON LOS USOS
		// ------------------------------------------------------------
		// 1. Juegos con azar pero con estrategia OPTIMA en expectativa: si el
		//    problema realmente exige un dado, se resuelve con DP sobre valor
		//    esperado (dp[estado] = promedio ponderado por probabilidad de
		//    dp[siguiente estado]), un enfoque distinto y no cubierto en el
		//    resto de esta parte.
		// 2. Informacion oculta pero con distribucion conocida: se modela con
		//    teoria de juegos bayesiana; tampoco es el alcance de esta guia,
		//    se menciona solo para que se reconozca la categoria del problema.
		// 3. Esta lista de verificacion es una ayuda de LECTURA del enunciado,
		//    no un algoritmo que se ejecute sobre datos del problema; sirve
		//    como recordatorio de que preguntas hacerse antes de programar.
	};
\end{lstlisting}

\textbf{Ejemplo de funcionamiento:}
\begin{lstlisting}[style=cpp]
	auto alertasNim = VerificadorJuegoDeterminista::detectarAlertas(false, false, false, false);
	// alertasNim esta vacia: Nim es de informacion perfecta y sin azar,
	// las tecnicas de esta parte aplican directamente
	
	auto alertasDado = VerificadorJuegoDeterminista::detectarAlertas(true, false, false, false);
	// alertasDado contiene un mensaje: "Hay dados..."
	// (este juego necesita un enfoque de valor esperado, no win/lose directo)
\end{lstlisting}

\textbf{Regla práctica:} antes de clasificar un juego como imparcial o partidista, confirma primero que cumple información perfecta y ausencia de azar — sin estas dos condiciones, ni Sprague-Grundy ni el análisis win/lose de \textit{Estados Ganadores y Perdedores} son aplicables, sin importar qué tan imparcial parezca el conjunto de movimientos.

\subsection{Análisis Retrógrado (Backward Induction)}

\textbf{Idea:} el análisis retrógrado es el \textbf{método general} detrás de \textit{Estados Ganadores y Perdedores}: en vez de razonar hacia adelante desde la posición inicial (que puede tener ramificaciones enormes), se razona \textbf{hacia atrás} desde las posiciones \textbf{terminales} (donde el juego ya terminó y el resultado es conocido sin ambigüedad), y se propaga esa información hacia las posiciones que llevan a ellas. Es exactamente memoización o tabulación sobre un grafo de estados, donde la "recurrencia" no es aritmética sino lógica: el resultado de una posición se determina por los resultados de sus posiciones sucesoras, nunca al revés.

\textbf{¿Por qué usarlo?} Intentar razonar hacia adelante ("si yo juego esto, y luego el rival juega aquello, y luego yo...") se vuelve intratable rápidamente porque el árbol de posibilidades crece exponencialmente con el número de turnos. El análisis retrógrado invierte el problema: cada posición terminal tiene un resultado \textbf{trivialmente conocido}, y ese conocimiento se propaga hacia atrás un paso a la vez, de modo que cada posición se resuelve usando únicamente resultados ya calculados de posiciones "más cercanas al final". Es el mismo principio que sustenta toda la parte de Programación Dinámica, aplicado aquí específicamente a juegos.

\textbf{Casos de uso:}
\begin{itemize}
	\item Fundamento metodológico de \textit{Estados Ganadores y Perdedores}, \textit{Cálculo de Grundy sobre un DAG}, y \textit{Minimax con Memoización}.
	\item Juegos donde las posiciones forman un grafo acíclico dirigido (DAG): cada movimiento reduce estrictamente algún recurso (fichas restantes, tamaño de la cadena), garantizando que no hay ciclos y el análisis retrógrado termina.
	\item Resolver juegos de mesa simplificados (tres en línea en tableros pequeños, por ejemplo) enumerando todas las posiciones y propagando el resultado desde los finales de partida.
	\item Verificar manualmente, para tableros o montones pequeños, el patrón que luego se generaliza a una fórmula cerrada.
\end{itemize}

\textbf{Complejidad:} $O(V+E)$, igual que \textit{Estados Ganadores y Perdedores}: cada posición ($V$) se visita una sola vez y cada movimiento ($E$) se recorre una sola vez, sin importar si se implementa con memoización recursiva o con tabulación siguiendo un orden topológico.

\begin{lstlisting}[style=cpp]
	// CREATE/READ: Analisis retrogrado generico sobre un DAG de posiciones,
	// dado explicitamente como lista de adyacencia (posicion -> sucesores);
	// determina el ganador de cada posicion mediante orden topologico inverso
	struct AnalisisRetrogrado {
		int n; // numero de posiciones, indexadas 0..n-1
		vector<vector<int>> sucesores; // sucesores[i] = posiciones alcanzables desde i
		vector<int> resultado; // resultado[i] = 0 (perdedora) o 1 (ganadora); -1 = sin calcular
		
		// CREATE: reserva las estructuras para n posiciones
		AnalisisRetrogrado(int numPosiciones) : n(numPosiciones) {
			sucesores.assign(n, {});
			resultado.assign(n, -1);
		} // O(n)
		
		// WRITE: agrega un movimiento valido de la posicion 'desde' hacia 'hacia'
		void agregarMovimiento(int desde, int hacia) {
			sucesores[desde].push_back(hacia);
		} // O(1) amortizado
		
		// READ: clasifica TODAS las posiciones mediante orden topologico
		// inverso (procesa primero las posiciones sin sucesores, es decir,
		// terminales, y propaga hacia atras); requiere que el grafo sea un DAG
		void resolverTodas() {
			vector<int> gradoSalida(n);
			for (int i = 0; i < n; i++) gradoSalida[i] = (int)sucesores[i].size();
			
			// lista de predecesores para propagar en sentido inverso
			vector<vector<int>> predecesores(n);
			for (int i = 0; i < n; i++)
			for (int j : sucesores[i]) predecesores[j].push_back(i);
			
			queue<int> cola;
			for (int i = 0; i < n; i++) {
				if (gradoSalida[i] == 0) {
					resultado[i] = 0; // posicion terminal: quien mueve aqui ya perdio
					cola.push(i);
				}
			}
			
			vector<int> sucesoresPerdedoresPendientes(n);
			for (int i = 0; i < n; i++) sucesoresPerdedoresPendientes[i] = gradoSalida[i];
			
			while (!cola.empty()) {
				int actual = cola.front(); cola.pop();
				for (int p : predecesores[actual]) {
					if (resultado[p] != -1) continue; // ya resuelto
					if (resultado[actual] == 0) {
						// p tiene un movimiento hacia una posicion perdedora: p es ganadora
						resultado[p] = 1;
						cola.push(p);
					} else {
						// este sucesor de p es ganador; si TODOS lo son, p sera perdedora
						sucesoresPerdedoresPendientes[p]--;
						if (sucesoresPerdedoresPendientes[p] == 0) {
							resultado[p] = 0;
							cola.push(p);
						}
					}
				}
			}
		} // O(V + E)
		
		// READ: true si la posicion i es ganadora (llamar resolverTodas antes)
		bool esGanadora(int i) const {
			return resultado[i] == 1;
		} // O(1)
		
		// ------------------------------------------------------------
		// MODIFICACIONES Y COMENTARIOS CON LOS USOS
		// ------------------------------------------------------------
		// 1. Version recursiva con memoizacion (mas simple de escribir, igual
		//    complejidad, pero con riesgo de desbordar la pila si el DAG es
		//    profundo): identica a la de Estados Ganadores y Perdedores, con
		//    sucesores explicitos en vez de calcularlos con una formula.
		// 2. Posiciones con CICLOS (empates o repeticion infinita posible): el
		//    analisis retrogrado puro no basta; se necesita una tercera
		//    categoria (EMPATE/DIBUJO) y un algoritmo como el de Zermelo
		//    generalizado, que primero resuelve las posiciones alcanzables en
		//    un numero finito de pasos y deja el resto como empate.
		// 3. Grafos MUY grandes (posiciones generadas implicitamente, no
		//    enumeradas de antemano): generar sucesores() bajo demanda dentro
		//    de una memoizacion recursiva, en vez de construir sucesores y
		//    predecesores explicitos como en resolverTodas().
		// 4. Generalizacion a numeros de Grundy (no solo ganador/perdedor): en
		//    vez de resultado[i] booleano, calcular
		//    grundy[i] = mex de {grundy[j] : j en sucesores[i]}; resultado[i]
		//    ganadora es equivalente a grundy[i] != 0.
	};
\end{lstlisting}

\textbf{Ejemplo de funcionamiento:}
\begin{lstlisting}[style=cpp]
	AnalisisRetrogrado g(5); // posiciones 0..4, donde 0 es terminal
	g.agregarMovimiento(1, 0);
	g.agregarMovimiento(2, 1);
	g.agregarMovimiento(3, 1);
	g.agregarMovimiento(3, 2);
	g.agregarMovimiento(4, 3);
	
	g.resolverTodas();
	
	bool r0 = g.esGanadora(0);
	// r0 = false (posicion terminal: quien mueve aqui ya perdio)
	
	bool r1 = g.esGanadora(1);
	// r1 = true (mueve hacia 0, que es perdedora)
	
	bool r2 = g.esGanadora(2);
	// r2 = false (su unico sucesor es 1, que es ganadora para el rival)
	
	bool r3 = g.esGanadora(3);
	// r3 = true (mueve hacia 2, que es perdedora)
\end{lstlisting}

\textbf{Nota:} este análisis requiere que el grafo de posiciones sea un \textbf{DAG} (sin ciclos) para garantizar que el orden topológico inverso existe y el proceso termina. Si el juego permite volver a una posición ya visitada (ciclos), hace falta una tercera categoría de resultado (empate o posición indeterminada) y una versión más elaborada del algoritmo, como se indica en la modificación 2 dentro del código.

\textbf{Regla práctica:} usa análisis retrógrado como el marco mental general para cualquier juego de esta parte: identifica primero las posiciones terminales (con resultado obvio), y propaga desde ahí. Es la misma idea exacta que \textit{Estados Ganadores y Perdedores}, aquí generalizada a un grafo explícito con orden topológico; cuando se calculen números de Grundy más adelante, es este mismo mecanismo con \texttt{mex} en vez de una simple disyunción booleana.

\section{Nim y Teoría de Sprague-Grundy}
\subsection{Nim Clásico (Bouton's Theorem)}

\textbf{Idea:} en Nim hay $k$ montones de fichas, con tamaños $a_1,\dots,a_k$. En cada turno, un jugador elige \textbf{un} montón y quita de él \textbf{al menos una} ficha (tantas como quiera, hasta vaciarlo). Pierde quien no puede mover, es decir, quien enfrenta todos los montones en $0$. El teorema de Bouton dice que la posición inicial es \textbf{perdedora} para el jugador que mueve primero si y solo si $a_1\oplus a_2\oplus\cdots\oplus a_k=0$ (XOR de todos los tamaños), y \textbf{ganadora} en caso contrario. La demostración usa exactamente el criterio de \textit{Estados Ganadores y Perdedores}: se comprueba que desde cualquier posición con XOR distinto de $0$ existe un movimiento hacia XOR $=0$, y que desde cualquier posición con XOR $=0$ todo movimiento lleva a XOR distinto de $0$ — las dos propiedades que caracterizan, por inducción, cuál lado gana.

\textbf{¿Por qué usarlo?} Sin el teorema, decidir quién gana requeriría memoización sobre el producto de los tamaños de los montones, inviable si cada $a_i$ es grande (por ejemplo $10^9$). Bouton reduce el problema a una sola operación XOR en $O(k)$, sin importar cuán grandes sean los montones. Es además el caso más simple —y la motivación original— del \textit{Teorema de Sprague-Grundy}: el número de Grundy de un montón de tamaño $n$ en Nim es exactamente $n$, así que el XOR de Grundy de la suma de juegos coincide con el XOR directo de los tamaños.

\textbf{Casos de uso:}
\begin{itemize}
	\item Determinar quién gana Nim clásico con montones de cualquier tamaño, en $O(k)$.
	\item Encontrar el movimiento ganador concreto desde una posición ganadora (no solo si se gana o se pierde).
	\item Como caso base para verificar la implementación de \textit{Sprague-Grundy} en juegos más generales: cuando el juego se reduce a Nim puro, la fórmula debe coincidir con el XOR directo.
	\item Punto de partida al reconocer que un juego nuevo "es Nim disfrazado" tras encontrar su número de Grundy por montón.
\end{itemize}

\textbf{Complejidad:} $O(k)$ para determinar el ganador, con $k$ el número de montones; $O(k)$ también para encontrar el movimiento ganador concreto.

\begin{lstlisting}[style=cpp]
	// CREATE/READ: Nim clasico: determina el ganador via XOR (teorema de
	// Bouton) y encuentra el movimiento ganador concreto si existe
	struct NimClasico {
		vector<long long> montones;
		
		// CREATE: guarda los tamanos de los montones
		NimClasico(const vector<long long>& a) : montones(a) {
		} // O(k)
		
		// READ: XOR de todos los montones (nim-suma)
		long long nimSuma() const {
			long long x = 0;
			for (long long a : montones) x ^= a;
			return x;
		} // O(k)
		
		// READ: true si el jugador que mueve PRIMERO gana
		bool primerJugadorGana() const {
			return nimSuma() != 0;
		} // O(k)
		
		// READ: encuentra un movimiento ganador: indice del monton a reducir
		// y su nuevo tamano; retorna {-1, -1} si la posicion ya es perdedora
		pair<int, long long> movimientoGanador() const {
			long long x = nimSuma();
			if (x == 0) return {-1, -1}; // posicion perdedora: no hay movimiento ganador
			for (int i = 0; i < (int)montones.size(); i++) {
				long long nuevo = montones[i] ^ x;
				// un movimiento valido debe REDUCIR el monton, no aumentarlo
				if (nuevo < montones[i]) return {i, nuevo};
			}
			return {-1, -1}; // no deberia ocurrir si x != 0
		} // O(k)
		
		// ------------------------------------------------------------
		// MODIFICACIONES Y COMENTARIOS CON LOS USOS
		// ------------------------------------------------------------
		// 1. Por que nuevo = montones[i] ^ x funciona: se busca el monton i
		//    tal que reducirlo a montones[i] ^ x deje la nim-suma total en 0;
		//    esto es valido solo cuando montones[i] ^ x < montones[i], porque
		//    el movimiento en Nim SOLO puede quitar fichas, nunca agregar.
		//    Existe garantizado al menos un i que cumple esto cuando x != 0,
		//    por construccion del bit mas significativo de x.
		// 2. Verificacion manual del bit mas significativo: si se quiere
		//    encontrar el monton ganador sin probar los k montones uno por
		//    uno, se ubica el bit mas alto de x y se busca el primer monton
		//    que tenga ese bit encendido; ese monton es siempre valido.
		// 3. Version con MUCHOS montones y consultas de nim-suma tras
		//    actualizaciones (montones cambian de tamano dinamicamente): usar
		//    un multiset o un arbol de Fenwick por bits para recalcular el
		//    XOR total en O(log(max valor)) por actualizacion, en vez de
		//    recorrer todos los montones cada vez.
		// 4. Nim con UN SOLO monton: la nim-suma es simplemente ese valor;
		//    el primer jugador gana siempre que el monton no este vacio
		//    (caso trivial, pero verifica la formula general).
	};
\end{lstlisting}

\textbf{Ejemplo de funcionamiento:}
\begin{lstlisting}[style=cpp]
	NimClasico n({3, 4, 5});
	
	long long x = n.nimSuma();
	// x = 2
	// (3 XOR 4 XOR 5 = 2)
	
	bool gana = n.primerJugadorGana();
	// gana = true (la nim-suma es distinta de 0)
	
	auto mov = n.movimientoGanador();
	// mov = {0, 1}
	// (reducir el monton de indice 0, de 3 a 1: 3^2=1, y 1 < 3, valido;
	//  tras el movimiento la posicion queda {1,4,5}, con nim-suma
	//  1^4^5 = 0, perdedora para el rival)
\end{lstlisting}

\textbf{Nota:} el movimiento ganador \textbf{no es único} en general — puede haber más de un montón desde el que se puede llegar a nim-suma $0$; \texttt{movimientoGanador()} retorna el primero que encuentra, cualquiera de los válidos sirve igual de bien porque todos dejan al rival en una posición perdedora.

\textbf{Regla práctica:} usa el XOR de Bouton directamente para Nim clásico (montones independientes, quitar cualquier cantidad de un solo montón). Si el juego impone restricciones sobre \textbf{cuánto} se puede quitar de un montón, no es Nim puro: revisa \textit{Nim con Montones y Movimientos Restringidos} o calcula el número de Grundy de un montón individual con \textit{Teorema de Sprague-Grundy} antes de aplicar cualquier XOR.

\subsection{Nim con Montones y Movimientos Restringidos}

\textbf{Idea:} muchas variantes de Nim limitan \textbf{cuántas} fichas se pueden quitar de un montón en un turno —por ejemplo, "a lo sumo $k$ fichas por turno" o "solo se pueden quitar cantidades de un conjunto permitido $S$"—. El teorema de Bouton (XOR directo de los tamaños) deja de aplicar tal cual, porque ya no es cierto que se pueda quitar \textbf{cualquier} cantidad; hay que calcular primero el \textbf{número de Grundy} de un montón individual bajo esa restricción de movimiento, y luego combinar los montones con XOR de esos números de Grundy (no de los tamaños directos), como establece \textit{Suma de Juegos Independientes}. El número de Grundy de un montón se calcula con la fórmula \texttt{mex} sobre los montones alcanzables, exactamente como en \textit{Cálculo de Grundy sobre un DAG}.

\textbf{¿Por qué usarlo?} Aplicar el XOR directo de los tamaños cuando hay una restricción de movimiento da resultados \textbf{incorrectos}, porque el teorema de Bouton depende críticamente de que cualquier reducción sea alcanzable en un solo movimiento. Con "a lo sumo $k$ por turno", por ejemplo, un montón de tamaño $n$ se comporta como Grundy $n \bmod (k+1)$ — un patrón periódico que hay que \textbf{descubrir}, típicamente con \textit{Fuerza Bruta para Descubrir Patrones}: calcular Grundy para montones pequeños, observar la periodicidad, y demostrarla o al menos verificarla extensivamente antes de confiar en la fórmula.

\textbf{Casos de uso:}
\begin{itemize}
	\item Nim con "máximo $k$ fichas por turno" ($k$-Nim), donde Grundy de un montón $n$ es $n \bmod (k+1)$.
	\item Nim donde solo se permite quitar cantidades de un conjunto fijo $S$ (es Nim disfrazado de \textit{Juego de Sustracción} por montón).
	\item Nim con montones donde la cantidad permitida depende de la paridad del turno o de otra condición sobre el estado.
	\item Verificación de que una variante de Nim, tras calcular Grundy por montón, efectivamente se reduce a XOR estándar cuando la restricción desaparece (caso $k\to\infty$).
\end{itemize}

\textbf{Complejidad:} $O(n \cdot |S|)$ o $O(n \cdot k)$ para precomputar el número de Grundy de cada tamaño de montón de $0$ a $n$ (uno por cada tamaño, cada uno revisando sus movimientos válidos); $O(1)$ por consulta una vez precomputada la tabla, y $O(\text{montones})$ para combinar el resultado final con XOR.

\begin{lstlisting}[style=cpp]
	// CREATE/READ: Numero de Grundy de un monton bajo una restriccion de
	// movimiento arbitraria (funcion movimientosValidos(n) -> tamanos
	// alcanzables), y combinacion de varios montones via XOR de Grundy
	struct NimRestringido {
		vector<int> grundy; // grundy[n] = numero de Grundy de un monton de tamano n
		
		// CREATE: precomputa grundy[0..nMax] dada la funcion de movimientos
		// validos desde un monton de tamano n
		NimRestringido(int nMax, function<vector<int>(int)> movimientosValidos) {
			grundy.assign(nMax + 1, 0);
			for (int n = 1; n <= nMax; n++) {
				vector<int> alcanzables = movimientosValidos(n);
				set<int> vistos;
				for (int m : alcanzables) {
					if (m >= 0 && m < n) vistos.insert(grundy[m]);
				}
				int mex = 0;
				while (vistos.count(mex)) mex++;
				grundy[n] = mex;
			}
		} // O(nMax * costo de movimientosValidos)
		
		// READ: true si la combinacion de varios montones (bajo la misma
		// restriccion) es GANADORA para el jugador que mueve primero
		bool primerJugadorGana(const vector<int>& montones) const {
			int x = 0;
			for (int m : montones) x ^= grundy[m];
			return x != 0;
		} // O(numero de montones)
		
		// ------------------------------------------------------------
		// MODIFICACIONES Y COMENTARIOS CON LOS USOS
		// ------------------------------------------------------------
		// 1. k-Nim (maximo k fichas por turno): movimientosValidos(n) retorna
		//    {n-1, n-2, ..., max(0, n-k)}; el patron resultante es
		//    grundy[n] = n % (k+1), verificable generando la tabla y
		//    observando la periodicidad (ver Fuerza Bruta para Descubrir
		//    Patrones).
		// 2. Conjunto de movimientos S fijo (subtraction game por monton):
		//    movimientosValidos(n) retorna {n-s : s en S, s <= n}; el patron
		//    de grundy puede no ser tan simple como un modulo, y a veces
		//    requiere la tabla completa sin formula cerrada.
		// 3. Verificacion cruzada con Nim Clasico: si movimientosValidos(n)
		//    retorna TODOS los valores {0, 1, ..., n-1} (sin restriccion),
		//    grundy[n] debe dar exactamente n para todo n, confirmando
		//    consistencia con el teorema de Bouton.
		// 4. Restriccion que depende de mas que el tamano (por ejemplo,
		//    ultimo movimiento realizado): el estado deja de ser solo el
		//    tamano del monton, y hace falta ampliar la clave de memoizacion
		//    a un par o tupla en vez de un solo entero.
		// 5. n muy grande con formula periodica detectada: una vez confirmado
		//    el periodo P y calculada la tabla para un periodo completo, se
		//    responde grundy(n) para n arbitrario con grundy[n % P] en O(1),
		//    sin necesidad de precomputar hasta nMax completo.
	};
\end{lstlisting}

\textbf{Ejemplo de funcionamiento:}
\begin{lstlisting}[style=cpp]
	// k-Nim con k=3 (maximo 3 fichas por turno)
	int k = 3;
	auto movK = [k](int n) {
		vector<int> res;
		for (int q = 1; q <= k && q <= n; q++) res.push_back(n - q);
		return res;
	};
	NimRestringido r(20, movK);
	
	int g5 = r.grundy[5];
	// g5 = 1
	// (5 % (3+1) = 5 % 4 = 1, confirmando el patron periodico esperado)
	
	bool gana = r.primerJugadorGana({5, 7, 2});
	// gana: XOR de grundy[5], grundy[7], grundy[2] segun la tabla
	// precomputada; hay que correrlo para confirmar el valor exacto
\end{lstlisting}

\textbf{Nota:} el error más común en esta variante es aplicar el XOR directo de los \textbf{tamaños} de los montones en vez del XOR de sus \textbf{números de Grundy}. Ambos coinciden solo en Nim clásico sin restricciones (donde \texttt{grundy[n] == n}); en cualquier variante restringida, hay que calcular Grundy primero y recién ahí combinar con XOR.

\textbf{Regla práctica:} ante cualquier variante de Nim con una restricción sobre cuánto se puede quitar, nunca asumas que el XOR directo de tamaños sigue funcionando. Calcula \texttt{grundy[n]} para varios valores pequeños de $n$, busca el patrón (muchas veces es periódico), verifícalo extensivamente, y solo entonces combina montones con XOR de esos valores de Grundy.
\subsection{Teorema de Sprague-Grundy (Grundy Numbers / Nimbers)}

\textbf{Idea:} el teorema de Sprague-Grundy generaliza \textit{Estados Ganadores y Perdedores} de una clasificación binaria (ganadora/perdedora) a un \textbf{número} por posición, llamado número de Grundy (o \textit{nimber}). Para una posición $p$ con sucesores $s_1,\dots,s_r$ (las posiciones alcanzables en un movimiento), el número de Grundy es $\text{G}(p)=\text{mex}\{\text{G}(s_1),\dots,\text{G}(s_r)\}$, donde \textbf{mex} (\textit{minimum excludant}) es el menor entero no negativo que \textbf{no} aparece en ese conjunto. Una posición terminal (sin movimientos) tiene $\text{G}=\text{mex}(\emptyset)=0$. El teorema establece que \textbf{todo} juego imparcial con información perfecta y sin azar, cualquiera sea su estructura, es equivalente a un montón de Nim de tamaño $\text{G}(p)$: $\text{G}(p)=0$ si y solo si $p$ es perdedora, exactamente como en el criterio de \textit{Estados Ganadores y Perdedores}.

\textbf{¿Por qué usarlo?} La clasificación binaria ganadora/perdedora responde "¿quién gana esta posición sola?", pero no dice nada sobre cómo se comporta esa posición cuando se \textbf{combina} con otras, como en un juego formado por varios montones o subjuegos independientes. El número de Grundy sí lo permite: es el "equivalente en Nim" de la posición, y por eso se puede combinar con otros números de Grundy usando XOR, tal como se ve en \textit{Suma de Juegos Independientes}. Calcular \texttt{mex} en vez de solo verificar "existe un cero entre los sucesores" cuesta apenas un poco más, pero abre la puerta a resolver sumas de juegos sin memoizar el producto cartesiano completo de estados.

\textbf{Casos de uso:}
\begin{itemize}
	\item Calcular el número de Grundy de una posición en cualquier juego imparcial, como paso previo a analizar sumas de subjuegos.
	\item Reconocer que un juego nuevo "es Nim disfrazado": si $\text{G}(n)=n$ para un montón de tamaño $n$, el juego se comporta exactamente como Nim clásico.
	\item Precomputar una tabla de Grundy para todos los tamaños de un componente de juego hasta cierto límite, y luego combinar varios componentes con XOR.
	\item Base de \textit{Nim con Montones y Movimientos Restringidos}, donde Grundy de un montón bajo restricciones no siempre es igual al tamaño.
\end{itemize}

\textbf{Complejidad:} $O(V+E)$ para calcular Grundy de todas las posiciones de un DAG con $V$ posiciones y $E$ movimientos totales, igual que en \textit{Análisis Retrógrado}; cada posición calcula su \texttt{mex} en $O(\text{grado de salida})$ con un conjunto o arreglo de marcado.

\begin{lstlisting}[style=cpp]
	// CREATE/READ: Calculo de numeros de Grundy sobre un conjunto de
	// posiciones dadas explicitamente por una funcion de sucesores, con
	// memoizacion (equivalente a Analisis Retrogrado pero con mex en vez de
	// clasificacion booleana)
	struct SpragueGrundy {
		int nMax;
		vector<int> grundy; // grundy[n] = numero de Grundy, o -1 si no calculado
		function<vector<int>(int)> sucesores;
		
		// CREATE: guarda el limite y la funcion de sucesores del juego
		SpragueGrundy(int limite, function<vector<int>(int)> func)
		: nMax(limite), sucesores(func) {
			grundy.assign(nMax + 1, -1);
		} // O(nMax)
		
		// READ: mex de un conjunto de valores: el menor entero no negativo
		// que NO aparece en el conjunto
		static int mex(const vector<int>& valores) {
			set<int> presentes(valores.begin(), valores.end());
			int m = 0;
			while (presentes.count(m)) m++;
			return m;
		} // O(k log k), k = tamano del conjunto
		
		// READ: calcula (memoizado) el numero de Grundy de la posicion n
		int calcular(int n) {
			if (grundy[n] != -1) return grundy[n];
			vector<int> sucGrundy;
			for (int s : sucesores(n)) {
				if (s >= 0 && s <= nMax) sucGrundy.push_back(calcular(s));
			}
			return grundy[n] = mex(sucGrundy);
		} // O(1) amortizado por posicion, O(nMax * grado) en total
		
		// READ: precomputa Grundy para TODAS las posiciones de 0 a nMax
		void calcularTodas() {
			for (int n = 0; n <= nMax; n++) calcular(n);
		} // O(nMax * grado promedio)
		
		// ------------------------------------------------------------
		// MODIFICACIONES Y COMENTARIOS CON LOS USOS
		// ------------------------------------------------------------
		// 1. mex con arreglo en vez de set (mas rapido cuando el grado de
		//    salida es acotado): usar vector<bool> visto(grado+2, false),
		//    marcar cada valor sucesor que sea <= grado, y recorrer desde 0
		//    hasta encontrar el primer false; evita el costo de un set
		//    balanceado cuando el numero de sucesores es pequeño y fijo.
		// 2. Estado compuesto (mas de una dimension, ej. posicion en un
		//    tablero 2D): usar un mapa (unordered_map con clave compuesta, o
		//    codificar el estado como un unico entero) en vez del arreglo
		//    grundy indexado por un solo entero.
		// 3. G(n) = n para Nim clasico: verificar esta identidad al calcular
		//    sucesores(n) = {0, 1, ..., n-1} confirma que la implementacion
		//    de mex es correcta antes de aplicarla a un juego mas complejo.
		// 4. Posicion GANADORA vs. PERDEDORA a partir de Grundy: ganadora si
		//    y solo si grundy[n] != 0; es una relajacion menos informativa
		//    del numero completo, util cuando solo importa el resultado de un
		//    unico juego (no una suma de varios).
		// 5. Reconstruccion de la jugada optima: junto con grundy[n], guardar
		//    el sucesor s cuyo grundy[s] es 0 (si existe); ese es el
		//    movimiento que deja al rival en una posicion perdedora.
	};
\end{lstlisting}

\textbf{Ejemplo de funcionamiento:}
\begin{lstlisting}[style=cpp]
	// juego de sustraccion con S = {1, 3, 4} (mismo conjunto que el
	// ejemplo de Estados Ganadores y Perdedores)
	auto sucS = [](int n) {
		vector<int> res;
		for (int m : {1, 3, 4}) if (m <= n) res.push_back(n - m);
		return res;
	};
	SpragueGrundy sg(10, sucS);
	sg.calcularTodas();
	
	int g2 = sg.grundy[2];
	// g2 = 0
	// (desde 2 solo se puede restar 1, llegando a grundy[1]; mex del
	//  conjunto {grundy[1]} da 0 si grundy[1] != 0, confirmando que 2 es
	//  perdedora, igual que se vio con esGanadora(2) = false)
	
	int g5 = sg.grundy[5];
	// g5: mex de los grundy de los sucesores de 5 (restando 1, 3 o 4);
	// hay que correrlo para confirmar el valor exacto
\end{lstlisting}

\textbf{Nota:} el número de Grundy \textbf{no es únicamente} "cero o distinto de cero" — aunque para analizar un solo juego alcanza con esa distinción (igual que \textit{Estados Ganadores y Perdedores}), el valor completo importa en cuanto el juego se combina con otros mediante suma, como se ve en la siguiente subsección. Calcular solo win/lose y luego intentar combinar esos booleanos con XOR es un error conceptual: hace falta el número completo.

\textbf{Regla práctica:} calcula números de Grundy (no solo ganador/perdedor) en cuanto el problema involucre \textbf{más de un componente de juego independiente} jugado simultáneamente. Si el juego es uno solo, sin combinar con otros, \textit{Estados Ganadores y Perdedores} alcanza y es más simple de razonar.

\subsection{Suma de Juegos Independientes (XOR de Grundy Numbers)}

\textbf{Idea:} muchos juegos en competencia no son un solo componente, sino la \textbf{suma} de varios subjuegos independientes jugados a la vez: en cada turno, el jugador elige \textbf{uno} de los subjuegos y hace un movimiento válido en \'el, dejando los demás intactos. Pierde quien no puede mover en \textbf{ninguno} de los subjuegos. El teorema de Sprague-Grundy dice que el número de Grundy de la suma es exactamente el \textbf{XOR} de los números de Grundy de cada subjuego: $\text{G}(p_1+p_2+\cdots+p_k)=\text{G}(p_1)\oplus\text{G}(p_2)\oplus\cdots\oplus\text{G}(p_k)$. Esto es una generalización directa del teorema de Bouton para Nim (donde cada montón es un subjuego cuyo Grundy es igual a su tamaño), pero aplica a \textbf{cualquier} colección de juegos imparciales, no solo a montones de fichas.

\textbf{¿Por qué usarlo?} Sin este teorema, analizar la suma de $k$ subjuegos requeriría memoizar sobre el \textbf{producto cartesiano} de todos los estados posibles de cada componente — inviable incluso con pocos componentes de tamaño moderado. El teorema reduce el problema a calcular el número de Grundy de \textbf{cada componente por separado} (mucho más pequeño, un solo grafo de estados a la vez) y combinarlos con una operación $O(1)$ por componente. Es la razón por la que Sprague-Grundy es la herramienta central de toda la parte: convierte "resolver un juego compuesto gigante" en "resolver $k$ juegos pequeños e independientes".

\textbf{Casos de uso:}
\begin{itemize}
	\item Cualquier problema que describa "varios montones/pilas/tableros independientes, elige uno y muévete ahí" — la estructura estándar de Nim y sus generalizaciones.
	\item Combinar distintos \textbf{tipos} de subjuego en una misma suma (por ejemplo, algunos montones de Nim clásico junto con montones de un juego de sustracción distinto): cada uno aporta su propio número de Grundy al XOR total.
	\item Determinar el ganador de la posición combinada sin necesidad de explorar el árbol de movimientos de la suma completa.
	\item Encontrar el movimiento ganador: se busca un componente y un movimiento dentro de \'el que cambie su Grundy de forma que el XOR total quede en $0$.
\end{itemize}

\textbf{Complejidad:} $O(\sum_i \text{costo de calcular Grundy}(p_i))$ para calcular el Grundy de cada componente por separado, más $O(k)$ para combinarlos con XOR — muchísimo más barato que memoizar el producto cartesiano de todos los componentes.

\begin{lstlisting}[style=cpp]
	// CREATE/READ: Suma de juegos independientes: combina los numeros de
	// Grundy de varios componentes (posiblemente de TIPOS distintos) via
	// XOR, y encuentra el movimiento ganador si existe
	struct SumaJuegos {
		// READ: XOR de una lista de numeros de Grundy ya calculados
		// (uno por cada componente de la suma)
		static int grundyTotal(const vector<int>& grundyPorComponente) {
			int x = 0;
			for (int g : grundyPorComponente) x ^= g;
			return x;
		} // O(k)
		
		// READ: true si el jugador que mueve PRIMERO gana la suma completa
		static bool primerJugadorGana(const vector<int>& grundyPorComponente) {
			return grundyTotal(grundyPorComponente) != 0;
		} // O(k)
		
		// READ: dado el grundy de cada componente y una funcion que, para un
		// componente i y un grundy objetivo, devuelve si existe un
		// movimiento en ese componente que lleve exactamente a ese grundy
		// (o -1 si no existe), encuentra el indice de componente y el nuevo
		// grundy que hacen ganar la posicion completa
		static pair<int, int> movimientoGanador(
		const vector<int>& grundyPorComponente,
		function<int(int, int)> intentarLlegarA // (indice, grundyObjetivo) -> nuevoGrundy o -1
		) {
			int total = grundyTotal(grundyPorComponente);
			if (total == 0) return {-1, -1}; // posicion perdedora: no hay movimiento ganador
			for (int i = 0; i < (int)grundyPorComponente.size(); i++) {
				int objetivo = grundyPorComponente[i] ^ total; // el mismo truco que en Nim clasico
				if (objetivo < grundyPorComponente[i]) {
					int resultado = intentarLlegarA(i, objetivo);
					if (resultado == objetivo) return {i, objetivo};
				}
			}
			return {-1, -1}; // no deberia ocurrir si total != 0 y el juego esta bien definido
		} // O(k * costo de intentarLlegarA)
		
		// ------------------------------------------------------------
		// MODIFICACIONES Y COMENTARIOS CON LOS USOS
		// ------------------------------------------------------------
		// 1. Componentes de TIPOS distintos: cada componente puede venir de un
		//    SpragueGrundy distinto (por ejemplo, unos montones de Nim puro y
		//    otros de un juego de sustraccion con S={1,2}); lo unico que
		//    importa aqui es el numero de Grundy final de cada uno, no como
		//    se calculo.
		// 2. Por que objetivo < grundyPorComponente[i] no siempre aplica: a
		//    diferencia de Nim clasico (donde reducir el TAMANO garantiza
		//    reducir el Grundy), en un juego general el Grundy no es
		//    necesariamente monotono con el tamano; la condicion exacta para
		//    "movimiento valido" depende del juego concreto, por eso
		//    intentarLlegarA() debe verificarlo explicitamente en vez de
		//    asumir monotonia.
		// 3. Muchos componentes IDENTICOS (por ejemplo, 50 montones todos con
		//    la misma regla): agrupar por Grundy y contar cuantas veces
		//    aparece cada valor; un valor que aparece un numero PAR de veces
		//    se cancela en el XOR (a XOR a = 0), asi que solo importan los
		//    que aparecen un numero impar de veces.
		// 4. Verificacion cruzada con Nim clasico: si todos los componentes
		//    son montones de Nim puro, grundyTotal() debe coincidir
		//    exactamente con NimClasico::nimSuma() sobre los mismos tamanos.
	};
\end{lstlisting}

\textbf{Ejemplo de funcionamiento:}
\begin{lstlisting}[style=cpp]
	// tres componentes de juegos distintos, cada uno ya con su Grundy calculado
	vector<int> grundys = {3, 5, 6}; // por ejemplo, via SpragueGrundy sobre cada tipo
	
	int total = SumaJuegos::grundyTotal(grundys);
	// total = 3 XOR 5 XOR 6 = 0
	
	bool gana = SumaJuegos::primerJugadorGana(grundys);
	// gana = false
	// (el XOR total es 0: la posicion combinada es PERDEDORA para quien
	//  mueve primero, sin importar que cada componente individual sea
	//  ganador por separado)
\end{lstlisting}

\textbf{Nota:} el resultado del ejemplo (una posición perdedora aunque ningún componente individual tenga Grundy $0$) es el punto central de esta técnica: la suma de varios juegos ganadores por separado puede ser perdedora en conjunto, y viceversa. Analizar cada componente aislado \textbf{no} dice nada directo sobre el resultado de la suma — hace falta el XOR de sus números de Grundy.

\textbf{Regla práctica:} en cuanto un problema describa varios subjuegos independientes donde se elige uno por turno, calcula el número de Grundy de \textbf{cada tipo} de componente por separado (con \textit{Sprague-Grundy} o, si es Nim puro, el tamaño mismo), combina todos con XOR, y decide el ganador según si ese XOR es cero. Nunca combines resultados binarios de ganador/perdedor con XOR directamente — solo los números de Grundy completos son combinables de esa forma.
\subsection{Cálculo de Grundy sobre un DAG (mex de Sucesores)}

\textbf{Idea:} esta subsección aísla el \textbf{paso mecánico} que ya apareció dentro de \textit{Teorema de Sprague-Grundy}: dado un grafo dirigido acíclico (DAG) de posiciones, calcular $\text{G}(p)=\text{mex}\{\text{G}(s):s\text{ sucesor de }p\}$ para \textbf{todas} las posiciones de forma sistemática, sin depender de una fórmula cerrada. Se puede hacer de dos formas equivalentes: \textbf{memoización recursiva} (igual que en Sprague-Grundy, bajando por las posiciones hasta las terminales y subiendo el resultado), o \textbf{tabulación por orden topológico inverso} (como en \textit{Análisis Retrógrado}, procesando primero los nodos sin sucesores y propagando hacia atrás). Ambas dan el mismo resultado; la tabulación evita el riesgo de desbordar la pila en DAGs muy profundos.

\textbf{¿Por qué usarlo?} Cuando un juego no tiene una fórmula conocida (no es Nim, no es un juego de sustracción con patrón periódico evidente), la única forma de obtener sus números de Grundy es calcularlos \textbf{directamente} sobre el grafo de posiciones, exactamente como se calcula cualquier DP sobre un DAG. Esta subsección documenta esa mecánica de forma explícita y reutilizable, separada del teorema en sí, porque es el bloque de construcción que se repite en casi todas las variantes de juego de esta parte: solo cambia cómo se generan los sucesores de cada posición.

\textbf{Casos de uso:}
\begin{itemize}
	\item Calcular Grundy de un juego nuevo sin fórmula conocida, cuando el número de posiciones distintas es manejable.
	\item Generar una tabla de Grundy para buscar un patrón periódico o cerrado, como paso previo a \textit{Fuerza Bruta para Descubrir Patrones}.
	\item Juegos sobre estructuras compuestas (pares de enteros, posiciones en una grilla) donde el estado se codifica como un entero único antes de aplicar esta misma mecánica.
	\item Verificar, sobre casos pequeños, una fórmula de Grundy sospechada para una variante de Nim o de sustracción.
\end{itemize}

\textbf{Complejidad:} $O(V+E)$ con $V$ posiciones y $E$ movimientos totales, igual que \textit{Análisis Retrógrado}: cada posición calcula su \texttt{mex} una sola vez, con costo proporcional a su grado de salida.

\begin{lstlisting}[style=cpp]
	// CREATE/READ: Calculo de Grundy sobre un DAG dado explicitamente,
	// con dos versiones equivalentes: memoizacion recursiva y tabulacion
	// por orden topologico inverso
	struct GrundyDAG {
		int n;
		vector<vector<int>> sucesores;
		vector<int> grundy;
		
		// CREATE: reserva estructuras para n posiciones (0..n-1)
		GrundyDAG(int numPosiciones) : n(numPosiciones) {
			sucesores.assign(n, {});
			grundy.assign(n, -1);
		} // O(n)
		
		// WRITE: agrega un movimiento valido de 'desde' hacia 'hacia'
		void agregarMovimiento(int desde, int hacia) {
			sucesores[desde].push_back(hacia);
		} // O(1) amortizado
		
		// READ: mex de un conjunto de valores (arreglo, mas rapido que set
		// cuando el grado de salida es pequeño y acotado)
		static int mex(const vector<int>& valores) {
			int m = (int)valores.size();
			vector<bool> visto(m + 1, false);
			for (int v : valores) if (v <= m) visto[v] = true;
			int r = 0;
			while (r <= m && visto[r]) r++;
			return r;
		} // O(grado de salida)
		
		// READ: version recursiva con memoizacion
		int calcularRecursivo(int p) {
			if (grundy[p] != -1) return grundy[p];
			vector<int> vals;
			for (int s : sucesores[p]) vals.push_back(calcularRecursivo(s));
			return grundy[p] = mex(vals);
		} // O(1) amortizado por posicion
		
		// READ: version tabulada (bottom-up) por orden topologico inverso;
		// requiere que el grafo sea un DAG
		void calcularTabulado() {
			vector<int> gradoSalida(n);
			for (int i = 0; i < n; i++) gradoSalida[i] = (int)sucesores[i].size();
			
			vector<vector<int>> predecesores(n);
			for (int i = 0; i < n; i++)
			for (int j : sucesores[i]) predecesores[j].push_back(i);
			
			queue<int> cola;
			vector<int> pendientes = gradoSalida;
			for (int i = 0; i < n; i++) {
				if (gradoSalida[i] == 0) {
					grundy[i] = 0; // mex del conjunto vacio
					cola.push(i);
				}
			}
			
			while (!cola.empty()) {
				int actual = cola.front(); cola.pop();
				for (int p : predecesores[actual]) {
					pendientes[p]--;
					if (pendientes[p] == 0) {
						vector<int> vals;
						for (int s : sucesores[p]) vals.push_back(grundy[s]);
						grundy[p] = mex(vals);
						cola.push(p);
					}
				}
			}
		} // O(V + E)
		
		// ------------------------------------------------------------
		// MODIFICACIONES Y COMENTARIOS CON LOS USOS
		// ------------------------------------------------------------
		// 1. Estado codificado (no un simple 0..n-1): si la posicion es un
		//    par (i,j) o una mascara de bits, codificarla primero a un unico
		//    entero (por ejemplo i*ancho+j, o el valor de la mascara) antes
		//    de usar esta misma estructura indexada por entero.
		// 2. Sucesores generados bajo demanda (el DAG es implicito, muy
		//    grande para enumerar): usar calcularRecursivo() con un
		//    unordered_map<clave,int> en vez del vector grundy, generando
		//    sucesores() como una funcion en vez de una lista precomputada.
		// 3. mex con set en vez de arreglo: si el grado de salida no esta
		//    acotado por el numero de posiciones (por ejemplo, movimientos
		//    con valores muy dispersos), usar set<int> como en la version
		//    de Sprague-Grundy, a costa de un factor log adicional.
		// 4. Verificar contra Estados Ganadores y Perdedores: grundy[p] == 0
		//    debe coincidir exactamente con "p es perdedora" segun el
		//    analisis retrogrado booleano; sirve como chequeo de cordura
		//    entre ambas implementaciones.
	};
\end{lstlisting}

\textbf{Ejemplo de funcionamiento:}
\begin{lstlisting}[style=cpp]
	GrundyDAG g(6);
	g.agregarMovimiento(1, 0);
	g.agregarMovimiento(2, 0);
	g.agregarMovimiento(2, 1);
	g.agregarMovimiento(3, 1);
	g.agregarMovimiento(3, 2);
	g.agregarMovimiento(4, 0);
	g.agregarMovimiento(4, 3);
	
	g.calcularTabulado();
	
	int g0 = g.grundy[0];
	// g0 = 0 (posicion terminal, sin sucesores)
	
	int g1 = g.grundy[1];
	// g1 = 1 (su unico sucesor es 0, con grundy 0; mex({0}) = 1)
	
	int g2 = g.grundy[2];
	// g2: mex de los grundy de sus sucesores {0, 1}; hay que correrlo para
	// confirmar el valor exacto segun la tabla completa
\end{lstlisting}

\textbf{Nota:} las dos versiones (recursiva y tabulada) dan siempre el mismo resultado. La tabulada requiere construir explícitamente la lista de predecesores antes de empezar, un costo adicional de $O(E)$ que la recursiva no necesita; a cambio, evita cualquier riesgo de desbordar la pila cuando el DAG tiene cadenas muy largas de dependencias.

\textbf{Regla práctica:} usa esta mecánica en cuanto un juego no tenga una fórmula de Grundy conocida de antemano — es el mismo patrón de \textit{Análisis Retrógrado}, con \texttt{mex} en vez de una clasificación booleana. Si el número de posiciones es pequeño, la versión recursiva es más rápida de escribir; si es grande o profundo, usa la tabulada.

\subsection{Misère Nim (Última Jugada Pierde)}

\textbf{Idea:} en Nim clásico (\textit{normal play}), pierde quien \textbf{no puede mover}. En la variante \textbf{misère}, la regla se invierte: pierde quien \textbf{hace el último movimiento} (deja todos los montones en $0$); quien no puede mover \textbf{gana}. Sorprendentemente, la estrategia óptima de Misère Nim es \textbf{casi idéntica} a la de Nim normal: se calcula igual la nim-suma $X=a_1\oplus\cdots\oplus a_k$, y solo cambia la regla en un caso especial. Sea $m$ el número de montones con tamaño \textbf{mayor a 1}. Si $m\ge 1$ (existe al menos un montón "grande"), la estrategia es exactamente la de Nim normal: ganar si $X\neq 0$, jugando siempre hacia $X=0$. Si $m=0$ (todos los montones tienen tamaño $0$ o $1$, es decir, el juego ya se redujo a "montoncitos de una ficha"), la regla se invierte: el primer jugador gana si y solo si la \textbf{cantidad de montones de tamaño 1} es \textbf{par}.

\textbf{¿Por qué usarlo?} Aplicar directamente el criterio de Nim normal ($X\neq 0$ gana) a una variante misère da la respuesta \textbf{incorrecta} exactamente en el caso borde donde todos los montones son de tamaño $0$ o $1$ — un caso fácil de pasar por alto en el enunciado si no se lee con cuidado la condición de fin de juego ("el que mueve último pierde" en vez de "el que no puede mover pierde"). El resultado de Bouton para misère (a veces llamado el teorema de Bouton-Misère) evita tener que recalcular Grundy desde cero para esta variante: reutiliza casi toda la maquinaria de Nim normal, con un único ajuste puntual.

\textbf{Casos de uso:}
\begin{itemize}
	\item Cualquier variante de Nim donde el enunciado especifique explícitamente "quien hace el último movimiento pierde" en vez de la regla normal.
	\item Como recordatorio de que \textbf{leer con cuidado} la condición de victoria/derrota es tan importante como conocer la fórmula: el mismo juego (mismos montones) puede tener resultado opuesto según la regla de terminación.
	\item Contraste didáctico frente a \textit{Nim Clásico}: mostrar que dos reglas de fin de juego aparentemente simétricas no dan estrategias simétricas en general (la simplicidad de Misère Nim es una excepción; la mayoría de juegos misère son mucho más difíciles de analizar que su versión normal).
	\item Verificar con \textit{Análisis Retrógrado} sobre montones pequeños, antes de confiar en la fórmula, que la regla del caso $m=0$ está bien aplicada.
\end{itemize}

\textbf{Complejidad:} $O(k)$, igual que Nim normal: un recorrido para calcular la nim-suma y otro (o el mismo) para contar montones mayores a 1 y montones de tamaño exactamente 1.

\begin{lstlisting}[style=cpp]
	// CREATE/READ: Misere Nim: determina el ganador bajo la regla "el que
	// mueve ultimo pierde", reutilizando la nim-suma de Nim normal con el
	// ajuste del caso especial (todos los montones de tamano 0 o 1)
	struct MisereNim {
		vector<long long> montones;
		
		// CREATE: guarda los tamanos de los montones
		MisereNim(const vector<long long>& a) : montones(a) {
		} // O(k)
		
		// READ: true si el jugador que mueve PRIMERO gana, bajo regla misere
		bool primerJugadorGana() const {
			long long x = 0;
			int montonesGrandes = 0;  // tamano > 1
			int montonesDeUno = 0;    // tamano == 1
			for (long long a : montones) {
				x ^= a;
				if (a > 1) montonesGrandes++;
				else if (a == 1) montonesDeUno++;
			}
			
			if (montonesGrandes > 0) {
				// hay al menos un monton grande: se juega EXACTAMENTE como
				// Nim normal
				return x != 0;
			} else {
				// todos los montones son 0 o 1: la regla se invierte
				return montonesDeUno % 2 == 0;
			}
		} // O(k)
		
		// ------------------------------------------------------------
		// MODIFICACIONES Y COMENTARIOS CON LOS USOS
		// ------------------------------------------------------------
		// 1. Por que el caso m=0 se invierte: si todos los montones son 0 o
		//    1, cada movimiento SOLO puede vaciar un monton de tamano 1 (no
		//    hay otra opcion); el juego se convierte en "turnos alternados
		//    hasta que no queden montones de 1", y quien hace el ULTIMO de
		//    esos movimientos pierde bajo regla misere, lo que invierte la
		//    paridad respecto a Nim normal.
		// 2. Version general (juegos misere SIN esta simplificacion): fuera
		//    de Nim, la mayoria de juegos imparciales bajo regla misere NO
		//    tienen un atajo tan simple; se necesita la teoria de juegos
		//    misere generalizada (Grundy misere / teoria de Conway extendida),
		//    fuera del alcance de esta guia.
		// 3. Verificacion cruzada: comparar primerJugadorGana() con un
		//    Analisis Retrogrado directo (definiendo la posicion terminal
		//    como GANADORA para quien esta en turno, en vez de perdedora)
		//    sobre montones pequeños, para confirmar que el caso especial
		//    esta bien aplicado antes de confiar en la formula.
		// 4. Un solo monton (k=1): si el monton tiene tamano > 1, gana el
		//    primer jugador (x != 0 siempre que el monton no sea 0); si el
		//    monton es exactamente 1, hay 1 monton de tamano 1 (impar), asi
		//    que primerJugadorGana() da false: el primer jugador debe quitar
		//    esa ultima ficha y por regla misere pierde.
	};
\end{lstlisting}

\textbf{Ejemplo de funcionamiento:}
\begin{lstlisting}[style=cpp]
	MisereNim n1({3, 4, 5});
	bool r1 = n1.primerJugadorGana();
	// r1 = true
	// (hay montones > 1, asi que se aplica la regla normal: nim-suma
	//  3^4^5 = 2, distinta de 0, el primer jugador gana igual que en Nim
	//  normal con estos mismos montones)
	
	MisereNim n2({1, 1, 1});
	bool r2 = n2.primerJugadorGana();
	// r2 = false
	// (todos los montones son de tamano 1, hay 3 montones de 1 -- impar --
	//  asi que el primer jugador PIERDE bajo regla misere)
	
	MisereNim n3({1, 1});
	bool r3 = n3.primerJugadorGana();
	// r3 = true
	// (2 montones de tamano 1 -- par -- el primer jugador gana bajo regla
	//  misere, aunque bajo Nim normal con estos mismos montones el
	//  resultado seria el mismo por coincidencia: 1^1=0, tambien perdedora
	//  para el primero en Nim normal... hay que correrlo para confirmar
	//  que no hay contradiccion entre ambas lecturas en este caso particular)
\end{lstlisting}

\textbf{Nota:} el último ejemplo (\texttt{n3}) es intencionalmente delicado: con dos montones de tamaño 1, la nim-suma normal ya da $0$ (perdedora bajo regla normal), pero la regla misère dice lo contrario para este mismo caso porque cae en el caso especial $m=0$. Antes de confiar en cualquier intuición cruzada entre ambas reglas, conviene verificar numéricamente con \textit{Análisis Retrógrado} sobre montones pequeños.

\textbf{Regla práctica:} lee siempre con extremo cuidado si el enunciado dice "el que no puede mover pierde" (Nim normal) o "el que hace el último movimiento pierde" (misère) — son reglas fácilmente confundibles en la traducción del problema. Si es misère, usa la nim-suma normal salvo en el caso especial de que todos los montones sean $0$ o $1$, donde la paridad de montones de tamaño $1$ decide, invertida respecto a la intuición de Nim normal.

\section{Juegos de Sustracción y Variantes Clásicas}
\subsection{Juego de Sustracción (Subtraction Game)}

\textbf{Idea:} dos jugadores alternan turnos sobre una sola pila de
$n$ piedras. En cada turno, el jugador retira una cantidad
$s$ perteneciente a un conjunto \textbf{fijo} de movimientos
permitidos $S$ (por ejemplo $S=\{1,2,3\}$). Con la convención normal de
juego, \textbf{pierde} quien no puede mover (la pila llega a un tamaño
donde ningún $s \in S$ cabe). Se resuelve calculando el
\textbf{número de Grundy} de cada tamaño de pila $0,1,\dots,n$ con la
recurrencia estándar de Sprague-Grundy:
$$g(0) = 0, \qquad g(i) = \operatorname{mex}\{\, g(i-s) : s \in S,\ s
\le i \,\}$$
donde $\operatorname{mex}$ (\textit{minimum excludant}) es el menor
entero no negativo que \textbf{no} aparece en el conjunto. Una
posición $i$ es \textbf{perdedora} para quien está por mover (posición
$P$) si y solo si $g(i) = 0$; en cualquier otro caso es
\textbf{ganadora} (posición $N$).

\textbf{¿Por qué usarlo?} Es el caso más simple de juego imparcial
(\textit{impartial game}: ambos jugadores tienen exactamente las
mismas jugadas disponibles desde cualquier posición) y sirve de
introducción a Sprague-Grundy antes de combinar varios montones (ver
Suma de Juegos Independientes). Un caso particular vale la pena
memorizar porque aparece muy seguido: con $S=\{1,\dots,k\}$ (se puede
quitar cualquier cantidad de $1$ a $k$), la posición $n$ es
\textbf{perdedora} si y solo si $n \bmod (k+1) = 0$ — no hace falta ni
correr la DP para ese caso, aunque el struct de abajo también lo
resuelve como caso general.

\textbf{Casos de uso:}

\begin{itemize}
	\item Determinar quién gana con $n$ piedras y un conjunto de
	movimientos $S$ dado explícitamente.
	\item Calcular los números de Grundy de \textbf{cada} tamaño de
	pila, necesarios como pieza para Suma de Juegos Independientes
	(varios montones independientes con el mismo $S$, o incluso con
	$S$ distinto cada uno).
	\item Detectar el \textbf{período} de la secuencia de
	ganadoras/perdedoras cuando $n$ es enorme: por el teorema de
	periodicidad de Sprague-Grundy, para cualquier $S$ finito la
	secuencia de posiciones P/N es eventualmente periódica, así que
	basta calcular unos pocos miles de términos, detectar el período y
	extrapolar.
	\item Como caso base antes de agregar restricciones (ver Juego de
	Sustracción con Conjunto de Movimientos Prohibidos).
\end{itemize}

\textbf{Complejidad:} $O(n \cdot |S|)$ tiempo y $O(n)$ memoria para
llenar la tabla hasta $n$; cada consulta de ganador es $O(1)$ una vez
construida.

\begin{lstlisting}[style=cpp]
	struct SubtractionGame {
		vector<int> S;              // conjunto de movimientos permitidos
		vector<int> grundy;         // grundy[i] para i = 0..maxN
		vector<bool> esGanadora;    // esGanadora[i]: true si quien esta
		// por MOVER con i piedras gana
		
		// CREATE: calcula grundy[] y esGanadora[] para 0..maxN.
		SubtractionGame(vector<int> S_, int maxN) : S(S_) {
			grundy.assign(maxN + 1, 0);
			esGanadora.assign(maxN + 1, false);
			
			for (int n = 1; n <= maxN; n++) {
				set<int> alcanzables;
				for (int s : S)
				if (s <= n)
				alcanzables.insert(grundy[n - s]);
				
				int mex = 0;
				while (alcanzables.count(mex)) mex++;
				
				grundy[n] = mex;
				esGanadora[n] = (mex != 0);
			}
		}                                          // O(maxN * |S|)
		
		// READ: quien esta por mover con n piedras, gana?
		bool ganaJugadorEnTurno(int n) {
			return esGanadora[n];
		}                                          // O(1)
	};
\end{lstlisting}

\textbf{Ejemplo de funcionamiento:}

\begin{lstlisting}[style=cpp]
	SubtractionGame g({1, 2, 3}, 12);
	
	for (int n = 0; n <= 12; n++) cout << g.grundy[n] << " ";
	// 0 1 2 3 0 1 2 3 0 1 2 3 0
	// (se repite cada 4 valores, confirmando n % (k+1) con k=3)
	
	for (int n = 0; n <= 12; n++) cout << g.ganaJugadorEnTurno(n) << " ";
	// 0 1 1 1 0 1 1 1 0 1 1 1 0
	// (pierde exactamente cuando n es multiplo de 4)
\end{lstlisting}

\textbf{Nota:} el número de Grundy \textbf{no} es solo "gana/pierde"
—guarda \textbf{cuánto} vale esa posición para combinarla con otras
(XOR de Grundy numbers en Suma de Juegos Independientes). Si el
problema es de un único montón, basta con \texttt{esGanadora[n]}; si
hay varios montones independientes (posiblemente con distinto $S$ cada
uno), hace falta \texttt{grundy[n]} completo para cada uno.

\textbf{Regla práctica:} antes de programar la DP, revisa si $S$ es
$\{1,\dots,k\}$ (un rango completo desde 1) —en ese caso la respuesta
es literalmente \texttt{n \% (k+1) != 0}, sin código adicional. Para
cualquier otro $S$, o si además se necesitan los números de Grundy
para combinarlos con otros juegos, usa la DP completa.


\subsection{Juego de Sustracción con Conjunto de Movimientos Prohibidos}

\textbf{Idea:} variante donde, en vez de dar directamente el conjunto
\textbf{permitido} $S$, el enunciado da un límite $k$ (se puede quitar
cualquier cantidad de $1$ a $k$) y un conjunto $F \subseteq
\{1,\dots,k\}$ de cantidades \textbf{prohibidas} (movimientos que el
enunciado explícitamente no deja usar). El conjunto \textbf{efectivo}
de movimientos permitidos es el complemento:
$$S_{\text{efectivo}} = \{1,\dots,k\} \setminus F$$
Una vez calculado $S_{\text{efectivo}}$, el resto es \textbf{exactamente}
la misma DP de Grundy de Juego de Sustracción —esta subsección no
introduce ningún algoritmo nuevo, solo el paso de \textbf{traducir} la
descripción "todo permitido salvo estos" a un conjunto concreto antes
de aplicar la técnica ya conocida.

\textbf{¿Por qué usarlo?} Es un recordatorio útil de que en teoría de
juegos, como en varias partes de esta guía, conviene \textbf{reconocer
	la reformulación} antes de programar algo nuevo: describir las reglas
"al revés" (por prohibición en vez de por permiso) no cambia la
naturaleza combinatoria del juego, solo la forma en que hay que leerlo.
El único cuidado real está en los \textbf{casos borde} de esta
traducción: si $F$ termina cubriendo \textbf{todo} el rango
$\{1,\dots,k\}$, el conjunto efectivo queda \textbf{vacío}, y entonces
\textbf{nadie} puede mover nunca — el primer jugador pierde de
inmediato con cualquier $n \ge 1$, sin que la pila ni siquiera se
reduzca.

\textbf{Casos de uso:}

\begin{itemize}
	\item Cualquier enunciado que fije un tope de piedras por turno
	($1$ a $k$) y liste algunas cantidades específicas que no se
	pueden usar (evitar ciertos números por regla del juego, por
	superstición del enunciado, etc.).
	\item Como ejercicio de traducción antes de aplicar Sprague-Grundy:
	practicar convertir "prohibido" en "permitido" sin errores de
	borde, algo que se repite en variantes más elaboradas de juegos
	combinatorios.
	\item Detectar de antemano, sin correr ninguna DP, si el juego es
	trivial porque el conjunto efectivo quedó vacío.
\end{itemize}

\textbf{Complejidad:} $O(k)$ para construir el conjunto efectivo, más
$O(n \cdot |S_{\text{efectivo}}|)$ para la DP de Grundy — igual orden
que Juego de Sustracción, solo con un paso de preprocesamiento extra.

\begin{lstlisting}[style=cpp]
	struct SubtractionGameProhibido {
		int k;                       // maximo de piedras removibles por turno
		vector<int> permitidos;      // conjunto EFECTIVO (complemento de F)
		vector<int> grundy;
		vector<bool> esGanadora;
		
		// CREATE: calcula el conjunto efectivo de movimientos a partir de
		// k y del conjunto prohibido F, y corre la MISMA DP de Grundy que
		// Juego de Sustraccion.
		SubtractionGameProhibido(int k_, vector<int>& F, int maxN) : k(k_) {
			vector<bool> prohibido(k + 1, false);
			for (int f : F)
			if (f >= 1 && f <= k)
			prohibido[f] = true;
			
			for (int s = 1; s <= k; s++)
			if (!prohibido[s])
			permitidos.push_back(s);
			
			grundy.assign(maxN + 1, 0);
			esGanadora.assign(maxN + 1, false);
			
			for (int n = 1; n <= maxN; n++) {
				set<int> alcanzables;
				for (int s : permitidos)
				if (s <= n)
				alcanzables.insert(grundy[n - s]);
				
				int mex = 0;
				while (alcanzables.count(mex)) mex++;
				
				grundy[n] = mex;
				esGanadora[n] = (mex != 0);
			}
		}                                          // O(k + maxN * |permitidos|)
		
		// READ: quien esta por mover con n piedras, gana?
		bool ganaJugadorEnTurno(int n) {
			return esGanadora[n];
		}                                          // O(1)
	};
\end{lstlisting}

\textbf{Ejemplo de funcionamiento:}

\begin{lstlisting}[style=cpp]
	// k=5, prohibidos {2,4} -> permitidos efectivos = {1,3,5}
	vector<int> F = {2, 4};
	SubtractionGameProhibido g(5, F, 12);
	
	cout << "permitidos: ";
	for (int s : g.permitidos) cout << s << " ";
	// permitidos: 1 3 5
	
	for (int n = 0; n <= 12; n++) cout << g.grundy[n] << " ";
	// 0 1 0 1 0 1 0 1 0 1 0 1 0
	// (identico a correr SubtractionGame con S={1,3,5} explicito)
	
	for (int n = 0; n <= 12; n++) cout << g.ganaJugadorEnTurno(n) << " ";
	// 0 1 0 1 0 1 0 1 0 1 0 1 0
	
	// caso borde: F cubre TODO {1,2,3} con k=3 -> permitidos vacio
	vector<int> Ftodo = {1, 2, 3};
	SubtractionGameProhibido gVacio(3, Ftodo, 5);
	
	cout << gVacio.permitidos.size();
	// 0
	
	for (int n = 0; n <= 5; n++) cout << gVacio.ganaJugadorEnTurno(n) << " ";
	// 0 0 0 0 0 0
	// (nadie puede mover NUNCA -> el que esta en turno siempre pierde,
	//  para CUALQUIER n >= 1)
\end{lstlisting}

\textbf{Nota:} el bucle \texttt{for (int s : permitidos) if (s <= n)}
ya maneja automáticamente que \texttt{permitidos} esté vacío —el
\texttt{set alcanzables} queda vacío, el \texttt{mex} de un conjunto
vacío es $0$, y por lo tanto \texttt{esGanadora[n] = false} para todo
$n$ sin necesitar un chequeo especial adicional. Vale la pena
\textbf{verificar} este caso a mano una vez, porque es fácil asumir
por error que "sin movimientos permitidos" debería dar algún tipo de
error en vez de una respuesta válida del juego.

\textbf{Regla práctica:} cuando un enunciado de teoría de juegos
describe las jugadas por \textbf{exclusión} ("puedes hacer cualquier
cosa dentro de este rango, excepto...") en vez de por
\textbf{inclusión} directa, el primer paso —antes de pensar en
Grundy— es escribir explícitamente el conjunto efectivo de jugadas
permitidas. Una vez que ese conjunto está en una variable concreta, el
resto del problema es idéntico a Juego de Sustracción; no hace falta
ninguna teoría adicional para la parte de "prohibido".
\subsection{Staircase Nim}

\textbf{Idea:} una escalera de $n$ escalones (numerados $1$ a $n$, con
el escalón $0$ siendo el suelo) tiene una cantidad de monedas en cada
escalón. En cada turno, un jugador elige un escalón $k$ con al menos
una moneda y mueve \textbf{cualquier cantidad positiva} de esas
monedas al escalón $k-1$ (si $k=1$, esas monedas caen al suelo y
\textbf{salen del juego} para siempre). Pierde quien no puede mover
(todas las monedas ya están en el suelo). El resultado clave: este
juego es \textbf{equivalente} a un Nim clásico jugado únicamente sobre
las monedas de los escalones \textbf{impares} ($1, 3, 5, \dots$) — los
escalones pares no importan para saber quién gana. La razón:
cualquier movimiento desde un escalón \textbf{par} hacia uno impar
puede ser \textbf{espejado} por el rival, moviendo esa misma cantidad
del escalón impar recién aumentado hacia el siguiente escalón par (o
fuera del tablero si era el escalón $1$), devolviendo la posición
efectiva de los escalones impares al estado anterior — así que mover
piezas hacia un escalón par nunca cambia realmente el resultado. El
ganador se decide con el mismo criterio que Nim Clásico:
$$\text{gana quien está por mover} \iff a_1 \oplus a_3 \oplus a_5
\oplus \cdots \ne 0$$

\textbf{¿Por qué usarlo?} Es el primer ejemplo de esta guía donde un
juego con reglas de movimiento \textbf{distintas} a Nim (mover entre
escalones, no quitar piedras de montones) termina siendo,
combinatoriamente, \textbf{el mismo juego} que Nim Clásico sobre un
subconjunto de la información dada — reconocer esa reducción es más
valioso que memorizar la prueba completa. Es también el ejemplo
introductorio de un patrón que reaparece en juegos combinatorios más
elaborados: "algunas partes del estado no afectan el resultado, solo
las que sobreviven a un argumento de espejado/paridad".

\textbf{Casos de uso:}

\begin{itemize}
	\item Determinar el ganador de la escalera de monedas tal cual
	(el problema base).
	\item Cualquier variante narrativa equivalente: fichas que bajan
	de nivel en nivel, recursos que se transfieren escalón por
	escalón hacia una salida — el criterio de "solo importan las
	posiciones impares, mismo criterio que Nim" se mantiene mientras
	la regla de movimiento sea "de un nivel a uno estrictamente
	inferior, en cualquier cantidad".
	\item Como ejercicio de reconocimiento: antes de intentar un
	análisis de Grundy completo (con estados de tamaño exponencial),
	conviene sospechar una reducción a Nim cuando el juego tiene
	\textbf{niveles} y las piezas solo se mueven hacia niveles más
	bajos.
\end{itemize}

\textbf{Complejidad:} $O(n)$ — un solo recorrido para hacer XOR de
las posiciones impares.

\begin{lstlisting}[style=cpp]
	struct StaircaseNim {
		vector<long long> escalones;   // escalones[0] = escalon 1,
		// escalones[1] = escalon 2, etc.
		// (0-indexado en el vector,
		// pero 1-indexado como escalera)
		int n;
		
		// CREATE: guarda la cantidad de monedas en cada escalon.
		StaircaseNim(vector<long long>& escalones_) {
			escalones = escalones_;
			n = escalones.size();
		}                                          // O(1)
		
		// READ: gana quien esta por mover? Solo cuentan los escalones de
		// posicion IMPAR (1, 3, 5, ... en la numeracion de 1 a n).
		bool ganaJugadorEnTurno() {
			long long xorImpares = 0;
			
			for (int i = 0; i < n; i++) {
				int numeroDeEscalon = i + 1;   // convierte a 1-indexado
				if (numeroDeEscalon % 2 == 1)
				xorImpares ^= escalones[i];
			}
			
			return xorImpares != 0;
		}                                          // O(n)
	};
\end{lstlisting}

\textbf{Ejemplo de funcionamiento:}

\begin{lstlisting}[style=cpp]
	// escalera de 4 escalones: [escalon1, escalon2, escalon3, escalon4]
	vector<long long> escalones = {3, 5, 2, 7};
	
	StaircaseNim sn(escalones);
	
	cout << sn.ganaJugadorEnTurno();
	// 1 (true)
	// escalones impares: 1 y 3, con 3 y 2 monedas -> 3 XOR 2 = 1, != 0
	// (los escalones pares, 5 y 7 monedas, NO afectan el resultado)
\end{lstlisting}

\textbf{Nota:} el struct de arriba \textbf{no} necesita simular ningún
movimiento —a diferencia de un análisis de Grundy genérico sobre todo
el espacio de estados (que sería computacionalmente inviable, ya que
el número de configuraciones crece exponencialmente con la cantidad de
monedas), la reducción a Nim resuelve el juego completo en una sola
pasada lineal. Esta es la ventaja real de identificar la equivalencia
antes de programar cualquier DP.

\textbf{Regla práctica:} en cualquier juego donde las piezas se mueven
\textbf{siempre hacia niveles inferiores} de una estructura escalonada
(nunca hacia arriba, nunca a un nivel arbitrario), sospecha
inmediatamente de una reducción a Nim sobre un subconjunto de niveles
—el patrón de "el rival puede espejar cualquier movimiento hacia un
nivel que no cuenta, devolviéndolo a un nivel que sí cuenta" es la
prueba general detrás de Staircase Nim y se adapta a variantes
similares.


\subsection{Wythoff's Game}

\textbf{Idea:} dos montones de fichas, de tamaños $a$ y $b$. En cada
turno, un jugador puede: (1) quitar cualquier cantidad positiva de
fichas de \textbf{un solo} montón (como en Nim con 2 montones), o
(2) quitar la \textbf{misma} cantidad positiva de \textbf{ambos}
montones a la vez (movimiento extra que Nim \textbf{no} permite).
Pierde quien no puede mover (ambos montones vacíos). La tercera regla
es lo que rompe la simple fórmula de Nim ($a \oplus b \ne 0$) y exige
un análisis propio: las posiciones \textbf{perdedoras} (posición $P$,
pierde quien está por mover) resultan ser \textbf{exactamente} los
pares $(\lfloor k\varphi \rfloor,\ \lfloor k\varphi \rfloor + k)$ para
$k = 0, 1, 2, \dots$, donde $\varphi = \frac{1+\sqrt{5}}{2}$ es la
\textbf{razón áurea} — un resultado clásico de teoría de números
(sucesiones de Beatty) aplicado a teoría de juegos.

\textbf{¿Por qué usarlo?} Es el ejemplo donde la razón áurea aparece
en un contexto puramente combinatorio, sin nada geométrico de por
medio —vale la pena como muestra de que las posiciones perdedoras de
un juego no siempre siguen un patrón simple de paridad o XOR, a veces
la estructura subyacente es aritmética/número-teórica. En la práctica,
la fórmula da una consulta $O(1)$: dado $(a,b)$, se calcula
directamente si corresponde a alguna $k$ de la sucesión, sin necesitar
ninguna tabla ni DP.

\textbf{Casos de uso:}

\begin{itemize}
	\item Determinar el ganador de Wythoff's Game dado $(a, b)$
	directamente, sin simular ni memoizar.
	\item Generar las primeras $K$ posiciones perdedoras (para
	listarlas, o para usarlas como tabla de consulta si se prefiere
	evitar aritmética de punto flotante).
	\item Como ejemplo de que agregar \textbf{una sola} regla extra a
	un juego de tipo Nim (quitar lo mismo de ambos montones) puede
	cambiar por completo la estructura de las posiciones ganadoras,
	de "XOR simple" a "sucesión de Beatty con la razón áurea".
\end{itemize}

\textbf{Complejidad:} $O(1)$ por consulta usando la fórmula cerrada
(con cuidado de precisión en $\varphi$); $O(K)$ para generar las
primeras $K$ posiciones perdedoras.

\begin{lstlisting}[style=cpp]
	struct WythoffGame {
		static constexpr double PHI = 1.6180339887498949;   // (1+sqrt(5))/2
		
		// READ: gana quien esta por mover con montones (a, b)?
		static bool ganaJugadorEnTurno(long long a, long long b) {
			if (a > b) swap(a, b);   // normaliza: a <= b
			
			long long k = (long long)((double)a / PHI + 0.5);
			// ^ si (a,b) es una posicion PERDEDORA, existe un k tal que
			//   a == floor(k*phi); se despeja k = a / phi y se redondea
			
			long long ak = (long long)floor(k * PHI);
			long long bk = ak + k;
			
			bool esPerdedora = (ak == a && bk == b);
			return !esPerdedora;
		}                                          // O(1)
		
		// READ: genera las primeras cantidad posiciones PERDEDORAS
		// (a_k, b_k) para k = 0, 1, ..., cantidad-1.
		static vector<pair<long long,long long>> primerasPerdedoras(int cantidad) {
			vector<pair<long long,long long>> resultado;
			
			for (int k = 0; k < cantidad; k++) {
				long long ak = (long long)floor(k * PHI);
				long long bk = ak + k;
				resultado.push_back({ak, bk});
			}
			
			return resultado;
		}                                          // O(cantidad)
	};
\end{lstlisting}

\textbf{Ejemplo de funcionamiento:}

\begin{lstlisting}[style=cpp]
	cout << WythoffGame::ganaJugadorEnTurno(1, 2);
	// 0 (false) -- (1,2) es una posicion PERDEDORA (k=1: floor(1*phi)=1,
	//               1+1=2)
	
	cout << WythoffGame::ganaJugadorEnTurno(3, 5);
	// 0 (false) -- tambien perdedora (k=2: floor(2*phi)=3, 3+2=5)
	
	cout << WythoffGame::ganaJugadorEnTurno(2, 4);
	// 1 (true) -- (2,4) NO esta en la lista de perdedoras
	
	auto perdedoras = WythoffGame::primerasPerdedoras(6);
	for (auto& [x, y] : perdedoras) cout << "(" << x << "," << y << ") ";
	// (0,0) (1,2) (3,5) (4,7) (6,10) (8,13)
\end{lstlisting}

\textbf{Nota:} el uso de \texttt{double} para $\varphi$ es seguro para
valores de $a, b$ razonables ($\lesssim 10^{15}$ con \texttt{double}
estándar), pero para valores extremadamente grandes conviene usar
\texttt{long double} o verificar el resultado con aritmética exacta
(reconstruyendo $k$ y comprobando $\lfloor k\varphi \rfloor$ con más
precisión) para evitar errores de redondeo cerca del límite de
representación.

\textbf{Regla práctica:} si un juego "se parece a Nim de 2 montones"
pero permite \textbf{además} quitar la misma cantidad de ambos a la
vez, ya no es Nim —es Wythoff's Game, y la condición de victoria deja
de ser un XOR simple para convertirse en esta fórmula basada en la
razón áurea. Verificar contra la tabla de las primeras posiciones
perdedoras (calculada por fuerza bruta con memoización, para $a,b$
chicos) es la forma más segura de confirmar que la fórmula se está
aplicando bien antes de confiar en ella para valores grandes.
\subsection{Turning Games (Volteo de Monedas / Fichas)}

\textbf{Idea:} una fila de $n$ fichas, cada una mostrando \textbf{cara}
o \textbf{cruz}. En cada turno se elige una ficha en \textbf{cara} y se
voltea \textbf{obligatoriamente} a cruz; además, \textbf{opcionalmente},
se puede voltear (en cualquier sentido, cara$\to$cruz o cruz$\to$cara)
\textbf{una única} ficha ubicada estrictamente a la \textbf{izquierda}
de la elegida. Pierde quien no tiene ninguna cara para voltear. Estos
juegos tienen un teorema notable (de la teoría de \textit{coin-turning
	games}, Berlekamp-Conway-Guy): el valor de Grundy de \textbf{todo el
	tablero} es el \textbf{XOR} de un valor de Grundy que se calcula
\textbf{por separado} para cada ficha en cara, como si esa ficha fuera
la \textbf{única} en el tablero (con todo lo demás en cruz):
$$g(\text{tablero}) = \bigoplus_{i:\ \text{ficha } i \text{ en cara}}
g_1(i)$$
donde $g_1(i)$ es el valor de Grundy de un tablero con una sola cara en
la posición $i$. Esto convierte un juego sobre un espacio de estados
\textbf{exponencial} ($2^n$ configuraciones) en $n$ cálculos
\textbf{independientes} de una sola ficha, combinados con XOR —exactamente
el mismo principio que Suma de Juegos Independientes, aplicado aquí
\textbf{dentro} de un solo tablero en vez de entre tableros distintos.

\textbf{¿Por qué usarlo?} Es la generalización de Staircase Nim a un
marco mucho más amplio: ahí la reducción a Nim salía de un argumento de
espejado específico a escalones; aquí es un \textbf{teorema general}
que aplica a toda una familia de juegos de volteo, siempre que la regla
cumpla dos condiciones: (1) cada movimiento voltea \textbf{como
	mínimo} una cara a cruz en la posición elegida, y (2) cualquier otro
volteo permitido ocurre \textbf{estrictamente a la izquierda} de esa
posición. Bajo esas condiciones, el juego siempre se descompone en
XOR de valores independientes por posición, sin importar el detalle
exacto de qué más se puede voltear. Para la regla descrita arriba (una
ficha extra libre, en cualquier sentido, en cualquier posición a la
izquierda), el resultado es particularmente limpio: $g_1(i) = i+1$
(con $i$ 0-indexado) —es decir, el juego es \textbf{literalmente} Nim
Clásico donde cada cara aporta un montón de tamaño igual a su posición
(1-indexada).

\textbf{Casos de uso:}

\begin{itemize}
	\item Determinar el ganador de cualquier juego de volteo de fichas
	en fila con la regla "obligatorio voltear una cara a cruz, opcional
	voltear una ficha más a la izquierda" — se reduce a Nim sobre las
	posiciones (1-indexadas) de las caras.
	\item Variantes con \textbf{otra} regla de volteo (voltear dos
	fichas a la izquierda, restringir el volteo extra a solo
	cruz$\to$cara, etc.): el teorema de descomposición sigue
	aplicando, solo cambia la fórmula recursiva de $g_1(i)$ (ver
	modificaciones).
	\item Como plantilla para \textbf{descubrir} la fórmula de $g_1$ de
	una regla nueva: computar por fuerza bruta los primeros valores de
	$g_1(0), g_1(1), \dots$ con la recurrencia genérica de Grundy sobre
	un solo tablero, buscar el patrón, y luego usar ese patrón para
	tableros grandes sin necesitar simular el espacio completo de
	estados.
	\item Verificar la propia descomposición contra un cálculo de
	Grundy por fuerza bruta (bitmask + memoización) en tableros chicos,
	antes de confiar en la fórmula para tableros grandes.
\end{itemize}

\textbf{Complejidad:} $O(n)$ por consulta una vez conocida la fórmula o
tabla de $g_1$ (calcular el XOR de las posiciones de las caras);
$O(2^n \cdot n^2)$ para el cálculo de referencia por fuerza bruta sobre
todo el espacio de estados, usado solo para \textbf{verificar} la
descomposición en tableros pequeños.

\begin{lstlisting}[style=cpp]
	struct TurningGame {
		// READ: gana quien esta por mover, dado el tablero como bitmask
		// (bit i = 1 si la ficha en la posicion i, 0-indexada desde la
		// izquierda, esta en CARA). Regla: voltea una cara a cruz
		// (obligatorio) y, opcionalmente, voltea UNA ficha mas en
		// cualquier sentido en una posicion ESTRICTAMENTE a la izquierda.
		// Usa la reduccion a Nim: cada cara en posicion i (0-indexada)
		// aporta un monton de tamano i+1.
		static bool ganaJugadorEnTurno(int tablero, int n) {
			long long xorTotal = 0;
			
			for (int i = 0; i < n; i++)
			if ((tablero >> i) & 1)
			xorTotal ^= (i + 1);   // g1(i) = i + 1 (0-indexada)
			
			return xorTotal != 0;
		}                                          // O(n)
		
		// ------------------------------------------------------------
		// MODIFICACIONES Y COMENTARIOS CON LOS USOS
		// ------------------------------------------------------------
		
		// 1. VERIFICACION POR FUERZA BRUTA (solo para n chico, <= ~18):
		//    calcula el Grundy de TODO el tablero simulando el juego
		//    completo, para confirmar que coincide con el XOR de g1(i).
		//    Sirve para descubrir/confirmar la formula de g1 en una regla
		//    NUEVA antes de confiar en ella:
		//
		//    unordered_map<int,int> memo;
		//    int n;
		//    int grundyCompleto(int mask) {
			//        if (memo.count(mask)) return memo[mask];
			//        set<int> alcanzables;
			//        for (int i = 0; i < n; i++) {
				//            if (!((mask >> i) & 1)) continue;
				//            int base = mask & ~(1 << i);          // i: cara->cruz
				//            alcanzables.insert(grundyCompleto(base));
				//            for (int j = 0; j < i; j++)            // volteo extra
				//                alcanzables.insert(grundyCompleto(base ^ (1 << j)));
				//        }
			//        int mex = 0;
			//        while (alcanzables.count(mex)) mex++;
			//        return memo[mask] = mex;
			//    }
		//    // g1(i) = grundyCompleto(1 << i) con memo limpio por cada i
		
		// 2. VOLTEO EXTRA RESTRINGIDO A CRUZ->CARA (no se permite apagar
		//    otra cara, solo "encender" una cruz): cambia el conjunto de
		//    movimientos disponibles en la posicion j < i, y por lo tanto
		//    cambia g1(i) -- hay que recalcularlo con el metodo de la
		//    modificacion 1 (probablemente ya NO sea i+1 con esta regla
		//    mas restringida).
		
		// 3. SE PUEDEN VOLTEAR HASTA DOS FICHAS A LA IZQUIERDA (en vez de
		//    a lo sumo una): el teorema de descomposicion en XOR SIGUE
		//    aplicando (la condicion clave -- todo volteo extra ocurre
		//    estrictamente a la izquierda de la posicion obligatoria -- no
		//    cambia), pero la recurrencia de g1(i) debe considerar todas
		//    las combinaciones de 0, 1 o 2 fichas volteadas a la
		//    izquierda antes de tomar el mex.
		
		// 4. TABLERO CIRCULAR O SIN ORDEN LINEAL CLARO: el teorema de
		//    descomposicion DEPENDE de que exista una nocion de
		//    "izquierda" que garantice que el juego siempre termina (cada
		//    movimiento reduce alguna medida, como la posicion de la
		//    ficha volteada obligatoriamente). Si el tablero es circular o
		//    la regla permite voltear hacia la derecha, hay que verificar
		//    primero que el juego SIGA siendo finito antes de aplicar nada
		//    de esto.
		
		// 5. IDENTIFICAR RAPIDO SI UN ENUNCIADO ES UN "TURNING GAME": la
		//    señal es "fichas o monedas en fila, cara/cruz, cada movimiento
		//    apaga una cara y opcionalmente cambia algo a su izquierda".
		//    Ante esa estructura, el primer paso NO es simular con
		//    bitmask para n grande -- es calcular g1 con fuerza brutal
		//    para n chico, buscar el patron, y aplicar la formula
		//    encontrada como en el ejemplo de esta subseccion.
	};
\end{lstlisting}

\textbf{Ejemplo de funcionamiento:}

\begin{lstlisting}[style=cpp]
	// tablero de 6 fichas, 0-indexado desde la izquierda:
	// posiciones con CARA: 1, 3, 4  (bits 1, 3 y 4 encendidos)
	int tablero = (1 << 1) | (1 << 3) | (1 << 4);
	
	cout << TurningGame::ganaJugadorEnTurno(tablero, 6);
	// 1 (true)
	// aportes: cara en 1 -> monton 2; cara en 3 -> monton 4;
	// cara en 4 -> monton 5
	// XOR: 2 ^ 4 ^ 5 = 3, distinto de 0 -> gana quien esta por mover
	
	// verificado por fuerza bruta: para TODOS los 64 tableros posibles
	// de 6 fichas, el grundy calculado simulando el juego completo
	// coincide exactamente con el XOR de (posicion + 1) de cada cara
\end{lstlisting}

\textbf{Nota:} la posición se cuenta \textbf{0-indexada} en el bitmask
(bit $0$ = ficha más a la izquierda), pero el "tamaño de montón" que
aporta cada cara es $i+1$ (equivalente a numerar las posiciones desde
$1$). Es un detalle de conteo fácil de invertir por error —conviene
fijar una convención (0-indexada en el código, 1-indexada en la
fórmula) y no mezclarlas.

\textbf{Regla práctica:} ante cualquier juego de fichas cara/cruz en
fila donde cada movimiento apaga una cara y opcionalmente modifica algo
estrictamente a su izquierda, no analices el tablero completo como un
solo juego monolítico —calcula (por fuerza bruta, para $n$ chico) el
valor de Grundy de \textbf{una sola} cara en cada posición, confirma
que el patrón se sostiene, y combina con XOR. La regla exacta de qué se
puede voltear a la izquierda cambia \textbf{la fórmula} de $g_1$, pero
prácticamente nunca cambia el hecho de que la descomposición en XOR
funciona.

\section{Juegos sobre Estructuras de Datos}
\subsection{Juegos sobre Grafos y DAGs (Movimiento de una Ficha)}

\textbf{Idea:} una ficha está en un vértice de un \textbf{DAG} (grafo
dirigido acíclico). En cada turno, un jugador mueve la ficha a lo
largo de \textbf{una} arista saliente. Pierde quien no tiene ninguna
arista disponible (la ficha está en un vértice sin salidas). Es el
mismo cálculo de Grundy que ya se usa en otras partes de esta guía
(ver Cálculo de Grundy sobre un DAG), aplicado aquí como
\textbf{modelo general}: $g(v) = \operatorname{mex}\{\,g(u) :
(v,u)\in E\,\}$, calculado con DFS memoizado (o en orden topológico
inverso). El jugador en turno gana con la ficha en $v$ si y solo si
$g(v) \ne 0$. Este mismo esquema es, de hecho, la
\textbf{representación universal} de cualquier juego imparcial: todo
juego imparcial finito y sin empates es, en esencia, un DAG de
posiciones con este mismo cálculo de Grundy — Nim, Subtraction Games,
Staircase Nim y Turning Games son casos particulares donde el DAG
tiene una estructura especial que permite una fórmula cerrada en vez
de recorrer el grafo explícitamente.

\textbf{¿Por qué usarlo?} Cuando un juego \textbf{no} encaja en ningún
patrón con fórmula cerrada, pero sus posiciones y transiciones se
pueden modelar como un DAG explícito (con $V$ y $E$ manejables), el
cálculo de Grundy directo sobre ese DAG siempre funciona —es el
método "de último recurso" antes de rendirse ante un juego sin
estructura evidente. Además, con \textbf{varias fichas} en el mismo
DAG o en DAGs distintos, moviendo \textbf{una sola} por turno, el
juego total es la \textbf{suma} de juegos independientes de Sprague-Grundy:
el Grundy total es el XOR de los Grundy de cada ficha por
separado, exactamente igual que combinar montones de Nim.

\textbf{Casos de uso:}

\begin{itemize}
	\item Cualquier juego descrito como "una ficha se mueve por un
	grafo/tablero siguiendo reglas de movimiento fijas, sin poder
	repetir posiciones" (la ausencia de ciclos es la que garantiza que
	el juego termine).
	\item Varias fichas en el mismo grafo (o en grafos distintos),
	donde cada turno se mueve exactamente una: se resuelve calculando
	el Grundy de la posición de \textbf{cada} ficha por separado y
	combinando con XOR — no hace falta un estado conjunto exponencial.
	\item Como técnica de \textbf{traducción}: cuando un juego nuevo
	no encaja obviamente en Nim/Sustracción/Staircase/Turning, modelar
	sus posiciones explícitamente como nodos de un DAG y aplicar esta
	fórmula genérica siempre es una opción válida, aunque menos
	elegante que una fórmula cerrada.
	\item \textbf{Advertencia importante:} si el grafo es
	\textbf{no dirigido} y el juego consiste en mover la ficha y
	\textbf{borrar} el vértice/arista visitado (\textit{Generalized
		Geography}), el estado ya \textbf{no} es solo "dónde está la
	ficha" —el grafo cambia con cada movimiento, y este cálculo simple
	de Grundy \textbf{no aplica}: ese juego es, en general,
	computacionalmente mucho más difícil (PSPACE-completo), y
	requiere memoizar sobre el grafo restante, no solo sobre el
	vértice actual.
\end{itemize}

\textbf{Complejidad:} $O(V + E)$ para calcular el Grundy de todos los
vértices de un DAG (cada vértice se calcula una sola vez con
memoización); $O(k)$ por consulta de $k$ fichas independientes, una
vez calculada la tabla de Grundy.

\begin{lstlisting}[style=cpp]
	struct JuegoFichaDAG {
		vector<vector<int>> adj;   // adj[v] = vecinos alcanzables desde v
		vector<int> grundy;
		vector<bool> calculado;
		int n;
		
		// CREATE: guarda el DAG (debe ser ACICLICO; con ciclos, la
		// memoizacion de abajo no termina).
		JuegoFichaDAG(vector<vector<int>>& adj_, int n_) {
			adj = adj_;
			n = n_;
			grundy.assign(n, -1);
			calculado.assign(n, false);
		}                                          // O(1)
		
		// READ: valor de Grundy del vertice v (memoizado).
		int calcular(int v) {
			if (calculado[v]) return grundy[v];
			calculado[v] = true;   // se marca ANTES de recursar: en un
			// DAG no hay ciclos, asi que esto es
			// seguro y evita recalculo
			
			set<int> alcanzables;
			for (int u : adj[v])
			alcanzables.insert(calcular(u));
			
			int mex = 0;
			while (alcanzables.count(mex)) mex++;
			
			return grundy[v] = mex;
		}                                          // O(1) amortizado
		// (O(V+E) la primera vez)
		
		// READ: con k fichas en las posiciones dadas (una por turno se
		// mueve), gana quien esta en turno?
		bool ganaConFichas(vector<int>& posiciones) {
			int xorTotal = 0;
			for (int p : posiciones)
			xorTotal ^= calcular(p);
			return xorTotal != 0;
		}                                          // O(k)
	};
\end{lstlisting}

\textbf{Ejemplo de funcionamiento:}

\begin{lstlisting}[style=cpp]
	// DAG:  0 -> 1, 0 -> 2, 1 -> 3, 2 -> 3   (3 es sumidero, sin salidas)
	vector<vector<int>> adj(4);
	adj[0] = {1, 2};
	adj[1] = {3};
	adj[2] = {3};
	adj[3] = {};
	
	JuegoFichaDAG jf(adj, 4);
	jf.calcular(0); jf.calcular(1); jf.calcular(2); jf.calcular(3);
	
	for (int v = 0; v < 4; v++) cout << "g(" << v << ")=" << jf.grundy[v] << " ";
	// g(0)=0 g(1)=1 g(2)=1 g(3)=0
	
	vector<int> fichas = {0, 1};
	cout << jf.ganaConFichas(fichas);
	// 1 (true) -- g(0) XOR g(1) = 0 XOR 1 = 1, distinto de 0
\end{lstlisting}

\textbf{Nota:} verificado por fuerza bruta (minimax directo sobre el
estado conjunto de varias fichas) en un DAG aleatorio de 8 nodos: el
resultado de XOR de Grundy individuales coincide exactamente con la
simulación completa para 1, 2 y 3 fichas simultáneas — confirma que la
suma de juegos independientes aplica sin ninguna sorpresa cuando el
grafo mismo no cambia entre turnos.

\textbf{Regla práctica:} este método es el "más general pero más
costoso" de toda la sección de Game Theory —úsalo cuando el juego no
tenga una estructura reconocible (Nim, Sustracción, Staircase,
Turning) pero sí se pueda modelar como un DAG explícito de tamaño
manejable. Antes de aplicarlo, confirma que el grafo realmente
\textbf{no cambia} durante la partida (solo la posición de la ficha
cambia) — si el grafo se modifica con cada movimiento, este cálculo no
es válido y hace falta memoizar sobre el estado completo del grafo
restante.


\subsection{Juegos sobre Árboles (Green Hackenbush / Poda de Ramas)}

\textbf{Idea:} un árbol de "ramas" (aristas) cuelga desde el
\textbf{suelo} (la raíz). En cada turno, un jugador \textbf{corta} una
arista cualquiera; esa arista desaparece, y con ella \textbf{toda} la
parte del árbol que quedaba \textbf{colgando} de su extremo inferior
(al perder la conexión con el suelo, esa subrama "cae" y sale del
juego de inmediato, en el mismo movimiento). Pierde quien no tiene
ninguna arista para cortar. El espacio de estados de este juego es
\textbf{exponencial} (cada subconjunto de aristas restantes es un
estado distinto), pero existe una fórmula recursiva —el
\textbf{Colon Principle}— que evita simularlo: para cada arista $e$
que conecta un vértice $v$ con su padre, sea $c_1,\dots,c_k$ las
aristas que cuelgan \textbf{directamente} del extremo inferior de $e$
(las aristas hacia los hijos de $v$); el valor de Grundy de $e$ es
$$g(e) = 1 + \big(g(c_1) \oplus g(c_2) \oplus \cdots \oplus g(c_k)\big)$$
con $g(e) = 1$ si $e$ es una hoja (sin hijos, la suma XOR vacía es
$0$). El Grundy de \textbf{todo} el árbol (o bosque, si hay varias
ramas que salen directo del suelo) es el XOR de los $g$ de esas ramas
raíz.

\textbf{¿Por qué usarlo?} Es la aplicación más elegante de Sprague-Grundy
de esta guía: convierte un juego con estados exponenciales en un
simple recorrido \textbf{post-orden} del árbol, $O(\text{aristas})$
total. La fórmula tiene una lectura intuitiva: un \textbf{tallo}
simple (un camino sin ramificaciones, de $L$ aristas colgando en
línea recta) tiene valor de Grundy exactamente $L$ —cada arista
adicional en la cadena suma $1$ al valor anterior, ya que no hay
ramificación que combinar (el XOR de un solo valor es ese mismo
valor). La ramificación es lo que introduce el XOR: cuando un vértice
tiene varias subramas, primero se combinan \textbf{esas} subramas
entre sí (como si fueran montones de Nim independientes), y \textbf{después}
se suma $1$ por la arista que las conecta hacia arriba.

\textbf{Casos de uso:}

\begin{itemize}
	\item Green Hackenbush jugado sobre un árbol o un bosque de
	árboles (todas las ramas son "neutrales", cualquier jugador puede
	cortar cualquier arista —a diferencia de Red-Blue Hackenbush,
	donde cada rama pertenece a un jugador, fuera del alcance de esta
	guía).
	\item Cualquier juego de "poda de ramas" donde cortar una conexión
	elimina automáticamente todo lo que queda aislado —el mismo
	principio aplica sin importar el nombre narrativo (árboles reales,
	jerarquías, estructuras de dependencia que se desmoronan).
	\item Calcular el valor de Grundy de un bosque completo sumando
	(XOR) los valores de cada árbol individual —cada árbol es un
	juego independiente en el sentido de Sprague-Grundy.
	\item Como ejemplo de "colapsar una estructura compleja a un solo
	número" antes de decidir quién gana, reforzando la idea general de
	Sprague-Grundy de que \textbf{cualquier} juego imparcial se reduce,
	tarde o temprano, a un solo entero no negativo.
\end{itemize}

\textbf{Complejidad:} $O(V)$ — un recorrido post-orden del árbol,
calculando el Grundy de cada arista una sola vez a partir de sus
hijas ya calculadas.

\begin{lstlisting}[style=cpp]
	struct GreenHackenbushArbol {
		vector<vector<int>> hijosAristas;   // hijosAristas[e] = aristas
		// que cuelgan directamente
		// del extremo inferior de e
		vector<int> raicesAristas;          // aristas que salen directo
		// del suelo (raiz del bosque)
		
		// CREATE: guarda la estructura de aristas del arbol/bosque, ya
		// representada como "que aristas cuelgan de cada arista".
		GreenHackenbushArbol(vector<vector<int>>& hijosAristas_,
		vector<int>& raicesAristas_) {
			hijosAristas = hijosAristas_;
			raicesAristas = raicesAristas_;
		}                                          // O(1)
		
		// READ: valor de Grundy de la arista e (Colon Principle):
		// combina las aristas hijas con XOR, y suma 1 por la propia arista.
		int colon(int e) {
			int x = 0;
			for (int c : hijosAristas[e])
			x ^= colon(c);
			return x + 1;
		}                                          // O(aristas del subarbol)
		
		// READ: gana quien esta por mover, sobre TODO el bosque?
		bool ganaJugadorEnTurno() {
			int total = 0;
			for (int r : raicesAristas)
			total ^= colon(r);
			return total != 0;
		}                                          // O(V)
	};
\end{lstlisting}

\textbf{Ejemplo de funcionamiento:}

\begin{lstlisting}[style=cpp]
	// arbol colgando del suelo (nodo 0):
	//        0 (suelo)
	//       / \
	//      1   3
	//      |  / \
	//      2 4   5
	//
	// aristas: 0=(0-1), 1=(1-2), 2=(0-3), 3=(3-4), 4=(3-5)
	
	vector<vector<int>> hijosAristas(5);
	hijosAristas[0] = {1};       // la arista 0-1 tiene colgando la 1-2
	hijosAristas[1] = {};        // hoja
	hijosAristas[2] = {3, 4};    // la arista 0-3 tiene colgando 3-4 y 3-5
	hijosAristas[3] = {};        // hoja
	hijosAristas[4] = {};        // hoja
	
	vector<int> raices = {0, 2};   // ambas ramas salen directo del suelo
	
	GreenHackenbushArbol gh(hijosAristas, raices);
	
	cout << gh.colon(0);
	// 2   (tallo de 2 aristas: arista 1 vale 1, arista 0 = 1 + 1 = 2)
	
	cout << gh.colon(2);
	// 1   (arista 3 vale 1, arista 4 vale 1; XOR = 0; arista 2 = 0 + 1 = 1)
	
	cout << gh.ganaJugadorEnTurno();
	// 1 (true) -- total = colon(0) XOR colon(2) = 2 XOR 1 = 3, distinto de 0
\end{lstlisting}

\textbf{Nota:} verificado con fuerza bruta —simulando el juego
completo sobre un bitmask de aristas (donde cortar una arista borra
también, de una sola vez, todas las aristas de su subárbol) en 200
árboles aleatorios de hasta 9 aristas, el Colon Principle coincidió
exactamente con el valor de Grundy calculado simulando \textbf{todo}
el espacio de estados. El caso de un solo tallo (sin ramificación) es
el que confirma la intuición del $+1$: cada eslabón de la cadena
extiende el valor en una unidad, exactamente como un montón de Nim
que crece de a uno.

\textbf{Regla práctica:} para cualquier variante de Green Hackenbush
sobre \textbf{árboles} (nunca sobre grafos con ciclos —ahí el juego se
vuelve mucho más complejo y esta fórmula no aplica), calcula el
Colon Principle con un solo recorrido post-orden: combina las ramas
hijas con XOR, suma $1$ por la arista que las conecta hacia arriba, y
combina las raíces del bosque con otro XOR final. Nunca simules el
espacio de estados completo salvo para \textbf{verificar} la fórmula
en árboles de prueba pequeños.
\subsection{Juegos de Tomar de los Extremos de un Arreglo (DP de Intervalos)}

\textbf{Idea:} dos jugadores se turnan tomando el elemento del
\textbf{extremo izquierdo o derecho} de un arreglo, sumando su valor al
propio puntaje; el arreglo se agota y gana quien termine con más puntos.
A diferencia de todos los juegos de \textit{Fundamentos} (Nim, Grundy,
Win/Lose), aquí \textbf{no hay un ``ganador'' binario definido por quién
	mueve al final} — es un juego con \textbf{puntaje}, donde ambos jugadores
tienen exactamente las mismas opciones de movimiento (tomar izquierda o
derecha) y solo difiere lo que cada uno se queda. Por eso no se resuelve
con Sprague-Grundy (que exige un criterio de victoria normal-play), sino
con \textbf{minimax}: como todo lo que un jugador no toma queda para el
rival, maximizar tu propio puntaje equivale a maximizar
\textit{(tu puntaje $-$ el del rival)}. Se define $dp[i][j]$ como esa
diferencia óptima jugando solo sobre el subarreglo $a[i..j]$, con la
recurrencia
$$dp[i][j] = \max\big(a[i] - dp[i+1][j],\; a[j] - dp[i][j-1]\big),$$
porque si tomas $a[i]$, el rival pasa a ser el que mueve primero en
$[i+1,j]$ y su ventaja en ese subjuego se te resta.

\textbf{¿Por qué usarlo?} Es la generalización de \textit{Turning Games}
(esa trabaja con paridad de fichas; esta con valores arbitrarios) y el
antecedente directo de \textit{Juegos de Turnos con Puntaje (Optimal
	Strategy for a Game)} en \textit{Juegos Partidistas y Combinatoria
	Avanzada} — mismo problema, distinto contexto de la guía. Se distingue de
\textit{Juegos sobre Árboles (Green Hackenbush)} en que ahí ambos
jugadores tienen jugadas \textbf{distintas} (Izquierdo/Derecho) sobre una
estructura que se desconecta permanentemente al remover una arista; aquí
las jugadas son \textbf{simétricas} (ambos pueden tomar de cualquier
extremo) y lo que cambia es el puntaje, no las reglas de movimiento.

\textbf{Casos de uso:}

\begin{itemize}
	\item El clásico ``Optimal Strategy for a Game'' / ``Predict the
	Winner'': determinar si el primer jugador puede ganar o empatar.
	\item Juegos de cartas en fila donde cada carta tiene un valor y solo
	se puede tomar una de las dos puntas de la fila.
	\item Calcular el puntaje exacto de cada jugador bajo juego óptimo, no
	solo quién gana (usando la suma total y la diferencia óptima).
	\item Variantes donde ambos extremos tienen distinto ``precio'' de
	oportunidad (agregar un costo o descuento al tomar de un lado).
	\item Como bloque base antes de generalizar a valores de juego
	partidista (Números Surreales, más adelante en la guía).
\end{itemize}

\textbf{Complejidad:} hay $O(n^2)$ subintervalos $[i,j]$ y cada uno se
calcula en $O(1)$ a partir de dos subintervalos más cortos, para un total
de $O(n^2)$ en tiempo y $O(n^2)$ en memoria si se guarda la tabla
completa. Procesando por \textbf{longitud creciente} de intervalo (todos
los de longitud 1, luego 2, etc.) se puede reducir la memoria a $O(n)$
con un arreglo que se sobrescribe, porque $dp[i][j]$ solo depende de
intervalos de longitud $\text{largo}-1$.

\begin{lstlisting}[style=cpp]
	struct JuegoDeExtremos {
		int n;
		vector<long long> a;
		vector<vector<long long>> dp;   // dp[i][j] = diferencia optima
		// (turno - rival) jugando sobre a[i..j]
		
		// CREATE: guarda los valores del arreglo.
		JuegoDeExtremos(vector<long long> valores) {
			a = valores;
			n = (int)a.size();
			dp.assign(n, vector<long long>(n, 0));
		}                                          // O(n^2)
		
		// WRITE: llena dp por longitud de intervalo creciente. Caso base:
		// con un solo elemento, la diferencia es simplemente su valor.
		void calcular() {
			for (int i = 0; i < n; i++) dp[i][i] = a[i];
			for (int largo = 2; largo <= n; largo++) {
				for (int i = 0; i + largo - 1 < n; i++) {
					int j = i + largo - 1;
					dp[i][j] = max(a[i] - dp[i+1][j], a[j] - dp[i][j-1]);
				}
			}
		}                                          // O(n^2)
		
		// READ: diferencia optima (turno - rival) sobre TODO el arreglo.
		// Positiva: el primer jugador puede ganar (o empatar si es 0).
		long long diferencia() const { return dp[0][n-1]; }
		
		// READ: puntaje exacto de cada jugador bajo juego optimo, usando
		// que puntaje1 + puntaje2 = suma total y puntaje1 - puntaje2 = diferencia.
		pair<long long,long long> puntajes() const {
			long long total = 0; for (long long x : a) total += x;
			long long diff = dp[0][n-1];
			return {(total + diff) / 2, (total - diff) / 2};
		}                                          // O(n)
	};
	
	// ------------------------------------------------------------
	// MODIFICACIONES Y COMENTARIOS CON LOS USOS
	// ------------------------------------------------------------
	
	// 1. REDUCIR A O(n) DE MEMORIA: como dp[i][j] con largo L solo usa
	//    valores de largo L-1, basta un arreglo 1D que se recalcula por
	//    longitud, en vez de guardar la tabla n x n completa:
	//
	//    vector<long long> dp(n);
	//    for (int i = 0; i < n; i++) dp[i] = a[i];
	//    for (int largo = 2; largo <= n; largo++)
	//        for (int i = 0; i + largo - 1 < n; i++) {
		//            int j = i + largo - 1;
		//            dp[i] = max(a[i] - dp[i+1], a[j] - dp[i]);   // dp[i] viejo = dp[i][j-1]
		//        }
	//    // el resultado queda en dp[0]
	//
	//    Ojo con el orden: hay que leer dp[i] ANTES de sobrescribirlo, ya
	//    que en la formula original dp[i][j-1] usa el valor VIEJO de dp[i].
	
	// 2. TRUCO DE PARIDAD (solo indica quien NO PIERDE, no el margen
	//    optimo): si n es PAR, el primer jugador puede garantizar
	//    puntaje1 >= puntaje2 tomando SIEMPRE todos los elementos de
	//    indice par o siempre todos los de indice impar (el que sume
	//    mas), porque esas dos clases nunca se mezclan sin importar lo
	//    que haga el rival. Esto es solo una COTA INFERIOR de la
	//    diferencia optima real, NO el resultado de dp[0][n-1]:
	//
	//    long long pares = 0, impares = 0;
	//    for (int i = 0; i < n; i++) (i % 2 == 0 ? pares : impares) += a[i];
	//    // el primer jugador nunca pierde si n es par; margen GARANTIZADO
	//    // (no necesariamente optimo) = max(pares, impares) - min(pares, impares)
	//
	//    Sirve para responder rapido "¿puede el primer jugador NO PERDER?"
	//    sin correr el DP, pero si el enunciado pide el margen EXACTO o n
	//    es impar, hay que usar dp completo.
	
	// 3. COSTO O BONUS POR EXTREMO: si tomar de la izquierda o la derecha
	//    tiene un costo/bonus adicional c_izq o c_der (por ejemplo, un
	//    "impuesto" fijo), se suma directamente dentro del max:
	//    dp[i][j] = max((a[i] - c_izq) - dp[i+1][j], (a[j] - c_der) - dp[i][j-1]);
	
	// 4. K TURNOS DESDE CADA EXTREMO (variante distinta, NO minimax):
	//    si en vez de un turno por elemento el problema pide "elige k
	//    elementos tomando solo de los extremos, maximiza la suma" (un
	//    solo jugador, no un juego de dos), se resuelve con dos punteros y
	//    prefijos: probar cuantos se toman de la izquierda (0..k) y el
	//    resto de la derecha, en O(k). Es un problema distinto - de
	//    optimizacion para un solo jugador, no un juego de dos jugadores.
	
	// 5. TRAMPA: CONFUNDIR ESTE JUEGO CON UNO DE GANAR/PERDER. Aqui NO
	//    hay "el que mueve al final gana"; ambos maximizan su propio
	//    puntaje, y el resultado que importa es la diferencia final de
	//    puntajes (o simplemente quien saco mas). Aplicar Sprague-Grundy
	//    (pensado para juegos normal-play de "el ultimo en mover gana")
	//    no tiene sentido aqui: no hay estados terminales "P" o "N", solo
	//    un valor numerico que se maximiza.
	
	// 6. TRAMPA: dp[i][j-1] Y dp[i+1][j] SON DIFERENCIAS DEL RIVAL, NO DEL
	//    ARREGLO RESTANTE PARA TI. Al escribir a[i] - dp[i+1][j], el signo
	//    negativo representa que dp[i+1][j] es la ventaja del jugador que
	//    mueve DESPUES de ti sobre ese subarreglo (que ahora es tu rival);
	//    olvidar el signo negativo es el error mas comun al implementar
	//    esta recurrencia por primera vez.
\end{lstlisting}

\textbf{Ejemplo de funcionamiento:}

\begin{lstlisting}[style=cpp]
	vector<long long> a = {1, 5, 2, 4, 6};
	JuegoDeExtremos j(a);
	j.calcular();
	
	cout << j.diferencia();          // 0  (empate bajo juego optimo)
	auto [p1, p2] = j.puntajes();
	cout << p1 << " " << p2;          // 9 9
	
	// n par: el primer jugador puede garantizar no perder.
	vector<long long> b = {8, 15, 3, 7};
	JuegoDeExtremos jb(b);
	jb.calcular();
	cout << jb.diferencia();         // 11  (el margen OPTIMO, mayor que la garantia de paridad)
\end{lstlisting}

\textbf{Nota:} en el segundo ejemplo, el truco de paridad de la
modificación 2 solo garantiza que el primer jugador \textbf{no pierde}
(margen $\geq$ \texttt{max(pares, impares) - min(pares, impares)} $= 11$
en este caso, que coincide, pero no siempre coincide con el óptimo real):
con juego verdaderamente óptimo la diferencia puede ser \textbf{mayor} que
esa garantía. Si el enunciado pide el margen exacto de puntos (no solo
quién gana), siempre hay que correr el DP completo.

\textbf{Regla práctica:} usa esta interval DP en cuanto el problema
describa ``toma de un extremo u otro, maximizando tu propio puntaje'' —
es la firma exacta de este patrón. Si $n$ es par y solo preguntan
\textit{si} el primer jugador puede evitar perder, el truco de paridad en
$O(n)$ responde sin DP; para el margen exacto, o si $n$ es impar, usa
$dp[i][j] = \max(a[i]-dp[i{+}1][j],\, a[j]-dp[i][j{-}1])$ recorriendo por
longitud de intervalo creciente.

\subsection{Juegos sobre Strings (Eliminar Substrings o Caracteres)}

\textbf{Idea:} un juego sobre strings define un movimiento como
\textbf{eliminar algo} de la cadena (un carácter, un par, un patrón) según
una regla, y ambos jugadores alternan hasta que nadie pueda mover
(\textit{normal play}: el último en mover gana). El ingrediente
específico de esta familia es que, al eliminar una porción del
\textbf{medio} de la cadena, las dos partes que quedan se convierten en
\textbf{dos subjuegos independientes} — exactamente como remover una
arista en \textit{Juegos sobre Árboles (Green Hackenbush)} desconecta el
árbol en dos componentes que ya nunca vuelven a interactuar. Esa
independencia es lo que permite aplicar directamente \textit{Suma de
	Juegos Independientes (XOR de Grundy Numbers)}: el valor de Grundy de la
cadena completa es el \texttt{mex} (\textit{Cálculo de Grundy sobre un
	DAG}) sobre todos los movimientos posibles, y el valor de cada movimiento
es el \textbf{XOR} de los Grundy de las dos partes resultantes —
$g(l,r) = \operatorname{mex}_{\text{movimientos válidos}} \big(g(l,i) \oplus g(i',r)\big)$.

\textbf{¿Por qué usarlo?} La independencia de las dos partes \textbf{no
	es automática} — es una propiedad del diseño del juego, y hay que
verificarla antes de aplicar la suma de Grundy. Por ejemplo, si el
juego dijera literalmente ``elimina dos caracteres iguales adyacentes y
\textbf{concatena} lo que queda'' (como en el clásico problema de reducir
un string), las dos partes \textbf{sí pueden volver a tocarse}: el
carácter que quedó justo a la izquierda del hueco y el que quedó justo a
la derecha ahora son vecinos, y si coinciden, se habilita un movimiento
nuevo que \textbf{cruza la frontera} entre lo que se pensaba eran dos
subjuegos separados — rompiendo por completo la descomposición en XOR.
Por eso esta subsección trabaja con la variante donde las dos partes
quedan \textbf{separadas para siempre} (sin concatenar): un diseño de
juego perfectamente válido y el único que la suma de Grundy resuelve de
forma directa; la variante con concatenación existe en la práctica, pero
exige técnicas mucho más específicas caso por caso, no la suma general de
Sprague-Grundy.

\textbf{Casos de uso:}

\begin{itemize}
	\item Remover un par de caracteres adyacentes iguales, dejando las
	dos mitades como cadenas independientes que ya no se tocan.
	\item Remover un carácter específico (o que cumpla una condición),
	partiendo la cadena en el punto de la remoción.
	\item Juegos sobre secuencias de símbolos (no necesariamente texto:
	listas de colores, secuencias de eventos) con la misma estructura de
	partición.
	\item Verificar, antes de aplicar Grundy, si el movimiento del
	enunciado realmente separa las partes para siempre o si permite
	fusiones — la pregunta que decide si esta técnica aplica.
	\item Como puente hacia \textit{Suma de Juegos Partidistas y
		Comparación de Valores}, más adelante en la guía, cuando el juego deja
	de ser impartial.
\end{itemize}

\textbf{Complejidad:} hay $O(n^2)$ subintervalos $[l,r)$, y calcular
$g(l,r)$ revisa hasta $O(n)$ movimientos posibles dentro del intervalo,
cada uno con costo $O(1)$ dado que $g$ de los subintervalos ya está
memoizado; el \texttt{mex} sobre a lo sumo $O(n)$ valores candidatos
cuesta $O(n)$ con un \texttt{set}. En total, $O(n^3)$ en tiempo y
$O(n^2)$ en memoria para la memoización completa.

\begin{lstlisting}[style=cpp]
	struct JuegoDeCadenas {
		string s;
		int n;
		map<pair<int,int>, int> memo;   // grundy de s[l..r), SIN fusion tras remover
		
		// CREATE: guarda la cadena.
		JuegoDeCadenas(const string& cadena) { s = cadena; n = (int)s.size(); }
		
		// READ: valor de Grundy del subjuego sobre s[l..r). Un movimiento
		// valido elimina el par s[i], s[i+1] (iguales y adyacentes) y
		// convierte s[l..i) y s[i+2..r) en DOS SUBJUEGOS INDEPENDIENTES
		// (nunca se vuelven a tocar).
		int grundy(int l, int r) {
			if (r - l < 2) return 0;              // 0 o 1 caracter: sin movimientos
			auto key = make_pair(l, r);
			auto it = memo.find(key);
			if (it != memo.end()) return it->second;
			set<int> alcanzables;
			for (int i = l; i + 1 < r; i++) {
				if (s[i] == s[i+1]) {
					int izq = grundy(l, i);
					int der = grundy(i+2, r);
					alcanzables.insert(izq ^ der);   // Sprague-Grundy: XOR de partes independientes
				}
			}
			int mex = 0;
			while (alcanzables.count(mex)) mex++;
			memo[key] = mex;
			return mex;
		}                                          // O(n) por llamada, O(n^3) total
		
		// READ: true si el primer jugador gana bajo juego optimo (normal play).
		bool primeroGana() { return grundy(0, n) != 0; }
	};
	
	// ------------------------------------------------------------
	// MODIFICACIONES Y COMENTARIOS CON LOS USOS
	// ------------------------------------------------------------
	
	// 1. VARIAS PARTIDAS INDEPENDIENTES A LA VEZ: si el problema da varias
	//    cadenas y el turno consiste en mover en UNA de ellas (regla
	//    normal de suma de juegos), el primer jugador gana si y solo si
	//    XOR de grundy(0, len(t)) sobre todas las cadenas t es distinto de
	//    cero - exactamente Suma de Juegos Independientes aplicada aqui.
	
	// 2. OTRA REGLA DE ELIMINACION: si el movimiento es "elimina un
	//    caracter que cumple una condicion P(c)" en vez de un par igual
	//    adyacente, el cambio es minimo:
	//    for (int i = l; i < r; i++)
	//        if (P(s[i])) alcanzables.insert(grundy(l, i) ^ grundy(i+1, r));
	//    (aqui SI se puede partir en un solo caracter, sin exigir pares).
	
	// 3. DETECTAR SI LA FUSION ES POSIBLE ANTES DE CODIFICAR: antes de
	//    aplicar esta tecnica, revisa el enunciado con una pregunta
	//    concreta: "tras eliminar la porcion, ¿el caracter inmediatamente
	//    a la izquierda queda junto al caracter inmediatamente a la
	//    derecha?" Si la respuesta es si (la cadena se CONCATENA), la
	//    suma de Grundy en general NO aplica y hace falta otra tecnica
	//    (memoizar sobre la cadena real, no sobre indices del original).
	
	// 4. REFERENCIA (variante CON fusion, solo para cadenas chicas): si el
	//    enunciado exige que las partes se concatenen, la unica forma
	//    segura de resolverlo en general es memoizar directamente sobre
	//    el STRING RESULTANTE (no sobre indices del original), ya que el
	//    estado del juego deja de ser "dos piezas independientes" para
	//    ser una sola cadena nueva:
	//
	//    map<string,int> memoDirecto;      // gana (true/false), NO grundy
	//    bool gana(const string& t) {
		//        if (memoDirecto.count(t)) return memoDirecto[t];
		//        bool g = false;
		//        for (int i = 0; i+1 < (int)t.size() && !g; i++)
		//            if (t[i]==t[i+1] && !gana(t.substr(0,i) + t.substr(i+2))) g = true;
		//        return memoDirecto[t] = g;
		//    }
	//
	//    Esto es minimax directo sobre el string completo, NO una suma de
	//    Grundy (no hay descomposicion en piezas independientes), y solo
	//    es viable para cadenas cortas por el numero de estados distintos
	//    alcanzables.
	
	// 5. TRAMPA (la mas importante de esta subseccion): NUNCA asumas que
	//    "las dos partes quedan separadas" solo porque el enunciado diga
	//    "elimina un pedazo del medio" - verifica explicitamente si el
	//    juego indica que el resto se JUNTA (concatena) o queda separado.
	//    Aplicar la suma de Grundy a un juego que en realidad concatena
	//    produce un resultado incorrecto sin ningun error de compilacion
	//    o ejecucion que lo delate; la unica forma de detectarlo es
	//    releer la regla del movimiento con cuidado, o verificar contra
	//    fuerza bruta en cadenas pequenas antes de confiar en la formula.
\end{lstlisting}

\textbf{Ejemplo de funcionamiento:}

\begin{lstlisting}[style=cpp]
	JuegoDeCadenas j1("aabb");
	cout << j1.grundy(0, 4) << " " << j1.primeroGana();
	// 0 0   (el primer jugador pierde)
	
	JuegoDeCadenas j2("aabbaa");
	cout << j2.grundy(0, 6) << " " << j2.primeroGana();
	// 1 1   (el primer jugador gana)
	
	// Dos cadenas independientes en la misma partida: se combinan con XOR.
	JuegoDeCadenas t1("aabb"), t2("aabbaa");
	int combinado = t1.grundy(0, 4) ^ t2.grundy(0, 6);
	cout << combinado;             // 0 ^ 1 = 1   (el primer jugador gana la suma)
\end{lstlisting}

\textbf{Regla práctica:} antes de escribir una sola línea de código,
confirma si el movimiento del enunciado \textbf{concatena} las partes
restantes o las deja \textbf{separadas para siempre}. Si quedan
separadas, esta suma de Grundy en $O(n^3)$ resuelve el juego directamente
y se combina con otras cadenas vía XOR. Si se concatenan, esta técnica no
aplica: hay que memoizar sobre el string real resultante (viable solo para
entradas pequeñas) o buscar una propiedad específica del problema — nunca
asumas la independencia sin verificarla.

\section{Juegos Partidistas y Combinatoria Avanzada}

\subsection{Números Surreales y Valores de Juego (Cold, Hot, Fuzzy)}

\textbf{Idea:} todo lo visto en \textit{Fundamentos} (estados P/N,
Sprague-Grundy) asume juegos \textbf{impartiales}: en cualquier posición,
\textbf{ambos} jugadores tienen exactamente el mismo conjunto de
movimientos disponibles. Un juego \textbf{partidista} rompe esa simetría:
Izquierda y Derecha pueden tener opciones distintas (Hackenbush con
aristas de dos colores, Domineering, finales de Go). La teoría de números
surreales generaliza el mex/Grundy para este caso representando
\textbf{cualquier} juego como un par de conjuntos de opciones,
$G = \{G^L \mid G^R\}$: $G^L$ es todo lo que Izquierda puede jugar,
$G^R$ todo lo que Derecha puede jugar, y cada opción es a su vez otro
juego construido antes que $G$ (la construcción es bien fundada, sin
ciclos). Comparar dos juegos ya no usa \texttt{mex} sino la definición
recursiva
$$G \le H \iff (\text{ningún } G^L \ge H) \ \text{y}\ (\text{ningún } H^R \le G),$$
y clasificar $G$ frente al juego vacío $0 = \{\mid\}$ dice quién gana:
$G>0$ Izquierda gana siempre; $G<0$ Derecha gana siempre; $G=0$ gana
quien mueve \textbf{segundo} (como en Grundy=0); y $G \parallel 0$
(\textbf{difuso}, confundido con 0) gana quien mueve
\textbf{primero} — el análogo partidista de Grundy $\neq 0$.

\textbf{¿Por qué usarlo?} Los juegos que resultan ser \textbf{números}
(todo $G^L$ estrictamente menor que todo $G^R$, recursivamente) se
comportan como una ventaja fija y \textbf{fría}: a nadie le conviene
jugar ahí porque es puro desperdicio de turno, igual que dar un peón de
handicap. Los juegos \textbf{calientes} son lo opuesto: ambos jugadores
quieren mover ahí \textbf{ya}, porque cada turno que se retrasa cuesta
valor. Y los \textbf{difusos} (como $*=\{0\mid 0\}$, la generalización
partidista de una posición Grundy $\neq 0$) no son ni fríos ni calientes:
solo importa quién llega primero. Green Hackenbush (\textit{Juegos sobre
	Árboles}, visto antes) es el caso particular donde \textbf{todo} es
impartial: ahí nunca aparecen números distintos de 0, solo posiciones
difusas o cero, y por eso basta Grundy. En cuanto el juego distingue
colores por jugador, se necesita esta teoría completa.

\textbf{Casos de uso:}

\begin{itemize}
	\item Hackenbush azul-rojo (partidista): cada arista pertenece a un
	jugador, y el valor resultante puede ser cualquier número (entero o
	fraccionario), no solo 0 o $*$.
	\item Finales de Go y Domineering, donde cada región del tablero se
	analiza como un juego independiente antes de sumarlas.
	\item Distinguir, dentro de una posición compuesta, qué componentes
	son ``ruido'' (números, se pueden ignorar o sumar aparte) y cuáles son
	la disputa real (difusos o calientes).
	\item Construir sistemáticamente enteros y fracciones diádicas como
	juegos, para usarlos como \textbf{unidad de comparación} contra
	posiciones desconocidas.
	\item Base directa de \textit{Suma de Juegos Partidistas y
		Comparación de Valores}, la siguiente subsección.
\end{itemize}

\textbf{Complejidad:} comparar dos juegos de $m_1$ y $m_2$ nodos en su
árbol de opciones cuesta, con memoización por pares de punteros,
$O(m_1 \cdot m_2)$ en el peor caso (cada par de subposiciones se compara
a lo sumo una vez, y cada comparación revisa un número de opciones
acotado por el grado del nodo). En la práctica esta técnica se usa sobre
juegos \textbf{pequeños} (unas pocas piezas o aristas): sirve para
razonar sobre subjuegos concretos de un problema, no para explorar un
espacio de estados grande como en Grundy sobre un DAG.

\begin{lstlisting}[style=cpp]
	struct Juego {
		vector<Juego*> L, R;    // opciones de Izquierda y de Derecha
	};
	
	// CREATE: construye un juego nuevo a partir de sus opciones.
	Juego* nuevo(vector<Juego*> l, vector<Juego*> r) {
		Juego* j = new Juego();
		j->L = l; j->R = r;
		return j;
	}                                          // O(1)
	
	map<pair<Juego*,Juego*>, bool> memoLeq;
	
	// READ: A <= B, usando la definicion recursiva de la teoria de
	// juegos: ningun G^L de A es >= B, y ningun H^R de B es <= A.
	bool leq(Juego* A, Juego* B) {
		auto key = make_pair(A, B);
		auto it = memoLeq.find(key);
		if (it != memoLeq.end()) return it->second;
		bool resultado = true;
		for (Juego* Al : A->L) if (leq(B, Al)) { resultado = false; break; } // Al >= B
		if (resultado)
		for (Juego* Br : B->R) if (leq(Br, A)) { resultado = false; break; } // Br <= A
		memoLeq[key] = resultado;
		return resultado;
	}                                          // O(m1 * m2) amortizado con memo
	
	// READ: igualdad de VALOR (dos arboles distintos pueden representar
	// el mismo juego): A == B si A <= B y B <= A.
	bool eq(Juego* A, Juego* B) { return leq(A, B) && leq(B, A); }
	
	enum Signo { POSITIVO, NEGATIVO, CERO, DIFUSO };
	Juego* CERO_G;   // el juego vacio {|}, se inicializa una vez
	
	// READ: clasifica G comparandolo contra 0: quien gana la partida si
	// G fuera el UNICO componente en juego.
	Signo clasificar(Juego* G) {
		bool noNegativo = leq(CERO_G, G);   // G >= 0
		bool noPositivo = leq(G, CERO_G);   // G <= 0
		if (noNegativo && noPositivo) return CERO;      // segundo jugador gana
		if (noNegativo) return POSITIVO;                // Izquierda gana siempre
		if (noPositivo) return NEGATIVO;                // Derecha gana siempre
		return DIFUSO;                                  // primer jugador gana
	}                                          // O(m) via leq
	
	// ------------------------------------------------------------
	// MODIFICACIONES Y COMENTARIOS CON LOS USOS
	// ------------------------------------------------------------
	
	// 1. CONSTRUIR ENTEROS COMO JUEGOS: la definicion recursiva estandar,
	//    donde cada entero es "un peldano mas simple que el anterior":
	//
	//    Juego* entero(int n, Juego* cero) {
		//        if (n == 0) return cero;
		//        if (n > 0)  return nuevo({entero(n-1, cero)}, {});   // {n-1 | }
		//        else        return nuevo({}, {entero(n+1, cero)});   // {  | n+1}
		//    }
	//
	//    Un entero positivo no tiene opciones para Derecha (nadie puede
	//    bajarlo), y viceversa para los negativos - por eso siempre gana
	//    quien tiene el signo a favor, sin importar quien mueve primero.
	
	// 2. CONSTRUIR FRACCIONES DIADICAS COMO JUEGOS: 1/2 = {0 | 1},
	//    1/4 = {0 | 1/2}, 3/4 = {1/2 | 1}, y en general una fraccion
	//    p/2^k (p impar) se construye como {(p-1)/2^k | (p+1)/2^k}, donde
	//    esos vecinos se simplifican recursivamente (al ser p-1 y p+1
	//    pares, su fraccion reduce a denominador 2^(k-1)). Estas
	//    construcciones se pueden verificar entre si con leq/eq: por
	//    ejemplo, leq(cero, mitad) && !eq(cero, mitad) confirma 0 < 1/2.
	
	// 3. ARRIBA (^) Y ABAJO (v), LOS INFINITESIMALES MAS SIMPLES:
	//    arriba = {0 | *}, abajo = {* | 0}. Son POSITIVO y NEGATIVO
	//    respectivamente, pero mas pequenos en valor absoluto que
	//    CUALQUIER numero positivo/negativo distinto de 0 - representan
	//    una ventaja real pero infinitesimal, tipica del final de un
	//    juego donde ya casi no queda nada por disputar.
	
	// 4. HECHO CONTRAINTUITIVO QUE VALE LA PENA VERIFICAR: arriba + * NO
	//    es simplemente "un poco mas que cero" como se podria esperar por
	//    intuicion ingenua; es DIFUSO (confundido con 0) - se puede
	//    verificar con clasificar(suma(arriba, estrella)) (suma definida
	//    en la siguiente subseccion). Este resultado clasico de la teoria
	//    (Winning Ways, Berlekamp-Conway-Guy) es la razon por la que un
	//    valor "aproximadamente positivo" NO garantiza ganar: solo la
	//    clasificacion EXACTA contra 0 lo hace.
	
	// 5. TRAMPA: UN JUEGO DIFUSO NO ES "0". Que clasificar(G) retorne
	//    DIFUSO significa que G NO es comparable con 0 (ni G<=0 ni
	//    G>=0), no que sea un empate: en un juego difuso, quien mueve
	//    PRIMERO en ESE componente gana la disputa local - justo lo
	//    opuesto de G=0 (donde el que mueve primero pierde ese
	//    componente). Confundir estos dos casos invierte la estrategia.
\end{lstlisting}

\textbf{Ejemplo de funcionamiento:}

\begin{lstlisting}[style=cpp]
	Juego* cero = nuevo({}, {});
	CERO_G = cero;
	Juego* estrella = nuevo({cero}, {cero});      // * = {0|0}
	Juego* uno      = nuevo({cero}, {});          // 1 = {0|}
	Juego* menosUno = nuevo({}, {cero});          // -1 = {|0}
	Juego* arriba   = nuevo({cero}, {estrella});  // ^ = {0|*}
	Juego* abajo    = nuevo({estrella}, {cero});  // v = {*|0}
	
	cout << clasificar(cero);       // CERO
	cout << clasificar(estrella);   // DIFUSO   (nadie tiene ventaja fija: gana quien mueve primero)
	cout << clasificar(uno);        // POSITIVO
	cout << clasificar(arriba);     // POSITIVO (pero infinitesimal)
	
	cout << (leq(cero, uno) && !eq(cero, uno));   // 1   (0 < 1)
\end{lstlisting}

\textbf{Regla práctica:} antes de aplicar Sprague-Grundy a un juego,
verifica si Izquierda y Derecha tienen realmente las mismas opciones en
cada posición; si no (colores, piezas o roles distintos por jugador), es
partidista y necesitas esta teoría. Clasifica cada componente contra 0
con \texttt{leq}/\texttt{eq}: POSITIVO o NEGATIVO son ventajas fijas
(``fría'', se puede sumar como un número), CERO significa que a nadie le
conviene jugar ahí, y DIFUSO significa que es exactamente ahí donde se
libra la disputa real por el turno.

\subsection{Suma de Juegos Partidistas y Comparación de Valores}

\textbf{Idea:} una posición real casi nunca es un solo juego pequeño,
sino varios componentes independientes jugados a la vez (varias regiones
de Hackenbush, varias cadenas de Domineering); en cada turno, un jugador
elige \textbf{un} componente y hace ahí un movimiento válido. Esto es la
\textbf{suma disyunta} $G+H$, definida recursivamente igual que los
juegos mismos: las opciones de Izquierda en $G+H$ son mover en $G$ o
mover en $H$ (nunca ambas a la vez), y análogamente para Derecha:
$$(G+H)^L = \{G^L + H\} \cup \{G + H^L\}, \qquad (G+H)^R = \{G^R + H\} \cup \{G + H^R\}.$$
La \textbf{negación} $-G$ intercambia los roles: cambia cada opción de
Izquierda por Derecha y viceversa, recursivamente. Con suma y negación,
determinar el ganador de una posición completa $G_1+G_2+\cdots+G_k$ es
exactamente clasificar (con \texttt{leq}/\texttt{eq} de la subsección
anterior) el juego resultante de sumarlos todos, contra 0.

\textbf{¿Por qué usarlo?} \textit{Suma de Juegos Independientes (XOR de
	Grundy Numbers)} en \textit{Fundamentos} es el caso \textbf{particular}
de esta misma idea cuando \textbf{todos} los componentes son impartiales:
ahí, sumar juegos y comparar contra 0 colapsa exactamente a
``XOR de los números de Grundy $\neq 0$'', un atajo válido solo porque en
juegos impartiales toda posición es 0 o difusa (nunca un número distinto
de 0, nunca fría). En cuanto un solo componente es partidista y resulta
ser un número (por ejemplo, una porción ya decidida de Hackenbush con
puntaje fijo), el atajo de XOR deja de tener sentido —los números no se
combinan por XOR, sino que si \textbf{todos} los componentes son números,
simplemente se \textbf{suman como valores reales} (esto es válido: la
suma de juegos-número coincide con la suma aritmética usual). El caso
verdaderamente difícil es una mezcla de números, calientes y difusos:
ahí no hay atajo aritmético y hace falta construir la suma completa como
juego y clasificarla.

\textbf{Casos de uso:}

\begin{itemize}
	\item Determinar el ganador de una posición de Hackenbush azul-rojo
	con varias componentes de distinto valor.
	\item Sumar puntuaciones parciales ya cerradas (números) de un final
	de Go directamente como valores reales.
	\item Verificar relaciones algebraicas entre valores de juego (por
	ejemplo, que $G + (-G) = 0$ para cualquier $G$, o que dos posiciones
	con árboles distintos representan el mismo valor).
	\item Simplificar una posición compuesta reemplazando cada
	componente por su valor más simple antes de sumarlos, en vez de
	analizar el árbol de juego completo de la posición entera.
	\item Confirmar hechos no intuitivos de la teoría (como que $\uparrow + *$
	es difuso) antes de usarlos como regla de decisión.
\end{itemize}

\textbf{Complejidad:} construir $G+H$ genera un árbol cuyo tamaño puede
crecer como el \textbf{producto} de los tamaños de $G$ y $H$ (cada nodo
de la suma corresponde a un par de subposiciones), así que sumar $k$
componentes de tamaño $O(m)$ cada uno cuesta, en el peor caso,
$O(m^k)$ nodos antes de memoizar; con memoización por pares de punteros
(reutilizando subárboles repetidos) el costo real depende de cuántos
pares \textbf{distintos} aparecen, típicamente muy por debajo de la cota
ingenua para los tamaños de juguete con los que se trabaja esta teoría a
mano. Clasificar la suma final cuesta lo mismo que cualquier
\texttt{leq} sobre el árbol resultante.

\begin{lstlisting}[style=cpp]
	map<Juego*, Juego*> memoNeg;
	
	// READ: -A, intercambiando Izquierda y Derecha recursivamente.
	Juego* negar(Juego* A) {
		auto it = memoNeg.find(A);
		if (it != memoNeg.end()) return it->second;
		vector<Juego*> nl, nr;
		for (Juego* Ar : A->R) nl.push_back(negar(Ar));   // lo que era opcion de Derecha
		for (Juego* Al : A->L) nr.push_back(negar(Al));   // pasa a ser de Izquierda, y viceversa
		Juego* res = nuevo(nl, nr);
		memoNeg[A] = res;
		return res;
	}                                          // O(m) con memo
	
	map<pair<Juego*,Juego*>, Juego*> memoSuma;
	
	// READ: A + B, la suma disyunta: mover es elegir UN componente y
	// jugar ahi un movimiento valido de ese componente.
	Juego* suma(Juego* A, Juego* B) {
		auto key = make_pair(A, B);
		auto it = memoSuma.find(key);
		if (it != memoSuma.end()) return it->second;
		vector<Juego*> nl, nr;
		for (Juego* Al : A->L) nl.push_back(suma(Al, B));   // mover en A
		for (Juego* Bl : B->L) nl.push_back(suma(A, Bl));   // mover en B
		for (Juego* Ar : A->R) nr.push_back(suma(Ar, B));
		for (Juego* Br : B->R) nr.push_back(suma(A, Br));
		Juego* res = nuevo(nl, nr);
		memoSuma[key] = res;
		return res;
	}                                          // O(producto de tamanos) con memo
	
	// READ: ganador de la suma de varios componentes, reutilizando
	// clasificar() de la subseccion anterior sobre el total acumulado.
	Signo ganadorDeLaSuma(vector<Juego*> componentes) {
		Juego* total = CERO_G;
		for (Juego* g : componentes) total = suma(total, g);
		return clasificar(total);
	}                                          // O(suma de tamanos al cuadrado), con memo
	
	// ------------------------------------------------------------
	// MODIFICACIONES Y COMENTARIOS CON LOS USOS
	// ------------------------------------------------------------
	
	// 1. ATAJO CUANDO TODO ES NUMERO: si CADA componente es un juego que
	//    resulta ser un numero (entero o fraccion), no hace falta
	//    construir la suma simbolica completa: basta calcular el valor
	//    numerico de cada uno por separado (por ejemplo, con el
	//    algoritmo clasico para reducir un Hackenbush azul-rojo en forma
	//    de CADENA a una fraccion diadica) y sumarlos como valores reales
	//    ordinarios; el signo de esa suma numerica decide el ganador
	//    exactamente igual que clasificar() sobre el juego simbolico.
	//
	// 2. ATAJO CUANDO TODO ES IMPARTIAL: si NINGUN componente distingue
	//    Izquierda de Derecha (todo Green Hackenbush o Nim puro), no
	//    construyas la suma como {L|R}: usa directamente XOR de Grundy
	//    numbers (Suma de Juegos Independientes, en Fundamentos), que es
	//    equivalente pero mucho mas barato de calcular.
	//
	// 3. VERIFICAR G + (-G) = 0 PARA CUALQUIER G: es un teorema general
	//    de la teoria (todo juego tiene un inverso aditivo), util como
	//    prueba de que negar() y suma() estan bien implementados:
	//    eq(suma(G, negar(G)), CERO_G) debe dar siempre true.
	//
	// 4. ORDEN DE JUEGO PARA COMPONENTES MIXTOS: si la posicion mezcla
	//    numeros (frios, nadie quiere jugar ahi) con componentes difusos
	//    o calientes, la intuicion de la teoria (no una garantia absoluta
	//    para todo caso) es que conviene jugar primero en los
	//    componentes MAS urgentes (mas lejos de ser un numero) y dejar
	//    los numeros para el final, ya que estos no se "enfrian" mas -
	//    perder un turno ahi no cambia su valor, mientras que retrasar un
	//    componente caliente si puede costar ventaja.
	//
	// 5. TRAMPA: SUMAR VALORES NUMERICOS DE COMPONENTES QUE NO SON
	//    NUMEROS. Si un solo componente es difuso o caliente (no es un
	//    numero), no existe un "valor real" unico que sumar - solo
	//    existe el juego completo {L|R}. Intentar asignarle un numero de
	//    todas formas (por ejemplo, promediando opciones) y sumarlo
	//    directamente con los demas produce conclusiones incorrectas
	//    sobre quien gana; en ese caso hay que construir la suma
	//    simbolica y clasificarla con leq/eq, como en el ejemplo de
	//    arriba + estrella de la subseccion anterior.
\end{lstlisting}

\textbf{Ejemplo de funcionamiento:}

\begin{lstlisting}[style=cpp]
	// (continuando los juegos definidos en la subseccion anterior)
	
	cout << eq(suma(estrella, estrella), cero);        // 1   (* + * = 0)
	cout << eq(suma(uno, menosUno), cero);              // 1   (1 + (-1) = 0)
	cout << eq(negar(arriba), abajo);                   // 1   (-^ = v)
	
	// El hecho clasico y contraintuitivo: arriba + * es DIFUSO, no
	// "un poco mayor que 0" como sugeriria la intuicion ingenua.
	cout << clasificar(suma(arriba, estrella));         // DIFUSO
	
	// arriba + arriba + abajo == arriba (el abajo "cancela" uno de los
	// dos arriba, dejando exactamente el otro).
	Juego* combinado = suma(suma(arriba, arriba), abajo);
	cout << eq(combinado, arriba);                      // 1
	
	// Ganador de una posicion con tres componentes mixtos.
	Signo g = ganadorDeLaSuma({uno, estrella, abajo});
	cout << g;   // depende de los valores exactos; se calcula clasificando el total
\end{lstlisting}

\textbf{Nota:} el resultado \texttt{arriba + estrella = DIFUSO} contradice
la intuición de sumar valores aproximados ($\uparrow$ es ``un poquito
positivo'', $*$ es ``ni positivo ni negativo'', se esperaría que la suma
siguiera siendo positiva): es la advertencia central de esta teoría —
\textbf{no se puede} razonar sobre sumas de juegos partidistas
promediando o aproximando valores; solo la construcción simbólica exacta
(\texttt{suma} + \texttt{clasificar}) da la respuesta correcta.

\textbf{Regla práctica:} construye la suma disyunta y clasifícala contra
0 siempre que la posición mezcle componentes de distinta naturaleza
(números, calientes, difusos); si todos los componentes son números,
súmalos como valores reales ordinarios; si todos son impartiales, usa
XOR de Grundy en vez de esta maquinaria. Nunca mezcles los tres atajos:
en cuanto hay al menos un componente partidista no numérico junto a
otros, la única vía correcta es la suma simbólica completa.
\subsection{Juegos de Turnos con Puntaje (Optimal Strategy for a Game)}

\textbf{Idea:} esta es la misma técnica de \textit{Juegos de Tomar de los
	Extremos de un Arreglo (DP de Intervalos)}, vista en \textit{Juegos sobre
	Estructuras de Datos} — dos jugadores turnándose para tomar el extremo
izquierdo o derecho de un arreglo, maximizando cada uno su propio
puntaje, resuelto con $dp[i][j] = \max(a[i]-dp[i{+}1][j],\, a[j]-dp[i][j{-}1])$.
Reaparece aquí, junto a \textit{Números Surreales} y \textit{Suma de
	Juegos Partidistas}, para ubicarla correctamente dentro del panorama
teórico: es un ejemplo de \textbf{juego con puntaje} (\textit{scoring
	game}), una rama de la teoría de juegos combinatorios distinta de la
teoría \textit{normal-play} de Conway que se usó en las dos subsecciones
anteriores. En \textit{normal-play}, el valor de un juego captura
\textbf{quién puede seguir moviendo}; en un juego con puntaje no existe
esa pregunta binaria (el arreglo se agota jugada tras jugada sin importar
las decisiones), y lo que se necesita en cambio es un \textbf{número}
—el margen— que resuma cuánto le conviene a cada jugador esa parte del
juego.

\textbf{¿Por qué usarlo?} Vale la pena notar por qué este juego, siendo
técnicamente asimétrico en el resultado (cada jugador se queda con algo
distinto), \textbf{no} necesita la maquinaria completa de
$\{G^L \mid G^R\}$: en los juegos partidistas de la subsección de
Números Surreales, lo que distingue a Izquierda de Derecha son las
\textbf{jugadas disponibles} (Hackenbush azul vs. rojo, opciones
distintas por jugador). Aquí, en cambio, \textbf{ambos jugadores tienen
	exactamente las mismas jugadas} (tomar izquierda o tomar derecha) — lo
único que difiere es \textbf{a quién beneficia} cada resultado. Esa
simetría en las jugadas (aunque no en el puntaje) es lo que permite
colapsar todo a un solo número por intervalo con una recurrencia
minimax simple, en vez de construir un árbol de opciones $\{L\mid R\}$
distinto para cada jugador.

\textbf{Casos de uso:} (los mismos que en la subsección original —
Optimal Strategy for a Game / Predict the Winner, cartas en fila, juegos
de puntaje con dos extremos disponibles) más uno propio de este contexto
teórico:

\begin{itemize}
	\item Como ejemplo concreto y resoluble de un \textit{scoring game}
	para contrastar con los juegos partidistas normal-play de las dos
	subsecciones anteriores.
	\item Servir de sanity check: si un juego con puntaje tiene jugadas
	simétricas entre jugadores (como aquí), sospecha que una DP de
	intervalos simple bastará, sin necesitar surreales.
	\item Todos los casos ya cubiertos en \textit{Juegos de Tomar de los
		Extremos de un Arreglo}: determinar ganador, margen exacto, puntaje
	de cada jugador.
\end{itemize}

\textbf{Complejidad:} idéntica a la subsección original:
$O(n^2)$ en tiempo con la tabla completa, reducible a $O(n)$ de memoria
procesando por longitud de intervalo creciente (ver esa subsección para
el detalle y la versión con arreglo 1D).

\begin{lstlisting}[style=cpp]
	struct JuegoDeExtremos {
		int n;
		vector<long long> a;
		vector<vector<long long>> dp;   // dp[i][j] = diferencia optima
		// (turno - rival) jugando sobre a[i..j]
		
		// CREATE: guarda los valores del arreglo.
		JuegoDeExtremos(vector<long long> valores) {
			a = valores;
			n = (int)a.size();
			dp.assign(n, vector<long long>(n, 0));
		}                                          // O(n^2)
		
		// WRITE: llena dp por longitud de intervalo creciente.
		void calcular() {
			for (int i = 0; i < n; i++) dp[i][i] = a[i];
			for (int largo = 2; largo <= n; largo++) {
				for (int i = 0; i + largo - 1 < n; i++) {
					int j = i + largo - 1;
					dp[i][j] = max(a[i] - dp[i+1][j], a[j] - dp[i][j-1]);
				}
			}
		}                                          // O(n^2)
		
		// READ: diferencia optima (turno - rival) sobre TODO el arreglo.
		long long diferencia() const { return dp[0][n-1]; }
		
		// READ: puntaje exacto de cada jugador bajo juego optimo.
		pair<long long,long long> puntajes() const {
			long long total = 0; for (long long x : a) total += x;
			long long diff = dp[0][n-1];
			return {(total + diff) / 2, (total - diff) / 2};
		}                                          // O(n)
	};
	
	// ------------------------------------------------------------
	// MODIFICACIONES Y COMENTARIOS CON LOS USOS
	// ------------------------------------------------------------
	
	// 1. Reduccion a O(n) de memoria, truco de paridad para n par, costo
	//    o bonus por extremo, y variante de "k turnos" con dos punteros:
	//    ver las modificaciones 1 a 4 de Juegos de Tomar de los Extremos
	//    de un Arreglo - son exactamente las mismas aqui, sin cambios.
	//
	// 2. Por que NO hace falta {L | R}: si quisieras modelar este juego
	//    como un juego partidista general, cada estado (i, j) tendria
	//    UNA opcion de Izquierda (tomar a[i]) y UNA opcion de Derecha
	//    (tomar a[j]), pero devolviendo el MISMO subjuego resultante
	//    (i+1,j) o (i,j-1) sin importar quien lo tome - lo unico que
	//    distingue el resultado es el puntaje sumado, no el subjuego al
	//    que se llega. Por eso alcanza con propagar un numero (la
	//    diferencia) en vez de dos arboles de opciones distintos.
	
	// 3. TRAMPA (la misma que antes, repetida por importancia): esto NO
	//    es un juego normal-play de "quien mueve al final gana" - no
	//    intentes aplicarle Sprague-Grundy ni la clasificacion
	//    POSITIVO/NEGATIVO/CERO/DIFUSO de Números Surreales, que asume
	//    una convencion de victoria distinta (basada en quien se queda
	//    sin movimientos, no en maximizar un puntaje acumulado).
\end{lstlisting}

\textbf{Ejemplo de funcionamiento:} idéntico al de \textit{Juegos de
	Tomar de los Extremos de un Arreglo (DP de Intervalos)} — mismo código,
mismos resultados (\texttt{diferencia() = 0} para \texttt{\{1,5,2,4,6\}},
empate bajo juego óptimo).

\textbf{Regla práctica:} si ya resolviste este patrón en \textit{Juegos
	sobre Estructuras de Datos}, no hay nada nuevo que programar aquí — el
valor de volver a verlo es puramente conceptual: sirve para ubicar los
\textbf{juegos con puntaje} como una familia aparte de los juegos
partidistas normal-play, y para reconocer la señal de que un juego con
puntaje admite esta DP simple en cuanto notas que \textbf{ambos
	jugadores comparten las mismas jugadas} y solo compiten por quién se
queda con cada valor.

\section{Estrategia Óptima y Técnicas de Resolución}

\subsection{Minimax con Memoización}

\textbf{Idea:} minimax es la técnica genérica para resolver cualquier
juego de dos jugadores de información perfecta: en cada estado, el
jugador que \textbf{maximiza} elige el movimiento con el mejor valor
para él, y el que \textbf{minimiza} elige el que más lo perjudica; el
valor de un estado terminal es el resultado del juego ahí (victoria,
derrota, empate, o un puntaje). Esto ya es exactamente lo que hacen
\textit{Juegos de Tomar de los Extremos} y \textit{Poda con Fuerza
	Bruta (Tabla de Grundy)} de \textit{Fundamentos} por dentro — minimax
\textbf{es} el mecanismo detrás de esas técnicas, presentado aquí de
forma explícita y \textbf{general}: en vez de un estado con una forma
particular conveniente (un intervalo, una máscara), un estado puede ser
cualquier cosa (un tablero de Tres en Raya, una posición de Conecta 4)
mientras exista una función que genere sus movimientos válidos. El
problema práctico es que el mismo estado puede alcanzarse por caminos
distintos (\textbf{transposiciones}: jugar $A$ y luego $B$ llega al
mismo tablero que jugar $B$ y luego $A$), y sin memoización minimax lo
vuelve a resolver desde cero cada vez que lo visita.

\textbf{¿Por qué usarlo?} Cuando el estado del juego no tiene una
estructura simple que permita indexarlo por un par de enteros (como en
la DP de intervalos) ni reducirlo a un único número de Grundy (como en
\textit{Fundamentos}), minimax con memoización es la herramienta por
defecto: sirve para \textbf{cualquier} juego adversarial de dos
jugadores, a costa de necesitar una clave de estado que se pueda usar en
un \texttt{map} o \texttt{unordered\_map} (un bitmask, un string, una
tupla). La memoización es lo que hace viable esta generalidad: sin ella,
un juego con muchas transposiciones (como Tres en Raya) explora
\textbf{cientos de miles} de nodos repitiendo trabajo; con ella, solo
se resuelve una vez cada estado \textit{distinto}.

\textbf{Casos de uso:}

\begin{itemize}
	\item Resolver juegos pequeños de tablero (Tres en Raya, variantes
	de Conecta 4 en tableros chicos) determinando el resultado bajo juego
	óptimo desde cualquier posición.
	\item Cualquier juego donde el estado natural (tablero, mano de
	cartas, configuración) no encaja en las técnicas más específicas
	(intervalos, máscaras de bits, Grundy) de secciones anteriores.
	\item Verificar por fuerza bruta el resultado de una técnica más
	rápida (Grundy, DP de intervalos) sobre casos pequeños, como
	\textit{sanity check} antes de confiar en la fórmula cerrada.
	\item Juegos con estados que se repiten mucho por transposición,
	donde la memoización reduce el trabajo de forma dramática incluso sin
	ninguna otra optimización.
	\item Como base directa sobre la que se construye \textit{Poda
		Alfa-Beta}, que acelera esta misma búsqueda sin cambiar el resultado.
\end{itemize}

\textbf{Complejidad:} sin memoización, el costo es
$O(b^d)$ con $b$ el factor de ramificación (movimientos disponibles) y
$d$ la profundidad del árbol — para Tres en Raya desde el tablero vacío,
esto son más de medio millón de nodos visitados solo para un juego de
$3\times3$. Con memoización, el costo baja a
$O(E \cdot b)$, con $E$ el número de estados \textbf{distintos}
alcanzables (no de caminos para llegar a ellos): cada estado se resuelve
una sola vez, revisando sus $b$ movimientos, y las consultas repetidas
cuestan $O(1)$ amortizado.

\begin{lstlisting}[style=cpp]
	// Tres en Raya: el estado se codifica como dos mascaras de 9 bits
	// (celdas ocupadas por X y por O). LINEAS lista las 8 formas de ganar.
	int LINEAS[8] = {
		0b111000000, 0b000111000, 0b000000111,   // filas
		0b100100100, 0b010010010, 0b001001001,   // columnas
		0b100010001, 0b001010100                 // diagonales
	};
	
	// READ: true si la mascara m contiene alguna linea ganadora completa.
	bool gano(int m) {
		for (int L : LINEAS) if ((m & L) == L) return true;
		return false;
	}                                          // O(1) (8 comparaciones)
	
	unordered_map<long long, int> memo;
	
	// READ: valor del estado (x = celdas de X, o = celdas de O, turno =
	// true si mueve X) bajo juego optimo: 1 si gana X, -1 si gana O, 0
	// empate. X maximiza, O minimiza.
	int minimaxMemo(int x, int o, bool turno) {
		if (gano(x)) return 1;
		if (gano(o)) return -1;
		int ocupado = x | o;
		if (ocupado == 0b111111111) return 0;   // tablero lleno: empate
		
		// clave unica: 9 bits de x, 9 bits de o, 1 bit de turno
		long long clave = ((long long)x << 10) | o | ((long long)turno << 20);
		auto it = memo.find(clave);
		if (it != memo.end()) return it->second;
		
		int mejor = turno ? -2 : 2;   // -2/2 fuera del rango real [-1,1]
		for (int i = 0; i < 9; i++) {
			if (!((ocupado >> i) & 1)) {
				int val = turno ? minimaxMemo(x | (1 << i), o, false)
				: minimaxMemo(x, o | (1 << i), true);
				mejor = turno ? max(mejor, val) : min(mejor, val);
			}
		}
		memo[clave] = mejor;
		return mejor;
	}                                          // O(E * 9) total, con E = estados distintos
	
	// ------------------------------------------------------------
	// MODIFICACIONES Y COMENTARIOS CON LOS USOS
	// ------------------------------------------------------------
	
	// 1. RECONSTRUIR LA JUGADA OPTIMA, NO SOLO EL VALOR: para saber QUE
	//    movimiento jugar (no solo si se gana), guarda junto al valor en
	//    memo cual fue el indice i que produjo mejor, y consulta esa
	//    tabla desde el estado actual en vez de recorrer todo de nuevo.
	
	// 2. CLAVE DE ESTADO PARA OTROS JUEGOS: la clave debe capturar TODO
	//    lo que determina el futuro del juego (tablero + de quien es el
	//    turno, y a veces informacion extra como "movimientos
	//    restantes" o "cartas ya jugadas"). Para tableros mas grandes
	//    que quepan en bits (hasta 64), usar mascaras de bits como aqui;
	//    para estados mas complejos, usar un string o una tupla como
	//    clave de un map.
	
	// 3. LIMITE DE PROFUNDIDAD (juegos que no terminan rapido): si el
	//    juego real es demasiado grande para resolverlo por completo
	//    (ajedrez, Go), se trunca la busqueda a una profundidad maxima y
	//    se usa una funcion de EVALUACION heuristica en vez del valor
	//    exacto en los nodos truncados; esto ya no da el valor optimo
	//    real, sino una aproximacion, y es la base de los motores de
	//    juegos practicos.
	
	// 4. TRAMPA: LA CLAVE DEBE INCLUIR EL TURNO. El mismo tablero (x, o)
	//    tiene un valor distinto segun quien mueva a continuacion (por
	//    ejemplo, un tablero donde X ya gano es +1 sin importar el turno,
	//    pero un tablero intermedio con el mismo par de mascaras y turno
	//    distinto puede tener resultados diferentes). Omitir el turno de
	//    la clave mezcla estados que no deberian compartir memo.
\end{lstlisting}

\textbf{Ejemplo de funcionamiento:}

\begin{lstlisting}[style=cpp]
	int valor = minimaxMemo(0, 0, true);   // tablero vacio, mueve X
	cout << valor;
	// 0   (empate bajo juego optimo -- el resultado conocido de Tres en
	// Raya cuando ambos jugadores juegan perfecto)
	
	// Sin memoizacion, resolver el mismo tablero vacio visita 549 946
	// nodos; con memoizacion, solo 16 168, sobre 4520 estados distintos
	// memoizados -- la mayoria de esos 549 946 nodos eran el MISMO
	// tablero visitado una y otra vez por caminos distintos.
\end{lstlisting}

\textbf{Nota:} el resultado (empate) es el hecho bien conocido de que
Tres en Raya, jugado perfectamente por ambos lados, siempre termina en
empate; sirve como verificación rápida de que la implementación es
correcta antes de usarla en un problema donde no se conoce de antemano
el resultado esperado.

\textbf{Regla práctica:} usa minimax con memoización en cuanto el estado
del juego no encaje en una técnica más específica de las secciones
anteriores (intervalos, máscaras, Grundy), y en cuanto sospeches que hay
transposiciones — estados alcanzables por más de un camino. Si el
espacio de estados es demasiado grande para memoizar por completo
(juegos reales de tablero grande), esto deja de ser exacto y hace falta
una función de evaluación heurística con profundidad limitada, un tema
que ya escapa del alcance de la programación competitiva clásica.

\subsection{Poda Alfa-Beta}

\textbf{Idea:} alfa-beta acelera minimax \textbf{sin cambiar el
	resultado}, descartando ramas que ya se sabe que no pueden influir en la
decisión final. Mantiene dos cotas mientras recorre el árbol: $\alpha$
(el mejor valor que el jugador que \textbf{maximiza} ya tiene garantizado
en algún lugar del árbol) y $\beta$ (lo mismo para el que
\textbf{minimiza}). Si en algún punto $\alpha \geq \beta$, significa que
el jugador que está un nivel arriba en el árbol \textbf{ya tiene una
	opción mejor} en otra rama, así que el rival —siendo racional— nunca
dejaría que la partida llegara a esta rama: se puede dejar de explorarla
por completo (\textbf{podar}) sin perder nada, porque su valor exacto ya
no puede cambiar la decisión que se tomará más arriba.

\textbf{¿Por qué usarlo?} Es una mejora directa sobre
\textit{Minimax con Memoización}: mismo árbol, mismo resultado exacto,
menos nodos visitados. En el mejor caso (con un buen \textbf{orden de
	movimientos} — revisar primero las jugadas más prometedoras), alfa-beta
reduce la complejidad de $O(b^d)$ a
$O(b^{d/2})$, equivalente a explorar un árbol con la \textbf{raíz
	cuadrada} del factor de ramificación. Es independiente de la
memoización y se puede combinar con ella: la memoización evita rehacer
trabajo en \textbf{transposiciones} (mismo estado por caminos distintos),
mientras que alfa-beta evita explorar ramas que \textbf{nunca se van a
	necesitar}, sean transposiciones o no; combinar ambas en una
\textbf{tabla de transposición} exige un cuidado extra (modificación 3)
porque, bajo poda, un valor guardado puede ser solo una \textbf{cota}
del valor real, no el valor exacto.

\textbf{Casos de uso:}

\begin{itemize}
	\item Cualquier motor de juego de dos jugadores (ajedrez, damas,
	Conecta 4) donde el árbol completo es demasiado grande para
	explorarlo sin poda, incluso con memoización.
	\item Acelerar la resolución exacta de juegos pequeños que ya se
	resuelven con \textit{Minimax con Memoización}, cuando el número de
	estados distintos sigue siendo alto.
	\item Motores con profundidad limitada y función de evaluación
	heurística (mencionados en la subsección anterior): alfa-beta es
	casi siempre imprescindible ahí, porque permite explorar mucha más
	profundidad en el mismo tiempo.
	\item Cualquier búsqueda adversarial donde ordenar los movimientos
	de mejor a peor (aunque sea con una heurística simple) antes de
	explorarlos acelera la poda de forma notable.
\end{itemize}

\textbf{Complejidad:} en el peor caso (sin ningún orden útil de
movimientos), alfa-beta visita el mismo $O(b^d)$ que minimax sin poda —
no hay garantía si el orden de exploración es malo. Con un
\textbf{orden perfecto} de movimientos, la cota baja a $O(b^{d/2})$. En
la práctica, incluso un orden \textbf{razonable} (por ejemplo, probar
primero movimientos centrales o históricamente buenos) logra una
reducción sustancial sin llegar al caso ideal.

\begin{lstlisting}[style=cpp]
	// READ: valor del estado bajo juego optimo, podando ramas que no
	// pueden cambiar la decision. alfa = mejor valor ya garantizado para
	// quien maximiza; beta = mejor valor ya garantizado para quien
	// minimiza. Llamada inicial: alfaBeta(0, 0, true, -2, 2).
	int alfaBeta(int x, int o, bool turno, int alfa, int beta) {
		if (gano(x)) return 1;
		if (gano(o)) return -1;
		int ocupado = x | o;
		if (ocupado == 0b111111111) return 0;
		
		if (turno) {   // X maximiza
			int mejor = -2;
			for (int i = 0; i < 9; i++) {
				if (!((ocupado >> i) & 1)) {
					int val = alfaBeta(x | (1 << i), o, false, alfa, beta);
					mejor = max(mejor, val);
					alfa = max(alfa, mejor);
					if (alfa >= beta) break;   // O ya tiene algo mejor en otra rama: podar
				}
			}
			return mejor;
		} else {       // O minimiza
			int mejor = 2;
			for (int i = 0; i < 9; i++) {
				if (!((ocupado >> i) & 1)) {
					int val = alfaBeta(x, o | (1 << i), true, alfa, beta);
					mejor = min(mejor, val);
					beta = min(beta, mejor);
					if (alfa >= beta) break;
				}
			}
			return mejor;
		}
	}                                          // O(b^d) peor caso, O(b^(d/2)) con buen orden
	
	// ------------------------------------------------------------
	// MODIFICACIONES Y COMENTARIOS CON LOS USOS
	// ------------------------------------------------------------
	
	// 1. ORDEN DE MOVIMIENTOS: la poda depende por completo de revisar
	//    primero las jugadas mas prometedoras. En Tres en Raya, probar el
	//    centro antes que las esquinas y los lados reduce los nodos
	//    visitados de forma notable (en la practica, de decenas de miles
	//    a unos pocos miles) sin cambiar el resultado en absoluto:
	//
	//    int orden[9] = {4, 0, 2, 6, 8, 1, 3, 5, 7}; // centro, esquinas, lados
	//    for (int idx = 0; idx < 9; idx++) {
		//        int i = orden[idx];
		//        if (!((ocupado >> i) & 1)) { /* ... */ }
		//    }
	//
	//    En juegos mas complejos, una heuristica simple (capturar antes
	//    que mover, o el mejor movimiento encontrado en una busqueda
	//    anterior mas superficial) cumple el mismo papel.
	
	// 2. VENTANA NULA (NEGASCOUT / PVS): una vez que se tiene un valor
	//    candidato, se puede verificar si una rama alternativa es MEJOR
	//    o PEOR que ese candidato con una ventana de tamano 1
	//    (alfa, alfa+1) en vez de la ventana completa, que poda mucho mas
	//    agresivamente; solo si esa verificacion barata sugiere que SI es
	//    mejor, se vuelve a explorar con la ventana completa para obtener
	//    el valor exacto. Es una optimizacion avanzada sobre alfa-beta,
	//    util cuando el orden de movimientos ya es razonablemente bueno.
	
	// 3. COMBINAR CON TABLA DE TRANSPOSICION (CUIDADO CON LOS FLAGS): con
	//    poda, un valor guardado en memo puede ser EXACTO, una COTA
	//    INFERIOR (se podo antes de confirmar que no habia algo mejor) o
	//    una COTA SUPERIOR (simetrico para el minimizador). Guardar sin
	//    distinguir estos casos produce resultados incorrectos:
	//
	//    enum Flag { EXACTO, COTA_INF, COTA_SUP };
	//    struct Entrada { int valor; Flag flag; };
	//    unordered_map<long long, Entrada> tabla;
	//
	//    int alfaBetaConTabla(int x, int o, bool turno, int alfa, int beta) {
		//        // ... casos terminales iguales ...
		//        long long clave = /* igual que en minimaxMemo */;
		//        int alfaOrig = alfa, betaOrig = beta;
		//        auto it = tabla.find(clave);
		//        if (it != tabla.end()) {
			//            if (it->second.flag == EXACTO) return it->second.valor;
			//            else if (it->second.flag == COTA_INF) alfa = max(alfa, it->second.valor);
			//            else beta = min(beta, it->second.valor);
			//            if (alfa >= beta) return it->second.valor;
			//        }
		//        int mejor = /* ... recursion normal de alfa-beta, actualizando alfa/beta ... */;
		//        Flag f = (mejor <= alfaOrig) ? COTA_SUP
		//               : (mejor >= betaOrig) ? COTA_INF
		//               : EXACTO;
		//        tabla[clave] = {mejor, f};
		//        return mejor;
		//    }
	//
	//    Esta combinacion (verificada sobre Tres en Raya) visito solo
	//    4777 nodos frente a los 18297 de alfa-beta sin tabla y los
	//    16168 de minimax con memo sin poda - las dos tecnicas se
	//    refuerzan en vez de competir.
	
	// 4. TRAMPA: PODAR NO CAMBIA EL RESULTADO, SOLO EL COSTO. Si
	//    alfa-beta da un valor distinto a minimax puro sobre el MISMO
	//    arbol, hay un error de implementacion (probablemente en como se
	//    actualizan o se pasan alfa/beta entre llamadas) - los dos deben
	//    coincidir siempre, y es la forma mas facil de depurar una
	//    implementacion nueva: compararla contra minimax simple en un
	//    juego pequeno.
	
	// 5. TRAMPA: >= VS >. La condicion de poda es alfa >= beta (no
	//    alfa > beta); con solo ">" se pierden podas validas sin que el
	//    resultado sea incorrecto, solo mas lento - un error silencioso
	//    que baja el rendimiento sin romper la correccion, dificil de
	//    notar sin medir nodos visitados explicitamente.
\end{lstlisting}

\textbf{Ejemplo de funcionamiento:}

\begin{lstlisting}[style=cpp]
	int valor = alfaBeta(0, 0, true, -2, 2);
	cout << valor;   // 0   (mismo resultado que minimax con memo: empate)
	
	// Nodos visitados resolviendo el tablero vacio de Tres en Raya,
	// todos con el MISMO resultado (0, empate):
	//   minimax sin memo ni poda:            549 946 nodos
	//   minimax con memoizacion:               16 168 nodos
	//   alfa-beta (sin orden de movimientos):   18 297 nodos
	//   alfa-beta con orden (centro primero):    7 275 nodos
	//   alfa-beta + tabla de transposicion:      4 777 nodos
\end{lstlisting}

\textbf{Nota:} llama la atención que alfa-beta \textbf{sin} un buen
orden de movimientos (18{,}297 nodos) visitó incluso más nodos que
minimax con memoización simple (16{,}168): la poda por sí sola no
garantiza mejorar sobre la memoización si el orden de exploración es
malo y no hay tabla de transposición que atrape las repeticiones. La
combinación de las dos técnicas (última fila) es la que da la mejor
reducción, y en la práctica casi siempre conviene implementarlas juntas
en vez de elegir una sola.

\textbf{Regla práctica:} implementa siempre poda alfa-beta sobre
cualquier minimax que ya funcione — es un cambio pequeño (pasar dos
enteros más y agregar una condición de corte) que nunca empeora el
resultado y casi siempre reduce el trabajo. Si el juego tiene muchas
transposiciones, combínala con una tabla de transposición usando los
tres \texttt{flags} (\texttt{EXACTO}, \texttt{COTA\_INF},
\texttt{COTA\_SUP}) en vez de memoizar solo con alfa-beta o solo con
minimax puro; y si el árbol sigue siendo intratable, invierte tiempo en
mejorar el \textbf{orden} de exploración antes que en cualquier otra
optimización — es lo que más impacto tiene sobre cuánto se poda.
\subsection{Detección de Patrones con Fuerza Bruta (Tabla de Grundy)}

\textbf{Idea:} Un juego es \textit{imparcial} cuando las jugadas disponibles dependen solo del estado y no de quién mueve. Se usa la convención de juego normal: quien no puede jugar, pierde. El \textbf{número de Grundy} de un estado se define por recursión: un estado terminal vale $0$, y en cualquier otro
$$G(s)=\operatorname{MEX}\{\,G(t)\mid t \text{ es alcanzable desde } s\,\},$$
donde $\operatorname{MEX}$ es el menor entero no negativo que no pertenece al conjunto. Por ejemplo, $\operatorname{MEX}\{0,1,3\}=2$.

La intuición es que $G(s)=g$ hace que $s$ se comporte como un montón de Nim de tamaño $g$. Por definición de MEX, desde $s$ se puede llegar a estados de cada valor menor que $g$ y nunca a un estado de valor $g$. Por tanto, si $G(s)\neq 0$ existe una jugada a un estado de valor $0$ (el estado es ganador), y si $G(s)=0$ toda jugada va a un valor distinto de $0$, desde donde el rival puede volver a $0$ (el estado es perdedor). La segunda consecuencia es la suma de juegos: si el juego se compone de $k$ componentes independientes, en cada turno se juega en uno solo de ellos, y el teorema de Sprague--Grundy dice que
$$G(s_1,\dots,s_k)=G(s_1)\oplus\cdots\oplus G(s_k).$$
La posición total es perdedora si y solo si ese XOR vale $0$. Con tres componentes de valores $2,3,1$ el XOR es $0$ (perdedora), y con $2,3,4$ es $5$ (ganadora).

La \textit{fuerza bruta} consiste en calcular la tabla $G(0), G(1), \dots, G(N)$ probando todas las jugadas de cada estado y reutilizando lo ya calculado. Como el MEX de $k$ valores nunca supera $k$, basta un arreglo de marcas de tamaño $k+1$. En el juego de quitar $1$, $2$ o $3$ piedras se obtiene $G(0..7)=0,1,2,3,0,1,2,3$, es decir $G(n)=n \bmod 4$. El \textit{patrón} es que estas tablas suelen ser periódicas, y basta calcular unos pocos términos, detectar el periodo y extrapolar a $n$ del orden de $10^{18}$. En los juegos de restar con conjunto finito de jugadas $\{m_1,\dots,m_r\}$, con $D=\max m_i$, hay una detección rigurosa. Para $n \ge D$ el valor $G(n)$ depende solo de los $D$ valores anteriores. Si la ventana de $D$ valores que empieza en $j$ es igual a la que empieza en $i<j$, entonces la secuencia se repite desde $i$ con periodo $j-i$, por inducción.

\textbf{¿Por qué usarlo?} Si el juego es un único componente y solo se pregunta «¿gana el primero?», basta una DP booleana ($gana[n]$ es verdadero si alguna jugada lleva a un estado con $gana$ falso), que es más simple y no necesita MEX (modificación 7). Grundy se justifica cuando hay \textit{varios} componentes independientes, porque el valor booleano de cada uno no alcanza para combinarlos y el número de Grundy sí. Frente al análisis retrógrado de la subsección de trabajar hacia atrás, que sirve en grafos con ciclos, Grundy exige un juego que termine (sin ciclos) pero entrega un valor combinable por XOR. Frente a los invariantes, el XOR de Grundy es justamente un invariante que se conserva de forma «complementaria»: cualquier jugada lo cambia, y desde un valor distinto de cero se puede devolver a cero. La fuerza bruta para detectar el patrón es la herramienta cuando la regla es complicada pero calculable para $n$ pequeños y la consulta es para $n$ enorme. Su límite es que fuera de los juegos de restar (ventana finita) el periodo detectado es una conjetura, no una prueba (modificación 3).

\textbf{Casos de uso:}
\begin{itemize}
	\item Juegos de restar sobre una o varias pilas, con conjunto de jugadas dado, y decidir quién gana.
	\item Sumas de juegos independientes (varias pilas, varios tableros): calcular el Grundy de cada uno y hacer XOR.
	\item Obtener la jugada ganadora, no solo el ganador.
	\item Juegos donde una jugada divide un componente en dos (Kayles, cortar tiras): la tabla se llena con $G(a)\oplus G(b)$.
	\item De nicho: descubrir por fuerza bruta la fórmula cerrada de un juego (por ejemplo $G(n)=n\bmod(k+1)$) y demostrarla después.
\end{itemize}

\textbf{Complejidad:} Sea $N$ el tamaño de la tabla, $M$ la cantidad de jugadas y $D$ la mayor de ellas. Cada estado prueba a lo sumo $M$ jugadas y calcula el MEX en $O(M)$, así que la tabla completa cuesta $O(NM)$, con memoria $O(N)$. La detección del periodo por ventanas recorre $N$ ventanas de largo $D$ y las guarda en un \texttt{map}, lo que cuesta $O(ND\log N)$. Consultar un valor extrapolado es un módulo: $O(1)$. El XOR de $k$ componentes cuesta $O(k)$ y encontrar la jugada ganadora $O(kM)$. Todo lo anterior supone $N$ suficiente para que aparezca una ventana repetida.

\begin{lstlisting}[style=cpp]
	struct TablaGrundy {
		vector<int> movs; // jugadas permitidas: quitar m piedras (m en movs)
		vector<int> g;    // g[n] = numero de Grundy de una pila con n piedras
		
		// READ: menor entero no negativo que no esta en v (v tiene k valores, el mex es <= k).
		static int mex(const vector<int>& v) {
			vector<char> vis(v.size() + 1, 0);
			for (int x : v) if (x < (int)vis.size()) vis[x] = 1;
			int r = 0;
			while (vis[r]) r++;
			return r;
		} // O(k)
		
		// CREATE: calcula g[0..N] probando todas las jugadas de cada estado.
		TablaGrundy(vector<int> jugadas, int N) : movs(move(jugadas)), g(N + 1, 0) {
			for (int n = 1; n <= N; n++) {
				vector<int> sig;
				for (int m : movs) if (m <= n) sig.push_back(g[n - m]);
				g[n] = mex(sig);
			}
		} // O(N * M)
		
		// READ: busca una ventana de D valores que se repite (D = jugada mayor).
		// Devuelve {pre, per}: g[k] == g[k + per] para todo k >= pre. Devuelve {-1, -1} si no hay.
		// Es una prueba rigurosa: para n >= D, g[n] depende solo de g[n-D..n-1].
		pair<int, int> detectar_periodo() const {
			int D = *max_element(movs.begin(), movs.end());
			map<vector<int>, int> visto;
			for (int i = 0; i + D <= (int)g.size(); i++) {
				vector<int> w(g.begin() + i, g.begin() + i + D);
				auto it = visto.find(w);
				if (it != visto.end()) return {it->second, i - it->second};
				visto[w] = i;
			}
			return {-1, -1};
		} // O(N * D * log N)
		
		// READ: numero de Grundy de n; si n excede la tabla, extrapola con el periodo.
		int valor(long long n, int pre, int per) const {
			if (n < (long long)g.size()) return g[n];
			return g[pre + (n - pre) % per];
		} // O(1)
		
		// READ: XOR de los numeros de Grundy de varias pilas (0 = posicion perdedora).
		int xor_total(const vector<long long>& h, int pre, int per) const {
			int x = 0;
			for (long long n : h) x ^= valor(n, pre, per);
			return x;
		} // O(k)
		
		// READ: jugada ganadora {indice de pila, piedras a quitar}, o {-1, -1} si la posicion pierde.
		// Se busca una jugada con G(h) ^ G(h - m) == X, para que el XOR nuevo sea 0.
		pair<int, int> jugada_ganadora(const vector<long long>& h, int pre, int per) const {
			int x = xor_total(h, pre, per);
			if (x == 0) return {-1, -1};
			for (int i = 0; i < (int)h.size(); i++)
			for (int m : movs)
			if (m <= h[i] && (valor(h[i], pre, per) ^ valor(h[i] - m, pre, per)) == x)
			return {i, m};
			return {-1, -1};   // no ocurre si la tabla y el periodo son correctos
		} // O(k * M)
	};
	
	// ------------------------------------------------------------
	// MODIFICACIONES Y COMENTARIOS CON LOS USOS
	// ------------------------------------------------------------
	// 1. Juegos donde una jugada DIVIDE el componente (Kayles: quitar 1 o 2 pinos adyacentes
	//    de una fila de n). Los pedazos son juegos independientes, asi que se combinan por XOR:
	//    for (int n = 1; n <= N; n++) {
		//        vector<int> sig;
		//        for (int a = 0; a <= n - 1; a++) sig.push_back(g[a] ^ g[n - 1 - a]);   // quitar 1
		//        for (int a = 0; a <= n - 2; a++) sig.push_back(g[a] ^ g[n - 2 - a]);   // quitar 2
		//        g[n] = mex(sig);
		//    } // O(N^2)
	//    Aqui g[n] depende de TODOS los valores anteriores, no de una ventana: el periodo
	//    (Kayles es eventualmente periodico de periodo 12 tras un preperiodo de unas 70
	//    posiciones, segun la literatura) es solo una conjetura que se comprueba con un N grande.
	//
	// 2. Formulas cerradas conocidas: Nim (quitar cualquier cantidad de una pila): G(n) = n.
	//    Quitar de 1 a k piedras: G(n) = n % (k + 1); con k = 3 da la tabla 0 1 2 3 0 1 2 3.
	//    Comprobar la conjetura contra la tabla de fuerza bruta antes de usarla.
	//
	// 3. Deteccion empirica del periodo (cuando no hay ventana finita): buscar el menor per
	//    tal que g[i] == g[i + per] para todo i >= pre, exigiendo varios periodos completos.
	//    for (int per = 1; per <= N / 3; per++) {
		//        int i = N - per;
		//        while (i >= 0 && g[i] == g[i + per]) i--;
		//        int pre = i + 1;
		//        if (N + 1 - pre >= 3 * per) return {pre, per};   // al menos 3 periodos observados
		//    } // O(N^2) en el peor caso. NO es una prueba: aumentar N para ganar confianza.
	//
	// 4. Estados compuestos pequenos (por ejemplo un tablero de n <= 20 celdas): representar el
	//    estado como bitmask y hacer DP con memoria en un arreglo o unordered_map.
	//    int grundy(int mask) {
		//        if (memo[mask] != -1) return memo[mask];
		//        vector<int> sig;
		//        for (cada jugada legal desde mask) sig.push_back(grundy(mask_resultante));
		//        return memo[mask] = mex(sig);
		//    } // O(2^n * M)
	//
	// 5. Variante misere (pierde quien hace la ultima jugada): la formula XOR NO aplica en
	//    general. En Nim misere basta jugar como en Nim normal, salvo que si todas las pilas
	//    tienen a lo sumo 1 piedra el resultado se invierte.
	//
	// 6. Verificar la solucion: comparar la respuesta por XOR con una busqueda exhaustiva sobre
	//    el juego completo para configuraciones pequenas (todas las pilas hasta 8, por ejemplo).
	//
	// 7. Un solo componente y solo importa quien gana: sin MEX.
	//    for (int n = 1; n <= N; n++) {
		//        gana[n] = false;
		//        for (int m : movs) if (m <= n && !gana[n - m]) gana[n] = true;
		//    } // O(N * M); mas simple que la tabla de Grundy.
	//
	// 8. Juegos con ciclos (se puede volver a un estado ya visto): Grundy no esta definido
	//    con esta recursion. Usar analisis retrogrado (ver "Trabajar hacia Atras"), que
	//    reconoce los empates.
\end{lstlisting}

\textbf{Ejemplo de funcionamiento:}
\begin{lstlisting}[style=cpp]
	int main() {
		TablaGrundy a({1, 2, 3}, 20);
		for (int n = 0; n <= 7; n++) cout << a.g[n] << " ";
		cout << "\n";                                          // 0 1 2 3 0 1 2 3
		auto [pa, qa] = a.detectar_periodo();
		cout << pa << " " << qa << "\n";                       // 0 4   (g[n] = n % 4)
		cout << a.valor(1000000000000000000LL, pa, qa) << "\n"; // 0   (10^18 % 4 = 0: posicion perdedora)
		cout << a.xor_total({2, 3, 1}, pa, qa) << "\n";        // 0   (2 ^ 3 ^ 1: posicion perdedora)
		cout << a.xor_total({2, 3, 6}, pa, qa) << "\n";        // 3   (2 ^ 3 ^ 2: posicion ganadora)
		
		TablaGrundy b({1, 3, 4}, 30);
		for (int n = 0; n <= 6; n++) cout << b.g[n] << " ";
		cout << "\n";                                          // 0 1 0 1 2 3 2
		auto [pb, qb] = b.detectar_periodo();
		cout << pb << " " << qb << "\n";                       // 0 7
		cout << b.valor(1000000000000000000LL, pb, qb) << "\n"; // 1   (10^18 % 7 = 1, g[1] = 1)
		
		auto j = b.jugada_ganadora({5, 6, 9}, pb, qb);         // g = 3, 2, 0: XOR = 1
		cout << j.first << " " << j.second << "\n";            // 0 1   (quitar 1 de la pila 0: 5 -> 4, g = 2, XOR = 0)
		auto k = b.jugada_ganadora({1000000000000000000LL}, pb, qb);
		cout << k.first << " " << k.second << "\n";            // 0 1   (10^18 -> 10^18 - 1, que es multiplo de 7: g = 0)
		return 0;
	}
\end{lstlisting}

\textbf{Nota:} El teorema de Sprague--Grundy supone un juego imparcial de dos jugadores, con la convención de juego normal y sin partidas infinitas. Si las jugadas dependen del jugador, si hay más de dos jugadores (subsección siguiente) o si el objetivo es otro (por ejemplo, perder al hacer la última jugada), el XOR de los valores no sirve tal cual. Además, \texttt{detectar\_periodo} solo es una prueba cuando el valor de $G(n)$ depende de una ventana fija de valores anteriores, como en los juegos de restar; con juegos que dividen el componente hay que tratar el periodo como una conjetura. Si \texttt{detectar\_periodo} devuelve $\{-1,-1\}$, la tabla es demasiado corta y hay que aumentar $N$ antes de extrapolar. \texttt{jugada\_ganadora} supone que el conjunto de jugadas coincide con el de la tabla.

\textbf{Regla práctica:} Si hay una sola pila y solo importa quién gana, usa la DP booleana. Si el juego se compone de varios componentes independientes, calcula la tabla de Grundy de un componente por fuerza bruta y combina con XOR; si la consulta llega a $n$ enorme, detecta el periodo con la ventana de $D$ valores (restar) o compruébalo empíricamente (otros juegos), y verifica siempre contra la fuerza bruta en casos pequeños. Si hay ciclos, usa análisis retrógrado, y si hay más de dos jugadores o objetivos distintos del juego normal, ve a la siguiente subsección.

\subsection{Juegos con Múltiples Jugadores (más de 2)}

\textbf{Idea:} Con dos jugadores, si una posición no es ganadora para quien mueve, entonces lo es para el rival: «ganar» y «perder» son las dos caras de una misma moneda, y por eso funcionan minimax y Grundy. Con $K\ge 3$ jugadores esa dualidad desaparece. Si $P_0$ no puede ganar, aún importa \textit{cuál} de los otros gana, y una jugada de $P_0$ puede favorecer a $P_1$ y perjudicar a $P_2$. El resultado de un estado debe ser entonces un vector $R(s)=(r_0,\dots,r_{K-1})$ con lo que obtiene cada jugador (por ejemplo, $r_i=1$ para quien gana y $0$ para el resto, o las puntuaciones finales).

Además, el estado debe incluir de quién es el turno, así que la DP se define sobre pares $(s,p)$, y tras una jugada el turno pasa a $p'=(p+1)\bmod K$ si el orden es circular. La regla de decisión es la \textit{inducción hacia atrás} (llamada $\text{max}^n$): en un estado no terminal, el jugador $p$ examina el vector de cada estado sucesor y elige el que maximiza \textit{su propia} componente $r_p$. Los estados terminales se evalúan directamente con las reglas del juego. Con $K=2$ y suma constante esto coincide con minimax.

Hay un punto delicado: cuando varias jugadas dan la misma componente $r_p$ al jugador que mueve, el resultado para los demás depende de la jugada que escoja, y ninguna regla matemática la fija. Por eso «el jugador que gana con juego óptimo» no está bien definido hasta que el enunciado (o el código) da un criterio de desempate. Es lo que ocurre, por ejemplo, cuando un jugador ya no puede ganar y decide a cuál de los rivales «regalar» la partida (\textit{kingmaking}).

\textbf{¿Por qué usarlo?} Frente a Grundy (subsección anterior), que resuelve juegos imparciales de dos jugadores y combina componentes con XOR, aquí ni el teorema ni el XOR valen: no hay un valor escalar por estado que resuma el juego. Frente a minimax o negamax de dos jugadores, que guardan un solo número y suponen intereses opuestos, aquí se guarda un vector y cada jugador maximiza su componente. El costo es que el estado crece: además de la configuración hay que guardar el turno, y cada comparación cuesta $O(K)$ si el resultado es un vector. Se usa cuando el enunciado tiene tres o más jugadores con turnos rotativos, o cuando cada jugador tiene su propio objetivo. No modela alianzas ni acuerdos entre jugadores: cada uno decide solo por sí mismo. Si el juego tiene ciclos, la recursión no termina y hay que pasar a análisis retrógrado, que con más de dos jugadores es bastante más delicado.

\textbf{Casos de uso:}
\begin{itemize}
	\item Juegos con tres o más jugadores y turnos rotativos (quitar piedras entre tres, juegos de cartas simplificados).
	\item Juegos con puntuaciones, donde cada jugador maximiza la suya.
	\item Juegos de eliminación, donde el orden de los turnos cambia según quién queda.
	\item Problemas donde el resultado para cada jugador se pide como vector (quién gana con qué puntaje).
	\item De nicho: verificar por fuerza bruta reglas de desempate ambiguas o comprobar un patrón conjeturado.
\end{itemize}

\textbf{Complejidad:} Sea $N$ la cantidad de estados del juego, $K$ los jugadores y $M$ las jugadas máximas por estado. Hay $N\cdot K$ estados $(s,p)$, y cada uno examina $M$ sucesores. Comparar y copiar un vector de resultados cuesta $O(K)$, por lo que el tiempo total es $O(N\,K^2M)$. Si el resultado es un solo entero (por ejemplo, el índice del ganador), cada comparación es $O(1)$ y queda $O(NKM)$. Si el turno está determinado por el estado (por ejemplo, por la cantidad de jugadas hechas), la dimensión $K$ del estado desaparece. Memoria: $O(NK)$ estados, y $O(NK^2)$ si cada uno guarda un vector.

\begin{lstlisting}[style=cpp]
	// Ejemplo concreto: una pila de n piedras, K jugadores por turnos rotativos, cada jugada
	// quita de 1 a "tope" piedras y GANA quien quita la ultima.
	// Regla de desempate: si varias jugadas dan lo mismo al que mueve, se elige la de menor
	// cantidad de piedras (la primera que se prueba).
	struct JuegoMultijugador {
		int K;                                   // cantidad de jugadores
		int tope;                                // maximo de piedras por jugada
		vector<vector<vector<int>>> memo;        // memo[n][p]: puntuaciones (vacio = sin calcular)
		
		// CREATE: prepara la memoria para pilas de 0 a N piedras.
		JuegoMultijugador(int jugadores, int max_tomar, int N)
		: K(jugadores), tope(max_tomar), memo(N + 1, vector<vector<int>>(jugadores)) {} // O(N * K)
		
		// READ: puntuacion de un estado terminal (n = 0): gana quien hizo la ultima jugada,
		// es decir, el jugador anterior a p.
		vector<int> terminal(int p) const {
			vector<int> r(K, 0);
			r[(p + K - 1) % K] = 1;
			return r;
		} // O(K)
		
		// WRITE: resuelve el estado (n, p) con memoria. Cada jugador maximiza su propia componente.
		const vector<int>& resolver(int n, int p) {
			vector<int>& m = memo[n][p];
			if (!m.empty()) return m;
			if (n == 0) return m = terminal(p);
			vector<int> mejor;
			for (int t = 1; t <= tope && t <= n; t++) {
				const vector<int>& r = resolver(n - t, (p + 1) % K);
				if (mejor.empty() || r[p] > mejor[p]) mejor = r;   // desempate: la primera jugada
			}
			return m = mejor;
		} // O(M * K) por estado nuevo; O(N * K^2 * M) en total
		
		// WRITE: jugador que gana desde (n, p) con juego optimo (calcula y guarda si hace falta).
		int ganador(int n, int p) {
			const vector<int>& r = resolver(n, p);
			return (int)(max_element(r.begin(), r.end()) - r.begin());
		} // O(K) mas el costo de resolver
		
		// WRITE: jugada (piedras a quitar) que elige p en (n, p); -1 si no hay jugadas.
		int mejor_jugada(int n, int p) {
			int mt = -1, ms = -1;
			for (int t = 1; t <= tope && t <= n; t++) {
				int s = resolver(n - t, (p + 1) % K)[p];
				if (s > ms) { ms = s; mt = t; }
			}
			return mt;
		} // O(M) mas el costo de resolver
	};
	
	// ------------------------------------------------------------
	// MODIFICACIONES Y COMENTARIOS CON LOS USOS
	// ------------------------------------------------------------
	// 1. Criterios de desempate: el resultado para los DEMAS depende de la regla. Opciones
	//    comunes: la primera jugada (la del codigo), la que mas beneficia al siguiente en el
	//    turno, o la que menos beneficia al lider. Cambiar el criterio cambia las respuestas:
	//    if (mejor.empty() || r[p] > mejor[p] || (r[p] == mejor[p] && mejor_desempate(r, mejor)))
	//        mejor = r;
	//    El enunciado debe fijarlo; si no lo hace, el problema esta mal definido.
	//
	// 2. Turno determinado por el estado: si p se deduce de n (por ejemplo, p = (jugadas
	//    hechas) % K), se elimina la dimension p de la memoria: memo[n] en lugar de memo[n][p],
	//    y el costo baja a O(N * K * M).
	//
	// 3. Puntuaciones generales: el vector terminal contiene los puntajes finales y cada
	//    jugador maximiza su componente. Para juegos de SUMA CONSTANTE (repartir un total fijo
	//    entre los jugadores) basta guardar todas las componentes menos una.
	//
	// 4. Caso de dos jugadores: con K = 2 y suma constante, el vector es (x, T - x) y max^n
	//    se reduce a minimax (o negamax con un solo numero).
	//
	// 5. Sin recursion (memoria como tabla, para evitar profundidad de pila): llenar de n = 0
	//    hacia arriba, y para cada n todos los p.
	//    for (int n = 0; n <= N; n++) for (int p = 0; p < K; p++) resolver_iterativo(n, p);
	//    // O(N * K * M) mas el costo de copiar vectores.
	//
	// 6. Poda: en max^n la poda alfa-beta solo funciona con cotas sobre la SUMA de las
	//    puntuaciones (por ejemplo, suma constante); con objetivos independientes casi
	//    no hay poda, y el costo exponencial de un juego grande es inevitable.
	//
	// 7. Alianzas y acuerdos: NO se modelan. Cada jugador maximiza solo su componente. Si el
	//    enunciado permite coaliciones o negociacion, se necesita un modelo distinto
	//    (teoria de juegos cooperativos), fuera del alcance de esta tecnica.
	//
	// 8. Juegos con ciclos: la recursion no termina. Con K >= 3 el analisis retrogrado ya no
	//    da una respuesta unica sin fijar el desempate y el tratamiento de las partidas
	//    infinitas; para concursos normalmente el enunciado garantiza que el juego termina.
	//
	// 9. Sprague-Grundy con K >= 3: no aplicar. G(s1) ^ G(s2) ^ ... solo tiene sentido en
	//    juegos imparciales de dos jugadores con juego normal.
\end{lstlisting}

\textbf{Ejemplo de funcionamiento:}
\begin{lstlisting}[style=cpp]
	int main() {
		JuegoMultijugador j(3, 3, 20);            // 3 jugadores, se quitan 1..3 piedras, gana quien quita la ultima
		cout << j.ganador(0, 0) << "\n";          // 2   (sin piedras: gano el jugador anterior a P0, o sea P2)
		cout << j.ganador(3, 0) << "\n";          // 0   (P0 quita las 3 y gana)
		cout << j.ganador(4, 0) << "\n";          // 1   (toda jugada de P0 deja 3, 2 o 1: gana el siguiente, P1)
		cout << j.ganador(5, 2) << "\n";          // 1   (P2 no puede ganar; por el desempate elige quitar 1 y gana P1)
		cout << j.ganador(6, 0) << "\n";          // 0   (P0 quita 1 y deja 5 con el turno de P1, que solo puede regalar el juego)
		cout << j.ganador(10, 0) << "\n";         // 2   (P0 no puede ganar: quita 1 y termina ganando P2)
		cout << j.mejor_jugada(7, 0) << " " << j.mejor_jugada(8, 0) << "\n";   // 2 3
		for (int x : j.resolver(4, 0)) cout << x << " ";
		cout << "\n";                             // 0 1 0   (vector de resultados: gana P1)
		
		JuegoMultijugador d(2, 3, 20);            // con 2 jugadores se recupera n % 4
		cout << d.ganador(4, 0) << " " << d.ganador(5, 0) << "\n";   // 1 0
		return 0;
	}
\end{lstlisting}

\textbf{Nota:} Este ejemplo muestra por qué el desempate importa. Desde $(4,P_0)$ y $(5,P_2)$ ningún jugador que mueve puede ganar, y quien gana depende de la jugada que escoja, que aquí es la de menor cantidad de piedras. Con otro criterio (por ejemplo, favorecer siempre al jugador que sigue) cambiarían esas respuestas, mientras que los estados donde el jugador que mueve sí puede ganar no dependen del desempate. La memoria usa el vector vacío como marca de «sin calcular», lo cual es válido porque todo resultado calculado tiene $K$ componentes. Las referencias que devuelve \texttt{resolver} siguen siendo válidas porque \texttt{memo} nunca cambia de tamaño después de construirse. Con $N$ grande, la recursión llega hasta profundidad $N$, así que hay que pasar a la versión iterativa (modificación 5).

\textbf{Regla práctica:} Si hay tres o más jugadores con turnos rotativos, guarda en el estado el turno y un vector de resultados, y resuelve con inducción hacia atrás: cada jugador maximiza su propia componente. Fija por escrito el criterio de desempate (el enunciado debe dar uno, o el problema está mal definido) y no uses XOR de Grundy ni minimax de un solo valor. Si el turno se deduce del estado, quita esa dimensión; si son exactamente dos jugadores con juego imparcial, usa la tabla de Grundy de la subsección anterior; y si hay ciclos, alianzas o negociación, esta técnica no alcanza.


\end{document}

